% Generated mechanically from frozen input p02/ko/tex/more-algebra.tex.
% Do not edit this generated file; edit the bound locale source instead.

% OK, start here.
%

\setcounter{chapter}{14}
\chapter{대수학의 추가 주제}



\phantomsection
\label{more-algebra-section-phantom}




\section{서론}
\label{more-algebra-section-introduction}

\noindent
이 장에서는 가환대수학에 관한 결과 가운데 첫 번째 가환대수학 장의 결과보다
덜 초등적인 몇 가지를 증명한다. Algebra, Section
\ref{algebra-section-introduction}을 보라.
참고문헌으로는 \cite{MatCA}가 있다.






\section{독자를 위한 안내}
\label{more-algebra-section-advice}

\noindent
가환대수학 장에서보다도 이 장의 각 절은 서로 독립적이다.
Section \ref{more-algebra-section-derived-modules}부터는 환 위 가군의 (비유계)
유도 범주와 그에 딸린 모든 장치를 자유롭게 사용한다.




\section{안정 자유 가군}
\label{more-algebra-section-stably-free}

\noindent
다음은 일반적으로 받아들여지는 정의인 듯하다.

\begin{definition}
\label{more-algebra-definition-stably-free}
$R$을 환이라 하자.
\begin{enumerate}
\item $R$ 위의 두 가군 $M$, $N$에 대하여 어떤 $n, m \geq 0$이 존재하여
$R$-가군으로서
$M \oplus R^{\oplus m} \cong N \oplus R^{\oplus n}$
이면 이 두 가군이 {\it 안정 동형}이라고 한다.
\item 가군 $M$이 자유 가군과 안정 동형이면 이를 {\it 안정 자유}라고 한다.
\end{enumerate}
\end{definition}

\noindent
안정 자유 가군은 사영 가군임에 유의하라.

\begin{lemma}
\label{more-algebra-lemma-exact-category-stably-free}
$R$을 환이라 하자. 유한 생성 사영 $R$-가군의 짧은 완전열
$0 \to P' \to P \to P'' \to 0$이 주어졌다고 하자.
이 세 가군 가운데 $2$개가 안정 자유이면 나머지 $3$번째 가군도 안정 자유이다.
\end{lemma}

\begin{proof}
가군들이 사영이므로 이 열은 분할된다. 따라서 동형
$P = P' \oplus P''$을 택할 수 있다. $P' \oplus R^{\oplus n}$과
$P'' \oplus R^{\oplus m}$이 자유이면
$P \oplus R^{\oplus n + m}$도 자유임을 알 수 있다.
$P'$과 $P$가 안정 자유라고 하자. 즉, $P \oplus R^{\oplus n}$과
$P' \oplus R^{\oplus m}$이 자유라고 하자. 그러면
$$
P'' \oplus (P' \oplus R^{\oplus m}) \oplus R^{\oplus n} =
(P'' \oplus P') \oplus R^{\oplus m} \oplus R^{\oplus n} =
(P \oplus R^{\oplus n}) \oplus R^{\oplus m}
$$
은 자유이다. 따라서 $P''$은 안정 자유이다. 대칭성에 의해 나머지 경우도 얻는다.
\end{proof}

\begin{lemma}
\label{more-algebra-lemma-lift-stably-free}
$R$을 환, $I \subset R$을 아이디얼이라 하자.
$1 + I$의 모든 원소가 단원이라고 가정하자
(다시 말해 $I$는 $R$의 야코브슨 근기에 포함된다).
모든 유한 생성 안정 자유 $R/I$-가군 $E$에 대하여
$M/IM \cong E$를 만족하는 유한 생성 안정 자유 $R$-가군 $M$이 존재한다.
\end{lemma}

\begin{proof}
어떤 $n$, $m$과 동형
$E \oplus (R/I)^{\oplus n} \cong (R/I)^{\oplus m}$을 택한다.
사영 $(R/I)^{\oplus m} \to (R/I)^{\oplus n}$과 단사
$(R/I)^{\oplus n} \to (R/I)^{\oplus m}$을 각각 올리는
$R$-선형 사상 $\varphi : R^{\oplus m} \to R^{\oplus n}$과
$\psi : R^{\oplus n} \to R^{\oplus m}$을 택한다.
그러면 $\varphi \circ \psi : R^{\oplus n} \to R^{\oplus n}$은
$I$를 법으로 항등사상으로 내려간다. 따라서 $I$에 대한 가정에 의해
이 사상의 행렬식은 가역이다. 그러므로
$P = \Ker(\varphi)$는 안정 자유이고 $E$를 올린다.
\end{proof}

\begin{lemma}
\label{more-algebra-lemma-lift-projective}
$R$을 환, $I \subset R$을 아이디얼이라 하자.
$1 + I$의 모든 원소가 단원이라고 가정하자
(다시 말해 $I$는 $R$의 야코브슨 근기에 포함된다).
$M$을 유한 생성 평탄 $R$-가군이라 하고
$M/IM$이 사영 $R/I$-가군이라고 하자.
그러면 $M$은 유한 생성 사영 $R$-가군이다.
\end{lemma}

\begin{proof}
Algebra, Lemma \ref{algebra-lemma-finite-flat-local}에 의해 모든 소아이디얼
$\mathfrak p \subset R$에 대하여 $M_\mathfrak p$는 유한 생성 자유 가군이다.
Algebra, Lemma \ref{algebra-lemma-finite-projective}에 의해 함수
$\rho_M : \mathfrak p \mapsto
\dim_{\kappa(\mathfrak p)} M \otimes_R \kappa(\mathfrak p)$가
$\Spec(R)$ 위에서 국소 상수임을 보이면 충분하다.
$M/IM$이 유한 생성 사영이므로 이는 $V(I) \subset \Spec(R)$ 위에서 참이다.
$\Spec(R)$의 모든 닫힌점이 $V(I)$에 있고, 소아이디얼
$\mathfrak p \subset \mathfrak q \subset R$에 대하여
$\rho_M(\mathfrak p) = \rho_M(\mathfrak q)$이므로,
여기서는 생략하는 위상공간에 관한 초등적인 논증으로 결론을 얻는다.
\end{proof}

\noindent
보조정리 \ref{more-algebra-lemma-lift-stably-free}의 올림은 다음 보조정리에 의해
동형을 제외하고 유일하다.

\begin{lemma}
\label{more-algebra-lemma-isomorphic-finite-projective-lifts}
$R$을 환, $I \subset R$을 아이디얼이라 하자.
$1 + I$의 모든 원소가 단원이라고 가정하자
(다시 말해 $I$는 $R$의 야코브슨 근기에 포함된다).
$P$와 $P'$을 유한 생성 사영 $R$-가군이라 하자. 그러면
\begin{enumerate}
\item $\varphi : P \to P'$이 동형
$\overline{\varphi} : P/IP \to P'/IP'$을 유도하는 $R$-가군 사상이면,
$\varphi$는 동형이다.
\item $P/IP \cong P'/IP'$이면 $P \cong P'$이다.
\end{enumerate}
\end{lemma}

\begin{proof}
(1)의 증명. $P'$은 $R$-가군으로서 사영이므로 사상
$P' \to P'/IP' \xrightarrow{\overline{\varphi}^{-1}} P/IP$의 올림
$\psi : P' \to P$를 택할 수 있다.
나카야마 보조정리(Algebra, Lemma \ref{algebra-lemma-NAK})에 의해
$\psi \circ \varphi$와 $\varphi \circ \psi$는 전사이다.
따라서 이 사상들은 동형이다
(Algebra, Lemma \ref{algebra-lemma-fun}). 그러므로 $\varphi$는 동형이다.

\medskip\noindent
(2)의 증명. 동형 $P/IP \cong P'/IP'$을 택한다.
$P$가 사영이므로 사상 $P \to P/IP \to P'/IP'$의 올림
$\varphi : P \to P'$을 택할 수 있다. (1)에 의해 $\varphi$는 동형이다.
\end{proof}




\section{Artin--Rees 성질에 관한 논평}
\label{more-algebra-section-artin-rees}

\noindent
이 내용의 일부는 \cite{conrad-dejong}에서 가져왔다.
추가 참고문헌을 포함한 일반적인 논의는 \cite[Section 1]{Eis}에서 찾을 수 있다.

\medskip\noindent
$A$를 뇌터 환, $I \subset A$를 아이디얼이라 하자.
유한 생성 $A$-가군의 준동형 $f : M \to N$이 주어지면
$$
f(M) \cap I^nN \subset f(I^{n - c}M)
$$
가 모든 $n \geq c$에 대하여 성립하는 $c \geq 0$이 존재한다.
Algebra, Lemma \ref{algebra-lemma-map-AR}을 보라.
이 상황에서는 {\it $c$가 Artin--Rees 보조정리에서 $f$에 대하여 유효하다}고 한다.

\begin{lemma}
\label{more-algebra-lemma-approximate-complex}
$A$를 뇌터 환이라 하자. $I \subset A$를 $A$의 야코브슨 근기에 포함되는
아이디얼이라 하자. 다음과 같은 유한 생성 $A$-가군의 두 복합체
$$
S : L \xrightarrow{f} M \xrightarrow{g} N
\quad\text{and}\quad
S' : L \xrightarrow{f'} M \xrightarrow{g'} N
$$
가 주어졌다고 하자. 다음을 가정하자.
\begin{enumerate}
\item $c$가 Artin--Rees 보조정리에서 $f$와 $g$에 대하여 유효하다.
\item 복합체 $S$는 완전하다.
\item $f' = f \bmod I^{c + 1}M$이고 $g' = g \bmod I^{c + 1}N$이다.
\end{enumerate}
그러면 $c$는 Artin--Rees 보조정리에서 $g'$에 대하여 유효하고,
복합체 $S'$은 완전하다.
\end{lemma}

\begin{proof}
먼저 $n \geq c$일 때
$g'(M) \cap I^nN \subset g'(I^{n - c}M)$임을 보인다.
$g'(a) \in I^nN$을 만족하는 $M$의 원소 $a$를 택한다.
$g'(a)$를 바꾸지 않으면서 $f'(L)$의 원소로 $a$를 조정하여
$a \in I^{n-c}M$이 되게 하려 한다.
$r < n - c$이고 $a \in I^rM$이라고 가정하자. 그러면
$$
g(a) = g'(a) + (g - g')(a) \in
I^n N + I^{r + c + 1}N = I^{r + c + 1}N.
$$
$g$에 대한 Artin--Rees 성질에 의해 $g(a) \in g(I^{r + 1}M)$이다.
$a_1 \in I^{r + 1}M$인 원소를 택하여 $g(a) = g(a_1)$이라 쓰자.
열 $S$가 완전하므로 $a - a_1 \in f(L)$이다.
따라서 어떤 $b \in L$에 대하여 $a = f(b) + a_1$이라 쓸 수 있다.
그러면 $f(b) = a - a_1 \in I^rM$이다.
$f$에 대한 Artin--Rees 성질에 의해 $r \geq c$이면
$b$를 $I^{r - c}L$의 원소로 바꿀 수 있다.
따라서 모든 경우에 $a = f'(b) + a_2$이고, 여기서
$a_2 = (f - f')(b) + a_1 \in I^{r + 1}M$이다.
(실제로 $c \geq r$이면 가정에 의해
$(f - f')(b) \in I^{r + 1}M$이고, $c < r$이면
$b \in I^{r - c}$이므로 다시
$(f - f')(b) \in I^{c + 1} I^{r - c} M = I^{r + 1}M$이다.)
그러므로 $f'(b) \in f'(L)$로 $a$를 조정하여 $r$을 $1$만큼 늘릴 수 있다.

\medskip\noindent
사실 위의 논증은 모든 $n \geq c$에 대하여
$(g')^{-1}(I^nN) \subset f'(L) + I^{n - c}M$임을 보인다.
따라서
$$
(g')^{-1}(0) = (g')^{-1}(\bigcap I^nN) \subset
\bigcap f'(L) + I^{n - c}M = f'(L)
$$
이다. 여기서는 $I$가 $A$의 야코브슨 근기에 포함됨을 사용했다.
Algebra, Lemma \ref{algebra-lemma-intersection-powers-ideal-module}을 보라.
따라서 $S'$은 완전하다.
\end{proof}

\noindent
환 $A$의 아이디얼 $I \subset A$와 $A$-가군 $M$이 주어지면
$$
\text{Gr}_I(M) = \bigoplus I^nM/I^{n + 1}M.
$$
으로 둔다. 이를 등급 $\text{Gr}_I(A)$-가군으로 생각한다.

\begin{lemma}
\label{more-algebra-lemma-approximate-complex-graded}
가정은 보조정리 \ref{more-algebra-lemma-approximate-complex}에서와 같다.
$Q = \Coker(g)$와 $Q' = \Coker(g')$이라 하자. 그러면
등급 $\text{Gr}_I(A)$-가군으로서
$\text{Gr}_I(Q) \cong \text{Gr}_I(Q')$이다.
\end{lemma}

\begin{proof}
차수 $n$에서
$\text{Gr}_I(Q)_n = I^nN/(I^{n + 1}N + g(M) \cap I^nN)$이고
$Q'$에 대해서도 마찬가지이다. 다음을 주장한다.
$$
g(M) \cap I^nN \subset I^{n + 1}N + g'(M) \cap I^nN.
$$
대칭성에 의해(이 주장의 증명은 $c$가 $g$에 대하여 유효하다는 사실만을
사용하고, 보조정리에 의해 이는 $g'$에 대해서도 성립한다) 이로부터
$$
I^{n + 1}N + g(M) \cap I^nN = I^{n + 1}N + g'(M) \cap I^nN
$$
가 따른다. 따라서 $\text{Gr}_I(Q)_n$과 $\text{Gr}_I(Q')_n$은
$N$의 부분몫으로서 일치하고, 보조정리가 따른다.
$g = g' \bmod I^{c + 1}N$이므로 $n \leq c$일 때 주장은 명백하다.
$n > c$이고 $b \in g(M) \cap I^nN$이라고 하자.
$a \in I^{n - c}M$인 원소를 택하여 $b = g(a)$라고 쓴다.
$b' = g'(a)$로 두면 $b - b' = (g - g')(a) \in I^{n + 1}N$이고,
원하는 바를 얻는다.
\end{proof}

\begin{lemma}
\label{more-algebra-lemma-works-flat-extension}
$A \to B$를 뇌터 환의 평탄 사상이라 하자.
$I \subset A$를 아이디얼, $f : M \to N$을 유한 생성 $A$-가군의 준동형이라 하자.
$c$가 Artin--Rees 보조정리에서 $f$에 대하여 유효하다고 가정하자.
그러면 $c$는 아이디얼 $IB$에 관한 Artin--Rees 보조정리에서
$f \otimes 1 : M \otimes_A B \to N \otimes_A B$에 대하여 유효하다.
\end{lemma}

\begin{proof}
다음에 유의하라.
$$
(f \otimes 1)(M) \cap I^n N \otimes_A B
= (f \otimes 1)\left((f \otimes 1)^{-1}(I^n N \otimes_A B)\right)
$$
한편,
\begin{align*}
(f \otimes 1)^{-1}(I^n N \otimes_A B) &
= \Ker(M \otimes_A B \to N \otimes_A B/(I^n N \otimes_A B)) \\
& =
\Ker(M \otimes_A B \to (N/I^nN) \otimes_A B)
\end{align*}
$A \to B$가 평탄하므로 핵과 여핵을 취하는 것은 $B$와의 텐서곱과 가환한다.
따라서 이는 $f^{-1}(I^nN) \otimes_A B$와 같다.
가정에 의해 $f^{-1}(I^nN)$은
$\Ker(f) + I^{n - c}M$에 포함된다. 그러므로 보조정리가 성립한다.
\end{proof}





\section{환의 섬유곱, I}
\label{more-algebra-section-fibre-products-rings}

\noindent
환의 섬유곱은 스킴의 밂과 관련된다. 스킴의 밂 가운데 몇몇 경우는
More on Morphisms, Section \ref{more-morphisms-section-pushouts}에서 다룬다.

\begin{lemma}
\label{more-algebra-lemma-fibre-product-finite-type}
$R$을 환이라 하자. $A \to B$와 $C \to B$를 $R$-대수 사상이라 하자.
다음을 가정하자.
\begin{enumerate}
\item $R$은 뇌터이다.
\item $A$, $B$, $C$는 $R$ 위 유한형이다.
\item $A \to B$는 전사이다.
\item $B$는 $C$ 위 유한이다.
\end{enumerate}
그러면 $A \times_B C$는 $R$ 위 유한형이다.
\end{lemma}

\begin{proof}
$D = A \times_B C$로 두자. 완전한 행들을 갖는 다음 가환 그림이 있다.
$$
\xymatrix{
0 &
B \ar[l] &
A \ar[l] &
I \ar[l] &
0 \ar[l] \\
0 &
C \ar[l] \ar[u] &
D \ar[l] \ar[u] &
I \ar[l] \ar[u] &
0 \ar[l]
}
$$
$C$-가군으로서 $B$의 생성원 $y_1, \ldots, y_n \in B$를 택한다.
$y_i$로 가는 $x_i \in A$를 택한다.
그러면 $1, x_1, \ldots, x_n$은 $D$-가군으로서 $A$의 생성원이다.
사상 $D \to A \times C$는 단사이고, 환 $A \times C$는
$D$-가군으로서 유한이다(유한 $D$-가군 $A$와 $C$의 직합이기 때문이다).
따라서 Artin--Tate 보조정리
(Algebra, Lemma \ref{algebra-lemma-Artin-Tate})에서 결론이 따른다.
\end{proof}

\begin{lemma}
\label{more-algebra-lemma-formal-consequence}
$R$을 뇌터 환이라 하고 $I$를 유한 집합이라 하자. $\psi_i$와
$\varphi_i$가 전사이고 $Q$, $A_i$, $B_i$가 $R$ 위 유한형인
다음 데카르트 그림이 주어졌다고 하자.
$$
\xymatrix{
\prod B_i &
\prod A_i \ar[l]^{\prod \varphi_i} \\
Q \ar[u]^{\prod \psi_i} &
P \ar[u] \ar[l]
}
$$
그러면 $P$는 $R$ 위 유한형이다.
\end{lemma}

\begin{proof}
보조정리 \ref{more-algebra-lemma-fibre-product-finite-type}와
$I$의 크기에 대한 귀납법에서 따른다.
실제로 $I = I' \amalg \{i_0\}$라 하자.
$I'$을 사용하여 보조정리의 그림으로 정의되는 환을 $P'$이라 하자.
귀납 가정에 의해 $P'$은 유한형이다. 마지막으로 $P$는 다음 섬유곱 그림에 놓인다.
$$
\xymatrix{
B_{i_0} &
A_{i_0} \ar[l] \\
P' \ar[u] &
P \ar[u] \ar[l] &
}
$$
여기에 보조정리를 적용할 수 있다.
\end{proof}

\begin{lemma}
\label{more-algebra-lemma-diagram-localize}
다음 환들의 데카르트 그림, 즉 $B' = B \times_R R'$이 주어졌다고 하자.
$$
\xymatrix{
R &
R' \ar[l]^t \\
B \ar[u]_s &
B'\ar[u] \ar[l]
}
$$
$h \in B'$가 $g \in B$와 $f \in R'$에 대응하고
$s(g) = t(f)$라고 하자. 그러면 다음 그림도 데카르트이다.
$$
\xymatrix{
R_{s(g)} = R_{t(f)} &
(R')_f \ar[l]^-t \\
B_g \ar[u]_s &
(B')_h \ar[u] \ar[l]
}
$$
\end{lemma}

\begin{proof}
등식 $B' = B \times_R R'$은 $B'$-가군의 열
$$
0 \to B' \to B \oplus R' \xrightarrow{s, -t} R
$$
이 완전함을 말한다. $B'$-가군으로서
$B_g = B_h$, $R'_f = R'_h$, $R_{s(g)} = R_{t(f)} = R_h$이다.
국소화의 완전성
(Algebra, Proposition \ref{algebra-proposition-localization-exact})에 의해
$$
0 \to B'_h \to B_g \oplus R'_f \xrightarrow{s, -t} R_{s(g)} = R_{t(f)}
$$
가 완전함을 얻는다. 이로써 보조정리가 증명된다.
\end{proof}

\noindent
다음 환들의 가환 그림을 생각하자.
$$
\xymatrix{
R &
R' \ar[l] \\
B \ar[u] &
B' \ar[u] \ar[l]
}
$$
다음 함자를 생각하자. 여기서 범주의 섬유곱은
Categories, Example \ref{categories-example-2-fibre-product-categories}에서
구성한 것이다.
\begin{equation}
\label{more-algebra-equation-modules}
\text{Mod}_{B'} \longrightarrow
\text{Mod}_B \times_{\text{Mod}_R} \text{Mod}_{R'},\quad
L' \longmapsto (L' \otimes_{B'} B, L' \otimes_{B'} R', can)
\end{equation}
여기서 $can$은 표준 동일시
$L' \otimes_{B'} B \otimes_B R = L' \otimes_{B'} R' \otimes_{R'} R$이다.
이하에서는 오른쪽 범주의 대상을 $(N, M', \varphi)$로 쓴다.
즉 $N$은 $B$-가군이고 $M'$은 $R'$-가군이며,
$\varphi : N \otimes_B R \to M' \otimes_{R'} R$은 동형이다.

\begin{lemma}
\label{more-algebra-lemma-modules}
다음 환들의 가환 그림이 주어졌다고 하자.
$$
\xymatrix{
R &
R' \ar[l] \\
B \ar[u] &
B' \ar[u] \ar[l]
}
$$
함자 (\ref{more-algebra-equation-modules})는 오른쪽 수반 함자를 갖는다.
구체적으로 그 함자는
$$
F : (N, M', \varphi) \longmapsto N \times_\varphi M'
$$
이다(설명은 증명을 보라).
\end{lemma}

\begin{proof}
범주
$\text{Mod}_B \times_{\text{Mod}_R} \text{Mod}_{R'}$의 대상
$(N, M', \varphi)$가 주어지면
$$
N \times_\varphi M' = \{(n, m') \in N \times M' \mid
\varphi(n \otimes 1) = m' \otimes 1\text{ in }M' \otimes_{R'} R\}
$$
로 두고 이를 $B'$-가군으로 본다.
수반성 명제는 $B'$-가군 $L'$과 삼중항 $(N, M', \varphi)$에 대하여
$$
\Hom_{B'}(L', N \times_\varphi M') =
\Hom_B(L' \otimes_{B'} B, N)
\times_{\Hom_R(L' \otimes_{B'} R, M' \otimes_{R'} R)}
\Hom_{R'}(L' \otimes_{B'} R', M')
$$
이 성립한다는 것이다.
Algebra, Lemma \ref{algebra-lemma-adjoint-tensor-restrict}에 의해
오른쪽 항은
$$
\Hom_{B'}(L', N) \times_{\Hom_{B'}(L', M' \otimes_{R'} R)} \Hom_{B'}(L', M')
$$
과 같다. 따라서 이 섬유곱의 원소쌍 $(g, f')$에 대하여
$B'$-선형 사상
$L' \to N \times_\varphi M'$, $l' \mapsto (g(l'), f'(l'))$을 얻음이 명백하다.
반대로 $B'$-선형 사상 $g' : L' \to N \times_\varphi M'$이 주어지면,
$g$를 합성 $L' \to N \times_\varphi M' \to N$으로 두고
$f'$을 합성 $L' \to N \times_\varphi M' \to M'$으로 둘 수 있다.
이 두 구성은 서로 역이고 원하는 동형을 정의한다.
\end{proof}








\section{환의 섬유곱, II}
\label{more-algebra-section-fibre-products-rings-II}

\noindent
이 절에서는 다음 상황에서의 섬유곱을 논의한다.

\begin{situation}
\label{more-algebra-situation-module-over-fibre-product}
이하에서는 환 사상
$$
\xymatrix{
B \ar[r] & A & A' \ar[l]
}
$$
을 생각하며, 여기서 $A' \to A$는 핵이 $I$인 전사라고 가정한다.
이 상황에서 $B' = B \times_A A'$로 두면 데카르트 사각형
$$
\xymatrix{
A & A' \ar[l] \\
B \ar[u] & B' \ar[l] \ar[u]
}
$$
을 얻는다.
\end{situation}

\begin{lemma}
\label{more-algebra-lemma-points-of-fibre-product}
상황 \ref{more-algebra-situation-module-over-fibre-product}에서 위상공간으로서
$$
\Spec(B') = \Spec(B) \amalg_{\Spec(A)} \Spec(A')
$$
이다.
\end{lemma}

\begin{proof}
$B' = B \times_A A'$이므로 스펙트럼의 가환 사각형을 얻는다.
그 원천이 위상공간의 범주에서의 밂이므로(이는
Topology, Section \ref{topology-section-colimits}에 의해 존재한다)
이 사각형은 연속 사상
$$
can : \Spec(B) \amalg_{\Spec(A)} \Spec(A') \longrightarrow \Spec(B')
$$
을 유도한다.

\medskip\noindent
사상 $can$이 전사임을 보이기 위하여 소아이디얼
$\mathfrak q' \subset B'$을 택하자. $I \subset \mathfrak q'$이면
(여기와 이하에서는 $I$를 $A'$의 아이디얼일 뿐 아니라 $B'$의
아이디얼로도 간주한다) $\mathfrak q'$은 $B$의 소아이디얼에 대응하므로 상에 있다.
그렇지 않다면 $h \in I$, $h \not \in \mathfrak q'$인 원소를 택한다.
이 경우 $B_h = A_h = 0$이고 환 사상
$B'_h \to A'_h$는 동형이다. 보조정리
\ref{more-algebra-lemma-diagram-localize}를 보라. 따라서 $\mathfrak q'$은
$I$를 포함하지 않는 유일한 소아이디얼 $\mathfrak p' \subset A'$에 대응한다.

\medskip\noindent
$B' \to B$가 전사이므로 $can$은 성분 $\Spec(B)$ 위에서 단사이다.
위에서 $\Spec(A') \to \Spec(B')$가
$V(I) \subset \Spec(A')$의 여집합 위에서 단사임을 보았다.
또 $V(I) \subset \Spec(A')$는 정확히
$\Spec(A) \to \Spec(A')$의 상이므로, 자명한 집합론적 논증에 의해
$can$은 단사이다.

\medskip\noindent
증명을 마치려면 $can$이 열린 사상임을 보여야 한다.
이를 위해 밂의 열린집합은 $V \amalg U'$ 꼴임에 유의하자.
여기서 $V \subset \Spec(B)$와 $U' \subset \Spec(A')$는 열린집합이고,
$\Spec(A)$에서의 역상들이 일치한다. $v \in V$라 하자.
$v \in D(g) \subset V$인 $g \in B$를 찾을 수 있다.
그 상을 $f \in A$라 하고 $f$로 가는 $f' \in A'$을 택한다. 그러면
$D(f') \cap U' \cap V(I) = D(f') \cap V(I)$이다.
따라서 $V(I) \cap D(f')$와 $D(f') \cap (U')^c$는
$D(f') = \Spec(A'_{f'})$의 서로소인 닫힌 부분집합이다.
어떤 아이디얼 $J \subset A'$에 대하여 $(U')^c = V(J)$라고 쓰자.
방금 보인 서로소성에 의해
$A'_{f'} \to A'_{f'}/IA'_{f'} \times A'_{f'}/JA'_{f'}$
는 전사이므로, $A'_{f'}/IA'_{f'}$에서 $1$로 가고
$A'_{f'}/JA'_{f'}$에서 영으로 가는 $a'' \in A'_{f'}$을 찾을 수 있다.
분모를 없애면 $A$에서 $f^n$으로 가는 원소 $a' \in J$를 얻는다.
그러면 $D(a'f') \subset U'$이다.
$h' = (g^{n + 1}, a'f') \in B'$로 두자.
앞서 인용한 보조정리에 의해
$B'_{h'} = B_{g^{n + 1}} \times_{A_{f^{n + 1}}} A'_{a'f'}$이므로,
$D(h')$의 역상은 밂에서 $v$의 열린 근방이다.
즉 $V \amalg U'$의 상은 $v$의 상의 열린 근방을 포함한다.
$u' \in U'$이고 $u' \not \in V(I)$인 경우에도 같은 명제가 성립하며,
그 더 쉬운 증명은 생략한다.
\end{proof}

\begin{lemma}
\label{more-algebra-lemma-fibre-product-integral}
상황 \ref{more-algebra-situation-module-over-fibre-product}에서
$B \to A$가 정수적이면 $B' \to A'$도 정수적이다.
\end{lemma}

\begin{proof}
$a' \in A'$이라 하고 그 상을 $a \in A$라 하자.
% P02-E0744: frozen source says “monical polynomial”; Korean renders the standard mathematical term while preserving the polynomial.
$a$가 만족하는, 계수가 $B$에 있는 모닉 다항식을
$x^d + b_1 x^{d - 1} + \ldots + b_d$라 하자.
$b_i \in B$로 가는 $b'_i \in B'$을 택한다(이는 가능하다).
그러면 $(a')^d + b'_1 (a')^{d - 1} + \ldots + b'_d$는
$A' \to A$의 핵에 속한다.
$\Ker(B' \to B) = \Ker(A' \to A)$이므로 $b'_d$의 선택을 바꾸어
원하는 등식
$(a')^d + b'_1 (a')^{d - 1} + \ldots + b'_d = 0$을 얻을 수 있다.
\end{proof}

\noindent
상황 \ref{more-algebra-situation-module-over-fibre-product}에서
$A'$, $A$, $B$ 위의 가군을 이용하여 $B'$-가군을 이해하고자 한다.
이를 위하여 다음 함자를 생각한다. 여기서 범주의 섬유곱은
Categories, Example \ref{categories-example-2-fibre-product-categories}에서
구성한 것이다.
\begin{equation}
\label{more-algebra-equation-functor}
\text{Mod}_{B'} \longrightarrow
\text{Mod}_B \times_{\text{Mod}_A} \text{Mod}_{A'},\quad
L' \longmapsto (L' \otimes_{B'} B, L' \otimes_{B'} A', can)
\end{equation}
여기서 $can$은 표준 동일시
$L' \otimes_{B'} B \otimes_B A = L' \otimes_{B'} A' \otimes_{A'} A$이다.
이하에서는 오른쪽 범주의 대상을 $(N, M', \varphi)$로 쓴다.
즉 $N$은 $B$-가군이고 $M'$은 $A'$-가군이며,
$\varphi : N \otimes_B A \to M' \otimes_{A'} A$는 동형이다.
% P02-E0745: frozen source omits “to” in “more convenient to think of”.
그러나 $\varphi$를 $B$-선형 사상
$\varphi : N \to M'/IM'$으로 생각하는 것이 흔히 더 편리하다.
이 사상은 동형
$N \otimes_B A \to M' \otimes_{A'} A = M'/IM'$을 유도한다.

\begin{lemma}
\label{more-algebra-lemma-module-over-fibre-product}
상황 \ref{more-algebra-situation-module-over-fibre-product}에서
함자 (\ref{more-algebra-equation-functor})는 오른쪽 수반 함자, 즉
$$
F : (N, M', \varphi) \longmapsto N \times_{\varphi, M} M'
$$
을 갖는다. 여기서 $M = M'/IM'$이다.
또한 $F$와 (\ref{more-algebra-equation-functor})의 합성은
$\text{Mod}_B \times_{\text{Mod}_A} \text{Mod}_{A'}$ 위의 항등함자이다.
다시 말해 $N' = N \times_{\varphi, M} M'$로 두면
$N' \otimes_{B'} B = N$이고 $N' \otimes_{B'} A' = M'$이다.
\end{lemma}

\begin{proof}
수반성 명제는 더 일반적인 보조정리 \ref{more-algebra-lemma-modules}에서 따른다.
마지막 명제를 증명하기 위하여
$B' = B \times_A A'$과 $N' = N \times_{\varphi, M} M'$을 상기하고,
이 등식들을 다음 두 그림으로 확장한다.
$$
\vcenter{
\xymatrix{
A & A' \ar[l] & I \ar[l] \\
B \ar[u] & B' \ar[l] \ar[u] & J \ar[l] \ar[u]
}
}
\quad\text{and}\quad
\vcenter{
\xymatrix{
M & M' \ar[l] & K \ar[l] \\
N \ar[u]_\varphi & N' \ar[l] \ar[u] & L \ar[l] \ar[u]
}
}
$$
여기서 $I, J, K, L$은 원래 그림의 수평 사상들의 핵이다.
증명을 일련의 관찰로 제시한다.
\begin{enumerate}
\item $K = IM'$이다(보조정리의 명제를 보라).
\item $B' \to B$는 핵이 $J$인 전사이고 $J \to I$는 전단사이다.
\item $N' \to N$은 핵이 $L$인 전사이고 $L \to K$는 전단사이다.
\item $JN' \subset L$이다.
\item $\Im(N \to M)$은 $A$-가군으로서 $M$을 생성한다
($N \otimes_B A = M$이기 때문이다).
\item $\Im(N' \to M')$은 $A'$-가군으로서 $M'$을 생성한다
(이를 $K$를 법으로 하면 성립하고, $L$은 $K$에 동형으로 사상되기 때문이다).
\item $JN' = L$이다(앞 명제에 의해 $L \cong K = I M'$은
$x \in I$, $n' \in N'$인 원소 $x n'$의 상들로 생성되기 때문이다).
\item $N' \otimes_{B'} B = N$이다
($N = N'/L$, $B = B'/J$이고 앞 명제가 성립하기 때문이다).
\item 사상 $\gamma : N' \otimes_{B'} A' \to M'$이 있다.
\item $\gamma$는 전사이다(위를 보라).
% P02-E0746: frozen source has l in K here although the tensor source requires L; preserve the frozen span.
\item 합성 $N' \otimes_{B'} A' \to M' \to M$의 핵은
$l \in K$, $n' \in N'$, $x \in I$인 원소
$l \otimes 1$과 $n' \otimes x$로 생성된다
(가정에 의해 $M = N \otimes_B A$이고
$N' \to N$과 $A' \to A$는 각각 핵이 $L$과 $I$인 전사이기 때문이다).
\item $N' \otimes_{B'} A'$ 안에서
$l \in L$, $n' \in N'$, $x \in I$인 원소
$l \otimes 1$과 $n' \otimes x$로 생성되는 부분가군의 모든 원소는
어떤 $l \in L$에 대하여 $l \otimes 1$로 쓸 수 있다
($J$가 $I$에 동형으로 사상되므로
$N' \otimes_{B'} A'$에서 $n' \otimes x = n'x \otimes 1$이고,
마찬가지로 $n' \in N'$, $x \in J$, $a' \in A'$일 때
$N' \otimes_{B'} A'$에서
$x n' \otimes a' = n' \otimes xa' = n'(xa') \otimes 1$이다.
이미 $JN' = L$임을 보였으므로 명제가 따른다).
\item $\gamma$의 핵은 영이다
((10)과 (11)에 의해 핵의 모든 원소는 $l \in L$인
$l \otimes 1$ 꼴이고, $\gamma$는 이를 $l \in K \subset M'$로 보내기 때문이다).
\end{enumerate}
이로써 증명이 끝난다.
\end{proof}

\begin{lemma}
\label{more-algebra-lemma-module-over-fibre-product-bis}
보조정리 \ref{more-algebra-lemma-module-over-fibre-product}의 상황에서
$B'$-가군 $L'$에 대한 수반 사상
$$
L' \longrightarrow 
(L' \otimes_{B'} B) \times_{(L' \otimes_{B'} A)} (L' \otimes_{B'} A')
$$
은 전사이지만 일반적으로 단사는 아니다.
\end{lemma}

\begin{proof}
보조정리 \ref{more-algebra-lemma-module-over-fibre-product}의 증명에서와 같이
$J \subset B'$를 사상 $B' \to B$의 핵이라 하자.
그러면 $L' \otimes_{B'} B = L'/JL'$이다.
따라서 전사성을 보이려면 섬유곱의 $(0, z)$ 꼴 원소들이
보조정리의 사상의 상에 있음을 보이면 충분하다.
사상 $L' \otimes_{B'} A' \to L' \otimes_{B'} A$의 핵은
$L' \otimes_{B'} I \to L' \otimes_{B'} A'$의 상이다.
$B' \to A'$가 유도하는 사상 $J \to I$가 동형이므로 합성
$$
L' \otimes_{B'} J \to L' \to
(L' \otimes_{B'} B) \times_{(L' \otimes_{B'} A)} (L' \otimes_{B'} A')
$$
은 $L' \otimes_{B'} J$에서 $(0, z)$ 꼴 원소들의 집합으로 전사이다.
이 사상이 일반적으로 단사가 아님을 보기 위해 간단한 예를 들자.
체 $k$를 택하고
$B' = k[x, y]/(xy)$, $A' = B'/(x)$, $B = B'/(y)$,
$A = B'/(x, y)$, $L' = B'/(x - y)$로 두자.
이 경우 $L'$에서 $x$의 류는 영이 아니지만 위에 표시한 화살표에서 영으로 간다.
\end{proof}

\begin{lemma}
\label{more-algebra-lemma-surjection-module-over-fibre-product}
상황 \ref{more-algebra-situation-module-over-fibre-product}에서
$(N_1, M'_1, \varphi_1) \to (N_2, M'_2, \varphi_2)$를
$\text{Mod}_B \times_{\text{Mod}_A} \text{Mod}_{A'}$의 사상이라 하자.
$N_1 \to N_2$와 $M'_1 \to M'_2$가 전사이면
$$
N_1 \times_{\varphi_1, M_1} M'_1 \to N_2 \times_{\varphi_2, M_2} M'_2
$$
도 전사이다. 여기서 $M_1 = M'_1/IM'_1$이고 $M_2 = M'_2/IM'_2$이다.
\end{lemma}

\begin{proof}
$(x_2, y_2) \in N_2 \times_{\varphi_2, M_2} M'_2$를 택한다.
$x_2$로 가는 $x_1 \in N_1$을 택한다.
$M'_1 \to M_1$이 전사이므로 $\varphi_1(x_1)$로 가는
$y_1 \in M'_1$을 찾을 수 있다.
그러면 $(x_1, y_1)$은
$N_2 \times_{\varphi_2, M_2} M'_2$에서 $(x_2, y'_2)$로 간다.
따라서 $(0, y_2)$ 꼴 원소들이 사상의 상에 있음을 보이면 충분하다.
여기서 $y_2 \in IM'_2$이다.
$t_i \in I$인 원소들을 사용하여
$y_2 = \sum t_i y_{2, i}$라고 쓰자.
$y_{2, i}$로 가는 $y_{1, i} \in M'_1$을 택한다.
그러면 $y_1 = \sum t_iy_{1, i} \in IM'_1$이고 원소 $(0, y_1)$이
원하는 일을 한다.
\end{proof}

\begin{lemma}
\label{more-algebra-lemma-finite-module-over-fibre-product}
$A, A', B, B', I, M, M', N, \varphi$가 보조정리
\ref{more-algebra-lemma-module-over-fibre-product}에서와 같다고 하자.
% P02-E0747: frozen source omits the copula in “If N finite over B”.
$N$이 $B$ 위 유한이고 $M'$이 $A'$ 위 유한이면
$N' = N \times_{\varphi, M} M'$은 $B'$ 위 유한이다.
\end{lemma}

\begin{proof}
보조정리 \ref{more-algebra-lemma-module-over-fibre-product}의 결과를
더 언급하지 않고 사용하겠다.
$B$ 위 $N$의 생성원 $y_1, \ldots, y_r$와
$A'$ 위 $M'$의 생성원 $x_1, \ldots, x_s$를 택한다.
$N = N' \otimes_{B'} B$이고 $B' \to B$가 전사이므로
$N$에서 $y_1, \ldots, y_r$로 가는
$u_1, \ldots, u_r \in N'$을 찾을 수 있다.
$M' = N' \otimes_{B'} A'$이므로 어떤 $a'_{ij} \in A'$에 대하여
$x_i = \sum v_j \otimes a'_{ij}$가 되도록 하는
$v_1, \ldots, v_t \in N'$을 찾을 수 있다.
특히 $v_j$의 상 $\overline{v}_j \in M'$들은
$A'$ 위에서 $M'$을 생성한다.
$u_1, \ldots, u_r, v_1, \ldots, v_t$가
$B'$-가군으로서 $N'$을 생성한다고 주장한다.
실제로 $\xi \in N'$을 택하자.
먼저 $\xi$와 $\sum b'_i u_i$가 $N$의 같은 원소로 가도록
$b'_1, \ldots, b'_r \in B'$을 택한다.
$B' \to B$가 전사이고 $y_1, \ldots, y_r$가 $B$ 위에서 $N$을
생성하므로 이는 가능하다.
차 $\xi - \sum b'_i u_i$는 $IM'$에 속하는 어떤 $\theta$에 대하여
$(0, \theta)$ 꼴이다.
$t_j \in I$를 사용하여 $\theta$를 $\sum t_j\overline{v}_j$라고 쓰자.
$J = \Ker(B' \to B)$가 $I$에 동형으로 사상되므로
$t_j$로 가는 $s_j \in J \subset B'$을 택할 수 있다.
$N' = N \times_{\varphi, M} M'$이므로 원하는 대로
$\xi = \sum b'_i u_i + \sum s_j v_j$가 따른다.
\end{proof}

\begin{lemma}
\label{more-algebra-lemma-flat-module-over-fibre-product}
$A, A', B, B', I$가 상황
\ref{more-algebra-situation-module-over-fibre-product}에서와 같다고 하자.
\begin{enumerate}
\item $(N, M', \varphi)$를
$\text{Mod}_B \times_{\text{Mod}_A} \text{Mod}_{A'}$의 대상이라 하자.
$M'$이 $A'$ 위 평탄이고 $N$이 $B$ 위 평탄이면
$N' = N \times_{\varphi, M} M'$은 $B'$ 위 평탄이다.
% P02-E0748: frozen source uses undefined L rather than L' in all three tensor factors; preserve every frozen span.
\item $L'$이 평탄 $B'$-가군이면
$L' = (L \otimes_{B'} B) \times_{(L \otimes_{B'} A)} (L \otimes_{B'} A')$이다.
\item 평탄 $B'$-가군의 범주는
$\text{Mod}_B \times_{\text{Mod}_A} \text{Mod}_{A'}$의 다음
충만한 부분범주와 동치이다. 이 부분범주의 대상은
$N$이 $B$ 위 평탄이고 $M'$이 $A'$ 위 평탄인 삼중항
$(N, M', \varphi)$이다.
\end{enumerate}
\end{lemma}

\begin{proof}
증명에서는 보조정리 \ref{more-algebra-lemma-module-over-fibre-product}를
더 언급하지 않고 사용한다.

\medskip\noindent
(1)의 증명. $J = \Ker(B' \to B)$로 두자.
이는 $I = \Ker(A' \to A)$에 동형으로 사상되는 $B'$의 아이디얼이다.
$\mathfrak b' \subset B'$을 아이디얼이라 하자.
Algebra, Lemma \ref{algebra-lemma-flat}에 의해
$\mathfrak b' \otimes_{B'} N' \to N'$이 단사임을 보여야 한다.
$N$이 $B$ 위 평탄이므로
$$
\mathfrak b'/(\mathfrak b' \cap J) \otimes_{B'} N' =
\mathfrak b'/(\mathfrak b' \cap J) \otimes_B N \to N
$$
은 단사이다. 열
$\mathfrak b' \cap J \to \mathfrak b' \to
\mathfrak b'/(\mathfrak b' \cap J) \to 0$이 완전하므로
$(\mathfrak b' \cap J) \otimes_{B'} N' \to N'$이 단사임을 보이면
충분하다고 결론 내린다. 따라서 $\mathfrak b' \subset J$라고 가정할 수 있다.
다음으로 $J \to I$가 동형이므로
$$
J \otimes_{B'} N' =
I \otimes_{A'} A' \otimes_{B'} N' =
I \otimes_{A'} M'
$$
이다. $M'$이 평탄 $A'$-가군이므로 이는 $M'$에 단사로 사상된다.
따라서 $J \otimes_{B'} N' \to N'$은 단사이고
$$\text{Tor}_1^{B'}(B'/J, N') = 0$$
이라고 결론 내린다.
Algebra, Remark \ref{algebra-remark-Tor-ring-mod-ideal}을 보라.
그러므로 $B'$ 위의 $N'$과 아이디얼 $J$에
Algebra, Lemma \ref{algebra-lemma-what-does-it-mean}을 적용할 수 있다.
아이디얼 $\mathfrak b' \subset J$로 돌아가서
$\mathfrak b' \subset \mathfrak b'' \subset J$를
$I$에서의 상이 $I$의 $A'$-부분가군인 가장 작은 아이디얼이라 하자.
다시 말해 $J = I$를 $A'$-가군으로 보면
$\mathfrak b'' = A' \mathfrak b'$이다.
그러면 $\mathfrak b''/\mathfrak b'$은 $J$에 의해 소멸되고,
보조정리의 적용으로 얻는
$\text{Tor}_1^{B'}(\mathfrak b''/\mathfrak b', N')$의 소멸에 의해
짧은 완전열
$$
0 \to \mathfrak b' \otimes_{B'} N' \to
\mathfrak b'' \otimes_{B'} N' \to
\mathfrak b''/\mathfrak b' \otimes_{B'} N' \to 0
$$
을 얻는다. 따라서 $\mathfrak b'$을 $\mathfrak b''$로 바꿀 수 있다.
특히 $\mathfrak b'$이 $A'$-가군이고 $A'$의 아이디얼로 사상된다고
가정할 수 있다. 그러면
$$
\mathfrak b' \otimes_{B'} N' =
\mathfrak b' \otimes_{A'} A' \otimes_{B'} N' =
\mathfrak b' \otimes_{A'} M'
$$
이다. $M'$이 $A'$ 위 평탄이라는 가정에 의해 이 텐서곱은
$M'$에 단사로 사상된다.
따라서 $\mathfrak b' \otimes_{B'} N' \to N' \to M'$은 단사이고,
원하는 대로 첫 번째 사상도 단사이다.

\medskip\noindent
(2)의 증명. 짧은 완전열
$0 \to B' \to B \oplus A' \to A \to 0$을
$B'$ 위에서 $L'$과 텐서하면 된다.

\medskip\noindent
(3)의 증명. (1)과 (2)의 직접적인 귀결이다.
\end{proof}

\begin{lemma}
\label{more-algebra-lemma-finitely-presented-module-over-fibre-product}
$A, A', B, B', I$가 상황
\ref{more-algebra-situation-module-over-fibre-product}에서와 같다고 하자.
유한 생성 사영 $B'$-가군의 범주는
$\text{Mod}_B \times_{\text{Mod}_A} \text{Mod}_{A'}$의 다음
충만한 부분범주와 동치이다. 이 부분범주의 대상은
$N$이 $B$ 위 유한 생성 사영이고 $M'$이 $A'$ 위 유한 생성 사영인
삼중항 $(N, M', \varphi)$이다.
\end{lemma}

\begin{proof}
가군이 유한 생성 사영일 필요충분조건은 유한 표시이고 평탄인 것이다.
Algebra, Lemma \ref{algebra-lemma-finite-projective}를 보라.
보조정리 \ref{more-algebra-lemma-flat-module-over-fibre-product}와
\ref{more-algebra-lemma-finite-module-over-fibre-product}를 사용하면 다음을 보이는
문제로 환원된다. $N$이 $B$ 위 유한 생성 사영이고
$M'$이 $A'$ 위 유한 생성 사영이면
$N' = N \times_{\varphi, M} M'$은 유한 표시 $B'$-가군이다.

\medskip\noindent
보조정리 \ref{more-algebra-lemma-finite-module-over-fibre-product}에 의해
가군 $N'$은 $B'$ 위 유한이다.
핵이 $K'$인 전사 $(B')^{\oplus n} \to N'$을 택한다.
기저변환으로 핵이 각각 $K_B$, $K_{A'}$, $K_A$인 사상
$B^{\oplus n} \to N$, $(A')^{\oplus n} \to M'$,
$A^{\oplus n} \to M$을 얻는다. 표준 사상
$$
K' \longrightarrow K_B \times_{K_A} K_{A'}
$$
이 있다. 한편 $N' = N \times_{\varphi, M} M'$이고
$B' = B \times_A A'$이므로, 표시된 화살표의 역인 표준 사상
$K_B \times_{K_A} K_{A'} \to K'$도 있다.
따라서 표시된 사상은 동형이다.
Algebra, Lemma \ref{algebra-lemma-extension}에 의해
가군 $K_B$와 $K_{A'}$은 유한이다.
보조정리 \ref{more-algebra-lemma-finite-module-over-fibre-product}에 의해
$K_B \to K_A$와 $K_{A'} \to K_A$가 동형
$K_B \otimes_B A = K_A = K_{A'} \otimes_{A'} A$를 유도한다면
$K'$은 유한 $B'$-가군이라고 결론 내린다.
% P02-E0749: frozen source has plural subject with singular “implies”.
이는 평탄성 가정에 의해 텐서한 뒤에도 열
$$
0 \to K_B \to B^{\oplus n} \to N \to 0
\quad\text{and}\quad
0 \to K_{A'} \to (A')^{\oplus n} \to M' \to 0
$$
이 완전하게 남기 때문에 성립한다.
Algebra, Lemma \ref{algebra-lemma-flat-tor-zero}를 보라.
\end{proof}

\section{환의 섬유곱, III}
\label{more-algebra-section-fibre-products-rings-III}

\noindent
이 절에서는 다음 상황에서의 섬유곱을 논의한다.

\begin{situation}
\label{more-algebra-situation-relative-module-over-fibre-product}
$A, A', B, B', I$가 상황
\ref{more-algebra-situation-module-over-fibre-product}에서와 같다고 하자.
$B' \to D'$을 환 사상이라 하자.
$D = D' \otimes_{B'} B$, $C' = D' \otimes_{B'} A'$,
$C = D' \otimes_{B'} A$로 두자. 그러면 환들의 큰 가환 그림
$$
\xymatrix{
C & & & C' \ar[lll] \\
& A \ar[ul] & A' \ar[l] \ar[ru] \\
& B \ar[u] \ar[ld] & B' \ar[l] \ar[u] \ar[rd] \\
D \ar[uuu] & & & D' \ar[lll] \ar[uuu]
}
$$
을 얻는다.
$D' \to D \times_C C'$가 동형이라고 {\bf 가정하지 않음}에 유의하라
\footnote{그러나 보조정리 \ref{more-algebra-lemma-module-over-fibre-product-bis}에 의해
$D' \to D \times_C C'$는 전사이다.}.
이 상황에서 (\ref{more-algebra-equation-functor})와 유사한 함자
\begin{equation}
\label{more-algebra-equation-relative-functor}
\text{Mod}_{D'} \longrightarrow
\text{Mod}_D \times_{\text{Mod}_C} \text{Mod}_{C'},\quad
L' \longmapsto (L' \otimes_{D'} D, L' \otimes_{D'} C', can)
\end{equation}
가 있다. 다음에 유의하라.
% P02-E0750: both frozen identities use undefined L on the right; preserve every frozen span.
$L' \otimes_{D'} D = L \otimes_{D'} (D' \otimes_{B'} B) = L \otimes_{B'} B$
이고, 마찬가지로
$L' \otimes_{D'} C' = L \otimes_{D'} (D' \otimes_{B'} A') = L \otimes_{B'} A'$이다.
따라서 그림
$$
\xymatrix{
\text{Mod}_{D'} \ar[r] \ar[d] &
\text{Mod}_D \times_{\text{Mod}_C} \text{Mod}_{C'} \ar[d] \\
\text{Mod}_{B'} \ar[r] &
\text{Mod}_B \times_{\text{Mod}_A} \text{Mod}_{A'}
}
$$
은 가환한다.
이하에서는 $\text{Mod}_D \times_{\text{Mod}_C} \text{Mod}_{C'}$의
대상을 $(N, M', \varphi)$로 쓴다.
즉 $N$은 $D$-가군이고 $M'$은 $C'$-가군이며,
% P02-E0751: frozen English uses the wrong article before C'-module.
$\varphi : N \otimes_B A \to M' \otimes_{A'} A$는 $C$-가군의 동형이다.
% P02-E0752: frozen source omits “to” in “more convenient to think of”.
그러나 $\varphi$를 $D$-선형 사상
$\varphi : N \to M'/IM'$으로 생각하는 것이 흔히 더 편리하다.
이 사상은 동형
$N \otimes_B A \to M' \otimes_{A'} A = M'/IM'$을 유도한다.
\end{situation}

\begin{lemma}
\label{more-algebra-lemma-relative-module-over-fibre-product}
상황 \ref{more-algebra-situation-relative-module-over-fibre-product}에서
함자 (\ref{more-algebra-equation-relative-functor})는 오른쪽 수반 함자, 즉
$$
F : (N, M', \varphi) \longmapsto N \times_{\varphi, M} M'
$$
을 갖는다. 여기서 $M = M'/IM'$이다.
또한 $F$와 (\ref{more-algebra-equation-relative-functor})의 합성은
$\text{Mod}_D \times_{\text{Mod}_C} \text{Mod}_{C'}$ 위의 항등함자이다.
다시 말해 $N' = N \times_{\varphi, M} M'$로 두면
$N' \otimes_{D'} D = N$이고 $N' \otimes_{D'} C' = M'$이다.
\end{lemma}

\begin{proof}
수반성 명제는 더 일반적인 보조정리 \ref{more-algebra-lemma-modules}에서 따른다.
마지막 명제는 보조정리 \ref{more-algebra-lemma-module-over-fibre-product}의
대응하는 명제에서 따른다. 실제로
$N' \otimes_{D'} D = N' \otimes_{D'} D' \otimes_{B'} B = N' \otimes_{B'} B$
이고
$N' \otimes_{D'} C' = N' \otimes_{D'} D' \otimes_{B'} A' = N' \otimes_{B'} A'$
이기 때문이다.
\end{proof}

\begin{lemma}
\label{more-algebra-lemma-relative-surjection-ideals}
상황 \ref{more-algebra-situation-relative-module-over-fibre-product}에서
$J = \Ker(B' \to B)$라 하면 사상 $JD' \to IC'$는 전사이다.
\end{lemma}

\begin{proof}
$C' = D' \otimes_{B'} A'$이므로 $IC'$는 사상
$D' \otimes_{B'} I = C' \otimes_{A'} I \to C'$의 상이다.
환 사상 $B' \to A'$가 동형 $J \to I$를 유도하므로 보조정리가 따른다.
\end{proof}

\begin{lemma}
\label{more-algebra-lemma-relative-finite-module-over-fibre-product}
$A, A', B, B', C, C', D, D', I, M', M, N, \varphi$가 보조정리
\ref{more-algebra-lemma-relative-module-over-fibre-product}에서와 같다고 하자.
% P02-E0753: frozen source omits the copula in “If N finite over D”.
$N$이 $D$ 위 유한이고 $M'$이 $C'$ 위 유한이면
$N' = N \times_{\varphi, M} M'$은 $D'$ 위 유한이다.
\end{lemma}

\begin{proof}
보조정리 \ref{more-algebra-lemma-module-over-fibre-product-bis}에 의해
$D' \to D \times_C C'$가 전사임을 상기하자.
$N' = N \times_{\varphi, M} M'$은
$D \times_C C'$ 위의 가군임에 유의하라.
자료 $C, C', D, D', IC', M', M, N, \varphi$에
보조정리 \ref{more-algebra-lemma-finite-module-over-fibre-product}를 적용하면
$N' = N \times_{\varphi, M} M'$이 $D \times_C C'$ 위 유한임을 알 수 있다.
따라서 이는 $D'$ 위 유한이다.
\end{proof}

\begin{lemma}
\label{more-algebra-lemma-relative-flat-module-over-fibre-product}
$A, A', B, B', C, C', D, D', I$가 상황
\ref{more-algebra-situation-relative-module-over-fibre-product}에서와 같다고 하자.
\begin{enumerate}
\item $(N, M', \varphi)$를
$\text{Mod}_D \times_{\text{Mod}_C} \text{Mod}_{C'}$의 대상이라 하자.
$M'$이 $A'$ 위 평탄이고 $N$이 $B$ 위 평탄이면
$N' = N \times_{\varphi, M} M'$은 $B'$ 위 평탄이다.
% P02-E0754: frozen source starts with L' but uses undefined L in all three tensor factors; preserve every frozen span.
\item $L'$이 $B'$ 위 평탄인 $D'$-가군이면
$L' = (L \otimes_{D'} D) \times_{(L \otimes_{D'} C)} (L \otimes_{D'} C')$이다.
\item $B'$ 위 평탄인 $D'$-가군의 범주는
$\text{Mod}_D \times_{\text{Mod}_C} \text{Mod}_{C'}$에서
$N$이 $B$ 위 평탄이고 $M'$이 $A'$ 위 평탄인 대상
$(N, M', \varphi)$들의 범주와 동치이다.
\end{enumerate}
\end{lemma}

\begin{proof}
(1)은 보조정리 \ref{more-algebra-lemma-flat-module-over-fibre-product}의 (1)에서 따른다.

\medskip\noindent
(2)는 보조정리 \ref{more-algebra-lemma-flat-module-over-fibre-product}의 (2)에서 따른다.
여기서는
$L' \otimes_{D'} D = L' \otimes_{B'} B$,
$L' \otimes_{D'} C' = L' \otimes_{B'} A'$,
$L' \otimes_{D'} C = L' \otimes_{B'} A$를 사용한다.
상황 \ref{more-algebra-situation-relative-module-over-fibre-product}의 논의를 보라.

\medskip\noindent
(3)은 (1)과 (2)의 직접적인 귀결이다.
\end{proof}

\noindent
다음 보조정리는 절대적 경우의 대응 명제
(보조정리 \ref{more-algebra-lemma-finitely-presented-module-over-fibre-product})보다
훨씬 흥미롭지만, 증명은 본질적으로 같다.

\begin{lemma}
\label{more-algebra-lemma-relative-finitely-presented-module-over-fibre-product}
$A, A', B, B', C, C', D, D', I, M', M, N, \varphi$가 보조정리
\ref{more-algebra-lemma-relative-module-over-fibre-product}에서와 같다고 하자.
다음을 가정하자.
\begin{enumerate}
\item $N$은 $D$ 위 유한 표시이고 $B$ 위 평탄이다.
% P02-E0755: frozen source omits the copula in this condition.
\item $M'$은 $C'$ 위 유한 표시이고 $A'$ 위 평탄이다.
\item 환 사상 $B' \to D'$은 $B' \to D'' \to D'$으로 분해되며,
$B' \to D''$은 평탄이고 $D'' \to D'$은 유한 표시이다.
\end{enumerate}
그러면 $N' = N \times_M M'$은 $D'$ 위 유한 표시이다.
\end{lemma}

\begin{proof}
핵이 유한 생성인 전사
$D''' = D''[x_1, \ldots, x_n] \to D'$을 택하고 그 핵을 $J$라 하자.
Algebra, Lemma \ref{algebra-lemma-finite-finitely-presented-extension}에 의해
$N'$이 $D'''$-가군으로서 유한 표시임을 보이면 충분하다.
또한
$D''' \otimes_{B'} B \to D' \otimes_{B'} B = D$와
$D''' \otimes_{B'} A' \to D' \otimes_{B'} A' = C'$은 전사이고,
그 핵들은 $J$의 상으로 생성된다.
따라서 다시 Algebra, Lemma
\ref{algebra-lemma-finite-finitely-presented-extension}에 의해
$N$은 유한 표시 $D''' \otimes_{B'} B$-가군이고
$M'$은 유한 표시 $D''' \otimes_{B'} A'$-가군이다.
그러므로 $D'$을 $D'''$으로, $D$를
$D''' \otimes_{B'} B$로, 나머지도 마찬가지로 바꿀 수 있다.
$D'''$은 $B'$ 위 평탄이므로 $B' \to D'$이 평탄이라고 가정할 수 있다.

\medskip\noindent
$B' \to D'$이 평탄하다고 가정하자.
보조정리 \ref{more-algebra-lemma-relative-finite-module-over-fibre-product}에 의해
가군 $N'$은 $D'$ 위 유한이다.
핵이 $K'$인 전사 $(D')^{\oplus n} \to N'$을 택한다.
기저변환으로 핵이 각각 $K_D$, $K_{C'}$, $K_C$인 사상
$D^{\oplus n} \to N$, $(C')^{\oplus n} \to M'$,
$C^{\oplus n} \to M$을 얻는다. 표준 사상
$$
K' \longrightarrow K_D \times_{K_C} K_{C'}
$$
이 있다. 한편 $N' = N \times_M M'$이고
$D' = D \times_C C'$이므로
(평탄 $B'$-가군 $D'$에 적용한 보조정리
\ref{more-algebra-lemma-flat-module-over-fibre-product}에 의해)
표시된 화살표의 역인 표준 사상
$K_D \times_{K_C} K_{C'} \to K'$도 있다.
따라서 표시된 사상은 동형이다.
Algebra, Lemma \ref{algebra-lemma-extension}에 의해
가군 $K_D$와 $K_{C'}$은 유한이다.
보조정리 \ref{more-algebra-lemma-relative-finite-module-over-fibre-product}에 의해
$K_D \to K_C$와 $K_{C'} \to K_C$가 동형
$K_D \otimes_B A = K_C = K_{C'} \otimes_{A'} A$를 유도한다면
$K'$은 유한 $D'$-가군이라고 결론 내린다.
% P02-E0756: frozen source has plural subject with singular “implies”.
이는 평탄성 가정에 의해 텐서한 뒤에도 열
$$
0 \to K_D \to D^{\oplus n} \to N \to 0
\quad\text{and}\quad
0 \to K_{C'} \to (C')^{\oplus n} \to M' \to 0
$$
이 완전하게 남기 때문에 성립한다.
Algebra, Lemma \ref{algebra-lemma-flat-tor-zero}를 보라.
\end{proof}

\begin{lemma}
\label{more-algebra-lemma-properties-algebras-over-fibre-product}
$A, A', B, B', I$가 상황
\ref{more-algebra-situation-module-over-fibre-product}에서와 같다고 하자.
% P02-E0757: frozen English uses the wrong articles before B-algebra and A'-algebra.
$(D, C', \varphi)$를 $B$-대수 $D$, $A'$-대수 $C'$,
동형 $D \otimes_B A \to C'/IC' = C$로 이루어진 계라 하자.
$D' = D \times_C C'$로 두자
(보조정리 \ref{more-algebra-lemma-module-over-fibre-product}에서와 같이).
그러면 다음이 성립한다.
\begin{enumerate}
% P02-E0758: frozen source omits “of” in the finite-type predicates.
\item $B' \to D'$이 유한형일 필요충분조건은
$B \to D$와 $A' \to C'$이 유한형인 것이다.
\item $B' \to D'$이 평탄일 필요충분조건은
$B \to D$와 $A' \to C'$이 평탄인 것이다.
\item $B' \to D'$이 평탄이고 유한 표시일 필요충분조건은
$B \to D$와 $A' \to C'$이 평탄이고 유한 표시인 것이다.
\item $B' \to D'$이 매끄러울 필요충분조건은
$B \to D$와 $A' \to C'$이 매끄러운 것이다.
\item $B' \to D'$이 에탈일 필요충분조건은
$B \to D$와 $A' \to C'$이 에탈인 것이다.
\end{enumerate}
또 $D'$이 평탄 $B'$-대수이면
$D' \to (D' \otimes_{B'} B) \times_{(D' \otimes_{B'} A)} (D' \otimes_{B'} A')$
는 동형이다. 이 방식으로 평탄 $B'$-대수의 범주는 위와 같은 계
$(D, C', \varphi)$ 가운데 $D$가 $B$ 위 평탄이고
$C'$이 $A'$ 위 평탄인 것들의 범주와 동치이다.
\end{lemma}

\begin{proof}
함의 “$\Rightarrow$”는 Algebra의 보조정리
\ref{algebra-lemma-base-change-finiteness},
\ref{algebra-lemma-flat-base-change},
\ref{algebra-lemma-base-change-smooth},
\ref{algebra-lemma-etale}에서 따른다.
실제로 보조정리 \ref{more-algebra-lemma-module-over-fibre-product}에 의해
$D' \otimes_{B'} B = D$이고 $D' \otimes_{B'} A' = C'$이기 때문이다.
따라서 반대 방향의 함의들을 증명하면 충분하다.

\medskip\noindent
(1)에 대하여.
% P02-E0759: frozen source omits both copulas in this assumption.
$D$가 $B$ 위 유한형이고 $C'$이 $A'$ 위 유한형이라고 가정하자.
보조정리 \ref{more-algebra-lemma-module-over-fibre-product}의 결과를
더 언급하지 않고 사용하겠다.
$B$ 위 $D$의 생성원 $x_1, \ldots, x_r$와
$A'$ 위 $C'$의 생성원 $y_1, \ldots, y_s$를 택한다.
$D = D' \otimes_{B'} B$이고 $B' \to B$가 전사이므로
$D$에서 $x_1, \ldots, x_r$로 가는
$u_1, \ldots, u_r \in D'$을 찾을 수 있다.
$C' = D' \otimes_{B'} A'$이므로 어떤 $a'_{ij} \in A'$에 대하여
$y_i = \sum v_j \otimes a'_{ij}$가 되도록 하는
$v_1, \ldots, v_t \in D'$을 찾을 수 있다.
특히 $C'$에서 $v_j$들의 상은 $A'$-대수로서 $C'$을 생성한다.
$N = r + t$로 두고 다음 환들의 정육면체를 생각하자.
$$
\xymatrix{
A[x_1, \ldots, x_N] & & A'[x_1, \ldots, x_N] \ar[ll] \\
& A \ar[lu] & & A' \ar[ll] \ar[lu] \\
B[x_1, \ldots, x_N] \ar[uu] & & B'[x_1, \ldots, x_N] \ar[uu] \ar[ll] \\
& B \ar[uu] \ar[lu] & & B' \ar[ll] \ar[uu] \ar[lu]
}
$$
뒤쪽 사각형도 데카르트임에 유의하라.
환 사상
$$
B'[x_1, \ldots, x_N] \to D',\quad
x_i \mapsto u_i \quad\text{and}\quad x_{r + j} \mapsto v_j.
$$
을 생각하자. 그러면 유도된 사상
$B[x_1, \ldots, x_N] \to D$와
$A'[x_1, \ldots, x_N] \to C'$은 전사이고, 특히 유한이다.
보조정리 \ref{more-algebra-lemma-relative-finite-module-over-fibre-product}에 의해
$B'[x_1, \ldots, x_N] \to D'$은 유한이다.
따라서 예를 들어 Algebra, Lemma
\ref{algebra-lemma-compose-finite-type}에 의해
$D'$은 $B'$ 위 유한형이다.

\medskip\noindent
(2)에 대하여. 함의 “$\Leftarrow$”는
보조정리 \ref{more-algebra-lemma-relative-flat-module-over-fibre-product}에서 따른다.
또 마지막 명제도 보조정리
\ref{more-algebra-lemma-relative-flat-module-over-fibre-product}의 마지막 명제에서 따른다.

\medskip\noindent
(3)에 대하여. $B \to D$와 $A' \to C'$이 평탄이고 유한 표시라고 가정하자.
$B' \to D'$의 평탄성은 (2)에서 보았다.
(1)에 의해 $B' \to D'$이 유한형임을 안다.
전사 $B'[x_1, \ldots, x_N] \to D'$을 택한다.
Algebra, Lemma \ref{algebra-lemma-finite-presentation-independent}에 의해
환 $D$는 $B[x_1, \ldots, x_N]$-가군으로서 유한 표시이고,
환 $C'$은 $A'[x_1, \ldots, x_N]$-가군으로서 유한 표시이다.
보조정리 \ref{more-algebra-lemma-relative-finitely-presented-module-over-fibre-product}에 의해
$D'$은 $B'[x_1, \ldots, x_N]$-가군으로서 유한 표시이다.
즉 $B' \to D'$은 유한 표시이다.

\medskip\noindent
(4)에 대하여. $B \to D$와 $A' \to C'$이 매끄럽다고 가정하자.
(3)에 의해 $B' \to D'$은 평탄이고 유한 표시이다.
Algebra, Lemma \ref{algebra-lemma-flat-fibre-smooth}에 의해
$B'$ 위의 모든 체 $k$에 대하여 $D' \otimes_{B'} k$가 매끄러움을
확인하면 충분하다. 합성 $J \to B' \to k$가 영이면
$B' \to k$는 $B' \to B \to k$로 분해되고,
$$
D' \otimes_{B'} k = D' \otimes_{B'} B \otimes_B k
= D \otimes_B k
$$
임을 알 수 있다. $B \to D$가 매끄러우므로 이는 매끄럽다.
합성 $J \to B' \to k$가 영이 아니면
$k$에서 영으로 가지 않는 $h \in J$가 존재한다.
그러면 $B' \to k$는 $B' \to B'_h \to k$로 분해된다.
$h$는 $B$에서 영으로 가므로 $B_h = 0$임에 유의하라.
따라서 보조정리 \ref{more-algebra-lemma-diagram-localize}에 의해 $B'_h = A'_h$이고,
$$
D' \otimes_{B'} k = D' \otimes_{B'} B'_h \otimes_{B'_h} k
= C'_h \otimes_{A'_h} k
$$
를 얻는다. $A' \to C'$이 매끄러우므로 이는 매끄럽다.

\medskip\noindent
(5)에 대하여. $B \to D$와 $A' \to C'$이 에탈이라고 가정하자.
(4)에 의해 $B' \to D'$은 매끄럽다.
매끄러운 사상이 에탈인지 여부는 올들의 차원에서 읽을 수 있으므로
(4)의 증명에서처럼 논증하여 올을 동일시하면 (5)가 따른다.
몇몇 세부사항은 생략한다.
\end{proof}

\begin{remark}
\label{more-algebra-remark-relative-modules-over-fibre-product}
상황 \ref{more-algebra-situation-relative-module-over-fibre-product}에서
$B' \to D'$이 유한 표시라고 가정하고 $D'$-가군 $L'$이 주어졌다고 하자.
다음 두 자료 사이에 전단사 대응이 있다고 주장한다.
\begin{enumerate}
\item $Q'$이 $D'$ 위 유한 표시이고 $B'$ 위 평탄인
$D'$-가군의 전사 $L' \to Q'$,
\item 다음을 만족하는 가군의 전사쌍
$(L' \otimes_{D'} D \to Q_1, L' \otimes_{D'} C' \to Q_2)$:
\begin{enumerate}
\item $Q_1$은 $D$ 위 유한 표시이고 $B$ 위 평탄이다.
\item $Q_2$는 $C'$ 위 유한 표시이고 $A'$ 위 평탄이다.
\item $L' \otimes_{D'} C$의 몫으로서
$Q_1 \otimes_D C = Q_2 \otimes_{C'} C$이다.
\end{enumerate}
\end{enumerate}
% P02-E0760: frozen source switches from Q' to undefined Q and omits a copula; preserve the formal switch.
이 대응은 $Q \mapsto (Q_1, Q_2)$로 주어진다.
여기서 $Q_1 = Q \otimes_{D'} D$이고 $Q_2 = Q \otimes_{D'} C'$이다.
반대 방향에는 $Q = Q_1 \times_{Q_{12}} Q_2$를 사용한다.
여기서 $Q_{12}$은 $L' \otimes_{D'} C$의 공통 몫
$Q_1 \otimes_D C = Q_2 \otimes_{C'} C$이다.
몫사상으로는 다음을 사용한다.
$$
L' \longrightarrow
(L' \otimes_{D'} D) \times_{(L' \otimes_{D'} C)} (L' \otimes_{D'} C')
\longrightarrow Q_1 \times_{Q_{12}} Q_2 = Q
$$
첫 번째 화살표는 보조정리
\ref{more-algebra-lemma-module-over-fibre-product-bis}에 의해 전사이고,
두 번째 화살표는 보조정리
\ref{more-algebra-lemma-surjection-module-over-fibre-product}에 의해 전사이다.
주장은 보조정리 \ref{more-algebra-lemma-relative-flat-module-over-fibre-product}와
\ref{more-algebra-lemma-relative-finitely-presented-module-over-fibre-product}에서 따른다.
\end{remark}





\section{Fitting 아이디얼}
\label{more-algebra-section-fitting-ideals}

\noindent
유한 가군의 Fitting 아이디얼은
보조정리 \ref{more-algebra-lemma-fitting-ideal}의 구성으로 정해지는 아이디얼이다.

\begin{lemma}
\label{more-algebra-lemma-ideals-generated-by-minors}
$R$을 환이라 하자. $A$를 $R$에 계수를 갖는 $n \times m$ 행렬이라 하자.
$I_r(A)$를 $A$의 $r \times r$ 소행렬식들이 생성하는 아이디얼이라 하고,
$I_0(A) = R$이며 $r > \min(n, m)$이면 $I_r(A) = 0$이라고 약속하자.
그러면
\begin{enumerate}
\item $I_0(A) \supset I_1(A) \supset I_2(A) \supset \ldots$,
\item $B$가 $(n + n') \times m$ 행렬이고 $A$가 $B$의 처음
$n$개 행이면 $I_{r + n'}(B) \subset I_r(A)$,
\item $C$가 $n \times n$ 행렬이면 $I_r(CA) \subset I_r(A)$.
\item $A$가 블록 행렬
$$
\left(
\begin{matrix}
A_1 & 0 \\
0 & A_2 
\end{matrix}
\right)
$$
이면 $I_r(A) = \sum_{r_1 + r_2 = r} I_{r_1}(A_1) I_{r_2}(A_2)$.
\item 여기에 더 추가할 것.
\end{enumerate}
\end{lemma}

\begin{proof}
생략한다. (힌트: 행렬식은 한 열이나 한 행을 따라 전개하여 계산할 수 있다.)
\end{proof}

\begin{lemma}
\label{more-algebra-lemma-fitting-ideal}
$R$을 환이라 하자. $M$을 유한 $R$-가군이라 하자. $M$의 표시
$$
\bigoplus\nolimits_{j \in J} R \longrightarrow R^{\oplus n}
\longrightarrow M \longrightarrow 0.
$$
를 택하자. $A = (a_{ij})_{i = 1, \ldots, n, j \in J}$를 사상
$\bigoplus_{j \in J} R \to R^{\oplus n}$의 행렬이라 하자.
$A$의 $(n - k) \times (n - k)$ 소행렬식들이 생성하는 아이디얼
$\text{Fit}_k(M)$은 표시의 선택과 무관하다.
\end{lemma}

\begin{proof}
$K \subset R^{\oplus n}$을 전사 $R^{\oplus n} \to M$의 핵이라 하자.
$z_1, \ldots, z_{n - k} \in K$를 택하고
$z_j = (z_{1j}, \ldots, z_{nj})$로 쓰자.
아이디얼 $\text{Fit}_k(M)$은 이와 같이 얻는 모든 행렬
$(z_{ij})$의 $(n - k) \times (n - k)$ 소행렬식들이 생성하는
아이디얼이라고도 기술할 수 있다.

\medskip\noindent
이 전사를 핵이 $K'$인 전사 $R^{\oplus n + n'} \to M$으로 바꾼다고 하자.
여기서 $R^{\oplus n + n'}$의 처음 $n$개 표준 기저 원소에는 원래 사상을,
마지막 $n'$개 기저 벡터에는 $0$을 사용한다. 그러면 대응하는 아이디얼들은 같다.
실제로 위와 같이 $z_1, \ldots, z_{n - k} \in K$이면
$j = 1, \ldots, n - k$에 대하여
$z'_j = (z_{1j}, \ldots, z_{nj}, 0, \ldots, 0) \in K'$라 놓고
$z'_{n + j'} = (0, \ldots, 0, 1, 0, \ldots, 0) \in K'$라 놓자.
그러면 $(z_{ij})$의 $(n - k) \times (n - k)$ 소행렬식들의 아이디얼이
$(z'_{ij})$의 $(n + n' - k) \times (n + n' - k)$ 소행렬식들의
아이디얼과 일치함을 알 수 있다. 이로써 한쪽 포함을 얻는다.
거꾸로 $K'$ 안의 $z'_1, \ldots, z'_{n + n' - k}$가 주어졌다고 하자.
이들을 $R^{\oplus n}$으로 사영하여 $K$ 안의
$z_1, \ldots, z_{n + n' - k}$를 얻을 수 있다.
보조정리 \ref{more-algebra-lemma-ideals-generated-by-minors}에 의해
$(z'_{ij})$의 $(n + n' - k) \times (n + n' - k)$ 소행렬식들이
생성하는 아이디얼은 $(z_{ij})$의 $(n - k) \times (n - k)$
소행렬식들이 생성하는 아이디얼에 포함된다. 이로써 다른 쪽 포함을 얻는다.

\medskip\noindent
$R^{\oplus m} \to M$을 핵이 $L$인 또 하나의 전사라 하자.
Schanuel 보조정리(Algebra, Lemma \ref{algebra-lemma-Schanuel})와
앞 문단의 결과에 의해 $m = n$이고 $M$으로의 전사들과 가환하는 동형
$R^{\oplus n} \to R^{\oplus m}$이 있다고 가정해도 된다.
$C = (c_{li})$를 이 사상의 (가역) 행렬이라 하자
($n = m$이므로 정사각행렬이다). 그러면 위와 같이
$z'_1, \ldots, z'_{n - k} \in L$이 주어졌을 때
$z_1, \ldots, z_{n - k} \in K$로서
$z_1' = Cz_1, \ldots, z'_{n - k} = Cz_{n - k}$인 것들을 찾을 수 있다.
보조정리 \ref{more-algebra-lemma-ideals-generated-by-minors}에 의해 한쪽 포함을 얻고,
대칭성에 의해 다른 쪽 포함을 얻는다.
\end{proof}

\begin{definition}
\label{more-algebra-definition-fitting-ideal}
$R$을 환이라 하자. $M$을 유한 $R$-가군이라 하자. $k \geq 0$라 하자.
$M$의 {\it 제 $k$ Fitting 아이디얼}은
보조정리 \ref{more-algebra-lemma-fitting-ideal}에서 구성한 아이디얼
$\text{Fit}_k(M)$이다. $\text{Fit}_{-1}(M) = 0$으로 놓는다.
\end{definition}

\noindent
Fitting 아이디얼은 큰 행렬의 소행렬식 아이디얼들이고
(보조정리 \ref{more-algebra-lemma-ideals-generated-by-minors}의 순서와 반대 순서로
번호가 붙는다), 어떤 $t \gg 0$에 대하여
$$
0 = \text{Fit}_{-1}(M) \subset \text{Fit}_0(M) \subset \text{Fit}_1(M)
\subset \ldots \subset \text{Fit}_t(M) = R
$$
임을 알 수 있다. 다음은 Fitting 아이디얼의 몇 가지 기본 성질이다.

\begin{lemma}
\label{more-algebra-lemma-fitting-ideal-basics}
$R$을 환이라 하자. $M$을 유한 $R$-가군이라 하자.
\begin{enumerate}
\item $M$이 $n$개의 원소로 생성될 수 있으면 $\text{Fit}_n(M) = R$이다.
\item 두 번째 유한 $R$-가군 $M'$이 주어지면
$\text{Fit}_0(M \oplus M') = \text{Fit}_0(M)\text{Fit}_0(M')$이고,
더 일반적으로
$$
\text{Fit}_l(M \oplus M') =
\sum\nolimits_{k + k' = l} \text{Fit}_k(M)\text{Fit}_{k'}(M')
$$
\item $R \to R'$이 환 사상이면 $\text{Fit}_k(M \otimes_R R')$은
$\text{Fit}_k(M)$의 상이 $R'$에서 생성하는 아이디얼이다.
\item $M$이 유한 표시이면 $\text{Fit}_k(M)$은 유한 생성 아이디얼이다.
\item $M \to M'$이 전사이면
$\text{Fit}_k(M) \subset \text{Fit}_k(M')$이다.
\item $\text{Fit}_0(M) \subset \text{Ann}_R(M)$이다.
\item $V(\text{Fit}_0(M)) = \text{Supp}(M)$이다.
\item $I$가 $R$의 아이디얼이면 $\text{Fit}_0(R/I) = I$이다.
\item 여기에 더 추가할 것.
\end{enumerate}
\end{lemma}

\begin{proof}
(1)은 보조정리 \ref{more-algebra-lemma-ideals-generated-by-minors}에서
$I_0(A) = R$이라는 사실로부터 따른다.

\medskip\noindent
% P02-E0761: frozen source says “follows form”; preserve source and render the intended attribution.
(2)는 보조정리 \ref{more-algebra-lemma-ideals-generated-by-minors}의 대응하는 명제에서 따른다.

\medskip\noindent
(3)은 $\otimes_R R'$가 오른쪽 완전이므로 따른다. 따라서 $M$의 표시를
기저변환한 것은 $M \otimes_R R'$의 표시이다.

\medskip\noindent
(4)의 증명. 표시
$R^{\oplus m} \xrightarrow{A} R^{\oplus n} \to M \to 0$를 택하자.
그러면 $\text{Fit}_k(M)$은 행렬 $A$의
$(n - k) \times (n - k)$ 소행렬식들이 생성하는 아이디얼이다.

\medskip\noindent
(5)는 정의에서 즉시 따른다.

\medskip\noindent
(6)의 증명. 보조정리 \ref{more-algebra-lemma-fitting-ideal}에서와 같이
행렬 $A$를 갖는 $M$의 표시를 택하자.
$J' \subset J$를 크기가 $n$인 부분집합이라 하자.
$f = \det(a_{ij})_{i = 1, \ldots, n, j \in J'}$가 $M$을
영으로 만듦을 보이면 충분하다. 이는 다음 사상의 여핵이
$$
R^{\oplus n} \xrightarrow{A' = (a_{ij})_{i = 1, \ldots, n, j \in J'}}
R^{\oplus n} \to M \to 0
$$
$f$에 의해 영이 되므로 분명하다. 실제로
$A' B = f1_{n \times n}$인 행렬 $B$가 있다.

\medskip\noindent
(7)의 증명. 보조정리 \ref{more-algebra-lemma-fitting-ideal}에서와 같이
행렬 $A$를 갖는 $M$의 표시를 택하자.
Nakayama 보조정리(Algebra, Lemma \ref{algebra-lemma-NAK})에 의해
$$
M_\mathfrak p \not = 0
\Leftrightarrow
M \otimes_R \kappa(\mathfrak p) \not = 0
\Leftrightarrow
\text{rank}(\text{image }A\text{ in }\kappa(\mathfrak p)) < n
$$
이다. $\text{Fit}_0(M)$이 이 성질을 갖는 소아이디얼들의 집합을
정확히 잘라낸다는 것은 분명하다.
\end{proof}

\begin{example}
\label{more-algebra-example-fitting-free}
$R$을 환이라 하자.
유한 자유 가군 $M = R^{\oplus n}$의 Fitting 아이디얼은
$k < n$이면 $\text{Fit}_k(M) = 0$이고 $k \geq n$이면
$\text{Fit}_k(M) = R$이다.
\end{example}

\begin{example}
\label{more-algebra-example-laurent}
$R$을 환이라 하고 $I, J\subset R$을 아이디얼이라 하자.
그러면 $\text{Fit}_0(R/I) = I$, $\text{Fit}_0(R/I \oplus R/J) = IJ$,
$\text{Fit}_1(R/I\oplus R/J) = I+J$이다.
\end{example}

\begin{lemma}
\label{more-algebra-lemma-fitting-ideal-generate-locally}
$R$을 환이라 하자. $M$을 유한 $R$-가군이라 하자. $k \geq 0$라 하자.
$\mathfrak p \subset R$을 소아이디얼이라 하자. 다음은 동치이다.
\begin{enumerate}
\item $\text{Fit}_k(M) \not \subset \mathfrak p$,
\item $\dim_{\kappa(\mathfrak p)} M \otimes_R \kappa(\mathfrak p) \leq k$,
\item $M_\mathfrak p$는 $R_\mathfrak p$ 위에서 $k$개의 원소로 생성될 수 있고,
\item 어떤 $f \in R$, $f \not \in \mathfrak p$에 대하여
$M_f$는 $R_f$ 위에서 $k$개의 원소로 생성될 수 있다.
\end{enumerate}
\end{lemma}

\begin{proof}
Nakayama 보조정리(Algebra, Lemma \ref{algebra-lemma-NAK})에 의해
$M \otimes_R \kappa(\mathfrak p)$가 $k$개의 원소로 생성될 수 있으면
어떤 $f \in R$, $f \not \in \mathfrak p$에 대하여 $M_f$는
$R_f$ 위에서 $k$개의 원소로 생성될 수 있다. 따라서 (2), (3), (4)는
동치이다. 보조정리 \ref{more-algebra-lemma-fitting-ideal-basics}의 (3)을 사용하면
문제는 $R$이 체이고 $\mathfrak p = (0)$인 경우로 귀착된다.
이 경우 결과는 예 \ref{more-algebra-example-fitting-free}에서 따른다.
\end{proof}

\begin{lemma}
\label{more-algebra-lemma-fitting-ideal-finite-locally-free}
$R$을 환이라 하자. $M$을 유한 $R$-가군이라 하자. $r \geq 0$라 하자.
다음은 동치이다.
\begin{enumerate}
\item $M$은 랭크 $r$인 유한 국소 자유 가군이고
(Algebra, Definition \ref{algebra-definition-locally-free}),
\item $\text{Fit}_{r - 1}(M) = 0$이고 $\text{Fit}_r(M) = R$이며,
\item $k < r$이면 $\text{Fit}_k(M) = 0$이고 $k \geq r$이면
$\text{Fit}_k(M) = R$이다.
\end{enumerate}
\end{lemma}

\begin{proof}
Fitting 아이디얼들이 증가 아이디얼열을 이루므로 (2)와 (3)이
동치임은 즉시 알 수 있다.
$\text{Fit}_k(M)$의 형성은 기저변환과 가환하므로
(보조정리 \ref{more-algebra-lemma-fitting-ideal-basics}),
예 \ref{more-algebra-example-fitting-free}와 붙이기 결과
(Algebra, Section \ref{algebra-section-more-glueing})에 의해
(1)은 (2)를 함의한다. 거꾸로 (2)를 가정하자.
보조정리 \ref{more-algebra-lemma-fitting-ideal-generate-locally}에 의해
$M$이 $r$개의 원소로 생성된다고 가정해도 된다.
% P02-E0762: frozen source has a sentence fragment; preserve source and supply the required Korean predicate.
따라서 표시 $\bigoplus_{j \in J} R \to R^{\oplus r} \to M \to 0$가 있다.
이제 $\text{Fit}_{r - 1}(M) = 0$이라는 가정은 사상
$\bigoplus_{j \in J} R \to R^{\oplus r}$의 행렬의 모든 성분이
영임을 함의한다. 따라서 $M$은 자유이다.
\end{proof}

\begin{lemma}
\label{more-algebra-lemma-principal-fitting-ideal}
$R$을 국소환이라 하자. $M$을 유한 $R$-가군이라 하자. $k \geq 0$라 하자.
어떤 $f \in R$에 대하여 $\text{Fit}_k(M) = (f)$라고 가정하자.
$M'$을 $M$의 $\{x \in M \mid fx = 0\}$에 의한 몫이라 하자.
그러면 $M'$은 $k$개의 원소로 생성될 수 있다.
\end{lemma}

\begin{proof}
전사 $R^{\oplus n} \to M$에 대응하는 생성원
$x_1, \ldots, x_n \in M$을 택하자. $R$이 국소환이므로
원소들의 집합 $E \subset (f)$가 $(f)$를 생성하면
어떤 $e \in E$가 $(f)$를 생성한다. Algebra, Lemma
\ref{algebra-lemma-NAK}를 보라. 따라서 $R^{\oplus n} \to M$의
핵에서 $z_1, \ldots, z_{n - k}$를 택하여
$n \times (n - k)$ 행렬 $A = (z_{ij})$의 어떤
$(n - k) \times (n - k)$ 소행렬식이 $(f)$를 생성하게 할 수 있다.
$x_i$의 번호를 다시 붙여 처음 소행렬식
$\det(z_{ij})_{1 \leq i, j \leq n - k}$이 $(f)$를 생성한다고
가정해도 된다. 즉 어떤 단원 $u \in R$에 대하여
$\det(z_{ij})_{1 \leq i, j \leq n - k} = uf$이다.
다른 모든 소행렬식은 $f$의 배수이다.
Algebra, Lemma \ref{algebra-lemma-matrix-right-inverse}에 의해
$n - k \times n - k$ 행렬 $B$가 존재하여
$$
AB = f
\left(
\begin{matrix}
u 1_{n - k \times n - k} \\
C
\end{matrix}
\right)
$$
이다. 여기서 $C$는 $R$에 계수를 갖는 어떤 행렬이다.
이는 모든 $i \leq n - k$에 대하여
$y_i = ux_i + \sum_j c_{ji}x_j$가 $f$에 의해 영이 됨을 함의한다.
$M/\sum Ry_i$가 $x_{n - k + 1}, \ldots, x_n$의 상들로
생성되므로 결론을 얻는다.
\end{proof}

\begin{lemma}
\label{more-algebra-lemma-fitting-ideals-and-pd1}
$R$을 환이라 하자. $M$을 유한 표시 $R$-가군이라 하자. $k \geq 0$라 하자.
어떤 영인자가 아닌 원소 $f \in R$에 대하여
$\text{Fit}_k(M) = (f)$이고 $\text{Fit}_{k - 1}(M) = 0$이라고 가정하자.
그러면
\begin{enumerate}
\item $M$의 사영 차원은 $\leq 1$이고,
\item $M' = \Ker(f : M \to M)$은 $M$의 $f$-멱 꼬임 부분가군이며,
\item $M'$의 사영 차원은 $\leq 1$이고,
\item $M/M'$은 랭크 $k$인 유한 국소 자유 가군이며,
\item $M \cong M/M' \oplus M'$이다.
\end{enumerate}
\end{lemma}

\begin{proof}
$R$에 계수를 갖는 어떤 행렬 $A$에 의한 표시
$$
R^{\oplus m} \xrightarrow{A} R^{\oplus n} \to M \to 0
$$
를 택하자.

\medskip\noindent
먼저 $R$이 국소환일 때 보조정리를 증명한다.
명제에서처럼 $M' = \{x \in M \mid fx = 0\}$으로 놓자.
보조정리 \ref{more-algebra-lemma-principal-fitting-ideal}에 의해
$M/M'$을 생성하는 $x_1, \ldots, x_k \in M$을 택할 수 있다.
그러면 $x_1, \ldots, x_k$는 $M_f = (M/M')_f$를 생성한다.
따라서 $M$ 안에 관계식 $\sum a_ix_i = 0$이 있으면
$a_1, \ldots, a_k$의 $R_f$에서의 상은 영이다. 그렇지 않으면
$\text{Fit}_{k - 1}(M) R_f = \text{Fit}_{k - 1}(M_f)$가
영이 아니기 때문이다.
$f$는 영인자가 아니므로 $a_1 = \ldots = a_k = 0$이라고 결론짓는다.
따라서 $M \cong R^{\oplus k} \oplus M'$이다.
위 표시의 기저를 바꾸어 $R^{\oplus n}$의 처음 $n - k$개
기저 벡터는 $M$의 직합인자 $M'$으로 가고,
% P02-E0763: frozen source says “last k-basis vectors”; preserve source and render the intended noun phrase.
$R^{\oplus n}$의 마지막 $k$개 기저 벡터는 $M$의 직합인자
$R^{\oplus k}$의 기저 원소들로 간다고 가정해도 된다.
그러면 행렬 $A$의 마지막 $k$개 행은 영이다.
이렇게 하여 $M$을 $M'$으로, $k$를 $0$으로,
$n$을 $n - k$로, $A$를 마지막 $k$개 행을 삭제한 부분행렬로
바꾸면 다음 문단에서 논할 경우로 귀착된다.

\medskip\noindent
$R$이 국소환이고 $k = 0$이며 $M$이 $f$에 의해 영이 된다고 가정하자.
이제 $M$의 제 $0$ Fitting 아이디얼은 $(f)$이고, 크기가
$n \times m$인 행렬 $A$의 $n \times n$ 소행렬식들이 생성한다.
(특히 이는 $m \geq n$을 함의한다.)
$R$이 국소환이므로 $A$의 어떤 $n \times n$ 소행렬식은
어떤 단원 $u \in R$에 대하여 $uf$이다.
번호를 다시 붙여 이 소행렬식이 첫 번째 것이라고 가정해도 된다.
또한 $A$의 다른 모든 $n \times n$ 소행렬식들은 $f$로 나누어진다.
$A = (A_1 A_2)$로 블록 형태로 쓰자. 여기서
$A_1$은 $n \times n$ 행렬이고 $A_2$는
$n \times (m - n)$ 행렬이다.
Algebra, Lemma \ref{algebra-lemma-matrix-right-inverse}를
$A$의 전치에 적용하면 (!) 어떤 $n \times n$ 행렬 $B$가 존재하여
$$
BA = B(A_1 A_2) = f
\left(
\begin{matrix}
u 1_{n \times n} & C
\end{matrix}
\right)
$$
이다. 여기서 $C$는 $R$에 계수를 갖는
$n \times (m - n)$ 행렬이다. 먼저 $BA_1 = fu 1_{n \times n}$을 얻는다.
따라서
$$
BA_2 = fC = u^{-1}fuC = u^{-1}BA_1C
$$
이다. $B$의 행렬식이 영인자가 아니므로
$A_2 = u^{-1}A_1C$라고 결론짓는다. 따라서 $A$의 상은
$A_1$의 상과 같고, 후자는 $A_1$의 행렬식이 영인자가 아니므로
$R^{\oplus n}$과 동형이다. 그러므로 $M$의 사영 차원은 $\leq 1$이다.

\medskip\noindent
일반적인 환 $R$의 경우로 돌아가자. 국소적 경우에 의해
$M/M'$은 랭크 $k$인 유한 국소 자유 가군이다.
Algebra, Lemma \ref{algebra-lemma-finite-projective}를 보라.
따라서 완전열 $0 \to M' \to M \to M/M' \to 0$은 분할된다.
그러므로 $M'$은 유한 표시 가군이다. 짧은 완전열
$0 \to K \to R^{\oplus a} \to M' \to 0$을 택하자.
그러면 $K$는 유한 $R$-가군이다. Algebra, Lemma
\ref{algebra-lemma-extension}을 보라. 국소적 경우에 의해
모든 소아이디얼에 대하여 $K_\mathfrak p \cong R_\mathfrak p^{\oplus a}$이다.
따라서 다시 Algebra, Lemma \ref{algebra-lemma-finite-projective}에 의해
$K$는 랭크 $a$인 유한 국소 자유 가군이다.
그러므로 $M'$의 사영 차원은 $\leq 1$이고 보조정리가 증명되었다.
\end{proof}

\section{올림}
\label{more-algebra-section-lifting}

\noindent
% P02-E0764: frozen source says “we collection”; preserve source and render the intended verb.
이 절에서는 다음과 같은 종류의 올림 명제에 관한 몇 가지 보조정리를 모은다.
$A$가 환이고 $I \subset A$가 아이디얼이며
$\overline{\xi}$가 $A/I$ 위의 어떤 구조라면,
$\overline{\xi}$를 $A$ 위 또는 $A$의 어떤 에탈 확대 위의
비슷한 종류의 구조 $\xi$로 올릴 수 있는지를 다룬다.
다음은 이러한 구조 가운데 이미 몇 가지 결과를 증명한 것들이다.
\begin{enumerate}
\item 멱등원. Algebra, Lemmas \ref{algebra-lemma-lift-idempotents}와
\ref{algebra-lemma-lift-idempotents-noncommutative}를 보라.
\item 사영 가군. Algebra, Lemmas \ref{algebra-lemma-lift-projective-module}과
\ref{algebra-lemma-lift-finite-projective-module}을 보라.
\item 유한 안정 자유 가군. 보조정리 \ref{more-algebra-lemma-lift-stably-free}를 보라.
\item 기저 원소. Algebra, Lemmas
\ref{algebra-lemma-local-artinian-basis-when-flat}와
\ref{algebra-lemma-lift-basis}를 보라.
\item 환 사상, 즉 어떤 대수들이 형식적으로 매끄러움을 증명하는 문제.
Algebra, Lemma \ref{algebra-lemma-polynomial-ring-formally-smooth},
Proposition \ref{algebra-proposition-smooth-formally-smooth}, 그리고
Lemma \ref{algebra-lemma-smooth-strong-lift}를 보라.
\item 신토믹 환 사상. Algebra, Lemma \ref{algebra-lemma-lift-syntomic}을 보라.
\item 매끄러운 환 사상. Algebra, Lemma
\ref{algebra-lemma-lift-smooth}를 보라.
\item 에탈 환 사상. Algebra, Lemma \ref{algebra-lemma-lift-etale}을 보라.
\item 다항식의 인수분해. Algebra, Lemma
\ref{algebra-lemma-factor-mod-lift-etale}을 보라.
\item Algebra, Section \ref{algebra-section-henselian}에서는
헨젤 국소환을 논한다.
\end{enumerate}
관심 있는 독자는 Smoothing Ring Maps, Section
\ref{smoothing-section-presentations}, 특히 Smoothing Ring Maps,
Proposition \ref{smoothing-proposition-lift-smooth}에서
이와 같은 결과를 더 찾을 수 있다.

\medskip\noindent
$A$를 환이라 하고 $I \subset A$를 아이디얼이라 하자.
$\overline{\xi}$를 $A/I$ 위의 어떤 구조라 하자.
다음 보조정리들에서는 동형 $A/I \to A'/IA'$을 유도하는
에탈 환 사상 $A \to A'$과 $\overline{\xi}$를 올리는
$A'$ 위의 대상 $\xi'$을 찾는다. 일반적으로 에탈 환 사상
$A \to A' \to A''$이 주어지고
$A/I \cong A'/IA'$이며 $A'/IA' \cong A''/IA''$이면,
합성 $A \to A''$도 에탈이고(Algebra, Lemma
\ref{algebra-lemma-etale}) 또한 $A/I \cong A''/IA''$을 만족한다.
이하의 보조정리들에서는 이를 더 언급하지 않고 자주 사용한다.
먼저 우리가 찾는 종류의 결과에 대한 자명한 예를 들자.

\begin{lemma}
\label{more-algebra-lemma-lift-invertible-element}
$A$를 환이라 하고 $I \subset A$를 아이디얼이라 하며,
$\overline{u} \in A/I$를 가역원이라 하자.
동형 $A/I \to A'/IA'$을 유도하는 에탈 환 사상 $A \to A'$과
$\overline{u}$를 올리는 가역원 $u' \in A'$이 존재한다.
\end{lemma}

\begin{proof}
% P02-E0765: frozen source constructs u although the statement names u'; preserve the formal switch.
$\overline{u}$의 임의의 올림 $f \in A$를 택하고 $A' = A_f$로 놓으며,
$u$를 $A'$에서 $f$의 상이라 하자.
\end{proof}

\begin{lemma}
\label{more-algebra-lemma-lift-idempotent}
$A$를 환이라 하고 $I \subset A$를 아이디얼이라 하며,
$\overline{e} \in A/I$를 멱등원이라 하자.
동형 $A/I \to A'/IA'$을 유도하는 에탈 환 사상 $A \to A'$과
$\overline{e}$를 올리는 멱등원 $e' \in A'$이 존재한다.
\end{lemma}

\begin{proof}
$\overline{e}$의 임의의 올림 $x \in A$를 택하자. 다음과 같이 놓자.
$$
A' = A[t]/(t^2 - t)\left[\frac{1}{t - 1 + x}\right].
$$
$(2t - 1)\text{d}t = 0$이고 $(2t - 1)(2t - 1) = 1$은 가역이므로
환 사상 $A \to A'$은 에탈이다.
$$
A'/IA' = A/I[t]/(t^2 - t)[\frac{1}{t - 1 + \overline{e}}] \cong A/I
$$
이며 마지막 사상은 $t$를 $\overline{e}$로 보낸다.
$\overline{e}$가 $t^2 - t$의 근이므로 이 사상이 잘 정의된다.
이는 $e'$을 $A'$에서 $t$의 류로 놓아도 됨을 보인다.
\end{proof}

\begin{lemma}
\label{more-algebra-lemma-lift-open-covering}
$A$를 환이라 하고 $I \subset A$를 아이디얼이라 하자.
$\Spec(A/I) = \coprod_{j \in J} \overline{U}_j$를 유한 서로소 열린 덮개라 하자.
그러면 동형 $A/I \to A'/IA'$을 유도하는 에탈 환 사상 $A \to A'$과
주어진 덮개를 올리는 유한 서로소 열린 덮개
$\Spec(A') = \coprod_{j \in J} U'_j$가 존재한다.
\end{lemma}

\begin{proof}
보조정리 \ref{more-algebra-lemma-lift-idempotent}와 스펙트럼의 열린닫힌 부분집합들이
멱등원과 대응한다는 사실에서 따른다.
Algebra, Lemma \ref{algebra-lemma-disjoint-decomposition}을 보라.
\end{proof}

\begin{lemma}
\label{more-algebra-lemma-localize-upstairs}
$A \to B$를 환 사상이라 하고 $J \subset B$를 아이디얼이라 하자.
$A \to B$가 $V(J)$의 모든 소아이디얼에서 에탈이면,
$B/J$에서 가역원으로 가는 $g \in B$가 존재하여
$A' = B_g$는 $A$ 위 에탈이다.
\end{lemma}

\begin{proof}
$A \to B$가 에탈이 아닌 $\Spec(B)$의 점들의 집합은
$\Spec(B)$의 닫힌 부분집합이다. Algebra, Definition
\ref{algebra-definition-etale}을 보라.
이를 어떤 아이디얼 $J' \subset B$에 대해 $V(J')$로 쓰자.
그러면 $V(J') \cap V(J) = \emptyset$이므로 Algebra, Lemma
\ref{algebra-lemma-Zariski-topology}에 의해 $J + J' = B$이다.
$f \in J$이고 $g \in J'$인 $1 = f + g$를 쓰자.
그러면 $g$가 원하는 원소이다.
\end{proof}

\noindent
다음 세 보조정리는 다항식의 인수분해를 올릴 수 있음을 말한다.

\begin{lemma}
\label{more-algebra-lemma-lift-factorization-monic}
$A$를 환이라 하고 $I \subset A$를 아이디얼이라 하자.
$f \in A[x]$를 모닉 다항식이라 하자.
$\overline{f} = \overline{g} \overline{h}$를 $A/I[x]$에서 $f$의
인수분해라 하되, $\overline{g}$와 $\overline{h}$는 모닉이고
$A/I[x]$에서 단위 아이디얼을 생성한다고 하자.
그러면 동형 $A/I \to A'/IA'$을 유도하는 에탈 환 사상
$A \to A'$과, 주어진 $A/I$ 위의 인수분해를 올리는
$g'$, $h'$가 모닉인 $A'[x]$에서의 인수분해 $f = g' h'$가 존재한다.
\end{lemma}

\begin{proof}
앞에서 증명한 보편 인수분해에 관한 결과로부터 이를 유도하겠다.
하지만 독자가 이 방법을 사용하지 않는 독자적인 증명을 찾아보기를 권한다.
$\deg(\overline{g}) = n$이고 $\deg(\overline{h}) = m$이라 하자.
그러면 $\deg(f) = n + m$이다. 어떤
$\alpha_1, \ldots, \alpha_{n + m} \in A$에 대하여
$f = x^{n + m} + \sum \alpha_i x^{n + m - i}$로 쓰자.
Algebra, Example \ref{algebra-example-factor-polynomials-etale}의 환 사상
$$
R = \mathbf{Z}[a_1, \ldots, a_{n + m}]
\longrightarrow
S = \mathbf{Z}[b_1, \ldots, b_n, c_1, \ldots, c_m]
$$
을 생각하자. $R \to A$를 $a_i$를 $\alpha_i$로 보내는 환 사상이라 하자.
다음과 같이 놓자.
$$
B = A \otimes_R S
$$
구성에 의해 $B[x]$에서 $f$의 상 $f_B$는 인수분해된다.
이를 $f_B = g_B h_B$라 쓰자. 여기서
$g_B = x^n + \sum (1 \otimes b_i) x^{n - i}$이고 $h_B$도 마찬가지이다.
$\overline{g} = x^n + \sum \overline{\beta}_i x^{n - i}$와
$\overline{h} = x^m + \sum \overline{\gamma}_i x^{m - i}$로 쓰자.
$A$-대수 사상
$$
B \longrightarrow A/I, \quad
1 \otimes b_i \mapsto \overline{\beta}_i, \quad
1 \otimes c_i \mapsto \overline{\gamma}_i
$$
은 $g_B$와 $h_B$를 $A/I[x]$의 $\overline{g}$와 $\overline{h}$로 보낸다.
표시된 사상은 전사이다. 그 핵을 $J \subset B$로 나타내자.
Algebra, Example \ref{algebra-example-factor-polynomials-etale}의 논의에서
% P02-E0766: frozen source writes unaccented “etale”; Korean uses the established 에탈 register.
$A \to B$가 $V(J) \subset \Spec(B)$의 모든 점에서 에탈임은 분명하다.
보조정리 \ref{more-algebra-lemma-localize-upstairs}에서와 같은 $g \in B$를 택하고
$A$-대수 $B_g$를 생각하자. $g$가 $B/J = A/I$에서 단원으로 가므로
$A/I$-대수의 사상 $B_g/I B_g \to A/I$도 얻는다.
$A/I \to B_g/I B_g$가 에탈이므로 $B_g/IB_g \to A/I$도 에탈이다
(Algebra, Lemma \ref{algebra-lemma-map-between-etale}).
따라서 $A/I = (B_g/I B_g)_e$인 멱등원 $e \in B_g/I B_g$가 존재한다
(Algebra, Lemma \ref{algebra-lemma-surjective-flat-finitely-presented}).
$e$의 올림 $h \in B_g$를 택하자. 그러면
$A \to A' = (B_g)_h$와 $A'$에서 $f_B = g_B h_B$의 상으로 주어지는
인수분해가 보조정리의 문제에 대한 해이다.
\end{proof}

\begin{lemma}
\label{more-algebra-lemma-lift-factorization-easy}
$A$를 환이라 하고 $I \subset A$를 아이디얼이라 하자.
$f \in A[x]$를 모닉 다항식이라 하자.
$\overline{f} = \overline{g} \overline{h}$를 $A/I[x]$에서 $f$의
인수분해라 하고 다음을 가정하자.
\begin{enumerate}
\item $\overline{g}$의 최고차항 계수는 $A/I$의 가역원이고,
\item $\overline{g}$, $\overline{h}$는 $A/I[x]$에서 단위 아이디얼을 생성한다.
\end{enumerate}
그러면 동형 $A/I \to A'/IA'$을 유도하는 에탈 환 사상
$A \to A'$과, 주어진 $A/I$ 위의 인수분해를 올리는
$A'[x]$에서의 인수분해 $f = g' h'$가 존재한다.
\end{lemma}

\begin{proof}
보조정리 \ref{more-algebra-lemma-lift-invertible-element}를 적용하여
$\overline{g}$의 최고차항 계수가 어떤 가역원 $u \in A$의
축소라고 가정해도 된다. 그러면 $\overline{g}$를
$\overline{u}^{-1}\overline{g}$로, $\overline{h}$를
$\overline{u}\overline{h}$로 바꿀 수 있다. 따라서
$\overline{g}$가 모닉이라고 가정해도 된다.
$f$가 모닉이므로 $\overline{h}$도 모닉이라고 결론짓는다.
이 경우 결과는 보조정리 \ref{more-algebra-lemma-lift-factorization-monic}에서 따른다.
\end{proof}

\begin{example}
\label{more-algebra-example-cannot-lift-factorization}
이 예는 보조정리 \ref{more-algebra-lemma-lift-factorization-easy}에서 최고차항
계수에 관한 조건을 제거할 수 없음을 보인다. 실제로
$A = \mathbf{Z}$, $I = 4\mathbf{Z}$, $f = 1$이고
$A/I[x]$에서 $\bar{g} = \bar{h} = 2x^2 + 2x + 1$이라 하자.
$A \to A'$이 에탈이고 $A'/I A' = A/I$이면
$A'$의 $2$-진 완비화는 $\mathbf{Z}_2$와 동형이다.
이제 $\mathbf{Z}_2[x]$에서 $4$를 법으로
$2x^2 + 2x + 1$과 합동인 임의의 다항식 $g'$은
$\mathbf{Z}_2[x]$에서 다항식 곱셈 역원을 갖지 않는다.
따라서 보조정리의 결론은 성립하지 않는다.
\end{example}

\begin{lemma}
\label{more-algebra-lemma-separate-image-closed-from-closed}
$R \to S$를 환 사상이라 하자. $I \subset R$을 $R$의 아이디얼이라 하고
$J \subset S$를 $S$의 아이디얼이라 하자.
$V(J)$의 $\Spec(R)$에서의 상의 폐포가 $V(I)$와 서로소이면,
$R/I$에서 $1$로 가고 $S$에서 $J$의 원소로 가는 원소
$f \in R$이 존재한다.
\end{lemma}

\begin{proof}
$V(I')$가 $V(J)$의 상의 폐포가 되는 아이디얼 $I' \subset R$을 택하자.
그러면 가정에 의해 $V(I) \cap V(I') = \emptyset$이고,
Algebra, Lemma \ref{algebra-lemma-Zariski-topology}에 의해
$I + I' = R$이다. $g \in I$이고 $f \in I'$인
$1 = g + f$를 쓰자. $f'$이 $S$에서 $f$의 상이면
$V(f') \supset V(J)$이다. 따라서 어떤 $n$에 대하여
$(f')^n \in J$이다. Algebra, Lemma
\ref{algebra-lemma-Zariski-topology}를 보라.
$f$를 $f^n$으로 바꾸면 결론을 얻는다.
\end{proof}

\begin{lemma}
\label{more-algebra-lemma-helper-integral}
$I$를 환 $A$의 아이디얼이라 하자. $A \to B$를 정수적 환 사상이라 하자.
$b \in B$가 $B/IB$에서 멱등원으로 간다고 하자.
그러면 어떤 $d \geq 1$에 대하여 $f(b) = 0$이고
$f \bmod I = x^d(x - 1)^d$인 모닉 $f \in A[x]$가 존재한다.
\end{lemma}

\begin{proof}
$z = b^2 - b$가 $IB$의 원소임을 관찰하자.
Algebra, Lemma \ref{algebra-lemma-integral-integral-over-ideal}에 의해,
$g(z) = 0$ in $B$이고 $a_j \in I$인 차수 $d$의 모닉 다항식
$g(x) = x^d + \sum a_j x^j$가 존재한다.
따라서 $f(x) = g(x^2 - x) \in A[x]$는
$f(x) \equiv x^d(x - 1)^d \bmod I$이고
$B$에서 $f(b) = 0$인 모닉 다항식이다.
\end{proof}

\begin{lemma}
\label{more-algebra-lemma-lift-idempotent-upstairs}
$A$를 환이라 하고 $I \subset A$를 아이디얼이라 하자.
$A \to B$를 정수적 환 사상이라 하자.
$\overline{e} \in B/IB$를 멱등원이라 하자.
그러면 동형 $A/I \to A'/IA'$을 유도하는 에탈 환 사상
$A \to A'$과 $\overline{e}$를 올리는 멱등원
$e' \in B \otimes_A A'$이 존재한다.
\end{lemma}

\begin{proof}
$\overline{e}$를 올리는 원소 $y \in B$를 택하자.
$y$에 대하여 보조정리 \ref{more-algebra-lemma-helper-integral}에서와 같은
$f \in A[x]$를 택하자.
보조정리 \ref{more-algebra-lemma-lift-factorization-easy}에 의해
동형 $A/I \to A'/IA'$을 유도하는 에탈 환 사상
$A \to A'$을 찾아 다음을 만족하게 할 수 있다.
% P02-E0767: frozen source locates the lifted factorization in A[x], not A'[x]; preserve the exact formal span.
$A[x]$에서 $f = gh$이고, $g(x) = x^d \bmod IA'$이며
$h(x) = (x - 1)^d \bmod IA'$이다.
$A$를 $A'$으로 바꾼 뒤 인수분해가 $A$ 위에서 정의된다고
가정해도 된다. 이 경우
$b_1 = g(y) \in B$는 $\overline{e}^d = \overline{e}$의 올림이고,
$b_2 = h(y) \in B$는
$(\overline{e} - 1)^d = (-1)^d (1 - \overline{e})^d = (-1)^d(1 - \overline{e})$
의 올림이며, 또한 $b_1b_2 = 0$이다.
따라서 $(b_1, b_2)B/IB = B/IB$이고
$V(b_1, b_2) \subset \Spec(B)$는 $V(IB)$와 서로소이다.
$\Spec(B) \to \Spec(A)$가 닫힌 사상이므로(Algebra, Lemmas
\ref{algebra-lemma-integral-going-up}과
\ref{algebra-lemma-going-up-closed}를 보라),
$A/I$에서 가역원으로 가고 $B$에서의 상이 $(b_1, b_2)$에
속하는 $a \in A$를 찾을 수 있다.
보조정리 \ref{more-algebra-lemma-separate-image-closed-from-closed}를 보라.
$A$를 국소화 $A_a$로 바꾼 뒤에는 $(b_1, b_2) = B$이다.
그러면 $\Spec(B) = D(b_1) \amalg D(b_2)$이다.
$b_1b_2 = 0$이므로 서로소 합이고,
$(b_1, b_2) = B$이므로 $\Spec(B)$를 덮는다.
$D(b_1)$이라는 열린닫힌 부분집합에 대응하는 멱등원을
$e \in B$라 하자. Algebra, Lemma
\ref{algebra-lemma-disjoint-decomposition}을 보라.
$b_1$은 $\overline{e}$의 올림이고 $b_2$는
$\pm (1 - \overline{e})$의 올림이므로, Algebra, Lemma
\ref{algebra-lemma-disjoint-decomposition}의 유일성 명제에 의해
$e$가 $\overline{e}$의 올림이라고 결론짓는다.
\end{proof}

\begin{lemma}
\label{more-algebra-lemma-lift-projective-module}
$A$를 환이라 하고 $I \subset A$를 아이디얼이라 하자.
$\overline{P}$를 유한 사영 $A/I$-가군이라 하자.
그러면 동형 $A/I \to A'/IA'$을 유도하는 에탈 환 사상
$A \to A'$과 $\overline{P}$를 올리는 유한 사영
$A'$-가군 $P'$이 존재한다.
\end{lemma}

\begin{proof}
어떤 정수 $n$과 어떤 $R/I$-가군 $\overline{K}$에 대한
% P02-E0768: frozen source says R/I although no R is defined; preserve the exact formal span.
직합 분해 $(A/I)^{\oplus n} = \overline{P} \oplus \overline{K}$를 택할 수 있다.
직합인자 $\overline{P}$에 대응하는 사영자 $\overline{p}$의 올림
$\varphi : A^{\oplus n} \to A^{\oplus n}$을 택하자.
$f \in A[x]$를 $\varphi$의 특성다항식이라 하자.
$B = A[x]/(f)$로 놓자. Cayley--Hamilton에 의해
(Algebra, Lemma \ref{algebra-lemma-charpoly})
$x$를 $\varphi$로 보내는 사상
$B \to \text{End}_A(A^{\oplus n})$이 있다.
모든 소아이디얼 $\mathfrak p \supset I$에 대하여
$\kappa(\mathfrak p)$에서 $f$의 상은 $(x - 1)^rx^{n - r}$이다.
여기서 $r$은 $\overline{P} \otimes_{A/I} \kappa(\mathfrak p)$의 차원이다.
따라서 모든 $\mathfrak p \supset I$에 대하여
$(x - 1)^nx^n$은 $B \otimes_A \kappa(\mathfrak p)$에서 영으로 간다.
그러므로 $x(1 - x)$는 $B/IB$의 모든 소아이디얼에 속한다.
따라서 어떤 $N \geq 1$에 대하여 $x^N(1 - x)^N$은 $IB$에 속한다.
이로부터 $x^N + (1 - x)^N$은 $B/IB$의 단원이고
$$
\overline{e} = \text{image of }\frac{x^N}{x^N + (1 - x)^N}\text{ in }B/IB
$$
는 멱등원이다. 두 명제 모두
$\mathbf{Z}[x]/(x^N(x - 1)^N)$에서 성립하기 때문이다.
$\overline{e}$의 $\text{End}_{A/I}((A/I)^{\oplus n})$에서의 상은
$$
\frac{\overline{p}^N}{\overline{p}^N + (1 - \overline{p})^N} = \overline{p}
$$
이다. 이는 $\overline{p}$가 멱등원이기 때문이다.
보조정리에서와 같이 $A$를 에탈 확대 $A'$으로 바꾼 뒤에는,
보조정리 \ref{more-algebra-lemma-lift-idempotent-upstairs}에 의해
$B/IB$에서 $\overline{e}$로 가는 멱등원 $e \in B$가
존재한다고 가정해도 된다. 그러면 사상
$$
B = A[x]/(f) \longrightarrow \text{End}_A(A^{\oplus n}).
$$
아래에서 $e$의 상은 $\overline{p}$를 올리는 멱등원 $p$이다.
$P = \Im(p)$로 놓으면 결론을 얻는다.
\end{proof}

\begin{lemma}
\label{more-algebra-lemma-cotangent-complex-symmetric-algebra}
% P02-E0769: frozen source calls the required exact presentation only a sequence; preserve that wording boundary.
$A$를 환이라 하자. $0 \to K \to A^{\oplus m} \to M \to 0$을
$A$-가군들의 열이라 하자. 전사 $A^{\oplus m} \to M$에서 오는 표시
$\alpha : A[y_1, \ldots, y_m] \to C$를 갖는 $A$-대수
$C = \text{Sym}^*_A(M)$을 생각하자. 그러면
$$
\NL(\alpha) =
(K \otimes_A C \to \bigoplus\nolimits_{j = 1, \ldots, m} C \text{d}y_j)
$$
이고(Algebra, Section \ref{algebra-section-netherlander}를 보라),
특히 $\Omega_{C/A} = M \otimes_A C$이다.
\end{lemma}

\begin{proof}
$J = \Ker(\alpha)$로 놓자. 보조정리의 주장은
$J/J^2 \cong K \otimes_A C$이다.
$\alpha$가 등급 대수들의 준동형임에 유의하자.
차수 $d$에서 $(J/J^2)_d = K \otimes_A C_{d - 1}$임을 증명하겠다.
다음에 유의하자.
$$
J_d = \Ker(\text{Sym}^d_A(A^{\oplus m}) \to \text{Sym}^d_A(M))
= \Im(K \otimes_A \text{Sym}^{d - 1}_A(A^{\oplus m})
\to \text{Sym}^d_A(A^{\oplus m})),
$$
Algebra, Lemma \ref{algebra-lemma-presentation-sym-exterior}를 보라.
이로부터 $(J^2)_d = \sum_{a + b = d} J_a \cdot J_b$는 다음 사상의 상이다.
% P02-E0770: frozen source uses A^{\otimes m} in four symmetric-power loci; preserve all exact spans.
$$
K \otimes_A K \otimes_A \text{Sym}^{d - 2}_A(A^{\otimes m})
\to \text{Sym}^d_A(A^{\oplus m}).
$$
위에서 인용한 보조정리에 의해 사상
$K \otimes_A \text{Sym}^{d - 2}_A(A^{\otimes m}) \to
\text{Sym}^{d - 1}_A(A^{\oplus m})$의 여핵은
$\text{Sym}^{d - 1}_A(M)$이다.
따라서 $(J/J^2)_d = J_d/(J^2)_d$는 분명히 다음과 같다.
\begin{align*}
\Coker(
K \otimes_A K \otimes_A \text{Sym}^{d - 2}_A(A^{\otimes m})
\to K \otimes_A \text{Sym}^{d - 1}_A(A^{\otimes m}))
& = K \otimes_A \text{Sym}^{d - 1}_A(M) \\
& = K \otimes_A C_{d -1}
\end{align*}
이는 원하는 결과이다.
\end{proof}

\begin{lemma}
\label{more-algebra-lemma-symmetric-algebra-smooth}
$A$를 환이라 하자. $M$을 $A$-가군이라 하자.
그러면 $C = \text{Sym}_A^*(M)$가 $A$ 위 매끄러운 것과
$M$이 유한 사영 $A$-가군인 것은 동치이다.
\end{lemma}

\begin{proof}
$\sigma : C \to A$를 $C$의 차수 $0$ 부분으로의 사영이라 하자.
그러면 $J = \Ker(\sigma)$는 차수 $> 0$인 부분이고,
$A$-가군으로서 $J/J^2 = M$임을 알 수 있다.
따라서 $A \to C$가 매끄러우면 Algebra, Lemma
\ref{algebra-lemma-section-smooth}에 의해 $M$은 유한 사영 $A$-가군이다.

\medskip\noindent
거꾸로 $M$이 유한 사영이라고 가정하고, 핵이 $K$인 전사
$A^{\oplus n} \to M$을 택하자.
$M$이 사영이므로 열 $0 \to K \to A^{\oplus n} \to M \to 0$은
물론 분할된다. 특히 $K$는 유한 $A$-가군이고, 따라서
$C$는 $A$ 위 유한 표시이다. 실제로 $C$는
$K \subset \bigoplus Ax_i$가 생성하는 아이디얼에 의한
$A[x_1, \ldots, x_n]$의 몫이다.
보조정리 \ref{more-algebra-lemma-cotangent-complex-symmetric-algebra}의 계산은
% P02-E0771: frozen source uses (K to M), not the presentation complex; preserve the exact formal span.
$\NL_{C/A}$가 $(K \to M) \otimes_A C$와 호모토피 동치임을 보인다.
따라서 $\NL_{C/A}$는 차수 $0$에 놓인 $C \otimes_A M$과 준동형이고,
이는 Algebra, Definition \ref{algebra-definition-smooth}에 의해
$C$가 $A$ 위 매끄럽다는 뜻이다.
\end{proof}

\begin{lemma}
\label{more-algebra-lemma-lift-section-smooth-morphism}
$A$를 환이라 하고 $I \subset A$를 아이디얼이라 하자.
가환도식
$$
\xymatrix{
B \ar[rd] \\
A \ar[u] \ar[r] & A/I
}
$$
을 생각하자. 여기서 $B$는 매끄러운 $A$-대수이다.
그러면 동형 $A/I \to A'/IA'$을 유도하는 에탈 환 사상
$A \to A'$과 환 사상 $B \to A/I$을 올리는
$A$-대수 사상 $B \to A'$이 존재한다.
\end{lemma}

\begin{proof}
$J \subset B$를 $B \to A/I$의 핵이라 하자. 그러면 $B/J = A/I$이다.
Algebra, Lemma \ref{algebra-lemma-application-NL-smooth}에 의해 열
$$
0 \to I/I^2 \to J/J^2 \to \Omega_{B/A} \otimes_B B/J \to 0
$$
은 분할 완전열이다. 따라서
$\overline{P} = J/(J^2 + IB) = \Omega_{B/A} \otimes_B B/J$는
유한 사영 $A/I$-가군이다. 어떤 정수 $n$과 직합 분해
% P02-E0772: frozen source omits parentheses in A/I^{\oplus n}; preserve the exact span.
$A/I^{\oplus n} = \overline{P} \oplus \overline{K}$를 택하자.
보조정리 \ref{more-algebra-lemma-lift-projective-module}에 의해
동형 $A/I \to A'/IA'$을 유도하는 에탈 환 사상
$A \to A'$과 $\overline{K}$를 올리는 유한 사영
% P02-E0773: frozen source calls the lifted module an A-module before renaming A'; preserve the exact span.
$A$-가군 $K$를 찾을 수 있다. 이제 $A$를 $A'$으로 바꾼다.
$B' = B \otimes_A \text{Sym}_A^*(K)$로 놓자.
보조정리 \ref{more-algebra-lemma-symmetric-algebra-smooth}에 의해
$A \to \text{Sym}_A^*(K)$가 매끄러우므로 $B \to B'$은 매끄럽고,
이는 다시 $A \to B'$이 매끄러움을 함의한다
(Algebra, Lemmas \ref{algebra-lemma-base-change-smooth}와
\ref{algebra-lemma-locally-smooth}를 보라).
또한 절단 $\text{Sym}^*_A(K) \to A$는 절단 $B' \to B$를 정한다.
$B' \to A/I$를 합성 $B' \to B \to A/I$이라 하자.
$J' \subset B'$을 $B' \to A/I$의 핵이라 하자.
$JB' \subset J'$이고 $B \otimes_A K \subset J'$이다.
이 사상들은 합쳐져 동형
$$
(A/I)^{\oplus n} \cong J/(J^2 + IB) \oplus \overline{K}
\longrightarrow
J'/((J')^2 + IB')
$$
을 준다. 따라서 $B$를 $B'$으로 바꾼 뒤에는
$J/(J^2 + IB) = \Omega_{B/A} \otimes_B B/J$가
랭크 $n$인 자유 $A/I$-가군이라고 가정해도 된다.

\medskip\noindent
이 경우 $J/(J^2 + IB)$의 기저로 가는
$f_1, \ldots, f_n \in J$를 택하자.
유한 표시 $A$-대수 $C = B/(f_1, \ldots, f_n)$을 생각하자.
완전열
$$
0 \to H_1(L_{C/A}) \to (f_1, \ldots, f_n)/(f_1, \ldots, f_n)^2
\to \Omega_{B/A} \otimes_B C \to \Omega_{C/A} \to 0
$$
이 있음에 유의하자. Algebra, Lemma
\ref{algebra-lemma-exact-sequence-NL}을 보라
($H_1(L_{B/A}) = 0$이고 $\Omega_{B/A}$가 유한 사영이며,
특히 평탄이므로 Tor 군이 사라짐에 유의하라).
$B$의 임의의 소아이디얼 $\mathfrak q \supset J$에 대하여
$\Omega_{B/A, \mathfrak q}$는 랭크 $n$인 자유 가군이다.
이는 $\Omega_{B/A}$가 유한 사영이고
$\Omega_{B/A} \otimes_B B/J$가 랭크 $n$인 자유 가군이기 때문이다
(Algebra, Lemma \ref{algebra-lemma-finite-projective}를 보라).
$f_1, \ldots, f_n$의 선택에 의해 사상
$$
\left((f_1, \ldots, f_n)/(f_1, \ldots, f_n)^2\right)_{\mathfrak q}
\to
\Omega_{B/A, \mathfrak q}
$$
은 $J$를 법으로 전사이다. 따라서 국소환
$C_{\mathfrak q}$ 위 가군들의 이 사상은 동형이어야 한다
(Algebra, Lemma \ref{algebra-lemma-NAK}에 의해 사상은 전사이고,
예를 들어 Algebra, Lemma \ref{algebra-lemma-fun}에 의해,
$((f_1, \ldots, f_n)/(f_1, \ldots, f_n)^2)_{\mathfrak q}$가
$n$개의 원소로 생성되므로 사상은 단사이다).
따라서 $H_1(L_{C/A})_{\mathfrak q} = 0$이고
$\Omega_{C/A, \mathfrak q} = 0$이다.
Algebra, Lemma \ref{algebra-lemma-smooth-at-point}에 의해
$A \to C$가 $\mathfrak q$에 대응하는 $C$의 소아이디얼
$\overline{\mathfrak q}$에서 매끄러움을 알 수 있다.
$\Omega_{C/A, \mathfrak q} = 0$이므로 실제로
$\overline{\mathfrak q}$에서 에탈이다.
따라서 $A \to C$는 $JC$를 포함하는 $C$의 모든 소아이디얼에서 에탈이다.
보조정리 \ref{more-algebra-lemma-localize-upstairs}에 의해
$C/JC$에서 가역원으로 가고 $A \to C_f$가 에탈이 되게 하는
$f \in C$를 찾을 수 있다. $f$의 선택에 의해 여전히
$C_f/JC_f = A/I$이다. 사상 $C_f/IC_f \to A/I$은 전사이고,
Algebra, Lemma \ref{algebra-lemma-map-between-etale}에 의해 에탈이다.
따라서 $A/I$은 $C_f/IC_f$를 어떤 원소 $g \in C$에서 국소화한 것과 동형이다.
Algebra, Lemma \ref{algebra-lemma-surjective-flat-finitely-presented}를 보라.
$A' = C_{fg}$로 놓으면 증명이 끝난다.
\end{proof}

\section{Zariski 쌍}
\label{more-algebra-section-zariski-pairs}

\noindent
이 절과 다음 절에서 {\it 쌍}은 환 $A$와 아이디얼 $I \subset A$로
이루어진 쌍 $(A, I)$을 뜻한다. {\it 쌍의 사상}
$(A, I) \to (B, J)$은 $\varphi(I) \subset J$를 만족하는
환 사상 $\varphi : A \to B$이다.

\begin{definition}
\label{more-algebra-definition-zariski-pair}
{\it Zariski 쌍}은 $I$가 $A$의 야코브슨 근기에 포함되는 쌍
$(A, I)$이다.
\end{definition}

\begin{lemma}
\label{more-algebra-lemma-idempotents-determined-modulo-radical}
$(A, I)$을 Zariski 쌍이라 하자.
$A$의 멱등원에서 $A/I$의 멱등원으로 가는 사상은 단사이다.
\end{lemma}

\begin{proof}
국소환의 멱등원은 $0$ 또는 $1$이다.
따라서 멱등원은 그것이 사라지는 극대 아이디얼들의 집합으로 정해진다.
Algebra, Lemma \ref{algebra-lemma-characterize-zero-local}을 보라.
\end{proof}

\begin{lemma}
\label{more-algebra-lemma-check-isomorphism-zariski}
$(A, I)$을 Zariski 쌍이라 하자.
$A/I \to B/IB$가 동형이 되게 하는 평탄하고 정수적이며
유한 표시인 환 사상 $A \to B$가 주어졌다고 하자.
그러면 $A \to B$는 동형이다.
\end{lemma}

\begin{proof}
Algebra, Lemma \ref{algebra-lemma-characterize-finite-in-terms-of-integral}에 의해
환 사상 $A \to B$는 유한이다. 따라서 Algebra, Lemma
\ref{algebra-lemma-finite-finitely-presented-extension}에 의해
$B$는 $A$-가군으로서 유한 표시이다.
그러므로 Algebra, Lemma \ref{algebra-lemma-finite-projective}에 의해
$B$는 유한 국소 자유 $A$-가군이다.
가군 $B$는 $V(I)$을 따라 랭크 $1$이고
(Algebra, Lemma \ref{algebra-lemma-finite-projective}에서 설명한
랭크 함수를 보라), $(A, I)$이 Zariski 쌍이므로
랭크가 모든 곳에서 $1$이라고 결론짓는다.
따라서 $A \to B$는 동형이다. 실제로 $S$가 랭크 $1$인
국소 자유 $R$-가군이 되게 하는 환 사상 $R \to S$가 동형임을
보이는 것은 흥미로운 연습문제이다(힌트: 국소환을 보라).
\end{proof}

\begin{lemma}
\label{more-algebra-lemma-helper-finite}
$(A, I)$을 Zariski 쌍이라 하자. $A \to B$를 유한 환 사상이라 하자.
다음을 가정하자.
\begin{enumerate}
% P02-E0774: frozen source omits terminal punctuation in item (1); preserve the list topology.
\item $B/IB = B_1 \times B_2$는 $A/I$-대수들의 곱이고
\item $A/I \to B_1/IB_1$은 전사이며,
\item $b \in B$는 이 곱에서 $(1, 0)$으로 간다.
\end{enumerate}
그러면 어떤 $d \geq 1$에 대하여 $f(b) = 0$이고
$f \bmod I = (x - 1)x^d$인 모닉 $f \in A[x]$가 존재한다.
\end{lemma}

\begin{proof}
보조정리 \ref{more-algebra-lemma-lift-idempotent-upstairs}에 의해
동형 $A/I \to A'/IA'$을 유도하는 에탈 환 사상
$A \to A'$을 찾아, $B' = B \otimes_A A'$이
$B'/IB'$에서 $b$의 상을 올리는 멱등원 $e'$을 포함하게 할 수 있다.
이에 대응하는 $A'$-대수 분해
$$
B' = B'_1 \times B'_2
$$
를 생각하자. 이는 $I$을 법으로 축소했을 때
보조정리에 주어진 분해와 양립한다.
사상 $A' \to B'_1$은 $IA'$를 법으로 전사이다.
Nakayama 보조정리(Algebra, Lemma \ref{algebra-lemma-NAK})에 의해
어떤 $i \in IA'$를 찾아 $A'$을 $A'_{1 + i}$로 바꾼 뒤
사상 $A' \to B'_1$이 전사가 되게 할 수 있다.
$b$의 상 $b'_1 \in B'_1$은 $b'_1 - 1 \in IB'_1$을 만족한다.
따라서 $b'_1 - 1$로 가는 $a' \in IA'$를 택할 수 있다.
한편 $b$의 상 $b'_2 \in B'_2$는 $IB'_2$에 속한다.
Algebra, Lemma \ref{algebra-lemma-integral-integral-over-ideal}에 의해
$a'_j \in IA'$인 차수 $d$의 모닉 다항식
$g(x) = x^d + \sum a'_j x^j$가 존재하여
$B'_2$에서 $g(b'_2) = 0$이다.
따라서 $b$의 상 $b' = (b'_1, b'_2) \in B'$은
다항식 $(x - 1 - a')g(x)$의 근이다. 그러므로
$$
(b' - 1)(b')^d \in \sum\nolimits_{j = 0, \ldots, d} IA' \cdot (b')^j
$$
이다. 이것이 $B$에서
$$
(b - 1)b^d \in \sum\nolimits_{j = 0, \ldots, d} I \cdot b^j
$$
를 함의한다고 주장한다. 이를 위해서는 환 사상
$A \to A'$이 충실 평탄임을 보면 충분하다.
실제로 조건은 $B/\sum_{j = 0, \ldots, d} Ib^j$에서
$(b - 1)b^d$의 상이 영이라는 것이다
(Algebra, Lemma \ref{algebra-lemma-faithfully-flat-universally-injective}를 사용하라).
% P02-E0775: frozen source omits the verb in “The map A to A' flat”; Korean supplies it.
사상 $A \to A'$은 에탈이므로 평탄이다
(Algebra, Lemma \ref{algebra-lemma-etale}).
한편 유도된 스펙트럼 사이의 사상은 열린 사상이고
(Algebra, Proposition \ref{algebra-proposition-fppf-open}과
앞에서 인용한 보조정리를 사용하라), 그 상은 $V(I)$를 포함한다.
$I$가 $A$의 야코브슨 근기에 포함되므로 결론을 얻는다.
\end{proof}

\begin{lemma}
\label{more-algebra-lemma-noetherian-zariski-jacobson-complement}
$(A, I)$을 Zariski 쌍이라 하고 $A$가 뇌터라고 하자.
$f \in I$라 하자. 그러면 $A_f$는 야코브슨 환이다.
\end{lemma}

\begin{proof}
Algebra, Lemma \ref{algebra-lemma-noetherian-dim-1-Jacobson}의
판정법을 사용하겠다. $\mathfrak p_f = \mathfrak p A_f$가
소아이디얼이지만 극대가 아니게 하는 소아이디얼
$\mathfrak p \subset A$를 택하자.
$A_f/\mathfrak p_f = (A/\mathfrak p)_f$가 무한히 많은
소아이디얼을 가짐을 보여야 한다.
$A$를 $A/\mathfrak p$로 바꾸어 $A$가 정역이고
$\dim A_f > 0$이라고 가정해도 되며,
목표는 $\Spec(A_f)$가 무한임을 보이는 것이다.
$\dim A_f > 0$이므로 $f$를 포함하지 않는 영이 아닌 소아이디얼
$\mathfrak q \subset A$를 찾을 수 있다.
$\mathfrak q$를 포함하는 극대 아이디얼 $\mathfrak m \subset A$를 택하자.
$(A, I)$이 Zariski 쌍이므로 $I \subset \mathfrak m$이다.
따라서 $\mathfrak m \not = \mathfrak q$이고
$\dim(A_\mathfrak m) > 1$이다. 그러므로 Algebra, Lemma
\ref{algebra-lemma-Noetherian-local-domain-dim-2-infinite-opens}에 의해
$\Spec((A_\mathfrak m)_f) \subset \Spec(A_f)$는 무한이고 결론을 얻는다.
\end{proof}

\section{헨젤 쌍}
\label{more-algebra-section-henselian-pairs}

\noindent
절 \ref{more-algebra-section-lifting}의 결과 가운데 일부는 헨젤 쌍에 관한
결과로 볼 수 있다. 이 절에서 {\it 쌍}은 환 $A$와 아이디얼
$I \subset A$로 이루어진 쌍 $(A, I)$을 뜻한다.
{\it 쌍의 사상} $(A, I) \to (B, J)$은
$\varphi(I) \subset J$를 만족하는 환 사상 $\varphi : A \to B$이다.
절 \ref{more-algebra-section-lifting}에서와 같이, $A$ 위의 대상 $\xi$가 주어지면
$\overline{\xi}$로 $\xi$를 $A/I$ 위의 대상으로 ‘기저변환’한 것을
나타낸다(이것이 의미를 갖는 경우에 한한다).

\begin{definition}
\label{more-algebra-definition-henselian-pair}
{\it 헨젤 쌍}은 다음을 만족하는 쌍 $(A, I)$이다.
\begin{enumerate}
\item $I$는 $A$의 야코브슨 근기에 포함되고,
\item 임의의 모닉 다항식 $f \in A[T]$와, 모닉이며
$A/I[T]$에서 단위 아이디얼을 생성하는 $g_0, h_0 \in A/I[T]$에 의한
인수분해 $\overline{f} = g_0h_0$에 대하여,
$g, h$가 모닉이고 $g_0 = \overline{g}$ 및
$h_0 = \overline{h}$인 $A[T]$에서의 인수분해 $f = gh$가 존재한다.
\end{enumerate}
\end{definition}

\noindent
$A$가 국소환이고 $I = \mathfrak m$이 그 극대 아이디얼이면,
$(A, I)$이 헨젤 쌍인 것과 $A$가 헨젤 국소환인 것은 동치이다.
Algebra, Lemma \ref{algebra-lemma-characterize-henselian}을 보라.
보조정리 \ref{more-algebra-lemma-characterize-henselian-pair}에서
헨젤 쌍의 여러 동치인 특성화를 주며
(시간이 지나면서 더 추가할 것이다).

\begin{lemma}
\label{more-algebra-lemma-locally-nilpotent-henselian}
$(A, I)$을 $I$가 국소 멱영인 쌍이라 하자.
그러면 함자 $B \mapsto B/IB$는 $A$ 위 에탈 대수들의 범주와
$A/I$ 위 에탈 대수들의 범주 사이의 동치를 유도한다.
또한 이 쌍은 헨젤이다.
\end{lemma}

\begin{proof}
본질적 전사성은 Algebra, Lemma \ref{algebra-lemma-lift-etale}에서 따른다.
$B$, $B'$이 $A$ 위 에탈이고
$B/IB \to B'/IB'$이 $A/I$-대수들의 사상이면,
Algebra, Lemma \ref{algebra-lemma-smooth-strong-lift}에 의해
이를 올릴 수 있다. 마지막으로 $f \bmod I = g \bmod I$인 두
$A$-대수 사상 $f, g : B \to B'$이 주어졌다고 하자.
곱셈 사상 $B \otimes_A B \to B$의 핵을 생성하는 멱등원
$e \in B \otimes_A B$를 택하자. Algebra, Lemmas
\ref{algebra-lemma-diagonal-unramified}와
\ref{algebra-lemma-unramified}을 보라
(에탈 사상이 비분기임을 사용한다).
그러면 $(f \otimes g)(e) \in IB'$이다.
$IB'$은 국소 멱영이므로(Algebra, Lemma
\ref{algebra-lemma-locally-nilpotent}),
Algebra, Lemma \ref{algebra-lemma-lift-idempotents}에 의해
이는 $(f \otimes g)(e) = 0$을 함의한다. 따라서 $f = g$이다.

\medskip\noindent
$I$가 $A$의 야코브슨 근기에 포함된다는 것은 분명하다.
$f \in A[T]$를 모닉 다항식이라 하고,
모닉이며 $A/I[T]$에서 단위 아이디얼을 생성하는
$g_0, h_0 \in A/I[T]$에 의한
$\overline{f} = f \bmod I$의 인수분해
$\overline{f} = g_0h_0$가 주어졌다고 하자.
보조정리 \ref{more-algebra-lemma-lift-factorization-monic}에 의해
동형 $A/I \to A'/IA'$을 유도하는 에탈 환 사상
$A \to A'$이 존재하여, 인수분해가 $A'$ 위 모닉 다항식들의
인수분해로 올라간다. 위 결과에 의해 $A = A'$이므로
인수분해는 $A$ 위에 있다.
\end{proof}

\begin{lemma}
\label{more-algebra-lemma-limit-henselian}
$A = \lim A_n$이고 $(A_n)$이 전이사상들이 전사이며
그 핵들이 국소 멱영인 환들의 역계라고 하자.
그러면 $I_n = \Ker(A \to A_n)$으로 놓을 때 $(A, I_n)$은 헨젤 쌍이다.
\end{lemma}

\begin{proof}
$n$을 고정하자. $A_n$에서 $1$로 가는 원소 $a \in A$를 택하자.
Algebra, Lemma \ref{algebra-lemma-locally-nilpotent-unit}에 의해
모든 $m \geq n$에 대하여 $a$는 $A_m$에서 단원으로 간다.
따라서 $a$는 $A$의 단원이다. 그러므로 Algebra, Lemma
\ref{algebra-lemma-contained-in-radical}에 의해
아이디얼 $I_n$은 $A$의 야코브슨 근기에 포함된다.
$f \in A[T]$를 모닉 다항식이라 하고,
모닉이며 $A_n[T]$에서 단위 아이디얼을 생성하는
$g_n, h_n \in A_n[T]$에 의한
$\overline{f} = f \bmod I_n$의 인수분해
$\overline{f} = g_nh_n$이 주어졌다고 하자.
보조정리 \ref{more-algebra-lemma-locally-nilpotent-henselian}에 의해
이를 모든 $m \geq n$에 대하여 $A_m[T]$의 모닉
$g_m, h_m$에 의한 $f \bmod I_m = g_m h_m$으로 차례로 올릴 수 있다.
각 단계에서는 올림 $g_m, h_m$이
% P02-E0776: frozen source says A_n[T] at the m-th stage; preserve the exact formal span.
$A_n[T]$에서 단위 아이디얼을 생성하는지 확인해야 한다.
이는 $g_n, h_n$에 대한 대응하는 사실과
$A_m \to A_n$의 핵이 국소 멱영이어서
$\Spec(A_n[T]) = \Spec(A_m[T])$이라는 사실에서 따른다.
$A = \lim A_m$이므로 증명이 끝난다.
\end{proof}

\begin{lemma}
\label{more-algebra-lemma-complete-henselian}
$(A, I)$을 쌍이라 하자. $A$가 $I$-진 완비이면 이 쌍은 헨젤이다.
\end{lemma}

\begin{proof}
Algebra, Lemma \ref{algebra-lemma-radical-completion}에 의해
아이디얼 $I$는 $A$의 야코브슨 근기에 포함된다.
$f \in A[T]$를 모닉 다항식이라 하고,
모닉이며 $A/I[T]$에서 단위 아이디얼을 생성하는
$g_0, h_0 \in A/I[T]$에 의한
$\overline{f} = f \bmod I$의 인수분해
$\overline{f} = g_0h_0$가 주어졌다고 하자.
보조정리 \ref{more-algebra-lemma-locally-nilpotent-henselian}에 의해
이를 모든 $n \geq 1$에 대하여 $A/I^n[T]$의 모닉
$g_n, h_n$에 의한 $f \bmod I^n = g_n h_n$으로 차례로 올릴 수 있다.
$A = \lim A/I^n$이므로 증명이 끝난다.
\end{proof}

\begin{lemma}
\label{more-algebra-lemma-helper-finite-type}
$(A, I)$을 쌍이라 하자. $A \to B$를 유한형 환 사상이라 하고,
$B/IB = C_1 \times C_2$이며 $A/I \to C_1$이 유한이라고 하자.
$B'$을 $B$ 안에서 $A$의 정수적 폐포라 하자.
그러면 곱 분해를 보존하는 사상 $B'/IB' \to B/IB$과
$B'/IB' = C_1 \times C'_2$로 쓸 수 있으며,
$C_1 \times C'_2$에서 $(1, 0)$으로 가는 $g \in B'$로서
$B'_g \to B_g$가 동형인 것이 존재한다.
\end{lemma}

\begin{proof}
닫힌 부분집합 $T = \Spec(C_1) \subset \Spec(B)$의 모든
소아이디얼에서 $A \to B$가 준유한임을 관찰하자
(올환을 보면 알 수 있다. Algebra, Definition
\ref{algebra-definition-quasi-finite}을 보라).
위상공간들의 도식
$$
\xymatrix{
\Spec(B) \ar[rr]_\phi \ar[rd]_\psi & & \Spec(B') \ar[ld]^{\psi'} \\
& \Spec(A)
}
$$
을 생각하자. Algebra, Theorem \ref{algebra-theorem-main-theorem}에 의해
모든 $\mathfrak p \in T$에 대하여
$h_\mathfrak p \not \in \mathfrak p$인 $h_\mathfrak p \in B'$가 존재하여
% P02-E0777: frozen source localizes at undefined h rather than h_p; preserve the exact formal span.
$B'_h \to B_h$가 동형이다.
합집합 $U = \bigcup D(h_\mathfrak p)$는 열린집합
$U \subset \Spec(B')$을 주고, $\phi^{-1}(U) \to U$는 위상동형이며
$T \subset \phi^{-1}(U)$이다. $T$가
$\psi^{-1}(V(I))$에서 열려 있으므로
$\phi(T)$는 $U \cap (\psi')^{-1}(V(I))$에서 열려 있다.
따라서 $\phi(T)$는 $(\psi')^{-1}(V(I))$에서 열려 있다.
한편 $C_1$은 $A/I$ 위 유한이므로 $B'$ 위 유한이다.
그러므로 Algebra, Lemmas \ref{algebra-lemma-going-up-closed}와
\ref{algebra-lemma-integral-going-up}에 의해
$\phi(T)$는 $\Spec(B')$의 닫힌 부분집합이다.
따라서 $\Spec(B'/IB') \supset \phi(T)$는 열린닫힌 부분집합이다.
Algebra, Lemma \ref{algebra-lemma-disjoint-implies-product}에 의해
이에 대응하는 곱 분해 $B'/IB' = C'_1 \times C'_2$를 얻는다.
$B'/IB' \to B/IB$는 $C'_1$을 $C_1$ 안으로,
$C'_2$를 $C_2$ 안으로 보낸다.
이는 스펙트럼에서 일어나는 일을 보면 알 수 있다
(힌트: $\phi(T)$의 역상은 정확히 $T$이다. 몇몇 세부사항은 생략한다).
$C'_1 \times C'_2$에서 $(1, 0)$으로 가고
$D(g) \subset U$인 $g \in B'$를 택하자.
$\Spec(C'_1)$과 $\Spec(C'_2)$가 $\Spec(B')$에서 서로소인
닫힌집합이고 $\Spec(C'_1)$이 $U$에 포함되므로 이는 가능하다.
그러면 $B'_g \to B_g$는 스펙트럼 위 위상동형과
국소환 위 동형을 정한다(위에서 $U$를 그렇게 택했기 때문이다).
따라서 예를 들어 Algebra, Lemma
\ref{algebra-lemma-characterize-zero-local}에 의해 이는 동형이다.
마지막으로 $C'_1 = C_1$이 따르고 증명이 끝난다.
\end{proof}

\begin{lemma}
\label{more-algebra-lemma-characterize-henselian-pair}
\begin{reference}
\cite[Chapter XI]{Henselian} and \cite[Proposition 1]{Gabber-henselian}
\end{reference}
$(A, I)$을 쌍이라 하자. 다음은 동치이다.
\begin{enumerate}
\item $(A, I)$은 헨젤 쌍이다.
\item 에탈 환 사상 $A \to A'$과 $A$-대수 사상
$\sigma : A' \to A/I$이 주어지면,
$\sigma$를 올리는 $A$-대수 사상 $A' \to A$이 존재한다.
\item 임의의 유한 $A$-대수 $B$에 대하여 사상
$B \to B/IB$는 멱등원들 위에서 전단사를 유도한다.
\item 임의의 정수적 $A$-대수 $B$에 대하여 사상
$B \to B/IB$는 멱등원들 위에서 전단사를 유도하며,
\item (Gabber) $I$는 $A$의 야코브슨 근기에 포함되고,
형태가
$$
f(T) = T^n(T - 1) + a_n T^n + \ldots + a_1 T + a_0
$$
이며 $a_n, \ldots, a_0 \in I$이고 $n \ge 1$인 모든
모닉 다항식 $f(T) \in A[T]$는 근 $\alpha \in 1 + I$를 갖는다.
\end{enumerate}
또한 (5)에서 이 근은 유일하다.
\end{lemma}

\begin{proof}
(2)가 성립한다고 가정하자. 그러면 $I$는 $A$의 야코브슨 근기에 포함된다.
그렇지 않으면 $I$를 법으로 $1$과 합동인 어떤 비단원
$f \in A$가 존재하고, 사상 $A \to A_f$가 (2)에 모순된다.
따라서 $A$ 위 정수적인 $B$에 대하여 $IB \subset B$는
$B$의 야코브슨 근기에 포함된다. 실제로 Algebra, Lemmas
\ref{algebra-lemma-going-up-closed}와
\ref{algebra-lemma-integral-going-up}에 의해
$\Spec(B) \to \Spec(A)$는 닫힌 사상이다.
그러므로 보조정리 \ref{more-algebra-lemma-idempotents-determined-modulo-radical}에 의해
$B$의 멱등원에서 $B/IB$의 멱등원으로 가는 사상은 단사이다.
한편 (2)가 성립하므로 보조정리 \ref{more-algebra-lemma-lift-idempotent-upstairs}에 의해
$B/IB$의 모든 멱등원은 $B$의 멱등원으로 올라간다.
이로써 (2)가 (4)를 함의함을 알 수 있다.

\medskip\noindent
함의 (4) $\Rightarrow$ (3)은 자명하다.

\medskip\noindent
(3)을 가정하자. $\mathfrak m$을 극대 아이디얼이라 하고
유한 사상 $A \to B = A/(I \cap \mathfrak m)$를 생각하자.
$B \to B/IB$가 멱등원들 위에서 전단사를 유도한다는 조건은
$I \subset \mathfrak m$을 함의한다
(그렇지 않으면 $B = A/I \times A/\mathfrak m$이고 $B/IB = A/I$이다).
따라서 $I$가 $A$의 야코브슨 근기에 포함됨을 알 수 있다.
$f \in A[T]$를 모닉이라 하고, 모닉이며 $A/I[T]$에서
단위 아이디얼을 생성하는 $g_0, h_0 \in A/I[T]$에 의한
인수분해 $\overline{f} = g_0h_0$가 주어졌다고 하자.
$B = A[T]/(f)$로 놓자. $B/IB$의 분해
$$
B/IB = A/I[T]/(g_0) \times A/I[T]/(h_0)
$$
에 대응하는 멱등원을 $\overline{e}$라 하자.
이는 $A$-대수들의 분해이다. (3)을 가정했으므로 존재하는
$\overline{e}$의 올림 멱등원 $e \in B$를 택하자.
이는 곱 분해
$$
B = eB \times (1 - e)B
$$
를 준다. $B$는 $A$-가군으로서 랭크 $\deg(f)$인 자유 가군임에 유의하자.
따라서 $eB$와 $(1 - e)B$는 유한 국소 자유 $A$-가군이다.
그런데 $eB$와 $(1 - e)B$는 $A/I$ 위에서 각각 상수 랭크
$\deg(g_0)$와 $\deg(h_0)$를 가지므로,
$\Spec(A)$ 위에서도 마찬가지임을 알 수 있다. 따라서
\begin{align*}
f & = \text{CharPol}_A(T : B \to B) \\
& =
\text{CharPol}_A(T : eB \to eB)
\text{CharPol}_A(T : (1 - e)B \to (1 - e)B)
\end{align*}
이며, 이는 $I$을 법으로 주어진 인수분해로 축소되는
모닉 다항식들의 인수분해이다.
여기서 $\text{CharPol}_A$는 $A$ 위 유한 국소 자유 가군의
자기준동형의 특성다항식을 나타낸다. 가군이 자유이면
$\text{CharPol}_A$는 대응하는 행렬의 특성다항식으로 정의하고,
일반적으로는 Algebra, Lemma \ref{algebra-lemma-standard-covering}을
사용해 붙인다. 세부사항은 생략한다. 따라서 (3)은 (1)을 함의한다.

\medskip\noindent
(1)을 가정하자. $f$를 (5)에서와 같다고 하자.
$f \bmod I$의 $T^n$과 $T - 1$로의 인수분해는
정의 \ref{more-algebra-definition-henselian-pair}에 의해 $g$와 $h$가 모닉인
인수분해 $f = gh$로 올라간다. 그러면 $h$의 차수는 $1$이어야 하고,
$f$가 $I$을 법으로 $1$로 축소되는 근을 가짐을 알 수 있다.
마지막으로 헨젤 쌍의 정의에 의해 $I$는 야코브슨 근기에 포함된다.
따라서 (1)은 (5)를 함의한다.

\medskip\noindent
마지막 단계의 증명을 주기 전에, (5)의 근 $\alpha$가 존재하면
유일함을 보이자. 실제로 $f(T)$의 명시적 형태로 인해
$f'(\alpha) \in 1 + I$이다. 여기서 $f'$은 $T$에 대한 $f$의 도함수이다.
초등적인 논증은
$$
f(T) = f(\alpha + T - \alpha) =
f(\alpha) + f'(\alpha) \cdot (T - \alpha) \bmod
(T - \alpha)^2 A[T]
$$
을 보인다. 이는 $f(T)$의 다른 임의의 근
$\alpha' \in 1 + I$가 어떤 $i \in I$에 대하여
$0 = f(\alpha') - f(\alpha) = (\alpha' - \alpha)(1 + i)$를
만족함을 보인다. $1 + i$는 $A$의 단원이므로
$\alpha = \alpha'$이다.

\medskip\noindent
(5)를 가정하자. (2)가 성립함을 보이겠다. 다시 말해,
모든 에탈 사상 $A \to A'$에 대하여 $I$을 법으로 한 모든 절단
$\sigma : A' \to A/I$은 절단 $A' \to A$으로 올라간다.
$A \to A'$이 에탈이므로 절단 $\sigma$는 $A/I$-대수의 분해
\begin{equation}
\label{more-algebra-equation-GCHP}
A'/IA' \cong A/I \times C
\end{equation}
을 정한다. 실제로 전사 환 사상 $A'/IA' \to A/I$은
Algebra, Lemma \ref{algebra-lemma-map-between-etale}에 의해 에탈이고,
그 뒤 Algebra, Lemma
\ref{algebra-lemma-surjective-flat-finitely-presented}에 의해
원하는 멱등원을 얻는다. 이 분해가 $A'_1$이 $A$ 위 정수적인
$A$-대수의 분해
\begin{equation}
\label{more-algebra-equation-GCHP-want}
A' \cong A'_1 \times A'_2
\end{equation}
로 올라감을 보이겠다. 그러면 $A \to A'_1$은 정수적이고 에탈이며
$A/I \to A'_1/IA'_1$은 동형이다. 따라서 보조정리
\ref{more-algebra-lemma-check-isomorphism-zariski}에 의해 $A \to A'_1$은 동형이다
(여기서는 에탈 환 사상이 평탄이고 유한 표시라는 사실도 사용한다.
Algebra, Lemma \ref{algebra-lemma-etale}을 보라).

\medskip\noindent
$B'$을 $A'$ 안에서 $A$의 정수적 폐포라 하자.
보조정리 \ref{more-algebra-lemma-helper-finite-type}에 의해,
(\ref{more-algebra-equation-GCHP})와 양립하는 $A/I$-대수들의 분해
\begin{equation}
\label{more-algebra-equation-dec-mod-I}
B'/IB' \cong A/I \times C'
\end{equation}
를 택하고, $(1, 0)$을 올리는 $b \in B'$로서
$B'_b \to A'_b$가 동형인 것을 찾을 수 있다.
분해 (\ref{more-algebra-equation-dec-mod-I})가 $A$-대수들의 분해
\begin{equation}
\label{more-algebra-equation-want-2}
B' \cong B'_1 \times B'_2
\end{equation}
로 올라가면, 유도된 분해 $A' = A'_1 \times A'_2$는
원하는 (\ref{more-algebra-equation-GCHP-want})를 준다.
실제로 $b$는 $B'_1$의 단원이므로(세부사항은 생략한다)
$B'_1 \cong A'_1$이고, 따라서 $A'_1$은 $A$ 위 정수적이다.

\medskip\noindent
$b$를 포함하는 유한 $A$-부분대수 $B'' \subset B'$을 택하자
($B'$의 모든 유한 생성 $A$-부분대수는 $A$ 위 유한임을 관찰하라).
$B''$을 확대하여 $b$가 $B''/IB''$의 멱등원으로 가고 분해
\begin{equation}
\label{more-algebra-equation-again-dec-mod-I}
B''/IB'' \cong C''_1 \times C''_2
\end{equation}
를 만든다고 가정해도 된다.
$B'_b \cong A'_b$이므로 $B'_b$는 $A$ 위 유한형임을 알 수 있다.
$B'_b$가 $A$ 위에서 $b_1/b^n, \ldots, b_t/b^n$로 생성된다고 하고,
$B''$을 확대하여 $b_1, \ldots, b_t \in B''$라 하자.
그러면 $B''_b \to B'_b$는 전사인 동시에 단사이고, 따라서 동형이다.
특히 $C''_1 = A/I$임을 알 수 있다!
따라서 $A/I \to C''_1$은 동형이고 특히 전사이다.
보조정리 \ref{more-algebra-lemma-helper-finite}에 의해 형태가
$$
f(T) = T^n(T - 1) + a_n T^n + \ldots + a_1 T + a_0
$$
이고 $a_n, \ldots, a_0 \in I$이며 $n \ge 1$이고
$f(b) = 0$인 $f(T) \in A[T]$를 찾을 수 있다.
% P02-E0778: frozen source has the wrong English article before A[T]/(f); Korean has no article.
특히 $B'$은 $A[T]/(f)$-대수임을 알 수 있다.
(5)에 의해 $f$의 근 $a \in 1 + I$가 존재한다.
이는 $B'/IB'$의 분해 (\ref{more-algebra-equation-dec-mod-I})와 양립하는 곱 분해
$A[T]/(f) = A[T]/(T - a) \times D$를 만든다.
이로부터 유도되는 $B'$의 분해가 원하는 (\ref{more-algebra-equation-want-2})이다.
\end{proof}

\begin{lemma}
\label{more-algebra-lemma-change-ideal-henselian-pair}
$A$를 환이라 하자. $V(I) = V(J)$인 아이디얼
$I, J \subset A$를 택하자. 그러면 $(A, I)$이 헨젤인 것과
$(A, J)$가 헨젤인 것은 동치이다.
\end{lemma}

\begin{proof}
임의의 정수적 환 사상 $A \to B$에 대하여 $V(IB) = V(JB)$이다.
따라서 $B/IB$와 $B/JB$의 멱등원들은 전단사로 대응한다
(Algebra, Lemma \ref{algebra-lemma-disjoint-decomposition}).
그러므로 $B \to B/IB$가 멱등원들의 집합 위에서 전단사를
유도하는 것과 $B \to B/JB$가 멱등원들의 집합 위에서
전단사를 유도하는 것은 동치이다.
보조정리 \ref{more-algebra-lemma-characterize-henselian-pair}에 의해 결론을 얻는다.
\end{proof}

\begin{lemma}
\label{more-algebra-lemma-integral-over-henselian-pair}
$(A, I)$을 헨젤 쌍이라 하고 $A \to B$를 정수적 환 사상이라 하자.
그러면 $(B, IB)$는 헨젤 쌍이다.
\end{lemma}

\begin{proof}
보조정리 \ref{more-algebra-lemma-characterize-henselian-pair}의 네 번째 특성화와
정수적 환 사상들의 합성이 정수적이라는 사실에서 즉시 따른다.
\end{proof}

\begin{lemma}
\label{more-algebra-lemma-henselian-henselian-pair}
$I \subset J \subset A$를 환 $A$의 아이디얼들이라 하자.
다음은 동치이다.
\begin{enumerate}
\item $(A, I)$과 $(A/I, J/I)$은 헨젤 쌍이고,
% P02-E0779: frozen source says “an henselian pair”; Korean has no article.
\item $(A, J)$는 헨젤 쌍이다.
\end{enumerate}
\end{lemma}

\begin{proof}
(1)을 가정하자. $B$를 정수적 $A$-대수라 하자.
환 사상들
$$
B \to B/IB \to B/JB
$$
을 생각하자. 보조정리 \ref{more-algebra-lemma-characterize-henselian-pair}에 의해
두 화살표 모두 멱등원들 위에서 전단사를 유도한다.
따라서 합성도 그러하다. 그러므로 보조정리
\ref{more-algebra-lemma-characterize-henselian-pair}에 의해 $(A, J)$는 헨젤 쌍이다.

\medskip\noindent
거꾸로 (2)가 성립한다고 가정하자.
보조정리 \ref{more-algebra-lemma-integral-over-henselian-pair}에 의해
$(A/I, J/I)$은 헨젤 쌍이다.
$B$를 정수적 $A$-대수라 하고 환 사상들
$$
B \to B/IB \to B/JB
$$
을 생각하자. 보조정리 \ref{more-algebra-lemma-characterize-henselian-pair}에 의해
합성과 두 번째 화살표는 멱등원들 위에서 전단사를 유도한다.
따라서 첫 번째 화살표도 그러하다.
그러므로 (다시 같은 보조정리에 의해) $(A, I)$은 헨젤 쌍이다.
\end{proof}

\begin{lemma}
\label{more-algebra-lemma-sum-henselian}
$A$를 환이라 하고 $(A, I)$과 $(A, I')$을 헨젤 쌍이라 하자.
% P02-E0780: frozen source says “an henselian pair”; Korean has no article.
그러면 $(A, I + I')$은 헨젤 쌍이다.
\end{lemma}

\begin{proof}
보조정리 \ref{more-algebra-lemma-integral-over-henselian-pair}에 의해
쌍 $(A/I, (I' + I)/I)$은 헨젤이다.
따라서 보조정리 \ref{more-algebra-lemma-henselian-henselian-pair}에서 결론을 얻는다.
\end{proof}

\begin{lemma}
\label{more-algebra-lemma-product-henselian-pairs}
$J$를 집합이라 하고 $\{ (A_j, I_j)\}_{j \in J}$를 쌍들의 모임이라 하자.
그러면 $(\prod_{j \in J} A_j, \prod_{j\in J} I_j)$가 헨젤인 것과
각 $(A_j, I_j)$가 헨젤인 것은 동치이다.
\end{lemma}
 
\begin{proof}
모든 $j \in J$에 대하여 사영
$\prod_{j \in J} A_j \rightarrow A_j$은 정수적 환 사상이다.
따라서 $(\prod_{j \in J} A_j, \prod_{j\in J} I_j)$가 헨젤이면,
보조정리 \ref{more-algebra-lemma-integral-over-henselian-pair}에 의해
각 $(A_j, I_j)$가 헨젤이다.

\medskip\noindent
거꾸로 각 $(A_j, I_j)$가 헨젤 쌍이라고 가정하자.
그러면 $x \in \prod_{j \in J} I_j$인 모든 $1 + x$는
$\prod_{j \in J} A_j$에서 단원이다.
실제로 Algebra, Lemma \ref{algebra-lemma-contained-in-radical}과
정의 \ref{more-algebra-definition-henselian-pair}에 의해 성분별로 그러하다.
따라서 다시 Algebra, Lemma
\ref{algebra-lemma-contained-in-radical}에 의해
$\prod_{j \in J} I_j$는 $\prod_{j \in J} A_j$의
야코브슨 근기에 포함된다. 계속 성분별로 논증하면,
모든 모닉 $f \in (\prod_{j \in J} A_j)[T]$와,
모닉이며
$g_0, h_0 \in (\prod_{j \in J} A_j / \prod_{j \in J} I_j)[T] =
(\prod_{j \in J} A_j/I_j)[T]$이고
$(\prod_{j \in J} A_j / \prod_{j \in J} I_j)[T]$에서
단위 아이디얼을 생성하는 모든 인수분해
$\overline{f} = g_0h_0$에 대하여,
$g$, $h$가 모닉이고 $g_0$, $h_0$로 축소되는
$(\prod_{j \in J} A_j)[T]$에서의 인수분해 $f = gh$가 존재한다.
따라서 정의 \ref{more-algebra-definition-henselian-pair}에 의해
$(\prod_{j \in J} A_j, \prod_{j\in J} I_j)$는 헨젤 쌍이다.
\end{proof}

\begin{lemma}
\label{more-algebra-lemma-limits-henselian}
헨젤이라는 성질은 쌍들의 극한에서 보존된다.
더 정확히, $J$를 전순서집합이라 하고 $(A_j, I_j)$를
$J$ 위의 헨젤 쌍들의 역계라 하자.
그러면 아이디얼 $I = \lim I_j$를 갖춘
$A = \lim A_j$는 헨젤 쌍 $(A, I)$이다.
\end{lemma}

\begin{proof}
Categories, Lemma \ref{categories-lemma-limits-products-equalizers}에 의해
곱과 등화자만 생각하면 된다.
곱의 경우 주장은 보조정리 \ref{more-algebra-lemma-product-henselian-pairs}에서 따른다.
따라서 등화자 도식
$$
\xymatrix{
(A, I) \ar[r] & (A', I')
\ar@<1ex>[r]^{\varphi} \ar@<-1ex>[r]_{\psi}
&
(A'', I'')  
}
$$
을 생각하자. 여기서 쌍 $(A', I')$과 $(A'', I'')$은 헨젤이다.
쌍 $(A, I)$도 헨젤임을 확인하기 위해 보조정리
\ref{more-algebra-lemma-characterize-henselian-pair}의 Gabber 판정법을 사용하겠다.
$1 + I$의 모든 원소는 $A$의 단원이다.
단원의 역원의 유일성 때문에 이를 $(A', I')$에서 확인할 수 있기 때문이다.
따라서 $I$는 $A$의 야코브슨 근기에 포함된다.
Algebra, Lemma \ref{algebra-lemma-contained-in-radical}을 보라.
이제
% P02-E0781: frozen source permits N >= 1 for exponent N-1 although Gabber's criterion uses exponent at least 1; preserve both spans.
$$
f(T) = T^{N - 1}(T - 1) + a_{N - 1} T^{N - 1} + \dotsb + a_1 T + a_0
$$
를 $a_{N - 1}, \dotsc, a_0 \in I$이고 $N \ge 1$인
$A[T]$의 다항식이라 하자. $A'[T]$에서 $f(T)$의 상은
유일한 근 $\alpha' \in 1 + I'$를 가지며,
$A''[T]$에서의 더 나아간 상도 마찬가지이다.
따라서 유일성에 의해 $\varphi(\alpha') = \psi(\alpha')$이고,
결국 $\alpha'$은 원하는 대로 $1 + I$에서 $f(T)$의 근을 정한다.
\end{proof}

\begin{lemma}
\label{more-algebra-lemma-filtered-colimits-henselian}
헨젤이라는 성질은 쌍들의 여과 쌍대극한에서 보존된다.
더 정확히, $J$를 유향집합이라 하고 $(A_j, I_j)$를
$J$ 위의 헨젤 쌍들의 계라 하자.
그러면 아이디얼 $I = \colim I_j$를 갖춘
$A = \colim A_j$는 헨젤 쌍 $(A, I)$이다.
\end{lemma}

\begin{proof}
$u \in 1 + I$이면 어떤 $j \in J$에 대하여
$u$는 어떤 $u_j \in 1 + I_j$의 상이다.
Algebra, Lemma \ref{algebra-lemma-contained-in-radical}과
$I_j$가 $A_j$의 야코브슨 근기에 포함된다는 가정에 의해
$u_j$는 $A_j$에서 가역이다. 따라서 $u$는 $A$에서 가역이다.
그러므로 (같은 보조정리에 의해) $I$는 $A$의 야코브슨 근기에 포함된다.

\medskip\noindent
$f \in A[T]$를 모닉 다항식이라 하고,
모닉이며 $A/I[T]$에서 단위 아이디얼을 생성하는
$g_0, h_0 \in A/I[T]$에 의한 인수분해
$\overline{f} = g_0 h_0$가 주어졌다고 하자.
어떤 $g'_0, h'_0 \in A/I[T]$에 대하여
$1 = g_0 g'_0 + h_0 h'_0$로 쓰자.
$A = \colim A_j$와 $A/I = \colim A_j/I_j$가 여과 쌍대극한이므로,
어떤 $j \in J$와
% P02-E0782: frozen source places polynomial f_j in A_j rather than A_j[T]; preserve the exact formal span.
$f_j \in A_j$ 및 인수분해
$\overline{f}_j = g_{j, 0} h_{j, 0}$을 찾을 수 있다.
여기서 $g_{j, 0}, h_{j, 0} \in A_j/I_j[T]$는 모닉이고,
어떤 $g'_{j, 0}, h'_{j, 0} \in A_j/I_j[T]$에 대하여
$1 = g_{j, 0} g'_{j, 0} + h_{j, 0} h'_{j, 0}$이며,
$f_j, g_{j, 0}, h_{j, 0}, g'_{j, 0}, h'_{j, 0}$는 각각
$f, g_0, h_0, g'_0, h'_0$로 간다.
$(A_j, I_j)$가 헨젤 쌍이므로
$\overline{f}_j = g_{j, 0} h_{j, 0}$을 $A_j$ 위의
인수분해로 올릴 수 있다. 그 상을 $A$에서 취하여
$A$의 대응하는 인수분해를 얻는다. 따라서 $(A, I)$은 헨젤이다.
\end{proof}

\begin{example}[Moret-Bailly]
\label{more-algebra-example-moret-bailly}
쌍대극한이 여과되지 않으면 보조정리
\ref{more-algebra-lemma-filtered-colimits-henselian}은 틀린다.
예를 들어 헨젤 쌍 $(\mathbf{Z}_p, (p))$와
$(\mathbf{Z}_p, (p))$의 쌍대곱을 취하면,
$A = \mathbf{Z}_p \otimes_\mathbf{Z} \mathbf{Z}_p$인 $(A, pA)$를 얻는다.
이는 헨젤 쌍이 아니다. $A/pA = \mathbf{F}_p$이므로,
% P02-E0784: frozen source says “where henselian” instead of “were henselian”; Korean supplies the conditional.
$(A, pA)$가 헨젤이면 $A$는 국소환이어야 한다.
그러나 $\Spec(A)$는 비연결이다. 예를 들어 홀수 소수 $p$에 대하여
자명하지 않은 멱등원
$$
(1/2 \otimes 1)
\left(1 \otimes 1 - (1 + p)^{-1}u \otimes u\right)
$$
이 있다. 여기서 $u \in \mathbf{Z}_p$는 $1 + p$의 제곱근이다.
몇몇 세부사항은 생략한다.
\end{example}

\begin{lemma}
\label{more-algebra-lemma-largest-ideal-henselian}
$A$를 환이라 하자. $(A, I)$이 헨젤 쌍이 되게 하는
가장 큰 아이디얼 $I \subset A$가 존재한다.
\end{lemma}

\begin{proof}
보조정리들 \ref{more-algebra-lemma-henselian-henselian-pair},
\ref{more-algebra-lemma-sum-henselian}, 그리고
\ref{more-algebra-lemma-filtered-colimits-henselian}을 결합하라.
\end{proof}

\begin{lemma}
\label{more-algebra-lemma-irreducible-henselian-pair-connected}
$(A, I)$을 헨젤 쌍이라 하자. $\mathfrak p \subset A$를
소아이디얼이라 하자. 그러면 $V(\mathfrak p + I)$는 연결이다.
\end{lemma}

\begin{proof}
보조정리 \ref{more-algebra-lemma-integral-over-henselian-pair}에 의해
% P02-E0783: frozen source omits parentheses around the image ideal; preserve the exact span.
$(A/\mathfrak p, I + \mathfrak p/\mathfrak p)$는 헨젤 쌍이다.
따라서 다음을 증명하면 충분하다. $(A, I)$이 헨젤 쌍이고
$A$가 정역이면 $\Spec(A/I) = V(I)$는 연결이다.
그렇지 않으면 Algebra, Lemma
\ref{algebra-lemma-characterize-spec-connected}에 의해
$A/I$는 자명하지 않은 멱등원을 갖는다.
보조정리 \ref{more-algebra-lemma-characterize-henselian-pair}에 의해
이는 $A$가 자명하지 않은 멱등원을 가짐을 함의한다. 모순이다.
\end{proof}

\section{쌍의 헨젤화}
\label{more-algebra-section-henselization}

\noindent
절 \ref{more-algebra-section-henselian-pairs}에서 시작한 논의를 계속한다.

\begin{lemma}
\label{more-algebra-lemma-henselization}
포함 함자
$$
\text{category of henselian pairs}
\longrightarrow
\text{category of pairs}
$$
는 왼쪽 수반 $(A, I) \mapsto (A^h, I^h)$을 갖는다.
\end{lemma}

\begin{proof}
$(A, I)$을 쌍이라 하자. $A/I \to B/IB$가 동형이 되게 하는
에탈 환 사상 $A \to B$들로 이루어진 범주 $\mathcal{C}$를 생각하자.
범주 $\mathcal{C}$가 유향이고, 아이디얼 $I^h = IA^h$를 갖는
$A^h = \colim_{B \in \mathcal{C}} B$가 원하는 수반을 줌을 보이겠다.

\medskip\noindent
먼저 $\mathcal{C}$가 유향임을 증명하자
(Categories, Definition \ref{categories-definition-directed}).
$\text{id} : A \to A$가 대상이므로 공집합이 아니다.
$B$와 $B'$이 $\mathcal{C}$의 두 대상이면
$B'' = B \otimes_A B'$도 $\mathcal{C}$의 대상이고
(Algebra, Lemma \ref{algebra-lemma-etale}을 사용하라),
사상 $B \to B''$과 $B' \to B''$이 있다.
$f, g : B \to B'$을 $\mathcal{C}$의 대상 사이의 두 사상이라 하자.
그러면 쌍대등화자는
$$
(B' \otimes_{f, B, g} B') \otimes_{(B' \otimes_A B')} B'
$$
이고, Algebra, Lemmas \ref{algebra-lemma-etale}과
\ref{algebra-lemma-map-between-etale}에 의해 $A$ 위 에탈이다.
따라서 범주 $\mathcal{C}$는 유향이다.

\medskip\noindent
$\mathcal{C}$의 모든 대상 $B$에 대하여 $B/IB = A/I$이므로
$A^h/I^h = A^h/IA^h = \colim B/IB = \colim A/I = A/I$임을 알 수 있다.

\medskip\noindent
다음으로 $I^h = IA^h$를 갖는
$A^h = \colim_{B \in \mathcal{C}} B$가 헨젤 쌍임을 보이자.
이를 위해 보조정리 \ref{more-algebra-lemma-characterize-henselian-pair}의
조건 (2)를 확인하겠다. 에탈 환 사상 $A^h \to A'$과
$A^h$-대수 사상 $\sigma : A' \to A^h/I^h$이 주어졌다고 하자.
그러면 어떤 $B \in \mathcal{C}$와 에탈 환 사상 $B \to B'$이 존재하여
$A' = B' \otimes_B A^h$이다. Algebra, Lemma
\ref{algebra-lemma-etale}을 보라.
$A^h/I^h = A/I$이므로 사상 $\sigma$는
$A$-대수 사상 $s : B' \to A/I$을 유도한다.
그러면 $A/I$-대수로서 $B'/IB' = A/I \times C$이다.
% P02-E0785: frozen source calls the unital factor C the kernel; preserve that description.
여기서 $C$는 $s$가 유도하는 사상 $B'/IB' \to A/I$의 핵이다.
$(1, 0) \in A/I \times C$로 가는
$g \in B'$를 택하자. 그러면 $B \to B'_g$는 에탈이고
$A/I \to B'_g/IB'_g$는 동형이다. 즉 $B'_g$는 $\mathcal{C}$의 대상이다.
따라서 다음 두 도식이 가환하게 하는 표준 사상 $B'_g \to A^h$를 얻는다.
$$
\vcenter{
\xymatrix{
B'_g \ar[r] & A^h \\
B \ar[u] \ar[ur]
}
}
\quad\text{and}\quad
\vcenter{
\xymatrix{
B' \ar[r] \ar[rrd]_s & B'_g \ar[r] & A^h \ar[d] \\
& & A/I
}
}
$$
이는 원하는 대로 $\sigma$와 양립하는 사상
$A' = B' \otimes_B A^h \to A^h$을 유도한다.

\medskip\noindent
$(A, I) \to (A', I')$을 $(A', I')$이 헨젤인 쌍들의 사상이라 하자.
증명을 끝내기 위해 유일한 인수분해
$A \to A^h \to A'$이 있음을 보이겠다.
실제로 $\mathcal{C}$ 안의 각 $A \to B$에 대하여
환 사상 $A' \to B' = A' \otimes_A B$는 에탈이고
동형 $A'/I' \to B'/I'B'$을 유도한다.
따라서 보조정리 \ref{more-algebra-lemma-characterize-henselian-pair}에 의해
절단 $\sigma_B : B' \to A'$이 있다.
$\mathcal{C}$ 안의 사상 $B_1 \to B_2$가 주어지면 도식
$$
\xymatrix{
B'_1 \ar[rr] \ar[rd]_{\sigma_{B_1}} & &
B'_2 \ar[ld]^{\sigma_{B_2}} \\
& A'
}
$$
이 가환한다고 주장한다.
이는 $\mathcal{C}$의 모든 $B$에 대하여 절단 $\sigma_B$가
유일한 $A'$-대수 사상 $B' \to A'$임을 증명하면 따른다.
Algebra, Lemma \ref{algebra-lemma-diagonal-unramified}에 의해
어떤 환 $R$에 대하여 $B' \otimes_{A'} B' = B' \times R$이다.
지금의 경우 $B'/I'B' = A'/I'$이므로 $R/I'R = 0$이다.
따라서 두 $A'$-대수 사상 $\sigma_B, \sigma_B' : B' \to A'$이
주어지면 $e = (\sigma_B \otimes \sigma_B')(0, 1) \in A'$은
$I'$에 포함되는 멱등원이다.
보조정리 \ref{more-algebra-lemma-idempotents-determined-modulo-radical}에 의해
$e = 0$이라고 결론짓는다. 따라서 원하는 대로
$\sigma_B = \sigma_B'$이다. 이 가환성을 사용하여
$$
A^h = \colim_{B \in \mathcal{C}} B \to
\colim_{B \in \mathcal{C}} A' \otimes_A B \xrightarrow{\colim \sigma_B} A'
$$
을 얻는다. 사상 $\sigma_B$들의 유일성은 이 사상도 유일함을 보장한다.
따라서 $(A, I) \mapsto (A^h, I^h)$은 원하는 수반이다.
\end{proof}

\begin{lemma}
\label{more-algebra-lemma-henselization-flat}
$(A, I)$을 쌍이라 하고, $(A^h, I^h)$를
보조정리 \ref{more-algebra-lemma-henselization}에서와 같이 두자.
그러면 $A \to A^h$는 평탄이고, $I^h = IA^h$이며,
모든 $n$에 대하여 $A/I^n \to A^h/I^nA^h$는 동형이다.
\end{lemma}

\begin{proof}
보조정리 \ref{more-algebra-lemma-henselization}의 증명에서
$A^h$가 $A/I \to B/IB$가 동형이 되게 하는 에탈 $A$-대수 $B$들의
여과 쌍대극한이고 $I^h = IA^h$임을 보았다.
에탈 환 사상은 평탄이므로(Algebra, Lemma
\ref{algebra-lemma-etale}), Algebra, Lemma
\ref{algebra-lemma-colimit-flat}에 의해 $A \to A^h$가 평탄이라고 결론짓는다.
각 $A \to B$가 평탄이므로 사상 $A/I^n \to B/I^nB$도 동형이다
(예를 들어 Algebra, Lemma \ref{algebra-lemma-lift-basis}에 의해).
쌍대극한을 취하면 원하는 대로 $A/I^n = A^h/I^nA^h$을 얻는다.
\end{proof}

\begin{lemma}
\label{more-algebra-lemma-henselization-local-ring}
\begin{slogan}
% P02-E0786: frozen source omits “of” in its English slogan; Korean supplies the relation.
쌍의 헨젤화와 국소환의 헨젤화 사이의 양립성.
\end{slogan}
보조정리 \ref{more-algebra-lemma-henselization}의 함자는 국소환
$(A, \mathfrak m)$에 그 헨젤화를 대응시킨다.
\end{lemma}

\begin{proof}
$(A^h, \mathfrak m^h)$를 보조정리 \ref{more-algebra-lemma-henselization}에서 구성한
쌍 $(A, \mathfrak m)$의 헨젤화라 하자.
보조정리 \ref{more-algebra-lemma-henselization-flat}에 의해
$\mathfrak m^h = \mathfrak m A^h$는 극대 아이디얼이고,
야코브슨 근기에 포함되므로 $A^h$가 극대 아이디얼
$\mathfrak m^h$를 갖는 국소환이라고 결론짓는다.
이제 두 가지 방법으로 증명을 끝낼 수 있다.

\medskip\noindent
첫 번째 증명: Algebra, Lemma \ref{algebra-lemma-henselization}의 증명에서
쌍대극한으로 주어지는 구성은 보조정리
\ref{more-algebra-lemma-henselization}에서 $A^h$를 구성하는 데 사용한
쌍대극한과 같음을 관찰하자.
두 번째 증명: 헨젤화 $A \to S$와 보조정리
\ref{more-algebra-lemma-henselization}의 $A \to A^h$는 모두 국소환 준동형이고,
$S$와 $A^h$는 모두 에탈 $A$-대수들의 여과 쌍대극한이며,
$S$와 $A^h$는 모두 헨젤 국소환이고,
$S$와 $A^h$는 모두 $\kappa(\mathfrak m)$와 같은 잉여체를 갖는다
(두 번째 경우에는 보조정리 \ref{more-algebra-lemma-henselization-flat}에 의해).
따라서 Algebra, Lemma \ref{algebra-lemma-uniqueness-henselian}에 의해
표준적으로 동형이다.
\end{proof}

\begin{lemma}
\label{more-algebra-lemma-henselization-Noetherian-pair}
\begin{slogan}
뇌터 쌍의 헨젤화는 같은 완비화를 갖는 뇌터 쌍이다.
\end{slogan}
$(A, I)$을 $A$가 뇌터인 쌍이라 하고,
$(A^h, I^h)$를 보조정리 \ref{more-algebra-lemma-henselization}에서와 같이 두자.
그러면 $I$-진 완비화들의 사상
$$
A^\wedge \to (A^h)^\wedge
$$
은 동형이다. 또한 $A^h$는 뇌터이고, 사상
$A \to A^h \to A^\wedge$은 평탄이며,
$A^h \to A^\wedge$은 충실 평탄이다.
\end{lemma}

\begin{proof}
첫 번째 명제는 보조정리 \ref{more-algebra-lemma-henselization-flat}의 즉각적인
결과이며, 사실 $A$가 뇌터라는 가정 없이도 성립한다.
보조정리 \ref{more-algebra-lemma-henselization}의 증명에서
$A^h$는 $A/I \to B/IB$가 동형이 되게 하는 에탈
$A$-대수 $B$들의 여과 쌍대극한임을 보았다.
각각의 그러한 $A \to B$에 대하여 유도된 사상
$A^\wedge \to B^\wedge$은 동형이다
(보조정리 \ref{more-algebra-lemma-henselization-flat}의 증명을 보라).
Algebra, Lemma \ref{algebra-lemma-completion-flat}에 의해
각 $B$에 대하여 환 사상
$B \to A^\wedge = B^\wedge = (A^h)^\wedge$은 평탄하다.
따라서 Algebra, Lemma \ref{algebra-lemma-colimit-rings-flat}에 의해
$A^h \to A^\wedge = (A^h)^\wedge$은 평탄하다.
$I^h = IA^h$가 $A^h$의 야코브슨 근기에 포함되고
$A^h \to A^\wedge$이 동형 $A^h/I^h \to A/I$을 유도하므로,
Algebra, Lemma \ref{algebra-lemma-ff}에 의해
$A^h \to A^\wedge$은 충실 평탄이다.
Algebra, Lemma
\ref{algebra-lemma-completion-Noetherian-Noetherian}에 의해
환 $A^\wedge$은 뇌터이다. 따라서 Algebra, Lemma
\ref{algebra-lemma-descent-Noetherian}에 의해
$A^h$가 뇌터라고 결론짓는다.
\end{proof}

\begin{lemma}
\label{more-algebra-lemma-henselization-colimit}
$(A, I) = \colim (A_i, I_i)$를 쌍들의 여과 쌍대극한이라 하자.
보조정리 \ref{more-algebra-lemma-henselization}의 함자는
$A^h = \colim A_i^h$와 $I^h = \colim I_i^h$를 준다.
\end{lemma}

\noindent
이 보조정리는 여과되지 않은 쌍대극한에는 거짓이다.
예 \ref{more-algebra-example-moret-bailly}를 보라.

\begin{proof}
Categories, Lemma \ref{categories-lemma-adjoint-exact}에 의해
$(A^h, I^h)$은 헨젤 쌍들의 범주에서
계 $(A_i^h, I_i^h)$의 쌍대극한이다.
따라서 헨젤 쌍 $(B, J)$에 대하여
$$
\Mor((A^h, I^h), (B, J)) =
\lim \Mor((A_i^h, I_i^h), (B, J)) =
\Mor(\colim (A_i^h, I_i^h), (B, J))
$$
이다. 여기서 쌍대극한은 쌍들의 범주에서 취한다.
쌍대극한이 여과되었으므로 쌍들의 범주에서
$\colim (A_i^h, I_i^h) = (\colim A_i^h, \colim I_i^h)$를 얻는다.
세부사항은 생략한다.
% P02-E0787: frozen source omits “that” in “using the colimit is filtered”; Korean supplies the relation.
다시 쌍대극한이 여과되었다는 사실을 사용하면,
이는 헨젤 쌍이다(보조정리
\ref{more-algebra-lemma-filtered-colimits-henselian}).
따라서 Yoneda 보조정리에 의해
$(A^h, I^h) = (\colim A_i^h, \colim I_i^h)$를 얻는다.
\end{proof}

\begin{lemma}
\label{more-algebra-lemma-henselization-change-ideal}
\begin{slogan}
쌍의 헨젤화는 아이디얼의 근기에만 의존한다.
\end{slogan}
$A$를 아이디얼 $I$와 $J$를 갖는 환이라 하자.
$V(I) = V(J)$이면 보조정리 \ref{more-algebra-lemma-henselization}의 함자는
쌍 $(A, I)$과 쌍 $(A, J)$에 같은 환을 만든다.
\end{lemma}

\begin{proof}
$(A', IA')$을 쌍 $(A, I)$에서 시작하여
보조정리 \ref{more-algebra-lemma-henselization}으로 만든 쌍이라 하자.
보조정리 \ref{more-algebra-lemma-henselization-flat}을 보라.
$(A'', JA'')$을 쌍 $(A, J)$에서 시작하여
보조정리 \ref{more-algebra-lemma-henselization}으로 만든 쌍이라 하자.
보조정리 \ref{more-algebra-lemma-change-ideal-henselian-pair}에 의해
$(A', JA')$는 헨젤 쌍이고 $(A'', IA'')$도 헨젤 쌍이다.
구성의 보편 성질에 의해 유일한 $A$-대수 사상
$A'' \to A'$과 $A' \to A''$을 얻는다.
유일성은 이들이 서로 역임을 보인다.
\end{proof}

\begin{lemma}
\label{more-algebra-lemma-henselization-integral}
\begin{slogan}
헨젤화는 정수적 기저변환과 가환한다.
\end{slogan}
$(A, I) \to (B, J)$를 $V(J) = V(IB)$인 쌍들의 사상이라 하자.
$(A^h , I^h) \to (B^h, J^h)$를 헨젤화 위에서 유도된 사상이라 하자
(보조정리 \ref{more-algebra-lemma-henselization}).
$A \to B$가 정수적이면 유도된 사상
$A^h \otimes_A B \to B^h$은 동형이다.
\end{lemma}

\begin{proof}
보조정리 \ref{more-algebra-lemma-henselization-change-ideal}에 의해
$J = IB$라고 가정해도 된다.
보조정리 \ref{more-algebra-lemma-integral-over-henselian-pair}에 의해
쌍 $(A^h \otimes_A B, I^h(A^h \otimes_A B))$는 헨젤이다.
$(B^h, IB^h)$의 보편 성질에 의해 사상
$B^h \to A^h \otimes_A B$를 얻는다.
이 사상이 보조정리의 사상의 역이라는 증명은 생략한다.
\end{proof}


\begin{lemma}
\label{more-algebra-lemma-CRT-henselization}
$I_1, I_2, \dots, I_n \subset A$를 환 $A$의 아이디얼들이라 하자.
$i \not = j$, $1 \leq i, j \leq n$이면 $I_i + I_j = A$라고 가정하자.
$A^h$로 쌍 $(A, I_1 \cap \ldots \cap I_n)$의 헨젤화를 나타내고,
$A_i^h$로 쌍 $(A, I_i)$의 헨젤화를 나타내자.
그러면 자연스러운 사상
$$
A^h \to \prod\nolimits_{i = 1, \ldots, n} A^h_i
$$
은 동형이다.
\end{lemma}

\begin{proof}
이 사상은 구성의 함자성에서 나온다.
$n = 1$이면 결과는 즉시 따른다. $n > 1$이면
아이디얼들에 대한 조건으로부터
$(I_1 \cap \ldots \cap I_{n - 1}) + I_n = A$임을 관찰하자.
따라서 $n = 2$에 대해 보조정리를 증명하면
모든 $n$에 대한 보조정리가 귀납법으로 따른다.

\medskip\noindent
$n = 2$인 경우. $f_1 + f_2 = 1$인 $f_i \in I_i$를 택하자.
$A^h$가 $A \to B$가 에탈이고
% P02-E0788: frozen source omits parentheses in A/I_1 cap I_2; preserve the exact formal span.
$A/I_1 \cap I_2 \to B/(I_1 \cap I_2)B$가 동형이 되게 하는
$A$-대수 $B$들의 쌍대극한으로 구성됨을 상기하자.
보조정리 \ref{more-algebra-lemma-henselization}의 증명을 보라.
그러한 $A$-대수 $B$들의 범주를 $Hens$로 나타내자.
마찬가지로 $i = 1, 2$에 대하여 $Hens_i$로
$A/I_i \to B/I_iB$가 동형이 되게 하는 에탈 $A$-대수 $B$들의
범주를 나타내자.

\medskip\noindent
$f_2$가 가역인 $B$들로 이루어진 부분범주
$Hens_1(f_2) \subset Hens_1$을 생각하자.
이는 공종 부분범주이다. 실제로 $B$가 $Hens_1$에 속하면
$B_{f_2}$도 $Hens_1$에 속한다
(힌트: $f_2$는 $I_1$을 법으로 단원이므로
$B_{f_2}/I_1B_{f_2} = (B/I_1B)_{f_2} = B/I_1B$이다).
마찬가지로 부분범주 $Hens_2(f_1) \subset Hens_2$도 공종이다.
따라서 $A_1^h$와 $A_2^h$를 계산할 때
% P02-E0789: frozen source says “this smaller categories”; Korean supplies plural agreement.
이 더 작은 범주들을 사용해도 된다.
Categories, Lemma \ref{categories-lemma-cofinal}을 보라.

\medskip\noindent
함자
$$
Hens_1(f_2) \times Hens_2(f_1) \to Hens,\quad (B_1, B_2) \mapsto B_1 \times B_2
$$
를 생각하자. 이 함자는 다음 계산 때문에 잘 정의된다.
\begin{align*}
(B_1 \times B_2)/(I_1 \cap I_2)(B_1 \times B_2)
& =
B_1/(I_1 \cap I_2)B_1 \times B_2/(I_1 \cap I_2)B_2 \\
& =
B_1/I_1B_1 \times B_2/I_2B_2 \\
& = A/I_1 \times A/I_2 \\
& = A/(I_1 \cap I_2)
\end{align*}
두 번째 등식은 $f_2 \in I_2$가 $B_1$에서 가역이어서
$(I_1 \cap I_2)B_1 = I_1B_1 \cap I_2B_1 = I_1B_1$이므로 성립한다
(Algebra, Lemma \ref{algebra-lemma-flat-intersect-ideals}도 사용했다).
마지막 등식은 Algebra, Lemma
\ref{algebra-lemma-chinese-remainder}에 의해 성립한다.
증명을 끝내려면 표시된 함자가 공종임을 증명하면 충분하다
(위에서 인용한 보조정리를 보라).
Categories, Definition \ref{categories-definition-cofinal}의
첫 번째 조건은 $Hens$의 대상 $B$가 함자의 상에 속하는
$B_{f_2} \times B_{f_1}$로 사상되므로 성립한다.
두 번째 조건은 항상 성립한다. 실제로
$B \to B_1 \times B_2$와 $B \to B'_1 \times B'_2$가 주어지고
$B_1, B'_1 \in Hens_1(f_2)$ 및 $B_2, B'_2 \in Hens_2(f_1)$이면,
$Hens_1(f_2)$에 속하는
$B''_1 = B_1 \otimes_B B'_1$로 놓을 수 있고, 마찬가지로
% P02-E0790: frozen source defines B''_2 self-referentially; preserve the exact formal span.
$Hens_2(f_1)$ 안에서 $B''_2 = B_2 \otimes_B B''_2$로 놓아
가환도식에 들어가는 사상 $B \to B''_1 \times B''_2$를 찾을 수 있다.
$$
\xymatrix{
& B \ar[ld] \ar[d] \ar[rd] \\
B_1 \times B_2 \ar[r] & B_1'' \times B_2'' & B_1' \times B_2' \ar[l]
}
$$
이로써 증명이 끝난다.
\end{proof}
\section{올림과 헨젤 쌍}
\label{more-algebra-section-lifting-henselian-pairs}

\noindent
이 절에서는 주로 절 \ref{more-algebra-section-lifting}과
\ref{more-algebra-section-henselian-pairs}의 결과를 결합한다.

\begin{lemma}
\label{more-algebra-lemma-lift-finite-projective-module}
$(R, I)$을 헨젤 쌍이라 하자. 사상
$$
P \longrightarrow P/IP
$$
은 유한 사영 $R$-가군들의 동형류 집합과 유한 사영
$R/I$-가군들의 동형류 집합 사이의 전단사를 유도한다.
특히 임의의 유한 사영 $R/I$-가군은 어떤 유한 사영
$R$-가군 $P$에 대하여 $P/IP$와 동형이다.
\end{lemma}

\begin{proof}
먼저 마지막 명제를 증명하자.
$\overline{P}$를 유한 사영 $R/I$-가군이라 하자.
$R$ 위 에탈이고 $R/I = R'/IR'$인 어떤 $R'$ 위의
유한 사영 가군 $P'$로서 $P'/IP'$가 $\overline{P}$와 동형인 것을
찾을 수 있다. 보조정리 \ref{more-algebra-lemma-lift-projective-module}을 보라.
그러면 $(R, I)$이 헨젤 쌍이므로 에탈 환 사상 $R \to R'$은
절단 $\tau : R' \to R$을 갖는다
(보조정리 \ref{more-algebra-lemma-characterize-henselian-pair}).
$P = P' \otimes_{R', \tau} R$로 놓으면
$P/IP$가 $\overline{P}$와 동형이라고 결론짓는다.
물론 이는 보조정리의 명제에 있는 사상이 전사임을 말한다.

\medskip\noindent
단사성. $P_1$과 $P_2$를
$R/I$-가군으로서 $P_1/IP_1 \cong P_2/IP_2$인
유한 사영 $R$-가군들이라 하자.
$P_1$이 사영이므로 주어진 동형을 올리는
$R$-가군 사상 $u : P_1 \to P_2$를 찾을 수 있다.
그러면 Nakayama 보조정리(Algebra, Lemma
\ref{algebra-lemma-NAK})에 의해 $u$는 전사이다.
마찬가지로 전사 $v : P_2 \to P_1$을 찾는다.
Algebra, Lemma \ref{algebra-lemma-fun}에 의해 사상
$v \circ u$는 동형이고, 따라서 $u$가 동형이라고 결론짓는다.
\end{proof}

\begin{lemma}
\label{more-algebra-lemma-finite-etale-equivalence}
$(A, I)$을 헨젤 쌍이라 하자.
함자 $B \to B/IB$는 유한 에탈 $A$-대수들과
유한 에탈 $A/I$-대수들 사이의 동치를 정한다.
\end{lemma}

\begin{proof}
$B, B'$을 $A$ 위 유한 에탈인 두 $A$-대수라 하자.
그러면 $B' \to B'' = B \otimes_A B'$도 유한 에탈이다
(Algebra, Lemmas \ref{algebra-lemma-etale}과
\ref{algebra-lemma-base-change-integral}).
이제 다음 사이에 $1$ 대 $1$ 대응이 있다.
\begin{enumerate}
\item $A$-대수 사상 $B \to B'$,
\item $B' \to B''$의 절단,
\item $B' \to B'' \to eB''$가 동형이 되게 하는
$B''$의 멱등원 $e$.
\end{enumerate}
(2)와 (3) 사이의 전단사는 $\sigma : B'' \to B'$을,
Algebra, Lemmas \ref{algebra-lemma-map-between-etale}과
\ref{algebra-lemma-surjective-flat-finitely-presented}에 의해 존재하는
$\sigma$의 핵을 생성하는 멱등원이 $(1 - e)$가 되게 하는 $e$로 보낸다.
$A/I$-대수 사상 $B/IB \to B'/IB'$과,
$B'/IB' \to B''/IB'' \to \overline{e}(B''/IB'')$가 동형이 되게 하는
$B''/IB''$의 멱등원 $\overline{e}$ 사이에도 비슷한 대응이 있다.
그런데 $B''/IB''$의 모든 멱등원 $\overline{e}$는
$B''$의 멱등원 $e$로 유일하게 올라간다
(보조정리 \ref{more-algebra-lemma-characterize-henselian-pair}).
또한 $B'/IB' \to \overline{e}(B''/IB'')$가 동형이면
Nakayama 보조정리(Algebra, Lemma \ref{algebra-lemma-NAK})에 의해
$B' \to eB''$도 동형이다.
이로부터 함자가 충만충실임을 알 수 있다.

\medskip\noindent
본질적 전사성. $A/I \to C$를 유한 에탈 사상이라 하자.
Algebra, Lemma \ref{algebra-lemma-lift-etale}에 의해
$B/IB \cong C$가 되게 하는 에탈 사상 $A \to B$가 존재한다.
$B'$을 $B$ 안에서 $A$의 정수적 폐포라 하자.
보조정리 \ref{more-algebra-lemma-helper-finite-type}에 의해 어떤 환 $C'$에 대해
$B'/IB' = C \times C'$이고, $(1, 0) \in C \times C'$로 가는
어떤 $g \in B'$에 대하여 $B'_g \cong B_g$이다.
멱등원들이 올라가므로(보조정리
\ref{more-algebra-lemma-characterize-henselian-pair}),
$C = B'_1/IB'_1$이고 $C' = B'_2/IB'_2$인
$B' = B'_1 \times B'_2$를 얻는다.
$B'_1$에서 $g$의 상은 가역이다.
% P02-E0791: frozen source localizes undefined B_2 rather than B'_2; preserve the exact formal span.
그러면 $B_g = B'_g = B'_1 \times (B_2)_g$이고,
이는 $A \to B'_1$이 에탈임을 함의한다.
따라서 $B'_1$은 $A$ 위 유한 에탈이라고 결론짓고 증명이 끝난다
(예를 들어 Algebra, Lemma
\ref{algebra-lemma-characterize-finite-in-terms-of-integral}에 의해
정수적 에탈은 유한 에탈이다).
\end{proof}

\begin{lemma}
\label{more-algebra-lemma-lift-smooth-henselian}
$R$을 환이라 하고 $S$를 매끄러운 $R$-대수라 하자.
$A$가 $R$-대수이고 $(A,I)$가 헨젤 쌍이라고 가정하자.
그러면 임의의 $R$-대수 사상 $S \to A/I$은
$R$-대수 사상 $S \to A$로 올라간다.
\end{lemma}

\begin{proof}
$\tau : S \to A/I$를 $R$-대수 사상이라 하자.
Algebra, Lemma \ref{algebra-lemma-base-change-smooth}에 의해
$S \otimes_R A$는 매끄러운 $A$-대수임을 관찰하자.
따라서 보조정리 \ref{more-algebra-lemma-lift-section-smooth-morphism}에 의해,
유도된 사상 $S \otimes_R A \to A/I$을
% P02-E0792: frozen source misspells “homomorphism”; Korean supplies the mathematical noun.
$A$-대수 준동형 $S \otimes_R A \to A'$으로 올릴 수 있다.
여기서 $A \to A'$은 에탈이고 동형 $A/I \to A'/IA'$을 유도한다.
$(A, I)$이 헨젤이므로 $A$-대수 사상 $A' \to A$가 있다.
보조정리 \ref{more-algebra-lemma-characterize-henselian-pair}를 보라.
합성 $S \to S \otimes_R A \to A' \to A$가 원하는 올림이다.
\end{proof}

\begin{lemma}
\label{more-algebra-lemma-lim-finite-projective-gives-finite-projective}
$A = \lim A_n$을 환들의 역계 $(A_n)$의 극한이라 하자.
$A_n$-가군 $M_n$들과 $A_{n + 1}$-가군 사상
$M_{n + 1} \to M_n$이 주어졌다고 하자. 다음을 가정하자.
\begin{enumerate}
\item 전이사상 $A_{n + 1} \to A_n$은 전사이고 그 핵은 국소 멱영이다.
\item $M_1$은 유한 사영 $A_1$-가군이다.
\item $M_n$은 유한 평탄 $A_n$-가군이다.
\item 사상들은 동형
$M_{n + 1} \otimes_{A_{n + 1}} A_n \to M_n$을 유도한다.
\end{enumerate}
그러면 $M = \lim M_n$은 유한 사영 $A$-가군이고,
모든 $n$에 대하여 $M \otimes_A A_n \to M_n$은 동형이다.
\end{lemma}

\begin{proof}
보조정리 \ref{more-algebra-lemma-limit-henselian}에 의해
쌍 $(A, \Ker(A \to A_1))$은 헨젤이다.
보조정리 \ref{more-algebra-lemma-lift-finite-projective-module}에 의해
유한 사영 $A$-가군 $P$와 동형
$P \otimes_A A_1 \to M_1$을 택할 수 있다.
$P$가 사영이므로 $A$-가군 사상 $P \to M_1$을 차례로
$A$-가군 사상 $P \to M_2$, $P \to M_3$ 등으로 올릴 수 있다.
따라서 사상
$$
P \longrightarrow M
$$
을 얻는다. $P$가 유한 사영이므로 어떤 $m \geq 0$과
$A$-가군 $Q$에 대하여 $A^{\oplus m} = P \oplus Q$로 쓸 수 있다.
$A = \lim A_n$이므로 $P = \lim P \otimes_A A_n$이라고 결론짓는다.
따라서 표시된 $A$-가군 사상이 동형임을 보이려면
사상 $P \otimes_A A_n \to M_n$들이 동형임을 보이면 충분하다.
보조정리 \ref{more-algebra-lemma-lift-projective}에 의해
$M_n$은 유한 사영 가군이다.
보조정리 \ref{more-algebra-lemma-isomorphic-finite-projective-lifts}에 의해
사상 $P \otimes_A A_n \to M_n$들은 동형이다.
\end{proof}


\section{절대 정수적 폐포}
\label{more-algebra-section-absolute-integral-closure}

\noindent
다음과 같이 정의한다.

\begin{definition}
\label{more-algebra-definition-absolutely-integrally-closed}
모든 모닉 다항식 $f \in A[T]$가 일차 인자들의 곱이면
환 $A$를 {\it 절대 정수적으로 닫혔다}고 한다.
\end{definition}

\noindent
주의하자. $f$를 일차 인자들의 곱으로 쓰는 방법은 여러 가지일 수 있다.

\begin{lemma}
\label{more-algebra-lemma-absolutely-integrally-closed}
$A$를 환이라 하자. 다음은 동치이다.
\begin{enumerate}
\item $A$는 절대 정수적으로 닫혀 있다.
\item 임의의 모닉 다항식 $f \in A[T]$는 $A$ 안에 근을 갖는다.
\end{enumerate}
\end{lemma}

\begin{proof}
생략한다.
\end{proof}

\begin{lemma}
\label{more-algebra-lemma-absolutely-integrally-closed-quotient-localization}
$A$가 절대 정수적으로 닫혀 있다고 하자.
\begin{enumerate}
\item $A$의 임의의 몫환 $A/I$는 절대 정수적으로 닫혀 있다.
\item 임의의 국소화 $S^{-1}A$는 절대 정수적으로 닫혀 있다.
\end{enumerate}
\end{lemma}

\begin{proof}
생략한다.
\end{proof}

\begin{lemma}
\label{more-algebra-lemma-integrally-closed-in-absolutely-integrally-closed}
$A$를 환이라 하자. $S \subset A$를 영인자가 아닌 원소들로 이루어진
곱셈적 부분집합이라 하자.
$S^{-1}A$가 절대 정수적으로 닫혀 있고
$A \subset S^{-1}A$가 $S^{-1}A$ 안에서 정수적으로 닫혀 있으면,
$A$도 절대 정수적으로 닫혀 있다.
\end{lemma}

\begin{proof}
생략한다.
\end{proof}

\begin{lemma}
\label{more-algebra-lemma-normal-domain-absolutely-integrally-closed}
$A$를 정규 정역이라 하자. 그러면 $A$가 절대 정수적으로 닫혀 있을
필요충분조건은 그 분수체가 대수적으로 닫혀 있는 것이다.
\end{lemma}

\begin{proof}
체가 대수적으로 닫혀 있을 필요충분조건은 환으로서 절대 정수적으로
닫혀 있는 것임을 관찰하자. 따라서 보조정리는 보조정리
\ref{more-algebra-lemma-absolutely-integrally-closed-quotient-localization}와
\ref{more-algebra-lemma-integrally-closed-in-absolutely-integrally-closed}에서 따른다.
\end{proof}

\begin{lemma}
\label{more-algebra-lemma-construct-absolute-integral-closure}
임의의 환 $A$에 대하여 다음을 만족하는 확대 $A \subset B$가 존재한다.
\begin{enumerate}
\item $B$는 유한 자유 $A$-대수들의 여과 쌍대극한이다.
\item $B$는 $A$-가군으로서 자유이다.
\item $B$는 절대 정수적으로 닫혀 있다.
\end{enumerate}
\end{lemma}

\begin{proof}
$I$를 $A$ 위의 모닉 다항식들의 집합이라 하자. 각 $i \in I$에 대하여
% P02-E0793: frozen source omits the grammatical relation in “denote x_i a variable”; Korean supplies it.
$x_i$를 변수라 하고 $P_i$를 변수 $x_i$에 대응하는 모닉 다항식이라 하자.
다음과 같이 놓는다.
$$
F(A) = A[x_i; i \in I]/(P_i; i \in I)
$$
기호가 시사하듯이 $F$는 환들의 범주에서 자기 자신으로 가는 함자이다.
$A \subset F(A)$이고, $F(A)$가 $A$-가군으로서 자유이며,
$F(A)$가 유한 자유 $A$-대수들의 여과 쌍대극한임을 주목하자.
이제
$$
B = \colim F^n(A)
$$
로 놓는다. 여기서 전이사상들은 포함
$F^n(A) \subset F(F^n(A)) = F^{n + 1}(A)$이다.
계수가 $B$에 속하는 임의의 모닉 다항식은 실제로 어떤 $n$에 대하여
$F^n(A)$에 속하는 계수를 가지므로 구성에 의해 $F^{n + 1}(A)$ 안에 해를 갖는다.
보조정리 \ref{more-algebra-lemma-absolutely-integrally-closed}에 의해
이는 $B$가 절대 정수적으로 닫혀 있음을 함의한다.
다른 성질들의 증명은 생략한다.
\end{proof}

\begin{lemma}
\label{more-algebra-lemma-absolutely-integrally-closed-strictly-henselian}
$A$가 절대 정수적으로 닫혀 있다고 하자. $\mathfrak p \subset A$를
소아이디얼이라 하자. 그러면 국소환 $A_\mathfrak p$는 강 헨젤 환이다.
\end{lemma}

\begin{proof}
보조정리 \ref{more-algebra-lemma-absolutely-integrally-closed-quotient-localization}에 의해
$A$가 국소환이고 $\mathfrak p$가 그 극대아이디얼이라고 가정해도 된다.
잉여체는 보조정리
\ref{more-algebra-lemma-absolutely-integrally-closed-quotient-localization}에 의해
대수적으로 닫혀 있다. 모든 모닉 다항식은 일차 인자들로 완전히 분해되므로
Algebra, Definition \ref{algebra-definition-henselian}이 직접 적용된다.
\end{proof}

\begin{lemma}
\label{more-algebra-lemma-absolutely-integrally-closed-henselian-pair}
$A$가 절대 정수적으로 닫혀 있다고 하자. $I \subset A$를 아이디얼이라 하자.
다음 조건들이 성립하면(또 그때에만) $(A, I)$은 헨젤 쌍이다.
\begin{enumerate}
\item $I$는 $A$의 야코브슨 근기에 포함된다.
\item $A \to A/I$는 멱등원들 사이의 전단사를 유도한다.
\end{enumerate}
\end{lemma}

\begin{proof}
$f \in A[T]$를 모닉 다항식이라 하고,
$f \bmod I = g_0 h_0$를 $A/I$ 위의 분해라 하자.
여기서 $g_0$, $h_0$는 모닉이고
$g_0$와 $h_0$는 $A/I[T]$의 단위아이디얼을 생성한다고 하자.
이는 다음을 뜻한다.
$$
A/I[T]/(f) = A/I[T]/(g_0) \times A/I[T]/(h_0)
$$
% P02-E0794: frozen source misspells “corresponding”; Korean supplies the intended word.
$e \in A/I[T]/(f)$를 오른쪽 환의 멱등원 $(1, 0)$에 대응하는 원소라 하자.
$a_i \in A$인 원소들을 사용하여 $f = (T - a_1) \ldots (T - a_d)$로 쓰자.
각 $i \in \{1, \ldots, d\}$에 대하여 $A$-대수 사상
$\varphi_i : A[T]/(f) \to A$, $T \mapsto a_i$를 얻으며, 이는 비슷한
$A/I$-대수 사상 $\overline{\varphi}_i : A/I[T]/(f) \to A/I$를 유도한다.
$e_i = \overline{\varphi}_i(e) \in A/I$로 놓자. 이들은 멱등원이다.
가정 (2)에 의해 $e_i$를 $A$의 멱등원으로 올릴 수 있다.
이는 $A = \prod A_j$를 환들의 유한곱으로 써서, 각 $A_j/IA_j$에서
각 $e_i$가 $0$ 또는 $1$이 되게 할 수 있음을 뜻한다.
몇몇 세부사항은 생략한다. $A_j$는 $A$의 몫환이므로
절대 정수적으로 닫혀 있음을 관찰하자. $f$의 $A_j/IA_j$ 위 분해를
$A_j$로 올리는 것으로 충분하다.
이로써 다음 단락에서 논의하는 상황으로 환원된다.

\medskip\noindent
$i = 1, \ldots, r$에 대하여 $e_i = 1$이고
$i = r + 1, \ldots, d$에 대하여 $e_i = 0$이라고 가정하자.
$(g_0, h_0) = A/I[T]$이므로 $g_0 k_0 + h_0 l_0 = 1$이 되게 하는
$k_0, l_0 \in A/I[T]$가 존재한다.
$e = h_0 l_0$이고 $e_i = h_0(a_i) l_0(a_i)$임을 알 수 있다.
따라서 $i = 1, \dots ,r$에 대하여 $h_0(a_i)$는 단원이라고 결론짓는다.
$f(a_i) = 0$이므로 $0 = h_0(a_i)g_0(a_i)$를 얻고,
$i = 1, \ldots, r$에 대하여 $g_0(a_i) = 0$이라고 결론짓는다.
따라서 $(T - a_1)$은 $A/I[T]$ 안에서 $g_0$를 나누며,
$g_0 = (T - a_1) g_0'$로 쓸 수 있다.
$f' = (T - a_2) \ldots (T - a_d)$와 $h'_0 = h_0$로 놓자.
% P02-E0795: frozen source repeats “over” in “over over A”; Korean states the relation once.
$d$에 관한 귀납법으로 분해 $f' \bmod I = g'_0 h'_0$을
$A$ 위의 분해 $f' = g' h'$로 올릴 수 있으며, 이는 원하는 대로
분해 $f \bmod I = g_0 h_0$를 올리는 분해
$f = (T - a_1) g' h'$을 준다.
\end{proof}

\section{자기 동반 환}
\label{more-algebra-section-auto-ass}

\noindent
이 내용의 일부는 \cite{Autour}에 있다.

\begin{definition}
\label{more-algebra-definition-auto-ass}
환 $R$이 국소환이고 그 극대아이디얼 $\mathfrak m$이 $R$에 약하게
동반되면 $R$을 {\it 자기 동반} 환이라 한다.
\end{definition}

\begin{lemma}
\label{more-algebra-lemma-auto-ass-implies-P}
자기 동반 환 $R$은 다음 성질을 갖는다. (P)
모든 진 유한 생성 아이디얼 $I \subset R$은 영이 아닌 소멸자를 갖는다.
\end{lemma}

\begin{proof}
가정에 의해 모든 $f \in \mathfrak m$에 대하여 $f^n x = 0$인
영이 아닌 원소 $x \in R$이 존재한다.
% P02-E0796: frozen source leaves exponent n unquantified; no rendered source correction is introduced.
$I = (f_1, \ldots, f_r)$로 쓰자.
그러면 $x$는 $R \to \bigoplus R_{f_i}$의 핵에 속한다.
따라서 모든 $i$에 대하여 $f_i y = 0$인 영이 아닌 $y \in R$이
존재함을 알 수 있다. Algebra, Lemma
\ref{algebra-lemma-when-injective-covering}을 보라.
$y \in \text{Ann}_R(I)$이므로 증명이 끝난다.
\end{proof}

\begin{lemma}
\label{more-algebra-lemma-P-universally-injective}
$R$을 보조정리 \ref{more-algebra-lemma-auto-ass-implies-P}의 성질 (P)를 갖는 환이라 하자.
$u : N \to M$을 사영 $R$-가군들의 준동형이라 하자.
그러면 $u$가 보편 단사일 필요충분조건은 $u$가 단사인 것이다.
\end{lemma}

\begin{proof}
$u$가 단사라고 가정하자. 목표는 $u$가 보편 단사임을 보이는 것이다.
먼저 $N \oplus Q$가 자유가 되게 하는 가군 $Q$를 택한다.
사상 $N \oplus Q \to M \oplus Q$를 생각하면 $N$이 자유인 경우에
보조정리를 증명하는 것으로 충분함을 알 수 있다.
이 경우 $N$은 유한 자유 $R$-가군들의 유향 쌍대극한이다.
따라서 $N$이 유한 자유 $R$-가군인 경우로 환원되며,
$N = R^{\oplus n}$으로 쓰자. $n$에 관한 귀납법으로 보조정리를 증명한다.
$n = 0$인 경우는 자명하다.

\medskip\noindent
$u : R^{\oplus n} \to M$을 $M$이 사영인 단사 가군 사상이라 하자.
$M \oplus Q$가 자유가 되게 하는 $R$-가군 $Q$를 택한다.
$u$를 합성 $R^{\oplus n} \to M \to M \oplus Q$로 바꾸면
$M$이 자유라고 가정해도 됨을 알 수 있다.
그러면 $u(R^{\oplus n}) \subset R^{\oplus m}$이 되게 하는
직접합인자 $R^{\oplus m} \subset M$을 찾을 수 있다.
따라서 $M = R^{\oplus m}$이라고 가정해도 된다.
이 경우 $u$는 행렬 $A = (a_{ij})$로 주어지며
$u(x_1, \ldots, x_n) = (\sum x_i a_{i1}, \ldots, \sum x_i a_{im})$이다.
$u$가 단사이므로 특히 $x \not = 0$이면
$u(x, 0, \ldots, 0) = (xa_{11}, xa_{12}, \ldots, xa_{1m}) \not = 0$이고,
$R$이 성질 (P)를 가지므로
$a_{11}R + a_{12}R + \ldots + a_{1m}R = R$임을 알 수 있다.
% P02-E0797: frozen source says “Hence see that”; Korean supplies grammatical phrasing.
따라서 $R(a_{11}, \ldots, a_{1m}) \subset R^{\oplus m}$이
$R^{\oplus m}$의 직접합인자임을 알 수 있으며, 특히
$R^{\oplus m}/R(a_{11}, \ldots, a_{1m})$은 사영 $R$-가군이다.
행들이 분할 완전인 가환도식
$$
\xymatrix{
0 \ar[r] &
R \ar[rr] \ar[d]^1 & & R^{\oplus n} \ar[r] \ar[d]^u &
R^{\oplus n - 1} \ar[r] \ar[d] & 0 \\
0 \ar[r] & R \ar[rr]^{(a_{11}, \ldots, a_{1m})} & &
R^{\oplus m} \ar[r] & R^{\oplus m}/R(a_{11}, \ldots, a_{1m}) \ar[r] & 0
}
$$
를 얻는다. 따라서 오른쪽 수직 화살표는 단사이며, 귀납가정을 적용하여
오른쪽 수직 화살표가 보편 단사라고 결론 내릴 수 있다.
그러므로 가운데 수직 화살표도 보편 단사이다.
\end{proof}

\begin{lemma}
\label{more-algebra-lemma-P-fPD-zero}
$R$을 환이라 하자. 다음은 동치이다.
\begin{enumerate}
\item $R$은 보조정리 \ref{more-algebra-lemma-auto-ass-implies-P}의 성질 (P)를 갖는다.
\item 사영 $R$-가군들의 임의의 단사 사상은 보편 단사이다.
\item $u : N \to M$이 단사이고 $N$, $M$이 유한 사영 $R$-가군이면
$\Coker(u)$는 유한 사영 $R$-가군이다.
\item $N \subset M$이고 $N$, $M$이 $R$-가군으로서 유한 사영이면
$N$은 $M$의 직접합인자이다.
\item 임의의 단사 사상 $R \to R^{\oplus n}$은 분할 단사이다.
\end{enumerate}
\end{lemma}

\begin{proof}
함의 (1) $\Rightarrow$ (2)는 보조정리
\ref{more-algebra-lemma-P-universally-injective}이다.
(3)과 (4)가 동치임은 명백하다.
Algebra, Lemma \ref{algebra-lemma-universally-exact-split}에 의해
(2) $\Rightarrow$ (3), (4)이다. (5)는 (4)의 특수한 경우이다.
(5)를 가정하자. $I = (a_1, \ldots, a_n)$을 $R$의 진 유한 생성
아이디얼이라 하자. $I \not = R$이므로
$R \to R^{\oplus n}$, $x \mapsto (xa_1, \ldots, xa_n)$은
분할 단사가 아님을 알 수 있다. 따라서 그 핵은 영이 아니며,
$\text{Ann}_R(I) \not = 0$이라고 결론짓는다. 그러므로 (1)이 성립한다.
\end{proof}

\begin{example}
\label{more-algebra-example-auto-ass-weird-flat}
보조정리 \ref{more-algebra-lemma-P-fPD-zero}의 동치인 조건들이 성립하더라도,
자유 $R$-가군들의 모든 단사 사상이 항상 분할 단사인 것은 아니다.
예를 들어 $R = k[x_1, x_2, x_3, \ldots]/(x_i^2)$라고 하자.
이는 자기 동반 환이다. 자유 $R$-가군들의 사상
$$
u :
\bigoplus\nolimits_{i \geq 1} Re_i
\longrightarrow
\bigoplus\nolimits_{i \geq 1} Rf_i, \quad
e_i \longmapsto f_i - x_if_{i + 1}.
$$
를 생각하자. 임의의 정수 $n$에 대하여 $u$를
$\bigoplus_{i = 1, \ldots, n} Re_i$에 제한하면 단사이다.
상 $u(e_1), \ldots, u(e_n)$이 $R$-선형 독립이기 때문이다.
따라서 $u$는 단사이고, 보조정리에 의해 보편 단사이다.
$u \otimes \text{id}_k$가 전단사이므로, $u$가 분할 단사라면
$u$는 전사여야 함을 알 수 있다. 그러나 $u$는 전사가 아니다.
$f_1$의 역상이 될 원소
$$
\sum\nolimits_{i \geq 0} x_1 \ldots x_ie_{i + 1} =
e_1 + x_1e_2 + x_1x_2e_3 + \ldots
$$
는 직접합의 원소가 아니기 때문이다.
덧붙여 $\Coker(u)$는 ($u$가 보편 단사이므로) 평탄하고,
($u$가 분할되지 않으므로) 사영이 아닌 가산 생성 $R$-가군이며,
따라서 Mittag-Leffler 가군이 아니다
(Algebra, Lemma \ref{algebra-lemma-countgen-projective}를 보라).
\end{example}

\noindent
다음 보조정리는 국소환이 뇌터 환인 경우에는
Algebra, Proposition \ref{algebra-proposition-what-exact}의 특수한 경우이다.

\begin{lemma}
\label{more-algebra-lemma-exact-length-1}
$(R, \mathfrak m)$을 국소환이라 하자.
$\varphi : R^m \to R^n$이 유한 자유 가군들의 사상이라고 하자.
다음은 동치이다.
\begin{enumerate}
\item $\varphi$는 단사이다.
\item $\varphi$의 계수는 $m$이고, $R$ 안에서 $I(\varphi)$의 소멸자는 영이다.
\end{enumerate}
$R$이 뇌터 환이면 이들은 다음과도 동치이다.
\begin{enumerate}
\item[(3)] $\varphi$의 계수는 $m$이고,
$I(\varphi) = R$이거나 이 아이디얼이 영인자가 아닌 원소를 포함한다.
\end{enumerate}
여기서 $\varphi$의 계수와 $I(\varphi)$는
Algebra, Definition \ref{algebra-definition-rank}에서와 같이 정의된다.
\end{lemma}

\begin{proof}
$\varphi$의 행렬 계수 중 하나가 $\mathfrak m$에 속하지 않으면,
Algebra, Lemma \ref{algebra-lemma-add-trivial-complex}를 적용하여
$\varphi$를 $1 : R \to R$과 사상 $\varphi' : R^{m-1} \to R^{n-1}$의
합으로 쓴다. $\varphi'$에 대한 보조정리가 $\varphi$에 대한 보조정리를
함의함은 쉽게 알 수 있다. 따라서 처음부터 $\varphi$의 모든 행렬 계수가
$\mathfrak m$에 속한다고 가정해도 된다.

\medskip\noindent
$\varphi$가 단사라고 하자. $m > 0$이라고 가정해도 된다.
$\mathfrak q \in \text{WeakAss}(R)$라 하면 $R_\mathfrak q$는 자기 동반 환이다.
그러면 $\varphi$는 단사 사상 $R_\mathfrak q^m \to R_\mathfrak q^n$을 유도하며,
이는 보조정리 \ref{more-algebra-lemma-auto-ass-implies-P}와
\ref{more-algebra-lemma-P-universally-injective}에 의해 보편 단사이다.
따라서 $\varphi : \kappa(\mathfrak q)^m \to \kappa(\mathfrak q)^n$은 단사이다.
그러므로 $\varphi \bmod \mathfrak q$의 계수는 $m$이고
$I(\varphi \otimes \kappa(\mathfrak q))$는 영아이디얼이 아니다.
$m$은 $\varphi$가 가질 수 있는 최대 계수이므로,
$\varphi$의 계수 역시 $m$이라고 결론짓는다
(행렬의 계수는 기저변환 뒤에 감소할 수만 있다).
따라서 $I(\varphi) \cdot \kappa(\mathfrak q) =
I(\varphi \otimes \kappa(\mathfrak q))$는 영이 아니다.
그러므로 $I(\varphi)$는 $\mathfrak q$에 포함되지 않는다.
따라서 $R$의 약한 동반 소아이디얼 중 어느 것도
$R$-가군 $\text{Ann}_R I(\varphi)$의 약한 동반 소아이디얼이 아니다.
그러므로 $\text{Ann}_R I(\varphi)$는 약한 동반 소아이디얼을 갖지 않는다.
Algebra, Lemma \ref{algebra-lemma-weakly-ass}를 보라.
Algebra, Lemma \ref{algebra-lemma-weakly-ass-zero}에 의해
$\text{Ann}_R I(\varphi)$는 영이다.

\medskip\noindent
거꾸로 (2)를 가정하자. 계수가 $m$이므로 $n \geq m$이다.
$I(\varphi)$는 유한 생성이므로 $I(\varphi) = (f_1, \ldots, f_r)$로 쓸 수 있다.
Algebra, Lemma \ref{algebra-lemma-matrix-left-inverse}에 의해
% P02-E0798: frozen source uses undefined psi rather than psi_i in the following identity; preserve that exact span.
$\psi \circ \varphi = f_i \text{id}_{R^m}$이 되게 하는 사상
$\psi_i : R^n \to R^m$을 찾을 수 있다.
따라서 $\varphi(x) = 0$이면 $i = 1, \ldots, r$에 대하여 $f_i x = 0$이다.
이는 $x = 0$을 함의하므로 $\varphi$는 단사이다.

\medskip\noindent
뇌터 국소환의 경우 (1)과 (3)의 동치는
Algebra, Proposition \ref{algebra-proposition-what-exact}를 보라.
환 $R$이 뇌터 환이지만 국소환은 아니면 독자는 국소적인 경우에서 이를
도출할 수 있다. 세부사항은 생략한다.
또 다른 방법은 동반 소아이디얼들이 유한 개라는 사실,
소아이디얼 회피 정리, 그리고 뇌터 환에서 영인자가 아닌 원소들은
어느 동반 소아이디얼에도 포함되지 않는 원소들이라는 특징짓기를 사용하여,
위 논증을 동반 소아이디얼들로 다시 수행하는 것이다.
\end{proof}

\begin{lemma}
\label{more-algebra-lemma-coker-injective-free}
$R$을 환이라 하자.
% P02-E0799: frozen source combines “Suppose that” with “be”; Korean uses one grammatical construction.
$\varphi : R^n \to R^n$이 계수가 같은 유한 자유 가군들의 단사 사상이라고 하자.
그러면 $\Hom_R(\Coker(\varphi), R) = 0$이다.
\end{lemma}

\begin{proof}
$\varphi^t : R^n \to R^n$을 $\varphi$의 전치라 하자.
보조정리는 $\varphi^t$가 단사임을 주장한다.
보조정리 \ref{more-algebra-lemma-exact-length-1}의 기호를 사용하면
$\varphi^t$의 계수는 $n$이고 $I(\varphi) = I(\varphi^t)$임을 알 수 있다.
따라서 그 보조정리의 (1)과 (2)의 동치에 의해 결론을 얻는다.
\end{proof}

\section{평탄화 층화}
\label{more-algebra-section-flattening}

\noindent
$R \to S$를 환 사상이라 하고 $M$을 $S$-가군이라 하자.
임의의 $R$-대수 $R'$에 대하여 기저변환
$S' = S \otimes_R R'$과 $M' = M \otimes_R R'$을 생각할 수 있다.
가군 $M'$이 $R'$ 위 평탄이면 $R \to R'$이 $M$을 {\it 평탄화한다}고 한다.
$M$을 평탄화하는 환 사상 $R \to R'$들의 모임의 구조를 이해하고자 한다.
특히 {\it $M$의 보편 평탄화 $R \to R_{univ}$}가 존재하는지 알고 싶다.
이는 $M$을 평탄화하는 환 사상 $R \to R_{univ}$로서,
$M$을 평탄화하는 임의의 환 사상 $R \to R'$이 $R \to R_{univ}$를
통하여 분해되는 성질을 갖는 것을 말한다.
이러한 보편해는 대개 존재하지 않는다.

\medskip\noindent
{\it 보편 평탄화}와 {\it 평탄화 층화}는 스킴론적 상황
$\mathcal{F}/X/S$에서 More on Flatness, Section
\ref{flat-section-flattening}에서 논의할 것이다.
보편 평탄화 $R \to R_{univ}$가 존재하면,
스킴 사상 $\Spec(R_{univ}) \to \Spec(R)$은
$\Spec(S)$ 위의 준연접 가군 $\widetilde{M}$의 보편 평탄화이다.

\medskip\noindent
이 절과 다음 몇 절에서는 이와 관련된 몇 가지 기본적인 대수적 사실을 증명한다.
아마도 가장 기본적인 결과는 다음일 것이다.

\begin{lemma}
\label{more-algebra-lemma-intersection-flat}
$R$을 환이라 하자. $M$을 $R$-가군이라 하고 $I_1$, $I_2$를 $R$의 아이디얼이라 하자.
$M/I_1M$이 $R/I_1$ 위 평탄이고 $M/I_2M$이 $R/I_2$ 위 평탄이면,
$M/(I_1 \cap I_2)M$은 $R/(I_1 \cap I_2)$ 위 평탄이다.
\end{lemma}

\begin{proof}
$R$을 $R/(I_1 \cap I_2)$로, $M$을 $M/(I_1 \cap I_2)M$으로 바꾸면
$I_1 \cap I_2 = 0$이라고 가정해도 된다. $J \subset R$을 아이디얼이라 하자.
$M$이 $R$ 위 평탄임을 증명하려면 $J \otimes_R M \to M$이 단사임을
보여야 한다. Algebra, Lemma \ref{algebra-lemma-flat}을 보라.
$M/I_1M$이 $R/I_1$ 위 평탄이므로 사상
$$
J/(J \cap I_1) \otimes_R M =
(J + I_1)/I_1 \otimes_{R/I_1} M/I_1M
\longrightarrow M/I_1M
$$
은 단사이다. $0 \to (J \cap I_1) \to J \to J/(J \cap I_1) \to 0$이
완전하므로 도식
$$
\xymatrix{
(J \cap I_1) \otimes_R M \ar[r] \ar[d] &
J \otimes_R M \ar[r] \ar[d] &
J/(J \cap I_1) \otimes_R M \ar[r] \ar[d] & 0 \\
M \ar@{=}[r] &
M \ar[r] &
M/I_1M
}
$$
을 얻는다. 따라서 $(J \cap I_1) \otimes_R M \to M$이 단사임을
보이면 충분하다. $I_1 \cap I_2 = 0$이므로 아이디얼 $J \cap I_1$은
아이디얼 $J' \subset R/I_2$에 동형으로 사상되며,
$(J \cap I_1) \otimes_R M = J' \otimes_{R/I_2} M/I_2M$임을 알 수 있다.
$M/I_2M$이 $R/I_2$ 위 평탄이므로 사상
$J' \otimes_{R/I_2} M/I_2M \to M/I_2M$은 단사이며,
이는 분명히 $(J \cap I_1) \otimes_R M \to M$이 단사임을 함의한다.
\end{proof}

\section{아르틴 환 위의 평탄화}
\label{more-algebra-section-flattening-artinian}

\noindent
기저환이 아르틴 국소환이면 보편 평탄화가 존재한다.
이는 임의의 가군에 대하여 존재한다.
% P02-E0800: frozen source misspells “flattening” as “flatting”; Korean uses the intended construction.
따라서 나중에 보겠지만, 기저 스킴이 아르틴 국소환의 스펙트럼이면
평탄화 층화가 존재한다.

\begin{lemma}
\label{more-algebra-lemma-flattening-artinian}
$R$을 아르틴 환이라 하자. $M$을 $R$-가군이라 하자.
그러면 $M/IM$이 $R/I$ 위 평탄이 되게 하는 가장 작은 아이디얼
$I \subset R$이 존재한다.
\end{lemma}

\begin{proof}
이는 보조정리 \ref{more-algebra-lemma-intersection-flat}과 아르틴 성질에서 직접 따른다.
\end{proof}

\noindent
이 아이디얼은 다음 보편 성질을 갖는다.

\begin{lemma}
\label{more-algebra-lemma-flattening-artinian-universal-property}
$R$을 아르틴 환이라 하고 $M$을 $R$-가군이라 하자.
$I \subset R$을 $M/IM$이 $R/I$ 위 평탄이 되게 하는
가장 작은 아이디얼 $I \subset R$이라 하자.
그러면 $I$는 다음 보편 성질을 갖는다.
모든 환 사상 $\varphi : R \to R'$에 대하여 다음이 성립한다.
$$
R' \otimes_R M\text{ is flat over }R'
\Leftrightarrow
\text{we have }\varphi(I) = 0.
$$
\end{lemma}

\begin{proof}
보조정리 \ref{more-algebra-lemma-flattening-artinian}에 의해 $I$가 존재함을 주목하자.
% P02-E0801: frozen source attributes the forward arrow to flat base change; retain its exact arrow.
함의 $\Rightarrow$는 Algebra, Lemma \ref{algebra-lemma-flat-base-change}에서 따른다.
$\varphi : R \to R'$이 $M \otimes_R R'$을 $R'$ 위 평탄하게 하는 사상이라 하자.
$J = \Ker(\varphi)$로 두자. Algebra, Lemma
\ref{algebra-lemma-descent-flatness-injective-map-artinian-rings}에 의해,
그리고 $R' \otimes_R M = R' \otimes_{R/J} M/JM$이 $R'$ 위 평탄이므로,
$M/JM$이 $R/J$ 위 평탄이라고 결론짓는다.
따라서 원하는 대로 $I \subset J$이다.
\end{proof}

\section{기저의 닫힌 부분집합 위의 평탄화}
\label{more-algebra-section-flattening-local-base}

\noindent
$R \to S$를 환 사상이라 하자. $I \subset R$을 아이디얼이라 하자.
$M$을 $S$-가군이라 하자. 이하에서는 다음 조건을 생각한다.
\begin{equation}
\label{more-algebra-equation-flat-at-primes-over}
\forall \mathfrak q \in V(IS) \subset \Spec(S) :
M_{\mathfrak q}\text{ is flat over }R.
\end{equation}
기하학적으로 이는 $M$이 $\Spec(S)$ 안에서 $V(I)$의 역상을 따라
$R$ 위 평탄임을 뜻한다. $R$과 $S$가 뇌터 환이고 $M$이 유한 $S$-가군이면,
(\ref{more-algebra-equation-flat-at-primes-over})는 모든 $n \geq 1$에 대하여
$M/I^nM$이 $R/I^n$ 위 평탄이라는 조건과 동치이다.
Algebra, Lemma \ref{algebra-lemma-flat-module-powers}를 보라.

\begin{lemma}
\label{more-algebra-lemma-base-change-flat-at-primes-over}
$R \to S$를 환 사상이라 하자. $I \subset R$을 아이디얼이라 하고
$M$을 $S$-가군이라 하자. $R \to R'$을 환 사상이라 하고
$IR' \subset I' \subset R'$을 아이디얼이라 하자.
$(R \to S, I, M)$에 대하여 (\ref{more-algebra-equation-flat-at-primes-over})가 성립하면,
$(R' \to S \otimes_R R', I', M \otimes_R R')$에 대해서도
(\ref{more-algebra-equation-flat-at-primes-over})가 성립한다.
\end{lemma}

\begin{proof}
$(R \to S, I \subset R, M)$에 대하여
(\ref{more-algebra-equation-flat-at-primes-over})가 성립한다고 가정하자.
$I'(S \otimes_R R') \subset \mathfrak q'$를 $S \otimes_R R'$의 소아이디얼이라 하자.
$\mathfrak q \subset S$를 이에 대응하는 $S$의 소아이디얼이라 하자.
그러면 $IS \subset \mathfrak q$이다.
$(M \otimes_R R')_{\mathfrak q'}$는 기저변환
$M_{\mathfrak q} \otimes_R R'$의 국소화임을 주목하자.
따라서 $(M \otimes_R R')_{\mathfrak q'}$는 평탄 가군의 국소화이므로
$R'$ 위 평탄이다. Algebra, Lemmas \ref{algebra-lemma-flat-base-change}와
\ref{algebra-lemma-flat-localization}을 보라.
\end{proof}

\begin{lemma}
\label{more-algebra-lemma-flat-descent-flat-at-primes-over}
$R \to S$를 환 사상이라 하자. $I \subset R$을 아이디얼이라 하고
$M$을 $S$-가군이라 하자. $R \to R'$을 환 사상이라 하고
$IR' \subset I' \subset R'$을 다음 조건을 만족하는 아이디얼이라 하자.
\begin{enumerate}
\item $\Spec(R') \to \Spec(R)$이 유도하는 사상 $V(I') \to V(I)$는 전사이다.
\item 모든 소아이디얼 $\mathfrak p' \in V(I')$에 대하여
$R'_{\mathfrak p'}$는 $R$ 위 평탄이다.
\end{enumerate}
$(R' \to S \otimes_R R', I', M \otimes_R R')$에 대하여
(\ref{more-algebra-equation-flat-at-primes-over})가 성립하면,
$(R \to S, I, M)$에 대하여 (\ref{more-algebra-equation-flat-at-primes-over})가 성립한다.
\end{lemma}

\begin{proof}
% P02-E0802: frozen source opens this proof with IR' instead of I'; preserve its exact tuple.
$(R' \to S \otimes_R R', IR', M \otimes_R R')$에 대하여
(\ref{more-algebra-equation-flat-at-primes-over})가 성립한다고 가정하자.
소아이디얼 $IS \subset \mathfrak q \subset S$를 택한다.
$I \subset \mathfrak p \subset R$을 이에 대응하는 $R$의 소아이디얼이라 하자.
가정에 의해 $\mathfrak p$ 위에 놓이는 $R'$의 소아이디얼
$\mathfrak p' \in V(I')$가 존재하고,
$R_{\mathfrak p} \to R'_{\mathfrak p'}$는 평탄이다.
소아이디얼 $\overline{\mathfrak q}' \subset
\kappa(\mathfrak q) \otimes_{\kappa(\mathfrak p)} \kappa(\mathfrak p')$를 택한다.
이는 $\mathfrak q$ 위이면서 $\mathfrak p'$ 위에 놓이는 소아이디얼
$\mathfrak q' \subset S \otimes_R R'$에 대응한다.
$(S \otimes_R R')_{\mathfrak q'}$는
$S_{\mathfrak q} \otimes_{R_{\mathfrak p}} R'_{\mathfrak p'}$의 국소화임을 주목하자.
가정에 의해 가군 $(M \otimes_R R')_{\mathfrak q'}$는
$R'_{\mathfrak p'}$ 위 평탄이다. 따라서 Algebra, Lemma
\ref{algebra-lemma-base-change-flat-up-down}은
$M_{\mathfrak q}$가 $R_{\mathfrak p}$ 위 평탄임을 함의한다.
이것이 증명하고자 한 바이다.
\end{proof}

\begin{lemma}
\label{more-algebra-lemma-limit-preserving-flat-at-primes-over}
$R \to S$를 유한 표시 환 사상이라 하자.
$M$을 유한 표시 $S$-가군이라 하자.
$R' = \colim_{\lambda \in \Lambda} R_\lambda$를
$R$-대수들의 유향 쌍대극한이라 하자.
$I_\lambda \subset R_\lambda$를 모든 $\mu \geq \lambda$에 대하여
$I_\lambda R_\mu \subset I_\mu$를 만족하는 아이디얼들이라 하고
$I' = \colim_\lambda I_\lambda$로 놓자.
$(R' \to S \otimes_R R', I', M \otimes_R R')$에 대하여
(\ref{more-algebra-equation-flat-at-primes-over})가 성립하면,
$(R_\lambda \to S \otimes_R R_\lambda, I_\lambda, M \otimes_R R_\lambda)$에 대하여
(\ref{more-algebra-equation-flat-at-primes-over})가 성립하는 $\lambda \in \Lambda$가 존재한다.
\end{lemma}

\begin{proof}
$S_\lambda = S \otimes_R R_\lambda$,
$S' = S \otimes_R R'$, $M_\lambda = M \otimes_R R_\lambda$,
$M' = M \otimes_R R'$로 쓰겠다.
기저변환 $S'$는 $R'$ 위 유한 표시이고 $M'$는 $S'$ 위 유한 표시이며,
첨자 $\lambda$가 있는 경우도 마찬가지이다.
Algebra, Lemma \ref{algebra-lemma-base-change-finiteness}를 보라.
Algebra, Theorem \ref{algebra-theorem-openness-flatness}에 의해 집합
$$
U' = \{\mathfrak q' \in \Spec(S') \mid
M'_{\mathfrak q'}\text{ is flat over }R'\}
$$
는 $\Spec(S')$에서 열려 있다.
$V(I'S')$는 준콤팩트 공간이며 가정에 의해 $U'$에 포함됨을 주목하자.
따라서 $D(g'_j) \subset U'$이고 $V(I'S') \subset \bigcup D(g'_j)$인
유한 개의 $g'_j \in S'$, $j = 1, \ldots, m$이 존재한다.
특히 $(M')_{g'_j}$는 $R'$ 위 평탄 가군임을 주목하자.

\medskip\noindent
점점 더 큰 원소 $\lambda \in \Lambda$를 택하겠다.
먼저 $g'_j$로 가는 $g_{j, \lambda} \in S_{\lambda}$를 찾을 수 있을 만큼
충분히 크게 택한다. 포함 $V(I'S') \subset \bigcup D(g'_j)$는
$I'S' + (g'_1, \ldots, g'_m) = S'$를 뜻하며,
이는 어떤 $z_s \in I'$, $h_s, f_j \in S'$를 사용하여
$1 = \sum z_sh_s + \sum f_jg'_j$로 나타낼 수 있다.
$\lambda$를 더 크게 하면 이러한 등식이 $S_\lambda$에서 성립한다고
가정해도 된다. 따라서
$V(I_\lambda S_\lambda) \subset \bigcup D(g_{j, \lambda})$라고 가정할 수 있다.
Algebra, Lemma \ref{algebra-lemma-flat-finite-presentation-limit-flat}에 의해
충분히 큰 어떤 $\lambda$에 대하여 가군 $(M_\lambda)_{g_{j, \lambda}}$들이
$R_\lambda$ 위 평탄임을 알 수 있다.
특히 가군 $M_\lambda$는 아이디얼 $I_\lambda$ 위에 놓이는 모든
소아이디얼에서 $R_\lambda$ 위 평탄이다.
\end{proof}

% P02-E0803: frozen section title mixes singular article and plural noun; Korean supplies grammatical number.
\section{정의역과 기저의 닫힌 부분집합 위의 평탄화}
\label{more-algebra-section-flattening-local-source-base}

\noindent
이 절에서는 절 \ref{more-algebra-section-flattening-local-base}의 논의를 조금 일반화한다.
독자는 먼저 그 절을 읽고 이해하기를 강력히 권한다.

\begin{situation}
\label{more-algebra-situation-flattening-general}
$R \to S$를 환 사상이라 하자. $J \subset S$를 아이디얼이라 하고
$M$을 $S$-가군이라 하자.
\end{situation}

\noindent
이 상황에서 $R$-대수 $R'$과 아이디얼 $I' \subset R'$이 주어지면
$S' = S \otimes_R R'$과 $M' = M \otimes_R R'$로 놓는다.
다음 조건을 생각하겠다.
\begin{equation}
\label{more-algebra-equation-flat-at-primes}
\forall \mathfrak q' \in V(I'S' + JS') \subset \Spec(S') :
M'_{\mathfrak q'}\text{ is flat over }R'.
\end{equation}
기하학적으로 이는 $M'$가 $V(I')$의 역상과 $V(J)$의 역상의 교집합을
따라 $R'$ 위 평탄임을 뜻한다.
$(R \to S, J, M)$이 고정되어 있으므로 조건
(\ref{more-algebra-equation-flat-at-primes})는 쌍 $(R', I')$에만 의존한다.
여기서 $R'$은 $R$-대수로 본다.

\begin{lemma}
\label{more-algebra-lemma-base-change-flat-at-primes}
상황 \ref{more-algebra-situation-flattening-general}에서 $R' \to R''$을 $R$-대수 사상이라 하자.
$I' \subset R'$과 $I'R'' \subset I'' \subset R''$을 아이디얼이라 하자.
$(R', I')$에 대하여 (\ref{more-algebra-equation-flat-at-primes})가 성립하면,
$(R'', I'')$에 대해서도 (\ref{more-algebra-equation-flat-at-primes})가 성립한다.
\end{lemma}

\begin{proof}
$(R', I')$에 대하여 (\ref{more-algebra-equation-flat-at-primes})가 성립한다고 가정하자.
$I''S'' + JS'' \subset \mathfrak q''$를 $S''$의 소아이디얼이라 하자.
$\mathfrak q' \subset S'$를 이에 대응하는 $S'$의 소아이디얼이라 하자.
그러면 $\mathfrak q''$에 대하여 대응하는 조건들이 성립하므로
$I'S' \subset \mathfrak q'$와 $JS' \subset \mathfrak q'$가 모두 성립한다.
% P02-E0804: frozen source tensors the R'-module over R rather than R'; retain its exact tensor base.
$(M'')_{\mathfrak q''}$는 기저변환
$M'_{\mathfrak q'} \otimes_R R''$의 국소화임을 주목하자.
따라서 $(M'')_{\mathfrak q''}$는 평탄 가군의 국소화이므로 $R''$ 위 평탄이다.
Algebra, Lemmas \ref{algebra-lemma-flat-base-change}와
\ref{algebra-lemma-flat-localization}을 보라.
\end{proof}

\begin{lemma}
\label{more-algebra-lemma-flat-descent-flat-at-primes}
상황 \ref{more-algebra-situation-flattening-general}에서 $R' \to R''$을 $R$-대수 사상이라 하자.
$I' \subset R'$과 $I'R'' \subset I'' \subset R''$을 아이디얼이라 하자.
다음을 가정하자.
\begin{enumerate}
\item $\Spec(R'') \to \Spec(R')$이 유도하는 사상 $V(I'') \to V(I')$는 전사이다.
\item 모든 소아이디얼 $\mathfrak p'' \in V(I'')$에 대하여
$R''_{\mathfrak p''}$는 $R'$ 위 평탄이다.
\end{enumerate}
$(R'', I'')$에 대하여 (\ref{more-algebra-equation-flat-at-primes})가 성립하면,
$(R', I')$에 대해서도 (\ref{more-algebra-equation-flat-at-primes})가 성립한다.
\end{lemma}

\begin{proof}
$(R'', I'')$에 대하여 (\ref{more-algebra-equation-flat-at-primes})가 성립한다고 가정하자.
소아이디얼 $I'S' + JS' \subset \mathfrak q' \subset S'$를 택한다.
$I' \subset \mathfrak p' \subset R'$을 이에 대응하는 $R'$의 소아이디얼이라 하자.
가정에 의해 $\mathfrak p'$ 위에 놓이는 $R''$의 소아이디얼
$\mathfrak p'' \in V(I'')$가 존재하고,
$R'_{\mathfrak p'} \to R''_{\mathfrak p''}$는 평탄이다.
소아이디얼 $\overline{\mathfrak q}'' \subset
\kappa(\mathfrak q') \otimes_{\kappa(\mathfrak p')} \kappa(\mathfrak p'')$를 택한다.
이는 $\mathfrak q'$ 위이면서 $\mathfrak p''$ 위에 놓이는
소아이디얼 $\mathfrak q'' \subset S'' = S' \otimes_{R'} R''$에 대응한다.
특히 $I''S'' \subset \mathfrak q''$이고 $JS'' \subset \mathfrak q''$임을 알 수 있다.
$(S' \otimes_{R'} R'')_{\mathfrak q''}$는
$S'_{\mathfrak q'} \otimes_{R'_{\mathfrak p'}} R''_{\mathfrak p''}$의 국소화임을 주목하자.
가정에 의해 가군 $(M' \otimes_{R'} R'')_{\mathfrak q''}$는
$R''_{\mathfrak p''}$ 위 평탄이다. 따라서 Algebra, Lemma
\ref{algebra-lemma-base-change-flat-up-down}은
$M'_{\mathfrak q'}$가 $R'_{\mathfrak p'}$ 위 평탄임을 함의한다.
이것이 증명하고자 한 바이다.
\end{proof}

\begin{lemma}
\label{more-algebra-lemma-limit-preserving-flat-at-primes}
상황 \ref{more-algebra-situation-flattening-general}에서 $R \to S$가 본질적으로 유한 표시이고
$M$이 유한 표시 $S$-가군이라고 가정하자.
$R' = \colim_{\lambda \in \Lambda} R_\lambda$를
$R$-대수들의 유향 쌍대극한이라 하자.
$I_\lambda \subset R_\lambda$를 모든 $\mu \geq \lambda$에 대하여
$I_\lambda R_\mu \subset I_\mu$를 만족하는 아이디얼들이라 하고
$I' = \colim_\lambda I_\lambda$로 놓자.
$(R', I')$에 대하여 (\ref{more-algebra-equation-flat-at-primes})가 성립하면,
$(R_\lambda, I_\lambda)$에 대하여 (\ref{more-algebra-equation-flat-at-primes})가 성립하는
$\lambda \in \Lambda$가 존재한다.
\end{lemma}

\begin{proof}
먼저 $R \to S$가 유한 표시인 경우에 보조정리를 증명한 다음,
일반적인 경우에 무엇을 바꿔야 하는지 설명하겠다.
$S_\lambda = S \otimes_R R_\lambda$,
$S' = S \otimes_R R'$, $M_\lambda = M \otimes_R R_\lambda$,
$M' = M \otimes_R R'$로 쓰겠다.
기저변환 $S'$는 $R'$ 위 유한 표시이고 $M'$는 $S'$ 위 유한 표시이며,
첨자 $\lambda$가 있는 경우도 마찬가지이다.
Algebra, Lemma \ref{algebra-lemma-base-change-finiteness}를 보라.
Algebra, Theorem \ref{algebra-theorem-openness-flatness}에 의해 집합
$$
U' = \{\mathfrak q' \in \Spec(S') \mid
M'_{\mathfrak q'}\text{ is flat over }R'\}
$$
는 $\Spec(S')$에서 열려 있다.
$V(I'S' + JS')$는 준콤팩트 공간이며 가정에 의해 $U'$에 포함됨을 주목하자.
따라서 $D(g'_j) \subset U'$이고
$V(I'S' + JS') \subset \bigcup D(g'_j)$인 유한 개의
$g'_j \in S'$, $j = 1, \ldots, m$이 존재한다.
특히 $(M')_{g'_j}$는 $R'$ 위 평탄 가군임을 주목하자.

\medskip\noindent
점점 더 큰 원소 $\lambda \in \Lambda$를 택하겠다.
먼저 $g'_j$로 가는 $g_{j, \lambda} \in S_{\lambda}$를 찾을 수 있을 만큼
충분히 크게 택한다. 포함 $V(I'S' + JS') \subset \bigcup D(g'_j)$는
$I'S' + JS' + (g'_1, \ldots, g'_m) = S'$를 뜻하며,
이는 어떤 $z_s \in I'$, $y_t \in J$, $k_t, h_s, f_j \in S'$를 사용하여
$$
1 = \sum y_tk_t + \sum z_sh_s + \sum f_jg'_j
$$
로 나타낼 수 있다. $\lambda$를 더 크게 하면 이러한 등식이
$S_\lambda$에서 성립한다고 가정해도 된다. 따라서
$V(I_\lambda S_\lambda + J S_\lambda) \subset \bigcup D(g_{j, \lambda})$라고
가정할 수 있다. Algebra, Lemma
\ref{algebra-lemma-flat-finite-presentation-limit-flat}에 의해 충분히 큰 어떤
$\lambda$에 대하여 가군 $(M_\lambda)_{g_{j, \lambda}}$들이
$R_\lambda$ 위 평탄임을 알 수 있다.
특히 가군 $M_\lambda$는 $V(I_\lambda S_\lambda + J S_\lambda)$의 점들에
대응하는 모든 소아이디얼에서 $R_\lambda$ 위 평탄이다.

\medskip\noindent
$S$가 본질적으로 유한 표시인 경우에는 $S = \Sigma^{-1}C$로 쓸 수 있다.
여기서 $R \to C$는 유한 표시이고 $\Sigma \subset C$는 곱셈적 부분집합이다.
또한 어떤 유한 표시 $C$-가군 $N$에 대하여 $M = \Sigma^{-1}N$으로 쓸 수 있다.
Algebra, Lemma \ref{algebra-lemma-construct-fp-module}을 보라.
이제 $C_\lambda$, $C'$, $N_\lambda$, $N'$를 도입한다.
그러면 위 논의에서 $N'$가 $R'$ 위 평탄이 되는 열린집합
$U' \subset \Spec(C')$를 얻는다.
(\ref{more-algebra-equation-flat-at-primes})가 성립한다는 가정은
$V(I'S' + JS')$가 $U'$ 안으로 사상됨을 뜻한다.
소아이디얼 $\mathfrak r' \subset C'$에 대응하는 소아이디얼
$\mathfrak q' \subset S'$에 대하여 $M'_{\mathfrak q'} = N'_{\mathfrak r'}$이기 때문이다.
따라서 $\bigcup D(g'_j)$가 $V(I'S' + JS')$의 상을 포함하도록 하는
$g'_j \in C'$를 찾을 수 있다. 나머지 증명은 앞의 경우와 완전히 같다.
\end{proof}

\begin{lemma}
\label{more-algebra-lemma-flat-module-powers-variant}
상황 \ref{more-algebra-situation-flattening-general}에서 $I \subset R$을 아이디얼이라 하자.
다음을 가정하자.
\begin{enumerate}
\item $R$은 뇌터 환이다.
\item $S$는 뇌터 환이다.
\item $M$은 유한 $S$-가군이다.
\item 각 $n \geq 1$과 임의의 소아이디얼 $\mathfrak q \in V(J + IS)$에 대하여
가군 $(M/I^n M)_{\mathfrak q}$는 $R/I^n$ 위 평탄이다.
\end{enumerate}
그러면 $(R, I)$에 대하여 (\ref{more-algebra-equation-flat-at-primes})가 성립한다.
즉 모든 소아이디얼 $\mathfrak q \in V(J + IS)$에 대하여
국소화 $M_{\mathfrak q}$는 $R$ 위 평탄이다.
\end{lemma}

\begin{proof}
$\mathfrak q \in V(J + IS)$라 하자.
그러면 Algebra, Lemma \ref{algebra-lemma-flat-module-powers}를
$R \to S_{\mathfrak q}$와 $M_{\mathfrak q}$에 적용하면
$M_{\mathfrak q}$가 $R$ 위 평탄임을 얻는다.
\end{proof}

\section{뇌터 완비 국소환 위의 평탄화}
\label{more-algebra-section-flattening-Noetherian-complete-local}

\noindent
다음 세 보조정리는 기저가 완비 뇌터 국소환 $R$이고 주어진 가군이
유한형 $R$-대수 $S$ 위 유한일 때, ``국소적'' 평탄화 층화가 존재한다는
사실의 순전히 대수적인 증명을 준다.

\begin{lemma}
\label{more-algebra-lemma-flattening-complete-local-noetherian}
$R \to S$를 환 사상이라 하고 $M$을 $S$-가군이라 하자.
다음을 가정하자.
\begin{enumerate}
\item $(R, \mathfrak m)$은 완비 뇌터 국소환이다.
\item $S$는 뇌터 환이다.
\item $M$은 $S$ 위 유한이다.
\end{enumerate}
그러면 다음을 만족하는 아이디얼 $I \subset \mathfrak m$이 존재한다.
\begin{enumerate}
\item $\mathfrak m$ 위에 놓이는 $S/IS$의 모든 소아이디얼 $\mathfrak q$에 대하여
$(M/IM)_{\mathfrak q}$는 $R/I$ 위 평탄이다.
\item $J \subset R$이 $\mathfrak m$ 위에 놓이는 모든 소아이디얼 $\mathfrak q$에
대하여 $(M/JM)_{\mathfrak q}$를 $R/J$ 위 평탄하게 하는 아이디얼이면,
$I \subset J$이다.
\end{enumerate}
다시 말해 $I$는
$(\overline{R} \to \overline{S}, \overline{\mathfrak m}, \overline{M})$에 대하여
(\ref{more-algebra-equation-flat-at-primes-over})가 성립하도록 하는 $R$의 가장 작은 아이디얼이다.
여기서 $\overline{R} = R/I$, $\overline{S} = S/IS$,
$\overline{\mathfrak m} = \mathfrak m/I$, $\overline{M} = M/IM$이다.
\end{lemma}

\begin{proof}
$J \subset R$을 아이디얼이라 하자.
Algebra, Lemma \ref{algebra-lemma-flat-module-powers}를
환 $R/J$ 위의 가군 $M/JM$에 적용한다.
% P02-E0805: frozen source quantifies over all primes of S/JS here; do not insert the missing restriction in rendered prose.
그러면 $S/JS$의 모든 소아이디얼 $\mathfrak q$에 대하여
$(M/JM)_{\mathfrak q}$가 $R/J$ 위 평탄일 필요충분조건은
모든 $n \geq 1$에 대하여 $M/(J + \mathfrak m^n)M$이
$R/(J + \mathfrak m^n)$ 위 평탄인 것임을 알 수 있다.
아래에서 이 관찰을 사용하겠다.

\medskip\noindent
모든 $n \geq 1$에 대하여 국소환 $R/\mathfrak m^n$은 아르틴 환이다.
따라서 보조정리 \ref{more-algebra-lemma-flattening-artinian}에 의해
$M/I_nM$이 $R/I_n$ 위 평탄이 되게 하는 가장 작은 아이디얼
$I_n \supset \mathfrak m^n$이 존재한다.
% P02-E0806: frozen source says “is contains”; Korean uses one verb.
$I_{n + 1} + \mathfrak m^n$이 $I_n$을 포함함은 명백하고,
보조정리 \ref{more-algebra-lemma-intersection-flat}을 적용하면
$I_n = I_{n + 1} + \mathfrak m^n$임을 알 수 있다.
$R = \lim_n\ R/\mathfrak m^n$이므로
$I = \lim_n\ I_n/\mathfrak m^n$은 모든 $n \geq 1$에 대하여
$I_n = I + \mathfrak m^n$을 만족하는 $R$의 아이디얼임을 알 수 있다.
증명 첫머리의 관찰에 의해 $I$가 (1)과 (2)를 만족함을 알 수 있다.
몇몇 세부사항은 생략한다.
\end{proof}

\begin{lemma}
\label{more-algebra-lemma-flattening-complete-local-noetherian-property-by-finite-type}
$R \to S$, $M$, $I$의 기호와 가정은 보조정리
\ref{more-algebra-lemma-flattening-complete-local-noetherian}에서와 같다고 하자.
$R'$이 뇌터 환인 국소환들의 국소 준동형
$\varphi : (R, \mathfrak m) \to (R', \mathfrak m')$를 생각하자.
그러면 다음은 동치이다.
\begin{enumerate}
\item $(R' \to S \otimes_R R', \mathfrak m', M \otimes_R R')$에 대하여
조건 (\ref{more-algebra-equation-flat-at-primes-over})가 성립한다.
\item $\varphi(I) = 0$이다.
\end{enumerate}
\end{lemma}

\begin{proof}
함의 (2) $\Rightarrow$ (1)은 보조정리
\ref{more-algebra-lemma-base-change-flat-at-primes-over}에서 따른다.
$\varphi : R \to R'$이 보조정리에서와 같이 (1)을 만족한다고 하자.
$\varphi(I) = 0$임을 보여야 한다.
이는 $\varphi(I)R' = 0$이라는 조건과 동치이다.
뇌터 국소환 $R'$에서의 Artin-Rees 정리
(Algebra, Lemma \ref{algebra-lemma-intersect-powers-ideal-module-zero}를 보라)에 의해
이는 모든 $n > 0$에 대하여
$\varphi(I)R' + (\mathfrak m')^n = (\mathfrak m')^n$이라는 조건과 동치이다.
따라서 이는 각 $n$에 대하여 합성
$\varphi_n : R \to R' \to R'/(\mathfrak m')^n$이 $I$를 소멸시키는 조건과 동치이다.
이제 $\varphi$에 대한 가정 (1)은 보조정리
\ref{more-algebra-lemma-base-change-flat-at-primes-over}에 의해 $\varphi_n$에 대한
가정 (1)을 함의한다. 이로써 $R'$이 아르틴 국소환인 경우로 환원된다.

\medskip\noindent
% P02-E0807: frozen source omits “is” in “Assume R' Artinian”; Korean supplies the copula.
$R'$이 아르틴 환이라고 가정하자. $J = \Ker(\varphi)$로 놓자.
$I \subset J$임을 보여야 한다.
보조정리 \ref{more-algebra-lemma-flattening-complete-local-noetherian}에서의 $I$의 구성에 의해,
$\mathfrak m$ 위에 놓이는 $S/JS$의 모든 소아이디얼 $\mathfrak q$에 대하여
$(M/JM)_{\mathfrak q}$가 $R/J$ 위 평탄임을 보이면 충분하다.
$R'$이 아르틴 환이므로 조건 (1)은 $M \otimes_R R'$이 $R'$ 위 평탄임을 뜻한다.
$R'$이 아르틴 환이고 $R/J \to R'$이 국소 단사 환 사상이므로
$R/J$도 아르틴 환이다. 따라서
$M \otimes_R R' = M/JM \otimes_{R/J} R'$이 $R'$ 위 평탄이라는 사실은
Algebra, Lemma
\ref{algebra-lemma-descent-flatness-injective-map-artinian-rings}에 의해
$M/JM$이 $R/J$ 위 평탄임을 함의한다. 이로써 증명이 끝난다.
\end{proof}

\begin{lemma}
\label{more-algebra-lemma-flattening-complete-local-universal-property}
$R \to S$, $M$, $I$의 기호와 가정은 보조정리
\ref{more-algebra-lemma-flattening-complete-local-noetherian}에서와 같다고 하자.
또한 $R \to S$가 유한형이라고 가정하자.
그러면 임의의 국소환들의 국소 준동형
$\varphi : (R, \mathfrak m) \to (R', \mathfrak m')$에 대하여 다음은 동치이다.
\begin{enumerate}
\item $(R' \to S \otimes_R R', \mathfrak m', M \otimes_R R')$에 대하여
조건 (\ref{more-algebra-equation-flat-at-primes-over})가 성립한다.
\item $\varphi(I) = 0$이다.
\end{enumerate}
\end{lemma}

\begin{proof}
함의 (2) $\Rightarrow$ (1)은 보조정리
\ref{more-algebra-lemma-base-change-flat-at-primes-over}에서 따른다.
$\varphi : R \to R'$이 보조정리에서와 같이 (1)을 만족한다고 하자.
$R$이 뇌터 환이므로 $R \to S$는 유한 표시이고 $M$은 유한 표시
$S$-가군임을 알 수 있다.
$R' = \colim_\lambda R_\lambda$를 국소 $R$-부분대수
$R_\lambda \subset R'$들의 유향 쌍대극한으로 쓰자.
이들의 극대아이디얼은 $\mathfrak m_\lambda = R_\lambda \cap \mathfrak m'$이고,
각 $R_\lambda$는 $R$ 위 본질적으로 유한형이라고 하자.
보조정리 \ref{more-algebra-lemma-limit-preserving-flat-at-primes-over}에 의해 어떤 $\lambda$에 대하여
$(R_\lambda \to S \otimes_R R_\lambda, \mathfrak m_\lambda,
M \otimes_R R_\lambda)$에 대한 조건 (\ref{more-algebra-equation-flat-at-primes-over})가
성립함을 알 수 있다.
따라서 보조정리 \ref{more-algebra-lemma-flattening-complete-local-noetherian-property-by-finite-type}을
환 사상 $R \to R_\lambda$에 적용하면 $I$가 $R_\lambda$에서 영으로
사상됨을 알 수 있고, 따라서 당연히 $R'$에서도 영으로 사상된다.
\end{proof}

\section{정수적 사상을 따른 평탄성의 내림}
\label{more-algebra-section-descent-flatness-integral}

\noindent
먼저 몇 가지 간단한 보조정리를 제시한다.

\begin{lemma}
\label{more-algebra-lemma-have-one-root}
$R$을 환이라 하자. $P(T)$를 계수가 $R$에 속하는 모닉 다항식이라 하자.
$\alpha \in R$이 $P(\alpha) = 0$을 만족한다고 하자.
그러면 어떤 모닉 다항식 $Q(T) \in R[T]$에 대하여
$P(T) = (T - \alpha)Q(T)$이다.
\end{lemma}

\begin{proof}
$P$의 차수에 관한 귀납법을 사용한다.
$\deg(P) = 1$이면 $P(T) = T - \alpha$이므로 결과가 성립한다.
$\deg(P) > 1$이면 다항식 나눗셈에 의해 차수가 $< \deg(P)$인
어떤 다항식 $Q \in R[T]$와 어떤 $r \in R$에 대하여
$P(T) = (T - \alpha)Q(T) + r$로 쓸 수 있다.
가정에 의해 $0 = P(\alpha) = (\alpha - \alpha)Q(\alpha) + r = r$이므로
원하는 대로 $r = 0$이라고 결론짓는다.
\end{proof}

\begin{lemma}
\label{more-algebra-lemma-adjoin-one-root}
$R$을 환이라 하자. $P(T)$를 계수가 $R$에 속하는 모닉 다항식이라 하자.
어떤 $\alpha \in R'$과 어떤 모닉 다항식 $Q(T) \in R'[T]$에 대하여
$P(T) = (T - \alpha)Q(T)$가 되게 하는 유한 자유 환 사상 $R \to R'$이 존재한다.
\end{lemma}

\begin{proof}
$P(T) = T^d + a_1T^{d - 1} + \ldots + a_0$로 쓰자.
$R' = R[x]/(x^d + a_1x^{d - 1} + \ldots + a_0)$로 놓는다.
$\alpha$를 $x$의 합동류로 놓는다.
그러면 $P(\alpha) = 0$임은 명백하다.
따라서 보조정리 \ref{more-algebra-lemma-have-one-root}에 의해 결론을 얻는다.
\end{proof}

\begin{lemma}
\label{more-algebra-lemma-finite-split}
$R \to S$를 유한 환 사상이라 하자.
$S \otimes_R R'$이 다음 꼴의 환의 몫환이 되게 하는
유한 자유 환 확대 $R \subset R'$이 존재한다.
$$
R'[T_1, \ldots, T_n]/(P_1(T_1), \ldots, P_n(T_n))
$$
여기서 어떤 $\alpha_{ij} \in R'$에 대하여
$P_i(T) = \prod_{j = 1, \ldots, d_i} (T - \alpha_{ij})$이다.
\end{lemma}

\begin{proof}
$x_1, \ldots, x_n \in S$를 $R$ 위에서 $S$의 생성원들이라 하자.
각 $i$에 대하여 $S$에서 $P_i(x_i) = 0$이 되게 하는
모닉 다항식 $P_i(T) \in R[T]$를 택할 수 있다.
Algebra, Lemma \ref{algebra-lemma-finite-is-integral}을 보라.
$\deg(P_i) = d_i$로 쓰자.
보조정리 \ref{more-algebra-lemma-adjoin-one-root}를 $\sum d_i$번 적용하면,
각 $P_i$가 완전히 분해되는 유한 자유 환 확대 $R \subset R'$이 존재한다.
$$
P_i(T) = \prod\nolimits_{j = 1, \ldots, d_i} (T - \alpha_{ij})
$$
% P02-E0808: frozen source switches alpha_ij to unbound alpha_ik; retain the exact symbol.
여기서 $\alpha_{ik} \in R'$이다.
$R'[T_1, \ldots, T_n] \to S \otimes_R R'$을 $T_i$를 $x_i \otimes 1$로
보내는 $R'$-대수 사상이라 하자.
이는 $P_i(T_i)$를 영으로 보내므로 원하는 전사를 유도한다.
\end{proof}

\begin{lemma}
\label{more-algebra-lemma-split-image}
$R$을 환이라 하자. $S = R[T_1, \ldots, T_n]/J$라 하자.
$J$가 $P_i(T_i)$ 꼴의 원소들을 포함한다고 가정하자.
여기서 어떤 $\alpha_{ij} \in R$에 대하여
$P_i(T) = \prod_{j = 1, \ldots, d_i} (T - \alpha_{ij})$이다.
$1 \leq k_i \leq d_i$인 $\underline{k} = (k_1, \ldots, k_n)$에 대하여
환 사상
$$
\Phi_{\underline{k}} : R[T_1, \ldots, T_n] \to R,
\quad
T_i \longmapsto \alpha_{ik_i}
$$
를 생각하자. $J_{\underline{k}} = \Phi_{\underline{k}}(J)$로 놓는다.
그러면 $\Spec(S) \to \Spec(R)$의 상은 $V(\bigcap J_{\underline{k}})$와 같다.
\end{lemma}

\begin{proof}
이 보조정리는 자명하다. 힌트:
$V(\bigcap J_{\underline{k}}) = \bigcup V(J_{\underline{k}})$.
\end{proof}

\noindent
다음 결과는 Ferrand의 것이다. \cite{Ferrand}를 보라.

\begin{lemma}
\label{more-algebra-lemma-descent-flatness-injective-finite-Noetherian-rings}
$R \to S$를 뇌터 환들의 유한 단사 준동형이라 하자.
$M$을 $R$-가군이라 하자. $M \otimes_R S$가 평탄 $S$-가군이면
$M$은 평탄 $R$-가군이다.
\end{lemma}

\begin{proof}
$M$을 $M \otimes_R S$가 $S$ 위 평탄인 $R$-가군이라 하자.
Algebra, Lemma \ref{algebra-lemma-flatness-descends}에 의해,
$M$이 평탄임을 증명하려면 $R$을 임의의 충실 평탄 환 확대로 바꾸어도 된다.
보조정리 \ref{more-algebra-lemma-finite-split}에 의해 유한 국소 자유 환 확대
$R \subset R'$로서 $S' = S \otimes_R R' = R'[T_1, \ldots, T_n]/J$가 되는 것을
찾을 수 있다. 여기서 어떤 $\alpha_{ij} \in R'$에 대하여
$P_i(T) = \prod_{j = 1, \ldots, d_i} (T - \alpha_{ij})$인
$P_i(T_i)$ 꼴의 원소들을 포함하는 아이디얼
$J \subset R'[T_1, \ldots, T_n]$이 있다.
% P02-E0809: frozen source has doubled spacing before “elements”; Korean normalizes prose spacing.
$R'$이 뇌터 환이고 $R' \subset S'$이 환들의 유한 확대임을 주목하자.
따라서 $R$을 $R'$로 바꾸고 $S$가 보조정리 \ref{more-algebra-lemma-split-image}에서와 같은
표시를 갖는다고 가정해도 된다.
$\Spec(S) \to \Spec(R)$이 전사임을 주목하자.
Algebra, Lemma \ref{algebra-lemma-integral-overring-surjective}를 보라.
따라서 보조정리 \ref{more-algebra-lemma-split-image}를 사용하면
$I = \bigcap J_{\underline{k}}$는 $V(I) = \Spec(R)$을 만족하는 아이디얼이라고
결론짓는다. 이는 $I \subset \sqrt{(0)}$임을 뜻하며,
$R$이 뇌터 환이므로 $I$는 멱영이다.
사상 $\Phi_{\underline{k}}$들은 가환도식
$$
\xymatrix{
S \ar[rr] & & R/J_{\underline{k}} \\
& R \ar[lu] \ar[ru]
}
$$
을 유도한다. 이로부터 $M/J_{\underline{k}}M$이 $R/J_{\underline{k}}$ 위
평탄이라고 결론짓는다.
보조정리 \ref{more-algebra-lemma-intersection-flat}에 의해 $M/IM$이 $R/I$ 위 평탄임을
알 수 있다. 마지막으로 Algebra, Lemma \ref{algebra-lemma-lift-flatness}를
적용하여 $M$이 $R$ 위 평탄이라고 결론짓는다.
\end{proof}

\begin{lemma}
\label{more-algebra-lemma-descent-flatness-injective-integral}
$R \to S$를 단사 정수적 환 사상이라 하자.
$M$을 $R[x_1, \ldots, x_n]$ 위의 유한 표시 가군이라 하자.
$M \otimes_R S$가 $S$ 위 평탄이면 $M$은 $R$ 위 평탄이다.
\end{lemma}

\begin{proof}
표시
$$
R[x_1, \ldots, x_n]^{\oplus t} \to R[x_1, \ldots, x_n]^{\oplus r} \to M \to 0.
$$
를 택한다. 첫 번째 사상이 $f_{ij} \in R[x_1, \ldots, x_n]$인
$r \times t$ 행렬 $T = (f_{ij})$로 주어진다고 하자.
$f_{ij, I} \in R$을 사용하여 $f_{ij} = \sum f_{ij, I} x^I$로 쓴다
(다중지표 표기). 도식
$$
\xymatrix{
R \ar[r] & S \\
R_\lambda \ar[u] \ar[r] & S_\lambda \ar[u]
}
$$
을 생각하자. 여기서 $R_\lambda$는 모든 $f_{ij, I}$를 포함하는
$R$의 유한 생성 $\mathbf{Z}$-부분대수이고,
$S_\lambda$는 $S$의 유한 $R_\lambda$-부분대수이다.
$M_\lambda$를 위와 같은 표시로 정의되는 유한
$R_\lambda[x_1, \ldots, x_n]$-가군이라 하자.
같은 행렬 $T$를 사용하되 이제 이를 $R_\lambda[x_1, \ldots, x_n]$ 위의
행렬로 본다. $S$가 $S_\lambda$들의 유향 쌍대극한임을 주목하자
(세부사항은 생략한다).
Algebra, Lemma \ref{algebra-lemma-flat-finite-presentation-limit-flat}에 의해
어떤 $\lambda$에 대하여 가군 $M_\lambda \otimes_{R_\lambda} S_\lambda$가
$S_\lambda$ 위 평탄임을 알 수 있다.
보조정리 \ref{more-algebra-lemma-descent-flatness-injective-finite-Noetherian-rings}에 의해
$M_\lambda$가 $R_\lambda$ 위 평탄이라고 결론짓는다.
$M = M_\lambda \otimes_{R_\lambda} R$이므로
Algebra, Lemma \ref{algebra-lemma-flat-base-change}에 의해 결론을 얻는다.
\end{proof}

\begin{lemma}
\label{more-algebra-lemma-descent-projective-injective-finite-Noetherian-rings}
$R \to S$를 뇌터 환들의 유한 단사 준동형이라 하자.
$P$를 $R$-가군이라 하자. $P \otimes_R S$가 사영 $S$-가군이면
$P$는 사영 $R$-가군이다.
\end{lemma}

\begin{proof}
$P$를 $P \otimes_R S$가 $S$ 위 사영인 $R$-가군이라 하자.
Algebra, Theorem \ref{algebra-theorem-ffdescent-projectivity}에 의해,
$P$가 사영임을 증명하려면 $R$을 임의의 충실 평탄 환 확대로 바꾸어도 된다.
보조정리 \ref{more-algebra-lemma-finite-split}에 의해 유한 국소 자유 환 확대
$R \subset R'$로서 $S' = S \otimes_R R' = R'[T_1, \ldots, T_n]/J$가 되는 것을
찾을 수 있다. 여기서 어떤 $\alpha_{ij} \in R'$에 대하여
$P_i(T) = \prod_{j = 1, \ldots, d_i} (T - \alpha_{ij})$인
$P_i(T_i)$ 꼴의 원소들을 포함하는 아이디얼
$J \subset R'[T_1, \ldots, T_n]$이 있다.
% P02-E0809: the parallel frozen proof repeats the same doubled spacing; Korean normalizes prose spacing.
$R'$이 뇌터 환이고 $R' \subset S'$이 환들의 유한 확대임을 주목하자.
따라서 $R$을 $R'$로 바꾸고 $S$가 보조정리 \ref{more-algebra-lemma-split-image}에서와 같은
표시를 갖는다고 가정해도 된다.
$\Spec(S) \to \Spec(R)$이 전사임을 주목하자.
Algebra, Lemma \ref{algebra-lemma-integral-overring-surjective}를 보라.
따라서 보조정리 \ref{more-algebra-lemma-split-image}를 사용하면
$I = \bigcap J_{\underline{k}}$는 $V(I) = \Spec(R)$을 만족하는 아이디얼이라고
결론짓는다. 이는 $I \subset \sqrt{(0)}$임을 뜻하며,
$R$이 뇌터 환이므로 $I$는 멱영이다.
사상 $\Phi_{\underline{k}}$들은 가환도식
$$
\xymatrix{
S \ar[rr] & & R/J_{\underline{k}} \\
& R \ar[lu] \ar[ru]
}
$$
을 유도한다. 이로부터 $P/J_{\underline{k}}P$가 $R/J_{\underline{k}}$ 위
사영이라고 결론짓는다.
Algebra, Lemma \ref{algebra-lemma-projective-modulo-two-ideals}에 의해
$P/IP$가 $R/I$ 위 사영임을 알 수 있다.
보조정리 \ref{more-algebra-lemma-descent-flatness-injective-finite-Noetherian-rings}에 의해
$P$가 $R$ 위 평탄이므로, Algebra, Lemma \ref{algebra-lemma-lift-projective}를
적용하여 $P$가 $R$ 위 사영이라고 결론 내릴 수 있다.
\end{proof}

\section{비틀림 없는 가군}
\label{more-algebra-section-torsion-flat}

\noindent
이 절에서는 비틀림 없는 가군과 평탄성 사이의 관계를,
특히 차원 1인 환 위에서 논의한다.

\begin{definition}
\label{more-algebra-definition-torsion}
$R$을 정역이라 하자. $M$을 $R$-가군이라 하자.
\begin{enumerate}
\item $fx = 0$인 영이 아닌 $f \in R$이 존재하면
원소 $x \in M$을 {\it 비틀림 원소}라 한다.
\item $M$의 비틀림 원소가 $0$뿐이면 $M$에 {\it 비틀림이 없다}고 한다.
\item $M$의 모든 원소가 비틀림 원소이면 $M$을 {\it 비틀림 가군}이라 한다.
\end{enumerate}
\end{definition}

\noindent
$R$을 정역이라 하고 $S = R \setminus \{0\}$을 $R$의 영이 아닌 원소들의
곱셈적 집합이라 하자. 그러면 $R$-가군 $M$에 비틀림이 없을 필요충분조건은
$M \to S^{-1}M$이 단사인 것이다. 다시 말해 이는
$K = S^{-1}R$을 $R$의 분수체라 할 때 사상 $M \to M \otimes_R K$가
단사인 것과 동치이다.

\begin{lemma}
\label{more-algebra-lemma-torsion}
$R$을 정역이라 하고 $M$을 $R$-가군이라 하자.
$M$의 비틀림 원소들의 집합 $M_{tors}$는 사상 $M \to M \otimes_R K$의 핵이다.
따라서 $M_{tors}$는 $M$의 $R$-부분가군이다.
몫가군 $M/M_{tors}$에는 비틀림이 없다.
\end{lemma}

\begin{proof}
위의 논의를 보라.
\end{proof}

\begin{lemma}
\label{more-algebra-lemma-localize-torsion}
$R$을 정역이라 하자. $M$을 비틀림 없는 $R$-가군이라 하자.
임의의 곱셈적 집합 $S \subset R$에 대하여
가군 $S^{-1}M$은 비틀림 없는 $S^{-1}R$-가군이다.
\end{lemma}

\begin{proof}
생략한다.
\end{proof}

\begin{lemma}
\label{more-algebra-lemma-flat-pullback-torsion}
$R \to R'$을 정역들의 평탄 준동형이라 하자.
$M$이 비틀림 없는 $R$-가군이면 $M \otimes_R R'$은 비틀림 없는 $R'$-가군이다.
\end{lemma}

\begin{proof}
$M$에 비틀림이 없으면 $M \subset M \otimes_R K$는 단사이다.
여기서 $K$는 $R$의 분수체이다.
$R'$이 $R$ 위 평탄이므로
$M \otimes_R R' \to (M \otimes_R K) \otimes_R R'$은 단사임을 알 수 있다.
$M \otimes_R K$는 $K$의 복사본들의 직접합과 동형이므로
$K \otimes_R R'$에 비틀림이 없음을 보이면 충분하다.
이는 $R'$의 국소화이므로 성립한다.
\end{proof}

\begin{lemma}
\label{more-algebra-lemma-extension-torsion-free}
$R$을 정역이라 하자. $0 \to M \to M' \to M'' \to 0$을
$R$-가군들의 짧은 완전열이라 하자.
$M$과 $M''$에 비틀림이 없으면 $M'$에도 비틀림이 없다.
\end{lemma}

\begin{proof}
생략한다.
\end{proof}

\begin{lemma}
\label{more-algebra-lemma-check-torsion}
$R$을 정역이라 하고 $M$을 $R$-가군이라 하자.
그러면 $M$에 비틀림이 없을 필요충분조건은 $R$의 모든 극대아이디얼
$\mathfrak m$에 대하여 $M_\mathfrak m$이 비틀림 없는
$R_\mathfrak m$-가군인 것이다.
\end{lemma}

\begin{proof}
생략한다. 힌트: 보조정리 \ref{more-algebra-lemma-localize-torsion}과
Algebra, Lemma \ref{algebra-lemma-characterize-zero-local}을 사용하라.
\end{proof}

\begin{lemma}
\label{more-algebra-lemma-finite-torsion-free-submodule-free}
$R$을 정역이라 하자. $M$을 유한 $R$-가군이라 하자.
그러면 $M$에 비틀림이 없을 필요충분조건은
$M$이 유한 자유 가군의 부분가군인 것이다.
\end{lemma}

\begin{proof}
$M$이 $R^{\oplus n}$의 부분가군이면 $M$에는 비틀림이 없다.
거꾸로 $M$에 비틀림이 없다고 가정하자.
$K$를 $R$의 분수체라 하자.
그러면 $M \otimes_R K$는 유한 차원 $K$-벡터 공간이다.
이 벡터 공간의 기저 $e_1, \ldots, e_r$를 택한다.
$x_1, \ldots, x_n$을 $M$의 생성원들이라 하자.
$b_{ij} \not = 0$인 어떤 $a_{ij}, b_{ij} \in R$을 사용하여
$x_i = \sum (a_{ij}/b_{ij}) e_j$로 쓰자.
$b = \prod_{i, j} b_{ij}$로 놓는다.
$M$에 비틀림이 없으므로 사상 $M \to M \otimes_R K$는 단사이고,
그 상은 $R^{\oplus r} = R e_1/b \oplus \ldots \oplus Re_r/b$에 포함된다.
\end{proof}

\begin{lemma}
\label{more-algebra-lemma-torsion-free-finite-noetherian-domain}
$R$을 뇌터 정역이라 하자. $M$을 영이 아닌 유한 $R$-가군이라 하자.
다음은 동치이다.
\begin{enumerate}
\item $M$에는 비틀림이 없다.
\item $M$은 유한 자유 가군의 부분가군이다.
\item $(0)$은 $M$의 유일한 동반 소아이디얼이다.
\item $(0)$이 $M$의 지지집합에 속하고 $M$은 성질 $(S_1)$을 갖는다.
\item $(0)$이 $M$의 지지집합에 속하고 $M$은 매장된 동반 소아이디얼을 갖지 않는다.
\end{enumerate}
\end{lemma}

\begin{proof}
(1)과 (2)의 동치는 보조정리 \ref{more-algebra-lemma-finite-torsion-free-submodule-free}에서 보았다.
(4)와 (5)의 동치는 Algebra, Lemma \ref{algebra-lemma-criterion-no-embedded-primes}에서 보았다.
(3)과 (5)의 동치는 정의에서 바로 따른다.
비틀림 없는 가군의 국소화에는 비틀림이 없으므로
(보조정리 \ref{more-algebra-lemma-localize-torsion}),
% P02-E0810: frozen source has a spurious article before M; Korean supplies the intended subject.
$M$이 $(0)$과 다른 동반 소아이디얼을 갖지 않음은 명백하다.
따라서 (1)은 (5)를 함의한다. 거꾸로 (5)를 가정하자.
$M$에 비틀림이 있으면 $R$의 어떤 영이 아닌 아이디얼 $I$에 대하여
매장 $R/I \subset M$이 존재한다.
따라서 $M$은 $(0)$과 다른 동반 소아이디얼을 갖는다
(Algebra, Lemmas \ref{algebra-lemma-ass}와 \ref{algebra-lemma-ass-zero}를 보라).
이는 매장된 동반 소아이디얼이므로 가정과 모순이다.
\end{proof}

\begin{lemma}
\label{more-algebra-lemma-flat-torsion-free}
$R$을 정역이라 하자. 임의의 평탄 $R$-가군에는 비틀림이 없다.
\end{lemma}

\begin{proof}
$x \in R$이 영이 아니면 $x : R \to R$은 단사이다.
따라서 $M$이 $R$ 위 평탄이면 $x : M \to M$은 단사이다.
그러므로 $M$이 $R$ 위 평탄이면 $M$에는 비틀림이 없다.
\end{proof}

\begin{lemma}
\label{more-algebra-lemma-valuation-ring-torsion-free-flat}
$A$를 값매김환이라 하자.
$A$-가군 $M$이 $A$ 위 평탄일 필요충분조건은 $M$에 비틀림이 없는 것이다.
\end{lemma}

\begin{proof}
함의 ``평탄 $\Rightarrow$ 비틀림 없음''은 보조정리 \ref{more-algebra-lemma-flat-torsion-free}이다.
거꾸로 $M$에 비틀림이 없다고 가정하자.
평탄성의 방정식 판정법(Algebra, Lemma \ref{algebra-lemma-flat-eq}를 보라)에 의해
$M$의 모든 관계식이 자명함을 보여야 한다.
이를 위해 $x_i \in M$, $a_i \in A$에 대하여
$\sum_{i = 1, \ldots, n} a_i x_i = 0$이라고 가정하자.
번호를 다시 매기면 모든 $i$에 대하여 $v(a_1) \leq v(a_i)$라고 가정해도 된다.
따라서 어떤 $a'_i \in A$에 대하여 $a_i = a'_i a_1$로 쓸 수 있다.
$a'_1 = 1$임을 주목하자.
$M$에 비틀림이 없으므로 $x_1 = - \sum_{i \geq 2} a'_i x_i$임을 알 수 있다.
따라서 $y_i = x_i$, $i = 2, \ldots, n$으로 택하면
$$
x_1 = \sum\nolimits_{j \geq 2} -a'_j y_j, \quad
x_i = y_i, (i \geq 2)\quad
0 = a_1 \cdot (-a'_j) + a_j \cdot 1 (j \geq 2)
$$
은 관계식이 자명했음을 보인다. 명시적으로 쓰면 원소 $a_{ij}$들은
$a_{11} = 0$, $j > 1$에 대하여 $a_{1j} = -a'_j$,
$i, j \geq 2$에 대하여 $a_{ij} = \delta_{ij}$로 정의한다.
\end{proof}

\begin{lemma}
\label{more-algebra-lemma-dedekind-torsion-free-flat}
$A$를 데데킨트 정역이라 하자
(예를 들어 이산 값매김환, 더 일반적으로는 주아이디얼 정역).
\begin{enumerate}
\item $A$-가군이 평탄일 필요충분조건은 비틀림이 없는 것이다.
\item 유한 비틀림 없는 $A$-가군은 유한 국소 자유이다.
\item $A$가 주아이디얼 정역이면 유한 비틀림 없는 $A$-가군은 유한 자유이다.
\end{enumerate}
\end{lemma}

\begin{proof}
(보조정리의 명제에서 괄호 안의 설명에 대해서는
Algebra, Lemma \ref{algebra-lemma-PID-dedekind}를 보라.) (1)의 증명.
보조정리 \ref{more-algebra-lemma-check-torsion}과
Algebra, Lemma \ref{algebra-lemma-flat-localization}에 의해
극대아이디얼 $\mathfrak m \subset A$에 대하여 $A_\mathfrak m$ 위에서
명제를 확인하면 충분하다.
$A_\mathfrak m$은 이산 값매김환이므로
(Algebra, Lemma \ref{algebra-lemma-characterize-Dedekind}),
보조정리 \ref{more-algebra-lemma-valuation-ring-torsion-free-flat}에 의해 결론을 얻는다.

\medskip\noindent
(2)의 증명. Algebra, Lemma \ref{algebra-lemma-finite-projective}와 (1)에서 따른다.

\medskip\noindent
(3)의 증명. $A$를 주아이디얼 정역이라 하고 $M$을 유한 비틀림 없는 가군이라 하자.
보조정리 \ref{more-algebra-lemma-finite-torsion-free-submodule-free}에 의해
어떤 $n$에 대하여 $M \subset A^{\oplus n}$임을 알 수 있다.
$n$에 관한 귀납법으로 $M$이 자유임을 보이자.
$n = 1$인 경우는 정확히 $A$가 주아이디얼 정역이라는 사실을 나타낸다.
$n > 1$이면 $M' \subset A^{\oplus n - 1}$을
$A^{\oplus n}$의 마지막 $n - 1$개 직접합인자로의 사영의 상이라 하자.
그러면 짧은 완전열 $0 \to I \to M \to M' \to 0$을 얻는다.
여기서 $I$는 $M$과 $A^{\oplus n}$의 첫 번째 직접합인자 $A$의 교집합이다.
귀납법에 의해 $M$이 유한 자유 $A$-가군들의 확대임을 알 수 있으므로
유한 자유이다.
\end{proof}

\begin{lemma}
\label{more-algebra-lemma-hom-into-torsion-free}
$R$을 정역이라 하자. $M$, $N$을 $R$-가군이라 하자.
$N$에 비틀림이 없으면 $\Hom_R(M, N)$에도 비틀림이 없다.
\end{lemma}

\begin{proof}
전사 $\bigoplus_{i \in I} R \to M$을 택한다.
그러면 $\Hom_R(M, N) \subset \prod_{i \in I} N$이다.
\end{proof}

\section{가군의 계수}
\label{more-algebra-section-rank}

\noindent
다음과 같이 정의한다.

\begin{definition}
\label{more-algebra-definition-rank}
$R$을 분수체가 $K$인 정역이라 하자.
$R$-가군 $M$의 {\it 계수}는 $\dim_K(M \otimes_R K)$이다.
\end{definition}

\noindent
두 정의를 모두 적용할 수 있는 경우, 이 정의는
Algebra, Definition \ref{algebra-definition-locally-free}에 나오는
계수 $r$의 국소 자유 가군이라는 개념과 충돌하지 않는다.

\begin{lemma}
\label{more-algebra-lemma-torsion-does-not-matter-rank}
$R$을 정역이라 하자. $M \to M'$이 핵과 여핵이 비틀림 가군인
$R$-가군들의 사상이면, $M$의 계수는 $M'$의 계수와 같다.
\end{lemma}

\begin{proof}
생략한다. 힌트: $K$가 $R$의 분수체이면
유도된 사상 $M \otimes_R K \to M' \otimes_R K$는 동형이다.
\end{proof}

\begin{lemma}
\label{more-algebra-lemma-rank-additive}
$R$을 정역이라 하자. $0 \to M \to M' \to M'' \to 0$을
$R$-가군들의 짧은 완전열이라 하자.
그러면 $M'$의 계수는 $M$의 계수와 $M''$의 계수의 합이다.
\end{lemma}

\begin{proof}
생략한다.
\end{proof}

\begin{lemma}
\label{more-algebra-lemma-rank-tensor}
$R$을 정역이라 하자. $M$과 $N$을 $R$-가군이라 하자.
\begin{enumerate}
\item $M \otimes_R N$의 계수는 $M$의 계수와 $N$의 계수의 곱이다.
\item $M$이 유한 표시 $R$-가군이면
$\Hom_R(M, N)$의 계수는 $M$의 계수와 $N$의 계수의 곱이다.
\end{enumerate}
\end{lemma}

\begin{proof}
(1)은 벡터 공간에 대한 같은 사실에서 따른다.
$M$이 $R$-가군으로서 유한 표시라고 가정하자.
그러면 Algebra, Lemma \ref{algebra-lemma-hom-from-finitely-presented}에 의해
$\Hom_R(M, N) \otimes_R K$는 $K$-벡터 공간으로서
$\Hom_K(M \otimes_R K, N \otimes_R K)$와 동형이다.
또한 $M \otimes_R K$는 유한 차원이므로
벡터 공간에 대한 대응하는 사실에서 결론을 얻는다.
\end{proof}

\begin{lemma}
\label{more-algebra-lemma-dominant-pullback-rank}
$R \subset R'$을 정역들의 확대라 하자.
$M$이 $R$-가군이면 $R$ 위에서 $M$의 계수는
$R'$ 위에서 $M \otimes_R R'$의 계수와 같다.
\end{lemma}

\begin{proof}
벡터 공간의 차원은 기초체의 확대에 대해 불변이므로 성립한다.
\end{proof}

\begin{lemma}
\label{more-algebra-lemma-finite-module-rank}
$R$을 정역이라 하자. $M$을 유한 $R$-가군이라 하자. 그러면 다음이 성립한다.
\begin{enumerate}
\item $M$은 유한 계수 $r \geq 0$을 갖는다.
\item 핵과 여핵이 비틀림 가군인 사상 $M \to R^{\oplus r}$이 존재한다.
\item $M_f$가 계수 $r$의 자유 가군이 되게 하는 $f \in R$, $f \not = 0$이 존재한다.
\item 여핵이 비틀림 가군인 단사 사상 $R^{\oplus r} \to M$이 존재한다.
\end{enumerate}
\end{lemma}

\begin{proof}
$M$이 유한이므로 $M \otimes_R K$는 $R$의 분수체 $K$ 위의 유한 차원
벡터 공간이다. 따라서 계수 $r$은 유한하다.

\medskip\noindent
$M$의 생성원 $x_1, \ldots, x_n$을 택한다.
$M \otimes_R K$의 기저 $e_1, \ldots, e_r$를 택한다.
% P02-E0811: frozen source doubles prose spacing and has a pre-comma space in the coefficient span; retain the formal span.
$b_{ij} \not = 0$인 어떤 $a_{ij} , b_{ij} \in R$을 사용하여
$x_i \otimes 1 = \sum (a_{ij}/b_{ij}) e_j$로 쓸 수 있다.
$b = \prod b_{ij}$로 놓고 $a_{ij}$를 $a_{ij}b/b_{ij}$로 바꾸면
$x_i \otimes 1 = \sum (a_{ij}/b) e_j$라고 가정할 수 있다.
$M$에서 관계식 $\sum r_i x_i = 0$이 성립하면 모든 $j = 1, \ldots, r$에
대하여 $R$에서 $\sum r_ia_{ij} = 0$이 성립함을 주목하자.

\medskip\noindent
(2)의 증명. 위의 마지막 관찰에 의해 $x_i$를 $\sum a_{ij} e_j$로 보내는
% P02-E0812: frozen source inserts a period between the subject and predicate; Korean joins the clause.
사상 $M \to \bigoplus Re_j$는 잘 정의된다.
$e_1, \ldots, e_r$가 $M \otimes_R K$의 기저를 이룬다는 사실은
이 사상이 $K$와 텐서하면 동형이 됨을 말해 준다.
따라서 이는 (2)에서와 같은 사상이다.

\medskip\noindent
(3)의 증명. (2)에서 구성한 사상 $M \to R^{\oplus r}$의 여핵은
유한 비틀림 가군이다. 따라서 이 여핵을 소멸시키는
$g \in R$, $g \not = 0$을 찾을 수 있다.
% P02-E0813: frozen source has a sentence fragment beginning “In other words, such that”; Korean supplies a complete sentence.
다시 말해 사상 $M_g \to R_g^{\oplus r}$은 전사이다.
그러면 분할 $M_g = N \oplus R_g^{\oplus r}$을 택할 수 있다.
이때 $N$은 $N \otimes_{R_g} K = 0$인 유한 $R_g$-가군이다.
다시 말해 $N$은 비틀림 가군이고 $R_g$ 위 유한이다.
따라서 $N$을 소멸시키는 $g' \in R$을 찾을 수 있다.
$f= gg'$로 놓으면 (3)이 성립함을 알 수 있다.

\medskip\noindent
(4)를 증명하기 위해 $M_f$의 기저 $y_1, \ldots, y_r$를 택한다.
여기서 $f$는 (3)에서와 같다.
그러면 어떤 $m_j \in M$과 $t_j \geq 0$에 대하여
$y_j = m_j/f^{t_j}$로 쓴다.
$j$번째 기저벡터를 $m_j$로 보내는 사상 $R^{\oplus r} \to M$은
(4)에서와 같은 사상이 된다. 세부사항은 생략한다.
\end{proof}

\section{반사 가군}
\label{more-algebra-section-reflexive}

\noindent
다음과 같이 정의한다.

\begin{definition}
\label{more-algebra-definition-reflexive}
$R$을 정역이라 하자. 자연스러운 사상
$$
j : M \longrightarrow \Hom_R(\Hom_R(M, R), R)
$$
이 동형이면 $R$-가군 $M$을 {\it 반사 가군}이라 한다.
이 사상은 $m \in M$을, 각 $\varphi \in \Hom_R(M, R)$을
$\varphi(m) \in R$로 보내는 사상에 대응시킨다.
\end{definition}

\noindent
더 일반적인 환에 대해서도 이 정의를 할 수 있지만 위 정의에도 이미 단점이 있다.
뇌터 정역과 유한 비틀림 없는 가군으로 제한하고,
(아마도) $R$에 몇 가지 정칙성 조건을 부과하는 것이 현명할 것이다
(예를 들어 $R$이 정규라는 조건).

\begin{lemma}
\label{more-algebra-lemma-reflexive-torsion-free}
$R$을 정역이라 하고 $M$을 $R$-가군이라 하자.
\begin{enumerate}
\item $M$이 반사 가군이면 $M$에는 비틀림이 없다.
\item $M$이 유한이면 $j : M \to \Hom_R(\Hom_R(M, R), R)$의
핵과 여핵은 비틀림 가군이다.
\item $M$이 유한이면 $j$가 단사일 필요충분조건은 $M$에 비틀림이 없는 것이다.
\end{enumerate}
\end{lemma}

\begin{proof}
(1)을 보려면 보조정리 \ref{more-algebra-lemma-hom-into-torsion-free}에 의해
$M$의 이중 쌍대에 비틀림이 없다는 사실을 사용한다.
$M$이 유한이라고 가정하자.
$R$ 위에서 $M$의 생성원 $x_1, \ldots, x_n$을 택한다.
$K$를 $R$의 분수체라 하자.
번호를 다시 매기면 어떤 $0 \leq r \leq n$에 대하여
원소 $x_1, \ldots, x_r$가 $K$ 위에서 $M \otimes_R K$의 기저로
사상된다고 가정해도 된다.
$r < i \leq n$에 대하여 $k_{ij} \in K$를 사용하여
$x_i \otimes 1 = \sum_{1 \leq j \leq r} x_j \otimes k_{ij}$로 쓰자.
$a_{ij} = ck_{ij} \in R$이고 $M$에서
% P02-E0814: retain the frozen double space after the equality sign in the formal span.
$cx_i =  \sum_{1 \leq j \leq r} a_{ij} x_j$가 되게 하는 $c \in R$을 택한다
(작은 세부사항은 생략한다). 그러면 사상들
$$
R^{\oplus r} \xrightarrow{\alpha} M \xrightarrow{\beta} R^{\oplus r}
\xrightarrow{\alpha} M
$$
을 얻으며, $\alpha \circ \beta$와 $\beta \circ \alpha$는 모두
$c$를 곱하는 사상으로 주어진다.
여기서 $\alpha$는 $x_1, \ldots, x_r$를 사용하고,
$\beta$는 $1 \leq i \leq r$일 때 $x_i$를 $i$번째 기저벡터의 $c$배로,
$r < i \leq n$일 때는 벡터 $(a_{i1}, \ldots, a_{ir})$로 보낸다.
가환도식
$$
\xymatrix{
R^{\oplus r} \ar[r] \ar[d]^{\text{id}} &
M \ar[r] \ar[d]^j &
R^{\oplus r} \ar[d]^{\text{id}} \ar[r] &
M \ar[d]^j \\
R^{\oplus r} \ar[r] &
\Hom_R(\Hom_R(M, R), R) \ar[r] &
R^{\oplus r} \ar[r] &
\Hom_R(\Hom_R(M, R), R)
}
$$
을 생각하자. 물론 아래 행에서 연속한 두 화살표의 합성은
$c$를 곱하는 사상과 같다. 이 도식은 $j$의 핵과 여핵이 $c$에 의해
소멸됨을 보이며, 이로써 (2)가 증명된다.
(3)에 대해서는 $M$에 비틀림이 없으면 $\beta$가 단사이고,
이로부터 $j$가 단사임을 주목하자.
거꾸로의 함의는 $M$의 이중 쌍대에 비틀림이 없음을 이미 보았으므로 명백하다.
\end{proof}

\begin{lemma}
\label{more-algebra-lemma-cokernel-map-double-dual-dvr}
$R$을 이산 값매김환이라 하고 $M$을 유한 $R$-가군이라 하자.
그러면 사상 $j : M \to \Hom_R(\Hom_R(M, R), R)$은 전사이다.
\end{lemma}

\begin{proof}
$M_{tors} \subset M$을 비틀림 부분가군이라 하자.
그러면 $\Hom_R(M, R) = \Hom_R(M/M_{tors}, R)$이다
(이는 임의의 정역 위에서 성립한다).
따라서 $M$에 비틀림이 없다고 가정해도 된다.
그러면 보조정리 \ref{more-algebra-lemma-dedekind-torsion-free-flat}에 의해
$M$은 자유 가군이므로 보조정리는 명백하다.
\end{proof}

\begin{lemma}
\label{more-algebra-lemma-check-reflexive}
$R$을 뇌터 정역이라 하자. $M$을 유한 $R$-가군이라 하자.
다음은 동치이다.
\begin{enumerate}
\item $M$은 반사 가군이다.
\item 모든 소아이디얼 $\mathfrak p \subset R$에 대하여
$M_\mathfrak p$는 반사 $R_\mathfrak p$-가군이다.
\item $R$의 모든 극대아이디얼 $\mathfrak m$에 대하여
$M_\mathfrak m$은 반사 $R_\mathfrak m$-가군이다.
\end{enumerate}
\end{lemma}

\begin{proof}
$j : M \to \Hom_R(\Hom_R(M, R), R)$을 소아이디얼 $\mathfrak p$에서
국소화하면, 뇌터 국소 정역 $R_\mathfrak p$ 위의 가군 $M_\mathfrak p$에
대한 대응하는 사상을 얻는다.
Algebra, Lemma \ref{algebra-lemma-hom-from-finitely-presented}를 보라.
따라서 Algebra, Lemma \ref{algebra-lemma-characterize-zero-local}에 의해
보조정리가 성립한다.
\end{proof}

\begin{lemma}
\label{more-algebra-lemma-sequence-reflexive}
$R$을 뇌터 정역이라 하자.
% P02-E0815: frozen source omits “be” in the sequence introduction; Korean supplies grammatical phrasing.
$0 \to M \to M' \to M''$을 유한 $R$-가군들의 완전열이라 하자.
$M'$이 반사 가군이고 $M''$에 비틀림이 없으면 $M$은 반사 가군이다.
\end{lemma}

\begin{proof}
임의의 유한 $R$-가군 $N$과 $N'$에 대하여 $\Hom_R(N, N')$가
유한 $R$-가군이라는 사실을 더 언급하지 않고 사용하겠다.
Algebra, Lemma \ref{algebra-lemma-ext-noetherian}을 보라.
쌍대를 취하면 열
$$
\Hom_R(M, R) \leftarrow \Hom_R(M', R) \leftarrow \Hom_R(M'', R)
$$
을 얻는다. 다시 쌍대를 취하면 가환도식
$$
\xymatrix{
\Hom_R(\Hom_R(M, R), R) \ar[r]_j &
\Hom_R(\Hom_R(M', R), R) \ar[r] &
\Hom_R(\Hom_R(M'', R), R) \\
M \ar[u] \ar[r] & M' \ar[u] \ar[r] & M'' \ar[u]
}
$$
을 얻는다. 위 행이 완전한지는 알 수 없다.
그러나 가정에 의해 가운데 수직 화살표는 동형이고 오른쪽 수직 화살표는
단사이다(보조정리 \ref{more-algebra-lemma-reflexive-torsion-free}).
$j$가 단사라고 주장한다. 이 주장을 가정하면 도식 추적으로
왼쪽 수직 화살표가 동형임을 알 수 있다. 즉 $M$은 반사 가군이다.

\medskip\noindent
주장의 증명. $Q$를 정의하는 완전열
$\Hom_R(M',R)\to \Hom_R(M,R)\to Q \to 0$을 생각하자.
Algebra, Lemma \ref{algebra-lemma-hom-from-finitely-presented}를 적용하면
$$
\Hom_K(M'\otimes_R K,K) \to
\Hom_K(M\otimes_R K,K) \to
Q\otimes_R K\to 0
$$
을 얻는다. 그런데 $M \otimes_R K \to M' \otimes_R K$는 벡터 공간들의
단사 사상이므로 분할 단사이다. 따라서 $Q \otimes_R K = 0$이다.
즉 $Q$는 비틀림 가군이다. 그러면 완전열
$$
0 \to \Hom_R(Q,R) \to
\Hom_R(\Hom_R(M,R),R) \to
\Hom_R(\Hom_R(M',R),R)
$$
을 얻으며, $Q$가 비틀림 가군이므로 $\Hom_R(Q,R)=0$이다.
\end{proof}

\begin{lemma}
\label{more-algebra-lemma-characterize-reflexive}
$R$을 뇌터 정역이라 하자. $M$을 유한 $R$-가군이라 하자.
다음은 동치이다.
\begin{enumerate}
\item $M$은 반사 가군이다.
\item $F$가 유한 자유이고 $N$에 비틀림이 없는
짧은 완전열 $0 \to M \to F \to N \to 0$이 존재한다.
\end{enumerate}
\end{lemma}

\begin{proof}
유한 자유 가군은 반사 가군임을 관찰하자.
보조정리 \ref{more-algebra-lemma-sequence-reflexive}에 의해 (2)가 (1)을 함의함을 알 수 있다.
$M$이 반사 가군이라고 가정하자.
표시 $R^{\oplus m} \to R^{\oplus n} \to \Hom_R(M, R) \to 0$을 택한다.
쌍대를 취하면 완전열
$$
0 \to \Hom_R(\Hom_R(M, R), R) \to R^{\oplus n} \to N \to 0
$$
을 얻으며, $N = \Im(R^{\oplus n} \to R^{\oplus m})$은 비틀림 없는 가군이다.
$M = \Hom_R(\Hom_R(M, R), R)$이므로 (2)에서와 같은 완전열을 얻는다.
\end{proof}

\begin{lemma}
\label{more-algebra-lemma-flat-pullback-reflexive}
$R \to R'$을 뇌터 정역들의 평탄 준동형이라 하자.
$M$이 유한 반사 $R$-가군이면 $M \otimes_R R'$은 유한 반사 $R'$-가군이다.
\end{lemma}

\begin{proof}
$F$가 유한 자유이고 $N$에 비틀림이 없는 짧은 완전열
$0 \to M \to F \to N \to 0$을 택한다.
보조정리 \ref{more-algebra-lemma-characterize-reflexive}를 보라.
$R \to R'$이 평탄이므로 짧은 완전열
$0 \to M \otimes_R R' \to F \otimes_R R' \to N \otimes_R R' \to 0$을 얻는다.
여기서 $F \otimes_R R'$은 유한 자유이고 $N \otimes_R R'$에는 비틀림이 없다
(보조정리 \ref{more-algebra-lemma-flat-pullback-torsion}).
따라서 보조정리 \ref{more-algebra-lemma-characterize-reflexive}에 의해
$M \otimes_R R'$은 반사 가군이다.
\end{proof}

\begin{lemma}
\label{more-algebra-lemma-dual-reflexive}
$R$을 뇌터 정역이라 하자. $M$을 유한 $R$-가군이라 하고
$N$을 유한 반사 $R$-가군이라 하자. 그러면 $\Hom_R(M, N)$은 반사 가군이다.
\end{lemma}

\begin{proof}
표시 $R^{\oplus m} \to R^{\oplus n} \to M \to 0$을 택한다.
그러면
$$
0 \to \Hom_R(M, N) \to N^{\oplus n} \to N' \to 0
$$
을 얻으며, $N' = \Im(N^{\oplus n} \to N^{\oplus m})$에는 비틀림이 없다.
보조정리 \ref{more-algebra-lemma-sequence-reflexive}에 의해 결론을 얻는다.
\end{proof}

\begin{definition}
\label{more-algebra-definition-reflexive-hull}
$R$을 뇌터 정역이라 하자. $M$을 유한 $R$-가군이라 하자.
가군 $M^{**} = \Hom_R(\Hom_R(M, R), R)$을 $M$의 {\it 반사 포락}이라 한다.
\end{definition}

\noindent
보조정리 \ref{more-algebra-lemma-dual-reflexive}에 의해 반사 포락은 반사 가군이므로
이 정의는 타당하다. 대응 $M \mapsto M^{**}$은 함자이다.
$\varphi : M \to N$이 반사 $R$-가군 $N$으로 가는 $R$-가군 사상이면,
$\varphi$는 $M$의 반사 포락을 통하여 $M \to M^{**} \to N$으로 분해된다.
달리 말하면 반사 포락을 취하는 것은 뇌터 정역 $R$ 위의 포함 함자
$$
\text{finite reflexive modules} \subset
\text{finite modules}
$$
의 왼쪽 수반이다.

\begin{lemma}
\label{more-algebra-lemma-hom-into-depth}
$R$을 뇌터 국소환이라 하자. $M$, $N$을 유한 $R$-가군이라 하자.
\begin{enumerate}
\item $N$의 깊이가 $\geq 1$이면 $\Hom_R(M, N)$의 깊이는 $\geq 1$이다.
\item $N$의 깊이가 $\geq 2$이면 $\Hom_R(M, N)$의 깊이는 $\geq 2$이다.
\end{enumerate}
\end{lemma}

\begin{proof}
표시 $R^{\oplus m} \to R^{\oplus n} \to M \to 0$을 택한다.
쌍대를 취하면 완전열
$$
0 \to \Hom_R(M, N) \to N^{\oplus n} \to N' \to 0
$$
을 얻으며, $N' = \Im(N^{\oplus n} \to N^{\oplus m})$이다.
깊이가 $\geq 1$인 가군의 부분가군은 깊이가 $\geq 1$이다.
이는 정의에서 바로 따른다. 따라서 (1)은 명백하다.
% P02-E0816: frozen compound subject takes a singular verb; Korean supplies grammatical agreement.
(2)에 대해서는 가정과 앞의 관찰로부터 $N'$의 깊이가 $\geq 1$임을 주목하자.
가군 $N^{\oplus n}$의 깊이는 $\geq 2$이다.
Algebra, Lemma \ref{algebra-lemma-depth-in-ses}에 의해
$\Hom_R(M, N)$의 깊이가 $\geq 2$라고 결론짓는다.
\end{proof}

\begin{lemma}
\label{more-algebra-lemma-hom-into-S2}
$R$을 뇌터 환이라 하자. $M$, $N$을 유한 $R$-가군이라 하자.
\begin{enumerate}
\item $N$이 성질 $(S_1)$을 가지면 $\Hom_R(M, N)$도 성질 $(S_1)$을 갖는다.
\item $N$이 성질 $(S_2)$를 가지면 $\Hom_R(M, N)$도 성질 $(S_2)$를 갖는다.
\item $R$이 정역이고 $N$에 비틀림이 없으며 $(S_2)$이면,
$\Hom_R(M, N)$에는 비틀림이 없고 성질 $(S_2)$를 갖는다.
\end{enumerate}
\end{lemma}

\begin{proof}
유한 $R$-가군에 대해서는 소아이디얼에서의 국소화가 $\Hom_R$을 취하는 것과
교환하므로(Algebra, Lemma \ref{algebra-lemma-ext-noetherian}),
(1)과 (2)는 보조정리 \ref{more-algebra-lemma-hom-into-depth}에서 바로 따른다.
(3)은 (2)와 보조정리 \ref{more-algebra-lemma-hom-into-torsion-free}에서 따른다.
\end{proof}

\begin{lemma}
\label{more-algebra-lemma-check-injective-on-ass}
$R$을 뇌터 환이라 하자. $\varphi : M \to N$을 $R$-가군들의 사상이라 하자.
$R$의 모든 소아이디얼 $\mathfrak p$에 대하여 다음 중 적어도 하나가
성립한다고 가정하자.
\begin{enumerate}
\item $M_\mathfrak p \to N_\mathfrak p$는 단사이다.
\item $\mathfrak p \not \in \text{Ass}(M)$이다.
\end{enumerate}
그러면 $\varphi$는 단사이다.
\end{lemma}

\begin{proof}
$\mathfrak p$를 $\Ker(\varphi)$의 동반 소아이디얼이라 하자.
그러면 $M_\mathfrak p \to N_\mathfrak p$의 핵에 속하고 소멸자가
$\mathfrak pR_\mathfrak p$인 원소 $x \in M_\mathfrak p$가 존재한다
(Algebra, Lemma \ref{algebra-lemma-associated-primes-localize}).
이는 두 경우 모두 불가능하다.
따라서 $\text{Ass}(\Ker(\varphi)) = \emptyset$이고,
Algebra, Lemma \ref{algebra-lemma-ass-zero}에 의해 $\Ker(\varphi) = 0$이라고
결론짓는다.
\end{proof}

\begin{lemma}
\label{more-algebra-lemma-check-isomorphism-via-depth-and-ass}
$R$을 뇌터 환이라 하자. $\varphi : M \to N$을 $R$-가군들의 사상이라 하자.
$M$이 유한이고 $R$의 모든 소아이디얼 $\mathfrak p$에 대하여
다음 중 하나가 성립한다고 가정하자.
\begin{enumerate}
\item $M_\mathfrak p \to N_\mathfrak p$는 동형이다.
\item $\text{depth}(M_\mathfrak p) \geq 2$이고
$\mathfrak p \not \in \text{Ass}(N)$이다.
\end{enumerate}
그러면 $\varphi$는 동형이다.
\end{lemma}

\begin{proof}
보조정리 \ref{more-algebra-lemma-check-injective-on-ass}에 의해 $\varphi$는 단사임을 알 수 있다.
% P02-E0817: frozen source uses “an finitely generated”; Korean supplies grammatical phrasing.
$N' \subset N$을 $M$의 상을 포함하는 유한 생성 $R$-가군이라 하자.
% P02-E0818: retain the frozen emptiness implication rather than replacing it with the intended prime-exclusion condition.
그러면 $\text{Ass}(N_\mathfrak p) = \emptyset$이면
$\text{Ass}(N'_\mathfrak p) = \emptyset$이다.
따라서 $M \to N'$에 대하여 보조정리의 가정들이 성립한다.
$\varphi$가 동형임을 증명하려면 그러한 모든 $N' \subset N$에 대하여
같은 사실을 증명하면 충분하다. 따라서 $N$이 유한 $R$-가군이라고 가정해도 된다.
이 경우 Algebra, Lemma \ref{algebra-lemma-ideal-nonzerodivisor}에 의해
$\mathfrak p \not \in \text{Ass}(N) \Rightarrow
\text{depth}(N_\mathfrak p) \geq 1$이다.
$Q$를 정의하는 짧은 완전열
$$
0 \to M \to N \to Q \to 0
$$
을 생각하자. 조건들을 보면 (1)의 경우에는 $Q_\mathfrak p = 0$이고,
(2)의 경우에는 Algebra, Lemma \ref{algebra-lemma-depth-in-ses}에 의해
$\text{depth}(Q_\mathfrak p) \geq 1$임을 알 수 있다.
이는 $Q$가 동반 소아이디얼을 갖지 않음을 함의하므로,
Algebra, Lemma \ref{algebra-lemma-ass-zero}에 의해 $Q = 0$이다.
\end{proof}

\begin{lemma}
\label{more-algebra-lemma-isom-depth-2-torsion-free}
$R$을 뇌터 정역이라 하자. $\varphi : M \to N$을 $R$-가군들의 사상이라 하자.
$M$이 유한이고 $N$에 비틀림이 없으며,
$R$의 모든 소아이디얼 $\mathfrak p$에 대하여 다음 중 하나가 성립한다고 가정하자.
\begin{enumerate}
\item $M_\mathfrak p \to N_\mathfrak p$는 동형이다.
\item $\text{depth}(M_\mathfrak p) \geq 2$이다.
\end{enumerate}
그러면 $\varphi$는 동형이다.
\end{lemma}

\begin{proof}
이는 보조정리 \ref{more-algebra-lemma-check-isomorphism-via-depth-and-ass}의 특수한 경우이다.
\end{proof}

\begin{lemma}
\label{more-algebra-lemma-reflexive-depth-2}
$R$을 뇌터 정역이라 하자. $M$을 유한 $R$-가군이라 하자.
다음은 동치이다.
\begin{enumerate}
\item $M$은 반사 가군이다.
\item $R$의 모든 소아이디얼 $\mathfrak p$에 대하여 다음 중 하나가 성립한다.
\begin{enumerate}
\item $M_\mathfrak p$는 반사 $R_\mathfrak p$-가군이다.
\item $\text{depth}(M_\mathfrak p) \geq 2$이다.
\end{enumerate}
\end{enumerate}
\end{lemma}

\begin{proof}
% P02-E0819: frozen source says “all primes of p”; Korean supplies the intended quantifier.
(1)이 성립하면 보조정리 \ref{more-algebra-lemma-check-reflexive}에 의해
모든 소아이디얼 $\mathfrak p$에 대하여 $M_\mathfrak p$는 반사 가군이다.
따라서 (1) $\Rightarrow$ (2)이다. (2)를 가정하자.
$N = \Hom_R(\Hom_R(M, R), R)$로 놓으면 $R$의 모든 소아이디얼
$\mathfrak p$에 대하여 다음이 성립한다.
$$
N_\mathfrak p =
\Hom_{R_\mathfrak p}(
\Hom_{R_\mathfrak p}(M_\mathfrak p, R_\mathfrak p), R_\mathfrak p)
$$
Algebra, Lemma \ref{algebra-lemma-hom-from-finitely-presented}를 보라.
보조정리 \ref{more-algebra-lemma-isom-depth-2-torsion-free}를 사상 $j : M \to N$에 적용한다.
$M$이 유한이고 보조정리 \ref{more-algebra-lemma-hom-into-torsion-free}에 의해
$N$에 비틀림이 없으므로 적용할 수 있다.
(2)(a)의 경우 사상 $M_\mathfrak p \to N_\mathfrak p$는 동형이고,
(2)(b)의 경우에는 $\text{depth}(M_\mathfrak p) \geq 2$이다.
\end{proof}

\begin{lemma}
\label{more-algebra-lemma-reflexive-S2}
$R$을 뇌터 정역이라 하자. $M$을 유한 반사 $R$-가군이라 하자.
$\mathfrak p \subset R$을 소아이디얼이라 하자.
\begin{enumerate}
\item $\text{depth}(R_\mathfrak p) \geq 2$이면
$\text{depth}(M_\mathfrak p) \geq 2$이다.
\item $R$이 $(S_2)$이면 $M$도 $(S_2)$이다.
\end{enumerate}
\end{lemma}

\begin{proof}
반사 포락 $\Hom_R(\Hom_R(M, R), R)$을 취하는 것이 국소화와 교환하므로
(Algebra, Lemma \ref{algebra-lemma-hom-from-finitely-presented}),
(1)은 보조정리 \ref{more-algebra-lemma-hom-into-depth}에서 따른다.
(2)는 보조정리 \ref{more-algebra-lemma-hom-into-S2}에서 바로 따른다.
\end{proof}

\begin{example}
\label{more-algebra-example-ring-not-S2}
위와 아래의 결과들은 반사성이 $(S_2)$ 조건과 관련되어 있음을 시사한다.
지나치게 낙관적인 추측을 막기 위해 예를 하나 제시한다.
$k$를 체라 하자. $R$을 $1, y, x^2, xy, x^3$으로 생성되는
$k[x, y]$의 $k$-부분대수라 하자. 그러면 $R$은 $(S_2)$가 아니다.
따라서 $R$을 $R$-가군으로 보면 이는 $(S_2)$가 아닌 반사 $R$-가군의 예이다.
$M = k[x, y]$를 $R$-가군으로 보자. 그러면
$$
\Hom_R(M, R) = \mathfrak m = (y, x^2, xy, x^3)
\quad\text{and}\quad
\Hom_R(\mathfrak m, R) = M
$$
이므로 $M$은 반사 $R$-가군이고,
$M$은 $R$-가군으로서 $(S_2)$이다(계산은 생략한다).
따라서 $R$은 반사 $(S_2)$ 가군을 가지는 뇌터 정역이지만
$R$ 자체는 $(S_2)$가 아니다.
\end{example}

\begin{lemma}
\label{more-algebra-lemma-reflexive-over-normal}
$R$을 분수체가 $K$인 뇌터 정규 정역이라 하자.
$M$을 유한 $R$-가군이라 하자. 다음은 동치이다.
\begin{enumerate}
\item $M$은 반사 가군이다.
\item $M$에는 비틀림이 없고 성질 $(S_2)$를 갖는다.
\item $M$에는 비틀림이 없고
$M = \bigcap_{\text{height}(\mathfrak p) = 1} M_{\mathfrak p}$이다.
여기서 교집합은 $M_K = M \otimes_R K$ 안에서 취한다.
\end{enumerate}
\end{lemma}

\begin{proof}
Algebra, Lemma \ref{algebra-lemma-criterion-normal}에 의해
$R$이 $(R_1)$과 $(S_2)$를 만족함을 알 수 있다.

\medskip\noindent
(1)을 가정하자. 그러면 보조정리 \ref{more-algebra-lemma-reflexive-torsion-free}에 의해
$M$에는 비틀림이 없고, 보조정리 \ref{more-algebra-lemma-reflexive-S2}에 의해
$(S_2)$를 만족한다. 따라서 (2)가 성립한다.

\medskip\noindent
(2)를 가정하자. 정의에 의해
$M' = \bigcap_{\text{height}(\mathfrak p) = 1} M_{\mathfrak p}$는 사상
$$
M_K
\longrightarrow
\bigoplus\nolimits_{\text{height}(\mathfrak p) = 1} M_K/M_\mathfrak p
\subset
\prod\nolimits_{\text{height}(\mathfrak p) = 1} M_K/M_\mathfrak p
$$
의 핵이다. 주어진 $a/b \in K$에 대하여 $b \in \mathfrak p$인
높이 $1$의 소아이디얼 $\mathfrak p$는 유한 개뿐이므로,
위 사상은 실제로 표시한 직접합을 통하여 분해됨을 관찰하자.
$\mathfrak p_0$를 높이 $1$의 소아이디얼이라 하자.
그러면 $\mathfrak p = \mathfrak p_0$인 경우를 제외하면
$(M_K/M_\mathfrak p)_{\mathfrak p_0} = 0$이고,
그 경우에는 $(M_K/M_\mathfrak p)_{\mathfrak p_0} = M_K/M_{\mathfrak p_0}$을 얻는다.
따라서 국소화의 완전성과 국소화가 직접합과 교환한다는 사실에 의해
$M'_{\mathfrak p_0} = M_{\mathfrak p_0}$임을 알 수 있다.
$M$은 높이가 $> 1$인 소아이디얼에서 깊이가 $\geq 2$이므로,
보조정리 \ref{more-algebra-lemma-isom-depth-2-torsion-free}에 의해
$M \to M'$은 동형임을 알 수 있다. 따라서 (3)이 성립한다.

\medskip\noindent
(3)을 가정하자. $\mathfrak p$를 높이 $1$의 소아이디얼이라 하자.
그러면 $(R_1)$에 의해 $R_\mathfrak p$는 이산 값매김환이다.
보조정리 \ref{more-algebra-lemma-dedekind-torsion-free-flat}에 의해
$M_\mathfrak p$는 유한 자유 가군이며, 특히 반사 가군임을 알 수 있다.
따라서 사상 $M \to M^{**}$은 높이 $1$의 모든 소아이디얼 $\mathfrak p$에서
동형을 유도한다. 그러므로 조건
$M = \bigcap_{\text{height}(\mathfrak p) = 1} M_{\mathfrak p}$는
$M = M^{**}$를 함의하고 (1)이 성립한다.
\end{proof}

\begin{lemma}
\label{more-algebra-lemma-describe-reflexive-hull}
$R$을 뇌터 정규 정역이라 하자. $M$을 유한 $R$-가군이라 하자.
그러면 $M$의 반사 포락은 $M \otimes_R K$ 안에서 취한 교집합
$$
M^{**} =
\bigcap\nolimits_{\text{height}(\mathfrak p) = 1}
M_{\mathfrak p}/(M_\mathfrak p)_{tors} =
\bigcap\nolimits_{\text{height}(\mathfrak p) = 1}
(M/M_{tors})_\mathfrak p
$$
이다.
\end{lemma}

\begin{proof}
$\mathfrak p$를 높이 $1$의 소아이디얼이라 하자.
$M_\mathfrak p \to M \otimes_R K$의 핵은 $M_\mathfrak p$의
비틀림 부분가군 $(M_\mathfrak p)_{tors}$이다.
또한 $(M/M_{tors})_\mathfrak p = M_\mathfrak p/(M_\mathfrak p)_{tors}$이고,
이는 이산 값매김환 $R_\mathfrak p$ 위의 유한 자유 가군이다
(보조정리 \ref{more-algebra-lemma-dedekind-torsion-free-flat}).
그러면 $M_\mathfrak p/(M_\mathfrak p)_{tors} \to
(M_\mathfrak p)^{**} = (M^{**})_\mathfrak p$는 동형이므로,
보조정리는 보조정리 \ref{more-algebra-lemma-reflexive-over-normal}의 귀결이다.
\end{proof}

\begin{lemma}
\label{more-algebra-lemma-integral-closure-reflexive}
$A$를 분수체가 $K$인 뇌터 정규 정역이라 하자.
$L$을 $K$의 유한 확대라 하자.
$L$ 안에서 $A$의 정수적 폐포 $B$가 $A$ 위 유한이면,
$B$는 $A$-가군으로서 반사 가군이다.
\end{lemma}

\begin{proof}
$B = \bigcap B_\mathfrak p$임을 보이면 충분하다.
여기서 교집합은 높이 $1$의 소아이디얼 $\mathfrak p \subset A$들에 걸쳐 취한다.
보조정리 \ref{more-algebra-lemma-reflexive-over-normal}을 보라.
$b \in \bigcap B_\mathfrak p$라 하자.
$x^d + a_1x^{d - 1} + \ldots + a_d$를 $K$ 위에서 $b$의 최소다항식이라 하자.
$a_i \in A$임을 보이고자 한다.
Algebra, Lemma \ref{algebra-lemma-minimal-polynomial-normal-domain}에 의해
모든 $i$와 높이 1의 모든 소아이디얼 $\mathfrak p$에 대하여
$a_i \in A_\mathfrak p$임을 알 수 있다.
따라서 Algebra, Lemma
\ref{algebra-lemma-normal-domain-intersection-localizations-height-1}에서
원하는 결론을 얻는다
(또는 $A$가 자기 자신 위의 반사 가군이므로 이미 인용한 보조정리를 사용해도 된다).
\end{proof}

\section{내용 아이디얼}
\label{more-algebra-section-content}

\noindent
정의가 독자가 예상하는 것과 다를 수도 있다.

\begin{definition}
\label{more-algebra-definition-content-ideal}
$A$를 환이라 하고 $M$을 평탄 $A$-가군이라 하자. $x \in M$이라 하자.
$x \in IM$인 $A$의 아이디얼 $I$들의 집합이 가장 작은 원소를 가지면,
이를 {$x$의 \it 내용 아이디얼}이라 한다.
\end{definition}

\noindent
$M$이 $A$ 위 평탄이므로 $A$의 두 아이디얼 $I, I'$에 대하여
$IM \cap I'M = (I \cap I')M$임을 주목하자.
이는 완전열 $0 \to I \cap I' \to I \oplus I' \to I + I' \to 0$을
$M$과 텐서하면 알 수 있다.

\begin{lemma}
\label{more-algebra-lemma-content-finitely-generated}
$A$를 환이라 하고 $M$을 평탄 $A$-가군이라 하자. $x \in M$이라 하자.
$x$의 내용 아이디얼이 존재하면 이는 유한 생성이다.
\end{lemma}

\begin{proof}
$x \in IM$이라 하자. 그러면 $f_i \in I$와 $x_i \in M$을 사용하여
$x = \sum_{i = 1, \ldots, n} f_i x_i$로 쓸 수 있다.
따라서 $I' = (f_1, \ldots, f_n)$에 대하여 $x \in I'M$이다.
\end{proof}

\begin{lemma}
\label{more-algebra-lemma-equal-content}
$(A, \mathfrak m)$을 국소환이라 하자.
$u : M \to N$을 $\overline{u} : M/\mathfrak m M \to N/\mathfrak m N$이
단사인 평탄 $A$-가군들의 사상이라 하자.
$x \in M$의 내용 아이디얼이 $I$이면 $u(x)$의 내용 아이디얼도 $I$이다.
\end{lemma}

\begin{proof}
$u(x) \in IN$임은 명백하다. $u(x) \in I'N$이면
$u(x) \in (I' \cap I)N$이다. 정의 \ref{more-algebra-definition-content-ideal} 뒤의 논의를 보라.
따라서 다음을 보이면 충분하다.
% P02-E0820: frozen source has the ill-typed premise x in I'N; preserve the exact mathematical claim.
% P02-E0821: preserve the frozen double space in the inequality span.
$x \in I'N$이고 $I' \subset I$, $I' \not =  I$이면 $u(x) \not \in I'N$이다.
$I/I'$는 영이 아닌 유한 $A$-가군이므로
(보조정리 \ref{more-algebra-lemma-content-finitely-generated}),
Nakayama 보조정리(Algebra, Lemma \ref{algebra-lemma-NAK})에 의해
영이 아닌 $A$-가군 사상 $\chi : I/I' \to A/\mathfrak m$이 존재한다.
$I$가 $x$의 내용 아이디얼이므로 $I'' = \Ker(\chi)$에 대하여
$x \not \in I''M$임을 알 수 있다.
따라서 $x$는 사상
$$
IM = I \otimes_A M \xrightarrow{\chi \otimes 1}
A/\mathfrak m \otimes M \cong M/\mathfrak m M
$$
의 핵에 속하지 않는다.
$\overline{u}$에 대한 가정을 적용하면 $u(x)$는 사상
$$
IN = I \otimes_A N \xrightarrow{\chi \otimes 1}
A/\mathfrak m \otimes N \cong N/\mathfrak m N
$$
아래에서 영으로 가지 않는다고 결론지으며, 증명이 끝난다.
\end{proof}

\begin{lemma}
\label{more-algebra-lemma-content-exists-flat-Mittag-Leffler}
$A$를 환이라 하고 $M$을 평탄 Mittag-Leffler 가군이라 하자.
그러면 $M$의 모든 원소는 내용 아이디얼을 갖는다.
\end{lemma}

\begin{proof}
이는 Algebra, Lemma \ref{algebra-lemma-flat-ML-criterion}의 특수한 경우이다.
\end{proof}

\section{평탄성과 유한성 조건}
\label{more-algebra-section-flat-finite}

\noindent
이 절에서는 ``평탄 $+$ 유한형 $\Rightarrow$ 유한 표시'' 꼴의 몇 가지 함의를 논의한다.
평탄성에 관한 장에서 이 결과를 다시 다룰 것이다.
More on Flatness, Section \ref{flat-section-introduction}을 보라.
이러한 유형의 첫 번째 결과는
Algebra, Lemma \ref{algebra-lemma-finite-flat-module-finitely-presented}에서 증명하였다.

\begin{lemma}
\label{more-algebra-lemma-flat-finite-type-finite-presentation-local-module}
$R$을 환이라 하자. $S = R[x_1, \ldots, x_n]$을 $R$ 위의 다항식환이라 하고
$M$을 $S$-가군이라 하자. 다음을 가정하자.
\begin{enumerate}
\item 사상 $R \to \prod R_{\mathfrak p_j}$가 단사가 되게 하는
$R$의 유한 개의 소아이디얼 $\mathfrak p_1, \ldots, \mathfrak p_m$이 존재한다.
\item $M$은 유한 $S$-가군이다.
% P02-E0822: frozen source omits “is” in this condition; Korean supplies the predicate.
\item $M$은 $R$ 위 평탄이다.
\item $R$의 모든 소아이디얼 $\mathfrak p$에 대하여
가군 $M_{\mathfrak p}$는 $S_{\mathfrak p}$ 위 유한 표시이다.
\end{enumerate}
그러면 $M$은 $S$ 위 유한 표시이다.
\end{lemma}

\begin{proof}
$M$의 $S$-가군으로서의 표시
$$
0 \to K \to S^{\oplus r} \to M \to 0
$$
를 택한다. $\mathfrak q$를 $R$의 소아이디얼 $\mathfrak p$ 위에 놓이는
$S$의 소아이디얼이라 하자. 가정에 의해 유한 개의 원소
$k_1, \ldots, k_t \in K$가 존재하여 $K' = \sum Sk_j \subset K$로 놓으면
$K'_{\mathfrak p} = K_{\mathfrak p}$이고 $j = 1, \ldots, m$에 대하여
$K'_{\mathfrak p_j} = K_{\mathfrak p_j}$이다.
$M' = S^{\oplus r}/K'$로 놓으면 특히
$M'_{\mathfrak q} = M_{\mathfrak q}$라고 결론 내릴 수 있다.
평탄성의 열린성(Algebra, Theorem \ref{algebra-theorem-openness-flatness}를 보라)에 의해
$M'_g$가 $R$ 위 평탄이 되게 하는 $g \in S$, $g \not \in \mathfrak q$가
존재한다고 결론짓는다.
따라서 $M'_g \to M_g$는 평탄 $R$-가군들의 전사 사상이다.
가환도식
$$
\xymatrix{
M'_g \ar[r] \ar[d] & M_g \ar[d] \\
\prod (M'_g)_{\mathfrak p_j} \ar[r] & \prod (M_g)_{\mathfrak p_j}
}
$$
을 생각하자. 아래 화살표는 $k_1, \ldots, k_t$의 선택에 의해 동형이다.
$R \to \prod R_{\mathfrak p_j}$가 단사이고 $M'_g$가 $R$ 위 평탄이므로
왼쪽 수직 화살표는 단사이다.
따라서 위 수평 화살표는 단사이며, 따라서 동형이다.
이로써 $M_g$가 $S_g$ 위 유한 표시임이 증명된다.
Algebra, Lemma \ref{algebra-lemma-cover}를 적용하여 결론을 얻는다.
\end{proof}

\begin{lemma}
\label{more-algebra-lemma-flat-finite-type-finite-presentation-local}
$R \to S$를 환 준동형이라 하자. 다음을 가정하자.
\begin{enumerate}
\item 사상 $R \to \prod R_{\mathfrak p_j}$가 단사가 되게 하는
$R$의 유한 개의 소아이디얼 $\mathfrak p_1, \ldots, \mathfrak p_m$이 존재한다.
\item $R \to S$는 유한형이다.
% P02-E0823: frozen source omits “is” in this condition; Korean supplies the predicate.
\item $S$는 $R$ 위 평탄이다.
\item $R$의 모든 소아이디얼 $\mathfrak p$에 대하여
환 $S_{\mathfrak p}$는 $R_{\mathfrak p}$ 위 유한 표시이다.
\end{enumerate}
그러면 $S$는 $R$ 위 유한 표시이다.
\end{lemma}

\begin{proof}
가정에 의해 $S$는 $R$ 위 다항식환의 몫환이다.
따라서 결과는 보조정리
\ref{more-algebra-lemma-flat-finite-type-finite-presentation-local-module}에서 직접 따른다.
\end{proof}

\begin{lemma}
\label{more-algebra-lemma-flat-graded-finite-type-finite-presentation-module}
$R$을 환이라 하자.
$S = R[x_1, \ldots, x_n]$을 $R$ 위의 등급 다항식 대수라 하자.
즉 $\deg(x_i) > 0$이지만 반드시 $1$일 필요는 없다.
$M$을 등급 $S$-가군이라 하자. 다음을 가정하자.
\begin{enumerate}
\item $R$은 국소환이다.
\item $M$은 유한 $S$-가군이다.
\item $M$은 $R$ 위 평탄이다.
\end{enumerate}
그러면 $M$은 $S$-가군으로서 유한 표시이다.
\end{lemma}

\begin{proof}
$M = \bigoplus M_d$를 $M$의 등급이라 하자.
$M$의 동차 생성원 $m_1, \ldots, m_r \in M$을 택한다.
$\deg(m_i) = d_i \in \mathbf{Z}$로 쓰자. 이로부터 표시
$$
0 \to K \to \bigoplus\nolimits_{i = 1, \ldots, r} S(-d_i) \to M \to 0
$$
를 얻으며, 각 차수 $d$에서 짧은 완전열
$$
0 \to K_d \to \bigoplus\nolimits_{i = 1, \ldots, r} S_{d - d_i} \to
M_d \to 0.
$$
을 얻는다. 가정에 의해 각 $M_d$는 유한 평탄 $R$-가군이다.
Algebra, Lemma \ref{algebra-lemma-finite-flat-local}에 의해
이는 각 $M_d$가 유한 자유 $R$-가군임을 함의한다.
따라서 각 $K_d$는 유한 $R$-가군임을 알 수 있다.
또한 Algebra, Lemma \ref{algebra-lemma-flat-ses}에 의해
각 $K_d$는 $R$ 위 평탄이다.
그러므로 다시 Algebra, Lemma \ref{algebra-lemma-finite-flat-local}에 의해
각 $K_d$는 유한 자유라고 결론짓는다.

\medskip\noindent
$\mathfrak m$을 $R$의 극대아이디얼이라 하자.
$M$이 $R$ 위 평탄이므로 위의 짧은 완전열들은
$\kappa = \kappa(\mathfrak m)$와 텐서한 뒤에도 짧은 완전열로 남는다.
환 $S \otimes_R \kappa$가 뇌터 환이므로,
그 상 $\overline{k}_j$들이 $S \otimes_R \kappa$ 위에서
$K \otimes_R \kappa$를 생성하는 동차 원소 $k_1, \ldots, k_t \in K$가
존재함을 알 수 있다. $\deg(k_j) = e_j$로 쓰자.
따라서 임의의 $d$에 대하여 사상
$$
\bigoplus\nolimits_{j = 1, \ldots, t} S_{d - e_j}
\longrightarrow
K_d
$$
은 $\kappa$와 텐서한 뒤 전사가 된다.
Nakayama 보조정리(Algebra, Lemma \ref{algebra-lemma-NAK})에 의해
이는 사상이 $R$ 위에서 전사임을 함의한다.
따라서 $K$는 $S$ 위에서 $k_1, \ldots, k_t$로 생성되므로 결론을 얻는다.
\end{proof}

\begin{lemma}
\label{more-algebra-lemma-flat-graded-finite-type-finite-presentation}
$R$을 환이라 하자. $S = \bigoplus_{n \geq 0} S_n$을 등급 $R$-대수라 하자.
$M = \bigoplus_{d \in \mathbf{Z}} M_d$를 등급 $S$-가군이라 하자.
$S$가 $R$-대수로서 유한 생성이고 $S_0$가 유한 $R$-대수이며,
% P02-E0824: frozen source indexes the primes by j but gives an i-range; preserve both spans.
$R \to \prod R_{\mathfrak p_j}$가 단사가 되게 하는 유한 개의 소아이디얼
$\mathfrak p_j$, $i = 1, \ldots, m$이 존재한다고 가정하자.
\begin{enumerate}
\item $S$가 $R$ 위 평탄이면 $S$는 유한 표시 $R$-대수이다.
\item $M$이 $R$-가군으로서 평탄이고 $S$-가군으로서 유한이면
$M$은 $S$-가군으로서 유한 표시이다.
\end{enumerate}
\end{lemma}

\begin{proof}
$S$가 $R$-대수로서 유한 생성이므로 $S_0$-대수로서 유한 생성이며,
차수가 $d_1, \ldots, d_n > 0$인 동차 원소
$t_1, \ldots, t_n \in S$로 생성된다고 하자.
$\deg(x_i) = d_i$로 두고 $P = R[x_1, \ldots, x_n]$으로 놓는다.
$S_0$가 유한 $R$-가군이므로 환 사상 $P \to S$, $x_i \to t_i$는 유한이다.
(1)을 증명하려면 $S$가 유한 표시 $P$-가군임을 보이면 충분하다.
% P02-E0825: frozen source doubles the intersentence space; Korean uses ordinary prose spacing.
(2)를 증명하려면 $M$이 유한 표시 $P$-가군임을 보이면 충분하다.
따라서 $S = P$가 등급 다항식환이고 $M$이 $R$ 위 평탄인 유한
$S$-가군이면 $M$이 $S$-가군으로서 유한 표시임을 보이면 충분하다.
보조정리 \ref{more-algebra-lemma-flat-graded-finite-type-finite-presentation-module}에 의해
$R$의 모든 소아이디얼 $\mathfrak p$에 대하여 $M_{\mathfrak p}$는
유한 표시 $S_{\mathfrak p}$-가군임을 알 수 있다.
따라서 결과는 보조정리
\ref{more-algebra-lemma-flat-finite-type-finite-presentation-local-module}에서 따른다.
\end{proof}

\begin{remark}
\label{more-algebra-remark-when-does-condition-hold}
$R$을 환이라 하자. $R$은 언제 보조정리
\ref{more-algebra-lemma-flat-finite-type-finite-presentation-local-module},
\ref{more-algebra-lemma-flat-finite-type-finite-presentation-local},
\ref{more-algebra-lemma-flat-graded-finite-type-finite-presentation}에 언급된 조건을 만족하는가?
다음의 경우에 성립한다.
\begin{enumerate}
\item $R$은 국소환이다.
\item $R$은 뇌터 환이다.
\item $R$은 정역이다.
\item $R$은 최소 소아이디얼이 유한 개인 축약환이다.
\item $R$은 약한 동반 소아이디얼이 유한 개이다.
Algebra, Lemma \ref{algebra-lemma-zero-at-weakly-ass-zero}를 보라.
\end{enumerate}
따라서 이 보조정리들은 위에 나열한 모든 경우에 성립한다.
\end{remark}

\noindent
다음 보조정리는 More on Flatness, Proposition
\ref{flat-proposition-flat-finite-type-finite-presentation-domain}에서 개선될 것이다.

\begin{lemma}
\label{more-algebra-lemma-flat-finite-type-valuation-ring-finite-presentation}
\begin{reference}
\cite[Theorem 3]{Nagata-Finitely}
\end{reference}
$A$를 값매김환이라 하자. $A \to B$를 유한형 환 사상이라 하고
$M$을 유한 $B$-가군이라 하자.
\begin{enumerate}
\item $B$가 $A$ 위 평탄이면 $B$는 유한 표시 $A$-대수이다.
\item $M$이 $A$-가군으로서 평탄이면 $M$은 $B$-가군으로서 유한 표시이다.
\end{enumerate}
\end{lemma}

\begin{proof}
$A$-가군이 평탄일 필요충분조건은 비틀림이 없는 것이라는 사실을 사용하겠다.
보조정리 \ref{more-algebra-lemma-valuation-ring-torsion-free-flat}을 보라.
Algebra, Lemma \ref{algebra-lemma-homogenize}에 의해,
$S_0 = A$이고 차수 $1$인 유한 개의 원소로 생성되는 등급 $A$-대수 $S$,
원소 $f \in S_1$, 그리고 $B \cong S_{(f)}$와 $M \cong N_{(f)}$가 되게 하는
유한 등급 $S$-가군 $N$을 찾을 수 있다.
$M$에 비틀림이 없으면 $N$을 $N/N_{tors}$로 바꾸어
비틀림이 없게 할 수 있다. 보조정리 \ref{more-algebra-lemma-torsion}을 보라.
마찬가지로 $B$에 비틀림이 없으면 $S$를 $S/S_{tors}$로 바꾸어
비틀림이 없게 할 수 있다.
따라서 (1)의 경우에는 보조정리
\ref{more-algebra-lemma-flat-graded-finite-type-finite-presentation}을 적용하여
$S$가 유한 표시 $A$-대수임을 알 수 있으며,
이는 $B = S_{(f)}$가 유한 표시 $A$-대수임을 함의한다.
(2)를 보려면 먼저 $S$를 등급 다항식환으로 바꾼 뒤,
보조정리 \ref{more-algebra-lemma-flat-graded-finite-type-finite-presentation-module}을
적용하여 결론을 얻으면 된다.
\end{proof}

\begin{lemma}
\label{more-algebra-lemma-valuation-ring-flat-essentially-finite-type}
$A$를 값매김환이라 하자. $A \to B$를 본질적으로 유한형인
국소 준동형이라 하고 $M$을 유한 $B$-가군이라 하자.
\begin{enumerate}
\item $B$가 $A$ 위 평탄이면 $B$는 $A$ 위 본질적으로 유한 표시이다.
\item $M$이 $A$-가군으로서 평탄이면 $M$은 $B$-가군으로서 유한 표시이다.
\end{enumerate}
\end{lemma}

\begin{proof}
가정에 의해 $B$는 다항식 대수 $P = A[x_1, \ldots, x_n]$를
소아이디얼 $\mathfrak q$에서 국소화한 환의 몫환으로 쓸 수 있다.
(1)의 경우에는 $M = B$를 $P_\mathfrak q$ 위의 유한 가군으로 보고,
(2)의 경우에는 $M$을 $P_\mathfrak q$ 위의 유한 가군으로 본다.
두 경우 모두 이것이 유한 표시 $P_\mathfrak q$-가군임을 보여야 한다.
(2)의 경우에 대해서는 Algebra, Lemma
\ref{algebra-lemma-finitely-presented-over-subring}을 보라.

\medskip\noindent
표시 $0 \to K \to P_\mathfrak q^{\oplus r} \to M \to 0$을 택한다.
$M$이 $P_\mathfrak q$ 위 유한이므로 가능하다.
$L = P^{\oplus r} \cap K$로 놓는다. 그러면 $K = L_\mathfrak q$이다.
Algebra, Lemma \ref{algebra-lemma-submodule-localization}을 보라.
이때 $N = P^{\oplus r}/L$은 $M$의 부분가군이므로
보조정리 \ref{more-algebra-lemma-valuation-ring-torsion-free-flat}에 의해 평탄이다.
또한 $N$은 유한 $P$-가군이므로 보조정리
\ref{more-algebra-lemma-flat-finite-type-valuation-ring-finite-presentation}에 의해
$N$은 $P$-가군으로서 유한 표시임을 알 수 있다.
국소화는 완전하므로(Algebra, Proposition \ref{algebra-proposition-localization-exact}),
$N_\mathfrak q = M$임을 알 수 있고 결론을 얻는다.
\end{proof}

\section{블로업과 평탄성}
\label{more-algebra-section-blowup-flat}

\noindent
이 절에서는 ``블로업 뒤에 엄밀 변환이 평탄해진다''는 형태의 결과들에 대한
논의를 시작한다. 이러한 유형의 다른 결과들은
Divisors, Section \ref{divisors-section-blowup-flat}과
More on Flatness, Section \ref{flat-section-blowup-flat}에서 찾을 수 있다.

\begin{definition}
\label{more-algebra-definition-strict-transform}
$R$을 환이라 하자. $I \subset R$을 아이디얼이라 하고 $a \in I$라 하자.
$R[\frac{I}{a}]$를 아핀 블로업 대수라 하자.
Algebra, Definition \ref{algebra-definition-blow-up}을 보라.
$M$을 $R$-가군이라 하자.
{$R \to R[\frac{I}{a}]$를 따른 $M$의 \it 엄밀 변환}은
다음 $R[\frac{I}{a}]$-가군이다.
$$
M' = \left(M \otimes_R R[\textstyle{\frac{I}{a}}]\right)/a\text{-power torsion}
$$
\end{definition}

\noindent
다음은 블로업에 의한 평탄화의 매우 약한 형태이지만,
이 정도만으로도 때로는 유용한 결과이다.

\begin{lemma}
\label{more-algebra-lemma-flatten-on-affine-blowup}
$(R, \mathfrak m)$을 분수체가 $K$인 국소 정역이라 하자.
$S$를 유한형 $R$-대수라 하고 $M$을 유한 $S$-가군이라 하자.
$R$을 지배하는 모든 값매김환 $A \subset K$에 대하여,
다음을 만족하는 아이디얼 $I \subset \mathfrak m$과 영이 아닌 원소
$a \in I$가 존재한다.
\begin{enumerate}
\item $I$는 유한 생성이다.
\item $A$는 $R[\frac{I}{a}]$ 위에 중심을 갖는다.
\item $\mathfrak m$에서의 $R \to R[\frac{I}{a}]$의 섬유환은 영이 아니다.
\item $R \to R[\frac{I}{a}]$를 따른 $S$의 엄밀 변환 $S_{I, a}$는
$R$ 위 평탄이고 유한 표시이며,
$R \to R[\frac{I}{a}]$를 따른 $M$의 엄밀 변환 $M_{I, a}$는
$R$ 위 평탄이고 $S_{I, a}$ 위 유한 표시이다.
\end{enumerate}
\end{lemma}

\begin{proof}
$S = R[x_1, \ldots, x_n]/J$로 쓰고 $N = S \oplus M$을
$P = R[x_1, \ldots, x_n]$ 위의 가군으로 보자.
$S$가 $R$ 위의 다항식 대수인 경우에 보조정리를 증명할 수 있다면,
(1), (2), (3)을 만족하고 $R \to R[\frac{I}{a}]$를 따른 $N$의 엄밀 변환
$N_{I, a}$가 $R$ 위 평탄이며 $P$의 엄밀 변환 $P_{I, a}$ 위의 가군으로서
유한 표시가 되게 하는 $I, a$를 찾을 수 있다.
$P_{I, a} = R[\frac{I}{a}][x_1, \ldots, x_n]$이므로
(작은 세부사항은 생략한다), 직접합인자 $S_{I, a} \subset N_{I, a}$가
$R$ 위 평탄이고 $R[\frac{I}{a}][x_1, \ldots, x_n]$ 위의 가군으로서
유한 표시임을 얻는다.
따라서 $S_{I, a}$는 $R[\frac{I}{a}]$-대수로서 유한 표시이다.
또한 직접합인자 $M_{I, a} \subset N_{I, a}$는 $R$ 위 평탄이고
$P_{I, a}$ 위의 가군으로서 유한 표시이므로,
$S_{I, a}$ 위의 가군으로서도 유한 표시이다.
Algebra, Lemma \ref{algebra-lemma-finitely-presented-over-subring}을 보라.
이로써 다음 단락에서 논의하는 경우로 환원된다.

\medskip\noindent
$S = R[x_1, \ldots, x_n]$이라고 가정하자. 표시
$$
0 \to K \to S^{\oplus r} \to M \to 0.
$$
를 택한다. $M_A$를 $M \otimes_R A$를 그 비틀림 부분가군으로 나눈
몫가군이라 하자. 보조정리 \ref{more-algebra-lemma-torsion}을 보라.
그러면 $M_A$는 $S_A = A[x_1, \ldots, x_n]$ 위의 유한 가군이다.
보조정리 \ref{more-algebra-lemma-valuation-ring-torsion-free-flat}에 의해
$M_A$는 $A$ 위 평탄임을 알 수 있다.
보조정리 \ref{more-algebra-lemma-flat-finite-type-valuation-ring-finite-presentation}에 의해
$M_A$는 유한 표시임을 알 수 있다.
따라서 표시 $S_A^{\oplus r} \to M_A$의 핵을 $S_A$-가군으로서 생성하는
유한 개의 원소 $k_1, \ldots, k_t \in S_A^{\oplus r}$가 존재한다.
(1), (2), (3)을 만족하는 임의의 $a \in I \subset \mathfrak m$에 대하여
$M_{I, a}$로 $R \to R[\frac{I}{a}]$를 따른 $M$의 엄밀 변환을 나타내자.
이는 $S_{I, a} = R[\frac{I}{a}][x_1, \ldots, x_n]$ 위의 유한 가군이다.
Algebra, Lemma \ref{algebra-lemma-valuation-ring-colimit-affine-blowups}에 의해
$A = \colim_{I, a} R[\frac{I}{a}]$이다.
이는 $S_A = \colim S_{I, a}$와 다음을 함의한다.
$$
\colim M \otimes_R R[\textstyle{\frac{I}{a}}] = M \otimes_R A
$$
$I, a$와 올림 $k_1, \ldots, k_t \in S_{I, a}^{\oplus r}$를 택한다.
$M_A$는 $M \otimes_R A$를 비틀림으로 나눈 몫이므로,
$M \otimes_R A$ 안에서 $k_1, \ldots, k_t$의 상들은 영이 아닌 원소
$\alpha \in A$에 의해 소멸됨을 알 수 있다.
$I, a$를 다른 쌍으로 바꾸면(쌍대극한이 여과되어 있음을 상기하라),
영이 아닌 어떤 $x \in I^n$에 대하여 $\alpha = x/a^n$이라고 가정해도 된다.
그러면 $x k_1, \ldots, x k_t$는 $M \otimes_R A$에서 영으로 간다.
따라서 $I, a$를 다른 쌍으로 바꾸면 영이 아닌 어떤 $x \in R$에 대하여
$x k_1, \ldots, x k_t$가 $M \otimes_R R[\frac{I}{a}]$에서 영으로 간다고
가정해도 된다. 마지막으로 $I, a$를 $xI, xa$로 바꾸면,
$k_1, \ldots, k_t$가 $M \otimes_R R[\frac{I}{a}]$의
$a$의 거듭제곱에 의해 소멸되는 비틀림 원소들로 간다고 가정할 수 있다.
그러한 모든 쌍 $(I, a)$에 대하여 다음과 같이 놓는다.
$$
M'_{I, a} = S_{I, a}^{\oplus r}/ \sum S_{I, a}k_j.
$$
$M_A = S_A^{\oplus r}/ \sum S_Ak_j$이므로
$M_A = \colim_{I, a} M'_{I, a}$임을 알 수 있다.
이제 마지막으로 Algebra, Lemma
\ref{algebra-lemma-flat-finite-presentation-limit-flat} (3)을 적용하여,
위와 같은 어떤 쌍 $(I, a)$에 대하여 $M'_{I, a}$가 평탄이라고 결론짓는다.
이 보조정리는 전이사상들이 보조정리의 가정들을 만족하지 않을 수 있으므로,
엄밀 변환들의 계
$$
M_{I, a} = (M \otimes_R R[\textstyle{\frac{I}{a}}])/a\text{-power torsion}
$$
에 처음부터 적용할 수는 없다.
그러나 이제 평탄성은 비틀림이 없음을 함의하고
(보조정리 \ref{more-algebra-lemma-flat-torsion-free}),
$M_{I, a}$는 $M'_{I, a}$를 $a$의 거듭제곱에 의해 소멸되는 비틀림
부분가군으로 나눈 몫이므로(원소 $k_1, \ldots, k_t$가 $M_{I, a}$에서
영으로 가도록 해 두었기 때문이다),
그러한 쌍에 대하여 $M'_{I, a} = M_{I, a}$라고도 결론짓고 증명이 끝난다.
\end{proof}

\begin{lemma}
\label{more-algebra-lemma-blowup-fitting-ideal}
$R$을 환이라 하고 $M$을 유한 $R$-가군이라 하자.
$k \geq 0$이라 하고 $I = \text{Fit}_k(M)$이라 하자.
모든 $a \in I$에 대하여 $R' = R[\frac{I}{a}]$로 놓으면 엄밀 변환
$$
M' = (M \otimes_R R')/a\text{-power torsion}
$$
은 $\text{Fit}_k(M') = R'$을 만족한다.
\end{lemma}

\begin{proof}
먼저 $\text{Fit}_k(M \otimes_R R') = IR' = aR'$임을 관찰하자.
첫 번째 등식은 보조정리 \ref{more-algebra-lemma-fitting-ideal-basics}의 (3)에 의해,
두 번째 등식은 Algebra, Lemma \ref{algebra-lemma-affine-blowup}에 의해 성립한다.
보조정리 \ref{more-algebra-lemma-principal-fitting-ideal}과 국소화의 완전성에 의해
$R'$의 모든 소아이디얼 $\mathfrak p'$에 대하여 $M'_{\mathfrak p'}$가
$\leq k$개의 원소로 생성될 수 있음을 알 수 있다.
그러면 예를 들어 보조정리 \ref{more-algebra-lemma-fitting-ideal-generate-locally}에 의해
$\text{Fit}_k(M') = R'$이다.
\end{proof}

\begin{lemma}
\label{more-algebra-lemma-blowup-fitting-ideal-locally-free}
$R$을 환이라 하고 $M$을 유한 $R$-가군이라 하자.
$k \geq 0$이라 하고 $I = \text{Fit}_k(M)$이라 하자.
모든 $\mathfrak p \not \in V(I)$에 대하여 $M_\mathfrak p$가
계수 $k$의 자유 가군이라고 가정하자.
그러면 모든 $a \in I$에 대하여 $R' = R[\frac{I}{a}]$로 놓으면 엄밀 변환
$$
M' = (M \otimes_R R')/a\text{-power torsion}
$$
은 계수 $k$의 국소 자유 가군이다.
\end{lemma}

\begin{proof}
보조정리 \ref{more-algebra-lemma-blowup-fitting-ideal}에 의해 $\text{Fit}_k(M') = R'$이다.
보조정리 \ref{more-algebra-lemma-fitting-ideal-finite-locally-free}에 의해
$\text{Fit}_{k - 1}(M') = 0$임을 보이면 충분하다.
$R' \subset R'_a = R_a$임을 상기하자.
Algebra, Lemma \ref{algebra-lemma-affine-blowup}을 보라.
따라서 $\text{Fit}_{k - 1}(M')$가 $R'_a = R_a$에서 영으로 감을 증명하면 충분하다.
명백히 $(M')_a = M_a$이므로 이는 $\text{Fit}_{k - 1}(M_a) = 0$임을 보이는
문제로 환원된다. 보조정리 \ref{more-algebra-lemma-fitting-ideal-basics}의 (3)에 의해
Fitting 아이디얼을 취하는 것이 기저변환과 교환하기 때문이다.
이는 $M_a$가 계수 $k$의 유한 국소 자유 가군이라는 가정
(Algebra, Lemma \ref{algebra-lemma-finite-projective}를 보라)과
이미 인용한 보조정리 \ref{more-algebra-lemma-fitting-ideal-finite-locally-free}에 의해 성립한다.
\end{proof}

\begin{lemma}
\label{more-algebra-lemma-blowup-module}
$R$을 환이라 하고 $M$을 유한 $R$-가군이라 하자.
$f \in R$이 $M_f$를 계수 $r$의 유한 국소 자유 가군이 되게 하는 원소라고 하자.
그러면 $V(f) = V(I)$인 유한 생성 아이디얼 $I \subset R$이 존재하여,
모든 $a \in I$에 대하여 $R' = R[\frac{I}{a}]$로 놓으면 엄밀 변환
$$
M' = (M \otimes_R R')/a\text{-power torsion}
$$
은 계수 $r$의 국소 자유 가군이다.
\end{lemma}

\begin{proof}
전사 $R^{\oplus n} \to M$을 택한다.
$R^{\oplus n}/K \to M$이 $f$를 가역화한 뒤 동형이 되게 하는
유한 부분가군 $K \subset \Ker(R^{\oplus n} \to M)$을 택한다.
예를 들어 Algebra, Lemma \ref{algebra-lemma-finite-projective}에 의해
$M_f$가 유한 표시이므로 가능하다.
$M_1 = R^{\oplus n}/K$로 놓고 $M_1$에 대하여 보조정리를 증명할 수 있다고 하자.
$I \subset R$을 이에 대응하는 아이디얼이라 하자.
그러면 $a \in I$에 대하여 사상
$$
M_1' = (M_1 \otimes_R R')/a\text{-power torsion}
\longrightarrow
M' = (M \otimes_R R')/a\text{-power torsion}
$$
은 전사이다. 또한 $R'$에서 $a$를 가역화한 뒤 동형이다.
% P02-E0826: frozen source gives R'_a = R_f; preserve this exact equality.
$R'_a = R_f$이기 때문이다. Algebra, Lemma \ref{algebra-lemma-blowup-in-principal}을 보라.
그러나 $a$는 $M'_1$ 위의 영인자가 아니므로 표시된 사상은 동형이다.
따라서 $M$이 유한 표시 $R$-가군인 경우에 보조정리를 증명하면 충분하다.

\medskip\noindent
$M$이 유한 표시 $R$-가군이라고 가정하자.
% P02-E0827: frozen source places this Fitting ideal in undefined S; preserve its exact ambient-ring span.
그러면 $J = \text{Fit}_r(M) \subset S$는 유한 생성 아이디얼이다.
$I = fJ$가 원하는 조건을 만족한다고 주장한다.

\medskip\noindent
먼저 $V(f) = V(I)$임을 확인하자.
포함 $V(f) \subset V(I)$는 명백하다.
거꾸로 $f \not \in \mathfrak p$이면 보조정리
\ref{more-algebra-lemma-fitting-ideal-generate-locally}에 의해
$\mathfrak p$는 $V(J)$의 원소가 아니다.
따라서 $\mathfrak p \not \in V(fJ) = V(I)$이다.

\medskip\noindent
$a \in I$라 하고 $R' = R[\frac{I}{a}]$로 놓자.
어떤 $b \in J$에 대하여 $a = fb$로 쓸 수 있다.
Algebra, Lemmas \ref{algebra-lemma-affine-blowup}과
\ref{algebra-lemma-blowup-add-principal}에 의해 $J R' = b R'$이고
$b$는 $R'$의 영인자가 아님을 알 수 있다.
$\mathfrak p' \subset R' = R[\frac{I}{a}]$를 소아이디얼이라 하자.
그러면 $JR'_{\mathfrak p'}$는 $b$로 생성된다.
보조정리 \ref{more-algebra-lemma-principal-fitting-ideal}에 의해
$M'_{\mathfrak p'}$는 $r$개의 원소로 생성될 수 있다.
$M'$이 유한이므로, 대응하는 사상 $(R')^{\oplus r} \to M'$이 $g$를 가역화한 뒤
전사가 되게 하는 $m_1, \ldots, m_r \in M'$과
$g \in R'$, $g \not \in \mathfrak p'$가 존재한다.

\medskip\noindent
% P02-E0828: frozen source uses unbound k rather than fixed r; preserve the exact index.
마지막으로 아이디얼 $J' = \text{Fit}_{k - 1}(M')$을 생각하자.
$J' R'_g$는 $m_1, \ldots, m_r$ 사이의 관계식들의 계수로 생성됨을 주목하자
(Fitting 아이디얼과 기저변환의 양립성).
따라서 $J' = 0$임을 보이면 충분하다.
보조정리 \ref{more-algebra-lemma-fitting-ideal-finite-locally-free}를 보라.
% P02-E0826: the second occurrence retains the same frozen localization equality.
$R'_a = R_f$이고(Algebra, Lemma \ref{algebra-lemma-blowup-in-principal}),
$M'_a = M_f$가 계수 $r$의 자유 가군이므로 $J'_a = 0$임을 알 수 있다.
$a$는 $R'$의 영인자가 아니므로 $J' = 0$이라고 결론짓고 증명이 끝난다.
\end{proof}

\section{완비화와 평탄성}
\label{more-algebra-section-completion-flat}

\noindent
이 절에서는 “큰” 평탄 가군의 완비화가 언제 평탄인지를 논의한다.

\begin{lemma}
\label{more-algebra-lemma-ui-completion-direct-sum-into-product}
$R$을 환이라 하자. $I \subset R$을 아이디얼이라 하고 $A$를 집합이라 하자.
$R$이 Noether 환이고 $I$에 관하여 완비라고 가정하자.
직접합의 $I$-진 완비화에서 곱으로 가는 표준 사상
$$
\left(\bigoplus\nolimits_{\alpha \in A} R\right)^\wedge
\longrightarrow
\prod\nolimits_{\alpha \in A} R
$$
이 존재하며, 이 사상은 보편적으로 단사이다.
\end{lemma}

\begin{proof}
정의에 의해 왼쪽의 원소 $x$는 $x = (x_n)$의 형태이고,
$x_n = (x_{n, \alpha}) \in \bigoplus\nolimits_{\alpha \in A} R/I^n$이며
$x_{n, \alpha} = x_{n + 1, \alpha} \bmod I^n$을 만족한다.
$R = R^\wedge$이므로 모든 $\alpha$에 대하여
$x_{n, \alpha} = y_\alpha \bmod I^n$인 $y_\alpha \in R$이 존재함을 알 수 있다.
각 $n$에 대하여 원소 $x_{n, \alpha}$가 영이 아닌 $\alpha$는
유한 개뿐임을 주목하자. 거꾸로 $(y_\alpha) \in \prod_\alpha R$이 주어지고,
각 $n$에 대하여 $y_{\alpha} \bmod I^n$이 영이 아닌 $\alpha$가 유한 개뿐이면,
이는 왼쪽의 원소를 정의한다.
따라서 왼쪽의 원소를 $y_\alpha \in R$로 이루어진 무한 “수렴합”
$\sum_\alpha y_\alpha$로 생각할 수 있다. 여기서 각 $n$에 대하여
$I^n$을 법으로 영이 아닌 $y_\alpha$는 유한 개뿐이다.
% P02-E0829: duplicated English destination preposition; Korean uses one destination particle.
표시된 사상은 이 원소를 곱 안의 원소 $(y_\alpha)$로 보낸다.
특히 이 사상은 단사이다.

\medskip\noindent
$Q$를 유한 $R$-가군이라 하자. 사상
$$
Q \otimes_R \left(\bigoplus\nolimits_{\alpha \in A} R\right)^\wedge
\longrightarrow
Q \otimes_R \left(\prod\nolimits_{\alpha \in A} R\right)
$$
이 단사임을 보여야 한다.
Algebra, Theorem \ref{algebra-theorem-universally-exact-criteria}를 보라.
표시 $R^{\oplus k} \to R^{\oplus m} \to Q \to 0$을 택하고,
이에 대응하는 $Q$의 생성원들을 $q_1, \ldots, q_m \in Q$로 나타내자.
Artin--Rees 보조정리
(Algebra, Lemma \ref{algebra-lemma-Artin-Rees})에 의해
$\Im(R^{\oplus k} \to R^{\oplus m}) \cap (I^N)^{\oplus m}
\subset \Im((I^{N - c})^{\oplus k} \to R^{\oplus m})$인 상수 $c$가 존재한다.
행들이 완전한 다음 도표를 생각하자.
$$
\xymatrix{
\bigoplus_{l = 1}^k \left(\bigoplus\nolimits_{\alpha \in A} R\right)^\wedge
\ar[r] \ar[d] &
\bigoplus_{j = 1}^m \left(\bigoplus\nolimits_{\alpha \in A} R\right)^\wedge
\ar[r] \ar[d] &
Q \otimes_R \left(\bigoplus\nolimits_{\alpha \in A} R\right)^\wedge
\ar[r] \ar[d] &
0 \\
\bigoplus_{l = 1}^k \left(\prod\nolimits_{\alpha \in A} R\right)
\ar[r] &
\bigoplus_{j = 1}^m \left(\prod\nolimits_{\alpha \in A} R\right)
\ar[r] &
Q \otimes_R \left(\prod\nolimits_{\alpha \in A} R\right)
\ar[r] &
0
}
$$
$\bigoplus_{j = 1, \ldots, m}
\left(\bigoplus\nolimits_{\alpha \in A} R\right)^\wedge$의 원소
$\sum_j \sum_\alpha y_{j, \alpha}$를 택하자.
이 원소가 가군
$Q \otimes_R \left(\prod\nolimits_{\alpha \in A} R\right)$에서 영으로 가면,
특히 각 $\alpha$에 대하여 $Q$에서
$\sum_j q_j \otimes y_{j, \alpha} = 0$임을 알 수 있다.
따라서
$(y_{1, \alpha}, \ldots, y_{m, \alpha}) \in \bigoplus_{j = 1, \ldots, m} R$로 가는 원소
$(z_{1, \alpha}, \ldots, z_{k, \alpha}) \in \bigoplus_{l = 1, \ldots, k} R$을
찾을 수 있다. 또한 $j = 1, \ldots, m$에 대하여
$y_{j, \alpha} \in I^{N_\alpha}$이면,
$l = 1, \ldots, k$에 대하여 $z_{l, \alpha} \in I^{N_\alpha - c}$라고
가정해도 된다.
따라서 합 $\sum_l \sum_\alpha z_{l, \alpha}$는 “수렴”하며,
$\bigoplus_{l = 1, \ldots, k}
\left(\bigoplus\nolimits_{\alpha \in A} R\right)^\wedge$의 원소를 정의한다.
이 원소는 처음에 택한 원소 $\sum_j \sum_\alpha y_{j, \alpha}$로 간다.
따라서 오른쪽 수직 화살표는 단사이고 증명이 끝난다.
\end{proof}

\noindent
다음 보조정리는 아래의 보조정리 \ref{more-algebra-lemma-limit-flat}에서도 유도할 수 있다.

\begin{lemma}
\label{more-algebra-lemma-completed-direct-sum-flat}
$R$을 환이라 하자. $I \subset R$을 아이디얼이라 하고 $A$를 집합이라 하자.
$R$이 Noether 환이라고 가정하자. 완비화
$(\bigoplus\nolimits_{\alpha \in A} R)^\wedge$는 평탄 $R$-가군이다.
\end{lemma}

\begin{proof}
$R$의 $I$에 관한 완비화를 $R^\wedge$로 나타내자.
Algebra, Lemma \ref{algebra-lemma-completion-flat}에 의해
$R \to R^\wedge$가 평탄이므로,
$(\bigoplus\nolimits_{\alpha \in A} R)^\wedge$가 평탄 $R^\wedge$-가군임을
증명하면 충분하다
(Algebra, Lemma \ref{algebra-lemma-composition-flat}을 사용한다).
다음 등식이 성립하므로
$$
(\bigoplus\nolimits_{\alpha \in A} R)^\wedge
=
(\bigoplus\nolimits_{\alpha \in A} R^\wedge)^\wedge
$$
$R$을 $R^\wedge$로 바꾸어 $R$이 $I$에 관하여 완비라고 가정해도 된다
(Algebra, Lemma \ref{algebra-lemma-completion-complete}를 보라).
이 경우 보조정리 \ref{more-algebra-lemma-ui-completion-direct-sum-into-product}에 의해
사상
$(\bigoplus\nolimits_{\alpha \in A} R)^\wedge \to \prod_{\alpha \in A} R$은
보편적으로 단사이다. 따라서 Algebra, Lemma \ref{algebra-lemma-ui-flat-domain}에 의해
$\prod_{\alpha \in A} R$이 평탄임을 보이면 충분하다.
Algebra, Proposition \ref{algebra-proposition-characterize-coherent}
(그리고 Algebra, Lemma \ref{algebra-lemma-Noetherian-coherent})에 의해
$\prod_{\alpha \in A} R$이 평탄임을 알 수 있다.
\end{proof}

\begin{lemma}
\label{more-algebra-lemma-tor-strictly-pro-zero}
\begin{reference}
이는 \cite[Lemma 9.9]{quillenhomology}이다. 저자가 명제에서 “strict”라는
말을 빠뜨렸지만, 그 조건을 의도했음은 명백하다는 점에 유의하라.
\end{reference}
$A$를 Noether 환이라 하고 $I$를 $A$의 아이디얼이라 하자.
$M$을 유한 $A$-가군이라 하자. 모든 $p > 0$에 대하여
$c > 0$이 존재하여
$\text{Tor}_p^A(M, A/I^n) \to \text{Tor}_p^A(M, A/I^{n - c})$는
모든 $n \geq c$에 대하여 영사상이다.
\end{lemma}

\begin{proof}
$p = 1$인 경우를 증명하자. 짧은 완전열
$0 \to K \to A^{\oplus t} \to M \to 0$을 택한다.
그러면 $\text{Tor}_1^A(M, A/I^n) = K \cap (I^n)^{\oplus t}/I^nK$이다.
Artin--Rees 보조정리(Algebra, Lemma \ref{algebra-lemma-Artin-Rees})에 의해
상수 $c \geq 0$이 존재하여 $n \geq c$이면
$K \cap (I^n)^{\oplus t} \subset I^{n - c}K$이다.
% P02-E0830: the English sentence fragment lacks a predicate; Korean supplies it.
따라서 $p = 1$인 경우의 결과를 얻는다.
$p > 1$이면
$\text{Tor}_p^A(M, A/I^n) = \text{Tor}^A_{p - 1}(K, A/I^n)$이다.
따라서 보조정리는 귀납법으로 따른다.
\end{proof}

\begin{lemma}
\label{more-algebra-lemma-limit-flat}
$A$를 Noether 환이라 하고 $I$를 $A$의 아이디얼이라 하자.
$(M_n)$을 다음을 만족하는 $A$-가군들의 역계라 하자.
\begin{enumerate}
\item $M_n$은 평탄 $A/I^n$-가군이다.
\item $M_{n + 1} \to M_n$은 전사이다.
\end{enumerate}
그러면 $M = \lim M_n$은 평탄 $A$-가군이고,
모든 유한 $A$-가군 $Q$에 대하여
$Q \otimes_A M = \lim Q \otimes_A M_n$이다.
\end{lemma}

\begin{proof}
먼저 모든 유한 $A$-가군 $Q$에 대하여
$Q \otimes_A M = \lim Q \otimes_A M_n$임을 보이자.
유한 자유 $A$-가군 $F_i$들로 이루어진 분해
$F_2 \to F_1 \to F_0 \to Q \to 0$을 택한다. 그러면
$$
F_2 \otimes_A M_n \to F_1 \otimes_A M_n \to F_0 \otimes_A M_n
$$
은 차수 $0$의 호몰로지가 $Q \otimes_A M_n$이고
차수 $1$의 호몰로지가
$$
\text{Tor}_1^A(Q, M_n) = \text{Tor}_1^A(Q, A/I^n) \otimes_{A/I^n} M_n
$$
인 사슬 복합체이다. $M_n$이 $A/I^n$ 위 평탄이기 때문이다.
보조정리 \ref{more-algebra-lemma-tor-strictly-pro-zero}에 의해
이 계는 본질적으로 상수이며 그 값은 $0$임을 알 수 있다.
Homology, Lemma \ref{homology-lemma-apply-Mittag-Leffler-again}에 의해
$\lim Q \otimes_A M_n =
\Coker(\lim F_1 \otimes_A M_n \to \lim F_0 \otimes_A M_n)$이다.
$F_i$가 유한 자유이므로 이는
$\Coker(F_1 \otimes_A M \to F_0 \otimes_A M) = Q \otimes_A M$과 같다.

\medskip\noindent
다음으로 $Q \to Q'$를 유한 $A$-가군들 사이의 단사 사상이라 하자.
$Q \otimes_A M \to Q' \otimes_A M$이 단사임을 보여야 한다
(Algebra, Lemma \ref{algebra-lemma-flat}).
위의 결과에 의해 다음을 얻는다.
$$
\Ker(Q \otimes_A M \to Q' \otimes_A M) =
\Ker(\lim Q \otimes_A M_n \to \lim Q' \otimes_A M_n).
$$
각 $n$에 대하여 완전열
$$
\text{Tor}_1^A(Q', M_n) \to \text{Tor}_1^A(Q'', M_n) \to
Q \otimes_A M_n \to Q' \otimes_A M_n
$$
이 존재한다. 여기서 $Q'' = \Coker(Q \to Q')$이다.
위에서 Tor의 역계들이 본질적으로 값 $0$의 상수임을 보았다.
Homology, Lemma \ref{homology-lemma-apply-Mittag-Leffler-again}에 의해
가장 오른쪽 사상들의 역극한은 단사이다.
\end{proof}

\begin{lemma}
\label{more-algebra-lemma-flat-after-completion}
$R$을 환이라 하자. $I \subset R$을 아이디얼이라 하고 $M$을 $R$-가군이라 하자.
다음을 가정하자.
\begin{enumerate}
\item $I$는 유한 생성이다.
\item $R/I$는 Noether 환이다.
\item $M/IM$은 $R/I$ 위 평탄이다.
\item $\text{Tor}_1^R(M, R/I) = 0$이다.
\end{enumerate}
그러면 $I$-진 완비화 $R^\wedge$는 Noether 환이고,
$M^\wedge$는 $R^\wedge$ 위 평탄이다.
\end{lemma}

\begin{proof}
Algebra, Lemma \ref{algebra-lemma-what-does-it-mean}에 의해
모든 $n$에 대하여 가군 $M/I^nM$은 $R/I^n$ 위 평탄이다.
Algebra, Lemma \ref{algebra-lemma-hathat-finitely-generated}에 의해
(a) $R^\wedge$와 $M^\wedge$는 $I$-진 완비이고,
(b) 모든 $n$에 대하여 $R/I^n = R^\wedge/I^nR^\wedge$이다.
Algebra, Lemma \ref{algebra-lemma-completion-Noetherian}에 의해
환 $R^\wedge$는 Noether 환이다.
보조정리 \ref{more-algebra-lemma-limit-flat}을 적용하면
$M^\wedge = \lim M/I^nM$은 $R^\wedge$-가군으로서 평탄이라고 결론짓는다.
\end{proof}

\section{Koszul 복합체}
\label{more-algebra-section-koszul}

\noindent
Koszul 복합체를 다음과 같이 정의한다.

\begin{definition}
\label{more-algebra-definition-koszul}
$R$을 환이라 하고 $\varphi : E \to R$을 $R$-가군 사상이라 하자.
$\varphi$에 대응하는 {\it Koszul 복합체} $K_\bullet(\varphi)$는
다음과 같이 정의되는 가환 미분 등급 대수이다.
\begin{enumerate}
\item 바탕 등급 대수는 외대수 $K_\bullet(\varphi) = \wedge(E)$이다.
\item 미분 $d : K_\bullet(\varphi) \to K_\bullet(\varphi)$는
모든 $e \in E = K_1(\varphi)$에 대하여 $d(e) = \varphi(e)$를 만족하는
유일한 미분자(derivation)이다.
\end{enumerate}
\end{definition}

\noindent
구체적으로, $e_1 \wedge \ldots \wedge e_n$이
$K_\bullet(\varphi)$의 차수 $n$인 생성원 중 하나이면
$$
d(e_1 \wedge \ldots \wedge e_n) =
\sum\nolimits_{i = 1, \ldots, n} (-1)^{i + 1}
\varphi(e_i)e_1 \wedge \ldots \wedge \widehat{e_i} \wedge \ldots \wedge e_n.
$$
이것이 텐서 대수 위의 잘 정의된 미분자를 주며, $e \otimes e$를 영으로 보내므로
외대수를 거쳐 분해됨을 쉽게 확인할 수 있다.

\medskip\noindent
흔히 $E$가 유한 자유 가군이라고 가정하며, $E = R^{\oplus n}$이라 하자.
이 경우 사상 $\varphi$는 원소들의 열 $f_1, \ldots, f_n \in R$로 주어진다.

\begin{definition}
\label{more-algebra-definition-koszul-complex}
$R$을 환이라 하고 $f_1, \ldots, f_r \in R$이라 하자.
{\it $f_1, \ldots, f_r$에 대한 Koszul 복합체}는 사상
$(f_1, \ldots, f_r) : R^{\oplus r} \to R$에 대응하는 Koszul 복합체이다.
이를 $K_\bullet(f_\bullet)$, $K_\bullet(f_1, \ldots, f_r)$,
$K_\bullet(R, f_1, \ldots, f_r)$ 또는 $K_\bullet(R, f_\bullet)$로 나타낸다.
\end{definition}

\noindent
물론 $E$가 유한 국소 자유이면, $K_\bullet(\varphi)$는 $\Spec(R)$ 위에서
국소적으로 Koszul 복합체 $K_\bullet(f_1, \ldots, f_r)$와 동형이다.
이 복합체에는 흥미로운 형식적 성질이 많다.

\begin{lemma}
\label{more-algebra-lemma-functorial}
$\varphi : E \to R$과 $\varphi' : E' \to R$을 $R$-가군 사상이라 하자.
$\psi : E \to E'$를 $\varphi' \circ \psi = \varphi$인 $R$-가군 사상이라 하자.
그러면 $\psi$는 미분 등급 대수의 준동형
$K_\bullet(\varphi) \to K_\bullet(\varphi')$를 유도한다.
\end{lemma}

\begin{proof}
정의에서 바로 따른다.
\end{proof}

\begin{lemma}
\label{more-algebra-lemma-change-basis}
$f_1, \ldots, f_r \in R$을 원소들의 열이라 하자.
$(x_{ij})$를 $R$에 계수를 갖는 가역 $r \times r$ 행렬이라 하자.
그러면 복합체 $K_\bullet(f_\bullet)$와
$$
K_\bullet(\sum x_{1j}f_j, \sum x_{2j}f_j, \ldots, \sum x_{rj}f_j)
$$
는 동형이다.
\end{lemma}

\begin{proof}
$g_i = \sum x_{ij}f_j$로 놓자.
행렬 $(x_{ji})$는 $(g_1, \ldots, g_r) = (f_1, \ldots, f_r) \circ x$인 동형
$x : R^{\oplus r} \to R^{\oplus r}$를 준다.
따라서 이는 보조정리 \ref{more-algebra-lemma-functorial}에서 설명한
Koszul 복합체의 함자성에서 따른다.
\end{proof}

\begin{lemma}
\label{more-algebra-lemma-homotopy-koszul-abstract}
$R$을 환이라 하고 $\varphi : E \to R$을 $R$-가군 사상이라 하자.
$R$에서의 상이 $f = \varphi(e)$인 $e \in E$를 택하자.
그러면 $K_\bullet(\varphi)$의 자기준동형으로서
$$
f = de + ed
$$
가 성립한다.
\end{lemma}

\begin{proof}
$d(ea) = d(e)a - ed(a) = fa - ed(a)$이므로 성립한다.
\end{proof}

\begin{lemma}
\label{more-algebra-lemma-homotopy-koszul}
$R$을 환이라 하고 $f_1, \ldots, f_r \in R$을 원소들의 열이라 하자.
$K_\bullet(f_\bullet)$ 위에서 $f_i$를 곱하는 사상은 영사상과 호모토픽이다.
% P02-E0831: frozen source says cohomology for homological H_i; preserve that claim.
특히 코호몰로지 가군 $H_i(K_\bullet(f_\bullet))$는
아이디얼 $(f_1, \ldots, f_r)$에 의해 소멸된다.
\end{lemma}

\begin{proof}
보조정리 \ref{more-algebra-lemma-homotopy-koszul-abstract}의 특수한 경우이다.
\end{proof}

\noindent
Derived Categories, Section \ref{derived-section-cones}에서
코사슬 복합체 사상의 원뿔을 정의하였다.
사슬 복합체 사상 $f : A_\bullet \to B_\bullet$의 원뿔 $C(f)_\bullet$는
$C(f)_n = B_n \oplus A_{n - 1}$과 다음 미분으로 주어지는 복합체 $C(f)_\bullet$이다.
\begin{equation}
\label{more-algebra-equation-differential-cone}
d_{C(f), n} =
\left(
\begin{matrix}
d_{B, n} & f_{n - 1} \\
0 & -d_{A, n - 1}
\end{matrix}
\right)
\end{equation}
이 복합체에는 자명한 사상들 $B_n \to C(f)_n \to A_{n - 1}$에서 유도되는
표준 복합체 사상 $i : B_\bullet \to C(f)_\bullet$와
$p : C(f)_\bullet \to A_\bullet[-1]$가 주어진다.

\begin{lemma}
\label{more-algebra-lemma-cone-koszul-abstract}
$R$을 환이라 하고 $\varphi : E \to R$을 $R$-가군 사상이라 하자.
$f \in R$이라 하자. $E' = E \oplus R$로 놓고,
$E$ 위에서는 $\varphi$로, $R$ 위에서는 $f$를 곱하는 사상으로
$\varphi' : E' \to R$을 정의하자.
복합체 $K_\bullet(\varphi')$는 다음 복합체 사상의 원뿔과 동형이다.
$$
f :
K_\bullet(\varphi)
\longrightarrow
K_\bullet(\varphi).
$$
\end{lemma}

\begin{proof}
% P02-E0832: retain the frozen reversed direct-sum order in the displayed inclusion.
원소 $1 \in R \subset R \oplus E$를 $e_0 \in E'$로 나타내자.
위의 원뿔 정의에 의해 다음을 얻는다.
$$
C(f)_n = K_n(\varphi) \oplus K_{n - 1}(\varphi) =
\wedge^n(E) \oplus \wedge^{n - 1}(E) = \wedge^n(E')
$$
여기서 마지막 $=$에서는 $(0, e_1 \wedge \ldots \wedge e_{n - 1})$을
$\wedge^n(E')$의 원소 $e_0 \wedge e_1 \wedge \ldots \wedge e_{n - 1}$로 보낸다.
계산하면 이 동형이 미분과 양립함을 알 수 있다.
첫 번째 직접합인자의 원소들에 대해서는,
$\varphi'|_E = \varphi$이고 $d_{C(f)}$를 첫 번째 직접합인자에 제한한 것이
바로 $d_{K_\bullet(\varphi)}$이므로 명백하다.
한편 $e_1 \wedge \ldots \wedge e_{n - 1}$이 두 번째 직접합인자에 속하면
$$
d_{C(f)}(0, e_1 \wedge \ldots \wedge e_{n - 1}) =
fe_1 \wedge \ldots \wedge e_{n - 1}
- d_{K_\bullet(\varphi)}(e_1 \wedge \ldots \wedge e_{n - 1})
$$
이고, 다른 한편으로는 다음이 성립한다.
% P02-E0833: retain the frozen pair-like argument of the Koszul differential.
\begin{align*}
& d_{K_\bullet(\varphi')}(0, e_0 \wedge e_1 \wedge \ldots \wedge e_{n - 1}) \\
& =
\sum\nolimits_{i = 0, \ldots, n - 1}
(-1)^i \varphi'(e_i)e_0 \wedge \ldots \wedge \widehat{e_i}
\wedge \ldots \wedge e_{n - 1} \\
& =
fe_1 \wedge \ldots \wedge e_{n - 1} +
\sum\nolimits_{i = 1, \ldots, n - 1}
(-1)^i \varphi(e_i)e_0 \wedge \ldots \wedge \widehat{e_i}
\wedge \ldots \wedge e_{n - 1} \\
& =
fe_1 \wedge \ldots \wedge e_{n - 1} -
e_0 \left(\sum\nolimits_{i = 1, \ldots, n - 1}
(-1)^{i + 1} \varphi(e_i)e_1 \wedge \ldots \wedge \widehat{e_i}
\wedge \ldots \wedge e_{n - 1}\right)
\end{align*}
이는 앞선 계산 결과의 상이다.
\end{proof}

\begin{lemma}
\label{more-algebra-lemma-cone-koszul}
$R$을 환이라 하고 $f_1, \ldots, f_r$을 $R$의 원소들의 열이라 하자.
복합체 $K_\bullet(f_1, \ldots, f_r)$는 다음 복합체 사상의 원뿔과 동형이다.
$$
f_r :
K_\bullet(f_1, \ldots, f_{r - 1})
\longrightarrow
K_\bullet(f_1, \ldots, f_{r - 1}).
$$
\end{lemma}

\begin{proof}
보조정리 \ref{more-algebra-lemma-cone-koszul-abstract}의 특수한 경우이다.
\end{proof}

\begin{lemma}
\label{more-algebra-lemma-cone-squared}
$R$을 환이라 하고 $A_\bullet$을 $R$-가군들의 복합체라 하자.
$f, g \in R$이라 하자. $C(f)_\bullet$를
$f : A_\bullet \to A_\bullet$의 원뿔이라 하자.
$C(g)_\bullet$와 $C(fg)_\bullet$도 마찬가지로 정의한다.
그러면 $C(fg)_\bullet$는 다음 사상의 원뿔과 호모토피 동치이다.
$$
C(f)_\bullet[1] \longrightarrow C(g)_\bullet
$$
\end{lemma}

\begin{proof}
먼저 $A_\bullet$이 차수 $0$에 $R$을 놓은 복합체인 경우를 증명하자.
이 경우 복합체 $C(f)_\bullet$는
$$
\ldots \to 0 \to R \xrightarrow{f} R \to 0 \to \ldots
$$
이며, $R$은 (호몰로지) 차수 $1$과 $0$에 놓인다.
사용할 복합체 사상은 다음과 같다.
$$
\xymatrix{
0 \ar[r] \ar[d] &
0 \ar[r] \ar[d] &
R \ar[r]^f \ar[d]^1 &
R \ar[r] \ar[d] & 0 \ar[d] \\
0 \ar[r] & R \ar[r]^g & R \ar[r] & 0 \ar[r] & 0
}
$$
이 사상의 원뿔은 차수 $1$과 $0$에 $R^{\oplus 2}$를 놓고
미분 (\ref{more-algebra-equation-differential-cone})이
$$
\left(
\begin{matrix}
g & 1 \\
0 & -f
\end{matrix}
\right) :
R^{\oplus 2} \longrightarrow R^{\oplus 2}
$$
인 사슬 복합체이다.
이 사슬 복합체가 $C(fg)_\bullet$, 즉 $R \xrightarrow{fg} R$과
호모토픽임을 보이기 위해 다음 복합체 사상들을 생각하자.
$$
\xymatrix{
R \ar[d]_{(1, -g)} \ar[r]_{fg} & R \ar[d]^{(0, 1)} \\
R^{\oplus 2} \ar[r] & R^{\oplus 2}
}
\quad\quad
\xymatrix{
R^{\oplus 2} \ar[d]_{(1, 0)} \ar[r] & R^{\oplus 2} \ar[d]^{(f, 1)} \\
R \ar[r]^{fg} & R
}
$$
표기의 의미는 자명하다.
이 두 사상의 한 방향 합성은 $C(fg)_\bullet$ 위의 항등사상이지만,
반대 방향의 합성은 항등사상이 아니다.
필요한 호모토피를 구체적으로 쓰는 것은 생략한다.

\medskip\noindent
일반적인 경우를 보이기 위해 복합체 범주 위의 함자
$K_\bullet \mapsto \text{Tot}(A_\bullet \otimes_R K_\bullet)$가
호모토피 및 원뿔과 양립함을 이용한다.
이제 $C(f)_\bullet$와 $C(g)_\bullet$를 각각 이중 복합체
$$
(R \xrightarrow{f} R) \otimes_R A_\bullet
\quad\text{and}\quad
(R \xrightarrow{g} R) \otimes_R A_\bullet
$$
의 총복합체로 쓰면, 위에서 논의한 특수한 경우로부터 결과를 얻는다.
일부 세부사항은 생략한다.
\end{proof}

\begin{lemma}
\label{more-algebra-lemma-koszul-mult-abstract}
$R$을 환이라 하고 $\varphi : E \to R$을 $R$-가군 사상이라 하자.
$f, g \in R$이라 하자. $E' = E \oplus R$로 놓고,
$E$ 위에서는 $\varphi$로, $R$ 위에서는 각각 $f, g, fg$를 곱하는 사상으로
$\varphi'_f, \varphi'_g, \varphi'_{fg} : E' \to R$을 정의하자.
복합체 $K_\bullet(\varphi'_{fg})$는 다음 복합체 사상의 원뿔과 호모토피 동치이다.
$$
K_\bullet(\varphi'_f)[1]
\longrightarrow
K_\bullet(\varphi'_g).
$$
\end{lemma}

\begin{proof}
보조정리 \ref{more-algebra-lemma-cone-koszul-abstract}에 의해
복합체 $K_\bullet(\varphi'_f)$는 $K_\bullet(\varphi)$ 위에서 $f$를 곱하는 사상의
원뿔과 동형이고, 다른 두 경우도 마찬가지이다.
따라서 보조정리는 보조정리 \ref{more-algebra-lemma-cone-squared}에서 따른다.
\end{proof}

\begin{lemma}
\label{more-algebra-lemma-koszul-mult}
$R$을 환이라 하고 $f_1, \ldots, f_{r - 1}$을 $R$의 원소들의 열이라 하자.
$f, g \in R$이라 하자.
복합체 $K_\bullet(f_1, \ldots, f_{r - 1}, fg)$는
다음 복합체 사상의 원뿔과 호모토피 동치이다.
$$
K_\bullet(f_1, \ldots, f_{r - 1}, f)[1]
\longrightarrow
K_\bullet(f_1, \ldots, f_{r - 1}, g)
$$
\end{lemma}

\begin{proof}
보조정리 \ref{more-algebra-lemma-koszul-mult-abstract}의 특수한 경우이다.
\end{proof}

\begin{lemma}
\label{more-algebra-lemma-join-sequences-koszul-complex}
$R$을 환이라 하자.
$f_1, \ldots, f_r$, $g_1, \ldots, g_s$를 $R$의 원소들이라 하자.
그러면 Koszul 복합체 사이에 다음 동형이 존재한다.
$$
K_\bullet(R, f_1, \ldots, f_r, g_1, \ldots, g_s) =
\text{Tot}(K_\bullet(R, f_1, \ldots, f_r) \otimes_R
K_\bullet(R, g_1, \ldots, g_s)).
$$
\end{lemma}

\begin{proof}
생략한다. 힌트: $K_\bullet(R, f_1, \ldots, f_r)$가 미분 등급 대수로서
$\text{d}(x_i) = f_i$인 $x_1, \ldots, x_r$로 생성되고,
$K_\bullet(R, g_1, \ldots, g_s)$가 미분 등급 대수로서
$\text{d}(y_j) = g_j$인 $y_1, \ldots, y_s$로 생성된다면,
$K_\bullet(R, f_1, \ldots, f_r, g_1, \ldots, g_s)$를
$\text{d}(x_i) = f_i$와 $\text{d}(y_j) = g_j$를 만족하는 원소들의 열
$x_1, \ldots, x_r, y_1, \ldots, y_s$로 생성되는 미분 등급 대수로 생각할 수 있다.
\end{proof}

\section{확장 교대 {\v C}ech 복합체}
\label{more-algebra-section-alternating-cech}

\noindent
$R$을 환이라 하고 $f_1, \ldots, f_r \in R$이라 하자.
$R$의 확장 교대 {\v C}ech 복합체는 다음 코사슬 복합체이다.
$$
R \to \bigoplus\nolimits_{i_0} R_{f_{i_0}} \to
\bigoplus\nolimits_{i_0 < i_1} R_{f_{i_0}f_{i_1}} \to
\ldots \to R_{f_1\ldots f_r}
$$
여기서 $R$은 차수 $0$에 놓이고,
% P02-E0834: the English source misspells degree; Korean supplies the standard word.
항 $\bigoplus_{i_0} R_{f_{i_0}}$는 차수 $1$에 놓이며, 이후도 마찬가지이다.
사상들은 다음과 같이 정의한다.
\begin{enumerate}
\item 사상 $R \to \bigoplus\nolimits_{i_0} R_{f_{i_0}}$는
표준 사상들 $R \to R_{f_{i_0}}$로 주어진다.
\item $1 \leq i_0 < \ldots < i_{p + 1} \leq r$이고
$0 \leq j \leq p + 1$이면 표준 국소화 사상
$$
R_{f_{i_0} \ldots \hat f_{i_j} \ldots f_{i_{p + 1}}} \to
R_{f_{i_0} \ldots f_{i_{p + 1}}}
$$
이 존재한다.
\item 미분은 (2)의 표준 사상들에 부호 $(-1)^j$를 붙여 사용한다.
\end{enumerate}
$M$이 임의의 $R$-가군이면, $M$의 확장 교대 {\v C}ech 복합체는
마찬가지로 구성되는 다음 코사슬 복합체이다.
$$
M \to \bigoplus\nolimits_{i_0} M_{f_{i_0}} \to
\bigoplus\nolimits_{i_0 < i_1} M_{f_{i_0}f_{i_1}} \to
\ldots \to M_{f_1\ldots f_r}
$$
여기서도 $M$은 차수 $0$에 놓인다.

\begin{lemma}
\label{more-algebra-lemma-extended-alternating-is-complex}
위에서 정의한 확장 교대 {\v C}ech 복합체들은 $R$-가군들의 복합체이다.
\end{lemma}

\begin{proof}
생략한다.
\end{proof}

\begin{lemma}
\label{more-algebra-lemma-extended-alternating-form-module}
$R$을 환이라 하고 $f_1, \ldots, f_r \in R$이라 하자.
$M$을 $R$-가군이라 하자. $M$의 확장 교대 {\v C}ech 복합체는
$M$과 $R$의 확장 교대 {\v C}ech 복합체를 $R$ 위에서 취한 텐서곱이다.
\end{lemma}

\begin{proof}
생략한다.
\end{proof}

\begin{lemma}
\label{more-algebra-lemma-extended-alternating-base-change}
$R$을 환이라 하고 $f_1, \ldots, f_r \in R$이라 하자.
$M$을 $R$-가군이라 하자. $R \to S$를 환 준동형이라 하고,
$f_1, \ldots, f_r$의 상을 $g_1, \ldots, g_r \in S$로 나타내며
$N = M \otimes_R S$로 놓자.
$S$, $g_1, \ldots, g_r$, $N$을 사용하여 구성한 확장 교대 {\v C}ech 복합체는
$M$의 확장 교대 {\v C}ech 복합체와 $S$를 $R$ 위에서 취한 텐서곱이다.
\end{lemma}

\begin{proof}
생략한다.
\end{proof}

\begin{lemma}
\label{more-algebra-lemma-extended-alternating-homotopy-zero}
$R$을 환이라 하고 $f_1, \ldots, f_r \in R$이라 하자.
$M$을 $R$-가군이라 하자. $f_i$가 단위원인
$i \in \{1, \ldots, r\}$가 존재하면,
$M$의 확장 교대 {\v C}ech 복합체는 $0$과 호모토피 동치이다.
\end{lemma}

\begin{proof}
다음 표기를 사용한다. $M$의 확장 교대 {\v C}ech 복합체에서
차수 $p + 1$인 코사슬 $x$는 $x = (x_{i_0 \ldots i_p})$로 쓰며,
$x_{i_0 \ldots i_p}$는 $M_{f_{i_0} \ldots f_{i_p}}$에 속한다.
이 표기로 다음을 얻는다.
$$
d(x)_{i_0 \ldots i_{p + 1}} =
\sum\nolimits_j (-1)^j x_{i_0 \ldots \hat i_j \ldots i_{p + 1}}
$$
호모토피로는 다음 규칙으로 주어지는 사상들을 사용한다.
$$
h : \text{cochains of degree }p + 2 \to \text{cochains of degree }p + 1
$$
$$
h(x)_{i_0 \ldots i_p} = 0 \text{ if } i \in \{i_0, \ldots, i_p\}
\text{ and }
h(x)_{i_0 \ldots i_p} = 
(-1)^j x_{i_0 \ldots i_j i i_{j + 1} \ldots i_p} \text{ if not}
$$
% P02-E0835: frozen rule omits endpoint conventions for the unique index j.
둘째 경우에 $j$는 $i_j < i < i_{j + 1}$을 만족하는 유일한 지표이다.
또한 $f_i$가 단위원이므로 다음 등식이 성립한다.
$$
M_{f_{i_0} \ldots f_{i_p}} =
M_{f_{i_0} \ldots f_{i_j} f_i f_{i_{j + 1}} \ldots f_{i_p}}
$$
이를 이용하면
$(-1)^j x_{i_0 \ldots i_j i i_{j + 1} \ldots i_p}$를
$M_{f_{i_0} \ldots f_{i_p}}$의 원소로 생각할 수 있다.
계산을 통해 $d h + h d$가 항등사상의 음수와 같음을 보일 것이며,
그러면 증명이 끝난다. 이를 위해 차수 $p + 1$인 코사슬 $x$를 고정하고
$1 \leq i_0 < \ldots < i_p \leq r$이라 하자.

\medskip\noindent
경우 I: $i \in \{i_0, \ldots, i_p\}$.
$i = i_t$라 하자. 그러면 $h(d(x))_{i_0 \ldots i_p} = 0$이다.
한편
$$
d(h(x))_{i_0 \ldots i_p} =
\sum (-1)^j h(x)_{i_0 \ldots \hat i_j \ldots i_p} =
(-1)^t h(x)_{i_0 \ldots \hat i \ldots i_p} =
(-1)^t (-1)^{t - 1} x_{i_0 \ldots i_p}
$$
이다. 따라서 원하는 대로
$(dh + hd)(x)_{i_0 \ldots i_p} = -x_{i_0 \ldots i_p}$이다.

\medskip\noindent
경우 II: $i \not \in \{i_0, \ldots, i_p\}$.
$i_j < i < i_{j + 1}$인 $j$를 택하자. 그러면
\begin{align*}
h(d(x))_{i_0 \ldots i_p}
& =
(-1)^j d(x)_{i_0 \ldots i_j i i_{j + 1} \ldots i_p} \\
& =
\sum\nolimits_{j' \leq j} (-1)^{j + j'}
x_{i_0 \ldots \hat i_{j'} \ldots i_j i i_{j + 1} \ldots i_p} -
x_{i_0 \ldots i_p} \\
&
+ \sum\nolimits_{j' > j} (-1)^{j + j' + 1}
x_{i_0 \ldots i_j i i_{j + 1} \ldots \hat i_{j'} \ldots i_p}
\end{align*}
이다. 한편
\begin{align*}
d(h(x))_{i_0 \ldots i_p}
& =
\sum\nolimits_{j'} (-1)^{j'} h(x)_{i_0 \ldots \hat i_{j'} \ldots i_p} \\
& =
\sum\nolimits_{j' \leq j} (-1)^{j' + j - 1}
x_{i_0 \ldots \hat i_{j'} \ldots i_j i i_{j + 1} \ldots i_p} \\
& +
\sum\nolimits_{j' > j} (-1)^{j' + j}
x_{i_0 \ldots i_j i i_{j + 1} \ldots \hat i_{j'} \ldots i_p}
\end{align*}
이다. 이들을 더하면 원하는 대로
$(dh + hd)(x)_{i_0 \ldots i_p} = - x_{i_0 \ldots i_p}$를 얻는다.
\end{proof}

\begin{lemma}
\label{more-algebra-lemma-extended-alternating-torsion}
$R$을 환이라 하고 $f_1, \ldots, f_r \in R$이라 하자.
$M$을 $R$-가군이라 하자.
% P02-E0836: the source uses alternation rather than its defined alternating name.
$H^q$를 $M$의 확장 교대 {\v C}ech 복합체의 $q$차 코호몰로지 가군이라 하자.
그러면 다음이 성립한다.
\begin{enumerate}
\item $q \not \in [0, r]$이면 $H^q = 0$이다.
% P02-E0837: retain the frozen H^i and generator-index reuse.
\item $x \in H^i$이면 $i = 1, \ldots, r$에 대하여
$f_i^n x = 0$인 $n \geq 1$이 존재한다.
\item $H^q$의 지지는 $V(f_1, \ldots, f_r)$에 포함된다.
\item $M$ 위에서 가역적으로 작용하는 $f \in (f_1, \ldots, f_r)$가
존재하면 $H^q = 0$이다.
\end{enumerate}
\end{lemma}

\begin{proof}
(1)은 확장 교대 {\v C}ech 복합체가 차수 $< 0$과 $> r$에서
영이라는 사실에서 따른다. (2)를 증명하려면 각 $i$에 대하여
$f_i^n x = 0$인 $n \geq 1$이 존재함을 보이면 충분하다.
이를 위해서는 $(H^q)_{f_i} = 0$임을 보이면 충분하다.
국소화는 완전하므로, $(H^q)_{f_i}$는 $M$의 확장 교대 복합체를
$f_i$에서 국소화한 복합체의 $q$차 코호몰로지 가군이다.
보조정리 \ref{more-algebra-lemma-extended-alternating-base-change}에 의해
이 국소화는 $R_{f_i}$ 위의 $M_{f_i}$에 대하여
$f_1, \ldots, f_r$의 $R_{f_i}$ 안에서의 상들을 사용한
확장 교대 {\v C}ech 복합체이다.
따라서 $f_i$가 가역이면 $H^q$가 영임을 보이는 문제로 환원되며,
이는 보조정리 \ref{more-algebra-lemma-extended-alternating-homotopy-zero}에서 따른다.
(3)은 방금 증명한 모든 $i$에 대한 $(H^q)_{f_i} = 0$에서 따른다.
(4)를 보려면, 이 경우 $f$가 $H^q$ 위에서 가역적으로 작용하고,
(3)에 의해 $H^q$가 $V(f)$ 위에 지지됨을 주목하라.
이는 $H^q$가 영임을 강제한다(작은 세부사항은 생략한다).
\end{proof}

\begin{lemma}
\label{more-algebra-lemma-extended-alternating-Cech-is-colimit-koszul}
$R$을 환이라 하고 $f_1, \ldots, f_r \in R$이라 하자.
확장 교대 {\v C}ech 복합체
$$
R \to \bigoplus\nolimits_{i_0} R_{f_{i_0}} \to
\bigoplus\nolimits_{i_0 < i_1} R_{f_{i_0}f_{i_1}} \to
\ldots \to R_{f_1\ldots f_r}
$$
는 Koszul 복합체들 $K(R, f_1^n, \ldots, f_r^n)$의 쌍대극한이다.
정확한 명제는 증명을 보라.
\end{lemma}

\begin{proof}
독자에게 직접 증명해 보기를 권한다.
정의 \ref{more-algebra-definition-koszul-complex}의 Koszul 복합체를
차수 $0, \ldots, r$에 놓인 코사슬 복합체로 보고
$K(R, f_1^n, \ldots, f_r^n)$으로 나타내자. 따라서 다음을 얻는다.
$$
K(R, f_1^n, \ldots, f_r^n) :
0 \to \wedge^r(R^{\oplus r}) \to
\wedge^{r - 1}(R^{\oplus r}) \to \ldots \to
R^{\oplus r} \to R \to 0
$$
여기서 항 $\wedge^r(R^{\oplus r})$는 차수 $0$에 놓인다.
$e^n_1, \ldots, e^n_r$을 $R^{\oplus r}$의 표준 기저라 하자.
그러면 $1 \leq j_1 < \ldots < j_{r - p} \leq r$에 대한 원소들
$e^n_{j_1} \wedge \ldots \wedge e^n_{j_{r - p}}$는
Koszul 복합체의 차수 $p$인 항의 기저를 이룬다. 또한
$$
d(e^n_{j_1} \wedge \ldots \wedge e^n_{j_{r - p}}) =
\sum (-1)^{a + 1} f_{j_a}^n e^n_{j_1} \wedge \ldots \wedge \hat e^n_{j_a} \wedge
\ldots \wedge e^n_{j_{r - p}}
$$
임을 주목하자. 이는 Section \ref{more-algebra-section-koszul}에서의
Koszul 복합체 구성에서 따른다.
우리 계의 전이사상
$$
K(R, f_1^n, \ldots, f_r^n) \to K(R, f_1^{n + 1}, \ldots, f_r^{n + 1})
$$
은 다음 규칙으로 주어진다.
$$
e^n_{j_1} \wedge \ldots \wedge e^n_{j_{r - p}}
\longmapsto
f_{i_0} \ldots f_{i_{p - 1}}
e^{n + 1}_{j_1} \wedge \ldots \wedge e^{n + 1}_{j_{r - p}}
$$
여기서 지표 $1 \leq i_0 < \ldots < i_{p - 1} \leq r$는 다음을 만족한다.
% P02-E0838: retain the first frozen missing comma in the finite-set expression.
$$
\{1, \ldots r\} =
\{i_0, \ldots, i_{p - 1}\} \amalg \{j_1, \ldots, j_{r - p}\}
$$
이것이 미분과 양립함을 보이는 짧은 계산은 생략한다.
전이사상은 차수 $0$에서는 항상 $1$이고
차수 $r$에서는 $f_1 \ldots f_r$와 같음을 주목하자.

\medskip\noindent
Koszul 복합체의 차수 $p$인 항을
$K^p(R, f_1^n, \ldots, f_r^n)$으로 나타내자.
모든 $f \in R$에 대하여
$$
R_f = \colim (R \xrightarrow{f} R \xrightarrow{f} R \to \ldots )
$$
임을 주목하자. 따라서 차수 $p$에서는 다음을 얻는다.
% P02-E0839: retain the frozen missing comma in the Koszul argument list.
$$
\colim K^p(R, f_1^n, \ldots f_r^n) =
\bigoplus\nolimits_{1 \leq i_0 < \ldots < i_{p - 1} \leq r}
R_{f_{i_0} \ldots f_{i_{p - 1}}}
$$
위 Koszul 복합체의 원소
$e^n_{j_1} \wedge \ldots \wedge e^n_{j_{r - p}}$는
쌍대극한에서 직접합인자 $R_{f_{i_0} \ldots f_{i_{p - 1}}}$의 원소
$(f_{i_0} \ldots f_{i_{p - 1}})^{-n}$으로 간다.
여기서 지표들은 다음을 만족하도록 택한다.
% P02-E0838: retain the second frozen missing comma in the finite-set expression.
$$
\{1, \ldots r\} = \{i_0, \ldots, i_{p - 1}\}
\amalg \{j_1, \ldots, j_{r - p}\}
$$
따라서 이 복합체 위의 미분은 다음과 같이 주어진다.
$$
d(1\text{ in }R_{f_{i_0} \ldots f_{i_{p - 1}}}) =
\sum\nolimits_{i \not \in \{i_0, \ldots, i_{p - 1}\}}
(-1)^{i - t}\text{ in }
R_{f_{i_0} \ldots f_{i_t} f_i f_{i_{t + 1}} \ldots f_{i_{p - 1}}}
$$
이제 차수 $p$에서 다음 사상으로 주어지는 복합체 사상을 생각하자.
$$
\bigoplus\nolimits_{1 \leq i_0 < \ldots < i_{p - 1} \leq r}
R_{f_{i_0} \ldots f_{i_{p - 1}}}
\longrightarrow
\bigoplus\nolimits_{1 \leq i_0 < \ldots < i_{p - 1} \leq r}
R_{f_{i_0} \ldots f_{i_{p - 1}}}
$$
이는 다음 규칙으로 정해진다.
$$
1\text{ in }R_{f_{i_0} \ldots f_{i_{p - 1}}}
\longmapsto
(-1)^{i_0 + \ldots + i_{p - 1} + p}\text{ in }R_{f_{i_0} \ldots f_{i_{p - 1}}}
$$
그러면 $\colim K(R, f_1^n, \ldots, f_r^n)$에서 이 절에서 정의한
확장 교대 {\v C}ech 복합체로 가는 복합체 동형을 얻는다.
부호가 맞아떨어지는지 확인하는 과정은 생략한다.
\end{proof}

\section{Koszul 정칙열}
\label{more-algebra-section-koszul-regular}

\noindent
이 절을 보기 전에 Algebra, Sections
\ref{algebra-section-regular-sequences},
\ref{algebra-section-quasi-regular}, \ref{algebra-section-depth}를 보라.

\begin{definition}
\label{more-algebra-definition-koszul-regular-sequence}
$R$을 환이라 하고 $r \geq 0$이라 하며,
$f_1, \ldots, f_r \in R$을 원소들의 열이라 하자.
$M$을 $R$-가군이라 하자. 열 $f_1, \ldots, f_r$을 다음과 같이 부른다.
\begin{enumerate}
\item 모든 $i \not = 0$에 대하여
$H_i(K_\bullet(f_1, \ldots, f_r) \otimes_R M) = 0$이면
{\it $M$-Koszul-정칙}이라 한다.
\item $H_1(K_\bullet(f_1, \ldots, f_r) \otimes_R M) = 0$이면
{\it $M$-$H_1$-정칙}이라 한다.
\item 모든 $i \not = 0$에 대하여
$H_i(K_\bullet(f_1, \ldots, f_r)) = 0$이면 {\it Koszul-정칙}이라 한다.
\item $H_1(K_\bullet(f_1, \ldots, f_r)) = 0$이면
{\it $H_1$-정칙}이라 한다.
\end{enumerate}
\end{definition}

\noindent
보조정리 \ref{more-algebra-lemma-regular-koszul-regular},
\ref{more-algebra-lemma-koszul-regular-H1-regular},
\ref{more-algebra-lemma-H1-regular-quasi-regular}에서 환 $R$의 원소
$f_1, \ldots, f_r$에 대하여 다음 함의들이 성립함을 볼 것이다.
\begin{align*}
f_1, \ldots, f_r\text{ is a regular sequence}
& \Rightarrow f_1, \ldots, f_r\text{ is a Koszul-regular sequence} \\
& \Rightarrow f_1, \ldots, f_r\text{ is an }H_1\text{-regular sequence} \\
& \Rightarrow f_1, \ldots, f_r\text{ is a quasi-regular sequence.}
\end{align*}
일반적으로 이 함의들 중 어느 것도 역으로 성립하지 않지만,
$R$이 Noether 국소환이고 $f_1, \ldots, f_r \in \mathfrak m_R$이면
네 조건은 모두 동치이다
(보조정리 \ref{more-algebra-lemma-noetherian-finite-all-equivalent}).
$f = f_1 \in R$이 길이 $1$인 열이고 $f$가 $R$의 단위원이 아니면,
다음 조건들은 모두 동치임이 명백하다.
\begin{enumerate}
\item $f$는 길이 하나인 정칙열이다.
\item $f$는 길이 하나인 Koszul-정칙열이다.
\item $f$는 길이 하나인 $H_1$-정칙열이다.
\end{enumerate}
이 조건들이 $f$가 길이 하나인 준정칙열임을 함의하는 것도 명백하다.
그러나 정칙열이 아닌 길이 $1$의 준정칙열도 존재한다. 실제로
$$
R = k[x, y_0, y_1, \ldots]/(xy_0, xy_1 - y_0, xy_2 - y_1, \ldots)
$$
로 놓고 $f$를 $R$ 안에서 $x$의 상이라 하자.
그러면 $f$는 영인자이지만
$\bigoplus_{n \geq 0} (f^n)/(f^{n + 1}) \cong k[x]$는 다항식환이다.

\begin{lemma}
\label{more-algebra-lemma-regular-koszul-regular}
$R$을 환, $M$을 $R$-가군이라 하고 $f_1, \ldots, f_r \in R$이라 하자.
$i = 1, \ldots, r$에 대하여 $f_i$를 곱하는 사상이
$M/(f_1, \ldots, f_{i - 1})M$ 위에서 단사라고 가정하자.
그러면 $f_1, \ldots, f_r$은 $M$-Koszul 정칙이다.
특히 $M$-정칙열은 $M$-Koszul-정칙이고,
모든 정칙열은 Koszul-정칙이다.
\end{lemma}

\begin{proof}
$R$, $M$, $f_1, \ldots, f_r$을 보조정리의 첫 문장에서와 같이 택하자.
$r = 1$이면 $f_1$이 $M$-Koszul-정칙임은 즉시 따른다.
$r > 1$이라고 가정하자. $f_1$이 $M$ 위의 영인자가 아니므로,
복합체들의 다음 짧은 완전열을 얻는다.
$$
0 \to K_\bullet(f_2, \ldots, f_r) \otimes M
\xrightarrow{f_1}
K_\bullet(f_2, \ldots, f_r) \otimes M \to
K_\bullet(\overline{f}_2, \ldots, \overline{f}_r) \otimes M/f_1M \to 0
$$
여기서 $\overline{f}_i$는 $R/(f_1)$ 안에서 $f_i$의 상이다.
보조정리 \ref{more-algebra-lemma-cone-koszul}에 의해
복합체 $K_\bullet(f_1, \ldots, f_r)$은
$K_\bullet(f_2, \ldots, f_r)$ 위에서 $f_1$을 곱하는 사상의 원뿔과 동형이다.
따라서 $K_\bullet(R, f_1, \ldots, f_r) \otimes M$은
첫 번째 사상의 원뿔과 동형이다.
그러므로 $K_\bullet(\overline{f}_2, \ldots, \overline{f}_r) \otimes M/f_1M$은
$K_\bullet(f_1, \ldots, f_r) \otimes M$과 준동형이다.
$R/(f_1)$, $M/f_1M$, $\overline{f}_2, \ldots, \overline{f}_r$이
보조정리의 조건들을 만족하므로 귀납법에 의해 이 복합체는
% P02-E0840: the English source misspells positive; Korean supplies the standard word.
양의 차수에서 비순환적이다. 이로써 첫 번째 명제가 증명되며,
두 번째 명제는 첫 번째 명제에서 즉시 따른다.
\end{proof}

\begin{lemma}
\label{more-algebra-lemma-koszul-regular-H1-regular}
$M$-Koszul-정칙열은 $M$-$H_1$-정칙이다.
Koszul-정칙열은 $H_1$-정칙이다.
\end{lemma}

\begin{proof}
정의에서 즉시 따른다.
\end{proof}

\begin{lemma}
\label{more-algebra-lemma-mult-koszul-regular}
$f_1, \ldots, f_{r - 1} \in R$을 열이라 하고 $f, g \in R$이라 하자.
$M$을 $R$-가군이라 하자.
\begin{enumerate}
\item $f_1, \ldots, f_{r - 1}, f$와 $f_1, \ldots, f_{r - 1}, g$가
$M$-$H_1$-정칙이면 $f_1, \ldots, f_{r - 1}, fg$도
$M$-$H_1$-정칙이다.
\item $f_1, \ldots, f_{r - 1}, f$와 $f_1, \ldots, f_{r - 1}, g$가
$M$-Koszul-정칙이면 $f_1, \ldots, f_{r - 1}, fg$도
$M$-Koszul-정칙이다.
\end{enumerate}
\end{lemma}

\begin{proof}
보조정리 \ref{more-algebra-lemma-koszul-mult}에 의해 모든 $i$에 대하여 완전열
$$
H_i(K_\bullet(f_1, \ldots, f_{r - 1}, f) \otimes M) \to
H_i(K_\bullet(f_1, \ldots, f_{r - 1}, fg) \otimes M) \to
H_i(K_\bullet(f_1, \ldots, f_{r - 1}, g) \otimes M)
$$
을 얻는다.
\end{proof}

\begin{lemma}
\label{more-algebra-lemma-koszul-regular-flat-base-change}
$\varphi : R \to S$를 평탄 환 준동형이라 하자.
$f_1, \ldots, f_r \in R$이라 하자. $M$을 $R$-가군이라 하고
$N = M \otimes_R S$로 놓자.
\begin{enumerate}
\item $R$ 안의 $f_1, \ldots, f_r$이 $M$-$H_1$-정칙열이면
$\varphi(f_1), \ldots, \varphi(f_r)$은 $S$ 안의 $N$-$H_1$-정칙열이다.
\item $R$ 안의 $f_1, \ldots, f_r$이 $M$-Koszul-정칙열이면
$\varphi(f_1), \ldots, \varphi(f_r)$은 $S$ 안의 $N$-Koszul-정칙열이다.
\end{enumerate}
\end{lemma}

\begin{proof}
다음 등식들이 성립하므로 참이다.
$K_\bullet(f_1, \ldots, f_r) \otimes_R S =
K_\bullet(\varphi(f_1), \ldots, \varphi(f_r))$
이고, 따라서
$(K_\bullet(f_1, \ldots, f_r) \otimes_R M) \otimes_R S =
K_\bullet(\varphi(f_1), \ldots, \varphi(f_r)) \otimes_S N$이다.
\end{proof}

\begin{lemma}
\label{more-algebra-lemma-H1-regular-quasi-regular}
$M$-$H_1$-정칙열은 $M$-준정칙이다.
\end{lemma}

\begin{proof}
$R$을 환이라 하고 $M$을 $R$-가군이라 하자.
$f_1, \ldots, f_r$을 $M$-$H_1$-정칙열이라 하자.
$J = (f_1, \ldots, f_r)$로 나타내자.
가정은 다음 완전열이 존재한다는 뜻이다.
$$
\wedge^2(R^r) \otimes M \to R^{\oplus r} \otimes M \to JM \to 0
$$
여기서 첫 번째 화살표는
$e_i \wedge e_j \otimes m \mapsto (f_ie_j - f_je_i) \otimes m$으로 주어진다.
이 열을 $R/J$와 텐서곱하면 다음을 얻는다.
$$
JM/J^2M = (R/J)^{\oplus r} \otimes_R M = (M/JM)^{\oplus r}
$$
% P02-E0841: retain the frozen finite-free assertion for arbitrary M.
이는 유한 자유 가군이다. 증명을 끝내기 위해 모든 $n \geq 2$에 대하여
다음을 증명해야 한다. 만일
$$
\xi = \sum\nolimits_{|I| = n, I = (i_1, \ldots, i_r)}
m_I f_1^{i_1} \ldots f_r^{i_r} \in J^{n + 1}M
$$
이면 모든 $I$에 대하여 $m_I \in JM$이다.
다음 단락에서는 $I = (0, \ldots, 0, n)$에 대하여 $m_I \in JM$임을
증명하고, 마지막 단락에서는 이 특수한 경우로부터 일반적인 경우를 유도한다.

\medskip\noindent
$I = (0, \ldots, 0, n)$이라 하자. $\xi$를 위와 같이 택하자.
$\xi = m_1 f_1 + \ldots + m_{r - 1}f_{r - 1} + m_I f_r^n$으로 쓸 수 있다.
$\xi \in J^{n + 1}M$이라고 가정했으므로 다음과 같이 쓸 수도 있다.
$\xi = \sum_{1 \leq i \leq j \leq r - 1} m_{ij}f_if_j +
\sum_{1 \leq i \leq r - 1}m'_i f_if_r^n + m'' f_r^{n + 1}$.
그러면 다음을 얻는다.
$$
\begin{matrix}
(m_1 - m_{11}f_1 - m'_1f_r^n)f_1 + \\
(m_2 - m_{12}f_1 - m_{22}f_2 - m'_2f_r^n)f_2 + \\
\ldots + \\
(m_{r - 1} - m_{1 r - 1}f_1 - \ldots - m_{r - 1 r - 1}f_{r - 1}
- m'_{r - 1}f_r^n)f_{r - 1} + \\
(m_I - m'' f_r)f_r^n = 0
\end{matrix}
$$
보조정리 \ref{more-algebra-lemma-mult-koszul-regular}에 의해
$f_1, \ldots, f_{r - 1}, f_r^n$이 $M$-$H_1$-정칙이므로,
$m_I - m'' f_r$은 부분가군
$f_1M + \ldots + f_{r - 1}M + f_r^nM$에 속한다.
따라서 $m_I \in f_1M + \ldots + f_rM$이다.

\medskip\noindent
$S = R[x_1, x_2, \ldots, x_r, 1/x_r]$로 놓자.
환 준동형 $R \to S$는 충실하게 평탄이므로,
보조정리 \ref{more-algebra-lemma-koszul-regular-flat-base-change}에 의해
$f_1, \ldots, f_r$은 $S$ 안에서 $M$-$H_1$-정칙열이다.
보조정리 \ref{more-algebra-lemma-change-basis}에 의해
$$
g_1 = f_1 - \frac{x_1}{x_r} f_r,
\ \ldots,
\ g_{r - 1} = f_{r - 1} - \frac{x_{r - 1}}{x_r} f_r,
\ g_r = \frac{1}{x_r}f_r
$$
이 $S$ 안에서 $M$-$H_1$-정칙열임을 알 수 있다.
마지막으로 원소 $\xi$는 다음과 같이 다시 쓸 수 있음을 주목하자.
$$
\xi = \sum\nolimits_{|I| = n, I = (i_1, \ldots, i_r)}
m_I (g_1 + x_1 g_r)^{i_1} \ldots (g_{r - 1} + x_{r - 1} g_r)^{i_{r - 1}}
(x_rg_r)^{i_r}
$$
이 식에서 $g_r^n$의 계수는
$$
\sum m_I x_1^{i_1} \ldots x_r^{i_r}
$$
이다. 앞 단락에서 논의한 경우에 의해 이 합은
$J(M \otimes_R S)$에 속한다.
단항식들 $x_1^{i_1} \ldots x_r^{i_r}$는 $R$ 위에서 $S$의
$R$-기저의 일부를 이루므로,
% P02-E0842: retain the frozen ill-typed conclusion m_I in J rather than JM.
원하는 대로 모든 $I$에 대하여 $m_I \in J$라고 결론짓는다.
\end{proof}

\noindent
Noether 국소환 위의 영이 아닌 유한 가군에 대해서는
지금까지 도입한 모든 종류의 정칙열이 서로 동치이다.

\begin{lemma}
\label{more-algebra-lemma-noetherian-finite-all-equivalent}
$(R, \mathfrak m)$을 Noether 국소환이라 하자.
$M$을 영이 아닌 유한 $R$-가군이라 하자.
$f_1, \ldots, f_r \in \mathfrak m$이라 하자.
다음 조건들은 동치이다.
\begin{enumerate}
\item $f_1, \ldots, f_r$은 $M$-정칙열이다.
\item $f_1, \ldots, f_r$은 $M$-Koszul-정칙열이다.
\item $f_1, \ldots, f_r$은 $M$-$H_1$-정칙열이다.
\item $f_1, \ldots, f_r$은 $M$-준정칙열이다.
\end{enumerate}
특히 열 $f_1, \ldots, f_r$이 $R$ 안의 정칙열이라는 것과
Koszul 정칙열이라는 것, $H_1$-정칙열이라는 것,
준정칙열이라는 것은 서로 동치이다.
\end{lemma}

\begin{proof}
함의 (1) $\Rightarrow$ (2)는 보조정리 \ref{more-algebra-lemma-regular-koszul-regular}이다.
함의 (2) $\Rightarrow$ (3)은 보조정리
\ref{more-algebra-lemma-koszul-regular-H1-regular}이다.
함의 (3) $\Rightarrow$ (4)는 보조정리
\ref{more-algebra-lemma-H1-regular-quasi-regular}이다.
함의 (4) $\Rightarrow$ (1)은 Algebra, Lemma
\ref{algebra-lemma-quasi-regular-regular}이다.
\end{proof}

\begin{lemma}
\label{more-algebra-lemma-H1-regular-in-quotient}
$A$를 환이라 하고 $I \subset A$를 아이디얼이라 하자.
$g_1, \ldots, g_m$을 $A$ 안의 열이라 하고,
$A/I$에서의 상이 $H_1$-정칙이라고 가정하자. 그러면
$I \cap (g_1, \ldots, g_m) = I(g_1, \ldots, g_m)$이다.
\end{lemma}

\begin{proof}
복합체들의 완전열
$$
0 \to I \otimes_A K_\bullet(A, g_1, \ldots, g_m)
\to K_\bullet(A, g_1, \ldots, g_m) \to
K_\bullet(A/I, g_1, \ldots, g_m) \to 0
$$
을 생각하자. 가정에 의해 오른쪽 복합체에 대해 $H_1 = 0$이므로
$$
\Coker(I^{\oplus m} \to I)
\longrightarrow
\Coker(A^{\oplus m} \to A)
$$
가 단사임을 알 수 있다. 이는 보조정리의 주장과 동치이다.
\end{proof}

\begin{lemma}
\label{more-algebra-lemma-conormal-sequence-H1-regular}
$A$를 환이라 하고 $I \subset J \subset A$를 아이디얼들이라 하자.
$J/I \subset A/I$가 $H_1$-정칙열로 생성된다고 가정하자.
그러면 $I \cap J^2 = IJ$이다.
\end{lemma}

\begin{proof}
이를 증명하기 위해 $A/I$에서의 상들이 $J/I$를 생성하는
$H_1$-정칙열을 이루는 $g_1, \ldots, g_m \in J$를 택하자.
특히 $J = I + (g_1, \ldots, g_m)$이다.
$x \in I \cap J^2$라고 하자.
% P02-E0843: the English clause lacks a subject and auxiliary; Korean supplies them.
$x \in J^2$이므로 다음과 같이 쓸 수 있다.
$$
x =
\sum a_{ij} g_ig_j +
\sum a_j g_j +
a
$$
여기서 $a_{ij} \in A$, $a_j \in I$, $a \in I^2$이다.
그러면 $\sum a_{ij}g_ig_j \in I \cap (g_1, \ldots, g_m)$이고,
보조정리 \ref{more-algebra-lemma-H1-regular-in-quotient}에 의해
$\sum a_{ij}g_ig_j \in I(g_1, \ldots, g_m)$임을 알 수 있다.
따라서 원하는 대로 $x \in IJ$이다.
\end{proof}

\begin{lemma}
\label{more-algebra-lemma-join-quasi-regular-H1-regular}
$A$를 환이라 하자. $I$를 $A$ 안의 준정칙열
$f_1, \ldots, f_n$으로 생성되는 아이디얼이라 하자.
$g_1, \ldots, g_m \in A$를 원소들이라 하고,
그 상 $\overline{g}_1, \ldots, \overline{g}_m$이
$A/I$ 안에서 $H_1$-정칙열을 이룬다고 하자.
그러면 $f_1, \ldots, f_n, g_1, \ldots, g_m$은 $A$ 안의 준정칙열이다.
\end{lemma}

\begin{proof}
모든 $d$에 대하여 $g_1, \ldots, g_m$이
$A/I^d$ 안에서 $H_1$-정칙열을 이룬다고 주장한다.
귀납적으로 이것이 $A/I^{d - 1}$에서 성립한다고 가정하자.
복합체들의 짧은 완전열
$$
0 \to K_\bullet(A, g_\bullet) \otimes_A I^{d - 1}/I^d
\to K_\bullet(A/I^d, g_\bullet) \to
K_\bullet(A/I^{d - 1}, g_\bullet) \to 0
$$
이 존재한다.
$f_1, \ldots, f_n$이 준정칙이므로 첫 번째 복합체는
$K_\bullet(A/I, g_1, \ldots, g_m)$의 복사본들의 직접합이고,
따라서 차수 $1$에서 비순환적이다.
귀납 가정에 의해 마지막 복합체도 차수 $1$에서 비순환적이다.
따라서 가운데 복합체도 그렇다.
특히 모든 $d \geq 1$에 대하여 $g_1, \ldots, g_m$은
$A/I^d$ 안에서 준정칙열을 이룬다.
보조정리 \ref{more-algebra-lemma-H1-regular-quasi-regular}를 보라.
이제 $f_1, \ldots, f_n, g_1, \ldots, g_m$이
$A$ 안에서 준정칙열임을 증명할 준비가 되었다.
$J = (f_1, \ldots, f_n, g_1, \ldots, g_m)$로 놓고,
다항 지표 표기를 사용하여 어떤 $a_{N, M} \in A$에 대하여
$$
\sum\nolimits_{|N| + |M| = d} a_{N, M} f^N g^M \in J^{d + 1}
$$
이라고 가정하자. 모든 $N, M$에 대하여 $a_{N, M} \in J$임을 보여야 한다.
$e \in \{0, 1, \ldots, d\}$라 하자. 그러면
$$
\sum\nolimits_{|N| = d - e, \ |M| = e} a_{N, M} f^N g^M \in
(g_1, \ldots, g_m)^{e + 1} + I^{d - e + 1}
$$
이다. $g_1, \ldots, g_m$이 $A/I^{d - e + 1}$ 안에서 준정칙열이므로,
$$
\sum\nolimits_{|N| = d - e} a_{N, M} f^N \in
(g_1, \ldots, g_m) + I^{d - e + 1}
$$
임을 얻는다. 이는 $|M| = e$인 각 $M$에 대하여 성립한다.
환 $A/I^{d - e + 1}$ 안의 $I^{d - e}/I^{d - e + 1}$에
보조정리 \ref{more-algebra-lemma-H1-regular-in-quotient}을 적용하면,
이는 $\sum_{|N| = d - e} a_{N, M} f^N
\in I^{d - e}(g_1, \ldots, g_m)$을 함의한다.
$f_1, \ldots, f_n$이 $A$ 안에서 준정칙이므로
$|N| = d - e$, $|M| = e$인 각 $N, M$에 대하여
$a_{N, M} \in J$임을 함의한다. 이로써 보조정리가 증명된다.
\end{proof}

\begin{lemma}
\label{more-algebra-lemma-join-H1-regular-sequences}
$A$를 환이라 하자. $I$를 $A$ 안의 $H_1$-정칙열
$f_1, \ldots, f_n$으로 생성되는 아이디얼이라 하자.
$g_1, \ldots, g_m \in A$를 원소들이라 하고,
그 상 $\overline{g}_1, \ldots, \overline{g}_m$이
$A/I$ 안에서 $H_1$-정칙열을 이룬다고 하자.
그러면 $f_1, \ldots, f_n, g_1, \ldots, g_m$은 $A$ 안의 $H_1$-정칙열이다.
\end{lemma}

\begin{proof}
$H_1(A, f_1, \ldots, f_n, g_1, \ldots, g_m) = 0$임을 보여야 한다.
이를 위해 다음 가환 도표를 생각하자.
% P02-E0844: retain the frozen unparenthesized A/I direct-sum notation.
$$
\xymatrix{
\wedge^2(A^{\oplus n + m}) \ar[r] \ar[d] &
A^{\oplus n + m} \ar[r] \ar[d] &
A \ar[r] \ar[d] & 0 \\
\wedge^2(A/I^{\oplus m}) \ar[r] &
A/I^{\oplus m} \ar[r] &
A/I \ar[r] & 0
}
$$
$A$에서 영으로 가는 원소
$(a_1, \ldots, a_{n + m}) \in A^{\oplus n + m}$을 생각하자.
$\overline{g}_1, \ldots, \overline{g}_m$이 $A/I$ 안에서
$H_1$-정칙열을 이루므로,
$(\overline{a}_{n + 1}, \ldots, \overline{a}_{n + m})$은
$\wedge^2(A/I^{\oplus m})$의 어떤 원소 $\overline{\alpha}$의 상이다.
$\overline{\alpha}$를 $\alpha \in \wedge^2(A^{\oplus n + m})$로 올리고,
$A^{\oplus n + m}$에서 그 상을
우리의 원소 $(a_1, \ldots, a_{n + m})$에서 뺄 수 있다.
따라서 $a_{n + 1}, \ldots, a_{n + m} \in I$라고 가정해도 된다.
$I = (f_1, \ldots, f_n)$이므로,
우리의 원소 $(a_1, \ldots, a_{n + m})$을 위쪽 왼쪽 수평 화살표의 상에 속하는
다음 원소들의 선형결합으로 바꿀 수 있다.
% P02-E0845: retain the frozen same-sign vector used as a trivial relation.
$$
(0, \ldots, g_j, 0, \ldots, 0, f_i, 0, \ldots, 0)
$$
그렇게 하여 $a_{n + 1}, \ldots, a_{n + m}$이 영인 경우로 환원할 수 있다.
이 경우 $(a_1, \ldots, a_n, 0, \ldots, 0)$은
$H_1(A, f_1, \ldots, f_n)$의 원소를 정의하며,
이는 가정에 의해 영이다.
\end{proof}

\begin{lemma}
\label{more-algebra-lemma-truncate-H1-regular}
$A$를 환이라 하자.
$f_1, \ldots, f_n, g_1, \ldots, g_m \in A$가 $H_1$-정칙열이라고 하자.
그러면 $A/(f_1, \ldots, f_n)$에서의 상
$\overline{g}_1, \ldots, \overline{g}_m$은 $H_1$-정칙열을 이룬다.
\end{lemma}

\begin{proof}
$I = (f_1, \ldots, f_n)$로 놓자.
% P02-E0846: retain the malformed frozen range subscripts in all three sums.
$A/I$에서의 임의의 관계
$\sum_{j = 1, \ldots, m} \overline{a}_j \overline{g}_j$가
자명한 관계들의 선형결합임을 보여야 한다.
$I = (f_1, \ldots, f_n)$이므로 이 관계를 $A$ 안의 관계
$$
\sum\nolimits_{j = 1, \ldots, m} a_j g_j +
\sum\nolimits_{i = 1, \ldots, n} b_if_i = 0
$$
로 올릴 수 있다. 가정에 의해 $A$ 안의 이 관계는
자명한 관계들의 선형결합이다. $A/I$에서 상을 취하면 원하는 결과를 얻는다.
\end{proof}

\begin{lemma}
\label{more-algebra-lemma-join-koszul-regular-sequences}
$A$를 환이라 하자. $I$를 $A$ 안의 Koszul-정칙열
$f_1, \ldots, f_n$으로 생성되는 아이디얼이라 하자.
$g_1, \ldots, g_m \in A$를 원소들이라 하고,
그 상 $\overline{g}_1, \ldots, \overline{g}_m$이
$A/I$ 안에서 Koszul-정칙열을 이룬다고 하자.
그러면 $f_1, \ldots, f_n, g_1, \ldots, g_m$은 $A$ 안의 Koszul-정칙열이다.
\end{lemma}

\begin{proof}
가정에 의해 $K_\bullet(A, f_1, \ldots, f_n)$은 $A/I$의
유한 자유 분해이고,
$K_\bullet(A/I, \overline{g}_1, \ldots, \overline{g}_m)$은
$A/I$ 위에서 $A/(f_i, g_j)$의 유한 자유 분해이다. 그러면
\begin{align*}
K_\bullet(A, f_1, \ldots, f_n, g_1, \ldots, g_m)
& = \text{Tot}(K_\bullet(A, f_1, \ldots, f_n) \otimes_A
K_\bullet(A, g_1, \ldots, g_m)) \\
& \cong A/I \otimes_A K_\bullet(A, g_1, \ldots, g_m) \\
& = K_\bullet(A/I, \overline{g}_1, \ldots, \overline{g}_m) \\
& \cong A/(f_i, g_j)
\end{align*}
% P02-E0847: the English justification fragments lack finite verbs; Korean supplies them.
첫 번째 등식은 보조정리 \ref{more-algebra-lemma-join-sequences-koszul-complex}에서 따른다.
첫 번째 준동형 $\cong$은 Homology, Lemma
\ref{homology-lemma-double-complex-gives-resolution}의 쌍대에서 따른다.
이중 복합체
$K_\bullet(A, f_1, \ldots, f_n) \otimes_A K_\bullet(A, g_1, \ldots, g_m)$의
$q$번째 행은 $A/I \otimes_A K_q(A, g_1, \ldots, g_m)$의 분해이기 때문이다.
두 번째 등식은 명백하고, 마지막 준동형은 가정에서 따른다.
따라서 결과를 얻는다.
\end{proof}

\noindent
다음 보조정리의 결론을 얻으려면
$f_1, \ldots, f_n$과 $f_1, \ldots, f_n, g_1, \ldots, g_m$이
모두 Koszul-정칙이라고 가정해야 한다.
$f_1, \ldots, f_n$이 Koszul-정칙이라는 가정을 제거할 수 없음을 보이는 반례는
Examples, Lemma \ref{examples-lemma-strange-regular-sequence}이다.

\begin{lemma}
\label{more-algebra-lemma-truncate-koszul-regular}
$A$를 환이라 하고
$f_1, \ldots, f_n, g_1, \ldots, g_m \in A$라 하자.
$f_1, \ldots, f_n$과 $f_1, \ldots, f_n, g_1, \ldots, g_m$이
모두 $A$ 안의 Koszul-정칙열이면,
$A/(f_1, \ldots, f_n)$ 안의
$\overline{g}_1, \ldots, \overline{g}_m$은 Koszul-정칙열을 이룬다.
\end{lemma}

\begin{proof}
$I = (f_1, \ldots, f_n)$로 놓자.
가정에 의해 $K_\bullet(A, f_1, \ldots, f_n)$은 $A/I$의
유한 자유 분해이고,
$K_\bullet(A, f_1, \ldots, f_n, g_1, \ldots, g_m)$은
$A$ 위에서 $A/(f_i, g_j)$의 유한 자유 분해이다. 그러면
\begin{align*}
A/(f_i, g_j) & \cong K_\bullet(A, f_1, \ldots, f_n, g_1, \ldots, g_m) \\
& = \text{Tot}(K_\bullet(A, f_1, \ldots, f_n) \otimes_A
K_\bullet(A, g_1, \ldots, g_m)) \\
& \cong A/I \otimes_A K_\bullet(A, g_1, \ldots, g_m) \\
& = K_\bullet(A/I, \overline{g}_1, \ldots, \overline{g}_m)
\end{align*}
% P02-E0847: the parallel English justification fragments lack finite verbs.
첫 번째 준동형 $\cong$은 가정에서 따른다.
첫 번째 등식은 보조정리 \ref{more-algebra-lemma-join-sequences-koszul-complex}에서 따른다.
두 번째 준동형은 Homology, Lemma
\ref{homology-lemma-double-complex-gives-resolution}의 쌍대에서 따른다.
이중 복합체
$K_\bullet(A, f_1, \ldots, f_n) \otimes_A K_\bullet(A, g_1, \ldots, g_m)$의
$q$번째 행은 $A/I \otimes_A K_q(A, g_1, \ldots, g_m)$의 분해이기 때문이다.
두 번째 등식은 명백하다. 따라서 결과를 얻는다.
\end{proof}

\begin{lemma}
\label{more-algebra-lemma-independence-of-generators}
$R$을 환이라 하고 $I$를 $f_1, \ldots, f_r \in R$로 생성되는 아이디얼이라 하자.
\begin{enumerate}
\item $I$가 길이 $r$의 준정칙열로 생성될 수 있으면,
$f_1, \ldots, f_r$은 준정칙열이다.
\item $I$가 길이 $r$의 $H_1$-정칙열로 생성될 수 있으면,
$f_1, \ldots, f_r$은 $H_1$-정칙열이다.
\item $I$가 길이 $r$의 Koszul-정칙열로 생성될 수 있으면,
$f_1, \ldots, f_r$은 Koszul-정칙열이다.
\end{enumerate}
\end{lemma}

\begin{proof}
$I$가 길이 $r$의 준정칙열로 생성될 수 있으면,
$I/I^2$는 $R/I$ 위의 계수 $r$인 자유 가군이다.
가정에 의해 $f_1, \ldots, f_r$이 생성하므로,
상들 $\overline{f}_i$가 $R/I$ 위에서 $I/I^2$의 기저를 이룸을 알 수 있다.
따라서 $f_1, \ldots, f_r$은 준정칙열이다.
$I/I^2$의 자유성 외에 이 명제가 뜻하는 바는 모든 사상
$\text{Sym}^n_{R/I}(I/I^2) \to I^n/I^{n + 1}$이 동형이라는 것이기 때문이다.

\medskip\noindent
계속하여 $I$가 준정칙열 $g_1, \ldots, g_r$로 생성된다고 가정하자.
$g_j = \sum a_{ij}f_i$로 쓰자.
앞 단락에 의해 $f_1, \ldots, f_r$이 준정칙이므로,
$\det(a_{ij})$는 $I$를 법으로 가역임을 알 수 있다.
행렬 $a_{ij}$는 Koszul 복합체의 사상을 유도하는 사상
$R^{\oplus r} \to R^{\oplus r}$를 준다.
$$
\alpha : K_\bullet(R, f_1, \ldots, f_r) \to K_\bullet(R, g_1, \ldots, g_r)
$$
보조정리 \ref{more-algebra-lemma-functorial}을 보라.
이 사상은 $\det(a_{ij})$를 가역화하면 동형이 된다.
% P02-E0848: retain the frozen cohomology terminology for homological Koszul H_i.
보조정리 \ref{more-algebra-lemma-homotopy-koszul}에 의해
$K_\bullet(R, f_1, \ldots, f_r)$과
$K_\bullet(R, g_1, \ldots, g_r)$의 코호몰로지 가군들이 모두
$I$에 의해 소멸되므로, $\alpha$는 준동형임을 알 수 있다.

\medskip\noindent
이제 $g_1, \ldots, g_r$이 $I$를 생성하는 $H_1$-정칙열이라고 가정하자.
그러면 보조정리 \ref{more-algebra-lemma-H1-regular-quasi-regular}에 의해
$g_1, \ldots, g_r$은 준정칙열이다.
앞 단락으로부터 $f_1, \ldots, f_r$이 $H_1$-정칙열임을 결론짓는다.
Koszul-정칙열에 대해서도 마찬가지이다.
\end{proof}

\begin{lemma}
\label{more-algebra-lemma-make-nonzero-divisor}
\begin{reference}
이는 \cite[Corollary]{McCoy}의 특수한 경우이다.
\end{reference}
$R$을 환이라 하자. 원소 $a_1, \ldots, a_n \in R$가
$R \to R^{\oplus n}$, $x \mapsto (xa_1, \ldots, xa_n)$을
단사 사상이 되게 한다고 하자.
그러면 다항식환 $R[t_1, \ldots, t_n]$의 원소
$\sum a_i t_i$는 영인자가 아니다.
\end{lemma}

\begin{proof}
$a_i$ 중 하나가 단위원이면, 이는 $R$ 위의 다항식환에서
$t_1 + a_2 t_2 + \ldots + a_n t_n$ 형태의 모든 원소가
영인자가 아니라는 명제일 뿐이다.

\medskip\noindent
경우 I: $R$이 Noether 환인 경우.
$R$의 동반 소아이디얼들을
$\mathfrak q_j$, $j = 1, \ldots, m$이라 하자.
각 사상
$$
\sum a_i t_i :
\text{Sym}^d(R^{\oplus n})
\longrightarrow
\text{Sym}^{d + 1}(R^{\oplus n})
$$
이 단사임을 보여야 한다.
$\text{Sym}^d(R^{\oplus n})$은 자유 $R$-가군이므로
그 동반 소아이디얼들은 $\mathfrak q_j$, $j = 1, \ldots, m$이다.
각 $j$에 대하여 $a_i \not \in \mathfrak q_j$인 $i = i(j)$가 존재한다.
실제로 가정에 의해 어떤 $i$에 대하여
$\mathfrak q_jx = 0$이지만 $a_i x \not = 0$인 $x \in R$이 존재한다.
따라서 $a_i$는 $R_{\mathfrak q_j}$에서 단위원이고,
사상은 $\mathfrak q_j$에서 국소화하면 단사이다.
그러므로 사상은 단사이다.
Algebra, Lemma \ref{algebra-lemma-zero-at-ass-zero}를 보라.

\medskip\noindent
경우 II: 일반적인 $R$.
$R$을 $a_1, \ldots, a_n \in R_\lambda$인 Noether 환들
$R_\lambda$의 합집합으로 쓸 수 있다.
각 $R_\lambda$에 대하여 결과가 성립하므로 $R$에 대해서도 성립한다.
\end{proof}

\begin{lemma}
\label{more-algebra-lemma-Koszul-regular-flat-locally-regular}
$R$을 환이라 하고 $f_1, \ldots, f_n$을
$(f_1, \ldots, f_n) \not = R$인 $R$ 안의 Koszul-정칙열이라 하자.
충실하게 평탄하고 매끄러운 환 준동형
$$
R \longrightarrow
S = R[\{t_{ij}\}_{i \leq j}, t_{11}^{-1}, t_{22}^{-1}, \ldots, t_{nn}^{-1}]
$$
을 생각하자. $1 \leq i \leq n$에 대하여
$$
g_i = \sum\nolimits_{i \leq j} t_{ij} f_j \in S.
$$
로 놓자. 그러면 $g_1, \ldots, g_n$은 $S$ 안의 정칙열이고
$(f_1, \ldots, f_n)S = (g_1, \ldots, g_n)$이다.
\end{lemma}

\begin{proof}
행렬
$$
\left(
\begin{matrix}
t_{11} & t_{12} & t_{13} & \ldots \\
0 & t_{22} & t_{23} & \ldots \\
0 & 0 & t_{33} & \ldots \\
\ldots & \ldots & \ldots & \ldots
\end{matrix}
\right)
$$
이 $S$에서 가역이므로 아이디얼들의 등식은 명백하다.
$f_1, \ldots, f_n$이 Koszul-정칙열이므로,
사상 $R \to R^{\oplus n}$, $x \mapsto (xf_1, \ldots, xf_n)$의
핵은 영임을 알 수 있다.
% P02-E0849: the English ordinal is malformed; Korean supplies the degree-n construction.
실제로 이는 $f_1, \ldots, f_n$에 관한 $R$의 $n$차 Koszul 호몰로지를 계산한다.
따라서 보조정리 \ref{more-algebra-lemma-make-nonzero-divisor}에 의해
$g_1 = f_1 t_{11} + \ldots + f_n t_{1n}$은
$S' = R[t_{11}, t_{12}, \ldots, t_{1n}, t_{11}^{-1}]$에서
영인자가 아님을 알 수 있다.
% P02-E0850: retain the frozen undefined Koszul-sequence term without the regular modifier.
보조정리 \ref{more-algebra-lemma-koszul-regular-flat-base-change}와
\ref{more-algebra-lemma-independence-of-generators}에 의해
$g_1, f_2, \ldots, f_n$은 $S'$ 안의 Koszul-열임을 알 수 있다.
보조정리 \ref{more-algebra-lemma-truncate-koszul-regular}에 의해
$\overline{f}_2, \ldots, \overline{f}_n$은
$S'/(g_1)$ 안의 Koszul-정칙열이라고 결론짓는다.
따라서 $n$에 대한 귀납법에 의해,
$S'/(g_1)[\{t_{ij}\}_{2 \leq i \leq j}, t_{22}^{-1}, \ldots, t_{nn}^{-1}]$에서
$g_2, \ldots, g_n$의 상들
$\overline{g}_2, \ldots, \overline{g}_n$이 정칙열을 이룸을 알 수 있다.
이는 다시 $g_1, \ldots, g_n$이 $S$ 안의 정칙열을 이룬다는 뜻이다.
\end{proof}

\section{Koszul 정칙열에 관한 추가 결과}
\label{more-algebra-section-more-koszul-regular}

\noindent
Section \ref{more-algebra-section-koszul-regular}의 논의를 계속한다.

\begin{lemma}
\label{more-algebra-lemma-vanishing-extended-alternating-koszul}
$R$을 환이라 하고 $f_1, \ldots, f_r \in R$을 Koszul-정칙열이라 하자.
Section \ref{more-algebra-section-alternating-cech}의 확장 교대 {\v C}ech 복합체
$R \to \bigoplus\nolimits_{i_0} R_{f_{i_0}} \to
\bigoplus\nolimits_{i_0 < i_1} R_{f_{i_0}f_{i_1}} \to
\ldots \to R_{f_1\ldots f_r}$는 차수 $r$에서만 코호몰로지를 갖는다.
\end{lemma}

\begin{proof}
보조정리 \ref{more-algebra-lemma-mult-koszul-regular}과 귀납법에 의해
모든 $n \geq 1$에 대하여 열 $f_1, \ldots, f_{r - 1}, f_r^n$은
Koszul 정칙이다. 보조정리 \ref{more-algebra-lemma-change-basis}에 의해
Koszul 정칙열의 임의의 순열은 Koszul 정칙열이다.
따라서 임의의(또는 모든) $f_i$를 그 $n$제곱으로 바꾸어도
Koszul 정칙열을 얻는다.
% P02-E0851: retain the frozen cohomology terminology for homologically graded K_bullet.
그러므로 $K_\bullet(R, f_1^n, \ldots, f_r^n)$은
호몰로지 차수 $0$에서만 영이 아닌 코호몰로지를 갖는다.
보조정리 \ref{more-algebra-lemma-extended-alternating-Cech-is-colimit-koszul}에 의해
이는 원하는 결과를 함의한다.
\end{proof}

\begin{lemma}
\label{more-algebra-lemma-blowup-regular-sequence}
$a, a_2, \ldots, a_r$을 환 $R$ 안의 $H_1$-정칙열이라 하자.
(예를 들어 Koszul 정칙열 또는 정칙열이면 된다.
보조정리 \ref{more-algebra-lemma-regular-koszul-regular}과
\ref{more-algebra-lemma-koszul-regular-H1-regular}을 보라.)
$I = (a, a_2, \ldots, a_r)$로 놓으면 블로업 대수
$R' = R[\frac{I}{a}]$는
$R'' = R[y_2, \ldots, y_r]/(a y_i - a_i)$와 동형이다.
\end{lemma}

\begin{proof}
Algebra, Lemma \ref{algebra-lemma-affine-blowup-quotient-description}에 의해
$R''$에 $a$-비틀림이 없음을 보이면 충분하다.

\medskip\noindent
% P02-E0852: retain all frozen ay_n endpoints instead of silently changing them to ay_r.
$a, ay_2 - a_2, \ldots, ay_n - a_r$이
$R[y_2, \ldots, y_r]$ 안의 $H_1$-정칙열이라고 주장한다.
실제로 Koszul 복합체를 정의하는 데 사용되는 사상
$$
(a, ay_2 - a_2, \ldots, ay_n - a_r) :
R[y_2, \ldots, y_r]^{\oplus r}
\longrightarrow
R[y_2, \ldots, y_r]
$$
은 $a, ay_2 - a_2, \ldots, ay_n - a_r$에 대한 Koszul 복합체를 정의하며,
다음 사상과 동형이다.
$$
(a, a_2, \ldots, a_r) :
R[y_2, \ldots, y_r]^{\oplus r} \longrightarrow
R[y_2, \ldots, y_r]
$$
% P02-E0853: the English phrase has an extraneous article; Korean supplies normal grammar.
이는 $a, a_2, \ldots, a_r$에 대한 Koszul 복합체를 정의하는 데 사용되며,
동형
$$
R[y_2, \ldots, y_r]^{\oplus r}
\longrightarrow
R[y_2, \ldots, y_r]^{\oplus r}
$$
을 통해 얻어진다. 이 동형은 $(b_1, \ldots, b_r)$을
% P02-E0854: retain the frozen formula missing an operator before the ellipsis.
$(b_1 - b_2y_2 \ldots - b_ry_r, -b_2, \ldots, - b_r)$로 보낸다.
보조정리 \ref{more-algebra-lemma-functorial}에 의해 이 Koszul 복합체들은 동형이다.
평탄 환 준동형 $R \to R[y_2, \ldots, y_r]$에
보조정리 \ref{more-algebra-lemma-koszul-regular-flat-base-change}을 적용하면
우리의 주장이 참임을 결론짓는다.
보조정리 \ref{more-algebra-lemma-cone-koszul}에 의해,
$a, ay_2 - a_2, \ldots, ay_n - a_r$에 대한 Koszul 복합체 $K$는
$ay_2 - a_2, \ldots, ay_n - a_r$에 대한 Koszul 복합체를 $L$이라 할 때
$a : L \to L$의 원뿔임을 알 수 있다.
주장에 의해 $H_1(K) = 0$이므로
$a : H_0(L) \to H_0(L)$이 단사라고 결론짓는다.
달리 말하면 원하는 대로
$R'' = R[y_2, \ldots, y_r]/(a y_i - a_i)$에는
영이 아닌 $a$-비틀림 원소가 없다.
\end{proof}

\begin{lemma}
\label{more-algebra-lemma-base-change-H1-regular}
$A \to B$를 환 준동형이라 하자.
$f_1, \ldots, f_r$을 $B$ 안의 열이라 하고
$B/(f_1, \ldots, f_r)$이 $A$-평탄이라고 하자.
$A \to A'$를 환 준동형이라 하자. 그러면 표준 사상
$$
H_1(K_\bullet(B, f_1, \ldots, f_r)) \otimes_A A'
\longrightarrow
H_1(K_\bullet(B', f'_1, \ldots, f'_r))
$$
은 전사이다. 여기서 $B' = B \otimes_A A'$이고
$f_i' \in B'$는 $f_i$의 상이다.
\end{lemma}

\begin{proof}
열
$$
\wedge^2(B^{\oplus r}) \to B^{\oplus r} \to B \to B/J \to 0
$$
은 $A$-가군들의 복합체이고, $B/J$는 $A$ 위 평탄이며,
% P02-E0855: retain both frozen cohomology-group claims for chain Koszul H_1.
$B^{\oplus r}$의 자리에서 코호몰로지 군
$H_1 = H_1(K_\bullet(B, f_1, \ldots, f_r))$을 갖는다.
이를 $A'$와 텐서곱하면 복합체
$$
\wedge^2((B')^{\oplus r}) \to (B')^{\oplus r} \to B' \to B'/J' \to 0
$$
을 얻으며, 이는 $B'$와 $B'/J'$에서 완전하다.
$(B')^{\oplus r}$에서의 코호몰로지 군
$H'_1 = H_1(K_\bullet(B', f'_1, \ldots, f'_r))$을 계산하기 위해,
위의 첫 번째 열을 완전열들
$0 \to J \to B \to B/J \to 0$,
$0 \to K \to B^{\oplus r} \to J \to 0$,
$\wedge^2(B^{\oplus r}) \to K \to H_1 \to 0$으로 분해한다.
$A$ 위에서 $A'$와 텐서곱하면 다음 완전열들을 얻는다.
$$
\begin{matrix}
0 \to J \otimes_A A' \to B \otimes_A A' \to (B/J) \otimes_A A' \to 0 \\
K \otimes_A A' \to B^{\oplus r} \otimes_A A' \to J \otimes_A A' \to 0 \\
\wedge^2(B^{\oplus r}) \otimes_A A' \to K \otimes_A A' \to H_1 \otimes_A A'
\to 0
\end{matrix}
$$
첫 번째 열은 $B/J$가 $A$ 위 평탄이므로 완전하다.
Algebra, Lemma \ref{algebra-lemma-flat-tor-zero}를 보라.
따라서 $J' = J \otimes_A A' \subset B'$이고
$K \otimes_A A' \to \Ker((B')^{\oplus r} \to B')$가 전사임을 결론짓는다.
그러므로
\begin{align*}
H_1 \otimes_A A'
& =
\Coker\left(\wedge^2(B^{\oplus r}) \otimes_A A' \to K \otimes_A A'\right) \\
& \to
\Coker\left(
\wedge^2((B')^{\oplus r})  \to \Ker((B')^{\oplus r} \to B')
\right) = H'_1
\end{align*}
도 전사이다.
\end{proof}

\begin{lemma}
\label{more-algebra-lemma-relative-regular-immersion-algebra}
$A \to B$와 $A \to A'$를 환 준동형이라 하고 $B' = B \otimes_A A'$로 놓자.
$f_1, \ldots, f_r \in B$라 하자.
% P02-E0856: the English source omits terminal punctuation; Korean supplies it.
$B/(f_1, \ldots, f_r)B$가 $A$ 위 평탄이라고 가정하자.
\begin{enumerate}
\item $f_1, \ldots, f_r$이 준정칙열이면, $B'$에서의 상도 준정칙열이다.
\item $f_1, \ldots, f_r$이 $H_1$-정칙열이면,
$B'$에서의 상도 $H_1$-정칙열이다.
\end{enumerate}
\end{lemma}

\begin{proof}
$f_1, \ldots, f_r$이 준정칙이라고 가정하자.
$J = (f_1, \ldots, f_r)$로 놓자.
가정에 의해 $J^n/J^{n + 1}$은 $B/J$의 복사본들의 직접합과 동형이고,
따라서 $A$ 위 평탄이다.
귀납법과 Algebra, Lemma \ref{algebra-lemma-flat-ses}에 의해
$B/J^n$이 $A$ 위 평탄이라고 결론짓는다.
Algebra, Lemma \ref{algebra-lemma-flat-tor-zero}에 의해
아이디얼 $(J')^n$은 $J^n \otimes_A A'$와 같다.
따라서
$(J')^n/(J')^{n + 1} = J^n/J^{n + 1} \otimes_A A'$이며,
이는 명백히 $f_1, \ldots, f_r$이 $B'$ 안의 준정칙열임을 함의한다.

\medskip\noindent
$f_1, \ldots, f_r$이 $H_1$-정칙이라고 가정하자.
보조정리 \ref{more-algebra-lemma-base-change-H1-regular}에 의해
Koszul 호몰로지 군 $H_1(K_\bullet(B, f_1, \ldots, f_r))$의 소멸은
$H_1(K_\bullet(B', f'_1, \ldots, f'_r))$의 소멸을 함의하며,
따라서 결과를 얻는다.
\end{proof}

\begin{lemma}
\label{more-algebra-lemma-cut-by-koszul-flat}
$A' \to B'$를 환 준동형이라 하고 $I \subset A'$를 아이디얼이라 하자.
$A = A'/I$와 $B = B'/IB'$로 놓자.
$f'_1, \ldots, f'_r \in B'$라 하자. 다음을 가정하자.
\begin{enumerate}
\item $A' \to B'$는 평탄이고 유한 표시이다.
\item $I$는 국소적으로 멱영이다.
\item 상들 $f_1, \ldots, f_r \in B$는 준정칙열을 이룬다.
\item $B/(f_1, \ldots, f_r)$은 $A$ 위 평탄이다.
\end{enumerate}
그러면 $B'/(f'_1, \ldots, f'_r)$은 $A'$ 위 평탄이다.
\end{lemma}

\begin{proof}
$C' = B'/(f'_1, \ldots, f'_r)$로 놓자.
$A' \to C'$가 평탄임을 보여야 한다.
$\mathfrak p' \subset A'$ 위에 놓이는 소아이디얼
$\mathfrak r' \subset C'$를 택하자.
$\mathfrak q' \subset B'$를 $\mathfrak r'$의 역상이라 하자.
Algebra, Lemma \ref{algebra-lemma-flat-localization}에 의해
$A'_{\mathfrak p'} \to C'_{\mathfrak q'}$가 평탄임을 보이면 충분하다.
Algebra, Lemma \ref{algebra-lemma-grothendieck-regular-sequence-general}에 의해
$f'_1, \ldots, f'_r$이 다음 환 안에서 정칙열로 가는지 보이면 충분하다.
$$
B'_{\mathfrak q'}/\mathfrak p'B'_{\mathfrak q'} =
B_\mathfrak q/\mathfrak p B_\mathfrak q =
(B \otimes_A \kappa(\mathfrak p))_\mathfrak q
$$
표기의 의미는 자명하다. 알고 있는 것은
$f_1, \ldots, f_r$이 $B$ 안의 준정칙열이고
$B/(f_1, \ldots, f_r)$이 $A$ 위 평탄이라는 사실이다.
보조정리 \ref{more-algebra-lemma-relative-regular-immersion-algebra}에 의해
$B \otimes_A \kappa(\mathfrak p)$에서
$f'_1, \ldots, f'_r$의 상들
$\overline{f}_1, \ldots, \overline{f}_r$은 준정칙열을 이룬다.
$(B \otimes_A \kappa(\mathfrak p))_\mathfrak q$는 Noether 국소환이므로
보조정리 \ref{more-algebra-lemma-noetherian-finite-all-equivalent}에 의해 결론을 얻는다.
\end{proof}

\begin{lemma}
\label{more-algebra-lemma-cut-by-koszul}
$A' \to B'$를 환 준동형이라 하고 $I \subset A'$를 아이디얼이라 하자.
$A = A'/I$와 $B = B'/IB'$로 놓자.
$f'_1, \ldots, f'_r \in B'$라 하자. 다음을 가정하자.
\begin{enumerate}
\item $A' \to B'$는 평탄이고 유한 표시이다(예를 들어 매끄럽다).
\item $I$는 국소적으로 멱영이다.
\item 상들 $f_1, \ldots, f_r \in B$는 준정칙열을 이룬다.
\item $B/(f_1, \ldots, f_r)$은 $A$ 위 매끄럽다.
\end{enumerate}
그러면 $B'/(f'_1, \ldots, f'_r)$은 $A'$ 위 매끄럽다.
\end{lemma}

\begin{proof}
$C' = B'/(f'_1, \ldots, f'_r)$와
$C = B/(f_1, \ldots, f_r)$로 놓자.
그러면 $A' \to C'$는 유한 표시이다.
보조정리 \ref{more-algebra-lemma-cut-by-koszul-flat}에 의해
$A' \to C'$가 평탄임을 알 수 있다.
$A' \to C'$의 섬유환들은 $A \to C$의 섬유환들과 같고,
따라서 가정 (4)에 의해 매끄럽다.
Algebra, Lemma \ref{algebra-lemma-flat-fibre-smooth}에 의해
$A' \to C'$는 매끄럽다.
\end{proof}

\section{정칙 아이디얼}
\label{more-algebra-section-ideals}

\noindent
Divisors, Section \ref{divisors-section-regular-ideal-sheaves}에서
정칙 아이디얼 층의 개념을 매우 일반적으로 논의할 것이다.
여기서는 아핀인 경우, 즉 환 안의 아이디얼에 대응하는 개념을 정의한다.

\begin{definition}
\label{more-algebra-definition-regular-ideal}
$R$을 환이라 하고 $I \subset R$을 아이디얼이라 하자.
\begin{enumerate}
\item 모든 $\mathfrak p \in V(I)$에 대하여
$g \in R$, $g \not \in \mathfrak p$와 정칙열
$f_1, \ldots, f_r \in R_g$가 존재하여
$I_g$가 $f_1, \ldots, f_r$로 생성되면,
$I$를 {\it 정칙 아이디얼}이라 한다.
\item 모든 $\mathfrak p \in V(I)$에 대하여
$g \in R$, $g \not \in \mathfrak p$와 Koszul-정칙열
$f_1, \ldots, f_r \in R_g$가 존재하여
$I_g$가 $f_1, \ldots, f_r$로 생성되면,
$I$를 {\it Koszul-정칙 아이디얼}이라 한다.
\item 모든 $\mathfrak p \in V(I)$에 대하여
$g \in R$, $g \not \in \mathfrak p$와 $H_1$-정칙열
$f_1, \ldots, f_r \in R_g$가 존재하여
$I_g$가 $f_1, \ldots, f_r$로 생성되면,
$I$를 {\it $H_1$-정칙 아이디얼}이라 한다.
\item 모든 $\mathfrak p \in V(I)$에 대하여
$g \in R$, $g \not \in \mathfrak p$와 준정칙열
$f_1, \ldots, f_r \in R_g$가 존재하여
$I_g$가 $f_1, \ldots, f_r$로 생성되면,
$I$를 {\it 준정칙 아이디얼}이라 한다.
\end{enumerate}
\end{definition}

\noindent
$I \subset R$이 주어지면 다음 함의들이 성립함은 명백하다.
\begin{align*}
I\text{ is a regular ideal}
& \Rightarrow
I\text{ is a Koszul-regular ideal} \\
& \Rightarrow
I\text{ is a }H_1\text{-regular ideal} \\
& \Rightarrow
I\text{ is a quasi-regular ideal}
\end{align*}
보조정리 \ref{more-algebra-lemma-regular-koszul-regular},
\ref{more-algebra-lemma-koszul-regular-H1-regular},
\ref{more-algebra-lemma-H1-regular-quasi-regular}을 보라.
그러한 아이디얼은 항상 유한 생성이다.

\begin{lemma}
\label{more-algebra-lemma-quasi-regular-ideal-finite}
준정칙 아이디얼은 유한 생성이다.
\end{lemma}

\begin{proof}
$I \subset R$을 준정칙 아이디얼이라 하자.
$V(I)$는 준콤팩트이므로 $g_1, \ldots, g_m \in R$가 존재하여
$V(I) \subset D(g_1) \cup \ldots \cup D(g_m)$이고,
$I_{g_j}$는 준정칙열
$g_{j1}, \ldots, g_{jr_j} \in R_{g_j}$로 생성된다.
% P02-E0857: retain the frozen numerator index swap g'_{ji}/g'_{ij}.
$g_{ji} = g'_{ji}/g_j^{e_{ij}}$로 쓰자. 여기서 $g'_{ij} \in I$이다.
$1 + x = \sum g_j h_j$로 쓰자. 여기서 $x \in I$이다.
$V(I) \subset D(g_1) \cup \ldots \cup D(g_m)$이므로 가능하다.
% P02-E0858: the English affine-cover identity lacks terminal punctuation.
$\Spec(R) = D(g_1) \cup \ldots \cup D(g_m) \bigcup D(x)$임을 주목하자.
그러면 $I$는 원소들 $g'_{ij}$와 $x$로 생성된다.
이들이 덮개의 각 조각 위에서 생성하기 때문이다.
Algebra, Lemma \ref{algebra-lemma-cover}를 보라.
\end{proof}

\begin{lemma}
\label{more-algebra-lemma-quasi-regular-ideal-finite-projective}
$I \subset R$을 환의 준정칙 아이디얼이라 하자.
그러면 $I/I^2$는 유한 사영 $R/I$-가군이다.
\end{lemma}

\begin{proof}
Algebra, Lemma \ref{algebra-lemma-finite-projective}와 정의들에서 따른다.
\end{proof}

\noindent
Koszul-정칙, $H_1$-정칙, 준정칙 아이디얼에 대한 평탄 내림을 증명한다.

\begin{lemma}
\label{more-algebra-lemma-flat-descent-regular-ideal}
$A \to B$를 충실하게 평탄인 환 준동형이라 하고
$I \subset A$를 아이디얼이라 하자.
$IB$가 $B$ 안의 Koszul-정칙
(각각 $H_1$-정칙, 준정칙) 아이디얼이면,
$I$는 $A$ 안의 Koszul-정칙
(각각 $H_1$-정칙, 준정칙) 아이디얼이다.
\end{lemma}

\begin{proof}
증명 전체에서 소아이디얼 $\mathfrak p \supset I$를 고정한다.
$IB$가 준정칙이라고 가정하자.
보조정리 \ref{more-algebra-lemma-quasi-regular-ideal-finite}에 의해
$IB$는 유한 가군이고, 따라서 Algebra, Lemma
\ref{algebra-lemma-descend-properties-modules}에 의해
$I$는 유한 $A$-가군이다.
$A \to B$가 평탄이므로 다음을 얻는다.
$$
I/I^2 \otimes_{A/I} B/IB = I/I^2 \otimes_A B = IB/(IB)^2.
$$
$IB$가 준정칙이므로 $B/IB$-가군 $IB/(IB)^2$는 유한 국소 자유이다.
따라서 $I/I^2$는 유한 사영이다.
Algebra, Proposition \ref{algebra-proposition-ffdescent-finite-projectivity}를 보라.
특히 $A$를 어떤 $f \in A$, $f \not \in \mathfrak p$에 대한 $A_f$로
바꾸면 $I/I^2$가 계수 $r$의 자유 가군이라고 가정해도 된다.
$I/I^2$의 기저를 주는 $f_1, \ldots, f_r \in I$를 택하자.
Nakayama 보조정리(Algebra, Lemma \ref{algebra-lemma-NAK})에 의해,
다시 위와 같은 $A \leadsto A_f$로 바꾸면
$I$가 $f_1, \ldots, f_r$로 생성된다고 가정해도 된다.

\medskip\noindent
“준정칙”인 경우의 증명.
위에서 $I/I^2$가 $r$개의 생성원 $f_1, \ldots, f_r$ 위에서 자유임을 보았다.
이 경우 증명을 끝내려면 각 $d \geq 2$에 대하여 사상
$\text{Sym}^d(I/I^2) \to I^d/I^{d + 1}$이 동형임을 보여야 한다.
충실하게 평탄인 기저변환
$\text{Sym}^d(IB/(IB)^2) \to (IB)^d/(IB)^{d + 1}$이
가정에 의해 $B$ 위에서 국소적으로 동형이므로 이는 명백하다.
세부사항은 생략한다.

\medskip\noindent
“$H_1$-정칙” 및 “Koszul-정칙”인 경우의 증명.
위에서 구성한 $I$의 생성원들의 열 $f_1, \ldots, f_r$을 생각하자.
보조정리 \ref{more-algebra-lemma-independence-of-generators}에 의해,
$IB$가 $H_1$-정칙열 또는 Koszul-정칙열로 생성되는 임의의
$g \in B$에 대하여 $f_1, \ldots, f_r$은 $B_g$에서
$H_1$-정칙열 또는 Koszul-정칙열로 간다.
따라서
$K_\bullet(A, f_1, \ldots, f_r) \otimes_A B_g$의
$H_1$ 또는 $H_i$, $i > 0$은 소멸한다.
$K_\bullet(B, f_1, \ldots, f_r) =
K_\bullet(A, f_1, \ldots, f_r) \otimes_A B$의 호몰로지는
$IB$에 의해 소멸되고
(보조정리 \ref{more-algebra-lemma-homotopy-koszul}을 보라),
$V(IB) \subset \bigcup_{g\text{ as above}} D(g)$이므로
$K_\bullet(A, f_1, \ldots, f_r) \otimes_A B$는
차수 $1$ 또는 모든 양의 차수에서 호몰로지가 소멸한다고 결론짓는다.
$A \to B$가 충실하게 평탄이라는 사실을 사용하면
$K_\bullet(A, f_1, \ldots, f_r)$에 대해서도 마찬가지임을 결론짓는다.
\end{proof}

\begin{lemma}
\label{more-algebra-lemma-conormal-sequence-H1-regular-ideal}
$A$를 환이라 하고 $I \subset J \subset A$를 아이디얼들이라 하자.
$J/I \subset A/I$가 $H_1$-정칙 아이디얼이라고 가정하자.
그러면 $I \cap J^2 = IJ$이다.
\end{lemma}

\begin{proof}
보조정리 \ref{more-algebra-lemma-conormal-sequence-H1-regular}를 국소화하면 즉시 따른다.
\end{proof}







\section{국소 완전교차 사상}
\label{more-algebra-section-lci}

\noindent
위의 내용을 사용하면 (유한) 다항식 대수에 의한 표시를 통해
환들 사이의 국소 완전교차 사상을 정의할 수 있다.

\begin{lemma}
\label{more-algebra-lemma-koszul-independence-presentation}
$A \to B$를 유한형 환 사상이라 하자. 어떤 표시에 대해
$\alpha : A[x_1, \ldots, x_n] \to B$의 핵 $I$가 Koszul-정칙 아이디얼이면,
임의의 표시 $\beta : A[y_1, \ldots, y_m] \to B$의 핵
$J$도 Koszul-정칙 아이디얼이다.
\end{lemma}

\begin{proof}
$\alpha(f_j) = \beta(y_j)$를 만족하는 $f_j \in A[x_1, \ldots, x_n]$와
$\beta(g_i) = \alpha(x_i)$를 만족하는 $g_i \in A[y_1, \ldots, y_m]$를 택하자.
그러면 다음 가환 그림을 얻는다.
$$
\xymatrix{
A[x_1, \ldots, x_n, y_1, \ldots, y_m]
\ar[d]^{x_i \mapsto g_i} \ar[rr]_-{y_j \mapsto f_j} & &
A[x_1, \ldots, x_n] \ar[d] \\
A[y_1, \ldots, y_m] \ar[rr] & & B
}
$$
$A[x_i, y_j] \to B$의 핵 $K$는
$K = (I, y_j - f_j) = (J, x_i - g_i)$와 같음에 유의하자. 특히
보조정리 \ref{more-algebra-lemma-quasi-regular-ideal-finite}에 의해
$I$가 유한 생성이므로
$J = K/(x_i - g_i)$도 유한 생성임을 알 수 있다.

\medskip\noindent
소아이디얼 $\mathfrak q \subset B$를 택하자.
$I/I^2 \oplus B^{\oplus m} = J/J^2 \oplus B^{\oplus n}$이므로
(Algebra, Lemma \ref{algebra-lemma-conormal-module})
다음을 알 수 있다.
$$
\dim J/J^2 \otimes_B \kappa(\mathfrak q) + n =
\dim I/I^2 \otimes_B \kappa(\mathfrak q) + m.
$$
$I/I^2 \otimes \kappa(\mathfrak q) = I \otimes_{A[x_i]} \kappa(\mathfrak q)$의
기저로 가는 $p_1, \ldots, p_t \in I$를 택하자.
$J/J^2 \otimes \kappa(\mathfrak q) = J \otimes_{A[y_j]} \kappa(\mathfrak q)$의
기저로 가는 $q_1, \ldots, q_s \in J$를 택하자.
따라서 $s + n = t + m$이다. Nakayama 보조정리에 의해
$\kappa(\mathfrak q)$의 영이 아닌 원소로 가는
$h \in A[x_i]$와 $h' \in A[y_j]$가 존재하여,
$A[x_i, 1/h]$ 안에서 $I_h = (p_1, \ldots, p_t)$이고
$A[y_j, 1/h']$ 안에서 $J_{h'} = (q_1, \ldots, q_s)$이다.
$I$가 Koszul-정칙이므로 $I_h$가 Koszul 정칙열로 생성된다고
추가로 가정할 수 있다. 이 열의 길이는 반드시
$t = \dim I/I^2 \otimes_B \kappa(\mathfrak q)$이어야 하므로,
보조정리 \ref{more-algebra-lemma-independence-of-generators}에 의해
$p_1, \ldots, p_t$가 Koszul-정칙열임을 알 수 있다.
또한 $y_1 - f_1, \ldots, y_m - f_m$도 정칙열이므로
다음을 결론짓는다.
$$
y_1 - f_1, \ldots, y_m - f_m, p_1, \ldots, p_t
$$
이는 $A[x_i, y_j, 1/h]$ 안의 Koszul-정칙열이다
(보조정리 \ref{more-algebra-lemma-join-koszul-regular-sequences}를 보라).
이 열은 아이디얼 $K_h$를 생성한다. 따라서
아이디얼 $K_{hh'}$은 길이 $m + t = n + s$인
Koszul-정칙열로 생성된다. 한편 이는 같은 길이의 열
$$
x_1 - g_1, \ldots, x_n - g_n, q_1, \ldots, q_s
$$
로도 생성되므로, 보조정리 \ref{more-algebra-lemma-independence-of-generators}에 의해
이 열도 Koszul-정칙열이다. 마지막으로
보조정리 \ref{more-algebra-lemma-truncate-koszul-regular}에 의해
다음 안에서 $q_1, \ldots, q_s$의 상들이
$$
A[x_i, y_j, 1/hh']/(x_1 - g_1, \ldots, x_n - g_n)
\cong A[y_j, 1/h'']
$$
$J_{h''}$을 생성하는 Koszul-정칙열을 이룬다고 결론짓는다.
$h''$은 $hh'$의 상이므로 $\kappa(\mathfrak q)$에서 영으로 가지 않으며,
이것으로 증명이 끝난다.
\end{proof}

\noindent
이 보조정리로 다음 정의를 내릴 수 있다.

\begin{definition}
\label{more-algebra-definition-local-complete-intersection}
환 사상 $A \to B$가 유한형이고 어떤 (동치로 임의의) 표시
$B = A[x_1, \ldots, x_n]/I$에 대해 아이디얼 $I$가 Koszul-정칙이면,
이 사상을 {\it 국소 완전교차}라 한다.
\end{definition}

\noindent
이 개념은 국소적이다.

\begin{lemma}
\label{more-algebra-lemma-lci-local}
$R \to S$를 환 사상이라 하자. $g_1, \ldots, g_m \in S$가
단위 아이디얼을 생성한다고 하자. 각 $R \to S_{g_j}$가
국소 완전교차이면 $R \to S$도 국소 완전교차이다.
\end{lemma}

\begin{proof}
$S = R[x_1, \ldots, x_n]/I$를 하나의 표시라 하자.
$S$에서 $g_j$로 가는 $h_j \in R[x_1, \ldots, x_n]$를 택하자.
그러면 $R[x_1, \ldots, x_n, x_{n + 1}]/(I, x_{n + 1}h_j - 1)$은
$S_{g_j}$의 한 표시이다. 따라서
$I_j = (I, x_{n + 1}h_j - 1)$은
$R[x_1, \ldots, x_n, x_{n + 1}]$ 안의 Koszul-정칙 아이디얼이다.
소아이디얼 $I \subset \mathfrak q \subset R[x_1, \ldots, x_n]$을 택하자.
그러면 어떤 $j$에 대해 $h_j \not \in \mathfrak q$이고,
$\mathfrak q_j = (\mathfrak q, x_{n + 1}h_j - 1)$은
$\mathfrak q$ 위에 놓이는 $V(I_j)$의 소아이디얼이다.
$I/I^2 \otimes \kappa(\mathfrak q)$의 기저로 가는
$f_1, \ldots, f_r \in I$를 택하자. 그러면
$x_{n + 1}h_j - 1, f_1, \ldots, f_r$은 $I_j$의 원소들의 열로서
$I_j \otimes \kappa(\mathfrak q_j)$의 기저로 간다.
Algebra, Lemma \ref{algebra-lemma-principal-localization-NL}를 보라.
Nakayama 보조정리에 의해 $(I_j)_h$가
$x_{n + 1}h_j - 1, f_1, \ldots, f_r$로 생성되게 하는
$h \in R[x_1, \ldots, x_n, x_{n + 1}]$가 존재한다.
또한 $(I_j)_h$가 어떤 길이 $e$의 Koszul 정칙열로 생성된다고
가정할 수 있다. $I_j \otimes \kappa(\mathfrak q_j)$의 차원을 보면
$e = r + 1$임을 알 수 있다. 따라서
보조정리 \ref{more-algebra-lemma-independence-of-generators}에 의해
$x_{n + 1}h_j - 1, f_1, \ldots, f_r$은 어떤
$h \in R[x_1, \ldots, x_n, x_{n + 1}]$, $h \not \in \mathfrak q_j$에 대해
$(I_j)_h$를 생성하는 Koszul-정칙열이다. 이제
보조정리 \ref{more-algebra-lemma-truncate-koszul-regular}에 의해, 원하는 대로 어떤
$h' \in R[x_1, \ldots, x_n]$, $h' \not \in \mathfrak q$에 대해
$I_{h'}$이 Koszul-정칙열로 생성됨을 알 수 있다.
\end{proof}

\begin{lemma}
\label{more-algebra-lemma-relative-global-complete-intersection-koszul}
$R$을 환이라 하자. $R[x_1, \ldots, x_n]/(f_1, \ldots, f_c)$가
상대 전역 완전교차이면 $f_1, \ldots, f_c$는 Koszul 정칙열이다.
\end{lemma}

\begin{proof}
호몰로지 군 $H_i(K_\bullet(f_\bullet))$은 아이디얼
$(f_1, \ldots, f_c)$에 의해 소멸됨을 상기하자. 따라서
$(f_1, \ldots, f_c)$를 포함하는 모든 소아이디얼
$\mathfrak q \subset R[x_1, \ldots, x_n]$에 대해
$H_i(K_\bullet(f_\bullet))_\mathfrak q$가 영임을 보이면 충분하다.
이는 Algebra, Lemma
\ref{algebra-lemma-relative-global-complete-intersection-conormal}와
정칙열은 Koszul 정칙이라는 사실
(보조정리 \ref{more-algebra-lemma-regular-koszul-regular})에서 따른다.
\end{proof}

\begin{lemma}
\label{more-algebra-lemma-syntomic-lci}
$R \to S$를 환 사상이라 하자. 다음은 서로 동치이다.
\begin{enumerate}
\item $R \to S$는 신토믹이다
(Algebra, Definition \ref{algebra-definition-lci}).
\item $R \to S$는 평탄이고 국소 완전교차이다.
\end{enumerate}
\end{lemma}

\begin{proof}
(1)을 가정하자. 그러면 정의에 의해 $R \to S$는 평탄이다.
Algebra, Lemma \ref{algebra-lemma-syntomic}과
보조정리 \ref{more-algebra-lemma-lci-local}에 의해 상대 전역 완전교차가
국소 완전교차 준동형임을 보이면 충분하며, 이는
보조정리 \ref{more-algebra-lemma-relative-global-complete-intersection-koszul}이다.

\medskip\noindent
(2)를 가정하자. Koszul-정칙 아이디얼은 유한 생성이므로
국소 완전교차는 유한 표시이다. $R \to k$를 체로 가는 사상이라 하자.
$S' = S \otimes_R k$가 $k$ 위의 국소 완전교차임을 보이면 충분하다.
Algebra, Definition \ref{algebra-definition-lci-field}를 보라.
소아이디얼 $\mathfrak q' \subset S'$를 택하자.
$S = R[x_1, \ldots, x_n]/I$라고 쓰자.
그러면 $S' = k[x_1, \ldots, x_n]/I'$이고,
$I' \subset k[x_1, \ldots, x_n]$은 $I$의 상이다.
이에 대응하는 소아이디얼들을 각각
$\mathfrak p' \subset k[x_1, \ldots, x_n]$,
$\mathfrak q \subset S$,
$\mathfrak p \subset R[x_1, \ldots, x_n]$이라 하자.
% P02-E0859: frozen source line 8437 says “exists an g”; Korean uses the intended existential grammar.
정의 \ref{more-algebra-definition-regular-ideal}에 의해
$g \in R[x_1, \ldots, x_n]$, $g \not \in \mathfrak p$와
$I_g$를 생성하는 Koszul-정칙열을 이루는
$f_1, \ldots, f_r \in R[x_1, \ldots, x_n]_g$가 존재한다.
$S$, 따라서 $S_g$가 $R$ 위에서 평탄이므로,
$k[x_1, \ldots, x_n]_g$ 안의 상 $f'_1, \ldots, f'_r$은
$I'_g$를 생성하는 $H_1$-정칙열을 이룬다.
보조정리 \ref{more-algebra-lemma-relative-regular-immersion-algebra}를 보라.
따라서 보조정리 \ref{more-algebra-lemma-noetherian-finite-all-equivalent}에 의해
$f'_1, \ldots, f'_r$은 $k[x_1, \ldots, x_n]_{\mathfrak p'}$에서
$I'_{\mathfrak p'}$을 생성하는 정칙열로 간다.
Algebra, Lemma \ref{algebra-lemma-lci}를 적용하면,
어떤 $g' \in S$, $g' \not \in \mathfrak q'$에 대해
$S'_{gg'}$가 $k$ 위의 전역 완전교차라고 결론짓는다.
\end{proof}

\noindent
국소 완전교차 $R \to S$에 대해서는 $n \geq 2$일 때
$H_n(L_{S/R}) = 0$이다. 아직 완전한 여접복합체를 정의하지 않았으므로
이를 진술하고 증명할 수는 없지만, 그 귀결 하나를 이끌어낼 수 있다.

\begin{lemma}
\label{more-algebra-lemma-transitive-lci-at-end}
$A \to B \to C$를 환 사상들이라 하자. $B \to C$가 국소 완전교차
준동형이라고 가정하자. 핵이 $I$인 표시
$\alpha : A[x_s, s \in S] \to B$를 택하자. 핵이 $J$인 표시
$\beta : B[y_1, \ldots, y_m] \to C$를 택하자.
$K$를 핵으로 갖는 $C$의 유도된 표시를
$\gamma : A[x_s, y_t] \to C$라 하자. 그러면 행들이 완전한
다음 표준 가환 그림을 얻는다.
$$
\xymatrix{
0 \ar[r] &
\Omega_{A[x_s]/A} \otimes C \ar[r] &
\Omega_{A[x_s, y_t]/A} \otimes C \ar[r] &
\Omega_{B[y_t]/B} \otimes C \ar[r] &
0 \\
0 \ar[r] &
I/I^2 \otimes C \ar[r] \ar[u] &
K/K^2 \ar[r] \ar[u] &
J/J^2 \ar[r] \ar[u] &
0
}
$$
특히 Algebra, Lemma \ref{algebra-lemma-exact-sequence-NL}의
여섯 항 완전열은 왼쪽에 영을 붙여 완성할 수 있다. 즉 다음 열은 완전하다.
$$
0 \to H_1(\NL_{B/A} \otimes_B C) \to
H_1(L_{C/A}) \to
H_1(L_{C/B}) \to
\Omega_{B/A} \otimes_B C \to
\Omega_{C/A} \to
\Omega_{C/B} \to 0
$$
\end{lemma}

\begin{proof}
증명할 것은 사상 $I/I^2 \otimes C \to K/K^2$의 단사성뿐이다.
가정에 의해 아이디얼 $J$는 Koszul-정칙이다. 따라서
보조정리 \ref{more-algebra-lemma-conormal-sequence-H1-regular-ideal}에 의해
$IA[x_s, y_j] \cap K^2 = IK$이다. 이는
$K/K^2 \to J/J^2$의 핵이 $IA[x_s, y_j]/IK$와 동형임을 뜻한다.
% P02-E0860: frozen source tensors the conormal module over A; preserve both exact _A spans.
텐서곱의 오른쪽 완전성에 의해
$I/I^2 \otimes_A C = IA[x_s, y_j]/IK$이므로,
이는 원하는 $I/I^2 \otimes_A C \to K/K^2$의 단사성을 준다.
\end{proof}

\begin{lemma}
\label{more-algebra-lemma-transitive-colimit-lci-at-end}
$A \to B \to C$를 환 사상들이라 하자.
$B \to C$가 국소 완전교차 준동형들의 여과 여극한이면,
보조정리 \ref{more-algebra-lemma-transitive-lci-at-end}의 결론은 여전히 성립한다.
\end{lemma}

\begin{proof}
보조정리 \ref{more-algebra-lemma-transitive-lci-at-end}와
Algebra, Lemma \ref{algebra-lemma-colimits-NL}에서 따른다.
\end{proof}

\begin{lemma}
\label{more-algebra-lemma-henselization-NL}
$A \to B$를 국소환들의 국소 준동형이라 하자.
$A^h \to B^h$, 각각 $A^{sh} \to B^{sh}$를 Hensel화, 각각 엄밀 Hensel화에서
유도되는 사상이라 하자
(Algebra, Lemma \ref{algebra-lemma-henselian-functorial},
각각 Lemma \ref{algebra-lemma-strictly-henselian-functorial}).
그러면 $\NL_{B/A} \otimes_B B^h \to \NL_{B^h/A^h}$와
$\NL_{B/A} \otimes_B B^{sh} \to \NL_{B^{sh}/A^{sh}}$는
코호몰로지 군들 위에서 동형을 유도한다.
\end{lemma}

\begin{proof}
$A^h$는 $A$ 위의 에탈 대수들의 여과 여극한이므로,
Algebra, Lemma \ref{algebra-lemma-colimits-NL}와
Algebra, Definition \ref{algebra-definition-etale}에 의해
$\NL_{A^h/A}$는 비순환 복합체이다. $B^h/B$에 대해서도 마찬가지이다.
$A \to A^h \to B^h$에 Jacobi--Zariski 열
(Algebra, Lemma \ref{algebra-lemma-exact-sequence-NL})을 사용하면
$\NL_{B^h/A} \to \NL_{B^h/A^h}$가 코호몰로지 군들 위에서
동형을 유도함을 알 수 있다. 더욱이 에탈 환 사상은 전역 완전교차이기까지 하므로
국소 완전교차이다. Algebra, Lemma
\ref{algebra-lemma-etale-standard-smooth}를 보라.
보조정리 \ref{more-algebra-lemma-transitive-colimit-lci-at-end}에 의해
$A \to B \to B^h$에 붙은 여섯 항 완전 Jacobi--Zariski 열을 얻으며,
이로부터 $\NL_{B/A} \otimes_B B^h \to \NL_{B^h/A}$가
코호몰로지 군들 위에서 동형을 유도함이 증명된다.
이로써 Hensel화 위의 사상에 대한 증명이 끝난다.
엄밀 Hensel화의 경우도 정확히 같은 방식으로 증명한다.
\end{proof}







\section{Cartier 등식과 기하학적 정칙성}
\label{more-algebra-section-cartier-equality}

\noindent
이 절과 다음 절의 참고문헌은 \cite[Section 39]{MatCA}이다.
이 절을 수월하게 읽으려면 순진한 여접복합체와 그 성질들을 알고 있어야 한다.
Algebra, Section \ref{algebra-section-netherlander}를 보라.

\begin{lemma}[Cartier 등식]
\label{more-algebra-lemma-cartier-equality}
$K/k$를 유한 생성 체 확장이라 하자.
그러면 $\Omega_{K/k}$와 $H_1(L_{K/k})$는 유한 차원이고,
$\text{trdeg}_k(K) = \dim_K \Omega_{K/k} - \dim_K H_1(L_{K/k})$이다.
\end{lemma}

\begin{proof}
Algebra, Lemma \ref{algebra-lemma-colimit-syntomic}과 그 증명에 의해,
$k$ 위의 전역 완전교차
$A = k[x_1, \ldots, x_n]/(f_1, \ldots, f_c)$를 택하여
$K$가 $A$의 분수체와 동형이 되게 할 수 있다. 이 경우
Algebra, Lemmas \ref{algebra-lemma-NL-homotopy}와
\ref{algebra-lemma-localize-NL}에 의해 $\NL_{K/k}$는 다음 복합체와
호모토피 동치임을 알 수 있다.
% P02-E0861: frozen source uses malformed bounded-index subscripts in both direct sums.
$$
\bigoplus\nolimits_{j = 1, \ldots, c} K \longrightarrow
\bigoplus\nolimits_{i = 1, \ldots, n} K\text{d}x_i
$$
$k$ 위에서 $K$의 초월차수는 $A$의 차원이며
(Algebra, Lemma \ref{algebra-lemma-dimension-prime-polynomial-ring}에 의해),
이는 $n - c$이다. 따라서 결론이 따른다.
\end{proof}

\begin{lemma}
\label{more-algebra-lemma-transitivity-gamma}
$M/L/K$를 체 확장들이라 하자. 그러면 Jacobi--Zariski 열
$$
0 \to H_1(L_{L/K}) \otimes_L M \to
H_1(L_{M/K}) \to
H_1(L_{M/L}) \to
\Omega_{L/K} \otimes_L M \to
\Omega_{M/K} \to
\Omega_{M/L} \to 0
$$
은 완전하다.
\end{lemma}

\begin{proof}
보조정리 \ref{more-algebra-lemma-transitive-colimit-lci-at-end}와
Algebra, Lemma \ref{algebra-lemma-colimit-syntomic}을 결합하면 된다.
\end{proof}

\begin{lemma}
\label{more-algebra-lemma-gamma-commutative-diagram}
다음 체들의 가환 그림이 주어졌다고 하자.
$$
\xymatrix{
K \ar[r] & K' \\
k \ar[u] \ar[r] & k' \ar[u]
}
$$
$k'/k$와 $K'/K$가 유한 생성 체 확장이면, 사상들
$$
\alpha : \Omega_{K/k} \otimes_K K' \to \Omega_{K'/k'}
\quad\text{and}\quad
\beta : H_1(L_{K/k}) \otimes_K K' \to H_1(L_{K'/k'})
$$
의 핵과 여핵은 유한 차원이고 다음이 성립한다.
$$
\dim \Ker(\alpha) - \dim \Coker(\alpha)
-\dim \Ker(\beta) + \dim \Coker(\beta)
=
\text{trdeg}_k(k') - \text{trdeg}_K(K')
$$
\end{lemma}

\begin{proof}
$k \subset k' \subset K'$와 $k \subset K \subset K'$에 대한
Jacobi--Zariski 열은 각각 다음과 같다.
$$
0 \to H_1(L_{k'/k}) \otimes K'  \to
H_1(L_{K'/k}) \to
H_1(L_{K'/k'}) \to
\Omega_{k'/k} \otimes K' \to
\Omega_{K'/k} \to
\Omega_{K'/k'} \to 0
$$
그리고
$$
0 \to H_1(L_{K/k}) \otimes K' \to
H_1(L_{K'/k}) \to
H_1(L_{K'/K}) \to
\Omega_{K/k} \otimes K' \to
\Omega_{K'/k} \to
\Omega_{K'/K} \to 0
$$
보조정리 \ref{more-algebra-lemma-cartier-equality}에 의해 벡터공간들
$\Omega_{k'/k}$, $\Omega_{K'/K}$, $H_1(L_{K'/K})$, $H_1(L_{k'/k})$는
유한 차원이고, 그 차원들의 교대합은
$\text{trdeg}_k(k') - \text{trdeg}_K(K')$이다. 이제 보조정리가 따른다.
\end{proof}





\section{기하학적 정칙성}
\label{more-algebra-section-geometrically-regular}

\noindent
$k$를 체라 하고, $(A, \mathfrak m, K)$를 Noether 국소
$k$-대수라 하자. Algebra, Lemma \ref{algebra-lemma-exact-sequence-NL}의
Jacobi--Zariski 열은 다음 표준 완전열이다.
$$
H_1(L_{K/k}) \to \mathfrak m/\mathfrak m^2 \to \Omega_{A/k} \otimes_A K \to
\Omega_{K/k} \to 0
$$
실제로 Algebra, Lemma \ref{algebra-lemma-NL-surjection}에 의해
$H_1(L_{K/A}) = \mathfrak m/\mathfrak m^2$이다.
이 열의 왼쪽에서의 완전성이 정칙 국소환 $A$가 $k$ 위에서
기하학적으로 정칙인지 여부를 특징짓는다는 것을 보일 것이다.
이를 절 \ref{more-algebra-section-regular-fs}에서 형식적 매끄러움의 개념과 연결할 것이다.

\begin{proposition}
\label{more-algebra-proposition-characterization-geometrically-regular}
$k$를 표수가 $p > 0$인 체라 하자.
$(A, \mathfrak m, K)$를 Noether 국소 $k$-대수라 하자.
다음은 서로 동치이다.
\begin{enumerate}
\item $A$는 $k$ 위에서 기하학적으로 정칙이다.
\item $k$ 위에서 유한인 모든 $k \subset k' \subset k^{1/p}$에 대해
환 $A \otimes_k k'$은 정칙이다.
\item $A$는 정칙이고 표준 사상
$H_1(L_{K/k}) \to \mathfrak m/\mathfrak m^2$는 단사이다.
\item $A$는 정칙이고 사상
$\Omega_{k/\mathbf{F}_p} \otimes_k K \to \Omega_{A/\mathbf{F}_p} \otimes_A K$는
단사이다.
\end{enumerate}
\end{proposition}

\begin{proof}
(3) $\Rightarrow$ (1)의 증명.
(3)을 가정하자. $k'/k$를 유한 순수 비분리 확장이라 하자.
$A' = A \otimes_k k'$이라 놓자. 이는 극대 아이디얼
$\mathfrak m'$을 갖는 국소환이다. $K' = A'/\mathfrak m'$이라 놓자.
그러면 행들이 완전한 다음 가환 그림을 얻는다.
$$
\xymatrix{
0 \ar[r] &
H_1(L_{K/k}) \otimes K' \ar[r] \ar[d]_\beta &
\mathfrak m/\mathfrak m^2 \otimes K' \ar[r] \ar[d] &
\Omega_{A/k} \otimes_A K' \ar[r] \ar[d]_{\cong} &
\Omega_{K/k} \otimes K' \ar[r] \ar[d]_\alpha &
0
\\
 &
H_1(L_{K'/k'}) \ar[r] &
\mathfrak m'/(\mathfrak m')^2 \ar[r] &
\Omega_{A'/k'} \otimes_{A'} K' \ar[r] &
\Omega_{K'/k'} \ar[r] &
0
}
$$
미분 가군의 기저변환
(Algebra, Lemma \ref{algebra-lemma-differentials-base-change})에 의해
세 번째 세로 화살표는 동형이다. 따라서 $\alpha$는 전사이다.
보조정리 \ref{more-algebra-lemma-gamma-commutative-diagram}에 의해
$$
\dim \Ker(\alpha) - \dim \Ker(\beta) + \dim \Coker(\beta) = 0
$$
이다(이 차원들은 모두 유한하다). 그림 추적으로 다음을 얻는다.
$\dim \mathfrak m'/(\mathfrak m')^2 \leq \dim \mathfrak m/\mathfrak m^2$.
한편 $A \to A'$는 유한 평탄이므로
$\dim(A) = \dim(A')$이다. Algebra, Lemma
\ref{algebra-lemma-dimension-base-fibre-total}을 보라.
따라서 정의에 의해 $A'$은 정칙이다.

\medskip\noindent
(3)과 (4)의 동치. 다음 가환 그림의 행들에 대한
Jacobi--Zariski 열들을 생각하자.
$$
\xymatrix{
\mathbf{F}_p \ar[r] & A \ar[r] & K \\
\mathbf{F}_p \ar[r] \ar[u] & k \ar[r] \ar[u] & K \ar[u]
}
$$
그러면 행들이 완전한 다음 가환 그림을 얻는다.
$$
\xymatrix{
0 \ar[r] &
\mathfrak m/\mathfrak m^2 \ar[r] &
\Omega_{A/\mathbf{F}_p} \otimes_A K \ar[r] &
\Omega_{K/\mathbf{F}_p} \ar[r] & 0 & \\
0 \ar[r] &
H_1(L_{K/k}) \ar[r] \ar[u] &
\Omega_{k/\mathbf{F}_p} \otimes_k K \ar[r] \ar[u] &
\Omega_{K/\mathbf{F}_p} \ar[r] \ar[u] &
\Omega_{K/k} \ar[r] \ar[u] &
0
}
$$
여기서 $H_1(L_{K/A}) = \mathfrak m/\mathfrak m^2$라는 사실과,
$K/\mathbf{F}_p$가 분리이므로 $H_1(L_{K/\mathbf{F}_p}) = 0$이라는 사실을
사용했다. Algebra, Proposition
\ref{algebra-proposition-characterize-separable-field-extensions}를 보라.
따라서
$H_1(L_{K/k}) \to \mathfrak m/\mathfrak m^2$의 핵과
$\Omega_{k/\mathbf{F}_p} \otimes_k K \to \Omega_{A/\mathbf{F}_p} \otimes_A K$의
핵이 같은 차원을 갖는다는 것은 명백하다.

\medskip\noindent
Faltings를 따르는 (2) $\Rightarrow$ (4)의 증명은
\cite{Faltings-einfacher}를 보라.
$\Omega_{k/\mathbf{F}_p}$에서 $\text{d}a_1, \ldots, \text{d}a_n$이
선형 독립이 되게 하는 원소들 $a_1, \ldots, a_n \in k$를 택하자.
체 확장 $k' = k(a_1^{1/p}, \ldots, a_n^{1/p})$을 생각하자.
Algebra, Lemma \ref{algebra-lemma-size-extension-pth-roots}에 의해
$k' = k[x_1, \ldots, x_n]/(x_1^p - a_1, \ldots, x_n^p - a_n)$임을 알 수 있다.
특히 $k'/k$의 순진한 여접복합체는 다음 영 미분 복합체와 호모토픽이다.
% P02-E0862: frozen source has four malformed bounded-index direct sums below.
$\bigoplus_{j = 1, \ldots, n} k' \rightarrow \bigoplus_{i = 1, \ldots, n} k'$.
실제로 $\Omega_{k[x_1, \ldots, x_n]/k}$에서
$\text{d}(x_j^p - a_j) = 0$이다.
위와 같이 $A' = A \otimes_k k'$ 및 $K' = A'/\mathfrak m'$이라 놓자.
Algebra, Lemma \ref{algebra-lemma-change-base-NL}에 의해
$\NL_{A'/A}$는 영 미분을 갖는 복합체
$\bigoplus_{j = 1, \ldots, n} A' \rightarrow \bigoplus_{i = 1, \ldots, n} A'$와
호모토피 동치이다. 즉 $H_1(L_{A'/A})$와 $\Omega_{A'/A}$는
계수 $n$의 자유 가군이다. $\mathbf{F}_p \to A \to A'$에 대한
Jacobi--Zariski 열은 다음과 같다.
$$
H_1(L_{A'/A}) \to \Omega_{A/\mathbf{F}_p} \otimes_A A'
\to \Omega_{A'/\mathbf{F}_p} \to \Omega_{A'/A} \to 0
$$
핵이 $(x_j^p - a_j)$인 표시 $A[x_1, \ldots, x_n] \to A'$를 사용하고
Algebra, Lemma \ref{algebra-lemma-exact-sequence-NL}의 사상들을 풀어 쓰면,
$H_1(L_{A'/A})$의 $j$번째 기저벡터가
$\Omega_{A/\mathbf{F}_p} \otimes A'$ 안의 $\text{d}a_j \otimes 1$로 감을 알 수 있다.
$\Omega_{A'/A}$가 자유이고 따라서 평탄이므로, $K'$와 텐서곱하여
다음 완전열을 얻는다.
$$
K'^{\oplus n} \to \Omega_{A/\mathbf{F}_p} \otimes_A K'
\xrightarrow{\beta} \Omega_{A'/\mathbf{F}_p} \otimes_{A'} K' \to
K'^{\oplus n} \to 0
$$
원소들 $\text{d}a_j \otimes 1$이 $\Ker(\beta)$를 생성한다고 결론짓는다.
% P02-E0863: frozen source omits the subject of “are linearly independent”; Korean supplies it.
이들이 선형 독립임을, 즉 $\dim(\Ker(\beta)) = n$임을 보여야 한다.
다음 큰 그림을 생각하자.
$$
\xymatrix{
0 \ar[r] &
\mathfrak m'/(\mathfrak m')^2 \ar[r] &
\Omega_{A'/\mathbf{F}_p} \otimes K' \ar[r] &
\Omega_{K'/\mathbf{F}_p} \ar[r] & 0 \\
0 \ar[r] &
\mathfrak m/\mathfrak m^2 \otimes K' \ar[r] \ar[u]^\alpha &
\Omega_{A/\mathbf{F}_p} \otimes K' \ar[r] \ar[u]^\beta &
\Omega_{K/\mathbf{F}_p} \otimes K' \ar[r] \ar[u]^\gamma & 0
}
$$
보조정리 \ref{more-algebra-lemma-cartier-equality}과
$\mathbf{F}_p \to K \to K'$에 대한 Jacobi--Zariski 열에 의해
$\gamma$의 핵과 여핵은 같은 유한 차원을 갖는다.
가정에 의해 $A'$은 정칙이고 (위에서 보았듯 $A$와 차원이 같으므로),
$\alpha$의 핵과 여핵도 같은 차원을 갖는다. 따라서 $\beta$의 핵과 여핵은
같은 차원을 가지며, 이것이 보이고자 한 바이다.

\medskip\noindent
(1) $\Rightarrow$ (2)는 자명하다. 이로써 명제의 증명이 끝난다.
\end{proof}

\begin{lemma}
\label{more-algebra-lemma-geometrically-regular-over-field}
$k$를 표수가 $p > 0$인 체라 하자. $(A, \mathfrak m, K)$를
Noether 국소 $k$-대수라 하고, $A$가 $k$ 위에서 기하학적으로 정칙이라고 하자.
$K/F/k$를 유한 생성 부분확장이라 하자.
$\varphi : k[y_1, \ldots, y_m] \to A$를 다음 성질을 갖는 $k$-대수 사상이라 하자.
$y_i$는 $K$ 안의 $F$의 원소로 가고,
$\text{d}y_1, \ldots, \text{d}y_m$은 $\Omega_{F/k}$의 기저로 간다.
$\mathfrak p = \varphi^{-1}(\mathfrak m)$이라 놓자. 그러면
$$
k[y_1, \ldots, y_m]_\mathfrak p \to A
$$
는 평탄이고 $A/\mathfrak pA$는 정칙이다.
\end{lemma}

\begin{proof}
$A_0 = k[y_1, \ldots, y_m]_\mathfrak p$라 놓고, 그 극대 아이디얼을
$\mathfrak m_0$, 잉여체를 $K_0$라 하자. $\Omega_{A_0/k}$는 계수 $m$의 자유 가군이고
$\Omega_{A_0/k} \otimes K_0 \to \Omega_{K_0/k}$는 동형임에 유의하자.
$A_0$가 $k$ 위에서 기하학적으로 정칙임은 명백하다. 따라서
명제 \ref{more-algebra-proposition-characterization-geometrically-regular}에 의해
$H_1(L_{K_0/k}) \to \mathfrak m_0/\mathfrak m_0^2$는 동형이다.
이제 다음을 생각하자.
$$
\xymatrix{
H_1(L_{K_0/k}) \otimes K \ar[d] \ar[r] &
\mathfrak m_0/\mathfrak m_0^2 \otimes K \ar[d] \\
H_1(L_{K/k}) \ar[r] & \mathfrak m/\mathfrak m^2
}
$$
왼쪽 세로 화살표는 보조정리 \ref{more-algebra-lemma-transitivity-gamma}에 의해 단사이고,
% P02-E0864: frozen source omits the injectivity predicate for the lower horizontal arrow.
아래 가로 화살표는 명제 \ref{more-algebra-proposition-characterization-geometrically-regular}에 의해
단사이므로, 오른쪽 세로 화살표도 단사라고 결론짓는다.
따라서 $A_0$의 정칙 매개변수계는 $A$의 정칙 매개변수계의 일부로 간다.
Algebra, Lemmas \ref{algebra-lemma-flat-over-regular}와
\ref{algebra-lemma-regular-ring-CM}에 의해 결론이 따른다.
\end{proof}












\section{위상환과 위상 가군}
\label{more-algebra-section-topological-ring}

\noindent
위상 아벨 군의 몇 가지 성질을 간단히 논의하자.
아벨 군 $M$에 위상을 주어 덧셈 $M \times M \to M$,
$(x, y) \mapsto x + y$와 역원 $M \to M$, $x \mapsto -x$가
연속이 되면 $M$을 {\it 위상 아벨 군}이라 한다.
{\it 위상 아벨 군의 준동형}은 연속인 아벨 군 준동형일 뿐이다.
가환 위상군들의 범주는 가법적이고 핵과 여핵을 갖지만,
공리 $\Im = \Coim$이 성립하지 않으므로 아벨 범주는 아니다.
$N \subset M$이 부분군이면, $N$에는 유도 위상을,
$M/N$에는 몫위상을 주어 (즉 $M \to M/N$이 내려보내기 사상이 되게 하여)
$N$과 $M/N$도 위상군으로 간주한다.
$N \subset M$이 열린 부분군이면 $M/N$의 위상은 이산임에 유의하자.

\medskip\noindent
$0$의 근방 기본계 중 부분군들로 이루어진 것이 존재하면
$M$의 위상을 {\it 선형}이라 한다. 이 경우 이 부분군들은 열린 부분군이기도 하다.
다음은 한 예이다. $I$를 유향집합이라 하고, $G_i$를 $I$ 위의
(이산) 아벨 군들의 역계를 이루게 하자. 그러면
$$
G = \lim_{i \in I} G_i
$$
에 역극한 위상을 주면 선형 위상화되며,
$\Ker(G \to G_i)$가 $0$의 근방 기본계를 이룬다.
반대로 $M$을 선형 위상화된 아벨 군이라 하자.
열린 부분군들의 임의의 기본계 $U_i \subset M$, $i \in I$를 택하자
(즉 $U_i$들은 열린 근방의 기본계를 이루고 각 $U_i$는 $M$의 부분군이다).
$i \geq i' \Leftrightarrow U_i \subset U_{i'}$로 놓으면
$I$는 유향집합이다. 선형 위상화된 아벨 군들의 준동형
$$
c : M \longrightarrow \lim_{i \in I} M/U_i.
$$
를 얻는다. $M$이 (위상공간으로서) {\it 분리}일 필요충분조건은
$c$가 단사인 것임이 명백하다. $c$가 동형이면 $M$을 {\it 완비}라 한다.
% P02-E0865: frozen footnote says “a unique limits”; Korean renders the intended grammar.
\footnote{완비 군에서 극한이 유일하기를 원하므로 분리성을 완비성의 일부로 포함한다.
임의의 위상군에 대한 완비성의 정의도 있으며, 분리성 문제를 제외하면
우리의 특수한 경우에는 이 정의와 일치한다.}
이 조건이 위에서 택한 열린 부분군의 기본계
$\{U_i\}_{i \in I}$의 선택과 무관함을 확인하는 일은 독자에게 맡긴다.
실제로 위상 아벨 군 $M^\wedge = \lim_{i \in I} M/U_i$는
이 선택과 무관하며, 때로는 $M$의 {\it 완비화}라 부른다.
위와 같은 임의의 $G = \lim G_i$는 완비이다. 특히 완비화
$M^\wedge$는 항상 완비이다.

\begin{definition}[위상환]
\label{more-algebra-definition-topological-ring}
\begin{reference}
\cite[Chapter 0, Sections 7.1 and 7.2]{EGA1}
\end{reference}
$R$을 환이라 하고 $M$을 $R$-가군이라 하자.
\begin{enumerate}
\item $R$에 위상을 주어 덧셈과 곱셈이 사상
$R \times R \to R$으로서 모두 연속이 되게 하되,
여기서 $R \times R$에는 곱위상을 준다면 $R$을 {\it 위상환}이라 한다.
이 경우 $M$에 위상을 주어 덧셈 $M \times M \to M$과
스칼라 곱셈 $R \times M \to M$이 연속이면 $M$을 {\it 위상 가군}이라 한다.
\item {\it 위상 가군의 준동형}은 연속인 $R$-가군 사상일 뿐이다.
{\it 위상환의 준동형}은 주어진 위상에 대해 연속인 환 준동형이다.
\item $0$이 부분가군들로 이루어진 근방 기본계를 가지면
$M$을 {\it 선형 위상화되었다}고 한다. $0$이 아이디얼들로 이루어진
근방 기본계를 가지면 $R$을 {\it 선형 위상화되었다}고 한다.
\item $R$이 선형 위상화되었다고 하자. $I \subset R$이 열려 있고,
$0$의 모든 근방이 어떤 $n$에 대해 $I^n$을 포함하면
$I$를 {\it 정의 아이디얼}이라 한다.
\item $R$이 선형 위상화되었다고 하자. $R$이 정의 아이디얼을 가지면
$R$을 {\it pre-admissible}이라 한다.
\item $R$이 선형 위상화되었다고 하자. 이것이 pre-admissible이고
완비이면 $R$을 {\it admissible}이라 한다
\footnote{우리의 규약에서 이는 분리성을 포함한다.}.
\item $R$이 선형 위상화되었다고 하자. 어떤 정의 아이디얼 $I$가 존재하여
$\{I^n\}_{n \geq 0}$이 $0$의 근방 기본계를 이루면
$R$을 {\it pre-adic}이라 한다.
\item $R$이 선형 위상화되었다고 하자. $R$이 pre-adic이고 완비이면
$R$을 {\it adic}이라 한다.
\end{enumerate}
(Pre)adic 위상환은 모든 $n \geq 1$에 대해 $I^n$이 열린
정의 아이디얼 $I$를 갖는 (pre)admissible 위상환과 같음에 유의하자.
\end{definition}

\noindent
$R$을 환, $M$을 $R$-가군, $I \subset R$을 아이디얼이라 하자.
그러면 $\{I^n\}_{n \geq 0}$을 $0$의 근방 기본계로 갖는
$R$ 위의 선형 위상을 생각할 수 있다. 이 위상을 $I$-진 위상이라 한다.
$I$-진 위상에서 $R$은 pre-adic 위상환이다
\footnote{따라서 $I$-진 위상을 때로 $I$-pre-adic 위상이라고도 한다.}.
더욱이 $\{I^nM\}_{n \geq 0}$을 $0$의 열린 근방 기본계로 갖는
$M$ 위의 선형 위상은 $M$을 위상 $R$-가군으로 만든다.
이를 $M$ 위의 {\it $I$-진 위상}이라 한다.
Algebra, Definition \ref{algebra-definition-complete}에서 정의한 의미로
$M$이 $I$-진 완비일 필요충분조건은 $I$-진 위상에서 $M$이 완비인 것이다.
\footnote{$I$-진 완비화 $M^\wedge$가 극한 위상에 대해서는 항상 완비이지만,
$M^\wedge$가 $I$-진 완비가 아닐 수도 있다. $I$가 유한 생성이면 $M^\wedge$ 위의
$I$-진 위상과 극한 위상은 일치한다. Algebra, Lemma
\ref{algebra-lemma-hathat-finitely-generated}와 그 증명을 보라.}
특히 $R$이 $I$-진 완비일 필요충분조건은 $I$-진 위상에서
$R$이 adic 위상환인 것이다.

\medskip\noindent
특수한 경우로, 이산 위상은 $0$-진 위상이고 이산 위상의 모든 환은
adic임에 유의하자.

\begin{lemma}
\label{more-algebra-lemma-continuous}
$\varphi : R \to S$를 환 사상이라 하자.
$I \subset R$과 $J \subset S$를 아이디얼들이라 하고,
$R$에는 $I$-진 위상을, $S$에는 $J$-진 위상을 주자.
그러면 $\varphi$가 위상환의 준동형일 필요충분조건은
어떤 $n \geq 1$에 대해 $\varphi(I^n) \subset J$인 것이다.
\end{lemma}

\begin{proof}
생략한다.
\end{proof}

\begin{lemma}[Baire 범주 정리]
\label{more-algebra-lemma-baire-category-complete-module}
$M$을 위상 아벨 군이라 하자. $M$이 선형 위상화되고 완비이며,
$0$의 가산 근방 기본계를 갖는다고 하자.
$U_n \subset M$, $n \geq 1$이 열린 조밀 부분집합들이면
$\bigcap_{n \geq 1} U_n$은 조밀하다.
\end{lemma}

\begin{proof}
$U_n$들을 보조정리의 진술과 같이 하자.
$U_n$을 $U_1 \cap \ldots \cap U_n$으로 바꾸어
$U_1 \supset U_2 \supset \ldots$라고 가정해도 된다.
$M_n$, $n \in \mathbf{N}$을 $0$의 근방 기본계라 하자.
$M_{n + 1} \subset M_n$이라고 가정해도 된다. $x \in M$을 택하자.
모든 $k \geq 1$에 대해 $x - y \in M_k$를 만족하는
$y \in \bigcap_{n \geq 1} U_n$이 존재함을 보이겠다.

\medskip\noindent
$y$를 다음과 같이 구성한다. 먼저 $y_1 \in x + M_k$인
$y_1 \in U_1$을 택한다. $U_1$이 조밀하고 $x + M_k$가 열려 있으므로 가능하다.
그런 다음 $y_1 + M_{k_1} \subset U_1$이 되게 하는 $k_1 > k$를 택한다.
$U_1$이 열려 있으므로 가능하다. 다음으로
$y_2 \in y_1 + M_{k_1}$인 $y_2 \in U_2$를 택한다.
% P02-E0866: frozen source justifies density using y_2 + M_{k_1}; preserve that exact span.
$U_2$가 조밀하고 $y_2 + M_{k_1}$가 열려 있으므로 가능하다.
그런 다음 $y_2 + M_{k_2} \subset U_2$가 되게 하는 $k_2 > k_1$을 택한다.
$U_2$가 열려 있으므로 가능하다.

\medskip\noindent
이와 같이 계속하면 $M$의 원소들로 이루어진 수렴열 $y_i$와
그 극한 $y$를 얻는다. 구성에 의해 $x - y \in M_k$이다. 또한
$$
y - y_i = (y_{i + 1} - y_i) + (y_{i + 2} - y_{i + 1}) + \ldots
$$
가 $M_{k_i}$에 속하므로, 원하는 대로 모든 $i$에 대해
$y \in y_i + M_{k_i} \subset U_i$임을 알 수 있다.
\end{proof}

\begin{lemma}
\label{more-algebra-lemma-consequence-baire-complete-module}
% P02-E0867: frozen source omits “the” in “With same assumptions”; Korean is grammatical.
보조정리 \ref{more-algebra-lemma-baire-category-complete-module}와 같은 가정 아래,
어떤 닫힌 부분군들 $N_n$에 대해 $M = \bigcup_{n \geq 1} N_n$이면,
어떤 $n$에 대해 $N_n$은 열려 있다.
\end{lemma}

\begin{proof}
그렇지 않다면 모든 $n$에 대해 $U_n = M \setminus N_n$이 조밀하므로,
보조정리 \ref{more-algebra-lemma-baire-category-complete-module}와 모순이다.
\end{proof}

\begin{lemma}[열린 사상 보조정리]
\label{more-algebra-lemma-open-mapping}
$u : N \to M$을 위상화된 아벨 군들의 연속 사상이라 하자.
% P02-E0868: frozen source omits the copula after M; Korean supplies it.
$M$이 분리되어 있고, $N$은 완비이고 선형 위상화되며
$0$의 가산 근방 기본계를 갖는다고 하자. 그러면 다음 중 정확히 하나가 성립한다.
\begin{enumerate}
\item $u$는 열린 사상이다.
\item 어떤 열린 부분군 $N' \subset N$에 대해 상 $u(N')$은
$M$ 안에서 어느 곳에서도 조밀하지 않다.
\end{enumerate}
\end{lemma}

\begin{proof}
$N_n$, $n \in \mathbf{N}$을 $0$의 근방 기본계라 하자.
$N_n$이 부분군이고 $N_{n + 1} \subset N_n$이라고 가정해도 된다.
(2)가 성립하지 않으면, $u(N_n)$의 폐포 $M_n$은
$n = 1, 2, 3, \ldots$에 대해 열린 부분군이다.
$u$가 연속이므로 $M_n$, $n \in \mathbf{N}$은
$M$의 $0$의 열린 근방 기본계이어야 한다. 또한 $M_n$은
$u(N_n)$의 폐포이므로 모든 $n \geq 1$에 대해
$$
u(N_n) + M_{n + 1} = M_n
$$
이다. $x_1 \in M_1$을 택하자. 다음을 만족하도록
$y_i \in N_i$와 $x_{i + 1} \in M_{i + 1}$을 귀납적으로 택할 수 있다.
$$
u(y_i) + x_{i + 1} = x_i
$$
$N$이 완비이므로 $N$의 원소 $y = y_1 + y_2 + y_3 + \ldots$가 존재한다.
그러면 $M$이 분리되어 있으므로 $x_1 = u(y)$임을 알 수 있다.
따라서 $M_1 = u(N_1)$이다. 정확히 같은 방식으로 독자는
모든 $i \geq 2$에 대해 $M_i = u(N_i)$임을 보일 수 있으며,
따라서 $u$는 열린 사상이다.
\end{proof}









\section{위상환의 형식적으로 매끄러운 사상}
\label{more-algebra-section-formally-smooth}

\noindent
위상환의 준동형에 적용되는 형식적 매끄러움의 판본이 있다.

\begin{definition}
\label{more-algebra-definition-formally-smooth}
$R \to S$를 선형 위상화된 $R$과 $S$ 사이의 위상환 준동형이라 하자.
위상환 준동형들의 모든 가환인 실선 그림
$$
\xymatrix{
S \ar[r] \ar@{-->}[rd] & A/J \\
R \ar[r] \ar[u] & A \ar[u]
}
$$
에서 $A$가 이산환이고 $J \subset A$가 제곱영 아이디얼일 때,
그림을 가환이 되게 하는 점선 화살표가 존재하면
$S$를 {$R$ 위에서 \it 형식적으로 매끄럽다}고 한다.
\end{definition}

\noindent
주로 아이디얼들 $\mathfrak m \subset R$과 $\mathfrak n \subset S$가 주어져
$R$에는 $\mathfrak m$-진 위상을, $S$에는 $\mathfrak n$-진 위상을 주는 경우에
이 개념을 사용할 것이다. 보조정리 \ref{more-algebra-lemma-continuous}에 의해
$\varphi : R \to S$가 연속일 필요충분조건은 어떤 $m \geq 1$에 대해
$\varphi(\mathfrak m^m) \subset  \mathfrak n$인 것이다.
이 경우에는 $S$ 위의 위상만 관련된다는 것이 밝혀진다.

\begin{lemma}
\label{more-algebra-lemma-formally-smooth}
$\varphi : R \to S$를 환 사상이라 하자.
\begin{enumerate}
\item Algebra, Definition \ref{algebra-definition-formally-smooth}의 의미로
$R \to S$가 형식적으로 매끄러우면, $R$ 위의 임의의 선형 위상과
$R \to S$가 연속이 되게 하는 $S$ 위의 임의의 pre-adic 위상에 대해서도
$R \to S$는 형식적으로 매끄럽다.
\item $\mathfrak n \subset S$와 $\mathfrak m \subset R$을,
$R$ 위의 $\mathfrak m$-진 위상과 $S$ 위의 $\mathfrak n$-진 위상에 대해
$\varphi$가 연속이 되게 하는 아이디얼들이라 하자. 그러면 다음은 서로 동치이다.
\begin{enumerate}
\item $R$ 위의 $\mathfrak m$-진 위상과 $S$ 위의 $\mathfrak n$-진 위상에 대해
$\varphi$는 형식적으로 매끄럽다.
\item $R$ 위의 이산 위상과 $S$ 위의 $\mathfrak n$-진 위상에 대해
$\varphi$는 형식적으로 매끄럽다.
\end{enumerate}
\end{enumerate}
\end{lemma}

\begin{proof}
Algebra, Definition \ref{algebra-definition-formally-smooth}의 의미로
$R \to S$가 형식적으로 매끄럽다고 가정하자.
$S$가 pre-adic 위상을 가지면, $S$가 $\mathfrak n$-진 위상을 갖게 하는
아이디얼 $\mathfrak n \subset S$가 존재한다.
정의 \ref{more-algebra-definition-formally-smooth}에서와 같은 가환인 실선 그림이
주어졌다고 하자. $S \to A/J$의 연속성은 어떤 $k \geq 1$에 대해
$\mathfrak n^k$가 $A/J$에서 영으로 감을 뜻한다.
보조정리 \ref{more-algebra-lemma-continuous}를 보라. $S$가 $R$ 위에서 형식적으로
매끄럽다는 가정으로부터 환 사상 $\psi : S \to A$를 얻는다.
그러면 $\psi(\mathfrak n^k) \subset J$이고 $J^2 = 0$이므로
$\psi(\mathfrak n^{2k}) = 0$이다. 따라서 보조정리
\ref{more-algebra-lemma-continuous}에 의해 $\psi$는 연속이다. 이로써 (1)이 증명된다.

\medskip\noindent
(2)(b) $\Rightarrow$ (2)(a)의 증명은 (1)의 증명과 같다.
(2)(a)를 가정하자. $R$ 위에는 이산 위상을 사용하는
정의 \ref{more-algebra-definition-formally-smooth}에서와 같은 가환인 실선 그림이
주어졌다고 하자. $\varphi$가 연속이므로 어떤 $n \geq 1$에 대해
$\varphi(\mathfrak m^n) \subset \mathfrak n$이다.
$S \to A/J$가 연속이므로 어떤 $k \geq 1$에 대해
$\mathfrak n^k$가 $A/J$에서 영으로 간다. 따라서 사상 $R \to A$ 아래에서
$\mathfrak m^{nk}$는 $J$ 안으로 간다. 그러므로
$\mathfrak m^{2nk}$는 $A$에서 영으로 가며, $R \to A$가
$\mathfrak m$-진 위상에서 연속임을 알 수 있다.
따라서 (2)(a)는 원하는 점선 화살표를 준다.
\end{proof}

\begin{definition}
\label{more-algebra-definition-formally-smooth-adic}
$R \to S$를 환 사상이라 하고, $\mathfrak n \subset S$를 아이디얼이라 하자.
보조정리 \ref{more-algebra-lemma-formally-smooth}의 동치 조건 (2)(a)와 (2)(b)가 성립하면,
$R \to S$가 {$\mathfrak n$-진 위상에 대해 \it 형식적으로 매끄럽다}고 한다.
\end{definition}

\noindent
이 성질은 완비화들에 유전된다.

\begin{lemma}
\label{more-algebra-lemma-formally-smooth-completion}
$(R, \mathfrak m)$과 $(S, \mathfrak n)$을 유한 생성 아이디얼이 주어진
환들이라 하자. $R$과 $S$에 각각 $\mathfrak m$-진 위상과
$\mathfrak n$-진 위상을 주자. $R \to S$를 위상환의 준동형이라 하자.
다음은 서로 동치이다.
\begin{enumerate}
\item $R \to S$는 $\mathfrak n$-진 위상에 대해 형식적으로 매끄럽다.
\item $R \to S^\wedge$는 $\mathfrak n^\wedge$-진 위상에 대해
형식적으로 매끄럽다.
\item $R^\wedge \to S^\wedge$는 $\mathfrak n^\wedge$-진 위상에 대해
형식적으로 매끄럽다.
\end{enumerate}
여기서 $R^\wedge$와 $S^\wedge$는 각각 $R$의 $\mathfrak m$-진 완비화와
$S$의 $\mathfrak n$-진 완비화이다.
\end{lemma}

\begin{proof}
$\mathfrak m$이 유한 생성이라는 가정은 $R^\wedge$가
$\mathfrak mR^\wedge$-진 완비이고,
$\mathfrak mR^\wedge = \mathfrak m^\wedge$이며,
$R^\wedge/\mathfrak m^nR^\wedge = R/\mathfrak m^n$임을 함의한다.
Algebra, Lemma \ref{algebra-lemma-hathat-finitely-generated}와 그 증명을 보라.
$(S, \mathfrak n)$에 대해서도 마찬가지이다. 따라서
정의 \ref{more-algebra-definition-formally-smooth}에서와 같은 (1), (2), (3)의 그림들이
일대일로 대응함은 명백하다.
\end{proof}

\noindent
Adic 환을 다루는 이점은 더 강한 올림 성질을 얻는다는 것이다.

\begin{lemma}
\label{more-algebra-lemma-lift-continuous}
$R \to S$를 환 사상이라 하고, $\mathfrak n$을 $S$의 아이디얼이라 하자.
$R \to S$가 $\mathfrak n$-진 위상에서 형식적으로 매끄럽다고 가정하자.
다음 가환인 실선 그림을 생각하자.
$$
\xymatrix{
S \ar[r]_\psi \ar@{-->}[rd] & A/J \\
R \ar[r] \ar[u] & A \ar[u]
}
$$
여기서 화살표들은 위상환의 준동형이고, $A$는 adic이며,
$A/J$는 $A$를 닫힌 아이디얼 $J \subset A$로 나눈 몫(위상환으로서)이고,
어떤 $t \geq 1$에 대해 $J^t$는 $A$의 한 정의 아이디얼에 포함된다고 하자.
그러면 위상환의 범주에서 그림을 가환이 되게 하는 점선 화살표가 존재한다.
\end{lemma}

\begin{proof}
$I \subset A$를 어떤 $t$에 대해 $I \supset J^t$인 정의 아이디얼이라 하자.
% P02-E0869: frozen source leaves every denominator below unparenthesized; preserve exact spans.
$J$가 닫혀 있다고 가정했으므로
$A = \lim A/I^n$이고 $A/J = \lim A/J + I^n$이다.
이산 $R$-대수들 $A_{n, m} = A/J^n + I^m$의 다음 그림을 생각하자.
$$
\xymatrix{
A/J^3 + I^3 \ar[r] \ar[d] &
A/J^2 + I^3 \ar[r] \ar[d] &
A/J + I^3 \ar[d] \\
A/J^3 + I^2 \ar[r] \ar[d] &
A/J^2 + I^2 \ar[r] \ar[d] &
A/J + I^2 \ar[d] \\
A/J^3 + I \ar[r] &
A/J^2 + I \ar[r] &
A/J + I
}
$$
각 가환 사각형은 핵이 제곱영인 $R$-대수의 전사
$$
A_{n + 1, m + 1} \longrightarrow A_{n + 1, m} \times_{A_{n, m}} A_{n, m + 1}
$$
를 정의함에 유의하자. $R$-대수 사상
$\varphi_{n, m} : S \to A_{n, m}$을 귀납적으로 구성하겠다.
먼저 사상들 $\varphi_{1, m} = \psi \bmod J + I^m$이 있다.
각 사상은 $\psi$가 연속이므로 연속이다. 선택한 $\varphi_{1, 1}$에서 시작하여
위 그림의 맨 아래 행을 따라 올리면서 $S$가 $R$ 위에서 형식적으로
매끄럽다는 사실을 사용하면 사상들 $\varphi_{n, 1}$을 귀납적으로 택할 수 있다.
나머지 사상들 $\varphi_{n, m}$은 $n + m$에 관한 귀납으로 구성한다.
즉 위에 표시한 전사를 따라 쌍
$(\varphi_{n + 1, m}, \varphi_{n, m + 1})$을 올려
$\varphi_{n + 1, m + 1}$을 택한다(다시 $S$가 $R$ 위에서
형식적으로 매끄럽다는 사실을 사용한다). 이렇게 하면 모든
$\varphi_{n, m}$은 계의 전이사상들과 양립한다.
$J^t \subset I$이므로, 예를 들어
$\varphi_n = \varphi_{nt, n} \bmod I^n$은 사상 $S \to A/I^n$을 유도한다.
극한 $\varphi = \lim \varphi_n$을 취하면 사상
$S \to A = \lim A/I^n$을 얻는다. $A/J = \lim A/J + I^n$임을
보았으므로 $A/J$로 합성하면 $\psi$와 일치한다.
마지막으로 $\varphi$가 연속임을 보이자. 즉 $\psi$가 위상환의 사상이라는
가정과 보조정리 \ref{more-algebra-lemma-continuous}에 의해 어떤 $r \geq 1$에 대해
% P02-E0870: frozen source has ill-typed J + I/J; preserve exact formal span.
$\psi(\mathfrak n^r) \subset J + I/J$임을 안다. 따라서
$\varphi(\mathfrak n^r) \subset J + I$이고,
따라서 원하는 대로 $\varphi(\mathfrak n^{rt}) \subset I$이다.
\end{proof}

\begin{lemma}
\label{more-algebra-lemma-increase-ideal}
$R \to S$를 환 사상이라 하고,
$\mathfrak n \subset \mathfrak n' \subset S$를 아이디얼들이라 하자.
$R \to S$가 $\mathfrak n$-진 위상에 대해 형식적으로 매끄러우면,
$R \to S$는 $\mathfrak n'$-진 위상에 대해서도 형식적으로 매끄럽다.
\end{lemma}

\begin{proof}
생략한다.
\end{proof}

\begin{lemma}
\label{more-algebra-lemma-compose-formally-smooth}
선형 위상화된 환들의 형식적으로 매끄러운 연속 준동형들의 합성은
형식적으로 매끄럽다.
\end{lemma}

\begin{proof}
생략한다. (힌트: 이는 완전히 형식적이며 적절한 그림을 생각하면 따른다.)
\end{proof}

\begin{lemma}
\label{more-algebra-lemma-base-change-fs}
$R$, $S$를 환들이라 하고, $\mathfrak n \subset S$를 아이디얼이라 하자.
$R \to S$가 $\mathfrak n$-진 위상에 대해 형식적으로 매끄럽다고 하자.
$R \to R'$를 임의의 환 사상이라 하자. 그러면
$R' \to S' = S \otimes_R R'$은
$\mathfrak n' = \mathfrak nS'$-진 위상에서 형식적으로 매끄럽다.
\end{lemma}

\begin{proof}
정의 \ref{more-algebra-definition-formally-smooth}에서와 같은 실선 그림
$$
\xymatrix{
S \ar[r] \ar@{-->}[rrd] & S' \ar[r] \ar@{-->}[rd] & A/J \\
R  \ar[u] \ar[r] & R' \ar[r] \ar[u] & A \ar[u]
}
$$
가 주어졌다고 하자. 그러면 합성 $S \to S' \to A/J$는 연속이다.
가정에 의해 더 긴 점선 화살표가 존재한다. 텐서곱의 보편성질로부터
더 짧은 점선 화살표를 얻는다.
\end{proof}

\noindent
Algebra, Lemma \ref{algebra-lemma-descent-formally-smooth}에서
충실하게 평탄인 환 사상을 따른 형식적 매끄러움의 내림을 보았다.
현재의 위상환 상황에서도 비슷한 결과가 성립한다. 다만 여기서는 이미
상당히 유용한, 증명이 매우 간단하고 쉬운 다음 판본만 증명한다.

\begin{lemma}
\label{more-algebra-lemma-descent-fs}
$R$, $S$를 환들이라 하고, $\mathfrak n \subset S$를 아이디얼이라 하자.
$R \to R'$를 환 사상이라 하자. $S' = S \otimes_R R'$ 및
% P02-E0871: frozen source defines n' as nS rather than nS'; preserve exact span.
$\mathfrak n' = \mathfrak nS$라 놓자. 다음을 가정하자.
\begin{enumerate}
\item $R \to R'$는 $R$-가군으로서 $R$을 $R'$의 직합인자로 매장한다.
\item $R' \to S'$는 $\mathfrak n'$-진 위상에 대해 형식적으로 매끄럽다.
\end{enumerate}
그러면 $R \to S$는 $\mathfrak n$-진 위상에서 형식적으로 매끄럽다.
\end{lemma}

\begin{proof}
정의 \ref{more-algebra-definition-formally-smooth}에서와 같은 실선 그림
$$
\xymatrix{
S \ar[r] & A/J \\
R \ar[u] \ar[r] & A \ar[u]
}
$$
가 주어졌다고 하자. $A' = A \otimes_R R'$ 및
$J' = \Im(J \otimes_R R' \to A')$이라 놓자.
위 그림의 기저변환은 연속 화살표들로 이루어진 다음 그림이다.
$$
\xymatrix{
S' \ar[r] \ar@{-->}[rd]^{\psi'} & A'/J' \\
R' \ar[u] \ar[r] & A' \ar[u]
}
$$
조건 (2)에 의해 점선 화살표 $\psi' : S' \to A'$를 얻는다.
조건 (1)을 사용하여 $R$-가군으로서의 직합인자 분해
$R' = R \oplus C$를 택하자. (주의: $C$는 $R'$의 아이디얼이 아니다.)
그러면 $A' = A \oplus A \otimes_R C$이다. 다음과 같이 놓자.
$$
J'' = \Im(J \otimes_R C \to A \otimes_R C) \subset J' \subset A'.
$$
그러면 $A$-가군으로서 $J' = J \oplus J''$이다.
$\psi'$와 $S \to S'$의 합성 $\psi : S \to A'$의 상은
$A + J' = A \oplus J''$에 포함된다. 그런데 환
$A + J' = A \oplus J''$ 안에서 $A$-부분가군 $J''$은 아이디얼이다!
($J^2 = 0$임을 사용하라.) 따라서 합성
$S \to A + J' \to (A + J')/J'' = A$가 우리가 찾던 화살표이다.
\end{proof}








\section{국소환의 형식적으로 매끄러운 사상}
\label{more-algebra-section-formally-smooth-local}

\noindent
국소환들의 국소 준동형인 경우에는 올림 성질을 확인해야 하는 그림들을
제한할 수 있다. Algebra, Lemma \ref{algebra-lemma-smooth-test-artinian}과 비교하라.

\begin{lemma}
\label{more-algebra-lemma-fs-local}
$(R, \mathfrak m) \to (S, \mathfrak n)$을 국소환들의 국소 준동형이라 하자.
다음은 서로 동치이다.
\begin{enumerate}
\item $R \to S$는 $\mathfrak n$-진 위상에서 형식적으로 매끄럽다.
\item 다음과 같은 국소환들의 국소 준동형으로 이루어진 모든 가환인 실선 그림에서
$$
\xymatrix{
S \ar[r] \ar@{-->}[rd] & A/J \\
R \ar[r] \ar[u] & A \ar[u]
}
$$
$J \subset A$가 제곱영 아이디얼이고, 어떤 $n > 0$에 대해
$\mathfrak m_A^n = 0$이며, $S \to A/J$가 잉여체들 위에서 동형을 유도하면,
그림을 가환이 되게 하는 점선 화살표가 존재한다.
\end{enumerate}
$S$가 Noether이면 이 조건들은 다음과도 동치이다.
\begin{enumerate}
\item[(3)] (2)와 같되, 추가로 $A \to A/J$가 작은 확장인 그림들만 고려한다
(Algebra, Definition \ref{algebra-definition-small-extension}).
\end{enumerate}
\end{lemma}

\begin{proof}
(1) $\Rightarrow$ (2)는 정의에서 따른다.
$R$ 위의 $\mathfrak m$-진 위상과 $S$ 위의 $\mathfrak n$-진 위상에 대한
정의 \ref{more-algebra-definition-formally-smooth}에서와 같은 그림
$$
\xymatrix{
S \ar[r] \ar@{-->}[rd] & A/J \\
R \ar[r] \ar[u] & A \ar[u]
}
$$
을 생각하자. $\mathfrak n^m(A/J) = 0$이 되게 하는 $m > 0$을 택하자
(그림 속 사상들의 연속성에 의해 가능하다).
$A$에서 $S$의 $A/J$ 안의 상의 역상이 되는 부분환을 $A'$이라 하자.
$J' = J$라 놓되 이를 $A'$의 아이디얼로 보자.
그러면 $J'$은 $A'$ 안의 제곱영 아이디얼이고 $A'/J'$은
$S/\mathfrak n^m$의 몫이다. 따라서 $A'$은 국소환이고
$\mathfrak m_{A'}^{2m} = 0$이다. 그러므로 (2)에서와 같은 그림
$$
\xymatrix{
S \ar[r] \ar@{-->}[rd] & A'/J' \\
R \ar[r] \ar[u] & A' \ar[u]
}
$$
을 얻는다. 이 그림에서 점선 화살표를 구성할 수 있으면 이를
$A' \to A$와 합성하여 원래 그림의 점선 화살표를 얻는다.
이렇게 하여 (2)가 (1)을 함의함을 알 수 있다.

\medskip\noindent
$S$가 Noether라고 가정하자. (1) $\Rightarrow$ (3)은 즉시 따른다.
(3)을 가정하고 (2)에서와 같은 그림이 주어졌다고 하자.
그러면 어떤 $n > 0$에 대해 $\mathfrak m_A^n J = 0$이다. 사상들
$$
A \to A/\mathfrak m_A^{n - 1}J \to \ldots \to A/\mathfrak mJ \to A/J
$$
을 생각하면, $\mathfrak m_A J = 0$일 때 올림을 만들면 충분함을 알 수 있다.
$\mathfrak m_A J = 0$이라고 가정하고, $A' \subset A$를 위에서 구성한 환이라 하자.
그러면 $A'/J'$은 Artin 국소환 $S/\mathfrak n^m$의 몫이므로 Artin 환이다.
따라서 (3)이 주어졌을 때 $A/J$가 Artin이고
$\mathfrak m_A J = 0$인 (2)의 그림들에서 점선 화살표를 찾을 수 있음을
보이면 충분하다. $\kappa$를 $A$, $A/J$, $S$의 공통 잉여체라 하자.
(3)에 의해, $J_0 \subset J$가 $\dim_\kappa(J/J_0) = 1$인 아이디얼이면
점선 화살표 $S \to A/J_0$를 만들 수 있다. 곱을 취하면 다음을 얻는다.
$$
S \longrightarrow \prod\nolimits_{J_0 \text{ as above}} A/J_0
$$
이 화살표의 상은 작은 대각선
$A/J \subset \prod_{J_0} A/J$ 안으로 가는 원소들로 이루어진
부분 $R$-대수 $A'$에 포함됨이 명백하다.
$A/J$에서 영으로 가는 원소들로 이루어진 부분집합을 $J' \subset A'$이라 하자.
그러면 $J'$은 제곱영 아이디얼이고, $\kappa$-벡터공간으로서 다음과 같다.
$$
J' = \prod\nolimits_{J_0 \text{ as above}} J/J_0
$$
따라서 사상 $J \to J'$은 단사이다. 벡터공간 이론에 의해 분할
$J' = J \oplus M$을 택할 수 있다. 그러면 $R$-대수로서
$$
A' = A \oplus M
$$
이다. 따라서 사상 $S \to A'$을 사영 $A' \to A$와 합성하면
원하는 점선 화살표를 얻으며, 이로써 보조정리의 증명이 끝난다.
\end{proof}

\noindent
다음 보조정리는 절 \ref{more-algebra-section-regular-fs}에서 개선될 것이다.

\begin{lemma}
\label{more-algebra-lemma-fs-implies-regular}
$k$를 체라 하고 $(A, \mathfrak m, K)$를 Noether 국소
$k$-대수라 하자. $k \to A$가 $\mathfrak m$-진 위상에 대해
형식적으로 매끄러우면, $A$는 정칙 국소환이다.
\end{lemma}

\begin{proof}
$k_0 \subset k$를 소체라 하자. 그러면 $k_0$는 완전체이므로
$k / k_0$는 분리이고, 따라서 Algebra, Lemma
\ref{algebra-lemma-formally-smooth-extensions-easy}에 의해 형식적으로 매끄럽다.
보조정리 \ref{more-algebra-lemma-formally-smooth}와 \ref{more-algebra-lemma-compose-formally-smooth}에 의해
$k_0 \to A$는 $A$ 위의 $\mathfrak m$-진 위상에 대해 형식적으로 매끄럽다.
따라서 $k = \mathbf{Q}$ 또는 $k = \mathbf{F}_p$라고 가정해도 된다.

\medskip\noindent
Algebra, Lemmas \ref{algebra-lemma-completion-faithfully-flat}와
\ref{algebra-lemma-flat-under-regular}에 의해 완비화 $A^\wedge$가
정칙임을 증명하면 충분하다. 보조정리
\ref{more-algebra-lemma-formally-smooth-completion}에 의해 $A$를 $A^\wedge$로 바꾸어도 된다.
따라서 $A$가 Noether 완비 국소환이라고 가정해도 된다.
% P02-E0872: frozen source says “there exist a map”; Korean supplies singular grammar.
Cohen 구조 정리(Algebra, Theorem
\ref{algebra-theorem-cohen-structure-theorem})에 의해 사상 $K \to A$가 존재한다.
$k$가 소체이므로 $K \to A$는 $k$-대수 사상이다.

\medskip\noindent
$\mathfrak m/\mathfrak m^2$의 기저를 이루는 상들을 갖는 원소들
$x_1, \ldots, x_n \in \mathfrak m$을 택하자.
$T = K[[X_1, \ldots, X_n]]$이라 놓자. 다음에 유의하자.
$$
A/\mathfrak m^2 \cong K[x_1, \ldots, x_n]/(x_ix_j)
$$
그리고
$$
T/\mathfrak m_T^2 \cong K[X_1, \ldots, X_n]/(X_iX_j).
$$
% P02-E0873: frozen source target T/m_T^2 omits mathfrak; preserve exact span.
$x_i$의 류를 $X_i$의 류로 보내어 주어지는 국소 $K$-대수 동형
$A/\mathfrak m^2 \to T/m_T^2$를 택하자.
이 동형과 몫사상 $A \to A/\mathfrak m^2$의 합성을
$f_1 : A \to T/\mathfrak m_T^2$로 나타내자.
$k \to A$가 $\mathfrak m$-진 위상에서 형식적으로 매끄럽다는 가정은
$f_1$을 사상 $f_2 : A \to T/\mathfrak{m}_T^3$으로,
그다음 사상 $f_3 : A \to T/\mathfrak{m}_T^4$로, 이어서 모든
$n \geq 1$에 대해 계속 올릴 수 있음을 뜻한다.
주의: 사상들 $f_n$은 연속 $k$-대수 사상이지만 $K$-대수 사상일 필요는 없다.
국소 $k$-대수의 유도된 사상
$f : A \to T = \lim T/\mathfrak m_T^n$을 얻는다.
$f_1$의 선택에 의해 $f$는 동형
$\mathfrak m/\mathfrak m^2 \to \mathfrak m_T/\mathfrak m_T^2$을 유도한다.
따라서 각 $f_n$은 전사이고, $A$가 완비이므로 $f$는 전사라고 결론짓는다.
이는 $\dim(A) \geq \dim(T) = n$을 함의한다. 그러므로 정의에 의해
$A$는 정칙이다. (또한 $f$가 동형임도 따른다.)
\end{proof}

\begin{lemma}
\label{more-algebra-lemma-lift-residue-field}
$k$를 체라 하고, $(A, \mathfrak m, \kappa)$를 완비 국소
$k$-대수라 하자. $\kappa/k$가 분리이면,
$\kappa \to A \to \kappa$가 $\text{id}_\kappa$가 되게 하는
$k$-대수 사상 $\kappa \to A$가 존재한다.
\end{lemma}

\begin{proof}
Algebra, Proposition
\ref{algebra-proposition-characterize-separable-field-extensions}에 의해
확장 $\kappa/k$는 형식적으로 매끄럽다. 보조정리
\ref{more-algebra-lemma-formally-smooth}에 의해 $k \to \kappa$는
정의 \ref{more-algebra-definition-formally-smooth}의 의미로 형식적으로 매끄럽다.
이제 보조정리 \ref{more-algebra-lemma-lift-continuous}에서 $\kappa \to A$를 얻는다.
\end{proof}

\begin{lemma}
\label{more-algebra-lemma-power-series-over-residue-field}
$k$를 체라 하고, $(A, \mathfrak m, \kappa)$를 완비 국소
$k$-대수라 하자. $\kappa/k$가 분리이고 $A$가 정칙이면,
% P02-E0874: frozen source says “an isomorphism of A”; Korean supplies natural syntax.
$k$-대수의 동형 $A \cong \kappa[[t_1, \ldots, t_d]]$가 존재한다.
\end{lemma}

\begin{proof}
보조정리 \ref{more-algebra-lemma-lift-residue-field}에서와 같은 $\kappa \to A$를 택하고
Algebra, Lemma
\ref{algebra-lemma-regular-complete-containing-coefficient-field}를 적용한다.
\end{proof}

\noindent
다음 결과는 절 \ref{more-algebra-section-regular-fs}에서 개선될 것이다.

\begin{lemma}
\label{more-algebra-lemma-regular-implies-fs}
$k$를 체라 하자. $(A, \mathfrak m, K)$를 $K/k$가 분리인
정칙 국소 $k$-대수라 하자. 그러면 $k \to A$는
$\mathfrak m$-진 위상에서 형식적으로 매끄럽다.
\end{lemma}

\begin{proof}
보조정리 \ref{more-algebra-lemma-formally-smooth-completion}에 의해 $A$의 완비화가
$k$ 위에서 형식적으로 매끄러움을 증명하면 충분하다.
따라서 $A$가 완비 정칙 국소 $k$-대수이고 그 잉여체 $K$가
$k$ 위에서 분리라고 가정해도 된다.
보조정리 \ref{more-algebra-lemma-power-series-over-residue-field}에 의해
$A = K[[x_1, \ldots, x_n]]$임을 알 수 있다.

\medskip\noindent
멱급수환 $K[[x_1, \ldots, x_n]]$은 $k$ 위에서 형식적으로 매끄럽다.
실제로 $K$는 $k$ 위에서 형식적으로 매끄럽고,
$K[x_1, \ldots, x_n]$은 다항식 대수이므로 $K$ 위에서 형식적으로 매끄럽다.
따라서 Algebra, Lemma \ref{algebra-lemma-compose-formally-smooth}에 의해
$K[x_1, \ldots, x_n]$은 $k$ 위에서 형식적으로 매끄럽다.
그러므로 보조정리 \ref{more-algebra-lemma-formally-smooth}에 의해
$k \to K[x_1, \ldots, x_n]$은 $(x_1, \ldots, x_n)$-진 위상에 대해
형식적으로 매끄럽다. 마지막으로 보조정리
\ref{more-algebra-lemma-formally-smooth-completion}에 의해
$k \to K[[x_1, \ldots, x_n]]$은 $(x_1, \ldots, x_n)$-진 위상에 대해
형식적으로 매끄럽다.
\end{proof}

\begin{lemma}
\label{more-algebra-lemma-formally-smooth-finite-type}
$A \to B$를 $A$가 Noether인 유한형 환 사상이라 하자.
$\mathfrak p \subset A$ 위에 놓이는 소아이디얼 $\mathfrak q \subset B$를 택하자.
다음은 서로 동치이다.
\begin{enumerate}
\item $A \to B$는 $\mathfrak q$에서 매끄럽다.
\item $A_\mathfrak p \to B_\mathfrak q$는 $\mathfrak q$-진 위상에서
형식적으로 매끄럽다.
\end{enumerate}
\end{lemma}

\begin{proof}
(2) $\Rightarrow$ (1)은 Algebra, Lemma
\ref{algebra-lemma-smooth-test-artinian}에서 따른다.
반대로 $A \to B$가 $\mathfrak q$에서 매끄러우면, 어떤
$g \in B$, $g \not \in \mathfrak q$에 대해 $A \to B_g$는 매끄럽다.
그러면 Algebra, Proposition
\ref{algebra-proposition-smooth-formally-smooth}에 의해
$A \to B_g$는 형식적으로 매끄럽다. 따라서 국소화는 형식적 매끄러움을
보존하므로 $A_\mathfrak p \to B_\mathfrak q$는 형식적으로 매끄럽다
(예를 들어 Algebra, Proposition
\ref{algebra-proposition-characterize-formally-smooth}의 판정법과
여접복합체가 국소화와 잘 양립한다는 사실을 사용하라.
Algebra, Lemmas \ref{algebra-lemma-NL-localize-bottom}와
\ref{algebra-lemma-localize-NL}을 보라).
마지막으로 보조정리 \ref{more-algebra-lemma-formally-smooth}는
$A_\mathfrak p \to B_\mathfrak q$가 $\mathfrak q$-진 위상에서
형식적으로 매끄러움을 함의한다.
\end{proof}



\section{멱급수환에 관한 몇 가지 결과}
\label{more-algebra-section-power-series}

\noindent
Noether 국소환 사이의 형식적으로 매끄러운 사상에 관한 문제는 흔히
멱급수환 사이의 사상에 관한 문제로 환원할 수 있다. 이 절에서는
이러한 논증을 돕는 몇 가지 보조정리를 증명한다.

\begin{lemma}
\label{more-algebra-lemma-power-series-ring-over-Cohen-fs}
$K$를 표수가 $0$인 체라 하고 $A = K[[x_1, \ldots, x_n]]$이라 하자.
$L$을 표수가 $p > 0$인 체라 하고 $B = L[[x_1, \ldots, x_n]]$이라 하자.
$\Lambda$를 Cohen 환이라 하자. $C = \Lambda[[x_1, \ldots, x_n]]$이라 하자.
\begin{enumerate}
\item $\mathbf{Q} \to A$는 $\mathfrak m_A$-진 위상에서
형식적으로 매끄럽다.
\item $\mathbf{F}_p \to B$는 $\mathfrak m_B$-진 위상에서
형식적으로 매끄럽다.
\item $\mathbf{Z} \to C$는 $\mathfrak m_C$-진 위상에서
형식적으로 매끄럽다.
\end{enumerate}
\end{lemma}

\begin{proof}
멱급수환의 보편 성질에 의해 다음을 증명하면 충분하다.
\begin{enumerate}
\item $\mathbf{Q} \to K$는 형식적으로 매끄럽다.
\item $\mathbf{F}_p \to L$은 형식적으로 매끄럽다.
\item $\mathbf{Z} \to \Lambda$는 $\mathfrak m_\Lambda$-진 위상에서
형식적으로 매끄럽다.
\end{enumerate}
앞의 두 명제는 Algebra, Proposition
\ref{algebra-proposition-characterize-separable-field-extensions}이다.
셋째 명제는 Algebra, Lemma
\ref{algebra-lemma-cohen-ring-formally-smooth}에서 따른다. 실제로 정의
\ref{more-algebra-definition-formally-smooth}에서와 같은 임의의 시험 도식에 대해
$p$의 어떤 거듭제곱이 $A/J$에서 영이므로 $p$의 어떤 거듭제곱은
$A$에서도 영이다.
\end{proof}

\begin{lemma}
\label{more-algebra-lemma-quotient-power-series-ring-over-Cohen}
$K$를 체라 하고 $A = K[[x_1, \ldots, x_n]]$이라 하자.
$\Lambda$를 Cohen 환이라 하고 $B = \Lambda[[x_1, \ldots, x_n]]$이라 하자.
\begin{enumerate}
% P02-E0875: frozen source has plural parameters with singular “is”; Korean uses natural agreement.
\item $y_1, \ldots, y_n \in A$가 정칙 매개변수계를 이루면
$K[[y_1, \ldots, y_n]] \to A$는 동형이다.
\item $z_1, \ldots, z_r \in A$가 $A$의 정칙 매개변수계의 일부를 이루면
$r \leq n$이고
$A/(z_1, \ldots, z_r) \cong K[[y_1, \ldots, y_{n - r}]]$이다.
% P02-E0876: frozen source has plural parameters with singular “is”; Korean uses natural agreement.
\item $p, y_1, \ldots, y_n \in B$가 정칙 매개변수계를 이루면
$\Lambda[[y_1, \ldots, y_n]] \to B$는 동형이다.
\item $p, z_1, \ldots, z_r \in B$가 $B$의 정칙 매개변수계의 일부를 이루면
$r \leq n$이고
$B/(z_1, \ldots, z_r) \cong \Lambda[[y_1, \ldots, y_{n - r}]]$이다.
\end{enumerate}
\end{lemma}

\begin{proof}
(1)의 증명. $A' = K[[y_1, \ldots, y_n]]$라 두자. 사상 $A' \to A$가
모든 $n \geq 1$에 대해 동형
$A'/\mathfrak m_{A'}^n \to A/\mathfrak m_A^n$을 유도함은 명백하다.
$A$와 $A'$가 모두 완비이므로 $A' \to A$가 동형임을 얻는다.
(2)의 증명. $z_1, \ldots, z_r$을 $A$의 정칙 매개변수계
$z_1, \ldots, z_r, y_1, \ldots, y_{n - r}$로 연장하자. 사상
$A' = K[[z_1, \ldots, z_r, y_1, \ldots, y_{n - r}]] \to A$를 생각하자.
(1)에 의해 이는 동형이다. 또한
$A'/(z_1, \ldots, z_r) \cong K[[y_1, \ldots, y_{n - r}]]$임이
명백하므로 (2)가 따른다.
(3)과 (4)의 증명은 각각 (1)과 (2)의 증명과 완전히 같다.
\end{proof}

\begin{lemma}
\label{more-algebra-lemma-embed-map-Noetherian-complete-local-rings}
$A \to B$를 Noether 완비 국소환의 국소 준동형이라 하자.
그러면 다음 성질을 갖는 가환 도식이 존재한다.
$$
\xymatrix{
S \ar[r] & B \\
R \ar[u] \ar[r] & A \ar[u]
}
$$
\begin{enumerate}
\item 수평 화살표들은 전사이다.
\item $A/\mathfrak m_A$의 표수가 영이면 $S$와 $R$은 체 위의
멱급수환이다.
\item $A/\mathfrak m_A$의 표수가 $p > 0$이면 $S$와 $R$은 Cohen 환 위의
멱급수환이다.
\item $R \to S$는 $R$의 한 정칙 매개변수계를 $S$의 한
정칙 매개변수계의 일부로 보낸다.
\end{enumerate}
특히 $R \to S$는 평탄하고(Algebra, Lemma
\ref{algebra-lemma-flat-over-regular}를 보라), 그 올
$S/\mathfrak m_R S$는 정칙이다(Algebra, Lemma
\ref{algebra-lemma-regular-ring-CM}을 보라).
\end{lemma}

\begin{proof}
Cohen 구조 정리(Algebra, Theorem
\ref{algebra-theorem-cohen-structure-theorem})를 사용하여 보조정리의
명제에서와 같은 전사 $S \to B$를 택하되, 잉여 표수가 $p > 0$이면
$S$를 Cohen 환 위의 멱급수환으로 택하고 그렇지 않으면 체 위의
멱급수환으로 택한다. $J \subset S$를 $S \to B$의 핵이라 하자.
다음으로 전사 $R = \Lambda[[x_1, \ldots, x_n]] \to A$를 택하되,
$A$의 잉여 표수가 $p > 0$이면 $\Lambda$를 Cohen 환으로 택하고,
그렇지 않으면 $\Lambda$를 $A$의 잉여체로 택한다.
합성 $\Lambda[[x_1, \ldots, x_n]] \to A \to B$를 사상
$\varphi : R \to S$로 올린다. 보조정리
\ref{more-algebra-lemma-power-series-ring-over-Cohen-fs}와 보조정리
\ref{more-algebra-lemma-lift-continuous}를 적용하면
$\Lambda[[x_1, \ldots, x_n]]$이 $\mathfrak m$-진 위상에서
$\mathbf{Z}$ 위 형식적으로 매끄럽기 때문에 이러한 올림이 가능하다.
마지막으로 $\varphi$를 사상
$\varphi' : R = \Lambda[[x_1, \ldots, x_n]] \to S' = S[[y_1, \ldots, y_n]]$로
바꾸되, $\varphi'|_\Lambda = \varphi|_\Lambda$이고
$\varphi'(x_i) = \varphi(x_i) + y_i$가 되게 한다. 또한 $S \to B$를
$y_i$를 영으로 보내는 사상 $S' \to B$로 바꾼다. 이렇게 바꾸면
$R$의 한 정칙 매개변수계가 $S'$의 한 정칙열의 일부로 감은 명백하고,
이것으로 증명이 끝난다.
% P02-E0877: the statement says “part of a regular system of parameters”; the frozen proof only establishes “part of a regular sequence”.
\end{proof}

\noindent
다음 보조정리에는 초등적인 증명이 있을 것이다.

\begin{lemma}
\label{more-algebra-lemma-dominate-two-surjections}
$S \to R$과 $S' \to R$을 완비 Noether 국소환의 전사라 하자.
그러면 $S \times_R S'$는 완비 Noether 국소환이다.
\end{lemma}

\begin{proof}
$k$를 $R$의 잉여체라 하자. $k$의 표수가 $p > 0$이면 잉여체가
$k$인 Cohen 환 $\Lambda$를 택한다(Algebra, Definition
\ref{algebra-definition-cohen-ring} 및 Algebra, Lemma
\ref{algebra-lemma-cohen-rings-exist}를 보라).
% P02-E0878: frozen source omits the preposition in “denote Lambda a Cohen ring”; Korean uses natural syntax.
$k$의 표수가 $0$이면 $\Lambda = k$라 둔다.
전사 $\Lambda[[x_1, \ldots, x_n]] \to R$을 택하고
(Cohen 구조 정리와 같은 방식이다. Algebra, Theorem
\ref{algebra-theorem-cohen-structure-theorem}을 보라), 보조정리
\ref{more-algebra-lemma-power-series-ring-over-Cohen-fs}와
\ref{more-algebra-lemma-lift-continuous}를 사용하여 이를 사상
$\Lambda[[x_1, \ldots, x_n]] \to S$ 및
$\varphi : \Lambda[[x_1, \ldots, x_n]] \to S$, 그리고
$\varphi' : \Lambda[[x_1, \ldots, x_n]] \to S'$로 올린다.
% P02-E0879: frozen source redundantly gives an unnamed map to S and the named map phi to S; all three formal spans are preserved.
다음으로 $S \to R$의 핵을 생성하는 $f_1, \ldots, f_m \in S$와
$S' \to R$의 핵을 생성하는 $f'_1, \ldots, f'_{m'} \in S'$를 택한다.
그러면 $x_i$를 $(\varphi(x_i), \varphi'(x_i))$로,
$y_j$를 $(f_j, 0)$로, $z_{j'}$를 $(0, f'_j)$로 보내는 사상
$$
\Lambda[[x_1, \ldots, x_n, y_1, \ldots, y_m, z_1, \ldots, z_{m'}]]
\longrightarrow S \times_R S,
$$
은 전사이다.
% P02-E0880: frozen codomain is S times_R S although the argument and conclusion concern S times_R S-prime; preserved exactly.
% P02-E0881: frozen source sends z with index j-prime to f-prime with index j; preserved exactly.
따라서 $S \times_R S'$는 완비 국소환의 몫이고, 그러므로 완비이다.
\end{proof}






\section{기하학적 정칙성과 형식적 매끄러움}
\label{more-algebra-section-regular-fs}

\noindent
이 절에서는 앞 절들의 결과를 결합하여 체 위의 기하학적으로 정칙인
국소환을 특징짓는다. 그런 다음 논증의 일부를 다시 사용하여
Noether 국소환 사이의 $\mathfrak m$-진 위상에서 형식적으로 매끄러운
사상을 특징짓는다.

\begin{theorem}
\label{more-algebra-theorem-regular-fs}
$k$를 체라 하자. $(A, \mathfrak m, K)$를 Noether 국소 $k$-대수라 하자.
$k$의 표수가 영이면 다음은 서로 동치이다.
\begin{enumerate}
\item $A$는 정칙 국소환이다.
\item $k \to A$는 $\mathfrak m$-진 위상에서 형식적으로 매끄럽다.
\end{enumerate}
$k$의 표수가 $p > 0$이면 다음은 서로 동치이다.
\begin{enumerate}
\item $A$는 $k$ 위에서 기하학적으로 정칙이다.
\item $k \to A$는 $\mathfrak m$-진 위상에서 형식적으로 매끄럽다.
\item $k$ 위에서 유한인 모든 $k \subset k' \subset k^{1/p}$에 대해
환 $A \otimes_k k'$은 정칙이다.
\item $A$는 정칙이고 표준 사상
$H_1(L_{K/k}) \to \mathfrak m/\mathfrak m^2$은 단사이다.
\item $A$는 정칙이고 사상
$\Omega_{k/\mathbf{F}_p} \otimes_k K \to \Omega_{A/\mathbf{F}_p} \otimes_A K$는
단사이다.
\end{enumerate}
\end{theorem}

\begin{proof}
$k$의 표수가 영이면 (1)과 (2)의 동치는 보조정리
\ref{more-algebra-lemma-fs-implies-regular}와 \ref{more-algebra-lemma-regular-implies-fs}에서 따른다.

\medskip\noindent
$k$의 표수가 $p > 0$이면 Proposition
\ref{more-algebra-proposition-characterization-geometrically-regular}에서
(1), (3), (4), (5)가 서로 동치임을 알 수 있다. (2)를 가정하자.
보조정리 \ref{more-algebra-lemma-base-change-fs}에 의해
$k' \to A' = A \otimes_k k'$은 $\mathfrak m' = \mathfrak mA'$-진 위상에서
형식적으로 매끄럽다. 따라서 $k \subset k'$가 유한 순수 비분리이면
보조정리 \ref{more-algebra-lemma-fs-implies-regular}에 의해 $A'$은 정칙 국소환이다.
그러므로 (1)이 성립한다.

\medskip\noindent
마지막으로 (5)가 (2)를 함의함을 증명하자. 정의
\ref{more-algebra-definition-formally-smooth}에서와 같은 실선 도식
$$
\xymatrix{
A \ar[r]_{\bar\psi} \ar@{-->}[rd] & B/J \\
k \ar[u]^i \ar[r]^\varphi & B \ar[u]_\pi
}
$$
을 택하자. $J^2 = 0$이므로 $J$에는 표준적인 $B/J$-가군 구조가 있고,
$\bar\psi$를 통해 $A$-가군 구조가 있다. $\bar\psi$가
$\mathfrak m$-진 위상에 대해 연속이므로 어떤 $n$에 대해
$\mathfrak m^nJ = 0$이다. 따라서 $J$를 $B/J$-부분가군으로 여과하여
$0 \subset J_1 \subset J_2 \subset \ldots \subset J_n = J$라 할 수 있고,
각 몫 $J_{t + 1}/J_t$는 $\mathfrak m$에 의해 소멸된다.
환 사상들의 열
$B \to B/J_1 \to B/J_2 \to \ldots \to B/J$를 생각하면,
$J$가 $\mathfrak m$에 의해 소멸되는 경우, 즉 $J$가 $K$-벡터 공간인
경우에 점선 화살표가 존재함을 증명하면 충분함을 알 수 있다.

\medskip\noindent
위와 같은 도식이 주어지고 $J$가 $\mathfrak m$에 의해 소멸된다고 하자.
보조정리 \ref{more-algebra-lemma-regular-implies-fs}에 의해 $\mathbf{F}_p \to A$는
$\mathfrak m$-진 위상에서 형식적으로 매끄럽다. 따라서
$\pi \circ \psi = \bar \psi$를 만족하는 환 사상
$\psi : A \to B$를 찾을 수 있다. 그러면
$\psi \circ i, \varphi : k \to B$는 $\pi$와 합성한 사상이 같은 두
사상이다. 따라서 $D = \psi \circ i - \varphi : k \to J$는 미분이다.
Algebra, Lemma \ref{algebra-lemma-universal-omega}에 의해 어떤 $k$-선형 사상
$\xi : \Omega_{k/\mathbf{F}_p} \to J$에 대해
$D = \xi \circ \text{d}$로 쓸 수 있다. $J$의 $K$-벡터 공간 구조를
사용하여 $\xi$를 $K$-선형 사상
$\xi' : \Omega_{k/\mathbf{F}_p} \otimes_k K \to J$로 연장한다.
(5)를 사용하면 $K$-선형 사상
$\xi'' : \Omega_{A/\mathbf{F}_p} \otimes_A K$를 찾을 수 있고, 이것의
$\Omega_{k/\mathbf{F}_p} \otimes_k K$로의 제한은 $\xi'$이다.
% P02-E1431: frozen source declares xi-double-prime without a codomain; the following display supplies its use as a map to J.
다음과 같이 쓰자.
$$
D' : A \xrightarrow{\text{d}} \Omega_{A/\mathbf{F}_p}
\to \Omega_{A/\mathbf{F}_p} \otimes_A K \xrightarrow{\xi''} J.
$$
마지막으로 $\psi' = \psi - D' : A \to B$라 두자. 독자는 $\psi'$가
$\pi \circ \psi' = \bar \psi$와 $\psi' \circ i = \varphi$를 만족하는
환 사상임을 확인할 수 있고, 이것이 원하는 바이다.
\end{proof}

\begin{example}
\label{more-algebra-example-fs}
$k$를 표수가 $p > 0$인 체라 하자. $a \in k$가 $p$제곱이 아닌
원소라고 하자. 잉여체가 $k$ 위에서 순수 비분리인 기하학적으로 정칙인
국소 $k$-대수의 표준적인 예는 환
$$
A = k[x, y]_{(x, y^p - a)}/(y^p - a - x)
$$
이다. 실제로 $A$는 $k$ 위 매끄러운 대수의 국소화이므로
$k \to A$는 형식적으로 매끄럽고, 따라서 $k \to A$는
$\mathfrak m$-진 위상에서 형식적으로 매끄럽다. 다음 예도 이와
밀접하게 관련된다. $k = \mathbf{F}_p(s)$ 및
$K = \mathbf{F}_p(t)^{perf}$라 하자. 환 사상
$$
k \longrightarrow A = K[[x]],\quad s \longmapsto t + x
$$
가 $A$의 $(x)$-진 위상에서 형식적으로 매끄러움을 주장한다.
실제로 $\Omega_{k/\mathbf{F}_p}$는 기저가 $\text{d}s$인 $1$차원이다.
이는 $\Omega_{A/\mathbf{F}_p}$의 원소
$\text{d}x + \text{d}t = \text{d}x$로 간다. 독자는
$\Omega_{A/\mathbf{F}_p}$가 $A$-가군으로서 $\text{d}x$ 위의
자유 가군임을 보일 수 있다. 그러므로 정리 \ref{more-algebra-theorem-regular-fs}의
조건 (5)가 성립하고, $k \to A$가 $(x)$-진 위상에서 형식적으로
매끄럽다는 결론을 얻는다.
\end{example}

\begin{lemma}
\label{more-algebra-lemma-formally-smooth-flat}
$A \to B$를 Noether 국소환의 국소 준동형이라 하자.
$A \to B$가 $\mathfrak m_B$-진 위상에서 형식적으로 매끄럽다고
가정하자. 그러면 $A \to B$는 평탄하다.
\end{lemma}

\begin{proof}
보조정리 \ref{more-algebra-lemma-formally-smooth-completion}과 Algebra, Lemma
\ref{algebra-lemma-completion-Noetherian-Noetherian}에 의해 $A$와 $B$가
Noether 완비 국소환이라고 가정해도 된다.
% P02-E0882: frozen source omits “are” in “A and B a Noetherian complete local rings”; Korean supplies the predicate.
(여기서는 완비화 위의 사상이 평탄하면 $A \to B$도 평탄함을 알기 위해
Algebra, Lemma \ref{algebra-lemma-flatness-descends-more-general}과
\ref{algebra-lemma-completion-faithfully-flat}도 사용한다.)
보조정리 \ref{more-algebra-lemma-embed-map-Noetherian-complete-local-rings}에서와 같이
$R \to S$가 평탄인 가환 도식
$$
\xymatrix{
S \ar[r] & B \\
R \ar[u] \ar[r] & A \ar[u]
}
$$
을 택하자. $I \subset R$을 $R \to A$의 핵이라 하자.
$B$가 $A$ 위에서 형식적으로 매끄러우므로 $A$-대수 사상
$$
S/IS \longrightarrow B
$$
에는 단면이 존재한다. 보조정리 \ref{more-algebra-lemma-lift-continuous}를 보라.
따라서 $B$는 평탄 $A$-가군 $S/IS$의 직합 인자이고
(평탄성의 밑변환에 대해서는 Algebra, Lemma
\ref{algebra-lemma-flat-base-change}를 보라), 그러므로 평탄하다.
\end{proof}

\begin{lemma}
\label{more-algebra-lemma-formally-smooth-JZ}
$A \to B$를 Noether 국소환의 국소 준동형이라 하자.
$A \to B$가 $\mathfrak m_B$-진 위상에서 형식적으로 매끄럽다고
가정하자. $K$를 $B$의 잉여체라 하자. 그러면 $A \to B \to K$에 대한
Jacobi--Zariski 열은 다음 완전열을 준다.
$$
0 \to H_1(\NL_{K/A}) \to \mathfrak m_B/\mathfrak m_B^2
\to \Omega_{B/A} \otimes_B K \to \Omega_{K/A} \to 0
$$
\end{lemma}

\begin{proof}
Algebra, Lemma \ref{algebra-lemma-NL-surjection}에 의해
$\mathfrak m_B/\mathfrak m_B^2 = H_1(\NL_{K/B})$임에 주의하자.
Algebra, Lemma \ref{algebra-lemma-exact-sequence-NL}에 의해
$H_1(\NL_{K/A}) \to \mathfrak m_B/\mathfrak m_B^2$의 단사성을
보이면 충분하다. $k$를 $A$의 잉여체라 하면, $A \to k \to K$에 대한
Jacobi--Zariski 열은 $\Omega_{K/A} = \Omega_{K/k}$와 완전열
$$
\mathfrak m_A/\mathfrak m_A^2 \otimes_k K \to
H_1(\NL_{K/A}) \to
H_1(\NL_{K/k}) \to 0
$$
을 준다. $\overline{B} = B \otimes_A k$라 두자. $\overline{B}$가
정칙이므로 아이디얼 $\mathfrak m_{\overline{B}}$는 정칙열에 의해
생성된다. $\mathfrak m_A B \subset \mathfrak m_B$에 보조정리
\ref{more-algebra-lemma-conormal-sequence-H1-regular}와
\ref{more-algebra-lemma-noetherian-finite-all-equivalent}를 적용하면
$\mathfrak m_A B / (\mathfrak m_AB \cap \mathfrak m_B^2) =
\mathfrak m_A B / \mathfrak m_A \mathfrak m_B$를 얻는다.
보조정리 \ref{more-algebra-lemma-formally-smooth-flat}에 의해 $A \to B$가 평탄하므로
이는 $\mathfrak m_A/\mathfrak m_A^2 \otimes_k K$와 같다.
따라서 짧은 완전열
$$
0 \to
\mathfrak m_A/\mathfrak m_A^2 \otimes_k K \to
\mathfrak m_B/\mathfrak m_B^2 \to
\mathfrak m_{\overline{B}}/\mathfrak m_{\overline{B}}^2 \to 0
$$
을 얻는다. Jacobi--Zariski 열의 함자성에 의해 다음 가환 도식을 얻는다.
$$
\xymatrix{
&
\mathfrak m_A/\mathfrak m_A^2 \otimes_k K \ar[d] \ar[r] &
H_1(\NL_{K/A}) \ar[d] \ar[r] &
H_1(\NL_{K/k}) \ar[d] \ar[r] & 0 \\
0 \ar[r] &
\mathfrak m_A/\mathfrak m_A^2 \otimes_k K \ar[r] &
\mathfrak m_B/\mathfrak m_B^2 \ar[r] &
\mathfrak m_{\overline{B}}/\mathfrak m_{\overline{B}}^2 \ar[r] & 0
}
$$
왼쪽 수직 화살표는 정리 \ref{more-algebra-theorem-regular-fs}에 의해 단사이다.
실제로 보조정리 \ref{more-algebra-lemma-base-change-fs}에 의해 $k \to \overline{B}$는
$\mathfrak m_{\overline{B}}$-진 위상에서 형식적으로 매끄럽다.
뱀 보조정리를 적용하면 증명이 끝난다.
\end{proof}

\begin{proposition}
\label{more-algebra-proposition-fs-flat-fibre-fs}
$A \to B$를 Noether 국소환의 국소 준동형이라 하자. $k$를 $A$의
잉여체라 하고 $\overline{B} = B \otimes_A k$를 특수 올이라 하자.
다음은 서로 동치이다.
\begin{enumerate}
\item $A \to B$는 평탄이고 $\overline{B}$는 $k$ 위에서
기하학적으로 정칙이다.
\item $A \to B$는 평탄이고 $k \to \overline{B}$는
$\mathfrak m_{\overline{B}}$-진 위상에서 형식적으로 매끄럽다.
\item $A \to B$는 $\mathfrak m_B$-진 위상에서 형식적으로 매끄럽다.
\end{enumerate}
\end{proposition}

\begin{proof}
(1)과 (2)의 동치는 정리 \ref{more-algebra-theorem-regular-fs}에서 따른다.

\medskip\noindent
(3)을 가정하자. 보조정리 \ref{more-algebra-lemma-formally-smooth-flat}에 의해
$A \to B$는 평탄하다. 보조정리 \ref{more-algebra-lemma-base-change-fs}에 의해
$k \to \overline{B}$는 $\mathfrak m_{\overline{B}}$-진 위상에서
형식적으로 매끄럽다. 따라서 (2)가 성립한다.

\medskip\noindent
(2)를 가정하자. 보조정리 \ref{more-algebra-lemma-formally-smooth-completion}에 의해
형식적 매끄러움은 완비화로 보존된다. 평탄성도 Algebra, Lemma
\ref{algebra-lemma-completion-faithfully-flat}에 의해 마찬가지이다.
따라서 $A$와 $B$를 각각의 완비화로 바꾸어 $A$와 $B$가 Noether 완비
국소환이라고 가정해도 된다. 이 경우 보조정리
\ref{more-algebra-lemma-embed-map-Noetherian-complete-local-rings}에서와 같은 도식
$$
\xymatrix{
S \ar[r] & B \\
R \ar[u] \ar[r] & A \ar[u]
}
$$
을 택하자. 앞으로 이 도식의 모든 성질을 별도의 언급 없이 사용한다.
$R$의 정칙 매개변수계 $t_1, \ldots, t_d$를 택하되, $k$의 표수가
$p > 0$인 경우에는 $t_1 = p$가 되게 한다.
$\overline{S} = S \otimes_R k$라 두자. 짧은 완전열
$$
0 \to J \to S \to B \to 0
$$
을 생각하자. $\overline{B}$와 $\overline{S}$가 정칙이므로
$\overline{S} \to \overline{B}$의 핵은 $\overline{S}$의 한 정칙
매개변수계의 일부를 이루는 원소
$\overline{x}_1, \ldots, \overline{x}_r$로 생성된다. Algebra, Lemma
\ref{algebra-lemma-regular-quotient-regular}를 보라. 이 원소들을
$x_1, \ldots, x_r \in J$로 올리자. 그러면
$t_1, \ldots, t_d, x_1, \ldots, x_r$은 $S$의 한 정칙 매개변수계의
일부이다. 따라서 $S/(x_1, \ldots, x_r)$은 $k$의 표수가 영이면
체 위의 멱급수환이고, $k$의 표수가 $p > 0$이면 Cohen 환 위의
멱급수환이다. 보조정리
\ref{more-algebra-lemma-quotient-power-series-ring-over-Cohen}을 보라.
더욱이 $R \to S/(x_1, \ldots, x_r)$은 여전히
$t_1, \ldots, t_d$를 $S/(x_1, \ldots, x_r)$의 한 정칙 매개변수계의
일부로 보낸다. 다시 말해 $S$를 $S/(x_1, \ldots, x_r)$로 바꾸어,
보조정리 \ref{more-algebra-lemma-embed-map-Noetherian-complete-local-rings}에서와 같은
도식
$$
\xymatrix{
S \ar[r] & B \\
R \ar[u] \ar[r] & A \ar[u]
}
$$
이 있고 더 나아가 $\overline{S} = \overline{B}$라고 가정해도 된다.
이 경우 사상
$$
S \otimes_R A \longrightarrow B
$$
은 동형이다. 실제로 이는 전사이고 특수 올 위에서 동형이며, 원천과
과녁은 모두 $A$ 위에서 평탄하다. 예를 들어 Algebra, Lemma
\ref{algebra-lemma-mod-injective}를 사용하거나, 짧은 완전열
$0 \to I \to S \otimes_R A \to B \to 0$을 $A$ 위에서 $k$와
텐서하면 $I \otimes_A k = 0$을 얻고 따라서 Nakayama 보조정리에 의해
$I = 0$임을 사용하라. 그러므로 보조정리 \ref{more-algebra-lemma-base-change-fs}에
의해 $R \to S$가 $\mathfrak m_S$-진 위상에서 형식적으로 매끄러움을
보이면 충분하다. 물론 $\overline{S} = \overline{B}$이므로
$\overline{S}$는 $k = R/\mathfrak m_R$ 위에서 형식적으로 매끄럽다.

\medskip\noindent
$t_1, \ldots, t_d, y_1, \ldots, y_m$이 $S$의 정칙 매개변수계가 되도록
원소 $y_1, \ldots, y_m \in S$를 택하자. $k$의 표수가 영이면
계수체 $K \subset S$를 택하고, $k$의 표수가 $p > 0$이면 잉여체가
$K$인 Cohen 환 $\Lambda \subset S$를 택한다. 이제 사상
$K[[t_1, \ldots, t_d, y_1, \ldots, y_m]] \to S$
(표수 영인 경우) 또는
$\Lambda[[t_2, \ldots, t_d, y_1, \ldots, y_m]] \to S$
(표수가 $p > 0$인 경우)는 동형이다. 보조정리
\ref{more-algebra-lemma-quotient-power-series-ring-over-Cohen}을 보라.
이제부터 $S$를 위 멱급수환으로 생각한다.

\medskip\noindent
나머지 증명은 정리 \ref{more-algebra-theorem-regular-fs}의 증명에서 사용한 논증과
비슷하다. 정의 \ref{more-algebra-definition-formally-smooth}에서와 같은 실선 도식
$$
\xymatrix{
S \ar[r]_{\bar\psi} \ar@{-->}[rd] & N/J \\
R \ar[u]^i \ar[r]^\varphi & N \ar[u]_\pi
}
$$
을 택하자. $J^2 = 0$이므로 $J$에는 표준적인 $N/J$-가군 구조가 있고,
$\bar\psi$를 통해 $S$-가군 구조가 있다. $\bar\psi$가
$\mathfrak m_S$-진 위상에 대해 연속이므로 어떤 $n$에 대해
$\mathfrak m_S^nJ = 0$이다. 따라서 $J$를 $N/J$-부분가군으로 여과하여
$0 \subset J_1 \subset J_2 \subset \ldots \subset J_n = J$라 할 수 있고,
각 몫 $J_{t + 1}/J_t$는 $\mathfrak m_S$에 의해 소멸된다.
환 사상들의 열 $N \to N/J_1 \to N/J_2 \to \ldots \to N/J$를
생각하면, $J$가 $\mathfrak m_S$에 의해 소멸되는 경우, 즉 $J$가
$K$-벡터 공간인 경우에 점선 화살표가 존재함을 증명하면 충분하다.

\medskip\noindent
위와 같은 도식이 주어지고 $J$가 $\mathfrak m_S$에 의해 소멸된다고
하자. $\mathbf{Q} \to S$(표수 영인 경우) 또는
$\mathbf{Z} \to S$(표수가 $p > 0$인 경우)는
$\mathfrak m_S$-진 위상에서 형식적으로 매끄러우므로(보조정리
\ref{more-algebra-lemma-power-series-ring-over-Cohen-fs}를 보라),
$\pi \circ \psi = \bar \psi$를 만족하는 환 사상
$\psi : S \to N$을 찾을 수 있다. $S$는 어떤 부분환 위에서
$t_1, \ldots, t_d$(표수 영인 경우) 또는
$t_2, \ldots, t_d$(표수가 $p > 0$인 경우)에 관한 멱급수환이므로,
멱급수환의 보편 성질에 의해 $\psi(t_i)$가 $\varphi(t_i)$와 같도록
$\psi$의 선택을 바꿀 수 있다. 표수가 $p$인 경우 $t_1 = p$에 대해서는
자동으로 성립한다. 그러면 $\psi \circ i$와 $\varphi : R \to N$은
$\pi$와 합성한 사상이 같고 $t_1, \ldots, t_d$ 위에서 일치하는 두
사상이다. 따라서 $D = \psi \circ i - \varphi : R \to J$는
$t_1, \ldots, t_d$를 소멸시키는 미분이다.
Algebra, Lemma \ref{algebra-lemma-universal-omega}에 의해
$D = \xi \circ \text{d}$로 쓸 수 있다. 여기서
$\xi : \Omega_{R/\mathbf{Z}} \to J$는 $R$-선형 사상이며,
구성에 의해 $\text{d}t_1, \ldots, \text{d}t_d$를 소멸시키고,
$J$가 $\mathfrak m_R$에 의해 소멸되므로
$\mathfrak m_R \Omega_{R/\mathbf{Z}}$도 소멸시킨다. 따라서 $\xi$는 합성
$$
\Omega_{R/\mathbf{Z}} \to \Omega_{k/\mathbf{Z}} \xrightarrow{\xi'} J
$$
을 통해 인수분해되며, 여기서 $\xi'$는 $k$-선형이다. $J$의
$K$-벡터 공간 구조를 사용하여 $\xi'$를 $K$-선형 사상
$$
\xi'' : \Omega_{k/\mathbf{Z}} \otimes_k K \longrightarrow J.
$$
로 연장한다. $\overline{S}/k$가 형식적으로 매끄럽다는 사실과 정리
\ref{more-algebra-theorem-regular-fs}에 의해
$$
\Omega_{k/\mathbf{Z}} \otimes_k K \to
\Omega_{\overline{S}/\mathbf{Z}} \otimes_S K
$$
가 단사임을 알 수 있다. 표수 영인 경우에는
$\Omega_{k/\mathbf{Z}} \to \Omega_{K/\mathbf{Z}}$도 단사이므로 이 명제는
그 경우에도 성립한다. Algebra, Proposition
\ref{algebra-proposition-characterize-separable-field-extensions}를 보라.
따라서 $K$-선형 사상
$\xi''' : \Omega_{\overline{S}/\mathbf{Z}} \otimes_S K \to J$를 찾을 수 있고,
이것의 $\Omega_{k/\mathbf{Z}} \otimes_k K$로의 제한은 $\xi''$이다.
다음과 같이 쓰자.
$$
D' : S \xrightarrow{\text{d}} \Omega_{S/\mathbf{Z}}
\to \Omega_{\overline{S}/\mathbf{Z}} \to
\Omega_{\overline{S}/\mathbf{Z}} \otimes_S K \xrightarrow{\xi'''} J.
$$
마지막으로 $\psi' = \psi - D' : S \to N$이라 두자. 독자는
$\psi'$가 $\pi \circ \psi' = \bar \psi$와
$\psi' \circ i = \varphi$를 만족하는 환 사상임을 확인할 수 있고,
이것이 원하는 바이다.
\end{proof}

\noindent
위 결과의 응용으로 형식적으로 매끄러운 대수의 변형에는 장애가
없음을 증명한다.

\begin{lemma}
\label{more-algebra-lemma-lift-fs}
$A$를 잉여체가 $k$인 Noether 완비 국소환이라 하자. $B$를 Noether
완비 국소 $k$-대수라 하자. $k \to B$가 $\mathfrak m_B$-진 위상에서
형식적으로 매끄럽다고 가정하자. 그러면 Noether 완비 국소환 $C$와
$\mathfrak m_C$-진 위상에서 형식적으로 매끄러운 국소 준동형
$A \to C$가 존재하여 $C \otimes_A k \cong B$이다.
\end{lemma}

\begin{proof}
보조정리 \ref{more-algebra-lemma-embed-map-Noetherian-complete-local-rings}에서와 같은
도식
$$
\xymatrix{
S \ar[r] & B \\
R \ar[u] \ar[r] & A \ar[u]
}
$$
을 택하자. $R$의 정칙 매개변수계를 $t_1, \ldots, t_d$라 하되,
$k$의 표수가 $p > 0$인 경우에는 $t_1 = p$라 하자. $B$와
$\overline{S} = S \otimes_R k$가 정칙이므로
$\Ker(\overline{S} \to B)$는 $\overline{S}$의 한 정칙 매개변수계의
일부를 이루는 원소 $\overline{x}_1, \ldots, \overline{x}_r$로
생성된다. Algebra, Lemma
\ref{algebra-lemma-regular-quotient-regular}를 보라. 이 원소들을
$x_1, \ldots, x_r \in S$로 올리자. 그러면
$t_1, \ldots, t_d, x_1, \ldots, x_r$은 $S$의 한 정칙 매개변수계의
일부이다. 따라서 $S/(x_1, \ldots, x_r)$은 $k$의 표수가 영이면
체 위의 멱급수환이고, $k$의 표수가 $p > 0$이면 Cohen 환 위의
멱급수환이다. 보조정리
\ref{more-algebra-lemma-quotient-power-series-ring-over-Cohen}을 보라.
더욱이 $R \to S/(x_1, \ldots, x_r)$은 여전히
$t_1, \ldots, t_d$를 $S/(x_1, \ldots, x_r)$의 한 정칙 매개변수계의
일부로 보낸다. 다시 말해 $S$를 $S/(x_1, \ldots, x_r)$로 바꾸어,
보조정리 \ref{more-algebra-lemma-embed-map-Noetherian-complete-local-rings}에서와 같은
도식
$$
\xymatrix{
S \ar[r] & B \\
R \ar[u] \ar[r] & A \ar[u]
}
$$
이 있고 더 나아가 $\overline{S} = B$라고 가정해도 된다. 이 경우
Proposition \ref{more-algebra-proposition-fs-flat-fibre-fs}에 의해 $R \to S$는
$\mathfrak m_S$-진 위상에서 형식적으로 매끄럽다. 따라서 밑변환
$C = S \otimes_R A$는 보조정리 \ref{more-algebra-lemma-base-change-fs}에 의해
$\mathfrak m_C$-진 위상에서 $A$ 위 형식적으로 매끄럽다.
\end{proof}

\begin{remark}
\label{more-algebra-remark-what-does-it-mean}
보조정리 \ref{more-algebra-lemma-lift-fs}의 주장은 상당히 강하다. 실제로 잉여체의
동형 $k = A/\mathfrak m_A = A'/\mathfrak m_{A'}$을 유도하는 사상
$A \to A'$을 갖고 $B \otimes_{A'} k$가 $k$ 위에서 형식적으로
매끄러운, Noether 완비 국소환의 국소 준동형 도식
$$
\xymatrix{
& B \\
A \ar[r] & A' \ar[u]
}
$$
이 주어졌다고 하자. 그러면 이를 Noether 완비 국소환의 국소 준동형의
가환 도식
$$
\xymatrix{
C \ar[r] & B \\
A \ar[r] \ar[u] & A' \ar[u]
}
$$
으로 연장할 수 있다. 여기서 $A \to C$는 $\mathfrak m_C$-진 위상에서
형식적으로 매끄럽고 $C \otimes_A k \cong B \otimes_{A'} k$이다.
실제로 보조정리 \ref{more-algebra-lemma-lift-fs}에서와 같이
$B \otimes_{A'} k$를 $k$ 위에서 올리는 $A \to C$를 택한다.
형식적 매끄러움에 의해 화살표 $C \to B$를 찾을 수 있다. 보조정리
\ref{more-algebra-lemma-lift-continuous}를 보라.
$C \otimes_A^\wedge A'$를 $C \otimes_A \mathfrak m_{A'}$에 관한
$C \otimes_A A'$의 완비화라 하자.
% P02-E1432: frozen source omits “by” in its notation-introduction sentence; Korean uses natural syntax.
$C \otimes_A^\wedge A'$는 Noether 완비 국소환이고(Algebra, Lemma
\ref{algebra-lemma-completion-Noetherian}을 보라), $A'$ 위에서
평탄이다(Algebra, Lemma \ref{algebra-lemma-flat-module-powers}를 보라).
더 나아가 다음이 성립한다.
\begin{enumerate}
\item $C \otimes_A^\wedge A' \to B$는 전사이다.
\item $A \to A'$이 전사이면 $C \to B$는 전사이다.
\item $A \to A'$이 유한이면 $C \to B$는 유한이다.
\item $A' \to B$가 평탄이면 $C \otimes_A^\wedge A' \cong B$이다.
\end{enumerate}
실제로 멱영 아이디얼에 대한 Nakayama 보조정리(Algebra, Lemma
\ref{algebra-lemma-NAK}를 보라)에 의해
$C \otimes_A k \cong B \otimes_{A'} k$이면 모든 $n$에 대해
$C \otimes_A A'/\mathfrak m_{A'}^n \to B/\mathfrak m_{A'}^nB$가
전사임을 알 수 있다. 이로써 (1)이 증명된다. (2)와 (3)은 (1)에서
따른다. (4)는 Algebra, Lemma
\ref{algebra-lemma-mod-injective}에서 따른다.
\end{remark}




\section{정칙 환 사상}
\label{more-algebra-section-regular}

\noindent
$k$를 체라 하자. Noether $k$-대수 $A$가 $k$ 위에서
{\it 기하학적으로 정칙}이라는 것은 $k$의 모든 유한 순수 비분리 확대
$k'$에 대해 $A \otimes_k k'$이 정칙이라는 뜻임을 상기하자.
Algebra, Definition \ref{algebra-definition-geometrically-regular}를 보라.
더욱이 이 조건이 성립하면 모든 유한 생성 체 확대 $k'/k$에 대해
$A \otimes_k k'$은 정칙이다. Algebra, Lemma
\ref{algebra-lemma-geometrically-regular}를 보라. 다음 정의에서 이
개념을 사용한다.

\begin{definition}
\label{more-algebra-definition-regular}
환 사상 $R \to \Lambda$가 평탄이고, 모든 소아이디얼
$\mathfrak p \subset R$에 대한 올환
$$
\Lambda \otimes_R \kappa(\mathfrak p) =
\Lambda_\mathfrak p/\mathfrak p\Lambda_\mathfrak p
$$
이 Noether이며 $\kappa(\mathfrak p)$ 위에서 기하학적으로 정칙이면,
그 사상을 {\it 정칙}이라 한다.
\end{definition}

\noindent
$R \to \Lambda$가 환 사상이고 $\Lambda$가 Noether이면 올환들은
항상 Noether이다.

\begin{lemma}[정칙성은 국소적 성질이다]
\label{more-algebra-lemma-regular-local}
$R \to \Lambda$를 $\Lambda$가 Noether인 환 사상이라 하자.
다음은 서로 동치이다.
\begin{enumerate}
\item $R \to \Lambda$는 정칙이다.
\item $\mathfrak p \subset R$ 위에 놓이는 모든
$\mathfrak q \subset \Lambda$에 대해
$R_\mathfrak p \to \Lambda_\mathfrak q$는 정칙이다.
\item $R$의 $\mathfrak m$ 위에 놓이는 모든 극대 아이디얼
$\mathfrak m' \subset \Lambda$에 대해
$R_\mathfrak m \to \Lambda_{\mathfrak m'}$은 정칙이다.
\end{enumerate}
\end{lemma}

\begin{proof}
이는 Noether 환이 정칙이라는 것과 그 모든 국소환이 정칙 국소환이라는
것이 동치이고(Algebra, Definition \ref{algebra-definition-regular}을
보라), 환 사상이 평탄이라는 것과 국소환 사이에 유도된 모든 사상이
평탄이라는 것이 동치이기 때문이다(Algebra, Lemma
\ref{algebra-lemma-flat-localization}을 보라).
\end{proof}

\begin{lemma}[정칙 사상과 밑변환]
\label{more-algebra-lemma-regular-base-change}
$R \to \Lambda$를 정칙 환 사상이라 하자. 임의의 유한형 환 사상
$R \to R'$에 대해 밑변환 $R' \to \Lambda \otimes_R R'$도 정칙이다.
\end{lemma}

\begin{proof}
평탄성은 임의의 밑변환으로 보존된다. Algebra, Lemma
\ref{algebra-lemma-flat-base-change}를 보라.
$\mathfrak p \subset R$ 위에 놓이는 소아이디얼
$\mathfrak p' \subset R'$을 생각하자. $R'$이 $R$ 위 유한형이므로
잉여체 확대 $\kappa(\mathfrak p')/\kappa(\mathfrak p)$는 유한 생성이다.
따라서 올환
$$
(\Lambda \otimes_R R') \otimes_{R'} \kappa(\mathfrak p') =
\Lambda \otimes_R \kappa(\mathfrak p) \otimes_{\kappa(\mathfrak p)} 
\kappa(\mathfrak p')
$$
은 Algebra, Lemma \ref{algebra-lemma-Noetherian-field-extension}와
$R \to \Lambda$의 올환에 대한 가정에 의해 Noether이다.
올들의 기하학적 정칙성은 Algebra, Lemma
\ref{algebra-lemma-geometrically-regular}에 의해 보존된다.
\end{proof}

\begin{lemma}[정칙 사상의 합성]
\label{more-algebra-lemma-regular-composition}
$A \to B$와 $B \to C$를 정칙 환 사상이라 하자.
$A \to C$의 올환들이 Noether이면 $A \to C$는 정칙이다.
\end{lemma}

\begin{proof}
$\mathfrak p \subset A$를 소아이디얼이라 하자.
$\kappa(\mathfrak p) \subset k$를 유한 순수 비분리 확대라 하자.
$C \otimes_A k$가 정칙임을 보여야 한다. 보조정리
\ref{more-algebra-lemma-regular-base-change}에 의해 $A = k$라고 가정해도 되며,
$C$가 정칙임을 보이는 문제로 환원된다. 가정에 의해 $B$는 정칙이고
$B \to C$는 정칙 올을 갖는 평탄 사상이다. 따라서 Algebra, Lemma
\ref{algebra-lemma-flat-over-regular-with-regular-fibre}에 의해 $C$는
정칙이다. 몇 가지 세부 사항은 생략한다.
\end{proof}

\begin{lemma}
\label{more-algebra-lemma-colimit-smooth-regular}
$R$을 환이라 하자. $(A_i, \varphi_{ii'})$를 매끄러운 $R$-대수의
유향계라 하자. $\Lambda = \colim A_i$라 두자. 모든
$\mathfrak p \subset R$에 대해 올환
$\Lambda \otimes_R \kappa(\mathfrak p)$가 Noether이면
$R \to \Lambda$는 정칙이다.
\end{lemma}

\begin{proof}
Algebra, Lemmas \ref{algebra-lemma-colimit-flat}와
\ref{algebra-lemma-smooth-syntomic}에 의해 $\Lambda$는 $R$ 위에서
평탄임에 주의하자. $\kappa(\mathfrak p) \subset k$를 유한 순수
비분리 확대라 하자. 등식
$$
\Lambda \otimes_R \kappa(\mathfrak p) \otimes_{\kappa(\mathfrak p)} k =
\Lambda \otimes_R k = \colim A_i \otimes_R k
$$
에 주의하자. 이는 매끄러운 $k$-대수의 여극한이다. Algebra, Lemma
\ref{algebra-lemma-base-change-smooth}를 보라. 매끄러운 $k$-대수의
각 국소환은 Algebra, Lemma
\ref{algebra-lemma-characterize-smooth-over-field}에 의해 정칙이므로,
Algebra, Lemma \ref{algebra-lemma-colimit-regular}에 의해
$\Lambda \otimes_R k$의 모든 국소환이 정칙이라는 결론을 얻는다.
이로써 보조정리가 증명된다.
\end{proof}

\noindent
이제 체 확대가 언제 정칙 환 사상을 정의하는지 살펴보자.

\begin{lemma}
\label{more-algebra-lemma-regular-field-extension}
$K/k$를 체 확대라 하자. 그러면 $k \to K$가 정칙 환 사상인 것과
$K$가 $k$의 분리 체 확대인 것은 동치이다.
\end{lemma}

\begin{proof}
$k \to K$가 정칙이면 $K$는 $k$ 위에서 기하학적으로 기약원이고,
따라서 Algebra, Proposition
\ref{algebra-proposition-characterize-separable-field-extensions}에 의해
$K$는 $k$ 위에서 분리이다. 반대로 $K/k$가 분리이면 Algebra, Lemma
\ref{algebra-lemma-colimit-syntomic}에 의해 $K$는 매끄러운
$k$-대수의 여극한이고, 따라서 보조정리
\ref{more-algebra-lemma-colimit-smooth-regular}에 의해 정칙이다.
\end{proof}

\begin{lemma}
\label{more-algebra-lemma-regular-permanence}
$A \to B \to C$를 환 사상이라 하자. $A \to C$가 정칙이고
$B \to C$가 평탄이며 스펙트럼 위에서 전사이면 $A \to B$는 정칙이다.
\end{lemma}

\begin{proof}
Algebra, Lemma \ref{algebra-lemma-flat-permanence}에 의해 $A \to B$는
평탄하다. $\mathfrak p \subset A$를 소아이디얼이라 하자. 환 사상
$B \otimes_A \kappa(\mathfrak p) \to C \otimes_A \kappa(\mathfrak p)$는
평탄이고 스펙트럼 위에서 전사이다. 따라서 Algebra, Lemma
\ref{algebra-lemma-geometrically-regular-descent}에 의해
$B \otimes_A \kappa(\mathfrak p)$는 기하학적으로 정칙이다.
\end{proof}




\section{정칙 환 사상을 따른 성질의 상승}
\label{more-algebra-section-ascending-properties}

\noindent
이 절은 환 사상 $R \to S$가 정칙인 경우의 Algebra, Section
\ref{algebra-section-ascending-properties}에 해당한다.

\begin{lemma}
\label{more-algebra-lemma-reduced-goes-up}
$\varphi : R \to S$를 환 사상이라 하자. 다음을 가정하자.
\begin{enumerate}
\item $\varphi$는 정칙이다.
\item $S$는 Noether이다.
\item $R$은 Noether이고 기약원이다.
\end{enumerate}
그러면 $S$는 기약원이다.
\end{lemma}

\begin{proof}
Noether 환이 기약원이라는 것은 성질 $(S_1)$과 $(R_0)$을 갖는 것과
같다. Algebra, Lemma \ref{algebra-lemma-criterion-reduced}를 보라.
따라서 Algebra, Lemmas \ref{algebra-lemma-Sk-goes-up}과
\ref{algebra-lemma-Rk-goes-up}을 적용할 수 있다.
\end{proof}

\begin{lemma}
\label{more-algebra-lemma-normal-goes-up}
$\varphi : R \to S$를 환 사상이라 하자. 다음을 가정하자.
\begin{enumerate}
\item $\varphi$는 정칙이다.
\item $S$는 Noether이다.
\item $R$은 Noether이고 정규이다.
\end{enumerate}
그러면 $S$는 정규이다.
\end{lemma}

\begin{proof}
Noether 환이 정규라는 것은 성질 $(S_2)$와 $(R_1)$을 갖는 것과 같다.
Algebra, Lemma \ref{algebra-lemma-criterion-normal}을 보라. 따라서
Algebra, Lemmas \ref{algebra-lemma-Sk-goes-up}과
\ref{algebra-lemma-Rk-goes-up}을 적용할 수 있다.
\end{proof}

\begin{lemma}
\label{more-algebra-lemma-regular-goes-up}
$\varphi : R \to S$를 환 사상이라 하자. 다음을 가정하자.
\begin{enumerate}
\item $\varphi$는 정칙이다.
\item $S$는 Noether이다.
\item $R$은 Noether이고 정칙이다.
\end{enumerate}
그러면 $S$는 정칙이다.
\end{lemma}

\begin{proof}
Noether 환이 정칙이라는 것은 모든 $k$에 대해 성질 $(R_k)$를 갖는 것과
같다. 따라서 Algebra, Lemma \ref{algebra-lemma-Rk-goes-up}을 적용할 수 있다.
\end{proof}

\begin{lemma}
\label{more-algebra-lemma-CM-goes-up}
$\varphi : R \to S$를 환 사상이라 하자. 다음을 가정하자.
\begin{enumerate}
\item $\varphi$는 정칙이다.
\item $S$는 Noether이다.
\item $R$은 Noether이고 Cohen--Macaulay이다.
\end{enumerate}
그러면 $S$는 Cohen--Macaulay이다.
\end{lemma}

\begin{proof}
Noether 환이 Cohen--Macaulay라는 것은 모든 $k$에 대해 성질 $(S_k)$를
갖는 것과 같다. 따라서 Algebra, Lemma
\ref{algebra-lemma-Sk-goes-up}을 적용할 수 있다.
\end{proof}




\section{완비화에 의한 성질의 보존}
\label{more-algebra-section-permanence-completion}

\noindent
Noether 국소환 $(A, \mathfrak m)$이 주어졌을 때, $\mathfrak m$에 관한
$A$의 완비화를 $A^\wedge$로 나타낸다.
% P02-E1433: frozen source omits “by” in “we denote A-wedge the completion”; Korean uses natural notation syntax.
$A^\wedge$가 극대 아이디얼 $\mathfrak m^\wedge = \mathfrak m A^\wedge$을
갖는 Noether 완비 국소환이고 $A \to A^\wedge$가 충실히 평탄이라는
사실을 별도의 언급 없이 사용한다. Algebra, Lemmas
\ref{algebra-lemma-completion-Noetherian-Noetherian},
\ref{algebra-lemma-completion-complete}, 및
\ref{algebra-lemma-completion-faithfully-flat}을 보라.

\begin{lemma}
\label{more-algebra-lemma-completion-dimension}
$A$를 Noether 국소환이라 하자. 그러면 $\dim(A) = \dim(A^\wedge)$이다.
\end{lemma}

\begin{proof}
Algebra, Lemma \ref{algebra-lemma-completion-complete}에 의해 사상
$A \to A^\wedge$는 모든 $n \geq 1$에 대해 동형
$A/\mathfrak m^n = A^\wedge/(\mathfrak m^\wedge)^n$을 유도한다.
Algebra, Lemma \ref{algebra-lemma-pushdown-module}에 의해 이는 모든
$n \geq 1$에 대해
$$
\text{length}_A(A/\mathfrak m^n) =
\text{length}_{A^\wedge}(A^\wedge/(\mathfrak m^\wedge)^n)
$$
임을 함의한다. 따라서 $d(A) = d(A^\wedge)$이고, Algebra, Proposition
\ref{algebra-proposition-dimension}에 의해 결론을 얻는다. 다른 증명으로
Algebra, Lemma \ref{algebra-lemma-dimension-base-fibre-equals-total}을
사용할 수도 있다.
\end{proof}

\begin{lemma}
\label{more-algebra-lemma-completion-depth}
$A$를 Noether 국소환이라 하자. 그러면
$\text{depth}(A) = \text{depth}(A^\wedge)$이다.
\end{lemma}

\begin{proof}
Algebra, Lemma \ref{algebra-lemma-apply-grothendieck}를 보라.
\end{proof}

\begin{lemma}
\label{more-algebra-lemma-completion-CM}
$A$를 Noether 국소환이라 하자. 그러면 $A$가 Cohen--Macaulay일
필요충분조건은 $A^\wedge$가 Cohen--Macaulay인 것이다.
\end{lemma}

\begin{proof}
국소환 $A$가 Cohen--Macaulay일 필요충분조건은
$\dim(A) = \text{depth}(A)$인 것이다. 이 두 불변량은 완비화로 보존되므로
(보조정리 \ref{more-algebra-lemma-completion-dimension}과
\ref{more-algebra-lemma-completion-depth}), 주장이 따른다.
\end{proof}

\begin{lemma}
\label{more-algebra-lemma-completion-regular}
$A$를 Noether 국소환이라 하자. 그러면 $A$가 정칙일 필요충분조건은
$A^\wedge$가 정칙인 것이다.
\end{lemma}

\begin{proof}
$A^\wedge$가 정칙이면 Algebra, Lemma
\ref{algebra-lemma-flat-under-regular}에 의해 $A$는 정칙이다.
$A$가 정칙이라고 가정하자. $\mathfrak m$을 $A$의 극대 아이디얼이라
하자. 그러면
$\dim_{\kappa(\mathfrak m)} \mathfrak m/\mathfrak m^2 =
\dim(A) = \dim(A^\wedge)$이다(보조정리
\ref{more-algebra-lemma-completion-dimension}). 한편 $\mathfrak mA^\wedge$은
$A^\wedge$의 극대 아이디얼이므로 $\mathfrak m_{A^\wedge}$은 많아야
$\dim(A^\wedge)$개의 원소로 생성된다. 따라서 $A^\wedge$는 정칙이다.
(Algebra, Lemma
\ref{algebra-lemma-flat-over-regular-with-regular-fibre}를 사용할 수도 있다.)
\end{proof}

\begin{lemma}
\label{more-algebra-lemma-completion-dvr}
$A$를 Noether 국소환이라 하자. 그러면 $A$가 이산 값매김환일
필요충분조건은 $A^\wedge$가 이산 값매김환인 것이다.
\end{lemma}

\begin{proof}
이는 보조정리 \ref{more-algebra-lemma-completion-dimension}과
\ref{more-algebra-lemma-completion-regular}, 그리고 Algebra, Lemma
\ref{algebra-lemma-characterize-dvr}에서 따른다.
\end{proof}

\begin{lemma}
\label{more-algebra-lemma-completion-reduced}
$A$를 Noether 국소환이라 하자.
\begin{enumerate}
\item $A^\wedge$가 기약원이면 $A$도 기약원이다.
\item 일반적으로 $A$가 기약원이라고 해서 $A^\wedge$가 기약원인 것은 아니다.
\item $A$가 Nagata이면 $A$가 기약원일 필요충분조건은 $A^\wedge$가
기약원인 것이다.
\end{enumerate}
\end{lemma}

\begin{proof}
$A \to A^\wedge$가 충실히 평탄이므로 Algebra, Lemma
\ref{algebra-lemma-descent-reduced}에 의해 (1)을 얻는다.
(2)에 대해서는 Algebra, Example \ref{algebra-example-bad-dvr-char-p}를
보라. 표수 영인 예도 있다. Algebra, Remark
\ref{algebra-remark-resolution-dim-1}을 보라.
(3)에 대해서는 Algebra, Lemmas
\ref{algebra-lemma-local-nagata-domain-analytically-unramified}와
\ref{algebra-lemma-analytically-unramified-easy}를 보라.
\end{proof}

\begin{lemma}
\label{more-algebra-lemma-completion-normal}
$A$를 Noether 국소환이라 하자. $A^\wedge$가 정규이면 $A$도 정규이다.
\end{lemma}

\begin{proof}
$A \to A^\wedge$가 충실히 평탄이므로 이는 Algebra, Lemma
\ref{algebra-lemma-descent-normal}에서 따른다.
\end{proof}

\begin{lemma}
\label{more-algebra-lemma-flat-completion}
$A \to B$를 Noether 국소환의 국소 준동형이라 하자. 그러면 완비화의
유도 사상 $A^\wedge \to B^\wedge$가 평탄일 필요충분조건은
$A \to B$가 평탄인 것이다.
\end{lemma}

\begin{proof}
가환 도식
$$
\xymatrix{
A^\wedge \ar[r] & B^\wedge \\
A \ar[r] \ar[u] & B \ar[u]
}
$$
을 생각하자. 수직 화살표들은 충실히 평탄이다.
$A^\wedge \to B^\wedge$가 평탄이라고 가정하자. 그러면
$A \to B^\wedge$은 평탄하다. 따라서 Algebra, Lemma
\ref{algebra-lemma-flatness-descends-more-general}에 의해 $B$는
$A$ 위에서 평탄하다.

\medskip\noindent
$A \to B$가 평탄이라고 가정하자. 그러면 $A \to B^\wedge$은
평탄하다. 따라서 모든 $n \geq 1$에 대해
$B^\wedge/\mathfrak m_A^n B^\wedge$은 $A/\mathfrak m_A^n$ 위에서
평탄하다. $\mathfrak m_A^n A^\wedge$은 $A^\wedge$의 극대 아이디얼
$\mathfrak m_A^\wedge$의 $n$제곱이고
$A/\mathfrak m_A^n = A^\wedge/(\mathfrak m_A^\wedge)^n$임에 주의하자.
따라서 Algebra, Lemma \ref{algebra-lemma-flat-module-powers}를 적용하면
($R = A^\wedge$, $I = \mathfrak m_A^\wedge$,
$S = B^\wedge$, $M = S$), $B^\wedge$이 $A^\wedge$ 위에서 평탄임을
알 수 있다.
\end{proof}

\begin{lemma}
\label{more-algebra-lemma-flat-unramified}
$A \to B$를 $\mathfrak m_A B = \mathfrak m_B$ 및
$\kappa(\mathfrak m_A) = \kappa(\mathfrak m_B)$를 만족하는 Noether
국소환의 평탄 국소 준동형이라 하자. 그러면 $A \to B$는 완비화의 동형
$A^\wedge \to B^\wedge$을 유도한다.
\end{lemma}

\begin{proof}
Algebra, Lemma \ref{algebra-lemma-finite-after-completion}에 의해
$B^\wedge$은 $B$의 $\mathfrak m_A$-진 완비화이고
$A^\wedge \to B^\wedge$은 유한이다. $A \to B$가 평탄이므로
$\text{Tor}_1^A(B, \kappa(\mathfrak m_A)) = 0$이다. 따라서 보조정리
\ref{more-algebra-lemma-flat-after-completion}에 의해 $B^\wedge$은 $A^\wedge$ 위에서
평탄하다. 그러므로 Algebra, Lemma
\ref{algebra-lemma-finite-flat-local}에 의해 $B^\wedge$은 자유
$A^\wedge$-가군이다. 가정과 Algebra, Lemma
\ref{algebra-lemma-hathat-finitely-generated}에 의해
$A^\wedge \to B^\wedge$은 동형
$\kappa(\mathfrak m_A) = A^\wedge/\mathfrak m_A A^\wedge \to
B^\wedge/\mathfrak m_A B^\wedge = B^\wedge/\mathfrak m_B B^\wedge =
\kappa(\mathfrak m_B)$를 유도한다. 따라서 $B^\wedge$은 랭크 $1$인
자유 가군이다. 그러므로 $A^\wedge \to B^\wedge$은 동형이다.
\end{proof}





\section{에탈 사상에 의한 성질의 보존}
\label{more-algebra-section-permanence-etale}

\noindent
이 절에서는 에탈 환 사상 $\varphi : A \to B$를 생각하여 $A$의 어떤
성질이 $B$에 유전되는지, 그리고 $\mathfrak q$에서 $B$의 국소환이 갖는
어떤 성질이 $\mathfrak p = \varphi^{-1}(\mathfrak q)$에서 $A$의
국소환에 유전되는지 살펴본다. 기본적으로 이 절은 앞의 결과들을
복습하여 모은 것으로 새로운 내용을 더하지 않는다.

\medskip\noindent
에탈 환 사상은 평탄이고(Algebra, Lemma
\ref{algebra-lemma-etale}), 국소환의 평탄 국소 준동형은 충실히
평탄이라는 사실(Algebra, Lemma \ref{algebra-lemma-local-flat-ff})을
별도의 언급 없이 사용한다.

\begin{lemma}
\label{more-algebra-lemma-Noetherian-etale-extension}
$A \to B$가 에탈 환 사상이고 $\mathfrak q$가
$\mathfrak p \subset A$ 위에 놓이는 $B$의 소아이디얼이면,
$A_{\mathfrak p}$가 Noether일 필요충분조건은
$B_{\mathfrak q}$가 Noether인 것이다.
\end{lemma}

\begin{proof}
$A_\mathfrak p \to B_\mathfrak q$가 충실히 평탄이므로
$B_\mathfrak q$가 Noether이면 $A_\mathfrak p$도 Noether이다.
Algebra, Lemma \ref{algebra-lemma-descent-Noetherian}을 보라.
반대로 $A_\mathfrak p$가 Noether이면 $B_\mathfrak q$는 유한형
$A_\mathfrak p$-대수의 국소화이므로 Noether이다.
\end{proof}

\begin{lemma}
\label{more-algebra-lemma-dimension-etale-extension}
$A \to B$가 에탈 환 사상이고 $\mathfrak q$가
$\mathfrak p \subset A$ 위에 놓이는 $B$의 소아이디얼이면
$\dim(A_{\mathfrak p}) = \dim(B_{\mathfrak q})$이다.
\end{lemma}

\begin{proof}
$A_{\mathfrak p} \to B_{\mathfrak q}$가 평탄이므로 하강 정리가 성립하여
부등식 $\dim(A_{\mathfrak p}) \leq \dim(B_{\mathfrak q})$를 얻는다.
Algebra, Lemma \ref{algebra-lemma-dimension-going-up}을 보라. 한편
$$
\mathfrak q_0 \subset \mathfrak q_1 \subset \ldots \subset \mathfrak q_n
$$
이 $B_{\mathfrak q}$의 소아이디얼 사슬이라고 하자. 그러면 대응하는
소아이디얼의 열
$$
\mathfrak p_0 \subset \mathfrak p_1 \subset \ldots \subset \mathfrak p_n
$$
도 사슬이다. 여기서 $\mathfrak p_i = \mathfrak q_i \cap A_{\mathfrak p}$이고,
열에는 등식이 없다는 뜻이다. 실제로 에탈 환 사상은 준유한이고
(Algebra, Lemma \ref{algebra-lemma-etale-quasi-finite}를 보라),
준유한 환 사상은 정의상 이산 올을 갖는 스펙트럼의 사상을 유도한다.
% P02-E1434: frozen source says “is chain also,” omitting the article; Korean supplies natural syntax.
따라서 원하는 대로
$\dim(A_{\mathfrak p}) \geq \dim(B_{\mathfrak q})$이다.
\end{proof}

\begin{lemma}
\label{more-algebra-lemma-regular-etale-extension}
$A \to B$가 에탈 환 사상이고 $\mathfrak q$가
$\mathfrak p \subset A$ 위에 놓이는 $B$의 소아이디얼이면,
$A_{\mathfrak p}$가 정칙일 필요충분조건은 $B_{\mathfrak q}$가
정칙인 것이다.
\end{lemma}

\begin{proof}
동치를 증명하기 위해 보조정리 \ref{more-algebra-lemma-Noetherian-etale-extension}에
의해 $A_\mathfrak p$와 $B_\mathfrak q$가 모두 Noether라고 가정해도 된다.
$x_1, \ldots, x_t \in \mathfrak pA_\mathfrak p$를 극소 생성계라 하자.
$A_\mathfrak p \to B_\mathfrak q$가 충실히 평탄이므로 그 상
$y_1, \ldots, y_t$는 $B_\mathfrak q$에서
$\mathfrak pB_\mathfrak q = \mathfrak q B_\mathfrak q$의 극소 생성계를
이룬다(Algebra, Lemma \ref{algebra-lemma-etale-at-prime}). 정의상
$A_\mathfrak p$의 정칙성은 $t = \dim(A_\mathfrak p)$를 뜻하며
$B_\mathfrak q$에 대해서도 마찬가지이다. 따라서 보조정리
\ref{more-algebra-lemma-dimension-etale-extension}의 등식
$\dim(A_\mathfrak p) = \dim(B_\mathfrak q)$에서 결과가 따른다.
\end{proof}

\begin{lemma}
\label{more-algebra-lemma-Dedekind-etale-extension}
$A \to B$가 에탈 환 사상이고 $A$가 Dedekind 정역이면, $B$는
Dedekind 정역들의 유한 곱이다. 특히 극대 아이디얼
$\mathfrak q \subset B$에 대한 국소화 $B_\mathfrak q$는 이산 값매김환이다.
\end{lemma}

\begin{proof}
국소환에 대한 명제는 보조정리
\ref{more-algebra-lemma-dimension-etale-extension}과
\ref{more-algebra-lemma-regular-etale-extension}, 그리고 Algebra, Lemma
\ref{algebra-lemma-characterize-dvr}에서 따른다. 따라서 $B$는 차원이
$1$인 Noether 정규환이다. Algebra, Lemma
\ref{algebra-lemma-characterize-reduced-ring-normal}에 의해 $B$는 차원이
$1$인 정규 정역들의 유한 곱이다. 이들은 Algebra, Lemma
\ref{algebra-lemma-characterize-Dedekind}에 의해 Dedekind 정역이다.
\end{proof}



\section{헨젤화에 의한 성질의 보존}
\label{more-algebra-section-permanence-henselization}

\noindent
국소환 $R$이 주어졌을 때 $R$의 헨젤화와 엄밀 헨젤화를 각각
$R^h$, $R^{sh}$로 나타낸다. Algebra, Definition
\ref{algebra-definition-henselization}을 보라.
% P02-E1435: frozen source omits “by” in its paired notation-introduction sentence; Korean uses natural syntax.
이 절에서 보겠지만 $R$의 많은 성질은 $R^h$와 $R^{sh}$에 반영된다.

\begin{lemma}
\label{more-algebra-lemma-dumb-properties-henselization}
$(R, \mathfrak m, \kappa)$를 국소환이라 하자. 그러면 다음이 성립한다.
\begin{enumerate}
\item $R \to R^h \to R^{sh}$는 충실히 평탄인 환 사상이다.
\item $\mathfrak m R^h = \mathfrak m^h$이고
$\mathfrak m R^{sh} = \mathfrak m^h R^{sh} = \mathfrak m^{sh}$이다.
\item 모든 $n$에 대해 $R/\mathfrak m^n = R^h/\mathfrak m^nR^h$이다.
\item $R^{sh}/\mathfrak m^nR^{sh}$가
$x_i \bmod \mathfrak m^nR^{sh}$ 위의 자유 $R/\mathfrak m^n$-가군이
되게 하는 원소 $x_i \in R^{sh}$가 존재한다.
\end{enumerate}
\end{lemma}

\begin{proof}
구성상 $R^h$는 에탈 $R$-대수의 국소화들의 여극한이다. Algebra, Lemma
\ref{algebra-lemma-henselization}을 보라. 에탈 환 사상은 평탄이므로
(Algebra, Lemma \ref{algebra-lemma-etale}), Algebra, Lemma
\ref{algebra-lemma-colimit-flat}에 의해 $R^h$는 $R$ 위에서 평탄이다.
국소환의 평탄 국소 준동형은 충실히 평탄이므로(Algebra, Lemma
\ref{algebra-lemma-local-flat-ff}), $R \to R^h$는 충실히 평탄이다.
환 사상 $R^h \to R^{sh}$는 유한 에탈 환 사상들의 여극한이다.
Algebra, Lemma \ref{algebra-lemma-strict-henselization}의 증명을 보라.
따라서 위와 같은 논증으로 $R^h \to R^{sh}$도 충실히 평탄이다.

\medskip\noindent
(2)는 Algebra, Lemmas \ref{algebra-lemma-henselization}과
\ref{algebra-lemma-strict-henselization}에서 따른다. (3)은 Algebra, Lemma
\ref{algebra-lemma-local-artinian-basis-when-flat}에서 따른다. 실제로
$R/\mathfrak m \to R^h/\mathfrak mR^h$가 동형이고
$R/\mathfrak m^n \to R^h/\mathfrak m^nR^h$가 평탄 환 사상
$R \to R^h$의 밑변환으로서 평탄이므로(Algebra, Lemma
\ref{algebra-lemma-flat-base-change}) 이 주장을 적용할 수 있다.
$\kappa^{sep}$를 $R^{sh}$의 잉여체라 하자. 이는 $\kappa$의 분리
대수적 폐포이다. $\kappa$-벡터 공간으로서 $\kappa^{sep}$의 한 기저로
가는 $x_i \in R^{sh}$를 택하자. 그러면 위와 완전히 같은 방식으로
Algebra, Lemma \ref{algebra-lemma-local-artinian-basis-when-flat}에서
(4)가 따른다.
\end{proof}

\begin{lemma}
\label{more-algebra-lemma-henselization-formally-smooth}
$(R, \mathfrak m, \kappa)$를 국소환이라 하자. 그러면
\begin{enumerate}
\item $R \to R^h$, $R^h \to R^{sh}$, 및 $R \to R^{sh}$는 형식적으로
에탈이다.
\item $R \to R^h$, $R^h \to R^{sh}$, 그리고 $R \to R^{sh}$는 각각
$\mathfrak m^h$, $\mathfrak m^{sh}$, 그리고 $\mathfrak m^{sh}$-위상에서
형식적으로 매끄럽다.
\end{enumerate}
\end{lemma}

\begin{proof}
(1)은 $R^h$와 $R^{sh}$가 구성상 에탈 대수의 유향 여극한이라는 사실,
에탈 대수는 형식적으로 에탈이라는 사실(Algebra, Lemma
\ref{algebra-lemma-formally-etale-etale}), 그리고 형식적으로 에탈인
대수의 여극한도 형식적으로 에탈이라는 사실(Algebra, Lemma
\ref{algebra-lemma-colimit-formally-etale})에서 따른다.
(2)는 형식적으로 에탈인 환 사상이 형식적으로 매끄럽다는 사실과
보조정리 \ref{more-algebra-lemma-formally-smooth}에서 따른다.
\end{proof}

\begin{lemma}
\label{more-algebra-lemma-henselization-noetherian}
\begin{reference}
\cite[IV, Theorem 18.6.6 and Proposition 18.8.8]{EGA}
\end{reference}
$R$을 국소환이라 하자. 다음은 서로 동치이다.
\begin{enumerate}
\item $R$은 Noether이다.
\item $R^h$는 Noether이다.
\item $R^{sh}$는 Noether이다.
\end{enumerate}
이 경우 다음이 성립한다.
\begin{enumerate}
\item[(a)] $(R^h)^\wedge$와 $(R^{sh})^\wedge$는 Noether 완비 국소환이다.
\item[(b)] $R^\wedge \to (R^h)^\wedge$은 동형이다.
\item[(c)] $R^h \to (R^h)^\wedge$과
$R^{sh} \to (R^{sh})^\wedge$은 평탄이다.
\item[(d)] $R^\wedge \to (R^{sh})^\wedge$은
$\mathfrak m_{(R^{sh})^\wedge}$-진 위상에서 형식적으로 매끄럽다.
\item[(e)] $(R^\wedge)^{sh} = R^\wedge \otimes_{R^h} R^{sh}$이다.
\item[(f)] $((R^\wedge)^{sh})^\wedge = (R^{sh})^\wedge$이다.
\end{enumerate}
\end{lemma}

\begin{proof}
$R \to R^h \to R^{sh}$가 충실히 평탄이므로(보조정리
\ref{more-algebra-lemma-dumb-properties-henselization}), $R^h$ 또는 $R^{sh}$가
Noether이면 $R$도 Noether이다. Algebra, Lemma
\ref{algebra-lemma-descent-Noetherian}을 보라. 나머지 증명에서는
$R$이 Noether라고 가정한다.

\medskip\noindent
$\mathfrak m \subset R$은 유한 생성이므로
$\mathfrak m^h = \mathfrak m R^h$와
$\mathfrak m^{sh} = \mathfrak mR^{sh}$도 유한 생성이다. 보조정리
\ref{more-algebra-lemma-dumb-properties-henselization}을 보라. 따라서 Algebra, Lemma
\ref{algebra-lemma-complete-local-ring-Noetherian}에 의해
$(R^h)^\wedge$와 $(R^{sh})^\wedge$은 Noether이다. 이로써 (a)가 증명된다.

\medskip\noindent
(b)는 보조정리 \ref{more-algebra-lemma-dumb-properties-henselization}에서 즉시 따른다.
특히 $(R^h)^\wedge$은 $R$ 위에서 평탄이다. Algebra, Lemma
\ref{algebra-lemma-completion-faithfully-flat}을 보라.

\medskip\noindent
다음으로 $R^h \to (R^h)^\wedge$이 평탄임을 보이자. 이를 에탈
$R$-대수의 국소화들의 유향 여극한 $R^h = \colim_i R_i$로 쓰자.
Algebra, Lemma \ref{algebra-lemma-colimit-rings-flat}에 의해
$(R^h)^\wedge$이 각 $R_i$ 위에서 평탄이면
$R^h \to (R^h)^\wedge$은 평탄이다. 구성상 $R^h = R_i^h$임에
주의하자. 따라서 (b)에 의해 $R_i^\wedge = (R^h)^\wedge$이고,
이는 원하는 대로 $R_i$ 위에서 평탄이다. (c)의 증명을 끝내기 위해
$R^{sh} \to (R^{sh})^\wedge$이 평탄임을 보이자. 이를 위해 위와 같은
극한 논증에 의해 $(R^{sh})^\wedge$이 $R$ 위에서 평탄임을 보이면
충분하다. 보조정리 \ref{more-algebra-lemma-dumb-properties-henselization}에 의해
$(R^{sh})^\wedge$은 자유 $R$-가군의 완비화이다. 보조정리
\ref{more-algebra-lemma-completed-direct-sum-flat}에 의해 이는 원하는 대로 $R$ 위에서
평탄이다. 이로써 (c)가 증명된다.

\medskip\noindent
이제 (c)가 성립하고 $(R^h)^\wedge$과 $(R^{sh})^\wedge$이 Noether임을
안다. 따라서 Algebra, Lemma \ref{algebra-lemma-descent-Noetherian}에
의해 $R^h$와 $R^{sh}$는 Noether이다.

\medskip\noindent
(d)는 보조정리 \ref{more-algebra-lemma-henselization-formally-smooth}와 보조정리
\ref{more-algebra-lemma-formally-smooth-completion}에서 따른다.

\medskip\noindent
(e)는 Algebra, Lemma \ref{algebra-lemma-sh-from-h-map}과
$R^\wedge$이 헨젤이라는 사실에서 따른다. Algebra, Lemma
\ref{algebra-lemma-complete-henselian}을 보라.

\medskip\noindent
(f)의 증명. (e)를 사용하면 완비화에 의해 사상
$R^{sh} \to (R^\wedge)^{sh}$가 사상
$(R^{sh})^\wedge \to ((R^\wedge)^{sh})^\wedge$을 유도한다.
(e)를 사용하면 사상 $R^\wedge \to (R^{sh})^\wedge$도 있다.
$(R^{sh})^\wedge$은 엄밀 헨젤이므로(위를 보라), 이 사상은 Algebra,
Lemma \ref{algebra-lemma-strictly-henselian-functorial}에 의해 사상
$(R^\wedge)^{sh} \to (R^{sh})^\wedge$을 유도한다. 완비화하면 사상
$((R^\wedge)^{sh})^\wedge \to (R^{sh})^\wedge$을 얻는다. 이 두 사상이
서로 역이라는 확인은 생략한다.
\end{proof}

\begin{lemma}
\label{more-algebra-lemma-henselization-reduced}
\begin{slogan}
기약원성은 (엄밀) 헨젤화로 전달된다.
\end{slogan}
$R$을 국소환이라 하자. 다음은 서로 동치이다. $R$은 기약원이다.
$R$의 헨젤화 $R^h$는 기약원이다. $R$의 엄밀 헨젤화 $R^{sh}$는
기약원이다.
\end{lemma}

\begin{proof}
환 사상 $R \to R^h \to R^{sh}$는 충실히 평탄이다. 따라서 함의의
한 방향은 Algebra, Lemma \ref{algebra-lemma-descent-reduced}에서 따른다.
반대로 $R$이 기약원이라고 하자. $R^h$와 $R^{sh}$는 에탈, 따라서
매끄러운 $R$-대수의 여과 여극한이므로 결과는 Algebra, Lemma
\ref{algebra-lemma-reduced-goes-up}에서 따른다.
\end{proof}

\begin{lemma}
\label{more-algebra-lemma-henselization-nil}
$R$을 국소환이라 하자. $nil(R)$로 $R$의 멱영원들로 이루어진
아이디얼을 나타내자. 그러면 $nil(R)R^h = nil(R^h)$이고
$nil(R)R^{sh} = nil(R^{sh})$이다.
\end{lemma}

\begin{proof}
$nil(R)$은 멱영원들로 이루어진 가장 큰 아이디얼이며 그 몫
$R/nil(R)$은 기약원임에 주의하자. Algebra, Lemma
\ref{algebra-lemma-locally-nilpotent}에 의해 $nil(R)R^h$는 멱영원들로
이루어진다. 또한 $R^h/nil(R) R^h$는 Algebra, Lemma
\ref{algebra-lemma-quotient-henselization}에 의해 $R/nil(R)$의
헨젤화이다. 따라서 보조정리 \ref{more-algebra-lemma-henselization-reduced}에 의해
$R^h/nil(R)R^h$는 기약원이다. 그러므로 원하는 대로
$nil(R) R^h = nil(R^h)$이다. 엄밀 헨젤화도 마찬가지이며, 이때에는
Algebra, Lemma \ref{algebra-lemma-quotient-strict-henselization}을 사용한다.
\end{proof}

\begin{lemma}
\label{more-algebra-lemma-henselization-normal}
$R$을 국소환이라 하자. 다음은 서로 동치이다. $R$은 정규 정역이다.
$R$의 헨젤화 $R^h$는 정규 정역이다. $R$의 엄밀 헨젤화 $R^{sh}$는
정규 정역이다.
\end{lemma}

\begin{proof}
먼저 국소환이 정규일 필요충분조건은 정규 정역인 것임에 주의하자
(Algebra, Definition \ref{algebra-definition-ring-normal}을 보라).
환 사상 $R \to R^h \to R^{sh}$는 충실히 평탄이다. 따라서 함의의
한 방향은 Algebra, Lemma \ref{algebra-lemma-descent-normal}에서 따른다.
반대로 $R$이 정규라고 하자. $R^h$와 $R^{sh}$는 에탈, 따라서 매끄러운
$R$-대수의 여과 여극한이므로 결과는 Algebra, Lemmas
\ref{algebra-lemma-normal-goes-up}과
\ref{algebra-lemma-colimit-normal-ring}에서 따른다.
\end{proof}

\begin{lemma}
\label{more-algebra-lemma-henselization-dimension}
임의의 국소환 $R$에 대해
$\dim(R) = \dim(R^h) = \dim(R^{sh})$이다.
\end{lemma}

\begin{proof}
$R \to R^{sh}$는 충실히 평탄이므로(보조정리
\ref{more-algebra-lemma-dumb-properties-henselization}), 하강 정리에 의해
$\dim(R^{sh}) \geq \dim(R)$이다. Algebra, Lemma
\ref{algebra-lemma-dimension-going-up}을 보라. 반대 부등식을 위해
$R^{sh} = \colim R_i$를 국소환들의 유향 여극한으로 쓰자. 각 $R_i$는
에탈 $R$-대수의 국소화이다. 이제
$\mathfrak q_0 \subset \mathfrak q_1 \subset \ldots \subset \mathfrak q_n$이
$R^{sh}$의 소아이디얼 사슬이면, 충분히 큰 어떤 $i$에 대해 열
$$
R_i \cap \mathfrak q_0 \subset
R_i \cap \mathfrak q_1 \subset \ldots \subset
R_i \cap \mathfrak q_n
$$
은 $R_i$의 소아이디얼 사슬이다. 따라서
$\dim(R^{sh}) \leq \sup_i \dim(R_i)$이다. 그러나 보조정리
\ref{more-algebra-lemma-dimension-etale-extension}의 결과에 의해 각 $i$에 대해
$\dim(R_i) = \dim(R)$이고, 이것으로 끝난다.
\end{proof}

\begin{lemma}
\label{more-algebra-lemma-henselization-depth}
임의의 Noether 국소환 $R$에 대해
$\text{depth}(R) = \text{depth}(R^h) = \text{depth}(R^{sh})$이다.
\end{lemma}

\begin{proof}
보조정리 \ref{more-algebra-lemma-henselization-noetherian}에 의해 $R^h$와 $R^{sh}$는
Noether이다. 따라서 보조정리는 Algebra, Lemma
\ref{algebra-lemma-apply-grothendieck}에서 따른다.
\end{proof}

\begin{lemma}
\label{more-algebra-lemma-henselization-CM}
$R$을 Noether 국소환이라 하자. 다음은 서로 동치이다. $R$은
Cohen--Macaulay이다. $R$의 헨젤화 $R^h$는 Cohen--Macaulay이다.
$R$의 엄밀 헨젤화 $R^{sh}$는 Cohen--Macaulay이다.
\end{lemma}

\begin{proof}
보조정리 \ref{more-algebra-lemma-henselization-noetherian}에 의해 $R^h$와 $R^{sh}$가
Noether임을 알므로 보조정리의 내용은 의미가 있다. 보조정리
\ref{more-algebra-lemma-henselization-depth}와 \ref{more-algebra-lemma-henselization-dimension}에
의해
$\text{depth}(R) = \text{depth}(R^h) = \text{depth}(R^{sh})$
이고
$\dim(R) = \dim(R^h) = \dim(R^{sh})$
이므로 결론을 얻는다.
\end{proof}

\begin{lemma}
\label{more-algebra-lemma-henselization-regular}
$R$을 Noether 국소환이라 하자. 다음은 서로 동치이다. $R$은 정칙
국소환이다. $R$의 헨젤화 $R^h$는 정칙 국소환이다. $R$의 엄밀
헨젤화 $R^{sh}$는 정칙 국소환이다.
\end{lemma}

\begin{proof}
보조정리 \ref{more-algebra-lemma-henselization-noetherian}에 의해 $R^h$와 $R^{sh}$가
Noether임을 알므로 보조정리의 내용은 의미가 있다. $\mathfrak m$을
$R$의 극대 아이디얼이라 하자. $x_1, \ldots, x_t \in \mathfrak m$을
$\mathfrak m$의 극소 생성계, 즉 $\mathfrak m/\mathfrak m^2$에서의 상이
$\kappa = R/\mathfrak m$ 위의 기저를 이루는 생성계라 하자.
$R \to R^h$와 $R \to R^{sh}$가 충실히 평탄이므로, $R^h$에서의 상
$x_1^h, \ldots, x_t^h$와 $R^{sh}$에서의 상
$x_1^{sh}, \ldots, x_t^{sh}$는 각각
$\mathfrak m^h = \mathfrak mR^h$와
$\mathfrak m^{sh} = \mathfrak mR^{sh}$의 극소 생성계를 이룬다.
정의상 $R$의 정칙성은 $t = \dim(R)$를 뜻하고 $R^h$와 $R^{sh}$에
대해서도 마찬가지이다. 따라서 보조정리
\ref{more-algebra-lemma-henselization-dimension}의 차원 등식
$\dim(R) = \dim(R^h) = \dim(R^{sh})$에서 결과가 따른다.
% P02-E1436: frozen proof has no sentence-ending punctuation before the proof environment closes; Korean supplies normal punctuation.
\end{proof}

\begin{lemma}
\label{more-algebra-lemma-henselization-dvr}
$R$을 Noether 국소환이라 하자. 그러면 $R$이 이산 값매김환일
필요충분조건은 $R^h$가 이산 값매김환인 것이고, 이는 다시
$R^{sh}$가 이산 값매김환인 것과 동치이다.
\end{lemma}

\begin{proof}
이는 보조정리 \ref{more-algebra-lemma-henselization-dimension}과
\ref{more-algebra-lemma-henselization-regular}, 그리고 Algebra, Lemma
\ref{algebra-lemma-characterize-dvr}에서 따른다.
\end{proof}

\begin{lemma}
\label{more-algebra-lemma-filtered-colimit-etale-noetherian-fibres}
$A$를 환이라 하자. $B$를 에탈 $A$-대수의 여과 여극한이라 하자.
$\mathfrak p$를 $A$의 소아이디얼이라 하자. $B$가 Noether이면
$\mathfrak p$ 위에 놓이는 소아이디얼
$\mathfrak q_1, \ldots, \mathfrak q_r$은 유한 개이고,
$B \otimes_A \kappa(\mathfrak p) = \prod \kappa(\mathfrak q_i)$이며,
각 체 확대 $\kappa(\mathfrak q_i)/\kappa(\mathfrak p)$는 분리 대수적이다.
\end{lemma}

\begin{proof}
$B$를 $A \to B_i$가 에탈인 여과 여극한 $B = \colim B_i$로 쓰자.
한편
$B \otimes_A \kappa(\mathfrak p) = \colim B_i \otimes_A \kappa(\mathfrak p)$는
에탈 $\kappa(\mathfrak p)$-대수의 여과 여극한이고, 다른 한편으로
Noether이다. 에탈 $\kappa(\mathfrak p)$-대수는 유한 분리 체 확대들의
유한 곱이다(Algebra, Lemma \ref{algebra-lemma-etale-over-field}).
따라서 대수 $B_i \otimes_A \kappa(\mathfrak p)$의 소아이디얼들
(모두 극대이면서 극소이다) 사이에는 자명하지 않은 특수화가 없고,
$B \otimes_A \kappa(\mathfrak p)$의 소아이디얼들 사이에도 그러한
특수화가 없다. 따라서 $B \otimes_A \kappa(\mathfrak p)$는 기약원이고,
그 모든 소아이디얼은 극소이며 개수는 유한하다.
% P02-E1437: frozen source says the finitely many primes “which all minimal,” omitting “are”; Korean supplies the predicate.
따라서 이는 체들의 유한 곱이다(Algebra, Lemma
\ref{algebra-lemma-total-ring-fractions-no-embedded-points} 또는 Algebra,
Proposition \ref{algebra-proposition-dimension-zero-ring}을 사용하라).
각 체는 유한 분리 확대들의 여극한이므로 보조정리의 마지막 명제가 따른다.
\end{proof}

\begin{lemma}
\label{more-algebra-lemma-fibres-henselization}
$R$을 Noether 국소환이라 하자. $\mathfrak p \subset R$을 소아이디얼이라
하자. 그러면
$$
R^h \otimes_R \kappa(\mathfrak p) =
\prod\nolimits_{i = 1, \ldots, t} \kappa(\mathfrak q_i)
\quad\text{resp.}\quad
R^{sh} \otimes_R \kappa(\mathfrak p) =
\prod\nolimits_{i = 1, \ldots, s} \kappa(\mathfrak r_i)
$$
이다. 여기서 $\mathfrak q_1, \ldots, \mathfrak q_t$와
$\mathfrak r_1, \ldots, \mathfrak r_s$는 각각 $\mathfrak p$ 위에 놓이는
$R^h$와 $R^{sh}$의 소아이디얼이다.
% P02-E1438: frozen source says “are the prime” for two plural lists; Korean uses plural agreement.
더욱이 체 확대 $\kappa(\mathfrak q_i)/\kappa(\mathfrak p)$와
$\kappa(\mathfrak r_i)/\kappa(\mathfrak p)$는 각각 분리 대수적이다.
\end{lemma}

\begin{proof}
이는 헨젤화와 엄밀 헨젤화가 Noether라는 사실(위에서 보았다)과 더
일반적인 보조정리 \ref{more-algebra-lemma-filtered-colimit-etale-noetherian-fibres}에서
얻을 수 있다. 그러나 다음과 같은 직접적인 증명도 제시한다.

\medskip\noindent
보조정리 \ref{more-algebra-lemma-dumb-properties-henselization}과
\ref{more-algebra-lemma-henselization-noetherian}의 결과를 별도의 언급 없이
사용한다. $R^h/\mathfrak pR^h$와 $R^{sh}/\mathfrak pR^{sh}$는 각각
$R/\mathfrak p$의 헨젤화와 엄밀 헨젤화이다. 각각 Algebra, Lemma
\ref{algebra-lemma-quotient-henselization}과 Algebra, Lemma
\ref{algebra-lemma-quotient-strict-henselization}을 보라.
따라서 $R$을 $R/\mathfrak p$로 바꾸어 $R$이 Noether 국소 정역이고
$\mathfrak p = (0)$이라고 가정해도 된다. $R^h$와 $R^{sh}$는
Noether이므로 각각 유한 개의 극소 소아이디얼
$\mathfrak q_1, \ldots, \mathfrak q_t$와
$\mathfrak r_1, \ldots, \mathfrak r_s$를 갖는다.
$R \to R^h$와 $R \to R^{sh}$는 각각 평탄이므로, 하강 정리에 의해
이것들은 정확히 $\mathfrak p = (0)$ 위에 놓이는 소아이디얼들이다.
마지막으로 $R$이 정역이므로 보조정리
\ref{more-algebra-lemma-henselization-reduced}에 의해 $R^h$와 $R^{sh}$는 각각
기약원이다. 따라서 $R^h \otimes_R \kappa(\mathfrak p)$와
$R^{sh} \otimes_R \kappa(\mathfrak p)$는 소아이디얼이 유한 개이고 모두
극소인 기약원 Noether 환이다. 따라서 이 환들은 Artin이고 각 극대
아이디얼에서의 국소화들의 곱이며, 각 국소화는 필연적으로 체이다
(Algebra, Proposition \ref{algebra-proposition-dimension-zero-ring}과
Algebra, Lemma \ref{algebra-lemma-minimal-prime-reduced-ring}을 보라).

\medskip\noindent
마지막 명제는 $R \to R^h$와 $R \to R^{sh}$가 각각 에탈 환 사상의
여극한이고, 따라서 유도된 잉여체 확대가 유한 분리 확대들의
여극한이라는 사실에서 따른다. Algebra, Lemma
\ref{algebra-lemma-etale-at-prime}을 보라.
\end{proof}









\section{체 확대 다시 보기}
\label{more-algebra-section-p-bases}

\noindent
이 절에서는 표수가 $p > 0$인 체 확대의 몇 가지 특이한 점을 연구한다.

\begin{definition}
\label{more-algebra-definition-p-basis}
$p$를 소수라 하자. $k \to K$를 표수가 $p$인 체의 확대라 하자.
$K$ 안에서 $k$와 $K^p$의 합성체를 $kK^p$로 나타내자.
% P02-E1439: frozen source omits “by” in “Denote kK^p the compositum”; Korean uses natural notation syntax.
\begin{enumerate}
\item 부분집합 $\{x_i\} \subset K$에 대하여, $0 \leq e_i < p$인
$x^E = \prod x_i^{e_i}$들이 $kK^p$ 위에서 선형 독립이면
이를 $k$ 위에서 {\it p-독립}이라 한다.
\item $K$의 부분집합 $\{x_i\}$에 대하여, $x^E$들이 $kK^p$ 위에서
$K$의 한 기저를 이루면 이를 $k$ 위에서 $K$의
{\it p-기저}라 한다.
\end{enumerate}
\end{definition}

\noindent
이는 나중에 논의할 $\mathbf{F}_p$-대수의 $p$-기저라는 개념과 관련된다
(여기에 미래 참조를 삽입할 것).
% P02-E1440: frozen source contains the unresolved editorial placeholder “insert future reference here”.

\begin{lemma}
\label{more-algebra-lemma-p-basis}
$K/k$를 체 확대라 하자. $k$의 표수가 $p > 0$이라고 가정하자.
$\{x_i\}$를 $K$의 부분집합이라 하자. 다음은 서로 동치이다.
\begin{enumerate}
\item 원소들 $\{x_i\}$는 $k$ 위에서 $p$-독립이다.
\item 원소들 $\text{d}x_i$는 $\Omega_{K/k}$에서 $K$-선형 독립이다.
\end{enumerate}
임의의 $p$-독립 집합은 $k$ 위에서 $K$의 한 $p$-기저로 연장할 수 있다.
특히 체 $K$는 $k$ 위의 $p$-기저를 갖는다. 더욱이 다음은 서로 동치이다.
\begin{enumerate}
\item[(a)] $\{x_i\}$는 $k$ 위에서 $K$의 한 $p$-기저이다.
\item[(b)] $\text{d}x_i$는 $K$-벡터 공간 $\Omega_{K/k}$의 한 기저이다.
\end{enumerate}
\end{lemma}

\begin{proof}
(2)를 가정하고 $a_E \in k K^p$인 선형 관계
$\sum a_E x^E = 0$이 있다고 하자. $\theta_i(x_j) = \delta_{ij}$
(Kronecker 델타)를 만족하는 $k$-미분 $\theta_i : K \to K$를 택하자.
$K$의 임의의 $k$-미분은 $kK^p$를 소멸시킴에 주의하자. 주어진 관계에
$\theta_i$를 적용하면 새로운 관계
$$
\sum\nolimits_{E, e_i > 0}
e_i a_E x_1^{e_1}\ldots x_i^{e_i - 1} \ldots x_n^{e_n} = 0
$$
을 얻는다. 따라서 어떤 $a_E \not = 0$에 대해 총차수
$|E| = \sum e_i$가 최소인 관계로 $\sum a_E x^E$를 택하면 모순을
얻는다. 그러므로 (1)이 성립한다.

\medskip\noindent
$\{x_i\}$가 $k$ 위에서 $K$의 한 $p$-기저이면
$K \cong kK^p[X_i]/(X_i^p - x_i^p)$이다. 따라서 $\text{d}x_i$가
$K$ 위에서 $\Omega_{K/k}$의 한 기저를 이룸을 알 수 있다.
그러므로 (a)는 (b)를 함의한다.

\medskip\noindent
$\{x_i\}$를 $k$ 위에서 $K$의 $p$-독립 부분집합이라 하자.
Zorn 보조정리를 적용하여 이를 $k$ 위에서 $K$의 극대 $p$-독립
부분집합으로 확대할 수 있다. $K$의 임의의 극대 $p$-독립 부분집합
$\{x_i\}$가 $k$ 위에서 $K$의 한 $p$-기저라고 주장한다. 이 주장에서
(1)이 (2)를 함의하고 $p$-기저의 존재도 따른다. 주장을 증명하기 위해
$L$을 $kK^p$와 $x_i$들이 생성하는 $K$의 부분체라 하자.
$L = K$임을 보여야 한다. $x \in K$이지만 $x \not \in L$이면
$x^p \in L$이고 $L(x) \cong L[z]/(z^p - x)$이다.
% P02-E1441: frozen quotient has coefficient x although x is assumed outside L; the proof requires z^p minus x^p.
따라서 $\{x_i\} \cup \{x\}$는 $k$ 위에서 $p$-독립이고, 이는 모순이다.

\medskip\noindent
마지막으로 (b)가 (a)를 함의함을 보여야 한다. (1)과 (2)의 동치에 의해
$\{x_i\}$는 $k$ 위에서 $K$의 극대 $p$-독립 부분집합이다. 따라서 위
주장에 의해 이는 $p$-기저이다.
\end{proof}

\begin{lemma}
\label{more-algebra-lemma-intersection-subfields-subspace}
$K/k$를 체 확대라 하자. $\{K_\alpha\}_{\alpha \in A}$를 다음 성질을
갖는 $K$의 부분체들의 모음이라 하자.
\begin{enumerate}
\item 모든 $\alpha \in A$에 대해 $k \subset K_\alpha$이다.
\item $k = \bigcap_{\alpha \in A} K_\alpha$이다.
\item $\alpha, \alpha' \in A$에 대해
$K_{\alpha''} \subset K_\alpha \cap K_{\alpha'}$를 만족하는
$\alpha'' \in A$가 존재한다.
\end{enumerate}
그러면 $n \geq 1$이고 $V \subset K^{\oplus n}$가 $K$-벡터 공간일 때,
$V \cap k^{\oplus n} \not = 0$일 필요충분조건은 모든 $\alpha \in A$에
대해 $V \cap K_\alpha^{\oplus n} \not = 0$인 것이다.
% P02-E1442: frozen source omits a linking verb before “a K-vector space”; Korean supplies natural syntax.
\end{lemma}

\begin{proof}
$n$에 대한 귀납법을 사용한다. $n = 1$인 경우는 가정에서 따른다.
$K^{\oplus n - 1}$의 부분공간에 대해 결과가 증명되었다고 하자.
모든 $\alpha \in A$에 대해 $V \subset K^{\oplus n}$와
$K_\alpha^{\oplus n}$의 교집합이 영이 아니라고 가정하자.
$V \cap 0 \oplus k^{\oplus n - 1}$이 영이 아니면 끝난다. 따라서 그렇지
않다고 가정해도 된다. 귀납 가정에 의해
$V \cap 0 \oplus K_\alpha^{\oplus n - 1}$이 영인 $\alpha$를 찾을 수 있다.
$v = (x_1, \ldots, x_n) \in V \cap K_\alpha^{\oplus n}$를 영이 아닌
원소라 하자. $\alpha$의 선택에 의해 $x_1$은 영이 아니다.
$v$를 $x_1^{-1}v$로 바꾸어 $v = (1, x_2, \ldots, x_n)$이라 하자.
$v' = (x_1', \ldots, x'_n) \in V \cap K_\alpha^{\oplus n}$이면
$\alpha$의 선택에 의해 $v' - x_1'v = 0$임에 주의하자. 따라서
$V \cap K_\alpha^{\oplus n} = K_\alpha v$이다.
$K_{\alpha'} \subset K_\alpha$를 만족하는 어떤 $\alpha'$를 택하면,
$\alpha'$에 같은 논증을 적용하여 반드시
$v \in V \cap K_{\alpha'}^{\oplus n}$임을 알 수 있다. 따라서
$$
x_2, \ldots, x_n \in
\bigcap\nolimits_{\alpha' \in A, K_{\alpha'} \subset K_\alpha} K_{\alpha'}
$$
이고, 이는 (2)와 (3)에 의해 $k$와 같다.
\end{proof}

\begin{lemma}
\label{more-algebra-lemma-intersection-subfields}
$K$를 표수가 $p$인 체라 하자. $\{K_\alpha\}_{\alpha \in A}$를 다음
성질을 갖는 $K$의 부분체들의 모음이라 하자.
\begin{enumerate}
\item 모든 $\alpha \in A$에 대해 $K^p \subset K_\alpha$이다.
\item $K^p = \bigcap_{\alpha \in A} K_\alpha$이다.
\item $\alpha, \alpha' \in A$에 대해
$K_{\alpha''} \subset K_\alpha \cap K_{\alpha'}$를 만족하는
$\alpha'' \in A$가 존재한다.
\end{enumerate}
그러면 다음이 성립한다.
\begin{enumerate}
\item 사상들 $\Omega_{K/\mathbf{F}_p} \to \Omega_{K/K_\alpha}$의 핵들의
교집합은 영이다.
\item 임의의 유한 확대 $L/K$에 대해
$L^p = \bigcap_{\alpha \in A} L^pK_\alpha$이다.
\end{enumerate}
\end{lemma}

\begin{proof}
(1)의 증명. $\mathbf{F}_p$ 위에서 $K$의 한 $p$-기저 $\{x_i\}$를 택하자.
$\eta = \sum_{i \in I'} y_i \text{d}x_i$가 모든 $\alpha \in A$에 대해
$\Omega_{K/K_\alpha}$에서 영으로 간다고 하자. 여기서 첨자 집합
$I'$은 유한하다. 보조정리 \ref{more-algebra-lemma-p-basis}에 의해 이는 모든
$\alpha$에 대해 관계
$$
\sum\nolimits_E a_{E, \alpha} x^E = 0,\quad a_{E, \alpha} \in K_\alpha
$$
가 존재한다는 뜻이다. 여기서 $E$는 $0 \leq e_i < p$인 다중 첨자
$E = (e_i)_{i \in I'}$를 돈다. 한편 보조정리
\ref{more-algebra-lemma-p-basis}에 의해 $a_E \in K^p$인 그러한 관계
$\sum a_E x^E = 0$은 없다. 이는 보조정리
\ref{more-algebra-lemma-intersection-subfields-subspace}에 모순이다.

\medskip\noindent
(2)의 증명. 유한 체 확대의 탑 $L/M/K$가 있다고 하자.
$M_\alpha = M^p K_\alpha$ 및
$L_\alpha = L^p K_\alpha = L^p M_\alpha$라 두자. 그러면 먼저
$M^p = \bigcap_{\alpha \in A} M_\alpha$를 증명한 다음
$L^p = \bigcap_{\alpha \in A} L_\alpha$를 증명할 수 있다.
따라서 자명하지 않은 부분체가 없는 단순 체 확대에 대해 (2)를
증명하면 충분하다. 먼저 $L = K(\theta)$가 $K$ 위에서 분리라고 하자.
그러면 $L$은 $K$ 위에서 $\theta^p$에 의해 생성되므로
$\theta \in L^p$라고 가정해도 된다. 이 경우
$$
L^p = K^p \oplus K^p\theta \oplus \ldots K^p\theta^{d - 1}
\quad\text{and}\quad
L^pK_\alpha =
K_\alpha \oplus K_\alpha \theta \oplus \ldots K_\alpha\theta^{d - 1}
$$
이고, 여기서 $d = [L : K]$이다. 따라서 이 경우 결론은 명백하다.
다른 경우에는 $L = K(\theta)$이고
$\theta^p = t \in K$, $t \not \in K^p$이다. 이 경우
$$
L^p = K^p \oplus K^pt \oplus \ldots K^pt^{p - 1}
\quad\text{and}\quad
L^pK_\alpha =
K_\alpha \oplus K_\alpha t \oplus \ldots K_\alpha t^{p - 1}
$$
이다. 다시 결론은 명백하다.
\end{proof}

\begin{lemma}
\label{more-algebra-lemma-power-series-ring-subfields}
$k$를 표수가 $p > 0$인 체라 하자. $\{x_i\}_{i \in I}$를 $k$의 한
$p$-기저라 하자. $n, m \geq 0$라 하자.
$K$를 $A = k[[x_1, \ldots, x_n]][y_1, \ldots, y_m]$의 분수체라 하자.
$J$를 $I$의 유한 부분집합이라 하자. $k^p$와 $i \in I \setminus J$인
$x_i$들이 생성하는 부분체 $k/k_J/k^p$를 생각하자. 다음의 분수체
$K_J$들은
$$
A_J = k_J[[x_1^p, \ldots, x_n^p]][y_1^p, \ldots, y_m^p]
$$
$K$의 부분체들의 족을 이루고 보조정리
\ref{more-algebra-lemma-intersection-subfields}의 조건을 만족한다. 더욱이 각 환 확대
$A_J \subset A$는 유한이다.
\end{lemma}

\begin{proof}
$k/k_J$가 유한이므로 환 확대
$k_J[[x_1^p, \ldots, x_d^p]] \subset k[[x_1, \ldots, x_d]]$는
Algebra, Lemma \ref{algebra-lemma-finite-after-completion}에 의해 유한이다.
% P02-E1443: frozen proof uses x_d although the lemma defines n power-series variables; formal span is preserved exactly.
이는 $A_J \to A$가 유한임을 함의한다.

\medskip\noindent
보조정리 \ref{more-algebra-lemma-intersection-subfields}의 성질 (1), (2), (3)을
확인하자. (1)의 증명. $a \in A$에 대해 $a^p \in A_J$이다. 따라서
$K^p \subset K_J$이다. (2)의 증명. 모든 $J$의 선택에 대해
$K_J$에 속하는 $f/g^p \in K$, $f, g \in A$, $g \not = 0$이 있다고 하자.
당분간 $J$를 고정하자. $f/g^p \in K_J$이므로 $a \in A_J$이고
$b \in A$가 영이 아닐 때 $f/g^p = a/b^p$로 쓸 수 있다. 따라서
$b^p f \in A_J$이다. 임의의 $A_J$-미분 $D : A \to A$에 대해
$0 = D(b^pf) = b^p D(f)$이고, $A$가 정역이므로 $D(f) = 0$이다.
$D = \partial_{x_i}$와 $D = \partial_{y_j}$를 취하면
$f \in k[[x_1^p, \ldots, x_n^p]][y_1^p, \ldots, y_d^p]$라는 결론을 얻는다.
% P02-E1444: frozen conclusion uses y_d although the polynomial variables are indexed through m; formal span is preserved exactly.
$k_J$-미분 $\theta : k \to k$를 적용해도 마찬가지로 $f$의 모든
계수가 $k_J$에 속함을 얻는다. 즉 $f \in A_J$이다. 보조정리에서와 같은
모든 부분체를 $J$가 돌 때 $A^p = \bigcap\nolimits_J A_J$임이 명백하므로,
원하는 대로 $f \in A^p$라는 결론을 얻는다.
% P02-E1445: frozen source says J ranges over subfields although J was defined as a finite subset of I.
(3)의 증명. $K_{J \cup J'} \subset K_J \cap K_{J'}$이므로 명백하다.
\end{proof}





\section{특이 자취}
\label{more-algebra-section-singular-locus}

\noindent
$R$을 Noether 환이라 하자. $X = \Spec(R)$의 {\it 정칙 자취}
$\text{Reg}(X)$는 $R_\mathfrak p$가 정칙 국소환인 소아이디얼
$\mathfrak p$들의 집합이다. $X = \Spec(R)$의 {\it 특이 자취}
$\text{Sing}(X)$는 여집합 $X \setminus \text{Reg}(X)$, 즉
$R_\mathfrak p$가 정칙 국소환이 아닌 소아이디얼 $\mathfrak p$들의
집합이다. Algebra, Definition \ref{algebra-definition-regular} 앞의
논의에 의해 $\text{Reg}(X)$는 일반화에 대해 안정적이다. 이 절에서는
$\text{Reg}(X)$가 열림임을 보장하는 조건을 연구한다.

\begin{definition}
\label{more-algebra-definition-J}
\begin{reference}
\cite[(32.B)]{MatCA}
\end{reference}
$R$을 Noether 환이라 하자. $X = \Spec(R)$라 하자.
\begin{enumerate}
\item $\text{Reg}(X)$가 공집합이 아닌 열린집합을 포함하면 $R$을
{\it J-0}이라 한다.
\item $\text{Reg}(X)$가 열림이면 $R$을 {\it J-1}이라 한다.
\item 임의의 유한형 $R$-대수가 J-1이면 $R$을 {\it J-2}라 한다.
\end{enumerate}
\end{definition}

\noindent
환 $\mathbf{Q}[x]/(x^2)$은 J-0을 만족하지 않지만 J-1은 만족한다.
한편 Noether 정역과 더 일반적으로 영이 아닌 기약원 Noether 환은
극소 소아이디얼에서 정칙이므로, 이러한 환에 대해서는 J-1이 J-0을
함의한다. 다음은 J-1 성질의 한 특징짓기이다.

\begin{lemma}
\label{more-algebra-lemma-J-1}
$R$을 Noether 환이라 하자. $X = \Spec(R)$라 하자. $R$이 J-1일
필요충분조건은 모든 $\mathfrak p \in \text{Reg}(X)$에 대해
$V(\mathfrak p) \cap \text{Reg}(X)$가 $V(\mathfrak p)$의 공집합이 아닌
열린 부분집합을 포함하는 것이다.
\end{lemma}

\begin{proof}
이는 Topology, Lemma \ref{topology-lemma-characterize-open-Noetherian}과,
Algebra, Lemma
\ref{algebra-lemma-localization-of-regular-local-is-regular}에 의해
$\text{Reg}(X)$가 일반화에 대해 안정적이라는 사실에서 따른다.
\end{proof}

\begin{lemma}
\label{more-algebra-lemma-intersection-regular-with-closed}
$R$을 Noether 환이라 하자. $X = \Spec(R)$라 하자. 모든 소아이디얼
$\mathfrak p \subset R$에 대해 환 $R/\mathfrak p$가 J-0이라고
가정하자. 그러면 $R$은 J-1이다.
\end{lemma}

\begin{proof}
보조정리 \ref{more-algebra-lemma-J-1}의 판정법을 적용할 수 있음을 보이자.
$\mathfrak p \in \text{Reg}(X)$를 높이가 $r$인 소아이디얼이라 하자.
$\mathfrak pR_\mathfrak p$의 생성계로 가는
$f_1, \ldots, f_r \in \mathfrak p$를 택하자.
$\mathfrak p \in \text{Reg}(X)$이므로 $f_1, \ldots, f_r$은
$R_\mathfrak p$에서 정칙열로 간다. Algebra, Lemma
\ref{algebra-lemma-regular-ring-CM}을 보라.
% P02-E1446: frozen source uses singular “maps” with the plural parameter list; Korean uses natural agreement.
따라서 Algebra, Lemma
\ref{algebra-lemma-regular-sequence-in-neighbourhood}에 의해 어떤
$g \in R$, $g \not \in \mathfrak p$에 대해 $R$을 $R_g$로 바꾼 뒤
$f_1, \ldots, f_r$이 $R$의 정칙열이라고 가정해도 된다. 한 번 더
바꾸어 $f_1, \ldots, f_r$이 $\mathfrak p$를 생성한다고 가정해도 된다.
다음으로 $\mathfrak p \subset \mathfrak q$가
$(R/\mathfrak p)_\mathfrak q$가 정칙 국소환인 소아이디얼이라고 하자.
보조정리의 가정에 의해 $V(\mathfrak p)$ 안에는 그러한 소아이디얼들로
이루어진 공집합이 아닌 열린 부분집합이 존재한다. 따라서
$R_\mathfrak q$가 정칙임을 증명하면 충분하다.
$f_1, \ldots, f_r$은 $R_\mathfrak q$의 정칙열이고
$R_\mathfrak q/(f_1, \ldots, f_r)R_\mathfrak q$는 정칙임에 주의하자.
따라서 Algebra, Lemma \ref{algebra-lemma-regular-mod-x}에 의해
$R_\mathfrak q$는 정칙이다.
\end{proof}

\begin{lemma}
\label{more-algebra-lemma-J-0-goes-down}
$R \to S$를 환 사상이라 하자. 다음을 가정하자.
\begin{enumerate}
\item $R$은 Noether 정역이다.
\item $R \to S$는 단사이고 유한형이다.
\item $S$는 J-0인 정역이다.
\end{enumerate}
그러면 $R$은 J-0이다.
\end{lemma}

\begin{proof}
영이 아닌 어떤 $g \in S$에 대해 $S$를 $S_g$로 바꾸어 $S$가 정칙환이라고
가정해도 된다. 일반적 평탄성에 의해 $R \to S$도 충실히 평탄이라고
가정해도 된다. Algebra, Lemma
\ref{algebra-lemma-generic-flatness-Noetherian}을 보라. 그러면 Algebra,
Lemma \ref{algebra-lemma-descent-regular}에 의해 $R$은 정칙이다.
\end{proof}

\begin{lemma}
\label{more-algebra-lemma-J-0-goes-up}
$R \to S$를 환 사상이라 하자. 다음을 가정하자.
\begin{enumerate}
\item $R$은 J-0인 Noether 정역이다.
\item $R \to S$는 단사이고 유한형이다.
\item $S$는 정역이다.
\item 유도된 분수체 확대는 분리이다.
\end{enumerate}
그러면 $S$는 J-0이다.
\end{lemma}

\begin{proof}
$R$을 주 국소화로 바꾸어 $R$이 정칙환이라고 가정해도 된다.
Algebra, Lemma \ref{algebra-lemma-smooth-at-generic-point}에 의해 환 사상
$R \to S$는 $(0)$에서 매끄럽다. 따라서 $S$를 주 국소화로 바꾸어
$S$가 $R$ 위에서 매끄럽다고 가정해도 된다. 그러면 Algebra, Lemma
\ref{algebra-lemma-regular-goes-up}에 의해 $S$도 정칙이다.
\end{proof}

\begin{lemma}
\label{more-algebra-lemma-J-2}
$R$을 Noether 환이라 하자. 다음은 서로 동치이다.
\begin{enumerate}
\item $R$은 J-2이다.
\item 정역인 모든 유한형 $R$-대수는 J-0이다.
\item 모든 유한 $R$-대수는 J-1이다.
\item 모든 소아이디얼 $\mathfrak p$와 모든 유한 순수 비분리 확대
$L/\kappa(\mathfrak p)$에 대해, 정역이고 J-0이며 분수체가 $L$인
유한 $R$-대수 $R'$이 존재한다.
\end{enumerate}
\end{lemma}

\begin{proof}
함의 (1) $\Rightarrow$ (2) 및 (2) $\Rightarrow$ (4)는 명백하다.
J-1인 정역은 J-0임을 상기하자. 따라서 함의
(1) $\Rightarrow$ (3) 및 (3) $\Rightarrow$ (4)도 성립한다.

\medskip\noindent
$R \to S$를 유한형 환 사상이라 하고 $S$가 J-1임을 보이자.
보조정리 \ref{more-algebra-lemma-intersection-regular-with-closed}에 의해 $S$의 모든
소아이디얼 $\mathfrak q$에 대해 $S/\mathfrak q$가 J-0임을 증명하면
충분하다. 이로부터 (2) $\Rightarrow$ (1)을 알 수 있다.

\medskip\noindent
(4)를 가정하자. (2)가 성립함을 보여 증명을 끝내자. $R \to S$를
$S$가 정역인 유한형 환 사상이라 하자. $\mathfrak p = \Ker(R \to S)$라
하자. $K$를 $S$의 분수체라 하자. 수평 화살표가 유한 순수 비분리
체 확대이고 $K'/L$이 분리인 체의 도식
$$
\xymatrix{
K \ar[r] & K' \\
\kappa(\mathfrak p) \ar[u] \ar[r] & L \ar[u]
}
$$
이 존재한다. Algebra, Lemma
\ref{algebra-lemma-make-separably-generated}를 보라. (4)에서와 같이
$R' \subset L$을 택하고 $S \otimes_R R' \to K'$의 상을 $S'$이라
하자. 그러면 $S'$은 분수체가 $K'$인 정역이다. 따라서 $S'$은 보조정리
\ref{more-algebra-lemma-J-0-goes-up}과 $R'$의 선택에 의해 J-0이다. 이제 보조정리
\ref{more-algebra-lemma-J-0-goes-down}을 적용하면 원하는 대로 $S$가 J-0임을 알 수 있다.
\end{proof}



\section{정칙성과 미분}
\label{more-algebra-section-regularity-derivations}

\noindent
$R \to S$를 환 사상이라 하자. $D : R \to R$을 미분이라 하자.
$D' : S \to S$인 미분이 존재하여 도식
$$
\xymatrix{
S \ar[r]_{D'} & S \\
R \ar[u] \ar[r]^D & R \ar[u]
}
$$
이 가환이면 $D$가 $S$로 {\it 연장된다}고 한다.

\begin{lemma}
\label{more-algebra-lemma-derivation-extends}
$R$을 환이라 하자. $D : R \to R$을 미분이라 하자.
\begin{enumerate}
\item 임의의 아이디얼 $I \subset R$에 대해 미분 $D$는 $I$-진 완비화
위의 미분 $D^\wedge : R^\wedge \to R^\wedge$으로 표준적으로 연장된다.
\item 임의의 곱셈적 부분집합 $S \subset R$에 대해 미분 $D$는
$R$의 국소화 $S^{-1}R$로 유일하게 연장된다.
\end{enumerate}
$R \subset R'$이 유한형 환 확대이고, $R'$에서 비영인자인 어떤
$g \in R$에 대해 $R_g \cong R'_g$이면, 어떤 $N \geq 0$에 대해
$g^ND$는 $R'$로 연장된다.
\end{lemma}

\begin{proof}
(1)의 증명. $n \geq 2$에 대해 Leibniz 법칙으로부터
$D(I^n) \subset I^{n - 1}$을 얻는다. 따라서 $D$는 사상
$D_n : R/I^n \to R/I^{n - 1}$을 유도한다. 극한을 취하면
$D^\wedge$을 얻는다. $D^\wedge$이 미분이라는 확인은 생략한다.

\medskip\noindent
(2)의 증명. $D$를 $S^{-1}R$로 연장하려면
$D(r/s) = D(r)/s - rD(s)/s^2$로 두고 공리들을 확인하면 된다.

\medskip\noindent
마지막 명제의 증명. $x_1, \ldots, x_n \in R'$을 $R$ 위에서 $R'$의
생성원들이라 하자. $g^Nx_i \in R$이 되게 하는 $N$을 택하자.
$g^{N + 1}D$를 생각하자. (2)에 의해 이는 $R_g$로 연장된다. 더욱이
Leibniz 법칙과 위 연장의 구성에 의해
$$
g^{N + 1}D(x_i) = g^{N + 1}D(g^{-N} g^Nx_i) = -Ng^Nx_iD(g) +
gD(g^Nx_i)
$$
이고 두 항은 모두 $R$에 속한다. 이로부터
$$
g^{N + 1}D(x_1^{e_1} \ldots x_n^{e_n}) =
\sum e_i x_1^{e_1} \ldots x_i^{e_i - 1} \ldots x_n^{e_n} g^{N + 1}D(x_i)
$$
가 $R'$의 원소임을 얻는다. 따라서 $R'$의 모든 원소는
($R$의 계수를 갖는 $x_i$들의 단항식의 합으로 쓸 수 있으므로)
$g^{N + 1}D$에 의해 $R'$의 원소로 가고, 증명이 끝난다.
\end{proof}

\begin{lemma}
\label{more-algebra-lemma-quotient-regular}
\begin{slogan}
초곡면에 대한 올바른 Jacobi 판정법.
\end{slogan}
$R$을 정칙환이라 하자. $f \in R$이라 하자. $D(f)$가 $R/(f)$의
가역원인 미분 $D : R \to R$이 존재한다고 가정하자. 그러면
$R/(f)$는 정칙이다.
\end{lemma}

\begin{proof}
$R$이 극대 아이디얼 $\mathfrak m$과 잉여체 $\kappa$를 갖는
국소환인 경우를 증명하면 충분하다. 이 경우
$f \not \in \mathfrak m^2$임을 증명하면 충분하다. Algebra, Lemma
\ref{algebra-lemma-regular-ring-CM}을 보라. 그러나
$f \in \mathfrak m^2$이면 Leibniz 법칙에 의해
$D(f) \in \mathfrak m$이고, 이는 모순이다.
\end{proof}

\begin{lemma}
\label{more-algebra-lemma-quotient-sequence-regular}
$(R, \mathfrak m, \kappa)$를 정칙 국소환이라 하자. $m \geq 1$이라 하자.
$f_1, \ldots, f_m \in \mathfrak m$이라 하자. 행렬식
$\det_{1 \leq i, j \leq m}(D_i(f_j))$가 $R$의 가역원이 되게 하는
미분 $D_1, \ldots, D_m : R \to R$이 존재한다고 가정하자.
그러면 $R/(f_1, \ldots, f_m)$은 정칙이고 $f_1, \ldots, f_m$은
정칙열이다.
\end{lemma}

\begin{proof}
$f_1, \ldots, f_m$이 $\mathfrak m/\mathfrak m^2$에서
$\kappa$-선형 독립임을 증명하면 충분하다. Algebra, Lemma
\ref{algebra-lemma-regular-ring-CM}을 보라. 그러나 자명하지 않은
선형 관계가 있으면, 어떤 $a_i \in R$에 대해
$\sum a_i f_i \in \mathfrak m^2$이지만 모든 $a_i \in \mathfrak m$인
것은 아니다. 미분의
Leibniz 법칙에 의해 $D_i(\mathfrak m^2) \subset \mathfrak m$이고
$D_i(a_j f_j) \equiv a_j D_i(f_j) \bmod \mathfrak m$임에 주의하자.
따라서 이는
$$
\sum a_j D_i(f_j) \in \mathfrak m
$$
을 함의하며, 행렬식에 대한 가정과 모순이다.
\end{proof}

\begin{lemma}
\label{more-algebra-lemma-degree-p-extension-regular}
$R$을 정칙환이라 하자. $f \in R$이라 하자. $D(f)$가 $R$의 가역원인
미분 $D : R \to R$이 존재한다고 가정하자. 그러면 임의의 정수
$n \geq 1$에 대해 $R[z]/(z^n - f)$는 정칙이다. 더 일반적으로 임의의
$p \in \mathbf{Z}[z]$에 대해 $R[z]/(p(z) - f)$는 정칙이다.
\end{lemma}

\begin{proof}
Algebra, Lemma \ref{algebra-lemma-regular-goes-up}에 의해 $R[z]$는
정칙환이다. $z$를 영으로 보내는 $D$의 $R[z]$로의 연장에 보조정리
\ref{more-algebra-lemma-quotient-regular}를 적용하자. $D$는 정수 계수를 갖는 모든
다항식을 소멸시키고 $f$를 가역원으로 보내므로 이 논증을 적용할 수 있다.
\end{proof}

\begin{lemma}
\label{more-algebra-lemma-find-D}
$p$를 소수라 하자. $B$를 $B$에서 $p = 0$인 정역이라 하자.
$f \in B$가 $B$의 분수체에서 $p$제곱이 아닌 원소라 하자. $B$가
Noether 완비 국소환 위에서 유한형이면, $D(f)$가 영이 아닌 미분
$D : B \to B$가 존재한다.
\end{lemma}

\begin{proof}
유한형 환 사상 $R \to B$가 존재하는 Noether 완비 국소환 $R$을
택하자. 물론 $R$을 $B$ 안에서의 상으로 바꾸어, $R$이 표수가
$p > 0$인 정역이며 Noether 완비 국소환이라고 가정해도 된다.
Algebra, Lemma
\ref{algebra-lemma-complete-local-Noetherian-domain-finite-over-regular}에
의해 어떤 체 $k$와 정수 $n$에 대해 $R$을
$k[[x_1, \ldots, x_n]]$의 유한 확대라고 쓸 수 있다. 따라서 $R$을
$k[[x_1, \ldots, x_n]]$로 바꾸어도 된다. 다음으로 Algebra, Lemma
\ref{algebra-lemma-Noether-normalization-over-a-domain}을 사용하여
$R \to B$를
$$
R \subset R[y_1, \ldots, y_d] \subset B' \subset B
$$
로 인수분해한다. 여기서 $B'$은 $R[y_1, \ldots, y_d]$ 위에서 유한이고,
영이 아닌 어떤 $g \in R$에 대해 $B'_g \cong B_g$이다. 충분히 큰 정수
$N$에 대해 $f' = g^{pN} f \in B'$임에 주의하자. $f'$이 $B'$의
분수체에서 $p$제곱이 아님은 명백하다. $D'(f') \not = 0$인 미분
$D' : B' \to B'$을 찾을 수 있다면, 보조정리
\ref{more-algebra-lemma-derivation-extends}에 의해 어떤 $M > 0$에 대해
$D = g^MD'$가 $B$로 연장된다. 그러면
$D(f) = g^MD'(f) = g^MD'(g^{-pN}f') = g^{M - pN}D'(f')$은 영이 아니다.
따라서 $B$가
$A = k[[x_1, \ldots, x_n]][y_1, \ldots, y_m]$의 유한 확대인 경우에
보조정리를 증명하면 충분하다.

\medskip\noindent
$B$가 $A = k[[x_1, \ldots, x_n]][y_1, \ldots, y_m]$의 유한 확대라고
가정하자. $B$의 분수체를 $L$로 나타내자.
% P02-E1447: frozen source omits “by” in “Denote L the fraction field”; Korean uses natural notation syntax.
$\text{d}f$는 $\Omega_{L/\mathbf{F}_p}$에서 영이 아님에 주의하자.
Algebra, Lemma \ref{algebra-lemma-derivative-zero-pth-power}를 보라.
보조정리 \ref{more-algebra-lemma-power-series-ring-subfields}를 적용하여 다음 성질을
갖는 유한 지수의 부분체 $k' \subset k$를 찾는다. 즉,
$A' = k'[[x_1^p, \ldots, x_n^p]][y_1^p, \ldots, y_m^p]$이고 $K'$이
$A'$의 분수체일 때, 원소 $\text{d}f$는 $\Omega_{L/K'}$에서 영으로
가지 않는다. 따라서 $D'(f) \not = 0$인 $K'$-미분
$D' : L \to L$을 택할 수 있다. 구성상 $A' \subset A$와
$A \subset B$가 유한이므로 $A' \subset B$도 유한이다.
$A'$-가군으로서 $B$를 생성하는 $b_1, \ldots, b_t \in B$를 택하자.
그러면 어떤 $f_i, g_i \in B$, $g_i \not = 0$에 대해
$D'(b_i) = f_i/g_i$이다. $D = g_1 \ldots g_t D'$로 두면 증명이 끝난다.
\end{proof}

\begin{lemma}
\label{more-algebra-lemma-complete-Noetherian-domain-J-0}
$A$를 Noether 완비 국소 정역이라 하자. 그러면 $A$는 J-0이다.
\end{lemma}

\begin{proof}
Algebra, Lemma
\ref{algebra-lemma-complete-local-Noetherian-domain-finite-over-regular}에
의해 $A$가 $A_0$ 위에서 유한인 정칙 부분환 $A_0 \subset A$를 찾을 수
있다. 유도된 분수체 확대 $K/K_0$는 유한이다. $K/K_0$가 분리이면
보조정리 \ref{more-algebra-lemma-J-0-goes-up}에 의해 끝난다. 그렇지 않으면
$A_0$와 $A$의 표수는 $p > 0$이다. 임의의 부분 확대 $K/M/K_0$에
대해 분수체가 $M$인 유한 부분 확대 $A_0 \subset B \subset A$가
존재한다. 따라서 $[K : K_0]$에 대한 귀납법을 사용하여 $B$가 J-0이고
$K/M$이 자명하지 않은 부분 확대를 갖지 않게 하는
$A_0 \subset B \subset A$가 존재한다고 가정해도 된다.
이 경우 $K/M$이 분리이면 보조정리 \ref{more-algebra-lemma-J-0-goes-up}에 의해
$A$가 J-0임을 알 수 있다. 그렇지 않으면 어떤
$b_1, b_2 \in B$, $b_2 \not = 0$에 대해
$K = M[z]/(z^p - b_1/b_2)$이고 $b_1/b_2$는 $M$에서 $p$제곱이 아니다.
$az \in A$가 되게 하는 영이 아닌 $a \in A$를 택하자. $z$를
$b_2 a^p z$로 바꾸면 $z \in A$이고 $b \in B$가 $M$에서 $p$제곱이
아닐 때 $K = M[z]/(z^p - b)$를 얻는다. 보조정리
\ref{more-algebra-lemma-find-D}에 의해 $D(b) \not = 0$인 미분
$D : B \to B$를 찾을 수 있다. 보조정리
\ref{more-algebra-lemma-degree-p-extension-regular}를 적용하면 $A$의 소아이디얼 중
$B$의 정칙 소아이디얼 위에 놓이고 $D(b)$를 포함하지 않는 임의의
$\mathfrak p$에 대해 $A_\mathfrak p$가 정칙임을 알 수 있다.
$B$가 J-0이므로 $A$도 J-0이라는 결론을 얻는다.
\end{proof}

\begin{proposition}
\label{more-algebra-proposition-ubiquity-J-2}
다음 종류의 환들은 J-2이다.
\begin{enumerate}
\item 체,
\item Noether 완비 국소환,
\item $\mathbf{Z}$,
\item 차원이 $1$인 Noether 국소환,
\item 차원이 $1$인 Nagata 환,
\item 분수체의 표수가 영인 Dedekind 정역,
\item 위 환들 중 하나의 유한형 환 확대.
\end{enumerate}
\end{proposition}

\begin{proof}
(1), (3), (5), (6)의 경우에는 보조정리 \ref{more-algebra-lemma-J-2}의 조건 (4)를
확인하면 된다. 여기서는 $R$이 차원이 $1$인 Nagata 환인 경우만
확인하자. $\mathfrak p \subset R$을 소아이디얼이라 하고
$L/\kappa(\mathfrak p)$를 유한 순수 비분리 확대라 하자.
$\mathfrak p \subset R$이 극대 아이디얼이면 $R \to L$은 유한이고
$L$은 정칙환이므로 조건을 확인했다. $\mathfrak p \subset R$이 극소
소아이디얼이면 Nagata 조건에 의해 $L$ 안에서 $R$의 정수적 폐포
$R' \subset L$은 $R$ 위에서 유한이다. 그러면 $R'$은 차원이 $1$인
정규 정역이고(Algebra, Lemma \ref{algebra-lemma-integral-dim-up}),
따라서 정칙이다(Algebra, Lemma \ref{algebra-lemma-criterion-normal}).
이 경우에도 조건을 확인했다.

\medskip\noindent
(2)의 경우에는 보조정리 \ref{more-algebra-lemma-J-2}의 조건 (3)을 사용한다.
$R$을 Noether 완비 국소환이라 하자. $R \to R'$이 유한이면 $R'$은
Noether 완비 국소환들의 곱임에 주의하자. Algebra, Lemma
\ref{algebra-lemma-quotient-complete-local}을 보라. 따라서 보조정리
\ref{more-algebra-lemma-intersection-regular-with-closed}에 의해 Noether 완비
국소 정역이 J-0임을 증명하면 충분하고, 이는 보조정리
\ref{more-algebra-lemma-complete-Noetherian-domain-J-0}이다.

\medskip\noindent
(4)의 경우에도 보조정리 \ref{more-algebra-lemma-J-2}의 조건 (3)을 사용한다.
실제로 $R$이 차원이 $1$인 국소 Noether 환이고 $R \to R'$이 유한이면
$\Spec(R')$은 유한하다. 정칙 자취는 일반화에 대해 안정적이므로
$R'$은 J-1이다.
\end{proof}





\section{형식적 매끄러움과 정칙성}
\label{more-algebra-section-fs-regular}

\noindent
이 절의 제목은 Andr\'e 정리 \ref{more-algebra-theorem-fs-regular}를 가리킨다.

\begin{lemma}
\label{more-algebra-lemma-lift-derivation-through-fs}
$A \to B$를 Noether 국소환의 국소 준동형이라 하자.
$D : A \to A$를 미분이라 하자. $B$가 완비이고 $A \to B$가
$\mathfrak m_B$-진 위상에서 형식적으로 매끄럽다고 가정하자.
그러면 $D$의 연장 $D' : B \to B$가 존재한다.
\end{lemma}

\begin{proof}
$B$ 위의 쌍대수환을 $B[\epsilon] = B[x]/(x^2)$로 나타내자.
% P02-E1448: frozen source omits “by” in “Denote B[epsilon] ... the ring”; Korean uses natural notation syntax.
환 사상 $\psi : A \to B[\epsilon]$, $a \mapsto a + \epsilon D(a)$를
생각하자. 가환 도식
$$
\xymatrix{
B \ar[r]_1 & B \\
A \ar[u] \ar[r]^\psi & B[\epsilon] \ar[u]
}
$$
을 생각하자. 보조정리 \ref{more-algebra-lemma-lift-continuous}와 $B/A$의 형식적
매끄러움에 대한 가정에 의해 이 도식에 들어맞는 사상
$\varphi : B \to B[\epsilon]$을 찾을 수 있다.
$\varphi(b) = b + \epsilon D'(b)$로 쓰자. 그러면 $D' : B \to B$는
원하는 연장이다.
\end{proof}

\begin{proposition}
\label{more-algebra-proposition-fs-regular}
$A \to B$를 Noether 완비 국소환의 국소 준동형이라 하자. $k$를
$A$의 잉여체라 하고 $\overline{B} = B \otimes_A k$를 특수 올이라 하자.
다음은 서로 동치이다.
\begin{enumerate}
\item $A \to B$는 정칙이다.
\item $A \to B$는 평탄이고 $\overline{B}$는 $k$ 위에서 기하학적으로
정칙이다.
\item $A \to B$는 평탄이고 $k \to \overline{B}$는
$\mathfrak m_{\overline{B}}$-진 위상에서 형식적으로 매끄럽다.
\item $A \to B$는 $\mathfrak m_B$-진 위상에서 형식적으로 매끄럽다.
\end{enumerate}
\end{proposition}

\begin{proof}
Proposition \ref{more-algebra-proposition-fs-flat-fibre-fs}에서 (2), (3), (4)의
동치를 보았다. (1)이 (2)를 함의함은 명백하다. 따라서 서로 동치인
조건 (2), (3), (4)가 성립한다고 가정하고 (1)을 증명한다.

\medskip\noindent
$\mathfrak p$를 $A$의 소아이디얼이라 하자.
$B \otimes_A \kappa(\mathfrak p)$가 $\kappa(\mathfrak p)$ 위에서
기하학적으로 정칙임을 보이자. 보조정리 \ref{more-algebra-lemma-base-change-fs}에
의해 $A$를 $A/\mathfrak p$로, $B$를 $B/\mathfrak pB$로 바꾸어도 된다.
따라서 $A$가 정역이고 $\mathfrak p = (0)$이라고 가정해도 된다.

\medskip\noindent
Algebra, Lemma
\ref{algebra-lemma-complete-local-Noetherian-domain-finite-over-regular}에서와
같이 $A_0 \subset A$를 택하자. 그 보조정리에서 명시한 모든 성질을
별도의 언급 없이 사용한다. $A_0 \to A$는 잉여체 위의 동형을 유도하고
$B/\mathfrak m_A B$는 $A/\mathfrak m_A$ 위에서 기하학적으로 정칙이므로,
$A_0 \to C$가 $\mathfrak m_C$-진 위상에서 형식적으로 매끄럽고
$B = C \otimes_{A_0} A$인 도식
$$
\xymatrix{
C \ar[r] & B \\
A_0 \ar[r] \ar[u] & A \ar[u]
}
$$
을 찾을 수 있다. Remark \ref{more-algebra-remark-what-does-it-mean}을 보라.
$A_0 \to A$가 유한이므로 텐서곱에서 완비화는 필요하지 않다.
Algebra, Lemma \ref{algebra-lemma-completion-tensor}를 보라.
따라서 $C \otimes_{A_0} K_0$가 $A_0$의 분수체 $K_0$ 위에서
기하학적으로 정칙인 대수임을 보이면 충분하다.

\medskip\noindent
앞 문단의 결론은 $k$가 체일 때
$A = k[[x_1, \ldots, x_n]]$이거나 $\Lambda$가 Cohen 환일 때
$A = \Lambda[[x_1, \ldots, x_n]]$이라고 가정해도 된다는 것이다.
이 경우 Algebra, Lemma
\ref{algebra-lemma-flat-over-regular-with-regular-fibre}에 의해 $B$는
정칙환이다. 따라서 $B \otimes_A K$도 정칙환이다. 여기서 $K$는
$A$의 분수체이다. 그러므로 $K$의 표수가 영이면 끝난다.

\medskip\noindent
이제 $A = k[[x_1, \ldots, x_n]]$이고 $k$가 표수가 $p > 0$인 체인
경우만 남는다. $L/K$를 유한 순수 비분리 체 확대라 하자.
$B \otimes_A L$이 정칙임을 $[L : K]$에 대한 귀납법으로 보이자.
기본 경우는 $L = K$이고, 이는 위에서 보았다. $K \subset M \subset L$을
$L$이 $M$의 어떤 원소 $f \in M$의 $p$제곱근을 첨가하여 얻는 차수
$p$의 확대가 되게 하는 부분체라 하자. $A'$을 $M$의 유한
$A$-부분대수로 택하되 그 분수체는 $M$이라 하자. 분모를 제거하여
$f \in A'$이라고 가정한다.
$A'' = A'[z]/(z^p -f)$라 두고, $A' \subset A''$이 유한이며 $A''$의
분수체가 $L$임에 주의하자. 귀납 가정에 의해 환
$B \otimes_A M$은 정칙이다.
% P02-E1449: frozen source says “B tensor M ring is regular,” with an extraneous “ring”; Korean uses natural syntax.
등식
$$
B \otimes_A L = B \otimes_A M[z]/(z^p - f)
$$
이 성립한다. 보조정리 \ref{more-algebra-lemma-find-D}에 의해
$D(f) \not = 0$인 미분 $D : A' \to A'$이 존재한다. 보조정리
\ref{more-algebra-lemma-descent-fs}에 의해 $A' \to B \otimes_A A'$은
$\mathfrak m$-진 위상에서 형식적으로 매끄러우므로, 보조정리
\ref{more-algebra-lemma-lift-derivation-through-fs}를 사용하여 $D$를 미분
$D' : B \otimes_A A' \to B \otimes_A A'$으로 연장할 수 있다.
$D(f)$가 $A' \subset M$에서 영이 아니므로
$D'(f) = D(f)$는 $B \otimes_A M$의 가역원임에 주의하자. 따라서
보조정리 \ref{more-algebra-lemma-degree-p-extension-regular}에 의해
$B \otimes_A L$은 정칙이고, 증명이 끝난다.
\end{proof}

\begin{theorem}[Andr\'e]
\label{more-algebra-theorem-fs-regular}
\begin{reference}
함의 (4) $\Rightarrow$ (1)은 \cite{Andre-smooth}의 주결과이다.
\end{reference}
$A \to B$를 Noether 국소환의 국소 준동형이라 하자. $k$를 $A$의
잉여체라 하고 $\overline{B} = B \otimes_A k$를 특수 올이라 하자.
$A \to A^\wedge$이 정칙이라고 가정하자\footnote{예를 들어 $A$가
G-환이거나 (준)우수환인 경우이다. 절 
\ref{more-algebra-section-G-ring}과
\ref{more-algebra-section-excellent}를 보라.}. 다음은 서로 동치이다.
\begin{enumerate}
\item $A \to B$는 정칙이다.
\item $A \to B$는 평탄이고 $\overline{B}$는 $k$ 위에서 기하학적으로
정칙이다.
\item $A \to B$는 평탄이고 $k \to \overline{B}$는
$\mathfrak m_{\overline{B}}$-진 위상에서 형식적으로 매끄럽다.
\item $A \to B$는 $\mathfrak m_B$-진 위상에서 형식적으로 매끄럽다.
\end{enumerate}
\end{theorem}

\begin{proof}
Proposition \ref{more-algebra-proposition-fs-flat-fibre-fs}에서 (2), (3), (4)의
동치를 보았다. (1)이 (2)를 함의함은 명백하다. 따라서 (4)가
성립한다고 가정하고 (1)을 증명한다.

\medskip\noindent
보조정리 \ref{more-algebra-lemma-formally-smooth-completion}에 의해
$A^\wedge \to B^\wedge$은 형식적으로 매끄럽다. Proposition
\ref{more-algebra-proposition-fs-regular}에 의해 $A^\wedge \to B^\wedge$은 정칙이다.
가정에 의해 $A \to A^\wedge$이 정칙이므로, 보조정리
\ref{more-algebra-lemma-regular-composition}에 의해 $A \to B^\wedge$은 정칙이다.
$B \to B^\wedge$은 충실히 평탄이므로 보조정리
\ref{more-algebra-lemma-regular-permanence}에 의해 $A \to B$가 정칙이라는 결론을
얻는다.
\end{proof}




\section{G-환}
\label{more-algebra-section-G-ring}

\noindent
$A$를 Noether 국소환이라 하자. 절
\ref{more-algebra-section-permanence-completion}에서 $A$의 일부 성질은 $A$의
완비화 $A^\wedge$에 반영되지만 전부가 그렇지는 않음을 보았다.
이를 더 연구하기 위해 몇 가지 용어를 도입한다. $A$의 소아이디얼
$\mathfrak q$에 대하여 올환
% P02-E0883: frozen source tensors over undefined A/q instead of A/mathfrak q; Korean preserves the display.
$$
A^\wedge \otimes_A \kappa(\mathfrak q) =
(A^\wedge)_\mathfrak q/\mathfrak q(A^\wedge)_\mathfrak q =
(A/\mathfrak q)^\wedge \otimes_{A/q} \kappa(\mathfrak q)
$$
을 $A$의 {\it 형식적 올}이라 한다. 형식적 올은
$\kappa(\mathfrak q)$ 위의 대수로 생각한다. 따라서 $A \to A^\wedge$이
정칙 환 준동형일 필요충분조건은 모든 형식적 올이 기하학적으로
정칙인 대수인 것이다.

\begin{definition}
\label{more-algebra-definition-G-ring}
환 $R$이 Noether이고 $R$의 모든 소아이디얼 $\mathfrak p$에 대하여
환 사상 $R_\mathfrak p \to (R_\mathfrak p)^\wedge$이 정칙이면
$R$을 {\it G-환}이라 한다.
\end{definition}

\noindent
위의 논의에 의해 $R$이 G-환일 필요충분조건은 모든 국소환
$R_\mathfrak p$의 형식적 올이 기하학적으로 정칙인 것이다.
$\mathbf{Q} \subset R$이면 형식적 올이 정칙인지 확인하는 것으로
충분함에 주의하자. 다음 보조정리는 G-환 조건의 또 다른 표현을 준다.

\begin{lemma}
\label{more-algebra-lemma-check-G-ring-easy}
$R$을 Noether 환이라 하자. $R$이 G-환일 필요충분조건은 모든
소아이디얼 쌍 $\mathfrak q \subset \mathfrak p \subset R$에 대하여 대수
$$
(R/\mathfrak q)_\mathfrak p^\wedge \otimes_{R/\mathfrak q} \kappa(\mathfrak q)
$$
가 $\kappa(\mathfrak q)$ 위에서 기하학적으로 정칙인 것이다.
\end{lemma}

\begin{proof}
$\kappa(\mathfrak q)$ 위의 대수로서
$$
R_\mathfrak p^\wedge \otimes_R \kappa(\mathfrak q) =
(R/\mathfrak q)_\mathfrak p^\wedge \otimes_{R/\mathfrak q} \kappa(\mathfrak q)
$$
가 성립한다는 사실에서 곧바로 따른다.
\end{proof}

\begin{lemma}
\label{more-algebra-lemma-G-ring-goes-up-quasi-finite}
$R \to R'$을 Noether 환의 유한형 사상이라 하고
$$
\xymatrix{
\mathfrak q' \ar[r] & \mathfrak p' \ar[r] & R' \\
\mathfrak q \ar[r] \ar@{-}[u] &
\mathfrak p \ar[r] \ar@{-}[u] & R \ar[u]
}
$$
가 소아이디얼들로 이루어져 있다고 하자. $R \to R'$이
$\mathfrak p'$에서 준유한이라고 가정하자.
\begin{enumerate}
\item 형식적 올 $R_\mathfrak p^\wedge \otimes_R \kappa(\mathfrak q)$가
$\kappa(\mathfrak q)$ 위에서 기하학적으로 정칙이면, 형식적 올
$(R'_{\mathfrak p'})^\wedge \otimes_{R'} \kappa(\mathfrak q')$은
$\kappa(\mathfrak q')$ 위에서 기하학적으로 정칙이다.
\item $R_\mathfrak p$의 형식적 올들이 기하학적으로 정칙이면,
$R'_{\mathfrak p'}$의 형식적 올들도 기하학적으로 정칙이다.
\item $R \to R'$이 준유한이고 $R$이 G-환이면 $R'$도 G-환이다.
\end{enumerate}
\end{lemma}

\begin{proof}
(1) $\Rightarrow$ (2) $\Rightarrow$ (3)은 명백하다.
$R_\mathfrak p^\wedge \otimes_R \kappa(\mathfrak q)$가
$\kappa(\mathfrak q)$ 위에서 기하학적으로 정칙이라고 가정하자.
Algebra, Lemma \ref{algebra-lemma-completion-at-quasi-finite-prime}에
의해 어떤 $R_\mathfrak p^\wedge$-대수 $B$에 대하여
$$
R_\mathfrak p^\wedge \otimes_R R'
=
(R'_{\mathfrak p'})^\wedge \times B
$$
임을 안다. 따라서 $R'_{\mathfrak p'} \to
(R'_{\mathfrak p'})^\wedge$은 사상
$R_\mathfrak p \to R_\mathfrak p^\wedge$의 기저변환의 한 인자이다.
그러므로 $(R'_{\mathfrak p'})^\wedge \otimes_{R'}
\kappa(\mathfrak q')$은
$$
R_\mathfrak p^\wedge \otimes_R R' \otimes_{R'} \kappa(\mathfrak q') =
R_\mathfrak p^\wedge \otimes_R \kappa(\mathfrak q)
\otimes_{\kappa(\mathfrak q)} \kappa(\mathfrak q').
$$
의 한 인자이다. 기저체의 확대가 기하학적 정칙성을 보존하므로
결론이 따른다. Algebra, Lemma
\ref{algebra-lemma-geometrically-regular}를 보라.
\end{proof}

\begin{lemma}
\label{more-algebra-lemma-check-G-ring}
$R$을 Noether 환이라 하자. $R$이 G-환일 필요충분조건은 모든 유한
자유 환 사상 $R \to S$에 대하여 $S$의 형식적 올들이 정칙환인 것이다.
\end{lemma}

\begin{proof}
모든 유한 자유 환 사상 $R \to S$에 대하여 환 $S$의 형식적 올들이
정칙이라고 가정하자. $\mathfrak q \subset \mathfrak p \subset R$을
소아이디얼들이라 하고, $\kappa(\mathfrak q) \subset L$을 유한 순수
비분리 확대라 하자. $R$이 G-환임을 보이려면 환
$$
R_\mathfrak p^\wedge \otimes_R \kappa(\mathfrak q)
\otimes_{\kappa(\mathfrak q)} L
$$
이 정칙임을 보이면 충분하다. $\mathfrak q' = \mathfrak qR'$이
소아이디얼이고 $\kappa(\mathfrak q')$이 $\kappa(\mathfrak q)$ 위에서
$L$과 동형이 되도록 하는 유한 자유 확대 $R \to R'$을 택하자.
Algebra, Lemma
\ref{algebra-lemma-finite-free-given-residue-field-extension}을 보라.
Algebra, Lemma \ref{algebra-lemma-completion-finite-extension}에 의해
$$
R_\mathfrak p^\wedge \otimes_R R' = \prod (R'_{\mathfrak p_i'})^\wedge
$$
이다. 여기서 $\mathfrak p_i'$들은 $\mathfrak p$ 위에 놓이는 $R'$의
소아이디얼들이다. 따라서
$$
R_\mathfrak p^\wedge \otimes_R \kappa(\mathfrak q)
\otimes_{\kappa(\mathfrak q)} L =
R_\mathfrak p^\wedge \otimes_R R'
\otimes_{R'} \kappa(\mathfrak q')
=
\prod (R'_{\mathfrak p_i'})^\wedge
\otimes_{R'_{\mathfrak p'_i}} \kappa(\mathfrak q')
$$
를 얻는다. 가정에 의해 우변의 환들은 정칙이므로 좌변의 환도
정칙이다. 따라서 $R$은 G-환이다. 역은 보조정리
\ref{more-algebra-lemma-G-ring-goes-up-quasi-finite}에서 따른다.
\end{proof}

\begin{lemma}
\label{more-algebra-lemma-helper-G-ring}
$k$를 표수가 $p$인 체라 하자.
$A = k[[x_1, \ldots, x_n]][y_1, \ldots, y_m]$라 하고, $A$의 분수체를
$K$로 나타내자.
% P02-E1450: frozen source says “denote K the fraction field,” omitting “by”; Korean uses natural notation syntax.
$\mathfrak p \subset A$를 소아이디얼이라 하자. 그러면
$A_\mathfrak p^\wedge \otimes_A K$는 $K$ 위에서 기하학적으로 정칙이다.
\end{lemma}

\begin{proof}
$L/K$를 유한 순수 비분리 체 확대라 하자.
$A_\mathfrak p^\wedge \otimes L$이 정칙임을 $[L : K]$에 대한 귀납법으로
보인다. 기본 경우는 $L = K$이다. $A$가 정칙이므로
$A_\mathfrak p^\wedge$은 정칙이고(보조정리
\ref{more-algebra-lemma-completion-regular}), 따라서 국소화
$A_\mathfrak p^\wedge \otimes K$도 정칙이다.
$K \subset M \subset L$을 부분체라 하되, $L$은 원소
$f \in M$의 $p$제곱근을 첨가하여 얻는 $M$의 차수 $p$인 확대라고 하자.
$B$를 $M$의 유한 $A$-부분대수로 택하되 그 분수체는 $M$이라 하자.
분모를 제거하여 $f \in B$라고 가정해도 된다.
$C = B[z]/(z^p -f)$라 두고, $B \subset C$가 유한이며 $C$의 분수체가
$L$임에 주의하자. $A \subset B \subset C$는 유한이고 $L/M/K$는
순수 비분리이므로, $B$ 또는 $C$의 모든 원소는 어떤 거듭제곱이
$A$에 속한다. 따라서 $\mathfrak p$ 위에 놓이는 소아이디얼
$\mathfrak r \subset B$와 $\mathfrak q \subset C$는 각각 유일하다.
다음이 성립함에 주의하자.
$$
A_\mathfrak p^\wedge \otimes_A M = B_\mathfrak r^\wedge \otimes_B M
$$
Algebra, Lemma \ref{algebra-lemma-completion-finite-extension}을 보라.
귀납 가정에 의해 이 환은 정칙이다. 같은 방식으로
% P02-E0884: frozen source completes C at r although r is the prime of B and q is the prime of C; Korean preserves the display.
$$
A_\mathfrak p^\wedge \otimes_A L =
C_\mathfrak r^\wedge \otimes_C L =
B_\mathfrak r^\wedge \otimes_B M[z]/(z^p - f)
$$
이다. 마지막 등식은 $C = B[z]/(z^p - f)$의 완비화가
$B_\mathfrak r^\wedge[z]/(z^p -f)$와 같기 때문이다.
보조정리 \ref{more-algebra-lemma-find-D}에 의해 $D(f) \not = 0$인 미분
$D : B \to B$가 존재한다. 다시 말해 $g = D(f)$는 $M$의 가역원이다!
보조정리 \ref{more-algebra-lemma-derivation-extends}에 의해 $D$는
$B_\mathfrak r$, $B_\mathfrak r^\wedge$, 그리고
$B_\mathfrak r^\wedge \otimes_B M$의 미분으로 차례로 연장된다
(각각 국소화, 완비화, 국소화를 통해 연장한다). 연장이므로 끝에서
얻은 $B_\mathfrak r^\wedge \otimes_B M$의 미분은 $f$를 $g$로 보내며,
$g$는 환 $B_\mathfrak r^\wedge \otimes_B M$의 가역원이다. 따라서
보조정리 \ref{more-algebra-lemma-degree-p-extension-regular}에 의해
$A_\mathfrak p^\wedge \otimes_A L$은 정칙이고, 증명이 끝난다.
\end{proof}

\begin{proposition}
\label{more-algebra-proposition-Noetherian-complete-G-ring}
Noether 완비 국소환은 G-환이다.
\end{proposition}

\begin{proof}
$A$를 Noether 완비 국소환이라 하자. 보조정리
\ref{more-algebra-lemma-check-G-ring-easy}에 의해 $B = A/\mathfrak q$의 최소
소아이디얼 $(0)$ 위, 즉 $B$의 최소 소아이디얼 위의 형식적 올들이
기하학적으로 정칙인지 확인하면
충분하다. 따라서 $A$가 정역이라고 가정해도 되고, $A$의 최소
소아이디얼 $(0)$ 위의 형식적 올에 대한 조건을 확인하면 충분하다.
$K$를 $A$의 분수체라 하자.

\medskip\noindent
Algebra, Lemma
\ref{algebra-lemma-complete-local-Noetherian-domain-finite-over-regular}에
의해 $A$가 어떤 정칙 완비 국소 부분환 $A_0 \subset A$ 위에서 유한이
되도록 $A_0$를 택할 수 있다. 더욱이 $A_0$가 체 또는 Cohen 환 위의
멱급수환이라고 가정할 수 있다. 보조정리
\ref{more-algebra-lemma-G-ring-goes-up-quasi-finite}에 의해 $A_0$에 대해 결과를
증명하는 것으로 충분하다.

\medskip\noindent
$A$가 체 또는 Cohen 환 위의 멱급수환이라고 가정하자. $A$가 정칙이므로
국소화 $A_\mathfrak p$들은 정칙이다(Algebra, Definition
\ref{algebra-definition-regular}과 그 앞의 논의를 보라).
따라서 보조정리 \ref{more-algebra-lemma-completion-regular}에 의해 완비화
$A_\mathfrak p^\wedge$들도 정칙이다. 그러므로 올
$A_{\mathfrak p}^\wedge \otimes_A K$는 $A_\mathfrak p^\wedge$의
국소화로서 역시 정칙이다. 따라서 $K$의 표수가 $0$이면 끝난다.
양의 표수인 경우는 $A = k[[x_1, \ldots, x_d]]$인 경우이며, 이는
보조정리 \ref{more-algebra-lemma-helper-G-ring}의 특수한 경우이다.
\end{proof}

\begin{lemma}
\label{more-algebra-lemma-check-G-ring-maximal-ideals}
$R$을 Noether 환이라 하자. $R$이 G-환일 필요충분조건은 $R$의 모든
극대 아이디얼 $\mathfrak m$에 대하여 $R_\mathfrak m$의 형식적 올들이
기하학적으로 정칙인 것이다.
\end{lemma}

\begin{proof}
$R$의 모든 극대 아이디얼 $\mathfrak m$에 대하여
$R_\mathfrak m \to R_\mathfrak m^\wedge$이 정칙이라고 가정하자.
$\mathfrak p$를 $R$의 소아이디얼이라 하고, 이를 포함하는 극대 아이디얼
$\mathfrak p \subset \mathfrak m$을 택하자. $R_\mathfrak m \to
R_\mathfrak m^\wedge$은 충실히 평탄이므로
$\mathfrak pR_\mathfrak m$ 위에 놓이는 $R_\mathfrak m^\wedge$의
소아이디얼 $\mathfrak p'$을 택할 수 있다. 가환도식
$$
\xymatrix{
R_\mathfrak m^\wedge \ar[r] &
(R_\mathfrak m^\wedge)_{\mathfrak p'} \ar[r] &
(R_\mathfrak m^\wedge)_{\mathfrak p'}^\wedge
\\
R_\mathfrak m \ar[u] \ar[r] & R_\mathfrak p \ar[u] \ar[r] &
R_\mathfrak p^\wedge \ar[u]
}
$$
을 생각하자. 가정에 의해 환 사상
$R_\mathfrak m \to R_\mathfrak m^\wedge$은 정칙이다. Proposition
\ref{more-algebra-proposition-Noetherian-complete-G-ring}에 의해
$(R_\mathfrak m^\wedge)_{\mathfrak p'} \to
(R_\mathfrak m^\wedge)_{\mathfrak p'}^\wedge$은 정칙이다. 국소화
$R_\mathfrak m^\wedge \to (R_\mathfrak m^\wedge)_{\mathfrak p'}$도
정칙이다. 따라서 보조정리 \ref{more-algebra-lemma-regular-composition}에 의해
$R_\mathfrak m \to (R_\mathfrak m^\wedge)_{\mathfrak p'}^\wedge$은
정칙이다. 이 사상은 국소화 $R_\mathfrak p$를 통해 분해되므로 환 사상
$R_\mathfrak p \to (R_\mathfrak m^\wedge)_{\mathfrak p'}^\wedge$도
정칙이다. 이제 보조정리 \ref{more-algebra-lemma-regular-permanence}를 적용하면
$R_\mathfrak p \to R_\mathfrak p^\wedge$이 정칙임을 알 수 있다.
\end{proof}

\begin{lemma}
\label{more-algebra-lemma-henselization-G-ring}
$R$을 G-환인 Noether 국소환이라 하자. 그러면 헨젤화 $R^h$와 엄밀
헨젤화 $R^{sh}$은 G-환이다.
\end{lemma}

\begin{proof}
보조정리 \ref{more-algebra-lemma-check-G-ring-maximal-ideals}의 판정법을 사용한다.
$\mathfrak q \subset R^h$를 소아이디얼이라 하고
$\mathfrak p = R \cap \mathfrak q$라 두자.
$\mathfrak q_1 = \mathfrak q$라 두고
$\mathfrak q_2, \ldots, \mathfrak q_t$를 $\mathfrak p$ 위에 놓이는
$R^h$의 나머지 소아이디얼들이라 하자. 보조정리
\ref{more-algebra-lemma-fibres-henselization}에 의해
$R^h \otimes_R \kappa(\mathfrak p) =
\prod\nolimits_{i = 1, \ldots, t} \kappa(\mathfrak q_i)$이다.
$(R^h)^\wedge = R^\wedge$이라는 사실(보조정리
\ref{more-algebra-lemma-henselization-noetherian})을 사용하면
$$
\prod\nolimits_{i = 1, \ldots, t}
(R^h)^\wedge \otimes_{R^h} \kappa(\mathfrak q_i) =
(R^h)^\wedge \otimes_{R^h} (R^h \otimes_R \kappa(\mathfrak p)) =
R^\wedge \otimes_R \kappa(\mathfrak p)
$$
를 얻는다. 따라서 가정에 의해
$(R^h)^\wedge \otimes_{R^h} \kappa(\mathfrak q_i)$는
$\kappa(\mathfrak p)$ 위에서 기하학적으로 정칙이다.
$\kappa(\mathfrak q_i)$는 $\kappa(\mathfrak p)$의 분리 대수적 확대이므로
Algebra, Lemma
\ref{algebra-lemma-geometrically-regular-over-separable-algebraic}에 의해
$(R^h)^\wedge \otimes_{R^h} \kappa(\mathfrak q_i)$는
$\kappa(\mathfrak q_i)$ 위에서 기하학적으로 정칙이다.

\medskip\noindent
$\mathfrak r \subset R^{sh}$을 소아이디얼이라 하고
$\mathfrak p = R \cap \mathfrak r$라 두자.
$\mathfrak r_1 = \mathfrak r$라 두고
$\mathfrak r_2, \ldots, \mathfrak r_s$를 $\mathfrak p$ 위에 놓이는
$R^{sh}$의 나머지 소아이디얼들이라 하자. 보조정리
\ref{more-algebra-lemma-fibres-henselization}에 의해
$R^{sh} \otimes_R \kappa(\mathfrak p) =
\prod\nolimits_{i = 1, \ldots, s} \kappa(\mathfrak r_i)$이다.
그러면
$$
\prod\nolimits_{i = 1, \ldots, s}
(R^{sh})^\wedge \otimes_{R^{sh}} \kappa(\mathfrak r_i) =
(R^{sh})^\wedge \otimes_{R^{sh}} (R^{sh} \otimes_R \kappa(\mathfrak p)) =
(R^{sh})^\wedge \otimes_R \kappa(\mathfrak p)
$$
를 얻는다. 보조정리 \ref{more-algebra-lemma-henselization-noetherian}에 의해
$R^\wedge \to (R^{sh})^\wedge$은
$\mathfrak m_{(R^{sh})^\wedge}$-진 위상에서 형식적으로 매끄럽다.
따라서 Proposition \ref{more-algebra-proposition-fs-regular}에 의해
$R^\wedge \to (R^{sh})^\wedge$은 정칙이다. 가정에 의해
$R^\wedge \otimes_R \kappa(\mathfrak p)$가
$\kappa(\mathfrak p)$ 위에서 정칙이므로, 보조정리
\ref{more-algebra-lemma-regular-composition}에 의해
$(R^{sh})^\wedge \otimes_{R^{sh}} \kappa(\mathfrak r_i)$는
$\kappa(\mathfrak p)$ 위에서 정칙이다.
$\kappa(\mathfrak r_i)$는 $\kappa(\mathfrak p)$의 분리 대수적
확대이므로 Algebra, Lemma
\ref{algebra-lemma-geometrically-regular-over-separable-algebraic}에 의해
$(R^{sh})^\wedge \otimes_{R^{sh}} \kappa(\mathfrak r_i)$는
$\kappa(\mathfrak r_i)$ 위에서 기하학적으로 정칙이다.
\end{proof}

\begin{lemma}
\label{more-algebra-lemma-another-helper-G-ring}
$p$를 소수라 하자. $A$를 분수체 $K$의 표수가 $p$인 Noether 완비
국소 정역이라 하자. $\mathfrak q \subset A[x]$를 $A$의 극대 아이디얼
위에 놓이는 극대 아이디얼이라 하고, $(0) \subset A$ 위에 놓이는
소아이디얼 $(0) \not = \mathfrak r \subset \mathfrak q$를 택하자. 그러면
$A[x]_\mathfrak q^\wedge \otimes_{A[x]} \kappa(\mathfrak r)$는
$\kappa(\mathfrak r)$ 위에서 기하학적으로 정칙이다.
\end{lemma}

\begin{proof}
$K \subset \kappa(\mathfrak r)$가 유한임에 주의하자. 따라서 유한 순수
비분리 확대 $L/\kappa(\mathfrak r)$가 주어지면, 사상
$\kappa(\mathfrak r) \otimes_A B$가 $L$ 위로 전사인 Noether 완비 국소
정역의 유한 확대 $A \subset B$가 존재한다. 실제로 분수체가 $L$인
유한 $A$-부분대수 $B \subset L$을 택하면 된다.
% P02-E1451: frozen source says “Denote r-prime ... the kernel,” omitting “by”; Korean uses natural notation syntax.
$\mathfrak r' \subset B[x]$를 사상
$B[x] = A[x] \otimes_A B \to \kappa(\mathfrak r) \otimes_A B \to L$의
핵이라 하자. 그러면 $\kappa(\mathfrak r') = L$이고
$$
A[x]_\mathfrak q^\wedge \otimes_{A[x]} L =
A[x]_\mathfrak q^\wedge \otimes_{A[x]} B[x]
\otimes_{B[x]} \kappa(\mathfrak r') =
\prod B[x]_{\mathfrak q_i}^\wedge \otimes_{B[x]} \kappa(\mathfrak r')
$$
이다. 여기서 $\mathfrak q_1, \ldots, \mathfrak q_t$는 $\mathfrak q$
위에 놓이는 $B[x]$의 소아이디얼들이다. Algebra, Lemma
\ref{algebra-lemma-completion-finite-extension}을 보라. 따라서 환들
$B[x]_{\mathfrak q_i}^\wedge \otimes_{B[x]} \kappa(\mathfrak r')$이
정칙임을 증명하면 충분하다. 이는 $K = \kappa(\mathfrak r)$인 특수한
경우에 $A[x]_\mathfrak q^\wedge \otimes_{A[x]} \kappa(\mathfrak r)$가
정칙임을 보이는 것으로 환원된다.

\medskip\noindent
$K = \kappa(\mathfrak r)$라고 가정하자. 이 경우 어떤 $f \in K$에 대하여
$\mathfrak r K[x]$가 $x - f$로 생성되고
$$
A[x]_\mathfrak q^\wedge \otimes_{A[x]} \kappa(\mathfrak r)
=
(A[x]_\mathfrak q^\wedge \otimes_A K)/(x - f)
$$
임을 알 수 있다. $A[x]$의 미분 $D = \text{d}/\text{d}x$는 $K[x]$로
연장되며 $x - f$를 $K[x]$의 가역원으로 보낸다. 또한 보조정리
\ref{more-algebra-lemma-derivation-extends}에 의해 $D$는
$A[x]_\mathfrak q^\wedge \otimes_A K$로 연장된다.
$A \to A[x]_\mathfrak q^\wedge$이 형식적으로 매끄러우므로(보조정리
\ref{more-algebra-lemma-formally-smooth}와
\ref{more-algebra-lemma-formally-smooth-completion}을 보라), Proposition
\ref{more-algebra-proposition-fs-regular}에 의해 환
$A[x]_\mathfrak q^\wedge \otimes_A K$는 정칙이다(이 특수한 경우에는
그 Proposition의 증명 논변이 상당히 단순해진다). 보조정리
\ref{more-algebra-lemma-quotient-regular}에 의해 결론이 따른다.
\end{proof}

\begin{proposition}
\label{more-algebra-proposition-finite-type-over-G-ring}
$R$을 G-환이라 하자. $R \to S$가 본질적으로 유한형이면 $S$는 G-환이다.
\end{proposition}

\begin{proof}
G-환이라는 것은 국소환들의 성질이므로 G-환의 국소화가 G-환임은
명백하다. 역으로 모든 소아이디얼에서의 국소화가 G-환이면 원래 환도
G-환이다. 따라서 모든 유한형 $R$-대수 $S$와 $S$의 모든 소아이디얼
$\mathfrak q$에 대하여 $S_\mathfrak q$가 G-환임을 보이면 충분하다.
$S$를 $R[x_1, \ldots, x_n]$의 몫으로 쓰면, 보조정리
\ref{more-algebra-lemma-G-ring-goes-up-quasi-finite}에 의해
$R[x_1, \ldots, x_n]$이 G-환임을 증명하면 충분하다. $n$에 대한
귀납법으로 $R[x]$가 G-환임을 증명하면 충분하다.
$\mathfrak q \subset R[x]$를 극대 아이디얼이라 하자. 보조정리
\ref{more-algebra-lemma-check-G-ring-maximal-ideals}에 의해
$$
R[x]_\mathfrak q \longrightarrow R[x]_\mathfrak q^\wedge
$$
이 정칙임을 보이면 충분하다. $\mathfrak q$가
$\mathfrak p \subset R$ 위에 놓이면 $R$을 $R_\mathfrak p$로 바꾸어도
된다. 따라서 $R$이 극대 아이디얼 $\mathfrak m$을 갖는 Noether 국소
G-환이고 $\mathfrak q \subset R[x]$가 $\mathfrak m$ 위에 놓인다고
가정해도 된다. $\mathfrak q$ 위에 놓이는 소아이디얼
$\mathfrak q' \subset R^\wedge[x]$은 유일함에 주의하자.
도식
$$
\xymatrix{
R[x]_\mathfrak q^\wedge \ar[r] &
(R^\wedge[x]_{\mathfrak q'})^\wedge \\
R[x]_\mathfrak q \ar[r] \ar[u] & R^\wedge[x]_{\mathfrak q'} \ar[u]
}
$$
을 생각하자. $R$이 G-환이므로 아래 수평 화살표는 정칙이다(정칙 환
사상 $R \to R^\wedge$의 기저변환의 국소화이기 때문이다). 오른쪽 수직
화살표가 정칙임을 증명할 수 있다고 하자. 그러면 합성
$R[x]_\mathfrak q \to (R^\wedge[x]_{\mathfrak q'})^\wedge$이 정칙이고,
보조정리 \ref{more-algebra-lemma-regular-permanence}에 의해 왼쪽 수직 화살표도
정칙이다. 그러므로 $R$이 Noether 완비 국소환이고 $\mathfrak q$가
$R$의 극대 아이디얼 위에 놓이는 소아이디얼이라고 가정해도 된다.

\medskip\noindent
$R$을 Noether 완비 국소환이라 하고, $\mathfrak q \subset R[x]$를
$R$의 극대 아이디얼 위에 놓이는 극대 아이디얼이라 하자.
$\mathfrak r \subset \mathfrak q$를 소아이디얼이라 하자. 우리는
$R[x]_\mathfrak q^\wedge \otimes_{R[x]} \kappa(\mathfrak r)$가
$\kappa(\mathfrak r)$ 위에서 기하학적으로 정칙인 대수임을 보이고자
한다. $\mathfrak p = R \cap \mathfrak r$라 두자. 그러면 $R$을
$R/\mathfrak p$로 바꾸고, $\mathfrak q$와 $\mathfrak r$를
% P02-E0186: frozen source writes the ill-grouped R/mathfrak p[x]; Korean preserves the formal span.
$R/\mathfrak p[x]$에서의 상으로 바꾸어도 된다. 보조정리
\ref{more-algebra-lemma-check-G-ring-easy}를 보라. 따라서 $R$이 정역이고
$\mathfrak r \cap R = (0)$이라고 가정해도 된다.

\medskip\noindent
Algebra, Lemma
\ref{algebra-lemma-complete-local-Noetherian-domain-finite-over-regular}에
의해 $R$이 정칙 부분환 $R_0 \subset R$ 위에서 유한이 되도록 하는
$R_0$를 찾을 수 있다. 보조정리
\ref{more-algebra-lemma-G-ring-goes-up-quasi-finite}를 적용하면, 위의 조건들에 더해
$R$이 정칙 완비 국소환인 경우
$R[x]_\mathfrak q^\wedge \otimes_{R[x]} \kappa(\mathfrak r)$가
$\kappa(\mathfrak r)$ 위에서 기하학적으로 정칙임을 증명하는 것으로
충분하다.

\medskip\noindent
이제 $R$은 정칙 완비 국소환이고,
$\mathfrak r \subset \mathfrak q \subset R[x]$이며,
$(0) = R \cap \mathfrak r$이고, $\mathfrak q$는 $R$의 극대 아이디얼
위에 놓이는 극대 아이디얼이다. $R$이 정칙이므로 환 $R[x]$도 정칙이다
(Algebra, Lemma \ref{algebra-lemma-regular-goes-up}). 따라서 국소화
$R[x]_\mathfrak q$는 정칙이다. 그러므로 보조정리
\ref{more-algebra-lemma-completion-regular}에 의해 완비화
$R[x]_\mathfrak q^\wedge$은 정칙이다. 따라서 올
$R[x]_{\mathfrak q}^\wedge \otimes_{R[x]} \kappa(\mathfrak r)$는
$R[x]_\mathfrak q^\wedge$의 국소화로서 역시 정칙이다.
그러므로 $R$의 분수체의 표수가 $0$이면 끝난다.

\medskip\noindent
$R$의 표수가 양수이면 $R = k[[x_1, \ldots, x_n]]$이다. 이 경우 논변을
두 경우로 나눈다.
\begin{enumerate}
\item $\mathfrak r = (0)$인 경우. 보조정리
\ref{more-algebra-lemma-helper-G-ring}에서 곧바로 따른다.
\item $\mathfrak r \not = (0)$인 경우. 이는 보조정리
\ref{more-algebra-lemma-another-helper-G-ring}이다.
\end{enumerate}
\end{proof}

\begin{remark}
\label{more-algebra-remark-G-does-not-survive-completion}
$R$을 G-환이라 하고 $I \subset R$을 아이디얼이라 하자. 일반적으로
$I$-진 완비화 $R^\wedge$은 G-환이 아니다. Nishimura가
\cite{Nishimura}에서 예를 제시하였다. 이 예의 일반화이자 어떤 의미의
해명은 \cite{Dumitrescu}의 마지막 절에서 찾을 수 있다.
\end{remark}

\begin{proposition}
\label{more-algebra-proposition-ubiquity-G-ring}
다음 종류의 환들은 G-환이다.
\begin{enumerate}
\item 체,
\item Noether 완비 국소환,
\item $\mathbf{Z}$,
\item 분수체의 표수가 영인 Dedekind 정역,
\item 위의 어느 환의 유한형 환 확대.
\end{enumerate}
\end{proposition}

\begin{proof}
체, $\mathbf{Z}$, 그리고 표수 영인 Dedekind 정역의 경우에는 정의와
이산 값매김환의 완비화가 이산 값매김환이라는 사실에서 곧바로 따른다.
Noether 완비 국소환은 Proposition
\ref{more-algebra-proposition-Noetherian-complete-G-ring}에 의해 G-환이다. 유한형
상위환에 대한 명제는 Proposition
\ref{more-algebra-proposition-finite-type-over-G-ring}이다.
\end{proof}

\begin{lemma}
\label{more-algebra-lemma-henselian-local-limit-G-rings}
$(A, \mathfrak m)$를 헨젤 국소환이라 하자. 그러면 $A$는 국소 전이
사상들을 갖는 헨젤 국소 G-환들의 계의 여과 여극한이다.
\end{lemma}

\begin{proof}
$A = \colim A_i$를 유한형 $\mathbf{Z}$-대수들의 여과 여극한으로 쓰자.
$\mathfrak p_i$를 $\mathfrak m$ 아래에 놓이는 $A_i$의 소아이디얼이라
하자. $A_i$를 $\mathfrak p_i$에서의 $A_i$의 국소화로 바꾸어도 된다. 그러면
$A_i$는 Noether 국소 G-환이다(Proposition
\ref{more-algebra-proposition-ubiquity-G-ring}). 보조정리
\ref{more-algebra-lemma-henselization-colimit}에 의해 $A = \colim A_i^h$이다.
보조정리 \ref{more-algebra-lemma-henselization-G-ring}에 의해 환 $A_i^h$들은
G-환이다.
\end{proof}

\begin{lemma}
\label{more-algebra-lemma-map-G-ring-to-completion-regular}
\begin{reference}
\cite[Theorem 79]{MatCA}
\end{reference}
$A$를 G-환이라 하자. $I \subset A$를 아이디얼이라 하고
$A^\wedge$을 $I$에 관한 $A$의 완비화라 하자. 그러면
$A \to A^\wedge$은 정칙이다.
\end{lemma}

\begin{proof}
환 사상 $A \to A^\wedge$은 Algebra, Lemma
\ref{algebra-lemma-completion-flat}에 의해 평탄이다. 환 $A^\wedge$은
Algebra, Lemma
\ref{algebra-lemma-completion-Noetherian-Noetherian}에 의해 Noether이다.
따라서 보조정리 \ref{more-algebra-lemma-regular-local}의 세 번째 조건을 확인하면
충분하다. $\mathfrak m' \subset A^\wedge$을 $\mathfrak m \subset A$
위에 놓이는 극대 아이디얼이라 하자. Algebra, Lemma
\ref{algebra-lemma-radical-completion}에 의해
$IA^\wedge \subset \mathfrak m'$이다. $A^\wedge/IA^\wedge = A/I$이므로
$I \subset \mathfrak m$, $\mathfrak m/I = \mathfrak m'/IA^\wedge$,
그리고 $A/\mathfrak m = A^\wedge/\mathfrak m'$이다.
$A^\wedge/\mathfrak m'$이 체이므로 $\mathfrak m$도 극대 아이디얼이다.
그러면 $A_\mathfrak m \to A^\wedge_{\mathfrak m'}$은 잉여체를
동일시하고
$\mathfrak m A^\wedge_{\mathfrak m'} =
\mathfrak m'A^\wedge_{\mathfrak m'}$을 만족하는 Noether 국소환의
평탄 국소 준동형이다. 따라서 보조정리 \ref{more-algebra-lemma-flat-unramified}에
의해 완비 국소환들 사이의 동형을 유도한다.
$(A_\mathfrak m)^\wedge$을 $A_\mathfrak m$의 극대 아이디얼에 관한
완비화라 하자. 환 사상
$$
(A^\wedge)_{\mathfrak m'} \to
((A^\wedge)_{\mathfrak m'})^\wedge = (A_\mathfrak m)^\wedge
$$
은 충실히 평탄이다(Algebra, Lemma
\ref{algebra-lemma-completion-faithfully-flat}). 따라서 환 사상들
$$
A_\mathfrak m \to (A^\wedge)_{\mathfrak m'} \to (A_\mathfrak m)^\wedge
$$
에 보조정리 \ref{more-algebra-lemma-regular-permanence}를 적용하여 결론을 얻는다.
실제로 $A$가 G-환이므로
$A_\mathfrak m \to (A_\mathfrak m)^\wedge$은 정칙이다.
\end{proof}

\begin{lemma}
\label{more-algebra-lemma-henselization-pair-G-ring}
\begin{slogan}
G-환인 성질은 아이디얼을 따른 헨젤화 아래에서 안정적이다
\end{slogan}
\begin{reference}
\cite[Theorem 5.3 i)]{Greco}
\end{reference}
$A$를 G-환이라 하고 $I \subset A$를 아이디얼이라 하자.
$(A^h, I^h)$를 쌍 $(A, I)$의 헨젤화라 하자. 보조정리
\ref{more-algebra-lemma-henselization}을 보라. 그러면 $A^h$는 G-환이다.
\end{lemma}

\begin{proof}
$\mathfrak m^h \subset A^h$를 극대 아이디얼이라 하자. 보조정리
\ref{more-algebra-lemma-check-G-ring-maximal-ideals}에 따라
$A^h_{\mathfrak m^h}$에서 그 완비화로 가는 사상의 올들이
기하학적으로 정칙임을 보여야 한다.
$\mathfrak m$을 $A$에서 $\mathfrak m^h$의 역상이라 하자.
$(A^h, I^h)$가 헨젤 쌍이므로 $I^h \subset \mathfrak m^h$이고 따라서
$I \subset \mathfrak m$임에 주의하자. $A^h$가 Noether이고,
$I^h = IA^h$이며, $A \to A^h$가 $I$-진 완비화들 사이의 동형을
유도함을 상기하자. 보조정리
\ref{more-algebra-lemma-henselization-Noetherian-pair}를 보라. 그러면 Noether
국소환들의 국소 준동형
$$
A_\mathfrak m \to A^h_{\mathfrak m^h}
$$
은 보조정리 \ref{more-algebra-lemma-flat-unramified}에 의해 극대 아이디얼에 관한
완비화들 사이의 동형을 유도한다(세부사항은 생략한다).
$\mathfrak q \subset A_\mathfrak m$ 위에 놓이는
$A^h_{\mathfrak m^h}$의 소아이디얼을 $\mathfrak q^h$라 하자.
$\mathfrak q_1 = \mathfrak q^h$라 두고,
$\mathfrak q_2, \ldots, \mathfrak q_t$를 $\mathfrak q$ 위에 놓이는
$A^h$의 나머지 소아이디얼들이라 하자. 보조정리
\ref{more-algebra-lemma-filtered-colimit-etale-noetherian-fibres}에 의해
$A^h \otimes_A \kappa(\mathfrak q) =
\prod\nolimits_{i = 1, \ldots, t} \kappa(\mathfrak q_i)$이다.
위에서 논의했듯이
$(A^h)_{\mathfrak m^h}^\wedge = (A_\mathfrak m)^\wedge$이므로
$$
\prod\nolimits_{i = 1, \ldots, t}
(A^h_{\mathfrak m^h})^\wedge \otimes_{A^h_{\mathfrak m^h}}
\kappa(\mathfrak q_i) =
(A^h_{\mathfrak m^h})^\wedge \otimes_{A^h_{\mathfrak m^h}}
(A^h_{\mathfrak m^h} \otimes_{A_{\mathfrak m}} \kappa(\mathfrak q)) =
(A_{\mathfrak m})^\wedge \otimes_{A_{\mathfrak m}} \kappa(\mathfrak q)
$$
를 얻는다. 따라서 한 성분인 환
$$
(A^h_{\mathfrak m^h})^\wedge \otimes_{A^h_{\mathfrak m^h}}
\kappa(\mathfrak q^h)
$$
은 $A$에 대한 가정에 의해 $\kappa(\mathfrak q)$ 위에서 기하학적으로
정칙이다. $\kappa(\mathfrak q^h)$는 $\kappa(\mathfrak q)$의 분리
대수적 확대이므로 Algebra, Lemma
\ref{algebra-lemma-geometrically-regular-over-separable-algebraic}에
의해
$$
(A^h_{\mathfrak m^h})^\wedge \otimes_{A^h_{\mathfrak m^h}}
\kappa(\mathfrak q^h)
$$
은 원하는 대로 $\kappa(\mathfrak q^h)$ 위에서 기하학적으로 정칙이다.
\end{proof}







\section{형식적 올의 성질}
\label{more-algebra-section-properties-formal-fibres}

\noindent
% P02-E0885: frozen source has the malformed “for to be able”; Korean renders the intended purpose clause.
이 절에서는 Noether 환의 형식적 올의 성질을 제대로 논의할 수 있도록
절 \ref{more-algebra-section-G-ring}의 몇 가지 논변을 다시 전개한다.

\medskip\noindent
$k$가 체이고 $R$이 Noether인 환 사상 $k \to R$의 성질을 $P$라 하자.
$B$가 Noether인 환 준동형 $A \to B$에 대하여, $A$의 모든 소아이디얼
$\mathfrak q$에 대해
$\kappa(\mathfrak q) \to B \otimes_A \kappa(\mathfrak q)$가 $P$를
가지면 $P$가 이 준동형의 올들에 대하여 성립한다고 말한다.
이하에서는 다음 명제들을 사용한다.
\begin{enumerate}
\item[(A)] 유한 생성 체 확대 $k'/k$에 대하여
$P(k \to R) \Rightarrow P(k' \to R \otimes_k k')$,
\item[(B)] $P(k \to R_\mathfrak p),\ \forall \mathfrak p \in \Spec(R)
\Leftrightarrow P(k \to R)$,
\item[(C)] Noether 환의 평탄 사상 $A \to B \to C$가 주어졌을 때,
$A \to B$의 올들이 $P$를 가지고 $B \to C$가 정칙이면
$A \to C$의 올들이 $P$를 가진다,
\item[(D)] Noether 환의 평탄 사상 $A \to B \to C$가 주어졌을 때,
$A \to C$의 올들이 $P$를 가지고 $B \to C$가 충실히 평탄이면
$A \to B$의 올들이 $P$를 가진다,
\item[(E)] $R$이 Noether인 $k \to k' \to R$이 주어졌을 때,
$k'/k$가 분리 대수적이고 $P(k \to R)$이면 $P(k' \to R)$이다,
% P02-E0886: frozen item (F) is the unresolved editorial placeholder “add more here.”
\item[(F)] 여기에 더 추가할 것.
\end{enumerate}
Noether 국소환 $A$가 주어졌을 때, $A \to A^\wedge$의 올들에 대하여
$P$가 성립하면 “{\it $A$의 형식적 올들이 $P$를 가진다}”고 말한다.
$R$이 Noether이고, $R$의 모든 소아이디얼 $\mathfrak p$에 대하여
$R_\mathfrak p$의 형식적 올들이 $P$를 가지면
{\it $R$은 $P$-환}이라고 말한다.

\begin{lemma}
\label{more-algebra-lemma-check-P-ring-easy}
$R$을 Noether 환이라 하고 $P$를 위와 같은 성질이라 하자. $R$이
$P$-환일 필요충분조건은 모든 소아이디얼 쌍
$\mathfrak q \subset \mathfrak p \subset R$에 대하여
$\kappa(\mathfrak q)$-대수
$$
(R/\mathfrak q)_\mathfrak p^\wedge \otimes_{R/\mathfrak q} \kappa(\mathfrak q)
$$
가 성질 $P$를 가지는 것이다.
\end{lemma}

\begin{proof}
$\kappa(\mathfrak q)$ 위의 대수로서
$$
R_\mathfrak p^\wedge \otimes_R \kappa(\mathfrak q) =
(R/\mathfrak q)_\mathfrak p^\wedge \otimes_{R/\mathfrak q} \kappa(\mathfrak q)
$$
가 성립한다는 사실에서 따른다.
\end{proof}

\begin{lemma}
\label{more-algebra-lemma-P-local}
$R \to \Lambda$를 Noether 환의 준동형이라 하자. $P$가 성질 (B)를
만족한다고 가정하자. 다음은 서로 동치이다.
\begin{enumerate}
\item $R \to \Lambda$의 올들이 $P$를 가진다.
\item $\mathfrak p \subset R$ 위에 놓이는 모든
$\mathfrak q \subset \Lambda$에 대하여
$R_\mathfrak p \to \Lambda_\mathfrak q$의 올들이 $P$를 가진다.
\item $R$의 $\mathfrak m$ 위에 놓이는 모든 극대 아이디얼
$\mathfrak m' \subset \Lambda$에 대하여
$R_\mathfrak m \to \Lambda_{\mathfrak m'}$의 올들이 $P$를 가진다.
\end{enumerate}
\end{lemma}

\begin{proof}
$\mathfrak p \subset R$을 소아이디얼이라 하자. $\mathfrak p$ 위의 올은
환 $\Lambda \otimes_R \kappa(\mathfrak p)$이고, 그 스펙트럼은
$\mathfrak p$ 위에 놓이는 소아이디얼 $\mathfrak q$들로 이루어진
$\Spec(\Lambda)$의 부분집합 위로 전단사적으로 사상된다. Algebra,
Remark \ref{algebra-remark-fundamental-diagram}을 보라.
그러한 소아이디얼 $\mathfrak q$에 대하여
$\mathfrak q \subset \mathfrak m'$인 극대 아이디얼을 택하고
$\mathfrak m = R \cap \mathfrak m'$이라 두자.
그러면 $\mathfrak p \subset \mathfrak m$이고
$$
(\Lambda \otimes_R \kappa(\mathfrak p))_\mathfrak q \cong
(\Lambda_{\mathfrak m'} \otimes_{R_\mathfrak m}
\kappa(\mathfrak p))_\mathfrak q
$$
이다.
% P02-E0189: frozen source calls these kappa(q)-algebras although the displayed fibre base is kappa(p).
$\kappa(\mathfrak q)$-대수로서 위 동형이 성립한다. 따라서 (B)에 의해
국소환들에서 성질 $P$를 확인할 수 있으므로 (1), (2), (3)은 동치이다.
\end{proof}

\begin{lemma}
\label{more-algebra-lemma-P-ring-goes-up-quasi-finite}
$R \to R'$을 Noether 환의 유한형 사상이라 하고
$$
\xymatrix{
\mathfrak q' \ar[r] & \mathfrak p' \ar[r] & R' \\
\mathfrak q \ar[r] \ar@{-}[u] &
\mathfrak p \ar[r] \ar@{-}[u] & R \ar[u]
}
$$
가 소아이디얼들로 이루어져 있다고 하자. $R \to R'$이
$\mathfrak p'$에서 준유한이라고 가정하자. $P$가 (A)와 (B)를
만족한다고 가정하자.
\begin{enumerate}
\item $\kappa(\mathfrak q) \to
R_\mathfrak p^\wedge \otimes_R \kappa(\mathfrak q)$가 $P$를 가지면
$\kappa(\mathfrak q') \to
(R'_{\mathfrak p'})^\wedge \otimes_{R'} \kappa(\mathfrak q')$도 $P$를 가진다.
\item $R_\mathfrak p$의 형식적 올들이 $P$를 가지면
$R'_{\mathfrak p'}$의 형식적 올들도 $P$를 가진다.
\item $R \to R'$이 준유한이고 $R$이 $P$-환이면 $R'$은 $P$-환이다.
\end{enumerate}
\end{lemma}

\begin{proof}
(1) $\Rightarrow$ (2) $\Rightarrow$ (3)은 명백하다. 사상
$\kappa(\mathfrak q) \to
R_\mathfrak p^\wedge \otimes_R \kappa(\mathfrak q)$에 대하여 $P$가
성립한다고 가정하자. Algebra, Lemma
\ref{algebra-lemma-completion-at-quasi-finite-prime}에 의해
$$
R_\mathfrak p^\wedge \otimes_R R'
=
(R'_{\mathfrak p'})^\wedge \times B
$$
이다. 여기서 $B$는 어떤 $R_\mathfrak p^\wedge$-대수이다. 따라서
$R'_{\mathfrak p'} \to (R'_{\mathfrak p'})^\wedge$은
$R_\mathfrak p \to R_\mathfrak p^\wedge$의 기저변환의 한 인자이다.
그러므로 $(R'_{\mathfrak p'})^\wedge \otimes_{R'}
\kappa(\mathfrak q')$은
$$
R_\mathfrak p^\wedge \otimes_R R' \otimes_{R'} \kappa(\mathfrak q') =
R_\mathfrak p^\wedge \otimes_R \kappa(\mathfrak q)
\otimes_{\kappa(\mathfrak q)} \kappa(\mathfrak q').
$$
의 한 인자이다. 따라서 $P$에 대한 가정들에서 결론이 따른다.
\end{proof}

\begin{lemma}
\label{more-algebra-lemma-check-P-ring-maximal-ideals}
$R$을 Noether 환이라 하자. $P$가 (C)와 (D)를 만족한다고 가정하자.
그러면 $R$이 $P$-환일 필요충분조건은 $R$의 모든 극대 아이디얼
$\mathfrak m$에 대하여 $R_\mathfrak m$의 형식적 올들이 $P$를
가지는 것이다.
\end{lemma}

\begin{proof}
$R$의 모든 극대 아이디얼 $\mathfrak m$에 대하여 $R_\mathfrak m$의
형식적 올들이 $P$를 가진다고 가정하자. $\mathfrak p$를 $R$의
소아이디얼이라 하고 극대 아이디얼 $\mathfrak p \subset \mathfrak m$을
택하자. $R_\mathfrak m \to R_\mathfrak m^\wedge$은 충실히 평탄이므로
$\mathfrak pR_\mathfrak m$ 위에 놓이는
$R_\mathfrak m^\wedge$의 소아이디얼 $\mathfrak p'$을 택할 수 있다.
가환도식
$$
\xymatrix{
R_\mathfrak m^\wedge \ar[r] &
(R_\mathfrak m^\wedge)_{\mathfrak p'} \ar[r] &
(R_\mathfrak m^\wedge)_{\mathfrak p'}^\wedge
\\
R_\mathfrak m \ar[u] \ar[r] & R_\mathfrak p \ar[u] \ar[r] &
R_\mathfrak p^\wedge \ar[u]
}
$$
을 생각하자. 가정에 의해 환 사상
$R_\mathfrak m \to R_\mathfrak m^\wedge$의 올들은 $P$를 가진다.
Proposition \ref{more-algebra-proposition-Noetherian-complete-G-ring}에 의해
$(R_\mathfrak m^\wedge)_{\mathfrak p'} \to
(R_\mathfrak m^\wedge)_{\mathfrak p'}^\wedge$은 정칙이다. 국소화
$R_\mathfrak m^\wedge \to (R_\mathfrak m^\wedge)_{\mathfrak p'}$도
정칙이다. 따라서 보조정리 \ref{more-algebra-lemma-regular-composition}에 의해
$R_\mathfrak m^\wedge \to
(R_\mathfrak m^\wedge)_{\mathfrak p'}^\wedge$은 정칙이다.
그러므로 (C)에 의해
$R_\mathfrak m \to (R_\mathfrak m^\wedge)_{\mathfrak p'}^\wedge$의
올들은 $P$를 가진다.
$R_\mathfrak m \to (R_\mathfrak m^\wedge)_{\mathfrak p'}^\wedge$은
국소화 $R_\mathfrak p$를 통해 분해되므로
$R_\mathfrak p \to (R_\mathfrak m^\wedge)_{\mathfrak p'}^\wedge$의
올들도 $P$를 가진다. 이제 (D)를 적용하면
$R_\mathfrak p \to R_\mathfrak p^\wedge$의 올들이 $P$를 가짐을
알 수 있다.
\end{proof}

\begin{proposition}
\label{more-algebra-proposition-finite-type-over-P-ring}
$R$을 $P$-환이라 하고, $P$가 (A), (B), (C), (D)를 만족한다고 하자.
$R \to S$가 본질적으로 유한형이면 $S$는 $P$-환이다.
\end{proposition}

\begin{proof}
$P$-환이라는 것은 국소환들의 성질이므로 $P$-환의 국소화가
$P$-환임은 명백하다. 역으로 모든 소아이디얼에서의 국소화가
$P$-환이면 원래 환도 $P$-환이다. 따라서 모든 유한형 $R$-대수
$S$와 $S$의 모든 소아이디얼 $\mathfrak q$에 대하여
$S_\mathfrak q$가 $P$-환임을 보이면 충분하다.
$S$를 $R[x_1, \ldots, x_n]$의 몫으로 쓰면 보조정리
\ref{more-algebra-lemma-P-ring-goes-up-quasi-finite}에 의해
$R[x_1, \ldots, x_n]$이 $P$-환임을 증명하면 충분하다. $n$에 대한
귀납법으로 $R[x]$가 $P$-환임을 증명하면 충분하다.
$\mathfrak q \subset R[x]$를 극대 아이디얼이라 하자. 보조정리
\ref{more-algebra-lemma-check-P-ring-maximal-ideals}에 의해
$$
R[x]_\mathfrak q \longrightarrow R[x]_\mathfrak q^\wedge
$$
의 올들이 $P$를 가짐을 보이면 충분하다. $\mathfrak q$가
$\mathfrak p \subset R$ 위에 놓이면 $R$을 $R_\mathfrak p$로
바꾸어도 된다. 따라서 $R$이 극대 아이디얼 $\mathfrak m$을 갖는
Noether 국소 $P$-환이고 $\mathfrak q \subset R[x]$가
$\mathfrak m$ 위에 놓인다고 가정해도 된다.
$\mathfrak q$ 위에 놓이는 소아이디얼
$\mathfrak q' \subset R^\wedge[x]$은 유일함에 주의하자. 도식
$$
\xymatrix{
R[x]_\mathfrak q^\wedge \ar[r] &
(R^\wedge[x]_{\mathfrak q'})^\wedge \\
R[x]_\mathfrak q \ar[r] \ar[u] & R^\wedge[x]_{\mathfrak q'} \ar[u]
}
$$
을 생각하자. $R$이 $P$-환이므로 $R[x] \to R^\wedge[x]$의 올들은
$P$를 가진다. 실제로 이들은 $R \to R^\wedge$의 올들을 유한 생성
체 확대하여 기저변환한 것이므로 (A)를 적용할 수 있다. 따라서 예를
들어 보조정리 \ref{more-algebra-lemma-P-local}에 의해 아래 수평 화살표의 올들은
$P$를 가진다. 오른쪽 수직 화살표는 $R^\wedge$이 G-환이므로 정칙이다
(Proposition \ref{more-algebra-proposition-Noetherian-complete-G-ring}과
\ref{more-algebra-proposition-finite-type-over-G-ring}). 따라서 (C)에 의해 합성
$R[x]_\mathfrak q \to (R^\wedge[x]_{\mathfrak q'})^\wedge$의 올들은
$P$를 가진다. 이제 (D)에 의해 왼쪽 수직 화살표의 올들도 $P$를
가지므로 증명이 끝난다.
\end{proof}

\begin{lemma}
\label{more-algebra-lemma-map-P-ring-to-completion-P}
$A$를 $P$-환이라 하고 $P$가 (B)와 (D)를 만족한다고 하자.
$I \subset A$를 아이디얼이라 하고 $A^\wedge$을 $I$에 관한 $A$의
완비화라 하자. 그러면 $A \to A^\wedge$의 올들은 $P$를 가진다.
\end{lemma}

\begin{proof}
환 사상 $A \to A^\wedge$은 Algebra, Lemma
\ref{algebra-lemma-completion-flat}에 의해 평탄이다. 환 $A^\wedge$은
Algebra, Lemma
\ref{algebra-lemma-completion-Noetherian-Noetherian}에 의해 Noether이다.
따라서 보조정리 \ref{more-algebra-lemma-P-local}의 세 번째 조건을 확인하면
충분하다. $\mathfrak m' \subset A^\wedge$을 $\mathfrak m \subset A$
위에 놓이는 극대 아이디얼이라 하자. Algebra, Lemma
\ref{algebra-lemma-radical-completion}에 의해
$IA^\wedge \subset \mathfrak m'$이다. $A^\wedge/IA^\wedge = A/I$이므로
$I \subset \mathfrak m$, $\mathfrak m/I = \mathfrak m'/IA^\wedge$,
그리고 $A/\mathfrak m = A^\wedge/\mathfrak m'$이다.
$A^\wedge/\mathfrak m'$이 체이므로 $\mathfrak m$도 극대 아이디얼이다.
그러면 $A_\mathfrak m \to A^\wedge_{\mathfrak m'}$은 잉여체를
동일시하고
$\mathfrak m A^\wedge_{\mathfrak m'} =
\mathfrak m'A^\wedge_{\mathfrak m'}$을 만족하는 Noether 국소환의
평탄 국소 준동형이다. 따라서 보조정리 \ref{more-algebra-lemma-flat-unramified}에
의해 완비 국소환들 사이의 동형을 유도한다.
$(A_\mathfrak m)^\wedge$을 $A_\mathfrak m$의 극대 아이디얼에 관한
완비화라 하자. 환 사상
$$
(A^\wedge)_{\mathfrak m'} \to
((A^\wedge)_{\mathfrak m'})^\wedge = (A_\mathfrak m)^\wedge
$$
은 충실히 평탄이다(Algebra, Lemma
\ref{algebra-lemma-completion-faithfully-flat}). 따라서 환 사상들
$$
A_\mathfrak m \to (A^\wedge)_{\mathfrak m'} \to (A_\mathfrak m)^\wedge
$$
에 (D)를 적용하여 결론을 얻는다. 실제로 $A$가 $P$-환이므로
$A_\mathfrak m \to (A_\mathfrak m)^\wedge$의 올들은 $P$를 가진다.
\end{proof}

\begin{lemma}
\label{more-algebra-lemma-henselization-pair-P-ring}
\begin{slogan}
환의 헨젤화는 형식적 올의 좋은 성질을 물려받는다
\end{slogan}
$A$를 $P$-환이라 하고 $P$가 (B), (C), (D), (E)를 만족한다고 하자.
$I \subset A$를 아이디얼이라 하자. $(A^h, I^h)$를 쌍 $(A, I)$의
헨젤화라 하자. 보조정리 \ref{more-algebra-lemma-henselization}을 보라.
그러면 $A^h$는 $P$-환이다.
\end{lemma}

\begin{proof}
$\mathfrak m^h \subset A^h$를 극대 아이디얼이라 하자.
보조정리 \ref{more-algebra-lemma-check-P-ring-maximal-ideals}에 따라
$A^h_{\mathfrak m^h} \to (A^h_{\mathfrak m^h})^\wedge$의 올들이
$P$를 가짐을 보여야 한다. $\mathfrak m$을 $A$에서
$\mathfrak m^h$의 역상이라 하자. $(A^h, I^h)$가 헨젤 쌍이므로
$I^h \subset \mathfrak m^h$이고 따라서 $I \subset \mathfrak m$임에
주의하자. $A^h$가 Noether이고, $I^h = IA^h$이며, $A \to A^h$가
$I$-진 완비화들 사이의 동형을 유도함을 상기하자. 보조정리
\ref{more-algebra-lemma-henselization-Noetherian-pair}를 보라. 그러면 Noether
국소환들의 국소 준동형
$$
A_\mathfrak m \to A^h_{\mathfrak m^h}
$$
은 보조정리 \ref{more-algebra-lemma-flat-unramified}에 의해 극대 아이디얼에 관한
완비화들 사이의 동형을 유도한다(세부사항은 생략한다).
$\mathfrak q \subset A_\mathfrak m$ 위에 놓이는
$A^h_{\mathfrak m^h}$의 소아이디얼을 $\mathfrak q^h$라 하자.
$\mathfrak q_1 = \mathfrak q^h$라 두고,
$\mathfrak q_2, \ldots, \mathfrak q_t$를 $\mathfrak q$ 위에 놓이는
$A^h$의 나머지 소아이디얼들이라 하자. 보조정리
\ref{more-algebra-lemma-filtered-colimit-etale-noetherian-fibres}에 의해
$A^h \otimes_A \kappa(\mathfrak q) =
\prod\nolimits_{i = 1, \ldots, t} \kappa(\mathfrak q_i)$이다.
위에서 논의했듯이
$(A^h)_{\mathfrak m^h}^\wedge = (A_\mathfrak m)^\wedge$이므로
$$
\prod\nolimits_{i = 1, \ldots, t}
(A^h_{\mathfrak m^h})^\wedge \otimes_{A^h_{\mathfrak m^h}}
\kappa(\mathfrak q_i) =
(A^h_{\mathfrak m^h})^\wedge \otimes_{A^h_{\mathfrak m^h}}
(A^h_{\mathfrak m^h} \otimes_{A_{\mathfrak m}} \kappa(\mathfrak q)) =
(A_{\mathfrak m})^\wedge \otimes_{A_{\mathfrak m}} \kappa(\mathfrak q)
$$
를 얻는다. 따라서 국소환들을 보고 (B)를 사용하면
$$
\kappa(\mathfrak q) \longrightarrow
(A^h_{\mathfrak m^h})^\wedge \otimes_{A^h_{\mathfrak m^h}}
\kappa(\mathfrak q^h)
$$
가 $P$를 가짐을 알 수 있다. 실제로 $A$에 대한 가정에 의해
$\kappa(\mathfrak q) \to
(A_\mathfrak m)^\wedge \otimes_{A_\mathfrak m} \kappa(\mathfrak q)$가
그 성질을 가진다. $\kappa(\mathfrak q^h)/\kappa(\mathfrak q)$는 분리
대수적이므로 (E)에 의해
$\kappa(\mathfrak q^h) \to
(A^h_{\mathfrak m^h})^\wedge \otimes_{A^h_{\mathfrak m^h}}
\kappa(\mathfrak q^h)$가 원하는 대로 $P$를 가진다.
\end{proof}

\begin{lemma}
\label{more-algebra-lemma-henselization-P-ring}
$R$을 $P$-환인 Noether 국소환이라 하고 $P$가 (B), (C), (D), (E)를
만족한다고 하자. 그러면 헨젤화 $R^h$와 엄밀 헨젤화 $R^{sh}$은
$P$-환이다.
\end{lemma}

\begin{proof}
헨젤화에 대해서는 보조정리
\ref{more-algebra-lemma-henselization-pair-P-ring}에서 보았다. 엄밀 헨젤화에 대해
증명하려면 보조정리 \ref{more-algebra-lemma-check-P-ring-maximal-ideals}에 따라
$R^{sh}$의 형식적 올들이 $P$를 가짐을 보이면 충분하다.
$\mathfrak r \subset R^{sh}$을 소아이디얼이라 하고
$\mathfrak p = R \cap \mathfrak r$라 두자.
$\mathfrak r_1 = \mathfrak r$라 두고,
$\mathfrak r_2, \ldots, \mathfrak r_s$를 $\mathfrak p$ 위에 놓이는
$R^{sh}$의 나머지 소아이디얼들이라 하자. 보조정리
\ref{more-algebra-lemma-fibres-henselization}에 의해
$R^{sh} \otimes_R \kappa(\mathfrak p) =
\prod\nolimits_{i = 1, \ldots, s} \kappa(\mathfrak r_i)$이다.
그러면 다음을 얻는다.
% P02-E0887: frozen source indexes primes through s but the product through t; Korean preserves the display.
$$
\prod\nolimits_{i = 1, \ldots, t}
(R^{sh})^\wedge \otimes_{R^{sh}} \kappa(\mathfrak r_i) =
(R^{sh})^\wedge \otimes_{R^{sh}} (R^{sh} \otimes_R \kappa(\mathfrak p)) =
(R^{sh})^\wedge \otimes_R \kappa(\mathfrak p)
$$
보조정리 \ref{more-algebra-lemma-henselization-noetherian}에 의해
$R^\wedge \to (R^{sh})^\wedge$은
$\mathfrak m_{(R^{sh})^\wedge}$-진 위상에서 형식적으로 매끄럽다.
따라서 Proposition \ref{more-algebra-proposition-fs-regular}에 의해
$R^\wedge \to (R^{sh})^\wedge$은 정칙이다. 그러므로 (C)와 $R$에 대한
가정에 의해 $P$는
$\kappa(\mathfrak p) \to (R^{sh})^\wedge \otimes_R
\kappa(\mathfrak p)$에 대하여 성립한다. 성질 (B)와 위의 분해를
사용하고 국소환들을 보면, $P$가
$\kappa(\mathfrak p) \to (R^{sh})^\wedge \otimes_{R^{sh}}
\kappa(\mathfrak r)$에 대하여 성립함을 얻는다.
$\kappa(\mathfrak r)/\kappa(\mathfrak p)$는 분리 대수적이므로
(E)에 의해 $P$는
$\kappa(\mathfrak r) \to (R^{sh})^\wedge \otimes_{R^{sh}}
\kappa(\mathfrak r)$에 대하여 성립한다.
\end{proof}

\begin{lemma}
\label{more-algebra-lemma-formal-fibres-reduced}
$P(k \to R) =$ “$R$은 $k$ 위에서 기하학적으로 기약이다”인 경우
성질 (A), (B), (C), (D), (E)가 성립한다.
\end{lemma}

\begin{proof}
(A)는 기하학적으로 기약인 대수의 정의에서 따른다(Algebra,
Definition \ref{algebra-definition-geometrically-reduced}).
(B)도 따른다. 환이 기약일 필요충분조건은 모든 국소환이 기약인 것이다.
(C)는 보조정리 \ref{more-algebra-lemma-reduced-goes-up}에서 따른다.
(D)는 Algebra, Lemma \ref{algebra-lemma-descent-reduced}에서 따른다.
(E)는 Algebra, Lemma
\ref{algebra-lemma-geometrically-reduced-over-separable-algebraic}에서
따른다.
\end{proof}


\begin{lemma}
\label{more-algebra-lemma-formal-fibres-normal}
$P(k \to R) =$ “$R$은 $k$ 위에서 기하학적으로 정규이다”인 경우
성질 (A), (B), (C), (D), (E)가 성립한다.
\end{lemma}

\begin{proof}
(A)는 기하학적으로 정규인 대수의 정의에서 따른다(Algebra,
Definition \ref{algebra-definition-geometrically-normal}).
(B)도 따른다. 환이 정규일 필요충분조건은 모든 국소환이 정규인 것이다.
(C)는 보조정리 \ref{more-algebra-lemma-normal-goes-up}에서 따른다.
(D)는 Algebra, Lemma \ref{algebra-lemma-descent-normal}에서 따른다.
(E)는 Algebra, Lemma
\ref{algebra-lemma-geometrically-normal-over-separable-algebraic}에서
따른다.
\end{proof}

\begin{lemma}
\label{more-algebra-lemma-formal-fibres-Sk}
$n \geq 1$을 고정하자. $P(k \to R) =$ “$R$이 $(S_n)$을 만족한다”인
경우 성질 (A), (B), (C), (D), (E)가 성립한다.
\end{lemma}

\begin{proof}
$k$가 체이고 $R$이 Noether인 환 사상 $k \to R$을 생각하자.
$k'/k$를 유한 생성 체 확대라 하자. 그러면 환 사상
$R \to R \otimes_k k'$의 올들은 Algebra, Lemma
\ref{algebra-lemma-tensor-fields-CM}에 의해 Cohen--Macaulay이다.
따라서 환 사상 $R \to R \otimes_k k'$에 Algebra, Lemma
\ref{algebra-lemma-Sk-goes-up}을 적용하면 $R$이 $(S_n)$을 만족할 때
$R \otimes_k k'$도 이를 만족함을 알 수 있다. 이로써 (A)가 증명된다.
% P02-E0888: frozen source says “a Noetherian rings has”; Korean uses the grammatical singular.
(B)도 따른다. Noether 환이 $(S_n)$을 만족할 필요충분조건은 모든
국소환이 $(S_n)$을 만족하는 것이다. (C)는 Algebra, Lemma
\ref{algebra-lemma-Sk-goes-up}에서 따른다. 실제로 정칙 준동형의 올들은
정칙이고, 특히 Cohen--Macaulay이다. (D)는 Algebra, Lemma
\ref{algebra-lemma-descent-Sk}에서 따른다. (E)는 이 조건이 기저체를
언급하지 않으므로 즉시 성립한다.
\end{proof}

\begin{lemma}
\label{more-algebra-lemma-formal-fibres-CM}
$P(k \to R) =$ “$R$이 Cohen--Macaulay이다”인 경우
성질 (A), (B), (C), (D), (E)가 성립한다.
\end{lemma}

\begin{proof}
보조정리 \ref{more-algebra-lemma-formal-fibres-Sk}와, Noether 환이 Cohen--Macaulay일
필요충분조건은 모든 $n$에 대하여 조건 $(S_n)$을 만족하는 것이라는
사실에서 곧바로 따른다.
\end{proof}

\begin{lemma}
\label{more-algebra-lemma-formal-fibres-Rk}
$n \geq 0$을 고정하자.
$P(k \to R) =$ “모든 유한 확대 $k'/k$에 대하여
$R \otimes_k k'$이 $(R_n)$을 만족한다”인 경우
성질 (A), (B), (C), (D), (E)가 성립한다.
\end{lemma}

\begin{proof}
$k$가 체이고 $R$이 Noether인 환 사상 $k \to R$을 생각하자.
$P(k \to R)$가 참이라고 가정하자. $K/k$를 유한 생성 체 확대라 하자.
Algebra, Lemma \ref{algebra-lemma-make-separable}에 의해 도식
$$
\xymatrix{
K \ar[r] & K' \\
k \ar[u] \ar[r] & k' \ar[u]
}
$$
을 찾을 수 있다. 여기서 $k'/k$와 $K'/K$는 유한 순수 비분리 체
확대이고 $K'/k'$은 분리이다. Algebra, Lemma
\ref{algebra-lemma-localization-smooth-separable}에 의해 매끄러운
$k'$-대수 $B$가 존재하여 $K'$이 $B$의 분수체가 된다.
이제 다음과 같이 논증한다.

1단계: $k \to R$에 대하여 $P$를 가정했으므로
$R \otimes_k k'$은 $(R_n)$을 만족한다.

2단계: $R \otimes_k k' \to R \otimes_k k' \otimes_{k'} B$는 매끄러운
환 사상이고(Algebra, Lemma \ref{algebra-lemma-base-change-smooth}),
Algebra, Lemma \ref{algebra-lemma-Rk-goes-up}에 의해
$R \otimes_k k' \otimes_{k'} B$는 $(R_n)$을 만족한다(가정들이
성립함을 보려면 Algebra, Lemma
\ref{algebra-lemma-characterize-smooth-over-field}도 사용한다).

3단계. $R \otimes_k k' \otimes_{k'} K' = R \otimes_k K'$은
$(R_n)$을 만족하는 환의 국소화이므로 $(R_n)$을 만족한다.

4단계. 끝으로, 충실히 평탄인 환 사상
$K \otimes_k R \to K' \otimes_k R$을 따른 $(R_n)$의 내림에 의해
$R \otimes_k K$가 $(R_n)$을 만족한다(Algebra, Lemma
\ref{algebra-lemma-descent-Rk}). 이로써 (A)가 증명된다.

(B)도 따른다. Noether 환이 $(R_n)$을 만족할 필요충분조건은 모든
국소환이 $(R_n)$을 만족하는 것이다. (C)는 Algebra, Lemma
\ref{algebra-lemma-Rk-goes-up}에서 따른다. 실제로 정칙 준동형의 올들은
정칙이다(작은 세부사항은 생략한다). (D)는 Algebra, Lemma
\ref{algebra-lemma-descent-Rk}에서 따른다(작은 세부사항은 생략한다).

\medskip\noindent
(E). $l/k$를 분리 대수적 체 확대라 하고, $R$이 Noether인 환 사상
$l \to R$을 생각하자. $k \to R$이 $P$를 가진다고 가정하자.
$l \to R$이 $P$를 가짐을 보여야 한다. $l'/l$을 유한 확대라 하자.
먼저 유한 부분 확대 $l/m/k$와 유한 확대 $m'/m$이 존재하여
$l' = l \otimes_m m'$이 됨을 관찰하자. 그러면
$R \otimes_l l' = R \otimes_m m'$이다. 따라서 $m \to R$이 성질
$P$를 가짐을 증명하면 충분하다. 즉, $l/k$가 유한이라고 가정해도 된다.
$l/k$가 유한이면 $l'/k$도 유한이고
$$
l' \otimes_l R = (l' \otimes_k R) \otimes_{l \otimes_k l} l
$$
은 Noether 환 $l' \otimes_k R$의 국소화이다(Algebra, Lemma
\ref{algebra-lemma-separable-algebraic-diagonal}). 이 환은
$k \to R$에 대한 가정 $P$에 의해 성질 $(R_n)$을 가진다. 따라서
$l' \otimes_l R$도 원하는 대로 성질 $(R_n)$을 가진다.
\end{proof}

\section{우수환}
\label{more-algebra-section-excellent}

\noindent
이 절에서는 Grothendieck의 우수환 개념을 논의한다.
G-환, J-2 환, 보편적 catenary 환의 정의는 각각 정의
\ref{more-algebra-definition-G-ring}, 정의 \ref{more-algebra-definition-J}, 그리고
Algebra, Definition \ref{algebra-definition-universally-catenary}를 보라.

\begin{definition}
\label{more-algebra-definition-excellent}
$R$을 환이라 하자.
\begin{enumerate}
\item $R$이 Noether 환이고 G-환이며 J-2이면 $R$을
{\it 준우수}라고 한다.
\item $R$이 준우수이고 보편적 catenary이면 $R$을
{\it 우수}하다고 한다.
\end{enumerate}
\end{definition}

\noindent
따라서 Noether 환이 준우수라는 것은 그 형식적 올들이 기하학적으로
정칙이고 그 위의 모든 유한형 대수의 특이 집합이 닫혀 있다는 뜻이다.
그러한 환이 우수하려면 이에 더하여 (국소적으로) 좋은 차원 함수가
존재해야 한다. 뒤에서(절 \ref{more-algebra-section-formally-catenary})
보편적 catenary 조건을 $R$의 극대아이디얼 $\mathfrak m$에 대한 사상
$R_\mathfrak m \to R_\mathfrak m^\wedge$의 조건으로 나타낼 수 있음을
볼 것이다.

\begin{lemma}
\label{more-algebra-lemma-finite-type-over-excellent}
(준)우수환 위의 유한형 환을 국소화한 모든 환은 (준)우수하다.
\end{lemma}

\begin{proof}
유한형 대수에 대해서는 J-2와 보편적 catenary 성질의 정의에서 따른다.
G-환에 대해서는 명제 \ref{more-algebra-proposition-finite-type-over-G-ring}을 보라.
국소화가 (준)우수성을 보존한다는 증명은 생략한다.
\end{proof}

\begin{proposition}
\label{more-algebra-proposition-ubiquity-excellent}
다음 종류의 환들은 우수하다.
\begin{enumerate}
\item 체,
\item Noether 완비 국소환,
\item $\mathbf{Z}$,
\item 분수체의 표수가 영인 Dedekind 정역,
\item 위 환들 중 하나의 유한형 환 확대.
\end{enumerate}
\end{proposition}

\begin{proof}
이 환들이 G-환이고 J-2를 가짐은 명제
\ref{more-algebra-proposition-ubiquity-G-ring}과
\ref{more-algebra-proposition-ubiquity-J-2}를 보라. 모든 Cohen--Macaulay 환은
보편적 catenary이다. Algebra, Lemma
\ref{algebra-lemma-CM-ring-catenary}를 보라. 특히 체, Dedekind 환,
더 일반적으로 정칙환은 보편적 catenary이다. Cohen 구조 정리를
통해 완비 국소환이 보편적 catenary임을 알 수 있다. Algebra, Remark
\ref{algebra-remark-Noetherian-complete-local-ring-universally-catenary}를
보라.
\end{proof}

\noindent
위에서 전개한 내용은 Nagata 환에 관하여 몇 가지 결과를 준다.

\begin{lemma}
\label{more-algebra-lemma-Nagata-local-ring}
$(A, \mathfrak m)$를 Noether 국소환이라 하자. 다음은 서로 동치이다.
\begin{enumerate}
\item $A$는 Nagata이다.
\item $A$의 형식적 올들은 기하학적으로 기약이다.
\end{enumerate}
\end{lemma}

\begin{proof}
(2)를 가정하자. Algebra, Lemma
\ref{algebra-lemma-local-nagata-and-analytically-unramified}에 의해,
$A \to B$가 유한이고 $B$가 정역이며 $\mathfrak m' \subset B$가
극대아이디얼이면 $B_{\mathfrak m'}$이 해석적으로 비분기임을 보이면 된다.
보조정리 \ref{more-algebra-lemma-formal-fibres-reduced}와
\ref{more-algebra-lemma-check-P-ring-maximal-ideals}, 그리고 명제
\ref{more-algebra-proposition-finite-type-over-P-ring}을 함께 적용하면
$B_{\mathfrak m'}$의 형식적 올들이 기하학적으로 기약임을 알 수 있다.
특히 $L$을 $B$의 분수체라 할 때
$B_{\mathfrak m'}^\wedge \otimes_B L$은 기약원이다. 따라서
$B_{\mathfrak m'}^\wedge$는 기약원이다. 즉,
$B_{\mathfrak m'}$은 해석적으로 비분기이다.

\medskip\noindent
(1)을 가정하자. $\mathfrak q \subset A$를 소아이디얼이라 하고
$K/\kappa(\mathfrak q)$를 유한 확대라 하자.
$A^\wedge \otimes_A K$가 기약원임을 보여야 한다.
$A/\mathfrak q \subset B \subset K$를 $A$ 위 유한이고 분수체가 $K$인
국소 부분환이라 하자. $B$를 구성하기 위해 $\kappa(\mathfrak q)$ 위에서
$K$를 생성하고 다음과 같은 일계수 다항식을 만족하는
$x_1, \ldots, x_n \in K$를 택한다.
$P_i(T) = T^{d_i} + a_{i, 1} T^{d_i - 1} + \ldots + a_{i, d_i} = 0$,
여기서 $a_{i, j} \in \mathfrak m$이다. 이제 $B$를
$x_1, \ldots, x_n$으로 생성되는 $K$의 $A$-부분대수로 잡는다.
(더 자세한 내용은 Algebra, Lemma
\ref{algebra-lemma-local-nagata-and-analytically-unramified}의 증명을 보라.)
그러면
$$
A^\wedge \otimes_A K =
(A^\wedge \otimes_A B)_\mathfrak q =
B^\wedge_\mathfrak q
$$
이다. Algebra, Lemma
\ref{algebra-lemma-local-nagata-and-analytically-unramified}에 의해
$B^\wedge$가 기약원이므로 증명이 끝난다.
\end{proof}

\begin{lemma}
\label{more-algebra-lemma-quasi-excellent-nagata}
준우수환은 Nagata 환이다.
\end{lemma}

\begin{proof}
$R$을 준우수환이라 하자. $R$ 위의 유한형 대수가 준우수라는 사실
(보조정리 \ref{more-algebra-lemma-finite-type-over-excellent})을 사용하면,
모든 준우수 정역이 N-1임을 보이는 것으로 충분하다. Algebra, Lemma
\ref{algebra-lemma-check-universally-japanese}를 보라.
Algebra, Lemma \ref{algebra-lemma-characterize-N-1}을 적용하고
(준우수환이 J-2임을 사용하여), 준우수 국소 정역 $R$이 N-1임을
보이는 문제로 환원한다. $R \to R^\wedge$가 정칙이므로 보조정리
\ref{more-algebra-lemma-reduced-goes-up}에 의해 $R^\wedge$는 기약원이다.
다시 말해 $R$은 해석적으로 비분기이다. 따라서 Algebra, Lemma
\ref{algebra-lemma-analytically-unramified-easy}에 의해 $R$은 N-1이다.
\end{proof}

\begin{lemma}
\label{more-algebra-lemma-completion-normal-local-ring}
$(A, \mathfrak m)$를 Noether 국소환이라 하자. $A$가 정규이고
$A$의 형식적 올들이 정규이면(예를 들어 $A$가 우수 또는 준우수이면),
$A^\wedge$는 정규이다.
\end{lemma}

\begin{proof}
Algebra, Lemma \ref{algebra-lemma-normal-goes-up-noetherian}에서
즉시 따른다.
\end{proof}





\section{가군의 아벨 범주}
\label{more-algebra-section-abelian-categories-modules}

\noindent
$R$을 환이라 하자. $R$-가군의 범주 $\text{Mod}_R$은 아벨 범주이다.
다음은 $\text{Mod}_R$의 아벨 부분범주의 몇 가지 예이다(여기서는
Homology, Definition \ref{homology-definition-serre-subcategory}와
Homology, Lemmas \ref{homology-lemma-characterize-serre-subcategory} 및
\ref{homology-lemma-characterize-weak-serre-subcategory}에서 도입한
용어를 사용한다).
\begin{enumerate}
\item 연접 $R$-가군의 범주는 $\text{Mod}_R$의 약한 Serre 부분범주이다.
이는 Algebra, Lemma \ref{algebra-lemma-coherent}에서 따른다.
\item $S \subset R$을 곱셈적 부분집합이라 하자. 각 $s \in S$의
곱셈이 $M$ 위의 동형사상인 $R$-가군 $M$들로 이루어진
충만한 부분범주는 $\text{Mod}_R$의 Serre 부분범주이다. 이는
Algebra, Lemma \ref{algebra-lemma-localization-and-modules}에서 따른다.
\item $I \subset R$을 유한 생성 아이디얼이라 하자. $I$-멱 비틀림
가군들의 충만한 부분범주는 $\text{Mod}_R$의 Serre 부분범주이다.
보조정리 \ref{more-algebra-lemma-I-power-torsion}을 보라.
\item 어떤 문헌에서는 {\it 비틀림 가군}을, 모든 $x \in M$에 대하여
$fx = 0$을 만족하는 영인자가 아닌 원소 $f \in R$가 존재하는 가군
$M$으로 정의한다. 비틀림 가군들의 충만한 부분범주는
$\text{Mod}_R$의 Serre 부분범주이다.
\item $R$이 Noether가 아니면 유한 생성 $R$-가군의 범주
$\text{Mod}^{fg}_R$은 {\bf 아벨 범주가 아니다}. 실제로
$I \subset R$이 유한 생성되지 않은 아이디얼이면 사상 $R \to R/I$은
$\text{Mod}^{fg}_R$에서 핵을 갖지 않는다.
\item $R$이 Noether이면 연접 $R$-가군은 유한 생성(즉, 유한)
$R$-가군과 일치한다. Algebra, Lemmas
\ref{algebra-lemma-Noetherian-coherent},
\ref{algebra-lemma-coherent-ring}, 그리고
\ref{algebra-lemma-Noetherian-finite-type-is-finite-presentation}를 보라.
따라서 위의 (1)에 의해 $\text{Mod}^{fg}_R$은 아벨 범주이다. 더구나
% P02-E0889: frozen source omits a space in ``in fact,in''.
이 경우 범주 $\text{Mod}_R^{fg}$는 $\text{Mod}_R$의 (강한) Serre
부분범주이다.
\end{enumerate}





\section{단사 아벨 군}
\label{more-algebra-section-abelian-groups}

\noindent
이 절에서는 아벨 군의 범주가 충분한 단사 대상을 가짐을 보인다.
아벨 군 $M$이 {\it 가분}이라는 것은, 모든 $x \in M$과 모든
$n \in \mathbf{N}$에 대하여 $n y = x$를 만족하는 $y \in M$이
존재한다는 것과 동치임을 상기하자.

\begin{lemma}
\label{more-algebra-lemma-injective-abelian}
아벨 군 $J$가 아벨 군의 범주에서 단사 대상일 필요충분조건은
$J$가 가분인 것이다.
\end{lemma}

\begin{proof}
$J$가 가분이 아니라고 하자. 그러면 어떤 $x \in J$와
$n \in \mathbf{N}$이 존재하여 $n y = x$를 만족하는 $y \in J$가
존재하지 않는다. 이때 준동형 $\mathbf{Z} \to J$, $m \mapsto mx$는
$\frac{1}{n}\mathbf{Z} \supset \mathbf{Z}$로 연장되지 않는다.
따라서 $J$는 단사 대상이 아니다.

\medskip\noindent
$A \subset B$를 아벨 군이라 하자. $J$가 가분 아벨 군이라고 가정하고
$\varphi : A \to J$를 준동형이라 하자. $A \subset A' \subset B$이고
$\varphi'|_A = \varphi$인 준동형 $\varphi' : A' \to J$들의 집합을
생각하자. $A' \supset A''$이고 $\varphi'|_{A''} = \varphi''$일
필요충분조건으로 $(A', \varphi') \geq (A'', \varphi'')$를 정의한다.
$(A_i, \varphi_i)_{i \in I}$가 그러한 쌍들의 전순서 집합이면,
$a \in A_i$를 $\varphi_i(a)$로 보내는 사상
$\bigcup_{i \in I} A_i \to J$를 얻는다. 따라서 Zorn 보조정리를
적용할 수 있다. 결론을 얻으려면 쌍 $(A', \varphi')$가 극대일 때
$A' = B$임을 보이면 된다. 다시 말해, 임의의 부분군 $A \subset B$,
$A \not = B$와 임의의 $\varphi : A \to J$가 주어졌을 때 다음을
만족하는 $\varphi' : A' \to J$, $A \subset A' \subset B$를 찾으면
충분하다. (a) 포함 $A \subset A'$는 진포함이고, (b) 준동형
$\varphi'$는 $\varphi$를 연장한다.

\medskip\noindent
이를 증명하기 위해 $x \in B$, $x \not \in A$를 택하자.
$nx \in A$를 만족하는 $n\in \mathbf{N}$이 존재하지 않으면
$A \oplus \mathbf{Z} \cong A + \mathbf{Z}x$이다. 따라서 예를 들어
$A$ 위에서는 $\varphi$를 사용하고 $x$를 영으로 보내어
$\varphi$를 $A' = A + \mathbf{Z}x$로 연장할 수 있다.
$nx \in A$를 만족하는 $n \in \mathbf{N}$이 존재하면, 그러한
최소 정수를 $n$이라 하자. $nz = \varphi(nx)$를 만족하는 원소
$z \in J$를 택한다. 다음과 같이 준동형
$\tilde\varphi : A \oplus \mathbf{Z} \to J$를 정의한다.
$(a, m) \mapsto \varphi(a) + mz$. $z$의 선택에 의해
$\tilde \varphi$의 핵은 사상 $A \oplus \mathbf{Z} \to B$,
$(a, m) \mapsto a + mx$의 핵을 포함한다. 따라서 $\tilde \varphi$는
상 $A' = A + \mathbf{Z}x$를 통해 인수분해되며, 이는 준동형
$\varphi$를 연장한다.
\end{proof}

\noindent
이 보조정리를 사용하면 모든 아벨 군이 단사 아벨 군에 매입됨을
보일 수 있다. 그러나 이는 다음 절의 결과의 특수한 경우이다.





\section{단사 가군}
\label{more-algebra-section-injectives-modules}

\noindent
단사 가군에 관한 몇 가지 보조정리를 제시한다.

\begin{definition}
% P02-E0890: frozen label names projectivity although this defines injectivity.
\label{more-algebra-definition-projective}
$R$을 환이라 하자. $R$-가군 $J$가 {\it 단사}일 필요충분조건은
함자 $\Hom_R(-, J) : \text{Mod}_R \to \text{Mod}_R$이 완전 함자인 것이다.
\end{definition}

\noindent
임의의 $R$-가군 $M$에 대하여 함자 $\Hom_R(- , M)$은 왼쪽 완전이다.
Algebra, Lemma \ref{algebra-lemma-hom-exact}를 보라. 따라서 $J$가
단사라는 조건은 실제로 $R$-가군의 단사 사상 $M \to M'$이 주어지면
사상 $\Hom_R(M', J) \to \Hom_R(M, J)$가 전사라는 뜻이다.

\medskip\noindent
이를 ${Ext}$-가군으로 다시 나타내기 전에, $\Ext^1_R(M, N)$과
Homology, Section \ref{homology-section-extensions}의 의미에서의
확장 사이의 관계를 논의하자.

\begin{lemma}
\label{more-algebra-lemma-relation-ext-ext}
$R$을 환이라 하고 $\mathcal{A}$를 $R$-가군의 아벨 범주라 하자.
다음 표준 동형사상이 존재한다.
$\Ext_\mathcal{A}(M, N) = \Ext^1_R(M, N)$. 이는 Algebra, Lemmas
\ref{algebra-lemma-long-exact-seq-ext}와
\ref{algebra-lemma-reverse-long-exact-seq-ext}의 긴 완전열 및
Homology, Lemma \ref{homology-lemma-six-term-sequence-ext}의
$6$-항 완전열과 양립한다.
\end{lemma}

\begin{proof}
생략한다.
\end{proof}

\begin{lemma}
\label{more-algebra-lemma-characterize-injective}
$R$을 환이라 하고 $J$를 $R$-가군이라 하자. 다음은 서로 동치이다.
\begin{enumerate}
\item $J$는 단사이다.
\item 모든 $R$-가군 $M$에 대하여 $\Ext^1_R(M, J) = 0$이다.
\end{enumerate}
\end{lemma}

\begin{proof}
$0 \to M'' \to M' \to M \to 0$을 $R$-가군의 짧은 완전열이라 하자.
Algebra, Lemma \ref{algebra-lemma-reverse-long-exact-seq-ext}의 긴 완전열
$$
\begin{matrix}
0
\to \Hom_R(M, J)
\to \Hom_R(M', J)
\to \Hom_R(M'', J)
\\
\phantom{0\ }
\to \Ext^1_R(M, J)
\to \Ext^1_R(M', J)
\to \Ext^1_R(M'', J)
\to \ldots
\end{matrix}
$$
을 생각하자. 이로부터 (2)가 (1)을 함의함을 알 수 있다. 거꾸로
$J$가 단사이면 Homology, Lemma
\ref{homology-lemma-characterize-injectives}와 보조정리
\ref{more-algebra-lemma-relation-ext-ext}에 의해 $\Ext$-군은 영이다.
\end{proof}

\begin{lemma}
\label{more-algebra-lemma-characterize-injective-bis}
$R$을 환이라 하고 $J$를 $R$-가군이라 하자. 다음은 서로 동치이다.
\begin{enumerate}
\item $J$는 단사이다.
\item 모든 아이디얼 $I \subset R$에 대하여 $\Ext^1_R(R/I, J) = 0$이다.
\item 아이디얼 $I \subset R$와 가군 사상 $I \to J$가 주어지면
그 연장 $R \to J$가 존재한다.
\end{enumerate}
\end{lemma}

\begin{proof}
$I \subset R$이 아이디얼이면 짧은 완전열
$0 \to I \to R \to R/I \to 0$은 다음 완전열을 준다.
$$
\Hom_R(R, J) \to
\Hom_R(I, J) \to
\Ext^1_R(R/I, J) \to 0
$$
이는 Algebra, Lemma \ref{algebra-lemma-reverse-long-exact-seq-ext}와,
$R$이 사영이어서 $\Ext^1_R(R, J) = 0$이라는 사실에서 따른다
(Algebra, Lemma \ref{algebra-lemma-characterize-projective}).
따라서 (2)와 (3)은 동치이다. 이 증명에서는 Baer 판정법이라고 알려진
(1) $\Leftrightarrow$ (3)을 보일 것이다.

\medskip\noindent
(1)을 가정하자. (3)과 같은 가군 사상 $I \to J$가 주어지면 정의에 의해
사상 $\Hom_R(R, J) \to \Hom_R(I, J)$가 전사이므로 연장
$R \to J$를 얻는다.

\medskip\noindent
(3)을 가정하자. $M \subset N$을 $R$-가군의 포함이라 하고
$\varphi : M \to J$를 준동형이라 하자. $\varphi$가 $N$으로 연장됨을
보일 것이며, 이로써 보조정리의 증명이 끝난다. $M \subset M' \subset N$이고
$\varphi'|_M = \varphi$인 준동형 $\varphi' : M' \to J$들의 집합을
생각하자. $M' \supset M''$이고 $\varphi'|_{M''} = \varphi''$일
필요충분조건으로 $(M', \varphi') \geq (M'', \varphi'')$를 정의한다.
$(M_i, \varphi_i)_{i \in I}$가 그러한 쌍들의 전순서 집합이면,
$a \in M_i$를 $\varphi_i(a)$로 보내는 사상
$\bigcup_{i \in I} M_i \to J$를 얻는다. 따라서 Zorn 보조정리를
적용할 수 있다. 결론을 얻으려면 쌍 $(M', \varphi')$가 극대일 때
$M' = N$임을 보이면 된다. 다시 말해, 임의의
% P02-E0891: frozen source says subgroup although M and N are R-modules.
$R$-부분가군 $M \subset N$, $M \not = N$과 임의의
$\varphi : M \to J$가 주어졌을 때 다음을 만족하는
$\varphi' : M' \to J$, $M \subset M' \subset N$를 찾으면 충분하다.
(a) 포함 $M \subset M'$는 진포함이고, (b) 준동형 $\varphi'$는
$\varphi$를 연장한다.

\medskip\noindent
이를 증명하기 위해 $x \in N$, $x \not \in M$을 택하자.
$I = \{f \in R \mid fx \in M\}$이라 두면 이는 $R$의 아이디얼이다.
$f \mapsto \varphi(fx)$로 준동형 $\psi : I \to J$를 정의한다.
가정 (3)에 의해 가능한 사상 $\tilde\psi : R \to J$로 이를 연장하자.
$I$의 선택에 의해 준동형
$M \oplus R \to J$, $(y, f) \mapsto \varphi(y) + \tilde\psi(f)$의 핵은
사상 $M \oplus R \to N$, $(y, f) \mapsto y + fx$의 핵을 포함한다.
따라서 이 준동형은 상 $M' = M + Rx$를 통해 인수분해되며, 이는
원하는 대로 주어진 준동형을 연장한다.
\end{proof}

\noindent
이 절의 나머지에서는 환 $R$ 위에 충분한 단사 가군이 있음을 증명한다.
먼저 $\mathbf{Q}/\mathbf{Z}$가 단사 아벨 군이라는 사실에서 시작한다.
이는 보조정리 \ref{more-algebra-lemma-injective-abelian}에서 따른다.

\begin{definition}
\label{more-algebra-definition-simple-functors}
$R$을 환이라 하자.
\begin{enumerate}
\item 임의의 $R$-가군 $M$에 대하여 그 자연스러운 $R$-가군 구조를
갖춘 $M^\vee = \Hom(M, \mathbf{Q}/\mathbf{Z})$를 나타낸다.
{\it $M \mapsto M^\vee$}를 $R$-가군의 범주에서 자기 자신으로 가는
반변 함자로 생각한다.
\item 임의의 $R$-가군 $M$에 대하여
$$
F(M) = \bigoplus\nolimits_{m \in M} R[m]
$$
으로 $m \in M$인 원소 $[m]$들이 주는 기저를 갖는
{\it 자유 가군}을 나타낸다. $R$-가군의 자연스러운 전사
$F(M)\to M$, $\sum f_i [m_i] \mapsto \sum f_i m_i$를 취한다.
$M \mapsto (F(M) \to M)$을 $R$-가군의 범주에서 $R$-가군의
화살표 범주로 가는 함자로 생각한다.
\end{enumerate}
\end{definition}

\begin{lemma}
\label{more-algebra-lemma-vee-exact}
$R$을 환이라 하자. 함자 $M \mapsto M^\vee$는 완전하다.
\end{lemma}

\begin{proof}
% P02-E0892: frozen source omits ``is'' in ``This because''.
이는 보조정리 \ref{more-algebra-lemma-injective-abelian}에 의해
$\mathbf{Q}/\mathbf{Z}$가 단사 아벨 군이기 때문이다.
\end{proof}

\noindent
평가로 주어지는 표준 사상 $ev : M \to (M^\vee)^\vee$가 존재한다.
즉, $x \in M$이 주어지면 $ev(x) \in (M^\vee)^\vee =
\Hom(M^\vee, \mathbf{Q}/\mathbf{Z})$를 사상
$\varphi \mapsto \varphi(x)$로 정의한다.

\begin{lemma}
\label{more-algebra-lemma-ev-injective}
임의의 $R$-가군 $M$에 대하여 평가 사상
$ev : M \to (M^\vee)^\vee$는 단사이다.
\end{lemma}

\begin{proof}
$\mathbf{Q}/\mathbf{Z}$가 단사 아벨 군이라는 사실을 사용하여 이를
확인할 수 있다. 실제로 $x \in M$이 영이 아니면 $M' \subset M$을
그 원소가 생성하는 순환군이라 하자. 반드시 $x$를 소멸시키지 않는
영이 아닌 사상 $M' \to \mathbf{Q}/\mathbf{Z}$가 존재한다. 이는 사상
$\varphi : M \to \mathbf{Q}/\mathbf{Z}$로 연장되고,
$ev(x)(\varphi) = \varphi(x) \not = 0$이다.
\end{proof}

\noindent
위에서 본 것처럼, $R$-가군의 표준 전사 $F(M) \to M$은
$R$-가군의 표준 단사 사상
$$
(M^\vee)^\vee \longrightarrow (F(M^\vee))^\vee.
$$
으로 바뀐다. $J(M) = (F(M^\vee))^\vee$라 두자.
% P02-E0893: frozen source has the stray phrase ``this the displayed map''.
$ev$와 표시된 사상의 합성은 $M$에 관하여 함자적으로 $M \to J(M)$을 준다.

\begin{lemma}
\label{more-algebra-lemma-JM-injective}
$R$을 환이라 하자. 모든 $R$-가군 $M$에 대하여 $R$-가군
$J(M)$은 단사이다.
\end{lemma}

\begin{proof}
$R$-가군으로서 $J(M) \cong \prod_{\varphi \in M^\vee} R^\vee$임에
주의하자. 단사 가군들의 곱은 단사이므로 $R^\vee$가 단사임을 보이면
충분하다. 이를 위해
$$
\Hom_R(N, R^\vee) =
\Hom_R(N, \Hom_{\mathbf{Z}}(R, \mathbf{Q}/\mathbf{Z})) =
N^\vee
$$
와 보조정리 \ref{more-algebra-lemma-vee-exact}에 의해 $(-)^\vee$가 완전 함자라는
사실을 사용한다.
\end{proof}

\begin{lemma}
\label{more-algebra-lemma-injectives-modules}
$R$을 환이라 하자. 위의 구성은 $R$-가군의 범주에서 $R$-가군의
화살표 범주로 가는 공변 함자 $M \mapsto (M \to J(M))$을 정의하며,
모든 가군 $M$에 대한 출력 $M \to J(M)$은 $M$에서 단사
$R$-가군 $J(M)$으로 가는 단사 사상이다.
\end{lemma}

\begin{proof}
위에서 따른다.
\end{proof}

\noindent
특히 $R$-가군의 임의의 사상 $M \to N$에는 다음 그림을 가환하게 하는
준동형 $J(M) \to J(N)$이 연관된다.
$$
\xymatrix{
M \ar[d] \ar[r] & N \ar[d] \\
J(M) \ar[r] & J(N) }
$$
이것이 일반적으로 갖고 싶은 종류의 구성이다. Homology, Section
\ref{homology-section-injectives}에서 이를 표현하는 용어를 도입하였다.
즉, 이것은 $R$-가군의 범주가 함자적 단사 매장을 가진다는 뜻이라고 한다.










\section{가군의 유도 범주}
\label{more-algebra-section-derived-modules}

\noindent
이 절에서는 환 위 가군의 유도 범주에 관한 몇 가지 일반론을 정리한다.

\medskip\noindent
$A$를 환이라 하자. $A$-가군의 범주를 $\text{Mod}_A$로 나타낸다.
$A$-가군 복합체의 호모토피 범주를 기호 $K(A)$로 나타낼 것이다.
즉, 범주로서 $K(A) = K(\text{Mod}_A)$로 둔다. Derived Categories,
Section \ref{derived-section-homotopy}를 보라. 유계 버전은
$K^+(A)$, $K^-(A)$, $K^b(A)$이다. $K(A)$를 Derived Categories,
Section \ref{derived-section-homotopy-triangulated}에서처럼 삼각 범주로
본다. $A$의 {\it 유도 범주}는 $D(A)$로 나타내며, $K(A)$에서
준동형사상들을 가역화하여 얻는 범주이다. 즉,
$D(A) = D(\text{Mod}_A)$로 둔다. Derived Categories, Section
\ref{derived-section-derived-categories}\footnote{Injectives, Remark
\ref{injectives-remark-existence-D}도 보라.}를 보라. 유계 버전은
$D^+(A)$, $D^-(A)$, $D^b(A)$이다.

\medskip\noindent
$A$를 환이라 하자. $A$-가군의 범주는 곱을 가지며 곱은 완전하다.
보조정리 \ref{more-algebra-lemma-injectives-modules}에 의해 $A$-가군의 범주는
충분한 단사 대상을 가진다. 따라서 모든 $A$-가군 복합체는 K-단사
복합체와 준동형이다(Derived Categories, Lemma
\ref{derived-lemma-enough-K-injectives-Ab4-star}). 이로부터 $D(A)$가
가산 곱(Derived Categories, Lemma \ref{derived-lemma-products})을,
실제로는 임의의 곱(Injectives, Lemma
\ref{injectives-lemma-derived-products})을 가짐을 알 수 있다. 따라서
$D(A)$의 대상들로 이루어진 모든 역계는 (동형을 제외하고 잘 정의되는)
유도 극한을 가진다. Derived Categories, Section
\ref{derived-section-derived-limit}을 보라.

\begin{lemma}
\label{more-algebra-lemma-K-injective-flat}
$R \to S$를 평탄 환 사상이라 하자. $I^\bullet$이 $S$-가군의 K-단사
복합체이면, $I^\bullet$은 $R$-가군의 복합체로서도 K-단사이다.
\end{lemma}

\begin{proof}
Algebra, Lemma \ref{algebra-lemma-adjoint-tensor-restrict}와
$S$와의 텐서곱이 완전하다는 사실에 의해
$\Hom_{K(R)}(M^\bullet, I^\bullet) =
\Hom_{K(S)}(M^\bullet \otimes_R S, I^\bullet)$이기 때문이다.
\end{proof}

\begin{lemma}
\label{more-algebra-lemma-K-injective-epimorphism}
$R \to S$를 환의 전사라 하고 $I^\bullet$을 $S$-가군의 복합체라 하자.
$I^\bullet$이 $R$-가군의 복합체로서 K-단사이면, $I^\bullet$은
$S$-가군의 K-단사 복합체이다.
\end{lemma}

\begin{proof}
임의의 $S$-가군 복합체 $N^\bullet$에 대하여
$\Hom_{K(R)}(N^\bullet, I^\bullet) =
\Hom_{K(S)}(N^\bullet, I^\bullet)$이기 때문이다. Algebra, Lemma
\ref{algebra-lemma-epimorphism-modules}를 보라.
\end{proof}

\begin{lemma}
\label{more-algebra-lemma-hom-K-injective}
$A \to B$를 환 사상이라 하자. $I^\bullet$이 $A$-가군의 K-단사
복합체이면 $\Hom_A(B, I^\bullet)$은 $B$-가군의 K-단사 복합체이다.
\end{lemma}

\begin{proof}
Algebra, Lemma \ref{algebra-lemma-adjoint-hom-restrict}에 의해
$\Hom_{K(B)}(N^\bullet, \Hom_A(B, I^\bullet)) =
\Hom_{K(A)}(N^\bullet, I^\bullet)$이기 때문이다.
\end{proof}









\section{Tor의 계산}
\label{more-algebra-section-computing-tor}

\noindent
$R$을 환이라 하자. $R$-가군의 아벨 범주 $\text{Mod}_R$의 유도 범주를
$D(R)$로 나타낸다. 모든 자유 $R$-가군이 사영이므로
$\text{Mod}_R$은 충분한 사영 대상을 가진다는 점에 주의하자. 따라서
$\text{Mod}_R$에서 임의의 아벨 범주로 가는 모든 가법 함자의 왼쪽
유도 함자를 정의할 수 있다.

\medskip\noindent
이는 특히 함자 $ - \otimes_R M : \text{Mod}_R \to \text{Mod}_R$에
적용된다. 이 함자의 왼쪽 유도 함자는 Tor 함자
$\text{Tor}_i^R(-, M)$이다. Algebra, Section
\ref{algebra-section-tor}를 보라. 다음 전체 왼쪽 유도 함자도 존재한다.
\begin{equation}
\label{more-algebra-equation-derived-tensor-module}
-\otimes_R^{\mathbf{L}} M :
D^{-}(R)
\longrightarrow
D^{-}(R)
\end{equation}
이는 $-\otimes_R^{\mathbf{L}} M$로 나타낸다. 그 위성 함자들은 Tor
가군이다. 즉,
$$
H^{-p}(N \otimes_R^{\mathbf{L}} M) = \text{Tor}_p^R(N, M).
$$

\medskip\noindent
$R$-대수 $A$와의 텐서곱을 생각할 때에는 특별한 상황이 생긴다.
이 경우 $- \otimes_R A$를 $\text{Mod}_R$에서 $\text{Mod}_A$로 가는
함자로 생각한다. 따라서 다음 전체 오른쪽 유도 함자가 존재한다.
% P02-E0894: frozen source says right derived although tensor and L indicate left.
\begin{equation}
\label{more-algebra-equation-derived-tensor-algebra}
-\otimes_R^{\mathbf{L}} A :
D^{-}(R)
\longrightarrow
D^{-}(A)
\end{equation}
이는 $-\otimes_R^{\mathbf{L}} A$로 나타낸다. 그 위성 함자들은 tor
군이다. 즉,
$$
H^{-p}(N \otimes_R^{\mathbf{L}} A) = \text{Tor}_p^R(N, A).
$$
특히 이 Tor 군들은 자연스럽게 $A$-가군 구조를 가진다.

\medskip\noindent
다음 몇 절에서는 이 절의 내용을 비유계 복합체로 일반화할 것이다.





\section{복합체의 텐서곱}
\label{more-algebra-section-symmetric-monoidal}

\noindent
$R$을 환이라 하자. $R$-가군 복합체의 범주 $\text{Comp}(R)$은
대칭 모노이드 구조를 가진다. 실제로 $R$-가군의 두 복합체
$L^\bullet$과 $M^\bullet$이 주어졌다고 하자. Homology, Example
\ref{homology-example-double-complex-as-tensor-product-of}와 Homology,
Definition \ref{homology-definition-associated-simple-complex}을 사용하면
세 번째 $R$-가군 복합체
$$
\text{Tot}(L^\bullet \otimes_R M^\bullet)
$$
를 얻는다. 이 구성은 명백히 $L^\bullet$과 $M^\bullet$ 모두에 대하여
함자적이다. 결합 제약은 Homology, Remark
\ref{homology-remark-triple-complex}에서 삼중 복합체
$K^\bullet \otimes_R L^\bullet \otimes_R M^\bullet$로부터 구성한
복합체의 표준 동형사상
$$
\text{Tot}(\text{Tot}(K^\bullet \otimes_R L^\bullet) \otimes_R M^\bullet)
\longrightarrow
\text{Tot}(K^\bullet \otimes_R \text{Tot}(L^\bullet \otimes_R M^\bullet))
$$
이다. 교환 제약은 직합 성분 $L^p \otimes_R M^q$ 위에서 부호
$(-1)^{pq}$를 사용하는 표준 동형사상
$$
\text{Tot}(L^\bullet \otimes_R M^\bullet) \to
\text{Tot}(M^\bullet \otimes_R L^\bullet)
$$
이다. 이것이 복합체의 사상임을 보기 위해 $x \in L^p$와
$y \in M^q$에 대하여 다음을 계산한다.
$$
\text{d}(x \otimes y) =
\text{d}_L(x) \otimes y + (-1)^px \otimes \text{d}_M(y)
$$
우리의 규칙에 따르면 우변은 다음으로 보내진다.
$$
(-1)^{(p + 1)q}y \otimes \text{d}_L(x) +
(-1)^{p + p(q + 1)} \text{d}_M(y) \otimes x
$$
한편,
$$
\text{d}((-1)^{pq}y \otimes x) =
(-1)^{pq} \text{d}_M(y) \otimes x +
(-1)^{pq + q} y \otimes \text{d}_L(x)
$$
임을 알 수 있다. 이 두 식은 바로 확인하면 일치하므로 원하는 결론을 얻는다.

\begin{lemma}
\label{more-algebra-lemma-symmetric-monoidal-cat-complexes}
$R$을 환이라 하자. 위의 함자
$(L^\bullet, M^\bullet) \mapsto \text{Tot}(L^\bullet \otimes_R M^\bullet)$와
결합 및 교환 제약을 갖춘 $R$-가군 복합체의 범주 $\text{Comp}(R)$은
대칭 모노이드 범주이다.
\end{lemma}

\begin{proof}
생략한다. 힌트: 단위원 $\mathbf{1}$로는 차수 $0$에 $R$을 가지고
다른 차수에서는 영인 복합체를 택하며, 다음 자명한 동형사상들을 사용한다.
$\text{Tot}(\mathbf{1} \otimes_R M^\bullet) = M^\bullet$ 및
$\text{Tot}(K^\bullet \otimes_R \mathbf{1}) = K^\bullet$.
% P02-E0895: frozen source begins the next sentence with lowercase ``to''.
보조정리를 증명하려면 여러 그림의 가환성을 확인해야 한다. Categories,
Definitions \ref{categories-definition-monoidal-category}와
\ref{categories-definition-symmetric-monoidal-category}를 보라.
각 경우의 확인은 직접적이다.
\end{proof}

\begin{lemma}
\label{more-algebra-lemma-derived-tor-homotopy}
$R$을 환이라 하고 $P^\bullet$을 $R$-가군의 복합체라 하자.
$\alpha, \beta : L^\bullet \to M^\bullet$를 호모토픽인 복합체 사상이라
하자. 그러면 $\alpha$와 $\beta$는 다음 호모토픽인 사상들을 유도한다.
$$
\text{Tot}(\alpha \otimes \text{id}_P),
\text{Tot}(\beta \otimes \text{id}_P) :
\text{Tot}(L^\bullet \otimes_R P^\bullet)
\longrightarrow
\text{Tot}(M^\bullet \otimes_R P^\bullet).
$$
특히 구성 $L^\bullet \mapsto \text{Tot}(L^\bullet \otimes_R P^\bullet)$은
복합체의 호모토피 범주의 자기함자를 정의한다.
\end{lemma}

\begin{proof}
$h^n : L^n \to M^{n - 1}$로 정의되는 어떤 호모토피 $h$에 대하여
$\alpha = \beta + dh + hd$라고 하자. 다음과 같이 둔다.
$$
H^n = \bigoplus\nolimits_{a + b = n} h^a \otimes \text{id}_{P^b} :
\bigoplus\nolimits_{a + b = n} L^a \otimes_R P^b
\longrightarrow
\bigoplus\nolimits_{a + b = n} M^{a - 1} \otimes_R P^b
$$
그러면 직접 계산하여 사상
$\text{Tot}(L^\bullet \otimes_R P^\bullet) \to
\text{Tot}(M^\bullet \otimes_R P^\bullet)$로서
$$
\text{Tot}(\alpha \otimes \text{id}_P) =
\text{Tot}(\beta \otimes \text{id}_P) + dH + Hd
$$
임을 알 수 있다.
\end{proof}

\begin{lemma}
\label{more-algebra-lemma-symmetric-monoidal-homotopy-category}
$R$을 환이라 하자. 위의 함자
$(L^\bullet, M^\bullet) \mapsto \text{Tot}(L^\bullet \otimes_R M^\bullet)$와
결합 및 교환 제약을 갖춘 $R$-가군 복합체의 호모토피 범주 $K(R)$은
대칭 모노이드 범주이다.
\end{lemma}

\begin{proof}
보조정리 \ref{more-algebra-lemma-symmetric-monoidal-cat-complexes}와
\ref{more-algebra-lemma-derived-tor-homotopy}에서 따른다. 자세한 내용은 생략한다.
\end{proof}

\begin{lemma}
\label{more-algebra-lemma-derived-tor-exact}
$R$을 환이라 하고 $P^\bullet$을 $R$-가군의 복합체라 하자. 함자
$$
K(R) \longrightarrow K(R), \quad
L^\bullet \longmapsto \text{Tot}(P^\bullet \otimes_R L^\bullet)
$$
와
$$
K(R) \longrightarrow K(R), \quad
L^\bullet \longmapsto \text{Tot}(L^\bullet \otimes_R P^\bullet)
$$
는 삼각 범주의 완전 함자이다.
\end{lemma}

\begin{proof}
Derived Categories, Remark
\ref{derived-remark-double-complex-as-tensor-product-of}에서 따른다.
\end{proof}




\section{유도 텐서곱}
\label{more-algebra-section-derived-tensor-product}

\noindent
유도 텐서곱을 더 일반적으로 구성할 수 있다. 실제로 K-평탄 분해를
택하면 유계성 가정이 필요하지 않은 것으로 밝혀진다.

\begin{definition}
\label{more-algebra-definition-K-flat}
$R$을 환이라 하자. 모든 비순환 복합체 $M^\bullet$에 대하여 전체 복합체
$\text{Tot}(M^\bullet \otimes_R K^\bullet)$가 비순환이면 복합체
$K^\bullet$을 {\it K-평탄}이라고 한다.
\end{definition}

\begin{lemma}
\label{more-algebra-lemma-K-flat-quasi-isomorphism}
$R$을 환이라 하고 $K^\bullet$을 K-평탄 복합체라 하자. 그러면 함자
$$
K(R) \longrightarrow K(R), \quad
L^\bullet \longmapsto \text{Tot}(L^\bullet \otimes_R K^\bullet)
$$
는 준동형사상을 준동형사상으로 보낸다.
\end{lemma}

\begin{proof}
보조정리 \ref{more-algebra-lemma-derived-tor-exact}와, $K(R)$의 준동형사상은 그
원뿔이 비순환이라는 것으로 특징지어진다는 사실에서 따른다.
\end{proof}

\begin{lemma}
\label{more-algebra-lemma-base-change-K-flat}
$R \to R'$을 환 사상이라 하자. $K^\bullet$이 $R$-가군의 K-평탄
복합체이면 $K^\bullet \otimes_R R'$은 $R'$-가군의 K-평탄 복합체이다.
\end{lemma}

\begin{proof}
정의와, 임의의 $R'$-가군 복합체 $L^\bullet$에 대한 등식
$(K^\bullet \otimes_R R') \otimes_{R'} L^\bullet =
K^\bullet \otimes_R L^\bullet$에서 따른다.
\end{proof}

\begin{lemma}
\label{more-algebra-lemma-tensor-product-K-flat}
$R$을 환이라 하자. $K^\bullet$, $L^\bullet$이 $R$-가군의 K-평탄
복합체이면 $\text{Tot}(K^\bullet \otimes_R L^\bullet)$은 $R$-가군의
K-평탄 복합체이다.
\end{lemma}

\begin{proof}
동형사상
$$
\text{Tot}(M^\bullet \otimes_R \text{Tot}(K^\bullet \otimes_R L^\bullet))
=
\text{Tot}(\text{Tot}(M^\bullet \otimes_R K^\bullet) \otimes_R L^\bullet)
$$
과 정의에서 따른다.
\end{proof}

\begin{lemma}
\label{more-algebra-lemma-K-flat-two-out-of-three}
$R$을 환이라 하자. $(K_1^\bullet, K_2^\bullet, K_3^\bullet)$이
$K(R)$의 특수 삼각형이라고 하자. $K_i^\bullet$ 가운데 둘이 K-평탄이면
나머지 하나도 K-평탄이다.
\end{lemma}

\begin{proof}
보조정리 \ref{more-algebra-lemma-derived-tor-exact}와, $K(R)$의 특수 삼각형에서
셋 가운데 둘이 비순환이면 나머지 하나도 비순환이라는 사실에서 따른다.
\end{proof}

\begin{lemma}
\label{more-algebra-lemma-K-flat-two-out-of-three-ses}
$R$을 환이라 하자. 다음을 복합체의 짧은 완전열이라 하자.
$0 \to K_1^\bullet \to K_2^\bullet \to K_3^\bullet \to 0$.
모든 $n \in \mathbf{Z}$에 대하여 $K_3^n$이 평탄이고
$K_i^\bullet$ 가운데 둘이 K-평탄이면 나머지 하나도 K-평탄이다.
\end{lemma}

\begin{proof}
$L^\bullet$을 $R$-가군의 복합체라 하자. 그러면
$$
0 \to
\text{Tot}(L^\bullet \otimes_R K_1^\bullet) \to
\text{Tot}(L^\bullet \otimes_R K_2^\bullet) \to
\text{Tot}(L^\bullet \otimes_R K_3^\bullet) \to 0
$$
은 복합체의 짧은 완전열이다. 실제로 각 $n, m$에 대하여 가군열
$0 \to L^n \otimes_R K_1^m \to
L^n \otimes_R K_2^m \to
L^n \otimes_R K_3^m \to 0$은 Algebra, Lemma
\ref{algebra-lemma-flat-tor-zero}에 의해 완전하고, 위 복합체열은
이 열들의 직합이다. 따라서 이 사실과, 복합체의 짧은 완전열에서
셋 가운데 둘이 비순환이면 나머지 하나도 비순환이라는 사실로부터
보조정리가 따른다.
\end{proof}

\begin{lemma}
\label{more-algebra-lemma-derived-tor-quasi-isomorphism}
$R$을 환이라 하고 $P^\bullet$을 평탄 $R$-가군의 위로 유계인
복합체라 하자. 그러면 $P^\bullet$은 K-평탄이다.
\end{lemma}

\begin{proof}
$L^\bullet$을 $R$-가군의 비순환 복합체라 하자.
$\xi \in H^n(\text{Tot}(L^\bullet \otimes_R P^\bullet))$라 하자.
$\xi = 0$임을 보여야 한다. $\text{Tot}^n(L^\bullet \otimes_R P^\bullet)$은
$L^a \otimes_R P^b$들을 항으로 하는 직합이므로, 어떤
$m \in \mathbf{Z}$에 대하여 $\xi$는
$H^n(\text{Tot}(\tau_{\leq m}L^\bullet \otimes_R P^\bullet))$의 원소에서
온다. $\tau_{\leq m}L^\bullet$도 비순환이므로 $L^\bullet$을
$\tau_{\leq m}L^\bullet$로 바꾸어도 된다. 따라서 $L^\bullet$이
위로 유계라고 가정해도 된다. 이 경우 Homology, Lemma
\ref{homology-lemma-first-quadrant-ss}의 스펙트럼 열은
$$
{}'E_1^{p, q} = H^p(L^\bullet \otimes_R P^q)
$$
를 가지며, $P^q$가 평탄이고 $L^\bullet$이 비순환이므로 이는 영이다.
따라서 $H^*(\text{Tot}(L^\bullet \otimes_R P^\bullet)) = 0$이다.
\end{proof}

\noindent
다음 보조정리에서 복합체 계의 여극한은 항별 여극한을 뜻한다.

\begin{lemma}
\label{more-algebra-lemma-colimit-K-flat}
$R$을 환이라 하고 $K_1^\bullet \to K_2^\bullet \to \ldots$를
K-평탄 복합체의 계라 하자. 그러면 $\colim_i K_i^\bullet$은
K-평탄이다. 더 일반적으로 K-평탄 복합체의 모든 여과 여극한은
K-평탄이다.
\end{lemma}

\begin{proof}
항별 여극한을 취하고 있으므로 Algebra, Lemma
\ref{algebra-lemma-tensor-products-commute-with-limits}에 의해
$$
\colim_i \text{Tot}(M^\bullet \otimes_R K_i^\bullet) =
\text{Tot}(M^\bullet \otimes_R \colim_i K_i^\bullet)
$$
이다. 따라서 여과 여극한이 완전하다는 사실에서 보조정리가 따른다.
Algebra, Lemma \ref{algebra-lemma-directed-colimit-exact}를 보라.
\end{proof}

\begin{lemma}
\label{more-algebra-lemma-universally-acyclic-K-flat}
$R$을 환이라 하고 $K^\bullet$을 $R$-가군의 복합체라 하자.
모든 유한 표시 $R$-가군 $M$에 대하여 $K^\bullet \otimes_R M$이
비순환이면 $K^\bullet$은 K-평탄이다.
\end{lemma}

\begin{proof}
% P02-E0896: frozen source has ``tensor product commute''.
텐서곱이 여극한과 교환한다는 사실을 반복해서 사용할 것이다
(Algebra, Lemma \ref{algebra-lemma-tensor-products-commute-with-limits}).
따라서 임의의 $R$-가군 $M$에 대하여 $K^\bullet \otimes_R M$은
비순환이다. 실제로 모든 $R$-가군은 유한 표시 $R$-가군 $M$들의
% P02-E0200: frozen source leaves a stray repeated formal variable M.
여과 여극한이다. Algebra, Lemma \ref{algebra-lemma-module-colimit-fp}를
보라. $M^\bullet$을 $R$-가군의 비순환 복합체라 하자.
$\text{Tot}(M^\bullet \otimes_R K^\bullet)$이 비순환임을 보여야 한다.
$M^\bullet = \colim \tau_{\leq n} M^\bullet$(항별 여극한)이므로
$$
\text{Tot}(M^\bullet \otimes_R K^\bullet) =
\colim \text{Tot}(\tau_{\leq n} M^\bullet \otimes_R K^\bullet)
$$
이다. 절단은 Homology, Section \ref{homology-section-truncations}에서와
같다. 여과 여극한이 완전하므로(Algebra, Lemma
\ref{algebra-lemma-directed-colimit-exact}), $M^\bullet$을
$\tau_{\leq n}M^\bullet$로 바꾸고 $M^\bullet$이 위로 유계라고
가정해도 된다. 위로 유계인 경우에는
$M^\bullet = \colim \sigma_{\geq -n} M^\bullet$로 쓸 수 있다.
여기서 복합체 $\sigma_{\geq -n} M^\bullet$들은 유계이지만 더 이상
비순환이 아닐 수 있다. 위와 같이 논증하면 $M^\bullet$이 유계 복합체인
경우로 환원된다. 끝으로 유계 복합체 $M^a \to \ldots \to M^b$에
대하여 복합체의 길이 $b - a$에 관한 귀납법을 쓸 수 있다.
% P02-E0897: frozen source starts the induction at b-a=1 although one term is b-a=0.
$b - a = 1$인 경우는 위에서 보았다. $b - a > 1$이면 복합체의 분할
짧은 완전열
$$
0 \to \sigma_{\geq a + 1}M^\bullet \to M^\bullet \to M^a[-a] \to 0
$$
을 생각하고 보조정리 \ref{more-algebra-lemma-derived-tor-exact}를 적용하여 귀납
단계를 수행한다. 일부 세부사항은 생략한다.
\end{proof}

\begin{lemma}
\label{more-algebra-lemma-K-flat-resolution}
$R$을 환이라 하자. 임의의 복합체 $M^\bullet$에 대하여, 각 항이 평탄
$R$-가군인 K-평탄 복합체 $K^\bullet$과 항별 전사인 준동형사상
$K^\bullet \to M^\bullet$이 존재한다.
\end{lemma}

\begin{proof}
$\mathcal{P} \subset \Ob(\text{Mod}_R)$를 평탄 $R$-가군의 모임이라 하자.
Derived Categories, Lemma \ref{derived-lemma-special-direct-system}에 의해
계 $K_1^\bullet \to K_2^\bullet \to \ldots$와 그림
$$
\xymatrix{
K_1^\bullet \ar[d] \ar[r] &
K_2^\bullet \ar[d] \ar[r] & \ldots \\
\tau_{\leq 1}M^\bullet \ar[r] &
\tau_{\leq 2}M^\bullet \ar[r] & \ldots
}
$$
이 존재하고, 그 보조정리에 열거한 성질 (1), (2), (3)을 만족한다.
이 성질들에 의해 각 복합체 $K_i^\bullet$은 평탄 가군의 위로 유계인
복합체이다. 따라서 보조정리 \ref{more-algebra-lemma-derived-tor-quasi-isomorphism}에
의해 $K_i^\bullet$은 K-평탄이다. 유도된 사상
$\colim_i K_i^\bullet \to M^\bullet$은 구성에 의해 준동형사상이고
항별 전사이다. 복합체 $\colim_i K_i^\bullet$은 보조정리
\ref{more-algebra-lemma-colimit-K-flat}에 의해 K-평탄이다. 여과 여극한의 평탄
가군은 평탄이므로 각 항 $\colim K_i^n$은 평탄이다. Algebra, Lemma
\ref{algebra-lemma-colimit-flat}을 보라.
\end{proof}

\begin{remark}
\label{more-algebra-remark-P-resolution}
사실 보조정리 \ref{more-algebra-lemma-K-flat-resolution}보다 더 나은 결과를 얻을 수
있다. 즉, $R$-가군의 복합체 $P^\bullet$이 부분복합체들에 의한 여과
$$
0 = F_{-1}P^\bullet \subset F_0P^\bullet \subset
F_1P^\bullet \subset \ldots \subset P^\bullet
$$
를 갖고 다음을 만족하도록 준동형사상 $P^\bullet \to M^\bullet$을
찾을 수 있다.
\begin{enumerate}
\item $P^\bullet = \bigcup F_pP^\bullet$이다.
\item 포함 $F_iP^\bullet \to F_{i + 1}P^\bullet$은 항별 분할 단사이다.
\item 몫 $F_{i + 1}P^\bullet/F_iP^\bullet$은 이동된 복합체 $R[k]$들의
직합과 동형이다(복합체로서이므로 미분은 영이다).
\end{enumerate}
이는 Differential Graded Algebra, Lemma \ref{dga-lemma-resolve}에서
보일 것이다(또는 보조정리 \ref{more-algebra-lemma-resolution-filtered-complex}의
증명에서처럼 논증할 수 있다). 더구나 그러한 복합체가 주어지면
$D(R)$에서 특수 삼각형
$$
\bigoplus F_iP^\bullet \to \bigoplus F_iP^\bullet \to M^\bullet
\to \bigoplus F_iP^\bullet[1]
$$
을 얻는다. 이를 사용하면 때때로 일반 복합체에 관한 명제를
$R[k]$에 관한 명제로 환원할 수 있다(물론 이 방법은 그 명제가 직합을
취해도 보존될 때에만 작동한다). 더 정확히, $T$를 $D(R)$의 대상들이
가질 수 있는 성질이라 하자. 다음을 가정하자.
\begin{enumerate}
\item $K_i \in D(R)$, $i \in I$가 모든 $i \in I$에 대하여
$T(K_i)$를 만족하는 대상족이면 $T(\bigoplus K_i)$가 성립한다.
\item $K \to L \to M \to K[1]$이 특수 삼각형이고 두 대상에 대하여
$T$가 성립하면 세 번째 대상에 대해서도 $T$가 성립한다.
\item 모든 $k$에 대하여 $T(R[k])$가 성립한다.
\end{enumerate}
그러면 $D(R)$의 모든 대상에 대하여 $T$가 성립한다.
\end{remark}

\begin{lemma}
\label{more-algebra-lemma-derived-tor-quasi-isomorphism-other-side}
$R$을 환이라 하자. $\alpha : P^\bullet \to Q^\bullet$를 $R$-가군의
K-평탄 복합체들 사이의 준동형사상이라 하자. 모든 $R$-가군 복합체
$L^\bullet$에 대하여 유도된 사상
$$
\text{Tot}(\text{id}_L \otimes \alpha) :
\text{Tot}(L^\bullet \otimes_R P^\bullet)
\longrightarrow
\text{Tot}(L^\bullet \otimes_R Q^\bullet)
$$
은 준동형사상이다.
\end{lemma}

\begin{proof}
보조정리 \ref{more-algebra-lemma-K-flat-resolution}에 따라 $K^\bullet$이 K-평탄인
준동형사상 $K^\bullet \to L^\bullet$을 택한다. 가환 그림
$$
\xymatrix{
\text{Tot}(K^\bullet \otimes_R P^\bullet) \ar[r] \ar[d] &
\text{Tot}(K^\bullet \otimes_R Q^\bullet) \ar[d] \\
\text{Tot}(L^\bullet \otimes_R P^\bullet) \ar[r] &
\text{Tot}(L^\bullet \otimes_R Q^\bullet)
}
$$
을 생각하자. 보조정리 \ref{more-algebra-lemma-K-flat-quasi-isomorphism}에 의해
두 수직 화살표와 위쪽 수평 화살표가 준동형사상이므로 결론이 따른다.
\end{proof}

\noindent
$R$을 환이라 하고 $M^\bullet$을 $D(R)$의 대상이라 하자.
보조정리 \ref{more-algebra-lemma-K-flat-resolution}에 따라 K-평탄 분해
$K^\bullet \to M^\bullet$을 택한다. 보조정리
\ref{more-algebra-lemma-derived-tor-homotopy}와 \ref{more-algebra-lemma-derived-tor-exact}에 의해
삼각 범주의 완전 함자
$$
K(R) \longrightarrow K(R), \quad
L^\bullet \longmapsto \text{Tot}(L^\bullet \otimes_R K^\bullet)
$$
를 얻는다. $D(R)$은 $K(R)$을 준동형사상들에 대하여 국소화한 것이므로,
보조정리 \ref{more-algebra-lemma-K-flat-quasi-isomorphism}에 의해 이 함자는 단순히
함자 $D(R) \to D(R)$를 유도한다. $M^\bullet$의 $K$-평탄 분해들의 범주는
쌍대여과이고 보조정리
\ref{more-algebra-lemma-derived-tor-quasi-isomorphism-other-side}가 성립하므로,
그 결과로 얻는 함자는 (동형을 제외하면) $K^\bullet$의 선택에 의존하지
않는다.

\begin{definition}
\label{more-algebra-definition-derived-tor}
$R$을 환이라 하고 $M^\bullet$을 $D(R)$의 대상이라 하자. {\it 유도 텐서곱}
$$
- \otimes_R^{\mathbf{L}} M^\bullet : D(R) \longrightarrow D(R)
$$
은 위에서 설명한 삼각 범주의 완전 함자이다.
\end{definition}

\noindent
이 함자는 함자 (\ref{more-algebra-equation-derived-tensor-module})를 연장한다.
우리의 명시적 구성으로부터, $L^\bullet$과 $M^\bullet$이 모두
$D(R)$에 속할 때마다 동형사상(부호의 선택이 개입한다. 아래를 보라)
$$
M^\bullet \otimes_R^{\mathbf{L}} L^\bullet
\cong
L^\bullet \otimes_R^{\mathbf{L}} M^\bullet
$$
이 존재함은 명백하다. 따라서 $M^\bullet \otimes_R^{\mathbf{L}} L^\bullet$이라고
쓸 때에는 보통 유도 텐서곱의 정의에 어느 변수를 사용하는지 구별하지
않는다.

\begin{lemma}
% P02-E0202: frozen permanent label misspells ``double'' as ``douoble''.
\label{more-algebra-lemma-flip-douoble-tensor-product}
$R$을 환이라 하고 $K^\bullet, L^\bullet$을 $R$-가군의 복합체라 하자.
두 복합체 모두에 대하여 함자적이고, 사상
$K^p \otimes_R L^q \to L^q \otimes_R K^p$에서 부호 $(-1)^{pq}$를
사용하는 표준 동형사상
$$
K^\bullet \otimes_R^\mathbf{L} L^\bullet \longrightarrow
L^\bullet \otimes_R^\mathbf{L} K^\bullet
$$
이 존재한다(설명은 증명을 보라).
\end{lemma}

\begin{proof}
복합체들을 K-평탄 복합체 $K^\bullet$과 $L^\bullet$로 바꾸어도 되고,
실제로 그렇게 한 뒤 절 \ref{more-algebra-section-symmetric-monoidal}에서 논의한
교환 제약을 사용한다.
\end{proof}

\begin{lemma}
\label{more-algebra-lemma-triple-tensor-product}
$R$을 환이라 하고 $K^\bullet, L^\bullet, M^\bullet$을 $R$-가군의
복합체라 하자. 세 복합체 모두에 대하여 함자적인 표준 동형사상
$$
(K^\bullet \otimes_R^\mathbf{L} L^\bullet) \otimes_R^\mathbf{L} M^\bullet
=
K^\bullet \otimes_R^\mathbf{L} (L^\bullet \otimes_R^\mathbf{L} M^\bullet)
$$
이 존재한다.
\end{lemma}

\begin{proof}
복합체들을 K-평탄 복합체로 바꾸고 절
\ref{more-algebra-section-symmetric-monoidal}의 결합 제약을 사용한다.
\end{proof}

\begin{lemma}
\label{more-algebra-lemma-factor-through-K-flat}
$R$을 환이라 하고 $a : K^\bullet \to L^\bullet$를 $R$-가군 복합체의
사상이라 하자. $K^\bullet$이 K-평탄이면 복합체 $N^\bullet$과 복합체
사상 $b : K^\bullet \to N^\bullet$ 및 $c : N^\bullet \to L^\bullet$이
존재하여 다음을 만족한다.
\begin{enumerate}
\item $N^\bullet$은 K-평탄이다.
\item $c$는 준동형사상이다.
\item $a$는 $c \circ b$와 호모토픽이다.
\end{enumerate}
$K^\bullet$의 항들이 평탄이면, $N^\bullet$에 대해서도 같은 성질이
성립하도록 $N^\bullet$, $b$, $c$를 택할 수 있다.
\end{lemma}

\begin{proof}
호모토피 범주 $K(R)$이 삼각 범주라는 사실을 사용할 것이다.
Derived Categories, Proposition
\ref{derived-proposition-homotopy-category-triangulated}를 보라.
특수 삼각형 $K^\bullet \to L^\bullet \to C^\bullet \to K^\bullet[1]$을
택하자. 보조정리 \ref{more-algebra-lemma-K-flat-resolution}에 따라, $M^\bullet$이
평탄한 항들을 갖는 K-평탄 복합체인 준동형사상
$M^\bullet \to C^\bullet$을 택하자. 삼각 범주의 공리들에 의해 합성
$M^\bullet \to C^\bullet \to K^\bullet[1]$을 특수 삼각형
$K^\bullet \to N^\bullet \to M^\bullet \to K^\bullet[1]$에 넣을 수 있다.
보조정리 \ref{more-algebra-lemma-K-flat-two-out-of-three}에 의해 $N^\bullet$은
K-평탄이다. 다시 삼각 범주의 공리들을 사용하면 다음 특수 삼각형의
사상에 들어가는 사상 $N^\bullet \to L^\bullet$을 택할 수 있다.
$$
\xymatrix{
K^\bullet \ar[r] \ar[d] &
N^\bullet \ar[r] \ar[d] &
M^\bullet \ar[r] \ar[d] &
K^\bullet[1] \ar[d] \\
K^\bullet \ar[r] &
L^\bullet \ar[r] &
C^\bullet \ar[r] &
K^\bullet[1]
}
$$
화살표들 가운데 둘이 준동형사상이므로, 이 특수 삼각형들에 연관된
코호몰로지의 긴 완전열에 의해 세 번째 화살표
$N^\bullet \to L^\bullet$도 준동형사상이다(원한다면 이 그림의
$D(R)$에서의 상을 보고 Derived Categories, Lemma
\ref{derived-lemma-third-isomorphism-triangle}을 사용해도 된다).
이로써 (1), (2), (3)의 증명이 끝난다. 마지막 명제를 증명하기 위해서는
Derived Categories, Lemma
\ref{derived-lemma-improve-distinguished-triangle-homotopy}에 따라
$N^n \cong M^n \oplus K^n$이 되도록 $N^\bullet$을 택할 수 있다.
따라서 $K^\bullet$의 항들이 평탄이면 원하는 평탄성을 얻는다.
\end{proof}








\section{환의 유도 변경}
\label{more-algebra-section-derived-base-change}

\noindent
$R \to A$를 환 사상이라 하고 $N^\bullet$을 $A$-가군의 복합체라 하자.
K-평탄 분해를 사용하여 함자
$$
- \otimes_R^{\mathbf{L}} N^\bullet : D(R) \to D(A)
$$
를 함자 $K(R) \to K(A)$,
$M^\bullet \mapsto \text{Tot}(M^\bullet \otimes_R N^\bullet)$의 왼쪽
유도 함자로 정의할 수도 있다. 특히 $N^\bullet = A[0]$으로 두면
함자 (\ref{more-algebra-equation-derived-tensor-algebra})를 연장하는 유도 밑변경 함자
$$
- \otimes_R^{\mathbf{L}} A : D(R) \to D(A)
$$
를 얻는다. 실제로 모든 $R$-가군 복합체 $M^\bullet$에 대하여
K-평탄 분해 $K^\bullet \to M^\bullet$을 택하고
$$
M^\bullet \otimes_R^{\mathbf{L}} N^\bullet =
\text{Tot}(K^\bullet \otimes_R N^\bullet).
$$
로 둘 수 있다. 이것이 잘 정의됨은 보조정리
\ref{more-algebra-lemma-K-flat-resolution}과
\ref{more-algebra-lemma-derived-tor-quasi-isomorphism-other-side}을 사용하여 확인할
수 있다. 그러나 모든 세부사항을 빠짐없이 확인하려면 약간의 일반 이론을
사용하는 편이 더 편리할 수 있다.

\begin{lemma}
\label{more-algebra-lemma-derived-base-change}
위 구성은 선택에 의존하지 않으며 삼각 범주의 완전 함자
$- \otimes_R^\mathbf{L} N^\bullet : D(R) \to D(A)$를 정의한다.
$D(R)$의 $E^\bullet$에 대하여 함자적 동형사상
$$
E^\bullet \otimes_R^\mathbf{L} N^\bullet =
(E^\bullet \otimes_R^\mathbf{L} A) \otimes_A^\mathbf{L} N^\bullet
$$
이 존재한다.
\end{lemma}

\begin{proof}
유도 함자 $- \otimes_R^\mathbf{L} N^\bullet$의 존재를 증명하기 위해
Derived Categories, Section \ref{derived-section-derived-functors}에서
전개한 일반 이론을 사용한다. $\mathcal{D} = K(R)$ 및
$\mathcal{D}' = D(A)$로 두자. 규칙
$F(M^\bullet) = \text{Tot}(M^\bullet \otimes_R N^\bullet)$로 정의되는
삼각 범주의 완전 함자를 $F : \mathcal{D} \to \mathcal{D}'$로 쓰자.
$F$가 서술한 성질들을 가짐을 증명하려면 보조정리
\ref{more-algebra-lemma-derived-tor-homotopy}와 \ref{more-algebra-lemma-derived-tor-exact}를 사용한다.
$S$를 $\mathcal{D} = K(R)$의 준동형사상들의 집합이라 하자. 그러면
Derived Categories, Situation \ref{derived-situation-derived-functor}의
상황이므로 Derived Categories, Definition
\ref{derived-definition-right-derived-functor-defined}을 적용할 수 있다.
$LF$가 모든 곳에서 정의된다고 주장한다. 이는
$\mathcal{P} \subset \Ob(\mathcal{D})$를 K-평탄 복합체의 모임으로
잡고 Derived Categories, Lemma
\ref{derived-lemma-find-existence-computes}를 적용하면 따른다. 즉,
(1)은 보조정리 \ref{more-algebra-lemma-K-flat-resolution}에서, (2)는 보조정리
\ref{more-algebra-lemma-derived-tor-quasi-isomorphism-other-side}에서 따른다.
따라서 유도 함자
$$
LF : D(R) = S^{-1}\mathcal{D} \longrightarrow \mathcal{D}' = D(A)
$$
를 얻는다. Derived Categories, Equation
(\ref{derived-equation-everywhere})을 보라. 끝으로 Derived Categories,
Lemma \ref{derived-lemma-find-existence-computes}는 $K^\bullet$이
K-평탄일 때
$LF(K^\bullet) = F(K^\bullet) =
\text{Tot}(K^\bullet \otimes_R N^\bullet)$임을 보장한다. 즉, 실제로
$LF$는 위에서 설명한 방식으로 계산된다. 더구나 보조정리
\ref{more-algebra-lemma-base-change-K-flat}에 의해 복합체
$K^\bullet \otimes_R A$는 $A$-가군의 K-평탄 복합체이다. 따라서
% P02-E0898: frozen source tensors the merely R-complex K over A in the last terms.
$$
(K^\bullet \otimes_R^\mathbf{L} A) \otimes_A^\mathbf{L} N^\bullet =
\text{Tot}((K^\bullet \otimes_R A) \otimes_A N^\bullet) =
\text{Tot}(K^\bullet \otimes_A N^\bullet) =
K^\bullet \otimes_A^\mathbf{L} N^\bullet
$$
이며, 이는 보조정리의 마지막 명제를 증명한다.
\end{proof}

\begin{lemma}
\label{more-algebra-lemma-functoriality-derived-base-change}
$R \to A$를 환 사상이라 하고 $f : L^\bullet \to N^\bullet$를
$A$-가군 복합체의 사상이라 하자. 그러면 $f$는 함자의 변환
% P02-E0899: frozen display uses tensor over A although the proof evaluates R-complexes.
$$
1 \otimes f :
- \otimes_A^\mathbf{L} L^\bullet
\longrightarrow
- \otimes_A^\mathbf{L} N^\bullet
$$
을 유도한다. $f$가 준동형사상이면 $1 \otimes f$는 함자의 동형사상이다.
\end{lemma}

\begin{proof}
% P02-E0900: frozen source says the functors ``are computing''.
함자들은 K-평탄 복합체 $K^\bullet$에서 값을 취하여 계산하므로, 함자성
$$
\text{Tot}(K^\bullet \otimes_R L^\bullet) \to
\text{Tot}(K^\bullet \otimes_R N^\bullet)
$$
을 그대로 사용하여 변환을 정의할 수 있다. 마지막 명제는 보조정리
\ref{more-algebra-lemma-K-flat-quasi-isomorphism}에서 따른다.
\end{proof}

\begin{lemma}
\label{more-algebra-lemma-tensor-hom-adjoint}
$R \to A$를 환 사상이라 하자. 보조정리
\ref{more-algebra-lemma-derived-base-change}의 함자
$D(R) \to D(A)$, $E \mapsto E \otimes_R^\mathbf{L} A$는 제한 함자
$D(A) \to D(R)$의 왼쪽 수반이다.
\end{lemma}

\begin{proof}
Derived Categories, Lemma
\ref{derived-lemma-pre-derived-adjoint-functors-general}과,
$- \otimes_R A$와 제한이 수반이라는 사실에서 따른다. Algebra, Lemma
\ref{algebra-lemma-adjoint-tensor-restrict}를 보라.
\end{proof}

\begin{remark}[주의]
\label{more-algebra-remark-warning-compute-base-change}
$R \to A$를 환 사상이라 하고 $N$과 $N'$을 $A$-가군이라 하자.
$N$과 $N'$을 $R$-가군으로 제한한 것을 $N_R$과 $N'_R$로 나타낸다.
Algebra, Section \ref{algebra-section-base-change}를 보라. 이 상황에서
$D(A)$의 대상 $N_R \otimes_R^\mathbf{L} N'$과
$N \otimes_R^\mathbf{L} N'_R$은 일반적으로 동형이 아니다! 다시 말해,
두 쪽 가운데 어느 쪽으로 $A$-가군 구조를 주는지 세심하게 살펴야 한다.

\medskip\noindent
구체적인 예로 $R = k[x, y]$, $A = R/(xy)$, $N = R/(x)$,
$N' = A = R/(xy)$로 두자. 분해
$0 \to R \xrightarrow{xy} R \to N'_R \to 0$은 $D(A)$에서
$N \otimes_R^\mathbf{L} N'_R = N[1] \oplus N$임을 보인다. 분해
$0 \to R \xrightarrow{x} R \to N_R \to 0$은
$N_R \otimes_R^\mathbf{L} N'$이 복합체 $A \xrightarrow{x} A$로
표현됨을 보인다. 이 두 복합체가 동형이 아님을 보이려면, 두 번째
복합체가 $D(A)$에서 그 코호몰로지 군들의 직합과 동형이 아님을 보이거나,
첫 번째 복합체는 $D(A)$의 완전 대상이 아니지만 두 번째 복합체는
완전 대상임을 보이면 된다. 일부 세부사항은 생략한다.
\end{remark}

\begin{lemma}
\label{more-algebra-lemma-double-base-change}
$A \to B \to C$를 환 사상이라 하고 $N^\bullet$을 $B$-가군의 복합체,
$K^\bullet$을 $C$-가군의 복합체라 하자. 함자들의 합성
% P02-E0901: frozen source uses plural ``compositions'' with singular ``is''.
$$
D(A) \xrightarrow{- \otimes_A^\mathbf{L} N^\bullet}
D(B) \xrightarrow{- \otimes_B^\mathbf{L} K^\bullet} D(C)
$$
은 함자 $- \otimes_A^\mathbf{L}
(N^\bullet \otimes_B^\mathbf{L} K^\bullet) : D(A) \to D(C)$이다.
$M$, $N$, $K$가 각각 $A$, $B$, $C$ 위의 가군이면 $D(C)$에서
$$
(M \otimes_A^\mathbf{L} N) \otimes_B^\mathbf{L} K =
M \otimes_A^\mathbf{L} (N \otimes_B^\mathbf{L} K) =
(M \otimes_A^\mathbf{L} C) \otimes_C^\mathbf{L} (N \otimes_B^\mathbf{L} K)
$$
이다. 또한 부호를 사용하는 표준 동형사상
$$
(M \otimes_A^\mathbf{L} N) \otimes_B^\mathbf{L} K \longrightarrow
(M \otimes_A^\mathbf{L} K) \otimes_C^\mathbf{L} (N \otimes_B^\mathbf{L} C)
$$
을 가진다. 복합체에 대해서도 비슷한 결과가 성립한다.
% P02-E0902: frozen source has subject-verb disagreement in this sentence.
\end{lemma}

\begin{proof}
$B$-가군의 K-평탄 복합체 $P^\bullet$과 준동형사상
$P^\bullet \to N^\bullet$을 택하자(보조정리
\ref{more-algebra-lemma-K-flat-resolution}). $M^\bullet$을 $D(A)$의 임의의 대상을
표현하는 $A$-가군의 K-평탄 복합체라 하자. 그러면 괄호 안의 내용에
보조정리 \ref{more-algebra-lemma-functoriality-derived-base-change}을 적용하여
$$
(M^\bullet \otimes_A^\mathbf{L} P^\bullet) \otimes_B^\mathbf{L} K^\bullet
\longrightarrow
(M^\bullet \otimes_A^\mathbf{L} N^\bullet) \otimes_B^\mathbf{L} K^\bullet
$$
이 동형사상임을 알 수 있다. 보조정리
\ref{more-algebra-lemma-base-change-K-flat}과 \ref{more-algebra-lemma-tensor-product-K-flat}에 의해
복합체
% P02-E0903: frozen display introduces undefined R, wrong bases, and lacks a literal closing parenthesis.
$$
\text{Tot}(M^\bullet \otimes_A P^\bullet) =
\text{Tot}((M^\bullet \otimes_R A) \otimes_A P^\bullet
$$
는 $B$-가군의 복합체로서 K-평탄이고, 구성에 의해 $D(B)$의 유도
텐서곱을 표현한다. 따라서
$(M^\bullet \otimes_A^\mathbf{L} P^\bullet) \otimes_B^\mathbf{L} K^\bullet$은
$C$-가군의 복합체
$$
\text{Tot}(\text{Tot}(M^\bullet \otimes_A P^\bullet)\otimes_B K^\bullet) =
\text{Tot}(M^\bullet \otimes_A \text{Tot}(P^\bullet \otimes_B K^\bullet))
$$
로 표현됨을 알 수 있다.
% P02-E0211: frozen source has the fragment ``Equality by ...''.
등식은 Homology, Remark \ref{homology-remark-triple-complex}에 의한 것이다.
왔던 길을 되짚으면 이는 다음과 같음을 알 수 있다.
$$
M^\bullet \otimes_A^\mathbf{L} (P^\bullet \otimes_B^\mathbf{L} K^\bullet)
\longleftarrow
M^\bullet \otimes_A^\mathbf{L} (N^\bullet \otimes_B^\mathbf{L} K^\bullet)
$$
화살표는 함자 $-\otimes_B^\mathbf{L} K^\bullet$의 정의에 의해
동형사상이다. 이 모든 구성은 복합체 $M^\bullet$에 대하여 함자적이므로
원하는 함자의 동형사상을 얻는다.

\medskip\noindent
위의 결과로 다음 식의 첫 번째 등식을 얻는다.
$$
(M \otimes_A^\mathbf{L} N) \otimes_B^\mathbf{L} K =
M \otimes_A^\mathbf{L} (N \otimes_B^\mathbf{L} K) =
(M \otimes_A^\mathbf{L} C) \otimes_C^\mathbf{L} (N \otimes_B^\mathbf{L} K)
$$
두 번째 등식은 보조정리 \ref{more-algebra-lemma-derived-base-change}의 마지막 명제에서
따른다. 같은 결과를 사용하면
$N \otimes_B^\mathbf{L} K = (N \otimes_B^\mathbf{L} C) \otimes_C^\mathbf{L} K$로
쓸 수 있고, 이를 대입하면
\begin{align*}
(M \otimes_A^\mathbf{L} N) \otimes_B^\mathbf{L} K
& =
(M \otimes_A^\mathbf{L} C) \otimes_C^\mathbf{L}
((N \otimes_B^\mathbf{L} C) \otimes_C^\mathbf{L} K) \\
& =
(M \otimes_A^\mathbf{L} C) \otimes_C^\mathbf{L}
(K \otimes_C^\mathbf{L} (N \otimes_B^\mathbf{L} C)) \\
& =
((M \otimes_A^\mathbf{L} C) \otimes_C^\mathbf{L} K)
\otimes_C^\mathbf{L} (N \otimes_B^\mathbf{L} C)) \\
% P02-E0904: frozen aligned line has one unmatched literal closing parenthesis.
& =
(M \otimes_C^\mathbf{L} K)
% P02-E0905: frozen source tensors M directly over C although M is only an A-module.
\otimes_C^\mathbf{L} (N \otimes_B^\mathbf{L} C)
\end{align*}
이는 보조정리 \ref{more-algebra-lemma-flip-douoble-tensor-product}와
\ref{more-algebra-lemma-triple-tensor-product}, 그리고 앞에서 언급한 보조정리에서
따른다.
\end{proof}






\section{Tor 독립성}
\label{more-algebra-section-tor-independence}

\noindent
환들의 가환 그림
$$
\xymatrix{
A \ar[r] & A' \\
R \ar[r] \ar[u] & R' \ar[u]
}
$$
을 생각하자. $D(A)$의 대상 $K$가 주어지면, 이를 $D(A')$의 대상으로
유도 밑변경한 $K \otimes_A^\mathbf{L} A'$을 생각할 수 있다. 또는
$K$를 $D(R)$의 대상으로 제한한 뒤 이를 $D(R')$의 대상으로 유도
밑변경할 수 있으며, 그 결과를 $K \otimes_R^\mathbf{L} R'$으로
나타낸다. $D(R')$에 함자적 비교 사상
\begin{equation}
\label{more-algebra-equation-comparison-map}
K \otimes_R^{\mathbf{L}} R'
\longrightarrow
K \otimes_A^{\mathbf{L}} A'
\end{equation}
이 존재한다고 주장한다. 이 비교 사상을 구성하기 위해 $K$를 표현하는
$A$-가군의 K-평탄 복합체 $K^\bullet$을 택한다. 다음으로
$E^\bullet$이 $R$-가군의 K-평탄 복합체인 준동형사상
$E^\bullet \to K^\bullet$을 택한다. 위 사상은
$$
K \otimes_R^{\mathbf{L}} R' =
E^\bullet \otimes_R R'
\longrightarrow
K^\bullet \otimes_A A' =
K \otimes_A^{\mathbf{L}} A'
$$
이다. 일반적으로 이 사상이 동형사상일 가능성은 없다.

\medskip\noindent
그러나 위 그림이 환들의 ``밑변경'' 그림인 상황, 즉
$A' = A \otimes_R R'$인 상황을 자주 만난다. 이 경우 임의의
$A$-가군 $M$에 대하여 $M \otimes_A A' = M \otimes_R R'$이다.
따라서 $- \otimes_R R'$은 함자 $\text{Mod}_A \to \text{Mod}_{A'}$로서
$- \otimes_A A'$과 같다. 일반적으로 이 등식은 유도 텐서곱으로
{\bf 연장되지 않는다}. 다시 말해 비교 사상은 동형사상이 아니다.
간단한 예로 $R = k[x]$, $A = R' = A' = k[x]/(x) = k$ 및
$K^\bullet = A[0]$을 취할 수 있다. 명백한 필요조건은 모든
$p > 0$에 대하여 $\text{Tor}_p^R(A, R') = 0$인 것이다.

\begin{definition}
\label{more-algebra-definition-tor-independent}
$R$을 환이라 하고 $A$, $B$를 $R$-대수라 하자. 모든 $p > 0$에 대하여
$\text{Tor}_p^R(A, B) = 0$이면 $A$와 $B$가 {\it $R$ 위에서
Tor-독립}이라고 한다.
\end{definition}

\begin{lemma}
\label{more-algebra-lemma-base-change-comparison}
$A' = A \otimes_R R'$이고 $A$와 $R'$이 $R$ 위에서 Tor-독립이면
비교 사상 (\ref{more-algebra-equation-comparison-map})은 동형사상이다.
\end{lemma}

\begin{proof}
$R'$의 $R$-가군으로서의 자유 분해 $F^\bullet \to R'$을 택한다.
$A$와 $R'$이 $R$ 위에서 Tor-독립이므로
$F^\bullet \otimes_R A$는 $A'$의 $A$ 위 자유 $A$-가군 분해이다.
위의 유도 텐서곱의 일반 구성에 의해
$$
K^\bullet \otimes_A A' \cong
\text{Tot}(K^\bullet \otimes_A (F^\bullet \otimes_R A)) =
\text{Tot}(K^\bullet \otimes_R F^\bullet) \cong
\text{Tot}(E^\bullet \otimes_R F^\bullet) \cong
E^\bullet \otimes_R R'
$$
임을 알 수 있으며, 이는 원하는 결과이다.
\end{proof}

\begin{lemma}
\label{more-algebra-lemma-tor-independent-flat}
환들의 가환 그림
$$
\xymatrix{
A' & R' \ar[r] \ar[l] & B' \\
A \ar[u] & R \ar[l] \ar[u] \ar[r] & B \ar[u]
}
$$
을 생각하자. $R'$이 $R$ 위 평탄이고, $A'$이 $A \otimes_R R'$ 위
평탄이며, $B'$이 $R' \otimes_R B$ 위 평탄이라고 가정하자. 그러면
$$
\text{Tor}_i^R(A, B) \otimes_{(A \otimes_R B)} (A' \otimes_{R'} B') =
\text{Tor}_i^{R'}(A', B')
$$
이다.
\end{lemma}

\begin{proof}
Algebra, Section \ref{algebra-section-functoriality-tor}에 의해 표준 사상
$$
\text{Tor}_i^R(A, B) \longrightarrow
\text{Tor}_i^{R'}(A \otimes_R R', B \otimes_R R') \longrightarrow
\text{Tor}_i^{R'}(A', B')
$$
들이 존재한다. 이들은 보조정리 식의 왼쪽에서 오른쪽으로 가는 사상을
유도한다.

\medskip\noindent
$A$의 $R$-가군으로서의 자유 분해 $F_\bullet \to A$를 택하자.
그러면 $F_\bullet \otimes_R R'$은 $A \otimes_R R'$의 분해이다.
따라서 $\text{Tor}_i^{R'}(A \otimes_R R', B \otimes_R R')$은
$F_\bullet \otimes_R B \otimes_R R'$으로 계산된다. $R'$이 $R$ 위
평탄이라는 가정에 의해 이는 $\text{Tor}_i^R(A, B) \otimes_R R'$을
계산한다. 따라서
$\text{Tor}_i^{R'}(A \otimes_R R', B \otimes_R R') =
\text{Tor}_i^R(A, B) \otimes_R R'$이다(여기에는 $R'$이 $R$ 위
평탄이라는 사실만 사용한다).

\medskip\noindent
Lazard 정리(Algebra, Theorem \ref{algebra-theorem-lazard})에 의해
$A'$ 및 $B'$을 각각 유한 자유 $A \otimes_R R'$-가군 및 유한 자유
$B \otimes_R R'$-가군의 여과 여극한으로 쓸 수 있다. 이를
$A' = \colim M_i$ 및 $B' = \colim N_j$로 쓰자. 위 결과는
$$
\text{Tor}_i^{R'}(M_i, N_j) =
\text{Tor}_i^R(A, B) \otimes_{A \otimes_R B} (M_i \otimes_{R'} N_j)
$$
를 준다. 이는 모든 것을 기저를 사용해 풀어 쓰면 알 수 있다.
여극한을 취하면 보조정리의 결과를 얻는다.
\end{proof}

\begin{lemma}
\label{more-algebra-lemma-flat-base-change-tor-independent}
$R \to A$와 $R \to B$를 환 사상이라 하자. $R \to R'$을 환 사상이라
하고 $A' = A \otimes_R R'$ 및 $B' = B \otimes_R R'$로 두자.
% P02-E0906: frozen source writes lowercase ``tor independent'' at this locus.
$A$와 $B$가 $R$ 위에서 tor-독립이고 $R \to R'$이 평탄이면,
$A'$과 $B'$은 $R'$ 위에서 tor-독립이다.
\end{lemma}

\begin{proof}
보조정리 \ref{more-algebra-lemma-tor-independent-flat}과 정의
\ref{more-algebra-definition-tor-independent}에서 즉시 따른다.
\end{proof}

\begin{lemma}
% P02-E0212: frozen label repeats the lemma namespace word.
\label{more-algebra-lemma-lemma-tor-independent-flat-compare}
보조정리 \ref{more-algebra-lemma-tor-independent-flat}과 같은 가정을 하자.
$M \in D(A)$에 대하여 $A' \otimes_{R'} B'$-가군의 표준 동형사상
$$
H^i((M \otimes_A^\mathbf{L} A') \otimes_{R'}^\mathbf{L} B') =
H^i(M \otimes_R^\mathbf{L} B) \otimes_{(A \otimes_R B)} (A' \otimes_{R'} B')
$$
이 존재한다.
\end{lemma}

\begin{proof}
등식의 양변을 명확히 설명하자. 좌변에는 함자
$D(A) \to D(A') \to D(R') \to D(B')$의 합성에 함자
$H^i : D(B') \to \text{Mod}_{B'}$를 합성한 것이 있다.
$D(B')$의 대상 $(M \otimes_A^\mathbf{L} A') \otimes_{R'}^\mathbf{L} B'$의
자기준동형들로 가는 $A'$의 사상이 있으므로, 좌변은 실제로
$A' \otimes_{R'} B'$-가군이다. 같은 논증으로
$H^i(M \otimes_R^\mathbf{L} B)$가 $A \otimes_R B$-가군 구조를
가짐을 알 수 있다.

\medskip\noindent
먼저 $B' = R' \otimes_R B$인 경우에 결과를 증명한다. 이 경우
자유 $R$-가군에 의한 분해 $F^\bullet \to B$를 택한다. 또한 $M$을
표현하는 $A$-가군의 K-평탄 복합체 $M^\bullet$을 택한다. 그러면
좌변은 다음으로 표현된다.
\begin{align*}
H^i(\text{Tot}((M^\bullet \otimes_A A') \otimes_{R'} (R' \otimes_R F^\bullet)))
& =
H^i(\text{Tot}(M^\bullet \otimes_A A' \otimes_R F^\bullet)) \\
& =
H^i(\text{Tot}(M^\bullet \otimes_R F^\bullet) \otimes_A A') \\
& =
H^i(M \otimes_R^\mathbf{L} B) \otimes_A A'
\end{align*}
% P02-E0907: frozen source has the fragment ``The final equality because ...''.
마지막 등식은 $A \to A'$이 평탄이기 때문이다. 이 문단에서
$B' = R' \otimes_R B$라고 가정했으므로
$A' \otimes_{R'} B' = A' \otimes_R B$이고, 따라서 마지막 가군이
원하는 가군이다.

\medskip\noindent
일반적인 경우를 보자. $B' \to B''$가 평탄 환 사상이라고 하자.
그러면 다음을 쉽게 알 수 있다.
$$
H^i((M \otimes_A^\mathbf{L} A') \otimes_{R'}^\mathbf{L} B'') =
H^i((M \otimes_A^\mathbf{L} A') \otimes_{R'}^\mathbf{L} B')
\otimes_{B'} B''
$$
및
$$
H^i(M \otimes_R^\mathbf{L} B) \otimes_{(A \otimes_R B)} (A' \otimes_{R'} B'')
=
\left(
H^i(M \otimes_R^\mathbf{L} B) \otimes_{(A \otimes_R B)} (A' \otimes_{R'} B')
\right) \otimes_{B'} B''
$$
따라서 $B'$에 대한 결과는 $B''$에 대한 결과를 함의한다. 앞 문단에서
$R' \otimes_R B$에 대한 결과를 증명했으므로 일반적인 결과가 따른다.
\end{proof}

\begin{lemma}
\label{more-algebra-lemma-tor-independent}
$R$을 환이라 하고 $A$, $B$를 $R$-대수라 하자. 다음은 서로 동치이다.
\begin{enumerate}
\item $A$와 $B$는 $R$ 위에서 Tor-독립이다.
\item 같은 소아이디얼 $\mathfrak r \subset R$ 위에 놓이는 모든
소아이디얼 쌍 $\mathfrak p \subset A$ 및 $\mathfrak q \subset B$에 대하여
$A_\mathfrak p$와 $B_\mathfrak q$는 $R_\mathfrak r$ 위에서 Tor-독립이다.
\item $A \otimes_R B$의 모든 소아이디얼 $\mathfrak s$에 대하여 가군
% P02-E0214: frozen condition never quantifies the free index i.
$$
\text{Tor}_i^R(A, B)_\mathfrak s =
\text{Tor}_i^{R_\mathfrak r}(A_\mathfrak p, B_\mathfrak q)_\mathfrak s
$$
(여기서 $\mathfrak p = A \cap \mathfrak s$, $\mathfrak q = B \cap \mathfrak s$,
$\mathfrak r = R \cap \mathfrak s$이다)은 영이다.
\end{enumerate}
\end{lemma}

\begin{proof}
$\mathfrak s$를 (3)에서와 같은 $A \otimes_R B$의 소아이디얼이라 하자.
등식
$$
\text{Tor}_i^R(A, B)_\mathfrak s =
\text{Tor}_i^{R_\mathfrak r}(A_\mathfrak p, B_\mathfrak q)_\mathfrak s
$$
은 $\mathfrak p = A \cap \mathfrak s$, $\mathfrak q = B \cap \mathfrak s$,
$\mathfrak r = R \cap \mathfrak s$일 때 보조정리
\ref{more-algebra-lemma-tor-independent-flat}에서 따른다. 따라서 (2)는 (3)을 함의한다.
가군이 영인지 여부는 소아이디얼에서 국소화하여 판정할 수 있으므로
(Algebra, Lemma \ref{algebra-lemma-characterize-zero-local}), (3)이 (1)을
함의한다고 결론 내린다. (1) $\Rightarrow$ (2)를 위해서는 다시
보조정리 \ref{more-algebra-lemma-tor-independent-flat}에 의한 등식
$$
\text{Tor}_i^{R_\mathfrak r}(A_\mathfrak p, B_\mathfrak q) =
\text{Tor}_i^R(A, B) \otimes_{(A \otimes_R B)}
(A_\mathfrak p \otimes_{R_{\mathfrak r}} B_\mathfrak q)
$$
을 사용한다.
\end{proof}








\section{Tor의 스펙트럼 열}
\label{more-algebra-section-spectral-sequence-tor}


\noindent
이 절에서는 Tor 함자를 생각할 때 나타나는 여러 스펙트럼 열을 모은다.

\begin{example}
\label{more-algebra-example-cohomology-complex-tensored}
$R$을 환이라 하자. $K_n = 0$ ($n \ll 0$)인 $R$-가군의 사슬 복합체
$K_\bullet$을 생각하자. $M$을 $R$-가군이라 하자. 자유 $R$-가군에 의한
$M$의 분해 $P_\bullet \to M$을 택한다. 이로부터 이중 사슬 복합체
$K_\bullet \otimes_R P_\bullet$을 얻는다. Homology, Section
\ref{homology-section-double-complex}의 내용(특히 Homology, Lemma
\ref{homology-lemma-first-quadrant-ss})을 사슬 복합체의 언어로 옮겨
적용하면 $H_*(K_\bullet \otimes_R^\mathbf{L} M)$으로 수렴하는 두
스펙트럼 열을 얻는다. 한편으로 $E_2$-쪽이
$$
(E_2)_{i, j} = \text{Tor}^R_j(H_i(K_\bullet), M)
\Rightarrow
H_{i + j}(K_\bullet \otimes^{\mathbf{L}}_R M)
$$
인 스펙트럼 열이 있고, 그 미분 $d_2$는 사상
$\text{Tor}^R_j(H_i(K_\bullet), M) \to
\text{Tor}^R_{j - 2}(H_{i + 1}(K_\bullet), M)$으로 주어진다.
다른 스펙트럼 열은 $E_1$-쪽
$$
(E_1)_{i, j} = \text{Tor}^R_j(K_i, M)
\Rightarrow
H_{i + j}(K_\bullet \otimes^{\mathbf{L}}_R M)
$$
을 가지며, 그 미분 $d_1$은 $K_i \to K_{i - 1}$이 유도하는 사상
$\text{Tor}^R_j(K_i, M) \to \text{Tor}^R_j(K_{i - 1}, M)$으로 주어진다.
\end{example}

\begin{example}
\label{more-algebra-example-tor-change-rings}
$R \to S$를 환 사상이라 하자. $M$을 $R$-가군, $N$을 $S$-가군이라
하자. 그러면 스펙트럼 열
$$
\text{Tor}^S_n(\text{Tor}^R_m(M, S), N) \Rightarrow
\text{Tor}^R_{n + m}(M, N).
$$
이 존재한다. 이를 구성하기 위해 $M$의 $R$-자유 분해
% P02-E0908: frozen source uses the wrong English indefinite article before R-free.
$P^\bullet$을 택한다. 그러면
$$
M \otimes_R^{\mathbf{L}} N = P^\bullet \otimes_R N =
(P^\bullet \otimes_R S) \otimes_S N
$$
이고, 이제 예 \ref{more-algebra-example-cohomology-complex-tensored}의 첫 번째
스펙트럼 열을 적용한다.
\end{example}

\begin{example}
\label{more-algebra-example-tor-base-change}
가환 그림
$$
\xymatrix{
B \ar[r] & B' = B \otimes_A A' \\
A \ar[r] \ar[u] & A' \ar[u]
}
$$
과 $B$-가군 $M, N$을 생각하자. $M' = M \otimes_A A' = M \otimes_B B'$,
$N' = N \otimes_A A' = N \otimes_B B'$으로 두자.
{\it $A \to B$가 평탄이고 $M$과 $N$이 $A$-평탄이라고 가정하자.}
그러면 스펙트럼 열
$$
\text{Tor}^A_i(\text{Tor}_j^B(M, N), A')
\Rightarrow
\text{Tor}^{B'}_{i + j}(M', N')
$$
이 존재한다. 그 이유는 다음과 같다. $B$-가군으로서의 자유 분해
% P02-E0909: frozen source omits an article before ``free resolution''.
$F_\bullet \to M$을 택한다. $B$와 $M$이 $A$-평탄이므로
$F_\bullet \otimes_A A'$은 $M'$의 자유 $B'$-분해이다. 따라서 군
$\text{Tor}^{B'}_n(M', N')$은 복합체
$$
(F_\bullet \otimes_A A') \otimes_{B'} N' =
(F_\bullet \otimes_B N) \otimes_A A' =
(F_\bullet \otimes_B N) \otimes^{\mathbf{L}}_A A'
$$
로 계산된다. 마지막 등식은 $N$이 $A$ 위 평탄이므로
% P02-E0910: frozen source has the fragment ``the last equality because''.
$F_\bullet \otimes_B N$이 평탄 $A$-가군의 복합체이기 때문이다.
따라서 예 \ref{more-algebra-example-cohomology-complex-tensored}의 스펙트럼 열을
적용하여 위 스펙트럼 열을 얻는다.
\end{example}

\begin{example}
\label{more-algebra-example-tor}
$K^\bullet, L^\bullet$을 $D^{-}(R)$의 대상이라 하자. 그러면
$$
E_2^{p, q} = H^p(K^\bullet \otimes_R^{\mathbf{L}} H^q(L^\bullet))
\Rightarrow H^{p + q}(K^\bullet \otimes_R^{\mathbf{L}} L^\bullet)
$$
을 갖는 스펙트럼 열과
$$
E_2^{p, q} =
H^p(H^q(K^\bullet) \otimes_R^{\mathbf{L}} L^\bullet)
\Rightarrow H^{p + q}(K^\bullet \otimes_R^{\mathbf{L}} L^\bullet)
$$
을 갖는 또 다른 스펙트럼 열이 존재한다. 두 스펙트럼 열 모두
$d_2^{p, q} : E_2^{p, q} \to E_2^{p + 2, q - 1}$를 가진다.
$K^\bullet$과 $L^\bullet$을 사영 가군의 위로 유계인 복합체로 바꾸면,
이 스펙트럼 열들은 Homology, Section
\ref{homology-section-double-complex}에서 논의한
$\text{Tot}(K^\bullet \otimes L^\bullet)$의 코호몰로지를 계산하는 두
스펙트럼 열일 뿐이다.
\end{example}









\section{곱과 Tor}
\label{more-algebra-section-products-tor}

\noindent
곱 사상의 가장 간단한 예는 다음 상황에서 나온다.
$K^\bullet, L^\bullet \in D(R)$이라고 하자. 그러면 사상
\begin{equation}
\label{more-algebra-equation-simple-tor-product}
H^i(K^\bullet) \otimes_R H^j(L^\bullet)
\longrightarrow
H^{i + j}(K^\bullet \otimes_R^{\mathbf{L}} L^\bullet)
\end{equation}
들이 존재한다. 이 사상들을 정의하기 위해 $K^\bullet, L^\bullet$ 가운데
하나가 $R$-가군의 K-평탄 복합체라고 가정해도 된다(예를 들어 자유
또는 사영 $R$-가군의 위로 유계인 복합체). 이 경우 절
\ref{more-algebra-section-derived-tensor-product}(또는 절
\ref{more-algebra-section-computing-tor})에 의해
$K^\bullet \otimes_R^{\mathbf{L}} L^\bullet$은 복합체
$\text{Tot}(K^\bullet \otimes_R L^\bullet)$로 표현된다. 이제
$\xi \in H^i(K^\bullet)$ 및 $\zeta \in H^j(L^\bullet)$라고 하자.
$\xi$와 $\zeta$를 표현하는
$k \in \Ker(K^i \to K^{i + 1})$ 및
$l \in \Ker(L^j \to L^{j + 1})$을 택한다. 다음과 같이 둔다.
$$
\xi \cup \zeta =
\text{class of }k \otimes l\text{ in }
H^{i + j}(\text{Tot}(K^\bullet \otimes_R L^\bullet)).
$$
이는 잘 정의된다. 실제로 전체 복합체의 미분 $\text{d}$에 대한 공식
(Homology, Definition \ref{homology-definition-associated-simple-complex}을
보라)은 $k \otimes l$이 코사이클임을 보인다.
% P02-E0911: frozen source has ``This make sense''.
더구나 어떤 $k'' \in K^{i - 1}$에 대하여 $k' = d_K(k'')$이면,
$l$이 코사이클이므로 $k' \otimes l = \text{d}(k'' \otimes l)$이다.
마찬가지로 $\zeta$를 표현하는 $l$의 선택을 바꾸어도 $k \otimes l$의
류는 변하지 않는다. $\cup$가 쌍선형임도 똑같이 명백하므로,
$H^i(K^\bullet) \otimes_R H^j(L^\bullet)$의 일반 원소에는
$$
\sum \xi_i \otimes \zeta_i \longmapsto \sum \xi_i \cup \zeta_i
$$
를 $H^{i + j}(\text{Tot}(K^\bullet \otimes_R L^\bullet))$의 원소로
대응시킨다.

\medskip\noindent
$R \to A$를 환 사상이라 하고 $K^\bullet, L^\bullet \in D(R)$이라고
하자. 그러면 $D(A)$에서 표준 식별
\begin{equation}
\label{more-algebra-equation-pullback-derived-tensor-product}
(K^\bullet \otimes_R^{\mathbf{L}} A)
\otimes_A^{\mathbf{L}}
(L^\bullet \otimes_R^{\mathbf{L}} A)
=
(K^\bullet \otimes_R^{\mathbf{L}} L^\bullet) \otimes_R^{\mathbf{L}} A
\end{equation}
이 존재한다. 이는 다음과 같이 구성된다. 먼저 $R$ 위의 K-평탄 분해
$P^\bullet \to K^\bullet$과 $Q^\bullet \to L^\bullet$을 택한다.
그러면 좌변은 복합체
$\text{Tot}((P^\bullet \otimes_R A) \otimes_A (Q^\bullet \otimes_R A))$로,
우변은 복합체 $\text{Tot}(P^\bullet \otimes_R Q^\bullet) \otimes_R A$로
표현된다. 이 복합체들은 표준적으로 동형이다. 따라서 위 구성은 때때로
유용한 곱
$$
\text{Tor}^R_n(K^\bullet, A) \otimes_A \text{Tor}^R_m(L^\bullet, A)
\longrightarrow
\text{Tor}_{n + m}^R(K^\bullet \otimes_R^\mathbf{L} L^\bullet, A)
$$
을 유도한다.

\medskip\noindent
$M$, $N$을 $R$-가군이라 하자. 위의 일반 구성, 표준 사상
$M \otimes_R^\mathbf{L} N \to M \otimes_R N$, 그리고 $\text{Tor}$의
함자성을 사용하면 표준 사상
\begin{equation}
\label{more-algebra-equation-tor-product}
\text{Tor}^R_n(M, A) \otimes_A \text{Tor}^R_m(N, A)
\longrightarrow \text{Tor}_{n + m}^R(M \otimes_R N, A)
\end{equation}
을 얻는다. 다음은 사영 분해를 사용한 직접 구성이다. 먼저 $R$ 위의
사영 분해들
$$
P_\bullet \to M, \quad Q_\bullet \to N, \quad T_\bullet \to M \otimes_R N
$$
을 택한다. $\otimes_R$의 오른쪽 완전성에 의해
$H_0(\text{Tot}(P_\bullet \otimes_R Q_\bullet)) = M \otimes_R N$이다.
따라서 Derived Categories, Lemmas
\ref{derived-lemma-morphisms-lift-projective}와
\ref{derived-lemma-morphisms-equal-up-to-homotopy-projective}는
$H_0(\mu) = \text{id}_{M \otimes_R N}$을 만족하는 복합체 사상
% P02-E0219: the second cited lemma gives uniqueness only up to homotopy.
$\mu : \text{Tot}(P_\bullet \otimes_R Q_\bullet) \to T_\bullet$의 존재성과
유일성을 보장한다. 이는 $D(A)$에서 표준 사상
\begin{align*}
(M \otimes_R^{\mathbf{L}} A) \otimes_A^{\mathbf{L}}
(N \otimes_R^{\mathbf{L}} A)
& =
\text{Tot}((P_\bullet \otimes_R A) \otimes_A (Q_\bullet \otimes_R A)) \\
& =
\text{Tot}(P_\bullet \otimes_R Q_\bullet) \otimes_R A \\
& \to
T_\bullet \otimes_R A \\
& = (M \otimes_R N) \otimes_R^{\mathbf{L}} A
\end{align*}
을 유도한다. 따라서 위의 곱 (\ref{more-algebra-equation-tor-product})은 먼저
$A$ 위의 (\ref{more-algebra-equation-simple-tor-product})을 사용하여
$$
\text{Tor}^R_n(M, A) \otimes_A \text{Tor}^R_m(N, A) \to
H^{-n-m}((M \otimes_R^{\mathbf{L}} A) \otimes_A^{\mathbf{L}}
(N \otimes_R^{\mathbf{L}} A))
$$
를 구성한 뒤, 위의 표시된 사상과 합성하여
$\text{Tor}_{n + m}^R(M \otimes_R N, A)$에 이르게 함으로써 구성된다.

\medskip\noindent
위 내용의 흥미로운 특수한 경우는 $B$가 $R$-대수이고
$M = N = B$일 때 생긴다. 이 경우 사상
$$
\text{Tor}_n^R(B, A) \otimes_A \text{Tor}_m^R(B, A)
\longrightarrow
\text{Tor}_{n + m}^R(B \otimes_R B, A)
\longrightarrow
\text{Tor}_{n + m}^R(B, A)
$$
을 얻는다. 두 번째 화살표는 $\text{Tor}$의 함자성을 통해 곱셈 사상
$B \otimes_R B \to B$가 유도한다. 다시 말해
$\text{Tor}^R_{\star}(B, A)$ 위에 $A$-대수 구조를 얻는다. 이 대수
구조는 다른 곳에서 논의할 여러 흥미로운 성질(결합성, 등급 가환성,
% P02-E0912: frozen source mixes a noun and adjective in the property list.
$B$-대수 구조, 어떤 경우의 분할 거듭제곱 등)을 가진다.
% P02-E0913: frozen source uses singular ``case'' after ``some''.
(향후 참조를 여기에 삽입).
% P02-E0914: frozen source contains this unresolved editorial placeholder.

\begin{lemma}
\label{more-algebra-lemma-functoriality-product-tor}
$R$을 환이라 하자. $A, B, C$를 $R$-대수라 하고 $B \to C$를
$R$-대수 사상이라 하자. 그러면 유도된 사상
$$
\text{Tor}^R_{\star}(B, A)
\longrightarrow
\text{Tor}^R_{\star}(C, A)
$$
은 $A$-대수 준동형이다.
\end{lemma}

\begin{proof}
생략한다. 힌트: 모든 복합체를 명시적으로 써서 정의를 따라가면 증명할
수 있다.
\end{proof}








\section{K\"unneth 스펙트럼 열}
\label{more-algebra-section-kunneth-spectral-sequence}

\noindent
$R$을 환이라 하자. $K^\bullet$과 $L^\bullet$을 $R$-가군의 여과 복합체라
하자(Homology, Definition \ref{homology-definition-filtered-complex}을
보라. 우리의 여과는 감소한다는 점에 주의하자). 그러면 복합체
$$
T^\bullet = \text{Tot}(K^\bullet \otimes_R L^\bullet)
$$
도 다음 공식으로 정의되는 감소 여과를 가진다.
$$
F^nT^\bullet =
\Im\left(
\bigoplus\nolimits_{i + j = n}
\text{Tot}(F^iK^\bullet \otimes_R F^jL^\bullet) \to
\text{Tot}(K^\bullet \otimes_R L^\bullet)
\right)
$$
여과 복합체들에 몇 가지 가정을 두고, Homology, Section
\ref{homology-section-filtered-complex}의 구성으로 여기서 생기는
스펙트럼 열을 결정할 것이다.

\medskip\noindent
각각의 $K^n$, $F^iK^n$, $\text{gr}^iK^n$, $L^m$, $F^jL^m$,
$\text{gr}^jL^m$이 평탄 $R$-가군이라고 가정하자. 이 경우 가군
$F^iK^n \otimes_R L^m$, $K^n \otimes_R F^jL^m$,
$F^iK^n \otimes_R F^jL^m$은 $K^n \otimes_R L^m$의 부분가군이다.
마찬가지로 가군 $\text{gr}^iK^n \otimes_R \text{gr}^jL^m$은
$K^n/F^{i + 1}K^n \otimes_R L^m/F^{j + 1}L^m$의 부분가군이다.
사상
% P02-E0915: frozen display has F^nT^n although the source must be a complex.
$$
F^nT^n \longrightarrow
\bigoplus\nolimits_{i + j = n}
\text{Tot}(K^\bullet / F^{i + 1}K^\bullet
\otimes_R L^\bullet / F^{j + 1}L^\bullet)
$$
을 생각하자. $F^nT^\bullet$의 정의 때문에 실제로 직곱이 아니라 직합에
도달한다는 확인은 독자에게 맡긴다. $a + b = n$일 때, 표시된 화살표를
$F^nT^\bullet$의 부분복합체
$\text{Tot}(F^aK^\bullet \otimes_R F^bL^\bullet)$에 제한하면
$i = a$, $j = b$인 직합 성분으로 사상된다. 더구나 우리의 평탄성
가정에 의해 그 상은
$\text{Tot}(\text{gr}^iK^\bullet \otimes_R \text{gr}^jL^\bullet)$과
동형이고, 사상
$\text{Tot}(F^iK^\bullet \otimes_R F^jL^\bullet) \to
\text{Tot}(\text{gr}^iK^\bullet \otimes_R \text{gr}^jL^\bullet)$의 핵은
$\text{Tot}(F^{i + 1}K^\bullet \otimes_R F^jL^\bullet) +
\text{Tot}(F^iK^\bullet \otimes_R F^{j + 1}L^\bullet)$이다.
따라서 우리의 가정 아래 복합체의 짧은 완전열
$$
0 \to F^{n + 1}T^\bullet \to F^nT^\bullet \to
\bigoplus\nolimits_{i + j = n}
\text{Tot}(\text{gr}^iK^\bullet \otimes_R \text{gr}^jL^\bullet) \to 0
$$
을 얻는다. 다시 말해, 이는 $\text{gr}^nT^\bullet$이 $i + j = n$인
복합체 $\text{Tot}(\text{gr}^iK^\bullet \otimes_R
\text{gr}^jL^\bullet)$들의 직합임을 말한다.

\medskip\noindent
이에 더하여 $R$-가군의 복합체 $K^\bullet$, $F^iK^\bullet$,
$\text{gr}^iK^\bullet$, $L^\bullet$, $F^jL^\bullet$,
$\text{gr}^jL^\bullet$이 K-평탄이라고 가정하자. 이 경우 여과 복합체
$T^\bullet$에 연관된 Homology, Section
\ref{homology-section-filtered-complex}의 스펙트럼 열은 항
$$
E_1^{p, q} =
\bigoplus\nolimits_{i + j = p}
H^{p + q}(\text{gr}^iK^\bullet \otimes_R^\mathbf{L} \text{gr}^jL^\bullet)
$$
을 가지며, 그 미분은 Homology, Lemma
\ref{homology-lemma-spectral-sequence-filtered-complex-d1}에서 설명한
짧은 완전열
$$
0 \to \text{gr}^{n + 1}T^\bullet \to
F^nT^\bullet/F^{n + 2}T^\bullet \to
\text{gr}^nT^\bullet \to 0
$$
에서 유도된다. 특히 독자는 $E_1^{p, q}$의 직합 성분
$H^{p + q}(\text{gr}^iK^\bullet \otimes_R^\mathbf{L}
\text{gr}^jL^\bullet)$이 $E_1^{p + 1, q}$ 안의 합
$$
H^{p + q + 1}(\text{gr}^{i + 1}K^\bullet \otimes_R^\mathbf{L}
\text{gr}^jL^\bullet) \oplus
H^{p + q + 1}(\text{gr}^iK^\bullet \otimes_R^\mathbf{L}
\text{gr}^{j + 1}L^\bullet)
$$
으로 사상된다는 뜻임을 보일 수 있다.

\begin{lemma}
\label{more-algebra-lemma-kunneth-converges}
위의 가정 아래 다음 중 하나를 추가로 가정하자.
\begin{enumerate}
\item $K^\bullet$의 여과와 $L^\bullet$의 여과가 유한이다.
\item 더 일반적으로 다음이 성립한다.
\begin{enumerate}
\item $i \gg 0$에 대하여 $F^iK^\bullet$은 비순환이다.
\item $i \ll 0$에 대하여 $F^iK^\bullet \to K^\bullet$은 준동형사상이다.
\item $j \gg 0$에 대하여 $F^jL^\bullet$은 비순환이다.
\item $j \ll 0$에 대하여 $F^jL^\bullet \to L^\bullet$은 준동형사상이다.
\end{enumerate}
\end{enumerate}
그러면 스펙트럼 열은 유계이고, 각
$H^n(T^\bullet) = H^n(K^\bullet \otimes_R^\mathbf{L} L^\bullet)$ 위의
연관된 여과는 유한이며, 다음 수렴이 성립한다.
% P02-E0916: frozen convergence target uses free n instead of p+q.
$$
E_1^{p, q} =
\bigoplus\nolimits_{i + j = p}
H^{p + q}(\text{gr}^iK^\bullet \otimes_R^\mathbf{L} \text{gr}^jL^\bullet)
\Rightarrow
H^n(K^\bullet \otimes_R^\mathbf{L} L^\bullet)
$$
\end{lemma}

\begin{proof}
(1)의 경우 $T^\bullet$의 여과가 유한이므로 Homology, Lemma
\ref{homology-lemma-biregular-ss-converges}에서 보조정리가 즉시 따른다.
(2)의 경우 $i, j > b$일 때 $F^iK^\bullet$과 $F^jL^\bullet$이
비순환이고, $i, j < a$일 때 $F^iK^\bullet \to K^\bullet$과
$F^jL^\bullet \to L^\bullet$이 준동형사상이 되도록 $a < b$를 택하자.
이 경우 $n > 2b$이면 복합체 $F^nT^\bullet$이 비순환이고,
$n < 2a - 1$이면 $F^nT^\bullet \to T^\bullet$이 준동형사상이라고
주장한다. 이 주장에 의해 Homology, Lemma
\ref{homology-lemma-ss-converges-trivial}을 적용할 수 있으므로 보조정리가
참임을 알 수 있다.

\medskip\noindent
주장의 증명. $n$과 정수 $i_1 \leq i_2$가 주어졌을 때
$$
S^{n, i_1, i_2} =
\Im\left(
\bigoplus\nolimits_{i_1 \leq i \leq i_2}
\text{Tot}(F^iK^\bullet \otimes_R F^{n - i}L^\bullet)
\to
\text{Tot}(K^\bullet \otimes_R L^\bullet)
\right)
$$
을 $\text{Tot}(K^\bullet \otimes_R L^\bullet)$의 부분복합체로 생각하자.
$F^nT^\bullet = \colim S^{n, i_1, i_2}$가 여과 여극한임에 주의하자.
따라서 주장을 보이려면 $n > 2b$일 때 $S^{n, i_1, i_2}$가 비순환이고,
$n < 2a - 1$일 때 쌍 $i_1, i_2$의 공종 집합에 대하여 사상
% P02-E0225: frozen comparison target omits the base ring R.
$S^{n, i_1, i_2} \to \text{Tot}(K^\bullet \otimes L^\bullet)$이
준동형사상임을 보이면 충분하다.

\medskip\noindent
$i_1 = i_2 = i$이고 $n > 2b$이면
$$
S^{n, i, i} =
\text{Tot}(F^iK^\bullet \otimes_R F^{n - i}L^\bullet)
$$
은 비순환이다. 실제로 $i > b$이면 $F^{n - i}L^\bullet$이 K-평탄이라는
가정에 의해 이 복합체가 비순환이고, $i \leq b$이면
$n - i \geq n - b > b$이므로 같은 결론이 성립한다. 우리의 평탄성
가정을 사용하면 복합체의 짧은 완전열
$$
0 \to
S^{n + 1, i_2 + 1, i_2 + 1}
\to
S^{n, i_1, i_2} \oplus S^{n, i_2 + 1, i_2 + 1}
\to
S^{n, i_1, i_2 + 1}
\to
0
$$
이 존재함을 독자가 보일 수 있다. 따라서 $i_2 - i_1$에 관한 귀납법으로
$n > 2b$일 때 모든 복합체 $S^{n, i_1, i_2}$가 비순환임을 알 수 있다.

\medskip\noindent
$n < 2a - 1$이라고 가정하자. 사상
$$
S^{n, a - 1, a - 1} =
\text{Tot}(F^{a - 1}K^\bullet \otimes_R F^{n - a + 1}L^\bullet)
\to \text{Tot}(K^\bullet \otimes_R L^\bullet)
$$
은 준동형사상임에 주의하자. 실제로 $a - 1 < a$와
$n - a + 1 < a$가 모두 성립하고, 우리의 모든 복합체가 K-평탄이므로
양변은 $D(R)$의 같은 대상을 계산한다. $i_2 \geq a - 1$ 및
$n < 2a - 1$일 때 사상
$$
S^{n + 1, i_2 + 1, i_2 + 1} \to S^{n, i_2 + 1, i_2 + 1}
$$
은 준동형사상이다. 실제로 좌변과 우변은
% P02-E0226: frozen prose uses singular ``side'' for both sides.
$\text{Tot}(F^{i_2 + 1}K^\bullet \otimes L^\bullet)$과 준동형이다.
% P02-E0227: frozen total complex omits the base ring R.
따라서 앞 문단의 복합체 짧은 완전열을 사용하면 모든
$t \geq a - 1$에 대하여
$S^{n, a - 1, a - 1} \to S^{n, a - 1, t}$가 준동형사상임을 알 수 있다.
한편 비슷하게 복합체의 짧은 완전열
$$
0 \to
S^{n + 1, i_1, i_1}
\to
S^{n, i_1 - 1, i_1 - 1} \oplus
S^{n, i_1, i_2}
\to
S^{n, i_1 - 1, i_2}
\to
0
$$
들이 존재한다. 위와 마찬가지로 $i_1 \leq a - 1$이면 사상
$$
S^{n + 1, i_1, i_1} \to S^{n, i_1 - 1, i_1 - 1}
$$
은 준동형사상이다. 실제로 좌변과 우변은
% P02-E0228: frozen total complex omits the base ring R.
$\text{Tot}(K^\bullet \otimes F^{n - i_1 + 1}L^\bullet)$과 준동형이다.
따라서 모든 $s \leq a - 1$에 대하여
% P02-E0229: frozen arrow points opposite to the natural inclusion.
$S^{n, s, t} \to S^{n, a - 1, t}$가 준동형사상이라고 결론 내린다.
이 둘을 합치면 $s \leq a - 1$ 및 $t \geq a - 1$인
$S^{n, s, t}$가 $\text{Tot}(K^\bullet \otimes_R L^\bullet)$로
준동형사상으로 사상됨을 알 수 있다. 이로써 보조정리의 증명이 끝난다.
\end{proof}

\begin{lemma}
\label{more-algebra-lemma-resolution-filtered-complex}
$R$을 환이라 하고 $K^\bullet$을 여과 복합체라 하자. 여과 복합체의 사상
$f : P^\bullet \to K^\bullet$이 존재하여 다음을 만족한다.
\begin{enumerate}
\item 각각의 $P^n$, $F^iP^n$, $\text{gr}^iP^n$은 자유 $R$-가군이다.
\item $R$-가군의 복합체 $P^\bullet$, $F^iP^\bullet$,
$\text{gr}^iP^\bullet$은 K-평탄이다.
\item $f$는 준동형사상 $P^\bullet \to K^\bullet$,
$F^iP^\bullet \to F^iK^\bullet$, $\text{gr}^iP^\bullet \to
\text{gr}^iK^\bullet$을 유도한다.
\end{enumerate}
\end{lemma}

\begin{proof}
각각의 $L^n$, $F^iL^n$, $\text{gr}^iL^n$이 자유 $R$-가군이고 모든
미분이 영일 때 여과 복합체 $L^\bullet$을 {\it 기본}이라고 하자.

\medskip\noindent
기본 여과 복합체 $P_0^\bullet$과 여과 복합체의 사상
$f_0 : P_0^\bullet \to K^\bullet$이 존재하여 $f_0$와
$F^if_0$, $i \in \mathbf{Z}$가 코호몰로지에서 전사이다. 이를 보려면
$$
P_0^n = \bigoplus\nolimits_{z \in \Ker(d_K^n)} R \xi_z
\oplus
\bigoplus\nolimits_{j \in \mathbf{Z}}
\bigoplus\nolimits_{z \in \Ker(d_{F^jK}^n)} R \xi_z
$$
로 두고 미분을 영으로 하며, $f_0(\xi_z) = z$로 사상 $f_0$를 정의한다.
여과는
% P02-E0917: frozen direct-sum subscript has malformed chained condition j >= i in Z.
$$
F^iP_0^n = \bigoplus\nolimits_{j \geq i \in \mathbf{Z}}
\bigoplus\nolimits_{z \in \Ker(d_{F^jK}^n)} R \xi_z
$$
로 둔다. 이것이 주장한 사상 $f_0 : P_0^\bullet \to K^\bullet$을
준다는 확인은 독자에게 맡긴다.

\medskip\noindent
$m \geq 0$에 관한 귀납법으로 여과 복합체의 매입
% P02-E0918: frozen display omits a relation sign before P_m.
$$
P_0^\bullet \subset \ldots P_m^\bullet \subset P_{m + 1}^\bullet
$$
과 사상 $f_m : P_m^\bullet \to K^\bullet$을 구성하여 다음 성질들을
만족시킬 것이다.
\begin{enumerate}
\item 여과 복합체 $P_{m + 1}^\bullet / P_m^\bullet$은 기본이다.
\item $H^n(f_m)$과 $H^n(F^if_m)$의 핵은
$H^n(P_{m + 1}^\bullet)$과 $H^n(F^iP_{m + 1}^\bullet)$에서 영으로
사상된다.
\item 사상 $f_{m + 1} : P_{m + 1}^\bullet \to K^\bullet$은 사상
$f_m$을 연장한다.
\end{enumerate}
% P02-E0919: frozen source begins ``To to this''.
이를 위해
$$
\Omega_{n, m} = \Ker\left(\Ker(d^n_{P_m}) \to H^n(K^\bullet)\right),
\text{ resp. }
\Omega_{n, i, m} = \Ker\left(\Ker(d^n_{F^iP_m}) \to H^n(F^iK^\bullet)\right),
$$
로 두자. $\Omega_{n, m}$은 $H^n(f_m)$의 핵으로 전사하고,
$\Omega_{n, i, m}$은 $H^n(F^if_n)$의 핵으로 전사함에 주의하자.
% P02-E0920: frozen second kernel statement uses f_n instead of stage f_m.
각 $z \in \Omega_{n, m}$ 및 각각의 $z \in \Omega_{n, i, m}$에 대하여
$d_K(y_z) = f_m(z)$를 만족하는 $y_z \in K^{n - 1}$ 및 각각
$y_z \in F^iK^{n - 1}$을 택한다. 이제
$$
P^n_{m + 1} = P_m^n \oplus
\bigoplus\nolimits_{z \in \Omega_{n + 1, m}} R \eta_z
\oplus
\bigoplus\nolimits_{j \in \mathbf{Z}}
\bigoplus\nolimits_{z \in \Omega_{n + 1, j, m}} R \eta_z
$$
로 둔다. $d_{P_m}(\eta_z) = z$ 및 $f_{m + 1}(\eta_z) = y_z$로 둔다.
% P02-E0921: frozen differential is indexed P_m although eta_z is adjoined to P_{m+1}.
끝으로
$$
F^iP_{m + 1}^n = F^iP_m^n \oplus
\bigoplus\nolimits_{j \geq i}
\bigoplus\nolimits_{z \in \Omega_{n + 1, j, m}} R \eta_z
$$
로 둔다. 이것이 주장한 $P_m^\bullet \subset P_{m + 1}^\bullet$과
$f_{m + 1} : P_{m + 1}^\bullet \to K^\bullet$을 준다는 확인은 독자에게
맡긴다.

\medskip\noindent
이제 단순히
$$
P^\bullet = \bigcup P_m^\bullet
$$
을 여과 복합체로 취하고, $\bigcup f_m$으로 주어지는 사상
$f : P^\bullet \to K^\bullet$을 취한다.

\medskip\noindent
보조정리 명제의 (1)은 여과 가군으로서 $P^n$이 $P_0^n$과
$m \geq 0$에 대한 $P_{m + 1}^n/P_m^n$들의 직합과 동형이므로 성립한다.
작은 세부사항은 생략한다.

\medskip\noindent
명제의 (3)을 보자. $H^n(P^\bullet) = \colim H^n(P_m^\bullet)$임에
주의하자. 사상 $f_0$는 코호몰로지에서 전사이므로 각 $f_m$도 그러하다.
각 $m$에 대하여 구성상 매입 $P_m^\bullet \subset P_{m + 1}^\bullet$은
$H^n(f_m)$의 핵을 영으로 만든다. 이 사실들을 합치면 독자는
$H^n(f)$가 동형사상임을 쉽게 알 수 있다. $H^n(F^if_n)$도 마찬가지이다.
% P02-E0922: frozen filtered statement retains an unrelated stage subscript f_n.
그러면 짧은 완전열 $0 \to F^{i + 1} \to F^i \to \text{gr}^i \to 0$
(이를 여과 복합체 범주 위의 함자열로 볼 수 있다)에 의해
$H^n(\text{gr}^if)$도 동형사상이어야 한다. 작은 세부사항은 생략한다.

\medskip\noindent
명제의 (2)를 보자. $P^\bullet$이 K-평탄임을 보이기 위해, 보조정리
\ref{more-algebra-lemma-colimit-K-flat}에 따라 $P_m^\bullet$이 K-평탄임을 보이면
충분하다. 보조정리 \ref{more-algebra-lemma-K-flat-two-out-of-three-ses}와 귀납법에
의해, 미분이 영이고 항들이 자유인 복합체가 K-평탄임에 주의하면
충분하다. 같은 논증으로 모든 $i \in \mathbf{Z}$에 대하여
$F^iP^\bullet$이 K-평탄임을 알 수 있다. 끝으로 보조정리
\ref{more-algebra-lemma-K-flat-two-out-of-three-ses}을 한 번 더 적용하면
$\text{gr}^iP^\bullet$이 K-평탄임을 알 수 있다.
\end{proof}

\begin{proposition}
\label{more-algebra-proposition-kunneth-filtered}
$R$을 환이라 하고 $K^\bullet$과 $L^\bullet$을 $R$-가군의 여과 복합체라
하자. $D(R)$에서 $K^\bullet \otimes_R^\mathbf{L} L^\bullet$을 표현하고,
연관된 스펙트럼 열의 $E_1$-쪽이
$$
E_1^{p, q} =
\bigoplus\nolimits_{i + j = p}
H^{p + q}(\text{gr}^iK^\bullet \otimes_R^\mathbf{L} \text{gr}^jL^\bullet)
$$
인 여과 복합체 $T^\bullet$이 존재한다. 다음을 가정하자.
\begin{enumerate}
\item $i \gg 0$에 대하여 $F^iK^\bullet$은 비순환이다.
\item $i \ll 0$에 대하여 $F^iK^\bullet \to K^\bullet$은 준동형사상이다.
\item $j \gg 0$에 대하여 $F^jL^\bullet$은 비순환이다.
\item $j \ll 0$에 대하여 $F^jL^\bullet \to L^\bullet$은 준동형사상이다.
\end{enumerate}
그러면 스펙트럼 열은 유계이고, 각각의
$H^n(K^\bullet \otimes_R^\mathbf{L} L^\bullet)$ 위의 연관된 여과는
유한이며, 스펙트럼 열은 수렴한다.
% P02-E0923: frozen source says ``the spectral sequence convergences''.
\end{proposition}

\begin{proof}
보조정리 \ref{more-algebra-lemma-resolution-filtered-complex}에서와 같은
$P^\bullet \to K^\bullet$과 $Q^\bullet \to L^\bullet$을 택한다.
그런 다음 이 절의 본문과 보조정리 \ref{more-algebra-lemma-kunneth-converges}에서
설명한 여과 복합체
$$
T^\bullet = \text{Tot}(P^\bullet \otimes_R Q^\bullet)
$$
의 스펙트럼 열을 사용한다.
\end{proof}

\begin{lemma}[K\"unneth 스펙트럼 열]
\label{more-algebra-lemma-kunneth-spectral-sequence}
$R$을 환이라 하고 $K$와 $L$을 $D^b(R)$의 대상이라 하자. 다음을 갖는
이중등급 유계 스펙트럼 열 $\{E_r\}_{r \geq 2}$가 존재한다.
$$
E_2^{p, q} = \bigoplus\nolimits_{i + j = q} \text{Tor}^R_{-p}(H^i(K), H^j(L))
$$
$d_r$은 이중차수 $(r, -r + 1)$이고, 스펙트럼 열은
$H^{p + q}(K \otimes_R^\mathbf{L} L)$로 수렴한다.
\end{lemma}

\begin{proof}
$K^\bullet$을 $K$를 표현하는 $R$-가군의 복합체라 하고,
$L^\bullet$을 $L$을 표현하는 $R$-가군의 복합체라 하자.
$F^iK^\bullet = \tau_{\leq -i}K^\bullet$로 두고 $L$에 대해서도
마찬가지로 둔다. Homology, Section
\ref{homology-section-truncations}를 보라. 명제
\ref{more-algebra-proposition-kunneth-filtered}을 적용하고
$\text{gr}^iK^\bullet = H^{-i}(K)[i]$임에 주의하면,
$$
(E')^{p, q}_1 =
\bigoplus\nolimits_{i + j = p} \text{Tor}^R_{-2p - q}(H^{-i}(K), H^{-j}(L))
$$
을 가지고 $K \otimes_R^\mathbf{L} L$의 코호몰로지로 수렴하는
유계 스펙트럼 열 $\{(E')_r\}_{r \geq 1}$을 얻는다. 구성에 의해 이
스펙트럼 열은 $r \geq 1$에 대하여 미분
$d_r^{p, q} : (E')_r^{p, q} \to (E')_r^{p + r, q - r + 1}$을 가진다.
보조정리의 스펙트럼 열을 얻기 위해 $r \geq 2$에 대하여
$$
E_r^{p, q} = (E')_{r - 1}^{-q, p + 2q}
$$
로 둔다. 이것이 성립한다는 확인은 독자에게 맡긴다.
\end{proof}

\begin{example}
\label{more-algebra-example-kunneth-dedekind}
$R = \mathbf{Z}$이거나, 더 일반적으로 $R$이 Dedekind 정역이면 모든
$R$-가군 $M$과 $N$에 대하여 $i \not \in \{0, 1\}$일 때
$\text{Tor}_i^R(M, N) = 0$이다. 따라서 보조정리
\ref{more-algebra-lemma-kunneth-spectral-sequence}의 스펙트럼 열은 $E_2$-쪽에서
퇴화하며, 모든 $n \in \mathbf{Z}$에 대하여 짧은 완전열
$$
0 \to \bigoplus_{i + j = n} H^i(K) \otimes_R H^j(L) \to
H^n(K \otimes_R^\mathbf{L} L) \to
\bigoplus_{i + j = n + 1} \text{Tor}^R_1(H^i(K), H^j(L)) \to 0
$$
을 얻는다.
\end{example}











































\section{유사연접 가군, I}
\label{more-algebra-section-pseudo-coherent}

\noindent
$R$이 환이라고 하자. $R$-가군 $M$에 전사 $R^{\oplus a} \to M$이
존재하면 유한형이고, 표시
$R^{\oplus a_1} \to R^{\oplus a_0} \to M \to 0$이 존재하면
유한 표시임을 상기하자. 마찬가지로 유한 자유 $R$-가군들에 의한 길이
$n$의 분해
\begin{equation}
\label{more-algebra-equation-pseudo-coherent}
R^{\oplus a_n} \to R^{\oplus a_{n - 1}} \to \ldots \to R^{\oplus a_0} \to
M \to 0
\end{equation}
가 존재하는 $R$-가군들을 생각할 수 있다. 모든 $n$에 대하여 이러한
분해를 찾을 수 있는 가군을 {\it 유사연접}이라고 한다. 형식적 정의는
다음과 같다.

\begin{definition}
\label{more-algebra-definition-pseudo-coherent}
$R$을 환이라 하자.
% P02-E0924: frozen source omits the preposition in “Denote D(R) its derived category”; Korean renders the intended notation.
그 유도 범주를 $D(R)$로 나타내자. $m \in \mathbf{Z}$라 하자.
\begin{enumerate}
\item $D(R)$의 대상 $K^\bullet$에 대하여, 유한 자유 $R$-가군들의
유계 복합체 $E^\bullet$과 사상 $\alpha : E^\bullet \to K^\bullet$이
존재하여 $i > m$이면 $H^i(\alpha)$가 동형사상이고 $H^m(\alpha)$가
전사이면 이를 {\it $m$-유사연접}이라고 한다.
\item $D(R)$의 대상 $K^\bullet$이 유한 자유 $R$-가군들의 위로 유계인
복합체와 준동형적이면 이를 {\it 유사연접}이라고 한다.
\item $R$-가군 $M$에 대하여 $M[0]$이 $D(R)$의 $m$-유사연접 대상이면
이를 {\it $m$-유사연접}이라고 한다.
\item $R$-가군 $M$에 대하여
{\it 유사연접}\footnote{이는 \cite{Bourbaki-CA}에서 유사연접 가군이
뜻하는 바와 충돌한다.}이라는 말은 $M[0]$이 $D(R)$의 유사연접
대상이라는 뜻이다.
\end{enumerate}
\end{definition}

\noindent
늘 그렇듯 이 용어를 $R$-가군의 복합체에도 적용한다. $D(R)$의 임의의
사상 $E^\bullet \to K^\bullet$은 실제 복합체 사상으로 표현되므로
(Derived Categories, Lemma
\ref{derived-lemma-morphisms-from-projective-complex} 참조) 모호함은 없다.
$K^\bullet$이 유사연접일 필요충분조건은 모든 $m \in \mathbf{Z}$에
대하여 $K^\bullet$이 $m$-유사연접인 것임이 밝혀진다. 보조정리
\ref{more-algebra-lemma-pseudo-coherent}를 보라. 또한 환이 Noether이면 이 조건은
코호몰로지의 유한 생성 조건으로 이해할 수 있다. 보조정리
\ref{more-algebra-lemma-Noetherian-pseudo-coherent}를 보라. 먼저 이를 위의 비형식적
논의와 연결하자.

\begin{lemma}
\label{more-algebra-lemma-cone-pseudo-coherent}
$R$을 환이라 하고 $m \in \mathbf{Z}$라 하자.
$(K^\bullet, L^\bullet, M^\bullet, f, g, h)$를 $D(R)$의 구별 삼각형이라
하자.
\begin{enumerate}
\item $K^\bullet$이 $(m + 1)$-유사연접이고 $L^\bullet$이
$m$-유사연접이면 $M^\bullet$은 $m$-유사연접이다.
\item $K^\bullet, M^\bullet$이 $m$-유사연접이면 $L^\bullet$도
$m$-유사연접이다.
\item $L^\bullet$이 $(m + 1)$-유사연접이고 $M^\bullet$이
$m$-유사연접이면 $K^\bullet$은 $(m + 1)$-유사연접이다.
\end{enumerate}
\end{lemma}

\begin{proof}
(1)의 증명. $P^\bullet$을 유한 자유 가군들의 유계 복합체라 하고,
$i > m + 1$이면 $H^i(\alpha)$가 동형사상이며 $i = m + 1$이면
전사인 $\alpha : P^\bullet \to K^\bullet$을 택한다.
$P^\bullet$을 $\sigma_{\geq m + 1}P^\bullet$로 바꿀 수 있으므로
$i < m + 1$이면 $P^i = 0$이라고 가정할 수 있다. $E^\bullet$을
유한 자유 가군들의 유계 복합체라 하고, $i > m$이면 $H^i(\beta)$가
동형사상이며 $i = m$이면 전사인
$\beta : E^\bullet \to L^\bullet$을 택한다. Derived Categories,
Lemma \ref{derived-lemma-lift-map}에 의해 그림
$$
\xymatrix{
K^\bullet \ar[r] & L^\bullet \\
P^\bullet \ar[u] \ar[r]^\gamma & E^\bullet \ar[u]_\beta
}
$$
이 $D(R)$에서 가환하도록 하는 사상
$\gamma : P^\bullet \to E^\bullet$을 찾을 수 있다.
원뿔 $C(\gamma)^\bullet$은 유한 자유 $R$-가군들의 유계 복합체이고,
그림의 가환성은 구별 삼각형의 사상
$$
(P^\bullet, E^\bullet, C(\gamma)^\bullet)
\longrightarrow
(K^\bullet, L^\bullet, M^\bullet).
$$
이 존재함을 뜻한다. 유도된 긴 완전 코호몰로지 열의 사상과 Homology,
Lemmas \ref{homology-lemma-four-lemma} 및
\ref{homology-lemma-five-lemma}로부터
$C(\gamma)^\bullet \to M^\bullet$이 차수 $> m$의 코호몰로지에서
동형사상을, 차수 $m$에서 전사를 유도함을 알 수 있다. 따라서
$M^\bullet$은 $m$-유사연접이다.

\medskip\noindent
(2)와 (3)은 구별 삼각형을 회전하여 (1)에서 따른다.
\end{proof}

\begin{lemma}
\label{more-algebra-lemma-finite-cohomology}
$R$을 환이라 하고 $K^\bullet$을 $R$-가군의 복합체라 하자.
$m \in \mathbf{Z}$라 하자.
\begin{enumerate}
\item $K^\bullet$이 $m$-유사연접이고 $i > m$이면
$H^i(K^\bullet) = 0$이라면 $H^m(K^\bullet)$은 유한형 $R$-가군이다.
\item $K^\bullet$이 $m$-유사연접이고 $i > m + 1$이면
$H^i(K^\bullet) = 0$이라면 $H^{m + 1}(K^\bullet)$은 유한 표시
$R$-가군이다.
\end{enumerate}
\end{lemma}

\begin{proof}
(1)의 증명. $E^\bullet$을 유한 사영 $R$-가군들의 유계 복합체라 하고,
차수 $> m$의 코호몰로지에서 동형사상을, 차수 $m$에서 전사를 유도하는
사상 $\alpha : E^\bullet \to K^\bullet$을 택한다.
$E^\bullet$에 대하여 결과를 증명하면 충분함은 명백하다.
$E^n \not = 0$인 가장 큰 정수를 $n$이라 하자. $n = m$이면 결과는
명백하다. $n > m$이면 $H^n(E^\bullet) = 0$이므로
$E^{n - 1} \to E^n$은 전사이다. $E^n$은 유한 사영이므로
$E^{n - 1} = E' \oplus E^n$이다. 따라서 차수 $n - 1$의 항이 $E'$이고
차수 $n$의 항이 $0$이라는 점을 제외하면 $E^\bullet$과 같은 복합체
% P02-E0925: frozen source says “except has”, omitting the subject; Korean renders the intended comparison.
$(E')^\bullet$에 대하여 결과를 증명하면 충분하다. $n$에 대한 귀납으로
결론을 얻는다.

\medskip\noindent
(2)의 증명. $E^\bullet$을 유한 사영 $R$-가군들의 유계 복합체라 하고,
차수 $> m$의 코호몰로지에서 동형사상을, 차수 $m$에서 전사를 유도하는
사상 $\alpha : E^\bullet \to K^\bullet$을 택한다. (1)의 증명에서처럼
$i > m + 1$이면 $E^i = 0$인 경우로 환원할 수 있다. 그러면
$$
H^{m + 1}(K^\bullet) \cong
H^{m + 1}(E^\bullet) = \Coker(E^m \to E^{m + 1})
$$
이고, 이는 유한 표시이다.
\end{proof}

\begin{lemma}
\label{more-algebra-lemma-n-pseudo-module}
$R$을 환이라 하고 $M$을 $R$-가군이라 하자. 그러면
\begin{enumerate}
\item $M$이 $0$-유사연접일 필요충분조건은 $M$이 유한
$R$-가군인 것이고,
\item $M$이 $(-1)$-유사연접일 필요충분조건은 $M$이 유한 표시
$R$-가군인 것이며,
\item $M$이 $(-d)$-유사연접일 필요충분조건은 길이 $d$인 분해
$$
R^{\oplus a_d} \to R^{\oplus a_{d - 1}} \to \ldots \to R^{\oplus a_0} \to
M \to 0
$$
가 존재하는 것이고,
\item $M$이 유사연접일 필요충분조건은 유한 자유 $R$-가군들에 의한
무한 분해
$$
\ldots \to R^{\oplus a_1} \to R^{\oplus a_0} \to M \to 0
$$
가 존재하는 것이다.
\end{enumerate}
\end{lemma}

\begin{proof}
$M$이 유한형(각각 유한 표시)이면 정의
\ref{more-algebra-definition-pseudo-coherent} 앞의 논의에 의해 $M$은
$0$-유사연접(각각 $(-1)$-유사연접)이다. 반대로 $M$이
$0$-유사연접이면 보조정리 \ref{more-algebra-lemma-finite-cohomology}에 의해
$M = H^0(M[0])$은 유한형이다. $M$이 $(-1)$-유사연접이면
$0$-유사연접이므로 유한형이다. 전사 $R^{\oplus a} \to M$을 택하고
$K = \Ker(R^{\oplus a} \to M)$라 쓰자. 보조정리
\ref{more-algebra-lemma-cone-pseudo-coherent}에 의해 $K$가 $0$-유사연접이고,
따라서 유한형임을 알 수 있으므로 $M$은 유한 표시이다.

\medskip\noindent
% P02-E0926: frozen source uses singular “statement” for items (3) and (4); Korean renders the intended plural.
세 번째와 네 번째 명제를 증명하려면 귀납법과 위와 비슷한 논증을
사용하면 된다(세부사항은 생략한다).
\end{proof}
\begin{lemma}
\label{more-algebra-lemma-pseudo-coherent}
$R$을 환이라 하고 $K^\bullet$을 $R$-가군의 복합체라 하자.
다음은 동치이다.
\begin{enumerate}
\item $K^\bullet$은 유사연접이다.
\item 모든 $m \in \mathbf{Z}$에 대하여 $K^\bullet$은 $m$-유사연접이다.
\item $K^\bullet$은 유한 사영 $R$-가군들의 위로 유계인 복합체와
준동형적이다.
\end{enumerate}
(1), (2), (3)이 성립하고 $i > b$이면 $H^i(K^\bullet) = 0$이라면,
$F^i$가 유한 자유 $R$-가군이고 $i > b$이면 $F^i = 0$인 준동형사상
$F^\bullet \to K^\bullet$을 찾을 수 있다.
\end{lemma}

\begin{proof}
유한 자유 가군은 유한 사영 $R$-가군이므로
(1) $\Rightarrow$ (3)임을 알 수 있다. 반대로 $P^\bullet$을 유한 사영
$R$-가군들의 위로 유계인 복합체라 하자. $i > n_0$이면
$P^i = 0$이라고 하자.
% P02-E0927: frozen source says “a direct sum decompositions”; Korean renders the intended singular construction.
$F^{n_0}$가 유한 자유 $R$-가군인 직합 분해
$F^{n_0} = P^{n_0} \oplus C^{n_0}$을 택하고, $n \leq n_0$에 대하여
귀납적으로
$$
F^{n - 1} = P^{n - 1} \oplus C^n \oplus C^{n - 1}
$$
을 택하되,
% P02-E0928: frozen source repeats F^{n_0} instead of the newly constructed F^{n-1}; Korean preserves the formal span.
$F^{n_0}$가 유한 자유 $R$-가군이도록 한다. 복합체로서
$F^\bullet$의 사상 $F^{n - 1} \to F^n$은
$P^{n - 1} \to P^n$과 일치하고, $C^n \to C^n$ 위에서는 항등사상을
유도하며, $C^{n - 1}$ 위에서는 영이다.
사상 $F^\bullet \to P^\bullet$은 준동형사상(실제로는 호모토피
동치)이므로 (3)은 (1)을 함의한다.

\medskip\noindent
(1)을 가정하자. $E^\bullet$을 유한 자유 $R$-가군들의 위로 유계인
복합체라 하고 $E^\bullet \to K^\bullet$을 준동형사상이라 하자.
그러면 $E^\bullet$의 ‘stupid’ 절단에서 $K^\bullet$로 가는 유도 사상
$\sigma_{\geq m}E^\bullet \to K^\bullet$은 $K^\bullet$이
$m$-유사연접임을 보인다. 따라서 (1)은 (2)를 함의한다.

\medskip\noindent
(2)를 가정하자. $K^\bullet$이 $0$-유사연접이므로 특히
$K^\bullet$은 위로 유계이다. $i > b$이면
$H^i(K^\bullet) = 0$이도록 정수 $b$를 택하자. $n \in \mathbf{Z}$에
대한 내림 귀납법으로, $i \geq n$에 대한 유한 자유 $R$-가군 $F^i$,
$i \geq n$에 대한 미분 $d^i : F^i \to F^{i + 1}$, 그리고 미분과
호환되는 사상 $\alpha : F^i \to K^i$를 다음과 같이 구성한다.
(1) $i > n$이면 $H^i(\alpha)$는 동형사상이고 $i = n$이면 전사이며,
(2) $i > b$이면 $F^i = 0$이다. 그림은
$$
\xymatrix{
& F^n \ar[r] \ar[d]^\alpha & F^{n + 1} \ar[d]^\alpha \ar[r] & \ldots \\
K^{n - 1} \ar[r] & K^n \ar[r] & K^{n + 1} \ar[r] & \ldots
}
$$
이다. 기초 단계는 $n = b + 1$이며, 여기서는 모든 $i$에 대하여
$F^i = 0$으로 둘 수 있다.

귀납 단계. $C^\bullet$을 $\alpha$의 원뿔이라 하자(Derived Categories,
Definition \ref{derived-definition-cone}). 긴 완전 코호몰로지 열
$$
0 \to H^{n - 1}(K^\bullet) \to H^{n - 1}(C^\bullet) \to
H^n(F^\bullet) \to H^n(K^\bullet) \to H^n(C^\bullet) \to \ldots
$$
은 $i \geq n$이면 $H^i(C^\bullet) = 0$임을 보인다. 보조정리
\ref{more-algebra-lemma-cone-pseudo-coherent}에 의해 $C^\bullet$은
$(n - 1)$-유사연접이다. 보조정리 \ref{more-algebra-lemma-finite-cohomology}에
의해 $H^{n - 1}(C^\bullet)$은 유한 $R$-가군이다. 특히
$H^n(F^\bullet) \to H^n(K^\bullet)$의 핵이 유한 $R$-가군임을 알 수
있다. 유한 자유 $R$-가군 $F^{n - 1}$과 사상
% P02-E0929: frozen source targets Ker(F^n -> F^{n-1}), reversing the cochain differential; Korean preserves the formal span.
$F^{n - 1} \to \Ker(F^n \to F^{n - 1})$을 택하되,
$F^{n - 1}$이 $H^n(F^\bullet) \to H^n(K^\bullet)$의 핵 위로
전사하도록 한다.
% P02-E0930: frozen source has the broken phrase “extend our of complexes”; Korean renders the intended map-of-complexes construction.
복합체의 사상을 다음과 같이 확장한다.
$$
\xymatrix{
& F^{n - 1} \ar[d]^{\alpha^{n - 1}} \ar[r] &
F^n \ar[r] \ar[d]^\alpha &
F^{n + 1} \ar[d]^\alpha \ar[r] & \ldots \\
\ldots \ar[r] &
K^{n - 1} \ar[r] & K^n \ar[r] & K^{n + 1} \ar[r] & \ldots
}
$$
구성에 의해 $F^{n - 1} \to F^n \to K^n$의 상은
$K^{n - 1} \to K^n$의 상에 들어가므로 $\alpha^{n - 1}$이 존재한다.
% P02-E0931: frozen source begins “Dnote again”; Korean renders the intended verb.
이 확장된 복합체 사상의 원뿔도 다시 $C^\bullet$로 나타내자.
이제 완전열
$$
H^{n - 1}(F^\bullet) \to H^{n - 1}(K^\bullet) \to H^{n - 1}(C^\bullet) \to
H^n(F^\bullet) \cong H^n(K^\bullet) \to H^n(C^\bullet) \to \ldots
$$
을 얻는다. 다시 말해 $H^{n - 1}(F^\bullet) \to
H^{n - 1}(K^\bullet)$의 여핵은 유한 $R$-가군이며, 이를
$\xi_1, \ldots, \xi_r \in \Ker(K^{n - 1} \to K^n)$의 류들이
생성한다고 하자. 이제 $F^{n - 1}$을
$F^{n - 1} \oplus R^{\oplus r}$로 바꾸되, 자유 직합항의 기저 원소들이
$F^n$에서는 영으로, $K^{n - 1}$에서는 $\xi_i$로 가게 한다.
% P02-E0932: frozen source contains the malformed phrase “and the to xi_i”; Korean renders the intended parallel maps.
이로써 귀납 단계의 증명이 끝난다.
\end{proof}

\begin{lemma}
\label{more-algebra-lemma-two-out-of-three-pseudo-coherent}
$R$을 환이라 하고 $(K^\bullet, L^\bullet, M^\bullet, f, g, h)$를
$D(R)$의 구별 삼각형이라 하자. $K^\bullet, L^\bullet, M^\bullet$ 중
둘이 유사연접이면 나머지도 유사연접이다.
\end{lemma}

\begin{proof}
보조정리 \ref{more-algebra-lemma-cone-pseudo-coherent}와
\ref{more-algebra-lemma-pseudo-coherent}를 결합하면 된다.
\end{proof}

\begin{lemma}
\label{more-algebra-lemma-recognize-pseudo-coherent}
$R$을 환이라 하고 $K^\bullet$을 $R$-가군의 복합체라 하자.
$m \in \mathbf{Z}$라 하자.
\begin{enumerate}
\item 모든 $i \geq m$에 대하여 $H^i(K^\bullet) = 0$이면
$K^\bullet$은 $m$-유사연접이다.
\item $i > m$이면 $H^i(K^\bullet) = 0$이고 $H^m(K^\bullet)$이
유한 $R$-가군이면 $K^\bullet$은 $m$-유사연접이다.
\item $i > m + 1$이면 $H^i(K^\bullet) = 0$이고,
$H^{m + 1}(K^\bullet)$이 유한 표시이며 $H^m(K^\bullet)$이
유한형이면 $K^\bullet$은 $m$-유사연접이다.
\end{enumerate}
\end{lemma}

\begin{proof}
(3)을 증명하면 충분하다. $M = H^{m + 1}(K^\bullet)$로 두자.
$\tau_{\geq m + 1}K^\bullet$은 $M[- m - 1]$과 준동형적이다.
보조정리 \ref{more-algebra-lemma-n-pseudo-module}에 의해
$M[- m - 1]$은 $m$-유사연접이다. 구별 삼각형
$$
(\tau_{\leq m}K^\bullet, K^\bullet, \tau_{\geq m + 1}K^\bullet)
$$
이 있으므로(Derived Categories, Remark
\ref{derived-remark-truncation-distinguished-triangle}),
보조정리 \ref{more-algebra-lemma-cone-pseudo-coherent}에 의해
% P02-E0933: frozen source says pseudo-coherent here although the construction proves only m-pseudo-coherence; Korean preserves the stated claim.
$\tau_{\leq m}K^\bullet$이 유사연접임을 증명하면 충분하다.
가정에 의해 $H^m(\tau_{\leq m}K^\bullet)$은 유한형 $R$-가군이다.
따라서 유한 자유 $R$-가군 $E$와 사상
$E \to \Ker(d_K^m)$을 찾아 합성
$E \to \Ker(d_K^m) \to H^m(\tau_{\leq m}K^\bullet)$이 전사이게 할
수 있다. 그러면 $E[-m] \to \tau_{\leq m}K^\bullet$은
$\tau_{\leq m}K^\bullet$이 $m$-유사연접임을 보인다.
\end{proof}

\begin{lemma}
\label{more-algebra-lemma-summands-pseudo-coherent}
$R$을 환이라 하고 $m \in \mathbf{Z}$라 하자.
$K^\bullet \oplus L^\bullet$이 $m$-유사연접(각각 유사연접)이면
$K^\bullet$과 $L^\bullet$도 그러하다.
\end{lemma}

\begin{proof}
이 증명에서는 위첨자 ${}^\bullet$을 생략한다.
$K \oplus L$이 $m$-유사연접이라고 가정하자.
$K, L \in D^{-}(R)$임은 명백하다. 구별 삼각형
$$
(K \oplus L, K \oplus L, L \oplus L[1]) =
(K, K, 0) \oplus (L, L, L \oplus L[1])
$$
이 있음에 유의하자. Derived Categories, Lemma
\ref{derived-lemma-direct-sum-triangles}를 보라. 보조정리
\ref{more-algebra-lemma-cone-pseudo-coherent}에 의해
$L \oplus L[1]$은 $m$-유사연접이다. 따라서
$L[1] \oplus L[2]$도 $m$-유사연접이다. 귀납적으로
$L[n] \oplus L[n + 1]$은 $m$-유사연접이다. 보조정리
\ref{more-algebra-lemma-recognize-pseudo-coherent}에 의해 충분히 큰 $n$에 대하여
$L[n]$은 $m$-유사연접이다. 이제 구별 삼각형
$$
(L[n], L[n] \oplus L[n - 1], L[n - 1])
$$
을 사용하여 거꾸로 진행하면 $L[n], L[n - 1], \ldots, L$이
$m$-유사연접이라는 원하는 결론을 얻는다. 유사연접인 경우는 이 결과와
보조정리 \ref{more-algebra-lemma-pseudo-coherent}에서 따른다.
\end{proof}

\begin{lemma}
\label{more-algebra-lemma-complex-pseudo-coherent-modules}
$R$을 환이라 하고 $m \in \mathbf{Z}$라 하자. $K^\bullet$을
위로 유계인 $R$-가군 복합체라 하되 모든 $i$에 대하여 $K^i$가
$(m - i)$-유사연접이라고 하자. 그러면 $K^\bullet$은
$m$-유사연접이다. 특히 $K^\bullet$이 유사연접 $R$-가군들의 위로
유계인 복합체이면 $K^\bullet$은 유사연접이다.
\end{lemma}

\begin{proof}
$K^\bullet$을 (예를 들어) $\sigma_{\geq m - 1}K^\bullet$로 바꿀 수
있으므로 $K^\bullet$이 유계라고 가정할 수 있다. 그러면 복합체의
길이에 대한 귀납법으로 $K^\bullet$이 $m$-유사연접임을 알 수 있다.
각 $K^i[-i]$가 $m$-유사연접임을 이용하여 보조정리
\ref{more-algebra-lemma-cone-pseudo-coherent}와 ‘stupid’ 절단을 적용하면 된다.
마지막 명제를 위해서는 모든 $m \in \mathbf{Z}$에 대하여
$K^\bullet$이 $m$-유사연접임을 증명하면 충분하다. 보조정리
\ref{more-algebra-lemma-pseudo-coherent}를 보라. 이는 첫 부분에서 따른다.
\end{proof}
\begin{lemma}
\label{more-algebra-lemma-cohomology-pseudo-coherent}
$R$을 환이라 하고 $m \in \mathbf{Z}$라 하자.
$K^\bullet \in D^{-}(R)$이고 모든 $i$에 대하여
$H^i(K^\bullet)$이 $(m - i)$-유사연접(각각 유사연접)이라고 하자.
그러면 $K^\bullet$은 $m$-유사연접(각각 유사연접)이다.
\end{lemma}

\begin{proof}
$K^\bullet$을 $D^{-}(R)$의 대상이라 하고, 각 $H^i(K^\bullet)$이
$(m - i)$-유사연접이라고 가정하자. $H^n(K^\bullet)$이 영이 아닌
가장 큰 정수를 $n$이라 하자. $n$에 대한 귀납법으로 보조정리를
증명한다. $n < m$이면 보조정리
\ref{more-algebra-lemma-recognize-pseudo-coherent}에 의해 $K^\bullet$은
$m$-유사연접이다. $n \geq m$이면 구별 삼각형
$$
(\tau_{\leq n - 1}K^\bullet, K^\bullet, H^n(K^\bullet)[-n])
$$
이 있다(Derived Categories, Remark
\ref{derived-remark-truncation-distinguished-triangle}).
가정에 의해 $H^n(K^\bullet)[-n]$은 $m$-유사연접이므로, 보조정리
\ref{more-algebra-lemma-cone-pseudo-coherent}를 사용하면
$\tau_{\leq n - 1}K^\bullet$이 $m$-유사연접임을 증명하는 것으로
충분하다. $n$에 대한 귀납법으로 결론을 얻는다. (유사연접인 경우는
이 결과와 보조정리 \ref{more-algebra-lemma-pseudo-coherent}에서 따른다.)
\end{proof}

\begin{lemma}
\label{more-algebra-lemma-finite-push-pseudo-coherent}
$A \to B$를 환 사상이라 하자. $B$가 $A$-가군으로서 유사연접이라고
가정하자. $K^\bullet$을 $B$-가군의 복합체라 하자. 다음은 동치이다.
\begin{enumerate}
\item $K^\bullet$은 $B$-가군 복합체로서 $m$-유사연접이다.
\item $K^\bullet$은 $A$-가군 복합체로서 $m$-유사연접이다.
\end{enumerate}
유사연접성에 대해서도 같은 동치가 성립한다.
\end{lemma}

\begin{proof}
(1)을 가정하자. $E^\bullet$을 유한 자유 $B$-가군들의 유계 복합체라
하고, 차수 $> m$의 코호몰로지에서는 동형사상이고 차수 $m$에서는
전사인 사상 $\alpha : E^\bullet \to K^\bullet$을 택한다.
구별 삼각형 $(E^\bullet, K^\bullet, C(\alpha)^\bullet)$을 생각하자.
보조정리 \ref{more-algebra-lemma-recognize-pseudo-coherent}에 의해
$C(\alpha)^\bullet$은 $A$-가군의 복합체로서 $m$-유사연접이다.
따라서 $E^\bullet$이 $A$-가군의 복합체로서 유사연접임을 증명하면
충분한데, 이는 보조정리
\ref{more-algebra-lemma-complex-pseudo-coherent-modules}에서 따른다.
유사연접인 경우의 (1) $\Rightarrow$ (2)는 이 결과와 보조정리
\ref{more-algebra-lemma-pseudo-coherent}에서 따른다.

\medskip\noindent
(2)를 가정하자. $H^n(K^\bullet) \not = 0$인 가장 큰 정수를 $n$이라
하자. $n - m$에 대한 귀납법으로 $K^\bullet$이 $B$-가군의
복합체로서 $m$-유사연접임을 증명한다. $n < m$인 경우는 보조정리
\ref{more-algebra-lemma-recognize-pseudo-coherent}에서 따른다.
$E^\bullet$을 유한 자유 $A$-가군들의 유계 복합체라 하고,
차수 $> m$의 코호몰로지에서는 동형사상이고 차수 $m$에서는 전사인
사상 $\alpha : E^\bullet \to K^\bullet$을 택한다. 유도된 복합체 사상
$$
\alpha \otimes 1 : E^\bullet \otimes_A B \to K^\bullet.
$$
을 생각하자.
% P02-E0934: frozen source begins this composite with H^n(E), not H^n(E^\bullet); Korean preserves the formal span.
구성에 의해 $H^n(E) \to H^n(E^\bullet \otimes_A B) \to H^n(K^\bullet)$가
전사이고,
Example \ref{more-algebra-example-tor}의 스펙트럼 열에 의해 $i > n$이면
$H^i(E^\bullet \otimes_A B) = 0$이므로
$C(\alpha \otimes 1)^\bullet$은 차수 $\geq n$에서 비순환이다.
한편 $K^\bullet$과 $E^\bullet \otimes_A B$가 모두
$A$-가군의 복합체로서 $m$-유사연접이므로(보조정리
\ref{more-algebra-lemma-complex-pseudo-coherent-modules} 참조),
보조정리 \ref{more-algebra-lemma-cone-pseudo-coherent}에 의해
$C(\alpha \otimes 1)^\bullet$도 그러하다. 따라서 귀납 가정으로
$C(\alpha \otimes 1)^\bullet$은 $B$-가군의 복합체로서
$m$-유사연접이다. 마지막으로 보조정리
\ref{more-algebra-lemma-cone-pseudo-coherent}를 다시 적용하면
$K^\bullet$이 $B$-가군의 복합체로서 $m$-유사연접임을 알 수 있다
($E^\bullet \otimes_A B$가 $B$-가군의 복합체로서 유사연접임은
명백하다). 유사연접인 경우의 (2) $\Rightarrow$ (1)는 이 결과와
보조정리 \ref{more-algebra-lemma-pseudo-coherent}에서 따른다.
\end{proof}
\begin{lemma}
\label{more-algebra-lemma-pull-pseudo-coherent}
$A \to B$를 환 사상이라 하자. $K^\bullet$을 $A$-가군의
$m$-유사연접(각각 유사연접) 복합체라 하자. 그러면
$K^\bullet \otimes_A^{\mathbf{L}} B$는 $B$-가군의
$m$-유사연접(각각 유사연접) 복합체이다.
\end{lemma}

\begin{proof}
먼저 $K^\bullet$은 위로 유계이므로
$K^\bullet \otimes_A^{\mathbf{L}} B$가 식
(\ref{more-algebra-equation-derived-tensor-algebra})으로 정의되어 보조정리의
명제가 의미를 가짐에 유의하자. 이제 $E^\bullet$을 유한 자유
$A$-가군들의 유계 복합체라 하고, $i > m$이면 $H^i(\alpha)$가
동형사상이고 $i = m$이면 전사인
$\alpha : E^\bullet \to K^\bullet$을 택한다. 그러면 원뿔
$C(\alpha)^\bullet$은 차수 $\geq m$에서 비순환이다.
$-\otimes_A^{\mathbf{L}} B$가 완전 함자이므로 $B$-가군 복합체의
구별 삼각형
$$
(E^\bullet \otimes_A^{\mathbf{L}} B, K^\bullet \otimes_A^{\mathbf{L}} B,
C(\alpha)^\bullet \otimes_A^{\mathbf{L}} B)
$$
을 얻는다. Derived Categories, Lemma
\ref{derived-lemma-negative-vanishing}의 쌍대에 의해
$i \geq m$이면
$H^i(C(\alpha)^\bullet \otimes_A^{\mathbf{L}} B) = 0$이다.
% P02-E0935: frozen source calls E^\bullet projective over an undefined R instead of A; Korean preserves the ring symbol.
$E^\bullet$은 사영 $R$-가군들의 복합체이므로
$E^\bullet \otimes_A^{\mathbf{L}} B = E^\bullet \otimes_A B$이고, 따라서
$$
E^\bullet \otimes_A B
\longrightarrow
K^\bullet \otimes_A^{\mathbf{L}} B
$$
는 $K^\bullet \otimes_A^{\mathbf{L}} B$가 $m$-유사연접임을 보이는
$B$-가군 복합체의 사상이다. 유사연접 복합체인 경우는
$m$-유사연접 복합체인 경우와 보조정리
\ref{more-algebra-lemma-pseudo-coherent}에서 따른다.
\end{proof}

\begin{lemma}
\label{more-algebra-lemma-flat-base-change-pseudo-coherent}
$A \to B$를 평탄 환 사상이라 하자.
$M$을 $m$-유사연접(각각 유사연접) $A$-가군이라 하자. 그러면
$M \otimes_A B$는 $m$-유사연접(각각 유사연접) $B$-가군이다.
\end{lemma}

\begin{proof}
$B$가 $A$ 위에서 평탄이므로
$M \otimes_A^{\mathbf{L}} B = M \otimes_A B$라는 사실과 보조정리
\ref{more-algebra-lemma-pull-pseudo-coherent}에서 즉시 따른다.
\end{proof}

\noindent
다음 보조정리는 더 강한 보조정리
\ref{more-algebra-lemma-flat-descent-pseudo-coherent}에서도 따른다.

\begin{lemma}
\label{more-algebra-lemma-glue-pseudo-coherent}
$R$을 환이라 하자. $f_1, \ldots, f_r \in R$을 단위 아이디얼을 생성하는
원소들이라 하자. $m \in \mathbf{Z}$라 하고 $K^\bullet$을
$R$-가군의 복합체라 하자. 각 $i$에 대하여 복합체
$K^\bullet \otimes_R R_{f_i}$가 $m$-유사연접(각각 유사연접)이면,
$K^\bullet$도 $m$-유사연접(각각 유사연접)이다.
\end{lemma}

\begin{proof}
$- \otimes_R R_{f_i}$가 완전 함자이고, 따라서
$$
H^i(K^\bullet)_{f_i} =
H^i(K^\bullet) \otimes_R R_{f_i} = H^i(K^\bullet \otimes_R R_{f_i}).
$$
임을 더 언급하지 않고 사용하겠다.
$i = 1, \ldots, r$에 대하여 $K^\bullet \otimes_R R_{f_i}$가
$m$-유사연접이라고 가정하자. 어떤 $i$에 대하여
$H^n(K^\bullet \otimes_R R_{f_i})$가 영이 아닌 가장 큰 정수
$n \in \mathbf{Z}$를 택하자. 이는 특히 $i > n$이면
$H^i(K^\bullet) = 0$이고 $H^n(K^\bullet) \not = 0$임을 뜻한다.
Algebra, Lemma \ref{algebra-lemma-cover}를 보라.
$n - m$에 대한 귀납법으로 보조정리를 증명한다.
$n < m$이면 보조정리 \ref{more-algebra-lemma-recognize-pseudo-coherent}에 의해
참이다. $n \geq m$이면 각 $i$에 대하여
$H^n(K^\bullet)_{f_i}$는 유한 $R_{f_i}$-가군이다. 보조정리
\ref{more-algebra-lemma-finite-cohomology}를 보라. 따라서
$H^n(K^\bullet)$은 유한 $R$-가군이다. Algebra, Lemma
\ref{algebra-lemma-cover}를 보라. 유한 자유 $R$-가군 $E$와 전사
$E \to H^n(K^\bullet)$을 택한다. $E$가 사영이므로 이를 복합체 사상
$\alpha : E[-n] \to K^\bullet$로 올릴 수 있다. 그러면 원뿔
$C(\alpha)^\bullet$의 코호몰로지는 차수 $\geq n$에서 소멸한다.
한편 각 $i$에 대하여 복합체
$C(\alpha)^\bullet \otimes_R R_{f_i}$는 $m$-유사연접이다.
보조정리 \ref{more-algebra-lemma-cone-pseudo-coherent}를 보라. 따라서 귀납
가정에 의해 $C(\alpha)^\bullet$은 $R$-가군의 복합체로서
$m$-유사연접이다. 보조정리 \ref{more-algebra-lemma-cone-pseudo-coherent}를
한 번 더 적용하면 결론을 얻는다.
\end{proof}

\begin{lemma}
\label{more-algebra-lemma-flat-descent-pseudo-coherent}
$R$을 환이라 하고 $m \in \mathbf{Z}$라 하자. $K^\bullet$을
$R$-가군의 복합체라 하고 $R \to R'$을 충실히 평탄한 환 사상이라
하자. 복합체 $K^\bullet \otimes_R R'$이
$m$-유사연접(각각 유사연접)이면 $K^\bullet$도
$m$-유사연접(각각 유사연접)이다.
\end{lemma}

\begin{proof}
$- \otimes_R R'$이 완전 함자이고, 따라서
$$
H^i(K^\bullet) \otimes_R R' = H^i(K^\bullet \otimes_R R').
$$
임을 더 언급하지 않고 사용하겠다. $K^\bullet \otimes_R R'$이
$m$-유사연접이라고 가정하자. $H^n(K^\bullet)$이 영이 아닌 가장
큰 정수를 $n \in \mathbf{Z}$라 하자. 그러면 $n$도
$H^n(K^\bullet \otimes_R R')$이 영이 아닌 가장 큰 정수이기도 하다.
$n - m$에 대한 귀납법으로 보조정리를 증명한다. $n < m$이면 보조정리
\ref{more-algebra-lemma-recognize-pseudo-coherent}에 의해 참이다.
$n \geq m$이면 $H^n(K^\bullet) \otimes_R R'$은 유한
$R'$-가군이다. 보조정리 \ref{more-algebra-lemma-finite-cohomology}를 보라.
따라서 $H^n(K^\bullet)$은 유한 $R$-가군이다. Algebra, Lemma
\ref{algebra-lemma-descend-properties-modules}를 보라. 유한 자유
$R$-가군 $E$와 전사 $E \to H^n(K^\bullet)$을 택한다.
$E$가 사영이므로 이를 복합체 사상
$\alpha : E[-n] \to K^\bullet$로 올릴 수 있다. 그러면 원뿔
$C(\alpha)^\bullet$의 코호몰로지는 차수 $\geq n$에서 소멸한다.
한편 복합체 $C(\alpha)^\bullet \otimes_R R'$은
$m$-유사연접이다. 보조정리 \ref{more-algebra-lemma-cone-pseudo-coherent}를
보라. 따라서 귀납 가정에 의해 $C(\alpha)^\bullet$은
$R$-가군의 복합체로서 $m$-유사연접이다. 보조정리
\ref{more-algebra-lemma-cone-pseudo-coherent}를 한 번 더 적용하면 결론을 얻는다.
\end{proof}

\begin{lemma}
\label{more-algebra-lemma-tensor-pseudo-coherent}
$R$을 환이라 하고 $K, L$을 $D(R)$의 대상이라 하자.
\begin{enumerate}
\item $K$가 $n$-유사연접이고 $i > a$이면 $H^i(K) = 0$이며,
$L$이 $m$-유사연접이고 $j > b$이면 $H^j(L) = 0$이라면
$K \otimes_R^\mathbf{L} L$은 $t = \max(m + a, n + b)$에 대하여
$t$-유사연접이다.
\item $K$와 $L$이 유사연접이면 $K \otimes_R^\mathbf{L} L$도
유사연접이다.
\end{enumerate}
\end{lemma}

\begin{proof}
(1)의 증명. 유한 자유 $R$-가군들의 유계 복합체 $K^\bullet$과
$L^\bullet$, 그리고 사상 $\alpha : K^\bullet \to K$ 및
$\beta : L^\bullet \to L$이 존재한다고 가정해도 된다.
% P02-E0936: frozen source twice says “and isomorphism”; Korean renders the intended article.
$i > n$이면 $H^i(\alpha)$가 동형사상이고 $i = n$이면 전사이며,
$i > m$이면 $H^i(\beta)$가 동형사상이고 $i = m$이면 전사라고 하자.
그러면 사상
$$
\alpha \otimes^\mathbf{L} \beta :
\text{Tot}(K^\bullet \otimes_R L^\bullet)
\to K \otimes_R^\mathbf{L} L
$$
은 차수 $i$의 코호몰로지에서 $i > t$이면 동형사상을 유도하고
$i = t$이면
전사를 유도한다. 이는 Tor의 스펙트럼 열에서 따른다(세부사항은
생략한다). (2)는 (1)과 보조정리
\ref{more-algebra-lemma-pseudo-coherent}에서 따른다.
\end{proof}

\begin{lemma}
\label{more-algebra-lemma-Noetherian-pseudo-coherent}
$R$을 Noether 환이라 하자. 그러면
\begin{enumerate}
\item $R$-가군 복합체 $K^\bullet$이 $m$-유사연접일 필요충분조건은
$K^\bullet \in D^{-}(R)$이고 모든 $i \geq m$에 대하여
$H^i(K^\bullet)$이 유한 $R$-가군인 것이다.
\item $R$-가군 복합체 $K^\bullet$이 유사연접일 필요충분조건은
$K^\bullet \in D^{-}(R)$이고 모든 $i$에 대하여
$H^i(K^\bullet)$이 유한 $R$-가군인 것이다.
\item $R$-가군이 유사연접일 필요충분조건은 그것이 유한인 것이다.
\end{enumerate}
\end{lemma}

\begin{proof}
Algebra, Lemma \ref{algebra-lemma-resolution-by-finite-free}에서 임의의
유한 $R$-가군이 유사연접임을 보였다. 한편 유사연접 가군은 유한이다.
보조정리 \ref{more-algebra-lemma-n-pseudo-module}를 보라. 따라서 (3)이 성립한다.
$K^\bullet$이 $m$-유사연접 복합체라고 하자. 그러면 유한 자유
$R$-가군들의 유계 복합체 $E^\bullet$이 존재하여
$i > m$이면 $H^i(K^\bullet)$은 $H^i(E^\bullet)$과 동형이고,
$H^m(K^\bullet)$은 $H^m(E^\bullet)$의 몫이다. 따라서 각
$i \geq m$에 대하여 $H^i(K^\bullet)$이 유한 가군임은 명백하다.
(1)의 역은 보조정리 \ref{more-algebra-lemma-cohomology-pseudo-coherent}와
(3)에서 따른다. (2)는 (1)과 보조정리
\ref{more-algebra-lemma-pseudo-coherent}에서 따른다.
\end{proof}

\begin{lemma}
\label{more-algebra-lemma-coherent-pseudo-coherent}
$R$을 연접환이라 하자(Algebra, Definition
\ref{algebra-definition-coherent}). $K \in D^-(R)$라 하자.
다음은 동치이다.
\begin{enumerate}
\item $K$는 $m$-유사연접이다.
\item $H^m(K)$는 유한 $R$-가군이고 $i > m$이면 $H^i(K)$는
연접이다.
\item $H^m(K)$는 유한 $R$-가군이고 $i > m$이면 $H^i(K)$는
유한 표시이다.
\end{enumerate}
따라서 $K$가 유사연접일 필요충분조건은 모든 $i$에 대하여
$H^i(K)$가 연접 가군인 것이다.
\end{lemma}

\begin{proof}
$R$-가군 $M$이 연접일 필요충분조건은 유한 표시인 것임을 상기하자
(Algebra, Lemma \ref{algebra-lemma-coherent-ring}).
이것이 (2)와 (3)의 동치를 설명한다. 이들이 성립하고 완전열
$0 \to N \to R^{\oplus m} \to M \to 0$을 택하면, Algebra, Lemma
\ref{algebra-lemma-coherent}에 의해 $N$은 연접이다. 따라서 이
경우에 $N$에 대하여 이 절차를 반복하면 유한 자유 $R$-가군들에 의한
분해
$$
\ldots \to R^{\oplus n} \to R^{\oplus m} \to M \to 0
$$
을 얻는다. 다시 말해 $M$은 유사연접이다. (1)과 (2)의 동치는 이
결과와 보조정리 \ref{more-algebra-lemma-cohomology-pseudo-coherent} 및
\ref{more-algebra-lemma-n-pseudo-module}에서 따른다. 마지막 명제는 (1)과 (2)의
동치와 보조정리 \ref{more-algebra-lemma-pseudo-coherent}를 결합하면 따른다.
\end{proof}
\section{유사연접 가군, II}
\label{more-algebra-section-pseudo-coherent-bis}

\noindent
절 \ref{more-algebra-section-pseudo-coherent}에서 시작한 논의를 계속한다.

\begin{lemma}
\label{more-algebra-lemma-pseudo-coherence-colimit-ext}
$R$을 환이라 하자. $M = \colim M_i$를 $R$-가군들의 여과 여극한이라
하자. $K \in D(R)$가 $m$-유사연접이라고 하자. 그러면 $n < -m$에
대하여 $\colim \Ext^n_R(K, M_i) = \Ext^n_R(K, M)$이고,
$\colim \Ext^{-m}_R(K, M_i) \to \Ext^{-m}_R(K, M)$은 단사이다.
\end{lemma}

\begin{proof}
정의에 의해 $D(R)$의 구별 삼각형
$$
E \to K \to L \to E[1]
$$
을 찾을 수 있는데, 여기서 $E$는 유한 자유 $R$-가군들의 유계
복합체로 표현되고 $i \geq m$이면 $H^i(L) = 0$이다. 그러면 임의의
$R$-가군 $N$과 $n \leq -m$에 대하여
$\Ext^n_R(L, N) = 0$이다. Derived Categories, Lemma
\ref{derived-lemma-negative-exts}를 보라. 이 구별 삼각형에 결부된
$\Ext$의 긴 완전열로부터, $n < -m$이면
$\Ext^n_R(K, N) \to \Ext^n_R(E, N)$이 동형사상이고 $n = -m$이면
단사임을 알 수 있다. 따라서 $E$가 유한 자유 $R$-가군들의 유계
복합체 $E^\bullet$로 표현될 수 있을 때
$M \mapsto \Ext_R^n(E, M)$이 여과 여극한과 가환함을 증명하면 충분하다.
가군 $\Ext^n_R(E, M)$은 복합체
$\Hom_R(E^\bullet, M)$으로 계산된다. Derived Categories, Lemma
\ref{derived-lemma-morphisms-from-projective-complex}를 보라.
$E^p$가 유한 자유이므로 함자 $M \mapsto \Hom_R(E^p, M)$은 여과
여극한과 가환한다. 따라서 복합체로서
$\Hom_R(E^\bullet, M) = \colim \Hom_R(E^\bullet, M_i)$이다.
여과 여극한은 완전하므로 결론을 얻는다(Algebra, Lemma
\ref{algebra-lemma-directed-colimit-exact}).
\end{proof}

\begin{lemma}
\label{more-algebra-lemma-characterize-pseudo-coherent-colimit-ext}
$R$을 환이라 하고 $K \in D^-(R)$라 하자. $m \in \mathbf{Z}$라 하자.
$K$가 $m$-유사연접일 필요충분조건은 $R$-가군들의 임의의 여과 여극한
$M = \colim M_i$에 대하여 $n < -m$이면
$\colim \Ext^n_R(K, M_i) = \Ext^n_R(K, M)$이고,
$\colim \Ext^{-m}_R(K, M_i) \to \Ext^{-m}_R(K, M)$가 단사인 것이다.
\end{lemma}

\begin{proof}
한쪽 함의는 보조정리 \ref{more-algebra-lemma-pseudo-coherence-colimit-ext}에서
증명했다. $R$-가군들의 임의의 여과 여극한 $M = \colim M_i$에 대하여
$n < -m$이면 $\colim \Ext^n_R(K, M_i) = \Ext^n_R(K, M)$이고
$\colim \Ext^{-m}_R(K, M_i) \to \Ext^{-m}_R(K, M)$가 단사라고
가정하자. $K$가 $m$-유사연접임을 보이겠다.

\medskip\noindent
$H^t(K)$가 영이 아닌 가장 큰 정수를 $t$라 하자. $t$에 대한
귀납법을 사용한다. $t < m$이면 보조정리
\ref{more-algebra-lemma-recognize-pseudo-coherent}에 의해 $K$는 $m$-유사연접이다.
$t \geq m$이면 $\Hom_R(H^t(K), M) = \Ext^{-t}_R(K, M)$이므로 임의의
여과 여극한 $M = \colim M_i$에 대하여
$\colim \Hom_R(H^t(K), M_i) \to \Hom_R(H^t(K), M)$가 단사임을
알 수 있다. Algebra, Lemma
\ref{algebra-lemma-characterize-finite-module-hom}에 의해 이는
$H^t(K)$가 유한 $R$-가군임을 뜻한다. 유한 자유 $R$-가군 $F$와 전사
$F \to H^t(K)$을 택한다. 이를 $D(R)$의 사상 $F[-t] \to K$로
올리고 $D(R)$의 구별 삼각형
$$
F[-t] \to K \to L \to F[-t + 1]
$$
을 택할 수 있다. 그러면 $i \geq t$이면 $H^i(L) = 0$이다.
또한 이 구별 삼각형에 결부된 $\Ext$의 긴 완전열에 대한 생략한
짧은 논증으로 $L$이 $K$에 둔 가정을 이어받음을 알 수 있다.
$t$에 대한 귀납법으로 $L$은 $m$-유사연접이다. 따라서 보조정리
\ref{more-algebra-lemma-cone-pseudo-coherent}에 의해 $K$도 $m$-유사연접이다.
\end{proof}

\begin{lemma}
\label{more-algebra-lemma-pseudo-coherence-and-ext}
$R$을 환이라 하고 $L$, $M$, $N$을 $R$-가군이라 하자.
\begin{enumerate}
\item $M$이 유한 표시이고 $L$이 평탄이면 표준 사상
$\Hom_R(M, N) \otimes_R L \to \Hom_R(M, N \otimes_R L)$은
동형사상이다.
\item $M$이 $(-m)$-유사연접이고 $L$이 평탄이면 표준 사상
$\Ext^i_R(M, N) \otimes_R L \to \Ext^i_R(M, N \otimes_R L)$은
$i < m$에 대하여 동형사상이다.
\end{enumerate}
\end{lemma}

\begin{proof}
항들이 자유 $R$-가군인 분해 $F_\bullet \to M$을 택하자.
Algebra, Lemma \ref{algebra-lemma-resolution-by-finite-free}를 보라.
복합체 $\Hom_R(F_\bullet, N)$은 $\Ext^i_R(M, N)$을 계산하고,
복합체 $\Hom_R(F_\bullet, N \otimes_R L)$은
$\Ext^i_R(M, N \otimes_R L)$을 계산한다. 항상 공사슬 복합체의 사상
$$
\Hom_R(F_\bullet, N) \otimes_R L
\longrightarrow
\Hom_R(F_\bullet, N \otimes_R L)
$$
이 존재하며, 이는 모든 $i \geq 0$에 대하여 표준 사상
$\Ext^i_R(M, N) \otimes_R L \to \Ext^i_R(M, N \otimes_R L)$을
유도한다(예를 들어 이 사상들이 분해 $F_\bullet$의 선택에 의존하지
않는다는 의미에서 표준이다). $L$이 평탄이면 $L$과의 텐서곱이
코호몰로지를 취하는 것과 가환하므로 복합체
$\Hom_R(F_\bullet, N) \otimes_R L$은
$\Ext^i_R(M, N) \otimes_R L$을 계산한다.

\medskip\noindent
이제 $M$이 $(-m)$-유사연접이면 $i = 0, \ldots, m$에 대하여
$F_i$가 유한 자유이도록 $F_\bullet$을 택할 수 있다. 그러면 위에
표시한 공사슬 복합체의 사상은 차수 $\leq m$에서 동형사상이고, 따라서
차수 $< m$의 코호몰로지 군들에서 동형사상이다. 이것이 (2)를
증명한다. $M$이 유한 표시이면 보조정리
\ref{more-algebra-lemma-n-pseudo-module}에 의해 $M$은 $(-1)$-유사연접이고,
$\Hom = \Ext^0$이므로 결과를 얻는다.
\end{proof}

\begin{lemma}
\label{more-algebra-lemma-pseudo-coherence-and-base-change-ext}
$R \to R'$을 평탄 환 사상이라 하고 $M$, $N$을 $R$-가군이라 하자.
\begin{enumerate}
\item $M$이 유한 표시 $R$-가군이면
$\Hom_R(M, N) \otimes_R R' = \Hom_{R'}(M \otimes_R R', N \otimes_R R')$이다.
\item $M$이 $(-m)$-유사연접이면 $i < m$에 대하여
$\Ext^i_R(M, N) \otimes_R R' = \Ext^i_{R'}(M \otimes_R R', N \otimes_R R')$이다.
\end{enumerate}
특히 $R$이 Noether이고 $M$이 유한 가군이면 이는 모든 $i$에 대하여
성립한다.
\end{lemma}

\begin{proof}
Algebra, Lemma \ref{algebra-lemma-flat-base-change-ext}에 의해
$\Ext^i_{R'}(M \otimes_R R', N \otimes_R R') =
\Ext^i_R(M, N \otimes_R R')$이다. 이를 보조정리
\ref{more-algebra-lemma-pseudo-coherence-and-ext}와 결합하면
% P02-E0937: frozen source says “(1) and (2) holds”; Korean renders the compound predicate correctly.
(1)과 (2)를 얻는다. 마지막 명제는 이 결과와 보조정리
\ref{more-algebra-lemma-Noetherian-pseudo-coherent}에서 따른다.
\end{proof}
\begin{lemma}
\label{more-algebra-lemma-pseudo-coherent-tensor}
$R$을 환이라 하고 $K \in D^-(R)$라 하자. 다음은 동치이다.
\begin{enumerate}
\item $K$는 유사연접이다.
\item $R$-가군의 모든 족 $(Q_{\alpha})_{\alpha \in A}$에 대하여 표준 사상
$$
\alpha :
K \otimes_R^\mathbf{L} \left( \prod\nolimits_\alpha Q_{\alpha} \right)
\longrightarrow
\prod\nolimits_\alpha (K \otimes_R^\mathbf{L} Q_{\alpha})
$$
는 $D(R)$에서 동형사상이다.
\item 모든 $R$-가군 $Q$와 모든 집합 $A$에 대하여 표준 사상
$$
\beta : K \otimes_R^\mathbf{L} Q^A \longrightarrow (K \otimes_R^\mathbf{L} Q)^A
$$
는 $D(R)$에서 동형사상이다.
\item 모든 집합 $A$에 대하여 표준 사상
$$
\gamma : K \otimes_R^\mathbf{L} R^A \longrightarrow K^A
$$
는 $D(R)$에서 동형사상이다.
\end{enumerate}
$m \in \mathbf{Z}$가 주어졌을 때 다음은 동치이다.
\begin{enumerate}
\item[(a)] $K$는 $m$-유사연접이다.
\item[(b)] $R$-가군의 모든 족 $(Q_{\alpha})_{\alpha \in A}$에 대하여
$\alpha$가 위와 같을 때 $i > m$이면 $H^i(\alpha)$는 동형사상이고
$i = m$이면 전사이다.
\item[(c)] 모든 $R$-가군 $Q$와 모든 집합 $A$에 대하여 $\beta$가 위와
같을 때 $i > m$이면 $H^i(\beta)$는 동형사상이고 $i = m$이면 전사이다.
\item[(d)] 모든 집합 $A$에 대하여 $\gamma$가 위와 같을 때
$i > m$이면 $H^i(\gamma)$는 동형사상이고 $i = m$이면 전사이다.
\end{enumerate}
\end{lemma}

\begin{proof}
$K$가 유사연접이면 $K$는 유한 자유 $R$-가군들의 위로 유계인 복합체로
표현될 수 있다. 그러면 유도 텐서곱은 이 복합체와 텐서곱하여 계산된다.
또한 $D(R)$의 곱은 대표 복합체들을 어떻게 택하든 그 곱을 취하여
주어진다. 따라서 가군에 대한 대응하는 사실에 의해 (1)은 (2), (3),
(4)를 함의한다. Algebra, Proposition
\ref{algebra-proposition-fp-tensor}를 보라.

\medskip\noindent
% P02-E0939: frozen source parenthetical omits the complementizer before “the tensor product is right exact”; Korean renders the intended dependency.
같은 방식으로(텐서곱이 오른쪽 완전함을 사용하여) (a)가 (b), (c),
(d)를 함의함을 독자가 보일 수 있다.

\medskip\noindent
(4)를 가정하자. $K$가 유사연접임을 보이려면 모든 $m$에 대하여
$K$가 $m$-유사연접임을 보이면 충분하다(보조정리
\ref{more-algebra-lemma-pseudo-coherent}).
% P02-E0938: frozen source says “finish then proof”; Korean renders the intended “finish the proof”.
따라서 증명을 끝내려면 (d)가 (a)를 함의함을 증명하면 충분하다.

\medskip\noindent
(d)를 가정하자. $H^i(K)$가 영이 아닌 가장 큰 정수를 $i$라 하자.
$i < m$이면 끝난다. 그렇지 않으면 (d)와 위에서 준 $D(R)$의 곱의
기술로부터 $H^i(K) \otimes_R R^A \to H^i(K)^A$가 전사임을 알 수
있다. 따라서 Algebra, Proposition
\ref{algebra-proposition-fg-tensor}에 의해 $H^i(K)$는 유한 생성
$R$-가군이다. 그러므로 차수 $i$에 놓인 하나의 유한 자유 가군으로
이루어진 복합체 $L$과 복합체 사상 $L \to K$를 택하여
$H^i(L) \to H^i(K)$가 전사이게 할 수 있다. 특히 $L$은 (1), (2),
(3), (4)를 만족한다. 구별 삼각형
$$
L \to K \to M \to L[1]
$$
을 택하자. 그러면 $j \geq i$이면 $H^j(M) = 0$임을 알 수 있다.
한편 생략하는 짧은 논증에 의해 $M$은 여전히 성질 (d)를 갖는다.
$i$에 대한 귀납법으로 $M$이 $m$-유사연접임을 알 수 있다. 따라서
보조정리 \ref{more-algebra-lemma-cone-pseudo-coherent}에 의해 $K$도
$m$-유사연접이다.
\end{proof}

\begin{lemma}
\label{more-algebra-lemma-detect-cohomology-pseudo-coherent}
$R$을 환이라 하고 $K \in D(R)$를 유사연접이라 하자.
$i \in \mathbf{Z}$라 하자. 유한 표시 $R$-가군 $M$과 $D(R)$의 사상
$K \to M[-i]$가 존재하여 단사 $H^i(K) \to M$을 유도한다.
\end{lemma}

\begin{proof}
정의 \ref{more-algebra-definition-pseudo-coherent}에 의해 $K$를 유한 자유
$R$-가군들의 복합체 $P^\bullet$로 표현할 수 있다.
$M = \Coker(P^{i - 1} \to P^i)$로 두자.
\end{proof}

\begin{lemma}
\label{more-algebra-lemma-detect-cohomology}
$A$를 Noether 환이라 하고 $K \in D(A)$를 유사연접이라 하자. 즉,
$K \in D^-(A)$이고 코호몰로지 가군들이 유한이라고 하자.
$\mathfrak m$을 $A$의 극대 아이디얼이라 하자.
$H^i(K)/\mathfrak m H^i(K) \not = 0$이면, $\mathfrak m$의 어떤
거듭제곱에 의해 소멸되는 유한 $A$-가군 $E$와
$H^i(K)$ 위에서 영이 아닌 사상 $K \to E[-i]$가 존재한다.
\end{lemma}

\begin{proof}
(보조정리의 명제에 있는 유사연접성의 동치인 정식화는 보조정리
\ref{more-algebra-lemma-Noetherian-pseudo-coherent}이다.) 보조정리
\ref{more-algebra-lemma-detect-cohomology-pseudo-coherent}에서처럼
$K \to M[-i]$를 택하자. Artin--Rees(Algebra, Lemma
\ref{algebra-lemma-Artin-Rees})에 의해
$H^i(K) \cap \mathfrak m^n M \subset \mathfrak m H^i(K)$가 되도록
$n$을 택할 수 있다. $E = M/\mathfrak m^n M$로 두자.
\end{proof}
\section{Tor 차원}
\label{more-algebra-section-tor}

\noindent
사영 가군으로 분해하는 대신 평탄 가군에 의한 분해를 살펴볼 수 있다.
이는 다음 개념으로 이어진다.

\begin{definition}
\label{more-algebra-definition-tor-amplitude}
$R$을 환이라 하자.
% P02-E0940: frozen source repeats the malformed “Denote D(R) its derived category”; Korean renders the intended notation.
그 유도 범주를 $D(R)$로 나타내자. $a, b \in \mathbf{Z}$라 하자.
\begin{enumerate}
\item $D(R)$의 대상 $K^\bullet$에 대하여, 모든 $R$-가군 $M$과 모든
$i \not \in [a, b]$에 대해
$H^i(K^\bullet \otimes_R^\mathbf{L} M) = 0$이면
{\it $[a, b]$에 Tor 진폭을 갖는다}고 한다.
\item $D(R)$의 대상 $K^\bullet$이 어떤 $a, b$에 대하여
$[a, b]$에 Tor 진폭을 가지면 {\it 유한 Tor 차원}을 갖는다고 한다.
\item $R$-가군 $M$에 대하여, $D(R)$의 대상으로서 $M[0]$이
$[-d, 0]$에 Tor 진폭을 가지면 {\it Tor 차원이 $\leq d$}라고 한다.
\item $R$-가군 $M$에 대하여, $D(R)$의 대상으로서 $M[0]$이
유한 Tor 차원을 가지면 {\it 유한 Tor 차원}을 갖는다고 한다.
\end{enumerate}
\end{definition}

\noindent
$K^\bullet$이 유한 Tor 차원을 가지면 $K^\bullet \in D^b(R)$임을
관찰할 수 있다.

\begin{lemma}
\label{more-algebra-lemma-last-one-flat}
$R$을 환이라 하고 $K^\bullet$을 $[a, b]$에 Tor 진폭을 갖는 평탄
$R$-가군들의 위로 유계인 복합체라 하자. 그러면
$\Coker(d_K^{a - 1})$은 평탄 $R$-가군이다.
\end{lemma}

\begin{proof}
$K^\bullet$은 평탄 가군들의 위로 유계인 복합체이므로
$K^\bullet \otimes_R M = K^\bullet \otimes_R^{\mathbf{L}} M$이다.
따라서 모든 $R$-가군 $M$에 대하여 열
$$
K^{a - 2} \otimes_R M \to K^{a - 1} \otimes_R M \to K^a \otimes_R M
$$
은 가운데에서 완전하다.
$K^{a - 2} \to K^{a - 1} \to K^a \to \Coker(d_K^{a - 1}) \to 0$은
평탄 분해이므로 이는 모든 $R$-가군 $M$에 대하여
$\text{Tor}_1^R(\Coker(d_K^{a - 1}), M) = 0$임을 뜻한다. 따라서
$\Coker(d_K^{a - 1})$은 평탄하다. Algebra, Lemma
\ref{algebra-lemma-characterize-flat}을 보라.
\end{proof}

\begin{lemma}
\label{more-algebra-lemma-tor-amplitude}
$R$을 환이라 하고 $K^\bullet$을 $D(R)$의 대상이라 하자.
$a, b \in \mathbf{Z}$라 하자. 다음은 동치이다.
\begin{enumerate}
\item $K^\bullet$은 $[a, b]$에 Tor 진폭을 갖는다.
\item $K^\bullet$은 평탄 $R$-가군들의 복합체 $E^\bullet$과
준동형적이며, $i \not \in [a, b]$이면 $E^i = 0$이다.
\end{enumerate}
\end{lemma}

\begin{proof}
(2)가 성립하면
$K^\bullet \otimes_R^\mathbf{L} M = E^\bullet \otimes_R M$으로 계산할
수 있으므로 (1)이 성립함은 명백하다. (1)을 가정하자.
$K^\bullet$을 $i > b$이면 $K^i = 0$인 사영 분해로 바꿀 수 있다.
Derived Categories, Lemma
\ref{derived-lemma-projective-resolutions-exist}를 보라.
$E^\bullet = \tau_{\geq a}K^\bullet$로 두자. $E^a$가 평탄이라는
점만 빼면 모두 명백한데, 이 점은 보조정리
\ref{more-algebra-lemma-last-one-flat}과 정의에서 즉시 따른다.
\end{proof}

\begin{lemma}
\label{more-algebra-lemma-bounded-below-tor-amplitude}
$R$을 환이라 하자. $a \in \mathbf{Z}$라 하고 $K$를 $D(R)$의
대상이라 하자. 다음은 동치이다.
\begin{enumerate}
\item $K$는 $[a, \infty]$에 Tor 진폭을 갖는다.
\item $K$는 항들이 평탄 $R$-가군이고
$i \not \in [a, \infty]$이면 $E^i = 0$인 K-평탄 복합체
$E^\bullet$과 준동형적이다.
\end{enumerate}
\end{lemma}

\begin{proof}
(2) $\Rightarrow$ (1)은 즉시 따른다. (1)을 가정하자. 먼저 $K$를
대표하는 평탄 항들의 K-평탄 복합체 $K^\bullet$을 택한다. 보조정리
\ref{more-algebra-lemma-K-flat-resolution}를 보라. 임의의 $R$-가군 $M$에 대하여
$$
K^{n - 1} \otimes_R M \to K^n \otimes_R M \to K^{n + 1} \otimes_R M
$$
의 코호몰로지는 $H^n(K \otimes_R^\mathbf{L} M)$을 계산한다.
이는 $n < a$이면 언제나 영이다. 따라서 복합체
$\ldots \to K^{a - 1} \to K^a \to K^{a + 1}$에 보조정리
\ref{more-algebra-lemma-last-one-flat}을 적용하면
$N = \Coker(K^{a - 1} \to K^a)$이 평탄 $R$-가군임을 알 수 있다.
다음과 같이 둔다.
$$
E^\bullet = \tau_{\geq a}K^\bullet =
(\ldots \to 0 \to N \to K^{a + 1} \to \ldots )
$$
$K^\bullet \to E^\bullet$의 핵 $L^\bullet$은 복합체
$$
L^\bullet = (\ldots \to K^{a - 1} \to I \to 0 \to \ldots)
$$
이고, 여기서 $I \subset K^a$는 $K^{a - 1} \to K^a$의 상이다.
짧은 완전열 $0 \to I \to K^a \to N \to 0$이 있으므로 $I$는 평탄
$R$-가군이다. 따라서 $L^\bullet$은 평탄 가군들의 위로 유계인
복합체이고, 보조정리
\ref{more-algebra-lemma-derived-tor-quasi-isomorphism}에 의해 K-평탄이다.
보조정리 \ref{more-algebra-lemma-K-flat-two-out-of-three-ses}에 의해
$E^\bullet$도 K-평탄임이 따른다.
\end{proof}

\begin{lemma}
\label{more-algebra-lemma-cone-tor-amplitude}
$R$을 환이라 하고 $(K^\bullet, L^\bullet, M^\bullet, f, g, h)$를
$D(R)$의 구별 삼각형이라 하자. $a, b \in \mathbf{Z}$라 하자.
\begin{enumerate}
\item $K^\bullet$이 $[a + 1, b + 1]$에 Tor 진폭을 갖고
$L^\bullet$이 $[a, b]$에 Tor 진폭을 가지면 $M^\bullet$은
$[a, b]$에 Tor 진폭을 갖는다.
\item $K^\bullet, M^\bullet$이 $[a, b]$에 Tor 진폭을 가지면
$L^\bullet$도 $[a, b]$에 Tor 진폭을 갖는다.
\item $L^\bullet$이 $[a + 1, b + 1]$에 Tor 진폭을 갖고
$M^\bullet$이 $[a, b]$에 Tor 진폭을 가지면 $K^\bullet$은
$[a + 1, b + 1]$에 Tor 진폭을 갖는다.
\end{enumerate}
\end{lemma}

\begin{proof}
생략한다. 힌트: 구별 삼각형에 결부된 긴 완전 코호몰로지 열과
$- \otimes_R^{\mathbf{L}} M$이 구별 삼각형을 보존한다는 사실에서
따를 뿐이다. (2)가 가장 쉽게 증명되며, 나머지는 평행이동으로
그것에서 따른다.
\end{proof}

\begin{lemma}
\label{more-algebra-lemma-tor-dimension}
$R$을 환이라 하고 $M$을 $R$-가군이라 하자. $d \geq 0$이라 하자.
다음은 동치이다.
\begin{enumerate}
\item $M$의 Tor 차원은 $\leq d$이다.
\item $F_i$가 평탄 $R$-가군인 분해
$$
0 \to F_d \to \ldots \to F_1 \to F_0 \to M \to 0
$$
이 존재한다.
\end{enumerate}
특히 $R$-가군의 Tor 차원이 $0$일 필요충분조건은 그것이 평탄
$R$-가군인 것이다.
\end{lemma}

\begin{proof}
(2)를 가정하자. $E^{-i} = F_i$인 복합체 $E^\bullet$은 $M$과
준동형적이다. 따라서 보조정리 \ref{more-algebra-lemma-tor-amplitude}에 의해
$M$의 Tor 차원은 $d$ 이하이다. 반대로 (1)을 가정하자.
$P^\bullet \to M$을 $M$의 사영 분해라 하자. 보조정리
\ref{more-algebra-lemma-last-one-flat}에 의해 $\tau_{\geq -d}P^\bullet$은
$M$의 길이 $d$인 평탄 분해이다. 즉 (2)가 성립한다.
\end{proof}

\begin{lemma}
\label{more-algebra-lemma-summands-tor-amplitude}
$R$을 환이라 하고 $a, b \in \mathbf{Z}$라 하자.
$K^\bullet \oplus L^\bullet$이 $[a, b]$에 Tor 진폭을 가지면
$K^\bullet$과 $L^\bullet$도 그러하다.
\end{lemma}

\begin{proof}
Tor 함자들이 가법적이라는 사실에서 명백하다.
\end{proof}

\begin{lemma}
\label{more-algebra-lemma-complex-finite-tor-dimension-modules}
$R$을 환이라 하고 $K^\bullet$을 $R$-가군들의 유계 복합체라 하되,
모든 $i$에 대하여 $K^i$가 $[a - i, b - i]$에 Tor 진폭을 갖는다고
하자. 그러면 $K^\bullet$은 $[a, b]$에 Tor 진폭을 갖는다.
특히 $K^\bullet$이 유한 Tor 차원의 $R$-가군들로 이루어진 유한
복합체이면 $K^\bullet$은 유한 Tor 차원을 갖는다.
\end{lemma}

\begin{proof}
유한 복합체의 길이에 대한 귀납으로 따른다. 보조정리
\ref{more-algebra-lemma-cone-tor-amplitude}와 ‘stupid’ 절단을 사용하라.
\end{proof}
\begin{lemma}
\label{more-algebra-lemma-cohomology-tor-amplitude}
$R$을 환이라 하고 $a, b \in \mathbf{Z}$라 하자.
$K^\bullet \in D^b(R)$이고 모든 $i$에 대하여 $H^i(K^\bullet)$이
$[a - i, b - i]$에 Tor 진폭을 갖는다고 하자. 그러면
$K^\bullet$은 $[a, b]$에 Tor 진폭을 갖는다. 특히
$K^\bullet \in D^b(R)$이고 그 모든 코호몰로지 군이 유한 Tor 차원을
가지면 $K^\bullet$은 유한 Tor 차원을 갖는다.
\end{lemma}

\begin{proof}
유한 복합체의 길이에 대한 귀납으로 따른다. 보조정리
\ref{more-algebra-lemma-cone-tor-amplitude}와 표준 절단을 사용하라.
\end{proof}

\begin{lemma}
\label{more-algebra-lemma-push-tor-amplitude}
$A \to B$를 환 사상이라 하고 $K^\bullet$과 $L^\bullet$을
$B$-가군의 복합체라 하자. $a, b, c, d \in \mathbf{Z}$라 하자.
다음을 가정하자.
\begin{enumerate}
\item $K^\bullet$은 $B$-가군의 복합체로서 $[a, b]$에 Tor 진폭을 갖는다.
\item $L^\bullet$은 $A$-가군의 복합체로서 $[c, d]$에 Tor 진폭을 갖는다.
\end{enumerate}
그러면 $K^\bullet \otimes^\mathbf{L}_B L^\bullet$은 $A$-가군의
복합체로서 $[a + c, b + d]$에 Tor 진폭을 갖는다.
\end{lemma}

\begin{proof}
보조정리 \ref{more-algebra-lemma-tor-amplitude}에 의해 $K^\bullet$이 평탄
$B$-가군들의 복합체이고 $i \not \in [a, b]$이면 $K^i = 0$이라고
가정할 수 있다. $M$을 $A$-가군이라 하고 자유 분해
$F^\bullet \to M$을 택하자. 그러면
$$
(K^\bullet \otimes_B^\mathbf{L} L^\bullet) \otimes_A^{\mathbf{L}} M =
\text{Tot}(\text{Tot}(K^\bullet \otimes_B L^\bullet) \otimes_A F^\bullet) =
\text{Tot}(K^\bullet \otimes_B \text{Tot}(L^\bullet \otimes_A F^\bullet))
$$
이다. 두 번째 등식은 Homology, Remark
\ref{homology-remark-triple-complex}를 보라. 가정 (2)에 의해 복합체
$\text{Tot}(L^\bullet \otimes_A F^\bullet)$의 코호몰로지는
차수 $[c, d]$에서만 영이 아니다. 따라서 이중 복합체
$K^\bullet \otimes_B \text{Tot}(L^\bullet \otimes_A F^\bullet)$에
대한 Homology, Lemma \ref{homology-lemma-ss-double-complex}의
스펙트럼 열은
$(K^\bullet \otimes_B^\mathbf{L} L^\bullet) \otimes_A^{\mathbf{L}} M$의
코호몰로지가 차수 $[a + c, b + d]$에서만 영이 아님을 증명한다.
\end{proof}

\begin{lemma}
\label{more-algebra-lemma-flat-push-tor-amplitude}
$A \to B$를 환 사상이라 하고 $B$가 $A$-가군으로서 평탄이라고
가정하자. $K^\bullet$을 $B$-가군의 복합체라 하자.
$a, b \in \mathbf{Z}$라 하자. $K^\bullet$이 $B$-가군의 복합체로서
$[a, b]$에 Tor 진폭을 가지면, $K^\bullet$은 $A$-가군의 복합체로서도
$[a, b]$에 Tor 진폭을 갖는다.
\end{lemma}

\begin{proof}
이는 보조정리 \ref{more-algebra-lemma-push-tor-amplitude}의 특수한 경우이지만
다음과 같이 직접 볼 수도 있다. 등식
$K^\bullet \otimes_A^{\mathbf{L}} M =
K^\bullet \otimes_B^{\mathbf{L}} (M \otimes_A B)$
이 성립한다. 실제로 $K^\bullet$의 $B$-가군 복합체로서의 임의의
사영 분해는 $K^\bullet$의 $A$-가군 복합체로서의 평탄 분해이고,
$K^\bullet \otimes_A^{\mathbf{L}} M$을 계산하는 데 쓸 수 있다.
\end{proof}

\begin{lemma}
\label{more-algebra-lemma-finite-tor-dimension-push-tor-amplitude}
$A \to B$를 환 사상이라 하고 $B$가 $A$-가군으로서 Tor 차원
$\leq d$를 갖는다고 가정하자. $K^\bullet$을 $B$-가군의 복합체라
하자. $a, b \in \mathbf{Z}$라 하자. $K^\bullet$이 $B$-가군의
복합체로서 $[a, b]$에 Tor 진폭을 가지면, $K^\bullet$은
$A$-가군의 복합체로서
$[a - d, b]$에 Tor 진폭을 갖는다.
\end{lemma}

\begin{proof}
이는 보조정리 \ref{more-algebra-lemma-push-tor-amplitude}의 특수한 경우이지만
다음과 같이 직접 볼 수도 있다. $M$을 $A$-가군이라 하고 자유 분해
$F^\bullet \to M$을 택하자. 그러면
$$
K^\bullet \otimes_A^{\mathbf{L}} M =
\text{Tot}(K^\bullet \otimes_A F^\bullet) =
\text{Tot}(K^\bullet \otimes_B (F^\bullet \otimes_A B)) =
K^\bullet \otimes_B^{\mathbf{L}} (M \otimes_A^{\mathbf{L}} B).
$$
$B$를 $A$-가군으로 본 것에 대한 가정에 의해
$M \otimes_A^{\mathbf{L}} B$의 코호몰로지는 차수
$-d, -d + 1, \ldots, 0$에서만 영이 아니다. $K^\bullet$은
$[a, b]$에 Tor 진폭을 가지므로 Example \ref{more-algebra-example-tor}의
스펙트럼 열에서
$K^\bullet \otimes_B^{\mathbf{L}} (M \otimes_A^{\mathbf{L}} B)$의
코호몰로지가 원하는 대로 차수 $[-d + a, b]$에서만 영이 아님을
알 수 있다.
\end{proof}

\begin{lemma}
\label{more-algebra-lemma-pull-tor-amplitude}
$A \to B$를 환 사상이라 하고 $a, b \in \mathbf{Z}$라 하자.
$K^\bullet$을 $[a, b]$에 Tor 진폭을 갖는 $A$-가군의 복합체라
하자. 그러면 $K^\bullet \otimes_A^{\mathbf{L}} B$는 $B$-가군의
복합체로서 $[a, b]$에 Tor 진폭을 갖는다.
\end{lemma}

\begin{proof}
보조정리 \ref{more-algebra-lemma-tor-amplitude}에 의해 준동형사상
$E^\bullet \to K^\bullet$을 찾을 수 있는데, 여기서 $E^\bullet$은
평탄 $A$-가군들의 복합체이고 $i \not \in [a, b]$이면 $E^i = 0$이다.
그러면 구성에 의해 $E^\bullet \otimes_A B$는
$K^\bullet \otimes_A ^{\mathbf{L}} B$를 계산하고, 각
$E^i \otimes_A B$는 평탄 $B$-가군이다. Algebra, Lemma
\ref{algebra-lemma-flat-base-change}를 보라. 따라서 보조정리
\ref{more-algebra-lemma-tor-amplitude}로 결론을 얻는다.
\end{proof}

\begin{lemma}
\label{more-algebra-lemma-flat-base-change-finite-tor-dimension}
$A \to B$를 평탄 환 사상이라 하고 $d \geq 0$이라 하자.
$M$을 Tor 차원 $\leq d$를 갖는 $A$-가군이라 하자. 그러면
$M \otimes_A B$는 Tor 차원 $\leq d$를 갖는 $B$-가군이다.
\end{lemma}

\begin{proof}
$B$가 $A$ 위에서 평탄이므로
$M \otimes_A^{\mathbf{L}} B = M \otimes_A B$라는 사실과 보조정리
\ref{more-algebra-lemma-pull-tor-amplitude}에서 즉시 따른다.
\end{proof}
\begin{lemma}
\label{more-algebra-lemma-tor-amplitude-localization}
$A  \to B$를 환 사상이라 하고 $K^\bullet$을 $B$-가군의 복합체라 하자.
$a, b \in \mathbf{Z}$라 하자. 다음은 동치이다.
\begin{enumerate}
\item $K^\bullet$은 $A$-가군의 복합체로서 $[a, b]$에 Tor 진폭을 갖는다.
\item 모든 소 아이디얼 $\mathfrak q \subset B$에 대하여
$\mathfrak p = A \cap \mathfrak q$라 두면
$K^\bullet_\mathfrak q$는 $A_\mathfrak p$-가군의 복합체로서
$[a, b]$에 Tor 진폭을 갖는다.
\item 모든 극대 아이디얼 $\mathfrak m \subset B$에 대하여
$\mathfrak p = A \cap \mathfrak m$라 두면
$K^\bullet_\mathfrak m$은 $A_\mathfrak p$-가군의 복합체로서
$[a, b]$에 Tor 진폭을 갖는다.
\end{enumerate}
\end{lemma}

\begin{proof}
(3)을 가정하고 $M$을 $A$-가군이라 하자. 그러면
$H^i = H^i(K^\bullet \otimes_A^\mathbf{L} M)$은 $B$-가군이고
$(H^i)_\mathfrak m =
H^i(K^\bullet_\mathfrak m \otimes_{A_\mathfrak p}^\mathbf{L} M_\mathfrak p)$이다.
따라서 Algebra, Lemma
\ref{algebra-lemma-characterize-zero-local}에 의해
$i \not \in [a, b]$이면 $H^i = 0$이다. 이로써
(3) $\Rightarrow$ (1)을 얻는다.
(1) $\Rightarrow$ (2)와 (2) $\Rightarrow$ (3)의 증명은 생략한다.
\end{proof}

\begin{lemma}
\label{more-algebra-lemma-glue-tor-amplitude}
$R$을 환이라 하고 $f_1, \ldots, f_r \in R$을 단위 아이디얼을 생성하는
원소들이라 하자. $a, b \in \mathbf{Z}$라 하고 $K^\bullet$을
$R$-가군의 복합체라 하자. 각 $i$에 대하여 복합체
$K^\bullet \otimes_R R_{f_i}$가 $[a, b]$에 Tor 진폭을 가지면
$K^\bullet$도 $[a, b]$에 Tor 진폭을 갖는다.
\end{lemma}

\begin{proof}
이는 보조정리 \ref{more-algebra-lemma-tor-amplitude-localization}에서 즉시
따르지만 다음과 같이 직접 볼 수도 있다.
$- \otimes_R R_{f_i}$가 완전 함자이고, 따라서
$$
H^i(K^\bullet)_{f_i} =
H^i(K^\bullet) \otimes_R R_{f_i} = H^i(K^\bullet \otimes_R R_{f_i}).
$$
임에 유의하자. 마찬가지로 모든 $R$-가군 $M$에 대하여
$$
H^i(K^\bullet \otimes_R^{\mathbf{L}} M)_{f_i} =
H^i(K^\bullet \otimes_R^{\mathbf{L}} M) \otimes_R R_{f_i} =
H^i(K^\bullet \otimes_R R_{f_i} \otimes_{R_{f_i}}^{\mathbf{L}} M_{f_i}).
$$
이다. 따라서 $R$-가군 $N$이 영일 필요충분조건은 각 $i$에 대하여
$N_{f_i}$가 영인 것이라는 사실에서 결과가 따른다. Algebra, Lemma
\ref{algebra-lemma-cover}를 보라.
\end{proof}

\begin{lemma}
\label{more-algebra-lemma-flat-descent-tor-amplitude}
$R$을 환이라 하고 $a, b \in \mathbf{Z}$라 하자. $K^\bullet$을
$R$-가군의 복합체라 하고 $R \to R'$을 충실히 평탄한 환 사상이라
하자. 복합체 $K^\bullet \otimes_R R'$이 $[a, b]$에 Tor 진폭을
가지면 $K^\bullet$도 $[a, b]$에 Tor 진폭을 갖는다.
\end{lemma}

\begin{proof}
$M$을 $R$-가군이라 하자. $R \to R'$은 평탄하므로
% P02-E0941: frozen source has an unmatched opening parenthesis on the right-hand side; Korean preserves the displayed formal span.
$$
(M \otimes_R^{\mathbf{L}} K^\bullet) \otimes_R R'
=
((M \otimes_R R') \otimes_{R'}^{\mathbf{L}} (K^\bullet \otimes_R R')
$$
이고, 코호몰로지를 취하는 것은 $R'$과 텐서곱하는 것과 가환한다.
따라서 $\text{Tor}_i^R(M, K^\bullet) \otimes_R R' =
\text{Tor}_i^{R'}(M \otimes_R R', K^\bullet \otimes_R R')$이다.
$R \to R'$은 충실히 평탄하므로, $i \not \in [a, b]$일 때
$\text{Tor}_i^{R'}(M \otimes_R R', K^\bullet \otimes_R R')$가
소멸하면 $\text{Tor}_i^R(M, K^\bullet)$도 소멸한다.
\end{proof}

\begin{lemma}
\label{more-algebra-lemma-no-change-tor-amplitude}
환 사상 $R \to A \to B$가 주어지고 $A \to B$가 충실히 평탄하며
$K \in D(A)$라고 하자. $R$ 위에서 $K$의 Tor 진폭은 $R$ 위에서
$K \otimes_A^\mathbf{L} B$의 Tor 진폭과 같다.
\end{lemma}

\begin{proof}
이는 $R$-가군 $M$에 대하여 모든 $i$에서
$H^i(K \otimes_R^\mathbf{L} M) \otimes_A B =
H^i((K \otimes_A^\mathbf{L} B) \otimes_R^\mathbf{L} M)$이기 때문이다.
실제로 $K$를 $A$-가군의 복합체 $K^\bullet$로 나타내고 자유 분해
$F^\bullet \to M$을 택하자. 그러면 등식
$$
\text{Tot}(K^\bullet \otimes_A B \otimes_R F^\bullet) =
\text{Tot}(K^\bullet \otimes_R F^\bullet) \otimes_A B
$$
이 성립한다. 왼쪽의 코호몰로지 군은
$H^i((K \otimes_A^\mathbf{L} B) \otimes_R^\mathbf{L} M)$이고,
오른쪽에서는 $H^i(K \otimes_R^\mathbf{L} M) \otimes_A B$를 얻는다.
\end{proof}

\begin{lemma}
\label{more-algebra-lemma-finite-gl-dim-tor-dimension}
$R$을 전역 차원 $d$가 유한인 환이라 하자. 그러면
\begin{enumerate}
\item 모든 가군의 Tor 차원은 $\leq d$이다.
\item $i \in [a, b]$일 때만 $H^i(K^\bullet) \not = 0$인
$R$-가군의 복합체 $K^\bullet$은 $[a - d, b]$에 Tor 진폭을 갖는다.
\item $R$-가군의 복합체 $K^\bullet$이 유한 Tor 차원을 가질
필요충분조건은 $K^\bullet \in D^b(R)$인 것이다.
\end{enumerate}
\end{lemma}

\begin{proof}
$R$에 대한 가정은 모든 가군이 길이 $d$ 이하의 유한 사영 분해를
가짐을 뜻한다. 특히 모든 가군의 Tor 차원은 $\leq d$이다.
두 번째 명제는 보조정리 \ref{more-algebra-lemma-cohomology-tor-amplitude}와
정의에서 따른다. 세 번째 명제는 두 번째 명제를 바꾸어 말한 것이다.
\end{proof}

\begin{lemma}
\label{more-algebra-lemma-check-tor-dimension-modulo-nilpotent}
$R' \to R$을 핵이 멱영 아이디얼인 전사 환 사상이라 하자.
$K' \in D(R')$라 하고 $K = K' \otimes_{R'}^\mathbf{L} R$로 두자.
$a, b \in \mathbf{Z}$라 하자. 그러면 $K$가 $[a, b]$에 Tor 진폭을
가질 필요충분조건은 $K'$이 그러한 것이다.
\end{lemma}

\begin{proof}
한 방향은 보조정리 \ref{more-algebra-lemma-pull-tor-amplitude}에서 따른다.
다른 방향을 위해 $K$가 $[a, b]$에 Tor 진폭을 갖는다고 가정하고
$M'$을 $R'$-가군이라 하자.
$K' \otimes_{R'}^\mathbf{L} M'$의 코호몰로지가 구간 $[a, b]$에
포함되는 차수에서만 영이 아님을 보여야 한다.

\medskip\noindent
$I = \Ker(R' \to R)$로 두자. 어떤 $n$에 대하여 $I^n = 0$이다.
$IM' = 0$이면 $M'$을 $R$-가군으로 보고 다음과 같이 논증할 수 있다.
$$
K' \otimes_{R'}^\mathbf{L} M' =
K' \otimes_{R'}^\mathbf{L} (R \otimes_R^\mathbf{L} M') =
(K' \otimes_{R'}^\mathbf{L} R) \otimes_R^\mathbf{L} M' =
K \otimes_R^\mathbf{L} M'
$$
이는 $K$에 대한 가정에 의해 구간 $[a, b]$에서만 코호몰로지가 영이
아니다. $I^{t + 1}M' = 0$이면 짧은 완전열
$$
0 \to IM' \to M' \to M'/IM' \to 0
$$
을 생각한다. $t$에 대한 귀납 가정에 의해
$K' \otimes_{R'}^\mathbf{L} IM'$과
$K' \otimes_{R'}^\mathbf{L} M'/IM'$의 코호몰로지는 구간
$[a, b]$의 차수에서만 영이 아니다. 그러면 구별 삼각형
$$
K' \otimes_{R'}^\mathbf{L} IM' \to
K' \otimes_{R'}^\mathbf{L} M' \to
K' \otimes_{R'}^\mathbf{L} M'/IM' \to
(K' \otimes_{R'}^\mathbf{L} IM')[1]
$$
은 원하는 대로 $K' \otimes_{R'}^\mathbf{L} M'$에도 같은 사실이
성립함을 증명한다.
\end{proof}
\section{Ext의 스펙트럼 열}
\label{more-algebra-section-spectral-sequence-ext}

\noindent
이 절에서는 Ext 함자들을 생각할 때 나타나는 여러 스펙트럼 열을
모은다. 환 $R$의 유도 범주 $D(R)$의 임의의 두 대상 $L$, $K$에
대하여 Derived Categories, Section \ref{derived-section-ext}의
일반적 규약에 따라
$$
\Ext^n_R(L, K) = \Hom_{D(R)}(L, K[n])
$$
로 나타낸다.

\medskip\noindent
$M$이 $R$-가군이고 $K \in D^+(R)$이면 스펙트럼 열
\begin{equation}
\label{more-algebra-equation-first-ss-ext}
E_2^{i, j} = \Ext_R^i(M, H^j(K)) \Rightarrow \Ext_R^{i + j}(M, K)
\end{equation}
이 있다. 또한 $K$가 $R$-가군들의 아래로 유계인 복합체
$K^\bullet$로 표현되면 스펙트럼 열
\begin{equation}
\label{more-algebra-equation-second-ss-ext}
E_1^{i, j} = \Ext_R^j(M, K^i) \Rightarrow \Ext_R^{i + j}(M, K)
\end{equation}
이 있다. 이 스펙트럼 열들은 함자 $\Hom_R(M, -)$에
Derived Categories, Lemma
\ref{derived-lemma-two-ss-complex-functor}를 적용하여 얻는다.

\section{사영 차원}
\label{more-algebra-section-projective-dimension}

\noindent
가군의 사영 차원은 Algebra, Definition
\ref{algebra-definition-finite-proj-dim}에서 정의했다.

\begin{definition}
\label{more-algebra-definition-projective-dimension}
$R$을 환이라 하고 $K$를 $D(R)$의 대상이라 하자. $K$가 사영 가군들의
유계 복합체로 표현될 수 있으면 $K$가 {\it 유한 사영 차원}을 갖는다고 한다.
$K$가 복합체
$$
\ldots \to 0 \to P^a \to P^{a + 1} \to \ldots \to
P^{b - 1} \to P^b \to 0 \to \ldots
$$
와 준동형적이고 모든 $i \in \mathbf{Z}$에 대하여 $P^i$가 사영
$R$-가군이면 $K$가 {\it $[a, b]$에 사영 진폭을 갖는다}고 한다.
\end{definition}

\noindent
$K$가 유한 사영 차원을 가질 필요충분조건은 어떤
$a, b \in \mathbf{Z}$에 대하여 $K$가 $[a, b]$에 사영 진폭을 갖는
것임은 명백하다. 또한 $K$가 유한 사영 차원을 가지면 $K$는 유계이다.
다음 보조정리는 $D(R)$의 이러한 대상들을 판별한다.

\begin{lemma}
\label{more-algebra-lemma-projective-amplitude}
$R$을 환이라 하고 $K$를 $D(R)$의 대상이라 하자.
$a, b \in \mathbf{Z}$라 하자. 다음은 동치이다.
\begin{enumerate}
\item $K$는 $[a, b]$에 사영 진폭을 갖는다.
\item 모든 $R$-가군 $N$과 모든 $i \not \in [-b, -a]$에 대하여
$\Ext^i_R(K, N) = 0$이다.
\item $n > b$이면 $H^n(K) = 0$이고, 모든 $R$-가군 $N$과 모든
$i > -a$에 대하여 $\Ext^i_R(K, N) = 0$이다.
\item $n \not \in [a - 1, b]$이면 $H^n(K) = 0$이고, 모든
$R$-가군 $N$에 대하여 $\Ext^{-a + 1}_R(K, N) = 0$이다.
\end{enumerate}
\end{lemma}

\begin{proof}
(1)을 가정하자. $K$가 복합체
$$
\ldots \to 0 \to P^a \to P^{a + 1} \to \ldots \to
P^{b - 1} \to P^b \to 0 \to \ldots
$$
이고 모든 $i \in \mathbf{Z}$에 대하여 $P^i$가 사영 $R$-가군이라고
가정할 수 있다. 이 경우 복합체
$$
\ldots \to 0 \to \Hom_R(P^b, N) \to \ldots \to
\Hom_R(P^a, N) \to 0 \to \ldots
$$
로 Ext 군들을 계산할 수 있으므로 (2)를 얻는다.

\medskip\noindent
(2)가 성립한다고 가정하자. $I$가 단사 $R$-가군인 단사
$H^n(K) \to I$를 택하자. $\Hom_R(-, I)$는 완전 함자이므로
$\Ext^{-n}(K, I) = \Hom_R(H^n(K), I)$이다. 특히 $n > b$이면
$H^n(K)$가 영임을 알 수 있다. 따라서 (2)는 (3)을 함의한다.

\medskip\noindent
% P02-E0942: frozen source has the malformed phrase “By the same argument as in (2) implies (3) gives”; Korean renders the intended implication.
(2)가 (3)을 함의한다는 데 사용한 것과 같은 논증으로 (3)이 (4)를
함의함을 알 수 있다.

\medskip\noindent
(4)를 가정하자. (2)가 (3)을 함의한다는 데 사용한 것과 같은 논증으로
$H^{a - 1}(K) = 0$임을 알 수 있다. 즉,
$i \in [a, b]$가 아니면 $H^i(K) = 0$이다. 특히 $K$는 위로
유계이므로, 모든 $i \in \mathbf{Z}$에 대하여 $P^i$가 사영이고
$i > b$이면 $P^i = 0$인, $K$를 표현하는 복합체 $P^\bullet$을
% P02-E0943: frozen source repeats the article in “choose a a complex”; Korean renders the intended phrase.
택할 수 있다(예를 들어 자유로 택할 수 있다). Derived Categories,
Lemma \ref{derived-lemma-subcategory-left-resolution}를 보라.
$Q = \Coker(P^{a - 1} \to P^a)$로 두자. $i < a$이면
$H^i(K) = 0$이므로 $K$는 복합체
$$
\ldots \to 0 \to Q \to P^{a + 1} \to \ldots \to P^b \to 0 \to \ldots
$$
와 준동형적이다.
% P02-E0944: frozen source uses the malformed “Denote K' = (...) the corresponding object”; Korean renders the intended notation.
$K' = (P^{a + 1} \to \ldots \to P^b)$로 두고 이를 $D(R)$의 대응하는
대상으로 보자. $D(R)$의 구별 삼각형
$$
K' \to K \to Q[-a] \to K'[1]
$$
을 얻는다.
% P02-E0945: frozen source introduces the exact sequence without a finite verb; Korean supplies the intended verb.
따라서 모든 $R$-가군 $N$에 대하여 완전열
$$
\Ext^{-a}(K', N) \to \Ext^1(Q, N) \to \Ext^{1 - a}(K, N)
$$
이 있다. 가정에 의해 오른쪽 항은 소멸한다. 함의
(1) $\Rightarrow$ (2)에 의해 왼쪽 항도 소멸한다. 따라서 Algebra,
Lemma \ref{algebra-lemma-characterize-projective}에 의해 $Q$는 사영
$R$-가군이다. 그러므로 (1)이 성립하고 증명이 끝난다.
\end{proof}

\begin{example}
\label{more-algebra-example-ext-not-bounded}
$k$를 체라 하고 $R$을 $k$ 위의 쌍대수 환, 즉
$R = k[x]/(x^2)$라 하자.
% P02-E0946: frozen source uses the malformed “Denote epsilon in R the class of x”; Korean renders the intended notation.
$x$의 류를 $\epsilon \in R$로 나타내자. $M = R/(\epsilon)$이라 하자.
그러면 $M$은 복합체
$$
R \xrightarrow{\epsilon} R \xrightarrow{\epsilon} R \to \ldots
$$
와 준동형적이지만, $M$은 Algebra, Definition
\ref{algebra-definition-finite-proj-dim}에서 정의한 유한 사영 차원을
갖지 않는다. 이는 $D(R)$의 대상의 유한 사영 차원을 정의할 때 왜
양방향으로 유계인 사영 가군의 복합체를 생각하는지 설명한다.
\end{example}
\section{단사 차원}
\label{more-algebra-section-injective-dimension}

\noindent
이 절은 사영 차원에 관한 절의 쌍대이다.

\begin{definition}
\label{more-algebra-definition-injective-dimension}
$R$을 환이라 하고 $K$를 $D(R)$의 대상이라 하자.
$K$가 단사 $R$-가군들의 유한 복합체로 표현될 수 있으면 $K$가
{\it 유한 단사 차원}을 갖는다고 한다. $K$가 복합체
$$
\ldots \to 0 \to I^a \to I^{a + 1} \to \ldots \to
I^{b - 1} \to I^b \to 0 \to \ldots
$$
와 동형이고 모든 $i \in \mathbf{Z}$에 대하여 $I^i$가 단사
$R$-가군이면 $K$가 {\it $[a, b]$에 단사 진폭을 갖는다}고 한다.
\end{definition}

\noindent
% P02-E0947: frozen source switches from “finite” to “bounded” injective dimension after the definition; Korean preserves that terminology drift.
$K$가 유계 단사 차원을 가질 필요충분조건은 어떤
$a, b \in \mathbf{Z}$에 대하여 $K$가 $[a, b]$에 단사 진폭을 갖는
것임은 명백하다. 또한 $K$가 유계 단사 차원을 가지면 $K$는
유계이다. 다음은 당연히 필요한 보조정리이다.

\begin{lemma}
\label{more-algebra-lemma-injective-amplitude}
$R$을 환이라 하고 $K$를 $D(R)$의 대상이라 하자.
$a, b \in \mathbf{Z}$라 하자. 다음은 동치이다.
\begin{enumerate}
\item $K$는 $[a, b]$에 단사 진폭을 갖는다.
\item 모든 $R$-가군 $N$과 모든 $i \not \in [a, b]$에 대하여
$\Ext^i_R(N, K) = 0$이다.
\item 모든 아이디얼 $I \subset R$과 모든 $i \not \in [a, b]$에 대하여
% P02-E0948: frozen criterion omits the base-ring subscript on Ext; Korean preserves the formal span.
$\Ext^i(R/I, K) = 0$이다.
\end{enumerate}
\end{lemma}

\begin{proof}
(1)을 가정하자. $K$가 복합체
$$
\ldots \to 0 \to I^a \to I^{a + 1} \to \ldots \to
I^{b - 1} \to I^b \to 0 \to \ldots
$$
이고 모든 $i \in \mathbf{Z}$에 대하여 $I^i$가 단사
$R$-가군이라고 가정할 수 있다. 이 경우 복합체
$$
\ldots \to 0 \to \Hom_R(N, I^a) \to \ldots \to
\Hom_R(N, I^b) \to 0 \to \ldots
$$
로 Ext 군들을 계산할 수 있으므로 (2)를 얻는다.
(2)가 (3)을 함의함은 명백하다.

\medskip\noindent
(3)이 성립한다고 가정하자. 영이 아닌 사상 $R \to H^n(K)$를 택하자.
$\Hom_R(R, -)$은 완전 함자이므로
$\Ext^n_R(R, K) = \Hom_R(R, H^n(K)) = H^n(K)$이다.
따라서 $n \not \in [a, b]$이면 $H^n(K)$가 영임을 알 수 있다.
특히 $K$는 아래로 유계이므로 모든 $i \in \mathbf{Z}$에 대하여
$I^i$가 단사이고 $i < a$이면 $I^i = 0$인 준동형사상
$$
K \to I^\bullet
$$
을 택할 수 있다. Derived Categories, Lemma
\ref{derived-lemma-subcategory-right-resolution}를 보라.
$J = \Ker(I^b \to I^{b + 1})$로 두자. 그러면 $K$는 복합체
$$
\ldots \to 0 \to I^a \to \ldots \to I^{b - 1} \to J \to 0 \to \ldots
$$
와 준동형적이다.
% P02-E0955: frozen source uses the malformed “Denote K' = (...) the corresponding object”; Korean renders the intended notation.
$K' = (I^a \to \ldots \to I^{b - 1})$로 두고 이를 $D(R)$의 대응하는
대상으로 보자. $D(R)$의 구별 삼각형
$$
J[-b] \to K \to K' \to J[1 - b]
$$
을 얻는다.
% P02-E0956: frozen source introduces the exact sequence without a finite verb; Korean supplies the intended verb.
따라서 모든 아이디얼 $I \subset R$에 대하여 완전열
$$
\Ext^b(R/I, K') \to \Ext^1(R/I, J) \to \Ext^{1 + b}(R/I, K)
$$
이 있다. 가정에 의해 오른쪽 항은 소멸한다. 함의
(1) $\Rightarrow$ (2)에 의해 왼쪽 항도 소멸한다. 따라서 보조정리
\ref{more-algebra-lemma-characterize-injective-bis}에 의해 $J$는 단사
$R$-가군이다.
\end{proof}

\begin{example}
\label{more-algebra-example-finite-injective-finite-global-dimension}
$R$을 Dedekind 정역이라 하자. 그러면 영이 아닌 모든 아이디얼 $I$는
유한 사영 가군이다. 보조정리
\ref{more-algebra-lemma-dedekind-torsion-free-flat}을 보라. 따라서 $R/I$의
사영 차원은 $1$이다. 그러므로 보조정리
\ref{more-algebra-lemma-injective-amplitude}에 의해 모든 $R$-가군 $M$의 단사
차원은 $\leq 1$이다. 따라서 임의의 두 $R$-가군 $M, N$과
$i \geq 2$에 대하여 $\Ext^i_R(M, N) = 0$이다. 이로부터
$D^b(R)$의 임의의 대상 $K$는 그 코호몰로지들의 직합과 동형임이
따른다. 즉 $K \cong \bigoplus H^i(K)[-i]$이다. Derived Categories,
Lemma \ref{derived-lemma-ext-2-zero}를 보라.
\end{example}

\begin{example}
\label{more-algebra-example-ext-not-bounded-reversed}
$k$를 체라 하고 $R$을 $k$ 위의 쌍대수 환, 즉
$R = k[x]/(x^2)$라 하자.
% P02-E0957: frozen source repeats the malformed “Denote epsilon in R the class of x”; Korean renders the intended notation.
$x$의 류를 $\epsilon \in R$로 나타내자. $M = R/(\epsilon)$이라 하자.
그러면 $M$은 복합체
$$
\ldots \to R \xrightarrow{\epsilon} R \xrightarrow{\epsilon} R
$$
와 준동형적이고 $R$은 단사 $R$-가군이다. 그러나 이 상황에서는
보통 $M$이 유한 단사 차원을 갖는다고 보지 않는다.
% P02-E0947: frozen source again says “bounded injective dimension” rather than the defined “finite” term.
이는 $D(R)$의 대상의 유계 단사 차원을 정의할 때 왜 양방향으로
유계인 단사 가군들의 복합체를 생각하는지 설명한다.
\end{example}

\begin{lemma}
\label{more-algebra-lemma-finite-injective-dimension}
$R$을 환이라 하고 $K \in D(R)$라 하자.
\begin{enumerate}
\item $K$가 $D^b(R)$에 속하고 모든 $i$에 대하여 $H^i(K)$가
유한 단사 차원을 가지면 $K$도 유한 단사 차원을 갖는다.
\item $K^\bullet$이 $K$를 표현하는 $R$-가군들의 유계 복합체이고
모든 $i$에 대하여 $K^i$가 유한 단사 차원을 가지면 $K$도 유한
단사 차원을 갖는다.
\end{enumerate}
\end{lemma}

\begin{proof}
생략한다. 힌트: 함자 $F = \Hom_R(N, -)$에 Derived Categories,
Lemma \ref{derived-lemma-two-ss-complex-functor}의 스펙트럼 열들을
적용하여
% P02-E0949: frozen proof writes Ext_A although the lemma is over R and A is undefined; Korean preserves the formal span.
$\Ext^i_A(N, K)$를 계산하고, 보조정리
\ref{more-algebra-lemma-injective-amplitude}의 판정법을 사용하라.
\end{proof}

\begin{lemma}
\label{more-algebra-lemma-finite-injective-dimension-Noetherian-radical}
$R$을 Noether 환이라 하고 $I \subset R$을 $R$의 Jacobson 근기에
포함되는 아이디얼이라 하자. $K \in D^+(R)$가 유한 코호몰로지
가군들을 갖는다고 하자. 다음은 동치이다.
\begin{enumerate}
\item $K$는 유한 단사 차원을 갖는다.
\item 어떤 $b$가 존재하여 임의의 아이디얼 $J \supset I$와
$i > b$에 대하여 $\Ext^i_R(R/J, K) = 0$이다.
\end{enumerate}
\end{lemma}

\begin{proof}
함의 (1) $\Rightarrow$ (2)는 즉시 따른다. (2)를 가정하자.
$i < a$이면 $H^i(K) = 0$이라고 하자. 그러면 모든 $R$-가군 $M$과
$i < a$에 대하여 $\Ext^i(M, K) = 0$이다. 따라서 보조정리
\ref{more-algebra-lemma-injective-amplitude}에 의해, 모든 유한 $R$-가군 $M$과
% P02-E0950: frozen source omits a conjunction in “for i > b any finite R-module M”; Korean renders the intended quantification.
$i > b$에 대하여 $\Ext^i(M, K) = 0$임을 보이면 충분하다.
Algebra, Lemma \ref{algebra-lemma-filter-Noetherian-module}에 의해
$M$은 연속 몫들이 $R/\mathfrak p$ 꼴인 유한 여과를 가지며, 여기서
$\mathfrak p$는 소 아이디얼이다. $0 \to M_1 \to M \to M_2 \to 0$이
짧은 완전열이고 $j = 1, 2$ 및 $i > b$에 대하여
$\Ext^i(M_j, K) = 0$이면, $i > b$에 대하여
$\Ext^i(M, K) = 0$이다. 따라서 $M = R/\mathfrak p$라고 가정할 수 있다.
$I \subset \mathfrak p$이면 소멸은 가정에서 따른다. 그렇지 않으면
$f \in I$, $f \not \in \mathfrak p$를 택하고 짧은 완전열
$$
0 \to R/\mathfrak p \xrightarrow{f} R/\mathfrak p \to R/(\mathfrak p, f) \to 0
$$
을 생각하자. $R$-가군 $R/(\mathfrak p, f)$는 연속 몫들이
$(\mathfrak p, f) \subset \mathfrak q$인 $R/\mathfrak q$ 꼴인 여과를
갖는다. 따라서 Noether 귀납법과 위 논증에 의해
$R/(\mathfrak p, f)$에 대해 소멸이 성립한다고 가정할 수 있다.
한편 가정에 의해 $K$는 아래로 유계이고 유한 코호몰로지 가군들을
가지므로, 스펙트럼 열 (\ref{more-algebra-equation-first-ss-ext})과 Algebra, Lemma
\ref{algebra-lemma-ext-noetherian}에 의해 가군
$E^i = \Ext^i(R/\mathfrak p, K)$는 유한이다. 따라서 $i > b$이면
$E^i$는 $E^i/fE^i = 0$인 유한 $R$-가군이다. Nakayama 보조정리
(Algebra, Lemma \ref{algebra-lemma-NAK})에 의해 $E^i$가 영임을
알 수 있다.
\end{proof}

\begin{lemma}
\label{more-algebra-lemma-finite-injective-dimension-Noetherian-local}
$(R, \mathfrak m, \kappa)$를 Noether 국소환이라 하자.
$K \in D^+(R)$가 유한 코호몰로지 가군들을 갖는다고 하자.
다음은 동치이다.
\begin{enumerate}
\item $K$는 유한 단사 차원을 갖는다.
\item $i \gg 0$이면 $\Ext^i_R(\kappa, K) = 0$이다.
\end{enumerate}
\end{lemma}

\begin{proof}
이는 보조정리
\ref{more-algebra-lemma-finite-injective-dimension-Noetherian-radical}의 특수한
경우이다.
\end{proof}
\section{사영에 가까운 가군들}
\label{more-algebra-section-near-projective}

\noindent
% P02-E0951: frozen source says “many different of definitions”; Korean renders the intended phrase.
문헌에는 “거의 사영인 가군”에 대한 서로 다른 정의가 많이 있는 듯하다.
이 절에서는 여러 가능성 중 하나만 논의한다.

\begin{lemma}
\label{more-algebra-lemma-ideal-factor-through-projective}
$R$을 환이라 하고 $M$, $N$을 $R$-가군이라 하자.
\begin{enumerate}
\item $R$-가군 사상 $\varphi : M \to N$이 주어졌을 때 다음은
동치이다. (a) $\varphi$는 사영 $R$-가군을 통하여 분해되고,
(b) $\varphi$는 자유 $R$-가군을 통하여 분해된다.
\item (1)의 동치 조건들을 만족하는 $\varphi : M \to N$들의 집합은
$\Hom_R(M, N)$의 $R$-부분가군이다.
\item 사상 $\psi : M' \to M$과 $\xi : N \to N'$이 주어졌을 때,
$\varphi : M \to N$이 (1)의 동치 조건들을 만족하면
$\xi \circ \varphi \circ \psi : M' \to N'$도 그러하다.
\end{enumerate}
\end{lemma}

\begin{proof}
(1)(a)와 (1)(b)의 동치는 Algebra, Lemma
\ref{algebra-lemma-characterize-projective}에서 따른다.
$\varphi : M \to N$과 $\varphi' : M \to N$이 각각 가군 $P$와
$P'$을 통하여 분해되면 $\varphi + \varphi'$는 $P \oplus P'$을
통하여 분해되고, 모든 $\lambda \in R$에 대하여 $\lambda \varphi$는
$P$를 통하여 분해된다. 이것이 (2)를 증명한다.
$\varphi : M \to N$이 가군 $P$를 통하여 분해되고 $\psi$와 $\xi$가
(3)에서와 같으면 $\xi \circ \varphi \circ \psi$는 $P$를 통하여
분해된다. 이것이 (3)을 증명한다.
\end{proof}

\begin{lemma}
\label{more-algebra-lemma-factor-through-finite-projective}
$R$을 환이라 하고 $\varphi : M \to N$을 $R$-가군 사상이라 하자.
$\varphi$가 사영 가군을 통하여 분해되고 $M$이 유한 $R$-가군이면,
$\varphi$는 유한 사영 가군을 통하여 분해된다.
\end{lemma}

\begin{proof}
보조정리 \ref{more-algebra-lemma-ideal-factor-through-projective}에 의해
$\varphi = \tau \circ \sigma$로 분해할 수 있는데, 여기서
$\sigma$의 공역은 어떤 집합 $I$에 대한 $\bigoplus_{i \in I} R$이다.
$M$의 생성원 $x_1, \ldots, x_n$을 택하고
$\sigma(x_j) = (a_{ji})_{i \in I}$로 쓰자.
% P02-E0952: frozen source changes a_{ji} to a_{ij} in the next sentence; Korean preserves both formal spans.
각 $j$에 대하여 영이 아닌 $a_{ij}$는 유한 개뿐이다.
따라서 $\sigma$의 상은 유한 자유 $R$-가군에 포함되며 결론을 얻는다.
\end{proof}

\noindent
$R$을 환이라 하자.
% P02-E0953: frozen source tests projectivity of an arbitrary R-module using the identity on R rather than on that module; Korean preserves R.
$R$ 위의 항등사상이 사영 가군을 통하여 분해될 필요충분조건은
$R$-가군이 사영인 것임을 관찰하자.

\begin{lemma}
\label{more-algebra-lemma-near-projective}
$R$을 환이라 하고 $I \subset R$을 아이디얼이라 하자.
$M$을 $R$-가군이라 하자. 다음 조건들은 동치이다.
\begin{enumerate}
\item 모든 $a \in I$에 대하여 사상 $a : M \to M$은 사영
$R$-가군을 통하여 분해된다.
\item 모든 $a \in I$에 대하여 사상 $a : M \to M$은 자유
$R$-가군을 통하여 분해된다.
\item 모든 $R$-가군 $N$에 대하여 $\Ext^1_R(M, N)$은 $I$에 의해
소멸된다.
\end{enumerate}
\end{lemma}

\begin{proof}
(1)과 (2)의 동치는 보조정리
\ref{more-algebra-lemma-ideal-factor-through-projective}에서 따른다.
(1)이 성립하면 (3)이 성립한다. 실제로 임의의 $N$과 임의의 사영
가군 $P$에 대하여
% P02-E0954: frozen source omits the predicate after Ext^1_R(P,N); Korean supplies the intended vanishing.
$\Ext^1_R(P, N)$은 영이다. 반대로 (3)이 성립한다고 가정하자.
$P$가 사영(나아가 자유)인 짧은 완전열
$0 \to N \to P \to M \to 0$을 택하자. 가정에 의해
$\Ext^1_R(M, N)$에서 이 완전열에 대응하는 원소는 $I$에 의해
소멸된다. 따라서 모든 $a \in I$에 대하여 사상 $a : M \to M$은
전사 $P \to M$을 통하여 분해될 수 있으므로 (1)을 얻는다.
\end{proof}

\noindent
보조정리 \ref{more-algebra-lemma-near-projective}의 동치 조건들을 만족하는
가군을 편하게 논의하기 위해 이 성질에 이름을 붙인다.

\begin{definition}
\label{more-algebra-definition-near-projective}
$R$을 환이라 하고 $I \subset R$을 아이디얼이라 하자.
$M$을 $R$-가군이라 하자. 보조정리 \ref{more-algebra-lemma-near-projective}의
동치 조건들이 성립하면 $M$을
{\it $I$-사영}\footnote{이는 표준적이지 않은 표기이다.}이라고 한다.
\end{definition}

\noindent
$I$에 의해 소멸되는 가군은 $I$-사영이다.

\begin{lemma}
\label{more-algebra-lemma-torsion-near-projective}
$R$을 환이라 하고 $I \subset R$을 아이디얼이라 하자.
$M$을 $R$-가군이라 하자. $M$이 $I$에 의해 소멸되면 $M$은
$I$-사영이다.
\end{lemma}

\begin{proof}
정의와 영가군이 사영이라는 사실에서 즉시 따른다.
\end{proof}

\begin{lemma}
\label{more-algebra-lemma-ses-near-projective}
$R$을 환이라 하고 $I \subset R$을 아이디얼이라 하자.
$R$-가군들의 짧은 완전열
$$
0 \to K \to P \to M \to 0
$$
을 생각하자. $M$이 $I$-사영이고 $P$가 사영이면 $K$는 $I$-사영이다.
\end{lemma}

\begin{proof}
원소 $\text{id}_K \in \Hom_R(K, K)$는 $\Ext^1_R(M, K)$에서 주어진
확장의 류로 간다. 가정에 의해 이 류는 임의의 $a \in I$에 의해
소멸되므로 사상 $a : K \to K$는 $K \to P$를 통하여 분해된다.
이로써 결론을 얻는다.
\end{proof}

\begin{lemma}
\label{more-algebra-lemma-dual-near-projective}
$R$을 환이라 하고 $I \subset R$을 아이디얼이라 하자.
$M$이 유한 $I$-사영 $R$-가군이면
$M^\vee = \Hom_R(M, R)$도 $I$-사영이다.
\end{lemma}

\begin{proof}
$M$이 유한이고 $I$-사영이라고 가정하자. 짧은 완전열
$0 \to K \to R^{\oplus r} \to M \to 0$을 택하자. 이는 단사
$M^\vee \to R^{\oplus r} = (R^{\oplus r})^\vee$를 준다.
이 짧은 완전열에 대응하는 $\Ext^1_R(M, K)$의 확장류는 $I$에 의해
소멸되므로, 임의의 $a \in I$에 대하여 사상
$M \to R^{\oplus r}$을 찾아 주어진 사상 $R^{\oplus r} \to M$과의
합성이 $a : M \to M$이 되게 할 수 있다. 쌍대를 취하면
$a : M^\vee \to M^\vee$가 위의 사상
$M^\vee \to R^{\oplus r}$을 통하여 분해됨을 알 수 있고 결론을
얻는다.
\end{proof}
\section{Hom 복합체}
\label{more-algebra-section-hom-complexes}

\noindent
$R$을 환이라 하자. $L^\bullet$과 $M^\bullet$을 두 $R$-가군 복합체라
하자. 복합체 $\Hom^\bullet(L^\bullet, M^\bullet)$을 구성하자.
구체적으로 각 $n$에 대하여
$$
\Hom^n(L^\bullet, M^\bullet) =
\prod\nolimits_{n = p + q} \Hom_R(L^{-q}, M^p)
$$
로 둔다. $\Hom^n$을, 등급 가군으로 본 $L^\bullet$에서
$M^\bullet$로 가는 차수 $n$의 모든 동차 $R$-선형 사상들의
$R$-가군으로 생각하면 편리하다. 이 용어를 사용하여 미분을
$$
\text{d}(f) = \text{d}_M \circ f - (-1)^n f \circ \text{d}_L
$$
이라는 규칙으로 정의한다. 여기서
$f \in \Hom^n(L^\bullet, M^\bullet)$이다.
$\text{d}^2 = 0$의 확인은 생략한다. 부호 규칙은
절 \ref{more-algebra-section-sign-rules}을 보라. 이 구성은
미분 등급 대수, 예 \ref{dga-example-category-complexes}의 특수한
경우이다. 구성에서 곧바로 모든 $n \in \mathbf{Z}$에 대하여
\begin{equation}
\label{more-algebra-equation-cohomology-hom-complex}
H^n(\Hom^\bullet(L^\bullet, M^\bullet)) = \Hom_{K(R)}(L^\bullet, M^\bullet[n])
\end{equation}
임을 얻는다.

\begin{lemma}
\label{more-algebra-lemma-compose}
$R$을 환이라 하자. $R$-가군들의 복합체
$K^\bullet, L^\bullet, M^\bullet$이 주어지면 $R$-가군 복합체들의
표준 동형사상
$$
\Hom^\bullet(K^\bullet, \Hom^\bullet(L^\bullet, M^\bullet))
=
\Hom^\bullet(\text{Tot}(K^\bullet \otimes_R L^\bullet), M^\bullet)
$$
이 존재한다.
\end{lemma}

\begin{proof}
왼쪽 변의 차수 $n$인 원소를 $\alpha$라 하자. 그러면
$$
\alpha = (\alpha^{p, q}) \in
\prod\nolimits_{p + q = n} \Hom_R(K^{-q}, \Hom^p(L^\bullet, M^\bullet))
$$
이다. 각 $\alpha^{p, q}$는
$$
\alpha^{p, q} = (\alpha^{r, s, q}) \in
\prod\nolimits_{r + s + q = n} \Hom_R(K^{-q}, \Hom_R(L^{-s}, M^r))
$$
의 원소이다. 식별
\begin{equation}
\label{more-algebra-equation-identification}
\Hom_R(K^{-q}, \Hom_R(L^{-s}, M^r)) = \Hom_R(K^{-q} \otimes_R L^{-s}, M^r)
\end{equation}
을 하면 우리의 부호 규칙에 의해
\begin{align*}
\text{d}(\alpha^{r, s, q})
& =
\text{d}_{\Hom^\bullet(L^\bullet, M^\bullet)} \circ \alpha^{r, s, q}
- (-1)^n \alpha^{r, s, q} \circ \text{d}_K \\
& =
\text{d}_M \circ \alpha^{r, s, q}
- (-1)^{r + s} \alpha^{r, s, q} \circ \text{d}_L
- (-1)^{r + s + q} \alpha^{r, s, q} \circ \text{d}_K
\end{align*}
를 얻는다. 한편 오른쪽 변의 차수 $n$인 원소를 $\beta$라 하면
$$
\beta = (\beta^{r, s, q}) \in \prod\nolimits_{r + s + q = n}
\Hom_R(K^{-q} \otimes_R L^{-s}, M^r)
$$
이고, 우리의 부호 규칙(호몰로지, 정의
\ref{homology-definition-associated-simple-complex})에 의해
\begin{align*}
\text{d}(\beta^{r, s, q})
& =
\text{d}_M \circ \beta^{r, s, q}
- (-1)^n \beta^{r, s, q} \circ
\text{d}_{\text{Tot}(K^\bullet \otimes L^\bullet)} \\
& =
\text{d}_M \circ \beta^{r, s, q}
- (-1)^{r + s + q} \left(
\beta^{r, s, q} \circ \text{d}_K + (-1)^{-q} \beta^{r, s, q} \circ \text{d}_L
\right)
\end{align*}
를 얻는다. 따라서 식별 (\ref{more-algebra-equation-identification})에 의해 유도되는
사상이 실제로 복합체의 사상임을 알 수 있다.
\end{proof}

\begin{remark}
\label{more-algebra-remark-internal-hom}
$R$을 환이라 하자. $R$-가군 복합체의 범주 $\text{Comp}(R)$는
$\text{Tot}(- \otimes_R -)$로 주어지는 텐서곱을 갖는 대칭 모노이드
범주이다. 보조정리 \ref{more-algebra-lemma-symmetric-monoidal-cat-complexes}를 보라.
$\text{Comp}(R)$의 $L^\bullet$과 $M^\bullet$이 주어졌을 때,
원소 $f \in \Hom^0(L^\bullet, M^\bullet)$가 복합체의 사상
$f : L^\bullet \to M^\bullet$을 정의할 필요충분조건은
$\text{d}(f) = 0$이다. 따라서 보조정리 \ref{more-algebra-lemma-compose}는 또한
$\text{Comp}(R)$의 $K^\bullet, L^\bullet, M^\bullet$에 관하여 함자적으로
$$
\Mor_{\text{Comp}(R)}(K^\bullet, \Hom^\bullet(L^\bullet, M^\bullet))
=
\Mor_{\text{Comp}(R)}(\text{Tot}(K^\bullet \otimes_R L^\bullet), M^\bullet)
$$
임을 말해 준다. 이는 범주론, 주석
\ref{categories-remark-internal-hom-monoidal}에서 논의한 바와 같이
$\Hom^\bullet( - , -)$이 대칭 모노이드 범주 $\text{Comp}(R)$의 내부
Hom임을 뜻한다.
\end{remark}

\begin{lemma}
\label{more-algebra-lemma-composition}
$R$을 환이라 하자. $R$-가군들의 복합체
$K^\bullet, L^\bullet, M^\bullet$이 주어지면 $R$-가군 복합체의 표준
사상
$$
\text{Tot}\left(
\Hom^\bullet(L^\bullet, M^\bullet) \otimes_R \Hom^\bullet(K^\bullet, L^\bullet)
\right)
\longrightarrow
\Hom^\bullet(K^\bullet, M^\bullet)
$$
이 존재한다.
\end{lemma}

\begin{proof}
주석 \ref{more-algebra-remark-internal-hom}의 논의를 통하면 그러한 표준 사상의
존재는 범주론, 주석 \ref{categories-remark-internal-hom-monoidal}에서
따른다. 직접적인 구성도 주겠다.

\medskip\noindent
왼쪽 변의 차수 $n$인 원소 $\alpha$는
$$
\alpha = (\alpha^{p, q}) \in \bigoplus\nolimits_{p + q = n}
\Hom^p(L^\bullet, M^\bullet) \otimes_R \Hom^q(K^\bullet, L^\bullet)
$$
이다. 원소 $\alpha^{p, q}$는
$\alpha^{p, q} = \sum \beta^p_i \otimes \gamma^q_i$ 꼴의 유한합이고,
여기서
$$
\beta^p_i = (\beta^{r, s}_i)
\in \prod\nolimits_{r + s = p} \Hom_R(L^{-s}, M^r)
$$
이며
$$
\gamma^q_i = (\gamma^{u, v}_i)
\in \prod\nolimits_{u + v = q} \Hom_R(K^{-v}, L^u)
$$
이다. 이 사상은 $\alpha$를 $\delta = (\delta^{r, v})$로 보내며,
$$
\delta^{r, v} =
\sum\nolimits_{i, s}  \beta^{r, s}_i
\circ \gamma^{-s, v}_i
\in \Hom_R(K^{-v}, M^r)
$$
로 주어진다. $r + v = n$이 주어졌을 때 영이 아닌
$\alpha^{p, q}$는 유한 개뿐이고, 따라서 영이 아닌 $\beta^p_i$와
$\gamma^q_i$도 유한 개뿐이므로 이 합은 유한하다. 우리의 부호
규칙에 의해
\begin{align*}
\text{d}(\alpha^{p, q})
& =
\text{d}_{\Hom^\bullet(L^\bullet, M^\bullet)}(\alpha^{p, q})
+ (-1)^p \text{d}_{\Hom^\bullet(K^\bullet, L^\bullet)}(\alpha^{p, q}) \\
& =
\sum \Big( \text{d}_M \circ \beta^p_i \circ \gamma^q_i
- (-1)^p \beta^p_i \circ \text{d}_L \circ \gamma^q_i \Big) \\
&
\quad + (-1)^p \sum \Big( \beta^p_i \circ \text{d}_L \circ \gamma^q_i
- (-1)^q \beta^p_i \circ \gamma^q_i \circ \text{d}_K \Big) \\
& =
\sum \Big( \text{d}_M \circ \beta^p_i \circ \gamma^q_i 
-(-1)^n \beta^p_i \circ \gamma^q_i \circ \text{d}_K \Big)
\end{align*}
를 얻는다.
% P02-E0958: frozen source says “the rules ... is compatible”; Korean renders the intended singular assignment rule.
따라서 규칙 $\alpha \mapsto \delta$는 미분과 양립하며, 보조정리가
증명되었다.
\end{proof}

\begin{lemma}
\label{more-algebra-lemma-diagonal-better}
$R$을 환이라 하자. $R$-가군들의 복합체
$K^\bullet, L^\bullet, M^\bullet$이 주어지면 세 복합체 모두에 관하여
함자적인 $R$-가군 복합체의 표준 사상
$$
\text{Tot}(K^\bullet \otimes_R \Hom^\bullet(M^\bullet, L^\bullet))
\longrightarrow
\Hom^\bullet(M^\bullet, \text{Tot}(K^\bullet \otimes_R L^\bullet))
$$
이 존재한다.
\end{lemma}

\begin{proof}
주석 \ref{more-algebra-remark-internal-hom}의 논의를 통하면 그러한 표준 사상의
존재는 범주론, 주석 \ref{categories-remark-internal-hom-monoidal}에서
따른다. 직접적인 구성도 주겠다.

\medskip\noindent
오른쪽 변의 차수 $n$인 원소를 $\alpha$라 하자. 그러면
$$
\alpha = (\alpha^{p, q}) \in \prod\nolimits_{p + q = n}
\Hom_R(M^{-q}, \text{Tot}^p(K^\bullet \otimes_R L^\bullet))
$$
이다. 각 $\alpha^{p, q}$는
$$
\alpha^{p, q} = (\alpha^{r, s, q}) \in
\Hom_R(M^{-q}, \bigoplus\nolimits_{r + s + q = n} K^r \otimes_R L^s)
$$
의 원소이다. 여기서 $\alpha^{r, s, q}$를, 모든
$x \in M^{-q}$에 대하여 영이 아닌 $\alpha^{r, s, q}(x)$가 유한 개뿐인
사상들의 족으로 생각한다. 우리의 부호 규칙에 의해
\begin{align*}
\text{d}(\alpha^{r, s, q})
& =
\text{d}_{\text{Tot}(K^\bullet \otimes_R L^\bullet)} \circ \alpha^{r, s, q}
- (-1)^n \alpha^{r, s, q} \circ \text{d}_M \\
& =
\text{d}_K \circ \alpha^{r, s, q} + (-1)^r \text{d}_L \circ \alpha^{r, s, q}
- (-1)^n \alpha^{r, s, q} \circ \text{d}_M
\end{align*}
를 얻는다. 한편 왼쪽 변의 차수 $n$인 원소를 $\beta$라 하면
$$
\beta = (\beta^{p, q}) \in
\bigoplus\nolimits_{p + q = n} K^p \otimes_R \Hom^q(M^\bullet, L^\bullet)
$$
이고, $\beta^{p, q} = \sum \gamma_i^p \otimes \delta_i^q$로 쓸 수 있다.
여기서 $\gamma_i^p \in K^p$이고
$$
\delta_i^q = (\delta_i^{r, s}) \in
\prod\nolimits_{r + s = q} \Hom_R(M^{-s}, L^r)
$$
이다. 우리의 부호 규칙에 의해
\begin{align*}
\text{d}(\beta^{p, q})
& =
\text{d}_K(\beta^{p, q}) +
(-1)^p \text{d}_{\Hom^\bullet(M^\bullet, L^\bullet)}(\beta^{p, q}) \\
& =
\sum \text{d}_K(\gamma_i^p) \otimes \delta_i^q +
(-1)^p \sum \gamma_i^p \otimes
(\text{d}_L \circ \delta_i^q - (-1)^q \delta_i^q \circ \text{d}_M)
\end{align*}
를 얻는다. 원소 $\beta$를 다음 $\alpha$로 보낸다.
$$
\alpha^{r, s, q} = c^{r, s, q}(\sum \gamma_i^r \otimes \delta_i^{s, q})
$$
여기서 $c^{r, s, q} : K^r \otimes_R \Hom_R(M^{-q}, L^s) \to
\Hom_R(M^{-q}, K^r \otimes_R L^s)$는 표준 사상이다. 주어진
$\beta$와 $r$에 대하여 영이 아닌 $\gamma_i^r$는 유한 개뿐이다.
% P02-E0959: frozen source redundantly says “only finitely many nonzero alpha are nonzero”; Korean states the single finiteness condition.
따라서 (주어진 $r$에 대하여) 영이 아닌 $\alpha^{r, s, q}$도 유한
개뿐이다. 그러므로 이 사상들의 족은 위의 조건들을 만족하고 사상은
잘 정의된다. 부호를 비교하면 이것이 미분과 양립함을 알 수 있다.
\end{proof}

\begin{lemma}
\label{more-algebra-lemma-diagonal}
$R$을 환이라 하자. $R$-가군들의 복합체 $K^\bullet, L^\bullet$이
주어지면 두 복합체에 관하여 함자적인 $R$-가군 복합체의 표준 사상
$$
K^\bullet
\longrightarrow
\Hom^\bullet(L^\bullet, \text{Tot}(K^\bullet \otimes_R L^\bullet))
$$
이 존재한다.
\end{lemma}

\begin{proof}
주석 \ref{more-algebra-remark-internal-hom}의 논의를 통하면 그러한 표준 사상의
존재는 범주론, 주석 \ref{categories-remark-internal-hom-monoidal}에서
따른다. 직접적인 구성도 주겠다.

\medskip\noindent
오른쪽 변의 차수 $n$인 원소를 $\alpha$라 하자. 그러면
$$
\alpha = (\alpha^{p, q}) \in \prod\nolimits_{p + q = n}
\Hom_R(L^{-q}, \text{Tot}^p(K^\bullet \otimes_R L^\bullet))
$$
이다. 각 $\alpha^{p, q}$는
$$
\alpha^{p, q} = (\alpha^{r, s, q}) \in
\Hom_R(L^{-q}, \bigoplus\nolimits_{r + s + q = n} K^r \otimes_R L^s)
$$
의 원소이다. 여기서 $\alpha^{r, s, q}$를, 모든
$x \in L^{-q}$에 대하여 영이 아닌 $\alpha^{r, s, q}(x)$가 유한 개뿐인
사상들의 족으로 생각한다. 우리의 부호 규칙에 의해
\begin{align*}
\text{d}(\alpha^{r, s, q})
& =
\text{d}_{\text{Tot}(K^\bullet \otimes_R L^\bullet)} \circ \alpha^{r, s, q}
- (-1)^n \alpha^{r, s, q} \circ \text{d}_L \\
& =
\text{d}_K \circ \alpha^{r, s, q} + (-1)^r \text{d}_L \circ \alpha^{r, s, q}
- (-1)^n \alpha^{r, s, q} \circ \text{d}_L
\end{align*}
를 얻는다. 이제 원소 $\beta \in K^n$를
$\alpha^{n, -q, q} = \beta \otimes \text{id}_{L^{-q}}$이고
$r \not = n$이면 $\alpha^{r, s, q} = 0$인 $\alpha$로 보낸다.
이는 실제로 위와 같은 원소이다. 고정된 $q$에 대하여 영이 아닌
$\alpha^{r, s, q}$가 하나뿐이기 때문이다. 미분에 대한 설명으로부터
이 구성이 미분과 양립함을 알 수 있다.
\end{proof}

\begin{lemma}
\label{more-algebra-lemma-evaluate-and-more}
$R$을 환이라 하자. $R$-가군들의 복합체
$K^\bullet, L^\bullet, M^\bullet$이 주어지면 세 복합체 모두에 관하여
함자적인 $R$-가군 복합체의 표준 사상
$$
\text{Tot}(\Hom^\bullet(L^\bullet, M^\bullet) \otimes_R K^\bullet)
\longrightarrow
\Hom^\bullet(\Hom^\bullet(K^\bullet, L^\bullet), M^\bullet)
$$
이 존재한다.
\end{lemma}

\begin{proof}
주석 \ref{more-algebra-remark-internal-hom}의 논의를 통하면 그러한 표준 사상의
존재는 범주론, 주석 \ref{categories-remark-internal-hom-monoidal}에서
따른다. 직접적인 구성도 주겠다.

\medskip\noindent
오른쪽 변의 차수 $n$인 원소 $\beta$를 생각하자. 그러면
$$
\beta = (\beta^{p, s}) \in
\prod\nolimits_{p + s = n} \Hom_R(\Hom^{-s}(K^\bullet, L^\bullet), M^p)
$$
이다. 우리의 부호 규칙에 의해
\begin{align*}
\text{d}(\beta^{p, s})
& =
\text{d}_M \circ \beta^{p, s}
- (-1)^n \beta^{p, s} \circ
\text{d}_{\Hom^\bullet(K^\bullet, L^\bullet)}
\end{align*}
이다. 마지막 항은 다음과 같이 설명할 수 있다.
$$
(\beta^{p, s} \circ \text{d}_{\Hom^\bullet(K^\bullet, L^\bullet)})(f) =
\beta^{p, s}(\text{d}_L \circ f - (-1)^{s + 1} f \circ \text{d}_K)
$$
여기서 $f \in \Hom^{-s - 1}(K^\bullet, L^\bullet)$이다. 따라서 엄밀히
특정하지 않은 어떤 의미에서 $\text{d}(\beta^{p, s})$는 다음 부호를
갖는 세 항의 합이다.
\begin{equation}
\label{more-algebra-equation-beta}
\text{d}(\beta^{p, s}) =
\text{d}_M(\beta^{p, s})
-(-1)^n\text{d}_L(\beta^{p, s}) +
(-1)^{p + 1}\text{d}_K(\beta^{p, s})
\end{equation}

\noindent
다음으로 왼쪽 변의 차수 $n$인 원소 $\alpha$를 생각하자. 이는
다음과 같이 쓸 수 있다.
$$
\alpha = (\alpha^{t, r}) \in
\bigoplus\nolimits_{t + r = n} \Hom^t(L^\bullet, M^\bullet) \otimes K^r
$$
각 $\alpha^{t, r}$는 다음 원소로 사상된다.
$$
\alpha^{t, r} \mapsto (\alpha^{p, q, r}) \in
\prod\nolimits_{p + q = t} \Hom_R(L^{-q}, M^p) \otimes_R K^r
$$
우리의 부호 규칙에 의해
\begin{align*}
\text{d}(\alpha^{p, q, r})
& =
\text{d}_{\Hom^\bullet(L^\bullet, M^\bullet)}(\alpha^{p, q, r})
+ (-1)^{p + q} \text{d}_K(\alpha^{p, q, r})
\end{align*}
이다. 여기서 더 나아가
$\alpha^{p, q, r} = \sum g_i^{p, q} \otimes k_i^r$로 쓰면
$$
\text{d}_{\Hom^\bullet(L^\bullet, M^\bullet)}(\alpha^{p, q, r}) =
\sum (\text{d}_M \circ g_i^{p, q}) \otimes k_i^r
- (-1)^{p + q} \sum (g_i^{p, q} \circ \text{d}_L) \otimes k_i^r
$$
이다. 따라서 엄밀히 특정하지 않은 어떤 의미에서
$\text{d}(\alpha^{p, q, r})$는 다음 부호를 갖는 세 항의 합이다.
\begin{equation}
\label{more-algebra-equation-alpha}
\text{d}(\alpha^{p, q, r}) =
\text{d}_M(\alpha^{p, q, r})
-(-1)^{p + q}\text{d}_L(\alpha^{p, q, r}) +
(-1)^{p + q}\text{d}_K(\alpha^{p, q, r})
\end{equation}

\noindent
우리의 사상을 정의하기 위하여 표준 사상
$$
c_{p, q, r} : 
\Hom_R(L^{-q}, M^p) \otimes_R K^r
\longrightarrow
\Hom_R(\Hom_R(K^r, L^{-q}), M^p)
$$
을 사용한다. 이는 $\varphi \otimes k$를 사상
$\psi \mapsto \varphi(\psi(k))$로 보낸다. 이 구성은 세 변수 모두에
관하여 함자적이다. $s = q + r$로 두면 사영
$\Hom^{-s}(K^\bullet, L^\bullet) \to \Hom_R(K^r, L^{-q})$에서 오는 포함
$$
\Hom_R(\Hom_R(K^r, L^{-q}), M^p) \subset
\Hom_R(\Hom^{-s}(K^\bullet, L^\bullet), M^p)
$$
이 있다. $\alpha^{p, q, r}$는 유한 개의 $r$에 대해서만 영이 아니므로,
주어진 $s$에 대하여 $q + r = s$인 $q, r$도 유한 개뿐이다. 따라서
$\alpha$를 다음 $\beta$로 보낼 수 있다.
$$
\beta^{p, s} =
\sum\nolimits_{q + r = s} \epsilon_{p, q, r} c_{p, q, r}(\alpha^{p, q, r})
$$
% P02-E0960: frozen source begins “where where the sum uses”; Korean renders the intended single connective.
여기서 합에는 위에서 주어진 포함들을 사용하며
$\epsilon_{p, q, r} \in \{\pm 1\}$이다. 식
(\ref{more-algebra-equation-beta})와 (\ref{more-algebra-equation-alpha})에서 부호를 비교하면
다음을 얻는다.
\begin{enumerate}
\item $\epsilon_{p, q, r} = \epsilon_{p + 1, q, r}$
\item $-(-1)^n\epsilon_{p, q, r} = -(-1)^{p + q}\epsilon_{p, q - 1, r}$
또는 동치로 $\epsilon_{p, q, r} = (-1)^r\epsilon_{p, q - 1, r}$
\item $(-1)^{p + 1}\epsilon_{p, q, r} = (-1)^{p + q}\epsilon_{p, q, r + 1}$
또는 동치로
$(-1)^{q + 1}\epsilon_{p, q, r} = \epsilon_{p, q, r + 1}$.
\end{enumerate}
좋은 해는 다음을 택하는 것이다.
% P02-E0961: frozen source changes the index tuple from epsilon_{p,q,r} to epsilon_{p,r,s}; Korean preserves the displayed source formula.
$$
\epsilon_{p, r, s} = (-1)^{r + qr}
$$
이 부호를 택하는 이유는 증명 다음의 주석에서 설명한다.
\end{proof}

\begin{remark}
\label{more-algebra-remark-sign-explanation}
보조정리 \ref{more-algebra-lemma-evaluate-and-more}의 증명에서 직접 구성에 사용한
부호가 주석 \ref{more-algebra-remark-internal-hom}과 범주론, 주석
\ref{categories-remark-internal-hom-monoidal}의 논의를 사용하는 구성에서
얻는 부호와 일치하는 이유를 설명하자.
$- \otimes - = \text{Tot}(- \otimes_R -)$와
$hom(-, -) = \Hom^\bullet(-, -)$로 쓰자. 모노이드 범주의 언어를
사용한 구성은 $\text{Comp}(R)$의 사상
$$
hom(L^\bullet, M^\bullet) \otimes K^\bullet
\longrightarrow
hom(hom(K^\bullet, L^\bullet), M^\bullet)
$$
을 사용하라고 한다. 이는 사상
$$
hom(L^\bullet, M^\bullet) \otimes K^\bullet \otimes hom(K^\bullet, L^\bullet)
\longrightarrow
M^\bullet
$$
에 대응하며, 이 사상은 마지막 두 텐서곱의 순서를 바꾼 뒤 평가 사상
$hom(K^\bullet, L^\bullet) \otimes K^\bullet \to L^\bullet$과
% P02-E0962: frozen source gives hom(L,K) in the second evaluation map although the codomain is M; Korean preserves the formal source span.
$hom(L^\bullet, K^\bullet) \otimes L^\bullet \to M^\bullet$을 사용하여
얻는다. 부호는 순서를 바꿀 때에만 개입한다. 즉 동형사상
$$
K^\bullet \otimes hom(K^\bullet, L^\bullet)
\to
hom(K^\bullet, L^\bullet) \otimes K^\bullet
$$
에서 $K^r \otimes_R \Hom_R(K^{-r'}, L^q)$ 위의 부호는
$(-1)^{r(q + r')}$이다. 절 \ref{more-algebra-section-sign-rules}의 항목
(\ref{more-algebra-item-flip-tensor-product})을 보라. 내부 Hom으로 들어가는 사상을
설명하기 위해 우리가 사용하는 대응 때문에 이 부호는 $r = r'$일
때에만 중요하며, 이 경우 직접 증명에서와 같이
$(-1)^{r(q + r)} = (-1)^{r + qr}$를 얻는다는 것을 독자가 확인할 수
있다.
\end{remark}

\section{부호 규칙}
\label{more-algebra-section-sign-rules}

\noindent
이 절에서는 지금까지 사용한 부호 규칙을 검토하고 그 귀결 몇 가지를
논의한다. 흥미로운 현상의 대부분이 이미 이 경우에 나타나므로,
환 위의 가군 복합체 범주라는 배경에서 이 문제들을 논의하는 것도
적절해 보인다. 독자가 이 절의 더 난해한 측면들을 사용할 필요가
없기를 진심으로 바란다.

\medskip\noindent
이 절의 나머지 부분에서는 환 $R$을 고정하고,
% P02-E0963: frozen source uses the malformed construction “we denote M a complex”; Korean renders the intended setup.
$M^\bullet$을 미분 $d^n_M : M^n \to M^{n + 1}$을 갖는
$R$-가군 복합체라 하자.
\begin{enumerate}
\item $k$번째 이동 복합체 $M^\bullet[k]$의 항은
$(M^\bullet[k])^n = M^{n + k}$이고 미분은
$d_{M[k]}^n = (-1)^kd^{n + k}_M$이다. 호몰로지, 정의
\ref{homology-definition-shift-cochain}을 보라.
\item 복합체의 사상 $f : M^\bullet \to N^\bullet$이 주어지면 부호의
개입 없이 $f[k] : M^\bullet[k] \to N^\bullet[k]$를 정의한다.
호몰로지, 정의 \ref{homology-definition-shift-cochain}을 보라.
\item 부호의 개입 없이 $H^n(M^\bullet[k])$를
$H^{n + k}(M^\bullet)$와 식별한다. 호몰로지, 정의
\ref{homology-definition-cohomology-shift}를 보라.
\item 복합체의 짧은 완전열의 경계 사상은 부호의 개입 없이 뱀
보조정리에서처럼 정의한다. 호몰로지, 보조정리
\ref{homology-lemma-long-exact-sequence-cochain}를 보라.
\item 항별로 분할되는 복합체의 짧은 완전열
$0 \to K^\bullet \to L^\bullet \to M^\bullet \to 0$에 대응하는 특수
삼각형은
$$
K^\bullet \to L^\bullet \to M^\bullet \to K^\bullet[1]
$$
로 주어진다. 여기서 $s$와 $\pi$가 항별로 양립하는 분할이면
$M^n \to K^{n + 1}$은 사상
$\pi^{n + 1} \circ d^n_L \circ s^n$이다. 다시 말해 부호는 개입하지
않는다. 유도 범주, 정의 \ref{derived-definition-distinguished-triangle}와
\ref{derived-definition-split-ses}를 보라.
\item 전체 복합체 $\text{Tot}(M^\bullet \otimes_R N^\bullet)$의 미분
$d$는 라이프니츠 규칙
$d(x \otimes y) = d(x) \otimes y + (-1)^{\deg(x)}x \otimes d(y)$를
만족한다. 호몰로지, 예
\ref{homology-example-double-complex-as-tensor-product-of}와 호몰로지,
정의 \ref{homology-definition-associated-simple-complex}를 보라.
\item
\label{more-algebra-item-shift-tensor}
표준 동형사상
$$
\text{Tot}(M^\bullet \otimes_R N^\bullet)[a + b] \to
\text{Tot}(M^\bullet[a] \otimes_R N^\bullet[b])
$$
이 있으며, 이는 직합항 $M^p \otimes_R N^q$ 위에서 부호
$(-1)^{pb}$를 사용한다. 호몰로지, 주석
\ref{homology-remark-shift-double-complex}를 보라. 대응하는 이동된 사상
$\text{Tot}(M^\bullet \otimes_R N^\bullet) \to
\text{Tot}(M^\bullet[a] \otimes_R N^\bullet[b])[-a - b]$를 생각하는 것이
흔히 더 편리하다.
\item 부호의 개입 없이 정의되는 복합체의 표준 동형사상
$$
\text{Tot}(
\text{Tot}(K^\bullet \otimes_R L^\bullet) \otimes_R M^\bullet) \to
\text{Tot}(K^\bullet \otimes_R \text{Tot}(L^\bullet \otimes_R M^\bullet))
$$
이 있다. 절 \ref{more-algebra-section-symmetric-monoidal}을 보라.
\item
\label{more-algebra-item-flip-tensor-product}
표준 동형사상
$$
\text{Tot}(L^\bullet \otimes_R M^\bullet) \to
\text{Tot}(M^\bullet \otimes_R L^\bullet)
$$
이 있으며, 이는 직합항 $L^p \otimes_R M^q$ 위에서 부호
$(-1)^{pq}$를 사용한다. 절 \ref{more-algebra-section-symmetric-monoidal}을 보라.
\end{enumerate}
Hom 복합체에 관한 부호 규약을 논의하기에 앞서, 위의 규약에 따라
복합체의 쌍대를 구성하자.

\begin{lemma}
\label{more-algebra-lemma-left-dual-module}
$R$을 환이라 하자. $M$을 $R$-가군이라 하자.
$N, \eta, \epsilon$을 $R$-가군의 모노이드 범주에서 $M$의 왼쪽
쌍대라 하자. 범주론, 정의 \ref{categories-definition-dual}을 보라.
그러면 다음이 성립한다.
\begin{enumerate}
\item $M$과 $N$은 유한 사영 $R$-가군이다.
\item 사상
$e : \Hom_R(M, R) \to N$, $\lambda \mapsto (\lambda \otimes 1)(\eta)$는
동형사상이다.
\item $n \in N$과 $m \in M$에 대하여
$\epsilon(n, m) = e^{-1}(n)(m)$이다.
\end{enumerate}
\end{lemma}

\begin{proof}
가정은
$$
M \xrightarrow{\eta \otimes 1} M \otimes_R N \otimes_R M
\xrightarrow{1 \otimes \epsilon} M
\quad\text{and}\quad
N \xrightarrow{1 \otimes \eta} N \otimes_R M \otimes_R N
\xrightarrow{\epsilon \otimes 1} N
$$
이 항등사상이라는 뜻이다. 유한 자유 가군 $F$, $R$-가군 사상
$F \to M$, 그리고 $\eta$를 올린 사상
$\tilde \eta : R \to F \otimes_R N$을 택할 수 있다. 가환 그림
$$
\xymatrix{
M \ar[rr]_-{\eta \otimes 1} \ar[rrd]_-{\tilde \eta \otimes 1} & &
M \otimes_R N \otimes_R M \ar[r]_-{1 \otimes \epsilon} &
M \\
& &
F \otimes_R N \otimes_R M \ar[u] \ar[r]^-{1 \otimes \epsilon} &
F \ar[u]
}
$$
을 얻는다. 이는 $M$ 위의 항등사상이 유한 자유 가군을 통하여
분해됨을 보이므로 $M$은 유한 사영이다. 대칭성에 의해 $N$도 유한
사영임을 알 수 있다. 이로써 (1)이 증명된다. (2)는 범주론, 보조정리
\ref{categories-lemma-left-dual}과 그 증명에서 따른다. (3)은 그 증명의
첫 번째 등식에서 따른다.
\end{proof}

\begin{lemma}
\label{more-algebra-lemma-left-dual-complex}
$R$을 환이라 하자. $M^\bullet$을 $R$-가군 복합체라 하자.
$N^\bullet, \eta, \epsilon$을 $R$-가군 복합체의 모노이드 범주에서
$M^\bullet$의 왼쪽 쌍대라 하자. 그러면 다음이 성립한다.
\begin{enumerate}
\item $M^\bullet$과 $N^\bullet$은 유계이다.
\item $M^n$과 $N^n$은 유한 사영 $R$-가군이다.
\item $\epsilon_n : N^{-n} \otimes_R M^n \to R$인
$\epsilon = \sum \epsilon_n$과 $\eta_n : R \to  M^n \otimes_R N^{-n}$인
$\eta = \sum \eta_n$으로 쓰면 $(N^{-n}, \eta_n, \epsilon_n)$은
보조정리 \ref{more-algebra-lemma-left-dual-module}에서와 같은 $M^n$의 왼쪽
쌍대이다.
\item 미분 $d_N^n : N^n \to N^{n + 1}$은 사상
$$
N^n = \Hom_R(M^{-n}, R)
\xrightarrow{d_M^{-n - 1}}
\Hom_R(M^{-n - 1}, R) = N^{n + 1}
$$
의 $-(-1)^n$배이다. 여기서 등호들은 보조정리
\ref{more-algebra-lemma-left-dual-module}의 (2)에서 오는 식별이다.
\end{enumerate}
역으로 유한 사영 $R$-가군들로 이루어진 유계 복합체 $M^\bullet$이
주어졌다고 하자. $N^n = \Hom_R(M^{-n}, R)$으로 두고 미분을 위와 같이
정의하고, $\epsilon_n : N^{-n} \otimes_R M^n \to R$이 평가로 주어지는
$\epsilon = \sum \epsilon_n$으로 두며,
% P02-E0964: frozen source maps 1 to id_{M_n}, switching from the defined M^n; Korean preserves the formal source token.
$\eta_n : R \to M^n \otimes_R N^{-n}$이 $1$을
$\text{id}_{M_n}$으로 보내는 $\eta = \sum \eta_n$으로 두면,
$R$-가군 복합체의 모노이드 범주에서 $M^\bullet$의 왼쪽 쌍대를
얻는다.
\end{lemma}

\begin{proof}
범주론, 정의 \ref{categories-definition-dual}에 의해
$(1 \otimes \epsilon) \circ (\eta \otimes 1) = \text{id}_{M^\bullet}$이고
$(\epsilon \otimes 1) \circ (1 \otimes \eta) = \text{id}_{N^\bullet}$이므로
곧바로
$(1 \otimes \epsilon_n) \circ (\eta_n \otimes 1) = \text{id}_{M^n}$과
$(\epsilon_n \otimes 1) \circ (1 \otimes \eta_n) = \text{id}_{N^{-n}}$을
얻는다. 이로써 (3)이 증명된다. 보조정리
\ref{more-algebra-lemma-left-dual-module}에 의해 (2)를 얻는다.
합 $\eta = \sum \eta_n$이 유한하므로 (1)을 얻는다.
$\eta = \sum \eta_n$은 복합체의 사상
$R \to \text{Tot}(M^\bullet \otimes_R N^\bullet)$이므로
$$
(d_M^{-n - 1} \otimes 1) \circ \eta_{-n - 1}  +
(-1)^n (1 \otimes d_N^{-n}) \circ \eta_{-n} = 0
$$
이다. 이는 $\text{Tot}(M^\bullet \otimes_R N^\bullet)$ 위의 미분에 대한
우리의 부호 선택에서 따른다. 정의를 풀어 쓰면 (4)가 증명된다.
보조정리의 마지막 명제는 위의 논의를 거꾸로 읽으면 된다.
\end{proof}

\noindent
$\Hom$ 복합체의 부호 규칙에 대한 동기로 보조정리
\ref{more-algebra-lemma-left-dual-complex}에 주어진 복합체의 왼쪽 쌍대에 대한 설명을
사용할 것이다.
% P02-E0965: frozen source misspells “Namely” as “Namelly”; Korean renders the intended connective.
즉, (\ref{more-algebra-item-compatible})이 성립하도록 부호를 택한다.
위에서처럼 여러 부호 규칙에 대한 논의를 계속하자.
\begin{enumerate}
\setcounter{enumi}{9}
\item 복합체 $K^\bullet$, $M^\bullet$이 주어지면
$\Hom^\bullet(M^\bullet, K^\bullet)$을 항
$$
\Hom^n(M^\bullet, K^\bullet) = \prod\nolimits_{n = p + q} \Hom_R(M^{-q}, K^p)
$$
과 규칙
$$
d(f) = d_K \circ f - (-1)^n f \circ d_M
$$
으로 주어지는 미분을 갖는 복합체로 정의한다.
\item
\label{more-algebra-item-compatible}
위의 선택은 다음 성질을 갖는다. $M^\bullet$이 보조정리
\ref{more-algebra-lemma-left-dual-complex}에서와 같은 왼쪽 쌍대 $N^\bullet$을 가지면,
부호의 개입 없이 정의되는 표준 동형사상
$$
\text{Tot}(K^\bullet \otimes_R N^\bullet)
\longrightarrow
\Hom^\bullet(M^\bullet, K^\bullet)
$$
이 있다. 이는 $N^q = \Hom_R(M^{-q}, R)$와 표준 사상
$K^p \otimes_R \Hom_R(M^{-q}, R) \to \Hom_R(M^{-q}, K^p)$를 통하여
직합항 $K^p \otimes_R N^q$를 직합항
$\Hom_R(M^{-q}, K^p)$로 보낸다.
\item 부호의 개입 없이 정의되는 합성
$$
\text{Tot}(
\Hom^\bullet(L^\bullet, K^\bullet) \otimes_R
\Hom^\bullet(M^\bullet, L^\bullet))
\longrightarrow
\Hom^\bullet(M^\bullet, K^\bullet)
$$
이 있다. 보조정리 \ref{more-algebra-lemma-composition}을 보라.
\item 부호의 개입 없이 정의되는 표준 동형사상
$$
\Hom^\bullet(K^\bullet, \Hom^\bullet(L^\bullet, M^\bullet))
=
\Hom^\bullet(\text{Tot}(K^\bullet \otimes_R L^\bullet), M^\bullet)
$$
이 있다. 보조정리 \ref{more-algebra-lemma-compose}를 보라.
\item 부호의 개입 없이 정의되는 표준 사상
$$
\text{Tot}(K^\bullet \otimes_R \Hom^\bullet(M^\bullet, L^\bullet))
\longrightarrow
\Hom^\bullet(M^\bullet, \text{Tot}(K^\bullet \otimes_R L^\bullet))
$$
이 있다. 보조정리 \ref{more-algebra-lemma-diagonal-better}를 보라.
\item 부호의 개입 없이 정의되는 표준 사상
$$
K^\bullet
\longrightarrow
\Hom^\bullet(L^\bullet, \text{Tot}(K^\bullet \otimes_R L^\bullet))
$$
이 있다. 보조정리 \ref{more-algebra-lemma-diagonal}을 보라.
\item
% P02-E0966: frozen source says “By Lemma ... is a canonical map”; Korean supplies the intended existential construction.
보조정리 \ref{more-algebra-lemma-evaluate-and-more}에 의해 표준 사상
$$
\text{Tot}(\Hom^\bullet(L^\bullet, M^\bullet) \otimes_R K^\bullet)
\longrightarrow
\Hom^\bullet(\Hom^\bullet(K^\bullet, L^\bullet), M^\bullet)
$$
이 있다. 이는 가군 $\Hom_R(L^{-q}, M^p) \otimes_R K^r$ 위에서 부호
$(-1)^{r + qr}$를 사용하며, 그 이유는 주석
\ref{more-algebra-remark-sign-explanation}에서 설명한다.
\item
\label{more-algebra-item-evaluation}
$L^\bullet = M^\bullet$로 두고 $R \to \Hom^\bullet(M^\bullet, M^\bullet)$을
사용하면 앞 항목의 사상은 평가 사상
$$
ev : K^\bullet \longrightarrow
\Hom^\bullet(\Hom^\bullet(K^\bullet, M^\bullet), M^\bullet)
$$
이 된다. 이는 $x \in K^n$을,
$f \in \Hom^m(K^\bullet, M^\bullet)$를 $(-1)^{nm}f(x)$로 보내는 사상으로
보낸다.
\item
\label{more-algebra-item-shift-hom}
부호를 사용하는 표준 식별
$$
\Hom^\bullet(M^\bullet, K^\bullet)[a - b] \to
\Hom^\bullet(M^\bullet[b], K^\bullet[a])
$$
이 있다. 이는 그에 대응하는 이동된 사상
$$
\Hom^\bullet(M^\bullet, K^\bullet) \to
\Hom^\bullet(M^\bullet[b], K^\bullet[a])[b - a]
$$
이 $p + q = n$인 가군 $\Hom_R(M^{-q}, K^p)$ 위에서 부호
$(-1)^{nb}$를 사용하도록 정의된다. 즉
$f \in \Hom^n(M^\bullet, K^\bullet)$이면 원천에서
$$
d(f) = d_K \circ f - (-1)^n f \circ d_M
$$
이다. 한편 목표에서는 $f$가
$\left(\Hom^\bullet(M^\bullet[b], K^\bullet[a])[b - a]\right)^n =
\Hom^{n + b -a}(M^\bullet[b], K^\bullet[a])$에 속하므로
\begin{align*}
d(f) & =
(-1)^{b - a}
\left(d_{K[a]} \circ f - (-1)^{n + b - a} f \circ d_{M[b]}\right) \\
& =
(-1)^{b - a}
\left((-1)^a d_K \circ f - (-1)^{n + b - a} f \circ (-1)^b d_M \right) \\
& =
(-1)^b d_K \circ f - (-1)^{n + b} f \circ d_M
\end{align*}
를 얻는다. 따라서 차수 $n$에서 택한 부호 $(-1)^{nb}$가 이 미분들에
대하여 복합체의 사상을 준다는 것을 알 수 있다.
\end{enumerate}

\section{유도 Hom}
\label{more-algebra-section-RHom}

\noindent
$R$을 환이라 하자. 이 절에서 정의할 유도 Hom은 함자
$$
D(R)^{opp} \times D(R) \longrightarrow D(R),\quad
(K, L) \longmapsto R\Hom_R(K, L)
$$
이다. 이는 $R$-가군의 유도 범주에서 내부 Hom이다. 즉, $D(R)$의
대상 $K, L, M$에 대하여 식
\begin{equation}
\label{more-algebra-equation-internal-hom}
\Hom_{D(R)}(K, R\Hom_R(L, M))
=
\Hom_{D(R)}(K \otimes_R^\mathbf{L} L, M)
\end{equation}
으로 특징지어진다. 요네다 보조정리에 의해 이 식은 대상을 유일한
동형사상까지 특징짓는다는 점에 유의하라. 다음과 같이 구성할 수
있다. $M$을 나타내는 $R$-가군의 K-단사 복합체 $I^\bullet$과 $L$을
나타내는 복합체 $L^\bullet$을 택하고
$$
R\Hom_R(L, M) = \Hom^\bullet(L^\bullet, I^\bullet)
$$
로 둔다. 표기는 절 \ref{more-algebra-section-hom-complexes}에서와 같다. 이 구성의
일반화는 미분 등급 대수, 절 \ref{dga-section-restriction}에서 논의한다.
% P02-E0967: frozen source says “From (...) and [the lemma] that we have”, omitting a finite verb; Korean supplies the intended implication.
(\ref{more-algebra-equation-cohomology-hom-complex})과 유도 범주, 보조정리
\ref{derived-lemma-K-injective}에서 모든 $n \in \mathbf{Z}$에 대하여
\begin{equation}
\label{more-algebra-equation-h0-RHom}
H^n(R\Hom_R(L, M)) = \Hom_{D(R)}(L, M[n])
\end{equation}
임을 얻는다. 특히 $D(R)$의 대상 $R\Hom_R(L, M)$은 잘 정의된다.
즉, K-단사 복합체 $I^\bullet$의 선택과 무관하다.

\begin{lemma}
\label{more-algebra-lemma-internal-hom}
$R$을 환이라 하자. $K, L, M$을 $D(R)$의 대상이라 하자.
$K, L, M$에 관하여 함자적인 $D(R)$의 표준 동형사상
$$
R\Hom_R(K, R\Hom_R(L, M)) = R\Hom_R(K \otimes_R^\mathbf{L} L, M)
$$
이 존재하며, 여기에 $H^0$을 취하면 (\ref{more-algebra-equation-internal-hom})을
되찾는다.
\end{lemma}

\begin{proof}
$M$을 나타내는 K-단사 복합체 $I^\bullet$과 $L$을 나타내는
$R$-가군의 K-평탄 복합체 $L^\bullet$을 택하자. 임의의 $R$-가군
복합체 $K^\bullet$에 대하여 보조정리 \ref{more-algebra-lemma-compose}에 의해
$$
\Hom^\bullet(K^\bullet, \Hom^\bullet(L^\bullet, I^\bullet)) =
\Hom^\bullet(\text{Tot}(K^\bullet \otimes_R L^\bullet), I^\bullet)
$$
이다. $R\Hom$의 정의와
$\text{Tot}(K^\bullet \otimes_R L^\bullet)$이 유도 텐서곱을 나타낸다는
사실로부터 보조정리가 따른다.
\end{proof}

\begin{lemma}
\label{more-algebra-lemma-RHom-out-of-projective}
$R$을 환이라 하자. $P^\bullet$을 사영 $R$-가군들로 이루어진 위로
유계인 복합체라 하고, $L^\bullet$을 $R$-가군 복합체라 하자. 그러면
$R\Hom_R(P^\bullet, L^\bullet)$은 복합체
$\Hom^\bullet(P^\bullet, L^\bullet)$로 나타내어진다.
\end{lemma}

\begin{proof}
(\ref{more-algebra-equation-cohomology-hom-complex})과 유도 범주, 보조정리
\ref{derived-lemma-morphisms-from-projective-complex}에 의해 이 복합체의
코호몰로지 군들은 “올바르다”. 따라서 $I^\bullet$이 $R$-가군의
K-단사 복합체인 준동형사상 $L^\bullet \to I^\bullet$을 택하면 유도된
사상
$$
\Hom^\bullet(P^\bullet, L^\bullet)
\longrightarrow
\Hom^\bullet(P^\bullet, I^\bullet)
$$
은 준동형사상이다. 오른쪽 변이 $R\Hom_R(P^\bullet, L^\bullet)$의
정의이므로 결론을 얻는다.
\end{proof}

\begin{lemma}
\label{more-algebra-lemma-internal-hom-evaluate}
$R$을 환이라 하자. $K, L, M$을 $D(R)$의 대상이라 하자.
$K, L, M$에 관하여 함자적인 $D(R)$의 표준 사상
$$
R\Hom_R(L, M) \otimes_R^\mathbf{L} K
\longrightarrow
R\Hom_R(R\Hom_R(K, L), M)
$$
이 존재한다.
\end{lemma}

\begin{proof}
$M$을 나타내는 K-단사 복합체 $I^\bullet$, $L$을 나타내는 K-단사
복합체 $J^\bullet$, 그리고 $K$를 나타내는 K-평탄 복합체
$K^\bullet$을 택한다. 이 사상은 보조정리
\ref{more-algebra-lemma-evaluate-and-more}의 사상
$$
\text{Tot}(\Hom^\bullet(J^\bullet, I^\bullet) \otimes_R K^\bullet)
\longrightarrow
\Hom^\bullet(\Hom^\bullet(K^\bullet, J^\bullet), I^\bullet)
$$
을 사용하여 정의한다. 이것이 $D(R)$의 세 대상 모두에 관하여
함자적이라는 증명은 생략한다.
\end{proof}

\begin{lemma}
\label{more-algebra-lemma-internal-hom-composition}
$R$을 환이라 하자. $D(R)$의 $K, L, M$이 주어지면
$K, L, M$에 관하여 함자적인 $D(R)$의 표준 사상
$$
R\Hom_R(L, M) \otimes_R^\mathbf{L} R\Hom_R(K, L) \longrightarrow R\Hom_R(K, M)
$$
이 존재한다.
\end{lemma}

\begin{proof}
$M$을 나타내는 K-단사 복합체 $I^\bullet$, $L$을 나타내는 K-단사
복합체 $J^\bullet$, 그리고 $K$를 나타내는 임의의 $R$-가군 복합체
$K^\bullet$을 택하자. 보조정리 \ref{more-algebra-lemma-composition}에 의해 복합체의
사상
$$
\text{Tot}\left(
\Hom^\bullet(J^\bullet, I^\bullet) \otimes_R \Hom^\bullet(K^\bullet, J^\bullet)
\right)
\longrightarrow
\Hom^\bullet(K^\bullet, I^\bullet)
$$
이 있다. $R$-가군 복합체
$\Hom^\bullet(J^\bullet, I^\bullet)$,
$\Hom^\bullet(K^\bullet, J^\bullet)$,
$\Hom^\bullet(K^\bullet, I^\bullet)$은 각각
$R\Hom_R(L, M)$, $R\Hom_R(K, L)$, $R\Hom_R(K, M)$을 나타낸다.
K-평탄 복합체 $H^\bullet$과 준동형사상
$H^\bullet \to \Hom^\bullet(K^\bullet, J^\bullet)$을 택하면 사상
$$
\text{Tot}\left(
\Hom^\bullet(J^\bullet, I^\bullet) \otimes_R H^\bullet
\right)
\longrightarrow
\text{Tot}\left(
\Hom^\bullet(J^\bullet, I^\bullet) \otimes_R \Hom^\bullet(K^\bullet, J^\bullet)
\right)
$$
이 있고, 그 원천은
$R\Hom_R(L, M) \otimes_R^\mathbf{L} R\Hom_R(K, L)$을 나타낸다.
표시된 두 화살표를 합성하면 원하는 사상을 얻는다. 이 구성이
함자적이라는 증명은 생략한다.
\end{proof}

\begin{lemma}
\label{more-algebra-lemma-internal-hom-diagonal-better}
$R$을 환이라 하자. $D(R)$의 복합체 $K, L, M$이 주어지면
$K$, $L$, $M$에 관하여 함자적인 $D(R)$의 표준 사상
$$
K \otimes_R^\mathbf{L} R\Hom_R(M, L)
\longrightarrow
R\Hom_R(M, K \otimes_R^\mathbf{L} L)
$$
이 존재한다.
\end{lemma}

\begin{proof}
$K$를 나타내는 K-평탄 복합체 $K^\bullet$, $L$을 나타내는 K-단사
복합체 $I^\bullet$, 그리고 $M$을 나타내는 임의의 복합체
$M^\bullet$을 택하자. $J^\bullet$이 K-단사인 준동형사상
$\text{Tot}(K^\bullet \otimes_R I^\bullet) \to J^\bullet$을 택한다.
그러면 사상
$$
\text{Tot}\left(
K^\bullet \otimes_R \Hom^\bullet(M^\bullet, I^\bullet)
\right)
\to
\Hom^\bullet(M^\bullet, \text{Tot}(K^\bullet \otimes_R I^\bullet))
\to
\Hom^\bullet(M^\bullet, J^\bullet)
$$
을 사용한다. 여기서 첫 번째 사상은 보조정리
\ref{more-algebra-lemma-diagonal-better}의 사상이다.
\end{proof}

\begin{lemma}
\label{more-algebra-lemma-internal-hom-diagonal}
$R$을 환이라 하자. $D(R)$의 복합체 $K, L$이 주어지면
$K$와 $L$ 모두에 관하여 함자적인 $D(R)$의 표준 사상
$$
K \longrightarrow R\Hom_R(L, K \otimes_R^\mathbf{L} L)
$$
이 존재한다.
\end{lemma}

\begin{proof}
이는 보조정리 \ref{more-algebra-lemma-internal-hom-diagonal-better}의 특수한
경우이지만 직접 증명도 하겠다. $K$를 나타내는 K-평탄 복합체
$K^\bullet$과 $L$을 나타내는 임의의 복합체 $L^\bullet$을 택하자.
$J^\bullet$이 K-단사인 준동형사상
$\text{Tot}(K^\bullet \otimes_R L^\bullet) \to J^\bullet$을 택한다.
그러면 사상
$$
K^\bullet \to
\Hom^\bullet(L^\bullet, \text{Tot}(K^\bullet \otimes_R L^\bullet))
\to \Hom^\bullet(L^\bullet, J^\bullet)
$$
을 사용한다. 여기서 첫 번째 사상은 보조정리
\ref{more-algebra-lemma-diagonal}의 사상이다.
\end{proof}












\section{퍼펙트 복합체}
\label{more-algebra-section-perfect}

\noindent
퍼펙트 복합체란 유한 Tor 차원을 갖는 유사연접 복합체이다. 이를
정의로 사용하지는 않고, 환 위의 퍼펙트 복합체를 다음과 같이 직접
정의하겠다.

\begin{definition}
\label{more-algebra-definition-perfect}
$R$을 환이라 하자.
% P02-E0968: frozen source uses the malformed construction “Denote D(R) the derived category”; Korean renders the intended notation.
$R$-가군의 아벨 범주의 유도 범주를 $D(R)$로 나타내자.
\begin{enumerate}
\item $D(R)$의 대상 $K$가 유한 사영 $R$-가군들로 이루어진 유계
복합체와 준동형이면 이를 {\it 퍼펙트}라 한다.
\item $M[0]$이 $D(R)$의 퍼펙트 대상이면 $R$-가군 $M$을
{\it 퍼펙트}라 한다.
\end{enumerate}
\end{definition}

\noindent
예를 들어 뇌터 환 위에서 유한 가군이 퍼펙트일 필요충분조건은 유한
사영 차원을 갖는 것이다. 보조정리 \ref{more-algebra-lemma-perfect-module}과
대수학, 정의 \ref{algebra-definition-finite-proj-dim}을 보라.

\begin{lemma}
\label{more-algebra-lemma-perfect}
$K^\bullet$을 $D(R)$의 대상이라 하자. 다음은 동치이다.
\begin{enumerate}
\item $K^\bullet$은 퍼펙트이다.
\item $K^\bullet$은 유사연접이고 유한 Tor 차원을 갖는다.
\end{enumerate}
(1)과 (2)가 성립하고 $K^\bullet$의 Tor 진폭이 $[a, b]$ 안에 있으면,
$i \not \in [a, b]$에 대하여 $E^i = 0$인 유한 사영 $R$-가군들의
복합체 $E^\bullet$과 $K^\bullet$은 준동형이다.
\end{lemma}

\begin{proof}
(1)이 (2)를 함의함은 분명하다. 보조정리
\ref{more-algebra-lemma-pseudo-coherent}와 \ref{more-algebra-lemma-tor-amplitude}를 보라.
(2)가 성립하고 $K^\bullet$의 Tor 진폭이 $[a, b]$ 안에 있다고
가정하자. 특히 $i > b$이면 $H^i(K^\bullet) = 0$이다.
$i > b$이면 $F^i = 0$인 유한 자유 $R$-가군들의 복합체
$F^\bullet$과 준동형사상 $F^\bullet \to K^\bullet$을 택하자
(보조정리 \ref{more-algebra-lemma-pseudo-coherent}).
$E^\bullet = \tau_{\geq a}F^\bullet$로 두자. $E^i$는 $E^a$를
제외하면 유한 자유이고, 예외인 것은 유한 표시 $R$-가군임에 유의하라.
보조정리 \ref{more-algebra-lemma-last-one-flat}에 의해 $E^a$는 평탄하다.
따라서 대수학, 보조정리 \ref{algebra-lemma-finite-projective}에 의해
$E^a$가 유한 사영임을 알 수 있다.
\end{proof}

\begin{lemma}
\label{more-algebra-lemma-perfect-module}
$M$을 환 $R$ 위의 가군이라 하자. 다음은 동치이다.
\begin{enumerate}
\item $M$은 퍼펙트 가군이다.
\item 각 $F_i$가 유한 사영 $R$-가군인 분해
$$
0 \to F_d \to \ldots \to F_1 \to F_0 \to M \to 0
$$
가 존재한다.
\end{enumerate}
\end{lemma}

\begin{proof}
(2)를 가정하자. $E^{-i} = F_i$인 복합체 $E^\bullet$은 $M[0]$과
준동형이다. 따라서 $M$은 퍼펙트이다. 역으로 (1)을 가정하자.
보조정리 \ref{more-algebra-lemma-perfect}와 \ref{more-algebra-lemma-n-pseudo-module}에 의해
% P02-E0969: frozen source says “find resolution” without an article; Korean renders the intended construction.
$E^{-i}$가 유한 자유 $R$-가군인 분해 $E^\bullet \to M$을 찾을 수
있다. 보조정리 \ref{more-algebra-lemma-last-one-flat}에 의해 충분히 큰 어떤
$d$에 대하여
% P02-E0970: frozen source uses positive E exponents despite the established E^{-i} indexing; Korean preserves the displayed source formula.
$F_d = \Coker(E^{d - 1} \to E^d)$가 평탄함을 알 수 있다. 대수학,
보조정리 \ref{algebra-lemma-finite-projective}에 의해 $F_d$는 유한
사영이다. 따라서
$$
0 \to F_d \to E^{-d+1} \to \ldots \to E^0 \to M \to 0
$$
은 원하는 분해이다.
\end{proof}

\begin{lemma}
\label{more-algebra-lemma-two-out-of-three-perfect}
$R$을 환이라 하자. $(K^\bullet, L^\bullet, M^\bullet, f, g, h)$를
$D(R)$의 특수 삼각형이라 하자. $K^\bullet, L^\bullet, M^\bullet$
중 둘이 퍼펙트이면 나머지 하나도 퍼펙트이다.
\end{lemma}

\begin{proof}
보조정리 \ref{more-algebra-lemma-perfect}, \ref{more-algebra-lemma-two-out-of-three-pseudo-coherent},
\ref{more-algebra-lemma-cone-tor-amplitude}를 결합한다.
\end{proof}

\begin{lemma}
\label{more-algebra-lemma-summands-perfect}
$R$을 환이라 하자. $K^\bullet \oplus L^\bullet$이 퍼펙트이면
$K^\bullet$과 $L^\bullet$도 퍼펙트이다.
\end{lemma}

\begin{proof}
보조정리 \ref{more-algebra-lemma-perfect}, \ref{more-algebra-lemma-summands-pseudo-coherent},
\ref{more-algebra-lemma-summands-tor-amplitude}에서 따른다.
\end{proof}

\begin{lemma}
\label{more-algebra-lemma-complex-perfect-modules}
$R$을 환이라 하자. $K^\bullet$을 퍼펙트 $R$-가군들로 이루어진 유계
복합체라 하자. 그러면 $K^\bullet$은 퍼펙트 복합체이다.
\end{lemma}

\begin{proof}
유한 복합체의 길이에 대한 귀납법으로 따른다. 보조정리
\ref{more-algebra-lemma-two-out-of-three-perfect}와 ‘stupid’ 절단을 사용하라.
\end{proof}

\begin{lemma}
\label{more-algebra-lemma-cohomology-perfect}
$R$을 환이라 하자. $K^\bullet \in D^b(R)$이고 그 모든 코호몰로지
가군이 퍼펙트이면 $K^\bullet$은 퍼펙트이다.
\end{lemma}

\begin{proof}
유한 복합체의 길이에 대한 귀납법으로 따른다. 보조정리
\ref{more-algebra-lemma-two-out-of-three-perfect}와 표준 절단을 사용하라.
\end{proof}

\begin{lemma}
\label{more-algebra-lemma-perfect-push-perfect}
$A \to B$를 환 사상이라 하자. $B$가 $A$-가군으로서 퍼펙트라고
가정하자. $K^\bullet$을 $B$-가군의 퍼펙트 복합체라 하자. 그러면
$K^\bullet$은 $A$-가군 복합체로서 퍼펙트이다.
\end{lemma}

\begin{proof}
보조정리 \ref{more-algebra-lemma-perfect}를 사용하면 이는 유사연접 가군과 유한
Tor 차원의 가군에 관한 대응하는 결과로 바뀐다. 그 결과들에 대해서는
보조정리 \ref{more-algebra-lemma-finite-tor-dimension-push-tor-amplitude}와
보조정리 \ref{more-algebra-lemma-finite-push-pseudo-coherent}를 보라.
\end{proof}

\begin{lemma}
\label{more-algebra-lemma-pull-perfect}
$A \to B$를 환 사상이라 하자. $K^\bullet$을 $A$-가군의 퍼펙트
복합체라 하자. 그러면 $K^\bullet \otimes_A^{\mathbf{L}} B$는
$B$-가군의 퍼펙트 복합체이다.
\end{lemma}

\begin{proof}
보조정리 \ref{more-algebra-lemma-perfect}를 사용하면 이는 유사연접 가군과 유한
Tor 차원의 가군에 관한 대응하는 결과로 바뀐다. 그 결과들에 대해서는
보조정리 \ref{more-algebra-lemma-pull-tor-amplitude}와
보조정리 \ref{more-algebra-lemma-pull-pseudo-coherent}를 보라.
\end{proof}

\begin{lemma}
\label{more-algebra-lemma-flat-base-change-perfect}
$A \to B$를 평탄한 환 사상이라 하자. $M$을 퍼펙트 $A$-가군이라
하자. 그러면 $M \otimes_A B$는 퍼펙트 $B$-가군이다.
\end{lemma}

\begin{proof}
보조정리 \ref{more-algebra-lemma-perfect-module}에 의해 가정은 $M$이 유한 사영
% P02-E0971: frozen source says the resolution consists of finite projective R-modules although only A and B are defined; Korean preserves R.
$R$-가군들에 의한 유한 분해 $F_\bullet$을 가짐을 함의한다.
$A \to B$가 평탄하므로 복합체 $F_\bullet \otimes_A B$는 유한 사영
$B$-가군들에 의한 $M \otimes_A B$의 유한 길이 분해이다. 따라서
$M \otimes_A B$는 퍼펙트이다.
\end{proof}

\begin{lemma}
\label{more-algebra-lemma-tensor-perfect}
$R$을 환이라 하자. $K$와 $L$이 $D(R)$의 퍼펙트 대상이면
$K \otimes_R^\mathbf{L} L$도 퍼펙트 대상이다.
\end{lemma}

\begin{proof}
다음과 같이 정의를 사용하여 증명할 수 있다. $K$, 각각 $L$을 유한
사영 $R$-가군들의 유계 복합체 $K^\bullet$, 각각 $L^\bullet$로
나타낼 수 있다. 그러면 $K \otimes_R^\mathbf{L} L$은 유계 복합체
$\text{Tot}(K^\bullet \otimes_R L^\bullet)$로 나타내어진다.
% P02-E0972: frozen source uses the undefined M^a three times after choosing K^bullet and L^bullet; Korean preserves every formal source span.
이 복합체의 항들은 가군 $M^a \otimes_R L^b$들의 직합이다.
$M^a$와 $L^b$는 유한 자유 $R$-가군들의 직합인자이므로
$M^a \otimes_R L^b$도 그러하다. 따라서 복합체
$\text{Tot}(K^\bullet \otimes_R L^\bullet)$의 항들은 유한
사영이라고 결론짓는다.

\medskip\noindent
또 다른 증명은 보조정리 \ref{more-algebra-lemma-perfect}의 퍼펙트 복합체에 대한
특징짓기와 유사연접 복합체에 대한 대응하는 보조정리
(보조정리 \ref{more-algebra-lemma-tensor-pseudo-coherent}), 그리고 Tor 진폭에
대한 대응하는 보조정리($A = B = R$로 둔 보조정리
\ref{more-algebra-lemma-push-tor-amplitude})를 사용하여 얻을 수 있다.
\end{proof}

\begin{lemma}
\label{more-algebra-lemma-glue-perfect}
$R$을 환이라 하자. $f_1, \ldots, f_r \in R$을 단위 아이디얼을
생성하는 원소들이라 하자. $K^\bullet$을 $R$-가군 복합체라 하자.
각 $i$에 대하여 복합체 $K^\bullet \otimes_R R_{f_i}$가 퍼펙트이면
$K^\bullet$은 퍼펙트이다.
\end{lemma}

\begin{proof}
보조정리 \ref{more-algebra-lemma-perfect}를 사용하면 이는 유사연접 가군과 유한
Tor 차원의 가군에 관한 대응하는 결과로 바뀐다. 그 결과들에 대해서는
보조정리 \ref{more-algebra-lemma-glue-tor-amplitude}와
보조정리 \ref{more-algebra-lemma-glue-pseudo-coherent}를 보라.
\end{proof}

\begin{lemma}
\label{more-algebra-lemma-flat-descent-perfect}
$R$을 환이라 하자. $K^\bullet$을 $R$-가군 복합체라 하자.
$R \to R'$을 충실하게 평탄한 환 사상이라 하자. 복합체
$K^\bullet \otimes_R R'$이 퍼펙트이면 $K^\bullet$은 퍼펙트이다.
\end{lemma}

\begin{proof}
보조정리 \ref{more-algebra-lemma-perfect}를 사용하면 이는 유사연접 가군과 유한
Tor 차원의 가군에 관한 대응하는 결과로 바뀐다. 그 결과들에 대해서는
보조정리 \ref{more-algebra-lemma-flat-descent-tor-amplitude}와
보조정리 \ref{more-algebra-lemma-flat-descent-pseudo-coherent}를 보라.
\end{proof}

\begin{lemma}
\label{more-algebra-lemma-regular-perfect}
$R$을 정칙환이라 하자. 그러면 다음이 성립한다.
\begin{enumerate}
\item $R$-가군이 퍼펙트일 필요충분조건은 유한 $R$-가군인 것이다.
\item $R$-가군 복합체 $K^\bullet$이 퍼펙트일 필요충분조건은
$K^\bullet \in D^b(R)$이고 각 $H^i(K^\bullet)$가 유한
$R$-가군인 것이다.
\end{enumerate}
\end{lemma}

\begin{proof}
정의에 의해 모든 퍼펙트 $R$-가군은 유한이다. 역으로 $M$을 유한
$R$-가군이라 하자. $F_i$가 유한 자유 $R$-가군인 분해
$$
\ldots \to F_2 \xrightarrow{d_2} F_1 \xrightarrow{d_1} F_0 \to M \to 0
$$
을 택하자(대수학, 보조정리
\ref{algebra-lemma-resolution-by-finite-free}).
$M_i = \Ker(d_i)$로 두자.
% P02-E0973: frozen source uses the malformed construction “Denote U_i ... the set”; Korean renders the intended notation.
$M_{i, \mathfrak p}$가 자유인 소아이디얼 $\mathfrak p$들의 집합을
$U_i \subset \Spec(R)$로 나타내자. 대수학, 보조정리
\ref{algebra-lemma-finitely-presented-localization-free}에 의해 $U_i$는
열린집합이다.
% P02-E0974: frozen source says “a exact sequence”; Korean renders the intended article-free statement.
완전열 $0 \to M_{i + 1} \to F_{i + 1} \to M_i \to 0$이 있다.
$\mathfrak p \in U_i$이면
$0 \to M_{i + 1, \mathfrak p} \to F_{i + 1, \mathfrak p} \to
M_{i, \mathfrak p} \to 0$은 분할된다. 따라서
$M_{i + 1, \mathfrak p}$는 유한 사영이고, 그러므로 자유이다
(대수학, 보조정리 \ref{algebra-lemma-finite-projective}).
이는 $U_i \subset U_{i + 1}$임을 보인다.
$\Spec(R) = \bigcup U_i$라고 주장한다. 실제로 모든 소아이디얼
$\mathfrak p$에 대하여 정칙 국소환 $R_\mathfrak p$는 대수학, 명제
\ref{algebra-proposition-regular-finite-gl-dim}에 의해 유한 대역 차원을
갖는다. 따라서 예를 들어 대수학, 보조정리
\ref{algebra-lemma-independent-resolution}에 의해 $i \gg 0$이면
$M_{i, \mathfrak p}$는 유한 사영이고, 그러므로 자유이다.
$R$의 스펙트럼은 뇌터 공간이므로(대수학, 보조정리
\ref{algebra-lemma-Noetherian-topology}), 어떤 $n$에 대하여
$U_n = \Spec(R)$이라고 결론짓는다. 그러면 대수학, 보조정리
\ref{algebra-lemma-finite-projective}에 의해 $M_n$은 사영
$R$-가군이다. 따라서
% P02-E0975: frozen source omits F_0 from the final bounded resolution although it appeared in the chosen resolution; Korean preserves the display.
$$
0 \to M_n \to F_n \to \ldots \to F_1 \to M \to 0
$$
은 유한 사영 가군들에 의한 유계 분해이므로 $M$은 퍼펙트이다.
이로써 (1)이 증명된다.

\medskip\noindent
$K^\bullet$을 $R$-가군 복합체라 하자. $K^\bullet$이 퍼펙트이면
$D^b(R)$에 속하고 유한 사영 $R$-가군들로 이루어진 유한 복합체와
준동형이므로, $R$이 뇌터 환인 것까지 고려하면 각
$H^i(K^\bullet)$는 분명히 유한 $R$-가군이다. 역으로
$K^\bullet$이 $D^b(R)$에 속하고 각 $H^i(K^\bullet)$가 유한
$R$-가군이라고 가정하자. 그러면 (1)에 의해 각 $H^i(K^\bullet)$는
퍼펙트 $R$-가군이고, 따라서 보조정리
\ref{more-algebra-lemma-cohomology-perfect}에 의해 $K^\bullet$은 퍼펙트이다.
\end{proof}

\begin{lemma}
\label{more-algebra-lemma-dual-perfect-complex}
$A$를 환이라 하자. $K \in D(A)$가 퍼펙트라고 하자. 그러면
$K^\vee = R\Hom_A(K, A)$는 퍼펙트 복합체이고
$K \cong (K^\vee)^\vee$이다. $L \in D(A)$에 대하여 함자적
동형사상
$$
L \otimes_A^\mathbf{L} K^\vee = R\Hom_A(K, L)
\quad\text{and}\quad
H^0(L \otimes_A^\mathbf{L} K^\vee) = \Ext_A^0(K, L)
$$
이 존재한다.
\end{lemma}

\begin{proof}
$K$를 유한 사영 $A$-가군들로 이루어진 유계 복합체 $K^\bullet$로
나타낼 수 있다. 보조정리 \ref{more-algebra-lemma-RHom-out-of-projective}에 의해
대상 $K^\vee$는 복합체
$E^\bullet = \Hom^\bullet(K^\bullet, A)$로 나타내어진다.
$E^n = \Hom_A(K^{-n}, A)$도 유한 사영임에 유의하라. 따라서
$E^\bullet$은 유한 사영 가군들로 이루어진 유계 복합체이고,
$K^\vee$가 퍼펙트라고 결론짓는다. 각 차수에서 부호를 제외하면 평가
사상을 사용하는 표준 사상
$$
K^\bullet = \text{Tot}(\Hom^\bullet(A, A) \otimes_A K^\bullet)
\longrightarrow
\Hom^\bullet(\Hom^\bullet(K^\bullet, A), A)
$$
이 있다. 보조정리 \ref{more-algebra-lemma-evaluate-and-more}를 보라. 부호 규칙은
절 \ref{more-algebra-section-sign-rules}을 보라. 유한 사영 가군의 이중 쌍대는
자기 자신이므로 이 사상은 표준 동형사상
$(K^\vee)^\vee \cong K$를 정의한다.

\medskip\noindent
두 번째 등식은 첫 번째 등식, 보조정리 \ref{more-algebra-lemma-internal-hom},
유도 범주, 보조정리
\ref{derived-lemma-morphisms-from-projective-complex}, 그리고 Ext 군의
정의로부터 따른다. 유도 범주, 절 \ref{derived-section-ext}를 보라.
$L$을 나타내는 $A$-가군 복합체를 $L^\bullet$이라 하자.
절 \ref{more-algebra-section-sign-rules}의 항목 (\ref{more-algebra-item-compatible})에 의해
$A$-가군 복합체의 표준 동형사상
$$
\text{Tot}(L^\bullet \otimes_A E^\bullet)
\longrightarrow
\Hom^\bullet(K^\bullet, L^\bullet)
$$
이 있다. 여기서는 $E^\bullet$이 복합체 범주에서 $K^\bullet$의
왼쪽 쌍대임을 알기 위해 보조정리 \ref{more-algebra-lemma-left-dual-complex}를
사용한다. 이로써 첫 번째 표시된 등식이 증명되고 증명이 끝난다.
\end{proof}

\begin{remark}
\label{more-algebra-remark-dual-perfect-complex}
$A$를 환이라 하자. $K \in D(A)$가 퍼펙트라고 하자. $K$의 Tor
진폭이 $[a, b]$ 안에 있으면 그 쌍대 $K^\vee$의 Tor 진폭은
$[-b, -a]$ 안에 있다. 이는 보조정리 \ref{more-algebra-lemma-perfect}의 마지막
명제와 보조정리 \ref{more-algebra-lemma-dual-perfect-complex}의 증명에 주어진
$K^\vee$의 설명에서 따른다.
\end{remark}

\begin{lemma}
\label{more-algebra-lemma-colim-and-lim-of-duals}
\begin{slogan}
퍼펙트 대상들의 계에 대한 자명한 쌍대성.
\end{slogan}
$A$를 환이라 하자. $(K_n)_{n \in \mathbf{N}}$을 $D(A)$의 퍼펙트
대상들의 계라 하자. $K = \text{hocolim} K_n$을 유도 여극한이라 하자
(유도 범주, 정의 \ref{derived-definition-derived-colimit}).
그러면 $D(A)$의 임의의 대상 $E$에 대하여
$$
R\Hom_A(K, E) = R\lim E \otimes^\mathbf{L}_A K_n^\vee
$$
이다. 여기서 $(K_n^\vee)$은 쌍대 퍼펙트 복합체들의 역계이다.
\end{lemma}

\begin{proof}
보조정리 \ref{more-algebra-lemma-dual-perfect-complex}에 의해
$R\lim E \otimes^\mathbf{L}_A K_n^\vee =
R\lim R\Hom_A(K_n, E)$이고, 이는 특수 삼각형
$$
R\lim R\Hom_A(K_n, E) \to
\prod R\Hom_A(K_n, E) \to
\prod R\Hom_A(K_n, E)
$$
에 들어간다. $K$도 마찬가지로 특수 삼각형
$\bigoplus K_n \to \bigoplus K_n \to K$에 들어가므로
$\prod R\Hom_A(K_n, E) = R\Hom_A(\bigoplus K_n, E)$임을 보이면
충분하다. 이는 (\ref{more-algebra-equation-internal-hom})과 유도 텐서곱이 직합과
교환한다는 사실의 형식적 귀결이다.
\end{proof}

\begin{lemma}
\label{more-algebra-lemma-colimit-perfect-complexes}
$R = \colim_{i \in I} R_i$를 환들의 여과 여극한이라 하자.
\begin{enumerate}
\item $D(R)$의 퍼펙트 $K$가 주어지면 어떤 $i \in I$와
$D(R_i)$의 퍼펙트 $K_i$가 존재하여 $D(R)$에서
$K \cong K_i \otimes_{R_i}^\mathbf{L} R$이다.
\item $0 \in I$와 $K_0, L_0 \in D(R_0)$가 주어지고 $K_0$가
퍼펙트이면
$$
\Hom_{D(R)}(K_0 \otimes_{R_0}^\mathbf{L} R, L_0 \otimes_{R_0}^\mathbf{L} R) =
\colim_{i \geq 0}
\Hom_{D(R_i)}(K_0 \otimes_{R_0}^\mathbf{L} R_i,
L_0 \otimes_{R_0}^\mathbf{L} R_i)
$$
이다.
\end{enumerate}
다시 말해 $R$ 위의 퍼펙트 복합체들의 삼각 범주는 $R_i$ 위의 퍼펙트
복합체들의 삼각 범주들의 여극한이다.
\end{lemma}

\begin{proof}
대수학, 보조정리
\ref{algebra-lemma-module-map-property-in-colimit}와
\ref{algebra-lemma-colimit-category-fp-modules}의 결과를 더 언급하지
않고 사용하겠다. 특히 이 보조정리들은 유한 표시 $R$-가군의 범주가
유한 표시 $R_i$-가군 범주들의 여극한임을 말한다. 유한 사영 가군은
유한 자유 가군의 직합인자로 특징지을 수 있으므로(대수학, 보조정리
\ref{algebra-lemma-finite-projective}), 유한 사영 가군의 범주에도
같은 사실이 성립함을 알 수 있다. $D(R)$의 퍼펙트 대상에 대한 우리의
정의로부터 (1)이 증명된다.

\medskip\noindent
(2)를 증명하기 위하여 $K_0$를 유한 사영 $R_0$-가군들로 이루어진
유계 복합체 $K_0^\bullet$로 나타낼 수 있다. $L_0$를 K-평탄 복합체
$L_0^\bullet$로 나타낼 수 있다(보조정리
\ref{more-algebra-lemma-K-flat-resolution}). 그러면 유도 범주, 보조정리
\ref{derived-lemma-morphisms-from-projective-complex}에 의해
$$
\Hom_{D(R)}(K_0 \otimes_{R_0}^\mathbf{L} R, L_0 \otimes_{R_0}^\mathbf{L} R) =
\Hom_{K(R)}(K_0^\bullet \otimes_{R_0} R, L_0^\bullet \otimes_{R_0} R)
$$
이다. $R$을 $R_i$로 바꾼 $\Hom$에 대해서도 마찬가지이다.
오른쪽 변에는 유한 개의 항만 관여하고, 증명 처음에 인용한
보조정리들에 의해
$$
\Hom_R(K_0^p \otimes_{R_0} R, L_0^q \otimes_{R_0} R) =
\colim_{i \geq 0}
\Hom_{R_i}(K_0^p \otimes_{R_0} R_i, L_0^q \otimes_{R_0} R_i)
$$
이며, 여과 여극한은 완전하므로(대수학, 보조정리
\ref{algebra-lemma-directed-colimit-exact}) (2)도 성립한다고
결론짓는다.
\end{proof}

\section{복합체의 올림}
\label{more-algebra-section-lifting-complexes}

\noindent
$R$을 환이라 하고 $I \subset R$을 아이디얼이라 하자. 우리가 생각할
올림 문제는 다음과 같다. $D(R)$의 대상 $K$와 $R/I$-가군 복합체
$E^\bullet$이 주어졌고, $E^\bullet$이
$K \otimes_R^\mathbf{L} R/I$를 $D(R)$에서 나타낸다고 하자.
질문은 다음과 같다.
$K$를 $D(R)$에서 나타내고 $E^\bullet$을 올리는 $R$-가군 복합체
$P^\bullet$이 존재하는가? 일반적으로 이 질문의 답은 아니지만,
좋은 경우에는 무언가를 할 수 있다. 먼저 비순환 복합체의 올림을
논의하자.

\begin{lemma}
\label{more-algebra-lemma-lift-acyclic-complex}
$R$을 환이라 하고 $I \subset R$을 아이디얼이라 하자.
$\mathcal{P}$를 $R$-가군들의 한 모임이라 하자. 다음을 가정하자.
\begin{enumerate}
\item 각 $P \in \mathcal{P}$는 사영 $R$-가군이다.
\item $P_1 \in \mathcal{P}$이고
$P_1 \oplus P_2 \in \mathcal{P}$이면 $P_2 \in \mathcal{P}$이다.
\item
% P02-E0976: frozen source attaches “is surjective modulo I” to a malformed compound subject; Korean renders the intended hypothesis.
$f : P_1 \to P_2$, $P_1, P_2 \in \mathcal{P}$가 주어지고 이
사상이 $I$로 나눈 뒤 전사이면 $f$는 전사이다.
\end{enumerate}
그러면 각 항이 어떤 $P \in \mathcal{P}$에 대한 $P/IP$ 꼴인 임의의
위로 유계인 비순환 복합체 $E^\bullet$에 대하여, 항들이
$\mathcal{P}$에 속하고 $E^\bullet$을 올리는 위로 유계인 비순환
복합체 $P^\bullet$이 존재한다.
\end{lemma}

\begin{proof}
$i > b$이면 $E^i = 0$이라고 하자. $n$과 복합체의 사상
$$
\xymatrix{
& & P^n \ar[r] \ar[d] & P^{n + 1} \ar[r] \ar[d] & \ldots \ar[r]
& P^b \ar[r] \ar[d] & 0 \ar[r] \ar[d] & \ldots \\
\ldots \ar[r] & E^{n - 1} \ar[r] &
E^n \ar[r] & E^{n + 1} \ar[r] & \ldots \ar[r] &
E^b \ar[r] & 0 \ar[r] & \ldots
}
$$
이 주어졌다고 가정하자. 여기서 $P^i \in \mathcal{P}$이고,
$P^n \to P^{n + 1} \to \ldots \to P^b$는 차수
$\geq n + 1$에서 비순환이며, 수직 사상들은 동형사상
$P^i/IP^i \to E^i$를 유도한다. 이 상황에서는 동형사상
$P^i = Z^i \oplus Z^{i + 1}$을 귀납적으로 택하여 사상
$P^i \to P^{i + 1}$이
$Z^i \oplus Z^{i + 1} \to Z^{i + 1} \to
Z^{i + 1} \oplus Z^{i + 2}$로 주어지게 할 수 있다. 성질 (2)와
귀납 논증에 의해 $Z^i \in \mathcal{P}$임을 알 수 있다.
$P^{n - 1} \in \mathcal{P}$와 동형사상
$P^{n - 1}/IP^{n - 1} \to E^{n - 1}$을 택하자.
$P^{n - 1}$은 사영이고
$Z^n/IZ^n = \Im(E^{n - 1} \to E^n)$이므로 사상
$P^{n - 1} \to E^{n - 1} \to E^n$을 사상
$P^{n - 1} \to Z^n$으로 올릴 수 있다. 성질 (3)에 의해
$P^{n - 1} \to Z^n$은 전사이다. 따라서 방금 구성한
$P^{n - 1}$과 사상들을 $P^n$의 왼쪽에 더하여 그림을 확장할 수
있다. $n > b$에 대해서는 원하는 꼴의 그림이 존재하므로 $n$에 대한
귀납법으로 결론을 얻는다.
\end{proof}

\begin{lemma}
\label{more-algebra-lemma-lift-complex}
$R$을 환이라 하고 $I \subset R$을 아이디얼이라 하자.
$\mathcal{P}$를 $R$-가군들의 한 모임이라 하자.
$K \in D(R)$이고, $K \otimes_R^\mathbf{L} R/I$를 나타내는
$R/I$-가군 복합체를 $E^\bullet$이라 하자. 다음을 가정하자.
\begin{enumerate}
\item 각 $P \in \mathcal{P}$는 사영 $R$-가군이다.
\item $P_1 \in \mathcal{P}$이고
$P_1 \oplus P_2 \in \mathcal{P}$일 필요충분조건은
$P_1, P_2 \in \mathcal{P}$인 것이다.
\item $f : P_1 \to P_2$, $P_1, P_2 \in \mathcal{P}$가 주어지고
이 사상이 $I$로 나눈 뒤 전사이면 $f$는 전사이다.
\item $E^\bullet$은 위로 유계이고, $E^i$는 어떤
$P \in \mathcal{P}$에 대한 $P/IP$ 꼴이다.
\item $K$는 항들이 $\mathcal{P}$에 속하는 위로 유계인 복합체로
나타낼 수 있다.
\end{enumerate}
그러면 항들이 $\mathcal{P}$에 속하고,
% P02-E0977: frozen source omits “such that” before P/IP isomorphic to E; Korean renders the intended relative clause.
$P^\bullet/IP^\bullet$이 $E^\bullet$과 동형이며 $K$를 $D(R)$에서
나타내는 위로 유계인 복합체 $P^\bullet$이 존재한다.
\end{lemma}

\begin{proof}
가정 (5)에 의해 $K$를 항들이 $\mathcal{P}$에 속하는 위로 유계인
복합체 $K^\bullet$로 나타낼 수 있다. 그러면
$K \otimes_R^\mathbf{L} R/I$는 $K^\bullet/IK^\bullet$로
나타내어진다. (4)에 의해 $E^\bullet$은 사영 $R/I$-가군들로 이루어진
위로 유계인 복합체이므로 준동형사상
$\delta : E^\bullet \to K^\bullet/IK^\bullet$을 택할 수 있다
(유도 범주, 보조정리
\ref{derived-lemma-morphisms-from-projective-complex}).
% P02-E0978: frozen source says “be cone on delta”, omitting the article; Korean renders the intended construction.
$C^\bullet$을 $\delta$의 원뿔이라 하자
(유도 범주, 정의 \ref{derived-definition-cone}).
가군 $C^i$는 직합 $K^i/IK^i \oplus E^{i + 1}$이므로 어떤
$P \in \mathcal{P}$에 대한 $P/IP$ 꼴이다. 실제로 (2)는 특히
$\mathcal{P}$가 직합을 취할 때 보존됨을 말한다.
$C^\bullet$은 비순환이므로 보조정리
\ref{more-algebra-lemma-lift-acyclic-complex}을 적용하여
% P02-E0979: frozen source says “find a acyclic lift”; Korean renders the intended phrase.
$C^\bullet$의 비순환 올림 $A^\bullet$을 찾을 수 있다.
복합체 $A^\bullet$은 위로 유계이고 그 항들은 $\mathcal{P}$에 속한다.
그림
$$
\xymatrix{
K^\bullet \ar@{..>}[r] \ar[d] & A^\bullet \ar[d] \\
K^\bullet/IK^\bullet \ar[r] & C^\bullet \ar[r] & E^\bullet[1]
}
$$
에서 그림을 가환이 되게 하는 점선 화살표를 유도 범주, 보조정리
\ref{derived-lemma-morphisms-lift-projective}에 의해 찾을 수 있다.
아래에서 (1), (2), (3)으로부터 모든 $i$에 대하여
$K^i \to A^i$가 직합인자의 포함이라는 사실이 따름을 보이겠다.
성질 (2)에 의해 $P^i = \Coker(K^i \to A^i)$가
$\mathcal{P}$에 속함을 알 수 있다. 따라서
$P^\bullet = \Coker(K^\bullet \to A^\bullet)[-1]$로 두어 결론을
얻을 수 있다.

\medskip\noindent
증명을 끝내려면 다음을 보여야 한다. $f : P_1 \to P_2$이고
$P_1, P_2 \in \mathcal{P}$이며
$P_1/IP_1 \to P_2/IP_2$가 분할 단사이고 그 여핵이 어떤
$P_3 \in \mathcal{P}$에 대한 $P_3/IP_3$ 꼴이면, $f$는 분할
단사이다. $E_i = P_i/IP_i$로 쓰자. 그러면
$E_2 = E_1 \oplus E_3$이다. $P_2$가 사영이므로 사상
$E_2 \to E_3$을 올리는 사상 $g : P_2 \to P_3$을 택할 수 있다.
조건 (3)에 의해 사상 $g$는 전사이고, $P_3$가 사영이므로 분할된다.
% P02-E0259: frozen source defines P_1' and immediately switches to P'_1 for the same kernel; Korean preserves both formal source forms.
$P_1' = \Ker(g)$로 두고 분할 $P_2 = P'_1 \oplus P_3$을 택하자.
(2)에 의해 $P'_1 \in \mathcal{P}$이다. $g \circ f = 0$인지는
알 수 없지만 사상
$$
P_1 \xrightarrow{f} P_2 \xrightarrow{projection} P'_1
$$
을 생각할 수 있다. $I$로 나눈 합성은 동형사상이다.
$P'_1$은 사영이므로 $P_1 = T \oplus P'_1$로 분할할 수 있다.
$T = 0$이면 끝이다. 이 경우 $P_2 \to P'_1$이 $f$의 분할이기
때문이다. (2)에 의해 $T \in \mathcal{P}$임을 알 수 있다.
$I$로 나누어 계산하면 $T/IT = 0$이다.
$0 \in \mathcal{P}$이므로($\mathcal{P}$의 임의의 $P$의 직합인자이기
때문이다) 사상 $0 \to T$는 전사임을 알 수 있고, 원하는 대로
$T = 0$이라고 결론짓는다.
\end{proof}

\begin{lemma}
\label{more-algebra-lemma-lift-complex-projectives}
$R$을 환이라 하고 $I \subset R$을 아이디얼이라 하자.
$E^\bullet$을 $R/I$-가군 복합체라 하고 $K$를 $D(R)$의 대상이라 하자.
다음을 가정하자.
\begin{enumerate}
\item $E^\bullet$은 사영 $R/I$-가군들로 이루어진 위로 유계인
복합체이다.
\item $K \otimes_R^\mathbf{L} R/I$는 $D(R/I)$에서
$E^\bullet$로 나타내어진다.
\item $I$는 멱영 아이디얼이다.
\end{enumerate}
그러면 $K$를 $D(R)$에서 나타내는 사영 $R$-가군들의 위로 유계인
복합체 $P^\bullet$이 존재하여 $P^\bullet \otimes_R R/I$는
$E^\bullet$과 동형이다.
\end{lemma}

\begin{proof}
모든 사영 $R$-가군의 모임 $\mathcal{P}$를 사용하여 보조정리
\ref{more-algebra-lemma-lift-complex}을 적용한다. 이 보조정리의 성질 (1)과 (2)는
즉시 성립한다. 성질 (3)은 나카야마 보조정리(대수학, 보조정리
\ref{algebra-lemma-NAK})에서 따른다. 성질 (4)는 사영 $R/I$-가군을
사영 $R$-가군으로 올릴 수 있다는 사실에서 따른다. 대수학, 보조정리
\ref{algebra-lemma-lift-projective-module}을 보라. (5)가 성립함을
보이려면 $K$가 $D^{-}(R)$에 속함을 보이면 충분하다.
$E^\bullet$이 위로 유계이므로
$K \otimes_R^\mathbf{L} R/I$가 $D^{-}(R/I)$에 속한다는 것이
주어져 있고, 따라서 이는 보조정리
\ref{more-algebra-lemma-check-tor-dimension-modulo-nilpotent}에서 따른다.
\end{proof}

\begin{lemma}
\label{more-algebra-lemma-check-pseudo-coherent-modulo-nilpotent}
$R' \to R$을 핵이 멱영 아이디얼인 전사 환 사상이라 하자.
$K' \in D(R')$이고 $K = K' \otimes_{R'}^\mathbf{L} R$로 두자.
그러면 $K$가 유사연접일 필요충분조건은 $K'$이 유사연접인 것이다.
\end{lemma}

\begin{proof}
한 방향은 보조정리 \ref{more-algebra-lemma-pull-pseudo-coherent}에서 따른다.
다른 방향을 위하여 $K$가 유사연접이라고 가정하자. 그러면 보조정리
\ref{more-algebra-lemma-pseudo-coherent}에 의해 $K$를 유한 자유 $R$-가군들로
이루어진 위로 유계인 복합체 $E^\bullet$로 나타낼 수 있다.
보조정리 \ref{more-algebra-lemma-lift-complex-projectives}에 의해
$P^n \otimes_{R'} R = E^n$을 만족하는 사영 $R'$-가군들의 위로
유계인 복합체 $P^\bullet$로 $K'$을 나타낼 수 있다. 나카야마
보조정리에 의해 $P^n$이 유한 자유임을 알 수 있고, 따라서 $K'$도
유사연접이라고 결론짓는다.
\end{proof}

\begin{lemma}
\label{more-algebra-lemma-lift-complex-stably-frees}
$R$을 환이라 하고 $I \subset R$을 아이디얼이라 하자.
$E^\bullet$을 $R/I$-가군 복합체라 하고 $K$를 $D(R)$의 대상이라 하자.
다음을 가정하자.
\begin{enumerate}
\item $E^\bullet$은 유한 안정 자유 $R/I$-가군들로 이루어진 위로
유계인 복합체이다.
\item $K \otimes_R^\mathbf{L} R/I$는 $D(R/I)$에서
$E^\bullet$로 나타내어진다.
\item
% P02-E0980: frozen source changes the introduced object K to undefined K^bullet in hypothesis (3); Korean preserves the formal source span.
$K^\bullet$은 유사연접이다.
\item $1 + I$의 모든 원소는 가역이다.
\end{enumerate}
그러면 $K$를 $D(R)$에서 나타내는 유한 안정 자유 $R$-가군들의 위로
유계인 복합체 $P^\bullet$이 존재하여
$P^\bullet \otimes_R R/I$는 $E^\bullet$과 동형이다. 더욱이
$E^i$가 자유이면 $P^i$도 자유이다.
\end{lemma}

\begin{proof}
모든 유한 안정 자유 $R$-가군의 모임 $\mathcal{P}$를 사용하여
보조정리 \ref{more-algebra-lemma-lift-complex}을 적용한다. 이 보조정리의 성질
(1)은 즉시 성립한다. 성질 (2)는 보조정리
\ref{more-algebra-lemma-exact-category-stably-free}에서 따른다. 성질 (3)은
나카야마 보조정리(대수학, 보조정리
\ref{algebra-lemma-NAK})에서 따른다. 성질 (4)는 유한 안정 자유
$R/I$-가군을 유한 안정 자유 $R$-가군으로 올릴 수 있다는 사실에서
따른다. 보조정리 \ref{more-algebra-lemma-lift-stably-free}를 보라.
% P02-E0981: frozen source switches from “Property (5)” to “Part (5)”; Korean uses the intended hypothesis terminology.
조건 (5)는 유사연접 복합체를 유한 자유 $R$-가군들로 이루어진 위로
유계인 복합체로 나타낼 수 있으므로 성립한다. 보조정리의 마지막
명제는 보조정리 \ref{more-algebra-lemma-isomorphic-finite-projective-lifts}에서
따른다.
\end{proof}

\begin{lemma}
\label{more-algebra-lemma-lift-pseudo-coherent-from-residue-field}
$(R, \mathfrak m, \kappa)$를 국소환이라 하고
$K \in D(R)$을 유사연접이라 하자.
$d_i = \dim_\kappa H^i(K \otimes_R^\mathbf{L} \kappa)$로 두자.
% P02-E0982: frozen source begins two consecutive conclusions with “Then”; Korean combines them without changing their content.
그러면 $d_i < \infty$이고 어떤 $b \in \mathbf{Z}$에 대하여
$i > b$이면 $d_i = 0$이다. 또한 $K$를 $D(R)$에서 나타내는 복합체
$$
\ldots \to
R^{\oplus d_{b - 2}} \to
R^{\oplus d_{b - 1}} \to
R^{\oplus d_b} \to 0 \to \ldots
$$
가 존재한다. 더욱이 이 복합체는 동형사상까지 유일하다(!).
\end{lemma}

\begin{proof}
$K \otimes_R^\mathbf{L} \kappa$는 $D(\kappa)$의 대상으로서
유사연접임에 유의하라. 보조정리
\ref{more-algebra-lemma-pull-pseudo-coherent}를 보라. 따라서 코호몰로지 공간들은
유한 차원이고 어떤 절단점 위에서 사라진다. $D(\kappa)$의 모든
대상은 $D(\kappa)$에서 미분이 영인 복합체 $E^\bullet$과 동형이다.
특히 $E^i \cong \kappa^{\oplus d_i}$는 유한 자유이다.
보조정리 \ref{more-algebra-lemma-lift-complex-stably-frees}를 적용하면 존재성을
얻는다.

\medskip\noindent
% P02-E0983: frozen source makes the “If we have two complexes ...” clause a sentence fragment; Korean renders one complete conditional sentence.
$F^i$와 $G^i$가 계수 $d_i$의 자유 가군이고 $K$를 나타내는 두
복합체 $F^\bullet$과 $G^\bullet$이 있다고 하자. 그러면 동형사상
$F^\bullet \cong K \cong G^\bullet$을 나타내는 복합체의 사상
$\beta : F^\bullet \to G^\bullet$를 택할 수 있다. 유도 범주,
보조정리 \ref{derived-lemma-morphisms-from-projective-complex}를 보라.
유도된 복합체의 사상
$\beta \otimes 1 : F^\bullet \otimes_R^\mathbf{L} \kappa \to
G^\bullet \otimes_R^\mathbf{L} \kappa$는 동형사상이어야 한다.
$F^\bullet \otimes_R^\mathbf{L} \kappa$와
$G^\bullet \otimes_R^\mathbf{L} \kappa$의 미분들이 영이기 때문이다.
따라서 $\beta^i : F^i \to G^i$는 $\mathfrak m$으로 나눈 축약이
동형사상인 유한 자유 $R$-가군의 사상이다. 그러므로 $\beta^i$는
동형사상이고 결론을 얻는다.
\end{proof}

\begin{lemma}
\label{more-algebra-lemma-lift-perfect-from-residue-field}
$R$을 환이라 하고 $\mathfrak p \subset R$을 소아이디얼이라 하자.
$K \in D(R)$이 퍼펙트라고 하자.
$d_i =
\dim_{\kappa(\mathfrak p)} H^i(K \otimes_R^\mathbf{L} \kappa(\mathfrak p))$로
두자.
% P02-E0984: frozen source begins two consecutive conclusions with “Then”; Korean combines them without changing their content.
그러면 $d_i < \infty$이고 영이 아닌 것은 유한 개뿐이다. 또한
$f \in R$, $f \not \in \mathfrak p$인 원소와
$K \otimes_R^\mathbf{L} R_f$를 $D(R_f)$에서 나타내는 복합체
$$
\ldots \to 0 \to R_f^{\oplus d_a} \to R_f^{\oplus d_{a + 1}} \to
\ldots \to
R_f^{\oplus d_{b - 1}} \to
R_f^{\oplus d_b} \to 0
\to \ldots
$$
가 존재한다.
\end{lemma}

\begin{proof}
$K \otimes_R^\mathbf{L} \kappa(\mathfrak p)$는
$D(\kappa(\mathfrak p))$의 대상으로서 퍼펙트임에 유의하라.
보조정리 \ref{more-algebra-lemma-pull-perfect}를 보라. 따라서 영이 아닌
$d_i$는 유한 개뿐이고 모두 유한하다. 보조정리
\ref{more-algebra-lemma-lift-pseudo-coherent-from-residue-field}를 적용하면 국소환
$R_\mathfrak p$ 위에서 원하는 꼴을 갖고 $K$를 나타내는 복합체를
얻는다. $f \in R$, $f \not \in \mathfrak p$에 대하여
$R_\mathfrak p = \colim R_f$이다(대수학, 보조정리
\ref{algebra-lemma-localization-colimit}). 보조정리
\ref{more-algebra-lemma-colimit-perfect-complexes}에 의해 결론을 얻는다.
몇 가지 세부사항은 생략한다.
\end{proof}

\begin{lemma}
\label{more-algebra-lemma-compare-representatives-perfect}
$R$을 환이라 하고 $\mathfrak p \subset R$을 소아이디얼이라 하자.
$M^\bullet$과 $N^\bullet$을 $D(R)$의 같은 대상을 나타내는 유한 사영
$R$-가군들의 유계 복합체라 하자. 그러면 어떤
$f \in R$, $f \not \in \mathfrak p$가 존재하여 복합체의
동형사상(!)
$$
M^\bullet_f \oplus P^\bullet \cong N^\bullet_f \oplus Q^\bullet
$$
이 존재한다. 여기서 $P^\bullet$과 $Q^\bullet$은 자명한 복합체,
즉 임의의 차수에 놓인
% P02-E0985: frozen source repeats “of the form the form”; Korean renders the intended single phrase.
$\ldots \to 0 \to R_f \xrightarrow{1} R_f \to 0 \to \ldots$ 꼴의
복합체들의 유한 직합이다.
\end{lemma}

\begin{proof}
국소화 $R_\mathfrak p$ 위에서 설명한 유형의 동형사상이 있다면,
$R_\mathfrak p = \colim R_f$라는 사실(대수학, 보조정리
\ref{algebra-lemma-localization-colimit})을 사용하여 어떤 $f$에 대해
이 동형사상을 $R_f$ 위의 동형사상으로 내릴 수 있다. 따라서
$R$이 국소환이고 $\mathfrak p$가 극대 아이디얼이라고 가정할 수 있다.
이 경우 결과는 퍼펙트 대상을 나타내는 “극소” 복합체의 유일성과
(보조정리 \ref{more-algebra-lemma-lift-pseudo-coherent-from-residue-field}), 모든
복합체가 자명한 복합체와 극소 복합체의 직합이라는 사실
(대수학, 보조정리 \ref{algebra-lemma-add-trivial-complex})에서 따른다.
\end{proof}

\begin{lemma}
\label{more-algebra-lemma-lift-complex-finite-projectives}
$R$을 환이라 하고 $I \subset R$을 아이디얼이라 하자.
$E^\bullet$을 $R/I$-가군 복합체라 하고 $K$를 $D(R)$의 대상이라 하자.
다음을 가정하자.
\begin{enumerate}
\item $E^\bullet$은 유한 사영 $R/I$-가군들로 이루어진 위로 유계인
복합체이다.
\item $K \otimes_R^\mathbf{L} R/I$는 $D(R/I)$에서
$E^\bullet$로 나타내어진다.
\item $K$는 유사연접이다.
\item $(R, I)$는 헨젤 쌍이다.
\end{enumerate}
그러면 $K$를 $D(R)$에서 나타내는 유한 사영 $R$-가군들의 위로 유계인
복합체 $P^\bullet$이 존재하여 $P^\bullet \otimes_R R/I$는
$E^\bullet$과 동형이다. 더욱이 $E^i$가 자유이면 $P^i$도 자유이다.
\end{lemma}

\begin{proof}
모든 유한 사영 $R$-가군의 모임 $\mathcal{P}$를 사용하여 보조정리
\ref{more-algebra-lemma-lift-complex}을 적용한다. 이 보조정리의 성질 (1)과 (2)는
즉시 성립한다. 성질 (3)은 나카야마 보조정리(대수학, 보조정리
\ref{algebra-lemma-NAK})에서 따른다. 성질 (4)는 유한 사영
$R/I$-가군을 유한 사영 $R$-가군으로 올릴 수 있다는 사실에서
따른다. 보조정리 \ref{more-algebra-lemma-lift-finite-projective-module}을 보라.
성질 (5)는 유사연접 복합체를 유한 자유 $R$-가군들로 이루어진 위로
유계인 복합체로 나타낼 수 있으므로 성립한다. 따라서 보조정리
\ref{more-algebra-lemma-lift-complex}을 적용할 수 있고 원하는 $P^\bullet$을
찾는다. 보조정리의 마지막 명제는 보조정리
\ref{more-algebra-lemma-isomorphic-finite-projective-lifts}에서 따른다.
\end{proof}

\section{복합체의 분할}
\label{more-algebra-section-splitting}

\noindent
이 절에서는 환의 유도 범주의 대상이 그 절단들의 직합이 되게 하는
조건들을 논의한다. 우리의 방법은 적절한 가정 아래 다음 보조정리를
사용하여 표준 특수 삼각형
$$
\tau_{\leq i}K \to K \to \tau_{\geq i + 1}K \to (\tau_{\leq i}K)[1]
$$
을 $D(R)$에서 분할하는 것이다. 유도 범주, 주석
\ref{derived-remark-truncation-distinguished-triangle}을 보라.

\begin{lemma}
\label{more-algebra-lemma-splitting-unique}
$R$을 환이라 하고 $K$와 $L$을 $D(R)$의 대상이라 하자.
$L$의 사영 진폭이 $[a, b]$ 안에 있다고 가정하자. 예를 들어
$L$이 Tor 진폭 $[a, b]$를 갖는 퍼펙트 대상이면 이 가정이 성립한다.
\begin{enumerate}
\item $i \geq a$에 대하여 $H^i(K) = 0$이면
$\Hom_{D(R)}(L, K) = 0$이다.
\item $i \geq a + 1$에 대하여 $H^i(K) = 0$이면, 임의의 특수 삼각형
$K \to M \to L \to K[1]$에 대하여 그 특수 삼각형의 사상들과
양립하는 $D(R)$의 동형사상 $M \cong K \oplus L$이 존재한다.
\item $i \geq a$에 대하여 $H^i(K) = 0$이면 (2)의 동형사상은
존재하며 유일하다.
\end{enumerate}
\end{lemma}

\begin{proof}
$L$의 사영 진폭이 $[a, b]$ 안에 있다는 가정은
$i \not \in [a, b]$에 대하여 $L^i = 0$인 사영 $R$-가군들의 복합체
$L^\bullet$로 $L$을 나타낼 수 있다는 뜻이다. 정의
\ref{more-algebra-definition-projective-dimension}을 보라. $L$이 Tor 진폭
$[a, b]$를 갖는 퍼펙트 대상이면 $i \not \in [a, b]$에 대하여
$L^i = 0$인 유한 사영 $R$-가군들의 복합체 $L^\bullet$로 $L$을
나타낼 수 있다. 보조정리 \ref{more-algebra-lemma-perfect}를 보라.
$i \geq a$에 대하여 $H^i(K) = 0$이면 $K$는
$\tau_{\leq a - 1}K$와 준동형이다. 따라서 $i \geq a$에 대하여
$K^i = 0$인 $R$-가군 복합체 $K^\bullet$로 $K$를 나타낼 수 있다.
그러면 유도 범주, 보조정리
\ref{derived-lemma-morphisms-from-projective-complex}에 의해
$$
\Hom_{D(R)}(L, K) = \Hom_{K(R)}(L^\bullet, K^\bullet) = 0
$$
을 얻는다. 이로써 (1)이 증명된다. (2)의 가정 아래에서는 (1)에 의해
$\Hom_{D(R)}(L, K[1]) = 0$이므로 유도 범주, 보조정리
\ref{derived-lemma-split}에 의해 특수 삼각형이 분할된다. (3)의
유일성은 (1)에서 따른다.
\end{proof}

\begin{lemma}
\label{more-algebra-lemma-better-cut-complex-in-two}
$R$을 환이라 하고 $\mathfrak p \subset R$을 소아이디얼이라 하자.
$K^\bullet$을 유사연접 $R$-가군 복합체라 하자. 어떤
$i \in \mathbf{Z}$에 대하여 사상
$$
H^i(K^\bullet) \otimes_R \kappa(\mathfrak p)
\longrightarrow
H^i(K^\bullet \otimes_R^{\mathbf{L}} \kappa(\mathfrak p))
$$
이 전사라고 가정하자. 그러면 어떤 $f \in R$,
$f \not \in \mathfrak p$가 존재하여
$\tau_{\geq i + 1}(K^\bullet \otimes_R R_f)$는 Tor 진폭
$[i + 1, \infty]$를 갖는 $D(R_f)$의 퍼펙트 대상이고, $D(R_f)$에서
표준 동형사상
$$
K^\bullet \otimes_R R_f \cong
\tau_{\leq i}(K^\bullet \otimes_R R_f) \oplus
\tau_{\geq i + 1}(K^\bullet \otimes_R R_f)
$$
이 존재한다.
\end{lemma}

\begin{proof}
이 증명에서 모든 텐서곱은 $R$ 위에서 취하며
$\kappa = \kappa(\mathfrak p)$로 쓴다. $K^\bullet$이 유한 자유
$R$-가군들로 이루어진 위로 유계인 복합체라고 가정할 수 있다.
차수 $i$에서 무슨 일이 일어나는지 살펴보자.
$$
\ldots \to K^{i - 1} \xrightarrow{d^{i - 1}} K^i \xrightarrow{d^i}
K^{i + 1} \to \ldots
$$
$0 \subset V \subset W \subset K^i \otimes \kappa$를 식
$$
V = \Im\left(
K^{i - 1} \otimes \kappa \to K^i \otimes \kappa
\right)
\quad\text{and}\quad
W = \Ker\left(
K^i \otimes \kappa \to K^{i + 1} \otimes \kappa
\right)
$$
으로 정의하자.
$\dim(V) = r$, $\dim(W/V) = s$,
$\dim(K^i \otimes \kappa/W) = t$로 두자. $d^{i - 1}$에 의해
$V$의 기저로 가는 $x_1, \ldots, x_r \in K^{i - 1}$을 택할 수 있다.
우리의 가정에 의해 $W/V$의 기저로 가는
$y_1, \ldots, y_s \in \Ker(d^i)$를 택할 수 있다. 마지막으로
$K^i \otimes \kappa/W$의 기저로 가는
$z_1, \ldots, z_t \in K^i$를 택하자. 그러면 원소
$d^i(z_1), \ldots, d^i(z_t) \in K^{i + 1}$이
$K^{i + 1} \otimes_R \kappa$에서 선형 독립임을 알 수 있다.
대수학, 보조정리 \ref{algebra-lemma-cokernel-flat}에 의해 어떤
$f \in R$, $f \not \in \mathfrak p$에 대하여 $R$을 $R_f$로
바꾼 뒤 다음을 가정할 수 있다.
\begin{enumerate}
\item
% P02-E0986: frozen source changes d^{i-1}(x_a) to d^i(x_a) in the claimed basis; Korean preserves the formal source span.
$d^i(x_a), y_b, z_c$는 $K^i$의 $R$-기저이다.
\item $d^i(z_1), \ldots, d^i(z_t)$는 $K^{i + 1}$에서
$R$-선형 독립이다.
\item 몫
$E^{i + 1} = K^{i + 1}/\sum Rd^i(z_c)$는 유한 사영이다.
\end{enumerate}
$d^i$는 $d^{i - 1}(x_a)$와 $y_b$를 소멸시키므로 조건 (2)로부터
$E^{i + 1} = \Coker(d^i : K^i \to K^{i + 1})$임을 얻는다.
따라서
$$
\tau_{\geq i + 1}K^\bullet =
(\ldots \to 0 \to E^{i + 1} \to K^{i + 2} \to \ldots)
$$
는 어떤 $b$에 대하여 차수 $[i + 1, b]$에 놓인 유한 사영 가군들의
유계 복합체이다. 따라서 $\tau_{\geq i + 1}K^\bullet$은 진폭
$[i + 1, b]$를 갖는 퍼펙트 대상이다.
$\tau_{\leq i}K^\bullet$은 차수 $> i$에서 코호몰로지를 갖지
않으므로 특수 삼각형
$$
\tau_{\leq i}K^\bullet \to K^\bullet \to \tau_{\geq i + 1}K^\bullet \to
(\tau_{\leq i}K^\bullet)[1]
$$
에 보조정리 \ref{more-algebra-lemma-splitting-unique}를 적용하여 결론을 얻는다
(유도 범주, 주석
\ref{derived-remark-truncation-distinguished-triangle}).
\end{proof}

\begin{lemma}
\label{more-algebra-lemma-isolate-a-cohomology-group}
$R$을 환이라 하고 $\mathfrak p \subset R$을 소아이디얼이라 하자.
$K^\bullet$을 유사연접 $R$-가군 복합체라 하자. 어떤
$i \in \mathbf{Z}$에 대하여 사상들
$$
H^i(K^\bullet) \otimes_R \kappa(\mathfrak p)
\longrightarrow
H^i(K^\bullet \otimes_R^{\mathbf{L}} \kappa(\mathfrak p))
\quad\text{and}\quad
H^{i - 1}(K^\bullet) \otimes_R \kappa(\mathfrak p)
\longrightarrow
H^{i - 1}(K^\bullet \otimes_R^{\mathbf{L}} \kappa(\mathfrak p))
$$
이 전사라고 가정하자. 그러면 어떤 $f \in R$,
$f \not \in \mathfrak p$가 존재하여 다음이 성립한다.
\begin{enumerate}
\item $\tau_{\geq i + 1}(K^\bullet \otimes_R R_f)$는 Tor 진폭
$[i + 1, \infty]$를 갖는 $D(R_f)$의 퍼펙트 대상이다.
\item $H^i(K^\bullet)_f$는 유한 자유 $R_f$-가군이다.
\item $D(R_f)$에서 표준 직합 분해
$$
K^\bullet \otimes_R R_f \cong
\tau_{\leq i - 1}(K^\bullet \otimes_R R_f) \oplus
H^i(K^\bullet)_f[-i] \oplus
\tau_{\geq i + 1}(K^\bullet \otimes_R R_f)
$$
가 존재한다.
\end{enumerate}
\end{lemma}

\begin{proof}
보조정리 \ref{more-algebra-lemma-better-cut-complex-in-two}에서 (1)과 $D(R_f)$의
분할
$K^\bullet \otimes_R R_f = \tau_{\leq i}K^\bullet \otimes_R R_f \oplus
\tau_{\geq i + 1}K^\bullet \otimes_R R_f$을 얻는다.
$\tau_{\leq i}K^\bullet \otimes_R R_f$에 보조정리
\ref{more-algebra-lemma-better-cut-complex-in-two}를 한 번 더 적용하면
($f$를 적절히 택한 뒤) $D(R_f)$의 분할
% P02-E0987: frozen source omits the [-i] shift and then changes H^i(K^bullet)_f to undefined H^i(K)_f; Korean preserves both formal spans.
$\tau_{\leq i}K^\bullet \otimes_R R_f =
\tau_{\leq i - 1}K^\bullet \otimes_R R_f \oplus H^i(K^\bullet)_f$을
얻는다. 또한 $H^i(K)_f$가 평탄한 퍼펙트 가군, 즉 유한 사영이라는
결론도 얻는다.
% P02-E0988: statement (2) claims finite free, while the frozen proof stops at finite projective; Korean preserves the proof's endpoint.
\end{proof}

\begin{lemma}
\label{more-algebra-lemma-cut-complex-in-two}
$R$을 환이라 하고 $\mathfrak p \subset R$을 소아이디얼이라 하자.
$i \in \mathbf{Z}$라 하자. $K^\bullet$을
$H^i(K^\bullet \otimes_R^{\mathbf{L}} \kappa(\mathfrak p)) = 0$을
만족하는 유사연접 $R$-가군 복합체라 하자. 그러면 어떤
$f \in R$, $f \not \in \mathfrak p$가 존재하고, $D(R_f)$에서
표준 직합 분해
$$
K^\bullet \otimes_R R_f =
\tau_{\geq i + 1}(K^\bullet \otimes_R R_f) \oplus
\tau_{\leq i - 1}(K^\bullet \otimes_R R_f)
$$
가 존재한다. 여기서 $\tau_{\geq i + 1}(K^\bullet \otimes_R R_f)$는
Tor 진폭 $[i + 1, \infty]$를 갖는 퍼펙트 복합체이다.
\end{lemma}

\begin{proof}
이는 자주 사용하는 보조정리 \ref{more-algebra-lemma-better-cut-complex-in-two}의
특수한 경우이다. 직접적인 증명은 다음과 같다. $K^\bullet$이 유한
자유 $R$-가군들로 이루어진 위로 유계인 복합체라고 가정할 수 있다.
차수 $i$에서 무슨 일이 일어나는지 살펴보자.
$$
\ldots \to K^{i - 2} \to R^{\oplus l}
\to R^{\oplus m} \to R^{\oplus n} \to K^{i + 2} \to \ldots
$$
$A$를 $K^{i - 1} \to K^i$에 대응하는 $m \times l$ 행렬이라 하고,
$B$를 $K^i \to K^{i + 1}$에 대응하는 $n \times m$ 행렬이라 하자.
가정은 $A \bmod \mathfrak p$의 계수가 $r$이고
$B \bmod \mathfrak p$의 계수가 $m - r$이라는 것이다. 다시 말해
$A$에는 $\mathfrak p$에 속하지 않는 어떤 $r \times r$ 소행렬식
$a$가 있고, $B$에는 $\mathfrak p$에 속하지 않는 어떤
$(m - r) \times (m - r)$-소행렬식 $b$가 있다.
$f = ab$로 두자. 그러면 $f$를 가역화한 뒤 직합 분해
$K^{i - 1} = R^{\oplus l - r} \oplus R^{\oplus r}$,
$K^i = R^{\oplus r} \oplus R^{\oplus m - r}$,
$K^{i + 1} = R^{\oplus m - r} \oplus R^{\oplus n - m + r}$를
찾아 다음이 성립하게 할 수 있다.
% P02-E0989: frozen source twice says a module map “kills of” a summand; Korean renders the intended annihilation statements.
가군 사상 $K^{i - 1} \to K^i$는 $R^{\oplus l - r}$을 소멸시키고
$R^{\oplus r}$을 $K^i$의 대응하는 직합인자 위로 동형으로 보내며,
가군 사상 $K^i \to K^{i + 1}$은 $R^{\oplus r}$을 소멸시키고
$R^{\oplus m - r}$을 $K^{i + 1}$의 대응하는 직합인자 위로 동형으로
보낸다. 따라서 $K^\bullet$은
$$
\ldots \to K^{i - 2} \to R^{\oplus l - r}
\to 0 \to R^{\oplus n - m + r} \to K^{i + 2} \to \ldots
$$
와 준동형이 되고, 이제 모든 것이 분명하다.
\end{proof}

\begin{lemma}
\label{more-algebra-lemma-split-using-ext-injective}
$R$을 환이라 하고 $K \in D^-(R)$라 하자. $a \in \mathbf{Z}$라 하자.
임의의 단사 $R$-가군 사상 $M \to M'$에 대하여 사상
$\Ext^{-a}_R(K, M) \to \Ext^{-a}_R(K, M')$이 단사라고 가정하자.
그러면 유일한 직합 분해
$K \cong \tau_{\leq a}K \oplus \tau_{\geq a + 1}K$가 존재하고,
어떤 $b$에 대하여 $\tau_{\geq a + 1}K$의 사영 진폭은
$[a + 1, b]$ 안에 있다.
\end{lemma}

\begin{proof}
$D(R)$의 특수 삼각형
$$
\tau_{\leq a}K \to K \to \tau_{\geq a + 1}K \to (\tau_{\leq a}K)[1]
$$
을 생각하자. 유도 범주, 주석
\ref{derived-remark-truncation-distinguished-triangle}을 보라.
유도 범주, 보조정리 \ref{derived-lemma-negative-exts}에 의해
$\Ext^{-a}_R(\tau_{\leq a}K, M) = \Hom_R(H^a(K), M)$이고
$\Ext^{-a - 1}_R(\tau_{\leq a}K, M) = 0$임에 유의하라. 따라서
$\Ext$의 긴 완전열은 $R$-가군 $M$에 관하여 함자적인 완전열
$$
0 \to
\Ext^{-a}_R(\tau_{\geq a + 1}K, M) \to
\Ext^{-a}_R(K, M) \to
\Hom_R(H^a(K), M)
$$
을 준다. 이제 $I$가 단사 $R$-가군이면 예를 들어 유도 범주,
보조정리 \ref{derived-lemma-compute-ext-resolutions}에 의해
$\Ext^{-a}_R(\tau_{\geq a + 1}K, I) = 0$이다. 모든 가군은 단사
가군으로 매장되므로 모든 $R$-가군 $M$에 대하여
$\Ext^{-a}_R(\tau_{\geq a + 1}K, M) = 0$이라고 결론짓는다.
보조정리 \ref{more-algebra-lemma-projective-amplitude}에 의해 어떤 $b$에 대하여
$\tau_{\geq a + 1}K$의 사영 진폭이 $[a + 1, b]$ 안에 있다고
결론짓는다. 여기서 $K$가 위로 유계라는 사실을 사용한다.
보조정리 \ref{more-algebra-lemma-splitting-unique}에 의해 분할을 얻는다.
\end{proof}

\begin{lemma}
\label{more-algebra-lemma-split-using-ext-zero}
$R$을 환이라 하고 $K \in D^-(R)$라 하자. $a \in \mathbf{Z}$라 하자.
임의의 $R$-가군 $M$에 대하여 $\Ext^{-a}_R(K, M) = 0$이라고
가정하자. 그러면 유일한 직합 분해
$K \cong \tau_{\leq a - 1}K \oplus \tau_{\geq a + 1}K$가
존재하고, 어떤 $b$에 대하여 $\tau_{\geq a + 1}K$의 사영 진폭은
$[a + 1, b]$ 안에 있다.
\end{lemma}

\begin{proof}
보조정리 \ref{more-algebra-lemma-split-using-ext-injective}에 의해 직합 분해
$K \cong \tau_{\leq a}K \oplus \tau_{\geq a + 1}K$가 있고,
어떤 $b$에 대하여 $\tau_{\geq a + 1}K$의 사영 진폭은
$[a + 1, b]$ 안에 있다. 분명히 $H^a(K) = 0$이어야 하므로
$D(R)$에서 $\tau_{\leq a}K = \tau_{\leq a - 1}K$라고 결론짓는다.
\end{proof}

\section{퍼펙트 복합체의 판별}
\label{more-algebra-section-recognizing-perfect}

\noindent
어떤 복합체들이 퍼펙트임을 증명하게 해 주는 몇 가지 보조정리이다.

\begin{lemma}
\label{more-algebra-lemma-lift-bounded-pseudo-coherent-to-perfect}
$R$을 환이라 하고 $\mathfrak p \subset R$을 소아이디얼이라 하자.
$K$가 유사연접이고 아래로 유계라고 하자.
$d_i = \dim_{\kappa(\mathfrak p)}
H^i(K \otimes_R^\mathbf{L} \kappa(\mathfrak p))$로 두자.
어떤 $a \in \mathbf{Z}$가 존재하여 $i < a$이면 $d_i = 0$이라고
하자. 그러면 어떤 $f \in R$, $f \not \in \mathfrak p$와
$K \otimes_R^\mathbf{L} R_f$를 $D(R_f)$에서 나타내는 복합체
$$
\ldots \to 0 \to R_f^{\oplus d_a} \to R_f^{\oplus d_{a + 1}} \to
\ldots \to
R_f^{\oplus d_{b - 1}} \to
R_f^{\oplus d_b} \to 0
\to \ldots
$$
가 존재한다. 특히 $K \otimes_R^\mathbf{L} R_f$는 퍼펙트이다.
\end{lemma}

\begin{proof}
% P02-E0267: the lemma introduces K, but the frozen proof repeatedly switches to undeclared K^bullet; Korean preserves the formal source objects.
$a$를 줄인 뒤 $i < a$이면 $H^i(K^\bullet) = 0$인 것도 가정할 수
있다. 보조정리 \ref{more-algebra-lemma-cut-complex-in-two}에 의해 어떤
$f \in R$, $f \not \in \mathfrak p$에 대하여 $R$을 $R_f$로
바꾼 뒤, $\tau_{\geq a}K^\bullet$이 퍼펙트인 $D(R)$의 분해
$K^\bullet = \tau_{\leq a - 1}K^\bullet \oplus
\tau_{\geq a}K^\bullet$을 쓸 수 있다. $i < a$이면
$H^i(K^\bullet) = 0$이므로 $D(R)$에서
$\tau_{\leq a - 1}K^\bullet = 0$임을 알 수 있다. 따라서
$K^\bullet$은 퍼펙트이다. 이제 보조정리
\ref{more-algebra-lemma-lift-perfect-from-residue-field}를 사용하여 결론을 얻는다.
\end{proof}

\begin{lemma}
\label{more-algebra-lemma-check-perfect-pointwise}
$R$을 환이라 하고 $a, b \in \mathbf{Z}$라 하자.
$K^\bullet$을 유사연접 $R$-가군 복합체라 하자. 다음은 동치이다.
\begin{enumerate}
\item $K^\bullet$은 Tor 진폭 $[a, b]$를 갖는 퍼펙트 복합체이다.
\item 모든 소아이디얼 $\mathfrak p$와 모든
$i \not \in [a, b]$에 대하여
$H^i(K^\bullet \otimes_R^{\mathbf{L}} \kappa(\mathfrak p)) = 0$이다.
\item 모든 극대 아이디얼 $\mathfrak m$과 모든
$i \not \in [a, b]$에 대하여
$H^i(K^\bullet \otimes_R^{\mathbf{L}} \kappa(\mathfrak m)) = 0$이다.
\end{enumerate}
\end{lemma}

\begin{proof}
함의 (1) $\Rightarrow$ (2) $\Rightarrow$ (3)의 증명은 생략한다.
(3)을 가정하자. $i \not \in [a, b]$인
$i \in \mathbf{Z}$를 택하자. 보조정리
\ref{more-algebra-lemma-cut-complex-in-two}에 의해 가정은 $R$의 모든 극대
아이디얼에 대하여 $H^i(K^\bullet)_{\mathfrak m} = 0$임을
함의한다. 따라서 $H^i(K^\bullet) = 0$이다. 대수학, 보조정리
\ref{algebra-lemma-characterize-zero-local}을 보라. 더욱이 이제
보조정리 \ref{more-algebra-lemma-cut-complex-in-two}는 각 극대 아이디얼
$\mathfrak m$에 대하여 어떤 $f \in R$, $f \not \in \mathfrak m$가
존재하여 $K^\bullet \otimes_R R_f$가 Tor 진폭 $[a, b]$를 갖는
퍼펙트 복합체임도 함의한다. 따라서 보조정리
\ref{more-algebra-lemma-glue-perfect}와 \ref{more-algebra-lemma-glue-tor-amplitude}를
적용하여 결론을 얻는다.
\end{proof}

\begin{lemma}
\label{more-algebra-lemma-check-perfect-stalks}
$R$을 환이라 하고 $K^\bullet$을 유사연접 $R$-가군 복합체라 하자.
다음 조건들을 생각하자.
\begin{enumerate}
\item $K^\bullet$은 퍼펙트이다.
\item 모든 소아이디얼 $\mathfrak p$에 대하여 복합체
$K^\bullet \otimes_R R_{\mathfrak p}$는 퍼펙트이다.
\item 모든 극대 아이디얼 $\mathfrak m$에 대하여 복합체
$K^\bullet \otimes_R R_{\mathfrak m}$은 퍼펙트이다.
\item 모든 소아이디얼 $\mathfrak p$에 대하여 모든 $i \ll 0$에서
$H^i(K^\bullet \otimes_R^{\mathbf{L}} \kappa(\mathfrak p)) = 0$이다.
\item 모든 극대 아이디얼 $\mathfrak m$에 대하여 모든 $i \ll 0$에서
$H^i(K^\bullet \otimes_R^{\mathbf{L}} \kappa(\mathfrak m)) = 0$이다.
\end{enumerate}
항상 다음 함의들이 성립한다.
% P02-E0990: frozen source duplicates condition (3) in the implication chain; Korean preserves the display exactly.
$$
(1) \Rightarrow (2) \Leftrightarrow (3) \Leftrightarrow (3)
\Leftrightarrow (4) \Leftrightarrow (5)
$$
$K^\bullet$이 아래로 유계이면 모든 조건은 동치이다.
\end{lemma}

\begin{proof}
보조정리 \ref{more-algebra-lemma-pull-perfect}에 의해 (1)은 (2)를 함의한다.
(2) $\Rightarrow$ (3)은 즉시 성립한다. 모든 소아이디얼
$\mathfrak p$는 어떤 극대 아이디얼 $\mathfrak m$에 포함되므로
사상 $R_\mathfrak m \to R_\mathfrak p$에 보조정리
\ref{more-algebra-lemma-pull-perfect}를 적용하면 (3)이 (2)를 함의함을 알 수 있다.
잉여 사상 $R_\mathfrak p \to \kappa(\mathfrak p)$와
$R_\mathfrak m \to \kappa(\mathfrak m)$에 보조정리
\ref{more-algebra-lemma-pull-perfect}를 적용하면 (2)가 (4)를 함의하고 (3)이
(5)를 함의함을 알 수 있다.

\medskip\noindent
$R$이 극대 아이디얼 $\mathfrak m$과 잉여체 $\kappa$를 갖는
국소환이라고 가정하자. 어떤 $a$에 대하여 $i < a$이면
$H^i(K^\bullet \otimes^\mathbf{L} \kappa) = 0$이라는 것으로부터
% P02-E0269: the frozen proof switches from K^bullet to undeclared K in its perfection and truncation claims; Korean preserves the formal source objects.
$K$가 퍼펙트임을 보이겠다. 이는 (4)가 (2)를 함의하고 (5)가 (3)을
함의함을 보여 보조정리의 첫 부분을 증명한다. 먼저
$i = a - 1$로 보조정리 \ref{more-algebra-lemma-cut-complex-in-two}를 적용하면,
$\tau_{\geq a}K^\bullet$이 $[a, \infty]$에 포함되는 Tor 진폭을 갖는
퍼펙트 대상인 $D(R)$의 분해
$K^\bullet = \tau_{\leq a - 1}K^\bullet \oplus
\tau_{\geq a}K^\bullet$을 얻는다. 증명을 끝내려면
$\tau_{\leq a - 1}K$가 영임을, 즉 그 코호몰로지 군들이 영임을
보여야 한다. 그렇지 않다면
$M = H^i(\tau_{\leq a - 1}K)$가 영이 아닌 가장 큰 지수를 $i$라
하자. $\tau_{\leq a - 1}K^\bullet$은 유사연접이므로
$M$은 유한 $R$-가군이다(보조정리
\ref{more-algebra-lemma-finite-cohomology}와
\ref{more-algebra-lemma-summands-pseudo-coherent}). 따라서 나카야마 보조정리
(대수학, 보조정리 \ref{algebra-lemma-NAK})에 의해
$M \otimes_R \kappa$가 영이 아님을 알 수 있다. 이는
$$
H^i((\tau_{\leq a - 1}K^\bullet) \otimes_R^\mathbf{L} \kappa) =
H^i(K^\bullet \otimes_R^\mathbf{L} \kappa)
$$
가 영이 아님을 함의하여 모순이다.

\medskip\noindent
동치인 조건 (2)--(5)가 성립하고 $K^\bullet$이 아래로 유계라고
가정하자. $i < a$이면 $H^i(K^\bullet) = 0$이라고 하자.
$R$의 극대 아이디얼 $\mathfrak m$을 택하자. 어떤
$f \in R$, $f \not \in \mathfrak m$가 존재하여
$K^\bullet \otimes_R^\mathbf{L} R_f$가 퍼펙트임을 보이면 충분하다
(보조정리 \ref{more-algebra-lemma-glue-perfect}와 대수학, 보조정리
\ref{algebra-lemma-quasi-compact}). 이는 보조정리
\ref{more-algebra-lemma-lift-bounded-pseudo-coherent-to-perfect}에서 따른다.
\end{proof}

\begin{lemma}
\label{more-algebra-lemma-projective-amplitude-pseudo-coherent}
$R$을 환이라 하고 $K$를 $D(R)$의 유사연접 대상이라 하자.
$a, b \in \mathbf{Z}$라 하자. 다음은 동치이다.
\begin{enumerate}
\item $K$의 사영 진폭은 $[a, b]$ 안에 있다.
\item $K$는 Tor 진폭 $[a, b]$를 갖는 퍼펙트 대상이다.
\item 모든 유한 표시 $R$-가군 $N$과 모든 $i \not \in [-b, -a]$에
대하여 $\Ext^i_R(K, N) = 0$이다.
\item $n > b$에 대하여 $H^n(K) = 0$이고, 모든 유한 표시 $R$-가군
$N$과 모든 $i > -a$에 대하여 $\Ext^i_R(K, N) = 0$이다.
\item $n \not \in [a - 1, b]$에 대하여 $H^n(K) = 0$이고, 모든 유한
표시 $R$-가군 $N$에 대하여 $\Ext^{-a + 1}_R(K, N) = 0$이다.
\end{enumerate}
\end{lemma}

\begin{proof}
보조정리 \ref{more-algebra-lemma-perfect}의 마지막 명제에서 (2)가 (1)을 함의함을
알 수 있다. (1)이 성립하면 $i \not \in [a, b]$에 대하여
$P^i = 0$인 사영 가군들의 복합체 $P^i$로 $K$를 나타낼 수 있다.
사영 가군은 자유 가군의 직합인자로서 평탄하므로 $K$의 Tor 진폭이
$[a, b]$ 안에 있음을 알 수 있다. 보조정리
\ref{more-algebra-lemma-tor-amplitude}를 보라. 따라서 보조정리
\ref{more-algebra-lemma-perfect}에 의해 (2)가 성립함을 알 수 있다.

\medskip\noindent
조건 (3), (4), (5)에서 유한 표시인 $M$에 대하여 가정한 Ext 군
$\Ext^i_R(K, M)$의 소멸은 모든 $R$-가군 $M$에 대한 소멸과
동치이다. 이는 보조정리
\ref{more-algebra-lemma-pseudo-coherence-colimit-ext}와 대수학, 보조정리
\ref{algebra-lemma-module-colimit-fp}에서 따른다. 따라서 (1), (3),
(4), (5)의 동치는 보조정리 \ref{more-algebra-lemma-projective-amplitude}에서
따른다.
\end{proof}

\noindent
% P02-E0991: frozen source says “The following lemma useful”, omitting the finite verb; Korean supplies it.
다음 보조정리는 다항식환 $B = A[x_1, \ldots, x_d]$ 위의 퍼펙트
복합체를 찾는 데 유용하다.

\begin{lemma}
\label{more-algebra-lemma-perfect-over-polynomial-ring}
$A \to B$를 환 사상이라 하자. $a, b \in \mathbf{Z}$라 하고
$d \geq 0$이라 하자. $K^\bullet$을 $B$-가군 복합체라 하자.
다음을 가정하자.
\begin{enumerate}
\item 환 사상 $A \to B$는 평탄하다.
\item 모든 소아이디얼 $\mathfrak p \subset A$에 대하여 환
$B \otimes_A \kappa(\mathfrak p)$의 대역 차원은 유한하고
$\leq d$이다.
\item $K^\bullet$은 $B$-가군 복합체로서 유사연접이다.
\item $K^\bullet$은 $A$-가군 복합체로서 Tor 진폭 $[a, b]$를 갖는다.
\end{enumerate}
그러면 $K^\bullet$은 Tor 진폭 $[a - d, b]$를 갖는 퍼펙트
$B$-가군 복합체이다.
\end{lemma}

\begin{proof}
$K^\bullet$이 유한 자유 $B$-가군들로 이루어진 위로 유계인
복합체라고 가정할 수 있다. 특히 $K^\bullet$은 $A$-가군 복합체로서
평탄하고, 임의의 $A$-가군 $M$에 대하여
$K^\bullet \otimes_A M = K^\bullet \otimes_A^{\mathbf{L}} M$이다.
$A$의 모든 소아이디얼 $\mathfrak p$에 대하여 복합체
$$
K^\bullet \otimes_A \kappa(\mathfrak p)
$$
는 $B \otimes_A \kappa(\mathfrak p)$ 위의 유한 자유 가군들로
이루어진 위로 유계인 복합체이고, $i \in [a, b]$인 경우를 제외하면
$H^i$가 사라진다. $B \otimes_A \kappa(\mathfrak p)$의 대역 차원이
$d$이므로 보조정리 \ref{more-algebra-lemma-finite-gl-dim-tor-dimension}에서
$K^\bullet \otimes_A \kappa(\mathfrak p)$의 Tor 진폭이
$[a - d, b]$ 안에 있음을 알 수 있다. $\mathfrak q$를
$\mathfrak p$ 위에 놓인 $B$의 소아이디얼이라 하자.
$K^\bullet \otimes_A \kappa(\mathfrak p)$는 자유
$B \otimes_A \kappa(\mathfrak p)$-가군들로 이루어진 위로 유계인
복합체이므로
\begin{align*}
K^\bullet \otimes_B^{\mathbf{L}} \kappa(\mathfrak q)
& = K^\bullet \otimes_B \kappa(\mathfrak q) \\
& = (K^\bullet \otimes_A \kappa(\mathfrak p))
\otimes_{B \otimes_A \kappa(\mathfrak p)} \kappa(\mathfrak q) \\
& = (K^\bullet \otimes_A \kappa(\mathfrak p))
\otimes^{\mathbf{L}}_{B \otimes_A \kappa(\mathfrak p)} \kappa(\mathfrak q)
\end{align*}
임을 알 수 있다. 따라서 위의 논증은 $i \not \in [a - d, b]$이면
$H^i(K^\bullet \otimes_B^{\mathbf{L}} \kappa(\mathfrak q)) = 0$임을
함의한다. 보조정리 \ref{more-algebra-lemma-check-perfect-pointwise}에 의해
결론을 얻는다.
\end{proof}

\noindent
다음 보조정리는 보조정리 \ref{more-algebra-lemma-perfect-over-polynomial-ring}의
국소 형태이다. 정칙 국소환 위의 퍼펙트 복합체를 찾는 데 사용할 수
있다.

\begin{lemma}
\label{more-algebra-lemma-perfect-over-regular-local-ring}
$A \to B$를 국소환 준동형이라 하자. $a, b \in \mathbf{Z}$라 하고
$d \geq 0$이라 하자. $K^\bullet$을 $B$-가군 복합체라 하자.
다음을 가정하자.
\begin{enumerate}
\item 환 사상 $A \to B$는 평탄하다.
\item 환 $B/\mathfrak m_AB$는 차원 $d$의 정칙환이다.
\item $K^\bullet$은 $B$-가군 복합체로서 유사연접이다.
\item $K^\bullet$은 $A$-가군 복합체로서 Tor 진폭 $[a, b]$를 갖는다.
사실 $H^i(K^\bullet \otimes_A^\mathbf{L} \kappa(\mathfrak m_A))$가
$i \in [a, b]$인 경우에만 영이 아니라고 가정해도 충분하다.
\end{enumerate}
그러면 $K^\bullet$은 Tor 진폭 $[a - d, b]$를 갖는 퍼펙트
$B$-가군 복합체이다.
\end{lemma}

\begin{proof}
(3)에 의해 $K^\bullet$이 유한 자유 $B$-가군들로 이루어진 위로
유계인 복합체라고 가정할 수 있다. 다음을 계산한다.
\begin{align*}
K^\bullet \otimes_B^{\mathbf{L}} \kappa(\mathfrak m_B)
& = K^\bullet \otimes_B \kappa(\mathfrak m_B) \\
& = (K^\bullet \otimes_A \kappa(\mathfrak m_A))
\otimes_{B/\mathfrak m_A B} \kappa(\mathfrak m_B) \\
& = (K^\bullet \otimes_A \kappa(\mathfrak m_A))
\otimes^{\mathbf{L}}_{B/\mathfrak m_A B} \kappa(\mathfrak m_B)
\end{align*}
% P02-E0992: frozen source contains the fragment “The first equality because ...”; Korean supplies the predicate.
첫 번째 등식은 $K^\bullet$이 평탄 $B$-가군들로 이루어진 위로 유계인
복합체이기 때문에 성립한다. 두 번째 등식은 텐서곱의 기본 성질에서
따른다. 세 번째 등식은
$K^\bullet \otimes_A \kappa(\mathfrak m_A) =
K^\bullet/ \mathfrak m_A K^\bullet$이 평탄
$B/\mathfrak m_A B$-가군들로 이루어진 위로 유계인 복합체이기 때문에
성립한다. (1)에 의해 $K^\bullet$은 평탄 $A$-가군들로 이루어진 위로
유계인 복합체이므로, 가정 (4)에 의해 복합체
$K^\bullet \otimes_A \kappa(\mathfrak m_A)$의 코호몰로지 가군
$H^i$는 $i \in [a, b]$인 경우에만 영이 아니다. 따라서 예
\ref{more-algebra-example-cohomology-complex-tensored}의 스펙트럼 열과
$B/\mathfrak m_AB$의 대역 차원이 유한하고 $d$라는 사실
(이는 (2)와 대수학, 명제
\ref{algebra-proposition-regular-finite-gl-dim}에서 따른다)로부터
% P02-E0272: frozen source gives the compound subject “the spectral sequence ... and the fact ...” a singular verb; Korean renders the intended conclusion.
$j \not \in [a - d, b]$이면
$H^j(K^\bullet \otimes_B^{\mathbf{L}} \kappa(\mathfrak m_B))$가
영임을 알 수 있다. 보조정리 \ref{more-algebra-lemma-check-perfect-pointwise}에
의해 증명이 끝난다.
\end{proof}



















\section{퍼펙트 복합체의 특성화}
\label{more-algebra-section-perfect-compact}

\noindent
이 절에서는 퍼펙트 복합체가 정확히 환의 유도 범주의 콤팩트
대상임을 증명한다. 먼저 다음을 보인다.

\begin{lemma}
\label{more-algebra-lemma-perfect-ring-classical-generator}
$R$을 환이라 하자. 퍼펙트 대상들의 충만 부분범주
$D_{perf}(R) \subset D(R)$은 $R = R[0]$을 포함하는 가장 작은 엄밀히
충만하고 포화된 삼각 부분범주이다. 다시 말해
$D_{perf}(R) = \langle R \rangle$이다. 특히 $R$은
$D_{perf}(R)$의 고전적 생성자이다.
\end{lemma}

\begin{proof}
명제의 뜻은 유도 범주론, 정의
\ref{derived-definition-saturated}와
\ref{derived-definition-generators}를 보라.
보조정리 \ref{more-algebra-lemma-two-out-of-three-perfect}와
\ref{more-algebra-lemma-summands-perfect}에서 $D_{perf}(R) \subset D(R)$이
$D(R)$의 엄밀히 충만하고 포화된 삼각 부분범주임을 보였다.
물론 $R \in D_{perf}(R)$이다.

\medskip\noindent
$\langle R \rangle = \bigcup \langle R \rangle_n$임을 상기하자.
증명을 끝내기 위해 $M \in D_{perf}(R)$이
$$
\ldots \to 0 \to M^a \to M^{a + 1} \to \ldots \to M^b \to 0 \to \ldots
$$
로 표시되고 $M^i$가 유한 사영이면
$M \in \langle R \rangle_{b - a + 1}$임을 보이겠다. 증명은
$b - a$에 대한 귀납법으로 한다. 정의에 의해
$\langle R \rangle_1$은 임의의 차수에 놓인 모든 유한 사영
$R$-가군을 포함하므로, 이는 귀납법의 기초 단계 $b - a = 0$을
처리한다. 일반적으로 다음 완전 삼각형을 생각한다.
% P02-E0273: frozen source uses the undefined subscripted M_b twice after defining the cochain terms with superscripts; Korean preserves the formulas.
$$
M_b[-b] \to M^\bullet \to \sigma_{\leq b - 1}M^\bullet \to M_b[-b + 1]
$$
귀납 가정에 의해 절단 복합체 $\sigma_{\leq b - 1}M^\bullet$은
$\langle R \rangle_{b - a}$에 속하고 $M_b[-b]$는
$\langle R \rangle_1$에 속한다. 따라서 정의에 의해
$M^\bullet \in \langle R \rangle_{b - a + 1}$이다.
\end{proof}

\noindent
$R$을 환이라 하자. $D(R)$에는 복합체들의 직합을 취해 주어지는
직합이 있음을 상기하자. 유도 범주론, 보조정리
\ref{derived-lemma-direct-sums}를 보라. 이 절의 보조정리들에서는
이를 더 언급하지 않고 사용하겠다.

\begin{lemma}
\label{more-algebra-lemma-commutes-with-countable-sums}
$R$을 환이라 하자. 모든 가산 대상족 $E_n \in D(R)$에 대해 표준 사상
$$
\bigoplus \Hom_{D(R)}(K, E_n) \longrightarrow \Hom_{D(R)}(K, \bigoplus E_n)
$$
이 전단사인 대상 $K \in D(R)$을 택하자. 그러면
$\mathbf{N}$ 위의 복합체계 $L_n^\bullet$이 주어질 때
$$
\colim \Hom_{D(R)}(K, L^\bullet_n) \longrightarrow \Hom_{D(R)}(K, L^\bullet)
$$
은 전단사이다. 여기서 $L^\bullet$은 항별 여극한, 즉 모든
$m \in \mathbf{Z}$에 대해 $L^m = \colim L_n^m$인 복합체이다.
\end{lemma}

\begin{proof}
복합체들의 짧은 완전열
$$
0 \to \bigoplus L_n^\bullet \to \bigoplus L_n^\bullet \to L^\bullet \to 0
$$
을 생각하자. 여기서 $t_n : L_n^\bullet \to L_{n + 1}^\bullet$은
전이 사상이고 첫 번째 사상은 차수 $n$에서 $1 - t_n$으로 주어진다.
% P02-E0994: frozen source says “in degree n” although n indexes system summands; Korean preserves that distinction rather than silently revising the source topology.
유도 범주론, 보조정리
\ref{derived-lemma-derived-canonical-delta-functor}에 의해 이는
$D(R)$의 완전 삼각형이다. 호몰로지 함자
$\Hom_{D(R)}(K, -)$를 적용하자. 유도 범주론, 보조정리
\ref{derived-lemma-representable-homological}를 보라.
% P02-E0274: frozen source lacks a governing verb before the displayed long exact sequence; Korean supplies the intended sentence.
그리하여 다음 긴 완전 코호몰로지 열을 얻는다.
$$
\xymatrix{
& \ldots \ar[r] & \Hom_{D(R)}(K, \colim L^\bullet_n[-1]) \ar[lld] \\
\Hom_{D(R)}(K, \bigoplus L^\bullet_n) \ar[r] &
\Hom_{D(R)}(K, \bigoplus L^\bullet_n) \ar[r] &
\Hom_{D(R)}(K, \colim L^\bullet_n) \ar[lld] \\
\Hom_{D(R)}(K, \bigoplus L^\bullet_n[1]) \ar[r] & \ldots
}
$$
$\Hom_{D(R)}(K, \bigoplus L^\bullet_n)$이
$\bigoplus \Hom_{D(R)}(K, L^\bullet_n)$과 같다고 가정했으므로,
그림의 각 행에서 첫 번째 사상은 단사임을 알 수 있다(이 사상이
$1 - t_n$으로 유도되는 사상들의 합이라는 명시적 기술을 사용한다).
따라서 $\Hom_{D(R)}(K, \colim L^\bullet_n)$은 위 그림의 가운데
행의 첫 번째 사상의 여핵이며, 이것이 보일 것이었다.
\end{proof}

\noindent
다음 명제는 퍼펙트 복합체를 유도 범주의 콤팩트 대상
(유도 범주론, 정의 \ref{derived-definition-compact-object})으로
특성화하며 여러 곳에 나타난다. 예를 들어
\cite[proof of Proposition 6.3]{Rickard}(유계인 경우를 다룬다),
\cite[Theorem 2.4.3]{TT}(진술이 정확히 일치하지는 않는다), 그리고
\cite[Proposition 6.4]{Bokstedt-Neeman}(끔찍한 표기 관례에 주의하라)를
보라.

\begin{proposition}
\label{more-algebra-proposition-perfect-is-compact}
$R$을 환이라 하자. $D(R)$의 대상 $K$에 대해 다음은 동치이다.
\begin{enumerate}
\item $K$는 퍼펙트이다.
\item $K$는 $D(R)$의 콤팩트 대상이다.
\end{enumerate}
\end{proposition}

\begin{proof}
$K$가 퍼펙트라고 가정하자. 즉 $K$는 유한 사영 가군들로 이루어진
유계 복합체 $P^\bullet$과 유사동형이다. 정의
\ref{more-algebra-definition-perfect}를 보라. $E_i$가 복합체 $E_i^\bullet$로
표시되면 $\bigoplus E_i$는 차수 $n$의 항이 $\bigoplus E_i^n$인
복합체로 표시된다. 한편 모든 $n$에 대해 $P^n$이 사영이므로, 임의의
$R$-가군 복합체 $K^\bullet$에 대해
$\Hom_{D(R)}(P^\bullet, K^\bullet) = \Hom_{K(R)}(P^\bullet, K^\bullet)$이다.
유도 범주론, 보조정리
\ref{derived-lemma-morphisms-from-projective-complex}를 보라. 따라서
$\Hom_{D(R)}(P^\bullet, E^\bullet)$은 복합체
$$
\prod \Hom_R(P^n, E^{n - 1}) \to
\prod \Hom_R(P^n, E^n) \to
\prod \Hom_R(P^n, E^{n + 1}).
$$
의 코호몰로지이다. $P^\bullet$이 유계이므로 위 복합체에서
$\prod$ 기호를 $\bigoplus$ 기호로 바꿀 수 있다. 각 $P^n$이 유한
$R$-가군이므로 모든 $n, m$에 대해
$\Hom_R(P^n, \bigoplus_i E_i^m) = \bigoplus_i \Hom_R(P^n, E_i^m)$이다.
이 관찰들을 합치면 유도 범주론, 정의
\ref{derived-definition-compact-object}의 사상이 전단사임을 알 수 있다.

\medskip\noindent
반대로 $K$가 콤팩트라고 가정하자. $K$를 복합체 $K^\bullet$로
표시하고 사상
$$
K^\bullet
\longrightarrow
\bigoplus\nolimits_{n \geq 0} \tau_{\geq n} K^\bullet
$$
을 생각하자. 여기서는 표준 절단을 사용했다. 호몰로지, 절
\ref{homology-section-truncations}를 보라. 각 차수에서 오른쪽의 직합이
유한하므로 이 사상은 의미가 있다. 가정에 의해 이 사상은 유한 직합을
통해 분해된다. 따라서 적어도 하나의 $n$에 대해
$K \to \tau_{\geq n} K$가 영, 즉 $K$가 $D^{-}(R)$에 속한다.

\medskip\noindent
$K \in D^{-}(R)$이고 모든 $R$-가군이 자유 가군의 몫이므로,
$K$를 자유 $R$-가군들로 이루어진 위로 유계인 복합체
$K^\bullet$로 표시할 수 있다. 유도 범주론, 보조정리
\ref{derived-lemma-subcategory-left-resolution}를 보라. ‘stupid’ 절단을
사용하면
$$
K^\bullet = \bigcup\nolimits_{n \leq 0} \sigma_{\geq n}K^\bullet
$$
이다. 호몰로지, 절 \ref{homology-section-truncations}를 보라. 따라서
보조정리 \ref{more-algebra-lemma-commutes-with-countable-sums}에 의해
$1 : K^\bullet \to K^\bullet$은 $D(R)$에서
$\sigma_{\geq n}K^\bullet \to K^\bullet$을 통해 분해된다. 즉
$1 : K^\bullet \to K^\bullet$은 $D(R)$에서 어떤 유계 복합체
$L^\bullet$을 통해
$$
K^\bullet \xrightarrow{\varphi} L^\bullet \xrightarrow{\psi} K^\bullet
$$
로 분해되며, 그 항들은 자유 $R$-가군이다.
$i \not \in [a, b]$이면 $L^i = 0$이라고 하자. 이제부터 $a, b$를
고정한다. 위와 같은 $1_{K^\bullet}$의 분해 중 $i < c$일 때
$L^i$가 유한 자유가 되도록 하는 가장 큰 정수 $c$를 택하자. 이 정수는
$\leq b + 1$이다.
% P02-E0275: frozen source has the malformed clause “with L^i is finite free”; Korean renders the intended condition.
귀납법으로 $c = b + 1$임을 보이겠다. 구체적으로
$L^c = \bigoplus_{\lambda \in \Lambda} R$라 쓰자. $L^{c - 1}$이
유한 자유이므로, $L^{c - 1} \to L^c$가
$\bigoplus_{\lambda \in \Lambda'} R \subset L^c$을 통해 분해되도록 하는
유한 부분집합 $\Lambda' \subset \Lambda$를 찾을 수 있다. 차수 $c$에서
$\Lambda \setminus \Lambda'$에 해당하는 인자들로의 사영으로 주어지는
복합체 사상
$$
\pi :
L^\bullet
\longrightarrow
(\bigoplus\nolimits_{\lambda \in \Lambda \setminus \Lambda'} R)[-c]
$$
을 생각하자. $K$에 대한 가정에 의해, 필요하면 $\Lambda'$을 더 큰 유한
부분집합으로 바꾸어 $D(R)$에서 $\pi \circ \varphi = 0$이라고 가정할
수 있다. $\pi$의 핵을 $(L')^\bullet \subset L^\bullet$이라 하자.
$\pi$가 전사이므로 복합체들의 짧은 완전열을 얻으며, 이는 $D(R)$에서
완전 삼각형을 준다(유도 범주론, 보조정리
\ref{derived-lemma-derived-canonical-delta-functor}를 보라).
$\Hom_{D(R)}(K, -)$가 호몰로지적이고(유도 범주론, 보조정리
\ref{derived-lemma-representable-homological}를 보라)
$\pi \circ \varphi = 0$이므로, $(L')^\bullet \to L^\bullet$과 합성하면
$\varphi$가 되는 $D(R)$의 사상
$\varphi' : K^\bullet \to (L')^\bullet$을 찾을 수 있다.
$\psi'$을 $(L')^\bullet \to L^\bullet$과 $\psi$의 합성으로 놓으면 새로운
분해를 얻는다. $(L')^\bullet$은 차수 $c$를 제외하면 $L^\bullet$과 같고
$(L')^c = \bigoplus_{\lambda \in \Lambda'} R$이므로 귀납 단계가 증명된다.

\medskip\noindent
앞 문단의 논의로부터 $D(R)$에서
$1_K : K \to K$가
$$
K \xrightarrow{\varphi} L \xrightarrow{\psi} K
$$
로 분해되며, 여기서 $L$은 자유 $R$-가군들로 이루어진 유한 복합체로
표시될 수 있다. 특히 $L$은 퍼펙트이다. 또한
$e = \varphi \circ \psi \in \text{End}_{D(R)}(L)$는 멱등원이다.
유도 범주론, 보조정리
\ref{derived-lemma-projectors-have-images-triangulated}에 의해
$L = \Ker(e) \oplus \Ker(1 - e)$이다. 사상 $\varphi : K \to L$은
$D(R)$에서 $\Ker(1 - e)$와의 동형을 유도한다. 마침내 보조정리
\ref{more-algebra-lemma-summands-perfect}에 의해 $K$는 퍼펙트이다.
\end{proof}

\begin{lemma}
\label{more-algebra-lemma-perfect-modulo-nilpotent-ideal}
$R$을 환이라 하고 $I \subset R$을 아이디얼이라 하자.
$K$를 $D(R)$의 대상이라 하자. 다음을 가정하자.
\begin{enumerate}
\item $K \otimes_R^\mathbf{L} R/I$는 $D(R/I)$에서 퍼펙트이다.
\item $I$는 멱영 아이디얼이다.
\end{enumerate}
그러면 $K$는 $D(R)$에서 퍼펙트이다.
\end{lemma}

\begin{proof}
$K \otimes_R^\mathbf{L} R/I$를 표시하는 유한 사영 $R/I$-가군들의
유한 복합체 $\overline{P}^\bullet$을 택하자. 정의
\ref{more-algebra-definition-perfect}를 보라. 보조정리
\ref{more-algebra-lemma-lift-complex-projectives}에 의해 $K$를 표시하고
$\overline{P}^\bullet = P^\bullet/IP^\bullet$을 만족하는 사영
$R$-가군들의 복합체 $P^\bullet$이 존재한다. Nakayama 보조정리
(대수학, 보조정리 \ref{algebra-lemma-NAK})로부터 $P^\bullet$은 유한
사영 $R$-가군들로 이루어진 유한 복합체이다.
\end{proof}

\begin{lemma}
\label{more-algebra-lemma-perfect-modulo-two-ideals}
$R$을 환이라 하고 $I, J \subset R$을 아이디얼이라 하자.
$K$를 $D(R)$의 대상이라 하자. 다음을 가정하자.
\begin{enumerate}
\item $K \otimes_R^\mathbf{L} R/I$는 $D(R/I)$에서 퍼펙트이다.
\item $K \otimes_R^\mathbf{L} R/J$는 $D(R/J)$에서 퍼펙트이다.
\end{enumerate}
그러면 $K \otimes_R^\mathbf{L} R/IJ$는 $D(R/IJ)$에서 퍼펙트이다.
\end{lemma}

\begin{proof}
% P02-E0276: frozen source says “may assume replace”; Korean renders the intended replacement while retaining every formal span.
$R$을 $R/IJ$로, $K$를 $K \otimes_R^\mathbf{L} R/IJ$로 바꾸어도
됨은 분명하다. 그러면 $R \to R/(I \cap J)$는 핵의 제곱이 영인
전사이다. 따라서 보조정리 \ref{more-algebra-lemma-perfect-modulo-nilpotent-ideal}에
의해 $K \otimes_R^\mathbf{L} R/(I \cap J)$가 퍼펙트임을 보이면
충분하다. 그러므로 $I \cap J = 0$이라고 가정할 수 있다.

\medskip\noindent
$I \cap J = 0$인 경우에 보조정리를 증명하자. 먼저 모든 $K^n$이
평탄인 K-평탄 복합체 $K^\bullet$로 $K$를 표시할 수 있다. 보조정리
\ref{more-algebra-lemma-K-flat-resolution}를 보라. 그러면 복합체들의 짧은 완전열
$$
0 \to
K^\bullet \to K^\bullet/IK^\bullet \oplus K^\bullet/JK^\bullet \to
K^\bullet/(I + J)K^\bullet \to 0
$$
을 얻는다. 유도 텐서곱의 구성에 의해 $K^\bullet/IK^\bullet$은
$K \otimes^\mathbf{L}_R R/I$를 표시한다. $K^\bullet/JK^\bullet$과
$K^\bullet/(I + J)K^\bullet$도 마찬가지이다.
$K^\bullet/(I + J)K^\bullet$은 퍼펙트 $R/(I + J)$-가군 복합체이다.
보조정리 \ref{more-algebra-lemma-pull-perfect}를 보라.
% P02-E0277: frozen source has a malformed three-item conjunction; Korean renders the intended list without changing its formal spans.
따라서 복합체들 $K^\bullet/IK^\bullet$,
$K^\bullet/JK^\bullet$, 그리고 $K^\bullet/(I + J)K^\bullet$의 영이 아닌
코호몰로지 군은 유한 개뿐이다(퍼펙트 복합체의 Tor 진폭이 유한하기
때문이다. 보조정리 \ref{more-algebra-lemma-perfect}를 보라).
% P02-E0278: frozen source omits the definite article before its uniquely identified short exact sequence; Korean identifies the displayed sequence naturally.
위에 표시한 복합체들의 짧은 완전열에 딸린 긴 완전 코호몰로지 열로부터
$K \in D^b(R)$임을 얻는다.
% P02-E0999: frozen source says “we assume” where the cited resolution lemma permits a representative choice; Korean renders the permitted choice.
특히 $K^\bullet$을 자유 $R$-가군들로 이루어진 위로 유계인 복합체라
가정할 수 있다(유도 범주론, 보조정리
\ref{derived-lemma-subcategory-left-resolution}를 보라).

\medskip\noindent
이제 명제 \ref{more-algebra-proposition-perfect-is-compact}의 판정법을 사용하여
$K$가 퍼펙트임을 보이겠다.
% P02-E0279: frozen source reuses J as both an ideal and the family’s parameter set; Korean preserves the source symbols and makes no silent rename.
따라서 집합 $J$로 매개화된 대상족 $E_j \in D(R)$을 택하자.
$E_j$를 표시하며 항들이 평탄인 복합체 $E_j^\bullet$을 택한다.
예를 들어 보조정리 \ref{more-algebra-lemma-K-flat-resolution}를 보라. 다음이
복합체들의 짧은 완전열임은 분명하다.
$$
0 \to
E_j^\bullet \to
E_j^\bullet/IE_j^\bullet \oplus E_j^\bullet/JE_j^\bullet \to
E_j^\bullet/(I + J)E_j^\bullet \to 0
$$
직합을 취하면 비슷한 짧은 완전열
$$
0 \to
\bigoplus E_j^\bullet \to
\bigoplus E_j^\bullet/IE_j^\bullet \oplus E_j^\bullet/JE_j^\bullet \to
\bigoplus E_j^\bullet/(I + J)E_j^\bullet \to 0
$$
을 얻는다. ($- \otimes_R R/I$가 직합과 가환함에 유의하라.) 이 짧은
완전열은 $D(R)$에서 완전 삼각형을 결정한다. 유도 범주론, 보조정리
\ref{derived-lemma-derived-canonical-delta-functor}를 보라. 호몰로지 함자
$\Hom_{D(R)}(K, -)$를 적용하면(유도 범주론, 보조정리
\ref{derived-lemma-representable-homological}를 보라) 다음 가환 그림을
얻는다.
% P02-E0280: frozen diagram places four [-1] shifts outside Hom rather than on its target complexes; Korean preserves the diagram exactly.
$$
\xymatrix{
\bigoplus \Hom_{D(R)}(K^\bullet, E_j^\bullet/(I + J))[-1] \ar[r] \ar[d] &
\Hom_{D(R)}(K^\bullet, \bigoplus E_j^\bullet/(I + J))[-1] \ar[d] \\
\bigoplus \Hom_{D(R)}(K^\bullet, E_j^\bullet/I \oplus E_j^\bullet/J)[-1]
\ar[r] \ar[d] &
\Hom_{D(R)}(K^\bullet, \bigoplus E_j^\bullet/I \oplus E_j^\bullet/J)[-1]
\ar[d] \\
\bigoplus \Hom_{D(R)}(K^\bullet, E_j^\bullet) \ar[r] \ar[d] &
\Hom_{D(R)}(K^\bullet, \bigoplus E_j^\bullet) \ar[d] \\
\bigoplus \Hom_{D(R)}(K^\bullet, E_j^\bullet/I \oplus E_j^\bullet/J)
\ar[r] \ar[d] &
\Hom_{D(R)}(K^\bullet, \bigoplus E_j^\bullet/I \oplus E_j^\bullet/J)
\ar[d] \\
\bigoplus \Hom_{D(R)}(K^\bullet, E_j^\bullet/(I + J)) \ar[r] &
\Hom_{D(R)}(K^\bullet, \bigoplus E_j^\bullet/(I + J))
}
$$
여기서 열들은 완전하다. 임의의 $R$-가군 복합체 $E^\bullet$에 대해
다음이 성립함은 분명하다.
\begin{align*}
\Hom_{D(R)}(K^\bullet, E^\bullet/I) & =
\Hom_{K(R)}(K^\bullet, E^\bullet/I) \\
& =
\Hom_{K(R/I)}(K^\bullet/IK^\bullet, E^\bullet/I) \\
& =
\Hom_{D(R/I)}(K^\bullet/IK^\bullet, E^\bullet/I)
\end{align*}
$J$나 $I + J$로 나누는 경우도 마찬가지이다. 유도 범주론, 보조정리
\ref{derived-lemma-morphisms-from-projective-complex}를 보라.
% P02-E0281: frozen source contains an orphan “Derived Categories.” fragment here; Korean omits only the accidental prose fragment.
복합체들 $K^\bullet/IK^\bullet$, $K^\bullet/JK^\bullet$,
$K^\bullet/(I + J)K^\bullet$은 각각 $R/I$, $R/J$, $R/(I + J)$-가군들의
퍼펙트 복합체이므로, 가운데 사상을 제외한 모든 수평 사상은
동형이다. 명제 \ref{more-algebra-proposition-perfect-is-compact}를 보라. $5$-보조정리
(호몰로지, 보조정리 \ref{homology-lemma-five-lemma})에 의해 가운데
사상도 동형이고, 다시 명제 \ref{more-algebra-proposition-perfect-is-compact}에 의해
보조정리가 따른다.
\end{proof}






\section{강한 생성자와 정칙환}
\label{more-algebra-section-strong-generators-regular}

\noindent
$R$을 환이라 하자.
% P02-E0282: frozen source omits “by” in introducing D(R)_c; Korean introduces the notation grammatically.
% P02-E1002: the same frozen sentence does not identify which saturated subcategory is intended; Korean preserves that under-specification rather than adding content.
$D(R)$의 포화된 충만 삼각 부분범주를 $D(R)_c$로 나타내자. 이미
$$
\langle R \rangle = D_{perf}(R) = D(R)_c
$$
임을 안다. 보조정리 \ref{more-algebra-lemma-perfect-ring-classical-generator}와
명제 \ref{more-algebra-proposition-perfect-is-compact}를 보라. $R$이 정칙이면
$R$이 강한 생성자임을 알 수 있다(유도 범주론, 정의
\ref{derived-definition-generators}).

\begin{lemma}
\label{more-algebra-lemma-ghost-lemma}
\begin{reference}
\cite{Kelly}
\end{reference}
$R$을 환이라 하자. $n \geq 1$이고
$K \in \langle R \rangle_n$이라 하자. 표기는 유도 범주론, 절
\ref{derived-section-generators}와 같다. $D(R)$의 사상들
$$
K \xrightarrow{f_1} K_1 \xrightarrow{f_2} K_2
\xrightarrow{f_3} \ldots \xrightarrow{f_n} K_n
$$
을 생각하자. 모든 $i, j$에 대해 $H^i(f_j) = 0$이면
$f_n \circ \ldots \circ f_1 = 0$이다.
\end{lemma}

\begin{proof}
$n = 1$이면 $K$는 $D(R)$에서 유계 복합체 $P^\bullet$의 직합인자이며,
그 항들은 유한 자유 $R$-가군이고 미분들은 영이다. 따라서
모든 $i$에 대해 $H^i(f) = 0$인 $D(R)$의 임의의 사상
$f : P^\bullet \to K_1$이 영임을 보이면 충분하다. $P^\bullet$은 유한
직합 $P^\bullet = \bigoplus R[m_j]$이므로, $D(R)$에서
$H^{-m}(g) = 0$인 임의의 사상 $g : R[m] \to K_1$이 영임을 보이면
충분하다. 이는 $\Hom_{D(R)}(R[-m], K) = H^m(K)$라는 사실에서 따른다.

\medskip\noindent
$n > 1$이면 $n$에 대한 귀납법을 사용한다. 즉 $K$는 $D(R)$에서
완전 삼각형
$$
P' \xrightarrow{i} P \xrightarrow{p} P'' \to P'[1]
$$
에 놓이는 대상 $P$의 직합인자이며, 여기서
$P' \in \langle R \rangle_1$이고
$P'' \in \langle R \rangle_{n - 1}$임을 안다. 위와 같이 $K$를
$P$로 바꾸어, 코호몰로지에서 영인 $f_j$를 갖는 $D(R)$의 사상들
$$
P \xrightarrow{f_1} K_1 \xrightarrow{f_2} K_2
\xrightarrow{f_3} \ldots \xrightarrow{f_n} K_n
$$
이 주어졌다고 가정할 수 있다. $n = 1$인 경우에 의해 합성
$f_1 \circ i$는 영이다. 따라서 유도 범주론, 보조정리
\ref{derived-lemma-representable-homological}에 의해
$f_1 = h \circ p$를 만족하는 사상 $h : P'' \to K_1$을 찾을 수 있다.
$f_2 \circ h$는 코호몰로지에서 영이다. 그러므로 귀납 가정에 의해
$f_n \circ \ldots \circ f_2 \circ h = 0$이고, 따라서 원하는 대로
$f_n \circ \ldots \circ f_1 =
f_n \circ \ldots \circ f_2 \circ h \circ p = 0$이다.
\end{proof}

\begin{lemma}
\label{more-algebra-lemma-not-regular-not-strong}
$R$을 뇌터 환이라 하자. $R$이 $D_{perf}(R)$의 강한 생성자이면
$R$은 차원이 유한한 정칙환이다.
\end{lemma}

\begin{proof}
어떤 $n \geq 1$에 대해
$D_{perf}(R) = \langle R \rangle_n$이라고 가정하자. 임의의 유한
$R$-가군 $M$에 대해, 유한 자유 $R$-가군들로 이루어지고
$i = -n, \ldots, - 1$일 때 $H^i(P) = 0$이며
$M \cong \Coker(d^{-1})$인 복합체
% P02-E0283: frozen display labels the differential leaving P^{-n+1} as d^1; Korean preserves the formula exactly.
$$
P = (
P^{-n - 1} \xrightarrow{d^{-n - 1}}
P^{-n} \xrightarrow{d^{-n}}
P^{-n + 1} \xrightarrow{d^1}
\ldots \xrightarrow{d^{-1}} P^0)
$$
를 택할 수 있다. $P$는 $D_{perf}(R)$에 속한다.
% P02-E0284: frozen source omits “using” before the finite free resolution; Korean states the intended computational relation.
임의의 $R$-가군 $N$에 대해 $M$의 유한 자유 분해 $P$를 사용하여
$\Ext^n_R(M, N)$을 계산할 수 있다. 대수학, 절
\ref{algebra-section-ext}를 보고 유도 범주론, 절
\ref{derived-section-ext}와 비교하라. 특히 위 열은 원소
$$
\xi \in \Ext^n_R(\Coker(d^{-1}), \Coker(d^{-n - 1})) =
\Ext^n_R(M, \Coker(d^{-n - 1}))
$$
를 정의한다.
% P02-E1005: frozen English uses “a R-module map”; Korean is article-free and preserves the intended map.
$\Ext^n_R(M, N)$의 임의의 원소 $\overline{\xi}$에 대해 $R$-가군 사상
$\varphi : \Coker(d^{-n - 1}) \to N$이 존재한다. 이때 $\varphi$는
$\xi$를 $\overline{\xi}$로 보낸다.
$j = 1, \ldots, n - 1$에 대해 복합체
$$
K_j = (\Coker(d^{-n - 1}) \to P^{-n + 1} \to \ldots \to P^{-j})
$$
를 생각하자. 여기서 $\Coker(d^{-n - 1})$은 차수 $-n$에 있고
$P^t$는 차수 $t$에 있다. 또한
$K_n = \Coker(d^{-n - 1})[n]$으로 둔다. 그러면 코호몰로지에서
영 사상을 유도하는 사상들
$$
P \to K_1 \to K_2 \to \ldots \to K_n
$$
을 얻는다.
% P02-E0285: frozen source says “the composition of this maps”; Korean uses the intended plural reference.
$P \in D_{perf}(R) = \langle R \rangle_n$이므로 보조정리
\ref{more-algebra-lemma-ghost-lemma}에 의해 이 사상들의 합성은 $D(R)$에서 영이다.
유도 범주론, 보조정리
\ref{derived-lemma-morphisms-from-projective-complex}에 의해
$\Hom_{D(R)}(P, K_n) = \Hom_{K(R)}(P, K_n)$이므로
$\xi = 0$이다. 따라서 위 논의에 의해 모든 $R$-가군 $N$에 대해
$\Ext^n_R(M, N) = 0$이다. 대수학, 보조정리
\ref{algebra-lemma-projective-dimension-ext}에 의해 $M$의 사영 차원은
$\leq n - 1$이다. 이것이 모든 유한 $R$-가군 $M$에 대해 성립하므로
$R$의 대역 차원은 유한하다. 대수학, 보조정리
\ref{algebra-lemma-finite-gl-dim}을 보라. 끝으로 대수학, 보조정리
\ref{algebra-lemma-finite-gl-dim-finite-dim-regular}에 의해 결론을 얻는다.
\end{proof}

\begin{lemma}
\label{more-algebra-lemma-ext-regular}
$R$을 차원 $d < \infty$인 뇌터 정칙환이라 하자.
$K, L \in D^-(R)$이라 하자.
% P02-E0286: frozen English says “there exists an k”; Korean renders the existential hypothesis naturally.
$i \leq k$이면 $H^i(K) = 0$이고 $i \geq k - d + 1$이면
$H^i(L) = 0$이 되게 하는 $k$가 존재한다고 가정하자. 그러면
$\Hom_{D(R)}(K, L) = 0$이다.
\end{lemma}

\begin{proof}
$K$를 표시하는 위로 유계인 복합체를 $K^\bullet$이라 하고,
$i \geq n + 1$이면 $K^i = 0$이라고 하자. $K^\bullet$을
$\tau_{\geq k + 1}K^\bullet$로 바꾸어 $i \leq k$이면
$K^i = 0$이라고 가정할 수 있다. 이제 완전 삼각형
$$
K^n[-n] \to K^\bullet \to \sigma_{\leq n - 1}K^\bullet
$$
을 사용하면 $K^n[-n]$과 $\sigma_{\leq n - 1}K^\bullet$에 대해
보조정리를 증명하면 충분함을 알 수 있다. $n$에 대한 귀납법으로,
어떤 $R$-가군 $M$과 어떤 $m \geq k + 1$에 대해 복합체
$M[-m]$으로 $K$가 표시되는 경우만 증명하면 충분하다. 대수학,
보조정리 \ref{algebra-lemma-finite-gl-dim-finite-dim-regular}에 의해
$R$의 대역 차원은 $d$이므로 $M$은 사영 분해
$0 \to P_d \to \ldots \to P_0 \to M \to 0$을 갖는다. 차수
$m - i$에 $P_i$를 갖는 복합체 $P^\bullet$은 $M[-m]$을 표시하는
사영 가군들의 유계 복합체이다. 한편 $i \geq k - d  + 1$이면
$L^i = 0$인, $L$을 표시하는 복합체 $L^\bullet$을 택할 수 있다.
따라서 복합체 사상 $P^\bullet \to L^\bullet$은 모두 영이다. 유도
범주론, 보조정리 \ref{derived-lemma-morphisms-from-projective-complex}에
의해 보조정리가 따른다.
\end{proof}

\begin{lemma}
\label{more-algebra-lemma-split-complex-regular}
$R$을 차원 $1 \leq d < \infty$인 뇌터 정칙환이라 하자.
$K \in D(R)$가 퍼펙트이고 $k \in \mathbf{Z}$라 하자.
$i = k - d + 2, \ldots, k$일 때 $H^i(K) = 0$이라고 가정하자
($d = 1$이면 빈 조건이다). 그러면
$K = \tau_{\leq k - d + 1}K \oplus \tau_{\geq k + 1}K$이다.
\end{lemma}

\begin{proof}
코호몰로지의 소멸에 의해 완전 삼각형
$$
\tau_{\leq k - d + 1}K \to K \to \tau_{\geq k + 1}K \to
(\tau_{\leq k - d + 1}K)[1]
$$
을 얻는다. 유도 범주론, 보조정리 \ref{derived-lemma-split}에 의해
세 번째 화살표가 영임을 보이면 충분하다. 따라서
$\Hom_{D(R)}(\tau_{\geq k + 1}K, (\tau_{\leq k - d + 1}K)[1]) = 0$을
보이면 충분하며, 이는 보조정리 \ref{more-algebra-lemma-ext-regular}에서 따른다.
\end{proof}

\begin{lemma}
\label{more-algebra-lemma-regular-strong-generator}
$R$을 차원이 유한한 뇌터 정칙환이라 하자. 그러면 $R$은 퍼펙트
대상들의 충만 부분범주 $D_{perf}(R) \subset D(R)$의 강한 생성자이다.
\end{lemma}

\begin{proof}
$D(R)$의 대상 $K$가 퍼펙트일 필요충분조건은 $K$가 유계이고 유한
코호몰로지 가군들을 갖는 것이라는 사실을 사용하겠다. 보조정리
\ref{more-algebra-lemma-regular-perfect}를 보라. 삼각 범주의 강한 생성자는 유도
범주론, 정의 \ref{derived-definition-generators}에 정의되어 있다.
$d = \dim(R)$라 하자.

\medskip\noindent
$K \in D_{perf}(R)$라 하자.
$K \in \langle R \rangle_{d + 1}$임을 보이겠다. 대수학, 보조정리
\ref{algebra-lemma-finite-gl-dim-finite-dim-regular}에 의해 모든 유한
$R$-가군의 사영 차원은 $\leq d$이다. 모든 $n \in \mathbf{Z}$에 대해
$H^n(K)$의 사영 차원이 $\leq i$이면
$K$가 $\langle R \rangle_{i + 1}$에 속한다는 것을
$0 \leq i \leq d$에 대한 귀납법으로 보이겠다.

\medskip\noindent
기초 단계 $i = 0$. 이 경우 $H^n(K)$는 사영 차원이 $0$인 유한
$R$-가군이다. 다시 말해 각 코호몰로지는 사영 $R$-가군이다. 따라서
모든 $i > 0$과 $m, n \in \mathbf{Z}$에 대해
$\Ext^i_R(H^n(K), H^m(K)) = 0$이다. 유도 범주론, 보조정리
\ref{derived-lemma-ext-2-zero-pre}에 의해 $K$는 그 코호몰로지 가군들의
옮김들의 직합과 동형이다. 각 코호몰로지 가군은 유한 사영
$R$-가군이므로 $R$의 복사본들의 직합의 직합인자이다. 따라서 정의에
의해 $K$는 $\langle R \rangle_1$에 포함된다.

\medskip\noindent
귀납 단계. $i < d$에 대해 주장이 성립한다고 가정하고, 모든
$n \in \mathbf{Z}$에 대해 $H^n(K)$의 사영 차원이 $\leq i + 1$인
$K \in D_{perf}(R)$를 택하자. $n \not \in [a, b]$이면
$H^n(K)$가 영이 되도록 $a \leq b$를 택한다. 각
$n \in [a, b]$에 대해 $F^n$이 유한 자유 $R$-가군인 전사
$F^n \to H^n(K)$를 택하자. $F^n$이 사영이므로
$F^n \to H^n(K)$를 $D(R)$의 사상 $F^n[-n] \to K$로 올릴 수 있다
(작은 세부 사항은 생략한다). 따라서 코호몰로지 가군들 위에서 전사인
사상 $\bigoplus_{a \leq n \leq b} F^n[-n] \to K$를 얻는다.
$D(R)$의 완전 삼각형
$$
K' \to \bigoplus\nolimits_{a \leq n \leq b} F^n[-n] \to K \to K'[1]
$$
을 택하자. 물론 대상 $K'$은 유계이고 유한 코호몰로지 가군들을
갖는다. 우리가 택한 사상들에 의해 긴 완전 코호몰로지 열은 짧은
완전열들
$$
0 \to H^n(K') \to F^n \to H^n(K) \to 0
$$
로 분해된다. 대수학, 보조정리
\ref{algebra-lemma-exact-sequence-projective-dimension}에 의해
$H^n(K')$의 사영 차원은 $\leq \max(0, i)$이다. 따라서
$K' \in \langle R \rangle_{i + 1}$이다. 정의에 의해 이는 원하는 대로
$K$가 $\langle R \rangle_{i + 1 + 1}$에 속한다는 뜻이다.
\end{proof}

\begin{proposition}
\label{more-algebra-proposition-regular-strong-generator}
$R$을 뇌터 환이라 하자. 다음은 동치이다.
\begin{enumerate}
\item $R$은 차원이 유한한 정칙환이다.
\item $D_{perf}(R)$은 강한 생성자를 갖는다.
\item $R$은 $D_{perf}(R)$의 강한 생성자이다.
\end{enumerate}
\end{proposition}

\begin{proof}
이는 보조정리 \ref{more-algebra-lemma-perfect-ring-classical-generator},
\ref{more-algebra-lemma-not-regular-not-strong},
\ref{more-algebra-lemma-regular-strong-generator}, 그리고 유도 범주론, 보조정리
\ref{derived-lemma-classical-generator-strong-generator}의 형식적 귀결이다.
\end{proof}







\section{상대 유한 표시 가군}
\label{more-algebra-section-relative-finite-presentation}

\noindent
$R$을 환이라 하자. $A \to B$를 유한형 $R$-대수들 사이의 유한
사상이라 하자. $M$을 유한 $B$-가군이라 하자. 이때 다음은
{\bf 참이 아니다}.
$$
M\text{ of finite presentation over }B
\Leftrightarrow
M\text{ of finite presentation over }A
$$
% P02-E0287: frozen English splits the standard noun “counterexample”; Korean uses the standard term 반례.
반례는 $R = k[x_1, x_2, x_3, \ldots]$, $A = R$, $B = R/(x_i)$,
$M = B$이다. 이를 ``바로잡기'' 위해 유한 표시의 상대적 개념을
도입한다.

\begin{lemma}
\label{more-algebra-lemma-relatively-finitely-presented}
$R \to A$를 유한형 환 사상이라 하고 $M$을 $A$-가군이라 하자.
다음은 동치이다.
\begin{enumerate}
\item 어떤 표시 $\alpha : R[x_1, \ldots, x_n] \to A$에 대해
$M$은 유한 표시 $R[x_1, \ldots, x_n]$-가군이다.
\item 모든 표시 $\alpha : R[x_1, \ldots, x_n] \to A$에 대해
$M$은 유한 표시 $R[x_1, \ldots, x_n]$-가군이다.
\item $A'$가 유한 표시 $R$-대수인 임의의 전사 $A' \to A$에 대해
$M$은 유한 표시 $A'$-가군이다.
\end{enumerate}
이 경우 $M$은 유한 표시 $A$-가군이다.
\end{lemma}

\begin{proof}
% P02-E0288: frozen proof begins with a consequent-free “If” fragment; Korean gives the intended complete setup.
$\alpha : R[x_1, \ldots, x_n] \to A$와
$\beta : R[y_1, \ldots, y_m] \to A$가 표시들이라고 하자.
$\alpha(f_j) = \beta(y_j)$인 $f_j \in R[x_1, \ldots, x_n]$와
$\beta(g_i) = \alpha(x_i)$인 $g_i \in R[y_1, \ldots, y_m]$를 택한다.
그러면 가환 그림
$$
\xymatrix{
R[x_1, \ldots, x_n, y_1, \ldots, y_m]
\ar[d]^{x_i \mapsto g_i} \ar[rr]_-{y_j \mapsto f_j} & &
R[x_1, \ldots, x_n] \ar[d] \\
R[y_1, \ldots, y_m] \ar[rr] & & A
}
$$
을 얻는다. 따라서 대수학, 보조정리
\ref{algebra-lemma-finitely-presented-over-subring}와
\ref{algebra-lemma-finite-finitely-presented-extension}를 적용하면
(1)과 (2)의 동치가 따른다. 표시
$A' = R[x_1, \ldots, x_n]/(f_1, \ldots, f_m)$을 택하고 대수학,
보조정리 \ref{algebra-lemma-finite-finitely-presented-extension}를
사용하면 $M$이 유한 표시 $A'$-가군일 필요충분조건은 $M$이 유한 표시
$R[x_1, \ldots, x_n]$-가군인 것임을 알 수 있다. 따라서 (2)와
(3)은 동치이다.
\end{proof}

\begin{definition}
\label{more-algebra-definition-relatively-finitely-presented}
$R \to A$를 유한형 환 사상이라 하고 $M$을 $A$-가군이라 하자.
보조정리 \ref{more-algebra-lemma-relatively-finitely-presented}의 동치 조건들이
성립하면 $M$을 {\it $R$에 대해 상대적으로 유한 표시인 $A$-가군}이라
한다.
\end{definition}

\noindent
$R \to A$가 유한 표시이면, $M$이 $R$에 대해 상대적으로 유한 표시인
$A$-가군일 필요충분조건은 $M$이 유한 표시 $A$-가군인 것이다. 또한
$A$가 $A$-가군으로서 $R$에 대해 상대적으로 유한 표시일 필요충분조건은
$A$가 $R$ 위에서 유한 표시인 것임도 분명하다. $R$이 뇌터이면 이
개념은 흥미롭지 않다. 이제 우리가 찾던 결과를 정식화할 수 있다.

\begin{lemma}
\label{more-algebra-lemma-finite-extension}
$R$을 환이라 하자. $A \to B$를 유한형 $R$-대수들 사이의 유한
사상이라 하고 $M$을 $B$-가군이라 하자. 그러면 $M$이 $R$에 대해
상대적으로 유한 표시인 $A$-가군일 필요충분조건은 $M$이 $R$에 대해
상대적으로 유한 표시인 $B$-가군인 것이다.
\end{lemma}

\begin{proof}
전사 $R[x_1, \ldots, x_n] \to A$를 택한다. $A$ 위에서 $B$를
생성하는 $y_1, \ldots, y_m \in B$를 택하자. $A \to B$가 유한이므로
각 $y_i$는 계수가 $A$에 있는 일계수 방정식을 만족한다. 따라서
$B$에서 $P_j(y_j) = 0$인 일계수 다항식
$P_j(T) \in R[x_1, \ldots, x_n][T]$를 찾을 수 있다. 그러면 가환 그림
$$
\xymatrix{
R[x_1, \ldots, x_n] \ar[d] \ar[r] &
R[x_1, \ldots, x_n, y_1, \ldots, y_m]/(P_j(y_j)) \ar[d] \\
A \ar[r] & B
}
$$
을 얻는다. 위쪽 화살표는 유한이고 유한 표시인 환 사상이므로,
대수학, 보조정리 \ref{algebra-lemma-finite-finitely-presented-extension}와
정의에 의해 결론을 얻는다.
\end{proof}

\noindent
이 결과에 의해 상대적 개념은 의미가 있고 $R$ 위의 유한형 환들
사이의 유한 사상에 대해 잘 작동한다. 또한 국소화와 밑변환에 대해
안정적이며 잘 ``붙는다''.

\begin{lemma}
\label{more-algebra-lemma-localize-relative-finite-presentation}
% P02-E1009: beyond the missing verb tracked by P02-E0289, the frozen let-list embeds “R_f to A is” without a valid governing construction; Korean supplies the required predicates.
$R$을 환, $f \in R$을 원소라 하고 $R_f \to A$를 유한형 환 사상이라
하자. 또한 $g \in A$이고 $M$이 $A$-가군이라 하자.
% P02-E0289: frozen statement omits the finite verb in its finite-presentation hypothesis; Korean renders the complete condition.
$M$이 $R_f$에 대해 상대적으로 유한 표시이면 $M_g$는 $R$에 대해
상대적으로 유한 표시인 $A_g$-가군이다.
\end{lemma}

\begin{proof}
표시 $R_f[x_1, \ldots, x_n] \to A$를 택한다.
$R_f = R[x_0]/(fx_0 - 1)$이라 쓰자. 주어진 사상을 확장하고 $x_0$를
$1/f$의 상으로, $x_{n + 1}$을 $1/g$로 보내는 표시
$R[x_0, x_1, \ldots, x_n, x_{n + 1}] \to A_g$를 생각하자. $g$로
가는 $g' \in R[x_0, x_1, \ldots, x_n]$를 택한다(이는 가능하다).
행렬 $(h_{ij})$로 주어지는 $M$의 표시가
$$
R_f[x_1, \ldots, x_n]^{\oplus s} \to
R_f[x_1, \ldots, x_n]^{\oplus t} \to M \to 0
$$
이라고 하자. $h_{ij}$로 가는
$h'_{ij} \in R[x_0, x_1, \ldots, x_n]$를 택한다. 그러면
$$
R[x_0, x_1, \ldots, x_n, x_{n + 1}]^{\oplus s + 2t} \to
R[x_0, x_1, \ldots, x_n, x_{n + 1}]^{\oplus t} \to M_g \to 0
$$
은
% P02-E0290: frozen proof calls the displayed presentation ending in M_g a presentation of M_f; Korean preserves the frozen claim.
$M_f$의 표시이다. 이 사상을 정의하는 $t \times (s + 2t)$ 행렬의 첫
$t \times s$ 블록은 행렬 $h'_{ij}$이고, 둘째 $t \times t$ 블록은
% P02-E0291: frozen matrix factor omits the constant term after the minus sign; Korean preserves the exact formula.
$(x_0f - )I_t$이며, 셋째 블록은 $(x_{n + 1}g' - 1)I_t$이다.
\end{proof}

\begin{lemma}
\label{more-algebra-lemma-base-change-relative-finite-presentation}
$R \to A$를 유한형 환 사상이라 하자. $M$을 $R$에 대해 상대적으로
유한 표시인 $A$-가군이라 하자. 임의의 환 사상 $R \to R'$에 대해
$A \otimes_R R'$-가군
% P02-E0292: frozen display uses undefined A' rather than the stated base-changed algebra; Korean preserves the formula.
$$
M \otimes_A A' = M \otimes_R R'
$$
은 $R'$에 대해 상대적으로 유한 표시이다.
\end{lemma}

\begin{proof}
전사 $R[x_1, \ldots, x_n] \to A$와 표시
$$
R[x_1, \ldots, x_n]^{\oplus s} \to
R[x_1, \ldots, x_n]^{\oplus t} \to M \to 0
$$
을 택하자. 그러면
$$
R'[x_1, \ldots, x_n]^{\oplus s} \to
R'[x_1, \ldots, x_n]^{\oplus t} \to M \otimes_R R' \to 0
$$
은 밑변환의 표시이므로 증명이 끝난다.
\end{proof}

\begin{lemma}
\label{more-algebra-lemma-pull-relative-finite-presentation}
$R \to A$를 유한형 환 사상이라 하자. $M$을 $R$에 대해 상대적으로
유한 표시인 $A$-가군이라 하고 $A \to A'$을 유한 표시 환 사상이라
하자. 그러면 $A'$-가군 $M \otimes_A A'$은 $R$에 대해 상대적으로
유한 표시이다.
\end{lemma}

\begin{proof}
전사 $R[x_1, \ldots, x_n] \to A$와 표시
$A' = A[y_1, \ldots, y_m]/(g_1, \ldots, g_l)$을 택한다.
$g_i$로 가는
$g'_i \in R[x_1, \ldots, x_n, y_1, \ldots, y_m]$를 택하자. 행렬
$(h_{ij})$로 주어지는 $M$의 표시가
$$
R[x_1, \ldots, x_n]^{\oplus s} \to
R[x_1, \ldots, x_n]^{\oplus t} \to M \to 0
$$
이라고 하자. 그러면
% P02-E0293: frozen display introduces undeclared x_0 only in the codomain polynomial ring; Korean preserves the exact asymmetric formula.
$$
R[x_1, \ldots, x_n, y_1, \ldots, y_m]^{\oplus s + tl} \to
R[x_0, x_1, \ldots, x_n, y_1, \ldots, y_m]^{\oplus t} \to M \otimes_A A' \to 0
$$
은 $M \otimes_A A'$의 표시이다. 이 사상을 정의하는
$t \times (s + lt)$ 행렬은 행렬 $h_{ij}$로 이루어진 첫
$t \times s$ 블록과, 그 뒤를 잇는 $g'_iI_t$인 크기
$t \times t$의 $l$개 블록으로 이루어진다.
\end{proof}

\begin{lemma}
\label{more-algebra-lemma-composition-relative-finite-presentation}
$R \to A \to B$를 유한형 환 사상들이라 하고 $M$을 $B$-가군이라
하자. $M$이 $A$에 대해 상대적으로 유한 표시이고 $A$가 $R$ 위에서
유한 표시이면 $M$은 $R$에 대해 상대적으로 유한 표시이다.
\end{lemma}

\begin{proof}
전사 $A[x_1, \ldots, x_n] \to B$를 택한다. 행렬 $(h_{ij})$로 주어지는
표시
$$
A[x_1, \ldots, x_n]^{\oplus s} \to
A[x_1, \ldots, x_n]^{\oplus t} \to M \to 0
$$
를 택하고, 표시
$$
A = R[y_1, \ldots, y_m]/(g_1, \ldots, g_u).
$$
를 택한다. $h_{ij}$로 가는
$h'_{ij} \in R[y_1, \ldots, y_m, x_1, \ldots, x_n]$를 택하자.
그러면 표시
$$
R[y_1, \ldots, y_m, x_1, \ldots, x_n]^{\oplus s + tu} \to
R[y_1, \ldots, y_m, x_1, \ldots, x_n]^{\oplus t} \to M \to 0
$$
을 얻는다. 여기서 $t \times (s + tu)$ 행렬은 $h'_{ij}$로 이루어진
첫 $t \times s$ 블록과, 그 뒤를 잇는 $g_iI_t$로 주어지는 크기
$t \times t$의 $u$개 블록으로 이루어지며, $i = 1, \ldots, u$이다.
\end{proof}

\begin{lemma}
\label{more-algebra-lemma-glue-relative-finite-presentation}
$R \to A$를 유한형 환 사상이라 하고 $M$을 $A$-가군이라 하자.
단위 아이디얼을 생성하는 $f_1, \ldots, f_r \in A$를 택하자.
다음은 동치이다.
\begin{enumerate}
\item 각 $M_{f_i}$는 $R$에 대해 상대적으로 유한 표시이다.
\item $M$은 $R$에 대해 상대적으로 유한 표시이다.
\end{enumerate}
\end{lemma}

\begin{proof}
함의 (2) $\Rightarrow$ (1)은 보조정리
\ref{more-algebra-lemma-localize-relative-finite-presentation}에 있다. (1)을
가정하자. $A$에서 $1 = \sum f_ig_i$라 쓰자.
% P02-E0294: frozen source separates a chosen surjection from its dependent “such that” clause; Korean joins the intended condition grammatically.
$y_i$가 $f_i$로 가고 $z_i$가 $g_i$로 가도록 전사
$R[x_1, \ldots, x_n, y_1, \ldots, y_r, z_1, \ldots, z_r] \to A$를
택한다. 그러면 전사
$$
P = R[x_1, \ldots, x_n, y_1, \ldots, y_r, z_1, \ldots, z_r]/(\sum y_iz_i - 1)
\longrightarrow
A.
$$
가 존재한다. 보조정리 \ref{more-algebra-lemma-relatively-finitely-presented}에 의해
$M_{f_i}$는 유한 표시 $A_{f_i}$-가군이고, 따라서 대수학, 보조정리
\ref{algebra-lemma-cover}에 의해 $M$은 유한 표시 $A$-가군이다.
그러므로 $M$은 유한 $P$-가군이다(여기서 $P$는 위와 같다). 전사
$P^{\oplus t} \to M$을 택하자. 이 사상의 핵 $K$가 유한 $P$-가군임을
보여야 한다. $P_{y_i}$가 $A_{f_i}$ 위로 전사하므로, 보조정리
\ref{more-algebra-lemma-relatively-finitely-presented}와 대수학, 보조정리
\ref{algebra-lemma-extension}에 의해 국소화 $K_{y_i}$는 유한 생성
$P_{y_i}$-가군이다. $K_{y_i}$에서 $k_{i, j}$의 상들이 생성하도록 원소들
$k_{i, j} \in K$, $i = 1, \ldots, r$, $j = 1, \ldots, s_i$를 택하자.
원소들 $k_{i, j}$가 생성하는 $P$-부분가군을 $K' \subset K$라 하자.
그러면 모든 $i$에 대해 $K/K'$의 $y_i$에서의 국소화는 영이다.
$(y_1, \ldots, y_r) = P$이므로 원하는 대로 $K/K' = 0$이다.
\end{proof}

\begin{lemma}
\label{more-algebra-lemma-ses-relatively-finite-presentation}
$R \to A$를 유한형 환 사상이라 하고
$0 \to M' \to M \to M'' \to 0$을 $A$-가군들의 짧은 완전열이라 하자.
\begin{enumerate}
\item $M', M''$이 $R$에 대해 상대적으로 유한 표시이면 $M$도
그러하다.
\item $M'$이 유한형 $A$-가군이고 $M$이 $R$에 대해 상대적으로
유한 표시이면 $M''$도 $R$에 대해 상대적으로 유한 표시이다.
\end{enumerate}
\end{lemma}

\begin{proof}
% P02-E0295: frozen proof is a subjectless sentence fragment; Korean supplies the intended complete proof sentence.
이는 대수학, 보조정리 \ref{algebra-lemma-extension}에서 바로 따른다.
\end{proof}

\begin{lemma}
\label{more-algebra-lemma-sum-relatively-finite-presentation}
$R \to A$를 유한형 환 사상이라 하자. $M, M'$을 $A$-가군들이라
하자. $M \oplus M'$이 $R$에 대해 상대적으로 유한 표시이면 $M$과
$M'$도 그러하다.
\end{lemma}

\begin{proof}
생략한다.
\end{proof}







\section{상대 유사연접 가군}
\label{more-algebra-section-relative-pseudo-coherent}

\noindent
이 절은 유사연접성에 대한 절
\ref{more-algebra-section-relative-finite-presentation}의 유사물이다.

\begin{lemma}
\label{more-algebra-lemma-pull-push}
$R$을 환이라 하고 $K^\bullet$을 $R$-가군들의 복합체라 하자.
$x$를 영으로 보내는 $R$-대수 사상 $R[x] \to R$을 생각하자. 그러면
$D(R)$에서
$$
K^\bullet \otimes_{R[x]}^{\mathbf{L}} R \cong K^\bullet \oplus K^\bullet[1]
$$
이다.
\end{lemma}

\begin{proof}
모든 $n$에 대해 $P^n$이 평탄 $R$-가군이 되도록 $R$ 위의 K-평탄
분해 $P^\bullet \to K^\bullet$을 택하자. 보조정리
\ref{more-algebra-lemma-K-flat-resolution}를 보라. 그러면
$P^\bullet \otimes_R R[x]$은 항들이 평탄 $R[x]$-가군인 K-평탄
$R[x]$-가군 복합체이다. 보조정리 \ref{more-algebra-lemma-base-change-K-flat}과
대수학, 보조정리 \ref{algebra-lemma-flat-base-change}를 보라. 특히
$x : P^n \otimes_R R[x] \to P^n \otimes_R R[x]$는 단사이고 그 여핵은
$P^n$과 동형이다. 따라서
$$
P^\bullet \otimes_R R[x] \xrightarrow{x} P^\bullet \otimes_R R[x]
$$
은 연관된 전체 복합체가 $P^\bullet$, 따라서 $K^\bullet$과
유사동형인 $R[x]$-가군들의 이중 복합체이다. 또한 이 연관 전체
복합체는 $R[x]$-가군들의 K-평탄 복합체이다. 예를 들어 보조정리
\ref{more-algebra-lemma-tensor-product-K-flat} 또는 보조정리
\ref{more-algebra-lemma-K-flat-two-out-of-three}를 보라. 따라서 원하는 대로
\begin{align*}
K^\bullet \otimes_{R[x]}^{\mathbf{L}} R
& \cong
\text{Tot}(P^\bullet \otimes_R R[x] \xrightarrow{x} P^\bullet \otimes_R R[x])
\otimes_{R[x]} R =
\text{Tot}(P^\bullet \xrightarrow{0} P^\bullet) \\
& = P^\bullet \oplus P^\bullet[1] \cong K^\bullet \oplus K^\bullet[1]
\end{align*}
이다.
\end{proof}

\begin{lemma}
\label{more-algebra-lemma-add-variable-pseudo-coherent}
$R$을 환, $K^\bullet$을 $R$-가군들의 복합체라 하자.
$m \in \mathbf{Z}$라 하자. $x$를 영으로 보내는 $R$-대수 사상
$R[x] \to R$을 생각하자. 그러면 $K^\bullet$이 $R$-가군 복합체로서
$m$-유사연접일 필요충분조건은 $K^\bullet$이 $R[x]$-가군 복합체로서
$m$-유사연접인 것이다.
\end{lemma}

\begin{proof}
이는 보조정리 \ref{more-algebra-lemma-finite-push-pseudo-coherent}의 특수한 경우이다.
다음과 같은 다른 방식으로도 증명하자.

\medskip\noindent
$0 \to R[x] \to R[x] \to R \to 0$이 완전임에 유의하라. 따라서
$R$은 유사연접 $R[x]$-가군이다. 그러므로 보조정리
\ref{more-algebra-lemma-finite-push-pseudo-coherent}에서 한쪽 함의가 따른다.
다른 함의를 증명하기 위해 $K^\bullet$이 $R[x]$-가군 복합체로서
$m$-유사연접이라고 가정하자. 보조정리
\ref{more-algebra-lemma-pull-pseudo-coherent}에 의해
$K^\bullet \otimes^{\mathbf{L}}_{R[x]} R$은 $R$-가군 복합체로서
$m$-유사연접이다. 보조정리 \ref{more-algebra-lemma-pull-push}에 의해
$K^\bullet \oplus K^\bullet[1]$은 $R$-가군 복합체로서
$m$-유사연접이다. 끝으로 보조정리
\ref{more-algebra-lemma-summands-pseudo-coherent}에 의해 $K^\bullet$은
$R$-가군 복합체로서 $m$-유사연접이다.
\end{proof}

\begin{lemma}
\label{more-algebra-lemma-relatively-pseudo-coherent}
$R \to A$를 유한형 환 사상, $K^\bullet$을 $A$-가군들의 복합체라
하고 $m \in \mathbf{Z}$라 하자. 다음은 동치이다.
\begin{enumerate}
\item 어떤 표시 $\alpha : R[x_1, \ldots, x_n] \to A$에 대해
$K^\bullet$은 $R[x_1, \ldots, x_n]$-가군들의 $m$-유사연접
복합체이다.
\item 모든 표시 $\alpha : R[x_1, \ldots, x_n] \to A$에 대해
$K^\bullet$은 $R[x_1, \ldots, x_n]$-가군들의 $m$-유사연접
복합체이다.
\end{enumerate}
특히 유사연접성에 대해서도 같은 동치가 성립한다.
\end{lemma}

\begin{proof}
% P02-E0296: frozen proof begins with a dangling conditional sentence; Korean gives the intended complete setup.
% P02-E0297: frozen proof reuses m both as the pseudo-coherence bound and as the y-variable count; Korean preserves the symbols.
$\alpha : R[x_1, \ldots, x_n] \to A$와
$\beta : R[y_1, \ldots, y_m] \to A$가 표시들이라고 하자.
$\alpha(f_j) = \beta(y_j)$인 $f_j \in R[x_1, \ldots, x_n]$와
$\beta(g_i) = \alpha(x_i)$인 $g_i \in R[y_1, \ldots, y_m]$를 택한다.
그러면 가환 그림
$$
\xymatrix{
R[x_1, \ldots, x_n, y_1, \ldots, y_m]
\ar[d]^{x_i \mapsto g_i} \ar[rr]_-{y_j \mapsto f_j} & &
R[x_1, \ldots, x_n] \ar[d] \\
R[y_1, \ldots, y_m] \ar[rr] & & A
}
$$
을 얻는다. 좌표를 바꾸면 환 준동형
$R[x_1, \ldots, x_n, y_1, \ldots, y_m] \to R[x_1, \ldots, x_n]$은
각 $y_i$를 영으로 보내는 환 준동형과 동형이다. 그림의 왼쪽 수직
사상도 마찬가지이다. 따라서 변수 개수에 대한 귀납법과 보조정리
\ref{more-algebra-lemma-add-variable-pseudo-coherent}에 의해 보조정리가 따른다.
유사연접인 경우는 이것과 보조정리 \ref{more-algebra-lemma-pseudo-coherent}에서
따른다.
\end{proof}

\begin{definition}
\label{more-algebra-definition-relatively-pseudo-coherent}
$R \to A$를 유한형 환 사상, $K^\bullet$을 $A$-가군들의 복합체,
$M$을 $A$-가군이라 하고 $m \in \mathbf{Z}$라 하자.
\begin{enumerate}
\item 보조정리 \ref{more-algebra-lemma-relatively-pseudo-coherent}의 동치 조건들이
성립하면 $K^\bullet$이 {\it $R$에 대해 상대적으로 $m$-유사연접}이라고
한다.
\item 모든 $m \in \mathbf{Z}$에 대해 $K^\bullet$이 $R$에 대해
상대적으로 $m$-유사연접이면 $K^\bullet$이 {\it $R$에 대해 상대적으로
유사연접}이라고 한다.
\item $M[0]$이 $R$에 대해 상대적으로 $m$-유사연접이면 $M$이
{\it $R$에 대해 상대적으로 $m$-유사연접}이라고 한다.
\item $M[0]$이 $R$에 대해 상대적으로 유사연접이면 $M$이
{\it $R$에 대해 상대적으로 유사연접}이라고 한다.
\end{enumerate}
\end{definition}

\noindent
% P02-E0298: frozen explanation names an x-polynomial module structure for a y-polynomial presentation; Korean preserves both formal spans.
(2)는 임의의 전사 $R[y_1, \ldots, y_m] \to A$에 대해
$K^\bullet$이 $R[x_1, \ldots, x_n]$-가군들의 복합체로서 유사연접이라는
뜻이다. 보조정리 \ref{more-algebra-lemma-pseudo-coherent}를 보라. 이 정의에는
다음과 같은 좋은 성질이 있다.

\begin{lemma}
\label{more-algebra-lemma-finite-extension-pseudo-coherent}
$R$을 환이라 하자. $A \to B$를 유한형 $R$-대수들 사이의 유한
사상이라 하자. $m \in \mathbf{Z}$이고 $K^\bullet$이 $B$-가군들의
복합체라 하자. 그러면 $K^\bullet$이 $R$에 대해 상대적으로
$m$-유사연접(각각 유사연접)일 필요충분조건은 $A$-가군들의 복합체로
본 $K^\bullet$이 $R$에 대해 상대적으로 $m$-유사연접(유사연접)인
것이다.
\end{lemma}

\begin{proof}
전사 $R[x_1, \ldots, x_n] \to A$를 택한다.
% P02-E0299: frozen proof reuses m as both the pseudo-coherence bound and the number of generators of B; Korean preserves the symbols.
$A$ 위에서 $B$를 생성하는 $y_1, \ldots, y_m \in B$를 택하자.
$A \to B$가 유한이므로 각 $y_i$는 계수가 $A$에 있는 일계수 방정식을
만족한다. 따라서 $B$에서 $P_j(y_j) = 0$인 일계수 다항식
$P_j(T) \in R[x_1, \ldots, x_n][T]$를 찾을 수 있다. 그러면 가환 그림
$$
\xymatrix{
& R[x_1, \ldots, x_n, y_1, \ldots, y_m] \ar[d] \\
R[x_1, \ldots, x_n] \ar[d] \ar[r] &
R[x_1, \ldots, x_n, y_1, \ldots, y_m]/(P_j(y_j)) \ar[d] \\
A \ar[r] & B
}
$$
을 얻는다. 위쪽 수평 화살표와 오른쪽 위 수직 화살표는 보조정리
\ref{more-algebra-lemma-finite-push-pseudo-coherent}의 가정을 만족한다. 따라서
$K^\bullet$이 $R[x_1, \ldots, x_n]$-가군들의 복합체로서
$m$-유사연접(각각 유사연접)일 필요충분조건은 $K^\bullet$이
$R[x_1, \ldots, x_n, y_1, \ldots, y_m]$-가군들의 복합체로서
$m$-유사연접(각각 유사연접)인 것이다.
\end{proof}

\begin{lemma}
\label{more-algebra-lemma-cone-relatively-pseudo-coherent}
$R$을 환, $R \to A$를 유한형 환 사상, $m \in \mathbf{Z}$라 하자.
$(K^\bullet, L^\bullet, M^\bullet, f, g, h)$를 $D(A)$의 완전
삼각형이라 하자.
\begin{enumerate}
\item $K^\bullet$이 $R$에 대해 상대적으로 $(m + 1)$-유사연접이고
$L^\bullet$이 $R$에 대해 상대적으로 $m$-유사연접이면
$M^\bullet$은 $R$에 대해 상대적으로 $m$-유사연접이다.
\item $K^\bullet, M^\bullet$이 $R$에 대해 상대적으로
$m$-유사연접이면 $L^\bullet$도 $R$에 대해 상대적으로
$m$-유사연접이다.
\item $L^\bullet$이 $R$에 대해 상대적으로 $(m + 1)$-유사연접이고
$M^\bullet$이 $R$에 대해 상대적으로 $m$-유사연접이면
$K^\bullet$은 $R$에 대해 상대적으로 $(m + 1)$-유사연접이다.
\end{enumerate}
% P02-E0300: frozen source says “the so is the third”; Korean renders the intended two-out-of-three conclusion.
또한 $K^\bullet, L^\bullet, M^\bullet$ 중 둘이 $R$에 대해 상대적으로
유사연접이면 셋째도 그러하다.
\end{lemma}

\begin{proof}
% P02-E0301: frozen proof is a subjectless fragment; Korean supplies a complete sentence.
이는 보조정리 \ref{more-algebra-lemma-cone-pseudo-coherent}와 정의들에서 바로 따른다.
\end{proof}

\begin{lemma}
\label{more-algebra-lemma-rel-n-pseudo-module}
$R \to A$를 유한형 환 사상이라 하고 $M$을 $A$-가군이라 하자.
그러면 다음이 성립한다.
\begin{enumerate}
\item $M$이 $R$에 대해 상대적으로 $0$-유사연접일 필요충분조건은
$M$이 유한형 $A$-가군인 것이다.
% P02-E0302: frozen item inserts an article before the relative finite-presentation predicate; Korean renders the defined predicate naturally.
\item $M$이 $R$에 대해 상대적으로 $(-1)$-유사연접일 필요충분조건은
$M$이 $R$에 대해 상대적으로 유한 표시인 것이다.
% P02-E0303: frozen statement uses d without introducing or quantifying it; Korean preserves the unbound formal symbol.
\item $M$이 $R$에 대해 상대적으로 $(-d)$-유사연접일 필요충분조건은
모든 전사 $R[x_1, \ldots, x_n] \to A$에 대해 길이가 $d$인 분해
$$
R[x_1, \ldots, x_n]^{\oplus a_d} \to R[x_1, \ldots, x_n]^{\oplus a_{d - 1}}
\to \ldots \to R[x_1, \ldots, x_n]^{\oplus a_0} \to M \to 0
$$
가 존재하는 것이다.
\item $M$이 $R$에 대해 상대적으로 유사연접일 필요충분조건은 모든
표시 $R[x_1, \ldots, x_n] \to A$에 대해 유한 자유
$R[x_1, \ldots, x_n]$-가군들에 의한 무한 분해
$$
\ldots \to R[x_1, \ldots, x_n]^{\oplus a_1} \to
R[x_1, \ldots, x_n]^{\oplus a_0} \to M \to 0
$$
가 존재하는 것이다.
\end{enumerate}
\end{lemma}

\begin{proof}
% P02-E0304: one of four consecutive frozen proofs is a subjectless fragment; Korean supplies a complete sentence.
이는 보조정리 \ref{more-algebra-lemma-n-pseudo-module}와 정의들에서 바로 따른다.
\end{proof}

\begin{lemma}
\label{more-algebra-lemma-summands-relative-pseudo-coherent}
$R \to A$를 유한형 환 사상, $m \in \mathbf{Z}$라 하고
$K^\bullet, L^\bullet \in D(A)$라 하자.
$K^\bullet \oplus L^\bullet$이 $R$에 대해 상대적으로
$m$-유사연접(각각 유사연접)이면 $K^\bullet$과 $L^\bullet$도 그러하다.
\end{lemma}

\begin{proof}
% P02-E0304: the second subjectless frozen proof is rendered as a complete Korean sentence.
이는 보조정리 \ref{more-algebra-lemma-summands-pseudo-coherent}와 정의들에서 바로
따른다.
\end{proof}

\begin{lemma}
\label{more-algebra-lemma-complex-relative-pseudo-coherent-modules}
$R \to A$를 유한형 환 사상, $m \in \mathbf{Z}$라 하자.
$K^\bullet$을 위로 유계인 $A$-가군 복합체라 하고, 모든 $i$에 대해
$K^i$가 $R$에 대해 상대적으로 $(m - i)$-유사연접이라고 하자.
그러면 $K^\bullet$은 $R$에 대해 상대적으로 $m$-유사연접이다. 특히
% P02-E0305: frozen prose omits the relative-clause link in the termwise hypothesis; Korean makes that hypothesis explicit.
$K^\bullet$이 각 항이 $R$에 대해 상대적으로 유사연접인 $A$-가군들의
위로 유계인 복합체이면 $K^\bullet$은 $R$에 대해 상대적으로
유사연접이다.
\end{lemma}

\begin{proof}
% P02-E0304: the third subjectless frozen proof is rendered as a complete Korean sentence.
이는 보조정리 \ref{more-algebra-lemma-complex-pseudo-coherent-modules}와 정의들에서
바로 따른다.
\end{proof}

\begin{lemma}
\label{more-algebra-lemma-cohomology-relative-pseudo-coherent}
$R \to A$를 유한형 환 사상, $m \in \mathbf{Z}$라 하자.
$K^\bullet \in D^{-}(A)$이고 모든 $i$에 대해 $H^i(K^\bullet)$이
$R$에 대해 상대적으로 $(m - i)$-유사연접(각각 유사연접)이라고 하자.
그러면 $K^\bullet$은 $R$에 대해 상대적으로
$m$-유사연접(각각 유사연접)이다.
\end{lemma}

\begin{proof}
% P02-E0304: the fourth subjectless frozen proof is rendered as a complete Korean sentence.
이는 보조정리 \ref{more-algebra-lemma-cohomology-pseudo-coherent}와 정의들에서 바로
따른다.
\end{proof}

\begin{lemma}
\label{more-algebra-lemma-localize-relative-pseudo-coherent}
% P02-E0306: frozen opening mixes let-clauses with an embedded finite verb and lacks parallel syntax; Korean introduces each datum grammatically.
% P02-E1018: the later candidate records the same broader grammatical failure; no additional hypothesis is introduced.
$R$을 환, $f \in R$을 원소라 하고 $R_f \to A$를 유한형 환 사상이라
하자. $g \in A$이고 $K^\bullet$이 $A$-가군들의 복합체라 하자.
$K^\bullet$이 $R_f$에 대해 상대적으로 $m$-유사연접(각각
유사연접)이면 $K^\bullet \otimes_A A_g$는 $R$에 대해 상대적으로
$m$-유사연접(각각 유사연접)이다.
\end{lemma}

\begin{proof}
먼저 $K^\bullet$이 $R$에 대해 상대적으로 $m$-유사연접임을 보이자.
구체적으로 $R_f[x_1, \ldots, x_n] \to A$가 전사라고 하자.
$R_f = R[x_0]/(fx_0 - 1)$이라 쓰자. 그러면
$R[x_0, x_1, \ldots, x_n] \to A$는 전사이고
$R_f[x_1, \ldots, x_n]$은 유사연접
$R[x_0, \ldots, x_n]$-가군이다. 따라서 보조정리
\ref{more-algebra-lemma-finite-push-pseudo-coherent}에 의해 $K^\bullet$은
$R[x_0, x_1, \ldots, x_n]$-가군 복합체로서 $m$-유사연접이다.

\medskip\noindent
$g \in A$로 가는 원소 $g' \in R[x_0, x_1, \ldots, x_n]$를 택하자.
보조정리 \ref{more-algebra-lemma-pull-pseudo-coherent}에 의해 다음이
$R[x_0, x_1, \ldots, x_n, \frac{1}{g'}]$-가군 복합체로서
$m$-유사연접이다.
% P02-E0307: frozen localization identity ends in A_f although g' maps to g and the lemma concerns A_g; Korean preserves the formula.
\begin{align*}
K^\bullet \otimes_{R[x_0, x_1, \ldots, x_n]}^{\mathbf{L}}
R[x_0, x_1, \ldots, x_n, \frac{1}{g'}] & =
K^\bullet \otimes_{R[x_0, x_1, \ldots, x_n]}
R[x_0, x_1, \ldots, x_n, \frac{1}{g'}] \\
& = K^\bullet \otimes_A A_f
\end{align*}
% P02-E0308: frozen source begins the next sentence with lowercase “write”; Korean starts it normally.
다음과 같이 쓰자.
$$
R[x_0, x_1, \ldots, x_n, \frac{1}{g'}] =
R[x_0, \ldots, x_n, x_{n + 1}]/(x_{n + 1}g' - 1).
$$
$R[x_0, x_1, \ldots, x_n, \frac{1}{g'}]$은 유사연접
$R[x_0, \ldots, x_n, x_{n + 1}]$-가군이므로, 보조정리
\ref{more-algebra-lemma-finite-push-pseudo-coherent}에 의해
$K^\bullet \otimes_A A_g$는 원하는 대로
$R[x_0, \ldots, x_n, x_{n + 1}]$-가군 복합체로서 $m$-유사연접이다.
\end{proof}

\begin{lemma}
\label{more-algebra-lemma-base-change-relative-pseudo-coherent}
$R \to A$를 유한형 환 사상, $m \in \mathbf{Z}$라 하자.
$K^\bullet$을 $R$에 대해 상대적으로 $m$-유사연접(각각 유사연접)인
$A$-가군 복합체라 하자. $A$와 $R'$이 $R$ 위에서 Tor-독립인 환 사상
$R \to R'$을 택하고 $A' = A \otimes_R R'$이라 하자. 그러면
$K^\bullet \otimes_A^{\mathbf{L}} A'$은 $R'$에 대해 상대적으로
$m$-유사연접(각각 유사연접)이다.
\end{lemma}

\begin{proof}
전사 $R[x_1, \ldots, x_n] \to A$를 택하자. 보조정리
\ref{more-algebra-lemma-base-change-comparison}를 두 번 적용하면
$$
K^\bullet \otimes_A^{\mathbf{L}} A' =
K^\bullet \otimes_R^{\mathbf{L}} R' =
K^\bullet \otimes_{R[x_1, \ldots, x_n]}^{\mathbf{L}} R'[x_1, \ldots, x_n]
$$
이다. 따라서 보조정리 \ref{more-algebra-lemma-pull-pseudo-coherent}에 의해 증명이
끝난다.
\end{proof}

\begin{lemma}
\label{more-algebra-lemma-pull-relative-pseudo-coherent}
$R \to A \to B$를 유한형 환 사상들이라 하고 $m \in \mathbf{Z}$라
하자. $K^\bullet$을 $A$-가군들의 복합체라 하자. $B$가 $B$-가군으로서
$A$에 대해 상대적으로 유사연접이라고 가정하자. $K^\bullet$이 $R$에
대해 상대적으로 $m$-유사연접(각각 유사연접)이면
$K^\bullet \otimes_A^{\mathbf{L}} B$도 $R$에 대해 상대적으로
$m$-유사연접(각각 유사연접)이다.
\end{lemma}

\begin{proof}
% P02-E0309: frozen proof reuses m as both the pseudo-coherence bound and the y-variable count; Korean preserves the formal symbols.
전사 $A[y_1, \ldots, y_m] \to B$와 전사
$R[x_1, \ldots, x_n] \to A$를 택하자. 이를 합치면 전사
% P02-E0310: frozen source omits the comma after \ldots in two polynomial-variable lists; Korean preserves the exact spans.
$R[x_1, \ldots, x_n, y_1, \ldots y_m] \to B$를 얻는다.
% P02-E0311: frozen resolution ring switches from y_m to unrelated y_n; Korean preserves the coefficient ring exactly.
$B$의 유한 자유 $A[y_1, \ldots, y_n]$-가군 복합체에 의한 분해
$E^\bullet \to B$를 택하자(환 사상 $A \to B$에 대한 가정으로
가능하다). $K^\bullet$이 평탄 $A$-가군들로 이루어진 위로 유계인
복합체라고 가정할 수 있다. 그러면
\begin{align*}
K^\bullet \otimes_A^{\mathbf{L}} B & =
\text{Tot}(K^\bullet \otimes_A B[0]) \\
& = \text{Tot}(K^\bullet \otimes_A A[y_1, \ldots, y_m]
\otimes_{A[y_1, \ldots, y_m]} B[0]) \\
& \cong
\text{Tot}\left(
(K^\bullet \otimes_A A[y_1, \ldots, y_m])
\otimes_{A[y_1, \ldots, y_m]} E^\bullet
\right) \\
& =
\text{Tot}(K^\bullet \otimes_A E^\bullet)
\end{align*}
이다. 이 $D(A[y_1, \ldots, y_m])$의 유사동형 $\cong$은 보조정리
\ref{more-algebra-lemma-derived-tor-quasi-isomorphism}을 적용하여 얻는다. 따라서
$\text{Tot}(K^\bullet \otimes_A E^\bullet)$이
$R[x_1, \ldots, x_n, y_1, \ldots y_m]$-가군 복합체로서
$m$-유사연접임을 보여야 한다.
$\text{Tot}(K^\bullet \otimes_A E^\bullet)$에는 연속 몫들이 복합체
$K^\bullet \otimes_A E^i[-i]$인 부분복합체들에 의한 여과가 있다.
% P02-E1023: frozen proof gives vanishing in degrees <= m, while its shifts and bounded-above hypothesis indicate the opposite direction; Korean preserves the inequality.
$i \ll 0$이면 복합체 $K^\bullet \otimes_A E^i[-i]$는 차수
$\leq m$에서 코호몰로지가 영이고, 따라서 (어느 환 위에서도)
$m$-유사연접이다.
% P02-E0312: frozen induction asks for full relative pseudo-coherence although the hypothesis and reduction provide only m-pseudo-coherence; Korean preserves the strength mismatch.
그러므로 보조정리 \ref{more-algebra-lemma-cone-relatively-pseudo-coherent}와 귀납법에
의해, 모든 $i$에 대해 $K^\bullet \otimes_A E^i[-i]$가 $R$에 대해
상대적으로 유사연접임을 보이면 충분하다. $i > 0$이면
$E^i = 0$임에 유의하라. 또한 $E^i$가 유한 자유이므로, 이는
$K^\bullet \otimes_A A[y_1, \ldots, y_m]$이 $R$에 대해 상대적으로
$m$-유사연접임을 보이는 것으로 귀착되며, 예를 들어 보조정리
\ref{more-algebra-lemma-base-change-relative-pseudo-coherent}에서 따른다.
\end{proof}

\begin{lemma}
\label{more-algebra-lemma-pull-relative-pseudo-coherent-module}
$R \to A \to B$를 유한형 환 사상들, $m \in \mathbf{Z}$라 하고
$M$을 $A$-가군이라 하자. $B$가 $A$ 위에서 평탄이고 $B$가
$B$-가군으로서 $A$에 대해 상대적으로 유사연접이라고 가정하자.
$M$이 $R$에 대해 상대적으로 $m$-유사연접(각각 유사연접)이면
$M \otimes_A B$도 $R$에 대해 상대적으로
$m$-유사연접(각각 유사연접)이다.
\end{lemma}

\begin{proof}
% P02-E0313: frozen proof is a subjectless fragment; Korean supplies a complete sentence.
이는 보조정리 \ref{more-algebra-lemma-pull-relative-pseudo-coherent}에서 바로 따른다.
\end{proof}

\begin{lemma}
\label{more-algebra-lemma-composition-relative-pseudo-coherent}
$R$을 환이라 하자. $A \to B$를 유한형 $R$-대수들의 사상이라 하고
$m \in \mathbf{Z}$라 하자. $K^\bullet$을 $B$-가군들의 복합체라 하자.
% P02-E0314: frozen proof below reuses m as both the pseudo-coherence bound and the y-variable count; Korean preserves the symbols.
$A$가 $R$에 대해 상대적으로 유사연접이라고 가정하자. 다음은
동치이다.
\begin{enumerate}
\item $K^\bullet$은 $A$에 대해 상대적으로
$m$-유사연접(각각 유사연접)이다.
\item $K^\bullet$은 $R$에 대해 상대적으로
$m$-유사연접(각각 유사연접)이다.
\end{enumerate}
\end{lemma}

\begin{proof}
전사 $R[x_1, \ldots, x_n] \to A$와 전사
$A[y_1, \ldots, y_m] \to B$를 택하자. 그러면
$R[x_1, \ldots, x_n] \to A$의 평탄 밑변환인 전사
$$
R[x_1, \ldots, x_n, y_1, \ldots, y_m] \to A[y_1, \ldots, y_m]
$$
을 얻는다. 가정에 의해 $A$는 유사연접
$R[x_1, \ldots, x_n]$-가군이므로 보조정리
\ref{more-algebra-lemma-flat-base-change-pseudo-coherent}에 의해
$A[y_1, \ldots, y_m]$은
$R[x_1, \ldots, x_n, y_1, \ldots, y_m]$ 위에서 유사연접이다. 따라서
보조정리 \ref{more-algebra-lemma-finite-push-pseudo-coherent}와 정의들에서 보조정리가
따른다.
\end{proof}

\begin{lemma}
\label{more-algebra-lemma-glue-relative-pseudo-coherent}
$R \to A$를 유한형 환 사상, $K^\bullet$을 $A$-가군들의 복합체,
$m \in \mathbf{Z}$라 하자. 단위 아이디얼을 생성하는
$f_1, \ldots, f_r \in A$를 택하자. 다음은 동치이다.
\begin{enumerate}
\item 각 $K^\bullet \otimes_A A_{f_i}$는 $R$에 대해 상대적으로
$m$-유사연접이다.
\item $K^\bullet$은 $R$에 대해 상대적으로 $m$-유사연접이다.
\end{enumerate}
$R$에 대한 상대 유사연접성에 대해서도 같은 동치가 성립한다.
\end{lemma}

\begin{proof}
함의 (2) $\Rightarrow$ (1)은 보조정리
\ref{more-algebra-lemma-localize-relative-pseudo-coherent}에 있다. (1)을 가정하자.
$A$에서 $1 = \sum f_ig_i$라 쓰자.
% P02-E0315: frozen source separates the chosen surjection from its dependent condition; Korean joins the intended sentence.
$y_i$가 $f_i$로 가고 $z_i$가 $g_i$로 가도록 전사
$R[x_1, \ldots, x_n, y_1, \ldots, y_r, z_1, \ldots, z_r] \to A$를
택하자. 그러면 전사
$$
P = R[x_1, \ldots, x_n, y_1, \ldots, y_r, z_1, \ldots, z_r]/(\sum y_iz_i - 1)
\longrightarrow
A.
$$
가 존재한다. $P$는 유사연접
$R[x_1, \ldots, x_n, y_1, \ldots, y_r, z_1, \ldots, z_r]$-가군이고
$P[1/y_i]$는 유사연접
$R[x_1, \ldots, x_n, y_1, \ldots, y_r, z_1, \ldots, z_r, 1/y_i]$-가군이다.
따라서 보조정리 \ref{more-algebra-lemma-finite-push-pseudo-coherent}에 의해 각
$i$에 대해 $K^\bullet \otimes_A A_{f_i}$는
$P[1/y_i]$-가군들의 $m$-유사연접 복합체이다.
% P02-E0316: frozen fixed-m proof twice concludes full pseudo-coherence instead of m-pseudo-coherence; Korean preserves both strength mismatches.
그러므로 보조정리 \ref{more-algebra-lemma-glue-pseudo-coherent}에 의해
$K^\bullet$은 $P$-가군 복합체로서 유사연접이고, 보조정리
\ref{more-algebra-lemma-finite-push-pseudo-coherent}에 의해 $K^\bullet$은
$R[x_1, \ldots, x_n, y_1, \ldots, y_r, z_1, \ldots, z_r]$-가군들의
복합체로서 유사연접이다.
\end{proof}

\begin{lemma}
\label{more-algebra-lemma-Noetherian-relative-pseudo-coherent}
$R$을 뇌터 환이라 하고 $R \to A$를 유한형 환 사상이라 하자. 그러면
다음이 성립한다.
\begin{enumerate}
\item $A$-가군들의 복합체 $K^\bullet$이 $R$에 대해 상대적으로
$m$-유사연접일 필요충분조건은 $K^\bullet \in D^{-}(A)$이고
$i \geq m$일 때 $H^i(K^\bullet)$이 유한 $A$-가군인 것이다.
\item $A$-가군들의 복합체 $K^\bullet$이 $R$에 대해 상대적으로
유사연접일 필요충분조건은 $K^\bullet \in D^{-}(A)$이고 모든 $i$에
대해 $H^i(K^\bullet)$이 유한 $A$-가군인 것이다.
\item $A$-가군이 $R$에 대해 상대적으로 유사연접일 필요충분조건은
그것이 유한인 것이다.
\end{enumerate}
\end{lemma}

\begin{proof}
% P02-E0317: frozen proof is a subjectless fragment; Korean supplies a complete sentence.
이는 보조정리 \ref{more-algebra-lemma-Noetherian-pseudo-coherent}와 정의들의 즉각적인
귀결이다.
\end{proof}










\section{유사연접 환 사상과 퍼펙트 환 사상}
\label{more-algebra-section-pseudo-coherent-perfect-ring-map}

\noindent
이러한 종류의 환 사상들을 다음과 같이 정의할 수 있다.

\begin{definition}
\label{more-algebra-definition-pseudo-coherent-perfect}
$A \to B$를 환 사상이라 하자.
\begin{enumerate}
\item $A \to B$가 유한형이고 $B$가 $B$-가군으로서 $A$에 대해
상대적으로 유사연접이면 이를 {\it 유사연접 환 사상}이라 한다.
\item $A \to B$가 유사연접 환 사상이고 $B$가 $A$-가군으로서 유한
Tor 차원을 가지면 이를 {\it 퍼펙트 환 사상}이라 한다.
\end{enumerate}
\end{definition}

\noindent
이 용어는 표준적이지 않을 수 있다. 보조정리
\ref{more-algebra-lemma-rel-n-pseudo-module}을 사용하면, $A \to B$가 유사연접일
필요충분조건은 $B = A[x_1, \ldots, x_n]/I$이고 $B$가
$A[x_1, \ldots, x_n]$-가군으로서 유한 자유
$A[x_1, \ldots, x_n]$-가군들에 의한 분해를 갖는 것임을 알 수 있다.
퍼펙트 환 사상의 정의는 보조정리 \ref{more-algebra-lemma-perfect}에서 동기를 얻었다.
% P02-E0318: frozen source says “The following lemmas gives” although one lemma follows; Korean uses the intended singular introduction.
다음 보조정리는 퍼펙트 환 사상에 대한 더 유용하고 직관적인 특성화를
준다.

\begin{lemma}
\label{more-algebra-lemma-perfect-ring-map}
환 사상 $A \to B$가 퍼펙트일 필요충분조건은
$B = A[x_1, \ldots, x_n]/I$이고 $B$가
$A[x_1, \ldots, x_n]$-가군으로서 유한 사영
$A[x_1, \ldots, x_n]$-가군들에 의한 유한 분해를 갖는 것이다.
\end{lemma}

\begin{proof}
$A \to B$가 퍼펙트이면 $B = A[x_1, \ldots, x_n]/I$이고, $B$는
$A[x_1, \ldots, x_n]$-가군으로서 유사연접이며 $A$-가군으로서 유한
Tor 차원을 갖는다. 따라서 보조정리
\ref{more-algebra-lemma-perfect-over-polynomial-ring}에 의해
% P02-E0319: frozen English uses the wrong article before A[x]; Korean is article-free and renders the intended module predicate.
$B$는 $A[x_1, \ldots, x_n]$-가군으로서 퍼펙트이다. 즉 유한 사영
$A[x_1, \ldots, x_n]$-가군들에 의한 유한 분해를 갖는다(보조정리
\ref{more-algebra-lemma-perfect-module}). 반대로 $B = A[x_1, \ldots, x_n]/I$이고
$B$가 $A[x_1, \ldots, x_n]$-가군으로서 유한 사영
$A[x_1, \ldots, x_n]$-가군들에 의한 유한 분해를 가지면, $B$는
$A[x_1, \ldots, x_n]$-가군으로서 유사연접이고 따라서 $A \to B$는
유사연접이다. 또한 $A[x_1, \ldots, x_n]$ 위의 주어진 분해는 평탄
$A$-가군들에 의한 유한 분해이므로 $B$는 $A$-가군으로서 유한 Tor
차원을 갖는다.
\end{proof}

\noindent
앞 절들의 많은 결과를 이 용어로 다시 정식화할 수 있다.
% P02-E0320: frozen source has duplicated interword spacing here; Korean uses normal spacing.
스킴 사상에 관한 대응하는 논의는 사상론의 심화, 절
\ref{more-morphisms-section-pseudo-coherent}와
\ref{more-morphisms-section-perfect}도 참고하라.

\begin{lemma}
\label{more-algebra-lemma-Noetherian-pseudo-coherent-ring-map}
뇌터 환들 사이의 유한형 환 사상은 유사연접이다.
\end{lemma}

\begin{proof}
보조정리 \ref{more-algebra-lemma-Noetherian-relative-pseudo-coherent}를 보라.
\end{proof}

\begin{lemma}
\label{more-algebra-lemma-flat-finite-presentation-perfect}
평탄하고 유한 표시인 환 사상은 퍼펙트이다.
\end{lemma}

\begin{proof}
$A \to B$를 평탄하고 유한 표시인 환 사상이라 하자. $B$가 유한 Tor
차원을 가짐은 분명하다. 대수학, 보조정리
\ref{algebra-lemma-flat-finite-presentation-limit-flat}에 의해 유한형
$\mathbf{Z}$-대수 $A_0 \subset A$와 평탄 유한형 환 사상
$A_0 \to B_0$가 존재하여 $B = B_0 \otimes_{A_0} A$이다. 보조정리
\ref{more-algebra-lemma-Noetherian-relative-pseudo-coherent}에 의해
$A_0 \to B_0$는 유사연접이다. $A_0 \to B_0$가 평탄하므로 $B_0$와
$A$는 $A_0$ 위에서 Tor 독립이고, 따라서 보조정리
\ref{more-algebra-lemma-base-change-relative-pseudo-coherent}를 사용하면
$A \to B$가 유사연접임을 얻는다.
\end{proof}

\begin{lemma}
\label{more-algebra-lemma-regular-perfect-ring-map}
$A$가 차원이 유한한 정칙환인 유한형 환 사상 $A \to B$를 택하자.
그러면 $A \to B$는 퍼펙트이다.
\end{lemma}

\begin{proof}
대수학, 보조정리
\ref{algebra-lemma-finite-gl-dim-finite-dim-regular}에 의해 $A$에 대한
가정은 $A$의 대역 차원이 유한하다는 뜻이다. 따라서 모든 가군은 유한
Tor 차원을 가지며, 특히 $B$도 그러하다. 보조정리
\ref{more-algebra-lemma-finite-gl-dim-tor-dimension}를 보라. 보조정리
\ref{more-algebra-lemma-Noetherian-pseudo-coherent-ring-map}에 의해 이 사상은
유사연접이다.
\end{proof}

\begin{lemma}
\label{more-algebra-lemma-lci-perfect}
국소 완전 교차 준동형은 퍼펙트이다.
\end{lemma}

\begin{proof}
$A \to B$를 국소 완전 교차 준동형이라 하자. 정의
\ref{more-algebra-definition-local-complete-intersection}에 의해 이는
$B = A[x_1, \ldots, x_n]/I$이고 $I$가
$A[x_1, \ldots, x_n]$의 Koszul 아이디얼이라는 뜻이다. 보조정리
\ref{more-algebra-lemma-perfect-ring-map}과 \ref{more-algebra-lemma-perfect-module}에 의해
$I$가 퍼펙트 $A[x_1, \ldots, x_n]$-가군임을 보이면 충분하다.
보조정리 \ref{more-algebra-lemma-glue-perfect}에 의해 이는 국소적 문제이다. 따라서
정의 \ref{more-algebra-definition-regular-ideal}에 의해 $I$가 Koszul-정칙 열로
생성된다고 가정할 수 있다. 물론 이는 $I$가 유한 자유 분해를 갖는다는
뜻이므로 증명이 끝난다.
\end{proof}

\begin{lemma}
\label{more-algebra-lemma-relative-pseudo-coherent-is-moot}
$R \to A$를 유사연접 환 사상이라 하고 $K \in D(A)$라 하자. 다음은
동치이다.
\begin{enumerate}
\item $K$는 $R$에 대해 상대적으로 $m$-유사연접(각각 유사연접)이다.
\item $K$는 $D(A)$에서 $m$-유사연접(각각 유사연접)이다.
\end{enumerate}
\end{lemma}

\begin{proof}
이는 보조정리 \ref{more-algebra-lemma-composition-relative-pseudo-coherent}의 특수한
경우를 다시 표현한 것이다.
\end{proof}

\begin{lemma}
\label{more-algebra-lemma-more-relative-pseudo-coherent-is-moot}
$R \to B \to A$를 환 사상들이라 하자. $\varphi : B \to A$는 전사이고
$R \to B$와 $R \to A$는 평탄하고 유한 표시라고 하자.
% P02-E0321: frozen source omits “by” when introducing the restriction notation; Korean introduces it grammatically.
$K \in D(A)$에 대해 그 제한을 $\varphi_*K \in D(B)$로 나타내자.
다음은 동치이다.
\begin{enumerate}
\item $K$는 유사연접이다.
\item $K$는 $R$에 대해 상대적으로 유사연접이다.
\item $K$는 $A$에 대해 상대적으로 유사연접이다.
\item $\varphi_*K$는 유사연접이다.
\item $\varphi_*K$는 $R$에 대해 상대적으로 유사연접이다.
\end{enumerate}
% P02-E0322: frozen source uses “Similar” as the subject; Korean renders the intended same-statement construction.
$m$-유사연접성에 대해서도 같은 진술이 성립한다.
\end{lemma}

\begin{proof}
$R \to A$와 $R \to B$는 퍼펙트 환 사상(보조정리
\ref{more-algebra-lemma-flat-finite-presentation-perfect})이고, 따라서 특히
유사연접 환 사상이다. 그러므로 보조정리
\ref{more-algebra-lemma-relative-pseudo-coherent-is-moot}에 의해
(1) $\Leftrightarrow$ (2)이고 (4) $\Leftrightarrow$ (5)이다.

\medskip\noindent
$A$가 $R$에 대해 상대적으로 유사연접임을 사용하고 보조정리
\ref{more-algebra-lemma-composition-relative-pseudo-coherent}를 적용하면
(2) $\Leftrightarrow$ (3)을 얻는다. 그러나
% P02-E0323: frozen proof reverses the assumed surjection B to A; Korean preserves the formal arrow A to B.
$A \to B$가 전사이므로 정의
\ref{more-algebra-definition-relatively-pseudo-coherent}에서 직접
% P02-E0324: frozen English says “equivalent with”; Korean uses the normal equivalence construction.
(3)과 (4)가 동치임을 알 수 있다.
\end{proof}


\section{상대 퍼펙트 가군}
\label{more-algebra-section-relatively-perfect}

\noindent
이 절은 유도 범주의 퍼펙트 대상에 대하여 절
\ref{more-algebra-section-relative-pseudo-coherent}와 유사한 이론을 다룬다.
% P02-E0325: frozen source begins the new sentence with lowercase “we”; Korean begins it normally.
일반적인 경우에 올바른 정의가 무엇인지 확실하지 않으므로 이 개념은
제한된 일반성에서만 정의한다. 이에 대한 논의는 Schemes의 유도 범주,
Remark \ref{perfect-remark-discuss-rel-perfect}을 보라.

\begin{definition}
\label{more-algebra-definition-relatively-perfect}
$R \to A$를 평탄하고 유한 표시인 환 사상이라 하자. $D(A)$의 대상
$K$가 유사연접이고(정의 \ref{more-algebra-definition-pseudo-coherent}) $R$ 위에서
유한 Tor 차원을 가지면(정의 \ref{more-algebra-definition-tor-amplitude}),
$K$를 {\it $R$-퍼펙트} 또는 {\it $R$에 대해 상대적으로 퍼펙트}라고
한다.
\end{definition}

\noindent
보조정리 \ref{more-algebra-lemma-more-relative-pseudo-coherent-is-moot}에 의해
$K$가 $R$에 대해 상대적으로 유사연접이라고 요구해도 같은 개념을
얻는다. 이제 몇 가지 필수적인 보조정리를 제시한다.

\begin{lemma}
\label{more-algebra-lemma-cone-relatively-perfect}
$R \to A$를 평탄하고 유한 표시인 환 사상이라 하자. $D(A)$의
$R$-퍼펙트 대상들은 엄밀히 충만하고 포화된\footnote{Derived
Categories, Definition \ref{derived-definition-saturated}.} 삼각
부분범주를 이룬다.
\end{lemma}

\begin{proof}
이는 보조정리 \ref{more-algebra-lemma-cone-pseudo-coherent},
\ref{more-algebra-lemma-summands-pseudo-coherent},
\ref{more-algebra-lemma-cone-tor-amplitude}, 그리고
\ref{more-algebra-lemma-summands-tor-amplitude}에서 따른다.
\end{proof}

\begin{lemma}
\label{more-algebra-lemma-perfect-relatively-perfect}
$R \to A$를 평탄하고 유한 표시인 환 사상이라 하자. $D(A)$의 퍼펙트
대상은 $R$-퍼펙트이다. $K, M \in D(A)$라 할 때 $K$가 퍼펙트이고
$M$이 $R$-퍼펙트이면 $K \otimes_A^\mathbf{L} M$은
$R$-퍼펙트이다.
\end{lemma}

\begin{proof}
두 번째 명제에서 $M = A$로 두면 첫 번째 명제가 따른다. 두 번째
명제는 보조정리 \ref{more-algebra-lemma-perfect},
\ref{more-algebra-lemma-push-tor-amplitude}, 그리고
\ref{more-algebra-lemma-tensor-pseudo-coherent}에서 따른다.
\end{proof}

\begin{lemma}
\label{more-algebra-lemma-structure-relatively-perfect}
$R \to A$를 평탄하고 유한 표시인 환 사상이라 하고 $K \in D(A)$라
하자. 다음은 동치이다.
\begin{enumerate}
\item $K$는 $R$-퍼펙트이다.
\item $K$는 $R$-평탄하고 유한 표시인 $A$-가군들로 이루어진 유한
복합체와 동형이다.
\end{enumerate}
\end{lemma}

\begin{proof}
(2)가 (1)을 함의함을 보이기 위해서는 보조정리
\ref{more-algebra-lemma-cone-relatively-perfect}에 의해 $R$-평탄하고 유한 표시인
$A$-가군 $M$이 $D(A)$의 $R$-퍼펙트 대상을 정의함을 보이면 충분하다.
$M$은 $R$ 위에서 유한 Tor 차원을 가지므로 $M$이 유사연접임을 보이면
충분하다. Algebra, Lemma
\ref{algebra-lemma-flat-finite-presentation-limit-flat}에 의해
유한형 $\mathbf{Z}$-대수 $R_0 \subset R$, 평탄한 유한형 환 사상
$R_0 \to A_0$, 그리고 $R_0$ 위에서 평탄한 유한 $A_0$-가군 $M_0$가
존재하여 $A = A_0 \otimes_{R_0} R$이고
$M = M_0 \otimes_{R_0} R$이다. 보조정리
\ref{more-algebra-lemma-Noetherian-pseudo-coherent}에 의해
% P02-E0326: frozen English omits the article before “pseudo-coherent A_0-module”; Korean renders the predicate naturally.
$M_0$는 유사연접 $A_0$-가군이다. 유한 자유 $A_0$-가군 $P_0^n$들에
의한 분해 $P_0^\bullet \to M_0$를 택하자. $A_0$는 $R_0$ 위에서
평탄하므로 이는 평탄 분해이다. $M_0$도 $R_0$ 위에서 평탄하므로
$P^\bullet = P_0^\bullet \otimes_{R_0} R$은 여전히
$M = M_0 \otimes_{R_0} R$의 분해이다. (이를 보려면 보조정리
\ref{more-algebra-lemma-base-change-comparison}을 사용할 수 있다.) 따라서
$P^\bullet$은 $M$의 $A$ 위 유한 자유 분해이고, 이에 따라 $M$은
유사연접이다.

\medskip\noindent
(1)을 가정하자. $K$를 유한 자유 $A$-가군들의 위로 유계인 복합체
$P^\bullet$로 나타낼 수 있다. $K$를 $D(R)$의 대상으로 보았을 때
$[a, b]$에 Tor 진폭을 갖는다고 하자. 보조정리
\ref{more-algebra-lemma-last-one-flat}에 의해 $\tau_{\geq a}P^\bullet$은 $K$를
나타내는 $R$-평탄하고 유한 표시인 $A$-가군들의 복합체이다.
\end{proof}

\begin{lemma}
\label{more-algebra-lemma-base-change-relatively-perfect}
$R \to A$를 평탄하고 유한 표시인 환 사상이라 하자. $R \to R'$을
환 사상이라 하고 $A' = A \otimes_R R'$로 두자. $K \in D(A)$가
$R$-퍼펙트이면 $K \otimes_A^\mathbf{L} A'$은 $R'$-퍼펙트이다.
\end{lemma}

\begin{proof}
보조정리 \ref{more-algebra-lemma-pull-pseudo-coherent}에 의해
$K \otimes_A^\mathbf{L} A'$은 유사연접이다. 보조정리
\ref{more-algebra-lemma-base-change-comparison}에 의해
$K \otimes_A^\mathbf{L} A'$은 $D(R')$에서
$K \otimes_R^\mathbf{L} R'$과 같다. 이제 보조정리
\ref{more-algebra-lemma-pull-tor-amplitude}을 적용하면 $D(R')$의
$K \otimes_R^\mathbf{L} R'$이 유한 Tor 차원을 가짐을 알 수 있다.
\end{proof}

\begin{lemma}
\label{more-algebra-lemma-compute-RHom-relatively-perfect}
$R \to A$를 평탄한 환 사상이라 하자. $K, L \in D(A)$이고 $K$는
유사연접이며 $L$은 $R$ 위에서 유한 Tor 차원을 갖는다고 하자. 다음을
택할 수 있다.
\begin{enumerate}
\item $K$를 나타내는 유한 자유 $A$-가군들의 위로 유계인 복합체
$P^\bullet$,
\item $L$을 나타내는 $R$-평탄 $A$-가군들의 유계 복합체
$F^\bullet$.
\end{enumerate}
이렇게 택하면 다음이 성립한다.
\begin{enumerate}
\item[(a)] $E^\bullet = \Hom^\bullet(P^\bullet, F^\bullet)$은
$R\Hom_A(K, L)$을 나타내는 $R$-평탄 $A$-가군들의 아래로 유계인
복합체이다.
\item[(b)] 임의의 환 사상 $R \to R'$과
$A' = A \otimes_R R'$에 대하여 복합체
$E^\bullet \otimes_R R'$은
$R\Hom_{A'}(K \otimes_A^\mathbf{L} A',
L \otimes_A^\mathbf{L} A')$을 나타낸다.
\end{enumerate}
또한 $R \to A$가 유한 표시이고 $L$이 $R$-퍼펙트이면, 각 $F^p$를
유한 표시 $A$-가군으로 택할 수 있으며 따라서 각 $E^n$도 유한 표시
$A$-가군이다.
\end{lemma}

\begin{proof}
$P^\bullet$의 존재는 유사연접 복합체의 정의이다. 먼저 $L$을 자유
$A$-가군들의 위로 유계인 복합체 $F^\bullet$로 나타낸다(유계 Tor
차원은 특히 유계임을 함의하므로 가능하다). 다음으로 $L$을 $D(R)$의
대상으로 보았을 때 $[a, b]$에 Tor 진폭을 갖는다고 하자.
$F^\bullet$을 $\tau_{\geq a}F^\bullet$로 바꾸면 (2)와 같은 복합체를
얻는다. 이는 보조정리 \ref{more-algebra-lemma-last-one-flat}에서 따른다.

\medskip\noindent
(a)의 증명.
% P02-E0327: frozen source has “bounded an since”; Korean supplies the intended conjunction.
$F^\bullet$이 유계이고 $P^\bullet$이 위로 유계이므로 $n \ll 0$일 때
$E^n = 0$이고, $E^n$은 다음 유한한(!) 직합이다.
$$
E^n = \bigoplus\nolimits_{p + q = n} \Hom_A(P^{-q}, F^p)
$$
$P^{-q}$가 유한 자유이므로 이는 실제로 $R$-평탄 $A$-가군이다.
$E^\bullet$이 $R\Hom_A(K, L)$을 나타낸다는 사실은 보조정리
\ref{more-algebra-lemma-RHom-out-of-projective}에서 따른다.

\medskip\noindent
(b)의 증명.
$R \to R'$을 환 사상이라 하고 $A' = A \otimes_R R'$라 하자.
보조정리 \ref{more-algebra-lemma-base-change-comparison}에 의해 대상
$L \otimes_A^\mathbf{L} A'$은 $A'$-가군들의 복합체로 본
$F^\bullet \otimes_R R'$에 의해 나타내어진다(각 $F^p$가 $R$ 위에서
평탄하기 때문이다). $P^\bullet \otimes_R R'$에 대해서도 마찬가지다.
위와 같이
$R\Hom_{A'}(K \otimes_A^\mathbf{L} A',
L \otimes_A^\mathbf{L} A')$은 다음으로 나타내어진다.
$$
\Hom^\bullet(P^\bullet \otimes_R R', F^\bullet \otimes_R R') =
E^\bullet \otimes_R R'
$$
이 등식은 복합체의 각 항을 따로 살펴보고
$\Hom_{A'}(P^{-q} \otimes_R R', F^p \otimes_R R') =
\Hom_A(P^{-q}, F^p) \otimes_R R'$임을 사용하면 얻는다.
\end{proof}

\begin{lemma}
\label{more-algebra-lemma-colimit-relatively-perfect}
$R = \colim_{i \in I} R_i$를 환들의 여과 여극한이라 하자.
$0 \in I$라 하고 $R_0 \to A_0$를 평탄하고 유한 표시인 환 사상이라
하자. $i \geq 0$에 대하여 $A_i = R_i \otimes_{R_0} A_0$로 두고
$A = R \otimes_{R_0} A_0$로 두자.
\begin{enumerate}
\item $D(A)$의 $R$-퍼펙트 대상 $K$가 주어지면, 어떤 $i \in I$와
$D(A_i)$의 $R_i$-퍼펙트 대상 $K_i$가 존재하여 $D(A)$에서
$K \cong K_i \otimes_{A_i}^\mathbf{L} A$이다.
\item $K_0, L_0 \in D(A_0)$이고 $K_0$가 유사연접이며 $L_0$가
$R_0$ 위에서 유한 Tor 차원을 가지면
$$
\Hom_{D(A)}(K_0 \otimes_{A_0}^\mathbf{L} A, L_0 \otimes_{A_0}^\mathbf{L} A) =
\colim_{i \geq 0}
\Hom_{D(A_i)}(K_0 \otimes_{A_0}^\mathbf{L} A_i,
L_0 \otimes_{A_0}^\mathbf{L} A_i)
$$
이다.
\end{enumerate}
특히 $A$ 위의 $R$-퍼펙트 복합체들의 삼각 범주는 $A_i$ 위의
$R_i$-퍼펙트 복합체들의 삼각 범주들의 여극한이다.
\end{lemma}

\begin{proof}
Algebra, Lemma \ref{algebra-lemma-colimit-category-fp-modules}에 의해
유한 표시 $A$-가군들의 범주는 유한 표시 $A_i$-가군들의 범주들의
여극한이다. 이를 바탕으로 Algebra, Lemma
\ref{algebra-lemma-flat-finite-presentation-limit-flat}에 의해
% P02-E0328: frozen source omits “the” before “category”; Korean renders the definite category naturally.
$R$-평탄하고 유한 표시인 $A$-가군들의 범주는 $R_i$-평탄하고 유한
표시인 $A_i$-가군들의 범주들의 여극한임을 알 수 있다. 따라서
보조정리 \ref{more-algebra-lemma-structure-relatively-perfect}의 특성화는 (1)을
증명한다.

\medskip\noindent
(2)를 증명하기 위해 보조정리
\ref{more-algebra-lemma-compute-RHom-relatively-perfect}에서와 같이 $K_0$를
나타내는 $P_0^\bullet$과 $L_0$를 나타내는 $F_0^\bullet$을 택하자.
그러면 $E_0^\bullet = \Hom^\bullet(P_0^\bullet, F_0^\bullet)$은
$$
H^0(E_0^\bullet \otimes_{R_0} R_i) =
\Hom_{D(A_i)}(K_0 \otimes_{A_0}^\mathbf{L} A_i,
L_0 \otimes_{A_0}^\mathbf{L} A_i)
$$
와
$$
H^0(E_0^\bullet \otimes_{R_0} R) =
\Hom_{D(A)}(K_0 \otimes_{A_0}^\mathbf{L} A, L_0 \otimes_{A_0}^\mathbf{L} A)
$$
를 만족한다.
% P02-E0329: frozen source omits the finite verb in “Thus the result because”; Korean states the intended inference.
따라서 텐서곱은 여극한과 가환하고 여과 여극한은 완전하므로 결론이
따른다(Algebra, Lemma \ref{algebra-lemma-directed-colimit-exact}).
\end{proof}

\begin{lemma}
\label{more-algebra-lemma-thickening-relatively-perfect}
$R' \to A'$을 평탄하고 유한 표시인 환 사상이라 하자. $R' \to R$을
그 핵이 멱영 아이디얼인 전사 환 사상이라 하자.
$A = A' \otimes_{R'} R$로 두자. $K' \in D(A')$라 하고
$K = K' \otimes_{A'}^\mathbf{L} A$를 $D(A)$의 대상으로 두자.
$K$가 $R$-퍼펙트이면 $K'$은 $R'$-퍼펙트이다.
\end{lemma}

\begin{proof}
$A' \to A$의 핵은 멱영이고, $R' \to A'$의 평탄성에 의해
$K = K' \otimes_{R'}^\mathbf{L} R$이다(절
\ref{more-algebra-section-tor-independence}을 보라). 따라서 보조정리
\ref{more-algebra-lemma-check-pseudo-coherent-modulo-nilpotent}와
\ref{more-algebra-lemma-check-tor-dimension-modulo-nilpotent}를 결합하면 결론을
얻는다.
\end{proof}

\begin{lemma}
\label{more-algebra-lemma-lift-from-fibre-relatively-perfect}
$R$을 환이라 하자. $A = R[x_1, \ldots, x_d]/I$가 $R$ 위에서 평탄하고
유한 표시라고 하자. $\mathfrak q \subset A$를
$\mathfrak p \subset R$ 위에 놓인 소 아이디얼이라 하자.
$K \in D(A)$는 유사연접이고 $a, b \in \mathbf{Z}$라 하자.
$H^i(K_\mathfrak q \otimes_{R_\mathfrak p}^\mathbf{L} \kappa(\mathfrak p))$가
$i \in [a, b]$일 때만 영이 아니면, $K_\mathfrak q$는 $R$ 위에서
$[a - d, b]$에 Tor 진폭을 갖는다.
\end{lemma}

\begin{proof}
보조정리 \ref{more-algebra-lemma-more-relative-pseudo-coherent-is-moot}에 의해
$K$는 $R[x_1, \ldots, x_d]$-가군들의 복합체로서 유사연접이다.
따라서 $A = R[x_1, \ldots, x_d]$라고 가정할 수 있다.
$R_\mathfrak p \to A_\mathfrak q$와 복합체 $K_\mathfrak q$에
보조정리 \ref{more-algebra-lemma-perfect-over-regular-local-ring}을 적용하고
가정을 사용하면, $K_\mathfrak q$는 $D(A_\mathfrak q)$에서
퍼펙트이고 $[a - d, b]$에 Tor 진폭을 갖는다.
$R_\mathfrak p \to A_\mathfrak q$가 평탄하므로 보조정리
\ref{more-algebra-lemma-flat-push-tor-amplitude}에 의해 결론을 얻는다.
\end{proof}

\begin{lemma}
\label{more-algebra-lemma-bounded-on-fibres-relatively-perfect}
$R \to A$를 평탄하고 유한 표시인 환 사상이라 하고
$K \in D(A)$가 유사연접이라고 하자. 다음은 동치이다.
\begin{enumerate}
\item $K$는 $R$-퍼펙트이다.
\item $K$는 아래로 유계이고, 모든 소 아이디얼
$\mathfrak p \subset R$에 대하여 대상
$K \otimes_R^\mathbf{L} \kappa(\mathfrak p)$는 아래로 유계이다.
\end{enumerate}
\end{lemma}

\begin{proof}
정의에 의해 $R$-퍼펙트 복합체는 $R$-가군들의 복합체로서 유한 Tor
차원을 가지므로 (1)은 (2)를 함의한다. 역함의를 증명하자.

\medskip\noindent
$A = R[x_1, \ldots, x_d]/I$로 쓰자.
% P02-E0330: frozen source omits “by” in the notation sentence; Korean introduces the restriction grammatically.
$D(R[x_1, \ldots, x_d])$에서 $K$의 제한을 $L$로 나타내자.
보조정리 \ref{more-algebra-lemma-more-relative-pseudo-coherent-is-moot}에 의해
$L$은 유사연접이다. $L$과 $K$는 $D(R)$에서 같은 상을 가지므로
$L$이 $R$-퍼펙트일 필요충분조건은 $K$가 $R$-퍼펙트인 것이다. 또한
$L \otimes_R^\mathbf{L} \kappa(\mathfrak p)$와
$K \otimes_R^\mathbf{L} \kappa(\mathfrak p)$는
$D(\kappa(\mathfrak p))$에서 같은 대상이다. 따라서
$A = R[x_1, \ldots, x_d]$인 경우로 귀착된다.

\medskip\noindent
$A = R[x_1, \ldots, x_d]$이고 $K$가 (2)를 만족한다고 하자.
$\mathfrak q \subset A$를 소 아이디얼
$\mathfrak p \subset R$ 위에 놓인 소 아이디얼이라 하자.
$R_\mathfrak p \to A_\mathfrak q$와 복합체 $K_\mathfrak q$에
보조정리 \ref{more-algebra-lemma-perfect-over-regular-local-ring}을 적용하고
가정을 사용하면 $K_\mathfrak q$는 $D(A_\mathfrak q)$에서
퍼펙트이다. $K$가 아래로 유계이므로 보조정리
\ref{more-algebra-lemma-check-perfect-stalks}에 의해 $K$는 $D(A)$에서
퍼펙트이다. 보조정리 \ref{more-algebra-lemma-perfect-relatively-perfect}에 의해
이는 $K$가 $R$-퍼펙트임을 함의하고, 증명이 끝난다.
\end{proof}









\section{두 항 복합체}
\label{more-algebra-section-two-term}

\noindent
이 절에서는 가군들의 두 항 복합체에 관한 몇 가지 결과를 증명한다.
이 결과들은 순진한 여접복합체에 대한 조건을 이해하는 데 도움을 줄
것이다.

\begin{lemma}
\label{more-algebra-lemma-ext-1-zero}
$R$을 환이라 하자. $K \in D(R)$이고 $i \not \in \{-1, 0\}$일 때
$H^i(K) = 0$이라고 하자. 다음은 동치이다.
\begin{enumerate}
\item $H^{-1}(K) = 0$이고 $H^0(K)$는 사영 가군이다.
\item 모든 $R$-가군 $M$에 대하여 $\Ext^1_R(K, M) = 0$이다.
\end{enumerate}
$R$이 뇌터 환이고 $i = -1, 0$일 때 $H^i(K)$가 유한
$R$-가군이면, 이들은 다음과도 동치이다.
\begin{enumerate}
\item[(3)] 모든 유한 $R$-가군 $M$에 대하여
$\Ext^1_R(K, M) = 0$이다.
\end{enumerate}
\end{lemma}

\begin{proof}
(1)과 (2)의 동치는 보조정리 \ref{more-algebra-lemma-projective-amplitude}에서
따른다. $R$이 뇌터 환이고 $i = -1, 0$일 때 $H^i(K)$가 유한
$R$-가군이면, 보조정리 \ref{more-algebra-lemma-Noetherian-pseudo-coherent}에
의해 $K$는 유사연접이다. 따라서 (1)과 (3)의 동치는 보조정리
\ref{more-algebra-lemma-projective-amplitude-pseudo-coherent}에서 따른다.
\end{proof}

\begin{remark}
\label{more-algebra-remark-smoothness-ext-1-zero}
다음 두 명제는 보조정리 \ref{more-algebra-lemma-ext-1-zero}, Algebra, Definition
\ref{algebra-definition-smooth}, 그리고 Algebra, Proposition
\ref{algebra-proposition-characterize-formally-smooth}에서 따른다.
\begin{enumerate}
\item 환 사상 $A \to B$가 매끄러울 필요충분조건은 $A \to B$가 유한
표시이고 모든 $B$-가군 $N$에 대하여
$\Ext^1_B(\NL_{B/A}, N) = 0$인 것이다.
\item 환 사상 $A \to B$가 형식적으로 매끄러울 필요충분조건은 모든
$B$-가군 $N$에 대하여 $\Ext^1_B(\NL_{B/A}, N) = 0$인 것이다.
\end{enumerate}
\end{remark}

\begin{lemma}
\label{more-algebra-lemma-represent-two-term-complex}
$R$을 환이라 하자. $K$를 $D(R)$의 대상이라 하고
$i \not \in \{-1, 0\}$일 때 $H^i(K) = 0$이라고 하자. 그러면
\begin{enumerate}
\item $K$는 $K^0$가 자유 가군인 두 항 복합체
$K^{-1} \to K^0$로 나타낼 수 있다.
\item $R$이 뇌터 환이고 $i = -1, 0$일 때 $H^i(K)$가 유한
$R$-가군이면, $K$는 $K^0$가 유한 자유 가군이고 $K^{-1}$이 유한인
두 항 복합체 $K^{-1} \to K^0$로 나타낼 수 있다.
\end{enumerate}
\end{lemma}

\begin{proof}
(1)의 증명. $K$가 가군들의 복합체 $M^\bullet$로 주어졌다고 하자.
먼저 $M^\bullet$을 $\tau_{\leq 0}M^\bullet$로 바꿀 수 있다. 따라서
$i > 0$일 때 $M^i = 0$이라고 가정할 수 있다.
% P02-E0331: frozen source joins this completed sentence to “Next” with a comma; Korean separates them.
다음으로 $i > 0$일 때 $P^i = 0$인 자유 분해
$P^\bullet \to M^\bullet$을 택할 수 있다. Derived Categories, Lemma
\ref{derived-lemma-subcategory-left-resolution}를 보라. 마지막으로
$K^\bullet = \tau_{\geq -1}P^\bullet$로 둘 수 있다.

\medskip\noindent
(2)의 증명. $R$이 뇌터 환이고 $i = -1, 0$일 때 $H^i(K)$가 유한
$R$-가군이라고 하자. 보조정리 \ref{more-algebra-lemma-pseudo-coherent}에 의해
$i > 0$일 때 $F^i = 0$이고 각 $F^i$가 유한 자유인 준동형
$F^\bullet \to M^\bullet$을 택할 수 있다. 이제
$K^\bullet = \tau_{\geq -1}F^\bullet$로 둘 수 있다.
\end{proof}

\noindent
순진한 여접복합체 $\NL_{B/A}$ 또는 $\NL(\alpha)$에서 나가는 유도
범주의 사상들은(Algebra, Section
\ref{algebra-section-netherlander}을 보라) 다음 보조정리의 결과로
쉽게 이해할 수 있다.

\begin{lemma}
\label{more-algebra-lemma-map-out-of-almost-free}
$R$을 환이라 하자. $M^\bullet$을 $i > 0$일 때 $M^i = 0$이고
$M^0$가 사영 $R$-가군인 $R$ 위의 가군 복합체라 하자.
$K^\bullet$을 또 하나의 복합체라 하자.
\begin{enumerate}
\item $i \leq -2$일 때 $K^i = 0$이라고 가정하자. 그러면
$\Hom_{D(R)}(M^\bullet, K^\bullet) =
\Hom_{K(R)}(M^\bullet, K^\bullet)$이다.
\item $i \not \in [-1, 0]$일 때 $K^i = 0$이고 $K^0$가 사영
$R$-가군이라고 가정하자. 복합체의 사상
$a^\bullet : M^\bullet \to K^\bullet$에 대하여 다음은 동치이다.
\begin{enumerate}
\item $a^\bullet$은 모든 $R$-가군 $N$에 대하여 영사상
$\Ext^1_R(K^\bullet, N) \to \Ext^1_R(M^\bullet, N)$을 유도한다.
\item 어떤 사상 $h^0 : M^0 \to K^{-1}$가 존재하여
% P02-E1039: frozen formula composes h^0 with d_K^{-1}, although the types require d_M^{-1}; Korean preserves the formal span.
$a^{-1} + h^0 \circ d^{-1}_K = 0$이다.
\end{enumerate}
\item $i \leq -3$일 때 $K^i = 0$이라고 가정하자.
$\alpha \in \Hom_{D(R)}(M^\bullet, K^\bullet)$라 하자. $\alpha$와
$K^\bullet \to K^{-2}[2]$의 합성이 $a \circ d_M^{-3} = 0$인
$R$-가군 사상 $a : M^{-2} \to K^{-2}$에서 오면, $\alpha$는
$a^{-2} = a$인 복합체의 사상
$a^\bullet : M^\bullet \to K^\bullet$로 나타낼 수 있다.
\item
% P02-E0332: frozen item says “In (2)” although its alpha and degree-minus-two map occur only in (3); Korean preserves the item number.
(2)에서 $\alpha$를 나타내고 $a = (a')^{-2}$인 임의의 두 번째
복합체 사상 $(a')^\bullet : M^\bullet \to K^\bullet$에 대하여,
$i = 0, -1$일 때 사상 $h^i : M^i \to K^{i - 1}$들이 존재하여
$$
h^{-1} \circ d_M^{-2} = 0, \quad
(a')^{-1} = a^{-1} + d_K^{-2} \circ h^{-1} + h^0 \circ d_M^{-1},\quad
(a')^0 = a^0 + d_K^{-1} \circ h^0
$$
이다.
\end{enumerate}
\end{lemma}

\begin{proof}
$F^0 = M^0$로 두자. 자유 $R$-가군 $F^{-1}$과 전사
$F^{-1} \to M^{-1}$을 택하자. 자유 $R$-가군 $F^{-2}$과 전사
$F^{-2} \to M^{-2} \times_{M^{-1}} F^{-1}$을 택하자. 이 과정을
계속하면 각 항에서 전사이고 모든 $i$에 대하여 $F^i$가 사영인 준동형
$p^\bullet : F^\bullet \to M^\bullet$을 얻는다.

\medskip\noindent
(1)의 증명. Derived Categories, Lemma
\ref{derived-lemma-morphisms-from-projective-complex}에 의해
$$
\Hom_{D(R)}(M^\bullet, K^\bullet) = \Hom_{K(R)}(F^\bullet, K^\bullet)
$$
이다. $i \leq -2$일 때 $K^i = 0$이면 복합체의 모든 사상
$F^\bullet \to K^\bullet$은 $p^\bullet$을 통하여 인수분해된다.
마찬가지로 모든 호모토피 $\{h^i : F^i \to K^{i - 1}\}$도
$p^\bullet$을 통하여 인수분해된다. 따라서 (1)이 성립한다.

\medskip\noindent
(2)의 증명. (2)(b)가 성립하면 $a^\bullet$은 차수 $-1$에서 영인
복합체 사상 $(a')^\bullet : M^\bullet \to K^\bullet$과 호모토픽이다.
한편 $N \to I^\bullet$을 단사 분해라 하자. Derived Categories,
Lemma \ref{derived-lemma-morphisms-into-injective-complex}에 의해
$$
\Ext^1_R(K^\bullet, N) = \Hom_{D(R)}(K^\bullet, I^\bullet[1]) =
\Hom_{K(R)}(K^\bullet, I^\bullet[1])
$$
이다.
$b^\bullet : K^\bullet \to I^\bullet[1]$을 복합체의 사상이라 하자.
$K^1 = 0$이므로 사상 $b^0 : K^0 \to I^1$은 $I^1 \to I^2$의 핵,
즉 $I^0 \to I^1$의 상 안으로 사상된다. $K^0$가 사영이므로
$b^0$를 사상 $h : K^0 \to I^0$로 올릴 수 있다. 따라서
$b^\bullet$은 $(b')^0 = 0$인 복합체의 사상 $(b')^\bullet$과
호모토픽이다. $i \not \in [-1, 0]$일 때 $K^i = 0$이므로 복합체의
사상으로서 $(b')^\bullet \circ (a')^\bullet = 0$이다. 그러므로
사상 $\Ext^1_R(K^\bullet, N) \to \Ext^1_R(M^\bullet, N)$은 영이다.
이로써 (2)(b)가 (2)(a)를 함의함을 알 수 있다. 역으로 (2)(a)를
가정하자. $\Ext^1_R(K^\bullet, K^{-1})$의 표준 원소는
$\Ext^1_R(M^\bullet, K^{-1})$에서 영으로 사상된다. (1)을 사용하면
(2)(b)의 사상 $h^0$를 즉시 얻는다.

\medskip\noindent
(3)의 증명. $\alpha$를 나타내는
$b^\bullet : F^\bullet \to K^\bullet$을 택하자. $\alpha$와
$K^\bullet \to K^{-2}[2]$의 합성은
$b^{-2} : F^{-2} \to K^{-2}$로 나타내어진다. 이는
$a \circ p^{-2} : F^{-2} \to M^{-2} \to K^{-2}$과 호모토픽이므로,
어떤 사상 $h : F^{-1} \to K^{-2}$가 존재하여
$b^{-2} = a \circ p^{-2} + h \circ d_F^{-2}$이다.
$h$를 $F^\bullet$에서 $K^\bullet$로의 호모토피로 보아
$b^\bullet$을 조정하면 $b^{-2} = a \circ p^{-2}$을 얻는다.
따라서 $b^{-2}$은 $p^{-2}$을 통하여 인수분해된다.
$F^0 = M^0$이므로 $p^{-2}$의 핵은 $p^{-1}$의 핵 위로 전사한다
(예를 들어 $p^\bullet$의 핵이 비순환 복합체라는 사실이나 도표 추적을
사용할 수 있다). 그러므로 $b^{-1}$도 반드시 $p^{-1}$을 통하여
인수분해되고, 이 인수분해들과 $a^0 = b^0$에 대하여 (3)이 성립함을
알 수 있다.

\medskip\noindent
(4)의 증명은 생략한다. 힌트: $a^\bullet \circ p^\bullet$과
$(a')^\bullet \circ p^\bullet$ 사이에 호모토피가 존재하고, 앞에서와
같이 이 호모토피가 $p^\bullet$을 통하여 인수분해됨을 보인다.
\end{proof}
\noindent
$A \to B$를 유한 표시인 환 사상이라 하자. 아이디얼 $I \subset B$가
주어지면 다음 조건을 생각할 수 있다.
\begin{enumerate}
\item[(*)] 모든 $B$-가군 $N$에 대하여
$\Ext^1_B(\NL_{B/A}, N)$은 $I$에 의해 소멸된다.
\end{enumerate}
% P02-E1041: frozen English leaves the scheme-theoretic compound open; Korean renders the intended containment.
이 조건은 “$A \to B$의 특이 자취가 스킴 이론적으로 $V(I)$에
포함된다”는 개념을 정밀하게 수학적으로 정식화하는 한 가지 방법이다.
Remark \ref{more-algebra-remark-smoothness-ext-1-zero}와 다음 보조정리들을 비교하라.

\begin{lemma}
\label{more-algebra-lemma-ext-1-annihilated-definite}
$R$을 환이라 하고 $I \subset R$을 아이디얼이라 하자.
$K \in D(R)$라 하고 $i \not \in \{-1, 0\}$일 때
$H^i(K) = 0$이라고 가정하자. 다음은 동치이다.
\begin{enumerate}
\item 모든 $R$-가군 $N$에 대하여 $\Ext^1_R(K, N)$은 $I$에 의해
소멸된다.
\item $K$는 $K^0$가 자유이고, 임의의 $a \in I$에 대하여 사상
$a : K^{-1} \to K^{-1}$이 $d_K^{-1} : K^{-1} \to K^0$을 통하여
인수분해되는 복합체 $K^{-1} \to K^0$로 나타낼 수 있다.
\item $K$가 $K^0$가 사영인 두 항 복합체
$K^{-1} \to K^0$로 나타내어질 때마다, 임의의 $a \in I$에 대하여
사상 $a : K^{-1} \to K^{-1}$은
$d_K^{-1} : K^{-1} \to K^0$을 통하여 인수분해된다.
\end{enumerate}
$R$이 뇌터 환이고 $i = -1, 0$일 때 $H^i(K)$가 유한 $R$-가군이면
이들은 다음과도 동치이다.
\begin{enumerate}
\item[(4)] 모든 유한 $R$-가군 $N$에 대하여 $\Ext^1_R(K, N)$은
$I$에 의해 소멸된다.
\item[(5)] $K$는 $K^0$가 유한 자유이고 $K^{-1}$이 유한이며, 임의의
$a \in I$에 대하여 사상 $a : K^{-1} \to K^{-1}$이
$d_K^{-1} : K^{-1} \to K^0$을 통하여 인수분해되는 복합체
$K^{-1} \to K^0$로 나타낼 수 있다.
\end{enumerate}
\end{lemma}
\begin{proof}

\medskip\noindent
(1)을 가정하고, $K$를 나타내며 $K^0$가 사영인 두 항 복합체
$K^{-1} \to K^0$를 택하자. 더 언급하지 않고 보조정리
\ref{more-algebra-lemma-map-out-of-almost-free}에 주어진 $K^\bullet$에서 나가는
$D(R)$의 사상에 대한 기술을 사용할 것이다.
$N = K^{-1}$로 택하고
$\text{id}_{K^{-1}} : K^{-1} \to K^{-1}$에 의해 주어지는
$\Ext^1_R(K, N)$의 원소 $\xi$를 생각하자.
% P02-E0333: frozen source omits the subject of “Since is annihilated”; Korean identifies xi.
$a \in I$가 이를 소멸시키므로 다음 가환 도표에 맞는 점선 화살표를
얻는다.
$$
\xymatrix{
K^{-1} \ar[d]_a \ar[r]_{d_K^{-1}} & K^0 \ar@{..>}[ld]^h \\
K^{-1}
}
$$
이는 (3)을 증명한다. 보조정리
\ref{more-algebra-lemma-represent-two-term-complex}의 (1)에 의해 (3)은 (2)를
함의한다. $K^\bullet$이 (2)에서와 같고 $N$이 임의의
$R$-가군이라고 하자. $\Ext^1_R(K, N)$의 임의의 원소 $\xi$는 사상
$\varphi : K^{-1} \to N$의 동치류로 주어진다. 그러면 $a \in I$에
대하여 가정에 의해 위 도표와 같은 사상 $h$를 택할 수 있고,
$a\varphi = \varphi \circ a =
\varphi \circ h \circ \text{d}_K^{-1}$임을 알 수 있다. 이는
$a \xi$가 $\Ext^1_R(K, N)$에서 영임을 증명한다. 따라서 (1), (2),
(3)은 동치이다.

\medskip\noindent
$R$이 뇌터 환이고 $i = -1, 0$일 때 $H^i(K)$가 유한 $R$-가군이라고
가정하자. 보조정리 \ref{more-algebra-lemma-represent-two-term-complex}의 (2)에
의해 (3)은 (5)를 함의한다. (5)가 (2)를 함의함은 명백하다. 자명하게
(1)은 (4)를 함의한다. 따라서 증명을 끝내려면 (4)가 다른 조건들 중
하나를 함의함을 보이면 충분하다. 보조정리
\ref{more-algebra-lemma-represent-two-term-complex}의 (2)에서와 같이 $K$를
나타내며 $K^0$가 유한 자유이고 $K^{-1}$이 유한인 복합체
$K^{-1} \to K^0$를 택하자. (2) $\Rightarrow$ (1)의 증명에 주어진
논증은 $\Ext^1_R(K, K^{-1})$이 $I$에 의해 소멸되면 (1)이
성립함을 보인다. 따라서 (4)는 (1)을 함의하고, 증명이 끝난다.
\end{proof}

\begin{lemma}
\label{more-algebra-lemma-two-term-base-change}
$R$을 환이라 하자. $K$를 $D(R)$의 대상이라 하고
$i \not \in \{-1, 0\}$일 때 $H^i(K) = 0$이라고 하자. $K$를 나타내는
$R$-가군들의 두 항 복합체 $K^{-1} \to K^0$에서 $K^0$가 평탄한
$R$-가군(예를 들어 사영 또는 자유)이라고 하자. $R \to R'$을 환
사상이라 하자. 그러면 복합체 $K^\bullet \otimes_R R'$은
$\tau_{\geq -1}(K \otimes_R^\mathbf{L} R')$을 나타낸다.
\end{lemma}

\begin{proof}
$D(R)$에서 다음 구별 삼각을 갖는다.
$$
K^0 \to K^\bullet \to K^{-1}[1] \to K^0[1]
$$
이는 다음 구별 삼각들의 사상을 결정한다.
$$
\xymatrix{
K^0 \otimes_R^\mathbf{L} R' \ar[d] \ar[r] &
K^\bullet \otimes_R^\mathbf{L} R' \ar[r] \ar[d] &
K^{-1} \otimes_R^\mathbf{L} R'[1] \ar[r] \ar[d] &
K^0 \otimes_R^\mathbf{L} R'[1] \ar[d] \\
K^0 \otimes_R R' \ar[r] &
K^\bullet \otimes_R R' \ar[r] &
K^{-1} \otimes_R R'[1] \ar[r] &
K^0 \otimes_R R'[1]
}
$$
$K^0$가 평탄하므로 왼쪽과 오른쪽 수직 화살표는 동형이다.
$K^{-1} \otimes_R^\mathbf{L} R' \to K^{-1} \otimes_R R'$은 차수
$0$의 코호몰로지에서 동형이므로 결론을 얻는다.
\end{proof}
\begin{lemma}
\label{more-algebra-lemma-base-change-property-ext-1-annihilated}
$I$를 환 $R$의 아이디얼이라 하자. $K$를 $D(R)$의 대상이라 하고
$i \not \in \{-1, 0\}$일 때 $H^i(K) = 0$이라고 하자.
$R \to R'$을 환 사상이라 하자. $K$가 $(R, I)$에 관하여 보조정리
\ref{more-algebra-lemma-ext-1-annihilated-definite}의 동치 조건 (1), (2), (3)을
만족하면, $\tau_{\geq -1}(K \otimes_R^\mathbf{L} R')$은
$(R', IR')$에 관하여 보조정리
\ref{more-algebra-lemma-ext-1-annihilated-definite}의 동치 조건 (1), (2), (3)을
만족한다.
% P02-E1043: frozen source has no terminal punctuation before the lemma ends; Korean closes the statement normally.
\end{lemma}

\begin{proof}
$K$가 $K^0$가 자유이고, 모든 $a \in I$에 대하여 어떤 사상
$h_a : K^0 \to K^{-1}$가 존재하여 사상
$a : K^{-1} \to K^{-1}$이 $h_a \circ d_K^{-1}$와 같은 두 항 복합체
$K^{-1} \to K^0$로 나타내어진다고 가정할 수 있다. 보조정리
\ref{more-algebra-lemma-two-term-base-change}에 의해
$\tau_{\geq -1}(K \otimes_R^\mathbf{L} R')$은
$K^\bullet \otimes_R R'$로 나타내어진다. 그러면 모든
$a \in I$에 대하여
$a \otimes 1 : K^{-1} \otimes_R R' \to K^{-1} \otimes_R R'$은
$(h_a \otimes 1) \circ (d_K^{-1} \otimes 1)$와 같다.
$\text{d}_K^{-1} \otimes 1$을 통하여 인수분해되는 사상
$K^{-1} \otimes_R R' \to K^{-1} \otimes_R R'$들의 모임은
$R'$-가군을 이루므로 결론을 얻는다.
\end{proof}

\begin{lemma}
\label{more-algebra-lemma-two-term-surjection-map-zero}
$R$을 환이라 하고 $\alpha : K \to K'$을 $D(R)$의 사상이라 하자.
다음을 가정하자.
\begin{enumerate}
\item $i \not \in \{-1, 0\}$일 때
$H^i(K) = H^i(K') = 0$이다.
\item $H^0(\alpha)$는 동형이고 $H^{-1}(\alpha)$는 전사이다.
\end{enumerate}
임의의 $f \in R$에 대하여 $f : K \to K$가 $0$이면
$f : K' \to K'$도 $0$이다.
\end{lemma}

\begin{proof}
$M = \Ker(H^{-1}(\alpha))$로 두자. 그러면 $\alpha$는 구별 삼각
$$
M[1] \to K \to K' \to M[2]
$$
에 들어간다. 가정에 의해 $K \to K' \xrightarrow{f} K'$가
영이므로 $f : K' \to K'$는 사상 $M[2] \to K'$을 통하여
인수분해된다.
% P02-E1044: frozen source says “factors over”; Korean uses the intended through-factorization without changing any formal span.
그러나 예를 들어 Derived Categories, Lemma
\ref{derived-lemma-negative-exts}에 의해
$\Hom(M[2], K') = 0$이다.
\end{proof}

\begin{lemma}
\label{more-algebra-lemma-surjection-property-ext-1-annihilated}
$I$를 환 $R$의 아이디얼이라 하자. $\alpha : K \to K'$을
$D(R)$의 사상이라 하고 다음을 가정하자.
\begin{enumerate}
\item $i \not \in \{-1, 0\}$일 때
$H^i(K) = H^i(K') = 0$이다.
\item $H^0(\alpha)$는 동형이고 $H^{-1}(\alpha)$는 전사이다.
\end{enumerate}
$K$가 보조정리 \ref{more-algebra-lemma-ext-1-annihilated-definite}의 동치 조건
(1), (2), (3)을 만족하면 $K'$도 그러하다.
\end{lemma}

\begin{proof}
$M = \Ker(H^{-1}(\alpha))$로 두자. 그러면 $\alpha$는 구별 삼각
$$
M[1] \to K \to K' \to M[2]
$$
에 들어간다. 임의의 $R$-가군 $N$에 대하여 이는 완전열
$$
\Ext^0_R(M[1], N) \to
\Ext^1_R(K', N) \to
\Ext^1_R(K, N)
$$
을 결정한다. $\Ext^0_R(M[1], N) = \Ext^{-1}_R(M, N) = 0$이므로
$\Ext^1_R(K', N)$은 $\Ext^1_R(K, N)$의 부분가군이다. 따라서
$\Ext^1_R(K, N)$이 $I$에 의해 소멸되면
$\Ext^1_R(K', N)$도 그러하다.
\end{proof}
\begin{lemma}
\label{more-algebra-lemma-ext-1-annihilated}
% P02-E0334: frozen source says “Let R be ring”; Korean renders the declaration normally.
$R$을 환이라 하고 $I \subset R$을 아이디얼이라 하자.
$K \in D(R)$이고 $i \not \in \{-1, 0\}$일 때
$H^i(K) = 0$이라고 하자. 다음은 동치이다.
\begin{enumerate}
\item 어떤 $c \geq 0$가 존재하여 $K$와 아이디얼 $I^c$에 대해
보조정리 \ref{more-algebra-lemma-ext-1-annihilated-definite}의 동치 조건
(1), (2), (3)이 성립한다.
\item 어떤 $c \geq 0$가 존재하여 (a) $I^c$가 $H^{-1}(K)$를
소멸시키고 (b) $H^0(K)$가 $I^c$-사영 가군이다(절
\ref{more-algebra-section-near-projective}를 보라).
\end{enumerate}
$R$이 뇌터 환이고 $i = -1, 0$일 때 $H^i(K)$가 유한
$R$-가군이면, 이들은 다음과도 동치이다.
\begin{enumerate}
\item[(3)] 어떤 $c \geq 0$가 존재하여 $K$와 아이디얼 $I^c$에 대해
보조정리 \ref{more-algebra-lemma-ext-1-annihilated-definite}의 동치 조건
(4), (5)가 성립한다.
\item[(4)] $H^{-1}(K)$는 $I$-거듭제곱 꼬임이고, 어떤
$f_1, \ldots, f_s \in R$가 존재하여
$V(f_1, \ldots, f_s) \subset V(I)$이고 국소화
$H^0(K)_{f_i}$들이 사영 $R_{f_i}$-가군이다.
\item[(5)] $H^{-1}(K)$는 $I$-거듭제곱 꼬임이고, 어떤
$f_1, \ldots, f_s \in I$가 존재하여
$V(f_1, \ldots, f_s) = V(I)$이고 국소화
$H^0(K)_{f_i}$들이 사영 $R_{f_i}$-가군이다.
\item[(6)] $H^{-1}(K)$는 $I$-거듭제곱 꼬임이고,
$V(f_1, \ldots, f_s) = V(I)$인 임의의
$f_1, \ldots, f_s \in I$에 대하여 국소화 $H^0(K)_{f_i}$들은
사영 $R_{f_i}$-가군이다.
\end{enumerate}
\end{lemma}

\begin{proof}
구별 삼각
$H^{-1}(K)[1] \to K \to H^0(K)[0] \to H^{-1}(K)[2]$는 완전열
$$
0 \to \Ext^1_R(H^0(K), N) \to \Ext^1_R(K, N) \to \Hom_R(H^{-1}(K), N) \to
\Ext^2_R(H^0(K), N)
$$
을 결정한다. 따라서 (2)는 모든 $R$-가군 $N$에 대하여 $I^{2c}$가
$\Ext^1_R(K, N)$을 소멸시킴을 함의한다. (1)을 가정하면
$H^0(K)$가 $I^c$-사영임을 즉시 알 수 있다. 한편 어떤 단사
$R$-가군 $N$에 대한 단사 사상 $H^{-1}(K) \to N$을 택할 수 있다.
그러면 완전열에서 $\Ext^2$가 영이므로 이 사상은
$\Ext^1_R(K, N)$의 한 원소의 상이고, 따라서 $H^{-1}(K)$가
$I^c$에 의해 소멸됨을 얻는다.

\medskip\noindent
$R$이 뇌터 환이고 $i = -1, 0$일 때 $H^i(K)$가 유한 $R$-가군이라고
가정하자. 보조정리 \ref{more-algebra-lemma-ext-1-annihilated-definite}에 의해
(3)은 (1), (2)와 동치이다. 또한 (3)이 성립하면
% P02-E0335: frozen fixed-power argument claims factorization by f rather than the guaranteed f^c; Korean preserves the missing exponent.
$f \in I$에 대하여 $H^0(K)$ 위의 $f$에 의한 곱셈은 사영 가군을
통하여 인수분해된다. 이는 $H^0(K)_f$가 사영 $R_f$-가군의 직합
인자이고 따라서 그 자체가 사영 $R_f$-가군임을 함의한다. 그러므로
동치 조건 (1), (2), (3)은 (6)을 함의한다. 물론 (6)은 (5)를
함의하고 (5)는 자명하게 (4)를 함의한다.

\medskip\noindent
(4)를 가정하자. $H^{-1}(K)$는 유한 $R$-가군이자
$I$-거듭제곱 꼬임이므로, 어떤 $c_1 \geq 0$에 대하여 $I^{c_1}$이
$H^{-1}(K)$를 소멸시킨다. 짧은 완전열
$$
0 \to M \to R^{\oplus r} \to H^0(K) \to 0
$$
을 택하자. 이는 원소 $\xi \in \Ext^1_R(H^0(K), M)$를 결정한다.
임의의 $f \in I$에 대하여 보조정리
\ref{more-algebra-lemma-pseudo-coherence-and-base-change-ext}에 의해
$\Ext^1_R(H^0(K), M)_f = \Ext^1_{R_f}(H^0(K)_f, M_f)$이다.
따라서 $H^0(K)_f$가 사영이면 $f$의 어떤 거듭제곱이 $\xi$를
소멸시킨다. 그러므로 어떤 $c_2 \geq 0$에 대하여
$(f_1, \ldots, f_s)^{c_2}$가 $\xi$를 소멸시킨다.
% P02-E0336: frozen formula omits the radical on the right; Korean preserves the exact inclusion.
$V(f_1, \ldots, f_s) \subset V(I)$이므로
$\sqrt{I} \subset (f_1, \ldots, f_s)$이다(Algebra, Lemma
\ref{algebra-lemma-Zariski-topology}). $R$이 뇌터 환이므로 어떤
$c_3 \geq 0$에 대하여
$I^{c_3} \subset (f_1, \ldots, f_s)$이다(Algebra, Lemma
\ref{algebra-lemma-Noetherian-power}).
% P02-E0337: frozen exponents lose both subscripts; Korean preserves I^{c2c3} and max(c_1,c_2c3).
따라서 $I^{c2c3}$가 $\xi$를 소멸시킨다.
% P02-E0338: frozen parenthesis uses a in I rather than the proved power ideal and has agreement errors; Korean preserves the formal locus.
이는 다시 $H^0(K)$가 $I^{c_2c_3}$-사영임을 뜻한다
($\xi$를 소멸시키는 $a \in I$에 의한 곱셈이
$R^{\oplus r}$을 통하여 인수분해되기 때문이다). 그러므로
$c = \max(c_1, c_2c_3)$로 택하면 (2)가 성립한다.
\end{proof}
\begin{lemma}
\label{more-algebra-lemma-zero-in-derived}
$R$을 환이라 하자. $j = 1, 2, 3$에 대하여 $K_j \in D(R)$이고
$i \not \in \{-1, 0\}$일 때 $H^i(K_j) = 0$이라고 하자.
$\varphi : K_1 \to K_2$와 $\psi : K_2 \to K_3$를 $D(R)$의
사상들이라 하자. $H^0(\varphi) = 0$이고 $H^{-1}(\psi) = 0$이면
% P02-E0339: frozen source reverses the composable order; Korean preserves the ill-typed formal span.
$\varphi \circ \psi = 0$이다.
\end{lemma}

\begin{proof}
Derived Categories, Lemma
\ref{derived-lemma-trick-vanishing-composition}을 적용하면
$\varphi \circ \psi$가 $\tau_{\leq -2}K_2 = 0$을 통하여
인수분해됨을 알 수 있다.
\end{proof}

\begin{lemma}
\label{more-algebra-lemma-silly}
$R$을 환이라 하자. $K \in D(R)$가
$R^{\oplus n} \to R^{\oplus n}$ 꼴의 두 항 복합체로 주어졌다고 하자.
% P02-E0340: frozen notation sentence omits “by”; Korean introduces the differential matrix grammatically.
미분의 행렬을 $A \in \text{Mat}(n \times n, R)$로 나타내자.
% P02-E0341: frozen conclusion uses undefined lowercase a although the matrix is uppercase A; Korean preserves det(a).
그러면 $D(R)$에서 $\det(a) : K \to K$는 영이다.
\end{lemma}

\begin{proof}
생략한다. 좋은 연습문제이다.
\end{proof}



\section{순진한 여접복합체}
\label{more-algebra-section-netherlander}

\noindent
이 절에서는 Algebra, Section \ref{algebra-section-netherlander}에서
시작한 논의를 계속한다. 먼저 밑변환을 논의한다. 첫 번째 보조정리는
순진한 여접복합체와 환 확대의 순진한 텐서곱을 취하는 일이 생각만큼
그렇게 순진한 것은 아님을 보여 준다.

\begin{lemma}
\label{more-algebra-lemma-tensor-NL}
$R \to S$와 $S \to S'$을 환 사상이라 하자. 표준 사상
$\NL_{S/R} \otimes_S^\mathbf{L} S' \to \NL_{S/R} \otimes_S S'$은
$D(S')$에서 동형
$\tau_{\geq -1}(\NL_{S/R} \otimes_S^\mathbf{L} S') \to
\NL_{S/R} \otimes_S S'$을 유도한다. 마찬가지로 $R$ 위의 $S$의
표시 $\alpha$가 주어지면, 표준 사상
$\NL(\alpha) \otimes_S^\mathbf{L} S' \to
\NL(\alpha) \otimes_S S'$은 $D(S')$에서 동형
$\tau_{\geq -1}(\NL(\alpha) \otimes_S^\mathbf{L} S') \to
\NL(\alpha) \otimes_S S'$을 유도한다.
\end{lemma}

\begin{proof}
% P02-E0342: frozen proof is a sentence fragment; Korean supplies a complete proof sentence.
이는 보조정리 \ref{more-algebra-lemma-two-term-base-change}의 특수한 경우이다.
\end{proof}

\begin{lemma}
\label{more-algebra-lemma-base-change-NL}
$R \to S$와 $R \to R'$을 환 사상이라 하자.
$\alpha : P \to S$를 $R$ 위의 $S$의 표시라 하자. 그러면
$\alpha' : P \otimes_R R' \to S \otimes_R R'$은 $R'$ 위의
$S' = S \otimes_R R'$의 표시이다. 표준 사상
% P02-E0343: frozen source omits the backslash on the first NL macro; Korean preserves plain NL(\alpha).
$$
NL(\alpha) \otimes_S S' \to \NL(\alpha')
$$
은 $H^0$에서 동형이고 $H^{-1}$에서 전사이다. 특히 표준 사상
$$
\NL_{S/R} \otimes_S S' \to \NL_{S'/R'}
$$
은 $H^0$에서 동형이고 $H^{-1}$에서 전사이다.
\end{lemma}

\begin{proof}
% P02-E0344: frozen notation clauses omit “by”; Korean introduces all three symbols grammatically.
$I = \Ker(P \to S)$로 나타내자. $P' = P \otimes_R R'$로 나타내고
$I' = \Ker(P' \to S')$로 나타내자. $P$가 $j \in J$에 대한
$x_j$들의 다항식 대수라고 하자. 보조정리에 표시된 사상은 다음이
된다.
$$
\xymatrix{
\bigoplus_{j \in J} S' \text{d}x_j \ar[r] &
\bigoplus_{j \in J} S' \text{d}x_j \\
I/I^2 \otimes_S S' \ar[r] \ar[u] &
I'/(I')^2 \ar[u]
}
$$
여기서 왼쪽 열은 $\NL(\alpha) \otimes_S S'$이고 오른쪽 열은
$\NL(\alpha')$이다. 텐서곱의 오른쪽 완전성에 의해
$I \otimes_R R' \to I'$은 전사이다. 따라서 아래쪽 화살표도
전사이다. 이는 보조정리의 첫 번째 명제를 증명한다.
$\NL_{S/R} \otimes_S S' \to \NL_{S'/R'}$에 대한 명제는 이 복합체들이
각각 $\NL(\alpha) \otimes_S S'$ 및 $\NL(\alpha')$와 호모토픽이라는
사실에서 따른다.
\end{proof}

\begin{lemma}
\label{more-algebra-lemma-base-change-NL-flat}
환들의 쌍대 데카르트 그림
$$
\xymatrix{
B \ar[r] & B' \\
A \ar[r] \ar[u] & A' \ar[u]
}
$$
을 생각하자. $B$가 $A$ 위에서 평탄하면 표준 사상
$\NL_{B/A} \otimes_B B' \to \NL_{B'/A'}$은 준동형이다. 또한
$\NL_{B/A}$가 $[-1, 0]$에 Tor 진폭을 가지면
$\NL_{B/A} \otimes_B^\mathbf{L} B' \to \NL_{B'/A'}$도
준동형이다.
\end{lemma}

\begin{proof}
Algebra, Section \ref{algebra-section-netherlander}에서와 같은 표시
$\alpha : P \to B$를 택하자. $I = \Ker(\alpha)$라 하자.
$P' = P \otimes_A A'$로 두고,
% P02-E0345: frozen denote-by construction omits “by”; Korean introduces alpha-prime grammatically.
$\alpha' : P' \to B'$로 $A'$ 위의 $B'$의 대응하는 표시를 나타내자.
$B$가 $A$ 위에서 평탄하므로 $I' = \Ker(\alpha')$은
$I \otimes_A A'$와 같다. 따라서
$$
I'/(I')^2 = \Coker(I^2 \otimes_A A' \to I \otimes_A A') =
I/I^2 \otimes_A A' = I/I^2 \otimes_B B'
$$
이다. 양변이 같은 기저를 가지므로
$\Omega_{P'/A'} = \Omega_{P/A} \otimes_A A'$이다. 따라서
$\Omega_{P'/A'} \otimes_{P'} B' =
\Omega_{P/A} \otimes_P B \otimes_B B'$이다. 이는
$\NL(\alpha) \otimes_B B' \to \NL(\alpha')$이 복합체의
동형임을 증명하므로 첫 번째 명제가 성립한다.

\medskip\noindent
$I/I^2$가 차수 $-1$에 놓인 $B$-가군들의 복합체로서
$$
\NL(\alpha) = I/I^2 \longrightarrow \Omega_{P/A} \otimes_P B
$$
를 갖는다. 차수 $0$의 항이 자유이므로 이 복합체가 $[-1, 0]$에
Tor 진폭을 가질 필요충분조건은 $I/I^2$가 평탄 $B$-가군인 것이다.
보조정리 \ref{more-algebra-lemma-last-one-flat}을 보라. 이 조건이 성립하면
$\NL(\alpha) \otimes_B^\mathbf{L} B' =
\NL(\alpha) \otimes_B B'$이고 두 번째 명제를 얻는다.
\end{proof}

\begin{lemma}
\label{more-algebra-lemma-lci-NL}
$A \to B$를 정의 \ref{more-algebra-definition-local-complete-intersection}에서와
같은 국소 완전교차라 하자. 그러면 $\NL_{B/A}$는 $[-1, 0]$에
Tor 진폭을 갖는 $D(B)$의 퍼펙트 대상이다.
\end{lemma}

\begin{proof}
$B = A[x_1, \ldots, x_n]/I$로 쓰자. 그러면 $\NL_{B/A}$는
$I/I^2$가 차수 $-1$에 놓인 $B$-가군들의 복합체
$$
I/I^2 \longrightarrow \bigoplus B \text{d}x_i
$$
로 나타내어진다. 차수 $0$의 항이 유한 자유이므로 이 복합체가
$[-1, 0]$에 Tor 진폭을 가질 필요충분조건은 $I/I^2$가 평탄
$B$-가군인 것이다. 보조정리 \ref{more-algebra-lemma-last-one-flat}을 보라.
정의에 의해 $I$는 Koszul 정칙 아이디얼이고 따라서 준정칙
아이디얼이다. 절 \ref{more-algebra-section-ideals}를 보라. 그러므로
$I/I^2$는 유한 사영 $B$-가군이고(보조정리
\ref{more-algebra-lemma-quasi-regular-ideal-finite-projective}), 이에 따라
$\NL_{B/A}$가 퍼펙트이고 $[-1, 0]$에 Tor 진폭을 가짐을 모두 얻는다.
\end{proof}

\begin{lemma}
\label{more-algebra-lemma-base-change-lci-bis}
환들의 쌍대 데카르트 그림
$$
\xymatrix{
B \ar[r] & B' \\
A \ar[r] \ar[u] & A' \ar[u]
}
$$
을 생각하자. $A \to B$와 $A' \to B'$이 정의
\ref{more-algebra-definition-local-complete-intersection}에서와 같은 국소
완전교차이면, 사상
$H^{-1}(\NL_{B/A} \otimes_B B') \to H^{-1}(\NL_{B'/A'})$의 핵은
유한 사영 $B'$-가군이다.
\end{lemma}

\begin{proof}
보조정리 \ref{more-algebra-lemma-lci-NL}에 의해 복합체 $\NL_{B/A}$와
$\NL_{B'/A'}$은 $[-1, 0]$에 Tor 진폭을 갖는 퍼펙트 복합체이다.
보조정리 \ref{more-algebra-lemma-tensor-NL}, \ref{more-algebra-lemma-pull-perfect}, 그리고
\ref{more-algebra-lemma-pull-tor-amplitude}을 결합하면
$\NL_{B/A} \otimes_B B' = \NL_{B/A} \otimes_B^\mathbf{L} B'$이고,
이 복합체 역시 $[-1, 0]$에 Tor 진폭을 갖는 퍼펙트 복합체이다.
$D(B')$에서 구별 삼각
$$
C \to \NL_{B/A} \otimes_B B'  \to \NL_{B'/A'} \to C[1]
$$
을 택하자. 보조정리 \ref{more-algebra-lemma-two-out-of-three-perfect}와
\ref{more-algebra-lemma-cone-tor-amplitude}에 의해 $C$는 $[-1, 1]$에
Tor 진폭을 갖는 퍼펙트 복합체이다. 보조정리
\ref{more-algebra-lemma-base-change-NL}에 의해 복합체 $C$의 영이 아닌
코호몰로지 가군은 하나뿐이며, 이는 보조정리의 가군이 차수 $-1$에
놓인 것이다. 이 가군은 유한 표시이고(보조정리
\ref{more-algebra-lemma-n-pseudo-module}) 평탄하다(보조정리
\ref{more-algebra-lemma-tor-dimension}). 따라서 Algebra, Lemma
\ref{algebra-lemma-finite-projective}에 의해 유한 사영이다.
\end{proof}











\section{아벨 군의 Rlim}
\label{more-algebra-section-Rlim}

\noindent
아벨 군에서의 $R\lim$을 간단히 논의한다. 이 절에서
$\textit{Ab}(\mathbf{N})$으로 아벨 군의 역계들로 이루어진 아벨
범주를 나타내겠다.
% P02-E0346: frozen source omits “by” in the Ab(N) notation sentence; Korean introduces it grammatically.
이 표기는 위의 아벨 군들의 층에 대한 표기와 양립한다. 실제로 아벨
군의 역계는 범주 $\mathbf{N}$ 위의 군의 층과 같은 것이다
($i \leq j$이면 유일한 사상 $i \to j$가 있다). Remark
\ref{more-algebra-remark-Rlim-cohomology}를 보라. 이 절의 많은 논증은 위의 아벨
군들의 층에 대한 코호몰로지 기계를 구성할 때 사용한 논증을
반복한다.

\begin{lemma}
\label{more-algebra-lemma-compute-Rlim}
함자 $\lim : \textit{Ab}(\mathbf{N}) \to \textit{Ab}$은 오른쪽
유도 함자
\begin{equation}
\label{more-algebra-equation-Rlim}
R\lim : D(\textit{Ab}(\mathbf{N})) \longrightarrow D(\textit{Ab})
\end{equation}
를 갖는다. 보통과 같이 $R^p\lim(K) = H^p(R\lim(K))$로 둔다. 또한
다음이 성립한다.
\begin{enumerate}
\item $\textit{Ab}(\mathbf{N})$의 임의의 $(A_n)$에 대하여
$p > 1$이면 $R^p\lim A_n = 0$이다.
\item $D(\textit{Ab})$의 대상 $R\lim A_n$은 차수 $0$과 $1$에 놓인
복합체
$$
\prod A_n \to \prod A_n,\quad (x_n) \mapsto (x_n - f_{n + 1}(x_{n + 1}))
$$
로 나타내어진다.
\item $(A_n)$이 ML이면 $R^1\lim A_n = 0$이다. 즉 $(A_n)$은
$\lim$에 대해 오른쪽 비순환이다.
\item 모든 $K^\bullet \in D(\textit{Ab}(\mathbf{N}))$은 그 각 항이
$\lim$에 대해 오른쪽 비순환인 복합체와 준동형이다.
\item 각 $K^p = (K^p_n)$가 $\lim$에 대해 오른쪽 비순환이면, 즉
% P02-E0347: frozen clause has a stray “of” before the vanishing condition; Korean states the intended condition.
$R^1\lim_n K^p_n = 0$이면 $R\lim K$는 차수 $p$의 항이
$\lim_n K_n^p$인 복합체로 나타내어진다.
\end{enumerate}
\end{lemma}

\begin{proof}
$(A_n)$을 임의의 역계라 하자. $(B_n)$을
$$
B_n = A_n \oplus A_{n - 1} \oplus \ldots \oplus A_1
$$
이고 전이 사상들이 사영으로 주어지는 역계라 하자.
$A_n \to B_n$을
$(1, f_n, f_{n - 1} \circ f_n, \ldots, f_2 \circ \ldots \circ f_n)$으로
주자. 여기서 $f_i : A_i \to A_{i - 1}$들은 전이 사상들이다.
따라서 모든 역계는 ML 계의 부분대상이다(Homology, Section
\ref{homology-section-inverse-systems}). Derived Categories, Lemma
\ref{derived-lemma-subcategory-right-acyclics}와 Homology, Lemma
\ref{homology-lemma-Mittag-Leffler}에 의해 모든 ML 계는 $\lim$에
대해 오른쪽 비순환이다. 즉 (3)이 성립한다. 이는 이미
% P02-E0348: frozen source names undefined RF instead of Rlim; Korean preserves the formal symbol.
$RF$가 $D^+(\textit{Ab}(\mathbf{N}))$ 위에서 정의됨을 함의한다.
Derived Categories, Proposition
\ref{derived-proposition-enough-acyclics}를 보라.
$n > 1$일 때 $C_n = A_{n - 1} \oplus \ldots \oplus A_1$이고
$C_1 = 0$이며 전이 사상들이 역시 사영인 계를 두자. 그러면 역계들의
짧은 완전열
$0 \to (A_n) \to (B_n) \to (C_n) \to 0$이 있으며,
$B_n \to C_n$은
$(x_i) \mapsto (x_i - f_{i + 1}(x_{i + 1}))$로 주어진다.
$(C_n)$도 ML이므로 (2)를 얻는다.
% P02-E1056: frozen parenthesis is an unresolved proposition-reference placeholder; Korean points to the already cited proposition.
(위에서 인용한 명제에 의한다.) 이는 (1)도 함의한다. 마지막으로
Derived Categories, Lemma
\ref{derived-lemma-unbounded-right-derived}에 의해 $R\lim$은 실제로
$D(\textit{Ab}(\mathbf{N}))$ 전체에서 정의된다. 실제로 Derived
Categories, Lemma \ref{derived-lemma-unbounded-right-derived}의 증명은
(4)와 (5)를 증명하는 방식으로 진행된다.
\end{proof}

\begin{lemma}
\label{more-algebra-lemma-six-term-Rlim}
아벨 군의 역계들의 짧은 완전열
$$
0 \to (A_i) \to (B_i) \to (C_i) \to 0
$$
이 주어졌다고 하자. 그러면 이에 연관된 $6$항 완전열
$0 \to \lim A_i \to \lim B_i \to \lim C_i \to
R^1\lim A_i \to R^1\lim B_i \to R^1\lim C_i \to 0$이 존재한다.
\end{lemma}

\begin{proof}
% P02-E0349: frozen proof is a sentence fragment; Korean supplies a complete sentence.
이는 보조정리 \ref{more-algebra-lemma-compute-Rlim}의 소멸에서 따른다.
\end{proof}

\noindent
다음은 Homology, Lemma
\ref{homology-lemma-apply-Mittag-Leffler-again}의 “올바른” 정식화이다.

\begin{lemma}
\label{more-algebra-lemma-apply-Mittag-Leffler-again}
아벨 군 복합체들의 역계
$$
(A^{-2}_n \to A^{-1}_n \to A^0_n \to A^1_n)
$$
이 주어졌다고 하고, 그 극한을
% P02-E0350: frozen source omits “by” in three notation clauses; Korean introduces each object grammatically.
$A^{-2} \to A^{-1} \to A^0 \to A^1$로 나타내자. 코호몰로지들의
역계를 $(H_n^{-1})$, $(H_n^0)$로 나타내고,
$A^{-2} \to A^{-1} \to A^0 \to A^1$의 코호몰로지들을
$H^{-1}$, $H^0$로 나타내자. 다음을 가정하자.
\begin{enumerate}
\item $(A^{-2}_n)$과 $(A^{-1}_n)$의 $R^1\lim$이 영이다.
\item $(H^{-1}_n)$의 $R^1\lim$이 영이다.
\end{enumerate}
그러면 $H^0 = \lim H_n^0$이다.
\end{lemma}

\begin{proof}
$K \in D(\textit{Ab}(\mathbf{N}))$를 그 $n$번째 성분이 복합체
$A^{-2}_n \to A^{-1}_n \to A^0_n \to A^1_n$인 복합체들의 계로
나타내어지는 대상이라 하자. Derived Categories, Lemma
\ref{derived-lemma-two-ss-complex-functor}의 두 스펙트럼 열을 모두
사용하여 $H^0(R\lim K)$를 계산할 것이다.\footnote{이 스펙트럼 열들을
사용하려면 $\textit{Ab}(\mathbf{N})$이 충분한 단사 대상을 가짐을
보여야 한다.
% P02-E0351: frozen footnote uses the wrong article before “inverse”; Korean is unaffected.
아벨 군의 역계 $(I_n)$이 단사일 필요충분조건은 각 $I_n$이 단사 아벨
군이고 전이 사상들이 분할 전사인 것이다. 모든 계는 그러한 계에
매장된다. 세부사항은 생략한다.}
첫 번째 것은 $E_1$-쪽
$$
\begin{matrix}
0 & 0 & R^1\lim A^0_n  & R^1\lim A^1_n \\
A^{-2} & A^{-1} & A^0  & A^1
\end{matrix}
$$
을 갖는다.
% P02-E0352: frozen sentence mixes a prepositional phrase and finite clause; Korean separates the predicates.
미분들은 수평이고 모든 고차 미분은 영이다. 두 번째 것은 $E_2$ 쪽
$$
\begin{matrix}
R^1\lim H^{-2}_n & 0 & R^1\lim H^0_n & R^1 \lim H^1_n \\
\lim H^{-2}_n &
\lim H^{-1}_n &
\lim H^0_n &
\lim H^1_n
\end{matrix}
$$
을 가지며 이 쪽에서 퇴화한다. 결론이 따른다.
\end{proof}
\begin{lemma}
\label{more-algebra-lemma-pro-isomorphism-iso-Rlim}
$(A_n)$과 $(B_n)$을 아벨 군의 역계라 하자. pro-계의 사상
$\varphi : (A_n) \to (B_n)$은 사상
$\lim A_n \to \lim B_n$과 $R^1\lim A_n \to R^1\lim B_n$을
결정한다. $\varphi$가 pro-동형이면 이 사상들은 동형이다.
\end{lemma}

\begin{proof}
pro-계의 사상에 관한 논의는 Categories, Example
\ref{categories-example-pro-morphism-inverse-systems}를 보라.
사상 $\varphi$는 $1 \leq m_1 < m_2 < m_3 < \ldots$와 전이
사상들과 양립하는 사상 $\varphi_n : A_{m_n} \to B_n$들로 주어진다.
$\varphi', \varphi'' : \prod A_n \to \prod B_n$을
$\varphi'((a_n)) = (\varphi_n(a_{m_n}))$ 및
$$
\varphi''((a_n)) =
(\varphi_n(a_{m_n} + f(a_{m_n + 1}) + \ldots + f(a_{m_{n + 1} - 1})))
$$
와 같게 두자.
% P02-E0353: frozen source misspells “occurrence”; Korean uses the intended term.
여기서 $f$의 각 출현은 역계 $(A_n)$의 알맞은 전이 사상을 뜻한다.
그러면 수평 화살표들이 보조정리 \ref{more-algebra-lemma-compute-Rlim}에서와 같은
도표
$$
\xymatrix{
\prod A_n \ar[r]_\delta \ar[d]^{\varphi'} &
\prod A_n \ar[d]^{\varphi''} \\
\prod B_n \ar[r]^\delta & \prod B_n
}
$$
는 가환한다. 이렇게 하여 원하는 사상들을 얻는다. 이 구성은 다음
의미에서 함자적이다. 역계 $(C_n)$과
$1 \leq m'_1 < m'_2 < m'_3 < \ldots$, 그리고
% P02-E0354: frozen source misspells “transition”; Korean uses the established term.
전이 사상들과 양립하는 사상 $\psi_n : B_{m'_n} \to C_n$들이 주어지면
$(\psi \circ \varphi)' = \psi' \circ \varphi'$이고
$(\psi \circ \varphi)'' = \psi'' \circ \varphi''$이다. 여기서 합성
$\psi \circ \varphi$는 정수들
$1 \leq m_{m'_1} < m_{m'_2} < \ldots$와 사상들
$\psi_n \circ \varphi_{m'_n} : A_{m_{m'_n}} \to C_n$을 가리킨다.
$B_n = A_n$이고 $\varphi_n$이 전이 사상이면 그 결과로 생기는
사상들이 항등사상임을 보이겠다. 이는 구성이 $\varphi$의 표시
$(m_n, \varphi_n)$의 선택에 무관함과 보조정리의 마지막 명제를 모두
증명할 것이다.

\medskip\noindent
따라서 $B_n = A_n$이고 $\varphi_n$이 전이 사상이라고 하자.
$(a_n) \in \lim A_n$을 원소라 하자. 그러면 $\varphi'((a_n))$의
차수 $n$ 성분은 $a_{m_n}$의 상이고, 이는 $a_n$과 같다. 이는
$\lim A_n$에 대한 명제를 증명한다.
$\xi \in R^1\lim A_n$을 $\prod A_n$의 원소 $(a_n)$의 동치류라 하자.
다음을 만족하는 원소 $(b_n)$을 생각하자.
$$
b_i = a_i + f(a_{i + 1}) + \ldots + f(a_{m_n - 1})
$$
$m_{n - 1} < i < m_n$이면 이와 같고, 어떤 $n$에 대해
$i = m_n$이면 $0$이라고 하자. 그러면
$(a'_n) = (a_n) - \delta((b_n))$은 $R^1\lim A_n$에서 같은 동치류를
정의하며, 계산하면 어떤 $n$에 대해 $i = m_n$인 경우를 제외하고
$a'_i = 0$이다. 따라서 어떤 $n$에 대해 $i = m_n$인 경우를 제외하면
$a_i = 0$이라고 가정할 수 있고 그렇게 하자. 영이 아닌 성분들이
자리 $j$와 $m_j$에 있는
$$
(0, \ldots, -f(a_{m_j}), 0, \ldots, 0, a_{m_j}, 0, \ldots)
$$
는 영이 아닌 성분들이 자리 $j + 1, \ldots, m_j$에 있는
$$
c_j = (0, \ldots, 0, f(a_{m_j}), f(a_{m_j}), \ldots, a_{m_j}, 0, \ldots)
$$
의 $\delta$ 아래 상이다. 합 $c = \sum c_j$는 $\prod A_n$에서
의미가 있다. $\varphi''((a_n)) = (a_{m_1}, a_{m_2}, \ldots)$임을
상기하면
$$
(a_n) - \varphi''((a_n)) = \delta(c)
$$
를 얻고 증명이 끝난다.
\end{proof}

\begin{lemma}
\label{more-algebra-lemma-map-into-Rlim}
$\mathcal{D}$를 삼각 범주라 하고 $(K_n)$을 $\mathcal{D}$의 대상들의
역계라 하자. $K$를 계 $(K_n)$의 유도 극한이라 하자. 그러면
$\mathcal{D}$의 모든 $L$에 대하여 짧은 완전열
$$
0 \to R^1\lim \Hom_\mathcal{D}(L, K_n[-1]) \to
\Hom_\mathcal{D}(L, K) \to
\lim \Hom_\mathcal{D}(L, K_n) \to 0
$$
이 존재한다.
\end{lemma}

\begin{proof}
이는 Derived Categories, Definition
\ref{derived-definition-derived-limit}, 보조정리
\ref{derived-lemma-representable-homological}, 그리고 위의 보조정리
\ref{more-algebra-lemma-compute-Rlim}에 주어진 $\lim$과 $R^1\lim$의 기술에서
따른다.
\end{proof}

\begin{lemma}
\label{more-algebra-lemma-pro-isomorphism-Rlim}
$\mathcal{D}$를 삼각 범주라 하자. $(K_n)$과 $(M_n)$을 유도 극한이
각각 $K$와 $M$인 $\mathcal{D}$의 대상들의 역계라 하자.
$a : (K_n) \to (M_n)$을 pro-대상들의 pro-동형이라 하자. 그러면
$a$를 사용하여 (비표준적인) 동형 $K \to M$을 만들 수 있다.
\end{lemma}

\begin{proof}
Derived Categories, Remark
\ref{derived-remark-functoriality-derived-limit}에 의해 정의하는 구별
삼각들의 사상
$$
\xymatrix{
K \ar[r] \ar[d] & \prod K_n \ar[r] \ar[d]^{a'} &
\prod K_n \ar[d]^{a''} \\
M \ar[r] & \prod M_n \ar[r] & \prod M_n
}
$$
에 들어가는 화살표 $K \to M$을 얻는다. 따라서 $\mathcal{D}$의 모든
대상 $L$에 대하여 짧은 완전열들의 사상
$$
\xymatrix{
0 \ar[r] &
R^1\lim \Hom_\mathcal{D}(L, K_n[-1]) \ar[r] \ar[d] &
\Hom_\mathcal{D}(L, K) \ar[r] \ar[d] &
\lim \Hom_\mathcal{D}(L, K_n) \ar[r] \ar[d] &
0 \\
0 \ar[r] &
R^1\lim \Hom_\mathcal{D}(L, M_n[-1]) \ar[r] &
\Hom_\mathcal{D}(L, M) \ar[r] &
\lim \Hom_\mathcal{D}(L, M_n) \ar[r] &
0
}
$$
을 얻는다. 보조정리 \ref{more-algebra-lemma-map-into-Rlim}과 그 증명을 보라.
보조정리 \ref{more-algebra-lemma-pro-isomorphism-iso-Rlim}에 의해 왼쪽과 오른쪽
수직 화살표는 동형이다.\footnote{pro-동형 $a$가 사상들
$a_n : K_{m_n} \to M_n$으로 주어지면,
% P02-E0355: frozen footnote has “the” where “then” is required; Korean states the consequence.
대응하는 사상 $\Hom_\mathcal{D}(L, K_{m_n}) \to
\Hom_\mathcal{D}(L, M_n)$을 사용하여 $\Hom$ 위의 pro-동형을 얻는다.
따라서 보조정리 \ref{more-algebra-lemma-pro-isomorphism-iso-Rlim}의 증명에서
$\lim$과 $R^1\lim$ 위의 화살표들을 구성하는 방식은 Derived
Categories, Remark \ref{derived-remark-functoriality-derived-limit}에서
$a'$와 $a''$을 구성하는 방식과 일치한다.} 따라서 가운데 화살표는
동형이다. Yoneda 보조정리에 의해 사상 $K \to M$은 동형이다.
\end{proof}

\begin{lemma}
\label{more-algebra-lemma-map-from-hocolim}
$\mathcal{D}$를 삼각 범주라 하자. $(K_n)$을 $\mathcal{D}$의 대상들의
계라 하고 $K$를 계 $(K_n)$의 유도 여극한이라 하자. 그러면
$\mathcal{D}$의 모든 $L$에 대하여 짧은 완전열
$$
0 \to R^1\lim \Hom_\mathcal{D}(K_n, L[-1]) \to
\Hom_\mathcal{D}(K, L) \to
\lim \Hom_\mathcal{D}(K_n, L) \to 0
$$
이 존재한다.
\end{lemma}

\begin{proof}
이는 Derived Categories, Definition
\ref{derived-definition-derived-colimit}, 보조정리
\ref{derived-lemma-representable-homological}, 그리고 위의 보조정리
\ref{more-algebra-lemma-compute-Rlim}에 주어진 $\lim$과 $R^1\lim$의 기술에서
따른다.
\end{proof}

\begin{remark}[코호몰로지로서의 Rlim]
\label{more-algebra-remark-Rlim-cohomology}
대상이 자연수이고 $j \geq i$이면 유일한 사상 $i \to j$를 갖는 범주
$\mathbf{N}$을 생각하자. $\mathbf{N}$에 카오스 위상(Sites, Example
\ref{sites-example-indiscrete})을 주면, 층 $\mathcal{F}$는
$\mathbf{N}$ 위의 집합들의 역계
$$
\mathcal{F}_1 \leftarrow \mathcal{F}_2 \leftarrow \mathcal{F}_3
\leftarrow \ldots
$$
와 같은 것이다. 다음에 유의하자.
$\Gamma(\mathbf{N}, \mathcal{F}) = \lim \mathcal{F}_n$.
아벨 군들의 역계 $\mathcal{F}_n$에 대해서는
$$
R^p\lim \mathcal{F}_n = H^p(\mathbf{N}, \mathcal{F})
$$
이다. 실제로 양변은 모두
$\mathcal{F} \mapsto \lim \mathcal{F}_n = H^0(\mathbf{N}, \mathcal{F})$의
고차 오른쪽 유도 함자이다. 따라서 $R\lim$의 존재는 Cohomology on
Sites, Sections \ref{sites-cohomology-section-cohomology-sheaves} and
\ref{sites-cohomology-section-unbounded}의 일반 이론에서도 따른다.
\end{remark}

\noindent
다음 보조정리의 곱들은 복합체들의 항별 곱으로 볼 수도 있고 유도 범주
$D(\textit{Ab})$에서의 곱으로 볼 수도 있다. Derived Categories, Lemma
\ref{derived-lemma-products}를 보라.

\begin{lemma}
\label{more-algebra-lemma-distinguished-triangle-Rlim}
$K = (K_n^\bullet)$를 $D(\textit{Ab}(\mathbf{N}))$의 대상이라 하자.
$D(\textit{Ab})$에는 표준적인 구별 삼각형
$$
R\lim K \to \prod\nolimits_n K_n^\bullet \to \prod\nolimits_n K_n^\bullet
\to R\lim K[1]
$$
이 존재한다. 다시 말해 $R\lim K$는 $D(\textit{Ab})$의 역계
$(K_n^\bullet)$의 유도 극한이다. Derived Categories, Definition
\ref{derived-definition-derived-limit}를 보라.
\end{lemma}

\begin{proof}
각 $p$에 대해 역계 $(K_n^p)$가 $\lim$에 대해 오른쪽 비순환이라고
가정하자. 보조정리 \ref{more-algebra-lemma-compute-Rlim}에 의해 각 $p$에 대해
짧은 완전열
$$
0 \to \lim_n K^p_n \to \prod\nolimits_n K^p_n \to \prod\nolimits_n K^p_n \to 0
$$
을 얻는다. $\lim_n K^p_n$들로 이루어진 복합체는 보조정리
\ref{more-algebra-lemma-compute-Rlim}에 의해 $R\lim K$를 계산하므로, 이 경우에는
보조정리가 성립한다.

\medskip\noindent
이제 $K = (K_n^\bullet)$가 일반적인 경우를 생각하자. 보조정리
\ref{more-algebra-lemma-compute-Rlim}에 의해 $D(\textit{Ab}(\mathbf{N}))$ 안에
유사동형 $K \to L$이 존재하며, 각 $p$에 대해 $(L_n^p)$는 비순환이다.
$\textit{Ab}$에서 곱은 완전하므로 $\prod K_n^\bullet$은
$\prod L_n^\bullet$과 유사동형이다. 따라서 위에서 증명한 $L$에 대한
결과로부터 $K$에 대한 결과가 따른다.
\end{proof}

\begin{lemma}
\label{more-algebra-lemma-break-long-exact-sequence}
보조정리 \ref{more-algebra-lemma-distinguished-triangle-Rlim}의 표기 아래, 구별
삼각형에 딸린 코호몰로지 긴 완전열은 짧은 완전열들
$$
0 \to R^1\lim_n H^{p - 1}(K_n^\bullet) \to
H^p(R\lim K) \to
\lim_n H^p(K_n^\bullet) \to 0
$$
로 분해된다.
\end{lemma}

\begin{proof}
구별 삼각형의 긴 완전열은
$$
\ldots \to H^p(R\lim K) \to \prod\nolimits_n H^p(K_n^\bullet)
\to \prod\nolimits_n H^p(K_n^\bullet) \to
H^{p + 1}(R\lim K) \to \ldots
$$
이다. 가운데 사상의 핵은 보조정리에 주어진 명시적 기술에 의해
$\lim_n H^p(K_n^\bullet)$이다. 이 사상의 여핵은 보조정리
\ref{more-algebra-lemma-compute-Rlim}에 의해 $R^1\lim_n H^p(K_n^\bullet)$이다.
\end{proof}

\noindent
{\bf 경고.} $D(\textit{Ab}(\mathbf{N}))$의 대상은 아벨 군들의 역계로
이루어진 복합체이다. 이를 복합체들의 역계 $(K_n^\bullet)$로 생각할
수도 있다. 그러나 이것은 $D(\textit{Ab})$의 대상들의 역계와
{\bf 같은 것이 아니다}. 다음 보조정리와 주석이 그 차이를 설명한다.

\begin{lemma}
\label{more-algebra-lemma-lift-to-system-complexes-Ab}
$(K_n)$을 $D(\textit{Ab})$의 대상들의 역계라 하자. 그러면
$D(\textit{Ab}(\mathbf{N}))$의 대상 $M = (M_n^\bullet)$과
$D(\textit{Ab})$에서의 동형들 $M_n^\bullet \to K_n$이 존재하여 도식
$$
\xymatrix{
M_{n + 1}^\bullet \ar[d] \ar[r] &
M_n^\bullet \ar[d] \\
K_{n + 1} \ar[r] & K_n
}
$$
은 $D(\textit{Ab})$에서 가환한다.
\end{lemma}

\begin{proof}
구체적으로, $M_1^\bullet$을 $K_1$을 나타내는 아벨 군들의 복합체라
하자. 이미
$M_e^\bullet \to M_{e - 1}^\bullet \to \ldots \to M_1^\bullet$
과 사상들 $\psi_i : M_i^\bullet \to K_i$를 구성하여, 보조정리의
도식들이 모든 $n < e$에 대해 가환한다고 가정하자. 그러면
$D(\textit{Ab})$에서 다음 도식을 생각한다.
% P02-E0356: frozen induction step switches from e to an undeclared n; formal indices are preserved.
$$
\xymatrix{
& M_n^\bullet \ar[d]^{\psi_n} \\
K_{n + 1} \ar[r] & K_n
}
$$
$D(\textit{Ab})$의 사상의 정의에 의해 아벨 군들의 복합체
$M_{n + 1}^\bullet$, $D(\textit{Ab})$에서의 동형
$M_{n + 1}^\bullet \to K_{n + 1}$, 그리고 합성
$$
K_{n + 1} \to K_n \xrightarrow{\psi_n^{-1}} M_n^\bullet
$$
을 $D(\textit{Ab})$에서 나타내는 복합체들의 사상
$M_{n + 1}^\bullet \to M_n^\bullet$을 찾을 수 있다. 따라서 귀납법으로
보조정리가 성립한다.
\end{proof}

\begin{remark}
\label{more-algebra-remark-compare-derived-limit}
$(K_n)$을 $D(\textit{Ab})$의 대상들의 역계라 하자. $K = R\lim K_n$을
이 계의 유도 극한이라 하자(Derived Categories, Section
\ref{derived-section-derived-limit}을 보라). $D(\textit{Ab})$는 가산
곱을 가지므로(Derived Categories, Lemma \ref{derived-lemma-products}),
그러한 유도 극한이 존재한다. 보조정리
\ref{more-algebra-lemma-lift-to-system-complexes-Ab}에 의해 $(K_n)$을
$D(\textit{Ab}(\mathbf{N}))$의 대상 $M$으로 들어 올릴 수도 있다.
그러면 $K \cong R\lim M$이다. 여기서 $R\lim$은 함자
(\ref{more-algebra-equation-Rlim})이다. 실제로 보조정리
\ref{more-algebra-lemma-distinguished-triangle-Rlim}에 의해 $R\lim M$도 계
$(K_n)$의 유도 극한이다. 따라서 계 $(K_n)$의 들어 올림 $M$에는 여러
동형류가 있을 수 있지만, $R\lim M$의 동형류는 선택과 무관하다.
그것이 이 계의 유도 극한 $K = R\lim K_n$과 동형이기 때문이다.
그러므로 이 절에서 $R\lim$에 대해 증명한 결과들을 유도 극한들에
적용할 수 있다. 예를 들어 모든 $p \in \mathbf{Z}$에 대해 표준적인
짧은 완전열
$$
0 \to R^1\lim H^{p - 1}(K_n) \to H^p(K) \to \lim H^p(K_n) \to 0
$$
이 존재한다. 이는 보조정리 \ref{more-algebra-lemma-break-long-exact-sequence}를
$M$에 적용할 수 있기 때문이다.
% P02-E0357: frozen source has “can also been seen”; Korean renders the intended construction.
$M$의 존재를 사용하지 않고도 이를 직접 알 수 있다. 실제로 보조정리
\ref{more-algebra-lemma-break-long-exact-sequence}의 증명에 쓰인 논증을 (정의에
나오는) 구별 삼각형 $K \to \prod K_n \to \prod K_n \to K[1]$에
적용하면 된다.
\end{remark}

\begin{lemma}
\label{more-algebra-lemma-Rlim-pro-equal}
$E \to D$를 $D(\textit{Ab}(\mathbf{N}))$의 사상이라 하자. $(E_n)$과
$(D_n)$을 각각 $E$와 $D$에 연관된 $D(\textit{Ab})$의 대상들의 계라
하자. $(E_n) \to (D_n)$이 pro-대상들의 동형이면
$R\lim E \to R\lim D$는 $D(\textit{Ab})$에서 동형이다.
\end{lemma}

\begin{proof}
가정은 특히 pro-대상들 $H^p(E_n)$과 $H^p(D_n)$이 동형임을 뜻한다.
결론은 보조정리 \ref{more-algebra-lemma-break-long-exact-sequence}의 짧은 완전열들과
보조정리 \ref{more-algebra-lemma-pro-isomorphism-iso-Rlim}에서 따른다.
\end{proof}

\begin{lemma}[Emmanouil]
\label{more-algebra-lemma-emmanouil}
\begin{reference}
\cite{Emmanouil}에서 가져왔다.
\end{reference}
$(A_n)$을 아벨 군들의 역계라 하자. 다음 조건들은 동치이다.
\begin{enumerate}
\item $(A_n)$은 Mittag-Leffler이다.
\item $R^1\lim A_n = 0$이고
$\bigoplus_{i \in \mathbf{N}} (A_n)$에 대해서도 같은 것이 성립한다.
\end{enumerate}
\end{lemma}

\begin{proof}
$B = \bigoplus_{i \in \mathbf{N}} (A_n)$라 두자. 따라서
$B = (B_n)$이고 $B_n = \bigoplus_{i \in \mathbf{N}} A_n$이다.
$(A_n)$이 ML이면 $B$도 ML이고, 따라서 보조정리
\ref{more-algebra-lemma-compute-Rlim}에 의해 $R^1\lim A_n = 0$이며
$R^1\lim B_n = 0$이다.

\medskip\noindent
반대로 $(A_n)$이 ML이 아니라고 가정하자. 그러면 어떤 $m$, 정수열
$m < m_1 < m_2 < \ldots$, 그리고 원소들 $x_i \in A_{m_i}$를 골라서
그 상 $y_i \in A_m$이 $A_{m_i + 1} \to A_m$의 상에 속하지 않게 할
수 있다. 원소들 $x_i$와 $y_i$를 사용하여
$R^1\lim B_n \not = 0$임을 두 가지 방법으로 보이겠다. 이로써
보조정리의 증명이 끝날 것이다.

\medskip\noindent
첫 번째 증명. $C_n = \prod_{i \in \mathbf{N}} A_n$으로
$C = (C_n)$라 두자. 표준적인 단사 사상 $B_n \to C_n$이 있고 그
여핵을 $Q_n$이라 하자. $Q = (Q_n)$이라 두자. $Q_n$의 원소 $q_n$을
$q_{n, i} \in A_n$인 원소들의 열
$q_n = (q_{n, 1}, q_{n, 2}, \ldots)$을 꼬리가 영인 열들로 나눈
것으로 생각해도 되고, 그렇게 하겠다(다시 말해 유한 개의 자리에서만
다른 열들을 같은 것으로 본다). 역계들의 짧은 완전열
$$
0 \to (B_n) \to (C_n) \to (Q_n) \to 0
$$
이 있다. 다음으로 주어지는 원소 $q_n \in Q_n$을 생각하자.
$$
q_{n, i} =
\left\{
\begin{matrix}
\text{image of }x_i &\text{if}& m_i \geq n \\
0 & \text{else}
\end{matrix}
\right.
$$
그러면 $q_{n + 1}$이 $q_n$으로 가는 것은 명백하다. 따라서
$q = (q_n) \in \lim Q_n$을 얻는다. 한편 $q$는
$\lim C_n \to \lim Q_n$의 상에 속하지 않는다고 주장한다. 실제로
$c = (c_n)$이 $q$로 간다고 하자. 그러면 $c_n = (c_{n, i})$라 쓸
수 있고, $n' \geq n$일 때 $c_{n', i} \mapsto c_{n, i}$이므로 모든
$n, i, n' \geq n$에 대해
$c_{n, i} \in \Im(C_{n'} \to C_n)$임을 알 수 있다. 특히
% P02-E0359: frozen source calls this the image of c_{m,i} in A_m although it already lies there.
$c_{m, i}$의 $A_m$ 안에서의 상은
$\Im(A_{m_i + 1} \to A_m)$에 속하므로 $y_i$와 같을 수 없다.
따라서 $c_m$과 $q_m = (y_1, y_2, y_3, \ldots)$은 무한히 많은
자리에서 다르며, 이는 모순이다. 코호몰로지 긴 완전열
$$
0 \to \lim B_n \to \lim C_n \to \lim Q_n \to R^1\lim B_n
$$
을 생각하면, 원하는 대로 마지막 군이 영이 아님을 결론내린다.

\medskip\noindent
두 번째 증명.
% P02-E0358: frozen source omits “by” in the notation sentence; Korean states the notation naturally.
$n' \geq n$에 대해 $A_{n, n'} = \Im(A_{n'} \to A_n)$이라 쓰자.
그러면 $y_i \in A_m$이고 $y_i \not \in A_{m, m_i + 1}$이다.
$\xi = (\xi_n) \in \prod B_n$을 다음과 같은 원소라 하자.
$n = m_i$가 아니면 $\xi_n = 0$이고,
$\xi_{m_i} = (0, \ldots, 0, x_i, 0, \ldots)$에서 $x_i$는 $i$번째
직화 성분에 놓인다. $\xi$가 보조정리 \ref{more-algebra-lemma-compute-Rlim}의 사상
$\prod B_n \to \prod B_n$의 상에 속하지 않는다고 주장한다. 이로부터
$R^1\lim B_n$이 영이 아님을 알 수 있고 증명이 끝난다. 실제로
$\xi$가 $\eta = (z_1, z_2, \ldots)$의 상이라고 가정하자. 여기서
$z_n = \sum z_{n, i} \in \bigoplus_i A_n$이다. 다음에 유의하자.
$x_i = z_{m_i, i} \bmod A_{m_i, m_i + 1}$.
그러면 $z_{m_i - 1, i}$는 $A_{m_i} \to A_{m_i - 1}$ 아래
$z_{m_i, i}$의 상이고, 같은 방식으로 계속하면 $z_{m, i}$는
$A_{m_i} \to A_m$ 아래 $z_{m_i, i}$의 상이다. 따라서 $z_{m, i}$는
$A_{m, m_i + 1}$을 법으로 $y_i$와 합동이다. 특히
$z_{m, i} \not = 0$이다. 이는 불가능하다. 실제로
$\sum z_{m, i} \in \bigoplus_i A_m$이므로 유한 개의
$z_{m, i}$만 영이 아닐 수 있다.
\end{proof}

\begin{lemma}
\label{more-algebra-lemma-Mittag-Leffler}
다음을 아벨 군들의 역계들의 짧은 완전열이라 하자.
$$
0 \to (A_i) \to (B_i) \to (C_i) \to 0
$$
$(A_i)$와 $(C_i)$가 ML이면 $(B_i)$도 ML이다.
\end{lemma}

\begin{proof}
이는 보조정리 \ref{more-algebra-lemma-emmanouil}, 무한 직합을 취하는 것이 완전하다는
사실, 그리고 $R\lim$에 딸린 코호몰로지 긴 완전열에서 따른다.
\end{proof}

\begin{lemma}
\label{more-algebra-lemma-Rlim-zero-of-direct-sums}
$(A_n)$을 아벨 군들의 역계라 하자. 다음 조건들은 동치이다.
\begin{enumerate}
\item $(A_n)$은 pro-대상으로서 영이다.
\item $\lim A_n = 0$이고 $R^1\lim A_n = 0$이며,
$\bigoplus_{i \in \mathbf{N}} (A_n)$에 대해서도 같은 것이 성립한다.
\end{enumerate}
\end{lemma}

\begin{proof}
보조정리 \ref{more-algebra-lemma-pro-isomorphism-iso-Rlim}에 의해 (1)은 (2)를
함의한다. (2)를 가정하자. 그러면 보조정리 \ref{more-algebra-lemma-emmanouil}에
의해 $(A_n)$은 ML이다. $m \geq n$에 대해
$A_{n, m} = \Im(A_m \to A_n)$이라 두면
$A_n = A_{n, n} \supset A_{n, n + 1} \supset \ldots$이다.
$(A_n)$이 pro-대상으로서 영이라는 것은 모든 $n$에 대해 어떤
$m \geq n$이 존재하여 $A_{n, m} = 0$이라는 것과 동치이다.
$(A_n)$이 ML이라는 것은 모든 $n$에 대해 어떤 $m_n \geq n$이
존재하여 $A_{n, m} = A_{n, m + 1} = \ldots$이라는 것과 동치이다.
ML인 경우 $\lim A_n = 0$이면 $A_{n, m_n} = 0$임이 명백하다. 실제로
% P02-E0360: frozen source leaves m free in the target of this map; the formal span is preserved.
사상들 $A_{n + 1, m_{n + 1}} \to A_{n, m}$은 전사이다. 이로써 증명이
끝난다.
\end{proof}


\section{가군의 Rlim}
\label{more-algebra-section-Rlim-modules}

\noindent
가군에 대한 $R\lim$을 간단히 논의한다. 이 절의 여러 논증은 환 달린
사이트 위의 가군에 대한 코호몰로지 장치를 구성할 때 쓴 논증을
되풀이한다.

\medskip\noindent
$(A_n)$을 환들의 역계라 하자.
% P02-E0361: frozen notation sentence omits “by”; Korean states the notation naturally.
$\textit{Mod}(\mathbf{N}, (A_n))$으로 다음 범주를 나타내겠다. 그 대상은
아벨 군들의 역계 $(M_n)$으로서, 각 $M_n$에는
% P02-E0362: frozen source uses the wrong English article before A_n-module.
$A_n$-가군 구조가 주어지고 전이 사상 $M_{n + 1} \to M_n$은
$A_{n + 1}$-가군 사상이다. 이는 아벨 범주이다. $A = \lim A_n$이라
두자. $\textit{Mod}(\mathbf{N}, (A_n))$의 대상 $(M_n)$이 주어지면
극한 $\lim M_n$은 $A$-가군이다.

\begin{lemma}
\label{more-algebra-lemma-compute-Rlim-modules}
% P02-E0363: frozen source begins with a detached sentence fragment; Korean integrates it.
위의 상황에서 함자
$\lim : \textit{Mod}(\mathbf{N}, (A_n)) \to \text{Mod}_A$는 오른쪽
유도 함자
$$
R\lim :
D(\textit{Mod}(\mathbf{N}, (A_n)))
\longrightarrow
D(A)
$$
를 가진다. 보통과 같이 $R^p\lim(K) = H^p(R\lim(K))$라 둔다.
더 나아가 다음이 성립한다.
\begin{enumerate}
\item $\textit{Mod}(\mathbf{N}, (A_n))$의 모든 $(M_n)$에 대해
$p > 1$이면 $R^p\lim M_n = 0$이다.
\item $D(\text{Mod}_A)$의 대상 $R\lim M_n$은 차수 $0$과 $1$에
놓인 복합체
$$
\prod M_n \to \prod M_n,\quad
(x_n) \mapsto (x_n - f_{n + 1}(x_{n + 1}))
$$
로 나타내어진다.
\item $(M_n)$이 ML이면 $R^1\lim M_n = 0$이다. 즉 $(M_n)$은
$\lim$에 대해 오른쪽 비순환이다.
\item 모든 $K^\bullet \in D(\textit{Mod}(\mathbf{N}, (A_n)))$은
각 항이 $\lim$에 대해 오른쪽 비순환인 복합체와 유사동형이다.
\item 각 $K^p = (K^p_n)$가 $\lim$에 대해 오른쪽 비순환이면, 즉
% P02-E0364: frozen source has a spurious “of” before the vanishing condition.
$R^1\lim_n K^p_n = 0$이면, $R\lim K$는 차수 $p$의 항이
$\lim_n K_n^p$인 복합체로 나타내어진다.
\end{enumerate}
\end{lemma}

\begin{proof}
이 증명은 보조정리 \ref{more-algebra-lemma-compute-Rlim}의 증명과 글자 그대로
같다.
\end{proof}

\begin{remark}
\label{more-algebra-remark-Rlim-cohomology-modules}
이 주석은 주석 \ref{more-algebra-remark-Rlim-cohomology}의 연속이다.
$\mathbf{N}$ 위의 환들의 층은 바로 환들의 역계 $(A_n)$이다.
$(A_n)$ 위의 가군들의 층은 위에서 정의한 범주
$\textit{Mod}(\mathbf{N}, (A_n))$의 대상과 정확히 같은 것이다.
보조정리 \ref{more-algebra-lemma-compute-Rlim-modules}의 유도 함자 $R\lim$은
가군들의 유도 범주에서 구조층의 전역 단면 환 위 가군들의 유도
범주로 가는 $R\Gamma(\mathbf{N}, -)$일 뿐이다. 일반적으로 군으로
계산한 코호몰로지와 가군으로 계산한 코호몰로지는 일치한다.
Cohomology on Sites, Lemma
\ref{sites-cohomology-lemma-cohomology-modules-abelian-agree}를 보라.
\end{remark}

\noindent
다음 보조정리의 곱들은 복합체들의 항별 곱으로 볼 수도 있고 유도
범주 $D(A)$에서의 곱으로 볼 수도 있다. Derived Categories, Lemma
\ref{derived-lemma-products}를 보라.

\begin{lemma}
\label{more-algebra-lemma-distinguished-triangle-Rlim-modules}
$K = (K_n^\bullet)$를
$D(\textit{Mod}(\mathbf{N}, (A_n)))$의 대상이라 하자.
$D(A)$에는 표준적인 구별 삼각형
$$
R\lim K \to \prod\nolimits_n K_n^\bullet \to \prod\nolimits_n K_n^\bullet
\to R\lim K[1]
$$
이 존재한다. 다시 말해 $R\lim K$는 $D(A)$의 역계
$(K_n^\bullet)$의 유도 극한이다. Derived Categories, Definition
\ref{derived-definition-derived-limit}을 보라.
\end{lemma}

\begin{proof}
증명은 보조정리 \ref{more-algebra-lemma-distinguished-triangle-Rlim}의 증명에서
% P02-E0365: frozen source spells “instead” as two words at both parallel proof loci.
보조정리 \ref{more-algebra-lemma-compute-Rlim-modules}를 보조정리
\ref{more-algebra-lemma-compute-Rlim} 대신 사용하면 정확히 같다.
\end{proof}

\begin{lemma}
\label{more-algebra-lemma-break-long-exact-sequence-modules}
보조정리 \ref{more-algebra-lemma-distinguished-triangle-Rlim-modules}의 표기 아래,
구별 삼각형에 딸린 코호몰로지 긴 완전열은 $A$-가군들의 짧은 완전열들
$$
0 \to R^1\lim_n H^{p - 1}(K_n^\bullet) \to
H^p(R\lim K) \to
\lim_n H^p(K_n^\bullet) \to 0
$$
로 분해된다.
\end{lemma}

\begin{proof}
증명은 보조정리 \ref{more-algebra-lemma-break-long-exact-sequence}의 증명에서
% P02-E0365: second occurrence of frozen “in stead”; the same source candidate applies.
보조정리 \ref{more-algebra-lemma-compute-Rlim-modules}를 보조정리
\ref{more-algebra-lemma-compute-Rlim} 대신 사용하면 정확히 같다.
\end{proof}

\noindent
{\bf 경고.} 아벨 군들의 경우와 마찬가지로
$D(\textit{Mod}(\mathbf{N}, (A_n)))$의 대상
$M = (M_n^\bullet)$은 가군 복합체들의 역계이며, 이는 유도 범주들의
대상들의 역계와 {\bf 같은 것이 아니다}. 다음 보조정리에서 유도
범주들의 대상들의 역계를 언제나
$D(\textit{Mod}(\mathbf{N}, (A_n)))$의 대상으로 들어 올릴 수 있음을
보인다.

\begin{lemma}
\label{more-algebra-lemma-lift-to-system-complexes}
$(A_n)$을 환들의 역계라 하자. 다음이 주어졌다고 가정하자.
\begin{enumerate}
\item 모든 $n$에 대해 $D(A_n)$의 대상 $K_n$.
\item 모든 $n$에 대해 $D(A_{n + 1})$의 사상
$\varphi_n : K_{n + 1} \to K_n$. 여기서 $A_{n + 1} \to A_n$을
통한 스칼라 제한으로 $K_n$을 $D(A_{n + 1})$의 대상으로 본다.
\end{enumerate}
그러면 대상
$M = (M_n^\bullet) \in D(\textit{Mod}(\mathbf{N}, (A_n)))$과
$D(A_n)$에서의 동형들 $\psi_n : M_n^\bullet \to K_n$이 존재하여
도식들
$$
\xymatrix{
M_{n + 1}^\bullet \ar[d]_{\psi_{n + 1}} \ar[r] & M_n^\bullet \ar[d]^{\psi_n} \\
K_{n + 1} \ar[r]^{\varphi_n} & K_n
}
$$
은 $D(A_{n + 1})$에서 가환한다.
\end{lemma}

\begin{proof}
증명을 자세히 쓰겠다. $A_n$-가군 $T$에 대해, 같은 가군을
% P02-E0366: frozen source misspells “viewed”; Korean states the intended relation.
$A_{n + 1}$-가군으로 본 것을 $T_{A_{n + 1}}$이라 쓰겠다.
$K_n^\bullet$을 $K_n$을 나타내는 $A_n$-가군들의 복합체라 하자.
그러면 $K_{n, A_{n + 1}}^\bullet$은 같은 복합체를
$A_{n + 1}$-가군들의 복합체로 본 것이다. 유도 범주의 구성에 의해
사상 $\psi_n$은 다음과 같이 주어질 수 있다.
% P02-E0367: frozen proof names this roof psi_n rather than the statement's phi_n; formal spans are preserved.
$$
\psi_n = \tau_n \circ \sigma_n^{-1}
$$
여기서 $\sigma_n : L_{n + 1}^\bullet \to K_{n + 1}^\bullet$은
$A_{n + 1}$-가군 복합체들의 유사동형이고,
$\tau_n : L_{n + 1}^\bullet \to K_{n, A_{n + 1}}^\bullet$은
$A_{n + 1}$-가군 복합체들의 사상이다.

\medskip\noindent
이제 복합체들 $M_n^\bullet$을 귀납적으로 구성한다. 기초 단계에서는
$M_1^\bullet = K_1^\bullet$이라 둔다. 이미
$M_e^\bullet \to M_{e - 1}^\bullet \to \ldots \to M_1^\bullet$과
복합체들의 사상들 $\psi_i : M_i^\bullet \to K_i^\bullet$을
구성하여, 위의 도식들
$$
\xymatrix{
M_{n + 1}^\bullet \ar[d]_{\psi_{n + 1}} \ar[rr] & &
M_{n, A_{n + 1}}^\bullet \ar[d]^{\psi_{n, A_{n + 1}}} \\
K_{n + 1}^\bullet & L_{n + 1}^\bullet \ar[l]_{\sigma_n} \ar[r]^{\tau_n} &
K_{n, A_{n + 1}}^\bullet
}
$$
이 모든 $n < e$에 대해 $D(A_{n + 1})$에서 가환한다고 가정하자.
그러면 도식
$$
\xymatrix{
& & M_{e, A_{e + 1}}^\bullet \ar[d]^{\psi_{e, A_{e + 1}}} \\
K_{e + 1}^\bullet & L_{e + 1}^\bullet \ar[r]^{\tau_e} \ar[l]_{\sigma_e} &
K_{e, A_{e + 1}}^\bullet
}
$$
을 $D(A_{e + 1})$에서 생각한다. $\psi_e$는 유사동형이므로
$\psi_{e, A_{e + 1}}$도 유사동형이다. $D(A_{e + 1})$에서의 사상의
정의에 의해, $A_{e + 1}$-가군 복합체들의 유사동형
$\psi_{e + 1} : M_{e + 1}^\bullet \to K_{e + 1}^\bullet$을 찾아서
$D(A_{e + 1})$에서 합성
$\psi_{e, A_{e + 1}}^{-1} \circ \tau_e \circ \sigma_e^{-1}$을
나타내는 $A_{e + 1}$-가군 복합체들의 사상
$M_{e + 1}^\bullet \to M_{e, A_{e + 1}}^\bullet$이 존재하게 할 수
있다. 따라서 귀납법으로 보조정리가 성립한다.
\end{proof}

\begin{remark}
\label{more-algebra-remark-how-unique}
% P02-E0368: frozen source starts with a detached assumptions fragment; Korean integrates it.
보조정리 \ref{more-algebra-lemma-lift-to-system-complexes}의 가정 아래, 선험적으로는
그 보조정리의 계 $(K_n, \varphi_n)$을 주는
$D(\textit{Mod}(\mathbf{N}, (A_n)))$의 대상 $M$의 동형류가 여럿
있다. 그러한 각 $M$에 대해 $A = \lim A_n$일 때 복합체
$R\lim M \in D(A)$를 생각할 수 있다. 보조정리
\ref{more-algebra-lemma-distinguished-triangle-Rlim-modules}에 의해 $R\lim M$은
$D(A)$의 역계 $(K_n)$의 유도 극한이다. 따라서 $D(A)$에서
$R\lim M$의 동형류는 $M$을 구성할 때 한 선택들과 무관하다.
특히 이 절에서 $R\lim$에 대해 증명한 결과들을 $D(A)$의 역계들의
유도 극한에 적용할 수 있다. 예를 들어 모든 $p \in \mathbf{Z}$에
대해 표준적인 짧은 완전열
$$
0 \to R^1\lim H^{p - 1}(K_n) \to H^p(R\lim K_n) \to \lim H^p(K_n) \to 0
$$
이 있다. 이는 보조정리 \ref{more-algebra-lemma-distinguished-triangle-Rlim-modules}를
$M$에 적용할 수 있기 때문이다.
% P02-E0369: frozen source has “can also been seen”; Korean renders the intended construction.
$M$의 존재를 사용하지 않고도 이를 직접 알 수 있다. 실제로 보조정리
\ref{more-algebra-lemma-distinguished-triangle-Rlim-modules}의 증명에 쓰인 논증을
유도 극한의 (정의에 나오는) 구별 삼각형
$R\lim K_n \to \prod K_n \to \prod K_n \to (R\lim K_n)[1]$에
적용하면 된다.
\end{remark}
\begin{lemma}
\label{more-algebra-lemma-get-ML-system}
$(A_n)$을 환들의 역계라 하자. 모든
$K \in D(\textit{Mod}(\mathbf{N}, (A_n)))$는 모든 전이 사상
$M_{n + 1}^\bullet \to M_n^\bullet$이 전사인 복합체들의 계
$(M_n^\bullet)$으로 나타낼 수 있다.
\end{lemma}

\begin{proof}
$K$가 계 $(K_n^\bullet)$으로 나타내어진다고 하자.
$M_1^\bullet = K_1^\bullet$이라 두자. 복합체들의 전사 사상들
$M_n^\bullet \to M_{n - 1}^\bullet \to \ldots \to M_1^\bullet$과
호모토피 동치들 $\psi_e : K_e^\bullet \to M_e^\bullet$을 이미
구성하여 도식들
$$
\xymatrix{
K_{e + 1}^\bullet \ar[d] \ar[r] & K_e^\bullet \ar[d] \\
M_{e + 1}^\bullet \ar[r] & M_e^\bullet
}
$$
이 모든 $e < n$에 대해 가환한다고 가정하자. 그러면 도식
$$
\xymatrix{
K_{n + 1}^\bullet \ar[r] & K_n^\bullet \ar[d] \\
& M_n^\bullet
}
$$
을 생각한다. Derived Categories, Lemma
\ref{derived-lemma-make-surjective}에 의해 합성
$K_{n + 1}^\bullet \to M_n^\bullet$을
$K_{n + 1}^\bullet \to M_{n + 1}^\bullet \to M_n^\bullet$로
인수분해하여, 첫째 화살표는 호모토피 동치이고 둘째 화살표는 항별
분할 전사가 되게 할 수 있다. 여기서 귀납법을 적용하면 보조정리가
따른다.
\end{proof}

\begin{lemma}
\label{more-algebra-lemma-get-K-flat-system}
$(A_n)$을 환들의 역계라 하자. 모든
$K \in D(\textit{Mod}(\mathbf{N}, (A_n)))$는 각
$K_n^\bullet$이 K-평탄인 복합체들의 계 $(K_n^\bullet)$으로 나타낼
수 있다.
\end{lemma}

\begin{proof}
먼저 보조정리 \ref{more-algebra-lemma-get-ML-system}을 사용하여 $K$를, 모든 전이
사상 $M_{n + 1}^\bullet \to M_n^\bullet$이 전사인 복합체들의 계
$(M_n^\bullet)$으로 나타낸다. 다음으로
$K_1^\bullet \to M_1^\bullet$을 $K_1^\bullet$이
$A_1$-가군들의 K-평탄 복합체인 유사동형으로 택한다(보조정리
\ref{more-algebra-lemma-K-flat-resolution}). 이미
$K_n^\bullet \to K_{n - 1}^\bullet \to \ldots \to K_1^\bullet$과
복합체들의 사상들 $\psi_e : K_e^\bullet \to M_e^\bullet$을
구성하여
$$
\xymatrix{
K_{e + 1}^\bullet \ar[d] \ar[r] & K_e^\bullet \ar[d] \\
M_{e + 1}^\bullet \ar[r] & M_e^\bullet
}
$$
이 모든 $e < n$에 대해 가환한다고 가정하자. 그러면
$D(A_{n + 1})$ 안의 도식
$$
\xymatrix{
C^\bullet \ar@{..>}[d] \ar@{..>}[r] & K_n^\bullet \ar[d]^{\psi_n} \\
M_{n + 1}^\bullet \ar[r]^{\varphi_n} & M_n^\bullet
}
$$
을 생각한다. $M_{n + 1}^\bullet \to M_n^\bullet$은 항별 전사이므로,
항들이
$$
C^p = M_{n + 1}^p \times_{M_n^p} K_n^p
$$
이고 왼쪽 위 모서리를 채우는 복합체 $C^\bullet$은
$M_{n + 1}^\bullet$과 유사동형이다(자세한 내용은 생략한다).
$K_{n +1}^\bullet$이 K-평탄인 유사동형
$K_{n + 1}^\bullet \to C^\bullet$을 택한다. 따라서 귀납법으로
보조정리가 성립한다.
\end{proof}

\begin{lemma}
\label{more-algebra-lemma-derived-tensor-product-systems}
$(A_n)$을 환들의 역계라 하자.
$K, L \in D(\textit{Mod}(\mathbf{N}, (A_n)))$이 주어지면
$D(A_n)$으로 가는 사상들과 양립하는 표준적인 유도 텐서곱
$K \otimes^\mathbf{L} L$이
% P02-E0370: frozen source omits Mod from the target category; the formal category is preserved.
$D(\mathbf{N}, (A_n))$ 안에 존재한다. 이 구성은 $K$와 $L$에 대해
대칭이며, 각 변수에 대해 삼각 범주들의 완전함자이다.
\end{lemma}

\begin{proof}
$K$의 대표 $(K_n^\bullet)$을 각 $K_n^\bullet$이 K-평탄 복합체가
되도록 택한다(보조정리 \ref{more-algebra-lemma-get-K-flat-system}).
그러면 $L$의 임의의 대표 $(L_n^\bullet)$에 대해 복합체들의 계
$$
(\text{Tot}(K_n^\bullet \otimes_{A_n} L_n^\bullet))
$$
으로 나타내어지는 대상을 $K \otimes^\mathbf{L} L$로 정의할 수 있다.
보조정리 \ref{more-algebra-lemma-K-flat-quasi-isomorphism}과
\ref{more-algebra-lemma-derived-tor-quasi-isomorphism-other-side}에 의해 이는 두
변수 모두에서 잘 정의된다. $D(A_n)$으로 가는 사상과의 양립성은
명백하다. 완전성은 보조정리 \ref{more-algebra-lemma-derived-tor-exact}에서와
정확히 같은 방식으로 따른다.
\end{proof}

\begin{remark}
\label{more-algebra-remark-constructing-tensor-with-limits-functorially}
$A$를 환이라 하자. $(E_n)$을 $D(A)$의 대상들의 역계라 하자.
위에서 유도 극한 $R\lim E_n$이 존재함을 보았다. 따라서
$D(A)$의 모든 대상 $K$에 대해 유도 극한
$R\lim( K \otimes_A^\mathbf{L} E_n )$도 존재한다.
사실 이러한 유도 극한들을 $K$에 대해 함자적으로 구성하여 삼각
범주들의 완전함자
$$
R\lim(- \otimes_A^\mathbf{L} E_n) : D(A) \longrightarrow D(A)
$$
를 얻을 수 있다. 구체적으로 먼저 $(E_n)$을
$D(\mathbf{N}, A)$의 대상 $E$로 들어 올린다. 보조정리
\ref{more-algebra-lemma-lift-to-system-complexes}를 보라. (이 함자는 이 들어
올림의 선택에 의존할 것이다.) 다음으로 ‘대각’ 또는 ‘상수’ 함자
$$
\Delta : D(A) \longrightarrow D(\mathbf{N}, A)
$$
가 있어서 복합체 $K^\bullet$을 값이 $K^\bullet$인 복합체들의 상수
역계로 보냄에 유의하자. 이제 단순히
$$
R\lim(K \otimes_A^\mathbf{L} E_n) = R\lim(\Delta(K)\otimes^\mathbf{L} E)
$$
로 정의한다. 우변에서는 보조정리
\ref{more-algebra-lemma-compute-Rlim-modules}의 함자 $R\lim$과 보조정리
\ref{more-algebra-lemma-derived-tensor-product-systems}의 함자
$- \otimes^\mathbf{L} -$를 사용한다.
\end{remark}

\begin{lemma}
\label{more-algebra-lemma-tensor-Rlim-exact}
$A$를 환이라 하자. $E \to D \to F \to E[1]$을
$D(\mathbf{N}, A)$의 구별 삼각형이라 하자. $(E_n)$, 각각 $(D_n)$,
각각 $(F_n)$을 $E$, 각각 $D$, 각각 $F$에 연관된 $D(A)$의 대상들의
계라 하자. 그러면 모든 $K \in D(A)$에 대해 표준적인 구별 삼각형
$$
R\lim (K \otimes^\mathbf{L}_A E_n) \to
R\lim (K \otimes^\mathbf{L}_A D_n) \to
R\lim (K \otimes^\mathbf{L}_A F_n) \to
R\lim (K \otimes^\mathbf{L}_A E_n)[1]
$$
이 $D(A)$ 안에 존재한다. 표기는 주석
\ref{more-algebra-remark-constructing-tensor-with-limits-functorially}에서와 같다.
\end{lemma}

\begin{proof}
이는 주석 \ref{more-algebra-remark-constructing-tensor-with-limits-functorially}의
구성과 $\Delta : D(A) \to D(\mathbf{N}, A)$,
$- \otimes^\mathbf{L} -$, 그리고 $R\lim$이 삼각 범주들의
완전함자라는 사실에서 명백하다.
\end{proof}

\begin{lemma}
\label{more-algebra-lemma-tensor-Rlim-pro-equal}
$A$를 환이라 하자. $E \to D$를 $D(\mathbf{N}, A)$의 사상이라 하자.
$(E_n)$과 $(D_n)$을 각각 $E$와 $D$에 연관된 $D(A)$의 대상들의
계라 하자. $(E_n) \to (D_n)$이 pro-대상들의 동형이면 모든
$K \in D(A)$에 대해 이에 대응하는 사상
$$
R\lim (K \otimes^\mathbf{L}_A E_n)
\longrightarrow
R\lim (K \otimes^\mathbf{L}_A D_n)
$$
은 $D(A)$에서 동형이다(표기는 주석
\ref{more-algebra-remark-constructing-tensor-with-limits-functorially}에서와 같다).
\end{lemma}

\begin{proof}
% P02-E0371: frozen proof is a subjectless sentence fragment; Korean supplies the subject.
이는 정의들과 보조정리 \ref{more-algebra-lemma-Rlim-pro-equal}에서 따른다.
\end{proof}




\section{꼬임 가군}
\label{more-algebra-section-torsion}

\noindent
이 절에서 ‘꼬임 가군’은 아핀 스킴 $\Spec(R)$의 주어진 닫힌 부분집합
$V(I)$ 위에 지지되는 가군을 뜻한다. 이는 정역 위의 꼬임 가군이라는
개념(정의 \ref{more-algebra-definition-torsion})과 다르지만 서로 유사하다.

\begin{definition}
\label{more-algebra-definition-f-power-torsion}
$R$을 환이라 하고 $M$을 $R$-가군이라 하자.
\begin{enumerate}
\item $I \subset R$을 아이디얼이라 하자. 모든 $m \in M$에 대해
어떤 $n > 0$이 존재하여 $I^n m = 0$이면, $M$을
{\it $I$-거듭제곱 꼬임 가군}이라고 한다.
\item $f \in R$이라 하자. 모든 $m \in M$에 대해 어떤 $n > 0$이
존재하여 $f^n m = 0$이면, $M$을
{\it $f$-거듭제곱 꼬임 가군}이라고 한다.
\end{enumerate}
\end{definition}

\noindent
따라서 $f$-거듭제곱 꼬임 가군은 $I = (f)$에 대한
$I$-거듭제곱 꼬임 가군과 같은 것이다. $R$-가군 $M$에 대해 표기
$$
M[I^n] = \{m \in M \mid I^nm = 0\}
$$
와
$$
M[I^\infty] = \bigcup M[I^n]
$$
를 사용하겠다. 따라서 $M$이 $I$-거듭제곱 꼬임이라는 것은
$M = M[I^\infty]$이라는 것과 동치이고, 이는 다시
$M = \bigcup M[I^n]$이라는 것과 동치이다.

\begin{lemma}
\label{more-algebra-lemma-I-power-torsion-presentation}
$R$을 환이라 하고 $I$를 $R$의 아이디얼이라 하자. $M$을
$I$-거듭제곱 꼬임 가군이라 하자. 그러면 $M$은 분해
$$
\ldots \to K_2 \to K_1 \to K_0 \to M \to 0
$$
를 가지며, 각 $K_i$는
% P02-E0372: frozen source's “for n variable” is nonidiomatic; Korean states that n varies.
$n$이 변하는 $R/I^n$의 사본들의 직합이다.
\end{lemma}

\begin{proof}
표준적인 전사
$$
\oplus_{m \in M} R/I^{n_m} \to M \to 0
$$
가 있다. 여기서 $n_m$은 $I^{n_m} \cdot m = 0$을 만족하는 가장 작은
양의 정수이다. 앞의 전사의 핵도 $I$-거듭제곱 꼬임 가군이다.
귀납적으로 계속하여 $M$의 원하는 분해를 구성한다.
\end{proof}

\begin{lemma}
\label{more-algebra-lemma-torsion-free}
$R$을 환이라 하고 $I$를 $R$의 아이디얼이라 하자. 모든 $R$-가군
$M$에 대해 $M[I^n] = \{m \in M \mid I^nm = 0\}$이라 두자.
$I$가 유한 생성이면 다음 조건들은 동치이다.
\begin{enumerate}
\item $M[I] = 0$.
\item 모든 $n \geq 1$에 대해 $M[I^n] = 0$이다.
\item $I = (f_1, \ldots, f_t)$이면 사상
$M \to \bigoplus M_{f_i}$는 단사이다.
\end{enumerate}
\end{lemma}

\begin{proof}
이는 Algebra, Lemma
\ref{algebra-lemma-when-injective-covering}에서 따른다.
\end{proof}

\begin{lemma}
\label{more-algebra-lemma-divide-by-torsion}
$R$을 환이라 하고 $I$를 $R$의 유한 생성 아이디얼이라 하자.
\begin{enumerate}
\item 모든 $R$-가군 $M$에 대해 $(M/M[I^\infty])[I] = 0$이다.
\item $I$-거듭제곱 꼬임 가군들의 확장은 $I$-거듭제곱 꼬임이다.
\end{enumerate}
\end{lemma}

\begin{proof}
$m \in M$이라 하자. $m$이 $(M/M[I^\infty])[I]$의 원소로 가면
$Im \subset M[I^\infty]$이다. $I = (f_1, \ldots, f_t)$라 쓰자.
그러면 $f_i m \in M[I^\infty]$이다. 즉 어떤 $n_i > 0$에 대해
$I^{n_i}f_i m = 0$이다. 따라서 $N = \sum n_i + 2$로 두면
$I^Nm = 0$이다. 그러므로 $m$은 $(M/M[I^\infty])$에서 영으로 간다.
이는 보조정리의 첫 번째 명제를 증명한다.

\medskip\noindent
두 번째 명제를 위해, $0 \to M' \to M \to M'' \to 0$을
$M'$과 $M''$이 모두 $I$-거듭제곱 꼬임 가군인 가군들의 짧은
완전열이라 하자. 그러면 $M[I^\infty] \supset M'$이고, 따라서
$M/M[I^\infty]$은 $M''$의 몫이므로 $I$-거듭제곱 꼬임이다.
이를 첫 번째 명제 및 보조정리 \ref{more-algebra-lemma-torsion-free}와 결합하면
그것이 영임을 알 수 있다.
\end{proof}

\begin{lemma}
\label{more-algebra-lemma-I-power-torsion}
$I$를 환 $R$의 유한 생성 아이디얼이라 하자.
$I$-거듭제곱 꼬임 가군들은 아벨 범주 $\text{Mod}_R$의 Serre
부분범주를 이룬다. Homology, Definition
\ref{homology-definition-serre-subcategory}를 보라.
\end{lemma}

\begin{proof}
% P02-E0373: frozen source uses a singular verb for the coordinated submodule/quotient subject.
$I$-거듭제곱 꼬임 가군의 부분가군과 몫가군은 모두
$I$-거듭제곱 꼬임임이 명백하다. 또한 보조정리
\ref{more-algebra-lemma-divide-by-torsion}에 의해 두 $I$-거듭제곱 꼬임 가군의
확장은 $I$-거듭제곱 꼬임이다.
% P02-E0374: frozen concluding sentence lacks a finite verb; Korean supplies the inference.
따라서 보조정리의 명제는 Homology, Lemma
\ref{homology-lemma-characterize-serre-subcategory}에서 따른다.
\end{proof}

\begin{lemma}
\label{more-algebra-lemma-local-cohomology-closed}
$R$을 환이라 하고 $I \subset R$을 유한 생성 아이디얼이라 하자.
부분범주 $I^\infty\text{-torsion} \subset \text{Mod}_R$은 닫힌
부분집합 $Z = V(I) \subset \Spec(R)$에만 의존한다. 실제로
$R$-가군 $M$이 $I$-거듭제곱 꼬임이라는 것은 그 지지가 $Z$에
포함된다는 것과 동치이다.
\end{lemma}

\begin{proof}
$M$을 $R$-가군이라 하고 $x \in M$이라 하자.
$x \in M[I^\infty]$이면 모든 $f \in I$에 대해 $x$는 $M_f$에서
영으로 간다. 따라서 모든 $\mathfrak p \not \supset I$에 대해
$x$는 $M_\mathfrak p$에서 영으로 간다. 반대로 모든
$\mathfrak p \not \supset I$에 대해 $x$가 $M_\mathfrak p$에서
영으로 간다면, 모든 $f \in I$에 대해 $x$는 $M_f$에서 영으로 간다.
따라서 $I = (f_1, \ldots, f_r)$이면 어떤 $n_i \geq 1$에 대해
$f_i^{n_i}x = 0$이다. 그러므로 $x \in M[I^{\sum n_i}]$이다.
따라서 $M[I^\infty]$은 사상
$M \to \prod_{\mathfrak p \not \in Z} M_\mathfrak p$의 핵이다.
보조정리의 두 번째 명제가 따르고, 그것은 첫 번째 명제를 함의한다.
\end{proof}

\noindent
다음 두 보조정리는 아마 다른 곳에 있어야 할 것이다.

\begin{lemma}
\label{more-algebra-lemma-bounded-above-mod-I}
$R$을 환, $I \subset R$을 아이디얼, $K$를 다음을 만족하는
$D(R)$의 대상이라 하자.
$i > 0$이면 $H^i(K \otimes_R^\mathbf{L} R/I) = 0$이다. 그러면
\begin{enumerate}
\item 모든 $n \geq 1$과 $i > 0$에 대해
$H^i(K \otimes_R^\mathbf{L} R/I^n) = 0$이다.
\item 모든 $I$-거듭제곱 꼬임 $R$-가군 $N$과 $i > 0$에 대해
$H^i(K \otimes_R^\mathbf{L} N) = 0$이다.
\item 코호몰로지 가군들 $H^i(M)$이 $I$-거듭제곱 꼬임이고
$i > 0$이면 $0$인 모든 $M \in D^b(R)$에 대해
$i > 0$이면 $H^i(K \otimes_R^\mathbf{L} M) = 0$이다.
\end{enumerate}
\end{lemma}

\begin{proof}
(2)의 증명. $N = \bigcup N[I^n]$이라 쓸 수 있다. 텐서곱은 여극한과
가환하므로
$K \otimes_R^\mathbf{L} N = \text{hocolim}_n K \otimes_R^\mathbf{L} N[I^n]$
이다(자세한 내용은 생략한다. 힌트: $K$를 K-평탄 복합체로 나타내고
직접 계산한다). 따라서 $N$이 $I^n$에 의해 소멸된다고 가정해도 된다.
$N$이 제곱이 영인 아이디얼인 $R$-대수
$R' = R/I^n \oplus N$을 생각하자. $D(R')$의 대상
$K' = K \otimes_R^\mathbf{L} R'$이 양의 차수에서 코호몰로지가
사라짐을 보이면 충분하다. $R$-대수들의 전사 $R' \to R/I$가 있고,
그 핵 $J$는 멱영이다.
% P02-E0375: frozen source claims n kernel factors vanish; the formal n-factor claim is preserved.
(그 핵의 임의의 $n$개 원소의 곱은 영이다.) 다음을 얻는다.
$$
K \otimes_R^\mathbf{L} R/I =
(K \otimes_R^\mathbf{L} R') \otimes_{R'}^\mathbf{L} R/I =
K' \otimes_{R'}^\mathbf{L} R/I
$$
가정에 의해 복합체 $K \otimes_R^\mathbf{L} R/I$는
$[-\infty, 0]$에 Tor 진폭을 갖는다.
% P02-E0376: frozen source's conclusion sentence lacks a finite verb; Korean supplies it.
따라서 결론은 보조정리
\ref{more-algebra-lemma-check-tor-dimension-modulo-nilpotent}에서 따른다.

\medskip\noindent
(1)은 (2)에서 자명하게 따른다. (3)은 (2), $M$의 영이 아닌
코호몰로지 가군들의 개수에 대한 귀납법, 그리고 Derived Categories,
Remark \ref{derived-remark-truncation-distinguished-triangle}의 절단
구별 삼각형들에서 따른다. 자세한 내용은 생략한다.
\end{proof}

\begin{lemma}
\label{more-algebra-lemma-derived-vanishing-mod-I}
$R$을 환, $I \subset R$을 아이디얼, $K$를 $D(R)$의 대상이라 하자.
$K \otimes_R^\mathbf{L} R/I = 0$이 $D(R)$에서 성립한다고 가정하자.
그러면
\begin{enumerate}
\item 모든 $n \geq 1$에 대해
$K \otimes_R^\mathbf{L} R/I^n = 0$이다.
\item 모든 $I$-거듭제곱 꼬임 $R$-가군 $N$에 대해
$K \otimes_R^\mathbf{L} N = 0$이다.
\item 코호몰로지 가군들이 $I$-거듭제곱 꼬임인 모든
$M \in D^b(R)$에 대해 $K \otimes_R^\mathbf{L} M = 0$이다.
\end{enumerate}
\end{lemma}

\begin{proof}
% P02-E0377: frozen proof is a sentence fragment; Korean supplies a complete sentence.
이는 보조정리 \ref{more-algebra-lemma-bounded-above-mod-I}의 귀결이다(일부 자세한
내용은 생략한다).
\end{proof}

\begin{lemma}
\label{more-algebra-lemma-torsion-completion}
\begin{reference}
\cite[Lemme~1]{Beauville-Laszlo}의 약간의 일반화이다.
\end{reference}
$R \to R'$을 환 사상이라 하자. $I \subset R$을 모든 $n > 0$에
대해 $R/I^n \to R'/I^nR'$이 동형인 아이디얼이라 하자. 모든
$I$-거듭제곱 꼬임 $R$-가군 $M$에 대해 사상
$M \to M \otimes_R R'$은 동형이다. 예를 들어 $I$가 유한 생성이고
$R^\wedge$가 $I$에 대한 $R$의 완비화이면
$M \cong M \otimes_R R^\wedge$이다.
\end{lemma}

\begin{proof}
$M$이 $I^n$에 의해 소멸되면
$$
M \otimes_R R' \cong
M \otimes_{R/I^n} R'/I^n R' \cong
M \otimes_{R/I^n} R/I^n \cong M.
$$
$M$이 $I$-거듭제곱 꼬임이면 $M = \bigcup M[I^n]$이다.
텐서곱은 직접 극한과 가환하므로(Algebra, Lemma
\ref{algebra-lemma-tensor-products-commute-with-limits}), 원하는
동형을 얻는다. 마지막 명제는 Algebra, Lemma
\ref{algebra-lemma-hathat-finitely-generated}에 의해 첫 번째 명제의
특수한 경우이다.
\end{proof}



\section{가군 범주들의 형식적 붙이기}
\label{more-algebra-section-formal-glueing}

\noindent
뇌터 스킴 $X$와 여집합이 $U$인 닫힌 부분스킴 $Z$를 고정하자.
우리의 목표는 $Z$를 따른 $X$의 형식적 완비화 위의 연접층과 $U$ 위의
연접층을, 교집합 위의 적절한 양립성과 함께 사용하여 $X$ 위의
연접층을 (유일하게) 구성하는 방법을 설명하는 것이다. 먼저 가환대수만
사용하여 이를 수행하고(이 절), 뒤에서 대수 공간의 문맥으로 설명한다
(Pushouts of Spaces, Section
\ref{spaces-pushouts-section-formal-glueing}).

\medskip\noindent
이 절의 내용 일부를 다루는 참고문헌은 다음과 같다.
\cite[Section 2]{ArtinII},
\cite[Appendix]{Ferrand-Raynaud},
\cite{Beauville-Laszlo},
\cite{MB}, 그리고
\cite[Section 4.6]{dJ-crystalline}.

\begin{lemma}
\label{more-algebra-lemma-characterize-flatness-on-torsion}
$\varphi : R \to S$를 환 사상, $I \subset R$을 아이디얼이라 하자.
다음 조건들은 동치이다.
\begin{enumerate}
\item $\varphi$는 평탄하고 $R/I \to S/IS$는 충실 평탄이다.
\item $\varphi$는 평탄하고 사상
$\Spec(S/IS) \to \Spec(R/I)$은 전사이다.
\item $\varphi$는 평탄하고 밑변환 함자
$M \mapsto M \otimes_R S$는 $I$에 의해 소멸되는 가군들 위에서
충실하다.
\item $\varphi$는 평탄하고 밑변환 함자
$M \mapsto M \otimes_R S$는 $I$-거듭제곱 꼬임 가군들 위에서
충실하다.
\end{enumerate}
\end{lemma}

\begin{proof}
$R \to S$가 평탄하면 모든 $n$에 대해
$R/I^n \to S/I^nS$는 평탄하다. Algebra, Lemma
\ref{algebra-lemma-flat-base-change}를 보라. 따라서 (1)과 (2)는
Algebra, Lemma \ref{algebra-lemma-ff-rings}에 의해 동치이다.
(1)과 (3)의 동치는 $I$-꼬임 $R$-가군을 $R/I$-가군과 동일시하고,
$I$에 의해 소멸되는 $R$-가군 $M$에 대해
$$
M \otimes_R S = M \otimes_{R/I} S/IS
$$
임과 Algebra, Lemma \ref{algebra-lemma-easy-ff}를 사용하면 따른다.
(4) $\Rightarrow$ (3)은 즉시 성립한다. (3)을 가정하자. 위에서
$R/I^n \to S/I^nS$가 평탄함을 보았고, 가정에 의해 이 사상은
스펙트럼 위에서 전사를 유도한다. 실제로
$\Spec(R/I^n) = \Spec(R/I)$이고 $S$에 대해서도 마찬가지이다.
따라서 밑변환 함자는 $I^n$에 의해 소멸되는 가군들 위에서 충실하다.
모든 $I$-거듭제곱 꼬임 가군 $M$은 $I^n$에 의해 소멸되는
$M_n$들의 합집합 $M = \bigcup M_n$이므로,
% P02-E0378: frozen source omits punctuation before its main clause; Korean separates it.
밑변환 함자는 모든 $I$-거듭제곱 꼬임 가군들의 범주 위에서 충실하다
(텐서곱은 여극한과 가환하기 때문이다).
\end{proof}

\begin{lemma}
\label{more-algebra-lemma-neighbourhood-isomorphism}
$(\varphi : R \to S, I)$가 보조정리
\ref{more-algebra-lemma-characterize-flatness-on-torsion}의 동치인 조건들을
만족한다고 가정하자. 다음 조건들은 동치이다.
\begin{enumerate}
\item 모든 $I$-거듭제곱 꼬임 가군 $M$에 대해 자연스러운 사상
$M \to M \otimes_R S$는 동형이다.
\item $R/I \to S/IS$는 동형이다.
\end{enumerate}
\end{lemma}

\begin{proof}
(1) $\Rightarrow$ (2)는 즉시 성립한다. (2)를 가정하자.
먼저 $M$이 $I$에 의해 소멸된다고 가정하자. 이 경우 $M$은
$R/I$-가군이다. 따라서 동형
$$
M \otimes_R S = M \otimes_{R/I} S/IS = M \otimes_{R/I} R/I = M
$$
을 얻고, 이는 주장을 증명한다. 다음으로 $I^n$에 의해 소멸되는
% P02-E0379: frozen source omits a relative pronoun; Korean states the restriction naturally.
모든 가군 $M$에 대해 $M \to M \otimes_R S$가 동형임을 귀납법으로
증명한다. 귀납 가정이 $n$에 대해 성립하고 $M$이 $I^{n + 1}$에 의해
소멸된다고 가정하자. 그러면 짧은 완전열
$$
0 \to I^nM \to M \to M/I^nM \to 0
$$
을 얻으며, $R \to S$가 평탄하므로 이는 짧은 완전열
$$
0 \to I^nM \otimes_R S \to M \otimes_R S \to M/I^nM \otimes_R S \to 0
$$
을 준다. $M' = I^nM$과 $M'' = M/I^nM$에 대해 표준 사상이
동형이라는 사실(귀납 가정)을 사용하면 $M$에 대해서도 같은 결론을
얻는다. 마지막으로 $M$을 일반적인 $I$-거듭제곱 꼬임 가군이라 하자.
그러면 $M = \bigcup M_n$이며 $M_n$은 $I^n$에 의해 소멸된다.
텐서곱이 여극한과 가환한다는 사실을 사용하여 결론을 얻는다.
\end{proof}

\begin{lemma}
\label{more-algebra-lemma-neighbourhood-equivalence}
$\varphi : R \to S$를 평탄한 환 사상, $I \subset R$을
$R/I \to S/IS$가 동형인 유한 생성 아이디얼이라 하자. 그러면
\begin{enumerate}
\item 모든 $R$-가군 $M$에 대해 사상 $M \to M \otimes_R S$는
$I$-거듭제곱 꼬임 부분가군들의 동형
$M[I^\infty] \to (M \otimes_R S)[(IS)^\infty]$을 유도한다.
\item $M$ 또는 $N$ 가운데 하나가 $I$-거듭제곱 꼬임이면 자연스러운
사상
$$
\Hom_R(M, N) \longrightarrow \Hom_S(M \otimes_R S, N \otimes_R S)
$$
은 동형이다.
\item 밑변환 함자 $M \mapsto M \otimes_R S$는
$I$-거듭제곱 꼬임 가군들과 $IS$-거듭제곱 꼬임 가군들 사이의
범주 동치를 정의한다.
\end{enumerate}
\end{lemma}

\begin{proof}
보조정리 \ref{more-algebra-lemma-characterize-flatness-on-torsion}과 보조정리
\ref{more-algebra-lemma-neighbourhood-isomorphism}의 동치인 조건들이 모두
만족됨에 유의하자. 앞으로는 이를 따로 언급하지 않고 사용하겠다.
먼저 (1)을 증명한다. $M$을 임의의 $R$-가군이라 하자.
$M' = M/M[I^\infty]$이라 두고 완전열
$$
0 \to M[I^\infty] \to M \to M' \to 0
$$
을 생각하자. $M[I^\infty] = M[I^\infty] \otimes_R S$이므로
$(M' \otimes_R S)[(IS)^\infty] = 0$임을 보이면 충분하다.
$I = (f_1, \ldots, f_t)$라 쓰자. 보조정리
\ref{more-algebra-lemma-divide-by-torsion}에 의해 $M'[I^\infty] = 0$이다.
따라서 모든 $n > 0$에 대해 사상
$$
M' \longrightarrow \bigoplus\nolimits_{i = 1, \ldots t} M',
\quad
x \longmapsto (f_1^n x, \ldots, f_t^n x)
$$
은 단사이다. $S$는 $R$ 위에서 평탄하므로 이에 대응하는 사상
$M' \otimes_R S \to \bigoplus_{i = 1, \ldots t} M' \otimes_R S$도
단사이다. 이는 원하는 대로 $(M' \otimes_R S)[I^n] = 0$임을 뜻한다.

\medskip\noindent
다음으로 (2)를 증명한다. $N$이 $I$-거듭제곱 꼬임이면
$N \otimes_R S = N$이고, (2)에 표시된 사상은 Algebra, Lemma
\ref{algebra-lemma-adjoint-tensor-restrict}에 의해 동형이다.
$M$이 $I$-거듭제곱 꼬임이면 모든 사상 $M \to N$의 상은
$N[I^\infty]$을 통해 인수분해되고, 모든 사상
$M \otimes_R S \to N \otimes_R S$의 상은
$(N \otimes_R S)[(IS)^\infty]$을 통해 인수분해된다. 따라서 이
경우에는 (1)에 의해 $N$을 $N[I^\infty]$으로 바꾸어도 된다.
결론은 방금 논의한 $N$이 $I$-거듭제곱 꼬임인 경우에서 따른다.

\medskip\noindent
다음으로 (3)을 증명한다. 이 함자는 (2)에 의해 충만충실하다.
본질적 전사성을 위해서는 모든 $IS$-거듭제곱 꼬임 $S$-가군 $N$에
대해 자연스러운 사상 $N \otimes_R S \to N$이 동형임을 유의하면
된다.
\end{proof}

\begin{lemma}
\label{more-algebra-lemma-map-identifies-koszul-and-cech-complexes}
$\varphi : R \to S$를 평탄한 환 사상, $I \subset R$을
$R/I \to S/IS$가 동형인 유한 생성 아이디얼이라 하자.
$V(f_1, \ldots, f_r) = V(I)$를 만족하는 모든
$f_1, \ldots, f_r \in R$에 대해 다음이 성립한다.
\begin{enumerate}
\item Koszul 복합체들의 사상
$K(R, f_1, \ldots, f_r) \to K(S, f_1, \ldots, f_r)$은
유사동형이다.
\item 확장 교대 {\v C}ech 복합체들의 사상
$$
\xymatrix{
R \to \prod_{i_0} R_{f_{i_0}} \to \prod_{i_0 < i_1} R_{f_{i_0}f_{i_1}}
\to \ldots \to R_{f_1\ldots f_r} \ar[d] \\
S \to \prod_{i_0} S_{f_{i_0}} \to \prod_{i_0 < i_1} S_{f_{i_0}f_{i_1}}
\to \ldots \to S_{f_1\ldots f_r}
}
$$
은 유사동형이다.
\end{enumerate}
\end{lemma}

\begin{proof}
두 경우 모두 $R$-가군들의 복합체 $K_\bullet$가 있고
% P02-E0380: frozen source omits the English compound hyphen; Korean is unaffected.
$K_\bullet \to K_\bullet \otimes_R S$가 유사동형임을 보이고자 한다.
보조정리 \ref{more-algebra-lemma-neighbourhood-isomorphism}과 $R \to S$의
평탄성에 의해, $K$의 모든 호몰로지 군이 $I$-거듭제곱 꼬임이면
충분하다. Koszul 복합체의 경우 이는 보조정리
\ref{more-algebra-lemma-homotopy-koszul}에 의해 성립하고, 확장 교대
{\v C}ech 복합체의 경우에는 보조정리
\ref{more-algebra-lemma-extended-alternating-torsion}에 의해 성립한다.
\end{proof}
\begin{lemma}
\label{more-algebra-lemma-naive-Koszul-complex}
$R$을 환, $I = (f_1, \ldots, f_n)$을 $R$의 유한 생성 아이디얼이라
하자. $M$을 원소들 $e_1, \ldots, e_n$으로 생성되고 관계식들
$f_i e_j - f_j e_i = 0$을 만족하는 $R$-가군이라 하자. 그러면
짧은 완전열
$$
0 \to K \to M \to I \to 0
$$
이 존재하며, $K$는 $I$에 의해 소멸된다.
\end{lemma}

\begin{proof}
이는 Koszul 복합체를 절단한 것일 뿐이다. 사상 $M \to I$는 규칙
$e_i \mapsto f_i$에 의해 결정된다. $m = \sum a_i e_i$가
$M \to I$의 핵에 속하면, 즉 $\sum a_i f_i = 0$이면
$f_j m = \sum f_j a_i e_i = (\sum f_i a_i) e_j = 0$이다.
\end{proof}

\begin{lemma}
\label{more-algebra-lemma-explicit-ext}
$R$을 환, $I = (f_1, \ldots, f_n)$을 $R$의 유한 생성 아이디얼이라
하자. 모든 $R$-가군 $N$에 대해
% P02-E0381: frozen denominator has mismatched tuple delimiters; formal span is preserved.
$$
H_1(N, f_\bullet) =
\frac{\{(x_1, \ldots, x_n) \in N^{\oplus n} \mid f_i x_j = f_j x_i \}}
{\{f_1x, \ldots, f_nx) \mid x \in N\}}
$$
이라 두자. 모든 $R$-가군 $N$에 대해 표준적인
% P02-E0382: frozen source calls this sequence short exact despite lacking terminal zero.
짧은 완전열
$$
0 \to \Ext_R(R/I, N) \to H_1(N, f_\bullet) \to \Hom_R(K, N)
$$
이 존재한다. 여기서 $K$는 보조정리
\ref{more-algebra-lemma-naive-Koszul-complex}에서와 같다.
\end{lemma}

\begin{proof}
위의 표기는 Homology, Section \ref{homology-section-extensions}에서
정의한 $\text{Mod}_R$ 안의 $\Ext$-군들을 뜻한다. 이들은
$\Ext_R(M, N)$으로 표기된다. 짧은 완전열
$0 \to I \to R \to R/I \to 0$에 딸린 Homology, Lemma
\ref{homology-lemma-six-term-sequence-ext}의 긴 완전열과
$\Ext_R(R, N) = 0$이라는 사실을 사용하면
$$
\Ext_R(R/I, N) =
\Coker(N \longrightarrow \Hom(I, N))
$$
을 얻는다. 보조정리 \ref{more-algebra-lemma-naive-Koszul-complex}의 짧은 완전열을
사용하면 복합체
$$
N \to \Hom(M, N) \to \Hom_R(K, N)
$$
를 얻으며, 그 가운데 호몰로지는 표준적으로
$\Ext_R(R/I, N)$과 동형이다. 이제 첫째 사상의 여핵은 표준적으로
$H_1(N, f_\bullet)$과 동형이므로 보조정리의 증명이 끝난다.
\end{proof}

\begin{lemma}
\label{more-algebra-lemma-koszul-homology-annihilated}
$R$을 환, $I = (f_1, \ldots, f_n)$을 $R$의 유한 생성 아이디얼이라
하자. 모든 $R$-가군 $N$에 대해 보조정리
\ref{more-algebra-lemma-explicit-ext}에서 정의한 Koszul 호몰로지 군
$H_1(N, f_\bullet)$은 $I$에 의해 소멸된다.
\end{lemma}

\begin{proof}
$(x_1, \ldots, x_n) \in N^{\oplus n}$이고
$f_i x_j = f_j x_i$라 하자. 그러면
% P02-E0383: frozen source's first coordinate pattern is erroneous; formal span is preserved.
$f_i(x_1, \ldots, x_n) = (f_i x_i, \ldots, f_i x_n)$이다.
다시 말해 $f_i$는 $H_1(N, f_\bullet)$을 소멸시킨다.
\end{proof}

\noindent
보조정리 \ref{more-algebra-lemma-neighbourhood-equivalence}의 충만충실성을
개선하여, 원천이 $I$-거듭제곱 꼬임인 $\Ext$-군도 $S$로 옮기는
것에 둔감함을 보일 수 있다. 다음 보조정리의 유도 버전은 Dualizing
Complexes, Lemma \ref{dualizing-lemma-neighbourhood-extensions}를 보라.

\begin{lemma}
\label{more-algebra-lemma-neighbourhood-extensions}
$\varphi : R \to S$를 평탄한 환 사상, $I \subset R$을
$R/I \to S/IS$가 동형인 유한 생성 아이디얼이라 하자. $M$, $N$을
$R$-가군이라 하고 $M$이 $I$-거듭제곱 꼬임이라고 가정하자.
% P02-E0384: frozen source uses the wrong English article before “short”.
짧은 완전열
$$
0 \to N \otimes_R S \to \tilde E \to M \otimes_R S \to 0
$$
이 주어지면 완전한 행들을 갖는 가환 도식
$$
\xymatrix{
0 \ar[r] &
N \ar[r] \ar[d] &
E \ar[r] \ar[d] &
M \ar[r] \ar[d] &
0 \\
0 \ar[r] &
N \otimes_R S \ar[r] &
\tilde E \ar[r] &
M \otimes_R S \ar[r] &
0
}
$$
이 존재한다.
\end{lemma}

\begin{proof}
$M$은 $I$-거듭제곱 꼬임이므로
$M \otimes_R S = M$이다. 보조정리
\ref{more-algebra-lemma-neighbourhood-isomorphism}을 보라. 앞으로 이 동일시를
따로 언급하지 않고 사용하겠다. $R \to S$가 평탄하므로 밑변환 함자는
완전하고, $\Ext$-군들의 함자적인 사상
$$
\Ext_R(M, N)
\longrightarrow
\Ext_S(M \otimes_R S, N \otimes_R S),
$$
을 얻는다. Homology, Lemma
\ref{homology-lemma-exact-functor-ext}를 보라. 보조정리의 주장은
$M$이 $I$-거듭제곱 꼬임일 때 이 사상이 전사라는 것이다. 실제로는
이 사상이 동형임을 보이겠다. 보조정리
\ref{more-algebra-lemma-I-power-torsion-presentation}에 의해, $R/I^n$ 꼴 가군들의
직합인 어떤 $M'$에서 오는 전사 $M' \to M$을 찾을 수 있다.
Homology, Lemma \ref{homology-lemma-six-term-sequence-ext}의 긴 완전열과
보조정리 \ref{more-algebra-lemma-neighbourhood-equivalence}를 사용하면 $M'$에
대해 보조정리를 증명하는 것으로 충분함을 알 수 있다. $\Ext$와
직합의 양립성(자세한 내용은 생략한다)을 사용하여 어떤 $n$에 대해
$M = R/I^n$인 경우로 환원한다.

\medskip\noindent
$f_1, \ldots, f_t$를 $I^n$의 생성원들이라 하자. 보조정리
\ref{more-algebra-lemma-explicit-ext}에 의해 가환 도식
% P02-E0385: frozen lower row omits every tensor base subscript; formal spans are preserved.
$$
\xymatrix{
0 \ar[r] &
\Ext_R(R/I^n, N) \ar[r] \ar[d] &
H_1(N, f_\bullet) \ar[r] \ar[d] &
\Hom_R(K, N) \ar[d] \\
0 \ar[r] &
\Ext_S(S/I^nS, N \otimes S) \ar[r] &
H_1(N \otimes S, f_\bullet) \ar[r] &
\Hom_S(K \otimes S, N \otimes S)
}
$$
을 얻는다. 그 행들은 완전하고, $K$는 보조정리
\ref{more-algebra-lemma-naive-Koszul-complex}에서와 같다. 따라서 오른쪽 두
수직 화살표가 동형임을 보이면 충분하다. $K$는 $I^n$에 의해
소멸되므로 보조정리 \ref{more-algebra-lemma-neighbourhood-equivalence}에 의해
$\Hom_R(K, N) = \Hom_S(K \otimes_R S, N \otimes_R S)$이다.
$R \to S$가 평탄하므로
$H_1(N, f_\bullet) \otimes_R S = H_1(N \otimes_R S, f_\bullet)$이다.
보조정리 \ref{more-algebra-lemma-koszul-homology-annihilated}에 의해
$H_1(N, f_\bullet)$은 $I^n$에 의해 소멸되므로, 보조정리
\ref{more-algebra-lemma-neighbourhood-isomorphism}에 의해
$H_1(N, f_\bullet) \otimes_R S = H_1(N, f_\bullet)$이다.
\end{proof}

\noindent
$R \to S$를 환 사상이라 하자. $f_1, \ldots, f_t \in R$이고
$I = (f_1, \ldots, f_t)$라 하자. 그러면 모든 $R$-가군 $M$에 대해
복합체
\begin{equation}
\label{more-algebra-equation-glueing-complex}
0 \to M \xrightarrow{\alpha}
M \otimes_R S \times \prod M_{f_i} \xrightarrow{\beta}
\prod (M \otimes_R S)_{f_i}
\times
% P02-E0387: frozen codomain leaves this product index unspecified; formal span is preserved.
\prod M_{f_if_j}
\end{equation}
을 정의할 수 있다. 여기서
$\alpha(m) = (m \otimes 1, m/1, \ldots, m/1)$이고
% P02-E0386: frozen formula omits the outer closing parenthesis; formal span is preserved.
$$
\beta(m', m_1, \ldots, m_t) =
((m'/1 - m_1 \otimes 1, \ldots, m'/1 - m_t \otimes 1),
(m_1 - m_2, \ldots, m_{t - 1} - m_t).
$$
이 복합체가 언제 완전한지 알고자 한다.

\begin{lemma}
\label{more-algebra-lemma-recover-module-from-glueing-data}
$\varphi : R \to S$를 평탄한 환 사상이라 하고
$I = (f_1, \ldots, f_t) \subset R$을 $R/I \to S/IS$가 동형인
아이디얼이라 하자. $M$을 $R$-가군이라 하자. 그러면 복합체
(\ref{more-algebra-equation-glueing-complex})는 완전하다.
\end{lemma}

\begin{proof}
첫 번째 증명. 보조정리
\ref{more-algebra-lemma-map-identifies-koszul-and-cech-complexes}의 확장 교대
{\v C}ech 복합체들의 유사동형
% P02-E1086: frozen imperative omits “by”; Korean introduces the notation naturally.
$\check{\mathcal{C}}_R \to \check{\mathcal{C}}_S$라 쓰자.
이 복합체들은 평탄한 항들을 갖는 유계 복합체들이므로
$M \otimes_R \check{\mathcal{C}}_R \to M \otimes_R \check{\mathcal{C}}_S$도
유사동형이다(보조정리
\ref{more-algebra-lemma-derived-tor-quasi-isomorphism}과
\ref{more-algebra-lemma-derived-tor-quasi-isomorphism-other-side}). 이제 복합체
(\ref{more-algebra-equation-glueing-complex})는 사상
$M \otimes_R \check{\mathcal{C}}_R \to M \otimes_R \check{\mathcal{C}}_S$의
원뿔을 절단한 것이므로 결론을 얻는다.

\medskip\noindent
두 번째 계산적 증명. $m \in M$이라 하자. $\alpha(m) = 0$이면
$m \in M[I^\infty]$이다. 보조정리 \ref{more-algebra-lemma-torsion-free}를 보라.
$I^n m = 0$인 $n$을 고르고 사상 $\varphi : R/I^n \to M$을
생각하자. $m \otimes 1 = 0$이면 $\varphi \otimes 1_S = 0$이고,
따라서 $\varphi = 0$이다(보조정리
\ref{more-algebra-lemma-neighbourhood-equivalence}를 보라). 그러므로 $m = 0$이다.
이렇게 하여 $\alpha$가 단사임을 알 수 있다.

\medskip\noindent
$(m', m'_1, \ldots, m'_t) \in \Ker(\beta)$라 하자. 어떤 $n > 0$과
$m_i \in M$에 대해 $m'_i = m_i/f_i^n$이라 쓰자.
% P02-E0388: frozen source omits the comma after its intervening clause.
$n$을 필요하면 더 크게 잡아서
$f_i^n m' = m_i \otimes 1$이 $M \otimes_R S$에서 성립하고
$f_j^nm_i - f_i^nm_j = 0$이 $M$에서 성립한다고 가정해도 된다.
특히 $(m_1, \ldots, m_t)$는
$H_1(M, (f_1^n, \ldots, f_t^n))$의 원소 $\xi$를 정의한다.
보조정리 \ref{more-algebra-lemma-koszul-homology-annihilated}에 의해
$H_1(M, (f_1^n, \ldots, f_t^n))$은 $I^{tn + 1}$에 의해 소멸되고,
$R \to S$는 평탄하므로
$$
H_1(M, (f_1^n, \ldots, f_t^n)) =
H_1(M, (f_1^n, \ldots, f_t^n)) \otimes_R S =
H_1(M \otimes_R S, (f_1^n, \ldots, f_t^n))
$$
이다. 이는 보조정리
% P02-E0389: frozen source omits sentence-ending punctuation after this reference.
\ref{more-algebra-lemma-neighbourhood-isomorphism}에 의한 것이다.
$m'$의 존재는 $\xi$가 마지막 군에서 영으로 감을 뜻한다. 즉 원소
$\xi$는 영이다. 따라서 어떤 $m \in M$이 존재하여
$m_i = f_i^n m$이다. 그러면 어떤
$m'' \in (M \otimes_R S)[(IS)^\infty]$에 대해
$(m', m'_1, \ldots, m'_t) - \alpha(m) = (m'', 0, \ldots, 0)$이다.
보조정리 \ref{more-algebra-lemma-neighbourhood-equivalence}에 의해
$m'' \in M[I^\infty]$임을 결론내리고 증명이 끝난다.
\end{proof}
\begin{remark}
\label{more-algebra-remark-glueing-data}
이 주석에서는 붙이기 자료의 범주를 정의한다. $R \to S$를 환
사상이라 하자. $f_1, \ldots, f_t \in R$이고
$I = (f_1, \ldots, f_t)$라 하자. 범주
$\text{Glue}(R \to S, f_1, \ldots, f_t)$를 다음과 같이 정의한다.
\begin{enumerate}
\item 대상은 계 $(M', M_i, \alpha_i, \alpha_{ij})$이다. 여기서
$M'$은 $S$-가군, $M_i$는 $R_{f_i}$-가군,
$\alpha_i : (M')_{f_i} \to M_i \otimes_R S$는 동형이고,
$\alpha_{ij} : (M_i)_{f_j} \to (M_j)_{f_i}$는 다음을 만족하는
동형들이다.
\begin{enumerate}
\item 사상 $(M')_{f_if_j} \to (M_j)_{f_i}$로서
$\alpha_{ij} \circ \alpha_i = \alpha_j$이다.
\item 사상 $(M_i)_{f_jf_k} \to (M_k)_{f_if_j}$로서
$\alpha_{jk} \circ \alpha_{ij} = \alpha_{ik}$이다(코사이클 조건).
\end{enumerate}
\item 사상
$(M', M_i, \alpha_i, \alpha_{ij}) \to (N', N_i, \beta_i, \beta_{ij})$은
주어진 사상들 $\alpha_i, \beta_i, \alpha_{ij}, \beta_{ij}$와 양립하는
사상들 $\varphi' : M' \to N'$과 $\varphi_i : M_i \to N_i$로
주어진다.
\end{enumerate}
표준적인 함자
$$
\text{Can} : \text{Mod}_R
\longrightarrow
\text{Glue}(R \to S, f_1, \ldots, f_t),
\quad
M \longmapsto (M \otimes_R S, M_{f_i}, \text{can}_i, \text{can}_{ij})
$$
가 있다. 여기서
$\text{can}_i : (M \otimes_R S)_{f_i} \to M_{f_i} \otimes_R S$와
$\text{can}_{ij} : (M_{f_i})_{f_j} \to (M_{f_j})_{f_i}$는 표준적인
동형들이다. 범주 $\text{Glue}(R \to S, f_1, \ldots, f_t)$의 모든
대상 $\mathbf{M} = (M', M_i, \alpha_i, \alpha_{ij})$에 대해
$$
H^0(\mathbf{M}) =
\{(m', m_i) \mid \alpha_i(m') = m_i \otimes 1, \alpha_{ij}(m_i) = m_j\}
$$
로 정의한다. 다시 말해 이는
(\ref{more-algebra-equation-glueing-complex})와 비슷한 완전열
$$
0 \to H^0(\mathbf{M}) \to
M' \times \prod M_i \to
\prod M'_{f_i}
\times
\prod (M_i)_{f_j}
$$
로 정의된다. $H^0(\mathbf{M})$을 $R$-가군으로 생각한다.
따라서 함자
$$
H^0 :
\text{Glue}(R \to S, f_1, \ldots, f_t)
\longrightarrow
\text{Mod}_R
$$
도 얻는다. 다음 목표는 함자들 $\text{Can}$과 $H^0$가 때때로 서로
준역임을 보이는 것이다.
\end{remark}

\begin{lemma}
\label{more-algebra-lemma-H0-Can-adjoint}
주석 \ref{more-algebra-remark-glueing-data}에서 함자
$H^0 : \text{Glue}(R \to S, f_1, \ldots, f_t) \to \text{Mod}_R$는
함자
$\text{Can} : \text{Mod}_R \to \text{Glue}(R \to S, f_1, \ldots, f_t)$의
오른쪽 수반이다.
\end{lemma}

\begin{proof}
$\mathbf{M} = (M', M_i, \alpha_i, \alpha_{ij})$을
$\text{Glue}(R \to S, f_1, \ldots, f_t)$의 대상이라 하자.
모든 $R$-가군 $N$에 대해 사상
$$
\Hom_{\text{Glue}(R \to S, f_1, \ldots, f_t)}(\text{Can}(N), \mathbf{M})
\to
\Hom_R(N, H^0(\mathbf{M}))
$$
이 있다. 이는 $\psi$를 $H^0(\psi)$와 명백한 사상
$N \to H^0(\text{Can}(N))$의 합성으로 보낸다. 구성에 의해 표시된
사상은 $N = R$일 때 동형이다(일반적으로
$R \to H^0(\text{Can}(R))$가 동형이 아니더라도 그렇다).
범주 $\text{Glue}(R \to S, f_1, \ldots, f_t)$는 직합과 여핵을
갖는다. 함자 $\text{Can}$은 직합 및 여핵과 가환한다. $N$을 자유
$R$-가군들 사이 사상의 여핵으로 쓰면, 이러한 관찰들로부터 표시된
사상이 전단사임을 알 수 있다. 자세한 내용은 생략한다.
\end{proof}

\begin{lemma}
\label{more-algebra-lemma-H0-inverse}
$\varphi : R \to S$를 평탄한 환 사상이라 하고
$I = (f_1, \ldots, f_t) \subset R$을 $R/I \to S/IS$가 동형인
아이디얼이라 하자. 그러면 함자 $H^0$는 주석
\ref{more-algebra-remark-glueing-data}의 함자 $\text{Can}$의 왼쪽 준역이다.
\end{lemma}

\begin{proof}
이는 보조정리 \ref{more-algebra-lemma-recover-module-from-glueing-data}를 다시
서술한 것이다.
\end{proof}

\begin{lemma}
\label{more-algebra-lemma-exact}
$\varphi : R \to S$를 평탄한 환 사상이라 하고
$I = (f_1, \ldots, f_t) \subset R$을 아이디얼이라 하자. 그러면
$\text{Glue}(R \to S, f_1, \ldots, f_t)$는 아벨 범주이고,
함자 $\text{Can}$은 완전하며 임의의 여극한과 가환한다.
\end{lemma}

\begin{proof}
범주 $\text{Glue}(R \to S, f_1, \ldots, f_t)$의 사상
$(\varphi', \varphi_i) :
(M', M_i, \alpha_i, \alpha_{ij})
\to
(N', N_i, \beta_i, \beta_{ij})$
이 주어지면, 그 핵은 존재하고 대상
$(\Ker(\varphi'), \Ker(\varphi_i), \alpha_i, \alpha_{ij})$와 같으며,
그 여핵은 존재하고 대상
$(\Coker(\varphi'), \Coker(\varphi_i), \beta_i, \beta_{ij})$와 같다.
이는 $R \to S$가 평탄하여 핵과 여핵을 취하는 것이
$- \otimes_R S$와 가환하기 때문에 성립한다. 자세한 내용은 생략한다.
완전성은 $R_{f_i}$와 $S$가 $R$-평탄이라는 사실에서 따르고,
여극한과의 가환성은 텐서곱이 여극한과 가환한다는 사실에서 따른다.
\end{proof}

\begin{lemma}
\label{more-algebra-lemma-equivalence-I-unit}
$\varphi : R \to S$를 평탄한 환 사상이라 하고
$(f_1, \ldots, f_t) = R$이라 하자. 그러면 $\text{Can}$과 $H^0$는
범주들의 서로 준역인 동치
$$
\text{Mod}_R = \text{Glue}(R \to S, f_1, \ldots, f_t)
$$
이다.
\end{lemma}

\begin{proof}
$\mathbf{M} = (M', M_i, \alpha_i, \alpha_{ij})$을
$\text{Glue}(R \to S, f_1, \ldots, f_t)$의 대상이라 하자.
Algebra, Lemma \ref{algebra-lemma-glue-modules}에 의해, 붙이기 자료
$\alpha_{ij}$를 복원하는 유일한 가군 $M$과 동형들
$M_{f_i} \to M_i$가 존재한다. 그러면 $M'$과
$M \otimes_R S$는 모두 $f_i$에서 국소화하면 가군들
$M_i \otimes_R S$를 복원하는 $S$-가군이다. 따라서 표준적인 동형
$M \otimes_R S \to M'$이 존재한다. 이는 $\mathbf{M}$이
$\text{Can}$의 본질적 상에 속함을 보인다. 이를 보조정리
\ref{more-algebra-lemma-H0-inverse}와 결합하면 보조정리가 따른다.
\end{proof}

\begin{lemma}
\label{more-algebra-lemma-base-change-glue}
$\varphi : R \to S$를 평탄한 환 사상이라 하고
% P02-E0390: frozen source has “and ideal” instead of “an ideal”; Korean states the intended phrase.
$I = (f_1, \ldots, f_t)$를 아이디얼이라 하자. $R \to R'$을 평탄한
환 사상이라 하고 $S' = S \otimes_R R'$이라 두자. 그러면 범주와
함자들의 가환 도식
$$
\xymatrix{
\text{Mod}_R \ar[r]_-{\text{Can}} \ar[d]_{-\otimes_R R'} &
\text{Glue}(R \to S, f_1, \ldots, f_t) \ar[r]_-{H^0} \ar[d]^{-\otimes_R R'} &
\text{Mod}_R \ar[d]^{-\otimes_R R'} \\
\text{Mod}_{R'} \ar[r]^-{\text{Can}} &
\text{Glue}(R' \to S', f_1, \ldots, f_t) \ar[r]^-{H^0} &
\text{Mod}_{R'}
}
$$
을 얻는다.
\end{lemma}

\begin{proof}
생략한다.
\end{proof}

\begin{proposition}
\label{more-algebra-proposition-equivalence}
$\varphi : R \to S$를 평탄한 환 사상이라 하고
$I = (f_1, \ldots, f_t) \subset R$을 $R/I \to S/IS$가 동형인
아이디얼이라 하자. 그러면 $\text{Can}$과 $H^0$는 범주들의 서로
준역인 동치
$$
\text{Mod}_R = \text{Glue}(R \to S, f_1, \ldots, f_t)
$$
이다.
\end{proposition}

\begin{proof}
$H^0 \circ \text{Can}$이 항등함자와 동형임은 이미 보았다.
보조정리 \ref{more-algebra-lemma-H0-inverse}를 보라. 대상
$\mathbf{M} = (M', M_i, \alpha_i, \alpha_{ij})$을
$\text{Glue}(R \to S, f_1, \ldots, f_t)$의 대상이라 하자.
자연스러운 사상
$$
\Psi :
(H^0(\mathbf{M}) \otimes_R S, H^0(\mathbf{M})_{f_i},
\text{can}_i, \text{can}_{ij})
\longrightarrow
(M', M_i, \alpha_i, \alpha_{ij}).
$$
을 얻는다. 실제로 정의에 의해 $H^0(\mathbf{M})$에는 서로 양립하는
$R$-가군 사상들 $H^0(\mathbf{M}) \to M'$과
$H^0(\mathbf{M}) \to M_i$가 주어진다. 이 사상이 동형임을 보여야
한다.

\medskip\noindent
지표 $i$를 하나 고르고 $R' = R_{f_i}$라 두자. 보조정리
\ref{more-algebra-lemma-base-change-glue}와 \ref{more-algebra-lemma-equivalence-I-unit}을
결합하면 $\Psi \otimes_R R'$이 동형임을 알 수 있다. 따라서
$\Psi$의 핵, 각각 여핵은 $(K, 0, 0, 0)$, 각각
$(Q, 0, 0, 0)$ 꼴의 계이다.
$H^0((K, 0, 0, 0)) = K$이고, $H^0$는 왼쪽 완전하며, 구성에 의해
$H^0(\Psi)$가 전단사임에 유의하자. 따라서 $K = 0$이다. 즉
$\Psi$의 핵은 영이다.

\medskip\noindent
위 결론에 의해 짧은 완전열
$$
0 \to H^0(\mathbf{M}) \otimes_R S \to M' \to Q \to 0
$$
을 얻고 $M_i = H^0(\mathbf{M})_{f_i}$이다. $Q$를
$Q = Q \otimes_R S$인 $I$-거듭제곱 꼬임 $R$-가군으로 생각할 수
있음에 유의하자. 보조정리 \ref{more-algebra-lemma-neighbourhood-extensions}에
의해 완전한 행들을 갖는 가환 도식
$$
\xymatrix{
0 \ar[r] &
H^0(\mathbf{M}) \ar[r] \ar[d] &
E \ar[r] \ar[d] &
Q \ar[r] \ar[d] &
0 \\
0 \ar[r] &
H^0(\mathbf{M}) \otimes_R S \ar[r] &
M' \ar[r] &
Q \ar[r] &
0
}
$$
이 존재한다. 이는 명백히 범주
$\text{Glue}(R \to S, f_1, \ldots, f_t)$ 안의 동형
$\text{Can}(E) \to (M', M_i, \alpha_i, \alpha_{ij})$을 결정하므로
결론을 얻는다. (물론 사후적으로 $Q = 0$이다.)
\end{proof}
\begin{lemma}
\label{more-algebra-lemma-application-formal-glueing}
$\varphi : R \to S$를 평탄한 환 사상, $I \subset R$을
$R/I \to S/IS$가 동형인 유한 생성 아이디얼이라 하자.
% P02-E0391: frozen statement reuses varphi for module maps after binding it as the ring map.
\begin{enumerate}
\item $R$-가군 $N$, $S$-가군 $M'$, 그리고 핵과 여핵이
$I$-거듭제곱 꼬임인 $S$-가군 사상
$\varphi : M' \to N \otimes_R S$가 주어지면, $R$-가군 사상
$\psi : M \to N$과 $\varphi$ 및 $\psi$와 양립하는 동형
$M \otimes_R S = M'$이 존재한다.
\item $R$-가군 $M$, $S$-가군 $N'$, 그리고 핵과 여핵이
$I$-거듭제곱 꼬임인 $S$-가군 사상
$\varphi : M \otimes_R S \to N'$이 주어지면, $R$-가군 사상
$\psi : M \to N$과 $\varphi$ 및 $\psi$와 양립하는 동형
$N \otimes_R S = N'$이 존재한다.
\end{enumerate}
두 경우 모두
$\Ker(\varphi) \cong \Ker(\psi)$이고
$\Coker(\varphi) \cong \Coker(\psi)$이다.
\end{lemma}

\begin{proof}
(1)의 증명. $I = (f_1, \ldots, f_t)$라 하자.
국소화 $\varphi_{f_i}$가 동형임은 명백하다. 따라서
% P02-E0392: frozen source writes these canonical-map names as bare can products.
$(M', N_{f_i}, \varphi_{f_i}, can_{ij})$는
$\text{Glue}(R \to S, f_1, \ldots, f_t)$의 대상이다. 주석
\ref{more-algebra-remark-glueing-data}를 보라. 명제
\ref{more-algebra-proposition-equivalence}에 의해, 동형들 $\varphi_{f_i}$ 및
$can_{ij}$와 양립하는 방식으로
$M' = M \otimes_R S$이고 $N_{f_i} = M_{f_i}$인 $R$-가군 $M$이
존재함을 결론내린다. $\text{Glue}(R \to S, f_1, \ldots, f_t)$의
사상
$$
(M \otimes_R S, M_{f_i}, can_i, can_{ij}) =
(M', N_{f_i}, \varphi_{f_i}, can_{ij})
\to
(N \otimes_R S, N_{f_i}, can_i, can_{ij})
$$
이 있으며, 그 첫 번째 성분에는 $\varphi$를 사용한다. 이는 명제
\ref{more-algebra-proposition-equivalence}의 범주 동치에 의해 $R$-가군 사상
$\psi : M \to N$에 대응한다. $M \to N$의 밑변환과 동형
$M' \cong M \otimes_R S$의 합성은 $\varphi$이다. 다시 말해
$M \to N$은 $\varphi$와 양립한다.

\medskip\noindent
(2)의 증명. 이는 위 논증의 쌍대일 뿐이다. 구체적으로 국소화
$\varphi_{f_i}$는 동형이다. 따라서
$(N', M_{f_i}, \varphi_{f_i}^{-1}, can_{ij})$는
$\text{Glue}(R \to S, f_1, \ldots, f_t)$의 대상이다. 주석
\ref{more-algebra-remark-glueing-data}를 보라. 명제
\ref{more-algebra-proposition-equivalence}에 의해, 동형들
$\varphi_{f_i}^{-1}$ 및 $can_{ij}$와 양립하는 방식으로
$N' = N \otimes_R S$이고 $N_{f_i} = M_{f_i}$인 $R$-가군 $N$이
존재함을 결론내린다. 다음은
$\text{Glue}(R \to S, f_1, \ldots, f_t)$의 사상이다.
% P02-E0393: frozen target tuple drops the inverse on varphi_{f_i}; formal span is preserved.
$$
(M \otimes_R S, M_{f_i}, can_i, can_{ij}) \to
(N', M_{f_i}, \varphi_{f_i}, can_{ij}) =
(N \otimes_R S, N_{f_i}, can_i, can_{ij})
$$
그 첫 번째 성분에는 $\varphi$를 사용한다. 이는 명제
\ref{more-algebra-proposition-equivalence}의 범주 동치에 의해 $R$-가군 사상
$\psi : M \to N$에 대응한다. $M \to N$의 밑변환과 동형
$N' \cong N \otimes_R S$의 합성은 $\varphi$이다. 다시 말해
$M \to N$은 $\varphi$와 양립한다.

\medskip\noindent
마지막 명제는 예를 들어 보조정리
\ref{more-algebra-lemma-neighbourhood-equivalence}에서 따른다.
\end{proof}

\noindent
이제 명제 \ref{more-algebra-proposition-equivalence}를 특수화하여 더 사용하기
쉬운 결과를 얻는다. 구체적으로 $I = (f)$가 주 아이디얼이면
$\text{Glue}(R \to S, f)$의 대상은 단순히 삼중항
$(M', M_1, \alpha_1)$이고, 확인할 코사이클 조건은 {\it 없다}!

\begin{theorem}
\label{more-algebra-theorem-formal-glueing}
$R$을 환, $f \in R$이라 하자. $\varphi : R \to S$를 동형
$R/fR \to S/fS$를 유도하는 평탄한 환 사상이라 하자. 그러면 함자
$$
\text{Mod}_R
\longrightarrow
\text{Mod}_S \times_{\text{Mod}_{S_f}} \text{Mod}_{R_f},
\quad
M
\longmapsto
(M \otimes_R S, M_f, \text{can})
$$
는 동치이다.
\end{theorem}

\begin{proof}
화살표의 오른쪽에 나타나는 범주는 다음 삼중항들의 범주이다.
$M'$은 $S$-가군이고,
% P02-E0394: frozen source uses the wrong English article before R_f-module.
$M_1$은 $R_f$-가군이며,
$\alpha_1 : M'_f \to M_1 \otimes_R S$는 $S_f$-동형인 삼중항
$(M', M_1, \alpha_1)$이다. Categories, Example
\ref{categories-example-2-fibre-product-categories}를 보라.
따라서 이 정리는 명제 \ref{more-algebra-proposition-equivalence}의 특수한
경우이다.
\end{proof}

\noindent
정리 \ref{more-algebra-theorem-formal-glueing}의 유용한 특수한 경우는 $R$이
뇌터 환이고 $S$가 원소 $f$에서의 $R$의 완비화일 때이다. 완비화
$R \to S$는 평탄하고, $M$이 유한 생성이면 함자
$M \mapsto M \otimes_R S$를 $f$-진 완비화 함자와 동일시할 수 있다.
이를 더 정확히 진술하기 위해 유한 생성 $R$-가군들의 범주를
$\text{Mod}^{fg}_R$로 표기하자.

\begin{proposition}
\label{more-algebra-proposition-formal-glueing}
$R$을 뇌터 환, $f \in R$을 원소라 하자. $R^\wedge$를 $R$의
$f$-진 완비화라 하자. 그러면 함자
$M \mapsto (M^\wedge, M_f, \text{can})$는 동치
$$
\text{Mod}^{fg}_R
\longrightarrow
\text{Mod}^{fg}_{R^\wedge}
\times_{\text{Mod}^{fg}_{(R^\wedge)_f}}
\text{Mod}^{fg}_{R_f}
$$
를 정의한다.
\end{proposition}

\begin{proof}
환 사상 $R \to R^\wedge$는 Algebra, Lemma
\ref{algebra-lemma-completion-flat}에 의해 평탄하다.
$R/fR = R^\wedge/fR^\wedge$임은 명백하다. Algebra, Lemma
\ref{algebra-lemma-completion-tensor}에 의해 유한 $R$-가군 $M$의
완비화는 $M \otimes_R R^\wedge$와 같다. 따라서 명제의 표시된
함자는 정리 \ref{more-algebra-theorem-formal-glueing}에 나타나는 함자와 같다.
특히 이는 충만충실하다. $(M_1, M_2, \psi)$를 우변의 대상이라 하자.
정리 \ref{more-algebra-theorem-formal-glueing}에 의해
$M_1 = M \otimes_R R^\wedge$이고 $M_2 = M_f$인 $R$-가군 $M$이
존재한다. $R \to R^\wedge \times R_f$는 충실 평탄하므로 Algebra,
Lemma \ref{algebra-lemma-cover}에 의해 $M$이 유한 생성임을, 즉
$M \in \text{Mod}^{fg}_R$임을 결론내린다. 이는 명제를 증명한다.
\end{proof}

\begin{remark}
\label{more-algebra-remark-formal-glueing-algebras}
명제 \ref{more-algebra-proposition-equivalence}, 정리
\ref{more-algebra-theorem-formal-glueing}, 그리고 명제
\ref{more-algebra-proposition-formal-glueing}의 동치들은 가군의 성질들을 보존한다.
예를 들어 $M$이
$\mathbf{M} = (M', M_i, \alpha_i, \alpha_{ij})$에 대응하면,
$M$이 $R$ 위에서 유한, 또는 유한 표시, 또는 평탄, 또는 사영이라는
것은 $M'$과 $M_i$가 각각 $S$와 $R_{f_i}$ 위에서 대응하는 성질을
갖는다는 것과 동치이다. 이는 $R \to S \times \prod R_{f_i}$가
충실 평탄이라는 사실과 충실 평탄 사상을 따른 이러한 성질들의
% P02-E0395: frozen source uses the verb “descend” where the noun “descent” is required.
하강 및 상승에서 따른다. Algebra, Lemma
\ref{algebra-lemma-descend-properties-modules}와 Theorem
\ref{algebra-theorem-ffdescent-projectivity}를 보라. 이 함자들은
양변의 $\otimes$-구조도 보존한다.
% P02-E0396: frozen source uses singular agreement after plural “These functors”.
따라서 이들은 대수의 범주와 같이 쌍
$(\text{Mod}_R, \otimes)$에서 만들어지는 여러 범주의 동치를
정의한다.
\end{remark}

\begin{remark}
\label{more-algebra-remark-topological-analogue}
% P02-E0397: frozen source uses nonstandard “differential manifold”; Korean uses the conventional term.
미분다양체 $X$와 그 안의 콤팩트한 닫힌 부분다양체 $Z$가 주어지고
그 여집합이 $U$라 하자. $X$ 위의 층을 지정하는 것은 $U$ 위의 층,
$X$ 안에서 $Z$의 지정되지 않은 관형 근방 $T$ 위의 층, 그리고
$T \cap U$를 따라 얻어지는 두 층 사이의 동형을 지정하는 것과 같다.
대수기하에서는 관형 근방 자체가 존재하지 않지만, 명제
\ref{more-algebra-proposition-equivalence}, 정리
\ref{more-algebra-theorem-formal-glueing}, 그리고 명제
\ref{more-algebra-proposition-formal-glueing} 같은 결과들을 사용하면 그 대신
형식적 근방으로 작업할 수 있다.
\end{remark}




\section{Beauville--Laszlo 정리}
\label{more-algebra-section-beauville-laszlo}

\noindent
$R$을 환, $f$를 $R$의 원소라 하자.
% P02-E0398: frozen source omits “by” in the construction introducing the notation.
$R^\wedge = \lim R/f^n R$을 $R$의 $f$-진 완비화라 하자.
이 절에서는 Beauville과 Laszlo의 정리를 논의하고 약간 일반화한다.
\cite{Beauville-Laszlo}를 보라. 이 정리는 적절한 조건 아래에서
$R^\wedge$ 위의 가군과 $R_f$ 위의 가군을
$(R^\wedge)_f$로 밑확대했을 때의 동형을 따라 ``서로 붙여서''
$R$ 위의 가군을 구성할 수 있다고 주장한다.

\medskip\noindent
\cite{Beauville-Laszlo}에서는 $f$가 $R$과 $M$ 모두에서
비영인자라고 가정한다. 사실은
$$
R[f^\infty] \longrightarrow R^\wedge[f^\infty]
$$
가 전단사이고
$$
M[f^\infty] \longrightarrow M \otimes_R R^\wedge
$$
가 단사라고만 가정하면 된다. 이러한 개선은 비아르키메데스
해석공간의 이론에 사용하기 위해 \cite[\S 1.3]{Kedlaya-Liu-I}에서
도입한 붙이기의 다른 접근법에서 일부 영감을 받았다.

\medskip\noindent
실제로 우리는 모든 $n > 0$에 대해 동형
$R/f^nR \to R'/f^nR'$과 동형
$R[f^\infty] \to R'[f^\infty]$을 유도하는 환 사상
$$
R \longrightarrow R'
$$
의 더 일반적인 설정에서 Beauville--Laszlo 정리를 증명할 것이다.
이 설정은 대역화에 더 적합하며, $R'$이 $R$의 완비화인 경우로부터
형식적으로 따르지 않는다. 예를 들어
$R[f^\infty] \to R'[f^\infty]$이 전단사라는 조건은
$R[f^\infty] \to R^\wedge[f^\infty]$이 전단사임을 함의하지 않는다.

\medskip\noindent
이 절에서 증명하는 Beauville과 Laszlo의 정리는 정리
\ref{more-algebra-theorem-formal-glueing}의 비평탄 판본으로 볼 수 있고,
$R' = R^\wedge$인 경우에는 명제
\ref{more-algebra-proposition-formal-glueing}의 비 Noether 판본으로 볼 수 있다.
평탄 하강과의 비교는 주석 \ref{more-algebra-remark-not-descent}을 보라.

\medskip\noindent
더 강한 결과들도 확립할 수 있지만(예를 들어 $M$에 제한을 두지 않고),
그러려면 유도범주의 수준에서 작업해야 한다. 자세한 내용은
\cite[\S 5]{Bhatt-Algebraize}를 보라.

\begin{lemma}
\label{more-algebra-lemma-same-quotients}
$R$을 환, $f \in R$이라 하자. 모든 양의 정수 $n$에 대해 사상
$R/f^nR \to R^\wedge/f^n R^\wedge$은 동형이다.
\end{lemma}

\begin{proof}
이는 Algebra, Lemma
\ref{algebra-lemma-hathat-finitely-generated}의 특수한 경우이다.
\end{proof}

\noindent
절 \ref{more-algebra-section-torsion}에서 도입한 표기를 사용하겠다. 즉,
% P02-E0399: frozen source omits “by” in the construction introducing M[f^n].
$R$-가군 $M$에 대해 $f^n$에 의해 소멸되는 $M$의 부분가군을
$M[f^n]$으로 나타내고
$$
M[f^\infty] = \bigcup\nolimits_{n = 1}^\infty M[f^n] = \Ker(M \to M_f).
$$
로 둔다. $M = M[f^\infty]$이면 $M$을 $f$-거듭제곱 꼬임 가군이라 한다.

\begin{lemma}
\label{more-algebra-lemma-BL-faithful}
$R$을 환, $f \in R$이라 하고, $n > 0$에 대해 동형
$R/f^nR \to R'/f^nR'$을 유도하는 환 사상 $R \to R'$을 잡자.
$R$-가군 $R' \oplus R_f$는 충실하다. 즉, 영이 아닌 모든
$R$-가군 $M$에 대해 가군 $M \otimes_R (R' \oplus R_f)$도 영이 아니다.
예를 들어 $M$이 영이 아니면
$M \otimes_R (R^\wedge \oplus R_f)$도 영이 아니다.
\end{lemma}

\noindent
그러나 사상 $M \to M \otimes_R (R' \oplus R_f)$가 반드시
단사인 것은 아니다. 예 \ref{more-algebra-example-not-glueing-pair}를 보라.

\begin{proof}
$M \neq 0$이지만 $M \otimes_R R_f = 0$이면 $M$은
$f$-거듭제곱 꼬임이다. 보조정리 \ref{more-algebra-lemma-torsion-completion}에 의해
$M \otimes_R R' \cong M \neq 0$을 얻는다. 마지막 명제는 보조정리
\ref{more-algebra-lemma-same-quotients}에 의해 첫 번째 명제의 특수한 경우이다.
\end{proof}

\begin{lemma}
\label{more-algebra-lemma-cover-spec}
$R$을 환, $f \in R$이라 하고, 동형 $R/fR \to R'/fR'$을
유도하는 환 사상 $R \to R'$을 잡자. 사상
$\Spec(R') \amalg \Spec(R_f) \to \Spec(R)$는 전사이다.
예를 들어 사상
$\Spec(R^\wedge) \amalg \Spec(R_f) \to \Spec(R)$는 전사이다.
\end{lemma}

\begin{proof}
$V(f) = \Spec(R/fR)$이고 $D(f) = \Spec(R_f)$인
$\Spec(R) = V(f) \amalg D(f)$임을 상기하자. Algebra, Section
\ref{algebra-section-spectrum-ring}, 특히 Lemmas
\ref{algebra-lemma-spec-closed}와
\ref{algebra-lemma-standard-open}을 보라. 따라서 환 사상
$R \to R/fR$이 $R'$을 통해 분해되므로 보조정리가 따른다.
마지막 명제는 보조정리 \ref{more-algebra-lemma-same-quotients}에 의해
첫 번째 명제의 특수한 경우이다.
\end{proof}

\begin{lemma}
\label{more-algebra-lemma-faithful-descent}
\begin{reference}
\cite[Lemme~2(a)]{Beauville-Laszlo}의 약간의 일반화이다.
\end{reference}
$R$을 환, $f \in R$이라 하고, $n > 0$에 대해 동형
$R/f^nR \to R'/f^nR'$을 유도하는 환 사상 $R \to R'$을 잡자.
$R$-가군 $M$이 유한 생성일 필요충분조건은
($R' \oplus R_f$)-가군 $M \otimes_R (R' \oplus R_f)$가
유한 생성인 것이다. 예를 들어
$M \otimes_R (R^\wedge \oplus R_f)$가
$R^\wedge \oplus R_f$ 위의 유한 생성 가군이면 $M$은 유한 생성
$R$-가군이다.
\end{lemma}

\begin{proof}
‘필요’ 방향은 명백하므로 $M \otimes_R (R' \oplus R_f)$가
유한 생성이라고 가정하자. 이 경우 각 생성원을 단순 텐서들의 합으로
쓰면 $M \otimes_R (R' \oplus R_f)$는 $M$의 원소들로 이루어진
유한 생성 집합을 갖는다. 다시 말해 어떤 유한 자유 $R$-가군에서
$M$으로 가는 사상이 있어서, 그 여핵은 $R' \oplus R_f$와 텐서하면
소멸한다.
% P02-E0400: frozen source says “finite generated” instead of “finitely generated.”
따라서 이 여핵에 보조정리 \ref{more-algebra-lemma-BL-faithful}을 적용하면
$M$이 유한 생성임을 알 수 있다. 마지막 명제는 보조정리
\ref{more-algebra-lemma-same-quotients}에 의해 첫 번째 명제의 특수한 경우이다.
\end{proof}

\begin{remark}
\label{more-algebra-remark-not-descent}
$R \to R_f$는 항상 평탄하지만, $R$이 Noether가 아니면
$R \to R^\wedge$은 보통 평탄하지 않다(Algebra, Lemma
\ref{algebra-lemma-completion-flat}과 Examples, Section
\ref{examples-section-nonflat}의 논의를 보라). 따라서 일반적으로
Descent, Section \ref{descent-section-descent-modules}에서 논의한
충실 평탄 하강을 사상 $R \to R^\wedge \oplus R_f$에 적용할 수 없다.
더욱이 Noether인 경우에도 이 사상에 대한 통상적인 하강 자료의 정의는
환 $R^\wedge \otimes_R R^\wedge$을 참조하는데, 이 절에서는 이 환을
다루지 않을 것이다.
\end{remark}

\noindent
{\bf 붙이기 쌍.} $R \to R'$을 $n > 0$에 대해 동형
$R/f^nR \to R'/f^nR'$을 유도하는 환 사상이라 하자. 오른쪽 사상이
두 표준적인 준동형의 차인 수열
\begin{equation}
\label{more-algebra-equation-BL-cech-re}
0 \to R \to R' \oplus R_f \to R'_f \to 0,
\end{equation}
을 생각하자. 이 수열이 완전이면 $(R \to R', f)$를
\emph{붙이기 쌍}이라고 한다. $(R \to R^\wedge, f)$가 붙이기 쌍이면
$(R, f)$를 \emph{붙이기 쌍}이라고 하겠다. 보조정리
\ref{more-algebra-lemma-same-quotients}에 의해 이는 잘 정의된다. 따라서
$(R, f)$가 붙이기 쌍일 필요충분조건은 수열
\begin{equation}
\label{more-algebra-equation-BL-cech}
0 \to R \to R^\wedge \oplus R_f \to (R^\wedge)_f \to 0,
\end{equation}
이 완전인 것이다.

\begin{lemma}
\label{more-algebra-lemma-same-f-torsion}
$R$을 환, $f \in R$이라 하고, $n > 0$에 대해 동형
$R/f^nR \to R'/f^nR'$을 유도하는 환 사상 $R \to R'$을 잡자.
수열 (\ref{more-algebra-equation-BL-cech-re})은 다음을 만족한다.
\begin{enumerate}
\item 오른쪽에서 완전하다.
\item 왼쪽에서 완전일 필요충분조건은
$R[f^\infty] \to R'[f^\infty]$이 단사인 것이다.
\item 중간에서 완전일 필요충분조건은
$R[f^\infty] \to R'[f^\infty]$이 전사인 것이다.
\end{enumerate}
특히 $(R \to R', f)$가 붙이기 쌍일 필요충분조건은
$R[f^\infty] \to R'[f^\infty]$이 전단사인 것이다. 예를 들어
$(R, f)$가 붙이기 쌍일 필요충분조건은
$R[f^\infty] \to R^\wedge[f^\infty]$이 전단사인 것이다.
\end{lemma}

\begin{proof}
$x \in R'_f$라 하자. $x' \in R'$인 $x = x'/f^n$으로 쓰고,
$x'' \in R$, $y \in R'$인 $x' = x'' + f^n y$로 쓰자. 그러면
$(y, -x''/f^n)$가 $x$로 사상됨을 알 수 있다. 따라서 (1)이 성립한다.

\medskip\noindent
(2)는 $\Ker(R \to R_f) = R[f^\infty]$이라는 사실에서 따른다.

\medskip\noindent
수열이 중간에서 완전하다면 $x \in R'[f^\infty]$인 꼴의 원소
$(x, 0)$은 첫 번째 화살표의 상에 속한다. 이는
$R[f^\infty] \to R'[f^\infty]$이 전사임을 함의한다. 역으로
$R[f^\infty] \to R'[f^\infty]$이 전사라고 가정하자. 중간에 있는
원소 $(x, y)$가 오른쪽에서 영으로 사상된다고 하자. 어떤
$y' \in R$에 대해 $y = y'/f^n$으로 쓰자. 그러면
$f^n x - y'$은 $R'$에서 $f$의 어떤 거듭제곱에 의해 소멸된다.
가정에 의해 어떤 $z \in R[f^\infty]$에 대해
$f^nx - y' = z$로 쓸 수 있다. 이때 $y'' = y' + z$라 두면
$y = y''/f^n$이고, 이 원소는 $R \to R/f^nR$의 핵에 속한다. 따라서
어떤 $y''' \in R$에 대해 $y$를 $y'''/1$로 나타낼 수 있다. 이제
$x - y'''$은 $R'[f^\infty]$에 속한다. 그러므로
$x - y''' = z' \in R[f^\infty]$이다. 따라서 원하는 대로
$(x, y'''/1) = (y''' + z', (y''' + z')/1)$이다.

\medskip\noindent
보조정리의 마지막 명제는 보조정리 \ref{more-algebra-lemma-same-quotients}에 의해
끝에서 두 번째 명제의 특수한 경우이다.
\end{proof}

\begin{remark}
\label{more-algebra-remark-BL-special-case}
$f$가 비영인자라고 가정하자. 그러면 Algebra, Lemma
\ref{algebra-lemma-completion-differ-by-torsion}에 의해 $f$는
$R^\wedge$에서도 비영인자이다. 따라서 $(R, f)$는 붙이기 쌍이다.
\end{remark}

\begin{remark}
\label{more-algebra-remark-noetherian-case}
$R \to R^\wedge$이 평탄하면, 각 양의 정수 $n$에 대해 수열
$0 \to R[f^n] \to R \to R$을 $R^\wedge$과 텐서하여 수열
$0 \to R[f^n] \otimes_R R^\wedge \to R^\wedge \to R^\wedge$을 얻는다.
이를 보조정리 \ref{more-algebra-lemma-torsion-completion}과 결합하면
$R[f^n] \to R^\wedge[f^n]$이 동형임을 알 수 있다. 따라서
$(R, f)$는 붙이기 쌍이다. 특히 $R$이 Noether이면 이것이 성립한다.
Algebra, Lemma \ref{algebra-lemma-completion-flat}을 보라.
\end{remark}

\begin{example}
\label{more-algebra-example-not-glueing-pair}
$k$를 체라 하고
$$
R = k[f, T_1, T_2, \ldots]/(fT_1, fT_2 - T_1, fT_3 - T_2, \ldots).
$$
로 두자. $R^\wedge$에서 $T_1$의 상이 $f$로 나누어지므로 사상
$R[f^\infty] \to R^\wedge[f^\infty]$은 단사가 아니다. 따라서
$(R, f)$는 붙이기 쌍이 아니다. 한편
$$
R = k[f, T_1, T_2, \ldots]/(fT_1, f^2T_2, \ldots),
$$
에 대해서는 원소 $T_1 + fT_2 + f^2 T_3 + \ldots$가 상에 속하지 않으므로
사상 $R[f^\infty] \to R^\wedge[f^\infty]$은 전사가 아니다.
특히 주석 \ref{more-algebra-remark-noetherian-case}에 의해 이 둘은 모두
$R \to R^\wedge$이 평탄하지 않은 예이다.
\end{example}

\noindent
{\bf 붙이기 가능한 가군.}
$R \to R'$을 $n > 0$에 대해 동형 $R/f^nR \to R'/f^nR'$을
유도하는 환 사상이라 하자. 임의의 $R$-가군 $M$에 대해
(\ref{more-algebra-equation-BL-cech-re})을 $M$과 텐서하여 수열
\begin{equation}
\label{more-algebra-equation-BL-cech-mod-re}
0 \to M \to (M \otimes_R R') \oplus (M \otimes_R R_f) \to
M \otimes_R R'_f \to 0
\end{equation}
을 얻는다. $M \otimes_R R_f = M_f$이고
$M \otimes_R R'_f = (M \otimes_R R')_f$임에 유의하라.
이 수열이 완전하면 $M$이
\emph{$(R \to R', f)$에 대해 붙이기 가능하다}고 한다.
% P02-E0401: frozen source leaves M without a grammatical antecedent in this definition.
$R$이 환이고 $f \in R$일 때, $R$-가군 $M$이
$(R \to R^\wedge, f)$에 대해 붙이기 가능하면 단순히
\emph{붙이기 가능하다}고 한다. 따라서 $M$이 붙이기 가능할
필요충분조건은 수열
\begin{equation}
\label{more-algebra-equation-BL-cech-mod}
0 \to M \to (M \otimes_R R^\wedge) \oplus (M \otimes_R R_f) \to
M \otimes_R (R^\wedge)_f \to 0
\end{equation}
이 완전인 것이다.

\begin{lemma}
\label{more-algebra-lemma-same-f-torsion-module}
$R$을 환, $f \in R$이라 하고, $n > 0$에 대해 동형
$R/f^nR \to R'/f^nR'$을 유도하는 환 사상 $R \to R'$을 잡자.
수열 (\ref{more-algebra-equation-BL-cech-mod-re})은 다음을 만족한다.
\begin{enumerate}
\item 오른쪽에서 완전하다.
\item 왼쪽에서 완전일 필요충분조건은
$M[f^\infty] \to (M \otimes_R R')[f^\infty]$이 단사인 것이다.
\item 중간에서 완전일 필요충분조건은
$M[f^\infty] \to (M \otimes_R R')[f^\infty]$이 전사인 것이다.
\end{enumerate}
따라서 $M$이 $(R \to R', f)$에 대해 붙이기 가능할 필요충분조건은
$M[f^\infty] \to (M \otimes_R R')[f^\infty]$이 전단사인 것이다.
$(R \to R', f)$가 붙이기 쌍이면, $M$이 $(R \to R', f)$에 대해
붙이기 가능할 필요충분조건은
$M[f^\infty] \to (M \otimes_R R')[f^\infty]$이 단사인 것이다.
예를 들어 $(R, f)$가 붙이기 쌍이면 $M$이 붙이기 가능할
필요충분조건은
$M[f^\infty] \to (M \otimes_R R^\wedge)[f^\infty]$이 단사인 것이다.
\end{lemma}

\begin{proof}
보조정리 \ref{more-algebra-lemma-same-f-torsion}의 결과들을 별도의 언급 없이
사용하겠다. 함자 $M \otimes_R -$는 오른쪽 완전이므로
(Algebra, Lemma \ref{algebra-lemma-tensor-product-exact}) (1)을 얻는다.

\medskip\noindent
$M \to M \otimes_R R_f = M_f$의 핵은 $M[f^\infty]$이다.
따라서 (2)가 따른다.

\medskip\noindent
수열이 중간에서 완전하다면
$x \in (M \otimes_R R')[f^\infty]$인 꼴의 원소 $(x, 0)$은
첫 번째 화살표의 상에 속한다. 이는
$M[f^\infty] \to (M \otimes_R R')[f^\infty]$이 전사임을 함의한다.
역으로 $M[f^\infty] \to (M \otimes_R R')[f^\infty]$이
전사라고 가정하자. 중간에 있는 원소 $(x, y)$가 오른쪽에서 영으로
사상된다고 하자. 어떤 $y' \in M$에 대해 $y = y'/f^n$으로 쓰자.
그러면 $f^n x - y'$은 $M \otimes_R R'$에서 $f$의 어떤 거듭제곱에
의해 소멸된다. 가정에 의해 어떤 $z \in M[f^\infty]$에 대해
$f^nx - y' = z$로 쓸 수 있다. 이때 $y'' = y' + z$라 두면
$y = y''/f^n$이고, 이 원소는 $M \to M/f^nM$의 핵에 속한다. 따라서
어떤 $y''' \in M$에 대해 $y$를 $y'''/1$로 나타낼 수 있다. 이제
$x - y'''$은 $(M \otimes_R R')[f^\infty]$에 속한다. 그러므로
$x - y''' = z' \in M[f^\infty]$이다. 따라서 원하는 대로
$(x, y'''/1) = (y''' + z', (y''' + z')/1)$이다.

\medskip\noindent
$(R \to R', f)$가 붙이기 쌍이면 임의의 $M$에 대해
(\ref{more-algebra-equation-BL-cech-mod-re})은 Algebra, Lemma
\ref{algebra-lemma-tensor-product-exact}에 의해 중간에서 완전하다. 이것으로
보조정리의 끝에서 두 번째 명제를 얻는다. 마지막 명제는 이 사실과
$(R, f)$가 붙이기 쌍일 필요충분조건이
$(R \to R^\wedge, f)$가 붙이기 쌍인 것이라는 사실에서 따른다.
\end{proof}

\begin{remark}
\label{more-algebra-remark-glueable}
$(R \to R', f)$를 붙이기 쌍, $M$을 $R$-가군이라 하자. 다음 관찰들은
$M$이 $(R \to R', f)$에 대해 붙이기 가능한지 판정하는 데 사용할 수 있다.
\begin{enumerate}
% P02-E1094: frozen item (1) invokes the completion pair without adding that it is a glueing pair.
\item 보조정리 \ref{more-algebra-lemma-same-f-torsion-module}에 의해 $M$이
$(R \to R^\wedge, f)$에 대해 붙이기 가능할 필요충분조건은
$M[f^\infty] \to M \otimes_R R^\wedge$이 단사인 것이다.
$M[f] \to M^\wedge$이 단사일 때, 즉
$M[f] \cap \bigcap_{n = 1}^\infty f^n M = 0$일 때 이는 성립한다.
\item $\text{Tor}_1^R(M, R'_f) = 0$이면 $M$은
$(R \to R', f)$에 대해 붙이기 가능하다
(Algebra, Lemma \ref{algebra-lemma-long-exact-sequence-tor}를 사용하라).
이는 $\text{Tor}_1^R(M, R')$가 $f$-거듭제곱 꼬임이라는 것과 동치이다.
특히 모든 평탄 $R$-가군은 $(R \to R', f)$에 대해 붙이기 가능하다.
\item $R \to R'$이 평탄하면 모든 $R$-가군에 대해
$\text{Tor}_1^R(M, R') = 0$이므로 모든 $R$-가군은
$(R \to R', f)$에 대해 붙이기 가능하다. 특히 $R$이 Noether이고
$R' = R^\wedge$일 때 이것이 성립한다. Algebra, Lemma
% P02-E0402: frozen source omits terminal punctuation after this reference.
\ref{algebra-lemma-completion-flat}
\end{enumerate}
\end{remark}

% P02-E0403: frozen English example title lacks a hyphen in “Non glueable”.
\begin{example}[붙이기 가능하지 않은 가군]
\label{more-algebra-example-not-glueable-module}
\begin{reference}
\cite[\S 4, Remarques]{Beauville-Laszlo}
\end{reference}
$R$을 $\mathbf{R}$ 위의 $0$에서의 $C^\infty$ 함수 싹들의 환이라 하자.
$f \in R$을 함수 $f(x) = x$라 하자. 그러면 $f$는 $R$에서
비영인자이므로 $(R, f)$는 붙이기 쌍이고
$R^\wedge \cong \mathbf{R}[[x]]$이다.
% P02-E0404: frozen source presents exp(-1/x^2) as a C-infinity germ without defining its value at zero.
$\varphi \in R$을 함수 $\varphi(x) = \text{exp}(-1/x^2)$라 하자.
그러면 $\varphi$의 Taylor 급수는 영이므로
$\varphi \in \Ker(R \to R^\wedge)$이다. $x \neq 0$에 대해
$\varphi(x) \neq 0$이므로 $\varphi$는 $R$에서 비영인자이다.
함수 $\varphi/f$ 역시 Taylor 급수가 영이므로
$M = R/\varphi R$에서의 그 상은 $M[f]$의 영이 아닌 원소이고,
$$
M \otimes_R R^\wedge = R^\wedge/\varphi R^\wedge = R^\wedge
$$
에서는 영으로 사상된다. 따라서 $M$은 붙이기 가능하지 않다.
\end{example}

\noindent
다음으로 몇몇 Tor 군을 계산한다.

\begin{lemma}
\label{more-algebra-lemma-first-tor}
$(R \to R', f)$를 붙이기 쌍이라 하자. 그러면 각 $n > 0$에 대해
$\text{Tor}^R_1(R', f^n R) = 0$이다.
\end{lemma}

\begin{proof}
완전열 $0 \to R[f^n] \to R \to f^n R \to 0$으로부터
$R[f^n] \otimes_R R' \to R'$이 단사임을 확인하면 충분함을 알 수 있다.
보조정리 \ref{more-algebra-lemma-torsion-completion}에 의해
$R[f^n] \otimes_R R' = R[f^n]$이고, $(R \to R', f)$가 붙이기
쌍이므로 보조정리 \ref{more-algebra-lemma-same-f-torsion}에 의해
$R[f^n] \to R'$이 단사이다.
\end{proof}

\begin{lemma}
\label{more-algebra-lemma-first-tor-total}
$(R \to R',f)$를 붙이기 쌍이라 하자. 그러면
$\text{Tor}^R_1(R', R/R[f^\infty]) = 0$이다.
\end{lemma}

\begin{proof}
$R/R[f^\infty] = \colim R/R[f^n] = \colim f^nR$이다. Tor 군의 형성은
여과 여극한과 가환하므로(Algebra, Lemma
\ref{algebra-lemma-tor-commutes-filtered-colimits}) 보조정리
\ref{more-algebra-lemma-first-tor}를 적용할 수 있다.
\end{proof}

\begin{lemma}
\label{more-algebra-lemma-BL3}
\begin{reference}
\cite[Lemme 3(a)]{Beauville-Laszlo}의 약간의 일반화이다.
\end{reference}
$(R \to R', f)$를 붙이기 쌍이라 하자. 모든 $R$-가군 $M$에 대해
$\text{Tor}^R_1(R', \Coker(M \to M_f)) = 0$이다.
\end{lemma}

\begin{proof}
$\overline{M} = M/M[f^\infty]$로 두자. 그러면
$\Coker(M \to M_f) \cong \Coker(\overline{M} \to \overline{M}_f)$이므로
$f$가 $M$에서 비영인자라고 가정해도 된다. 이 경우
$M \subset M_f$이고 $M_f/M = \colim M/f^nM$이며 전이사상들은
$f$에 의한 곱셈으로 주어진다. Tor 군의 형성은 여극한과 가환하므로
(Algebra, Lemma \ref{algebra-lemma-tor-commutes-filtered-colimits})
$\text{Tor}^R_1(R', M/f^n M) = 0$임을 보이면 충분하다.

\medskip\noindent
먼저 $M = R/R[f^\infty]$인 경우를 다루자. 보조정리
\ref{more-algebra-lemma-same-f-torsion}에 의해
$M \otimes_R R' = R'/R'[f^\infty]$이다. 짧은 완전열
$0 \to M \to M \to M/f^nM \to 0$으로부터 완전열
$$
\xymatrix{
\text{Tor}_1^R(R', R/R[f^\infty]) \ar[r] &
\text{Tor}_1^R(R', M/f^n M) \ar[r] &
R'/R'[f^\infty] \ar[dll]_{f^n} \\
R'/R'[f^\infty] \ar[r] &
(R'/R'[f^\infty])/(f^n
(R'/R'[f^\infty])) \ar[r] & 0
}
$$
을 얻는다. Algebra, Lemma
\ref{algebra-lemma-long-exact-sequence-tor}를 보라. 여기서 대각 화살표는
단사이다. 첫 번째 군 $\text{Tor}_1^R(R', R/R[f^\infty])$은 보조정리
\ref{more-algebra-lemma-first-tor-total}에 의해 영이므로, 원하는 대로
$\text{Tor}_1^R(R', M/f^nM) = 0$임을 얻는다.

\medskip\noindent
일반적인 경우를 다루기 위해 $F$가 자유 $R/R[f^\infty]$-가군인 전사
$F \to M$을 택하고 완전열
$$
0 \to N \to F/f^n F \to M/f^n M \to 0.
$$
을 만들자. 보조정리 \ref{more-algebra-lemma-torsion-completion}에 의해 이 수열은
$R'$과 텐서한 뒤에도 변하지 않으며, 따라서 완전하다. 앞 문단에 의해
$\text{Tor}^R_1(R', F/f^n F) = 0$이므로 원하는 대로
$\text{Tor}^R_1(R', M/f^n M) = 0$을 얻는다.
\end{proof}

\noindent
$(R \to R', f)$를 붙이기 쌍이라 하자. 이는 $n > 0$에 대해
$R/f^nR \to R'/f^nR'$이 동형이고 수열
$$
0 \to R \to R' \oplus R_f \to R_f' \to 0
$$
이 완전하다는 뜻이다. 주석 \ref{more-algebra-remark-glueing-data}에서 도입한 범주
$\text{Glue}(R \to R', f)$를 생각하자.
$\text{Glue}(R \to R', f)$의 대상 $(M', M_1, \alpha_1)$을
\emph{붙이기 자료}라고 부르겠다. 이는 $R'$-가군 $M'$,
$R_f$-가군 $M_1$, 그리고 동형
$\alpha_1 : (M')_f \to M_1 \otimes_R R'$로 이루어진다.
자명한 함자
$$
\text{Can} : \text{Mod}_R \longrightarrow \text{Glue}(R \to R', f),\quad
M \longmapsto (M \otimes_R R', M_f, \text{can}),
$$
가 있고, 반대 방향으로 함자
$$
H^0 : \text{Glue}(R \to R', f) \longrightarrow \text{Mod}_R,\quad
(M', M_1, \alpha_1) \longmapsto \Ker(M' \oplus M_1 \to (M')_f)
$$
가 있다. 정확한 정의는 주석 \ref{more-algebra-remark-glueing-data}를 보라.

\begin{theorem}
\label{more-algebra-theorem-BL-glueing}
\begin{reference}
\cite{Beauville-Laszlo}의 주정리의 약간의 일반화이다.
\end{reference}
$(R \to R',f)$를 붙이기 쌍이라 하자. 함자
$\text{Can} : \text{Mod}_R \longrightarrow \text{Glue}(R \to R', f)$는
% P02-E0407: frozen source lacks the relational construction needed between the two categories.
$(R \to R', f)$에 대해 붙이기 가능한 $R$-가군들의 범주와
붙이기 자료의 범주 $\text{Glue}(R \to R', f)$ 사이의 동치를 정한다.
\end{theorem}

\begin{proof}
$(M', M_1, \alpha_1)$을 붙이기 자료라 하자.
$M = H^0((M', M_1, \alpha_1))$이 $(R \to R', f)$에 대해
붙이기 가능한 가군이고
$(M', M_1, \alpha_1) \cong \text{Can}(M)$임을 보이겠다.

\medskip\noindent
먼저 함자 $H^0$의 정의에 사용되는 사상
$\text{d} : M' \oplus M_1 \to (M')_f$가 전사임을 확인하자.
$(x, y) \in M' \oplus M_1$은 $(M')_f$에서
$\text{d}(x, y) = x/1 - \alpha_1^{-1}(y \otimes 1)$로 사상된다.
$z \in (M')_f$이면 $g_i \in R'$, $y_i \in M_1$인
$\alpha_1(z) = \sum y_i \otimes g_i$로 쓸 수 있다. 모든 $i$에 대해
같은 $n$을 택할 수 있으므로, 어떤 $y'_i \in M'$과 $n \geq 0$에 대해
$\alpha_1^{-1}(y_i \otimes 1) = y_i'/f^n$으로 쓰자. 또한
$a_i \in R$, $b_i \in R'$인 $g_i = a_i + f^n b_i$로 쓰자. 그러면
$y = \sum a_i y_i \in M_1$과 $x = \sum b_i y'_i \in M'$에 대해
원하는 대로 $\text{d}(x, -y) = z$이다.

\medskip\noindent
$M = H^0((M', M_1, \alpha_1)) = \Ker(\text{d})$이므로
$R$-가군의 완전열
\begin{equation}
\label{more-algebra-equation-define-M}
0 \to M \to M' \oplus M_1 \to (M')_f \to 0.
\end{equation}
을 얻는다. 사상 $M \to M'$과 $M \to M_1$이 동형
$M \otimes_R R' \to M'$과 $M \otimes_R R_f \to M_1$을 유도함을
보이겠다. 그러면 원하는 대로 $M$이 $(R \to R', f)$에 대해 붙이기
가능하고 $\text{Can}(M) \cong (M', M_1, \alpha_1)$임이 따른다.

\medskip\noindent
$f$가 $M_1$에서 비영인자이므로
$M[f^\infty] \cong M'[f^\infty]$이다. 이로부터 완전열
\begin{equation}
\label{more-algebra-equation-exact-mod-torsion}
0 \to M/M[f^\infty] \to M_1 \to (M')_f/M' \to 0.
\end{equation}
을 얻는다. $R \to R_f$가 평탄하므로 이 완전열을 $R_f$와 텐서하여
$M \otimes_R R_f = (M/M[f^\infty]) \otimes_R R_f \to M_1$이
동형임을 얻는다.

\medskip\noindent
보조정리 \ref{more-algebra-lemma-BL3}에 의해
$\text{Tor}_1^R(R', \Coker(M' \to (M')_f)) = 0$이다. 따라서 수열
(\ref{more-algebra-equation-exact-mod-torsion})은 $R$ 위에서 $R'$과 텐서한 뒤에도
완전하다. $\alpha_1$과 보조정리 \ref{more-algebra-lemma-torsion-completion}을
사용하면, 그 결과로 얻는 완전열을 다음과 같이 쓸 수 있다.
\begin{equation}
\label{more-algebra-equation-mod-torsion-sequence}
0 \to (M/M[f^\infty]) \otimes_R R' \to
(M')_f \to (M')_f/M' \to 0
\end{equation}
이로부터 동형
$(M/M[f^\infty]) \otimes_R R' \cong M'/M'[f^\infty]$을 얻는다.
따라서 도표
$$
\xymatrix{
& M[f^\infty] \otimes_R R' \ar[r] \ar[d] &
M \otimes_R R'  \ar[r] \ar[d] &
(M/M[f^\infty]) \otimes_R R' \ar[r] \ar[d] & 0 \\
0 \ar[r] &
M'[f^\infty] \ar[r] &
M' \ar[r] &
M'/M'[f^\infty] \ar[r] & 0,
}
$$
에서 세 번째 수직 화살표는 동형이다. 두 행은 완전하고, 보조정리
\ref{more-algebra-lemma-torsion-completion} 및
$M[f^\infty] = M'[f^\infty]$에 의해 첫 번째 수직 화살표가 동형이므로,
오 보조정리에 의해 $M \otimes_R R' \to M'$은 동형이다.

\medskip\noindent
위 논의로 $\text{Can}$은 본질적 전사이고 함자 $H^0$는 붙이기 가능한
가군들의 범주로 사상됨을 알 수 있다. 붙이기 가능한 가군에 대해
(\ref{more-algebra-equation-BL-cech-mod-re})이 완전하므로, 붙이기 가능한 가군들의
범주 위에서 $H^0 \circ \text{Can} = \text{id}$이다. 보조정리
\ref{more-algebra-lemma-H0-Can-adjoint}와 Categories, Lemma
\ref{categories-lemma-adjoint-fully-faithful}를 결합하면 이는
$\text{Can}$이 충만충실임을 함의한다. 이로써 증명이 끝난다.
\end{proof}

\begin{remark}
\label{more-algebra-remark-what-you-get-for-general-modules}
$(R \to R', f)$를 붙이기 쌍이라 하자. $M$을 $(R \to R', f)$에 대해
붙이기 가능하다고 반드시 가정하지 않는 $R$-가군이라 하자.
$M' = M \otimes_R R'$과 $M_1 = M_f$로 두면 붙이기 자료
$\text{Can}(M) = (M', M_1, \text{can})$을 얻는다. 그러면
$\tilde M = H^0(M', M_1, \text{can})$은 $(R \to R', f)$에 대해
붙이기 가능한 $R$-가군이고, 표준적인 사상 $M \to \tilde M$은 동형
$M \otimes_R R' \to \tilde M \otimes_R R'$과
$M_f \to \tilde M_f$를 준다. 정리 \ref{more-algebra-theorem-BL-glueing}을 보라.
수열
$$
M \to (M \otimes_R R' )\oplus M_f \to M \otimes_R (R')_f  \to 0
$$
과
$$
0 \to \tilde M \to (\tilde M \otimes_R R') \oplus \tilde M_f \to
\tilde M \otimes_R  (R')_f \to 0
$$
의 완전성으로부터 사상 $M \to \tilde M$이 전사임을 얻는다.
\end{remark}

\noindent
붙이기 쌍 $(R \to R', f)$ 위의 평탄 $R$-가군은 붙이기 가능함을
상기하자(주석 \ref{more-algebra-remark-glueable}). 따라서 다음 보조정리는 정리
\ref{more-algebra-theorem-BL-glueing}이 평탄 $R$-가군들의 범주와,
$M'$과 $M_1$이 각각 $R'$과 $R_f$ 위에서 평탄한 붙이기 자료
$(M', M_1, \alpha_1)$의 범주 사이의 동치를 정함을 보여준다.

\begin{lemma}
\label{more-algebra-lemma-BL-flat}
$(R \to R', f)$를 붙이기 쌍이라 하자. $M$을 $(R \to R', f)$에 대해
붙이기 가능하다고 반드시 가정하지 않는 $R$-가군이라 하자. 그러면
$M$이 $R$ 위에서 평탄일 필요충분조건은 $M \otimes_R R'$이
$R'$ 위에서 평탄이고 $M_f$가 $R_f$ 위에서 평탄인 것이다.
\end{lemma}

\begin{proof}
한 방향은 Algebra, Lemma
\ref{algebra-lemma-flat-base-change}에서 따른다. 다른 방향을 위해
$M \otimes_R R'$이 $R'$ 위에서 평탄이고 $M_f$가 $R_f$ 위에서
평탄이라고 가정하자. $\tilde M$을 주석
\ref{more-algebra-remark-what-you-get-for-general-modules}에서와 같이 잡자.
$\tilde M$이 $R$ 위에서 평탄이면 Algebra, Lemma
\ref{algebra-lemma-flat-tor-zero}을 짧은 완전열
$0 \to \Ker(M \to \tilde M) \to M \to \tilde M \to 0$에 적용하여
$\Ker(M \to \tilde M) \otimes_R (R' \oplus R_f)$이 영임을 얻는다.
따라서 보조정리 \ref{more-algebra-lemma-BL-faithful}에 의해 $M = \tilde M$이고
결론을 얻는다. 달리 말하면 $M$을 $\tilde M$으로 바꾸어
$M$이 $(R \to R', f)$에 대해 붙이기 가능하다고 가정해도 된다.
$N$을 두 번째 $R$-가군이라 하자. Algebra, Lemma
\ref{algebra-lemma-characterize-flat}에 의해
$\text{Tor}_1^R(M, N) = 0$임을 증명하면 충분하다.

\medskip\noindent
짧은 완전열 $0 \to R \to R' \oplus R_f \to (R')_f \to 0$과
$N$에 연관된 Tor 긴 완전열은 완전열
$$
0 \to \text{Tor}_1^R(R', N) \to \text{Tor}_1^R((R')_f, N)
$$
과 $i \geq 2$에 대한 동형
$\text{Tor}_i^R(R', N) = \text{Tor}_i^R((R')_f, N)$을 준다.
$\text{Tor}_i^R((R')_f, N) = \text{Tor}_i^R(R', N)_f$이므로
$f$는 $\text{Tor}_1^R(R', N)$에서 비영인자이고,
$i \geq 2$에 대해 $\text{Tor}_i^R(R', N)$에서 가역이다.
$M \otimes_R R'$이 $R'$ 위에서 평탄하므로, 예
\ref{more-algebra-example-tor-change-rings}의 스펙트럼 열에 의해
$$
\text{Tor}_i^R(M \otimes_R R', N) =
(M \otimes_R R') \otimes_{R'} \text{Tor}_i^R(R', N)
$$
이다. $M \otimes_R R'$을 유한 자유 $R'$-가군들의 여과 여극한으로 쓰면
(Algebra, Theorem \ref{algebra-theorem-lazard}), $f$가
$\text{Tor}_1^R(M \otimes_R R', N)$에서 비영인자이고
% P02-E0405: frozen source drops the necessary qualifier i >= 2 from the invertibility clause.
$\text{Tor}_i^R(M \otimes_R R', N)$에서 가역임을 얻는다. 다음으로
$M$이 붙이기 가능하다는 사실에서 오는 완전열
$0 \to M \to M \otimes_R R' \oplus M_f \to M \otimes_R (R')_f \to 0$과
그에 연관된 $\text{Tor}$ 긴 완전열을 생각하자. 필요한 부분은
$$
\xymatrix{
\text{Tor}_1^R(M, N) \ar[r] &
\text{Tor}_1^R(M \otimes_R R', N) \ar[r] &
\text{Tor}_1^R(M \otimes_R (R')_f, N) \\
& \text{Tor}_2^R(M \otimes_R R', N) \ar[r] &
\text{Tor}_2^R(M \otimes_R (R')_f, N) \ar[llu]
}
$$
이다. 위에서 논의한
% P02-E0406: frozen source says “the action on f on” instead of “the action of f on”.
$\text{Tor}_i^R(M \otimes_R R', N)$ 위에서의 $f$의 작용으로부터
$\text{Tor}_1^R(M, N) = 0$임을 얻는다.
\end{proof}

\noindent
다음 보조정리의 결과 가운데 “유한 생성”인 경우는 이미 보조정리
\ref{more-algebra-lemma-faithful-descent}에서 보았음에 유의하라.

\begin{lemma}
\label{more-algebra-lemma-BL-properties}
$(R \to R', f)$를 붙이기 쌍이라 하자. $M$을 $(R \to R', f)$에 대해
붙이기 가능하다고 반드시 가정하지 않는 $R$-가군이라 하자. 그러면
$M$이 유한 사영 $R$-가군일 필요충분조건은 $M \otimes_R R'$이
$R'$ 위에서 유한 사영이고 $M_f$가 $R_f$ 위에서 유한 사영인 것이다.
\end{lemma}

\begin{proof}
$M \otimes_R R'$이 $R'$ 위의 유한 사영 가군이고 $M_f$가 $R_f$ 위의
유한 사영 가군이라고 가정하자. 우리의 목표는 $M$이 $R$ 위에서
유한 사영임을 증명하는 것이다. Algebra, Lemma
\ref{algebra-lemma-finite-projective}를 별도의 언급 없이 사용하겠다.
보조정리 \ref{more-algebra-lemma-BL-flat}에 의해 $M$은 평탄하다. 보조정리
\ref{more-algebra-lemma-faithful-descent}에 의해 $M$은 유한이다. 짧은 완전열
$0 \to K \to R^{\oplus n} \to M \to 0$을 택하자. 유한 사영 가군은
유한 표시이고, $R'$과 텐서한 뒤에도(Algebra, Lemma
\ref{algebra-lemma-flat-tor-zero}에 의해), 그리고 $R_f$와 텐서한
뒤에도 이 수열은 완전하므로, $K \otimes_R R'$과 $K_f$가 유한
가군임을 얻는다. 위 보조정리를 사용하면 $K$는 유한 생성이다.
따라서 $M$은 유한 표시이고, 그러므로 유한 사영이다.
\end{proof}

\begin{remark}
\label{more-algebra-remark-compare-BL}
\cite{Beauville-Laszlo}에서는 $f$가 $R$에서 비영인자이고
$R' = R^\wedge$이라고 가정한다. 이는 보조정리
\ref{more-algebra-lemma-same-f-torsion}에 의해 붙이기 쌍을 준다. 이 설정에서도
정리 \ref{more-algebra-theorem-BL-glueing}은 새로운 것을 말한다. 즉,
\cite{Beauville-Laszlo}의 결과들은 $f$가 비영인자인 가군에만 적용된다
(그러한 가군은 우리의 의미에서 붙이기 가능하다. 보조정리
\ref{more-algebra-lemma-same-f-torsion-module}을 보라). 보조정리
\ref{more-algebra-lemma-BL-properties}도 \cite{Beauville-Laszlo}의 결과를 약간
확장한다. $M$이 영이 아닌 $f$-거듭제곱 꼬임을 갖도록 허용할 수 있을
뿐 아니라, 그것이 붙이기 가능하다고 요구하지도 않는다.
\end{remark}

\section{유도 완비화}
\label{more-algebra-section-derived-completion}

\noindent
이 절의 내용에 대한 참고문헌으로는
\cite{Dwyer-Greenlees}, \cite{Greenlees-May}, \cite{PSY}, \cite{dag12}
(특히 4장)이 있다. 우리의 설명은 \cite{BS}를 따른다. 꼬임 가군에
대한 이 절의 유사물(또는 ``쌍대'')은 Dualizing Complexes, Section
\ref{dualizing-section-local-cohomology}이다. 꼬임 코호몰로지를 갖는
복합체들의 유도범주와 유도 완비 복합체들 사이의 관계는
Dualizing Complexes, Section
\ref{dualizing-section-torsion-and-complete}에서 찾을 수 있다.

\medskip\noindent
$K \in D(A)$, $f \in A$라 하자.
% P02-E0408: frozen source omits “by” when introducing T(K, f).
$D(A)$에서 계
$$
\ldots \to K \xrightarrow{f} K \xrightarrow{f} K
$$
의 한 유도 극한을 $T(K, f)$로 나타낸다.

\begin{lemma}
\label{more-algebra-lemma-hom-from-Af}
$A$를 환이라 하고, $f \in A$, $K \in D(A)$라 하자. 다음 조건들은
동치이다.
\begin{enumerate}
\item 모든 $n$에 대해 $\Ext^n_A(A_f, K) = 0$이다.
\item $D(A_f)$의 모든 $E$에 대해 $\Hom_{D(A)}(E, K) = 0$이다.
\item $T(K, f) = 0$이다.
\item 모든 $p \in \mathbf{Z}$에 대해 $T(H^p(K), f) = 0$이다.
\item 모든 $p \in \mathbf{Z}$에 대해
$\Hom_A(A_f, H^p(K)) = 0$이고 $\Ext^1_A(A_f, H^p(K)) = 0$이다.
\item $R\Hom_A(A_f, K) = 0$이다.
\item 사상 $\prod_{n \geq 0} K \to \prod_{n \geq 0} K$,
$(x_0, x_1, \ldots) \mapsto (x_0 - fx_1, x_1 - fx_2, \ldots)$는
$D(A)$에서 동형이다.
% P02-E0409: frozen source contains the unfinished editorial placeholder “add more here.”
\item 여기에 더 추가하라.
\end{enumerate}
\end{lemma}

\begin{proof}
(2)가 (1)을 함의하고 (1)과 (6)이 동치임은 명백하다. (1)을
가정하자. $I^\bullet$을 $K$를 나타내는 $A$-가군의 K-단사 복합체라
하자. 조건 (1)은 $\Hom_A(A_f, I^\bullet)$이 비순환임을 뜻한다.
$M^\bullet$을 $E$를 나타내는 $A_f$-가군의 복합체라 하자. 그러면
$$
\Hom_{D(A)}(E, K) =
\Hom_{K(A)}(M^\bullet, I^\bullet) =
\Hom_{K(A_f)}(M^\bullet, \Hom_A(A_f, I^\bullet))
$$
이다. Algebra, Lemma \ref{algebra-lemma-adjoint-hom-restrict}를 보라.
% P02-E0410: frozen source lacks punctuation between the introductory reference clause and its main clause.
보조정리 \ref{more-algebra-lemma-hom-K-injective}에 의해
$\Hom_A(A_f, I^\bullet)$은 $A_f$-가군의 K-단사 복합체이다. 따라서
이 복합체가 비순환이라는 사실은 그것이 영과 호모토피 동치임을 함의한다
(Derived Categories, Lemma \ref{derived-lemma-K-injective}). 이로써
(2)를 얻는다.

\medskip\noindent
$A$-가군 $A_f$의 자유 분해는
$$
0 \to \bigoplus\nolimits_{n \in \mathbf{N}} A \to
\bigoplus\nolimits_{n \in \mathbf{N}} A
\to A_f \to 0
$$
으로 주어진다. 여기서 첫 번째 사상은
% P02-E0411: frozen source has an extraneous definite article before this tuple.
$(a_0, a_1, a_2, \ldots)$를
$(a_0, a_1 - fa_0, a_2 - fa_1, \ldots)$로 보내고, 두 번째 사상은
$(a_0, a_1, a_2, \ldots)$를 $a_0 + a_1/f + a_2/f^2 + \ldots$로 보낸다.
$\Hom_A(-, I^\bullet)$을 적용하면
$$
0 \to \Hom_A(A_f, I^\bullet) \to \prod I^\bullet \to \prod I^\bullet \to 0
$$
을 얻는다. $\prod I^\bullet$이 $\prod_{n \geq 0} K$를 나타내므로
이는 (1)과 (7)의 동치를 증명한다. 한편 Derived Categories, Section
\ref{derived-section-derived-limit}의 유도 극한 구성에 의해, 표시된
완전열은 대상 $T(K, f)$가 $D(A)$에서 $R\Hom_A(A_f, K)$의 대표임을
보여준다.
% P02-E0412: frozen source's conclusion is a sentence fragment without a finite verb.
따라서 (1)과 (3)은 동치이다.

\medskip\noindent
스펙트럼 열
$$
E_2^{p, q} = \Ext^p_A(A_f, H^q(K)) \Rightarrow
\Ext^{p + q}_A(A_f, K)
$$
이 있다. 식 (\ref{more-algebra-equation-first-ss-ext})을 보라\footnote{$K$가
아래로 유계가 아니면 주어진 참고문헌을 사용할 수 없다. 그 경우에는
% P02-E0413: frozen footnote says “we chooses a complex.”
$K$를 나타내는 복합체 $K^\bullet$과 자유 두 항 분해
$0 \to F^{-1} \to F^0 \to A_f \to 0$을 택한다(위를 보라). 그러면
항이 $\Hom_A(F^{-q}, K^p)$인 이중 복합체를 얻으며, 그 전체화는
$R\Hom(A_f, K)$를 계산한다. 이 이중 복합체에 연관된 두 번째
스펙트럼 열(Homology, Section
\ref{homology-section-double-complex})이 원하는 것을 준다. 세부사항은
생략한다.}. $A_f$는 사영 $A$-가군들에 의한 길이 $1$의 분해를
가지므로(위를 보라) 이 스펙트럼 열은 $E_2$에서 퇴화한다. 따라서
$E_2$-면에는 영이 아닌 열이 두 개뿐이다. 그러므로 짧은 완전열
$$
0 \to \Ext^1_A(A_f, H^{p - 1}(K)) \to
\Ext^p_A(A_f, K) \to
\Hom_A(A_f, H^p(K)) \to 0
$$
을 얻는다.
% P02-E0414: frozen source omits the complementizer in “This proves (4) and (5) are ...”.
이는 (4)와 (5)가 (1)과 동치임을 증명한다.
\end{proof}

\begin{lemma}
\label{more-algebra-lemma-ideal-of-elements-complete-wrt}
$A$를 환, $K \in D(A)$라 하자. $T(K, f) = 0$을 만족하는
$f \in A$들의 집합 $I$는 $A$의 근기 아이디얼이다.
\end{lemma}

\begin{proof}
보조정리 \ref{more-algebra-lemma-hom-from-Af}의 결과들을 별도의 언급 없이
사용하겠다. $f \in I$, $g \in A$이면 $A_{gf}$는 $A_f$-가군이므로
모든 $n$에 대해 $\Ext^n_A(A_{gf}, K) = 0$이고, 따라서 $gf \in I$이다.
$f, g \in I$라고 가정하자. $f, g$가 $A_{f + g}$에서 단위
아이디얼을 생성하므로 짧은 완전열
$$
0 \to A_{f + g} \to A_{f(f + g)} \oplus A_{g(f + g)} \to A_{gf(f + g)} \to 0
$$
이 있다. 이는 Algebra, Lemma
\ref{algebra-lemma-standard-covering}과 마지막 화살표가 전사라는
쉬운 사실에서 따른다. $\Ext$의 긴 완전열과 모든 $n$에 대한
$\Ext^n_A(A_{f(f + g)}, K)$,
$\Ext^n_A(A_{g(f + g)}, K)$ 및
$\Ext^n_A(A_{gf(f + g)}, K)$의 소멸로부터, 모든 $n$에 대해
$\Ext^n_A(A_{f + g}, K)$가 소멸함을 얻는다. 마지막으로 어떤
$n > 0$에 대해 $f^n \in I$이면 $T(K, f) = T(K, f^n)$이거나
$A_f \cong A_{f^n}$이므로 $f \in I$이다.
\end{proof}

\begin{lemma}
\label{more-algebra-lemma-complete-derived-complete}
$A$를 환, $I \subset A$를 아이디얼, $M$을 $A$-가군이라 하자.
\begin{enumerate}
\item $M$이 $I$-진 완비이면 모든 $f \in I$에 대해
$T(M, f) = 0$이다.
\item 역으로 모든 $f \in I$에 대해 $T(M, f) = 0$이고 $I$가
유한 생성이면 $M \to \lim M/I^nM$은 전사이다.
\end{enumerate}
\end{lemma}

\begin{proof}
(1)의 증명. $M$이 $I$-진 완비라고 가정하자. 보조정리
\ref{more-algebra-lemma-hom-from-Af}에 의해 $\Ext^1_A(A_f, M) = 0$과
$\Hom_A(A_f, M) = 0$을 증명하면 충분하다. $M = \lim M/I^nM$이고
$\Hom_A(A_f, M/I^nM) = 0$이므로 $\Hom_A(A_f, M) = 0$이다. 확대
$$
0 \to M \to E \to A_f \to 0
$$
가 주어졌다고 하자. $n \geq 0$에 대해 $1/f^n$으로 사상되는
$e_n \in E$를 택하고, $n \geq 0$에 대해
$\delta_n = fe_{n + 1} - e_n \in M$으로 두자. $e_n$을
$$
e'_n = e_n + \delta_n + f\delta_{n + 1} + f^2 \delta_{n + 2} + \ldots
$$
로 바꾸자. $M$은 $I$에 관해 완비이고 $f \in I$이므로 이 무한합은
존재한다. 간단한 계산으로 $fe'_{n + 1} = e'_n$임을 알 수 있다.
따라서 $1/f^n$을 $e'_n$으로 보내어 이 확대의 분할을 얻는다.

\medskip\noindent
(2)의 증명. $I = (f_1, \ldots, f_r)$이고 $i = 1, \ldots, r$에 대해
$T(M, f_i) = 0$이라고 가정하자. Algebra, Lemma
\ref{algebra-lemma-when-surjective-to-completion}에 의해
$I = (f)$이고 $T(M, f) = 0$이라고 가정해도 된다. $n \geq 0$에 대해
$x_n \in M$이라 하자. 확대
$$
0 \to M \to E \to A_f \to 0
$$
를 생각하자. 이는
$$
E = M \oplus \bigoplus Ae_n\Big/\langle x_n - fe_{n + 1} + e_n\rangle
$$
로 주어지며 $e_n$을 $A_f$의 $1/f^n$으로 보낸다(위를 보라).
가정과 보조정리 \ref{more-algebra-lemma-hom-from-Af}에 의해 이 확대는 분할되므로
$E$에서 $A_f$의 한 복사본을 생성하는 원소 $x + e_0$을 얻는다. 그러면
$$
x + e_0 = x - x_0 + fe_1 = x - x_0 - f x_1 + f^2 e_2 = \ldots
$$
이다. 뱀 보조정리에 의해 $M/f^nM = E/f^nE$이므로
$x = x_0 + fx_1 + \ldots + f^{n - 1}x_{n - 1}$가 $f^nM$을 법으로
성립함을 알 수 있다. 다시 말해 원하는 대로 사상
$M \to \lim M/f^nM$은 전사이다.
\end{proof}

\noindent
위 결과들에 동기를 받아 다음과 같이 정의한다.

\begin{definition}
\label{more-algebra-definition-derived-complete}
$A$를 환, $K \in D(A)$라 하고, $I \subset A$를 아이디얼이라 하자.
모든 $f \in I$에 대해 $T(K, f) = 0$이면 $K$가
{\it $I$에 관해 유도 완비}라고 한다. $M$이 $A$-가군일 때,
$M[0] \in D(A)$가 $I$에 관해 유도 완비이면 $M$이
{\it $I$에 관해 유도 완비}라고 한다.
\end{definition}

\noindent
유도 완비 대상들로 이루어진 충만 부분범주
$D_{comp}(A) = D_{comp}(A, I) \subset D(A)$는 엄밀히 충만하고 포화된
삼각 부분범주이다. Derived Categories, Definitions
\ref{derived-definition-triangulated-subcategory}와
\ref{derived-definition-saturated}를 보라. 보조정리
\ref{more-algebra-lemma-ideal-of-elements-complete-wrt}에 의해 부분범주
$D_{comp}(A, I)$는 $I$의 근기 $\sqrt{I}$에만 의존한다. 달리 말해
$\Spec(A)$의 닫힌 부분집합 $Z = V(I)$에만 의존한다. 부분범주
$D_{comp}(A, I)$는 $D(A)$에서 곱과 호모토피 극한으로 보존된다.
그러나 일반적으로 가산 직합으로는 보존되지 않는다. 아이디얼 $I$의
선택이 문맥상 명백하면 흔히 $M$을 단순히 유도 완비 가군이라고 하겠다.

\begin{proposition}
\label{more-algebra-proposition-derived-complete-modules}
$I \subset A$를 환 $A$의 유한 생성 아이디얼, $M$을 $A$-가군이라 하자.
다음 조건들은 동치이다.
\begin{enumerate}
\item $M$은 $I$-진 완비이다.
\item $M$은 $I$에 관해 유도 완비이고 $\bigcap I^nM = 0$이다.
\end{enumerate}
\end{proposition}

\begin{proof}
이는 보조정리 \ref{more-algebra-lemma-complete-derived-complete}의 결과들에서
명백하다.
\end{proof}

\noindent
다음 보조정리는 유도 완비 가군들의 범주 $\mathcal{C}$가 아벨임을
보여준다. $\mathcal{C}$는 Grothendieck 아벨 범주가 아님이 밝혀진다.
Examples, Section \ref{examples-section-derived-complete-modules}를 보라.

\begin{lemma}
\label{more-algebra-lemma-serre-subcategory}
$I$를 환 $A$의 아이디얼이라 하자.
\begin{enumerate}
\item 유도 완비 $A$-가군들은 $\text{Mod}_A$의 약한 Serre 부분범주
$\mathcal{C}$를 이룬다.
\item $D_\mathcal{C}(A) \subset D(A)$는 유도 완비 대상들의 충만
부분범주이다.
\end{enumerate}
\end{lemma}

\begin{proof}
(2)는 보조정리 \ref{more-algebra-lemma-hom-from-Af}와 정의들에서 바로 따른다.
(1)을 위해 $M \to N$을 유도 완비 가군들의 사상이라 하자.
% P02-E0415: frozen source omits “by” when introducing K = (M -> N).
$D(A)$의 대응하는 대상을 $K = (M \to N)$으로 나타내자.
$f \in I$를 택하자. 모든 $n$에 대해 $\Ext_A^n(A_f, M)$과
$\Ext_A^n(A_f, N)$이 영이므로 모든 $n$에 대해
$\Ext_A^n(A_f, K)$도 영이다. 따라서 $K$는 유도 완비이다. (2)에 의해
$\Ker(M \to N)$과 $\Coker(M \to N)$은 $\mathcal{C}$의 대상이다.
마지막으로 $0 \to M_1 \to M_2 \to M_3 \to 0$이 $A$-가군들의
짧은 완전열이고 $M_1$, $M_3$가 유도 완비라고 하자. 그러면
$\Ext$들의 긴 완전열로부터 $M_2$가 유도 완비임이 따른다. 따라서
Homology, Lemma
\ref{homology-lemma-characterize-weak-serre-subcategory}에 의해
$\mathcal{C}$는 약한 Serre 부분범주이다.
\end{proof}

\noindent
다음 보조정리는 보조정리 \ref{more-algebra-lemma-derived-complete-zero-bis}에서
일반화할 것이다.

\begin{lemma}
\label{more-algebra-lemma-derived-complete-zero}
$I$를 환 $A$의 유한 생성 아이디얼, $M$을 유도 완비 $A$-가군이라 하자.
$M/IM = 0$이면 $M = 0$이다.
\end{lemma}

\begin{proof}
$M/IM$이 영이라고 가정하자. $I = (f_1, \ldots, f_r)$라 하자.
$N = M/(f_1, \ldots, f_i)M$이 영이 아니게 하는 가장 큰 정수
$i < r$을 잡자. 그러한 $i$가 존재하지 않으면 $M = 0$이며, 이것이
보이고자 하는 바이다. 그러면 보조정리
\ref{more-algebra-lemma-serre-subcategory}에 의해 $N$은 유도 완비 가군들 사이의
사상의 여핵으로서 유도 완비이다. $i$의 선택에 의해
$f_{i + 1} : N \to N$은 전사이다. 따라서
$$
\lim (\ldots \to N \xrightarrow{f_{i + 1}} N \xrightarrow{f_{i + 1}} N)
$$
은 영이 아니며, 이는 $N$의 유도 완비성과 모순이다.
\end{proof}

\noindent
환이 $I$-진 완비이면 유도 완비 복합체들을 풍부하게 얻는다.

\begin{lemma}
\label{more-algebra-lemma-pseudo-coherent-is-derived-complete}
$A$를 환, $I \subset A$를 아이디얼이라 하자. $A$가 유도 완비이면
(예를 들어 $I$-진 완비이면) $D(A)$의 모든 유사연접 대상은
유도 완비이다.
\end{lemma}

\begin{proof}
(보조정리 \ref{more-algebra-lemma-complete-derived-complete}이 보조정리의 괄호 안
명제를 설명한다.) $K$를 $D(A)$의 유사연접 대상이라 하자. 정의에 의해
이는 $K$가 유한 자유 $A$-가군들로 이루어진 위로 유계인 복합체
$K^\bullet$로 나타내어진다는 뜻이다. $A$가 유도 완비이므로 보조정리
\ref{more-algebra-lemma-serre-subcategory}의 (1)에 의해 모든 $n$에 대해
$H^n(K)$가 유도 완비이다. 다시 같은 보조정리의 (2)에 의해 이는
$K$가 유도 완비임을 함의한다.
\end{proof}

\begin{lemma}
\label{more-algebra-lemma-double-localize}
$A$를 환, $f, g \in A$라 하자. 그러면 $K \in D(A)$에 대해
$R\Hom_A(A_f, R\Hom_A(A_g, K)) = R\Hom_A(A_{fg}, K)$이다.
\end{lemma}

\begin{proof}
이는 보조정리 \ref{more-algebra-lemma-internal-hom}에서 따른다.
\end{proof}

\begin{lemma}
\label{more-algebra-lemma-derived-completion}
\begin{slogan}
유한 생성 아이디얼을 따른 유도 완비화는 존재하며, {\v C}ech 절차로
계산할 수 있다.
\end{slogan}
$I$를 환 $A$의 유한 생성 아이디얼이라 하자. 포함 함자
$D_{comp}(A, I) \to D(A)$는 왼쪽 수반을 갖는다. 즉, $D(A)$의 임의의
대상 $K$에 대해 $K$에서 $D(A)$의 유도 완비 대상으로 가는 사상
% P02-E0416: frozen source redundantly adds “of K” after a map already having source K.
$K \to K^\wedge$이 존재하여, $E$가 $D(A)$의 유도 완비 대상일 때마다
사상
$$
\Hom_{D(A)}(K^\wedge, E) \longrightarrow \Hom_{D(A)}(K, E)
$$
은 전단사이다. 실제로 $I$가 $f_1, \ldots, f_r \in A$에 의해
생성되면 $K$에 대해 함자적으로
$$
K^\wedge = R\Hom\left((A \to \prod\nolimits_{i_0} A_{f_{i_0}} \to
\prod\nolimits_{i_0 < i_1} A_{f_{i_0}f_{i_1}}
\to \ldots \to A_{f_1\ldots f_r}), K\right)
$$
이다.
\end{lemma}

\begin{proof}
보조정리의 마지막 표시식으로 $K^\wedge$을 정의하자. 복합체의 사상
$$
(A \to \prod\nolimits_{i_0} A_{f_{i_0}} \to
\prod\nolimits_{i_0 < i_1} A_{f_{i_0}f_{i_1}} \to
\ldots \to A_{f_1\ldots f_r}) \longrightarrow A
$$
가 있으며, 이는 사상 $K \to K^\wedge$을 유도한다.
$K^\wedge$이 유도 완비이고, $K$가 유도 완비일 때
$K \to K^\wedge$이 동형임을 증명하면 충분하다\footnote{실제로
$E$가 유도 완비이고 $a : K \to E$가 사상이면 가환도표
$$
\xymatrix{
K \ar[r]_a \ar[d] & E \ar[d]^{\cong} \\
K^\wedge \ar[r]^{a^\wedge} & E^\wedge
}
$$
는 $E$로 가는 모든 사상이 $K^\wedge$을 통해 분해됨을 보여준다.
구별 삼각형 $K \to K^\wedge \to C$를 택하고 ${}^\wedge$을 적용하자.
$C^\wedge = 0$을 얻는다. 위 논의에 의해 이는 유도 완비인 $E$에 대해
$\Hom(C, E) = 0$임을 뜻하며, 원하는 대로
$\Hom(K, E) = \Hom(K^\wedge, E)$임을 증명한다.}.

\medskip\noindent
$f \in A$라 하자. 보조정리 \ref{more-algebra-lemma-double-localize}에 의해 대상
$R\Hom_A(A_f, K^\wedge)$은
$$
R\Hom\left((A_f \to \prod\nolimits_{i_0} A_{ff_{i_0}} \to
\prod\nolimits_{i_0 < i_1} A_{ff_{i_0}f_{i_1}} \to
\ldots \to A_{ff_1\ldots f_r}), K\right)
$$
과 같다. $f \in I$이면 $f_1, \ldots, f_r$은 $A_f$에서 단위
아이디얼을 생성하므로, 확장 교대 {\v C}ech 복합체
$$
A_f \to \prod\nolimits_{i_0} A_{ff_{i_0}} \to
\prod\nolimits_{i_0 < i_1} A_{ff_{i_0}f_{i_1}} \to
\ldots \to A_{ff_1\ldots f_r}
$$
는 보조정리 \ref{more-algebra-lemma-extended-alternating-torsion}에 의해 $D(A)$에서
영이다. (실제로 어떤 $i$에 대해 $f = f_i$이면 이 복합체는 보조정리
\ref{more-algebra-lemma-extended-alternating-homotopy-zero}에 의해 영과 호모토피
동치이다. 이것이 우리가 필요한 유일한 경우이다.) 따라서
$R\Hom_A(A_f, K^\wedge) = 0$이고, 보조정리
\ref{more-algebra-lemma-hom-from-Af}에 의해 $K^\wedge$이 유도 완비임을 얻는다.

\medskip\noindent
역으로 $K$가 유도 완비이면 모든 $f = f_{i_0} \ldots f_{i_p}$,
$p \geq 0$에 대해 $R\Hom_A(A_f, K)$는 영이다. 따라서
$K \to K^\wedge$은 $D(A)$에서 동형이다.
\end{proof}

\begin{remark}
\label{more-algebra-remark-derived-completion}
$A$를 환, $I \subset A$를 유한 생성 아이디얼이라 하자. 보조정리
\ref{more-algebra-lemma-derived-completion}에 의해 존재하는 포함 함자
$D_{comp}(A, I) \to D(A)$의 왼쪽 수반을 {\it 유도 완비화}라고 한다.
이를 나타내기 위해 ``$K^\wedge$을 $K$의 유도 완비화라 하자''라고
말하겠다. 수반의 단위원은 함자적인 사상 $K \to K^\wedge$임을
명심하라.
\end{remark}

\begin{lemma}
\label{more-algebra-lemma-derived-completion-vanishes}
$A$를 환, $I \subset A$를 유한 생성 아이디얼이라 하자.
$K^\bullet$을 어떤 $f \in I$에 대해
$f : K^\bullet \to K^\bullet$이 동형인 $A$-가군 복합체, 즉
$K^\bullet$이 $A_f$-가군들의 복합체라 하자. 그러면 $K^\bullet$의 유도 완비화는
영이다.
\end{lemma}

\begin{proof}
% P02-E0417: frozen proof switches from K^bullet to undefined K and adds an unnecessary article before RHom.
실제로 이 경우 임의의 유도 완비 복합체 $L$에 대해
$R\Hom_A(K, L)$은 영이다. 보조정리 \ref{more-algebra-lemma-hom-from-Af}를 보라.
따라서 보조정리 \ref{more-algebra-lemma-derived-completion}의 보편 성질에 의해
$K^\wedge$은 영이다.
\end{proof}

\begin{lemma}
\label{more-algebra-lemma-completion-RHom}
$A$를 환, $I \subset A$를 유한 생성 아이디얼이라 하고,
$K, L \in D(A)$라 하자. 그러면
$$
R\Hom_A(K, L)^\wedge = R\Hom_A(K, L^\wedge) = R\Hom_A(K^\wedge, L^\wedge)
$$
이다.
\end{lemma}

\begin{proof}
보조정리 \ref{more-algebra-lemma-derived-completion}에 의해 유도 완비화는 어떤
$C \in D(A)$에 대한 $R\Hom_A(C, -)$로 주어진다. 그러면
\begin{align*}
R\Hom_A(C, R\Hom_A(K, L))
& =
R\Hom_A(C \otimes_A^\mathbf{L} K, L) \\
& =
R\Hom_A(K, R\Hom_A(C, L))
\end{align*}
이다. 보조정리 \ref{more-algebra-lemma-internal-hom}을 보라. 이것으로 첫 번째
등식을 증명한다. 사상 $K \to K^\wedge$은 사상
$$
R\Hom_A(K^\wedge, L^\wedge) \to R\Hom_A(K, L^\wedge)
$$
을 유도한다. 포함 함자의 왼쪽 수반으로서 유도 완비화의 정의에 의해
이 사상은 $D(A)$에서 동형이다.
\end{proof}

\begin{lemma}
\label{more-algebra-lemma-naive-derived-completion}
$A$를 환, $I \subset A$를 아이디얼이라 하자. $(K_n)$을 $D(A)$의
대상들의 역계라 하되, 모든 $f \in I$와 $n$에 대해 어떤
$e = e(n, f)$가 존재하여 $f^e$가 $K_n$ 위에서 영이라고 하자.
그러면 $K \in D(A)$에 대해 대상
$K' = R\lim (K \otimes_A^\mathbf{L} K_n)$은 $I$에 관해 유도 완비이다.
\end{lemma}

\begin{proof}
유도 완비 대상들의 범주는 $R\lim$으로 보존되므로 각
$K \otimes_A^\mathbf{L} K_n$이 유도 완비임을 보이면 충분하다.
가정에 의해 모든 $f \in I$에 대해 어떤 $e$가 존재하여
$f^e$가 $K \otimes_A^\mathbf{L} K_n$ 위에서 영이다. 물론 이는
$T(K \otimes_A^\mathbf{L} K_n, f) = 0$임을 함의하며, 결론을 얻는다.
\end{proof}

\begin{situation}
\label{more-algebra-situation-koszul}
$A$를 환이라 하고 $I = (f_1, \ldots, f_r) \subset A$라 하자.
$K_n^\bullet = K_\bullet(A, f_1^n, \ldots, f_r^n)$을
$f_1^n, \ldots, f_r^n$에 대한 Koszul 복합체라 하되, 차수
$-r, -r + 1, \ldots, 0$에 놓인 공사슬 복합체로 보자. 보조정리
\ref{more-algebra-lemma-functorial}의 함자성을 사용하면 역계
$$
\ldots \to K_3^\bullet \to K_2^\bullet \to K_1^\bullet
$$
을 얻는다. 이는 역계
$H^0(K_n^\bullet) = A/(f_1^n, \ldots, f_r^n)$ 및 사상
$A \to K_n^\bullet$들과 양립한다.
\end{situation}

\noindent
% P02-E0418: frozen source illogically says that a “feature ... will use” the following fact.
아래 논의의 핵심적인 특징은 $m > n$일 때 사상
$$
K_m^{-p} = \wedge^p(A^{\oplus r}) \to \wedge^p(A^{\oplus r}) = K_n^{-p}
$$
이 기저 원소 $e_{i_1} \wedge \ldots \wedge e_{i_p}$ 위에서
$f_{i_1}^{m - n} \ldots f_{i_p}^{m - n}$에 의한 곱셈으로
주어진다는 사실을 사용한다는 것이다.

\begin{lemma}
\label{more-algebra-lemma-koszul-derived-completion-complete}
상황 \ref{more-algebra-situation-koszul}에서, $K \in D(A)$에 대해 대상
$K' = R\lim (K \otimes_A^\mathbf{L} K_n^\bullet)$은 $I$에 관해
유도 완비이다.
\end{lemma}

\begin{proof}
이는 보조정리 \ref{more-algebra-lemma-naive-derived-completion}의 특수한 경우이다.
실제로 보조정리 \ref{more-algebra-lemma-homotopy-koszul}에 의해 $f_i^n$은
$K_n^\bullet$ 위에서 영과 호모토피인 자기준동형으로 작용한다.
\end{proof}

\begin{lemma}
\label{more-algebra-lemma-characterize-derived-complete-Koszul}
상황 \ref{more-algebra-situation-koszul}에서 $K \in D(A)$라 하자. 다음 조건들은
동치이다.
\begin{enumerate}
\item $K$는 $I$에 관해 유도 완비이다.
\item 표준적인 사상
$K \to R\lim (K \otimes_A^\mathbf{L} K_n^\bullet)$은 $D(A)$에서
동형이다.
\end{enumerate}
\end{lemma}

\begin{proof}
(2)가 성립하면 보조정리
\ref{more-algebra-lemma-koszul-derived-completion-complete}에 의해 $K$는 $I$에 관해
유도 완비이다. 역으로 $K$가 $I$에 관해 유도 완비라고 가정하자.
우직한 절단에 의한 여과
$$
K_n^\bullet \supset
\sigma_{\geq -r + 1}K_n^\bullet \supset
\sigma_{\geq -r + 2}K_n^\bullet \supset \ldots \supset
\sigma_{\geq -1}K_n^\bullet \supset
\sigma_{\geq 0}K_n^\bullet = A
$$
를 생각하자(Homology, Section \ref{homology-section-truncations}).
구성 $R\lim(K \otimes E)$은 두 번째 변수에서 완전하므로(보조정리
\ref{more-algebra-lemma-tensor-Rlim-exact}), $p < 0$에 대해
$$
R\lim \left(
K \otimes_A^\mathbf{L}
(\sigma_{\geq p}K_n^\bullet/ \sigma_{\geq p + 1}K_n^\bullet)
\right) = 0
$$
임을 보이면 충분하다. 위의 Koszul 복합체들의 명시적 기술에 의해
$$
R\lim \left(
K \otimes_A^\mathbf{L}
(\sigma_{\geq p}K_n^\bullet/ \sigma_{\geq p + 1}K_n^\bullet)
\right) =
\bigoplus\nolimits_{i_1, \ldots, i_{-p}}
T(K, f_{i_1}\ldots f_{i_{-p}})
$$
이다. 이는 $K$에 대한 가정으로부터 $p < 0$에 대해 영이다.
\end{proof}

\begin{lemma}
\label{more-algebra-lemma-derived-completion-koszul}
상황 \ref{more-algebra-situation-koszul}에서, $K \in D(A)$를 유도 극한
$K' = R\lim( K \otimes_A^\mathbf{L} K_n^\bullet )$으로 보내는 함자는
보조정리 \ref{more-algebra-lemma-derived-completion}에서 구성한 포함 함자
$D_{comp}(A) \to D(A)$의 왼쪽 수반이다.
\end{lemma}

\begin{proof}[첫 번째 증명]
대응 $K \leadsto K'$은 함자이고, 보조정리
\ref{more-algebra-lemma-koszul-derived-completion-complete}에 의해 $K'$은 $I$에
관해 유도 완비이다. 형식적 논증(생략한다)에 의해 $K$가 $I$에 관해
유도 완비일 때 $K \to K'$이 동형임을 보이면 충분하다. 이는 보조정리
\ref{more-algebra-lemma-characterize-derived-complete-Koszul}이다.
\end{proof}

\begin{proof}[두 번째 증명]
% P02-E0419: frozen source omits “by” when introducing the adjoint K -> K^wedge.
보조정리 \ref{more-algebra-lemma-derived-completion}에서 구성한 수반을
$K \mapsto K^\wedge$으로 나타내자. 그 보조정리에 의해
$$
K^\wedge = R\Hom\left((A \to \prod\nolimits_{i_0} A_{f_{i_0}} \to
\prod\nolimits_{i_0 < i_1} A_{f_{i_0}f_{i_1}}
\to \ldots \to A_{f_1\ldots f_r}), K\right)
$$
이다. 보조정리 \ref{more-algebra-lemma-extended-alternating-Cech-is-colimit-koszul}에서
확장 교대 {\v C}ech 복합체
$$
A \to \prod\nolimits_{i_0} A_{f_{i_0}} \to
\prod\nolimits_{i_0 < i_1} A_{f_{i_0}f_{i_1}}
\to \ldots \to A_{f_1\ldots f_r}
$$
가 차수 $0, \ldots, r$에 놓인 Koszul 복합체
$K^n = K(A, f_1^n, \ldots, f_r^n)$들의 여극한임을 보았다.
$K^n$은 유한 자유 $A$-가군들로 이루어진 유한 사슬 복합체이며,
그 쌍대는(보조정리 \ref{more-algebra-lemma-dual-perfect-complex}에서와 같이)
$R\Hom_A(K^n, A) = K_n$이다. 여기서 $K_n$은 (통상과 같이) 차수
$-r, \ldots, 0$에 놓인 Koszul 공사슬 복합체이다. 따라서
$$
R\Hom_A(\text{hocolim} K^n, K) =  R\lim (K \otimes_A^\mathbf{L} K_n)
$$
임을 보이면 충분하다. 이는 보조정리
\ref{more-algebra-lemma-colim-and-lim-of-duals}에서 따른다.
\end{proof}

\begin{lemma}
\label{more-algebra-lemma-derived-complete-bounded}
$I = (f_1, \ldots, f_r)$를 환 $A$의 유한 생성 아이디얼이라 하고,
$K$를 $D(A)$의 유도 완비 대상이라 하자. 다음 조건들은 동치이다.
\begin{enumerate}
\item $i > 0$에 대해 $H^i(K) = 0$이다.
\item $i > 0$에 대해 $H^i(K \otimes_A^\mathbf{L} A/I) = 0$이다.
\item $i > 0$에 대해 $H^i(K \otimes_A^\mathbf{L} K_1^\bullet) = 0$이다.
여기서 $K_1^\bullet$은 상황 \ref{more-algebra-situation-koszul}에서와 같다.
\end{enumerate}
\end{lemma}

\begin{proof}
함의 (1) $\Rightarrow$ (2)는 항상 참이다. 함의 (2) $\Rightarrow$ (3)은
보조정리 \ref{more-algebra-lemma-bounded-above-mod-I}에서 따른다. (3)을 가정하자.
$s = 0, \ldots, r$에 대해 복합체
$$
K(s) =  K \otimes_A^\mathbf{L} K_\bullet(A, f_1, \ldots, f_s)
$$
를 생각하자. 여기서 $K_\bullet(A, f_1, \ldots, f_s)$는 코호몰로지
차수 $-s, \ldots, 0$에 놓인 Koszul 복합체이다.
% P02-E0420: frozen source says “a distinguished triangles.”
보조정리 \ref{more-algebra-lemma-cone-koszul}에 의해 구별 삼각형
$$
K(s - 1) \xrightarrow{f_s} K(s - 1) \to K(s) \to K(s - 1)[1]
$$
이 있다. $K(0) = K$는 유도 완비이므로 상승 귀납법에 의해 각
$K(s)$는 $D(A)$의 유도 완비 대상이다. 하강 귀납법으로 $K(s)$가
차수 $> 0$에서 영이 아닌 코호몰로지를 갖지 않음을 보이자. 실제로
가정 (3)에 의해 $s = r$에 대해서는 성립한다. $K(s)$에 대해 성립하고
$s > 0$이면, $i > 0$에 대해
$f_s : H^i(K(s - 1)) \to H^i(K(s - 1))$은 전사이다. 보조정리들
\ref{more-algebra-lemma-serre-subcategory}와 \ref{more-algebra-lemma-derived-complete-zero}에 의해
$H^i(K(s - 1))$이 소멸한다.
\end{proof}

\begin{lemma}
\label{more-algebra-lemma-derived-complete-zero-bis}
\begin{slogan}
유도 Nakayama
\end{slogan}
\begin{reference}
관련 결과로 \cite[Proposition 6.5]{Dwyer-Greenlees}가 있다. 유도
Nakayama 보조정리는 예를 들어 2018년 가을 Columbia University에서
열린 Bhatt의 Prismatic 코호몰로지 세 번째 강의의 절 2 성질 (2)에서
찾을 수 있다. Leonid Positselski는
\url{https://mathoverflow.net/a/331501}에서 한 증명을 제안했다.
그러나 우리는 댓글에서 Anonymous가 제안한 증명을 따른다.
\end{reference}
$I$를 환 $A$의 유한 생성 아이디얼, $K$를 $D(A)$의 유도 완비
대상이라 하자. $K \otimes_A^\mathbf{L} A/I = 0$이면 $K = 0$이다.
\end{lemma}

\begin{proof}
이는 보조정리 \ref{more-algebra-lemma-derived-complete-bounded}의 특수한 경우이지만,
직접적인 증명도 주겠다. $I$의 생성원 $f_1, \ldots, f_r$을 택하자.
% P02-E0421: frozen source omits “by” when introducing K_n.
$A$ 위에서 $f_1^n, \ldots, f_r^n$에 대한 Koszul 복합체를
$K_n$으로 나타내자. $K_n$은 유계이고 $K_n$의 코호몰로지 가군들은
$f_1^n, \ldots, f_r^n$에 의해, 따라서 $I^{nr}$에 의해 소멸됨을
상기하라. 보조정리 \ref{more-algebra-lemma-derived-vanishing-mod-I}에 의해
$K \otimes_A^\mathbf{L} K_n = 0$이다. 보조정리
\ref{more-algebra-lemma-derived-completion-koszul}에 의해 $K$는 유도 완비이므로,
원하는 대로 $K = R\lim K \otimes_A^\mathbf{L} K_n = 0$이다.
\end{proof}

\noindent
Koszul 복합체와의 관계를 응용하면 유도 완비화가 유한 코호몰로지
차원을 가짐을 얻는다.

\begin{lemma}
\label{more-algebra-lemma-derived-completion-finite-cohomological-dimension}
$A$를 환, $I \subset A$를 $r$개의 원소로 생성될 수 있는
아이디얼이라 하자. 그러면 유도 완비화는 다음과 같이 유한
코호몰로지 차원을 갖는다.
\begin{enumerate}
\item $K \to L$을 $D(A)$의 사상이라 하되, $i \geq 1$에 대해
$H^i(K) \to H^i(L)$이 동형이고 $i = 0$에 대해 전사라고 하자.
그러면 $i \geq 1$에 대해 $H^i(K^\wedge) \to H^i(L^\wedge)$은
동형이고 $i = 0$에 대해 전사이다.
\item $K \to L$을 $D(A)$의 사상이라 하되, $i \leq -1$에 대해
$H^i(K) \to H^i(L)$이 동형이고 $i = 0$에 대해 단사라고 하자.
그러면 $i \leq -r - 1$에 대해
$H^i(K^\wedge) \to H^i(L^\wedge)$은 동형이고 $i = -r$에 대해
단사이다.
\end{enumerate}
\end{lemma}

\begin{proof}
$I$가 $f_1, \ldots, f_r$에 의해 생성된다고 하자. 임의의
$K \in D(A)$에 대해 보조정리 \ref{more-algebra-lemma-derived-completion-koszul}에 의해
$K^\wedge = R\lim K \otimes_A^\mathbf{L} K_n$이다. 여기서 $K_n$은
$f_1^n, \ldots, f_r^n$에 대한 Koszul 복합체이다. 따라서 보조정리
\ref{more-algebra-lemma-break-long-exact-sequence-modules}에 의해 짧은 완전열
$$
0 \to R^1\lim H^{i - 1}(K \otimes_A^\mathbf{L} K_n)
\to H^i(K^\wedge) \to \lim H^i(K \otimes_A^\mathbf{L} K_n) \to 0
$$
을 얻는다.

\medskip\noindent
(1)의 증명. 구별 삼각형 $K \to L \to C \to K[1]$을 택하자.
그러면 $i \geq 0$에 대해 $H^i(C) = 0$이다. $K_n$은 차수
$\leq 0$에 놓여 있으므로 $i \geq 0$에 대해
$H^i(C \otimes_A^\mathbf{L} K_n) = 0$이고,
$H^{-1}(C \otimes_A^\mathbf{L} K_n) =
H^{-1}(C) \otimes_A A/(f_1^n, \ldots, f_r^n)$은 전이사상들이 전사인
계이다. 위 표시식은 $i \geq 0$에 대해 $H^i(C^\wedge) = 0$임을
보여준다. 구별 삼각형
$K^\wedge \to L^\wedge \to C^\wedge \to K^\wedge[1]$을 적용하면
(1)을 얻는다.

\medskip\noindent
(2)의 증명. 구별 삼각형 $K \to L \to C \to K[1]$을 택하자.
그러면 $i < 0$에 대해 $H^i(C) = 0$이다. $K_n$은 차수
$-r, \ldots, 0$에 놓여 있으므로 $i < -r$에 대해
$H^i(C \otimes_A^\mathbf{L} K_n) = 0$이다. 위 표시식은
% P02-E0422: frozen source concludes i < r although the stated bound requires i < -r.
$i < r$에 대해 $H^i(C^\wedge) = 0$임을 보여준다. 구별 삼각형
$K^\wedge \to L^\wedge \to C^\wedge \to K^\wedge[1]$을 적용하면
(2)를 얻는다.
\end{proof}

\begin{lemma}
\label{more-algebra-lemma-derived-completion-spectral-sequence}
$A$를 환, $I \subset A$를 유한 생성 아이디얼이라 하고,
$K^\bullet$을 $A$-가군의 여과 복합체라 하자. 쌍차수 유도 완비
$A$-가군들의 표준적인 스펙트럼 열
% P02-E0423: frozen source switches from \text{d}_r here to plain d_r immediately below.
$(E_r, \text{d}_r)_{r \geq 1}$이 존재하며, $d_r$의 쌍차수는
$(r, -r + 1)$이고
$$
E_1^{p, q} = H^{p + q}((\text{gr}^pK^\bullet)^\wedge)
$$
이다. 각 $K^n$ 위의 여과가 유한이면 이 스펙트럼 열은 유계이고
$H^*((K^\bullet)^\wedge)$으로 수렴한다.
\end{lemma}

\begin{proof}
보조정리 \ref{more-algebra-lemma-derived-completion}에 의해 유도 완비화는 어떤
$C \in D^b(A)$에 대한 $R\Hom_A(C, -)$로 주어진다. 보조정리들
\ref{more-algebra-lemma-derived-completion-finite-cohomological-dimension}과
\ref{more-algebra-lemma-projective-amplitude}에 의해 $C$는 유한 사영 차원을 갖는다.
따라서 $C$를 나타내는 사영 가군들의 유계 복합체 $P^\bullet$을
택할 수 있다. 그러면
$$
M^\bullet = \Hom^\bullet(P^\bullet, K^\bullet)
$$
은 $(K^\bullet)^\wedge$을 나타내는 $A$-가군 복합체이다. 이 복합체는
$F^pM^\bullet = \Hom^\bullet(P^\bullet, F^pK^\bullet)$로 주어지는
여과를 갖는다. $F^pM^\bullet$은 $(F^pK^\bullet)^\wedge$을 나타내며,
% P02-E0424: frozen source drops the filtration index p from the associated-graded target.
따라서 $\text{gr}^pM^\bullet$은 $(\text{gr}K^\bullet)^\wedge$을 나타낸다.
그러므로 여과 복합체 $M^\bullet$의 스펙트럼 열을 취하여 우리의
스펙트럼 열을 얻는다. Homology, Section
\ref{homology-section-filtered-complex}을 보라. 각 $K^n$ 위의 여과가
유한이면 $P^\bullet$이 유계 복합체이므로 각 $M^n$ 위의 여과도
유한이다. 따라서 마지막 명제는 Homology, Lemma
\ref{homology-lemma-biregular-ss-converges}에서 따른다.
\end{proof}

\begin{example}
\label{more-algebra-example-derived-completion-spectral-sequence}
$A$를 환, $I \subset A$를 유한 생성 아이디얼이라 하고,
$K^\bullet$을 $A$-가군의 복합체라 하자. 보조정리
\ref{more-algebra-lemma-derived-completion-spectral-sequence}를
$F^pK^\bullet = \tau_{\leq -p}K^\bullet$에 적용할 수 있다. 그러면
유계 스펙트럼 열
$$
E_1^{p, q} = H^{p + q}(H^{-p}(K^\bullet)^\wedge[p]) =
H^{2p + q}(H^{-p}(K^\bullet)^\wedge)
$$
을 얻으며, 이는 $H^{p + q}((K^\bullet)^\wedge)$으로 수렴한다.
$p = -j$, $q = i + 2j$로 다시 번호를 매기면, 임의의 $K \in D(A)$에
대해 쌍차수 유도 완비 가군들의 유계 스펙트럼 열
$(E'_r, d'_r)_{r \geq 2}$이 존재하고, $d'_r$의 쌍차수는
$(r, -r + 1)$이며,
$$
(E'_2)^{i, j} = H^i(H^j(K)^\wedge)
$$
이고 $H^{i + j}(K^\wedge)$으로 수렴함을 얻는다.
\end{example}

\begin{lemma}
\label{more-algebra-lemma-restriction-derived-complete}
$A \to B$를 환 사상, $I \subset A$를 아이디얼이라 하자. 제한 함자
$D(B) \to D(A)$에 의한 $D_{comp}(A, I)$의 역상은
$D_{comp}(B, IB)$이다.
\end{lemma}

\begin{proof}
보조정리 \ref{more-algebra-lemma-ideal-of-elements-complete-wrt}에 의해
$L \in D(B)$가 $D_{comp}(B, IB)$에 속할 필요충분조건은
% P02-E0425: frozen source calls an ordinary ring element a “local section.”
모든 국소 단면 $f \in I$에 대해 $T(L, f)$가 영인 것이다.
$T(L, f)$의 코호몰로지는 아벨 군의 범주에서 계산되므로, $f$를
$A$의 원소로 생각하든 $B$에서의 $f$의 상을 취하든 상관없다.
이 사실과 유도 완비 대상의 정의로부터 보조정리가 바로 따른다.
\end{proof}

\begin{lemma}
\label{more-algebra-lemma-restriction-derived-complete-equivalence}
$A \to B$를 환 사상, $I \subset A$를 유한 생성 아이디얼이라 하자.
$A \to B$가 평탄하고 $A/I \cong B/IB$이면 제한 함자
$D(B) \to D(A)$는 동치
$D_{comp}(B, IB) \to D_{comp}(A, I)$를 유도한다.
\end{lemma}

\begin{proof}
$I$의 생성원 $f_1, \ldots, f_r$을 택하자.
% P02-E0426: frozen source omits “by” when introducing the quasi-isomorphism.
보조정리 \ref{more-algebra-lemma-map-identifies-koszul-and-cech-complexes}의 확장 교대
{\v C}ech 복합체들의 준동형사상을
$\check{\mathcal{C}}^\bullet_A \to \check{\mathcal{C}}^\bullet_B$로
나타내자. $K \in D_{comp}(A, I)$라 하자. $I^\bullet$을 $K$를
나타내는 $A$-가군의 K-단사 복합체라 하자. 모든 $f \in I$과
$n \in \mathbf{Z}$에 대해 $\Ext^n_A(A_f, K)$과
$\Ext^n_A(B_f, K)$이 영이므로(보조정리 \ref{more-algebra-lemma-hom-from-Af}),
$\check{\mathcal{C}}^\bullet_A \to A$와
$\check{\mathcal{C}}^\bullet_B \to B$는 준동형사상
$$
I^\bullet = \Hom_A(A, I^\bullet) \longrightarrow
\text{Tot}(\Hom_A(\check{\mathcal{C}}^\bullet_A, I^\bullet))
$$
과
$$
\Hom_A(B, I^\bullet) \longrightarrow
\text{Tot}(\Hom_A(\check{\mathcal{C}}^\bullet_B, I^\bullet))
$$
을 유도한다. 세부사항 일부는 생략한다.
$\check{\mathcal{C}}^\bullet_A \to \check{\mathcal{C}}^\bullet_B$가
준동형사상이고 $I^\bullet$이 K-단사이므로
$\Hom_A(B, I^\bullet) \to I^\bullet$이 준동형사상임을 얻는다.
복합체 $\Hom_A(B, I^\bullet)$은 $B$-가군들의 복합체이므로 $K$가
제한 사상의 상에 속함을, 즉 함자가 본질적 전사임을 얻는다
% P02-E0427: frozen source omits terminal punctuation before the next paragraph.

\medskip\noindent
실제로 이 논증은
$F : D_{comp}(A, I) \to D_{comp}(B, IB)$,
$K \mapsto \Hom_A(B, I^\bullet)$가 제한의 왼쪽 역임을 보여준다.
마지막으로 $L \in D_{comp}(B, IB)$라 하자. $L$을 $B$-가군의
K-단사 복합체 $J^\bullet$로 나타내자. 보조정리
\ref{more-algebra-lemma-K-injective-flat}에 의해 $J^\bullet$은 $A$-가군의
복합체로서도 K-단사이므로
$F(\text{restriction of }L) = \Hom_A(B, J^\bullet)$이다.
$B$-가군 복합체의 사상 $J^\bullet \to \Hom_A(B, J^\bullet)$이 있고,
이를 $\Hom_A(B, J^\bullet) \to J^\bullet$과 합성하면 항등이다.
따라서 $F$는 제한의 오른쪽 역이기도 하며, 증명이 끝난다.
\end{proof}

\section{유도 완비 가군들의 범주}
\label{more-algebra-section-derived-complete-modules}

\noindent
$A$를 환, $I$를 아이디얼이라 하자.
% P02-E0428: frozen source omits “by” when introducing the category C.
유도 완비 가군들의 범주를 $\mathcal{C}$로 나타내자. 정의
\ref{more-algebra-definition-derived-complete}를 보라. 이 절에서는 이 범주의 몇몇
성질을 논의한다. Examples, Section
\ref{examples-section-derived-complete-modules}에서 일반적으로
$\mathcal{C}$가 Grothendieck 아벨 범주가 아님을 보인다.

\medskip\noindent
보조정리 \ref{more-algebra-lemma-serre-subcategory}에 의해 범주 $\mathcal{C}$는
아벨이고 포함 함자 $\mathcal{C} \to \text{Mod}_A$는 완전하다.

\medskip\noindent
$D_{comp}(A) \subset D(A)$는 곱에 대해 닫혀 있고(정의
\ref{more-algebra-definition-derived-complete} 뒤의 논의를 보라), $D(A)$의 곱은
복합체 수준에서 계산되므로, $\mathcal{C}$는 $\text{Mod}_A$의 곱과
일치하는 곱을 가짐을 알 수 있다. 따라서 $\mathcal{C}$는 사실 모든
극한을 가지며 포함 함자 $\mathcal{C} \to \text{Mod}_A$는 이들과
가환한다. Categories, Lemma
\ref{categories-lemma-limits-products-equalizers}를 보라.

\medskip\noindent
$I$가 유한 생성이라고 가정하자. ${}^\wedge : D(A) \to D(A)$를
보조정리 \ref{more-algebra-lemma-derived-completion}의 유도 완비화 함자라 하자.
함자
$$
\text{Mod}_A \longrightarrow \mathcal{C},\quad
M \longmapsto H^0(M^\wedge)
$$
가 포함 함자 $\mathcal{C} \to \text{Mod}_A$의 왼쪽 수반임을 보이자.
예를 들어 보조정리
\ref{more-algebra-lemma-derived-completion-finite-cohomological-dimension}에 의해
$i > 0$에 대해 $H^i(M^\wedge) = 0$임에 유의하라. 따라서 $N$이
유도 완비 $A$-가군이면 원하는 대로
\begin{align*}
\Hom_\mathcal{C}(H^0(M^\wedge), N)
& =
\Hom_{D_{comp}(A)}(M^\wedge, N)\\
& =
\Hom_{D(A)}(M, N) \\
& =
\Hom_A(M, N)
\end{align*}
이다.

\medskip\noindent
$T$를 전순서 집합이라 하고 $t \mapsto M_t$를 유도 완비
$A$-가군들의 계, 즉 $\mathcal{C}$에서 $T$ 위의 계라 하자. Categories,
Section \ref{categories-section-posets-limits}를 보라.
% P02-E0429: frozen source omits “by” when introducing the colimit notation.
$\text{Mod}_A$에서 이 계의 여극한을 $\colim_{t \in T} M_t$로 나타내자.
위 논의로부터
$$
H^0((\colim_{t \in T} M_t)^\wedge)
$$
이 $\mathcal{C}$에서 이 계의 여극한임이 형식적으로 따른다. 이로써
$\mathcal{C}$가 모든 여극한을 가짐을 알 수 있다. 일반적으로 포함
함자 $\mathcal{C} \to \text{Mod}_A$는 여극한과 가환하지 않는다.
Examples, Section \ref{examples-section-derived-complete-modules}를 보라.

\begin{lemma}
\label{more-algebra-lemma-derived-complete-modules}
$A$를 환, $I \subset A$를 아이디얼이라 하자. 유도 완비 가군들의 범주
$\mathcal{C}$는 아벨이고 모든 극한을 가지며, 포함 함자
$F : \mathcal{C} \to \text{Mod}_A$는 완전하고 극한과 가환한다.
$I$가 유한 생성이면 $\mathcal{C}$는 모든 여극한을 가지고
% P02-E0430: frozen source omits terminal punctuation after this clause.
$F$는 왼쪽 수반을 갖는다
\end{lemma}

\begin{proof}
이는 위의 논의를 요약한 것이다.
\end{proof}

\section{주 아이디얼에 대한 유도 완비화}
\label{more-algebra-section-derived-completion-principal}

\noindent
이 절에서는 아이디얼이 하나의 원소로 생성될 때 유도 완비화에 어떤
일이 일어나는지 논의한다.

\begin{lemma}
\label{more-algebra-lemma-lift-universally}
$A$를 환, $f \in A$라 하자. 어떤 정수 $c \geq 1$이 존재하여
$A[f^c] = A[f^{c + 1}] = A[f^{c + 2}] = \ldots$라고 하자(예를 들어
$A$가 Noether이면 그렇다). 그러면 모든 $n \geq 1$에 대해 $D(A)$에서
사상
$$
(A \xrightarrow{f^n} A) \longrightarrow A/(f^n),
\quad\text{and}\quad
A/(f^{n + c}) \longrightarrow (A \xrightarrow{f^n} A)
$$
이 존재하고, $D(A)$에서 pro-대상들 $\{A/(f^n)\}$과
$\{(f^n : A \to A)\}$의 동형을 유도한다.
\end{lemma}

\begin{proof}
첫 번째 표시된 화살표는 명백하다. 보조정리의 두 번째 화살표는 도표
$$
\xymatrix{
A/A[f^c] \ar[r]_-{f^{n + c}} \ar[d]_{f^c} & A \ar[d]^1 \\
A \ar[r]^{f^n} & A
}
$$
로 정의할 수 있다. 위쪽 가로 화살표가 단사이므로 위 행의 복합체는
$A/f^{n + c}A$와 준동형이다. pro-대상들에 관한 명제를 보이는 데
필요한 합성의 계산은 생략한다.
\end{proof}

\begin{lemma}
\label{more-algebra-lemma-when-does-it-work}
$A$를 환, $f \in A$라 하고 $I = (f)$로 두자. 이 상황에는 순진한
유도 완비화 $K \mapsto K' = R\lim (K \otimes_A^\mathbf{L} A/f^nA)$와
보조정리 \ref{more-algebra-lemma-derived-completion-koszul}의 유도 완비화
$$
K \mapsto K^\wedge = R\lim (K \otimes_A^\mathbf{L} (A \xrightarrow{f^n} A))
$$
이 있다. 함자들의 자연변환 $K^\wedge \to K'$이 동형일 필요충분조건은
$A$의 $f$-거듭제곱 꼬임이 유계인 것이다.
\end{lemma}

\begin{proof}
$f$-거듭제곱 꼬임이 유계이면 보조정리
\ref{more-algebra-lemma-lift-universally}에 의해 pro-대상들
$\{(f^n : A \to A)\}$과 $\{A/f^nA\}$은 동형이다. 따라서 보조정리
\ref{more-algebra-lemma-Rlim-pro-equal}에 의해 두 함자는 동형이다. 역으로
보조정리 \ref{more-algebra-lemma-tensor-Rlim-exact}에 의해 그 조건은 정확히
$$
R\lim (K \otimes_A^\mathbf{L} A[f^n])
$$
이 모든 $K \in D(A)$에 대해 영이라는 것임을 알 수 있다. 여기서 계
$(A[f^n])$의 사상들은 $f$에 의한 곱셈으로 주어진다. $K = A$와
$K = \bigoplus_{i \in \mathbf{N}} A$를 취하면 보조정리
\ref{more-algebra-lemma-Rlim-zero-of-direct-sums}로부터 이는 $(A[f^n])$가
pro-대상으로서 영임을 함의함을 알 수 있다. 즉 어떤 $n$에 대해
$f^{n - 1}A[f^n] = 0$이고, 즉 $A[f^{n - 1}] = A[f^n]$이며, 즉
$f$-거듭제곱 꼬임이 유계이다.
\end{proof}

\begin{example}
\label{more-algebra-example-derived-complete-modules}
$A$를 환, $f \in A$를 비영인자라 하자. 염두에 둘 예는
$A = \mathbf{Z}_p$와 $f = p$이다. $M$을 $A$-가군이라 하자. 주장:
$M$이 $f$에 관해 유도 완비일 필요충분조건은 짧은 완전열
$$
0 \to K \to L \to M \to 0
$$
이 존재하여 $K, L$이 $f$-꼬임이 영인 $f$-진 완비 가군인 것이다.
% P02-E0431: frozen source says “a such a short exact sequence.”
실제로 그러한 짧은 완전열이 있으면
$$
M \otimes_A^\mathbf{L} (A \xrightarrow{f^n} A) = (K/f^nK \to L/f^nL)
$$
이다. 이는 $f$가 $K$와 $L$에서 비영인자이기 때문이다. 따라서
$R\lim (M \otimes_A^\mathbf{L} (A \xrightarrow{f^n} A))$은
$K \to L$, 즉 $M$과 준동형임을 얻는다. 이는 보조정리
\ref{more-algebra-lemma-characterize-derived-complete-Koszul}에 의해 $M$이 유도
완비임을 보여준다. 역으로 $M$이 유도 완비라고 가정하자. $F$가 자유
$A$-가군인 전사 $F \to M$을 택하자. $f$는 $F$에서 비영인자이므로
$F$의 유도 완비화는 $L = \lim F/f^nF$이다.
% P02-E1118: frozen source's compound “f-torsion free” lacks the second hyphen.
$L$에는 $f$-꼬임이 없음에 유의하라. 실제로 $x_n \in F$인 $(x_n)$이
$L$의 원소 $\xi$를 나타내고 $f\xi = 0$이면, 어떤
$z_n, y_n \in F$에 대해 $x_n = x_{n + 1} + f^nz_n$이고
$fx_n = f^ny_n$이다. 그러면
$f^n y_n = fx_n = fx_{n + 1} + f^{n + 1}z_n = f^{n + 1}y_{n + 1} + f^{n + 1}z_n$이고,
$f$가 $F$에서 비영인자이므로 $y_n \in fF$이다. 이는
$x_n \in f^nF$, 즉 $\xi = 0$임을 함의한다. $L$은 유도 완비화이므로
보편 성질은 $F \to M$을 분해하는 사상 $L \to M$을 준다.
$K = \Ker(L \to M)$를 그 핵이라 하자. 다시 $K$에는 $f$-꼬임이
없으므로 $K$의 유도 완비화는 $\lim K/f^nK$이다. 한편 $M$과 $L$은
둘 다 유도 완비이므로 보조정리 \ref{more-algebra-lemma-serre-subcategory}에 의해
$K$도 유도 완비이다. 따라서 $K = \lim K/f^nK$이고 주장이 증명된다.
\end{example}

\begin{example}
\label{more-algebra-example-derived-complete-not-complete}
$p$를 소수라 하자. $x$를 $py$로 보내는 다항식 대수의 사상
$\mathbf{Z}_p[x] \to \mathbf{Z}_p[y]$를 생각하자. 유도된 (통상적인)
$p$-진 완비화 사상의 여핵
$M = \Coker(\mathbf{Z}_p[x]^\wedge \to \mathbf{Z}_p[y]^\wedge)$을
생각하자. 그러면 명제 \ref{more-algebra-proposition-derived-complete-modules}와
보조정리 \ref{more-algebra-lemma-serre-subcategory}에 의해 $M$은 유도 완비
$\mathbf{Z}_p$-가군이다. 예
\ref{more-algebra-example-derived-complete-modules}의 논의도 보라. 그러나
$1 + py + p^2 y^2 + \ldots$가 $\bigcap p^nM$에 들어 있는 $M$의
영이 아닌 원소로 사상되므로 $M$은 $p$-진 완비가 아니다.
\end{example}

\begin{example}
\label{more-algebra-example-spectral-sequence-principal}
$A$를 환, $f \in A$라 하자.
% P02-E0432: frozen source omits “by” when introducing K -> K^wedge.
$(f)$에 관한 유도 완비화를 $K \mapsto K^\wedge$으로 나타내자.
$M$을 $A$-가군이라 하자. 보조정리
\ref{more-algebra-lemma-derived-completion-koszul}에 의해
$$
M^\wedge = R\lim (M \xrightarrow{f^n} M)
$$
임을 사용하고 보조정리 \ref{more-algebra-lemma-break-long-exact-sequence-modules}를
사용하면
$$
H^{-1}(M^\wedge) = \lim M[f^n] = T_f(M)
$$
을 얻는다. 이는 {\it $M$의 $f$-진 Tate 가군}이다. 여기서 사상
$M[f^n] \to M[f^{n - 1}]$은 $f$에 의한 곱셈으로 주어진다. 또한
$H^0(M^\wedge)$을 기술하는 짧은 완전열
$$
0 \to R^1\lim M[f^n] \to H^0(M^\wedge) \to \lim M/f^n M \to 0
$$
이 있다. 전이사상들이 전사이므로(보조정리
\ref{more-algebra-lemma-compute-Rlim-modules})
$H^1(M^\wedge) = R^1\lim M/f^nM = 0$이다. $M^\wedge$의 다른 모든
코호몰로지는 자명한 이유로 영이다. 마지막으로 $K \in D(A)$와
$p \in \mathbf{Z}$에 대해 짧은 완전열
$$
0 \to H^0(H^p(K)^\wedge) \to H^p(K^\wedge) \to T_f(H^{p + 1}(K)) \to 0
$$
이 있다. 이는 예 \ref{more-algebra-example-derived-completion-spectral-sequence}의
스펙트럼 열에서 따른다. 영이 아닌 항을 주는 지표가 $i = -1, 0$뿐이므로
이 스펙트럼 열은 $E_2$에서 퇴화한다. 다음 보조정리는 더 많은 정보를
준다.
\end{example}

\begin{lemma}
\label{more-algebra-lemma-compare-completion-derived-completion}
$A$를 환, $f \in A$라 하고 $K$를 $D(A)$의 대상이라 하자.
% P02-E0433: frozen source uses a compressed denote construction without “by” or “let”.
$K_n = K \otimes_A^\mathbf{L} (A \xrightarrow{f^n} A)$로 두자.
모든 $p \in \mathbf{Z}$에 대해 행과 열이 완전인 가환도표
$$
\xymatrix{
& 0 & 0 \\
0 \ar[r] &
\widehat{H^p(K)} \ar[r] \ar[u] &
\lim H^p(K_n) \ar[r] \ar[u] &
T_f(H^{p + 1}(K)) \ar[r] &
0 \\
0 \ar[r] &
H^0(H^p(K)^\wedge) \ar[r] \ar[u] &
H^p(K^\wedge) \ar[r] \ar[u] &
T_f(H^{p + 1}(K)) \ar[r] \ar@{=}[u] &
0 \\
&
R^1\lim H^p(K)[f^n] \ar[u] \ar[r]^\cong &
R^1\lim H^{p - 1}(K_n) \ar[u] \\
& 0 \ar[u] & 0 \ar[u]
}
$$
가 있다. 여기서
$\widehat{H^p(K)} = \lim H^p(K)/f^nH^p(K)$는 통상적인 $f$-진
완비화이다. 왼쪽 수직 짧은 완전열과 가운데 수평 짧은 완전열은 예
% P02-E0434: frozen source omits a sentence-ending period after this reference.
\ref{more-algebra-example-spectral-sequence-principal}에서 가져온 것이다
가운데 수직 짧은 완전열은 보조정리
\ref{more-algebra-lemma-break-long-exact-sequence-modules}에서 온 것이다.
\end{lemma}

\begin{proof}
위쪽 수평 짧은 완전열을 구성하기 위해, $K_n$을
$f^n : K \to K$의 원뿔의 옮김으로 구성한 데서 오는 역계의
짧은 완전열
$$
0 \to H^p(K)/f^nH^p(K) \to H^p(K_n) \to H^{p + 1}(K)[f^n] \to 0
$$
이 있음을 관찰하자. 이들의 역극한을 취하면 위쪽 수평 짧은 완전열을
얻는다. Homology, Lemma \ref{homology-lemma-Mittag-Leffler}를 보라.

\medskip\noindent
보조정리에서와 같은 가환도표가 있음을 증명하자. 사상
$L = \tau_{\leq p}K \to K$를 생각하자.
$L_n = L \otimes_A^\mathbf{L} (A \xrightarrow{f^n} A)$로 두면 짧은
완전열들의 사상을 유도하는 역계의 사상 $(L_n) \to (K_n)$을 얻는다.
$$
\xymatrix{
0 & 0 \\
\lim H^p(L_n) \ar[r] \ar[u] &
\lim H^p(K_n) \ar[u] \\
H^p(L^\wedge) \ar[r] \ar[u] &
H^p(K^\wedge) \ar[u] \\
R^1\lim H^{p - 1}(L_n) \ar[r] \ar[u] &
R^1\lim H^{p - 1}(K_n) \ar[u] \\
0 \ar[u] & 0 \ar[u]
}
$$
$i > p$에 대해 $H^i(L) = 0$이고 $H^p(L) = H^p(K)$이므로,
보조정리의 명제에 있는 참고문헌들을 사용한 계산은
$H^p(L^\wedge) = H^0(H^p(K)^\wedge)$이고
$H^p(L_n) = H^p(K)/f^nH^p(K)$임을 보여준다. 한편
$H^{p - 1}(L_n) = H^{p - 1}(K_n)$이고, 우리는 이미
$H^0(H^p(K)^\wedge) \to \widehat{H^p(K)}$의 핵이
$R^1\lim H^p(K)[f^n]$과 같음을 알고 있으므로, 보조정리의 명제에
표시된 동형을 얻는다. 증명의 첫 문단의 구성으로 위 행을 정의했을 때
도표의 가장 오른쪽 정사각형이 가환한다는 확인은 생략한다.
\end{proof}

\begin{remark}
\label{more-algebra-remark-ML-condition}
보조정리 \ref{more-algebra-lemma-compare-completion-derived-completion}에서와 같은
표기 아래, 역계 $H^p(K_n)$이 ML을 가질 필요충분조건은 역계
$H^{p + 1}(K)[f^n]$이 ML을 갖는 것임도 알 수 있다. 이는 역계의 짧은
완전열 $0 \to H^p(K)/f^nH^p(K) \to H^p(K_n) \to H^{p + 1}(K)[f^n] \to 0$
(보조정리의 증명을 보라)을 Homology, Lemma
\ref{homology-lemma-Mittag-Leffler} 및 보조정리
\ref{more-algebra-lemma-Mittag-Leffler}와 결합하면 따른다.
\end{remark}

\begin{lemma}[Bhatt]
\label{more-algebra-lemma-torsion-and-derived-complete}
$I$를 환 $A$의 유한 생성 아이디얼이라 하고 $M$을 유도 완비
$A$-가군이라 하자. $M$이 $I$-거듭제곱 꼬임 가군이면 어떤 $n$에 대해
$I^nM = 0$이다.
\end{lemma}

\begin{proof}
$I = (f_1, \ldots, f_r)$라 하자. 각 $i$에 대해 어떤 $n_i$가 존재하여
$f_i^{n_i}M = 0$임을 보이면 충분하다. 따라서 $I = (f)$가 주
아이디얼이라고 가정해도 된다. $x$를 $f$로 보내는 환 사상
$B = \mathbf{Z}[x] \to A$를 잡자. 보조정리
\ref{more-algebra-lemma-restriction-derived-complete}에 의해 $M$은 아이디얼 $(x)$에
관해 $B$-가군으로서 유도 완비이다. $A$를 $B$로 바꾸어 $f$가
$A$에서 비영인자라고 가정해도 된다.

\medskip\noindent
$I = (f)$이고 $f \in A$가 비영인자라고 가정하자. 예
\ref{more-algebra-example-derived-complete-modules}에 의해 짧은 완전열
$$
0 \to K \xrightarrow{u} L \to M \to 0
$$
이 존재하여 $K$와 $L$은 $f$-꼬임이 영인 $I$-진 완비
$A$-가군이다\footnote{증명을 위해서는 $K$와 $L$이 $I$-진 완비
$A$-가군인 수열 $K \xrightarrow{u} L \to M \to 0$이 존재함을
보이는 것으로 충분하다. 이는 $F_i$가 자유인 표시
$F_1 \to F_0 \to M \to 0$을 택한 다음, $K$와 $L$을 각각
$F_1$과 $F_0$의 $f$-진 완비화로 두면 보일 수 있다. 실제로 $f$가
비영인자이므로 이 완비화들은 유도 완비화이고 수열은 계속 완전하다.}.
$K$와 $L$을 $I$-진 위상을 갖는 위상 가군으로 보자. 그러면 $u$는
연속이다. 다음과 같이 두자.
$$
L_n = \{x \in L \mid f^n x \in u(K)\}
$$
$M$이 $f$-거듭제곱 꼬임이므로 $L = \bigcup L_n$이다. $N_n$을
$L$에서 $L_n$의 폐포라 하자. 보조정리
\ref{more-algebra-lemma-consequence-baire-complete-module}에 의해 어떤 $n$에 대해
$N_n$은 $L$에서 열린집합이다. 그런 $n$을 고정하자.
$f^{n + m} : L \to L$은 연속 열린 사상이고
$f^{n + m} L_n \subset u(f^m K)$이므로, 모든 $m \geq 1$에 대해
$u(f^mK)$의 폐포가 열려 있음을 얻는다. 따라서 보조정리
\ref{more-algebra-lemma-open-mapping}에 의해 $u$가 열린 사상임을 얻는다.
그러므로 어떤 $t$에 대해 $f^tL \subset \Im(u)$이고, 원하는 대로
$f^t$가 $M$을 소멸시킴을 얻는다.
\end{proof}

\begin{lemma}
\label{more-algebra-lemma-kernel-to-completion-square-zero}
% P02-E1119: frozen source uses A before introducing the surrounding ring.
$f \in A$를 한 환의 원소라 하고 $J = \bigcap f^nA$로 두자.
$M$을 $f$에 관해 유도 완비인 $A$-가군이라 하자. 그러면
$M' = \Ker(M \to \lim M/f^nM)$일 때 $JM' = 0$이다. 특히 $A$가
유도 완비이면 $J$는 제곱이 영인 아이디얼이다.
\end{lemma}

\begin{proof}
$x \in M'$, $g \in J$를 잡자. 모든 $n \geq 1$에 대해
$x = f^n x_n$으로 쓸 수 있다. $g$가 $f^nA$에 속하므로 $M'$의 원소
$y_n = gx_n$은 $x_n$의 선택과 무관하다. 특히
$x_n = fx_{n + 1}$로 택할 수 있고, 그러면 $y_n = fy_{n + 1}$이다.
따라서 $1/f^n$을 $y_n$으로 보내는 사상 $A_f \to M$을 얻는다.
$M$이 유도 완비이므로 이 사상은 영이어야 한다(보조정리
% P02-E0435: frozen source lacks punctuation before the new consequent after this reference.
\ref{more-algebra-lemma-hom-from-Af}). 따라서 모든 $n$에 대해 $y_n = 0$이다.
$gx = gfx_1 = fy_1$이므로 증명이 끝난다.
\end{proof}

\begin{lemma}
\label{more-algebra-lemma-derived-complete-henselian}
$A$를 아이디얼 $I$에 관해 유도 완비인 환이라 하자. 그러면
$(A, I)$는 Hensel 쌍이다.
\end{lemma}

\begin{proof}
$f \in I$라 하자. 보조정리 \ref{more-algebra-lemma-largest-ideal-henselian}에 의해
$(A, fA)$가 Hensel 쌍임을 보이면 충분하다. $A$가 $fA$에 관해 유도
완비임에 유의하라(정의 \ref{more-algebra-definition-derived-complete}에서 바로
따른다). 보조정리 \ref{more-algebra-lemma-complete-derived-complete}에 의해
$A$에서 $A$의 $f$-진 완비화 $A'$으로 가는 사상은 전사이다.
보조정리 \ref{more-algebra-lemma-complete-henselian}에 의해 쌍 $(A', fA')$은
Hensel이다. 따라서 보조정리 \ref{more-algebra-lemma-henselian-henselian-pair}에 의해
$(A, \bigcap f^nA)$가 Hensel 쌍임을 보이면 충분하다. 이는 보조정리들
\ref{more-algebra-lemma-kernel-to-completion-square-zero}와
\ref{more-algebra-lemma-locally-nilpotent-henselian}에서 따른다.
\end{proof}

\begin{lemma}
\label{more-algebra-lemma-kernel-to-completion-nilpotent}
$A$를 아이디얼 $I$에 관해 유도 완비인 환이라 하고
$J = \bigcap I^n$으로 두자. $I$가 $r$개의 원소로 생성될 수 있으면
$N = 2^r$일 때 $J^N = 0$이다.
\end{lemma}

\begin{proof}
$r = 1$일 때는 보조정리
\ref{more-algebra-lemma-kernel-to-completion-square-zero}이다.
$r > 1$이고 $I = (f_1, \ldots, f_r)$이라 하자. 보조정리
\ref{more-algebra-lemma-serre-subcategory}에 의해 환 $A_t = A/f_r^tA$는 $I$에
관해 유도 완비이고, 따라서 더구나
$I_t = (f_1, \ldots, f_{r - 1})A_t$에 관해 유도 완비이다. 사상
$A \to A_t$가 $J$를 $J_t = \bigcap I_t^n$ 안으로 보냄에 유의하라.
귀납법에 의해 $N = 2^r$일 때 $J_t^{N/2} = 0$이다. 아이디얼
$\bigcap \Ker(A \to A_t) = \bigcap f_r^t A$는 $r = 1$인 경우에 의해
제곱이 영이다. 이로써 증명이 끝난다.
\end{proof}

\begin{lemma}
\label{more-algebra-lemma-reduced-derived-complete-complete}
$A$를 유한 생성 아이디얼 $I$에 관해 유도 완비인 기약환이라 하자.
그러면 $A$는 $I$-진 완비이다.
\end{lemma}

\begin{proof}
% P02-E0436: frozen source proof is the sentence fragment “Follows from ...”.
이는 보조정리 \ref{more-algebra-lemma-kernel-to-completion-nilpotent}와 명제
\ref{more-algebra-proposition-derived-complete-modules}에서 따른다.
\end{proof}
\section{Noether 환에 대한 유도 완비화}
\label{more-algebra-section-derived-completion-noetherian}

\noindent
$A$를 환이라 하고 $I \subset A$를 아이디얼이라 하자. 임의의
$K \in D(A)$에 대해 유도 극한
$$
K' = R\lim (K \otimes_A^\mathbf{L} A/I^n)
$$
을 생각할 수 있다. 이것은 $K$에 대한 함자이다. 주석
\ref{more-algebra-remark-constructing-tensor-with-limits-functorially}을 보라.
사상들의 계 $A \to A/I^n$은 사상 $K \to K'$을 유도하고, $K'$은
$I$에 관해 유도 완비이다(보조정리
\ref{more-algebra-lemma-naive-derived-completion}). 이 ``순진한'' 유도 완비화 구성은
일반적으로 보조정리 \ref{more-algebra-lemma-derived-completion}의 수반과 일치하지
않는다. 예를 들어 둘째 직합 인자가 제곱 영 아이디얼인
$A = \mathbf{Z}_p \oplus \mathbf{Q}_p/\mathbf{Z}_p$와 $K = A[0]$ 및
$I = (p)$를 취하면, 순진한 유도 완비화는 $\mathbf{Z}_p[0]$을 주지만
보조정리 \ref{more-algebra-lemma-derived-completion}의 구성은
$K^\wedge \cong \mathbf{Z}_p[1] \oplus \mathbf{Z}_p[0]$을 준다(계산은
생략한다). 보조정리 \ref{more-algebra-lemma-when-does-it-work}은 $I$가 한 원소로
생성되는 경우 두 함자가 일치하는 때를 특징짓는다.

\medskip\noindent
% P02-E0437: frozen source says "is the show" instead of "is to show".
이 절의 주된 목표는 $A$가 Noether 환이면 순진한 유도 완비화가 유도
완비화와 같음을 보이는 것이다.

\begin{lemma}
\label{more-algebra-lemma-sequence-Koszul-complexes}
% P02-E1123: frozen source begins with the sentence fragment "In Situation ... .".
상황 \ref{more-algebra-situation-koszul}에서 $A$가 Noether 환이면, $D(A)$의
pro-대상들 $\{K_n^\bullet\}$과
$\{A/(f_1^n, \ldots, f_r^n)\}$은 동형이다\footnote{특히 모든
$n$에 대해, $K_m^\bullet \to K_n^\bullet$이 사상
$K_m^\bullet \to A/(f_1^m, \ldots, f_r^m)$을 통해 분해되도록 하는
$m \geq n$이 존재한다.}.
\end{lemma}

\begin{proof}
다음과 같은 특수 삼각형들의 역계가 있다.
% P02-E0438: frozen source uses m in this quotient although the surrounding triangle is indexed by n.
$$
\tau_{\leq -1}K_n^\bullet \to K_n^\bullet \to A/(f_1^m, \ldots, f_r^m) \to
(\tau_{\leq -1}K_n^\bullet)[1]
$$
Derived Categories, 주석
\ref{derived-remark-truncation-distinguished-triangle}을 보라. Derived
Categories, 보조정리 \ref{derived-lemma-pro-isomorphism}에 의해 역계
$\tau_{\leq -1}K_n^\bullet$이 pro-영임을 보이면 충분하다.
$K_n^\bullet$은 $-r \leq i \leq 0$인 차수 $i$에서만 영이 아닌 항을
가짐을 상기하라. 따라서 Derived Categories, 보조정리
\ref{derived-lemma-essentially-constant-cohomology}에 의해 $p \leq -1$에
대해 $H^p(K_n^\bullet)$이 pro-영임을 보이면 충분하다. 다시 말해 모든
$n \in \mathbf{N}$에 대해, $H^p(K_m^\bullet) \to H^p(K_n^\bullet)$이
영이 되도록 하는 $m \geq n$이 존재함을 보여야 한다. $A$가 Noether
환이므로
$$
H^p(K_n^\bullet) =
\frac{\Ker(K_n^p \to K_n^{p + 1})}{\Im(K_n^{p - 1} \to K_n^p)}
$$
은 유한 $A$-가군이다. 더욱이 $p < 0$일 때 사상
$K_m^p \to K_n^p$은 그 성분들이 아이디얼
$(f_1^{m - n}, \ldots, f_r^{m - n})$에 속하는 대각행렬로 주어진다.
$H^p(K_n^\bullet)$은 $J = (f_1^n, \ldots, f_r^n)$에 의해 소멸됨에
유의하라. 보조정리 \ref{more-algebra-lemma-homotopy-koszul}을 보라. 이제
$m - n \geq tn$이면
$(f_1^{m - n}, \ldots, f_r^{m - n}) \subset J^t$이다.
% P02-E0439: frozen source runs the Artin--Rees data into its conclusion without a clause boundary.
따라서 아이디얼 $J$, 가군 $M = K_n^p$, 부분가군
$N = \Ker(K_n^p \to K_n^{p + 1})$에 Algebra, 보조정리
\ref{algebra-lemma-Artin-Rees}(Artin--Rees)를 적용하면, 충분히 큰
$m$에 대해 $K_m^p \to K_n^p$의 상과
$\Ker(K_n^p \to K_n^{p + 1})$의 교집합은
$J \Ker(K_n^p \to K_n^{p + 1})$에 포함된다. 그러한 $m$에 대해 영사상을
얻는다.
\end{proof}

\begin{proposition}
\label{more-algebra-proposition-noetherian-naive-completion-is-completion}
$A$를 Noether 환이라 하고 $I \subset A$를 아이디얼이라 하자.
$K \in D(A)$를 유도 극한
$K' = R\lim( K \otimes_A^\mathbf{L} A/I^n )$으로 보내는 함자는 보조정리
\ref{more-algebra-lemma-derived-completion}에서 구성한 포함 함자
$D_{comp}(A) \to D(A)$의 왼쪽 수반이다.
\end{proposition}

\begin{proof}
$(f_1, \ldots, f_r) = I$라 하고 $K_n^\bullet$을
$f_1^n, \ldots, f_r^n$에 대한 Koszul 복합체라 하자. 보조정리
\ref{more-algebra-lemma-derived-completion-koszul}에 의해 다음을 증명하면 충분하다.
$$
R\lim (K \otimes_A^\mathbf{L} K_n^\bullet) =
R\lim (K \otimes_A^\mathbf{L} A/(f_1^n, \ldots, f_r^n) ) =
R\lim (K \otimes_A^\mathbf{L} A/I^n ).
$$
보조정리 \ref{more-algebra-lemma-sequence-Koszul-complexes}에 의해 $D(A)$의
pro-대상들 $\{K_n^\bullet\}$과
$\{A/(f_1^n, \ldots, f_r^n)\}$은 동형이다. pro-대상들
$\{A/(f_1^n, \ldots, f_r^n)\}$과 $\{A/I^n\}$이 동형임은 명백하다.
따라서 보조정리 \ref{more-algebra-lemma-tensor-Rlim-pro-equal}에 의해 왼쪽에서
오른쪽으로 가는 사상은 동형이다.
\end{proof}

\begin{lemma}
\label{more-algebra-lemma-noetherian-calculate}
$I$를 Noether 환 $A$의 아이디얼이라 하자. $M$을 $A$-가군이라 하고
그 유도 완비화를 $M^\wedge$이라 하자. 그러면 짧은 완전열
$$
0 \to R^1\lim \text{Tor}_{i + 1}^A(M, A/I^n) \to
H^{-i}(M^\wedge) \to \lim \text{Tor}_i^A(M, A/I^n) \to 0
$$
이 있다. $M \in D^-(A)$에 대해서도 비슷한 결과가 성립한다.
\end{lemma}

\begin{proof}
% P02-E0440: frozen source proof is the sentence fragment "Immediate consequence of ...".
이는 명제 \ref{more-algebra-proposition-noetherian-naive-completion-is-completion}과
보조정리 \ref{more-algebra-lemma-break-long-exact-sequence-modules}의 즉각적인
귀결이다.
\end{proof}

\noindent
위 명제의 응용으로 Noether 경우의 유사연접 복합체에 대한 유도
완비화를 식별한다.

\begin{lemma}
\label{more-algebra-lemma-derived-completion-pseudo-coherent}
$A$를 Noether 환이라 하고 $I \subset A$를 아이디얼이라 하자. 모든
% P02-E0441: frozen source omits "is" after H^n(K).
$n \in \mathbf{Z}$에 대해 $H^n(K)$가 유한 $A$-가군인 $D(A)$의 대상
$K$를 잡자. 그러면 유도 완비화의 코호몰로지 가군 $H^n(K^\wedge)$은
코호몰로지 가군 $H^n(K)$의 $I$-진 완비화이다.
\end{lemma}

\begin{proof}
보조정리 \ref{more-algebra-lemma-Noetherian-pseudo-coherent}에 의해 모든 $m$에 대해
복합체 $\tau_{\leq m}K$는 유사연접이다. 따라서
$\tau_{\leq m}K$는 유한 자유 $A$-가군들의 위로 유계인 복합체
$P^\bullet$로 표현된다. 그러면
$\tau_{\leq m}K \otimes_A^\mathbf{L} A/I^n = P^\bullet/I^nP^\bullet$이다.
그러므로
$(\tau_{\leq m}K)^\wedge = R\lim P^\bullet/I^nP^\bullet$이다(명제
\ref{more-algebra-proposition-noetherian-naive-completion-is-completion}). 또한
$R\lim$은 단지 항별 $\lim$으로 주어지고(보조정리
\ref{more-algebra-lemma-compute-Rlim-modules}), $I$-진 완비화는 유한 $A$-가군들에
대한 완전 함자이므로(Algebra, 보조정리
\ref{algebra-lemma-completion-flat}), 결과가 $\tau_{\leq m}K$에 대해
성립한다고 결론짓는다. 따라서 유도 완비화가 유한 코호몰로지 차원을
가지므로 결과는 $K$에 대해 성립한다. 보조정리
\ref{more-algebra-lemma-derived-completion-finite-cohomological-dimension}을 보라.
\end{proof}

\begin{lemma}
\label{more-algebra-lemma-derived-complete-finite}
$I$를 Noether 환 $A$의 아이디얼이라 하자. $M$을 유도 완비인
$A$-가군이라 하자. $M/IM$이 유한 $A/I$-가군이면
$M = \lim M/I^nM$이고 $M$은 유한 $A^\wedge$-가군이다.
\end{lemma}

\begin{proof}
$M/IM$이 유한이라 하자. $M/IM$의 생성원들로 가는 원소들
$x_1, \ldots, x_t \in M$을 택하자. 제$i$ 기저벡터를 $x_i$로 보내는
사상 $A^{\oplus t} \to M$을 얻는다. 명제
\ref{more-algebra-proposition-noetherian-naive-completion-is-completion}에 의해 $A$의
유도 완비화는 $A^\wedge = \lim A/I^n$이다. $M$이 유도 완비이므로
우리의 사상은 사상 $q : (A^\wedge)^{\oplus t} \to M$을 통해
분해된다. 보조정리 \ref{more-algebra-lemma-derived-complete-zero}에 의해 가군
$\Coker(q)$은 영이다. 따라서 $M$은 유한 $A^\wedge$-가군이다.
$A^\wedge$은 Noether이고 $IA^\wedge$에 관해 완비이므로 $M$은
$I$-진 완비이다(Algebra, 보조정리들
\ref{algebra-lemma-completion-Noetherian},
\ref{algebra-lemma-when-finite-module-complete-over-complete-ring},
\ref{algebra-lemma-Artin-Rees}를 사용하라).
\end{proof}

\begin{lemma}
\label{more-algebra-lemma-when-derived-completion-is-completion}
$I$를 Noether 환 $A$의 아이디얼이라 하자.
\begin{enumerate}
\item $M$이 유한 $A$-가군이고 $N$이 평탄 $A$-가군이면,
$M \otimes_A N$의 유도 $I$-진 완비화는 $M \otimes_A N$의 보통
$I$-진 완비화이다.
\item $M$이 유한 $A$-가군이고 $f \in A$이면, $M_f$의 유도
$I$-진 완비화는 $M_f$의 보통 $I$-진 완비화이다.
\end{enumerate}
\end{lemma}

\begin{proof}
% P02-E0442: frozen source's notation sentence omits "by".
$A$-가군 $M$에 대해 $M^\wedge$으로 유도 완비화를 나타내고
$\lim M/I^nM$으로 보통 완비화를 나타내자. $M$이 유한이라 하자.
$i > 0$에 대해 계 $\text{Tor}^A_i(M, A/I^n)$은 pro-영이다. 보조정리
\ref{more-algebra-lemma-tor-strictly-pro-zero}를 보라. $N$이 평탄이므로
$\text{Tor}_i^A(M \otimes_A N, A/I^n) =
\text{Tor}_i^A(M, A/I^n) \otimes_A N$이고, 계
$\text{Tor}^A_i(M \otimes_A N, A/I^n)$도 마찬가지이다. 보조정리
\ref{more-algebra-lemma-noetherian-calculate}에 의해 다음 유도 극한은 오직 차수
$0$에서만 코호몰로지를 가지며, 그 코호몰로지는 보통 완비화
% P02-E1127: frozen source leaves the n-dependent tensor factor outside Rlim's parentheses and leaves the quotient ungrouped.
$R\lim (M \otimes_A N) \otimes_A^\mathbf{L} A/I^n$
$\lim M \otimes_A N/ I^n(M \otimes_A N)$으로 주어짐을 얻는다. 이것은
(1)을 증명한다. (2)는 (1)과 $M_f = M \otimes_A A_f$에서 따른다.
\end{proof}

\begin{lemma}
\label{more-algebra-lemma-derived-completion-tensor-finite}
$I$를 Noether 환 $A$의 아이디얼이라 하자. ${}^\wedge$이 $I$에 대한
유도 완비화를 나타낸다고 하자. $K \in D^-(A)$라 하자.
\begin{enumerate}
\item $M$이 유한 $A$-가군이면
$(K \otimes_A^\mathbf{L} M)^\wedge = K^\wedge \otimes_A^\mathbf{L} M$이다.
\item $L \in D(A)$가 유사연접이면
$(K \otimes_A^\mathbf{L} L)^\wedge = K^\wedge \otimes_A^\mathbf{L} L$이다.
\end{enumerate}
\end{lemma}

\begin{proof}
$L$을 (2)와 같다고 하자. $K$를 자유 $A$-가군들의 위로 유계인
복합체 $P^\bullet$로 표현할 수 있다. $L$을 유한 자유 $A$-가군들의
위로 유계인 복합체 $F^\bullet$로 표현할 수 있다.
$\text{Tot}(P^\bullet \otimes_A F^\bullet)$이
$K \otimes_A^\mathbf{L} L$을 표현하므로
$(K \otimes_A^\mathbf{L} L)^\wedge$은 다음 복합체로 표현된다.
$$
\text{Tot}((P^\bullet)^\wedge \otimes_A F^\bullet)
$$
여기서 $(P^\bullet)^\wedge$은 그 항들이 보통의
% P02-E0443: frozen source contains the stray standalone equality span in "usual = derived completions".
$=$ 유도 완비화들 $(P^n)^\wedge$인 복합체이다. 예를 들어 명제
\ref{more-algebra-proposition-noetherian-naive-completion-is-completion}과 보조정리
\ref{more-algebra-lemma-when-derived-completion-is-completion}을 보라. 이것은 (2)를
증명한다. (1)은 (2)의 특수한 경우이다.
\end{proof}
\section{Berthelot과 Ogus가 도입한 작용소}
\label{more-algebra-section-eta}

\noindent
이 절에서는 \cite[Section 8]{Berthelot-Ogus}에서 도입되고
\cite[Section 6]{BMS}에서 일반화된 구성을 논의한다. 독자에게 이 개념을
논의하는 원 논문들을 살펴보기를 강력히 권한다.

\medskip\noindent
$A$를 환이라 하고 $f \in A$를 비영인자라 하자. $M$이 $A$-가군이면
% P02-E0444: frozen source omits "the" before "following are equivalent".
보조정리 \ref{more-algebra-lemma-torsion-free}에 의해 다음 조건들은 동치이다.
\begin{enumerate}
\item $f$는 $M$에서 비영인자이다.
\item $M[f] = 0$이다.
\item 모든 $n \geq 1$에 대해 $M[f^n] = 0$이다.
\item 사상 $M \to M_f$는 단사이다.
\end{enumerate}
% P02-E1129: frozen source systematically leaves torsion-free unhyphenated as an attributive compound.
이 동치인 조건들이 성립하면 이 절에서는 {\it $M$은 $f$-꼬임이 없다}고
말하겠다. 이 경우
% P02-E0445: frozen source's notation sentence omits "by".
$f^iM \subset M_f$로 $x \in M$인 $f^ix$ 꼴의 원소들로 이루어진
부분가군을 나타내자. 물론 $f^iM$은 $A$-가군으로서 $M$과 동형이다.
$M^\bullet$을 미분 $d^i : M^i \to M^{i + 1}$을 갖는, $f$-꼬임이 없는
$A$-가군들의 복합체라 하자. 이 경우 $\eta_fM^\bullet$을 다음 항들을
갖는 복합체로 정의한다.
$$
(\eta_fM)^i = \{x \in f^iM^i \mid d^i(x) \in f^{i + 1}M^{i + 1}\}
$$
미분은 $d^i$에 의해 유도된다. $\eta_fM^\bullet$도 $f$-꼬임이 없는
$A$-가군들의 복합체임을 관찰하라. $a^\bullet : M^\bullet \to N^\bullet$이
$f$-꼬임이 없는 $A$-가군들의 복합체들의 사상이면, 사상들
$f^iM^i \to f^iN^i$에 의해 유도되는 복합체의 사상
$$
\eta_fa^\bullet : \eta_fM^\bullet \longrightarrow \eta_fN^\bullet
$$
을 얻는다. 독자는 이로부터 $f$-꼬임이 없는 $A$-가군들의 복합체 범주
위의 자기함자를 얻음을 확인할 수 있다.
$a^\bullet, b^\bullet : M^\bullet \to N^\bullet$이 $f$-꼬임이 없는
$A$-가군들의 복합체들의 두 사상이고
$h = \{h^i : M^i \to N^{i - 1}\}$가 $a^\bullet$과 $b^\bullet$ 사이의
호모토피라 하자. $\eta_fh$를, $x$를 $h^i(x)$로 보내는 사상들의 족
$(\eta_fh)^i : (\eta_fM)^i \to (\eta_fN)^{i - 1}$로 정의한다.
$x \in f^iM^i$이면 $h^i(x) \in f^iN^{i - 1}$이고 이는 분명
$(\eta_fN)^{i - 1}$에 포함되므로 이 정의는 의미가 있다. 독자는
$\eta_fh$가 $\eta_fa^\bullet$과 $\eta_fb^\bullet$ 사이의 호모토피임을
확인할 수 있다. 종합하면, $f$-꼬임이 없는 $A$-가군들의 가법 범주의
호모토피 범주(Derived Categories, 절
\ref{derived-section-homotopy}) 위의 함자
$$
\eta_f :
K(f\text{-torsion free }A\text{-modules})
\longrightarrow
K(f\text{-torsion free }A\text{-modules})
$$
를 얻는다. $\eta_f$가 삼각 범주들의 완전 함자가 되는 의미는 없다.
예 \ref{more-algebra-example-eta-not-distinguished}를 보라.

\begin{example}
\label{more-algebra-example-eta-not-distinguished}
$A$를 환이라 하고 $f \in A$를 비영인자라 하자. 함자
$\eta_f : K(f\text{-torsion free }A\text{-modules})
\to K(f\text{-torsion free }A\text{-modules})$를 생각하자.
$M^\bullet$을 $f$-꼬임이 없는 $A$-가군들의 복합체라 하자. $f$를
곱하는 사상은 동형
$\eta_f(M^\bullet[1]) \to (\eta_fM^\bullet)[1]$을 정의하므로, 이
의미에서 $\eta_f$는 이동과 양립한다. 그러나 다음 도표를 생각하자.
$$
\xymatrix{
A \ar[r]_f &
A \ar[r]_1 &
A \ar[r] &
0 \\
0 \ar[r] \ar[u] &
0 \ar[r] \ar[u] &
A \ar[r]^{-1} \ar[u]^f &
A \ar[u]
}
$$
각 열을 $f$-꼬임이 없는 $A$-가군들의 복합체로 보되, 위의 가군은
차수 $1$에, 그 아래 가군은 차수 $0$에 있다고 하자. 그러면 이 도표는
Derived Categories, 절
\ref{derived-section-homotopy-triangulated}에서 주어진 삼각 구조를 갖는
$K(f\text{-torsion free }A\text{-modules})$에서 특수 삼각형을 준다.
즉, 셋째 복합체는 첫 두 복합체 사이 사상의 원뿔이다. 그러나 각 열에
$\eta_f$를 적용하면 다음을 얻는다.
$$
\xymatrix{
fA \ar[r]_f &
fA \ar[r]_1 &
fA \ar[r] &
0 \\
0 \ar[r] \ar[u] &
0 \ar[r] \ar[u] &
A \ar[r]^{-1} \ar[u]^f &
A \ar[u]
}
$$
그런데 셋째 복합체는 비순환이고 심지어 영과 호모토피 동치이다.
따라서 이것이 특수 삼각형이면 첫 화살표는 호모토피 범주에서 동형이어야
하지만, 이는 $f$가 단원인 경우가 아니면 참이 아니다.
\end{example}

\begin{lemma}
\label{more-algebra-lemma-eta-first-property}
$A$를 환이라 하고 $f \in A$를 비영인자라 하자. $M^\bullet$을
$f$-꼬임이 없는 $A$-가군들의 복합체라 하자. $f^i$를 곱하여 주어지는
표준 동형
$$
f^i : H^i(M^\bullet)/H^i(M^\bullet)[f] \longrightarrow H^i(\eta_fM^\bullet)
$$
이 있다.
\end{lemma}

\begin{proof}
$\Ker(d^i : (\eta_fM)^i \to (\eta_fM)^{i + 1})$은
$\Ker(d^i : f^iM^i \to f^iM^{i + 1}) = f^i\Ker(d^i : M^i \to M^{i + 1})$과
같음을 관찰하라.
% P02-E0446: frozen source says "This we get a surjection" rather than "Thus".
따라서 $z \in \Ker(d^i : M^i \to M^{i + 1})$의 류를 $f^iz$의 류로
보내는 전사 $f^i : H^i(M^\bullet) \to H^i(\eta_fM^\bullet)$을 얻는다.
$H^i(\eta_fM^\bullet)$에서 영류를 얻는다면, 어떤
$y \in M^{i - 1}$에 대해 $f^i z = d^{i - 1}(f^{i - 1}y)$임을 알 수
있다. 관련된 모든 가군에서 $f$가 비영인자이므로 이는
$f z = d^{i - 1}(y)$를 뜻하고, 이는 정확히 $z$의 류가 $f$-꼬임임을
뜻한다. 이것이 원하는 바이다.
\end{proof}

\begin{lemma}
\label{more-algebra-lemma-eta-qis}
$A$를 환이라 하고 $f \in A$를 비영인자라 하자. $M^\bullet \to N^\bullet$이
$f$-꼬임이 없는 $A$-가군들의 복합체들의 준동형사상이면, 유도된 사상
$\eta_fM^\bullet \to \eta_fN^\bullet$도 준동형사상이다.
\end{lemma}

\begin{proof}
이는 보조정리 \ref{more-algebra-lemma-eta-first-property}의 동형들이 복합체의 사상과
양립하기 때문이다.
\end{proof}

\begin{lemma}
\label{more-algebra-lemma-Leta}
$A$를 환이라 하고 $f \in A$를 비영인자라 하자. 가법 함자
\footnote{이 함자는 완전하지 않음, 즉 특수 삼각형들을 특수 삼각형들로
보내지 않음에 주의하라. 예 \ref{more-algebra-example-eta-not-distinguished}를 보라.}
$L\eta_f : D(A) \to D(A)$가 존재하여, $M \in D(A)$가 $f$-꼬임이 없는
$A$-가군들의 복합체 $M^\bullet$로 표현되면
$L\eta_fM = \eta_fM^\bullet$이고 사상들에 대해서도 마찬가지이다.
\end{lemma}

\begin{proof}
% P02-E0447: frozen source's notation sentence omits "by".
$\mathcal{T} \subset \text{Mod}_A$로 $f$-꼬임이 없는 $A$-가군들의 충만한
부분범주를 나타내자. $K(\mathcal{T})$를 $K(A)$의 충만한 삼각
부분범주로 포함하는 대응하는 포함
$$
K(\mathcal{T}) \quad\subset\quad K(\text{Mod}_A) = K(A)
$$
이 있다. $S \subset \text{Arrows}(K(\mathcal{T}))$를 준동형사상들의
모임이라 하자. 사상
$$
S^{-1}K(\mathcal{T}) \longrightarrow D(A)
$$
이 삼각 범주들의 동치임을 보이기 위해 Derived Categories, 보조정리
\ref{derived-lemma-localization-subcategory}를 적용하겠다. 그 보조정리에
의해 다음을 증명하면 충분하다. $A$-가군들의 복합체 $M^\bullet$이
주어지면, $K^\bullet$이 $f$-꼬임이 없는 가군들의 복합체가 되도록 하는
준동형사상 $K^\bullet \to M^\bullet$이 존재한다. 보조정리
\ref{more-algebra-lemma-K-flat-resolution}에 의해 복합체 $K^\bullet$이 K-평탄이고
(이 성질은 사용하지 않을 것이다) 평탄 $A$-가군들 $K^i$로 이루어지도록
하는 준동형사상 $K^\bullet \to M^\bullet$을 찾을 수 있다. 특히 원하는
대로 모든 $i$에 대해 $f$는 $K^i$에서 비영인자이다.

\medskip\noindent
이 예비 작업을 마쳤으므로 $L\eta_f$를 정의할 수 있다. 즉, 이 절의
시작 부분의 논의로부터 잘 정의된 함자
$$
K(\mathcal{T}) \xrightarrow{\eta_f} K(\mathcal{T}) \to K(A) \to D(A)
$$
를 이미 얻었고, 보조정리 \ref{more-algebra-lemma-eta-qis}에 따라 이 함자는
준동형사상들을 준동형사상들로 보낸다. 따라서
% P02-E0448: frozen source says "factors over" instead of "factors through".
Categories, 보조정리 \ref{categories-lemma-properties-left-localization}에
의해 이 함자는 $S^{-1}K(\mathcal{T}) = D(A)$를 통해 분해된다.
\end{proof}

\begin{remark}
\label{more-algebra-remark-eta-BZ}
$A$를 환이라 하고 $f \in A$를 비영인자라 하자. $M^\bullet$을
$f$-꼬임이 없는 $A$-가군들의 복합체라 하자. 모든 $i$에 대해
$\overline{M}^i = M^i/fM^i$로 두자.
% P02-E0449: frozen source's notation sentence omits "by".
$B^i \subset Z^i \subset \overline{M}^i$로 복합체
$\overline{M}^\bullet = M^\bullet \otimes_A A/fA$의 미분에 대한 경계와
코사이클을 나타내자. 완전한 행들을 갖는 가환 도표
$$
\xymatrix{
0 \ar[r] &
B^{i + 1} \ar[r] \ar@{=}[d] &
B^{i + 1} \oplus B^i \ar[r] \ar[d]^{s, s'} &
B^i \ar[r] \ar[d] & 0 \\
0 \ar[r] &
B^{i + 1} \ar[r]^-s &
(\eta_fM)^i /f(\eta_fM)^i \ar[r]^-t &
Z^i \ar[r] & 0
}
$$
가 존재한다고 주장한다. 사상들의 구성은 다음과 같다.
\begin{enumerate}
\item $x \in (\eta_fM)^i$이면 $d^i(x') = 0$이
$\overline{M}^{i + 1}$에서 성립하는 $x = f^ix'$으로 쓸 수 있다. 따라서
$x$를 $x'$의 류로 보내어 사상 $t$를 정의할 수 있다.
\item $y \in M^{i + 1}$의 류 $\overline{y}$가
$B^{i + 1} \subset \overline{M}^{i + 1}$에 속하면, 어떤
$y' \in M^{i + 1}$과 $x \in M^i$에 대해
$y = fy' + d^i(x)$로 쓸 수 있다. 따라서 $\overline{y}$를
$(\eta_fM)^i /f(\eta_fM)^i$에서 $f^{i + 1}x$의 류로 보내는 사상 $s$를
정의할 수 있다. 이것이 잘 정의됨의 확인은 생략한다.
\item $x \in M^i$의 류 $\overline{x}$가
$B^i \subset \overline{M}^i$에 속하면, 어떤 $x' \in M^i$와
$z \in M^{i - 1}$에 대해 $x = fx' + d^{i - 1}(z)$로 쓸 수 있다.
$\overline{x}$를 $(\eta_fM)^i/f(\eta_fM)^i$에서
$f^i d^{i - 1}(z)$의 류로 보내어 사상 $s'$을 정의한다. 만일
$fx' + d^{i - 1}(z) = 0$이면 $f^ix'$은 $(\eta_fM)^i$에 속하고, 따라서
$f^id^{i - 1}(z)$는 $f(\eta_fM)^i$에 속하므로 이것은 잘 정의된다.
\end{enumerate}
위에 표시한 도표의 아래 행이 가군들의 짧은 완전열임의 확인은 생략한다.
이 구성들로부터 가환 도표들
$$
\xymatrix{
B^{i + 1} \oplus B^i \ar[d]^{s, s'} \ar[r] &
B^{i + 2} \oplus B^{i + 1} \ar[d]^{s, s'} \\
(\eta_fM)^i /f(\eta_fM)^i \ar[r] &
(\eta_fM)^{i + 1} /f(\eta_fM)^{i + 1}
}
$$
을 얻음은 즉시 명백하다. 여기서 위 수평 화살표는 정의역과 공역의
직합 인자들 $B^{i + 1}$을 동일시하여 주어진다. 다시 말해,
우리는
$\eta_fM^\bullet / f(\eta_fM^\bullet) = \eta_fM^\bullet \otimes_A A/fA$의
비순환 부분복합체를 찾았고, 이 부분복합체로 나눈 몫은 그 항들
$Z^i/B^i$가 복합체
$\overline{M}^\bullet = M^\bullet \otimes_A A/fA$의 코호몰로지 가군들인
복합체이다.
\end{remark}

\noindent
주석 \ref{more-algebra-remark-eta-BZ}에서 관찰한 현상을 더 표준적인 방식으로
설명하기 위해 Bockstein 작용소들을 구성하겠다. $A$를 환이라 하고
$f \in A$를 비영인자라 하자. $M^\bullet$을 $f$-꼬임이 없는
$A$-가군들의 복합체라 하자. 모든 $i \in \mathbf{Z}$에 대해 다음
가환 도표가 있다(텐서곱은 $A$ 위에서 취한다).
$$
\xymatrix{
0 \ar[r] &
M^\bullet \otimes f^{i + 1}A \ar[r] \ar[d] &
M^\bullet \otimes f^iA \ar[r] \ar[d] &
M^\bullet \otimes f^iA/f^{i + 1}A \ar[r] \ar@{=}[d] &
0 \\
0 \ar[r] &
M^\bullet \otimes f^{i + 1}A/f^{i + 2}A \ar[r] &
M^\bullet \otimes f^iA/f^{i + 2}A \ar[r] &
M^\bullet \otimes f^iA/f^{i + 1}A \ar[r] &
0
}
$$
그 행들은 복합체들의 짧은 완전열이다. 물론 서로 다른 $i$에 대한 이
짧은 완전열들은 $f$의 적당한 거듭제곱을 곱함으로써 모두 서로
동형이다. 특히 아래 열의 코호몰로지 긴 완전열은 모든
$i \in \mathbf{Z}$에 대해 {\it Bockstein 작용소}
$$
\beta = \beta^i : H^i(M^\bullet \otimes f^iA/f^{i + 1}A) \to
H^{i + 1}(M^\bullet \otimes f^{i + 1}A/f^{i + 2}A)
$$
를 결정한다. 나중에 사용하기 위해, 위 가환 도표에 의해 Bockstein
작용소가 다음과 같이 분해됨을 기록한다.
\begin{equation}
\label{more-algebra-equation-factorization-bockstein}
\vcenter{
\xymatrix{
H^i(M^\bullet \otimes f^iA/f^{i + 1}A)
\ar[r]_\delta \ar[rd]_\beta &
H^{i + 1}(M^\bullet \otimes f^{i + 1}A) \ar[d] \\
&
H^{i + 1}(M^\bullet \otimes f^{i + 1}A/f^{i + 2}A)
}
}
\end{equation}
여기서 $\delta$는 위 가환 도표의 위 행에서 오는 경계 작용소이다.
이제 복합체
\begin{equation}
\label{more-algebra-equation-complex-bocksteins}
H^\bullet(M^\bullet/f) =
\left[
\begin{matrix}
\ldots \\
\downarrow \\
H^{i - 1}(M^\bullet \otimes f^{i - 1}A/f^iA) \\
\downarrow \beta \\
H^i(M^\bullet \otimes f^iA/f^{i + 1}A) \\
\downarrow \beta \\
H^{i + 1}(M^\bullet \otimes f^{i + 1}A/f^{i + 2}A) \\
\downarrow \\
\ldots
\end{matrix}
\right]
\end{equation}
를 얻음, 즉 $\beta \circ \beta = 0$임을 보이자\footnote{다른 방법은
$f^iM^\bullet$이라는 부분복합체들로 여과된 복합체 $(M^\bullet)_f$에
대한 스펙트럼 열의 미분으로 $\beta$가 나타난다고 논증하는 것이다.
더 강한 결과를 증명하는 또 다른 논증은 먼저 $M^\bullet = A$인 경우를
생각하는 것이다. 여기서 짧은 완전열들
$0 \to f^{i + 1}A/f^{i + 2}A \to f^iA/f^{i + 2}A \to f^iA/f^{i + 1}A \to 0$은
$D(A)$에서 사상들
$\beta^i : f^iA/f^{i + 1}A \to f^{i + 1}A/f^{i + 2}A[1]$을 정의한다.
그러면 본문과 비슷하게 논증하여 합성
$f^iA/f^{i + 1}A \to f^{i + 1}A/f^{i + 2}A[1] \to f^{i + 2}A/f^{i + 3}A[2]$이
$D(A)$에서 영임을 계산한다. $M^\bullet$의 항들이 $f$-꼬임이 없다는
가정에 의해
$M^\bullet \otimes f^iA/f^{i + 1}A =
M^\bullet \otimes^\mathbf{L} f^iA/f^{i + 1}A$이므로,
% P02-E0450: frozen source misspells "composition" as "compostion".
합성
$$
(M^\bullet \otimes f^iA/f^{i + 1}A)
\to
(M^\bullet \otimes f^{i + 1}A/f^{i + 2}A)[1]
\to
(M^\bullet \otimes f^{i + 2}A/f^{i + 3}A)[2]
$$
도 $D(A)$에서 영이라고 결론짓는다.}. 즉, 분해
(\ref{more-algebra-equation-factorization-bockstein})를 사용하면 사상
$$
H^{i + 1}(M^\bullet \otimes f^{i + 1}A)
\to
H^{i + 1}(M^\bullet \otimes f^{i + 1}A/f^{i + 2}A)
\xrightarrow{\beta^{i + 1}}
H^{i + 2}(M^\bullet \otimes f^{i + 2}A/f^{i + 3}A)
$$
이 영임을 보이면 충분하다. 이는 $\beta^{i + 1}$의 핵이
$H^{i + 1}(M^\bullet \otimes f^{i + 1}A/f^{i + 3}A)$로 올릴 수 있는
코호몰로지 류들로 이루어지고, 첫 사상의 상에 있는 것들은 분명 올릴
수 있기 때문이다!

\begin{lemma}
\label{more-algebra-lemma-eta-second-property}
$A$를 환이라 하고 $f \in A$를 비영인자라 하자. $M^\bullet$을
$f$-꼬임이 없는 $A$-가군들의 복합체라 하자. 복합체들의 표준 사상
$$
\eta_fM^\bullet \otimes_A A/fA
\longrightarrow
H^\bullet(M^\bullet/f)
$$
이 있으며, 이는 준동형사상이다. 여기서 오른쪽은 복합체
(\ref{more-algebra-equation-complex-bocksteins})이다.
\end{lemma}

\begin{proof}
$x \in (\eta_fM)^i$라 하자. 그러면
% P02-E0451: frozen source drops the degree superscript in f^iM here.
$x = f^ix' \in f^iM$이고
$d^i(x) = f^{i + 1}y \in f^{i + 1}M^{i + 1}$이다. 따라서 $d^i$는
$x' \otimes f^i$를 $M^{i + 1} \otimes f^iA/f^{i + 1}A$에서 영으로
보낸다. 이 증명에서 모든 텐서곱은 $A$ 위에서 취한다. 따라서 $x$를
$H^i(M^\bullet \otimes f^iA/f^{i + 1}A)$에서 $x' \otimes f^i$의 류로
보낼 수 있다. 이 규칙이 $A/fA$-가군들의 사상
$$
(\eta_fM)^i \otimes A/fA
\longrightarrow
H^i(M^\bullet \otimes f^iA/f^{i + 1}A)
$$
을 정의함은 명백하다. 위 상황에서 $x' \otimes f^i$를
$M^i \otimes f^iA/f^{i + 2}A$의 원소로 볼 수 있으며, 그 미분은
$d^i(x' \otimes f^i) = y \otimes f^{i + 1}$임을 관찰하라. 위의
$\beta$의 구성에 의해
$\beta(x' \otimes f^i) = y \otimes f^{i + 1}$임을 얻는다. 따라서
우리의 사상들이 미분과 양립함, 즉 복합체들의 사상이 있음을 얻는다.

\medskip\noindent
증명을 끝내기 위해, 앞 문단에서 주어진 구성이 주석
\ref{more-algebra-remark-eta-BZ}에서 논의한 사상들
$(\eta_fM)^i \otimes A/fA \to Z^i/B^i$과 일치함을 관찰한다. 이
사상들의 핵이 $\eta_fM^\bullet \otimes A/fA$의 비순환 부분복합체임을
보았으므로 보조정리가 증명되었다.
\end{proof}

\begin{lemma}
\label{more-algebra-lemma-vanishing-beta}
$A$를 환이라 하고 $f \in A$를 비영인자라 하자. $M^\bullet$을
$f$-꼬임이 없는 $A$-가군들의 복합체라 하자. $i \in \mathbf{Z}$에
대해 다음은 동치이다.
\begin{enumerate}
\item $\Ker(d^i \bmod f^2)$는 $\Ker(d^i \bmod f)$ 위로 전사한다.
\item $\beta : H^i(M^\bullet \otimes_A f^iA/f^{i + 1}A) \to
H^{i + 1}(M^\bullet \otimes_A f^{i + 1}A/f^{i + 2}A)$는 영이다.
\end{enumerate}
조건 $H^{i + 1}(M^\bullet)[f] = 0$은 이 동치인 조건들을 함의한다.
\end{lemma}

\begin{proof}
(1)과 (2)의 동치는 $\beta$가 복합체들의 짧은 완전열
$0 \to fM^\bullet/f^2M^\bullet \to M^\bullet/f^2M^\bullet \to
M^\bullet/fM^\bullet \to 0$과 동형인 복합체들의 짧은 완전열의
코호몰로지 위 경계 사상으로 정의된다는 사실에서 따른다.
$\beta \not = 0$이면 분해
(\ref{more-algebra-equation-factorization-bockstein}) 때문에
$H^{i + 1}(M^\bullet)[f] \not = 0$이다.
\end{proof}

\begin{lemma}
\label{more-algebra-lemma-eta-vanishing-beta}
$A$를 환이라 하고 $f \in A$를 비영인자라 하자. $M^\bullet$을
$f$-꼬임이 없는 $A$-가군들의 복합체라 하자.
$\Ker(d^i \bmod f^2)$가 $\Ker(d^i \bmod f)$ 위로 전사하면 표준 사상
$$
(1, d^i) :
(\eta_fM)^i / f(\eta_fM)^i \longrightarrow
f^iM^i/f^{i + 1}M^i \oplus f^{i + 1}M^{i + 1}/f^{i + 2}M^{i + 1}
$$
은 왼쪽을 오른쪽의 부분가군들의 직합과 동일시한다.
\end{lemma}

\begin{proof}
주석 \ref{more-algebra-remark-eta-BZ}의 표기를 사용하여 사상
$t^{-1} : Z^i \to (\eta_fM)^i / f(\eta_fM)^i$을 정의한다. 즉,
$d^i(x) = f^2y$인 $x \in M^i$에 대해 $Z^i$에서 $x$의 류를
$(\eta_fM)^i / f(\eta_fM)^i$에서 $f^ix$의 류로 보낸다. 이것이 잘
정의됨의 확인은 생략한다. 보조정리의 가정은 이 연산의 정의역이 정확히
$Z^i$ 전체임을 뜻한다. 그러면
$t \circ t^{-1} = \text{id}_{Z^i}$이다. 따라서 $t^{-1}$은 주석
\ref{more-algebra-remark-eta-BZ}의 짧은 완전열의 분할을 정의하고, 그로부터 얻는
직합 분해
$$
(\eta_fM)^i / f(\eta_fM)^i = Z^i \oplus B^{i + 1}
$$
는 보조정리에 표시된 사상과 양립한다.
\end{proof}

\begin{lemma}
\label{more-algebra-lemma-eta-third-property}
$A$를 환이라 하고 $f, g \in A$를 비영인자들이라 하자. $M^\bullet$을
모든 $M^i$에서 $fg$가 비영인자인 $A$-가군들의 복합체라 하자. 그러면
$\eta_f\eta_gM^\bullet = \eta_{fg}M^\bullet$이다.
\end{lemma}

\begin{proof}
이 명제는 차수 $i$에서 국소화 $M^i_{fg} = (M^i_g)_f$의 같은
부분가군을 얻는다는 뜻이다. 세부 사항은 생략한다.
\end{proof}

\begin{lemma}
\label{more-algebra-lemma-eta-flat-base-change}
$A$를 환이라 하고 $f \in A$를 비영인자라 하자. $A \to B$를 평탄한
환 사상이라 하고,
% P02-E0452: frozen source omits "be" in "let g ... the image of f".
$g \in B$를 $f$의 상이라 하자. $M^\bullet$을 $f$-꼬임이 없는
$A$-가군들의 복합체라 하자. 그러면 $g$는 비영인자이고,
$M^\bullet \otimes_A B$는 $g$-꼬임이 없는 가군들의 복합체이며,
$\eta_fM^\bullet \otimes_A B = \eta_g(M^\bullet \otimes_A B)$이다.
\end{lemma}

\begin{proof}
생략한다.
\end{proof}
\section{퍼펙트 복합체와 eta 작용소}
\label{more-algebra-section-perfect-eta}

\noindent
이 절에서는 Macpherson의 그래프 구성을 우리의 형태로 전개하기 위한
대수학을 준비한다. More on Flatness, 절
\ref{flat-section-blowup-complexes-III}을 보라. 절 \ref{more-algebra-section-eta}에서
도입한 $\eta_f$ 작용소를 사용할 것이다.

\medskip\noindent
$A$를 환이라 하고 $f \in A$를 비영인자라 하자. $M^\bullet$을 유한
자유 $A$-가군들의 유계 복합체라 하자. 각 $i$에 대해 $r_i$를 $M^i$의
계수라 하고 다음과 같이 두자.
$$
I_i(M^\bullet, f) = \text{ideal generated by the }
r_i \times r_i\text{-minors of }
(f, d^i) : M^i \to M^i \oplus M^{i + 1}
$$
$f^{r_i} \in I_i(M^\bullet, f)$임을 관찰하라.

\begin{lemma}
\label{more-algebra-lemma-ideal-well-defined}
$A$를 환이라 하고 $f \in A$를 비영인자라 하자. $M^\bullet$과
$N^\bullet$을 $D(A)$의 같은 대상을 표현하는 유한 자유 $A$-가군들의
두 유계 복합체라 하자. 그러면
$$
f^m I_i(M^\bullet, f) = f^n I_i(N^\bullet, f)
$$
이며, 여기서 $n, m \geq 0$은 다음을 만족하는 정수들이다.
$$
m + \sum\nolimits_{j \geq i} (-1)^{j - i}rk(M^j) =
n + \sum\nolimits_{j \geq i} (-1)^{j - i}rk(N^j)
$$
위 등식은 $A$의 아이디얼들의 등식이다.
\end{lemma}

\begin{proof}
$A$의 모든 소 아이디얼에서 국소화한 뒤 등식을 증명하면 충분하다.
% P02-E0453: frozen source does not delimit the omitted induction argument from the main clause.
따라서 보조정리 \ref{more-algebra-lemma-compare-representatives-perfect}와, 그 논증은
생략하는 귀납법에 의해, 어떤 자명한 복합체 $Q^\bullet$에 대해
$N^\bullet = M^\bullet \oplus Q^\bullet$이라 가정할 수 있다. 즉,
$$
Q^\bullet = \ldots \to 0 \to A \xrightarrow{1} A \to 0 \to \ldots
$$
% P02-E1136: frozen source uses singular "A" for the two displayed copies in degrees j and j+1.
여기서 $A$는 차수 $j$와 $j + 1$에 놓여 있다.
$j \not = i - 1, i, i + 1$이면 분명
$I_i(M^\bullet, f) = I_i(N^\bullet, f)$이고 $m = n$이므로 원하는
등식을 얻는다. $j = i + 1$이면 사상들
$$
(f, d^i) : M^i \to M^i \oplus M^{i + 1}
\quad\text{and}\quad
(f, d^i, 0) : M^i \to M^i \oplus M^{i + 1} \oplus A
$$
은 같은 영 아닌 소행렬식들을 가지므로 이 경우에도
$I_i(M^\bullet, f) = I_i(N^\bullet, f)$이고 $m = n$이다.
$j = i$이면 $I_i(M^\bullet, f)$는 사상
$$
(f, d^i) : M^i \to M^i \oplus M^{i + 1}
$$
의 $r_i \times r_i$ 소행렬식들이 생성하는 아이디얼이고,
$I_i(N^\bullet, f)$는 사상
$$
(f \oplus f, d^i \oplus 1) :
(M^i \oplus A) \to (M^i \oplus A) \oplus (M^{i + 1} \oplus A)
$$
의 $(r_i + 1) \times (r_i + 1)$ 소행렬식들이 생성하는 아이디얼이다.
% P02-E0454: frozen source repeatedly omits articles in the coordinate and matrix descriptions.
좌표를 적절히 택하면 둘째 사상의 행렬은 다음 블록 형태임을 알 수 있다.
$$
T =
\left(
\begin{matrix}
T_1 & 0 \\
0 & T_2
\end{matrix}
\right),
\quad
T_1 = \text{matrix of first map},
\quad
T_2 =
\left(
\begin{matrix}
f \\
1
\end{matrix}
\right)
$$
보조정리 \ref{more-algebra-lemma-ideals-generated-by-minors}의 표기를 쓰면
$I_0(T_2) = A$, $I_1(T_2) = A$, 그리고 $p \geq 2$에 대해
$I_p(T_2) = 0$이다. 따라서
$I_{r_i + 1}(T) = I_{r_i + 1}(T_1) + I_{r_i}(T_1) =
I_{r_i}(T_1)$이고, 이는
$I_i(M^\bullet, f) = I_i(N^\bullet, f)$이라는 뜻이다. 또한
$m = n$이므로 $j = i$인 경우가 끝났다. 마지막으로
$j = i - 1$이라 하자. 그러면 $m = n + 1$이므로
$fI_i(M^\bullet, f) = I_i(N^\bullet, f)$임을 보여야 한다. 이 경우
$I_i(M^\bullet, f)$는 사상
$$
(f, d^i) : M^i \to M^i \oplus M^{i + 1}
$$
의 $r_i \times r_i$ 소행렬식들이 생성하는 아이디얼이고,
$I_i(N^\bullet, f)$는 사상
$$
(f \oplus f, d^i) : (M^i \oplus A) \to (M^i \oplus A) \oplus M^{i + 1}
$$
의 $(r_i + 1) \times (r_i + 1)$ 소행렬식들이 생성하는 아이디얼이다.
좌표를 적절히 택하면 둘째 사상의 행렬은 다음 블록 형태임을 알 수 있다.
$$
T =
\left(
\begin{matrix}
T_1 & 0 \\
0 & T_2
\end{matrix}
\right),
\quad
T_1 = \text{matrix of first map},
\quad
T_2 =
\left(
\begin{matrix}
f
\end{matrix}
\right)
$$
위와 같이 논증하면 실제로
$fI_i(M^\bullet, f) = I_i(N^\bullet, f)$임을 얻는다.
\end{proof}

\begin{lemma}
\label{more-algebra-lemma-eta-change-unit}
$f \in A$를 환 $A$의 비영인자라 하고 $u \in A$를 단원이라 하자.
$M^\bullet$을 유한 자유 $A$-가군들의 유계 복합체라 하자. 그러면
$I_i(M^\bullet, f) = I_i(M^\bullet, uf)$이다.
\end{lemma}

\begin{proof}
생략한다.
\end{proof}

\begin{lemma}
\label{more-algebra-lemma-eta-base-change-pre}
$A \to B$를 환 사상이라 하고 $f \in A$를 비영인자라 하자.
$M^\bullet$을 유한 자유 $A$-가군들의 유계 복합체라 하자. $f$의
$B$에서의 상 $g$가 비영인자라고 가정하자. 그러면
$I_i(M^\bullet, f)B = I_i(M^\bullet \otimes_A B, g)$이다.
\end{lemma}

\begin{proof}
$(f, d^i) : M^i \to M^i \oplus M^{i + 1}$의 소행렬식들은
$(g, d^i) : M^i \otimes_A B \to M^i \otimes_A B \oplus M^{i + 1} \otimes_A B$의
대응하는 소행렬식들로 간다.
\end{proof}

\begin{lemma}
\label{more-algebra-lemma-eta-cohomology-free}
$A$를 환이라 하고 $\mathfrak p \subset A$를 소 아이디얼이라 하며
$f \in A$를 비영인자라 하자. $M^\bullet$을 유한 자유 $A$-가군들의
유계 복합체라 하자. 모든 $i$에 대해 $H^i(M^\bullet)_\mathfrak p$가
자유이면, 모든 $i$에 대해 $I_i(M^\bullet, f)_\mathfrak p$는 주
아이디얼이고 실제로 $f$의 거듭제곱으로 생성된다.
\end{lemma}

\begin{proof}
보조정리 \ref{more-algebra-lemma-eta-base-change-pre}에 의해 $A$가 극대 아이디얼
$\mathfrak p$를 갖는 국소환이라고 가정해도 된다. 보조정리
\ref{more-algebra-lemma-ideal-well-defined}에 의해 $M^\bullet$을 준동형인 복합체로
바꾸어도 된다. 코호몰로지 가군들이 자유라는 가정에 의해
$M^\bullet$은 차수 $i$의 항이 $H^i(M^\bullet)$이고 미분들이 영인
복합체와 준동형사상임을 알 수 있다. 예를 들어 Derived Categories,
보조정리 \ref{derived-lemma-ext-2-zero-pre}를 보라. 다시 말해 복합체
$M^\bullet$의 모든 미분이 영이라고 가정해도 된다. 이 경우
$I_i(M^\bullet, f) = (f^{r_i})$가 주 아이디얼임은 명백하다.
\end{proof}

\begin{lemma}
\label{more-algebra-lemma-eta-locally-free}
$A$를 환이라 하고 $f \in A$를 비영인자라 하자. $M^\bullet$을 유한
자유 $A$-가군들의 유계 복합체라 하자. $I_i(M^\bullet, f)$가 주
아이디얼이라고 가정하자. 그러면 $(\eta_fM)^i$는 계수 $r_i$인 국소
자유 가군이고, 사상
$(1, d^i) : (\eta_fM)^i \to f^iM^i \oplus f^{i + 1}M^{i + 1}$은 직합
인자의 포함이다.
\end{lemma}

\begin{proof}
$I_i(M^\bullet, f)$의 생성원 $g$를 택하자.
$f^{r_i} \in I_i(M^\bullet, f)$이므로 $g$는 $f$의 한 거듭제곱을
나눈다. 특히 $g$는 $A$에서 비영인자이다. 사상
$(f, d^i) : M^i \to M^i \oplus M^{i + 1}$의
$r_i \times r_i$ 소행렬식들은 아이디얼 $I_i(M^\bullet, f)$를 생성하고,
$(f, d^i)$의 $(r_i + 1) \times (r_i + 1)$ 소행렬식들은 영이다. 이는
$f$에서 국소화한 뒤 확인할 수 있는데, 거기서는 사상의 계수가
$r_i$와 같기 때문이다. 전사
$$
M^i \oplus M^{i + 1}
\longrightarrow
Q = \Coker(f, d^i)/g\text{-torsion}
$$
를 생각하자. 보조정리 \ref{more-algebra-lemma-fitting-ideals-and-pd1}에 의해 가군
$Q$는 계수 $r_{i + 1}$인 유한 국소 자유 가군이다. 따라서 $Q$는
$f$-꼬임이 없고, $(f, d^i)$의 여핵을 $f$-거듭제곱 꼬임으로 나눈 것도
$Q$라고 결론짓는다.

\medskip\noindent
유한 자유 $A$-가군들의 복합체
$$
0 \to f^{i + 1}M^i \xrightarrow{1, d^i} f^iM^i \oplus f^{i + 1}M^{i + 1}
\xrightarrow{d^i, -1} f^iM^{i + 1} \to 0
$$
를 생각하자. 이는 $f$에서 국소화한 뒤 분할 완전열이 된다. 사상
$(1, d^i) : f^{i + 1}M^i \to f^iM^i \oplus f^{i + 1}M^{i + 1}$은 위에서
연구한 사상 $(f, d^i) : M^i \to M^i \oplus M^{i + 1}$과 동형이다.
따라서 상
$$
Q' = \Im(f^iM^i \oplus f^{i + 1}M^{i + 1} \xrightarrow{d^i, -1} f^iM^{i + 1})
$$
은 $Q$와 동형이고,
% P02-E0455: frozen source runs the conclusion "in particular projective" into the isomorphism without a connector.
특히 사영이다. 한편 절 \ref{more-algebra-section-eta}의 $\eta_f$ 구성에 의해 단사
사상 $(1, d^i) : (\eta_fM)^i \to f^iM^i \oplus f^{i + 1}M^{i + 1}$의 상은
$(d^i, -1)$의 핵이다. 따라서 동형
$(\eta_fM)^i \oplus Q' = f^iM^i \oplus f^{i + 1}M^{i + 1}$을 얻고,
% P02-E0456: frozen source writes eta_f M^i instead of the degree-i term (eta_f M)^i.
실제로 $\eta_fM^i$가 계수 $r_i$인 유한 국소 자유 가군이며
$(1, d^i)$가 직합 인자의 포함임을 알 수 있다.
\end{proof}

\begin{lemma}
\label{more-algebra-lemma-eta-base-change}
$A \to B$를 환 사상이라 하고 $f \in A$를 비영인자라 하자.
$M^\bullet$을 유한 자유 $A$-가군들의 유계 복합체라 하자. $f$가
$B$의 비영인자 $g$로 가고 모든 $i \in \mathbf{Z}$에 대해
$I_i(M^\bullet, f)$가 주 아이디얼이라고 가정하자. 그러면 표준 동형
$\eta_fM^\bullet \otimes_A B = \eta_g(M^\bullet \otimes_A B)$이 있다.
\end{lemma}

\begin{proof}
$N^i = M^i \otimes_A B$로 두자. $(N^i)_g$의 부분가군으로서
$f^iM^i \otimes_A B = g^iN^i$임을 관찰하라. 사상들
% P02-E0457: frozen source uses tensor products in both codomains where the direct sums from the preceding lemma are required.
$$
(\eta_fM)^i \otimes_A B \to g^iN^i \otimes g^{i + 1}N^{i + 1}
\quad\text{and}\quad
(\eta_gN)^i \to g^iN^i \otimes g^{i + 1}N^{i + 1}
$$
은 보조정리 \ref{more-algebra-lemma-eta-locally-free}에 의해 직합 인자들의 포함이다.
$g$에서 국소화한 뒤 그 상들이 일치하므로 결론을 얻는다.
\end{proof}

\begin{lemma}
\label{more-algebra-lemma-ideal-direct-summand}
$A$를 환이라 하고 $M$, $N_1$, $N_2$를 유한 사영 $A$-가군들이라 하자.
$s : M \to N_1 \oplus N_2$를 분할 단사 사상이라 하자. 다음 성질을 갖는
유한 생성 아이디얼 $J \subset A$가 존재한다. 환 사상 $A \to B$가
$A/J$를 통해 분해될 필요충분조건은 $s \otimes \text{id}_B$가
$M \otimes_A B$를
$N_1 \otimes_A B \oplus N_2 \otimes_A B$의 부분가군들의 직합과
동일시하는 것이다.
\end{lemma}

\begin{proof}
$s$의 분할 $\pi : N_1 \oplus N_2 \to M$을 택하자.
% P02-E0458: frozen source's notation sentence omits "by".
$q_i : N_1 \oplus N_2 \to N_1 \oplus N_2$로 $N_i$ 위로의 사영자를
나타내자. $p_i = \pi \circ q_i \circ s$로 두자.
$p_1 + p_2 = \text{id}_M$임을 관찰하라. $M$이 $N_1 \oplus N_2$의
부분가군들의 직합일 필요충분조건은 $p_1$과 $p_2$가 직교 사영자인
것이라고 주장한다. 따라서 $J$는
$p_1 \circ p_1 - p_1$, $p_2 \circ p_2 - p_2$, $p_1 \circ p_2$, 그리고
$p_2 \circ p_1$이 $J \otimes_A \text{End}_A(M)$에 포함되도록 하는
$A$의 가장 작은 아이디얼이다. 일부 세부 사항은 생략한다.
\end{proof}

\noindent
$A$를 환이라 하고 $f \in A$를 비영인자라 하자. $M^\bullet$을 유한
자유 $A$-가군들의 유계 복합체라 하자. 모든 $i \in \mathbf{Z}$에 대해
아이디얼들 $I_i(M^\bullet, f)$가 주 아이디얼이라고 가정하자. 그러면
사상들
$$
(1, d^i) : (\eta_fM)^i / f(\eta_fM)^i \longrightarrow
f^iM^i/f^{i + 1}M^i \oplus f^{i + 1}M^{i + 1}/f^{i + 2}M^{i + 1}
$$
은 보조정리 \ref{more-algebra-lemma-eta-locally-free}에 의해 분할 단사 사상이다.
% P02-E0459: frozen source's notation sentence omits "by".
$J_i(M^\bullet, f) \subset A/fA$로 위에 표시한 분할 단사 사상
$(1, d^i)$에 대응하는 보조정리 \ref{more-algebra-lemma-ideal-direct-summand}의 유한
생성 아이디얼을 나타내자.

\begin{lemma}
\label{more-algebra-lemma-eta-vanishing-beta-plus-pre}
$A$를 환이라 하고 $f \in A$를 비영인자라 하자. $M^\bullet$과
$N^\bullet$을 $D(A)$의 같은 대상을 표현하는 유한 자유 $A$-가군들의
두 유계 복합체라 하자. 모든 $i \in \mathbf{Z}$에 대해
$I_i(M^\bullet, f)$가 주 아이디얼이라고 가정하자. 그러면 $A/fA$의
아이디얼로서 $J_i(M^\bullet, f) = J_i(N^\bullet, f)$이다.
\end{lemma}

\begin{proof}
$I_i(M^\bullet, f)$가 주 아이디얼이라는 사실은 보조정리
% P02-E0460: frozen source repeats I_i(M,f) where the second representative N is required.
\ref{more-algebra-lemma-ideal-well-defined}에 의해 $I_i(M^\bullet, f)$가 주
아이디얼임을 함의하므로 명제는 의미가 있다. 보조정리
\ref{more-algebra-lemma-ideal-well-defined}의 증명에서와 같이, 어떤 자명한 복합체
$Q^\bullet$에 대해 $N^\bullet = M^\bullet \oplus Q^\bullet$이라 가정할
수 있다. 즉,
$$
Q^\bullet = \ldots \to 0 \to A \xrightarrow{1} A \to 0 \to \ldots
$$
여기서 $A$는 차수 $j$와 $j + 1$에 놓여 있다. $\eta_f$가 직합과
양립하므로 사상
$$
(1, d^i) : (\eta_fN)^i / f(\eta_fN)^i \longrightarrow
f^iN^i/f^{i + 1}N^i \oplus f^{i + 1}N^{i + 1}/f^{i + 2}N^{i + 1}
$$
은 $M^\bullet$에 대한 대응하는 사상과 $Q^\bullet$에 대한 대응하는
사상의 직합임을 알 수 있다. 해당 아이디얼들을 정의하는 보편 성질에
의해
$J_i(N^\bullet, f) = J_i(M^\bullet, f) + J_i(Q^\bullet, f)$라고
결론짓는다. 따라서 모든 $i$에 대해 $J_i(Q^\bullet, f) = 0$임을 보이면
충분하다. 이 계산은 생략한다.
\end{proof}

\begin{lemma}
\label{more-algebra-lemma-eta-vanishing-beta-plus}
$A$를 환이라 하고 $f \in A$를 비영인자라 하자. $M^\bullet$을 유한
자유 $A$-가군들의 유계 복합체라 하자. 모든 $i \in \mathbf{Z}$에 대해
$I_i(M^\bullet, f)$가 주 아이디얼이라고 가정하자. $A/fA$의 아이디얼
$J(M^\bullet, f) = \sum_i J_i(M^\bullet, f)$를 생각하자. 소 아이디얼들의
집합
\begin{align*}
E
& =
\{f \in \mathfrak p \subset A \mid
\Ker(d^i \bmod f^2)_\mathfrak p \text{ surjects onto }
\Ker(d^i \bmod f)_\mathfrak p \text{ for all }i \in \mathbf{Z}\} \\
& =
\{f \in \mathfrak p \subset A \mid
\text{the localizations }\beta_\mathfrak p
\text{ of the Bockstein operators are zero}\}
\end{align*}
을 생각하자. 그러면 다음이 성립한다.
\begin{enumerate}
\item $J(M^\bullet, f)$는 유한 생성이다.
\item $A/fA \to C = (A/fA)/J(M^\bullet, f)$는 유한 표시인 전사이다.
\item $\mathfrak p \in E$이면 $J(M^\bullet, f)_\mathfrak p = 0$이다.
\item $f \in \mathfrak p$이고 모든 $i \in \mathbf{Z}$에 대해
$H^i(M^\bullet)_\mathfrak p$가 자유이면 $\mathfrak p \in E$이다.
\item $\eta_f M^\bullet \otimes_A C$의 코호몰로지 가군들은 유한 국소
자유 $C$-가군들이다.
\end{enumerate}
\end{lemma}

\begin{proof}
$E$의 정의에 나오는 등식은 보조정리 \ref{more-algebra-lemma-vanishing-beta}에서
따르고, 또한 그 보조정리의 마지막 명제가 (4)를 함의한다.

\medskip\noindent
(1)은 아이디얼들 $J_i(M^\bullet, f)$가 유한 생성이고 $M^\bullet$이
유계이므로 거의 모든 $i$에 대해 $J_i(M^\bullet, f)$가 영이기 때문에
참이다. (2)는 (1)을 다시 서술한 것일 뿐이다.

\medskip\noindent
(3)의 증명. 보조정리 \ref{more-algebra-lemma-eta-locally-free}에 의해 모든 $i$에
대해 $(\eta_fM)^i$가 계수 $r_i$인 유한 국소 자유 가군임을 얻는다.
사상
$$
(1, d^i) :
(\eta_fM)^i / f(\eta_fM)^i
\longrightarrow
f^iM^i/f^{i + 1}M^i \oplus f^{i + 1}M^{i + 1}/f^{i + 2}M^{i + 1}
$$
을 생각하자. $\mathfrak p \in E$를 택하자. 보조정리
\ref{more-algebra-lemma-eta-vanishing-beta}와 가군들 $(\eta_fM)^i$의 국소 자유성에
의해, 위 화살표 $(1, d^i)$와 양립하도록
$$
\left((\eta_fM)^i / f(\eta_fM)^i\right)_\mathfrak p =
(A/fA)_\mathfrak p^{\oplus m_i} \oplus (A/fA)_\mathfrak p^{\oplus n_i}
$$
로 쓸 수 있다. 아이디얼 $J_i(M^\bullet, f)$의 보편 성질에 의해
$J_i(M^\bullet, f)_\mathfrak p = 0$이라고 결론짓는다.
% P02-E0461: frozen source concludes with an undefined ideal I instead of the stated J(M,f).
따라서 $\mathfrak p \in E$에 대해 $I_\mathfrak p = fA_\mathfrak p$이다.

\medskip\noindent


\medskip\noindent
(5)의 증명. $\eta_fM^\bullet$의 미분이 가환 도표
% P02-E0462: both frozen diagrams use f^{i+1}M^i where degree consistency requires M^{i+1}.
$$
\xymatrix{
(\eta_fM)^i \ar[d] \ar[r] &
f^iM^i \oplus f^{i + 1}M^{i + 1}
\ar[d]^{\left(
\begin{matrix}
0 & 1 \\
0 & 0
\end{matrix}
\right)} \\
(\eta_fM)^{i + 1} \ar[r] &
f^{i + 1}M^i \oplus f^{i + 2}M^{i + 2}
}
$$
에 들어감을 관찰하라. 구성에 의해 $C$와 텐서곱한 뒤 왼쪽 가군들은
오른쪽 직합 인자들의 직합 인자들의 직합이다. 다음을 그려 보자.
$$
\xymatrix{
(\eta_fM)^i \otimes_A C \ar[d] \ar@{=}[r] &
K^i \oplus L^i \ar[r] \ar[d] &
f^iM^i \otimes_A C \oplus f^{i + 1}M^{i + 1} \otimes_A C
\ar[d]^{\left(
\begin{matrix}
0 & 1 \\
0 & 0
\end{matrix}
\right)} \\
(\eta_fM)^{i + 1} \otimes_A C \ar@{=}[r] &
K^{i + 1} \oplus L^{i + 1} \ar[r] &
f^{i + 1}M^i \otimes_A C \oplus f^{i + 2}M^{i + 2} \otimes_A C
}
$$
여기서 수평 화살표들은 직합 분해들과 양립할 뿐 아니라 직합 인자들의
포함들과도 양립한다. 따라서 미분은 $L^i$를 $K^{i + 1}$의 직합
인자와 동일시하고, $\eta_fM^\bullet \otimes_A C$의 차수 $i$
% P02-E0463: frozen source shifts the degree-i cohomology quotient to K^{i+1}/L^i.
코호몰로지는 원하는 대로 유한 사영인 가군 $K^{i + 1}/L^i$라고
결론짓는다.
\end{proof}
\section{복합체들의 극한을 취하기}
\label{more-algebra-section-limits}

\noindent
이 절에서는 복합체의 ``형식적 변형''이 주어졌을 때 그 극한을 취하면
무슨 일이 일어나는지 논의한다. 다음 두 경우를 생각할 것이다.
\begin{enumerate}
\item 전이 사상들이 전사이고 그 핵들이 국소 멱영인 환들의 역계의 극한
$A = \lim A_n$과, 다음 의미에서 서로 맞물리는 대상들
$K_n \in D(A_n)$이 주어진 경우,
$K_n = K_{n + 1} \otimes_{A_{n + 1}}^\mathbf{L} A_n$이다.
\item 환 $A$, 아이디얼 $I$, 그리고 다음 의미에서 서로 맞물리는 대상들
$K_n \in D(A/I^n)$이 주어진 경우,
$K_n = K_{n + 1} \otimes_{A/I^{n + 1}}^\mathbf{L} A/I^n$이다.
\end{enumerate}
추가 가정 아래에서는 $K = R\lim K_n$이
$K_n = K \otimes_A^\mathbf{L} A_n$ 또는
$K_n = K \otimes_A^\mathbf{L} A/I^n$이라는 의미에서 그 계를 복원함을
보일 수 있다.

\begin{lemma}
\label{more-algebra-lemma-Rlim-pseudo-coherent-gives-pseudo-coherent}
$A = \lim A_n$을 환들의 역계 $(A_n)$의 극한이라 하자.
$K_n \in D(A_n)$과 $D(A_{n + 1})$에서의 사상들
$K_{n + 1} \to K_n$이 주어졌다고 하자. 다음을 가정하자.
\begin{enumerate}
\item 전이 사상들 $A_{n + 1} \to A_n$은 전사이고 그 핵들은 국소
멱영이다.
\item 모든 $K_n$이 유사연접이거나, 어떤 정수 $n_0$가 존재하여
$K_{n_0}$가 유사연접이고 $n \geq n_0$에 대해
$A_{n + 1} \to A_n$의 핵들이 멱영 아이디얼이다.
\item 이 사상들은 동형
$K_{n + 1} \otimes_{A_{n + 1}}^\mathbf{L} A_n \to K_n$을 유도한다.
\end{enumerate}
그러면 $K = R\lim K_n$은 $D(A)$의 유사연접 대상이고, 모든 $n$에 대해
$K \otimes_A^\mathbf{L} A_n \to K_n$은 동형이다.
\end{lemma}

\begin{proof}
가정에 의해 $K_1$을 표현하는 유한 자유 $A_1$-가군들의 위로 유계인
복합체 $P_1^\bullet$을 찾을 수 있다. 정의
\ref{more-algebra-definition-pseudo-coherent}를 보라. 보조정리
\ref{more-algebra-lemma-check-pseudo-coherent-modulo-nilpotent}에 의해 모든 $n$에
대해 $K_n$이 유사연접이라고 결론짓는다. 그러면 보조정리
\ref{more-algebra-lemma-lift-complex-stably-frees}에 의해 $n > 1$에 대한 귀납으로
$K_n$을 표현하는 유한 자유 $A_n$-가군들의 복합체들 $P_n^\bullet$과,
사상들 $K_n \to K_{n - 1}$을 표현하면서 복합체들의 동형(!)
$P_n^\bullet \otimes_{A_n} A_{n - 1} \to P_{n - 1}^\bullet$을 유도하는
사상들 $P_n^\bullet \to P_{n - 1}^\bullet$을 찾을 수 있다. 따라서
$K = R\lim K_n$은 $P^\bullet = \lim P_n^\bullet$로 표현된다. 보조정리
\ref{more-algebra-lemma-compute-Rlim-modules}와 주석 \ref{more-algebra-remark-how-unique}를 보라.
각 $n$에 대해 $P_n^i$가 유한 자유 $A_n$-가군이고
$A = \lim A_n$이므로, $P^i$는 각 $i$에 대해 $P_1^i$와 같은 계수의
유한 자유 가군임을 알 수 있다. 이는 $K$가 유사연접이라는 뜻이다.
또한 $K \otimes_A^\mathbf{L} A_n$은
$P^\bullet \otimes_A A_n = P_n^\bullet$로 표현되며, 이것이 마지막
명제를 증명한다.
\end{proof}

\begin{lemma}
\label{more-algebra-lemma-Rlim-pseudo-coherent-gives-complete-pseudo-coherent}
$A$를 환이라 하고 $I \subset A$를 아이디얼이라 하자.
$K_n \in D(A/I^n)$과 $D(A/I^{n + 1})$에서의 사상들
$K_{n + 1} \to K_n$이 주어졌다고 하자. 다음을 가정하자.
\begin{enumerate}
\item $A$는 $I$-진 완비이다.
\item $K_1$은 유사연접이다.
\item 이 사상들은 동형
$K_{n + 1} \otimes_{A/I^{n + 1}}^\mathbf{L} A/I^n \to K_n$을 유도한다.
\end{enumerate}
그러면 $K = R\lim K_n$은 $D(A)$의 유사연접이고 유도 완비인 대상이며,
모든 $n$에 대해 $K \otimes_A^\mathbf{L} A/I^n \to K_n$은 동형이다.
\end{lemma}

\begin{proof}
보조정리 \ref{more-algebra-lemma-Rlim-pseudo-coherent-gives-pseudo-coherent}에 의해
$K$가 유사연접이고 모든 $n$에 대해
$K \otimes_A^\mathbf{L} A/I^n \to K_n$이 동형임을 이미 알고 있다.
마지막으로 보조정리 \ref{more-algebra-lemma-naive-derived-completion}에 의해 $K$는
유도 완비이다.
\end{proof}

\begin{lemma}
\label{more-algebra-lemma-Rlim-perfect-gives-perfect}
\begin{reference}
\cite[Lemma 4.2]{Bhatt-Algebraize}
\end{reference}
$A = \lim A_n$을 환들의 역계 $(A_n)$의 극한이라 하자.
$K_n \in D(A_n)$과 $D(A_{n + 1})$에서의 사상들
$K_{n + 1} \to K_n$이 주어졌다고 하자. 다음을 가정하자.
\begin{enumerate}
\item 전이 사상들 $A_{n + 1} \to A_n$은 전사이고 그 핵들은 국소
멱영이다.
\item 모든 $K_n$이 퍼펙트이거나, 어떤 정수 $n_0$가 존재하여
$K_{n_0}$가 퍼펙트이고 $n \geq n_0$에 대해
$A_{n + 1} \to A_n$의 핵들이 멱영이다.
\item 이 사상들은 동형
$K_{n + 1} \otimes_{A_{n + 1}}^\mathbf{L} A_n \to K_n$을 유도한다.
\end{enumerate}
그러면 $K = R\lim K_n$은 $D(A)$의 퍼펙트 대상이고, 모든 $n$에 대해
$K \otimes_A^\mathbf{L} A_n \to K_n$은 동형이다.
\end{lemma}

\begin{proof}
보조정리 \ref{more-algebra-lemma-Rlim-pseudo-coherent-gives-pseudo-coherent}에 의해
$K$가 유사연접이고 모든 $n$에 대해
$K \otimes_A^\mathbf{L} A_n \to K_n$이 동형임을 이미 알고 있다.
그 핵이 극대 아이디얼 $\mathfrak m$인 전사 $A \to \kappa$를 생각하자.
$A_1$에서 단원으로 가는 $A$의 모든 원소는 Algebra, 보조정리
\ref{algebra-lemma-locally-nilpotent-unit}에 의해 $A$에서 단원이다.
따라서 Algebra, 보조정리 \ref{algebra-lemma-contained-in-radical}에 의해
$\Ker(A \to A_1)$은 $A$의 야코브슨 근기에 포함된다. 그러므로
$A \to \kappa$는 $A \to A_1 \to \kappa$로 분해된다. 따라서
% P02-E0464: frozen source omits terminal punctuation after this display and the later Rlim exact-sequence display.
$$
K \otimes_A^\mathbf{L} \kappa =
K \otimes_A^\mathbf{L} A_1 \otimes_{A_1}^\mathbf{L} \kappa =
K_1 \otimes_{A_1}^\mathbf{L} \kappa
$$
이다. 가정 (2)에 의해 $K_1$은 퍼펙트임에 유의하라. 따라서 보조정리
\ref{more-algebra-lemma-perfect}에 의해 $K_1$은 유한 Tor 차원을 갖는다. 그러므로
모든 $i \not \in [a, b]$에 대해
$H^i(K \otimes_A^\mathbf{L} \kappa) = 0$이 되도록 하는
$a, b \in \mathbf{Z}$가 존재한다. 보조정리
\ref{more-algebra-lemma-check-perfect-pointwise}에 의해 $K$가 퍼펙트라고 결론짓는다.
\end{proof}

\begin{lemma}
\label{more-algebra-lemma-Rlim-perfect-gives-complete}
$A$를 환이라 하고 $I \subset A$를 아이디얼이라 하자.
$K_n \in D(A/I^n)$과 $D(A/I^{n + 1})$에서의 사상들
$K_{n + 1} \to K_n$이 주어졌다고 하자. 다음을 가정하자.
\begin{enumerate}
\item $A$는 $I$-진 완비이다.
\item $K_1$은 퍼펙트 대상이다.
\item 이 사상들은 동형
$K_{n + 1} \otimes_{A/I^{n + 1}}^\mathbf{L} A/I^n \to K_n$을 유도한다.
\end{enumerate}
그러면 $K = R\lim K_n$은 $D(A)$의 퍼펙트이고 유도 완비인 대상이며,
모든 $n$에 대해 $K \otimes_A^\mathbf{L} A/I^n \to K_n$은 동형이다.
\end{lemma}

\begin{proof}
보조정리들 \ref{more-algebra-lemma-Rlim-perfect-gives-perfect}와
\ref{more-algebra-lemma-Rlim-pseudo-coherent-gives-complete-pseudo-coherent}를
결합하라(후자는 유도 완비성을 준다).
\end{proof}

\noindent
다음 보조정리가 비유계 복합체들에 대해서도 성립하는지는 알지 못한다.

\begin{lemma}
\label{more-algebra-lemma-Rlim-gives-complete}
$A$를 환이라 하고 $I \subset A$를 아이디얼이라 하자.
$K_n \in D(A/I^n)$과 $D(A/I^{n + 1})$에서의 사상들
$K_{n + 1} \to K_n$이 주어졌다고 하자. 다음을 가정하자.
\begin{enumerate}
\item $A$는 Noether 환이다.
\item $K_1$은 위로 유계이다.
\item 이 사상들은 동형
$K_{n + 1} \otimes_{A/I^{n + 1}}^\mathbf{L} A/I^n \to K_n$을 유도한다.
\end{enumerate}
그러면 $K = R\lim K_n$은 $D^-(A)$의 유도 완비인 대상이고, 모든
$n$에 대해 $K \otimes_A^\mathbf{L} A/I^n \to K_n$은 동형이다.
\end{lemma}

\begin{proof}
$D(A)$의 대상 $K$는 보조정리 \ref{more-algebra-lemma-naive-derived-completion}에 의해
유도 완비이다.

\medskip\noindent
$i > b$에 대해 $H^i(K_1) = 0$이라고 하자. $K_1$을 표현하면서
$i > b$에 대해 $P_1^i = 0$인 자유 $A/I$-가군들의 복합체
$P_1^\bullet$을 찾을 수 있다. 보조정리
\ref{more-algebra-lemma-lift-complex-projectives}에 의해 $n > 1$에 대한 귀납으로
$K_n$을 표현하는 자유 $A/I^n$-가군들의 복합체들 $P_n^\bullet$과,
사상들 $K_n \to K_{n - 1}$을 표현하면서 복합체들의 동형(!)
$P_n^\bullet/I^{n - 1}P_n^\bullet \to P_{n - 1}^\bullet$을 유도하는
사상들 $P_n^\bullet \to P_{n - 1}^\bullet$을 찾을 수 있다.

\medskip\noindent
따라서 $R\lim K_n$이 $P^\bullet = \lim P_n^\bullet$로 표현되는 상황에
이르렀다. 보조정리 \ref{more-algebra-lemma-compute-Rlim-modules}와 주석
\ref{more-algebra-remark-how-unique}를 보라. 복합체들 $P_n^\bullet$은 평탄한
$A/I^n$-가군들의 균일하게 위로 유계인 복합체들이며 전이 사상들은
항별로 전사이다. 그러면 보조정리 \ref{more-algebra-lemma-limit-flat}에 의해
$P^\bullet$은 평탄한 $A$-가군들의 위로 유계인 복합체이다. 따라서
$K \otimes_A^\mathbf{L} A/I^t$는 $P^\bullet \otimes_A A/I^t$로
표현된다. 보조정리 \ref{more-algebra-lemma-limit-flat}에 의해 항별로
$P^\bullet \otimes_A A/I^t = \lim P_n^\bullet \otimes_A A/I^t$이다.
$P_n^\bullet$의 선택에 의해 $n \geq t$에 대해 전이 사상들
$P_{n + 1}^\bullet \otimes_A A/I^t \to P_n^\bullet \otimes_A A/I^t$은
동형이다. 따라서
$\lim P_n^\bullet \otimes_A A/I^t = P_t^\bullet \otimes_A A/I^t = P_t^\bullet$이다.
$P_t^\bullet$이 $K_t$를 표현하므로
$K \otimes_A^\mathbf{L} A/I^t \to K_t$가 동형임을 알 수 있다.
\end{proof}

\noindent
다음은 다른 유형의 결과이다.

\begin{lemma}[Koll\'ar-Kov\'acs]
\label{more-algebra-lemma-kollar-kovacs}
\begin{reference}
% P02-E0465: frozen source drops the accent from Kovacs in this reference line.
2018년 23/02의 Kovacs 이메일.
\end{reference}
$I$를 Noether 환 $A$의 아이디얼이라 하고 $K \in D(A)$라 하자.
$K_n = K \otimes_A^\mathbf{L} A/I^n$으로 두자. 모든
$i \in \mathbf{Z}$에 대해 다음을 가정하자.
\begin{enumerate}
\item $H^i(K)$는 유한 $A$-가군이다.
\item 계 $H^i(K_n)$은 Mittag--Leffler 조건을 만족한다.
\end{enumerate}
그러면 모든 $i \in \mathbf{Z}$에 대해
$\lim H^i(K)/I^nH^i(K)$은 $\lim H^i(K_n)$과 같다.
\end{lemma}

\begin{proof}
$K^\wedge = R\lim K_n$은 $K$의 유도 완비화임을 상기하라. 명제
\ref{more-algebra-proposition-noetherian-naive-completion-is-completion}을 보라.
보조정리 \ref{more-algebra-lemma-derived-completion-pseudo-coherent}에 의해
$H^i(K^\wedge) = \lim H^i(K)/I^nH^i(K)$이다. 보조정리
\ref{more-algebra-lemma-break-long-exact-sequence-modules}에 의해 짧은 완전열들
% P02-E0464: second frozen display with no terminal punctuation before a new sentence.
$$
0 \to R^1\lim H^{i - 1}(K_n) \to H^i(K^\wedge) \to \lim H^i(K_n) \to 0
$$
을 얻는다. Mittag--Leffler 조건은 왼쪽 항들이 영임을 보장한다(보조정리
\ref{more-algebra-lemma-compute-Rlim-modules}). 따라서 보조정리가 참이라고 결론짓는다.
\end{proof}











\section{몇 가지 평가 사상}
\label{more-algebra-section-evaluation}

\noindent
이 절에서는 적절한 종류의 복합체에 대해 $R\Hom$의 어떤 표준 사상들이
동형사상임을 증명한다.

\begin{lemma}
\label{more-algebra-lemma-internal-hom-evaluate-isomorphism}
$R$을 환이라 하자. $K, L, M$을 $D(R)$의 대상이라 하자.
% P02-E0466: the frozen source begins this sentence and its technical analogue with lowercase “the”.
보조정리 \ref{more-algebra-lemma-internal-hom-evaluate}의 사상
$$
R\Hom_R(L, M) \otimes_R^\mathbf{L} K \longrightarrow R\Hom_R(R\Hom_R(K, L), M)
$$
은 다음 두 경우에 동형사상이다.
\begin{enumerate}
\item $K$가 퍼펙트이거나,
\item $K$가 유사연접이고 $L \in D^+(R)$이며,
% P02-E0467: the frozen source omits the finite verb “has” before finite injective dimension.
$M$이 유한 단사 차원을 갖는다.
\end{enumerate}
\end{lemma}

\begin{proof}
$M$을 나타내는 K-단사 복합체 $I^\bullet$, $L$을 나타내는 K-단사
복합체 $J^\bullet$, 그리고 $K$를 나타내는 유한 사영 가군들로 이루어진
위로 유계인 복합체 $K^\bullet$을 택하자. 보조정리
\ref{more-algebra-lemma-evaluate-and-more}의 복합체 사상
$$
\text{Tot}(\Hom^\bullet(J^\bullet, I^\bullet) \otimes_R K^\bullet)
\longrightarrow
\Hom^\bullet(\Hom^\bullet(K^\bullet, J^\bullet), I^\bullet)
$$
을 생각하자. $K^s$가 유한 사영이므로
$$
\left(\prod\nolimits_{p + r = t} \Hom_R(J^{-r}, I^p)\right) \otimes_R K^s =
\prod\nolimits_{p + r = t} \Hom_R(J^{-r}, I^p) \otimes_R K^s
$$
임에 유의하라. 이 사상은 사상들
$$
c_{p, r, s} :
\Hom_R(J^{-r}, I^p) \otimes_R K^s
\longrightarrow
\Hom_R(\Hom_R(K^s, J^{-r}), I^p)
$$
로 주어지며, $K^s$가 유한 사영이므로 이들은 동형사상이다.
왼쪽의 차수 $n$인 임의의 원소 $\alpha = (\alpha^{p, r, s})$에 대해,
$\alpha^{p, r, s}$가 영이 아닌 $s$의 값은 유한 개뿐이다
($n = p + r + s$인 어떤 $p, r$에 대하여). 따라서 오른쪽 원소
$\beta = (\beta^{p, r, s})$에도 같은 소멸 조건이 강제되면 우리의
사상은 동형사상이다. $K^\bullet$이 유한 사영 가군들로 이루어진 유계
복합체이면 이는 명백하다. 한편 $I^\bullet$을 유계로, $J^\bullet$을
아래로 유계로 택할 수 있다면, $\beta^{p, r, s}$는 $p$가 어떤 고정된
범위 밖에 있을 때, $s \gg 0$일 때, 그리고 $r \gg 0$일 때 영이다.
따라서 $\beta^{p, r, s}$가 영이 아닌 $n = p + r + s$의 해들 중에는
$s$의 값이 유한 개만 나타난다.
\end{proof}

\begin{lemma}
\label{more-algebra-lemma-internal-hom-evaluate-isomorphism-technical}
$R$을 환이라 하자. $K, L, M$을 $D(R)$의 대상이라 하자.
% P02-E0466: second exact locus of the same lowercase-sentence defect.
보조정리 \ref{more-algebra-lemma-internal-hom-evaluate}의 사상
$$
R\Hom_R(L, M) \otimes_R^\mathbf{L} K \longrightarrow R\Hom_R(R\Hom_R(K, L), M)
$$
은 다음 세 조건이 성립하면 동형사상이다.
\begin{enumerate}
\item $L, M$은 유한 단사 차원을 갖고,
\item $R\Hom_R(L, M)$은 유한 Tor 차원을 가지며,
% P02-E0468: the frozen three-condition list lacks terminal punctuation.
\item 모든 $n \in \mathbf{Z}$에 대하여 절단 $\tau_{\leq n}K$는
유사연접이다.
\end{enumerate}
\end{lemma}

\begin{proof}
정수 $n$을 택하고 구별 삼각형
% P02-E0469: first frozen display lacks punctuation before the following “see”.
$$
\tau_{\leq n}K \to K \to \tau_{\geq n + 1}K \to \tau_{\leq n}K[1]
$$
을 생각하자. 유도 범주론, 주석
\ref{derived-remark-truncation-distinguished-triangle}을 보라.
가정 (3)과 보조정리 \ref{more-algebra-lemma-internal-hom-evaluate-isomorphism}에 의해
이 사상은 $\tau_{\leq n}K$에 대해 동형사상이다. 따라서 어떤 $c$에
대해 다음 두 대상의 상동군이 차수 $\leq n - c$에서 소멸함을 보이면
충분하다.
$$
R\Hom_R(L, M) \otimes_R^\mathbf{L} \tau_{\geq n + 1}K
\quad\text{and}\quad
R\Hom_R(R\Hom_R(\tau_{\geq n + 1}K, L), M)
$$
이는 가정 (2)와 (1)에서 바로 따른다.
\end{proof}

\begin{lemma}
\label{more-algebra-lemma-internal-hom-evaluate-tensor-isomorphism}
$R$을 환이라 하자. $K, L, M$을 $D(R)$의 대상이라 하자. 보조정리
\ref{more-algebra-lemma-internal-hom-diagonal-better}의 사상
$$
K \otimes_R^\mathbf{L} R\Hom_R(M, L) \longrightarrow
R\Hom_R(M, K \otimes_R^\mathbf{L} L)
$$
은 다음 경우에 동형사상이다.
\begin{enumerate}
\item $M$이 퍼펙트이거나,
\item $K$가 퍼펙트이거나,
\item $M$이 유사연접이고 $L \in D^+(R)$이며 $K$가
% P02-E0470: frozen formal interval incorrectly closes at infinity; preserved for source fidelity.
$[a, \infty]$에 Tor 진폭을 갖는다.
\end{enumerate}
\end{lemma}

\begin{proof}
$M$이 퍼펙트인 경우의 증명. 화살표의 양변은 $M$에 관한 구별
삼각형을 구별 삼각형으로 보내며 직합과 가환함에 유의하라. 따라서
$M = R[n]$일 때 성립하는지만 확인하면 충분하다. 유도 범주론, 주석
\ref{derived-remark-check-on-generator}와 보조정리
\ref{more-algebra-lemma-perfect-ring-classical-generator}를 보라. 이 경우 결과는
명백하다.

\medskip\noindent
$K$가 퍼펙트인 경우의 증명. 앞의 경우와 같은 논증을 쓰면 된다.

\medskip\noindent
(3)의 증명. $K$와 $L$을 각각 $R$-가군의 아래로 유계인 복합체
$K^\bullet$과 $L^\bullet$로 나타낼 수 있다. 보조정리
\ref{more-algebra-lemma-bounded-below-tor-amplitude}에 의해 $K^\bullet$은 평탄
$R$-가군들로 이루어진 K-평탄 복합체라고 가정할 수 있다.
정의 \ref{more-algebra-definition-pseudo-coherent}에 의해 $M$을 유한 자유
$R$-가군들로 이루어진 위로 유계인 복합체 $M^\bullet$로 나타낼 수
있다. 그러면 왼쪽 대상은
$$
\text{Tot}(K^\bullet \otimes_R \Hom^\bullet(M^\bullet, L^\bullet))
$$
로 나타내어지고, 오른쪽 대상은
% P02-E0469: second frozen display lacks terminal punctuation before “This uses”.
$$
\Hom^\bullet(M^\bullet, \text{Tot}(K^\bullet \otimes_R L^\bullet))
$$
로 나타내어진다. 여기서는 보조정리
\ref{more-algebra-lemma-RHom-out-of-projective}를 사용했다. 두 복합체의 차수
$n$인 가군은 모두
$$
\bigoplus\nolimits_{p + q + r = n} K^p \otimes \Hom_R(M^{-r}, L^q) =
\bigoplus\nolimits_{p + q + r = n} \Hom_R(M^{-r}, K^p \otimes_R L^q)
$$
이다. 이는 $M^{-r}$가 유한 자유이기 때문이다
% P02-E0471: frozen parenthetical has malformed English word order.
(또한 이들은 유한 직합이다). 보조정리
\ref{more-algebra-lemma-internal-hom-diagonal-better}에서 정의한 사상은 보조정리
\ref{more-algebra-lemma-diagonal-better}에서 정의한 복합체의 사상에서 오며,
후자는 이 가군들 사이의 표준 동형사상들을 사용한다.
\end{proof}

\begin{lemma}
\label{more-algebra-lemma-hom-complex-K-flat}
$R$을 환이라 하자. $P^\bullet$을 사영 $R$-가군들로 이루어진 위로
유계인 복합체라 하고, $K^\bullet$을 $R$-가군의 K-평탄 복합체라 하자.
$P^\bullet$이 $D(R)$의 퍼펙트 대상이면
$\Hom^\bullet(P^\bullet, K^\bullet)$은 K-평탄이고
$R\Hom_R(P^\bullet, K^\bullet)$을 나타낸다.
\end{lemma}

\begin{proof}
마지막 진술은 보조정리 \ref{more-algebra-lemma-RHom-out-of-projective}이다.
$P^\bullet$이 퍼펙트 대상을 나타내므로 유한 사영 $R$-가군들로
이루어진 유한 복합체 $F^\bullet$이 존재하여 $P^\bullet$과
$F^\bullet$이 $D(R)$에서 동형이다. 정의
\ref{more-algebra-definition-perfect}를 보라. 그러면 $P^\bullet$과 $F^\bullet$은
호모토피 동치이다. 유도 범주론, 보조정리
\ref{derived-lemma-morphisms-from-projective-complex}을 보라.
따라서 $\Hom^\bullet(P^\bullet, K^\bullet)$과
$\Hom^\bullet(F^\bullet, K^\bullet)$은 호모토피 동치이다. 그러므로
전자가 K-평탄인 것과 후자가 K-평탄인 것은 동치이다(정의
\ref{more-algebra-definition-K-flat}과 보조정리
\ref{more-algebra-lemma-derived-tor-homotopy}에서 따른다). 다음은 명백하다.
% P02-E0469: third frozen display needs a comma before its defining where-clause.
$$
\Hom^\bullet(F^\bullet, K^\bullet) =
\text{Tot}(E^\bullet \otimes_R K^\bullet)
$$
여기서 $E^\bullet$은 $F^\bullet$의 쌍대 복합체이며 그 항들은
$E^n = \Hom_R(F^{-n}, R)$이다. 보조정리
\ref{more-algebra-lemma-dual-perfect-complex}과 그 증명을 보라. $E^\bullet$은
사영 가군들로 이루어진 유계 복합체이므로 보조정리
\ref{more-algebra-lemma-derived-tor-quasi-isomorphism}에 의해 K-평탄이다. 이제
보조정리 \ref{more-algebra-lemma-tensor-product-K-flat}에 의해 결론을 얻는다.
\end{proof}

\section{유도 Hom의 밑변환}
\label{more-algebra-section-base-change-RHom}

\noindent
이에 관한 내용 일부는 이미 보조정리
\ref{more-algebra-lemma-pseudo-coherence-and-base-change-ext}와 대수학, 절
\ref{algebra-section-functoriality-ext}에서 보았다.

\begin{lemma}
\label{more-algebra-lemma-upgrade-adjoint-tensor-RHom}
$R \to R'$을 환 사상이라 하자. $K \in D(R)$와 $M \in D(R')$에
대하여 표준 동형사상
% P02-E0472: first of four frozen displays lacking its surrounding sentence punctuation.
$$
R\Hom_R(K, M) = R\Hom_{R'}(K \otimes_R^\mathbf{L} R', M)
$$
이 있다.
\end{lemma}

\begin{proof}
$M$을 나타내는 $R'$-가군의 K-단사 복합체 $J^\bullet$을 택하자.
$I^\bullet$이 $R$-가군의 K-단사 복합체인 유사동형사상
$J^\bullet \to I^\bullet$을 택하자. $K$를 나타내는 $R$-가군의
K-평탄 복합체 $K^\bullet$을 택하자. 사상
% P02-E0472: second frozen sentence-ending display without punctuation.
$$
\Hom^\bullet(K^\bullet \otimes_R R', J^\bullet)
\longrightarrow
\Hom^\bullet(K^\bullet, I^\bullet)
$$
을 생각하자. 차수 $n$의 항들 위의 사상은 사상
$$
\prod\nolimits_{n = p + q} \Hom_{R'}(K^{-q} \otimes_R R', J^p)
\longrightarrow
\prod\nolimits_{n = p + q} \Hom_R(K^{-q}, I^p)
$$
으로 주어지며, 이는 $K^{-q} \to K^{-q} \otimes_R R'$로 전합성하고
$J^p \to I^p$로 후합성하여 얻는다. 증명을 끝내려면 상동군 위에서
동형사상
% P02-E0472: third frozen sentence-ending display without punctuation.
$$
\Hom_{D(R)}(K, M) = \Hom_{D(R')}(K \otimes_R^\mathbf{L} R', M)
$$
을 얻는다는 것을 보이면 충분하다. 이는 보조정리
\ref{more-algebra-lemma-tensor-hom-adjoint}에 의해 밑변환
$- \otimes_R^\mathbf{L} R' : D(R) \to D(R')$이 제한 함자
$D(R') \to D(R)$의 왼쪽 수반이기 때문이다.
\end{proof}

\noindent
$R \to R'$을 환 사상이라 하자. 밑변환 사상
% P02-E0472: fourth frozen display requires a comma before the continuation in D(R').
\begin{equation}
\label{more-algebra-equation-base-change-RHom}
R\Hom_R(K, M) \otimes_R^\mathbf{L} R'
\longrightarrow
R\Hom_{R'}(K \otimes_R^\mathbf{L} R', M \otimes_R^\mathbf{L} R')
\end{equation}
이 $D(R')$에 있으며 $K, M \in D(R)$에 관하여 함자적이다. 실제로
$- \otimes_R^\mathbf{L} R' : D(R) \to D(R')$과 제한 함자
$D(R') \to D(R)$의 수반성에 의해, 이는 사상
$$
R\Hom_R(K, M)
\longrightarrow
R\Hom_{R'}(K \otimes_R^\mathbf{L} R', M \otimes_R^\mathbf{L} R') =
R\Hom_R(K, M \otimes_R^\mathbf{L} R')
$$
과 같은 것이다(등식은 보조정리
\ref{more-algebra-lemma-upgrade-adjoint-tensor-RHom}에 의한다). 이 사상에는 표준
사상 $M \to M \otimes_R^\mathbf{L} R'$(수반의 단위원)을 사용할 수 있다.

\begin{lemma}
\label{more-algebra-lemma-base-change-RHom}
$R \to R'$을 환 사상이라 하자. $K, M \in D(R)$이라 하자. 사상
(\ref{more-algebra-equation-base-change-RHom})
$$
R\Hom_R(K, M) \otimes_R^\mathbf{L} R'
\longrightarrow
R\Hom_{R'}(K \otimes_R^\mathbf{L} R', M \otimes_R^\mathbf{L} R')
$$
은 다음 경우에 $D(R')$에서 동형사상이다.
\begin{enumerate}
\item $K$가 퍼펙트인 경우,
\item $R'$이 $R$-가군으로서 퍼펙트인 경우,
\item $R \to R'$이 평탄하고 $K$가 유사연접이며
$M \in D^{+}(R)$인 경우, 또는
\item $R'$이 $R$-가군으로서 유한 Tor 차원을 갖고 $K$가
유사연접이며
% P02-E0473: the frozen final case has no terminal punctuation.
$M \in D^{+}(R)$인 경우.
\end{enumerate}
\end{lemma}

\begin{proof}
제한 함자 $D(R') \to D(R)$을 적용한 뒤에 이 사상이 동형사상인지
확인해도 된다. 이 함자를 적용하면 우리의 사상은 보조정리
\ref{more-algebra-lemma-internal-hom-diagonal-better}의 사상
% P02-E0474: frozen source uses undefined L twice here instead of the lemma's M; formal spans preserved.
$$
R\Hom_R(K, L) \otimes_R^\mathbf{L} R'
\longrightarrow
R\Hom_R(K, L \otimes_R^\mathbf{L} R')
$$
이 된다. 왼쪽과 오른쪽을 맞추려면 이 보조정리 앞의 논의를 보라.
특히 여기서는 보조정리
\ref{more-algebra-lemma-upgrade-adjoint-tensor-RHom}을 사용한다. 따라서 보조정리
\ref{more-algebra-lemma-internal-hom-evaluate-tensor-isomorphism}에 의해 결론을
얻는다.
\end{proof}

\section{가군의 계}
\label{more-algebra-section-systems}

\noindent
$I$를 Noether 환 $A$의 아이디얼이라 하자. 이 절에서는 $A$ 위의 유한
가군과 유한 $A/I^n$-가군의 계 사이의 관계에 대한 지식을 보충한다.

\begin{lemma}
\label{more-algebra-lemma-consequence-Artin-Rees}
$I$를 Noether 환 $A$의 아이디얼이라 하자.
$
K \xrightarrow{\alpha} L \xrightarrow{\beta} M
$
을 유한 $A$-가군의 복합체라 하자. $H = \Ker(\beta)/\Im(\alpha)$로
두자. $n \geq 0$에 대해
$$
K/I^nK \xrightarrow{\alpha_n} L/I^nL \xrightarrow{\beta_n} M/I^nM
$$
을 유도된 복합체라 하자. $H_n = \Ker(\beta_n)/\Im(\alpha_n)$으로
두자. 그러면 가환 그림
$$
\xymatrix{
& & & H \ar[lld] \ar[ld] \ar[d] \\
\ldots \ar[r] & H_3 \ar[r] & H_2 \ar[r] & H_1
}
$$
을 주는 표준 $A$-가군 사상들이 있다. 또한 $c > 0$과 $n \geq c$에
대한 표준 $A$-가군 사상 $H_n \to H/I^{n - c}H$가 존재하여 합성들
$$
H/I^n H \to H_n \to  H/I^{n - c}H
\quad\text{and}\quad
H_n \to H/I^{n - c}H \to H_{n - c}
$$
은 표준 사상들이다. 더 나아가 다음이 성립한다.
\begin{enumerate}
\item $(H_n)$과 $(H/I^nH)$은 $\text{Mod}_A$의 pro-대상으로서 동형이다.
\item $\lim H_n = \lim H/I^n H$이다.
\item 역계 $(H_n)$은 Mittag--Leffler 조건을 만족한다.
\item $H_{n + c} \to H_n$의 상은 $H \to H_n$의 상과 같다.
\item 합성 $I^cH_n \to H_n \to H/I^{n - c}H \to H_n/I^{n - c}H_n$은
포함 $I^cH_n \to H_n$에 몫사상
$H_n \to H_n/I^{n - c}H_n$을 뒤이어 합성한 것이다.
\item
% P02-E0478: frozen English uses singular agreement for the compound kernel-and-cokernel subject.
$H/I^nH \to H_n$의 핵과 여핵은 $I^c$에 의해 소멸된다.
\end{enumerate}
\end{lemma}

\begin{proof}
% P02-E0475: first frozen quotient omits parentheses around the denominator sum.
$H_n = \beta^{-1}(I^nM)/\Im(\alpha) + I^nL$임에 유의하라.
$n \geq 2$에 대해
$\beta^{-1}(I^nM) \subset \beta^{-1}(I^{n - 1}M)$이고
$\Im(\alpha) + I^nL \subset \Im(\alpha) + I^{n - 1}L$이다. 따라서 표준
사상 $H_n \to H_{n - 1}$을 얻는다. 마찬가지로
$\Ker(\beta) \subset \beta^{-1}(I^nM)$이고
$\Im(\alpha) \subset \Im(\alpha) + I^nL$이므로 표준 사상
$H \to H_n$을 얻는다. 그림이 가환한다는 확인은 생략한다.

\medskip\noindent
Artin--Rees에 의해 $n \geq c_1$이면
$\beta^{-1}(I^nM) \subset \Ker(\beta) + I^{n - c_1}L$이고,
$n \geq c_2$이면
$\Ker(\beta) \cap I^nL \subset I^{n - c_2}\Ker(\beta)$인
$c_1, c_2 \geq 0$을 택할 수 있다. 대수학, 보조정리
\ref{algebra-lemma-map-AR}과 \ref{algebra-lemma-Artin-Rees}를 보라.
$c = c_1 + c_2$로 두자.

\medskip\noindent
$n \geq c$라 하자. $\psi_n : H_n \to H/I^{n - c}H$를 다음과 같이
정의한다. $x \in H_n$이라 하자. $x$를 나타내는
$y \in \beta^{-1}(I^nM)$을 택하자. $z \in \Ker(\beta)$와
$w \in I^{n - c_1}L$에 대해 $y = z + w$로 쓰자(이는 $c_1$의 선택에
의해 가능하다). $\psi_n(x)$를 $H/I^{n - c}H$에서 $z$의 동치류로
둔다. 이것이 잘 정의됨을 보이기 위해 같은 $x$에 대하여 위와 같은
두 번째 선택 $y', z', w'$을 했다고 하자. 표기는 자명하다. 그러면
$y' - y \in \Im(\alpha) + I^nL$이므로 어떤 $v \in K$와
$u \in I^nL$에 대해 $y' - y = \alpha(v) + u$이다. 따라서
$$
y' = z' + w' = \alpha(v) + u + z + w
\Rightarrow
z' = z + \alpha(v) + u + w - w'
$$
이다. $\beta(z' - z - \alpha(v)) = 0$이므로
$u + w - w' \in \Ker(\beta) \cap I^{n - c_1}L$이다. 이는 $c_2$의
선택에 의해
$I^{n - c_1 - c_2}\Ker(\beta) = I^{n - c}\Ker(\beta)$에 포함된다.
따라서 원하는 대로 $z'$과 $z$는 $H/I^{n - c}H$에서 같은 상을 갖는다.

\medskip\noindent
합성 $H/I^n H \to H_n \to  H/I^{n - c}H$은 표준 사상이다. 실제로
$z \in \Ker(\beta)$가
% P02-E0475: second frozen quotient omits parentheses around the denominator sum.
$H/I^nH = \Ker(\beta)/\Im(\alpha) + I^n\Ker(\beta)$의 원소 $x$를
나타내면, 위에서 구성한 사상들 아래 $x$가 $H/I^{n - c}H$에서 $z$의
동치류로 가는 것은 명백하다.

\medskip\noindent
합성 $H_n \to H/I^{n - c}H \to H_{n - c}$를 생각하자. 위의
$\psi_n$ 구성에서와 같은 $x, y, z, w$가 주어지면,
% P02-E0476: frozen source misspells “class” as “cass”.
$x$는 $H_{n - c}$에서 $z$의 동치류로 보내진다. 한편 증명의 첫 문단의
표준 사상 $H_n \to H_{n - c}$는 $x$를 $y$의 동치류로 보낸다.
따라서 $y - z \in \Im(\alpha) + I^{n - c}L$임을 보이면 되는데, 이는
$y - z = w \in I^{n - c_1}L \subset I^{n - c}L$이기 때문이다.

\medskip\noindent
(1)--(4)는 방금 증명한 것의 형식적 귀결이다. 즉, (1)은 사상들의
존재와 범주론, 주석 \ref{categories-remark-pro-category}의 pro-대상
사상의 정의에서 따른다. (2)는 동형인 pro-대상들이 동형인 극한을
갖기 때문에 성립한다. (3)은 (4)에서 바로 따른다. (4)는 표준 사상
$H_{n + c} \to H_n$의 분해
$H_{n + c} \to H/I^nH \to H_n$에서 따른다.

\medskip\noindent
(5)의 증명. $x \in I^cH_n$이라 하자.
$x_i \in H_n$과 $f_i \in I^c$에 대해 $x = \sum f_i x_i$로 쓰자.
$x_i$에 대한 $\psi_n$의 구성에서와 같이 $y_i, z_i, w_i$를 택하자.
그러면 $x$에 대한 $\psi_n$의 계산에서는
$y = \sum f_iy_i$, $z = \sum f_i z_i$, $w = \sum f_i w_i$를 택할 수
있고, $\psi_n(x)$가 $z$의 동치류로 주어짐을 알 수 있다.
$w = \sum f_i w_i$가 $I^nL$에 속하므로 이것의
$H_n/I^{n - c}H_n$에서의 상은 $y$의 동치류와 같다. 이로써 (5)가
증명된다.

\medskip\noindent
(6)의 증명. $H$에서 동치류가 $x$인 $y \in \Ker(\beta)$를 택하자.
$x$가 $H_n$에서 영으로 가면 $y \in I^nL + \Im(\alpha)$이다. 따라서
어떤 $v \in K$에 대해
$y - \alpha(v) \in \Ker(\beta) \cap I^nL$이다. 그러면
$y - \alpha(v) \in I^{n - c_2}\Ker(\beta)$이므로 $H/I^nH$에서
$y$의 동치류는 $I^{c_2}$에 의해 소멸된다. 마지막으로
$x \in H_n$이 $y \in \beta^{-1}(I^nM)$의 동치류라고 하자. 위와 같이
$z \in \Ker(\beta)$와 $w \in I^{n - c_1}L$에 대해 $y = z + w$로
쓴다. 명백히 $f \in I^{c_1}$이면 $fx$는
% P02-E0477: frozen calculation uses fy+fw≡fy and then the class of fy; preserved exactly.
$\Im(\alpha) + I^nL$을 법으로 한 $fy + fw \equiv fy$의 동치류이고,
따라서 원하는 대로 $fx$는 $H$에서 $fy$의 동치류의 상이다.
\end{proof}

\begin{lemma}
\label{more-algebra-lemma-kollar-kovacs-pseudo-coherent}
\begin{reference}
% P02-E0479: first frozen reference locus omits the established accent in Kovács.
2018년 23/02의 Kovacs 이메일.
\end{reference}
$I$를 Noether 환 $A$의 아이디얼이라 하자. $K \in D(A)$가
유사연접이라고 하자. $K_n = K \otimes_A^\mathbf{L} A/I^n$으로 두자.
그러면 모든 $i \in \mathbf{Z}$에 대해 계 $H^i(K_n)$은
Mittag--Leffler 조건을 만족하며,
$\lim H^i(K)/I^nH^i(K)$는 $\lim H^i(K_n)$와 같다.
\end{lemma}

\begin{proof}
$K$를 유한 자유 $A$-가군들로 이루어진 위로 유계인 복합체
$P^\bullet$로 나타낼 수 있다. 그러면 $K_n$은
$P^\bullet/I^nP^\bullet$로 나타내어진다.
% P02-E0480: frozen proof fragment lacks a finite verb; Korean supplies the intended conclusion.
따라서 보조정리 \ref{more-algebra-lemma-consequence-Artin-Rees}에 의해
Mittag--Leffler 조건이 성립한다. 마지막 진술은 이제 보조정리
\ref{more-algebra-lemma-kollar-kovacs}에서 따른다.
\end{proof}

\begin{lemma}
\label{more-algebra-lemma-derived-completion-plain-completion}
$A$를 Noether 환이라 하자. $I \subset A$를 아이디얼이라 하자.
$M^\bullet$을 유한 $A$-가군들로 이루어진 유계 복합체라 하자. 사상들의
역계
% P02-E0481: first marked display in this section lacks its surrounding sentence punctuation.
$$
M^\bullet \otimes_A^\mathbf{L} A/I^n \longrightarrow M^\bullet/I^nM^\bullet
$$
은 $D(A)$의 pro-대상들의 동형사상을 정의한다.
\end{lemma}

\begin{proof}
$I = (f_1, \ldots, f_r)$이라 하자. $K_n \in D(A)$를
$f_1^n, \ldots, f_r^n$ 위의 Koszul 복합체로 나타내어지는 대상이라 하자.
역계들의 pro-동형사상을 유도하는 사상 $K_n \to A/I^n$이 있음을
상기하라. 보조정리 \ref{more-algebra-lemma-sequence-Koszul-complexes}를 보라.
따라서
$$
M^\bullet \otimes_A^\mathbf{L} K_n \longrightarrow M^\bullet/I^nM^\bullet
$$
이 $D(A)$의 pro-대상들의 동형사상을 정의함을 보이면 충분하다.
$K_n$은 차수 $-r, \ldots, 0$에 놓인 유한 자유 $A$-가군들의
복합체로 나타내어지므로, 모든 $n$에 대하여 표시된 화살표의 원천과
목표의 상동군이 $[a, b]$ 밖에서 소멸하는
$a, b \in \mathbf{Z}$가 존재한다. 따라서 유도 범주론, 보조정리
\ref{derived-lemma-pro-isomorphism-bis}를 적용하면, 임의의
$i \in \mathbf{Z}$에 대하여 사상들
$$
H^i(M^\bullet \otimes_A^\mathbf{L} A/I^n) \to H^i(M^\bullet/I^nM^\bullet)
$$
이 $A$-가군의 pro-계들의 동형사상을 정의함을 보이는 것으로 충분함을
알 수 있다. 이를 보기 위해 $P^\bullet$이 유한 자유 $A$-가군들로
이루어진 위로 유계인 복합체인 유사동형사상
$P^\bullet \to M^\bullet$을 택하자. 위의 화살표들은 사상들
% P02-E0481: second marked display lacks terminal punctuation before “These define”.
$$
H^i(P^\bullet/I^nP^\bullet) \to H^i(M^\bullet/I^nM^\bullet)
$$
로 주어진다. 이들은 보조정리 \ref{more-algebra-lemma-consequence-Artin-Rees}에 의해
pro-계들의 동형사상을 정의한다. 실제로 그 보조정리는 양쪽이
$H^i = H^i(M^\bullet) = H^i(P^\bullet)$인 pro-계
$H^i/I^nH^i$와 동형임을 보인다.
\end{proof}

\begin{lemma}
\label{more-algebra-lemma-hom-systems-ML}
$A$를 Noether 환이라 하자. $I \subset A$를 아이디얼이라 하자.
$M$, $N$을 유한 $A$-가군이라 하자. $M_n = M/I^nM$과
$N_n = N/I^nN$으로 두자. 그러면 다음이 성립한다.
\begin{enumerate}
\item 계 $(\Hom_A(M_n, N_n))$과 $(\text{Isom}_A(M_n, N_n))$은
Mittag--Leffler 조건을 만족한다.
\item 모든 $n$에 대해
$$
\Hom_A(M, N)/I^n\Hom_A(M, N) \to \Hom_A(M_n, N_n)
$$
의 핵과 여핵이 $I^c$에 의해 소멸되는 $c \geq 0$이 존재한다.
\item
% P02-E0482: frozen item (3) lacks punctuation or coordination before the final item.
$\lim \Hom_A(M_n, N_n) =\Hom_A(M, N)^\wedge =
\Hom_{A^\wedge}(M^\wedge, N^\wedge)$이다.
\item $\lim \text{Isom}_A(M_n, N_n) =
\text{Isom}_{A^\wedge}(M^\wedge, N^\wedge)$이다.
\end{enumerate}
여기서 ${}^\wedge$는 보통의 $I$-진 완비화를 뜻한다.
\end{lemma}

\begin{proof}
$\Hom_A(M_n, N_n) = \Hom_A(M, N_n)$임에 유의하라. 표시
$$
A^{\oplus t} \to A^{\oplus s} \to M \to 0
$$
를 택하자.
% P02-E0483: frozen source incorrectly calls this Hom functor right exact.
오른쪽 완전 함자 $\Hom_A(-, N)$을 적용하여 복합체
$$
0 \xrightarrow{\alpha} N^{\oplus s} \xrightarrow{\beta} N^{\oplus t}
$$
를 얻는다. 그 중간 상동군은 $\Hom_A(M, N)$이고, $n \geq 0$에
대하여 복합체
$$
0 \xrightarrow{\alpha_n} N_n^{\oplus s} \xrightarrow{\beta_n} N_n^{\oplus t}
$$
의 상동군은 $\Hom_A(M_n, N_n)$이다. 이 $A$, $I$, $\alpha$, $\beta$에
대하여 보조정리 \ref{more-algebra-lemma-consequence-Artin-Rees}에서와 같은
$c \geq 0$을 택하자. 그 보조정리의 (3)에 의해
$(\Hom_A(M_n, N_n))$의 Mittag--Leffler 조건을 얻는다. 사상들
$\Hom_A(M, N)/I^n\Hom_A(M, N) \to \Hom_A(M_n, N_n)$의 핵과 여핵은
% P02-E0484: frozen source contains the stray text “[art” before its part reference.
그 보조정리의 (6)에 의해 $I^c$에 의해 소멸된다. 그 보조정리의
(2)에 의해
$\lim \Hom_A(M_n, N_n) = \Hom_A(M, N)^\wedge$임을 얻는다. 등식
% P02-E0481: third marked display lacks terminal punctuation before the next sentence.
$$
\Hom_{A^\wedge}(M^\wedge, N^\wedge) = \lim \Hom_A(M_n, N_n)
$$
은 $M^\wedge = \lim M_n$ 및
$M_n = M^\wedge/I^nM^\wedge$이라는 사실과 $N$에 대한 대응하는
사실들에서 형식적으로 따른다. 대수학, 보조정리
\ref{algebra-lemma-completion-complete}를 보라.

\medskip\noindent
동형사상들에 대한 결과는 $(\Hom(M_n, N_n))$과
$(\Hom(N_n, M_n))$ 양쪽에 준동형사상에 대한 경우를 적용하고 다음
사실을 사용하면 따른다. $n > m > 0$에 대하여 사상
$\alpha : M_n \to N_n$과 $\beta : N_n \to M_n$이 있고
% P02-E0485: frozen source says “induce an isomorphisms”; Korean uses grammatical number.
동형사상 $M_m \to N_m$과 $N_m \to M_m$을 유도한다면,
$\alpha$와 $\beta$는 동형사상이다. 실제로
% P02-E1156: frozen proof checks only one composite before concluding both maps are isomorphisms.
$\alpha \circ \beta$는 Nakayama 보조정리(대수학, 보조정리
\ref{algebra-lemma-NAK})에 의해 전사이고, 따라서 대수학, 보조정리
\ref{algebra-lemma-fun}에 의해 $\alpha \circ \beta$는 동형사상이다.
\end{proof}

\begin{lemma}
\label{more-algebra-lemma-isomorphic-completions}
$A$를 Noether 환이라 하자. $I \subset A$를 아이디얼이라 하자.
$M$, $N$을 유한 $A$-가군이라 하자. $M_n = M/I^nM$과
$N_n = N/I^nN$으로 두자. 모든 $n$에 대해 $M_n \cong N_n$이면
$A^\wedge$-가군으로서 $M^\wedge \cong N^\wedge$이다.
\end{lemma}

\begin{proof}
보조정리 \ref{more-algebra-lemma-hom-systems-ML}에 의해 계
$(\text{Isom}_A(M_n, N_n))$은 Mittag--Leffler 조건을 만족한다. 가정에
의해 각 집합 $\text{Isom}_A(M_n, N_n)$은 공집합이 아니다. 따라서
$\lim \text{Isom}_A(M_n, N_n)$은 공집합이 아니다.
$\lim \text{Isom}_A(M_n, N_n) = \text{Isom}_{A^\wedge}(M^\wedge, N^\wedge)$
이므로 동형사상을 얻는다.
\end{proof}

\begin{remark}
\label{more-algebra-remark-weird-systems}
$I$를 Noether 환 $A$의 아이디얼이라 하자. $n \geq 1$에 대해
$A_n = A/I^n$으로 두자. 다음 범주를 생각하자.
\begin{enumerate}
\item 대상은 수열 $\{E_n\}_{n \geq 1}$이며, 여기서 $E_n$은 유한
$A_n$-가군이다.
\item 사상 $\{E_n\} \to \{E'_n\}$은 사상들
$$
\varphi_n : I^cE_n \longrightarrow E'_n/E'_n[I^c]
\quad\text{for }n \geq c
$$
로 주어진다. 여기서 $E'_n[I^c]$는 꼬임 부분가군(절
\ref{more-algebra-section-torsion})이고, 다음 동치까지 고려한다. $(c, \varphi_n)$과
$(c + 1, \overline{\varphi}_n)$이 같다고 하는데, 여기서
$\overline{\varphi}_n : I^{c + 1}E_n \longrightarrow E'_n/E'_n[I^{c + 1}]$
은 유도된 사상이다.
\end{enumerate}
$(c, \varphi_n) : \{E_n\} \to \{E'_n\}$과
$(c', \varphi'_n) : \{E'_n\} \to \{E''_n\}$의 합성은 자명한 합성들
% P02-E0486: frozen middle quotient need not be defined because its denominator need not lie in its numerator.
$$
I^{c + c'}E_n \to I^{c'}E'_n/E'_n[I^{c}] \to E''_n/E''_n[I^{c + c'}]
$$
로 정의되며, 이는 $n \geq c + c'$에 대한 것이다. 이것이 범주라는
확인은 생략한다.
\end{remark}

\begin{lemma}
\label{more-algebra-lemma-iso}
주석 \ref{more-algebra-remark-weird-systems}의 범주의 사상 $(c, \varphi_n)$이
동형사상일 필요충분조건은 모든 $n \gg 0$에 대하여
$\Ker(\varphi_n)$과 $\Coker(\varphi_n)$이 $I^{c'}$-꼬임이 되는
$c' \geq 0$이 존재하는 것이다.
\end{lemma}

\begin{proof}
$c' \geq c$이고 모든 $n$에 대해 $\Ker(\varphi_n)$과
$\Coker(\varphi_n)$이 $I^{c'}$-꼬임이라고 가정해도 좋으며 그렇게
한다. $n \geq c'$이고 $x \in I^{c'}E'_n$이면
$\Coker(\varphi_n)$이 $I^{c'}$에 의해 소멸되므로
$x = \varphi_n(y) \bmod E'_n[I^c]$인 $y \in I^cE_n$을 택할 수 있다.
$\psi_n(x)$를 $E_n/E_n[I^{c'}]$에서 $y$의 동치류로 두자.
$x = \varphi_n(y') \bmod E'_n[I^c]$인 다른 선택 $y' \in I^cE_n$에
대해 차 $y - y'$은 $E'_n/E'_n[I^c]$에서 영으로 가므로
$I^cE_n$에서 $I^{c'}$에 의해 소멸된다. 따라서 사상들
$\psi_n : I^{c'}E'_n \to E_n/E_n[I^{c'}]$은 잘 정의된다.
$(c', \psi_n)$이 이 범주에서 $(c, \varphi_n)$의 역이라는 확인은
생략한다.
\end{proof}

\begin{lemma}
\label{more-algebra-lemma-dejong-kollar-kovacs}
\begin{reference}
% P02-E0479: second frozen reference locus omits the established accents in the three names.
2018년 23/02의 Janos Kollar, Sandor Kovacs, Johan de Jong 사이의 이메일
서신.
\end{reference}
$I$를 Noether 환 $A$의 아이디얼이라 하자. $M$과 $N$을 유한
$A$-가군이라 하자. $A_n = A/I^n$, $M_n = M/I^nM$,
$N_n = N/I^nN$으로 쓰자. 모든 $i \geq 0$에 대하여 대상들
$$
\{\Ext^i_A(M, N)/I^n\Ext^i_A(M, N)\}_{n \geq 1}
\quad\text{and}\quad
\{\Ext^i_{A_n}(M_n, N_n)\}_{n \geq 1}
$$
은 주석 \ref{more-algebra-remark-weird-systems}의 범주 $\mathcal{C}$에서
동형이다.
\end{lemma}

\begin{proof}
짧은 완전열
$$
0 \to K \to A^{\oplus r} \to M \to 0
$$
을 택하고 $K_n = K/I^nK$로 두자. $n \geq 1$에 대해
$K(n) = \Ker(A_n^{\oplus r} \to M_n)$으로 정의하면 완전열들
$$
0 \to K(n) \to A_n^{\oplus r} \to M_n \to 0
$$
과 전사 $K_n \to K(n)$이 있다. 실제로 보조정리
\ref{more-algebra-lemma-consequence-Artin-Rees}에 의해 $c \geq 0$과 서로
‘거의 역’인 사상들 $K(n) \to K_n/I^{n - c}K_n$이 있다.
% P02-E0487: frozen sentence lacks a finite verb before “which witness”.
$I^{n - c}K_n \subset K_n[I^c]$이므로 이 사상들은 계
$\{K(n)\}_{n \geq 1}$과 $\{K_n\}_{n \geq 1}$이
$\mathcal{C}$에서 동형이라는 사실을 증명한다.

\medskip\noindent
계들
$$
\{\Ext^i_{A_n}(K(n), N_n)\}_{n \geq 1}
\quad\text{and}\quad
\{\Ext^i_{A_n}(K_n, N_n)\}_{n \geq 1}
$$
이 범주 $\mathcal{C}$에서 동형임을 주장한다. 실제로 전사
$K_n \to K(n)$의 핵들은 $I^c$에 의해 소멸되므로 사상들
$$
\Ext^i_{A_n}(K(n), N_n) \to \Ext^i_{A_n}(K_n, N_n)
$$
을 정한다. 이들의 핵과 여핵은 $I^c$에 의해 소멸된다. 따라서
보조정리 \ref{more-algebra-lemma-iso}에 의해 주장을 얻는다.

\medskip\noindent
$i \geq 2$에 대해 동형사상들
$$
\Ext^{i - 1}_A(K, N) = \Ext^i_A(M, N)
\quad\text{and}\quad
\Ext^{i - 1}_{A_n}(K(n), N_n) = \Ext^i_{A_n}(M_n, N_n)
$$
이 있다. 이로써 $i = 0, 1$인 경우에 보조정리를 증명하면 충분함을
알 수 있다.

\medskip\noindent
$i = 0, 1$에 대해서는 가환 그림
$$
\xymatrix{
0 \ar[r] &
\Hom(M, N) \ar[r] \ar[dd] &
N^{\oplus r} \ar[r]_-\varphi \ar[dd] &
\Hom(K, N) \ar[r] \ar[d] &
\Ext^1(M, N) \ar[r] &
0 \\
& & &
\Hom(K_n, N_n)
\\
0 \ar[r] &
\Hom(M_n, N_n) \ar[r] &
N_n^{\oplus r} \ar[r] &
\Hom(K(n), N_n) \ar[r] \ar[u] &
\Ext^1(M_n, N_n) \ar[r] &
0
}
$$
을 생각하자. 보조정리 \ref{more-algebra-lemma-hom-systems-ML}에 의해
% P02-E0488: frozen prose has an extra conjunction and inconsistent cokernel agreement.
$\Hom(M, N)/I^n \Hom(M, N) \to \Hom(M_n, N_n)$ 및
$\Hom(K, N)/I^n \Hom(K, N) \to \Hom(K_n, N_n)$의 핵과 여핵은
$n$과 무관한 어떤 $c \geq 0$에 대해 $I^c$-꼬임임을 알 수 있다.
위에서 단사 $\Hom(K(n), N_n) \to \Hom(K_n, N_n)$의 여핵은
필요하다면 $c$를 늘렸을 때 $I^c$에 의해 소멸됨을 보았다. 그런
$c$에 대해 그림에 들어가는 사상
$\delta_n : I^c\Hom(K, N)/I^n\Hom(K, N) \to \Hom(K(n), N_n)$을
얻는다(정확한 정식화는 생략한다).
$\Hom(K, N)/I^n \Hom(K, N) \to \Hom(K_n, N_n)$과
$\Hom(K(n), N_n) \to \Hom(K_n, N_n)$에 대해 같은 사실이 성립함을
알고 있으므로, 필요하다면 $c$를 늘렸을 때 $\delta_n$의 핵과 여핵은
$I^c$에 의해 소멸된다. 이제 실선 그림의 가환성
$$
\xymatrix{
\varphi^{-1}(I^c\Hom(K, N)) \ar[r]_-\varphi \ar[d] &
I^c\Hom(K, N)/I^n\Hom(K, N) \ar[r] \ar[d]^{\delta_n} &
I^c\Ext^1(M, N)/I^n\Ext^1(M, N) \ar[r] \ar@{..>}[d] & 0 \\
N_n^{\oplus r} \ar[r] &
\Hom(K(n), N_n) \ar[r] &
\Ext^1(M_n, N_n) \ar[r] & 0
}
$$
을 사용하여 점선 화살표를 정의할 수 있다. 간단한 그림 추적(생략)을
하면, 마지막으로 한 번 더 $c$를 늘렸을 때 점선 화살표의 핵과 여핵이
% P02-E0489: frozen source misspells “by” as “buy”.
$I^c$에 의해 소멸됨을 알 수 있다.
\end{proof}

\begin{remark}
\label{more-algebra-remark-awkward}
보조정리 \ref{more-algebra-lemma-dejong-kollar-kovacs}의 진술이 어색한 까닭 중
일부는 변하는 $n$에 대해 가군 $\Ext^i_{A_n}(M_n, N_n)$들 사이에
자명한 사상들이 없다는 사실이다. 그 보조정리에서 다음을 결론지을
수 있다. $c \geq 0$이 존재하여 $m \gg n \gg 0$일 때 (표준) 사상들
% P02-E0490: frozen source forms the left Ext groups over A_n although their arguments are M_m and N_m.
$$
I^c\Ext^i_{A_n}(M_m, N_m)/I^n\Ext^i_{A_n}(M_m, N_m) \to
\Ext^i_{A_n}(M_n, N_n)/\Ext^i_{A_n}(M_n, N_n)[I^c]
$$
이 있으며, 그 핵과 여핵은 $I^c$에 의해 소멸된다. 이것이 가군의 계를
얻는 (약한) 의미이다.
\end{remark}

\begin{example}
\label{more-algebra-example-has-to-be-awkward}
$k$를 체라 하자. $A = k[[x, y]]/(xy)$라 하자. 표기를 약간 남용하여
% P02-E0491: frozen naming sentence omits “by”; Korean supplies normal naming syntax.
$x$와 $y$로 $A$에서 $x$와 $y$의 상을 나타내자. $I = (x)$라 하고
$M = A/(y)$라 하자. 자유 분해
% P02-E0481: fourth marked display group lacks surrounding punctuation in the frozen prose.
$$
\ldots \to
A \xrightarrow{y}
A \xrightarrow{x}
A \xrightarrow{y} A \to M \to 0
$$
가 있다. 따라서
$$
\Ext^2_A(M, N) = N[y]/xN
$$
을 얻는다. 여기서 $N[y] = \Ker(y : N \to N)$이다.
$A_n = A/I^n$, $M_n = M/I^nM$, $N_n = N/I^nN$으로 나타내자.
각 $n$에 대해 자유 분해
$$
\ldots \to
A_n^{\oplus 2} \xrightarrow{y, x^{n - 1}}
A_n \xrightarrow{x}
A_n \xrightarrow{y} A_n \to M_n \to 0
$$
가 있다. 따라서
$$
\Ext^2_{A_n}(M_n, N_n) = (N_n[y] \cap N_n[x^{n - 1}])/xN_n
$$
을 얻는다. 여기서 $N_n[y] = \Ker(y : N_n \to N_n)$이고
% P02-E0492: frozen second torsion notation drops the subscript n; formal span preserved.
$N[x^{n - 1}] = \Ker(x^{n - 1} : N_n \to N_n)$이다.
$N = A/(y)$로 택하자. 그러면
$$
\Ext^2_A(M, N) = N[y]/xN = N/xN \cong k
$$
이지만
$$
\Ext^2_{A_n}(M_n, N_n) = (N_n[y] \cap N_n[x^{n - 1}])/xN_n =
N_n[x^{n - 1}]/xN_n = 0
$$
은 모든
% P02-E0493: frozen conclusion uses r although the running index is n.
$r$에 대해 성립한다. 실제로 $N_n = k[x]/(x^n)$이고 열
$$
N_n \xrightarrow{x} N_n \xrightarrow{x^{n - 1}} N_n
$$
은 완전하다. 따라서 보조정리 \ref{more-algebra-lemma-dejong-kollar-kovacs}와
같은 결과를 얻으려면 어떤 종류의 $I$-거듭제곱 꼬임을 무시해야 한다.
\end{example}

\begin{lemma}
\label{more-algebra-lemma-not-awkward}
\begin{reference}
% P02-E0479: third frozen reference locus omits the established accents in the three names.
2018년 23/02의 Janos Kollar, Sandor Kovacs, Johan de Jong 사이의 이메일
서신.
\end{reference}
$A \to B$를 Noether 환의 평탄 준동형이라 하자. $I \subset A$를
아이디얼이라 하자.
% P02-E0494: frozen statement declares only A-module structures although it then forms Ext_B and B-resolutions.
$M, N$을 $A$-가군이라 하자. $B_n = B/I^nB$, $M_n = M/I^nM$,
$N_n = N/I^nN$으로 두자. $M$이 $A$ 위에서 평탄이면, 모든
$i \in \mathbf{Z}$에 대해
$$
\lim \Ext^i_B(M, N)/I^n \Ext^i_B(M, N) =
\lim \Ext^i_{B_n}(M_n, N_n)
$$
이다.
\end{lemma}

\begin{proof}
분해
% P02-E0495: frozen proof misspells “modules”; Korean uses the intended noun.
% P02-E0481: fifth marked display group lacks surrounding punctuation in the frozen prose.
$$
\ldots \to P_2 \to P_1 \to P_0 \to M \to 0
$$
를 유한 자유 $B$-가군 $P_i$들로 택하자.
$P_{i, n} = P_i/I^nP_i$로 두자. $M$과 $B$가 $A$ 위에서 평탄이므로 열
$$
\ldots \to P_{2, n} \to P_{1, n} \to P_{0, n} \to M_n \to 0
$$
은 완전하다. 한편으로 복합체
$$
\Hom_B(P_0, N) \to \Hom_B(P_1, N) \to \Hom_B(P_2, N) \to \ldots
$$
는 가군 $\Ext^i_B(M, N)$들을 계산하고, 다른 한편으로 복합체
$$
\Hom_{B_n}(P_{0, n}, N_n) \to \Hom_{B_n}(P_{1, n}, N_n) \to
\Hom_{B_n}(P_{2, n}, N_n) \to \ldots
$$
는 가군 $\Ext^i_{B_n}(M_n, N_n)$들을 계산한다. 등식
$$
\Hom_{B_n}(P_{i, n}, N_n) = \Hom_B(P_i, N)/I^n \Hom_B(P_i, N)
$$
이 성립하므로 보조정리 \ref{more-algebra-lemma-consequence-Artin-Rees}의 (2)에서
결과를 얻는다.
\end{proof}

\section{가군의 계, II}
\label{more-algebra-section-systems-bis}

\noindent
$I$를 Noether 환 $A$의 아이디얼이라 하자. 절
\ref{more-algebra-section-systems}에서는 유한 $A$-가군 $M$에 대한 $M/I^nM$
꼴의 계를 생각할 때 어떤 일이 일어나는지 살펴보았다. 이 절에서는
그 대신 계 $I^nM$을 생각한다.

\begin{lemma}
\label{more-algebra-lemma-consequence-Artin-Rees-bis}
$I$를 Noether 환 $A$의 아이디얼이라 하자.
$
K \xrightarrow{\alpha} L \xrightarrow{\beta} M
$
을 유한 $A$-가군의 복합체라 하자. $H = \Ker(\beta)/\Im(\alpha)$로
두자. $n \geq 0$에 대해
$$
I^nK \xrightarrow{\alpha_n} I^nL \xrightarrow{\beta_n} I^nM
$$
을 유도된 복합체라 하자. $H_n = \Ker(\beta_n)/\Im(\alpha_n)$으로
두자. 그러면 표준 $A$-가군 사상들
$$
\ldots \to H_3 \to H_2 \to H_1 \to H
$$
이 있다. $c > 0$이 존재하여 $n \geq c$이면 $H_n \to H$의 상은
$I^{n - c}H$에 포함되고, 표준 $A$-가군 사상
$I^nH \to H_{n - c}$가 있어 합성들
$$
I^n H \to H_{n - c} \to  I^{n - 2c}H
\quad\text{and}\quad
H_n \to I^{n - c}H \to H_{n - 2c}
$$
은 표준 사상들이다. 특히 역계 $(H_n)$과 $(I^nH)$은
$\text{Mod}_A$의 pro-대상으로서 동형이다.
\end{lemma}

\begin{proof}
% P02-E0496: frozen quotient omits parentheses around its intersected numerator.
$H_n = \Ker(\beta) \cap I^nL/\alpha(I^nK)$이다.
$\Ker(\beta) \cap I^nL \subset \Ker(\beta) \cap I^{n - 1}L$이고
$\alpha(I^nK) \subset \alpha(I^{n - 1}K)$이므로 사상
$H_n \to H_{n - 1}$을 얻는다. $H_1 \to H$의 사상도 마찬가지이다.

\medskip\noindent
Artin--Rees에 의해, $n \geq c_1$이면
$\Im(\alpha) \cap I^nL \subset \alpha(I^{n - c_1}K)$이고
$n \geq c_2$이면
$\Ker(\beta) \cap I^nL \subset I^{n - c_2}\Ker(\beta)$인
$c_1, c_2 \geq 0$을 택할 수 있다. 대수학, 보조정리
\ref{algebra-lemma-map-AR}과 \ref{algebra-lemma-Artin-Rees}를 보라.
$c = c_1 + c_2$로 두자.

\medskip\noindent
$c \geq c_2$의 선택에서, $n \geq c$이면 $H_n \to H$의 상이
$I^{n - c}H$에 포함된다는 것이 바로 따른다.

\medskip\noindent
$n \geq c$라 하자. $\psi_n : I^nH \to H_{n - c}$를 다음과 같이
정의한다. $x \in I^nH$라 하자. $x$를 나타내는
$y \in I^n\Ker(\beta)$를 택하자. $\psi_n(x)$를 $H_{n - c}$에서
$y$의 동치류로 둔다. 이것이 잘 정의됨을 보기 위해 같은 $x$에 대해
위와 같은 두 번째 선택 $y'$을 했다고 하자. 그러면
$y' - y \in \Im(\alpha)$이다. $c \geq c_1$의 선택에 의해
$y' - y \in \alpha(I^{n - c}K)$라고 결론지을 수 있고, 이는 $y$와
$y'$이 $H_{n - c}$의 같은 원소를 나타냄을 뜻한다. 따라서
$\psi_n$은 잘 정의된다.

\medskip\noindent
합성 $I^n H \to H_{n - c} \to  I^{n - 2c}H$과
$H_n \to I^{n - c}H \to H_{n - 2c}$에 대한 진술은 우리의 정의에서
바로 따른다.
\end{proof}

\begin{lemma}
\label{more-algebra-lemma-ext-factors}
$A$를 Noether 환이라 하자. $I \subset A$를 아이디얼이라 하자.
$M$, $N$을 $A$-가군이라 하고 $M$은 유한이라고 하자. 각 $p > 0$에
대해 $c \geq 0$이 존재하여, $n \geq c$이면 사상
$\Ext_A^p(M, N) \to \Ext_A^p(I^nM, N)$은
$\Ext^p_A(I^nM, I^{n - c}N) \to \Ext_A^p(I^nM, N)$을 통하여
분해된다.
\end{lemma}

\begin{proof}
$p = 0$일 때 $\varphi : M \to N$이 $A$-선형 사상이면,
% P02-E0497: frozen p=0 argument quantifies f_i only in A although membership in I^n is needed.
$f_i \in A$와 $m_i \in M$에 대해
$\varphi(\sum f_i m_i) = \sum f_i \varphi(m_i)$이다. 따라서 모든
$n$에 대해 $\varphi$는 사상 $I^nM \to I^nN$을 유도하고 결과는
$c = 0$으로 참이다.

\medskip\noindent
짧은 완전열 $0 \to K \to A^{\oplus t} \to M \to 0$을 택하자.
각 $n$에 대해 짧은 완전열
$0 \to L_n \to A^{\oplus s_n} \to I^nM \to 0$을 택한다.
짧은 완전열들의 사상
$$
\xymatrix{
0 \ar[r] &
L_n \ar[r] \ar[d] &
A^{\oplus s_n} \ar[r] \ar[d] &
I^nM \ar[r] \ar[d] & 0 \\
0 \ar[r] &
K \ar[r] &
A^{\oplus t} \ar[r] &
M \ar[r] & 0
}
$$
을 구성하여 $A^{\oplus s_n} \to A^{\oplus t}$의 상이
$(I^n)^{\oplus t}$에 놓이게 할 수 있음은 명백하다.
Artin--Rees(대수학, 보조정리 \ref{algebra-lemma-Artin-Rees})에 의해
$c \geq 0$이 존재하여, $n \geq c$이면 $L_n \to K$는
$I^{n - c}K$를 통하여 분해된다.

\medskip\noindent
$p = 1$에 대해서는 위의 선택들이 실선 가환 그림
$$
\xymatrix{
\Hom_A(A^{\oplus s_n}, N) \ar[r] &
\Hom_A(L_n, N) \ar[r] &
\Ext_A^1(I^nM, N) \ar[r] & 0 \\
\Hom_A((I^n)^{\oplus t}, I^{n - c}N) \ar[r] \ar[u] &
\Hom_A(K \cap (I^n)^{\oplus t}, I^{n - c}N) \ar[r] \ar[u] &
\Ext_A^1(I^nM, I^{n - c}N) \ar[u] \\
\Hom_A(A^{\oplus t}, N) \ar[r] \ar[u] &
\Hom_A(K, N) \ar[r] \ar[u] &
\Ext_A^1(M, N) \ar@{..>}[u] \ar[r] & 0
}
$$
을 유도하며, 그 가로 화살표들은 완전하다. 아래 가운데 세로 화살표가
생기는 것은 $K \cap (I^n)^{\oplus t} \subset I^{n - c}K$이고,
따라서 첫 문단의 논증에 의해 임의의 $A$-선형 사상 $K \to N$이
% P02-E0498: frozen prose claims an induced map on the whole ambient power instead of its intersection with K.
$A$-선형 사상 $(I^n)^{\oplus t} \to I^{n - c}N$을 유도하기 때문이다.
따라서 원하는 점선 화살표를 얻는다.

\medskip\noindent
$p > 1$에 대해서는 가환 그림
$$
\xymatrix{
\Ext^{p - 1}_A(I^{n - c}K, N) \ar[r] &
\Ext^{p - 1}_A(L_n, N) \ar[r] &
\Ext_A^p(I^nM, N) \\
\Ext^{p - 1}_A(K, N) \ar[rr] \ar[u] & &
\Ext_A^p(M, N) \ar[u]
}
$$
을 얻으며, 아래 가로 화살표는 동형사상이다. $p$에 관한 귀납법에
의해 어떤 $c' \geq 0$과 모든 $n \geq c + c'$에 대해 왼쪽 세로 사상은
$\Ext^{p - 1}_A(I^{n - c}K, I^{n - c - c'}N)$을 통하여 분해된다.
합성
$\Ext^{p - 1}_A(I^{n - c}K, I^{n - c - c'}N) \to
\Ext^{p - 1}_A(L_n, I^{n - c - c'}N) \to \Ext^p_A(I^nM, I^{n - c - c'}N)$
을 사용하여 원하는 분해를 얻는다($n \geq c + c'$에 대해 $c$를
$c + c'$으로 바꾼다).
\end{proof}

\begin{lemma}
\label{more-algebra-lemma-ext-annihilated}
$A$를 Noether 환이라 하자. $I \subset A$를 아이디얼이라 하자.
$M$, $N$을 $A$-가군이라 하고, $M$은 유한이며 $N$은 $I$의 어떤
거듭제곱에 의해 소멸된다고 하자. 각 $p > 0$에 대해 사상
$\Ext_A^p(M, N) \to \Ext_A^p(I^nM, N)$이 영이 되는 $n$이 존재한다.
\end{lemma}

\begin{proof}
보조정리 \ref{more-algebra-lemma-ext-factors}와 어떤 $m > 0$에 대해
$I^mN = 0$이라는 사실의 바로 된 귀결이다.
\end{proof}

\begin{lemma}
% P02-E0499: frozen label misspells topology and its quantifier says “an c”; label preserved.
\label{more-algebra-lemma-ext-induced-toplogy}
$A$를 Noether 환이라 하자. $I \subset A$를 아이디얼이라 하자.
$K \in D(A)$가 유사연접이고 $M$이 유한 $A$-가군이라고 하자.
각 $p \in \mathbf{Z}$에 대해 $c$가 존재하여, $n \geq c$이면
$\Ext_A^p(K, I^nM) \to \Ext_A^p(K, M)$의 상은
$I^{n - c}\Ext_A^p(K, M)$에 포함된다.
\end{lemma}

\begin{proof}
$K$를 나타내는 유한 자유 $A$-가군들로 이루어진 위로 유계인 복합체
$P^\bullet$을 택하자. 그러면 $\Ext_A^p(K, M)$은 다음 복합체의
상동군이다.
% P02-E0500: frozen display uses undefined F in all three terms after choosing P; formal spans preserved.
$$
\Hom_A(F^{-p + 1}, M) \xrightarrow{a}
\Hom_A(F^{-p}, M) \xrightarrow{b}
\Hom_A(F^{-p - 1}, M)
$$
그리고 $\Ext_A^p(K, I^nM)$은 이 유한 $A$-가군들을 각각 그것의
$I^n$배로 바꾸어 계산한다.
% P02-E0501: frozen conclusion lacks its governing verb.
따라서 결과는 보조정리 \ref{more-algebra-lemma-consequence-Artin-Rees-bis}에서
따른다(실제로는 훨씬 더 많은 것이 참이다).
\end{proof}

\noindent
상황 \ref{more-algebra-situation-koszul}에서 복합체 $I_n^\bullet$을 정의하여
구별 삼각형
% P02-E0502: first affected Koszul display lacks sentence-ending punctuation in the frozen source.
$$
I_n^\bullet \to A \to K_n^\bullet \to I_n^\bullet[1]
$$
을 호모토피를 법으로 한 $A$-가군 복합체의 삼각 범주 $K(A)$에서
얻는다. 실제로 $I_n^\bullet = \sigma_{\leq -1}K_n^\bullet[-1]$로
둔다. 항별로 분할되는 복합체의 짧은 완전열
% P02-E0502: second affected Koszul display lacks sentence-ending punctuation in the frozen source.
$$
0 \to A \to K_n^\bullet \to I_n^\bullet[1] \to 0
$$
이 있고, 이들은 $K(A)$의 삼각 구조의 정의에 의해 구별 삼각형들을
정의한다. 이들을 회전하면 위의 원하는 구별 삼각형들이 정해진다.
$I_n^0 = A^{\oplus r} \to A$는 $i$번째 인자 위에서 $f_i^n$을 곱하는
사상으로 주어짐에 유의하라. 따라서 $I_n^\bullet \to A$는
$$
I_n^\bullet \to (f_1^n, \ldots, f_r^n) \to A
$$
로 분해된다. 실제로 짧은 완전열
$$
0 \to H^{-1}(K_n^\bullet) \to H^0(I_n^\bullet) \to (f_1^n, \ldots, f_r^n) \to 0
$$
이 있고, 모든 $i < 0$에 대해
$H^i(I_n^\bullet) = H^{i - 1}(K_n^\bullet)$이다.
사상 $K_{n + 1}^\bullet \to K_n^\bullet$은 사상
$I_{n + 1}^\bullet \to I_n^\bullet$을 유도하고 가환 그림
$$
\xymatrix{
\ldots \ar[r] &
I_3^\bullet \ar[d] \ar[r] &
I_2^\bullet \ar[d] \ar[r] &
I_1^\bullet \ar[d] \\
\ldots \ar[r] &
(f_1^3, \ldots, f_r^3) \ar[r] &
(f_1^2, \ldots, f_r^2) \ar[r] &
(f_1, \ldots, f_r)
}
$$
을 $K(A)$에서 얻는다.

\begin{lemma}
\label{more-algebra-lemma-sequence-powers-pro-bounded}
상황 \ref{more-algebra-situation-koszul}에서 $A$가 Noether라고 가정하자. 위의
표기 아래 역계 $(I^n)$은 $D(A)$에서 역계 $(I_n^\bullet)$과
pro-동형이다.
\end{lemma}

\begin{proof}
역계 $I^n$이 $A$-가군의 범주에서 역계
$(f_1^n, \ldots, f_r^n)$과 pro-동형임을 보이는 것은 초등적이다.
구별 삼각형의 역계
$$
I_n^\bullet \to (f_1^n, \ldots, f_r^n) \to C_n^\bullet \to I_n^\bullet[1]
$$
을 생각하자. 여기서 $C_n^\bullet$은 첫 화살표의 원뿔이다. 유도
범주론, 보조정리 \ref{derived-lemma-pro-isomorphism}에 의해 역계
$C_n^\bullet$이 pro-영임을 보이면 충분하다. 복합체
$I_n^\bullet$은 $-r + 1 \leq i \leq 0$인 차수 $i$에서만 영이 아닌
항을 가지므로 $C_n^\bullet$도 마찬가지로 유계이다. 따라서 유도
범주론, 보조정리
\ref{derived-lemma-essentially-constant-cohomology}에 의해
$H^p(C_n^\bullet)$이 pro-영임을 보이면 충분하다. 위의 논의에 의해
$p \leq -1$이면 $H^p(C_n^\bullet) = H^p(K_n^\bullet)$이고
$p \geq 0$이면 $H^p(C_n^\bullet) = 0$이다. 역계
$H^p(K_n^\bullet)$이 pro-영이라는 사실은 보조정리
\ref{more-algebra-lemma-sequence-Koszul-complexes}의 증명에서 보였다(그리고 바로
여기서 $A$가 Noether라는 가정을 사용한다).
\end{proof}

\begin{lemma}
\label{more-algebra-lemma-tensoring-Deligne-system}
$A$를 Noether 환이라 하자. $I \subset A$를 아이디얼이라 하자.
$M^\bullet$을 유한 $A$-가군들로 이루어진 유계 복합체라 하자.
사상들의 역계
$$
I^n \otimes_A^\mathbf{L} M^\bullet \longrightarrow I^nM^\bullet
$$
은 $D(A)$의 pro-대상들의 동형사상을 정의한다.
\end{lemma}

\begin{proof}
$I$의 생성원 $f_1, \ldots, f_r \in I$를 택하자. 역계 $I^n$은
$A$-가군의 범주에서 역계 $(f_1^n, \ldots, f_r^n)$과 pro-동형이다.
보조정리 \ref{more-algebra-lemma-sequence-powers-pro-bounded}에서와 같은 표기
아래 사상들의 역계
$$
I_n^\bullet \otimes_A^\mathbf{L} M^\bullet
\longrightarrow
(f_1^n, \ldots, f_r^n)M^\bullet
$$
이 $D(A)$의 pro-대상들의 동형사상을 정의함을 증명하면 충분하다.
$i \not \in [a, b]$이면 $M^i = 0$이 되는 $a \leq b$가 있다고 하자.
그러면 위 화살표들의 원천과 목표는 차수 $[-r + a, b]$에서만 상동군을
갖는다. 따라서 임의의 $p \in \mathbf{Z}$에 대해 사상들의 역계
$$
H^p(I_n^\bullet \otimes_A^\mathbf{L} M^\bullet)
\longrightarrow
H^p((f_1^n, \ldots, f_r^n)M^\bullet)
$$
이 $A$-가군의 pro-대상들의 동형사상을 정의함을 보이면 충분하다.
유도 범주론, 보조정리
\ref{derived-lemma-pro-isomorphism-bis}를 보라.
$I_n^\bullet \otimes_A^\mathbf{L} M^\bullet$과
$I^n \otimes_A^\mathbf{L} M^\bullet$ 사이의 pro-동형사상 및
$(f_1^n, \ldots, f_r^n)M^\bullet$과 $I^nM^\bullet$ 사이의
pro-동형사상을 사용하면, 이는 사상들의 역계
% P02-E0502: affected pro-object conclusion lacks a period before “Choose” in the frozen source.
$$
H^p(I^n \otimes_A^\mathbf{L} M^\bullet)
\longrightarrow
H^p(I^nM^\bullet)
$$
이 $A$-가군의 pro-대상들의 동형사상을 정의함을 보이는 것과
동치이다. 유한 자유 $A$-가군들로 이루어진 위로 유계인 복합체
$P^\bullet$과 유사동형사상 $P^\bullet \to M^\bullet$을 택하자.
그러면 사상들의 역계
$$
H^p(I^nP^\bullet)
\longrightarrow
H^p(I^nM^\bullet)
$$
이 pro-동형사상임을 보이면 충분하다. 이는
$H^p(P^\bullet) = H^p(M^\bullet)$이므로 보조정리
\ref{more-algebra-lemma-consequence-Artin-Rees-bis}에서 따른다.
\end{proof}

\begin{lemma}
\label{more-algebra-lemma-factor-through-derived-tensor-product}
$A$를 Noether 환이라 하자. $I \subset A$를 아이디얼이라 하자.
$M$을 유한 $A$-가군이라 하자. $D(A)$에서 $I^nM \to M$이 사상
$I \otimes_A^\mathbf{L} M \to M$을 통하여 분해되는 정수 $n > 0$이
존재한다.
\end{lemma}

\begin{proof}
이는 보조정리 \ref{more-algebra-lemma-tensoring-Deligne-system}에서 따른다.
% P02-E0503: frozen prose says “can also been seen”; Korean supplies the intended grammar.
다음과 같이 직접 볼 수도 있다. 구별 삼각형
% P02-E0502: affected distinguished-triangle display lacks a period before the next sentence.
$$
I \otimes_A^\mathbf{L} M \to M \to A/I \otimes_A^\mathbf{L} M \to
I \otimes_A^\mathbf{L} M[1]
$$
을 생각하자. 삼각 범주의 공리들에 의해 어떤 $n$에 대해
$I^nM \to A/I \otimes_A^\mathbf{L} M$이 $D(A)$에서 영임을 증명하면
충분하다. $I$의 생성원 $f_1, \ldots, f_r$를 택하고
$K = K_\bullet(A, f_1, \ldots, f_r)$를 Koszul 복합체라 하며, 몫사상의
분해 $A \to K \to A/I$를 생각하자. 그러면 어떤 $n > 0$에 대해
$I^nM \to K \otimes_A M$이 $D(A)$에서 영임을 보이면 충분함을 알 수
있다. 어떤 $n > 0$을 찾아 $I^nM \to K \otimes_A M$이 $D(A)$에서
$\tau_{\geq t}(K \otimes_A M)$을 통하여 분해된다고 하자. 그러면
$\tau_{\geq t + 1}(K \otimes_A M)$을 통하여 분해하는 데 대한 장애는
$\Ext^t(I^nM, H_t(K \otimes_A M))$의 원소이다. 유한 $A$-가군
$H_t(K \otimes_A M)$은 $I$에 의해 소멸된다. 그러면 보조정리
\ref{more-algebra-lemma-ext-annihilated}에 의해 $n$을 늘려 이 장애 원소가 영이라고
가정할 수 있다. 이를 유한 번 반복하면 $I^nM \to K \otimes_A M$이
$D(A)$에서 $0 = \tau_{\geq r + 1}(K \otimes_A M)$을 통하여 분해되는
$n$을 얻고, 원하는 결론이 따른다.
\end{proof}

\section{잡록}
\label{more-algebra-section-misc}

\noindent
다른 어디에도 들어맞지 않는 몇 가지 결과를 모은다.

\begin{lemma}
\label{more-algebra-lemma-ext-annihilated-into}
$A$를 Noether 환이라 하자. $I \subset A$를 아이디얼이라 하자.
$K \in D(A)$가 유사연접이라 하고 $a \in \mathbf{Z}$라 하자.
모든 유한 $A$-가군 $M$에 대하여 $i \geq a$이면 가군
$\Ext^i_A(K, M)$이 $I$-멱 꼬임이라고 가정하자. 그러면 $i \geq a$이고
$M$이 유한일 때, 계 $\Ext^i_A(K, M/I^nM)$은 다음 값을 갖는
본질적으로 상수인 계이다.
% P02-E0504: frozen display lacks terminal punctuation before following prose.
$$
\Ext^i_A(K, M) = \lim \Ext^i_A(K, M/I^nM)
$$
\end{lemma}

\begin{proof}
$M$을 유한 $A$-가군이라 하자. $K$가 유사연접이므로
$\Ext^i_A(K, M)$은 유한 $A$-가군이다. 따라서 $i \geq a$이면 어떤
$t \geq 0$에 대하여 $I^t$에 의해 소멸된다. 보조정리
\ref{more-algebra-lemma-ext-induced-toplogy}에 의해 어떤 $n > 0$에 대하여
$\Ext^i_A(K, I^nM) \to \Ext^i_A(K, M)$의 상은 영이다.
짧은 완전열 $0 \to I^nM \to M \to M/I^n M \to 0$은 긴 완전열
$$
\Ext^i_A(K, I^nM) \to \Ext^i_A(K, M) \to
\Ext^i_A(K, M/I^nM) \to \Ext^{i + 1}_A(K, I^nM)
$$
을 준다. 방금 말한 사실을 유한 $A$-가군 $I^mM$에 적용하면 계
$\Ext^i_A(K, I^nM)$과 $\Ext^{i + 1}_A(K, I^nM)$은 값이 $0$인
본질적으로 상수인 계이다. 그림 추적을 하면
$\Ext^i_A(K, M/I^nM)$은 값이 $\Ext^i_A(K, M)$인 본질적으로 상수인
계임을 알 수 있다.
\end{proof}

\begin{lemma}
\label{more-algebra-lemma-tor-annihilated}
$A$를 Noether 환이라 하자. $I \subset A$를 아이디얼이라 하자. $M$을
유한 $A$-가군이라 하고 $N$을 $I$에 의해 소멸되는 $A$-가군이라 하자.
모든 $p \geq 0$에 대하여
$\text{Tor}^A_p(I^nM, N) \to \text{Tor}^A_p(M, N)$이 영이 되는 정수
$n > 0$이 존재한다.
\end{lemma}

\begin{proof}
보조정리 \ref{more-algebra-lemma-factor-through-derived-tensor-product}에 의해
$I^nM \to M$을 $I^nM \to M \otimes_A^\mathbf{L} I \to M$으로 분해할
수 있다. 우리는 합성
$$
I^nM \otimes_A^\mathbf{L} N \to
(M \otimes_A^\mathbf{L} I) \otimes_A^\mathbf{L} N
\to M \otimes_A^\mathbf{L} N
$$
이 영이라고 주장한다. 실제로 그림
$$
\xymatrix{
(M \otimes_A^\mathbf{L} I) \otimes_A^\mathbf{L} N \ar[rr] \ar[rd] & &
M \otimes_A^\mathbf{L} (I \otimes_A^\mathbf{L} N) \ar[ld] \\
& M \otimes_A^\mathbf{L} N
}
$$
은 가환하고(세부사항은 생략한다), $N$이 $I$에 의해 소멸되므로 사상
$I \otimes_A^\mathbf{L} N \to N$은 영이다.
\end{proof}

\begin{lemma}
\label{more-algebra-lemma-pseudo-coherent-tensor-limit}
$R$을 환이라 하자. $K \in D(R)$가 유사연접이라 하자.
$(M_n)$을 $R$-가군의 역계라 하자. 그러면
$R\lim K \otimes_R^\mathbf{L} M_n = K \otimes_R^\mathbf{L} R\lim M_n$이다.
\end{lemma}

\begin{proof}
정의에 쓰이는 구별 삼각형
$$
R\lim M_n \to \prod M_n \to \prod M_n \to R\lim M_n[1]
$$
을 생각하고 보조정리 \ref{more-algebra-lemma-pseudo-coherent-tensor}를 적용하자.
\end{proof}

\begin{lemma}
\label{more-algebra-lemma-additivity-of-pd}
$R$을 Noether 국소환이라 하자. $I \subset R$을 아이디얼이라 하고
$E$를 영이 아닌 유한 $R/I$-가군이라 하자. $R/I$의 사영 차원이
유한하고 $E$의 $R/I$ 위 사영 차원이 유한하면, $E$의 $R$ 위 사영
차원도 유한하고 다음이 성립한다.
% P02-E0504: frozen display lacks terminal punctuation before following prose.
$$
\text{pd}_R(E) = \text{pd}_R(R/I) + \text{pd}_{R/I}(E)
$$
\end{lemma}

\begin{proof}
유한 가군에 대하여 $R$, resp.\ $R/I$ 위 사영 차원이 유한한 것은
퍼펙트 가군인 것과 같다는 사실을 사용한다. 정의
\ref{more-algebra-definition-perfect} 뒤의 논의를 보라. 보조정리
\ref{more-algebra-lemma-cohomology-perfect}에 의해 $E$의 $R$ 위 사영 차원은
유한하다. 따라서 Auslander--Buchsbaum(Algebra, Proposition
\ref{algebra-proposition-Auslander-Buchsbaum})를 적용하면
$$
\text{pd}_R(E) + \text{depth}(E) = \text{depth}(R),\quad
\text{pd}_{R/I}(E) + \text{depth}(E) = \text{depth}(R/I),
$$
이고
% P02-E0504: frozen display lacks terminal punctuation before following prose.
$$
\text{pd}_R(R/I) + \text{depth}(R/I) = \text{depth}(R)
$$
이다. 첫 번째 등식에서는 $E$의 깊이를 $R$-가군으로서 취하고 두 번째
등식에서는 $R/I$-가군으로서 취한다는 데 주의하자. 그러나 두 깊이는
같다(이는 자명하지만 Algebra, Lemma
\ref{algebra-lemma-depth-goes-down-finite}에서도 따른다). 이로써 증명이
끝난다.
\end{proof}

\begin{lemma}
\label{more-algebra-lemma-enlarge}
$A \to B$를 환 사상이라 하자. 다음 성질을 갖는 기수
$\kappa = \kappa(A \to B)$가 존재한다. $M^\bullet$, resp.\ $N^\bullet$을
$A$-가군, resp.\ $B$-가군의 복합체라 하자.
$a : M^\bullet \to N^\bullet$을 $A$-가군 복합체의 사상이라 하고,
$D(B)$에서 동형사상
$M^\bullet \otimes_A^\mathbf{L} B \to N^\bullet$을 유도한다고 하자.
$M_1^\bullet \subset M^\bullet$, resp.\ $N_1^\bullet \subset N^\bullet$을
$A$-가군, resp.\ $B$-가군의 부분복합체라 하고
$a(M_1^\bullet) \subset N_1^\bullet$이라 하자. 그러면 부분복합체
$$
M_1^\bullet \subset M_2^\bullet \subset M^\bullet
\quad\text{and}\quad
N_1^\bullet \subset N_2^\bullet \subset N^\bullet
$$
로서 $a(M_2^\bullet) \subset N_2^\bullet$이고 다음 성질들을 갖는 것들이
존재한다.
\begin{enumerate}
\item $\Ker(H^i(M_1^\bullet \otimes_A^\mathbf{L} B) \to H^i(N_1^\bullet))$
은 $H^i(M_2^\bullet \otimes_A^\mathbf{L} B)$에서 영으로 사상된다.
\item $\Im(H^i(N_1^\bullet) \to H^i(N_2^\bullet))$은
% P02-E0505: frozen target of the degree-i image map uses H^2.
$\Im(H^i(M_2^\bullet \otimes_A^\mathbf{L} B) \to H^2(N_2^\bullet))$에
포함된다.
\item $|\bigcup M_2^i \cup \bigcup N_2^i| \leq
\max(\kappa, |\bigcup M_1^i \cup \bigcup N_1^i|)$이다.
\end{enumerate}
\end{lemma}

\begin{proof}
$\kappa = \max(|A|, |B|, \aleph_0)$라 하자.
$|M^\bullet| = |\bigcup M^i|$라 놓고 다른 복합체에도 같은 방식으로
표기한다. 이 표기법으로 보조정리의 명제에 쓰인 양에 대하여
% P02-E0506: frozen cardinal identity uses the not-yet-constructed M_2 in place of the N_1-side quantity.
$$
\max(\kappa, |\bigcup M_1^i \cup \bigcup N_1^i|) =
\max(\kappa, |M_1^\bullet|, |M_2^\bullet|)
$$
가 성립한다. 이하에서는 이 사실과 기수 산술에서 나오는 다른 관찰들을
별도의 언급 없이 사용한다.

\medskip\noindent
먼저 ``작은'' 부분복합체가 충분히 많음을 보이자. 유한 부분집합들의 족
$E = \{E^i\}$와 $F = \{F^i\}$의 모든 쌍에 대하여, 여기서
$E^i \subset M^i$, $i \in \mathbf{Z}$이고
$F^i \subset N^i$, $i \in \mathbf{Z}$이며, 다음을 택할 수 있다.
$$
M_1^\bullet \subset M_1(E, F)^\bullet \subset M^\bullet
\quad\text{and}\quad
N_1^\bullet \subset N_1(E, F)^\bullet \subset N^\bullet
$$
이들은 각각 $A$-가군과 $B$-가군의 가장 작은 부분복합체로서
$a(M_1(E, F)^\bullet) \subset N_1(E, F)^\bullet$이고
$E^i \subset M_1(E, F)^i$이며
% P02-E0507: frozen F-side containment uses undefined M_2(E,F).
$F^i \subset M_2(E, F)^i$이다. 그러면 다음을 쉽게 알 수 있다.
% P02-E0508: frozen bounds use undefined M_2(E,F) and M_2 instead of the N_1-side quantities.
% P02-E0504: frozen display lacks terminal punctuation before following prose.
$$
|M_1(E, F)^\bullet| \leq \max(\kappa, |M_1^\bullet|)
\quad\text{and}\quad
|M_2(E, F)^\bullet| \leq \max(\kappa, |M_2^\bullet|)
$$
세부사항은 생략한다. 다음은 명백하다.
$$
M^\bullet = \colim_{(E, F)} M_1(E, F)^\bullet
\quad\text{and}\quad
N^\bullet = \colim_{(E, F)} N_1(E, F)^\bullet
$$
그리고 이 여극한들은 (항별로) 여과 여극한이다.

\medskip\noindent
자유 $A$-가군 $F_i$들에 의한 분해가 존재하여
% P02-E0509: frozen resolution uses superscripted F but names the free modules with subscripts.
$\ldots \to F^{-1} \to F^0 \to B$이고 $|F_i| \leq \kappa$이다
(세부사항은 생략한다). $M_1^\bullet \otimes_A^\mathbf{L} B$의 상동
가군은 $\text{Tot}(M_1^\bullet \otimes_A F^\bullet)$로 계산된다.
따라서
$|H^i(M_1^\bullet \otimes_A^\mathbf{L} B)| \leq \max(\kappa, |M_1^\bullet|)$이다.

\medskip\noindent
$i \in \mathbf{Z}$라 하고
$\xi \in H^i(M_1^\bullet \otimes_A^\mathbf{L} B)$를
$H^i(N_1^\bullet)$에서 영으로 사상되는 원소라 하자. 그러면 $\xi$는
$H^i(N^\bullet)$에서 영으로 사상되며, 따라서 $\xi$는
$H^i(M^\bullet \otimes_A^\mathbf{L} B)$에서도 영으로 사상된다.
유도 텐서곱은 여과 여극한과 가환하므로, 위와 같은 유한 족
$E_\xi$와 $F_\xi$를 찾아 $\xi$가
$H^i(M_1(E_\xi, F_\xi)^\bullet \otimes_A^\mathbf{L} B)$에서 영으로
사상되게 할 수 있다.

\medskip\noindent
$i \in \mathbf{Z}$라 하고 $\eta \in H^i(N_1^\bullet)$라 하자. 그러면
$H^i(N^\bullet)$에서 $\eta$의 상은 사상
$H^i(M^\bullet \otimes_A^\mathbf{L} B) \to H^i(N^\bullet)$의 상에 속한다.
따라서 앞에서와 같이 위와 같은 유한 족 $E_\eta$와 $F_\eta$를 찾아
$\eta$가 다음 상동군의 원소로 사상되고
% P02-E0510: frozen cohomology terms omit the closing parenthesis and complex superscript.
$H^i(N_1(E_\eta, F_\eta)$ 그 원소가 사상
$H^i(M_1(E_\eta, F_\eta)^\bullet \otimes_A^\mathbf{L} B) \to
H^i(N_1(E_\eta, F_\eta)$의 상에 속하게 할 수 있다.

\medskip\noindent
이제 단순히
$$
M_2^\bullet =
\sum\nolimits_\xi M_1(E_\xi, F_\xi)^\bullet +
\sum\nolimits_\eta M_1(E_\eta, F_\eta)^\bullet
$$
로 정의하자. 여기서 합은 앞의 두 문단에 나온 $\xi$와 $\eta$에 걸쳐
취하며 $M^\bullet$ 안에서 취한다. 마찬가지로
$$
N_2^\bullet =
\sum\nolimits_\xi N_1(E_\xi, F_\xi)^\bullet +
\sum\nolimits_\eta N_1(E_\eta, F_\eta)^\bullet
$$
로 놓고, 합은 $N^\bullet$ 안에서 취한다. 구성에 의해 이 선택들은
성질 (1)과 (2)를 만족한다. 집합을 이루는 $\xi$와 $\eta$의 기수가
기껏해야 $\max(\kappa, |M_1^\bullet|, |N_1^\bullet|)$이므로 상계 (3)도
따른다.
\end{proof}

\begin{lemma}
\label{more-algebra-lemma-f-bounded}
$R$을 환이라 하고 $f \in R$이라 하자. $M \in D(R)$이라 하고
$C$를 $f : M \to M$의 원뿔이라 하자. $i < 0$이면
$H^i(M)_f = 0$이고 $i < -1$이면 $H^i(C) = 0$이라고 하자. 그러면
$i < 0$이면 $H^i(M) = 0$이다.
\end{lemma}

\begin{proof}
$M_f = M \otimes_R^\mathbf{L} R_f$라 쓰고 구별 삼각형
$F \to M \to M_f$를 택하자. 가정에 의해 $i < 0$이면
$H^i(M_f) = H^i(M)_f = 0$이다. 따라서 $i < 0$이면 $H^i(F) = 0$임을
보이면 충분하다. 모든 $i \in \mathbf{Z}$에 대하여 $H^i(F)$가
$f$-멱 꼬임임에 주의하자. 다른 한편으로
$f : M_f \to M_f$가 동형사상이므로 $C$는 $f : F \to F$의 원뿔과
동형이다(Derived Categories, Proposition
\ref{derived-proposition-9}을 사용하라). 이제 $H^i(F) \not = 0$이면
$f : H^i(F) \to H^i(F)$의 핵은 영이 아니며, 이는
$H^{i - 1}(C)$가 영이 아님을 함의한다. 우리의 가정에 의하면
$i - 1 < -1$일 때에는 이런 일이 일어날 수 없으므로 증명이 끝난다.
\end{proof}

\begin{lemma}
\label{more-algebra-lemma-f-flat}
$R$을 환이라 하고 $f \in R$이라 하자. $M \in D(R)$이라 하자. 다음을
가정하자.
\begin{enumerate}
\item $i > 0$이면 $H^i(M) = 0$이다.
\item $M \otimes_R^\mathbf{L} R_f$는 차수 $0$에 놓인 평탄
$R_f$-가군과 동형이다.
\item $M \otimes_R^\mathbf{L} R/fR$는 차수 $0$에 놓인 평탄
$R/fR$-가군과 동형이다.
\end{enumerate}
그러면 $M$은 차수 $0$에 놓인 평탄 $R$-가군과 동형이다.
\end{lemma}

\begin{proof}
$K$를 $R$-가군이라 하자. 절 \ref{more-algebra-section-tor}를 보라. 다음
$N = M \otimes_R^\mathbf{L} K$가 차수 $0$에 놓인 $R$-가군과 동형임을
보이면 충분하다. 가정 (1)에 의해 $N$은 차수 $\leq 0$에서만 영이 아닌
상동군을 갖는다. 따라서 보조정리 \ref{more-algebra-lemma-f-bounded}에 의해
$i < 0$이면 $H^i(N)_f = 0$임과, $C$를 $f : N \to N$의 원뿔이라 할 때
$i < -1$이면 $H^i(C) = 0$임을 보이면 충분하다. 전자에 대해서는 다음에
주의하자.
% P02-E0504: frozen display lacks terminal punctuation before following prose.
$$
H^i(N)_f = H^i(N \otimes_R^\mathbf{L} R_f) =
H^i(M \otimes_R^\mathbf{L} K \otimes_R^\mathbf{L} R_f ) =
H^i((M \otimes_R^\mathbf{L} R_f) \otimes_{R_f}^\mathbf{L}
(K \otimes_R^\mathbf{L} R_f))
$$
우리의 가정 (2)에 의해 이는 $i = 0$인 경우를 제외하면 영이고, 그
경우에는 $H^0(M)_f \otimes_R K$를 얻는다. 후자에 대해서는 $C'$를
$f : K \to K$의 원뿔이라 하자. 그러면
% P02-E0504: frozen display lacks terminal punctuation before following prose.
$$
C = M \otimes_R^\mathbf{L} C'
$$
복합체 $C'$은 차수 $0$과 $-1$에서만 영이 아닌 상동군을 가지며, 이들은
각각 $K/fK$와 $K[f]$이다. 다시 말해 구별 삼각형
% P02-E0504: frozen display lacks terminal punctuation before following prose.
$$
(K[f])[1] \to C' \to K/fK
$$
이 있다. 다른 한편으로 $f$에 의해 소멸되는 임의의 $R$-가군 $K'$에
대하여 다음이 성립한다.
% P02-E0504: frozen display lacks terminal punctuation before following prose.
$$
M \otimes_R^\mathbf{L} K' =
M \otimes_R^\mathbf{L} R/fR \otimes_{R/fR}^\mathbf{L} K'
$$
우리의 가정 (3)에 의해 이는 차수 $0$에 놓인 가군
$H^0(M \otimes_R^\mathbf{L} R/fR) \otimes_{R/fR} K'$와 같다. 위의
내용을 종합하면 $C$는 차수 $0$과 $-1$에서만 영이 아닌 상동군을 가지며,
증명이 완결된다.
\end{proof}

\section{이중복합체에 관한 기법}
\label{more-algebra-section-tricks}

\noindent
이 절에서는 Homology, Section
\ref{homology-section-double-complexes-abelian-groups}의 논의를 이어간다.

\begin{lemma}
\label{more-algebra-lemma-vanishing-coh-prod-totalization}
$A_0^\bullet \to A_1^\bullet \to A_2^\bullet \to \ldots$을 아벨 군의
복합체들로 이루어진 복합체라 하자. 모든 $p \geq 0$에 대하여
$H^{-p}(A_p^\bullet) = 0$이라고 가정하자. $A^{p, q} = A_p^q$라 놓고
$A^{\bullet, \bullet}$을 이중복합체로 보자. 그러면
$H^0(\text{Tot}_\pi(A^{\bullet, \bullet})) = 0$이다.
\end{lemma}

\begin{proof}
% P02-E1178: frozen English omits the preposition in “Denote ... by”.
$f_p : A_p^\bullet \to A_{p + 1}^\bullet$로 주어진 복합체 사상들을
나타내자. $\text{Tot}_\pi(A^{\bullet, \bullet})$의 미분은 차수 $n$의
원소에서 다음으로 주어짐을 상기하자.
$$
\prod\nolimits_{p + q = n} (f^q_p + (-1)^p\text{d}^q_{A_p^\bullet})
$$
$\xi \in H^0(\text{Tot}_\pi(A^{\bullet, \bullet}))$를 상동류라 하자.
$\xi$가 영임을 보이겠다. $\xi$를
% P02-E0511: frozen English says “an cocycle”; Korean uses the grammatical mathematical term.
공순환 $x = (x_p) \in \prod A^{p, -p}$의 동치류로 나타내자.
$\text{d}(x) = 0$이므로
$\text{d}_{A_0^\bullet}(x_0) = 0$이다.
$H^0(A_0^\bullet) = 0$이므로
$\text{d}_{A_0^\bullet}(y_{-1}) = x_0$을 만족하는
$y_{-1} \in A^{0, -1}$가 존재한다. 그러면
$\text{d}_{A_1^\bullet}(x_1 + f_0(y_{-1})) = 0$이다.
$H^{-1}(A_1^\bullet) = 0$이므로
$-\text{d}_{A_1^\bullet}(y_{-2}) = x_1 + f_0(y_{-1})$을 만족하는
$y_{-2} \in A^{1, -2}$를 찾을 수 있다. 귀납법으로
$y_{-p - 1} \in A^{p, -p - 1}$를 찾아 다음을 만족하게 할 수 있다.
% P02-E0513: frozen display ending a complete sentence lacks terminal punctuation.
$$
(-1)^p\text{d}_{A_p^\bullet}(y_{-p - 1}) = x_p + f_{p - 1}(y_{-p})
$$
이는 $y = (y_{-p - 1})$일 때 $\text{d}(y) = x$임을 뜻한다.
\end{proof}

\begin{lemma}
\label{more-algebra-lemma-prod-qis-gives-qis}
다음이 아벨 군의 복합체들로 이루어진 두 복합체 사이의 사상이라고 하자.
$$
(A_0^\bullet \to A_1^\bullet \to A_2^\bullet \to \ldots)
\longrightarrow
(B_0^\bullet \to B_1^\bullet \to B_2^\bullet \to \ldots)
$$
$A^{p, q} = A_p^q$, $B^{p, q} = B_p^q$라 놓아 이중복합체들을 얻자.
$\text{Tot}_\pi(A^{\bullet, \bullet})$와
$\text{Tot}_\pi(B^{\bullet, \bullet})$를 이 이중복합체들에 딸린 곱
총복합체라 하자. 각 $A_p^\bullet \to B_p^\bullet$이 준동형사상이면
$\text{Tot}_\pi(A^{\bullet, \bullet}) \to \text{Tot}_\pi(B^{\bullet, \bullet})$
도 준동형사상이다.
\end{lemma}

\begin{proof}
$\text{Tot}_\pi(A^{\bullet, \bullet})$의 차수 $n$인 항은 다음으로
주어짐을 상기하자.
% P02-E0512: frozen second product index uses p+1=n instead of p+q=n.
$\prod_{p + q = n} A^{p, q} = \prod_{p + 1 = n} A^q_p$.
$C_p^\bullet$을 사상 $A_p^\bullet \to B_p^\bullet$의 원뿔이라 하자.
Derived Categories, Section \ref{derived-section-cones}를 보라.
원뿔 구성의 함자성에 의해 복합체들로 이루어진 복합체
% P02-E0513: frozen display ending a complete sentence lacks terminal punctuation.
$$
C_0^\bullet \to C_1^\bullet \to C_2^\bullet \to \ldots
$$
를 얻는다. 그러면 $\text{Tot}_\pi(C^{\bullet, \bullet})$의 차수 $n$인
항은 다음으로 주어진다.
% P02-E0513: frozen display ending a complete sentence lacks terminal punctuation.
$$
\prod_{p + q = n} C^{p, q} = \prod_{p + q = n} C^q_p =
\prod_{p + q = n} (B^q_p \oplus A^{q + 1}_p) =
\prod_{p + q = n} B^q_p \oplus \prod_{p + q = n} A^{q + 1}_p
$$
따라서 $\text{Tot}_\pi(C^{\bullet, \bullet})$은 사상
$\text{Tot}_\pi(A^{\bullet, \bullet}) \to \text{Tot}_\pi(B^{\bullet, \bullet})$
의 원뿔이다.
% P02-E0514: frozen mid-sentence parenthetical begins with an inappropriate capital letter.
(미분들이 일치한다는 검증은 생략한다.) 따라서 각 $A_p^\bullet$이
비순환이면 $\text{Tot}_\pi(A^{\bullet, \bullet})$이 비순환임을
보이면 충분하다. 이는 보조정리
\ref{more-algebra-lemma-vanishing-coh-prod-totalization}에서 따른다.
\end{proof}

\section{약 에탈 환 사상}
\label{more-algebra-section-weakly-etale}

\noindent
이 절의 결과 대부분은 Olivier의 논문 \cite{Olivier-AF}에서 왔다.
관련 논문 \cite{Ferrand-epi}도 보라.

\begin{definition}
\label{more-algebra-definition-weakly-etale}
환 $A$의 모든 $A$-가군이 $A$ 위 평탄이면 이를
{\it 절대 평탄}이라고 한다. $A \to B$와
$B \otimes_A B \to B$가 모두 평탄이면 환 사상 $A \to B$를
{\it 약 에탈} 또는 {\it 절대 평탄}이라고 한다.
\end{definition}

\noindent
절대 평탄환은 (흔히 비가환환을 다루는 맥락에서) von Neumann
정칙환이라고도 한다. 국소화는 약 에탈 환 사상이다. 에탈 환 사상은
약 에탈이다. 다음은 간단하지만 핵심적인 성질이다.

\begin{lemma}
\label{more-algebra-lemma-key}
$A \to B$를 $B \otimes_A B \to B$가 평탄인 환 사상이라 하자.
$N$을 $B$-가군이라 하자. $N$이 $A$-가군으로서 평탄이면
$N$은 $B$-가군으로서 평탄이다.
\end{lemma}

\begin{proof}
% P02-E0515: frozen English contains a spurious article before “flat”.
$N$이 $A$-가군으로서 평탄이라고 가정하자. 그러면 함자
$$
\text{Mod}_B \longrightarrow \text{Mod}_{B \otimes_A B},\quad
N' \mapsto N \otimes_A N'
$$
는 완전하다. $B \otimes_A B \to B$가 평탄이므로 함자
$$
\text{Mod}_B \longrightarrow \text{Mod}_B,\quad
N' \mapsto (N \otimes_A N') \otimes_{B \otimes_A B} B = N \otimes_B N'
$$
가 완전임을 알 수 있고, 따라서 $N$은 $B$ 위 평탄이다.
\end{proof}

\begin{definition}
\label{more-algebra-definition-weak-dimension}
$A$를 환이라 하자. $d \geq 0$을 정수라 하자. 모든 $A$-가군의
Tor 차원이 $\leq d$이면 $A$의 {\it 약한 차원이 $\leq d$}라고 한다.
\end{definition}

\begin{lemma}
\label{more-algebra-lemma-weak-dimension-goes-up}
$A \to B$를 약 에탈 환 사상이라 하자. $A$의 약한 차원이 기껏해야
$d$이면 $B$도 그러하다.
\end{lemma}

\begin{proof}
$N$을 $B$-가군이라 하자. $d = 0$이면 $N$은 $A$-가군으로서 평탄이고,
따라서 보조정리 \ref{more-algebra-lemma-key}에 의해 $B$-가군으로서 평탄이다.
$d > 0$이라고 가정하자. 자유 $B$-가군들에 의한 분해
$F_\bullet \to N$을 택하자. 우리의 가정과 보조정리
\ref{more-algebra-lemma-last-one-flat}에 의해
$K = \Im(F_d \to F_{d - 1})$는 $A$-평탄이다. 따라서 보조정리
\ref{more-algebra-lemma-key}에 의해 $B$-평탄이다. 그러므로
$0 \to K \to F_{d - 1} \to \ldots \to F_0 \to N \to 0$은 길이가
$d$인 평탄 분해이고, $N$의 Tor 차원이 기껏해야 $d$임을 알 수 있다.
\end{proof}

\begin{lemma}
\label{more-algebra-lemma-absolutely-flat}
$A$를 환이라 하자. 다음은 동치이다.
\begin{enumerate}
\item $A$의 약한 차원이 $\leq 0$이다.
\item $A$는 절대 평탄이다.
\item $A$는 기약이고 모든 소아이디얼은 극대아이디얼이다.
\item $A$의 모든 국소환은 체이다.
\end{enumerate}
\end{lemma}

\begin{proof}
(1)과 (2)의 동치는 즉시 알 수 있다.

\medskip\noindent
(1) $\Rightarrow$ (3)의 증명.
$A$가 절대 평탄이라고 가정하자. 그러면 $A$의 모든 아이디얼은
순수하다. Algebra, Definition \ref{algebra-definition-pure-ideal}을
보라. 따라서 Algebra, Lemma
\ref{algebra-lemma-finitely-generated-pure-ideal}에 의해 모든 유한 생성
아이디얼은 멱등원 하나로 생성된다. $f \in A$이면 어떤 멱등원
$e \in A$에 대하여 $(f) = (e)$이고, $D(f) = D(e)$는 열린닫힌집합이다
(Algebra, Lemma \ref{algebra-lemma-idempotent-spec}). 이로부터 Algebra,
Lemma \ref{algebra-lemma-ring-with-only-minimal-primes}에 의해 $A$의 모든
소아이디얼은 극대아이디얼임을 알 수 있다. 더 나아가 $f$가 멱영이면
$e = 0$이고 따라서 $f = 0$이다. 그러므로 $A$는 기약이다.

\medskip\noindent
(3) $\Rightarrow$ (4)의 증명. $A$가 기약이고 모든 소아이디얼이
극대아이디얼이면 Algebra, Lemma
\ref{algebra-lemma-minimal-prime-reduced-ring}에 의해 모든 국소환은
체이다.

\medskip\noindent
(4) $\Rightarrow$ (2)의 증명. $A$의 모든 국소환이 체이면 Algebra,
Lemma \ref{algebra-lemma-flat-localization}에 의해 모든 $A$-가군은
평탄이다.
\end{proof}

\begin{lemma}
\label{more-algebra-lemma-product-fields-absolutely-flat}
체들의 곱은 절대 평탄환이다.
\end{lemma}

\begin{proof}
$K_i$를 체들의 족이라 하자. $f = (f_i) \in \prod K_i$이면 $f$가
생성하는 아이디얼은, $f_i$가 $0$인지 아닌지에 따라 $e_i = 0, 1$인
멱등원 $e = (e_i)$가 생성하는 아이디얼과 같다. 따라서
$D(f) = D(e)$는 열린닫힌집합이고, 보조정리
\ref{more-algebra-lemma-absolutely-flat}과 Algebra, Lemma
\ref{algebra-lemma-ring-with-only-minimal-primes}에 의해 결론을 얻는다.
\end{proof}

\begin{lemma}
\label{more-algebra-lemma-base-change-weakly-etale}
$A \to B$와 $A \to A'$을 환 사상이라 하자.
$B' = B \otimes_A A'$을 $B$의 밑변환이라 하자.
\begin{enumerate}
\item $B \otimes_A B \to B$가 평탄이면
$B' \otimes_{A'} B' \to B'$도 평탄이다.
\item $A \to B$가 약 에탈이면 $A' \to B'$도 약 에탈이다.
\end{enumerate}
\end{lemma}

\begin{proof}
$B \otimes_A B \to B$가 평탄이라고 가정하자. 환 사상
$B' \otimes_{A'} B' \to B'$은 $A \to A'$에 의한
$B \otimes_A B \to B$의 밑변환이다. 따라서 Algebra, Lemma
\ref{algebra-lemma-flat-base-change}에 의해 평탄이다. 이로써 (1)이
증명된다. (2)는 (1)과 방금 사용한, 평탄 환 사상의 밑변환은 평탄이라는
사실에서 따른다.
\end{proof}

\begin{lemma}
\label{more-algebra-lemma-absolutely-flat-over-absolutely-flat}
$A \to B$를 $B \otimes_A B \to B$가 평탄인 환 사상이라 하자.
\begin{enumerate}
\item $A$가 절대 평탄환이면 $B$도 그러하다.
\item $A$가 기약이고 $A \to B$가 약 에탈이면 $B$는 기약이다.
\end{enumerate}
\end{lemma}

\begin{proof}
(1)은 보조정리 \ref{more-algebra-lemma-key}와 정의들에서 즉시 따른다.
$A$가 기약이면 단사 사상
$A \to A' = \prod_{\mathfrak p \subset A\text{ minimal}} A_\mathfrak p$
가 존재하여 $A$를 절대 평탄환에 넣는다(Algebra, Lemma
\ref{algebra-lemma-reduced-ring-sub-product-fields}와 보조정리
\ref{more-algebra-lemma-product-fields-absolutely-flat}). $A \to B$가 평탄이면
유도된 사상 $B \to B' = B \otimes_A A'$도 단사이다. 보조정리
\ref{more-algebra-lemma-base-change-weakly-etale}에 의해 환 사상
$A' \to B'$은 약 에탈이다. (1)에 의해 $B'$은 절대 평탄이다.
보조정리 \ref{more-algebra-lemma-absolutely-flat}에 의해 환 $B'$은 기약이다.
따라서 $B$는 기약이다.
\end{proof}

\begin{lemma}
\label{more-algebra-lemma-composition-weakly-etale}
$A \to B$와 $B \to C$를 환 사상이라 하자.
\begin{enumerate}
\item $B \otimes_A B \to B$와 $C \otimes_B C \to C$가 평탄이면
$C \otimes_A C \to C$는 평탄이다.
\item $A \to B$와 $B \to C$가 약 에탈이면 $A \to C$도 약 에탈이다.
\end{enumerate}
\end{lemma}

\begin{proof}
(1)은 곱셈 사상의 분해
$$
C \otimes_A C \longrightarrow C \otimes_B C \longrightarrow C
$$
와 등식
$$
C \otimes_B C = (C \otimes_A C) \otimes_{B \otimes_A B} B,
$$
그리고 평탄 사상의 밑변환은 평탄이라는 사실 및 평탄 환 사상들의 합성은
평탄이라는 사실에서 따른다. Algebra, Lemmas
\ref{algebra-lemma-flat-base-change}와
\ref{algebra-lemma-composition-flat}을 보라. (2)는 (1)과 방금 사용한,
평탄 환 사상들의 합성은 평탄이라는 사실에서 따른다.
\end{proof}

\begin{lemma}
\label{more-algebra-lemma-go-down}
$A \to B \to C$를 환 사상이라 하자.
\begin{enumerate}
\item $B \to C$가 충실 평탄이고 $C \otimes_A C \to C$가 평탄이면
$B \otimes_A B \to B$는 평탄이다.
\item $B \to C$가 충실 평탄이고 $A \to C$가 약 에탈이면
$A \to B$는 약 에탈이다.
\end{enumerate}
\end{lemma}

\begin{proof}
$B \to C$가 충실 평탄이고 $C \otimes_A C \to C$가 평탄이라고
가정하자. 가환 그림을 생각하자.
% P02-E0522: frozen commutative-diagram display ends a sentence without punctuation.
$$
\xymatrix{
C \otimes_A C \ar[r] & C \\
B \otimes_A B \ar[r] \ar[u] & B \ar[u]
}
$$
수직 화살표들은 평탄이고 위쪽 수평 화살표도 평탄이다. 따라서 $C$는
$B \otimes_A B$-가군으로서 평탄이다. 사상 $B \to C$는 충실 평탄이고
$C = B \otimes_B C$이다. 따라서 Algebra, Lemma
\ref{algebra-lemma-flatness-descends-more-general}에 의해 $B$는
$B \otimes_A B$-가군으로서 평탄이다. 이로써 (1)이 증명된다. (2)는
(1)과, $A \to C$가 평탄이고 $B \to C$가 충실 평탄이면
$A \to B$가 평탄이라는 사실(Algebra, Lemma
\ref{algebra-lemma-flatness-descends-more-general})에서 따른다.
\end{proof}

\begin{lemma}
\label{more-algebra-lemma-weakly-etale-permanence}
$A$를 환이라 하자. $B \to C$를 약 에탈 $A$-대수들 사이의
$A$-대수 사상이라 하자. 그러면 $B \to C$는 약 에탈이다.
\end{lemma}

\begin{proof}
환 사상 $B \to C$는 보조정리 \ref{more-algebra-lemma-key}에 의해 평탄이다.
환 사상 $C \otimes_A C \to C \otimes_B C$는 전사이고, 따라서
환의 전사이다.
% P02-E0516: frozen English places a spurious comma between “implies” and its clause.
그러므로 보조정리 \ref{more-algebra-lemma-key}에 의해 $C$가 $C \otimes_A C$ 위
평탄인 것으로부터 $C$가 $C \otimes_B C$ 위 평탄임을 알 수 있다.
\end{proof}

\begin{lemma}
\label{more-algebra-lemma-formally-unramified}
$A \to B$를 $B \otimes_A B \to B$가 평탄인 환 사상이라 하자.
그러면 $\Omega_{B/A} = 0$이다. 즉, $B$는 $A$ 위 형식적으로 비분기이다.
\end{lemma}

\begin{proof}
$I \subset B \otimes_A B$를 평탄 전사 사상
$B \otimes_A B \to B$의 핵이라 하자. 그러면 $I$는 순수 아이디얼이고
(Algebra, Definition \ref{algebra-definition-pure-ideal}), 따라서
$I^2 = I$이다(Algebra, Lemma \ref{algebra-lemma-pure}).
$\Omega_{B/A} = I/I^2$이므로
(Algebra, Lemma \ref{algebra-lemma-differentials-diagonal}) 소멸을 얻는다.
이는 Algebra, Lemma
\ref{algebra-lemma-characterize-formally-unramified}에 의해 $B$가
$A$ 위 형식적으로 비분기임을 뜻한다.
\end{proof}

\begin{lemma}
\label{more-algebra-lemma-weakly-etale-finite-type}
$A \to B$를 $B \otimes_A B \to B$가 평탄인 환 사상이라 하자.
\begin{enumerate}
\item $A \to B$가 유한형이면 $A \to B$는 비분기이다.
\item $A \to B$가 유한 표시이고 평탄이면 $A \to B$는 에탈이다.
\end{enumerate}
특히 유한 표시인 약 에탈 환 사상은 에탈이다.
\end{lemma}

\begin{proof}
(1)은 보조정리 \ref{more-algebra-lemma-formally-unramified}와 Algebra, Definition
\ref{algebra-definition-unramified}에서 따른다. (2)는 (1)과 Algebra,
Lemma \ref{algebra-lemma-etale-flat-unramified-finite-presentation}에서
따른다.
\end{proof}

\begin{lemma}
\label{more-algebra-lemma-when-weakly-etale}
$A \to B$를 환 사상이라 하자. 다음 각 경우에 $A \to B$는 약 에탈이다.
\begin{enumerate}
\item $B = S^{-1}A$는 $A$의 국소화이다.
\item $A \to B$는 에탈이다.
\item $B$는 약 에탈 $A$-대수들의 여과 여극한이다.
\end{enumerate}
\end{lemma}

\begin{proof}
에탈 환 사상은 평탄이고, 사상 $B \otimes_A B \to B$도 에탈
$A$-대수들 사이의 사상으로서 에탈이다(Algebra, Lemma
\ref{algebra-lemma-map-between-etale}). 이로써 (2)가 증명된다.

\medskip\noindent
$B_i$를 약 에탈 $A$-대수들의 유향계라 하자. 그러면 Algebra, Lemma
\ref{algebra-lemma-colimit-flat}에 의해 $B = \colim B_i$는 $A$ 위
평탄이다. 보조정리 \ref{more-algebra-lemma-weakly-etale-permanence}에 의해 전이사상
$B_i \to B_{i'}$가 평탄임을 주목하자. 따라서 각 $i$에 대하여 $B$는
$B_i$ 위 평탄이고, Algebra, Lemma \ref{algebra-lemma-composition-flat}에
의해 $B$는 $B_i \otimes_A B_i$ 위 평탄이다. 따라서 Algebra, Lemma
\ref{algebra-lemma-colimit-rings-flat}에 의해 $B$는
$B \otimes_A B = \colim B_i \otimes_A B_i$ 위 평탄이다.

\medskip\noindent
(1)은 직접 증명할 수도 있지만 (2)와 (3)을 결합해도 따른다.
\end{proof}

\begin{lemma}
\label{more-algebra-lemma-absolutely-flat-fields}
$L/K$를 체 확대라 하자. $L \otimes_K L \to L$이 평탄이면 $L$은
$K$의 대수적 분리가능 확대이다.
\end{lemma}

\begin{proof}
% P02-E0517: frozen English omits the preposition governing the quantified subfield.
보조정리 \ref{more-algebra-lemma-go-down}에 의해 모든 부분체
$K \subset L' \subset L$에 대하여 사상 $L' \otimes_K L' \to L'$이
평탄임을 알 수 있다. 따라서 $L$이 $K$의 유한 생성 체 확대라고 가정해도
된다. 이 경우 $L/K$가 형식적으로 비분기라는 사실(보조정리
\ref{more-algebra-lemma-formally-unramified})로부터 $L/K$가 유한 분리가능임을 알 수
있다. Algebra, Lemma
\ref{algebra-lemma-characterize-separable-algebraic-field-extensions}를 보라.
\end{proof}

\begin{lemma}
\label{more-algebra-lemma-absolutely-flat-over-field}
$B$를 체 $K$ 위의 대수라 하자. 다음은 동치이다.
\begin{enumerate}
\item $B \otimes_K B \to B$가 평탄이다.
\item $K \to B$는 약 에탈이다.
\item $B$는 에탈 $K$-대수들의 여과 여극한이다.
\end{enumerate}
더 나아가 $B$의 모든 유한 생성 $K$-부분대수는 $K$ 위 에탈이다.
\end{lemma}

\begin{proof}
모든 $K$-대수는 $K$ 위 평탄이므로 (1)과 (2)는 동치이다.
% P02-E0518: frozen English omits terminal punctuation before the paragraph break.
(3)은 보조정리 \ref{more-algebra-lemma-when-weakly-etale}에 의해 (1)과 (2)를
함의한다.

\medskip\noindent
(1)과 (2)가 성립한다고 가정하자. (3)과 보조정리의 유한성 명제를
증명하겠다.
% P02-E0519: frozen English contains two malformed article phrases for absolutely flat rings.
체는 절대 평탄환이므로 보조정리
\ref{more-algebra-lemma-absolutely-flat-over-absolutely-flat}에 의해 $B$도 절대
평탄환이다. 따라서 보조정리 \ref{more-algebra-lemma-absolutely-flat}에 의해 $B$는
기약이고 모든 국소환은 체이다.

\medskip\noindent
$\mathfrak q \subset B$를 소아이디얼이라 하자. 환 사상
$B \to B_\mathfrak q$는 약 에탈이고, 따라서 $B_\mathfrak q$는
$K$ 위 약 에탈이다(보조정리 \ref{more-algebra-lemma-composition-weakly-etale}).
그러므로 보조정리 \ref{more-algebra-lemma-absolutely-flat-fields}에 의해
$B_\mathfrak q$는 $K$의 분리가능 대수적 확대이다.

\medskip\noindent
% P02-E0520: frozen English splits the compound noun “subalgebra”.
$K \subset A \subset B$를 유한 생성 $K$-부분대수라 하자. $A$가 $K$ 위
에탈임을 보이면 보조정리의 증명이 끝난다. 그러면 모든 극소소아이디얼
$\mathfrak p \subset A$는 $B$의 어떤 소아이디얼 $\mathfrak q$의 상이다.
Algebra, Lemma \ref{algebra-lemma-injective-minimal-primes-in-image}를 보라.
따라서 $B_\mathfrak q = \kappa(\mathfrak q)$의 부분체인
$\kappa(\mathfrak p)$는 $K$ 위 분리가능 대수적이다. 그러므로
$\Spec(A)$의 모든 일반점은 닫힌점이다(Algebra, Lemma
\ref{algebra-lemma-finite-residue-extension-closed}). 따라서
$\dim(A) = 0$이다. 그러면 예를 들어 Algebra, Proposition
\ref{algebra-proposition-dimension-zero-ring}에 의해 $A$는 그 국소환들의
곱이다. 더 나아가 $A$가 기약이므로 모든 국소환은 각자의 잉여체와 같고,
이 잉여체들은 $K$의 유한 분리가능 확대이다. 이는 Algebra, Lemma
\ref{algebra-lemma-etale-over-field}에 의해 $A$가 $K$ 위 에탈이라는
뜻이며, 증명이 끝난다.
\end{proof}

\begin{lemma}
\label{more-algebra-lemma-weakly-etale-residue-field-extensions}
$A \to B$를 환 사상이라 하자. $A \to B$가 약 에탈이면
$A \to B$는 분리가능 대수적 잉여체 확대들을 유도한다.
\end{lemma}

\begin{proof}
$\mathfrak p$를 $A$의 소아이디얼이라 하자. 그러면 보조정리
\ref{more-algebra-lemma-base-change-weakly-etale}에 의해
$\kappa(\mathfrak p) \to B \otimes_A \kappa(\mathfrak p)$는 약
에탈이다. 따라서 보조정리 \ref{more-algebra-lemma-absolutely-flat-over-field}에 의해
$B \otimes_A \kappa(\mathfrak p)$는 에탈
$\kappa(\mathfrak p)$-대수들의 여과 여극한이다. 그러므로
$\mathfrak p$ 위에 놓인 $\mathfrak q \subset B$에 대하여 확대
$\kappa(\mathfrak q)/\kappa(\mathfrak p)$는 Algebra, Lemma
\ref{algebra-lemma-etale-over-field}에 의해 유한 분리가능 확대들의
여과 여극한이다.
\end{proof}

\begin{lemma}
\label{more-algebra-lemma-weak-dimension-at-most-1}
$A$를 환이라 하자. 다음은 동치이다.
\begin{enumerate}
\item $A$의 약한 차원이 $\leq 1$이다.
\item $A$의 모든 아이디얼은 평탄이다.
\item $A$의 모든 유한 생성 아이디얼은 평탄이다.
\item 평탄 $A$-가군의 모든 부분가군은 평탄이다.
\item $A$의 모든 국소환은 값매김환이다.
\end{enumerate}
\end{lemma}

\begin{proof}
$A$의 약한 차원이 $\leq 1$이면 분해
$0 \to I \to A \to A/I \to 0$으로부터 보조정리
\ref{more-algebra-lemma-last-one-flat}에 의해 모든 아이디얼 $I$가 평탄임을 알 수
있다. 따라서 (1) $\Rightarrow$ (2)이다.

\medskip\noindent
(4)를 가정하자. $M$을 $A$-가군이라 하자. $F$가 자유 $A$-가군인
전사 사상 $F \to M$을 택하자. 그러면 가정에 의해
$\Ker(F \to M)$은 평탄이고, 보조정리 \ref{more-algebra-lemma-tor-dimension}에 의해
$M$의 Tor 차원이 $\leq 1$임을 알 수 있다. 따라서
(4) $\Rightarrow$ (1)이다.

\medskip\noindent
모든 아이디얼은 그 안에 포함되는 유한 생성 아이디얼들의 합집합이다.
따라서 Algebra, Lemma \ref{algebra-lemma-colimit-flat}에 의해 (3)은
(2)를 함의한다. 그러므로 (3) $\Leftrightarrow$ (2)이다.

\medskip\noindent
(2)를 가정하자. $N \subset M$이고 $M$이 평탄 $A$-가군이라고 하자.
$N$이 평탄임을 증명하겠다. Algebra, Theorem
\ref{algebra-theorem-lazard}를 사용하여 각 $M_i$가 유한 자유인
$M = \colim M_i$로 쓸 수 있다. $N_i \subset M_i$를 $N$의 역상으로
놓으면 $N = \colim N_i$임을 알 수 있다.
% P02-E0521: frozen punctuation breaks the “By ... it suffices” sentence.
Algebra, Lemma \ref{algebra-lemma-colimit-flat}에 의해 $N_i$가 평탄임을
증명하면 충분하며, 경우 $M = R^{\oplus n}$로 환원된다.
% P02-E0523: frozen proof switches from ambient A to undefined R in both formulas.
이 경우 가군 $N$은 부분가군 $R^{\oplus j} \cap N$들에 의한 유한
여과를 가지며, 그 부분몫들은 아이디얼이다. (2)에 의해 이 아이디얼들은
평탄이고, 따라서 Algebra, Lemma \ref{algebra-lemma-flat-ses}에 의해
$N$은 평탄이다. 그러므로 (2) $\Rightarrow$ (4)이다.

\medskip\noindent
$A$가 (1)을 만족한다고 가정하고 $\mathfrak p \subset A$를 소아이디얼이라
하자. 보조정리들 \ref{more-algebra-lemma-when-weakly-etale}와
\ref{more-algebra-lemma-weak-dimension-goes-up}에 의해 $A_\mathfrak p$가 (1)을
만족함을 알 수 있다. $A$가 (3)을 만족하는 국소환이면 $A$가
값매김환임을 보이자.
$f \in \mathfrak m$을 영이 아닌 원소라 하자. 그러면 $(f)$는 한
원소로 생성되는 영이 아닌 평탄 가군이다. 따라서 Algebra, Lemma
\ref{algebra-lemma-finite-flat-local}에 의해 자유 $A$-가군이다. 이로부터
$f$는 비영인자이고 $A$는 정역임을 알 수 있다. $I \subset A$가 유한
생성 아이디얼이면, 마찬가지로 $I$가 유한 자유 $A$-가군이고, 따라서
% P02-E1189: frozen rank-1 claim overlooks the zero finitely generated ideal.
(계수를 생각하면) 계수 $1$인 자유 가군이며 $I$는 주아이디얼임을 알 수
있다. 그러므로 Algebra, Lemma
\ref{algebra-lemma-characterize-valuation-ring}에 의해 $A$는
값매김환이다. 따라서 (1) $\Rightarrow$ (5)이다.

\medskip\noindent
(5)를 가정하자. $I \subset A$를 유한 생성 아이디얼이라 하자. 그러면
$I_\mathfrak p \subset A_\mathfrak p$는 값매김환의 유한 생성
아이디얼이므로 주아이디얼이고(Algebra, Lemma
\ref{algebra-lemma-characterize-valuation-ring}), 따라서 평탄이다.
그러므로 Algebra, Lemma \ref{algebra-lemma-flat-localization}에 의해
$I$는 평탄이다. 따라서 (5) $\Rightarrow$ (3)이다. 이로써 보조정리의
증명이 끝난다.
\end{proof}

\begin{lemma}
\label{more-algebra-lemma-product-weak-dimension-at-most-1}
$J$를 집합이라 하자. 각 $j \in J$에 대하여 $A_j$를 분수체가
$K_j$인 값매김환이라 하자. $A = \prod A_j$와 $K = \prod K_j$라 놓자.
그러면 $A$의 약한 차원은 기껏해야 $1$이고 $A \to K$는 국소화이다.
\end{lemma}

\begin{proof}
$I \subset A$를 유한 생성 아이디얼이라 하자. 보조정리
\ref{more-algebra-lemma-weak-dimension-at-most-1}에 의해 $I$가 평탄 $A$-가군임을
보이면 충분하다. $I_j \subset A_j$를 $I$의 상이라 하자.
$I_j = I \otimes_A A_j$임을 관찰하자. 따라서 Algebra, Proposition
\ref{algebra-proposition-fg-tensor}에 의해 $I \to \prod I_j$는
전사이다. 그러므로 $I = \prod I_j$이다. $A_j$가 값매김환이므로
아이디얼 $I_j$는 한 원소로 생성된다(Algebra, Lemma
\ref{algebra-lemma-characterize-valuation-ring}). $I_j = (f_j)$라 하자.
그러면 $I$는 원소 $f = (f_j)$로 생성된다. $f_j$가 $0$인지 아닌지에
따라 $A_j$에서 $0$ 또는 $1$인 멱등원 $e \in A$를 택하자. 그러면
$f = g e$인 비영인자 $g \in A$가 존재한다.
% P02-E0524: frozen English uses programming-style “else” in the componentwise definition.
$g = (g_j)$로서 $f_j = 0$이면 $g_j = 1$이고 그렇지 않으면
$g_j = f_j$인 것을 택하면 된다. 따라서 가군으로서 $I \cong (e)$이다.
$(e)$는 $A$의 직합인자이므로 $I$가 평탄임을 알 수 있다. 마지막 명제는
$S = \prod (A_j \setminus \{0\})$일 때 $K = S^{-1}A$이므로 참이다.
\end{proof}

\begin{lemma}
\label{more-algebra-lemma-product-found-valuation-rings}
$A$를 분수체가 $K$인 정규 정역이라 하자. 환들의 Cartesian 그림
$$
\xymatrix{
A \ar[d] \ar[r] & K \ar[d] \\
V \ar[r] & L
}
$$
로서 $V$의 약한 차원이 기껏해야 $1$이고 $V \to L$이 평탄 단사
환 전사인 것이 존재한다.
\end{lemma}

\begin{proof}
% P02-E0525: frozen English has a malformed quantifier; the intended index set is K minus A.
각 $x \in K$, $x \not \in A$에 대하여 Algebra, Lemma
\ref{algebra-lemma-find-valuation-rings}에서와 같은 $V_x \subset K$를
택하자. $V = \prod_{x \in K \setminus A} V_x$와
$L = \prod_{x \in K \setminus A} K$라 놓자. 보조정리
\ref{more-algebra-lemma-product-weak-dimension-at-most-1}에 의해 환 $V$의 약한
차원은 기껏해야 $1$이고, 같은 보조정리에 의해 $V \to L$은
국소화이다. 국소화는 평탄이고 환 전사이다. Algebra, Lemmas
\ref{algebra-lemma-flat-localization}와
\ref{algebra-lemma-epimorphism-local}을 보라.
\end{proof}

\begin{lemma}
\label{more-algebra-lemma-weak-dimension-at-most-1-integrally-closed}
$A$를 약한 차원이 기껏해야 $1$인 환이라 하자.
$A \to B$가 평탄 단사 환 전사이면 $A$는 $B$ 안에서 정수적으로
닫혀 있다.
\end{lemma}

\begin{proof}
$x \in B$가 $A$ 위 정수적이라고 하자. $A' = A[x] \subset B$라 놓자.
그러면 Algebra, Lemma
\ref{algebra-lemma-characterize-finite-in-terms-of-integral}에 의해
$A'$은 $A$의 유한 환 확대이다. $A = A'$임을 보이기 위해서는 Algebra,
Lemma \ref{algebra-lemma-finite-epimorphism-surjective}에 의해
$A \to A'$이 환 전사임을 보이면 충분하다. $A$에 대한 가정과 $B$가
$A$ 위 평탄이라는 사실(보조정리
\ref{more-algebra-lemma-weak-dimension-at-most-1})에 의해 $A'$은 $A$ 위 평탄임을
주목하자. 따라서 합성
$$
A' \otimes_A A' \to B \otimes_A A' \to B \otimes_A B \to B
$$
은 단사이다. 즉, $A' \otimes_A A' \cong A'$이고 보조정리가
증명되었다.
\end{proof}

\begin{lemma}
\label{more-algebra-lemma-normality-goes-up}
$A$를 분수체가 $K$인 정규 정역이라 하자. $A \to B$가 약 에탈이라고
하자. 그러면 $B$는 $B \otimes_A K$ 안에서 정수적으로 닫혀 있다.
\end{lemma}

\begin{proof}
보조정리 \ref{more-algebra-lemma-product-found-valuation-rings}에서와 같은 그림을
택하자. $A \to B$가 평탄이므로 밑변환하여 환들의 Cartesian 그림
$$
\xymatrix{
B \ar[d] \ar[r] & B \otimes_A K \ar[d] \\
B \otimes_A V \ar[r] & B \otimes_A L
}
$$
을 얻는다. $V \to B \otimes_A V$가 약 에탈임에 주의하자(보조정리
\ref{more-algebra-lemma-base-change-weakly-etale}). 따라서 보조정리
\ref{more-algebra-lemma-weak-dimension-goes-up}에 의해 $B \otimes_A V$의 약한
차원은 기껏해야 $1$이다. $B \otimes_A V \to B \otimes_A L$은
그러한 사상의 평탄 밑변환이므로 평탄 단사 환 전사이다(Algebra, Lemmas
\ref{algebra-lemma-flat-base-change}와
\ref{algebra-lemma-base-change-epimorphism}). 보조정리
\ref{more-algebra-lemma-weak-dimension-at-most-1-integrally-closed}에 의해
$B \otimes_A V$가 $B \otimes_A L$ 안에서 정수적으로 닫혀 있음을
알 수 있다. 그림의 Cartesian 성질로부터 $B$가 $B \otimes_A K$ 안에서
정수적으로 닫혀 있음을 얻는다.
\end{proof}

\begin{lemma}
\label{more-algebra-lemma-integral-over-henselian}
$A \to B$를 환 준동형이라 하자. 다음을 가정하자.
\begin{enumerate}
\item $A$는 헨젤 국소환이다.
\item $A \to B$는 정수적이다.
\item $B$는 정역이다.
\end{enumerate}
그러면 $B$는 헨젤 국소환이고 $A \to B$는 국소 준동형이다.
$A$가 엄밀 헨젤이면 $B$는 엄밀 헨젤 국소환이고 잉여체의 확대
$\kappa(\mathfrak m_B)/\kappa(\mathfrak m_A)$는 순수 비분리이다.
\end{lemma}

\begin{proof}
% P02-E0520: frozen English splits the compound noun “subalgebras”.
$B$를 유한 $A$-부분대수들의 여과 여극한 $B = \colim B_i$로 쓰자.
각 $B_i$에 대하여 결과를 증명하면 $B$에 대한 결과가 따른다. Algebra,
Lemma \ref{algebra-lemma-colimit-henselian}을 보라. $A \to B$가
유한이면 Algebra, Lemma \ref{algebra-lemma-finite-over-henselian}에
의해 $B$는 헨젤 국소환들의 곱이다. $B$가 정역이므로 $B$는
국소환이다. $B$의 극대아이디얼은 $A \to B$에 대한 올라가기 정리에
의해 $A$의 극대아이디얼 위에 놓인다(Algebra, Lemma
\ref{algebra-lemma-integral-going-up}). $A$가 엄밀 헨젤이면
% P02-E1191: frozen English uses a dangling “being algebraic” construction.
체 확대 $\kappa(\mathfrak m_B)/\kappa(\mathfrak m_A)$가 대수적이므로
순수 비분리이어야 한다. 물론 그러면 $\kappa(\mathfrak m_B)$는 분리가능
대수적으로 닫혀 있고 $B$는 엄밀 헨젤이다.
\end{proof}

\begin{theorem}[Olivier]
\label{more-algebra-theorem-olivier}
$A \to B$를 국소환들의 국소 준동형이라 하자. $A$가 엄밀 헨젤이고
$A \to B$가 약 에탈이면 $A = B$이다.
\end{theorem}

\begin{proof}
모든 $\mathfrak p \subset A$에 대하여 $\mathfrak p$ 위에 놓이는
유일한 소아이디얼 $\mathfrak q \subset B$가 존재하고
$\kappa(\mathfrak p) = \kappa(\mathfrak q)$임을 보이겠다. 이는
$B \otimes_A B \to B$가 스펙트럼 위에서 전단사이고 전사이며
평탄임을 함의한다. 따라서 예를 들어 Algebra, Lemma
\ref{algebra-lemma-pure-open-closed-specializations}의 순수 아이디얼에
대한 기술에 의해 이는 동형사상이다. 그러므로 $A \to B$는 충실 평탄
환 전사이다. Algebra, Lemma
\ref{algebra-lemma-faithfully-flat-epimorphism}에 의해 $A = B$를 얻는다.

\medskip\noindent
보조정리들 \ref{more-algebra-lemma-base-change-weakly-etale}와
\ref{more-algebra-lemma-absolutely-flat-over-field}에 의해 올환
$B \otimes_A \kappa(\mathfrak p)$는 $\kappa(\mathfrak p)$의 에탈
확대들의 여극한임에 주의하자. 따라서 $\mathfrak p$ 위에 놓이는
소아이디얼이 둘 이상 존재하거나, 어떤 $\mathfrak q$에 대하여
$\kappa(\mathfrak p) \not = \kappa(\mathfrak q)$이면, 어떤
(분리가능) 대수적 체 확대 $L/\kappa(\mathfrak p)$에 대하여
$B \otimes_A L$은 자명하지 않은 멱등원을 갖는다.

\medskip\noindent
$L/\kappa(\mathfrak p)$을 대수적 체 확대라 하자. $A' \subset L$을
$L$ 안에서 $A/\mathfrak p$의 정수적 폐포라 하자. 보조정리
\ref{more-algebra-lemma-integral-over-henselian}에 의해 $A'$은 그 잉여체가 $A$의
잉여체의 순수 비분리 확대인 엄밀 헨젤 국소환이다. 따라서 Algebra,
Lemma \ref{algebra-lemma-local-tensor-with-integral}에 의해
$B \otimes_A A'$은 국소환이다. 다른 한편으로 보조정리
\ref{more-algebra-lemma-normality-goes-up}에 의해 $B \otimes_A A'$은
$B \otimes_A L$ 안에서 정수적으로 닫혀 있다.
$B \otimes_A A'$이 국소이므로 환 $B \otimes_A L$은 자명하지 않은
멱등원을 갖지 않는다. 이것이 증명하고자 한 바이다.
\end{proof}

\section{체 위의 약 에탈 대수}
\label{more-algebra-section-weakly-etale-over-field}

\noindent
$K$가 체이면, 대수 $B$가 $K$ 위 약 에탈일 필요충분조건은 에탈
$K$-대수들의 여과 여극한인 것이다. 이는 보조정리
\ref{more-algebra-lemma-absolutely-flat-over-field}이다.

\begin{lemma}
\label{more-algebra-lemma-class-weakly-etale-over-field}
$K$를 체라 하자. $B$가 $K$ 위 약 에탈이면 다음이 성립한다.
\begin{enumerate}
\item $B$는 기약이다.
\item $B$는 $K$ 위 정수적이다.
\item $B$의 모든 유한 생성 $K$-부분대수는 $K$의 유한 분리가능
확대들의 유한 곱이다.
\item $B$가 체일 필요충분조건은 $B$가 자명하지 않은 멱등원을 갖지
않는 것이며, 이 경우 이는 $K$의 분리가능 대수적 확대이다.
% P02-E0526: frozen English ungrammatically coordinates “sub” with “quotient algebra”.
\item $B$의 임의의 부분대수 또는 몫 $K$-대수는 $K$ 위 약 에탈이다.
\item $B'$가 $K$ 위 약 에탈이면 $B \otimes_K B'$도 $K$ 위 약 에탈이다.
\end{enumerate}
\end{lemma}

\begin{proof}
(1)은 보조정리 \ref{more-algebra-lemma-absolutely-flat-over-absolutely-flat}에서
따르지만, 물론 (3)에서도 따른다. (3)은 보조정리
\ref{more-algebra-lemma-absolutely-flat-over-field}와, 에탈 $K$-대수는 $K$의 유한
분리가능 확대들의 유한 곱이라는 사실에서 따른다. Algebra, Lemma
\ref{algebra-lemma-etale-over-field}을 보라. (3)은 (2)를 함의한다.
체들의 곱이 체일 필요충분조건은 자명하지 않은 멱등원을 갖지 않는
것이므로 (4)는 (3)에서 따른다.

\medskip\noindent
$S \subset B$가 부분대수이면, 이는 그 유한 생성 부분대수들의 여과
여극한이고 위에서 보았듯 이들은 모두 $K$ 위 에탈이다. 따라서 보조정리
\ref{more-algebra-lemma-absolutely-flat-over-field}에 의해 $S$는 $K$ 위 약
에탈이다. $B \to Q$가 몫대수이면 $Q$는 유한 분리가능 확대
$L_i/K$들의 유한 곱 $\prod_{i \in I} L_i$의 $K$-대수 몫들의 여과
여극한이다. 그러한 몫은 어떤 부분집합 $J \subset I$에 대하여
$\prod_{i \in J} L_i$의 꼴이다. 따라서 같은 논증으로 몫에 대해서도
결과가 성립한다.

\medskip\noindent
텐서곱에 관한 명제도 같은 방식으로 따르거나 보조정리들
\ref{more-algebra-lemma-base-change-weakly-etale}와
\ref{more-algebra-lemma-composition-weakly-etale}를 결합하면 따른다.
\end{proof}

\begin{lemma}
\label{more-algebra-lemma-max-weakly-etale-subalgebra}
$K$를 체라 하자. $A$를 $K$-대수라 하자. 극대 약 에탈 $K$-부분대수
$B_{max} \subset A$가 존재한다.
\end{lemma}

\begin{proof}
$B_1, B_2 \subset A$를 약 에탈 $K$-부분대수라 하자. 그러면
$B_1 \otimes_K B_2$는 $K$ 위 약 에탈이고, 그 사상
$B_1 \otimes_K B_2 \to A$의 상도 그러하다(보조정리
\ref{more-algebra-lemma-class-weakly-etale-over-field}). 따라서 약 에탈
$K$-부분대수 $B \subset A$들의 모임 $\mathcal{B}$는 유향적이고,
보조정리 \ref{more-algebra-lemma-when-weakly-etale}에 의해 여극한
$B_{max} = \colim_{B \in \mathcal{B}} B$는 약 에탈 $K$-대수이다.
따라서 $B_{max} \to A$의 상은 $K$ 위 약 에탈이다(앞의 보조정리).
그러므로 이 상은 $\mathcal{B}$에 속하고 $\mathcal{B}$는 극대원소를
갖는다(그리고 이 상은 $B_{max}$와 같다).
\end{proof}

\begin{lemma}
\label{more-algebra-lemma-properties-of-max-weakly-etale-subalgebra}
$K$를 체라 하자. $K$-대수 $A$에 대하여, 보조정리
\ref{more-algebra-lemma-max-weakly-etale-subalgebra}에서와 같은 $A$의 극대 약 에탈
% P02-E0527: frozen English notation clause omits the required “by”.
$K$-부분대수를 $B_{max}(A)$로 나타내자. 그러면 다음이 성립한다.
\begin{enumerate}
\item 임의의 $K$-대수 사상 $A' \to A$는 $K$-대수 사상
$B_{max}(A') \to B_{max}(A)$를 유도한다.
\item $A' \subset A$이면 $B_{max}(A') = B_{max}(A) \cap A'$이다.
\item $A = \colim A_i$가 여과 여극한이면
$B_{max}(A) = \colim B_{max}(A_i)$이다.
\item 사상 $B_{max}(A) \to B_{max}(A_{red})$는 동형사상이다.
\item $B_{max}(A_1 \times \ldots \times A_n) =
B_{max}(A_1) \times \ldots \times B_{max}(A_n)$이다.
\item $A$가 자명하지 않은 멱등원을 갖지 않으면 $B_{max}(A)$는 체이며
$K$의 분리가능 대수적 확대이다.
% P02-E0528: frozen item (7) is an unresolved editorial placeholder.
\item 여기에 더 추가할 것.
\end{enumerate}
\end{lemma}

\begin{proof}
(1)의 증명. 보조정리 \ref{more-algebra-lemma-class-weakly-etale-over-field}에 의해
$B_{max}(A') \to A$의 상이 $K$ 위 약 에탈이므로 이는 참이다.

\medskip\noindent
(2)의 증명. (1)에 의해 $B_{max}(A') \subset B_{max}(A)$이다.
반대로 보조정리 \ref{more-algebra-lemma-class-weakly-etale-over-field}에 의해
$B_{max}(A) \cap A'$은 약 에탈 $K$-대수이고, 따라서
$B_{max}(A')$에 포함된다.

\medskip\noindent
(3)의 증명. (1)에 의해 사상 $\colim B_{max}(A_i) \to A$가 존재한다.
계가 여과되어 있고 $B_{max}(A_i) \subset A_i$이므로 이 사상은
단사이다. 보조정리 \ref{more-algebra-lemma-when-weakly-etale}에 의해 여극한
$\colim B_{max}(A_i)$는 $K$ 위 약 에탈이다. 따라서 단사 사상
$\colim B_{max}(A_i) \to B_{max}(A)$를 얻는다.
$a \in B_{max}(A)$라고 하자. 그러면 $a$는 유한 표시
$K$-부분대수 $B \subset B_{max}(A)$를 생성한다. Algebra, Lemma
\ref{algebra-lemma-characterize-finite-presentation}에 의해 어떤
$i$와 $K$-대수 사상 $f : B \to A_i$가 존재하여 주어진 사상
$B \to A$를 올린다. 보조정리
\ref{more-algebra-lemma-class-weakly-etale-over-field}에 의해 $B$가 약 에탈이므로
$f(B) \subset B_{max}(A_i)$임을 알 수 있다. 따라서 $a$는
$\colim B_{max}(A_i) \to B_{max}(A)$의 상에 속한다.

\medskip\noindent
(4)의 증명. 보조정리 \ref{more-algebra-lemma-absolutely-flat-over-field}을 사용하여
$B_{max}(A_{red}) = \colim B_i$를 에탈 $K$-대수들의 여과 여극한으로
쓰자. Algebra, Lemma \ref{algebra-lemma-smooth-strong-lift}에 의해 각
$i$에 대하여 주어진 사상 $B_i \to A_{red}$을 올리는 $K$-대수 사상
$f_i : B_i \to A$가 존재한다. 따라서
% P02-E0529: frozen proof reverses the canonical map from assertion (4).
정준사상 $B_{max}(A_{red}) \to B_{max}(A)$는 전사이다. 그 핵은 멱영
원소들로 이루어지고,
% P02-E0530: frozen proof invokes reducedness of the wrong ring to kill that kernel.
$B_{max}(A_{red})$가 기약이므로 영이다(보조정리
\ref{more-algebra-lemma-class-weakly-etale-over-field}).

\medskip\noindent
(5)의 증명.
% P02-E1195: the nontrivial product assertion is left with only “Omitted.” in the frozen source.
생략한다.

\medskip\noindent
(6)의 증명. 보조정리 \ref{more-algebra-lemma-class-weakly-etale-over-field}의
(4)에서 따른다.
\end{proof}

\begin{lemma}
\label{more-algebra-lemma-change-fields-max-weakly-etale-subalgebra}
$L/K$를 체 확대라 하자. $A$를 $K$-대수라 하자. $B \subset A$를
보조정리 \ref{more-algebra-lemma-max-weakly-etale-subalgebra}에서와 같은 $A$의 극대
약 에탈 $K$-부분대수라 하자. 그러면 $B \otimes_K L$은
$A \otimes_K L$의 극대 약 에탈 $L$-부분대수이다.
\end{lemma}

\begin{proof}
$K$ 위 대수 $A$에 대하여 $A$의 극대 약 에탈 $K$-부분대수를
$B_{max}(A/K)$로 쓰자. 마찬가지로 $A'$이 $L$-대수이면 $A'$의 극대
약 에탈 $L$-부분대수를 $B_{max}(A'/L)$로 쓰자.
$B_{max}(A/K) \otimes_K L$은 $L$ 위 약 에탈이고(보조정리
\ref{more-algebra-lemma-base-change-weakly-etale}),
$B_{max}(A/K) \otimes_K L \subset A \otimes_K L$이므로 정준 단사
사상을 얻는다.
% P02-E0533: frozen display ending a complete sentence lacks terminal punctuation.
$$
B_{max}(A/K) \otimes_K L \to B_{max}((A \otimes_K L)/L)
$$
보조정리의 명제는 이 사상이 동형사상이라는 것이다.

\medskip\noindent
$L$과 우리의 $K$-대수 $A$에 대하여 보조정리를 증명하려면 $L$의
임의의 체 확대 $L'$에 대하여 보조정리를 증명하면 충분하다. 실제로
다음 분해가 있다.
% P02-E0533: frozen factorization display lacks punctuation before the result clause.
$$
B_{max}(A/K) \otimes_K L' \to
B_{max}((A \otimes_K L)/L) \otimes_L L' \to
B_{max}((A \otimes_K L')/L')
$$
따라서 합성이 전사이면서
$B_{max}(A/K) \otimes_K L \to B_{max}((A \otimes_K L)/L)$이 전사가
아닐 수는 없다. 그러므로 $L$이 대수적으로 닫혀 있다고 가정해도 된다.

\medskip\noindent
유한형 $K$-대수로의 환원.
% P02-E0531: frozen English says “write A is” instead of “write A as”.
$A$를 그 유한형 $K$-부분대수들의 여과 여극한으로 쓸 수 있다. 보조정리
\ref{more-algebra-lemma-properties-of-max-weakly-etale-subalgebra}를 사용하면 유한형
$K$-대수에 대하여 보조정리를 증명하면 충분함을 알 수 있다.

\medskip\noindent
$A$가 유한형 $K$-대수라고 가정하자. $A \to A_{red}$의 핵은 멱영이므로
$A \otimes_K L \to A_{red} \otimes_K L$의 핵도 그러하다. 그러면
$$
B_{max}((A \otimes_K L)/L) \to B_{max}((A_{red} \otimes_K L)/L)
$$
은 단사이다. 실제로 그 핵은 멱영이고, 약 에탈 $L$-대수
$B_{max}((A \otimes_K L)/L)$은 기약이다(보조정리
\ref{more-algebra-lemma-class-weakly-etale-over-field}). 보조정리
\ref{more-algebra-lemma-properties-of-max-weakly-etale-subalgebra}에 의해
$B_{max}(A/K) = B_{max}(A_{red}/K)$이므로 $A_{red}$에 대하여 보조정리를
증명하면 충분하다고 결론짓는다.

\medskip\noindent
$A$를 기약인 유한형 $K$-대수라 하자. $Q = Q(A)$를 $A$의 전체
몫환이라 하자. 그러면 $A \subset Q$이고
% P02-E0532: frozen inclusion tensors Q with L over A although L is only a K-algebra.
$A \otimes_K L \subset Q \otimes_A L$이다. 따라서
$$
B_{max}(A/K) = A \cap B_{max}(Q/K)
$$
이고
$$
B_{max}((A \otimes_K L)/L) =
(A \otimes_K L) \cap B_{max}((Q \otimes_K L)/L)
$$
이다. 이는 보조정리 \ref{more-algebra-lemma-properties-of-max-weakly-etale-subalgebra}에서
따른다. $-\otimes_K L$이 완전함자이므로 $Q$에 대하여 결과를 증명하면
$A$에 대한 결과가 따른다. $Q$는 체들의 유한 곱이고(Algebra, Lemmas
\ref{algebra-lemma-total-ring-fractions-no-embedded-points},
\ref{algebra-lemma-minimal-prime-reduced-ring},
\ref{algebra-lemma-Noetherian-irreducible-components}, and
\ref{algebra-lemma-Noetherian-permanence}), $B_{max}$는 곱과 가환하므로
(보조정리 \ref{more-algebra-lemma-properties-of-max-weakly-etale-subalgebra}),
$A$가 체인 경우에 보조정리를 증명하면 충분하다.

\medskip\noindent
$A$가 체라고 가정하자. 증명의 세 번째 문단의 논증을 사용하여 $A$가
$K$ 위 유한 생성이라고 환원한다. (사실 체인 경우로 환원한 방식은
$K$의 유한 생성 체 확대를 만든다.)

\medskip\noindent
$A$를 $K$의 유한 생성 체 확대라 하자. 그러면 보조정리
\ref{more-algebra-lemma-properties-of-max-weakly-etale-subalgebra}의 (6)에 의해
$K' = B_{max}(A/K)$는 $K$ 위 분리가능 대수적인 체이다. 따라서
$K'$은 $K$의 유한 분리가능 체 확대이고, Algebra, Lemma
\ref{algebra-lemma-make-geometrically-irreducible}에 의해 $A$는 $K'$ 위
기하적으로 기약이다. $L$이 대수적으로 닫혀 있고 $K'/K$가 유한
분리가능이므로 다음 사상은 동형사상이다.
$$
K' \otimes_K L \to \prod\nolimits_{\sigma \in \Hom_K(K', L)} L,\quad
\alpha \otimes \beta \mapsto (\sigma(\alpha)\beta)_\sigma
$$
(Fields, Lemma \ref{fields-lemma-finite-separable-tensor-alg-closed}).
따라서
% P02-E0533: frozen display ending a complete sentence lacks terminal punctuation.
$$
A \otimes_K L = A \otimes_{K'} (K' \otimes_K L) =
\prod\nolimits_{\sigma \in \Hom_K(K', L)} A \otimes_{K', \sigma} L
$$
이다. $A$가 $K'$ 위 기하적으로 기약이므로
$A \otimes_{K', \sigma} L$은 유일한 극소소아이디얼을 갖는다. $L$이
대수적으로 닫혀 있으므로
$B_{max}((A \otimes_{K', \sigma} L)/L) = L$이다. 실제로 보조정리
\ref{more-algebra-lemma-properties-of-max-weakly-etale-subalgebra}의 (6)에 의해 이
$L$-대수는 $L$ 위 대수적인 체이다. 위와 같이 이 부분대수들을 곱으로
분해할 수 있으므로 $A \otimes_K L$의 극대 약 에탈
$K' \otimes_K L$-부분대수는 $K' \otimes_K L$이다. 따라서 포함관계
$K' \otimes_K L \subset B_{max}((A \otimes_K L)/L)$는 등식이다.
실제로 보조정리 \ref{more-algebra-lemma-weakly-etale-permanence}에 의해 환 사상
$K' \otimes_K L \to B_{max}((A \otimes_K L)/L)$은 약 에탈이다.
\end{proof}

\section{국소 기약성}
\label{more-algebra-section-unibranch}

\noindent
다음 정의가 일반적으로 받아들여지는 정의인 듯하다. 이를 해석하기 위해,
국소 정역들의 정수적 확대가 $A \subset B$이면 going up에 의해
$A \to B$가 국소환 준동형임을 관찰하자(Algebra, Lemma
\ref{algebra-lemma-integral-going-up}).

\begin{definition}
\label{more-algebra-definition-unibranch}
\begin{reference}
\cite[Chapter 0 (23.2.1)]{EGA4}
\end{reference}
$A$를 국소환이라 하자. 축약 $A_{red}$가 정역이고 그 분수체 안에서
$A_{red}$의 정수적 폐포 $A'$이 국소환이면 $A$를 {\it 단일분지}라고
한다. $A$가 단일분지이고, 더 나아가 $A'$의 잉여체가 $A$의 잉여체
위 순수 비분리이면 $A$를 {\it 기하적 단일분지}라고 한다.
\end{definition}

\noindent
$A$를 국소환이라 하자. 다음은 동치인 정식화이다.
\begin{enumerate}
\item $A$가 단일분지라는 것은 $A$가 유일한 극소소아이디얼
$\mathfrak p$를 갖고, 그 분수체 안에서 $A/\mathfrak p$의 정수적 폐포가
국소환이라는 것이며,
\item $A$가 기하적 단일분지라는 것은 $A$가 유일한 극소소아이디얼
$\mathfrak p$를 갖고, 그 분수체 안에서 $A/\mathfrak p$의 정수적 폐포가
국소환이며 그 잉여체가 $A$의 잉여체 위 순수 비분리라는 것이다.
\end{enumerate}
정규인 국소환은 기하적 단일분지이다(정의
\ref{more-algebra-definition-unibranch}와 Algebra, Definition
\ref{algebra-definition-ring-normal}에서 따른다). 보조정리
\ref{more-algebra-lemma-unibranch}와 \ref{more-algebra-lemma-geometrically-unibranch}는 (기하적)
단일분지라는 것이 살펴볼 만한 합리적인 성질임을 시사한다.

\begin{lemma}
\label{more-algebra-lemma-branches}
$A$를 국소환이라 하자. $A$가 유한 개의 극소소아이디얼을 갖는다고
가정하자. $A'$을 $A_{red}$의 전체 몫환 안에서 $A$의 정수적 폐포라
하자. $A^h$를 $A$의 헨젤화라 하자. 다음 사상들을 생각하자.
$$
\Spec(A') \leftarrow \Spec((A')^h) \rightarrow \Spec(A^h)
$$
여기서 $(A')^h = A' \otimes_A A^h$이다. 그러면
\begin{enumerate}
\item 왼쪽 화살표는 극대아이디얼 위에서 전단사이고,
\item 오른쪽 화살표는 극소소아이디얼 위에서 전단사이며,
\item $(A')^h$의 모든 극소소아이디얼은 유일한 극대아이디얼에 포함되고
모든 극대아이디얼은 정확히 하나의 극소소아이디얼을 포함한다.
\end{enumerate}
\end{lemma}

\begin{proof}
$I \subset A$를 멱영원들의 아이디얼이라 하자. Algebra, Lemma
\ref{algebra-lemma-quotient-henselization}에 의해
$(A/I)^h = A^h/IA^h$이다. $A$, $A^h$, $A'$, $(A')^h$의 스펙트럼은 각각
$A/I$, $A^h/IA^h$, $A'$, 그리고
$(A')^h = A' \otimes_{A/I} A^h/IA^h$의 스펙트럼과 같다. 따라서 $A$를
$A_{red} = A/I$로 바꾸어 $A$가 기약이라고 가정해도 된다. 그러면
$A \subset A'$이며, 아래에서는 더 언급하지 않고 이를 사용한다.

\medskip\noindent
(1)의 증명. $A'$은 $A$ 위 정수적이므로 $(A')^h$는 $A^h$ 위 정수적이다.
Going up(Algebra, Lemma \ref{algebra-lemma-integral-going-up})에 의해
$A'$, resp.\ $(A')^h$의 모든 극대아이디얼은 $A$, resp.\ $A^h$의
극대아이디얼 $\mathfrak m$, resp.\ $\mathfrak m^h$ 위에 놓인다. 따라서
다음 동형사상에서 (1)이 따른다.
$$
(A')^h \otimes_{A^h} \kappa^h =
A' \otimes_A A^h \otimes_{A^h} \kappa^h = A' \otimes_A \kappa
$$
이는 $A \to A^h$가 유도하는 잉여체 확대 $\kappa^h/\kappa$가
자명하기 때문이다.
% P02-E0534: frozen English omits an article and punctuation in the profinite-spectrum consequence.
아래에서 이 표시된 환이 체 위 정수적이고, 따라서 그 스펙트럼이
profinite 공간이라는 사실을 사용할 것이다. Algebra, Lemmas
\ref{algebra-lemma-integral-over-field}와
\ref{algebra-lemma-ring-with-only-minimal-primes}를 보라.

\medskip\noindent
(3)의 증명. 환 $A'$은 정규환이며, 실제로 정규 정역들의 유한 곱이다.
Algebra, Lemma \ref{algebra-lemma-characterize-reduced-ring-normal}를 보라.
$A^h$는 에탈 $A$-대수들의 여과 여극한이므로,
% P02-E0535: frozen English omits the article before “filtered colimit” here and in the strict analogue.
$(A')^h$는 에탈 $A'$-대수들의 여과 여극한이다. 따라서 Algebra, Lemmas
\ref{algebra-lemma-normal-goes-up}와
\ref{algebra-lemma-colimit-normal-ring}에 의해 $(A')^h$는 정규환이다.
그러므로 $(A')^h$의 모든 국소환은 정규 정역이고, 모든 극대아이디얼이
유일한 극소소아이디얼을 포함함을 알 수 있다. 보조정리
\ref{more-algebra-lemma-integral-over-henselian-pair}를 $A^h \to (A')^h$에 적용하면
$((A')^h, \mathfrak m(A')^h)$가 헨젤 쌍임을 알 수 있다.
$\mathfrak q \subset (A')^h$가 극소소아이디얼(또는 임의의
소아이디얼)이면, 보조정리
\ref{more-algebra-lemma-irreducible-henselian-pair-connected}에 의해
$V(\mathfrak q)$와 $V(\mathfrak m (A')^h)$의 교집합은 연결이다.
% P02-E0536: frozen English drops punctuation and contains the malformed connective “by we see”.
$V(\mathfrak m (A')^h) = \Spec((A')^h \otimes \kappa^h)$는 profinite
공간이므로 $\mathfrak q$를 포함하는 극대아이디얼은 유일하다.

\medskip\noindent
(2)의 증명. $A'$의 극소소아이디얼들은 정확히 $A$의 극소소아이디얼
위에 놓이는 소아이디얼들이다(구성에 의한다). $A' \to (A')^h$는
평탄하므로 going down(Algebra, Lemma
\ref{algebra-lemma-flat-going-down})에 의해 $(A')^h$의 모든
극소소아이디얼은 $A'$의 극소소아이디얼 위에 놓인다. 반대로
$(A')^h$의 소아이디얼이 $A'$의 극소소아이디얼 위에 놓이면 극소이다.
이는 $(A')^h$가 $A'$ 위 에탈, 따라서 준유한인 대수들의 여과
% P02-E1202: the frozen source explicitly omits the converse fibre detail.
여극한이기 때문이다(작은 세부사항은 생략한다). 따라서 $(A')^h$의
극소소아이디얼들은 정확히 $A$의 극소소아이디얼 위에 놓이는
소아이디얼들이다. 마찬가지로 $A^h$의 극소소아이디얼들은 정확히
$A$의 극소소아이디얼 위에 놓이는 소아이디얼들이다. 구성에 의해
$A' \otimes_A Q(A) = Q(A)$이고, 여기서 $Q(A)$는 기약 국소환 $A$의
전체 몫환이다. 물론 $Q(A)$는 $A$의 극소소아이디얼들의 잉여체들의
유한 곱이다. 따라서
% P02-E0537: frozen display lacks punctuation before the independent conclusion.
$$
(A')^h \otimes_A Q(A) = A^h \otimes_A A' \otimes_A Q(A) = A^h \otimes_A Q(A)
$$
이다. 위 논의에 따르면 왼쪽 환의 스펙트럼은 $(A')^h$의
극소소아이디얼들의 집합이고,
% P02-E0538: frozen English duplicates “is the” in the parallel spectrum sentence.
오른쪽 환의 스펙트럼은 $A^h$의 극소소아이디얼들의 집합이다. 이로써
증명이 끝난다.
\end{proof}

\begin{lemma}
\label{more-algebra-lemma-unibranch}
\begin{reference}
\cite[Chapter IV Proposition 18.6.12]{EGA4}
\end{reference}
$A$를 국소환이라 하고 $A^h$를 $A$의 헨젤화라 하자.
% P02-E1204: frozen English lacks punctuation before the equivalence list.
다음은 서로 동치이다.
\begin{enumerate}
\item $A$는 단일분지이고,
\item $A^h$는 유일한 극소소아이디얼을 갖는다.
\end{enumerate}
\end{lemma}

\begin{proof}
이는 보조정리 \ref{more-algebra-lemma-branches}에서 따르지만 직접증명도 제시하겠다.
% P02-E0539: frozen notation assignments omit “by”; the strict analogue also has number disagreement.
환 $A$의 극대아이디얼을 $\mathfrak m$으로 나타내자. 잉여체
$\kappa = A/\mathfrak m$은 $A^h$의 잉여체와 같음을 상기하자.

\medskip\noindent
(2)를 가정하자. $\mathfrak p^h$를 $A^h$의 유일한 극소소아이디얼이라
하자. $A \to A^h$의 평탄성에 의해
$\mathfrak p = A \cap \mathfrak p^h$는 $A$의 유일한 극소소아이디얼이다
(going down에 의한다. Algebra, Lemma
\ref{algebra-lemma-flat-going-down}을 보라). 또한 Algebra, Lemma
\ref{algebra-lemma-quotient-henselization}에서 보듯
$A^h/\mathfrak pA^h = (A/\mathfrak p)^h$이고, 이는 보조정리
\ref{more-algebra-lemma-henselization-reduced}에 의해 기약이므로
$\mathfrak p^h = \mathfrak pA^h$이다. $A'$을 그 분수체 안에서
$A/\mathfrak p$의 정수적 폐포라 하자. $A'$이 국소환임을 보여야 한다.
$A \to A'$은 정수적이므로 $A'$의 모든 극대아이디얼은 $\mathfrak m$
위에 놓인다(정수적 환 사상에 대한 going up에 의한다. Algebra, Lemma
\ref{algebra-lemma-integral-going-up}을 보라). $A'$이 국소환이 아니면
서로 다른 극대아이디얼 $\mathfrak m_1$, $\mathfrak m_2$를 찾을 수 있다.
$f_i \in \mathfrak m_i$이고 $f_i \not \in \mathfrak m_{3 - i}$인 원소
$f_1, f_2 \in A'$을 택하자.
% P02-E0540: frozen formal span has ambiguous/wrong quotient-adjoining precedence.
서로 다른 극대아이디얼 $B \cap \mathfrak m_i$, $i = 1, 2$를 갖는 유한
부분대수 $B = A/\mathfrak p[f_1, f_2] \subset A'$을 얻는다. 포함관계
$$
A/\mathfrak p \subset B \subset \kappa(\mathfrak p)
$$
를 평탄한 환 사상 $A \to A^h$로 텐서하면 다음 포함관계를 얻는다.
$$
A^h/\mathfrak p^h \subset
B \otimes_A A^h \subset
\kappa(\mathfrak p) \otimes_A A^h \subset
\kappa(\mathfrak p^h)
$$
마지막 포함관계는
$\kappa(\mathfrak p) \otimes_A A^h =
\kappa(\mathfrak p) \otimes_{A/\mathfrak p} A^h/\mathfrak p^h$가 정역
$A^h/\mathfrak p^h$의 국소화이기 때문이다. $B/\mathfrak mB$가 두 개의
극대아이디얼을 가지므로 $B \otimes_A \kappa$는 적어도 두 개의
극대아이디얼을 갖는다. 따라서 $A^h$가 헨젤환이므로
$B \otimes_A A^h$는 $\geq 2$개의 국소환의 곱이다. Algebra, Lemma
\ref{algebra-lemma-mop-up}을 보라. 그러나 방금 $B \otimes_A A^h$가
정역의 부분환임을 보았으므로 모순이다.

\medskip\noindent
(1)을 가정하자. $\mathfrak p \subset A$를 유일한 극소소아이디얼이라
하고, $A'$을 그 분수체 안에서 $A/\mathfrak p$의 정수적 폐포라 하자.
$A \to B$를 잉여체들의 동형사상을 유도하며 에탈 $A$-대수의 국소화인
국소환들의 국소사상이라 하자. 특히 $\mathfrak m_B$는
$\mathfrak m B$를 포함하는 유일한 소아이디얼이다. 그러면
$B' = A' \otimes_A B$는 $B$ 위 정수적이고, $A \to A'$이 국소사상이라는
가정에 의해 $B'$은 국소환이다(Algebra, Lemma
\ref{algebra-lemma-local-tensor-with-integral}). 한편 $A' \to B'$은
에탈 환 사상의 국소화이므로 $B'$은 정규이다. Algebra, Lemma
\ref{algebra-lemma-normal-goes-up}을 보라. 따라서 $B'$은 (국소) 정규
정역이다. 마지막으로 다음 포함관계가 있다.
% P02-E0537: frozen display lacks punctuation before the following conclusion.
$$
B/\mathfrak pB \subset B \otimes_A \kappa(\mathfrak p)
= B' \otimes_{A'} (\text{fraction field of }A') \subset
\text{fraction field of }B'
$$
따라서 $B/\mathfrak pB$는 정역이고, 이는 $B$가 유일한
극소소아이디얼을 가짐을 함의한다($A \to B$의 평탄성에 의해 이들은
모두 $\mathfrak p$ 위에 놓여야 하기 때문이다). $A^h$는 국소환
$B$들의 여과 여극한이므로 $A^h$는 유일한 극소소아이디얼을 갖는다.
실제로 어떤 멱영이 아닌 원소 $f, g$에 대하여 $A^h$에서 $fg = 0$이면,
$f$와 $g$를 모두 포함하는 위와 같은 $B$를 찾을 수 있고 이는 모순이다.
\end{proof}

\begin{lemma}
\label{more-algebra-lemma-geometric-branches}
$(A, \mathfrak m, \kappa)$를 국소환이라 하자. $A$가 유한 개의
극소소아이디얼을 갖는다고 가정하자. $A'$을 $A_{red}$의 전체 몫환
안에서 $A$의 정수적 폐포라 하자. $\kappa$의 대수적 폐포
$\overline{\kappa}$를 택하고,
% P02-E1206: frozen notation assignment omits the required “by”.
$\kappa$의 분리가능 대수적 폐포를
$\kappa^{sep} \subset \overline{\kappa}$로 나타내자. $\kappa^{sep}$에
대한 $A$의 엄밀 헨젤화를 $A^{sh}$라 하자. 다음 사상들을 생각하자.
$$
\Spec(A') \xleftarrow{c} \Spec((A')^{sh}) \xrightarrow{e} \Spec(A^{sh})
$$
여기서 $(A')^{sh} = A' \otimes_A A^{sh}$이다. 그러면
\begin{enumerate}
\item 극대아이디얼 $\mathfrak m' \subset A'$에 대하여 잉여체
$\kappa'$은 $\kappa$ 위 대수적이고, $\mathfrak m'$ 위에서 $c$의 섬유는
$\Hom_\kappa(\kappa', \overline{\kappa})$와 정준적으로 동일시할 수 있다.
\item 오른쪽 화살표는 극소소아이디얼 위에서 전단사이다.
\item $(A')^{sh}$의 모든 극소소아이디얼은 유일한 극대아이디얼에
포함되고 모든 극대아이디얼은 유일한 극소소아이디얼을 포함한다.
\end{enumerate}
\end{lemma}

\begin{proof}
증명은 보조정리 \ref{more-algebra-lemma-branches}의 증명과 거의 똑같다.
$I \subset A$를 멱영원들의 아이디얼이라 하자. Algebra, Lemma
\ref{algebra-lemma-quotient-henselization}에 의해
$(A/I)^{sh} = A^{sh}/IA^{sh}$이다.
% P02-E0541: frozen strict-henselization proof lists (A')^h among otherwise strict objects.
$A$, $A^{sh}$, $A'$, $(A')^h$의 스펙트럼은 각각 $A/I$,
$A^{sh}/IA^{sh}$, $A'$, 그리고
$(A')^{sh} = A' \otimes_{A/I} A^{sh}/IA^{sh}$의 스펙트럼과 같다. 따라서
$A$를 $A_{red} = A/I$로 바꾸어 $A$가 기약이라고 가정해도 된다. 그러면
$A \subset A'$이며, 아래에서는 더 언급하지 않고 이를 사용한다.

\medskip\noindent
(1)의 증명. $A'$이 $A$ 위 정수적이므로 체 확대
$\kappa'/\kappa$는 대수적이다. $A'$이 $A$ 위 정수적이므로
$(A')^{sh}$는 $A^{sh}$ 위 정수적이다. Going up(Algebra, Lemma
\ref{algebra-lemma-integral-going-up})에 의해 $A'$, resp.\ $(A')^{sh}$의
모든 극대아이디얼은
% P02-E0542: frozen strict proof switches to ordinary-henselization objects here and in the display.
$A$, resp.\ $A^h$의 극대아이디얼 $\mathfrak m$, resp.\
$\mathfrak m^{sh}$ 위에 놓인다. 다음 등식이 있다.
$$
(A')^{sh} \otimes_{A^{sh}} \kappa^{sep} =
A' \otimes_A A^h \otimes_{A^h} \kappa^{sep} =
(A' \otimes_A \kappa) \otimes_{\kappa} \kappa^{sep}
$$
이는 $A^{sh}$의 잉여체가 $\kappa^{sep}$이기 때문이다. 따라서
$\mathfrak m'$ 위에서 $c$의 섬유는
$\kappa' \otimes_\kappa \kappa^{sep}$의 스펙트럼이다. 다음 전단사가
있으므로 (1)이 성립한다.
% P02-E0537: frozen display lacks punctuation before the following sentence.
$$
\Hom_\kappa(\kappa', \overline{\kappa}) \to
\Spec(\kappa' \otimes_\kappa \kappa^{sep}),\quad
\sigma \mapsto \Ker(
\sigma \otimes 1 : \kappa' \otimes_\kappa \kappa^{sep} \to \overline{\kappa}
)
$$
% P02-E1209: the parallel frozen profinite-spectrum sentence omits its article.
아래에서 이 표시된 환이 체 위 정수적이고, 따라서 그 스펙트럼이
profinite 공간이라는 사실을 사용할 것이다. Algebra, Lemmas
\ref{algebra-lemma-integral-over-field}와
\ref{algebra-lemma-ring-with-only-minimal-primes}를 보라.

\medskip\noindent
(3)의 증명. 환 $A'$은 정규환이며, 실제로 정규 정역들의 유한 곱이다.
Algebra, Lemma \ref{algebra-lemma-characterize-reduced-ring-normal}를 보라.
$A^{sh}$는 에탈 $A$-대수들의 여과 여극한이므로,
% P02-E0535: second frozen omission of the article before “filtered colimit”.
$(A')^{sh}$는 에탈 $A'$-대수들의 여과 여극한이다. 따라서 Algebra,
Lemmas \ref{algebra-lemma-normal-goes-up}와
\ref{algebra-lemma-colimit-normal-ring}에 의해 $(A')^{sh}$는 정규환이다.
그러므로 $(A')^{sh}$의 모든 국소환은 정규 정역이고, 모든
극대아이디얼이 유일한 극소소아이디얼을 포함함을 알 수 있다.
% P02-E0543: frozen paragraph has a malformed application clause, undefined kappa^{sh}, and broken punctuation/connective.
보조정리 \ref{more-algebra-lemma-integral-over-henselian-pair}를
$A^{sh} \to (A')^{sh}$에 적용하면
$((A')^{sh}, \mathfrak m(A')^{sh})$가 헨젤 쌍임을 알 수 있다.
$\mathfrak q \subset (A')^{sh}$가 극소소아이디얼(또는 임의의
소아이디얼)이면, 보조정리
\ref{more-algebra-lemma-irreducible-henselian-pair-connected}에 의해
$V(\mathfrak q)$와 $V(\mathfrak m (A')^{sh})$의 교집합은 연결이다.
$V(\mathfrak m (A')^{sh}) = \Spec((A')^{sh} \otimes \kappa^{sh})$는
profinite 공간이므로 $\mathfrak q$를 포함하는 극대아이디얼은 유일하다.

\medskip\noindent
(2)의 증명. $A'$의 극소소아이디얼들은 정확히 $A$의 극소소아이디얼
위에 놓이는 소아이디얼들이다(구성에 의한다). $A' \to (A')^{sh}$는
평탄하므로 going down(Algebra, Lemma
\ref{algebra-lemma-flat-going-down})에 의해 $(A')^{sh}$의 모든
극소소아이디얼은 $A'$의 극소소아이디얼 위에 놓인다. 반대로
$(A')^{sh}$의 소아이디얼이 $A'$의 극소소아이디얼 위에 놓이면 극소이다.
이는 $(A')^{sh}$가 $A'$ 위 에탈, 따라서 준유한인 대수들의 여과
% P02-E1214: the frozen source explicitly omits the parallel converse fibre detail.
여극한이기 때문이다(작은 세부사항은 생략한다). 따라서 $(A')^{sh}$의
극소소아이디얼들은 정확히 $A$의 극소소아이디얼 위에 놓이는
소아이디얼들이다. 마찬가지로 $A^{sh}$의 극소소아이디얼들은 정확히
$A$의 극소소아이디얼 위에 놓이는 소아이디얼들이다. 구성에 의해
$A' \otimes_A Q(A) = Q(A)$이고, 여기서 $Q(A)$는 기약 국소환 $A$의
전체 몫환이다. 물론 $Q(A)$는 $A$의 극소소아이디얼들의 잉여체들의
유한 곱이다. 따라서
% P02-E0537: frozen parallel display lacks punctuation before the conclusion.
$$
(A')^{sh} \otimes_A Q(A) =
A^{sh} \otimes_A A' \otimes_A Q(A) = A^{sh} \otimes_A Q(A)
$$
이다. 위 논의에 따르면 왼쪽 환의 스펙트럼은 $(A')^{sh}$의
극소소아이디얼들의 집합이고,
% P02-E0538: frozen parallel English sentence duplicates “is the”.
오른쪽 환의 스펙트럼은 $A^{sh}$의 극소소아이디얼들의 집합이다. 이로써
증명이 끝난다.
\end{proof}

\begin{lemma}
\label{more-algebra-lemma-geometrically-unibranch}
\begin{reference}
\cite[Lemma 2.2]{Etale-coverings} and
\cite[Chapter IV Proposition 18.8.15]{EGA4}
\end{reference}
$A$를 국소환이라 하고 $A^{sh}$를 $A$의 엄밀 헨젤화라 하자.
% P02-E1216: frozen English lacks punctuation before the equivalence list.
다음은 서로 동치이다.
\begin{enumerate}
\item $A$는 기하적 단일분지이고,
\item $A^{sh}$는 유일한 극소소아이디얼을 갖는다.
\end{enumerate}
\end{lemma}

\begin{proof}
이는 보조정리 \ref{more-algebra-lemma-geometric-branches}에서 따르지만 직접증명도
제시하겠다. 이 직접증명은 보조정리 \ref{more-algebra-lemma-unibranch}의 직접증명과
거의 똑같다.
% P02-E0539: frozen paired notation assignments omit “by” and use singular “field”.
환 $A$의 극대아이디얼을 $\mathfrak m$으로 나타내자. $A$, $A^{sh}$의
잉여체를 각각 $\kappa$, $\kappa^{sh}$로 나타내자.

\medskip\noindent
(2)를 가정하자. $\mathfrak p^{sh}$를 $A^{sh}$의 유일한
극소소아이디얼이라 하자. $A \to A^{sh}$의 평탄성에 의해
$\mathfrak p = A \cap \mathfrak p^{sh}$는 $A$의 유일한
극소소아이디얼이다(going down에 의한다. Algebra, Lemma
\ref{algebra-lemma-flat-going-down}을 보라). 또한 Algebra, Lemma
\ref{algebra-lemma-quotient-strict-henselization}에서 보듯
$A^{sh}/\mathfrak pA^{sh} = (A/\mathfrak p)^{sh}$이고, 이는 보조정리
\ref{more-algebra-lemma-henselization-reduced}에 의해 기약이므로
$\mathfrak p^{sh} = \mathfrak pA^{sh}$이다. $A'$을 그 분수체 안에서
$A/\mathfrak p$의 정수적 폐포라 하자. $A'$이 국소환이고 그 잉여체가
$\kappa$ 위 순수 비분리임을 보여야 한다. $A \to A'$은 정수적이므로
$A'$의 모든 극대아이디얼은 $\mathfrak m$ 위에 놓인다(정수적 환 사상에
대한 going up에 의한다. Algebra, Lemma
\ref{algebra-lemma-integral-going-up}을 보라). $A'$이 국소환이 아니면
서로 다른 극대아이디얼 $\mathfrak m_1$, $\mathfrak m_2$를 찾을 수 있다.
% P02-E0544: frozen participial clause lacks a conjunction or separating punctuation.
$f_i \in \mathfrak m_i, f_i \not \in \mathfrak m_{3 - i}$인 원소
$f_1, f_2 \in A'$을 택하면, 서로 다른 극대아이디얼
$B \cap \mathfrak m_i$, $i = 1, 2$를 갖는 유한 부분대수
$B = A[f_1, f_2] \subset A'$을 얻는다. $A'$이 극대아이디얼
$\mathfrak m'$을 갖는 국소환이지만
$A/\mathfrak m \subset A'/\mathfrak m'$이 순수 비분리가 아니면,
$A'/\mathfrak m'$ 안에서 상이 $A/\mathfrak m$의 유한이고 순수
비분리가 아닌 확대를 생성하는 $f \in A'$을 찾을 수 있다. 그리고
잉여체가 $A/\mathfrak m$의 순수 비분리 확대가 아닌 유한 국소
부분대수 $B = A[f] \subset A'$을 얻는다. 포함관계
$$
A/\mathfrak p \subset B \subset \kappa(\mathfrak p)
$$
를 평탄한 환 사상 $A \to A^{sh}$로 텐서하면 다음 포함관계를 얻는다.
$$
A^{sh}/\mathfrak p^{sh} \subset
B \otimes_A A^{sh} \subset
\kappa(\mathfrak p) \otimes_A A^{sh} \subset
\kappa(\mathfrak p^{sh})
$$
마지막 포함관계는
$\kappa(\mathfrak p) \otimes_A A^{sh} =
\kappa(\mathfrak p) \otimes_{A/\mathfrak p} A^{sh}/\mathfrak p^{sh}$가
정역 $A^{sh}/\mathfrak p^{sh}$의 국소화이기 때문이다.
$B/\mathfrak mB$가 두 개의 극대아이디얼을 갖거나 하나를 갖되 그
잉여체가 $\kappa$ 위 순수 비분리가 아니고, $\kappa^{sh}$가 분리가능
대수적으로 닫혀 있으므로 $B \otimes_A \kappa^{sh}$는 적어도 두 개의
극대아이디얼을 갖는다. 따라서 $A^{sh}$가 엄밀 헨젤환이므로
$B \otimes_A A^{sh}$는 $\geq 2$개의 국소환의 곱이다. Algebra, Lemma
\ref{algebra-lemma-mop-up-strictly-henselian}을 보라. 그러나 방금
$B \otimes_A A^{sh}$가 정역의 부분환임을 보았으므로 모순이다.

\medskip\noindent
(1)을 가정하자. $\mathfrak p \subset A$를 유일한 극소소아이디얼이라
하고, $A'$을 그 분수체 안에서 $A/\mathfrak p$의 정수적 폐포라 하자.
$A \to B$를 에탈 $A$-대수의 국소화인 국소환들의 국소사상이라 하자.
특히 $\mathfrak m_B$는 $\mathfrak m_AB$를 포함하는 유일한
소아이디얼이다. 그러면 $B' = A' \otimes_A B$는 $B$ 위 정수적이고,
$A \to A'$이 순수 비분리 잉여체 확대를 갖는 국소사상이라는 가정에
의해 $B'$은 국소환이다(Algebra, Lemma
\ref{algebra-lemma-local-tensor-with-integral}). 한편 $A' \to B'$은
에탈 환 사상의 국소화이므로 $B'$은 정규이다. Algebra, Lemma
\ref{algebra-lemma-normal-goes-up}을 보라. 따라서 $B'$은 (국소) 정규
정역이다. 마지막으로 다음 포함관계가 있다.
% P02-E0537: frozen display lacks punctuation before the following conclusion.
$$
B/\mathfrak pB \subset B \otimes_A \kappa(\mathfrak p)
= B' \otimes_{A'} (\text{fraction field of }A') \subset
\text{fraction field of }B'
$$
따라서 $B/\mathfrak pB$는 정역이고, 이는 $B$가 유일한
극소소아이디얼을 가짐을 함의한다($A \to B$의 평탄성에 의해 이들은
모두 $\mathfrak p$ 위에 놓여야 하기 때문이다). $A^{sh}$는 국소환
$B$들의 여과 여극한이므로 $A^{sh}$는 유일한 극소소아이디얼을 갖는다.
실제로 어떤 멱영이 아닌 원소 $f, g$에 대하여 $A^{sh}$에서 $fg = 0$이면,
$f$와 $g$를 모두 포함하는 위와 같은 $B$를 찾을 수 있고 이는 모순이다.
\end{proof}

\begin{definition}
\label{more-algebra-definition-number-of-branches}
$A$를 헨젤화 $A^h$와 엄밀 헨젤화 $A^{sh}$를 갖는 국소환이라 하자.
$A^h$의 극소소아이디얼의 수가 유한하면 그 수를, 그렇지 않으면
$\infty$를 {\it $A$의 분지의 수}라 한다. $A^{sh}$의 극소소아이디얼의
수가 유한하면 그 수를, 그렇지 않으면 $\infty$를 {\it $A$의 기하적
분지의 수}라 한다.
\end{definition}

\noindent
정의 \ref{more-algebra-definition-unibranch}와의 관계를 자세히 적어 두자.

\begin{lemma}
\label{more-algebra-lemma-number-of-branches-1}
$(A, \mathfrak m, \kappa)$를 국소환이라 하자.
\begin{enumerate}
\item $A$가 무한히 많은 극소소아이디얼을 가지면 $A$의 (기하적)
분지의 수는 $\infty$이다.
\item $A$의 분지의 수가 $1$일 필요충분조건은 $A$가 단일분지인 것이다.
\item $A$의 기하적 분지의 수가 $1$일 필요충분조건은 $A$가 기하적
단일분지인 것이다.
\end{enumerate}
$A$가 유한 개의 극소소아이디얼을 갖는다고 가정하고, $A'$을
$A_{red}$의 전체 몫환 안에서 $A$의 정수적 폐포라 하자. 그러면
\begin{enumerate}
\item[(4)] $A$의 분지의 수는 $A'$의 극대아이디얼 $\mathfrak m'$의
수이고,
\item[(5)] $A$의 기하적 분지의 수를 얻으려면 $A'$의 각 극대아이디얼
$\mathfrak m'$을 $\kappa(\mathfrak m')/\kappa$의 분리가능 차수를
중복도로 하여 세어야 한다.
\end{enumerate}
\end{lemma}

\begin{proof}
이 보조정리는 정의들, 보조정리 \ref{more-algebra-lemma-branches}, 보조정리
\ref{more-algebra-lemma-geometric-branches}, 그리고 Fields, Lemma
\ref{fields-lemma-separable-degree}에서 바로 따른다.
\end{proof}

\begin{lemma}
\label{more-algebra-lemma-invariance-number-branches-smooth}
$A \to B$를 매끄러운 환 사상의 국소화인 국소환들의 국소 준동형이라
하자.
\begin{enumerate}
\item $A$의 기하적 분지의 수는 $B$의 기하적 분지의 수와 같다.
\item $A \to B$가 잉여체들의 순수 비분리 확대를 유도하면 $A$의
분지의 수는 $B$의 분지의 수와 같다.
\end{enumerate}
\end{lemma}

\begin{proof}
다음 사실들을 사용할 것이다. 매끄러운 환 사상은 평탄하고(Algebra,
Lemma \ref{algebra-lemma-smooth-syntomic}), 국소화는 평탄하며(Algebra,
Lemma \ref{algebra-lemma-flat-localization}), 평탄한 환 사상들의 합성은
평탄하고(Algebra, Lemma \ref{algebra-lemma-composition-flat}), 평탄한 환
사상의 밑변환은 평탄하며(Algebra, Lemma
\ref{algebra-lemma-flat-base-change}), 평탄한 국소 준동형은 충실히
평탄하고(Algebra, Lemma \ref{algebra-lemma-local-flat-ff}), (엄밀)
헨젤화는 평탄하며(보조정리
\ref{more-algebra-lemma-dumb-properties-henselization}),
% P02-E0545: frozen English capitalizes “Going” in the middle of this list.
평탄한 환 사상에는 going down이 성립한다(Algebra, Lemma
\ref{algebra-lemma-flat-going-down}).

\medskip\noindent
(2)의 증명. $A^h$, $B^h$를 $A$, $B$의 헨젤화라 하자. 그러면 $B^h$는
$\mathfrak m_B$ 위에 놓이는 유일한 극대아이디얼에서
$A^h \otimes_A B$의 헨젤화이다. Algebra, Lemma
\ref{algebra-lemma-henselian-functorial-improve}를 보라. 따라서 $A$가
헨젤환이라고 가정해도 되며 그렇게 한다. $A \to B \to B^h$가
평탄하므로 $B^h$의 모든 극소소아이디얼은 $A$의 극소소아이디얼 위에
놓이고, $A \to B^h$가 충실히 평탄하므로 $A$의 모든
극소소아이디얼은 실제로 $B^h$의 어떤 극소소아이디얼 아래에 놓인다.
두 경우 모두 평탄한 환 사상에 대한 going down을 사용한다. 따라서
주어진 극소소아이디얼 $\mathfrak p \subset A$에 대하여, $\mathfrak p$
위에 놓이는 $B^h$의 극소소아이디얼이 많아야 하나임을 보이면 충분하다. $A$를
$A/\mathfrak p$로, $B$를 $B/\mathfrak p B$로 바꾸어 $A$가 정역이라고
가정해도 된다.
% P02-E0546: frozen English contains the spurious article in “the A is still henselian”.
이때도 Algebra, Lemma \ref{algebra-lemma-quotient-henselization}에 의해
$A$는 헨젤환이다. 보조정리 \ref{more-algebra-lemma-unibranch}에 의해 그 분수체
안에서 $A$의 정수적 폐포 $A'$은 국소 정역이다. 물론 $A'$은 정규
정역이다. Algebra, Lemma \ref{algebra-lemma-normal-goes-up}에 의해
$A' \otimes_A B^h$는 정규환이다(그 보조정리는
$A' \otimes_A B$에 대해서만 이를 주므로, $A' \otimes_A B^h$로
올라가기 위해 $B^h$가 에탈 $B$-대수들의 여극한임과 Algebra, Lemma
\ref{algebra-lemma-colimit-normal-ring}을 사용한다). Algebra, Lemma
\ref{algebra-lemma-local-tensor-with-integral}에 의해
$A' \otimes_A B^h$는 국소환이다(여기서 $A$와 $B$의 잉여체들에 관한
가정을 사용한다). 따라서 $A' \otimes_A B^h$는 국소 정규환이고,
따라서 국소 정역이다. $A \to B^h$의 평탄성에 의해
$B^h \subset A' \otimes_A B^h$이므로 원하는 대로 $B^h$가 정역이라고
결론짓는다.

\medskip\noindent
(1)의 증명. $A^{sh}$, $B^{sh}$를 $A$, $B$의 엄밀 헨젤화라 하자.
% P02-E0547: frozen strict proof copies A^h and its maximal ideal from the ordinary-henselization argument.
그러면 $B^{sh}$는 $\mathfrak m_B$와 $\mathfrak m_{A^h}$ 위에 놓이는
극대아이디얼에서 $A^h \otimes_A B$의 엄밀 헨젤화이다. Algebra, Lemma
\ref{algebra-lemma-strictly-henselian-functorial-improve}를 보라. 따라서
$A$가 엄밀 헨젤환이라고 가정해도 되며 그렇게 한다.
$A \to B \to B^{sh}$가 평탄하므로 $B^{sh}$의 모든
극소소아이디얼은 $A$의 극소소아이디얼 위에 놓이고,
$A \to B^{sh}$가 충실히 평탄하므로 $A$의 모든 극소소아이디얼은
실제로 $B^{sh}$의 어떤 극소소아이디얼 아래에 놓인다. 두 경우 모두
평탄한 환 사상에 대한 going down을 사용한다. 따라서 주어진
극소소아이디얼 $\mathfrak p \subset A$에 대하여, $\mathfrak p$ 위에
놓이는 $B^{sh}$의 극소소아이디얼이 많아야 하나임을 보이면 충분하다. $A$를
$A/\mathfrak p$로, $B$를 $B/\mathfrak p B$로 바꾸어 $A$가 정역이라고
가정해도 된다. 이때도 Algebra, Lemma
\ref{algebra-lemma-quotient-strict-henselization}에 의해 $A$는 엄밀
헨젤환이다. 보조정리 \ref{more-algebra-lemma-geometrically-unibranch}에 의해 그
분수체 안에서 $A$의 정수적 폐포 $A'$은 국소 정역이고, 그 잉여체는
$A$의 잉여체의 순수 비분리 확대이다. 물론 $A'$은 정규 정역이다.
Algebra, Lemma \ref{algebra-lemma-normal-goes-up}에 의해
$A' \otimes_A B^{sh}$는 정규환이다(그 보조정리는
$A' \otimes_A B$에 대해서만 이를 주므로, $A' \otimes_A B^{sh}$로
올라가기 위해 $B^{sh}$가 에탈 $B$-대수들의 여극한임과 Algebra, Lemma
\ref{algebra-lemma-colimit-normal-ring}을 사용한다). Algebra, Lemma
\ref{algebra-lemma-local-tensor-with-integral}에 의해
$A' \otimes_A B^{sh}$는 국소환이다($A \subset A'$이 순수 비분리
잉여체 확대를 유도하기 때문이다). 따라서 $A' \otimes_A B^{sh}$는
국소 정규환이고, 따라서 국소 정역이다. $A \to B^{sh}$의 평탄성에
의해 $B^{sh} \subset A' \otimes_A B^{sh}$이므로 원하는 대로
$B^{sh}$가 정역이라고 결론짓는다.
\end{proof}

\section{분지에 관한 잡록}
\label{more-algebra-section-branches}

\noindent
절 \ref{more-algebra-section-unibranch}에서 정의한 국소환의 분지와 관련된 몇 가지
결과를 제시한다.

\begin{lemma}
\label{more-algebra-lemma-ideal-inverse-nonzero}
$A$와 $B$를 정역이라 하고 $A \to B$를 환 사상이라 하자.
% P02-E0548: frozen English lacks the colon introducing the seven properties.
$A \to B$가 다음 성질 중 적어도 하나도 가진다고 가정하자.
\begin{enumerate}
\item 에탈 환 사상의 국소화이다.
\item 평탄하고 비분기 환 사상의 국소화이다.
\item 평탄하고 준유한 환 사상의 국소화이다.
\item 평탄하고 정수적 환 사상의 국소화이다.
\item 평탄하고 $\Spec(B) \to \Spec(A)$의 섬유들의 점들 사이에
자명하지 않은 특수화가 없다.
\item $\Spec(B) \to \Spec(A)$가 일반점을 일반점으로 보내고 섬유들의
점들 사이에 자명하지 않은 특수화가 없다.
\item $\Spec(B)$의 정확히 한 점이 $\Spec(A)$의 일반점으로 보내진다.
\end{enumerate}
그러면 $B$의 모든 영이 아닌 아이디얼 $J$에 대하여 $A \cap J$는 영이
아니다.
\end{lemma}

\begin{proof}
(7)의 경우의 증명. $K$, resp.\ $L$을 $A$, resp.\ $B$의 분수체라 하자.
Algebra, Lemma
\ref{algebra-lemma-minimal-prime-image-minimal-prime}에 의해
$\Spec(B)$에서 일반점 $(0) \in \Spec(A)$으로 보내지는
유일한 점은 $(0) \in \Spec(B)$임을 알 수 있다. 따라서
$B \otimes_A K$는 유일한 소아이디얼을 가지며 그 잉여체가 $L$인
환이다(실제로 이는 $L$과 같지만 필요하지는 않다). 영이 아닌
$b \in J$를 택하자. 그러면 $b$는 $L$의 단원으로 보내진다. 따라서
$b$는 $B \otimes_A K$의 단원으로 보내진다(Algebra, Lemma
\ref{algebra-lemma-surjective-on-spec-units}).
$B \otimes_A K = \colim_{f \in A \setminus \{0\}} B_f$이므로 어떤 영이
아닌 $f \in A$에 대하여 $b$가 $B_f$의 단원으로 보내짐을 알 수 있다.
이는 어떤 $b' \in B$와 $n \geq 1$에 대하여 $b b' = f^n$임을 뜻한다.
따라서 원하는 대로 $f^n \in A \cap J$이다.

\medskip\noindent
% P02-E0549: frozen English uses plural “imply” with singular “each”.
증명의 나머지에서는 다른 각 가정이 (7)을 함의함을 보인다. 가정
(1) -- (5) 아래에서 환 사상 $A \to B$는 평탄하고,
% P02-E1220: frozen proof invokes faithful flatness for a map not assumed local.
따라서 $A \to B$는 단사이다(Algebra, Lemma
\ref{algebra-lemma-local-flat-ff}에 의해 평탄한 국소 준동형은 충실히
평탄하기 때문이다). 그러므로 $\Spec(B)$의 일반점은 $\Spec(A)$의
일반점으로 보내진다. 이제 $\Spec(B) \to \Spec(A)$의 섬유들의 점들
사이에 자명하지 않은 특수화가 없다면, 물론 $\Spec(B)$의 이 일반점은
$\Spec(A)$의 일반점으로 보내지는 유일한 점이어야 한다. 따라서 (6)은
(7)을 함의한다. 마지막으로 증명을 끝내기 위해 (1) -- (5)의 경우에
$\Spec(B) \to \Spec(A)$의 섬유들의 점들 사이에 자명하지 않은
특수화가 없음을 보인다. 즉 정수적 경우에는 Algebra, Lemma
\ref{algebra-lemma-integral-no-inclusion}를, 준유한 경우에는 Algebra,
Definition \ref{algebra-definition-quasi-finite}를 보라. 그리고 비분기
환 사상과 에탈 환 사상이 준유한이라는 사실을 사용하라(Algebra,
Lemmas \ref{algebra-lemma-unramified-quasi-finite}와
\ref{algebra-lemma-etale-quasi-finite}).
\end{proof}

\begin{lemma}
\label{more-algebra-lemma-local-unramified-extension-unibranch-domain-is-etale}
$A \to B$를 환 사상이라 하자. $\mathfrak q \subset B$를
$\mathfrak p \subset A$ 위에 놓이는 소아이디얼이라 하자. 다음을
가정하자.
\begin{enumerate}
\item $A$는 정역이다.
\item $A_\mathfrak p$는 기하적 단일분지이다.
\item $A \to B$는 $\mathfrak q$에서 비분기이다.
\item $A_\mathfrak p \to B_\mathfrak q$는 단사이다.
\end{enumerate}
그러면 $g \not \in \mathfrak q$인 $g \in B$가 존재하여 $B_g$는 $A$ 위
에탈이다.
\end{lemma}

\begin{proof}
Algebra, Proposition \ref{algebra-proposition-unramified-locally-standard}에
의해 $B$를 주 국소화로 바꾸면 표준 에탈 환 사상 $A \to B'$과 전사
$B' \to B$를 찾을 수 있다.
% P02-E0550: frozen notation assignment omits “by”.
$\mathfrak q$의 역상을 $\mathfrak q' \subset B'$으로 나타내자. 필요하면
$B'$을 주 국소화로 바꾼 뒤 $B' \to B$가 단사임을 보이겠다.

\medskip\noindent
이 문단에서는 $B'$이 정역인 경우로 환원한다. $A$가 정역이므로 환
$B'$은 기약이다. Algebra, Lemma \ref{more-algebra-lemma-reduced-goes-up}을 보라.
$K$를 $A$의 분수체라 하자. 그러면 $B' \otimes_A K$는 체 위 에탈이고,
따라서 체들의 유한 곱이다. Algebra, Lemma
\ref{algebra-lemma-etale-over-field}을 보라. $A \to B'$은 에탈이고
(따라서 평탄하고),
% P02-E0551: frozen English combines “are” and “lie” in one predicate.
$B'$의 극소소아이디얼들은 $(0) \subset A$ 위에 놓인다(평탄한 환
사상에 대한 going down에 의한다). 따라서 $B'$은 유한 개의
극소소아이디얼 $\mathfrak r_1, \ldots, \mathfrak r_r \subset B'$을
갖는다. $A_\mathfrak p$는 기하적 단일분지이고 $A \to B'$은
에탈이므로 보조정리들
\ref{more-algebra-lemma-invariance-number-branches-smooth}와
\ref{more-algebra-lemma-number-of-branches-1}에 의해 환 $B'_{\mathfrak q'}$은
정역이다. 따라서 정확히 하나의 $i = i_0$에 대하여
$\mathfrak q' \supset \mathfrak r_i$이다. $i \not = i_0$에 대해서는
$g' \in \mathfrak r_i$이고 $g' \not \in \mathfrak r_{i_0}$인
$g' \in B'$을 택하자. Algebra, Lemma \ref{algebra-lemma-silly}를 보라.
$B'$과 $B$를 $B'_{g'}$와 $B_{g'}$로 바꾸면 $B'$이 정역이 된다.

\medskip\noindent
$B'$이 정역이라고 가정하자. 특히 $B' \subset B'_{\mathfrak q'}$이다.
$B' \to B$가 단사가 아니면
$J = \Ker(B'_{\mathfrak q'} \to B_\mathfrak q)$는 영이 아니다. 보조정리
\ref{more-algebra-lemma-ideal-inverse-nonzero}를
$A_\mathfrak p \to B'_{\mathfrak q'}$에 적용하면 $B_\mathfrak q$에서
영으로 보내지는 영이 아닌 원소 $a \in A_\mathfrak p$를 얻는데, 이는
가정 (4)에 모순이다. 이로써 증명이 끝난다.
\end{proof}

\begin{lemma}
\label{more-algebra-lemma-unramified-extension-unibranch-domain-is-etale}
\begin{reference}
Generalization of \cite[Expose I, Theorem 9.5 part (ii)]{SGA1}
\end{reference}
$(A, \mathfrak m)$을 기하적 단일분지 국소 정역이라 하자.
$A \to B$를 본질적으로 유한형인 국소환들의 단사 국소 준동형이라 하자.
$\mathfrak m B$가 $B$의 극대아이디얼이고 유도된 잉여체 확대가
분리가능이면, $A \to B$는 에탈 환 사상의 국소화이다.
\end{lemma}

\begin{proof}
$A \to C$가 유한형 환 사상이고 $\mathfrak q \subset C$가
$\mathfrak m$ 위에 놓이는 소아이디얼이 되도록
$B = C_\mathfrak q$로 쓸 수 있다. Algebra, Lemma
\ref{algebra-lemma-characterize-unramified}에 의해 환 사상 $A \to C$는
$\mathfrak q$에서 비분기이다. Algebra, Proposition
\ref{algebra-proposition-unramified-locally-standard}에 의해 $C$를 주
국소화로 바꾸면 표준 에탈 환 사상 $A \to C'$과 전사 $C' \to C$를
찾을 수 있다.
% P02-E1222: second frozen inverse-image notation assignment omits “by”.
$\mathfrak q$의 역상을 $\mathfrak q' \subset C'$으로 나타내고
$B' = C'_{\mathfrak q'}$로 두자. 그러면 $B' \to B$는 전사이다.
$B' \to B$도 단사임을 보이면 충분하다.

\medskip\noindent
$A$가 정역이므로 환들 $C'$과 $B'$은 기약이다. Algebra, Lemma
\ref{more-algebra-lemma-reduced-goes-up}을 보라. $A$가 기하적 단일분지이므로 환
$B'$은 정역이다.
% P02-E0552: frozen citation clause redundantly combines “see” and “by”.
보조정리들 \ref{more-algebra-lemma-invariance-number-branches-smooth}와
\ref{more-algebra-lemma-number-of-branches-1}을 보라. $B' \to B$가 단사가 아니면
보조정리 \ref{more-algebra-lemma-ideal-inverse-nonzero}에 의해
$A \cap \Ker(B' \to B)$는 영이 아닌데, 이는 $A \to B$가 단사라는
가정에 모순이다.
\end{proof}

\begin{lemma}
\label{more-algebra-lemma-minimal-primes-tensor-strictly-henselian}
$k$를 대수적으로 닫힌 체라 하자. $A$, $B$를 잉여체가 $k$와 같은 엄밀
헨젤 국소 $k$-대수라 하자. $C$를 극대아이디얼
$\mathfrak m_A \otimes_k B + A \otimes_k \mathfrak m_B$에서
$A \otimes_k B$의 엄밀 헨젤화라 하자.
% P02-E0553: frozen statement omits the construction needed for one-to-one correspondence.
그러면 $C$의 극소소아이디얼들은 $A$와 $B$의 극소소아이디얼들의
순서쌍들과 $1$-대-$1$로 대응한다.
\end{lemma}

\begin{proof}
먼저 $C$의 극소소아이디얼 $\mathfrak r$은 $A$의 극소소아이디얼
$\mathfrak p$로, $B$의 극소소아이디얼 $\mathfrak q$로 보내짐을
관찰하자. 이는 환 사상들 $A \to C$와 $B \to C$가 평탄하기 때문이다.
% P02-E0554: frozen parenthetical uses singular “map” for two maps and lacks citation punctuation.
(평탄한 환 사상에 대한 going down, Algebra, Lemma
\ref{algebra-lemma-flat-going-down}에 의한다.) 따라서
% P02-E0555: frozen quotient tensor factors lack disambiguating parentheses.
$(A/\mathfrak p \otimes_k B/\mathfrak q)_{
\mathfrak m_A \otimes_k B + A \otimes_k \mathfrak m_B}$의 엄밀 헨젤화가
유일한 극소소아이디얼을 가짐을 보이면 충분하다. Algebra, Lemma
\ref{algebra-lemma-quotient-strict-henselization}에 의해 환들
$A/\mathfrak p$, $B/\mathfrak q$는 엄밀 헨젤환이다. 따라서 $A$와
$B$가 엄밀 헨젤 국소 정역이라고 가정해도 되고, 우리의 목표는 $C$가
유일한 극소소아이디얼을 가짐을 보이는 것이다. 보조정리
\ref{more-algebra-lemma-geometrically-unibranch}에 의해 그 분수체 안에서 $A$의
정수적 폐포 $A'$은 잉여체가 $k$인 정규 국소 정역이다. 마찬가지로
% P02-E0556: frozen parallel sentence uses “into” rather than “in” the fraction field.
그 분수체 안에서 $B$의 정수적 폐포 $B'$도 그러하다. Algebra, Lemma
\ref{algebra-lemma-geometrically-normal-tensor-normal}에 의해
$A' \otimes_k B'$은 정규환이다. 따라서 그 국소화
% P02-E0558: first of three frozen display boundaries has broken punctuation.
$$
R = (A' \otimes_k B')_{
\mathfrak m_{A'} \otimes_k B' + A' \otimes_k \mathfrak m_{B'}}
$$
는 정규 국소 정역이다.
% P02-E0557: frozen parenthetical misspells “going up” as “gong up”.
$A \otimes_k B \to A' \otimes_k B'$은 정수적임을 관찰하자(따라서
going up이 성립한다---Algebra, Lemma
\ref{algebra-lemma-integral-going-up}). 그리고
$\mathfrak m_{A'} \otimes_k B' + A' \otimes_k \mathfrak m_{B'}$은
$\mathfrak m_A \otimes_k B + A \otimes_k \mathfrak m_B$ 위에 놓이는
$A' \otimes_k B'$의 유일한 극대아이디얼이다. 따라서
% P02-E0558: second frozen display is followed by a detached “by”.
$$
R = (A' \otimes_k B')_{
\mathfrak m_A \otimes_k B + A \otimes_k \mathfrak m_B}
$$
임을 알 수 있다. 이는 Algebra, Lemma
\ref{algebra-lemma-unique-prime-over-localize-below}에서 따른다. 따라서
% P02-E0558: third frozen display lacks punctuation before the conclusion.
$$
(A \otimes_k B)_{
\mathfrak m_A \otimes_k B + A \otimes_k \mathfrak m_B}
\longrightarrow
R
$$
은 정수적이다. 그러므로 $R$은
$(A \otimes_k B)_{
\mathfrak m_A \otimes_k B + A \otimes_k \mathfrak m_B}$의 분수체 안에서
그 정수적 폐포라고 결론짓는다. 그리고 다시 보조정리
\ref{more-algebra-lemma-geometrically-unibranch}에 의해
% P02-E0559: frozen conclusion says unique prime rather than unique minimal prime.
$C$가 유일한 소아이디얼을 갖는다고 결론짓는다.
\end{proof}

\section{완비화의 분지}
\label{more-algebra-section-branches-completion}

\noindent
$(A, \mathfrak m)$을 뇌터 국소환이라 하자. 사상들
$A \to A^h \to A^\wedge$를 생각하자. 일반적으로 사상
$A^h \to A^\wedge$는 극소소아이디얼 위의 전단사를 유도할 필요가
없다. Examples, Section \ref{examples-section-bad}를 보라. 다시 말해,
정의 \ref{more-algebra-definition-number-of-branches}에서 정의한 $A$의 분지의 수는
$A^\wedge$의 분지의 수와 다를 수 있다. 그러나 몇몇 조건 아래에서는
분지의 수가 같다. 예를 들어 $A$의 차원이 $1$이면 그러하다.

\begin{lemma}
\label{more-algebra-lemma-nr-branches-completion}
$(A, \mathfrak m)$을 뇌터 국소환이라 하자.
\begin{enumerate}
\item 사상 $A^h \to A^\wedge$는 $A^\wedge$의 극소소아이디얼들에서
$A^h$의 극소소아이디얼들로 가는 전사 사상을 정의한다.
\item $A$의 분지의 수는 $A^\wedge$의 분지의 수 이하이다.
\item $A$의 기하적 분지의 수는 $A^\wedge$의 기하적 분지의 수 이하이다.
\end{enumerate}
\end{lemma}

\begin{proof}
보조정리 \ref{more-algebra-lemma-henselization-noetherian}에 의해 사상
$A^h \to A^\wedge$는 평탄하고 단사이다. Going down(Algebra, Lemma
\ref{algebra-lemma-flat-going-down})과 Algebra, Lemma
\ref{algebra-lemma-injective-minimal-primes-in-image}를 결합하면 (1)이
성립함을 알 수 있다. (2)는 이 사실, 정의
\ref{more-algebra-definition-number-of-branches}, 그리고 $A^\wedge$가 헨젤환이라는
사실에서 따른다(Algebra, Lemma
\ref{algebra-lemma-complete-henselian}). 보조정리
\ref{more-algebra-lemma-henselization-noetherian}에 의해
$(A^\wedge)^{sh} = A^{sh} \otimes_{A^h} A^\wedge$이다. 따라서 평탄 단사
사상 $A^{sh} \to (A^\wedge)^{sh}$를 사용하여 위 논증을 반복하면 (3)을
증명할 수 있다.
\end{proof}

\begin{lemma}
\label{more-algebra-lemma-equal-nr-branches-completion}
$(A, \mathfrak m)$을 뇌터 국소환이라 하자. 완비화의 모든
극소소아이디얼 $\mathfrak q \subset A^h$에 대하여
$\sqrt{\mathfrak qA^\wedge}$가 소아이디얼일 필요충분조건은 $A$의
분지의 수와 $A^\wedge$의 분지의 수가 같은 것이다.
\end{lemma}

\begin{proof}
% P02-E0560: frozen proof begins with the sentence fragment “Follows from”.
이는 보조정리 \ref{more-algebra-lemma-nr-branches-completion}, Algebra, Lemma
\ref{algebra-lemma-Noetherian-irreducible-components}에 의해 $A$와
$A^\wedge$ 모두 분지의 수가 유한하다는 사실, 그리고
$A^h$와 $A^\wedge$가 뇌터라는 사실에서 따른다(보조정리
\ref{more-algebra-lemma-henselization-noetherian}).
\end{proof}

\noindent
간단한 접합 보조정리이다.

\begin{lemma}
\label{more-algebra-lemma-glueing-sum-components-open}
$A$를 환이라 하고 $I$를 유한 생성 아이디얼이라 하자. $A \to C$를
모든 $f \in I$에 대하여 환 사상 $A_f \to C_f$가 멱등원에서의
국소화가 되는 환 사상이라 하자. 그러면 전사 $A \to C'$이 존재하여
모든 $f \in I$에 대하여 $A_f \to (C \times C')_f$는 동형사상이다.
\end{lemma}

\begin{proof}
$I$의 생성원들 $f_1, \ldots, f_r$을 택하자. 어떤 멱등원
$e_i \in A_{f_i}$에 대하여
$$
C_{f_i} = (A_{f_i})_{e_i}
$$
로 쓰자. 어떤 $a_i \in A$와 $n \geq 0$에 대하여
$e_i = a_i/f_i^n$로 쓰자. 모든 $i = 1, \ldots, r$에 같은 $n$을 사용할
수 있다. 적당한 $m \gg 0$에 대하여 $a_i$를 $f_i^ma_i$로, $n$을
$n + m$으로 바꾸면, 모든 $i$에 대하여 $a_i^2 = f_i^n a_i$라고
가정해도 된다. $e_i$가
$C_{f_if_j} = (A_{f_if_j})_{e_j} = A_{f_if_ja_j}$에서 $1$로 보내지므로,
어떤 $N$에 대하여
$$
(f_if_ja_j)^N(f_j^n a_i  - f_i^na_j) = 0
$$
이다(모든 순서쌍 $i, j$에 대하여 같은 $N$을 택할 수 있다).
$a_j^2 = f_j^na_j$을 사용하면 다음을 얻는다.
% P02-E0561: frozen display lacks terminal punctuation before the next sentence.
$$
f_i^{N + n} f_j^{N + nN} a_j = f_i^N f_j^{N + n} a_ia_j^N
$$
$n$을 $n + N + nN$으로 늘리고 $a_i$를 $f_i^{N + nN}a_i$로 바꾸면,
모든 순서쌍 $i, j$에 대하여 $f_i^n a_j$가 $a_i$의 아이디얼에 속함을
알 수 있다. $C' = A/(a_1, \ldots, a_r)$로 두자. 그러면
$$
C'_{f_i} = A_{f_i}/(a_i) = A_{f_i}/(e_i)
$$
이다. 이는 $f_i$를 가역화하면 $a_j$가 $a_i$로 생성된 아이디얼에
속하기 때문이다. 환 $B$의 멱등원 $e$에 대하여
$B = B_e \times B/(e)$이므로, $f$가 $f_1, \ldots, f_r$ 중 하나인 경우
보조정리의 결론이 성립함을 알 수 있다. Algebra, Lemma
\ref{algebra-lemma-cover}의 함수 접합을 사용하면 모든 $f \in I$에
대하여 결과가 성립한다고 결론짓는다. 실제로 $f \in I$에 대하여
원소들 $f_1, \ldots, f_r$은 $A_f$의 단위아이디얼을 생성하므로,
$A_f \to (C \times C')_f$가 동형사상일 필요충분조건은
$f_1, \ldots, f_r$에서 국소화한 뒤에도 그러한 것이다.
\end{proof}

\noindent
보조정리 \ref{more-algebra-lemma-quotient-by-idempotent}를 사용하면 형식적 완비화의
주어진 유한형 확대에서 유한형 확대를 구성할 수 있다. Algebraization
of Formal Spaces, Lemma
\ref{restricted-lemma-approximate-by-etale-over-complement}에서 이
보조정리를 일반화할 것이다.

\begin{lemma}
\label{more-algebra-lemma-quotient-by-idempotent}
$A$를 뇌터환, $I$를 아이디얼이라 하자. $B$를 유한형 $A$-대수라 하자.
$B^\wedge$가 $I$-진 완비화일 때, 핵이 $J$인 전사 환 사상
$B^\wedge \to C$를 생각하자. 어떤 $c \geq 0$에 대하여 $J/J^2$가
$I^c$에 의해 소멸되면, $C$는 유한형 $A$-대수의 완비화와 동형이다.
\end{lemma}

\begin{proof}
$f \in I$라 하자. $B^\wedge$는 뇌터이므로(Algebra, Lemma
\ref{algebra-lemma-completion-Noetherian-Noetherian}), $J$는 유한 생성
아이디얼이다. 따라서 Algebra, Lemma
\ref{algebra-lemma-ideal-is-squared-union-connected}에서 어떤 멱등원
$e \in (B^\wedge)_f$에 대하여
$$
C_f = ((B^\wedge)_f)_e
$$
라고 결론짓는다. 보조정리 \ref{more-algebra-lemma-glueing-sum-components-open}에 의해
전사 $B^\wedge \to C'$을 찾을 수 있고, 임의의 $f \in I$를 가역화하면
$B^\wedge \to C \times C'$은 동형사상이 된다. $C \times C'$은 유한
$B^\wedge$-대수임을 관찰하자.

\medskip\noindent
$f_1, \ldots, f_r \in I$를 생성원들로 택하자.
% P02-E0562: both frozen notation assignments omit “by”.
위에서 얻은 $(B^\wedge)_{f_i}$-대수의 동형사상의 역을
$\alpha_i : (C \times C')_{f_i} \to B_{f_i} \otimes_B B^\wedge$로
나타내자. 명백한 $B$-대수 동형사상을
$\alpha_{ij} : (B_{f_i})_{f_j} \to (B_{f_j})_{f_i}$로 나타내자. Remark
\ref{more-algebra-remark-glueing-data}에서 도입한 범주
$\text{Glue}(B \to B^\wedge, f_1, \ldots, f_r)$의 대상
$$
(C \times C', B_{f_i}, \alpha_i, \alpha_{ij})
$$
을 생각하자.
% P02-E1229: frozen proof explicitly omits verification of the gluing conditions.
조건 (1)(a)와 (1)(b)의 확인은 생략한다. $B \to B^\wedge$는 평탄하고
(Algebra, Lemma \ref{algebra-lemma-completion-flat}) 동형사상
$B/IB \to B^\wedge/IB^\wedge$를 유도하므로, Proposition
\ref{more-algebra-proposition-equivalence}와 Remark
\ref{more-algebra-remark-formal-glueing-algebras}를 적용할 수 있다. 따라서 어떤 유한
$B$-대수 $D$에 대하여 $C \times C'$은 $D \otimes_B B^\wedge$와
동형이라고 결론짓는다. 그러면
$D/ID \cong C/IC \times C'/IC'$이다. $\overline{e} \in D/ID$를 인자
$C/IC$에 대응하는 멱등원이라 하자. 보조정리
\ref{more-algebra-lemma-lift-idempotent-upstairs}에 의해, 동형사상
$B/IB \to B'/IB'$을 유도하는 에탈 환 사상 $B \to B'$이 존재하여
$D' = D \otimes_B B'$은 $\overline{e}$를 올리는 멱등원 $e$를 포함한다.
$C \times C'$은 $I$-진 완비이므로 쌍
$(C \times C', IC \times IC')$은 헨젤 쌍이다(보조정리
\ref{more-algebra-lemma-complete-henselian}). 따라서 사상 $B \to C \times C'$을
$B'$을 통하여 분해할 수 있다(보조정리
\ref{more-algebra-lemma-characterize-henselian-pair}를 보라). 그렇게 한 뒤 $B$를
$B'$으로, $D$를 $D'$으로 바꾸어도 된다.
% P02-E0563: frozen parenthetical misspells “completions”.
(이렇게 해도 $I$-진 완비화들은 바뀌지 않기 때문이다.) 그러면
$D = D_e \times D_{1 - e} = D/(1 - e) \times D/(e)$가 유한형
$A$-대수들의 곱이고, 첫째 부분의 완비화는 $C$이며 둘째 부분의
완비화는 $C'$임을 알 수 있다.
\end{proof}

\begin{lemma}
\label{more-algebra-lemma-one-dimensional-formal-branch}
$(A, \mathfrak m)$을 헨젤화가 $A^h$인 뇌터 국소환이라 하자.
$\mathfrak q \subset A^\wedge$를
$\dim(A^\wedge/\mathfrak q) = 1$인 극소소아이디얼이라 하자. 그러면
$\mathfrak q = \sqrt{\mathfrak q^hA^\wedge}$가 되는 $A^h$의
극소소아이디얼 $\mathfrak q^h$가 존재한다.
\end{lemma}

\begin{proof}
$A$와 $A^h$의 완비화가 같으므로 $A$가 헨젤환이라고 가정해도 된다
(보조정리 \ref{more-algebra-lemma-henselization-noetherian}).
$J = \Ker(A^\wedge \to (A^\wedge)_{\mathfrak q})$일 때
$A^\wedge \to A^\wedge/J$에 보조정리
\ref{more-algebra-lemma-quotient-by-idempotent}를 적용하겠다.
$\dim((A^\wedge)_\mathfrak q) = 0$이므로 어떤 $n$에 대하여
$\mathfrak q^n \subset J$이다. 따라서 $J/J^2$는 $\mathfrak q^n$에 의해
소멸된다. 한편 $J_\mathfrak q = 0$이므로
$(J/J^2)_\mathfrak q = 0$이다. 따라서 $\mathfrak m$은 $J/J^2$의 유일한
연관소아이디얼이고, $\mathfrak m$의 어떤 거듭제곱이 $J/J^2$를
소멸함을 알 수 있다. 그러므로 보조정리가 적용되어 어떤 유한형
$A$-대수 $C$에 대하여 $A^\wedge/J = C^\wedge$임을 얻는다.

\medskip\noindent
$A^\wedge/J$가 같은 성질을 가지므로
$C/\mathfrak m C = A/\mathfrak m$이다. 따라서
$\mathfrak m_C = \mathfrak m C$는 극대아이디얼이고, $A \to C$는
$\mathfrak m_C$에서 비분기이다(Algebra, Lemma
\ref{algebra-lemma-characterize-unramified}). $C$를 주 국소화로 바꾸어
$C$가 에탈 $A$-대수 $B$의 몫이라고 가정해도 된다. Algebra,
Proposition \ref{algebra-proposition-unramified-locally-standard}를 보라.
그러나 $A \to C_{\mathfrak m_C}$의 잉여체 확대는 자명하고 $A$는
헨젤환이므로, 다시 국소화한 뒤 $B = A$라고 결론짓는다. 따라서 어떤
아이디얼 $I \subset A$에 대하여 $C = A/I$이고, 이로부터
$J = IA^\wedge$가 따른다(우리 상황에서 완비화가 완전하기 때문이다.
Algebra, Lemma \ref{algebra-lemma-completion-flat}을 보라). 또한
$I = J \cap A$이다($A \to A^\wedge$의 평탄성에 의한다).
$\mathfrak q^n \subset J \subset \mathfrak q$이므로
$\mathfrak p = \mathfrak q \cap A$는
$\mathfrak p^n \subset I \subset \mathfrak p$를 만족한다. 그러면
$\sqrt{\mathfrak p A^\wedge} = \mathfrak q$이고 증명이 끝난다.
\end{proof}

\begin{lemma}
\label{more-algebra-lemma-completion-disconnected}
$(A, \mathfrak m)$을 뇌터 국소환이라 하자. $A^\wedge$의 천공
스펙트럼이 비연결일 필요충분조건은 $A^h$의 천공 스펙트럼이
비연결인 것이다.
\end{lemma}

\begin{proof}
$A$와 $A^h$의 완비화가 같으므로 $A$가 헨젤환이라고 가정해도 된다
(보조정리 \ref{more-algebra-lemma-henselization-noetherian}).

\medskip\noindent
$A \to A^\wedge$는 충실히 평탄하므로(방금 제시한 참고문헌을 보라),
$A^\wedge$의 천공 스펙트럼에서 $A$의 천공 스펙트럼으로 가는 사상은
전사이다(Algebra, Lemma \ref{algebra-lemma-ff-rings}를 보라). 따라서
$A$의 천공 스펙트럼이 비연결이면 $A^\wedge$에 대해서도 그러하다.

\medskip\noindent
$A^\wedge$의 천공 스펙트럼이 비연결이라고 가정하자. 이는
% P02-E0564: frozen disjoint-union display lacks punctuation before the following with-clause.
$$
\Spec(A^\wedge) \setminus \{\mathfrak m^\wedge\} = Z \amalg Z'
$$
이고 $Z$와 $Z'$이 닫혀 있음을 뜻한다.
% P02-E0565: frozen closure notation leaves the prime outside overline at both occurrences.
$\overline{Z}, \overline{Z}' \subset \Spec(A^\wedge)$를 폐포들이라 하자.
어떤 아이디얼들 $J, J' \subset A^\wedge$에 대하여
$\overline{Z} = V(J)$, $\overline{Z}' = V(J')$라고 하자. 그러면
$V(J + J') = \{\mathfrak m^\wedge\}$이고
$V(JJ') = \Spec(A^\wedge)$이다. 첫째 등식은
$\mathfrak m^\wedge = \sqrt{J + J'}$임을 뜻하고, 이는 어떤
$e \geq 1$에 대하여 $(\mathfrak m^\wedge)^e \subset J + J'$임을
함의한다. 둘째 등식은 $JJ'$의 모든 원소가 멱영임을 함의하고, 따라서
어떤 $n \geq 1$에 대하여 $(JJ')^n = 0$이다.
% P02-E0566: frozen transition “Combined this means” is ungrammatical.
이들을 결합하면 $J^n/J^{2n}$은 $J^n$과 $(J')^n$에 의해, 따라서
$(\mathfrak m^\wedge)^{2en}$에 의해 소멸된다.
% P02-E0567: frozen proof later reuses the original exponent e after replacing J and J' by nth powers.
그러므로 보조정리 \ref{more-algebra-lemma-quotient-by-idempotent}를 적용하여 유한형
$A$-대수 $C$와 동형사상 $A^\wedge/J^n = C^\wedge$가 존재함을 알 수
있다.

\medskip\noindent
증명의 나머지는 보조정리
\ref{more-algebra-lemma-one-dimensional-formal-branch}의 증명의 둘째 부분과 정확히
같다. 물론 그 보조정리는 이 보조정리의 특수한 경우이다!
$A^\wedge/J^n$이 같은 성질을 가지므로
$C/\mathfrak m C = A/\mathfrak m$이다. 따라서
$\mathfrak m_C = \mathfrak m C$는 극대아이디얼이고, $A \to C$는
$\mathfrak m_C$에서 비분기이다(Algebra, Lemma
\ref{algebra-lemma-characterize-unramified}). $C$를 주 국소화로 바꾸어
$C$가 에탈 $A$-대수 $B$의 몫이라고 가정해도 된다. Algebra,
Proposition \ref{algebra-proposition-unramified-locally-standard}를 보라.
그러나 $A \to C_{\mathfrak m_C}$의 잉여체 확대는 자명하고 $A$는
헨젤환이므로, 다시 국소화한 뒤 $B = A$라고 결론짓는다. 따라서 어떤
아이디얼 $I \subset A$에 대하여 $C = A/I$이고, 이로부터
$J^n = IA^\wedge$가 따른다(우리 상황에서 완비화가 완전하기 때문이다.
Algebra, Lemma \ref{algebra-lemma-completion-flat}을 보라). 또한
$I = J^n \cap A$이다($A \to A^\wedge$의 평탄성에 의한다). 대칭성에
의해 $I' = (J')^n \cap A$는 $(J')^n = I'A^\wedge$를 만족한다.
% P02-E0567: frozen conclusion retains exponent e after the nth-power replacement.
그러면 $\mathfrak m^e \subset I + I'$이고 $II' = 0$이다. 따라서
$V(I)$와 $V(I')$은 $A$의 천공 스펙트럼의 원하는 서로소 합 분해를
주는 닫힌 부분스킴들이라고 결론짓는다.
\end{proof}

\begin{lemma}
\label{more-algebra-lemma-one-dimensional-number-of-branches}
$(A, \mathfrak m)$을 차원이 $1$인 뇌터 국소환이라 하자. 그러면 $A$와
$A^\wedge$의 (기하적) 분지의 수는 같다.
\end{lemma}

\begin{proof}
분지의 수에 대하여 이를 보려면 보조정리들
\ref{more-algebra-lemma-nr-branches-completion},
\ref{more-algebra-lemma-equal-nr-branches-completion},
\ref{more-algebra-lemma-one-dimensional-formal-branch}를 결합하고 $A^\wedge$의 차원이
일이라는 사실을 사용하라. 보조정리 \ref{more-algebra-lemma-completion-dimension}을
보라. 기하적 분지의 수에 대해서도 참임을 보이기 위해, 분지에 대한
결과, 엄밀 헨젤화에서 차원이 바뀌지 않는다는 사실(보조정리
\ref{more-algebra-lemma-henselization-dimension}), 그리고 보조정리
\ref{more-algebra-lemma-henselization-noetherian}에 의한 등식
$(A^{sh})^\wedge = ((A^\wedge)^{sh})^\wedge$를 사용한다.
\end{proof}

\begin{lemma}
\label{more-algebra-lemma-geometrically-normal-formal-fibres-number-of-branches}
\begin{reference}
\cite[Theorem 2.3]{Beddani}
\end{reference}
$(A, \mathfrak m)$을 뇌터 국소환이라 하자. $A$의 형식적 섬유들이
기하적으로 정규이면(예를 들어 $A$가 훌륭하거나 준훌륭하면), $A$는
나가타환이고 $A$와 $A^\wedge$의 (기하적) 분지의 수는 같다.
\end{lemma}

\begin{proof}
정규환은 기약이므로 보조정리 \ref{more-algebra-lemma-Nagata-local-ring}에 의해
$A$는 나가타환임을 알 수 있다. 증명의 나머지에서는 보조정리
\ref{more-algebra-lemma-formal-fibres-normal}, Proposition
\ref{more-algebra-proposition-finite-type-over-P-ring}, 보조정리
\ref{more-algebra-lemma-check-P-ring-maximal-ideals}를 사용할 것이다. 이들에 따르면
$P(k \to R) = $``$R$가 $k$ 위 기하적으로 정규이다''로 둘 때 $A$는
P-환이고, 임의의 (본질적으로) 유한형 $A$-대수도 그러하다.

\medskip\noindent
$\mathfrak q \subset A$를 극소소아이디얼이라 하자. 그러면
$A^\wedge/\mathfrak q A^\wedge = (A/\mathfrak q)^\wedge$이고
$A^h/\mathfrak qA^h = (A/\mathfrak q)^h$이다(Algebra, Lemma
\ref{algebra-lemma-quotient-henselization}). 따라서 $A$의 분지의 수는
환들 $A/\mathfrak q$의 분지의 수의 합이고, $A^\wedge$에 대해서도
마찬가지이다. 이 방식으로 $A$가 정역인 경우로 환원한다.

\medskip\noindent
$A$가 정역이라고 가정하자. $A'$을 $A$의 분수체 $K$ 안에서 $A$의
정수적 폐포라 하자. $A$는 나가타환이므로 $A \to A'$은
유한이다. $A$의 분지의 수는 $A'$의 극대아이디얼 $\mathfrak m'$의
수임을 상기하자(보조정리 \ref{more-algebra-lemma-branches}). 또한 다음을 상기하자.
% P02-E0568: frozen decomposition display is followed by a detached justification.
$$
(A')^\wedge = A' \otimes_A A^\wedge =
\prod\nolimits_{\mathfrak m' \subset A'} (A'_{\mathfrak m'})^\wedge
$$
이는 Algebra, Lemma
\ref{algebra-lemma-completion-finite-extension}에서 따른다.
$A'_{\mathfrak m'}$은 형식적 섬유들이 기하적으로 정규인 국소환이므로
$(A'_{\mathfrak m'})^\wedge$는 정규이다(보조정리
\ref{more-algebra-lemma-completion-normal-local-ring}). 따라서
$A' \otimes_A A^\wedge$의 극소소아이디얼들은 위 분해의 인자들과
$1$-대-$1$로 대응한다. $A \to A^\wedge$의 평탄성에 의해
% P02-E0568: frozen inclusion display lacks punctuation before the next sentence.
$$
A^\wedge \subset A' \otimes_A A^\wedge \subset K \otimes_A A^\wedge
$$
이다.
% P02-E0569: frozen English uses singular “ring” for the two endpoint rings.
왼쪽 환과 오른쪽 환은 같은 극소소아이디얼 집합을 가지므로 가운데
환도 그러하다.
% P02-E1232: frozen source explicitly omits the key intermediate minimal-prime argument.
(작은 세부사항은 생략한다.) 이로써 증명이 끝난다.

\medskip\noindent
기하적 분지의 수에 대해서도 참임을 보이기 위해, 분지에 대한 결과,
$A^{sh}$의 형식적 섬유들이 기하적으로 정규라는 사실(보조정리들
\ref{more-algebra-lemma-formal-fibres-normal}와 \ref{more-algebra-lemma-henselization-P-ring}),
그리고 보조정리 \ref{more-algebra-lemma-henselization-noetherian}에 의한 등식
$(A^{sh})^\wedge = ((A^\wedge)^{sh})^\wedge$를 사용한다.
\end{proof}

\section{형식적 catenary 환}
\label{more-algebra-section-formally-catenary}

\noindent
이 절에서는 뇌터 국소환이 보편적 catenary일 필요충분조건이 형식적
catenary라는 Ratliff의 정리 \cite{Ratliff}를 증명한다.

\begin{definition}
\label{more-algebra-definition-formally-catenary}
뇌터 국소환 $A$가 {\it 형식적 catenary}라는 것은 모든 극소소아이디얼
$\mathfrak p \subset A$에 대하여
$A^\wedge/\mathfrak p A^\wedge$의 스펙트럼이 등차원이라는 뜻이다.
\end{definition}

\noindent
$A$를 형식적 catenary인 뇌터 국소환이라 하자. Ratliff의 결과
(Proposition \ref{more-algebra-proposition-ratliff})에 의해 $A$의 모든 몫도 형식적
catenary이다(보편적 catenary 환의 모임은 몫 아래에서 안정적이기
때문이다). 따라서 $A$의 모든 소아이디얼 $\mathfrak p$에 대하여
$A^\wedge/\mathfrak p A^\wedge$의 스펙트럼은 등차원이다.

\begin{lemma}
\label{more-algebra-lemma-not-formally-catenary}
$(A, \mathfrak m)$을 형식적 catenary가 아닌 뇌터 국소환이라 하자.
그러면 $A$는 보편적 catenary가 아니다.
\end{lemma}

\begin{proof}
가정에 의해 어떤 극소소아이디얼 $\mathfrak p \subset A$가 존재하여
$A^\wedge /\mathfrak p A^\wedge$의 스펙트럼은 등차원이 아니다.
$A$를 $A/\mathfrak p$로 바꾸어 $A$가 정역이고 $A^\wedge$의
스펙트럼이 등차원이 아니라고 가정할 수 있다. $\mathfrak q$를
$A^\wedge$의 극소소아이디얼 가운데
$d = \dim(A^\wedge/\mathfrak q)$가 최소인 것으로 택하자. 그러면
$0 < d < \dim(A)$이다. $d$에 대한 귀납법으로 보조정리를 증명한다.

\medskip\noindent
$d = 1$인 경우를 보자. 이 경우
$\dim(A^\wedge_\mathfrak q) = 0$이다. 따라서
$A^\wedge_\mathfrak q$는 아르틴 국소환이고, 어떤 $n > 0$에 대하여
아이디얼 $J = \mathfrak q^n$의 $A^\wedge_\mathfrak q$에서의 상은
영이다. 그러므로 $J/J^2$의 유일한 결합소아이디얼은 $\mathfrak m$이고,
어떤 $m > 0$에 대하여 $\mathfrak m^m$이 $J/J^2$을 소멸시킨다.
따라서 보조정리 \ref{more-algebra-lemma-quotient-by-idempotent}를 사용하여
% P02-E0570: frozen English omits the noun after “find” in the finite-type map clause.
$B^\wedge \cong A^\wedge/J$인 유한형 환 사상 $A \to B$를 얻는다.
그러면 $\mathfrak m_B = \sqrt{\mathfrak mB}$는 $\mathfrak m$과
같은 잉여체를 갖는 극대아이디얼이고, $B^\wedge$은
$\mathfrak m_B$-진 완비화이다(Algebra, Lemma
\ref{algebra-lemma-finite-after-completion}). 따라서
$$
\dim(B_{\mathfrak m_B}) = \dim(B^\wedge) = 1 = d.
$$
분해 $A \to B \to A^\wedge/J$가 있으므로 $\mathfrak q/J$의 역상은
$A$의 $(0)$ 위에 놓이는 소아이디얼
$\mathfrak q' \subset \mathfrak m_B$이다. 따라서 $A$가 보편적
catenary라면 차원 공식(Algebra, Lemma
\ref{algebra-lemma-dimension-formula})에 의해
\begin{align*}
\dim(B_{\mathfrak m_B})
& \geq
\dim((B/\mathfrak q')_{\mathfrak m_B}) \\
& =
\dim(A) + \text{trdeg}_A(B/\mathfrak q') -
\text{trdeg}_{\kappa(\mathfrak m)}(\kappa(\mathfrak m_B)) \\
& =
\dim(A) + \text{trdeg}_A(B/\mathfrak q')
\end{align*}
% P02-E0571: frozen dimension-formula display lacks terminal punctuation.
이 모순으로 $d = 1$인 경우의 논증이 끝난다.

\medskip\noindent
$d > 1$이라고 가정하자. $Z \subset \Spec(A^\wedge)$를
$V(\mathfrak q)$와 다른 기약 성분들의 합집합이라 하자.
$\mathfrak r_1, \ldots, \mathfrak r_m \subset A^\wedge$를
$V(\mathfrak q) \cap Z$의 차원이 $> 0$인 기약 성분들에 대응하는
소아이디얼들이라 하자. 소아이디얼 회피(Algebra, Lemma
\ref{algebra-lemma-silly})를 사용하여
$f \in \mathfrak m$, $f \not \in A \cap \mathfrak r_j$를 택한다.
그러면 $\dim(A/fA) = \dim(A) - 1$이고 $V(\mathfrak q, f)$에는
차원 $d - 1$인 기약 성분이 있다. 따라서 $A/fA$는 형식적
catenary가 아니고 불변량 $d$가 감소하였다. 귀납 가정에 의해
$A/fA$는 보편적 catenary가 아니므로 $A$도 보편적 catenary가 아니다.
\end{proof}

\begin{lemma}
\label{more-algebra-lemma-flat-under-catenary-equidimensional}
$A \to B$를 국소 뇌터환들의 평탄 국소 환 사상이라 하자.
% P02-E0572: frozen hypothesis contains a spurious second “is”.
$B$가 catenary이고 $\Spec(B)$가 등차원이라고 가정하자. 그러면
\begin{enumerate}
\item 모든 $\mathfrak p \subset A$에 대하여
$\Spec(B/\mathfrak p B)$는 등차원이고,
\item $A$는 catenary이며 $\Spec(A)$는 등차원이다.
\end{enumerate}
\end{lemma}

\begin{proof}
$\mathfrak p \subset A$를 소아이디얼이라 하고,
$\mathfrak q \subset B$를 $\mathfrak pB$ 위의 극소소아이디얼이라 하자.
$A \to B$에 대한 내려감(Algebra, Lemma
\ref{algebra-lemma-flat-going-down})에 의해
$\mathfrak q \cap A = \mathfrak p$이다. 따라서
$A_\mathfrak p \to B_\mathfrak q$는 특수 섬유의 차원이 $0$인 평탄
국소 환 사상이고, 그러므로
$$
\dim(A_\mathfrak p) = \dim(B_\mathfrak q) = \dim(B) - \dim(B/\mathfrak q)
$$
이다(Algebra, Lemma
\ref{algebra-lemma-dimension-base-fibre-equals-total}).
% P02-E0573: frozen English explanation of the second equality lacks a finite verb.
두 번째 등식은 $\Spec(B)$가 등차원이고 $B$가 catenary이기 때문에
성립한다. 따라서 $\dim(B/\mathfrak q)$는 $\mathfrak q$의 선택과
무관하고, $\Spec(B/\mathfrak p B)$는 차원
$\dim(B) - \dim(A_\mathfrak p)$인 등차원 스펙트럼이다. 한편
평탄성에 의해(위에서 인용한 보조정리를 보라)
$\dim(B/\mathfrak p B) = \dim(A/\mathfrak p) + \dim(B/\mathfrak m_A B)$
이고
$\dim(B) = \dim(A) + \dim(B/\mathfrak m_A B)$
이므로
$$
\dim(A_\mathfrak p) = \dim(A) - \dim(A/\mathfrak p)
$$
이다.
% P02-E0574: frozen display lacks punctuation and describes a prime ideal as “in” A.
이는 $A$의 모든 $\mathfrak p$에 대해 성립한다. 이를 $A$의 모든
극소소아이디얼에 적용하면 $A$가 등차원임을 알 수 있다.
$\mathfrak p \subset \mathfrak p'$이 그 사이에 다른 소아이디얼이 없는
진포함이면, $A/\mathfrak p \to B/\mathfrak p B$가 평탄하고
$\Spec(B/\mathfrak p B)$가 등차원이므로 위의 논의를
$A/\mathfrak p$의 소아이디얼 $\mathfrak p'/\mathfrak p$에 적용하여
$$
1 = \dim((A/\mathfrak p)_{\mathfrak p'}) =
\dim(A/\mathfrak p) - \dim(A/\mathfrak p')
$$
을 얻는다.
% P02-E0575: frozen dimension-drop display lacks terminal punctuation.
따라서 $\mathfrak p \mapsto \dim(A/\mathfrak p)$는 차원 함수이고
$A$는 catenary이다.
\end{proof}

\begin{lemma}
\label{more-algebra-lemma-formally-catenary}
$A$를 형식적 catenary인 뇌터 국소환이라 하자.
그러면 $A$는 보편적 catenary이다.
\end{lemma}

\begin{proof}
$\mathfrak p$가 $A$의 극소소아이디얼일 때 $A$를
$A/\mathfrak p$로 바꾸어도 된다. Algebra, Lemma
\ref{algebra-lemma-catenary-check-irreducible}를 보라. 따라서
$A^\wedge$의 스펙트럼이 등차원이라고 가정할 수 있다.
$A$ 위에서 본질적으로 유한형인 모든 국소환이 catenary임을 보이면
충분하다(예를 들어 Algebra, Lemma
\ref{algebra-lemma-catenary-check-local}을 보라). 그러므로
$\mathfrak m \subset A[x_1, \ldots, x_n]$이
$\mathfrak m_A$ 위에 놓이는 극대아이디얼일 때
$A[x_1, \ldots, x_n]_\mathfrak m$이 catenary임을 보이면 충분하다.
Algebra, Lemma
\ref{algebra-lemma-localization-at-closed-point-special-fibre}
(그리고 Algebra, Lemmas \ref{algebra-lemma-quotient-catenary}와
\ref{algebra-lemma-localization-catenary})를 보라.
$\mathfrak m' \subset A^\wedge[x_1, \ldots, x_n]$을
$\mathfrak m$ 위에 놓이는 유일한 극대아이디얼이라 하자. 그러면
$$
A[x_1, \ldots, x_n]_\mathfrak m \to A^\wedge[x_1, \ldots, x_n]_{\mathfrak m'}
$$
은 국소이고 평탄하다(Algebra, Lemma
\ref{algebra-lemma-completion-flat}).
% P02-E0576: frozen map display lacks punctuation and the next clause omits “is”.
따라서 보조정리 \ref{more-algebra-lemma-flat-under-catenary-equidimensional}에 의해
오른쪽 환이 catenary이고 스펙트럼이 등차원임을 보이면 충분하다.
완비 국소환은 보편적 catenary이므로 이 환은 catenary이다
(Algebra, Remark
\ref{algebra-remark-Noetherian-complete-local-ring-universally-catenary}).
$A^\wedge[x_1, \ldots, x_n]_{\mathfrak m'}$의 임의의 극소소아이디얼
$\mathfrak q$를 택하자.
% P02-E1236: frozen source explicitly omits the minimal-prime extension argument.
그러면 어떤 $A^\wedge$의 극소소아이디얼 $\mathfrak p$에 대하여
$\mathfrak q = \mathfrak p A^\wedge[x_1, \ldots, x_n]_{\mathfrak m'}$
이다(작은 세부사항은 생략한다). 따라서
$$
\dim(A^\wedge[x_1, \ldots, x_n]_{\mathfrak m'}/\mathfrak q) =
\dim(A^\wedge/\mathfrak p) + n = \dim(A^\wedge) + n
$$
이다.
% P02-E0577: frozen dimension display lacks punctuation and its explanation lacks a finite verb.
첫 번째 등식은 Algebra, Lemma
\ref{algebra-lemma-dimension-base-fibre-equals-total}에서 따르고,
두 번째 등식은 $A^\wedge$의 스펙트럼이 등차원이기 때문에 성립한다.
이로써 증명이 끝난다.
\end{proof}

\begin{proposition}[Ratliff]
\label{more-algebra-proposition-ratliff}
\begin{reference}
\cite{Ratliff}
\end{reference}
뇌터 국소환이 보편적 catenary일 필요충분조건은 형식적
catenary라는 것이다.
\end{proposition}

\begin{proof}
보조정리 \ref{more-algebra-lemma-not-formally-catenary}와
\ref{more-algebra-lemma-formally-catenary}를 결합하라.
\end{proof}

\begin{lemma}
\label{more-algebra-lemma-geometrically-normal-fibres-universally-catenary}
\begin{reference}
\cite[Corollary 2.3]{Heinzer-Rotthaus-Wiegand}
\end{reference}
$(A, \mathfrak m)$을 형식적 섬유들이 기하적으로 정규인 뇌터
국소환이라 하자. 그러면
\begin{enumerate}
\item $A^h$는 보편적 catenary이고,
\item $A$가 단일분지이면(예를 들어 정규이면) $A$는 보편적
catenary이다.
\end{enumerate}
\end{lemma}

\begin{proof}
보조정리 \ref{more-algebra-lemma-geometrically-normal-formal-fibres-number-of-branches}에
의해 $A$와 $A^\wedge$의 분지의 수는 같으므로 보조정리
\ref{more-algebra-lemma-equal-nr-branches-completion}를 적용할 수 있다. 그러면
임의의 극소소아이디얼 $\mathfrak q \subset A^h$에 대하여
$A^\wedge/\mathfrak q A^\wedge$은 유일한 극소소아이디얼을 갖는다.
따라서 $A^h$는 정의에 의해 형식적 catenary이고, 명제
\ref{more-algebra-proposition-ratliff}에 의해 보편적 catenary이다. $A$가
단일분지이면 $A^h$는 유일한 극소소아이디얼을 가지므로
$A^\wedge$도 유일한 극소소아이디얼을 갖는다. 따라서 $A$는 형식적
catenary이고 같은 방식으로 결론을 얻는다.
\end{proof}

\section{군 작용과 정수적 폐포}
\label{more-algebra-section-group-actions-integral}

\noindent
이 절은 어떤 의미에서 Algebra, Section
\ref{algebra-section-going-down-integral-over-normal}의 연속이다.
비슷한 종류의 내용은 Fundamental Groups, Section
\ref{pione-section-group-actions-integral}에도 더 있다.
% P02-E0578: frozen section-opening cross-reference sentence lacks terminal punctuation.

\begin{lemma}
\label{more-algebra-lemma-pol-lifting}
$\varphi : A \to B$를 환의 전사라 하자. 위수가 $n$인 유한군 $G$가
$\varphi : A \to B$에 작용한다고 하자. $b \in B^G$이면
$B^G[T]$에서 $(T - b)^n$으로 가는 일계수 다항식
$P \in A^G[T]$가 존재한다.
\end{lemma}

\begin{proof}
$b$의 리프트 $a \in A$를 택하고
$P = \prod_{\sigma \in G} (T - \sigma(a))$로 둔다.
\end{proof}

\begin{lemma}
\label{more-algebra-lemma-invariants-modulo}
$R$을 환이라 하고 유한군 $G$가 $R$에 작용한다고 하자.
$I \subset R$을 모든 $\sigma \in G$에 대하여
$\sigma(I) \subset I$인 아이디얼이라 하자. 그러면
$R^G/I^G \subset (R/I)^G$는 스펙트럼 위 위상동형과 잉여체의 순수
비분리 확대들을 유도하는 정수적 환 확대이다.
\end{lemma}

\begin{proof}
$I^G = R^G \cap I$이므로 사상이 단사임은 명백하다. 보조정리
\ref{more-algebra-lemma-pol-lifting}에 의해 Algebra, Lemma
\ref{algebra-lemma-universally-bijective}를 적용할 수 있다.
\end{proof}

\begin{lemma}
\label{more-algebra-lemma-pol-in-ideal}
위수가 $n$인 유한군 $G$가 환 $R$에 작용한다고 하자.
$J \subset R^G$를 아이디얼이라 하자. $x \in JR$에 대하여
$\prod_{\sigma \in G} (T - \sigma(x)) =
T^n + a_1 T^{n - 1} + \ldots + a_n$이고 $a_i \in J$이다.
\end{lemma}

\begin{proof}
이 다항식은 실제로 일계수이고 그 계수들은 $R^G$에 속한다.
$f_j \in J$이고 $b_j \in R$인
$x = f_1 b_1 + \ldots + f_m b_m$으로 쓸 수 있다. 따라서 $m$에 대한
귀납법을 사용하여 $f \in J$, $b \in R$이고, 결론이 $y$에 대해
성립하는 $y \in JR$에 대하여 $x = y - fb$라고 가정할 수 있다.
그러면
% P02-E0579: both frozen coefficient-expansion displays lack their surrounding punctuation.
$$
\prod\nolimits_{\sigma \in G} (T - \sigma(x)) =
\prod\nolimits_{\sigma \in G} (T - \sigma(y) + f\sigma(b)) =
\prod\nolimits_{\sigma \in G} (T - \sigma(y)) +
\sum_{i = 1, \ldots, n} f^i a_i
$$
이고, 여기서
$$
a_i =
\sum\nolimits_{S \subset G,\ |S| = i}
\prod\nolimits_{\sigma \in S} \sigma(b)
\prod\nolimits_{\sigma \not \in S} (T - \sigma(y))
$$
이다.
% P02-E1238: frozen source omits the reindexing computation proving coefficient invariance.
생략하는 계산에 의해 $a_i \in R^G$임을 알 수 있다(힌트: 주어진 식은
대칭적이다). 따라서 보조정리의 명제에 나타나는 $x$에 대한 다항식은
$J$를 법으로 $y$에 대한 다항식과 합동이고, 이로써 귀납 단계가
증명된다.
\end{proof}

\begin{lemma}
\label{more-algebra-lemma-invariants-modulo-bis}
$R$을 환이라 하자. 위수가 $n$인 유한군 $G$가 $R$에 작용한다고 하자.
$J \subset R^G$를 아이디얼이라 하자.
% P02-E0580: frozen statement omits the article and list-introducing punctuation.
그러면 환 사상 $R^G/J \to (R/JR)^G$는 다음 성질들을 갖는다.
\begin{enumerate}
\item $b \in (R/JR)^G$에 대하여, $(R/JR)^G[T]$에서의 상이
$(T - b)^n$인 일계수 다항식 $P \in R^G/J[T]$가 존재한다.
\item $a \in \Ker(R^G/J \to (R/JR)^G)$에 대하여
$R^G/J[T]$에서 $(T - a)^n = T^n$이다.
\end{enumerate}
특히 $R^G/J \to (R/JR)^G$는 스펙트럼 위 위상동형과 잉여체의 순수
비분리 확대들을 유도하는 정수적 환 사상이다.
\end{lemma}

\begin{proof}
% P02-E0581: frozen English uses “Part (1) follow” with singular subject.
(1)은 $I = JR$로 둔 보조정리 \ref{more-algebra-lemma-pol-lifting}에서 따른다.
$a$가 (2)에서와 같으면 $a$는 어떤 $x \in R^G \cap JR$의 상이다.
따라서 보조정리 \ref{more-algebra-lemma-pol-in-ideal}에 의해
$(T - x)^n = \prod_{\sigma \in G} (T - \sigma(x))$는 $J$를 법으로
$T^n$과 합동이다. 이로써 (2)가 증명된다. 마지막 명제를 보이려면
Algebra, Lemma \ref{algebra-lemma-universally-bijective}를 적용하면 된다.
\end{proof}

\begin{remark}
\label{more-algebra-remark-when-iso}
보조정리 \ref{more-algebra-lemma-invariants-modulo-bis}에서 $n$이 $R$에서
가역이면 사상 $R^G/J \to (R/JR)^G$가 동형사상임을 알 수 있다.
\end{remark}

\begin{lemma}
\label{more-algebra-lemma-functor-invariants-tensor}
$R$을 환이라 하자. 위수가 $n$인 유한군 $G$가 $R$에 작용한다고 하자.
$A$를 $R^G$-대수라 하자.
\begin{enumerate}
\item $b \in (A \otimes_{R^G} R)^G$에 대하여,
$(A \otimes_{R^G} R)^G[T]$에서의 상이 $(T - b)^n$인 일계수 다항식
$P \in A[T]$가 존재한다.
\item $a \in \Ker(A \to (A \otimes_{R^G} R)^G)$에 대하여
$A[T]$에서 $(T - a)^n = T^n$이다.
\end{enumerate}
\end{lemma}

\begin{proof}
$E$가 $R^G$ 위의 다항식 대수인 전사 $E \to A$를 택한다.
$E$는 $R^G$-가군으로서 자유이므로
$(E \otimes_{R^G} R)^G = E$이다.
% P02-E0582: frozen notation assignment “Denote J = ...” is malformed.
$J = \Ker(E \to A)$로 둔다. 텐서곱은 오른쪽 완전이므로
$A \otimes_{R^G} R$은 $J$가 생성하는 아이디얼에 의한
$E \otimes_{R^G} R$의 몫이다. 따라서 이 보조정리는 보조정리
\ref{more-algebra-lemma-invariants-modulo-bis}의 특수한 경우이다.
\end{proof}

\begin{lemma}
\label{more-algebra-lemma-base-change-invariants}
$R$을 환이라 하고 유한군 $G$가 $R$에 작용한다고 하자.
$R^G \to A$를 환 사상이라 하자. 사상
$$
A \to (A \otimes_{R^G} R)^G
$$
은 $R^G \to A$가 평탄이면 동형사상이다. 일반적으로 이 사상은
정수적이고, 스펙트럼 위 위상동형과 잉여체의 순수 비분리 확대들을
유도한다.
\end{lemma}

\begin{proof}
첫 번째 명제를 보이기 위해 완전열
$0 \to R^G \to R \to \bigoplus_{\sigma \in G} R$을 생각하자.
두 번째 사상은 $x$를 $(\sigma(x) - x)_{\sigma \in G}$로 보낸다.
$A$와 텐서곱하면 $R^G \to A$가 평탄일 때 이 열은 여전히 완전하다.
% P02-E0583: frozen conclusion redundantly takes invariants inside an invariant ring.
따라서 $A$는 $(A \otimes_{R^G} R)^G$ 안의 $G$-불변원들이다.

\medskip\noindent
두 번째 명제는 보조정리 \ref{more-algebra-lemma-functor-invariants-tensor}와
Algebra, Lemma \ref{algebra-lemma-universally-bijective}에서 따른다.
\end{proof}

\begin{lemma}
\label{more-algebra-lemma-one-orbit}
유한군 $G$가 환 $R$에 작용한다고 하자. $R^G$의 같은 소아이디얼 위에
놓이는 임의의 두 소아이디얼 $\mathfrak q, \mathfrak q' \subset R$에
대하여 $\sigma(\mathfrak q) = \mathfrak q'$인
$\sigma \in G$가 존재한다.
\end{lemma}

\begin{proof}
모든 $x \in R$은 $R^G[T]$의 일계수 다항식
$\prod_{\sigma \in G}(T - \sigma(x))$의 근이므로
$R^G \subset R$은 정수적 확대이다. 따라서 주어진 소아이디얼
$\mathfrak p$ 위에 놓이는 소아이디얼들 사이에는 포함관계가 없다
(Algebra, Lemma \ref{algebra-lemma-integral-no-inclusion}).
보조정리가 거짓이라면 모든 $\sigma \in G$에 대하여
$x \in \mathfrak q'$, $x \not \in \sigma(\mathfrak q)$인 원소를
택할 수 있다. Algebra, Lemma \ref{algebra-lemma-silly}를 보라.
그러면 $y = \prod_{\sigma \in G} \sigma(x)$는 $R^G$에 속하고
$\mathfrak p = R^G \cap \mathfrak q'$에도 속한다. 한편 모든
$\sigma$에 대해 $x \not \in \sigma(\mathfrak q)$라는 것은 모든
$\sigma$에 대해 $\sigma(x) \not \in \mathfrak q$라는 뜻이다.
$\mathfrak q$는 소아이디얼이므로 $y \not \in \mathfrak q$이다.
이는 $y \in \mathfrak p \subset \mathfrak q$와 모순이다.
\end{proof}

\begin{lemma}
\label{more-algebra-lemma-one-orbit-geometric}
유한군 $G$가 환 $R$에 작용한다고 하자.
$\mathfrak p \subset R^G$ 위에 놓이는 소아이디얼
$\mathfrak q \subset R$를 택하자. 그러면
$\kappa(\mathfrak q)/\kappa(\mathfrak p)$는 대수적 정규 확대이고 사상
$$
D = \{\sigma \in G \mid \sigma(\mathfrak q) = \mathfrak q\}
\longrightarrow
\text{Aut}(\kappa(\mathfrak q)/\kappa(\mathfrak p))
$$
은 전사이다\footnote{갈루아 확대인 경우에만
$\text{Gal}$이라는 표기를 사용함을 상기하라.}.
\end{lemma}

\begin{proof}
$A = (R^G)_\mathfrak p$와 $B = A \otimes_{R^G} R$로 두면, 국소화가
평탄이므로 보조정리 \ref{more-algebra-lemma-base-change-invariants}에 의해
$A = B^G$이다. $\mathfrak pA$와 $\mathfrak qB$는 소아이디얼이고,
$D$는 $\mathfrak qB$의 안정자이며,
$\kappa(\mathfrak p) = \kappa(\mathfrak pA)$와
$\kappa(\mathfrak q) = \kappa(\mathfrak qB)$임을 관찰하자. 따라서
$R$을 $B$로 바꾸어 $\mathfrak p$가 극대아이디얼이라고 가정할 수
있다. $R^G \subset R$은 정수적 환 확대이므로 $R$의 극대아이디얼들은
정확히 $\mathfrak p$ 위에 놓이는 소아이디얼들이다(Algebra, Lemmas
\ref{algebra-lemma-integral-no-inclusion}와
\ref{algebra-lemma-integral-going-up}에서 따른다). 보조정리
\ref{more-algebra-lemma-one-orbit}에 의해 이들은 유한개이고,
$\mathfrak q = \mathfrak q_1, \mathfrak q_2, \ldots, \mathfrak q_m$라
쓸 수 있으며 $G$의 한 궤도를 이룬다. 중국인의 나머지 정리
(Algebra, Lemma \ref{algebra-lemma-chinese-remainder})에 의해 사상
$R \to \prod_{j = 1, \ldots, m} R/\mathfrak q_j$은 전사이다.

\medskip\noindent
먼저 확대가 정규임을 증명한다. 원소
$\alpha \in \kappa(\mathfrak q)$를 택하자. $\kappa(\mathfrak p)$ 위에서
$\alpha$의 최소다항식 $P$가 완전히 분해됨을 보이면 된다. 위의
결과에 의해 $\kappa(\mathfrak q)$에서 $\alpha$로 가는
$a \in \mathfrak q_2 \cap \ldots \cap \mathfrak q_m$을 택할 수 있다.
$R^G[T]$의 다항식
$Q = \prod_{\sigma \in G} (T - \sigma(a))$를 생각하자. $Q$의
$R[T]$에서의 상은 일차 인자들로 완전히 분해되므로
$\kappa(\mathfrak q)[T]$에서의 상도 그러하다. $P$는
$\kappa(\mathfrak p)[T]$에서 $Q$의 상을 나누므로, 원하는 대로
$P$는 $\kappa(\mathfrak q)$ 위에서 일차 인자들로 완전히 분해된다.

\medskip\noindent
$\kappa(\mathfrak q)/\kappa(\mathfrak p)$가 정규이므로 Fields, Lemma
\ref{fields-lemma-normal-case}에 의해
$\kappa(\mathfrak q) = \kappa_1 \otimes_{\kappa(\mathfrak p)} \kappa_2$이고
$\kappa_1/\kappa(\mathfrak p)$가 순수 비분리이며
$\kappa_2/\kappa(\mathfrak p)$가 갈루아라고 가정할 수 있다.
$\kappa_2$가 유한이면 $\kappa(\mathfrak p)$ 위에서 이를 생성하는
$\alpha \in \kappa_2$를 택하고, 무한이면 차수가 $> |G|$인 부분체를
생성하는 것을 택한다(모순을 얻기 위해서이다).
% P02-E0584: frozen infinite-case wording leaves the generated subfield and base implicit.
이는 Fields, Lemma \ref{fields-lemma-primitive-element}에 의해 가능하다.
앞 단락에서와 같이 $a$, $P$, $Q$를 택한다.
$\alpha' \in \kappa_2$가 $\kappa(\mathfrak p)$ 위에서 $\alpha$의
갈루아 켤레이면,
% P02-E0585: frozen source says P divides the image of P, but the argument requires Q.
$\kappa(\mathfrak p)[T]$에서 $P$가 $P$의 상을 나눈다는 사실에 의해
$\sigma(a)$가 $\alpha'$로 가는 $\sigma \in G$가 존재한다.
$a$의 선택(다른 극대아이디얼들에서 영이 됨)에 의해 이는
$\sigma \in D$임을 뜻하고,
$\text{Aut}(\kappa(\mathfrak q)/\kappa(\mathfrak p))$에서의
$\sigma$의 상은 $\alpha$를 $\alpha'$로 보낸다.
% P02-E0586: frozen paragraph ends with an incomplete two-case conclusion.
따라서 전사성을 얻거나, $\alpha$가 $\kappa(\mathfrak p)$ 위에서
차수 $> |G|$인 경우에는 원하는 모순을 얻는다.
\end{proof}

\begin{lemma}
\label{more-algebra-lemma-one-orbit-geometric-galois}
$A$를 분수체가 $K$인 정규 정역이라 하자. $L/K$를 (무한일 수도 있는)
갈루아 확대라 하자. $G = \text{Gal}(L/K)$라 하고, $B$를 $L$ 안에서
$A$의 정수적 폐포라 하자.
\begin{enumerate}
\item $A$의 같은 소아이디얼 위에 놓이는 임의의 두 소아이디얼
$\mathfrak q, \mathfrak q' \subset B$에 대하여
$\sigma(\mathfrak q) = \mathfrak q'$인 $\sigma \in G$가 존재한다.
\item $\mathfrak p \subset A$ 위에 놓이는 소아이디얼
$\mathfrak q \subset B$를 택하자. 그러면
$\kappa(\mathfrak q)/\kappa(\mathfrak p)$는 대수적 정규 확대이고 사상
$$
D = \{\sigma \in G \mid \sigma(\mathfrak q) = \mathfrak q\}
\longrightarrow
\text{Aut}(\kappa(\mathfrak q)/\kappa(\mathfrak p))
$$
은 전사이다.
\end{enumerate}
\end{lemma}

\begin{proof}
(1)의 증명. $K \subset M \subset L$이고 $M/K$가 갈루아가 되는
부분체와, $\sigma \in \text{Gal}(M/K)$가
$\sigma(\mathfrak q \cap M) = \mathfrak q' \cap M$을 만족하게 하는
쌍 $(M, \sigma)$를 생각하자.
$(M', \sigma') \geq (M, \sigma)$라는 것은 $M \subset M'$이고
$\sigma'|_M = \sigma$라는 뜻으로 정한다. $A$가 정규 정역이어서
$A = K \cap B$이므로 $(K, \text{id}_K)$는 이러한 쌍이다.
이 쌍들의 모임은 초른 보조정리의 가정들을 만족하므로 극대 쌍
$(M, \sigma)$가 존재한다. $M \not = L$이면
$M \subset M' \subset L$이고 $M'/M$은 자명하지 않은 유한 확대이며
$M'/K$는 갈루아가 되도록 택할 수 있다(Fields, Lemma
\ref{fields-lemma-normal-closure-inside-normal}).
$M$에 대한 제한이 $\sigma$인
$\sigma' \in \text{Gal}(M'/K)$를 택한다(Fields, Lemma
\ref{fields-lemma-galois-infinite}). 그러면 소아이디얼들
$\sigma'(\mathfrak q \cap M')$과 $\mathfrak q' \cap M'$은
$B \cap M$의 같은 소아이디얼로 제한된다.
$B \cap M = (B \cap M')^{\text{Gal}(M'/M)}$이므로 보조정리
\ref{more-algebra-lemma-one-orbit}를 사용하여
$\tau(\sigma'(\mathfrak q \cap M')) = \mathfrak q' \cap M'$인
$\tau \in \text{Gal}(M'/M)$를 찾을 수 있다. 따라서
$(M', \tau \circ \sigma') > (M, \sigma)$이고, 이는
$(M, \sigma)$의 극대성과 모순이다.

\medskip\noindent
(2)도 (1)과 정확히 같은 방식으로 증명된다. 세부사항을 적어 보자.
$\overline{\sigma} \in
\text{Aut}(\kappa(\mathfrak q)/\kappa(\mathfrak p))$를 택한다.
$K \subset M \subset L$이고 $M/K$가 갈루아가 되는 부분체와,
$\sigma \in \text{Gal}(M/K)$가
$\sigma(\mathfrak q \cap M) = \mathfrak q \cap M$을 만족시키며
$$
\xymatrix{
\kappa(\mathfrak q \cap M) \ar[r] \ar[d]_\sigma &
\kappa(\mathfrak q) \ar[d]_{\overline{\sigma}} \\
\kappa(\mathfrak q \cap M) \ar[r] & \kappa(\mathfrak q)
}
$$
가 가환하게 하는 쌍 $(M, \sigma)$를 생각하자.
$(M', \sigma') \geq (M, \sigma)$라는 것은 $M \subset M'$이고
$\sigma'|_M = \sigma$라는 뜻으로 정한다. 위와 마찬가지로
$(K, \text{id}_K)$는 이러한 쌍이다. 이 쌍들의 모임은 초른
보조정리의 가정들을 만족하므로 극대 쌍 $(M, \sigma)$가 존재한다.
$M \not = L$이면 $M \subset M' \subset L$이고 $M'/M$은 유한이며
$M'/K$는 갈루아가 되도록 택할 수 있다(Fields, Lemma
\ref{fields-lemma-normal-closure-inside-normal}).
$M$에 대한 제한이 $\sigma$인
$\sigma' \in \text{Gal}(M'/K)$를 택한다(Fields, Lemma
\ref{fields-lemma-galois-infinite}). 그러면 소아이디얼들
$\sigma'(\mathfrak q \cap M')$과 $\mathfrak q \cap M'$은
$B \cap M$의 같은 소아이디얼로 제한된다. 첫 단락에서와 같이
$\sigma'$의 선택을 조정하여
$\sigma'(\mathfrak q \cap M') = \mathfrak q \cap M'$이라고 가정할
수 있다. 그러면 $\sigma'$와 $\overline{\sigma}$는
$\kappa(\mathfrak q \cap M') \to \kappa(\mathfrak q)$인 사상들을
정의하고, 이들은 $\kappa(\mathfrak q \cap M)$ 위에서 일치한다.
$B \cap M = (B \cap M')^{\text{Gal}(M'/M)}$이므로 보조정리
\ref{more-algebra-lemma-one-orbit-geometric}를 사용하여
$\tau(\mathfrak q \cap M') = \mathfrak q \cap M'$이고
% P02-E0587: frozen source composes tau with sigma, although only sigma-prime acts on M-prime.
$\tau \circ \sigma$와 $\overline{\sigma}$가
$\kappa(\mathfrak q \cap M')$ 위에서 같은 사상을 유도하는
$\tau \in \text{Gal}(M'/M)$를 찾을 수 있다. 여기에는 작은
세부사항이 있다. 보조정리는 먼저
$\kappa(\mathfrak q \cap M')/\kappa(\mathfrak q \cap M)$이
정규임을 보장하고, 이로부터 두 사상의 차이가 이 확대의
자기동형사상임을 알 수 있다(Fields, Lemma
\ref{fields-lemma-normal-embeddings-differ-by-aut}). 여기에 보조정리를
적용하여 $\tau$를 얻는다. 따라서
$(M', \tau \circ \sigma') > (M, \sigma)$이고, 이는
$(M, \sigma)$의 극대성과 모순이다.
\end{proof}

\begin{lemma}
\label{more-algebra-lemma-one-orbit-geometric-galois-compare}
$A$를 분수체가 $K$인 정규 정역이라 하자.
$M/L/K$를 $K$의 (무한일 수도 있는) 갈루아 확대들의 탑이라 하자.
$H = \text{Gal}(M/K)$와 $G = \text{Gal}(L/K)$라 하고,
% P02-E0588: frozen English uses singular “closure” for the two respective rings.
$C$와 $B$를 각각 $M$과 $L$ 안에서의 $A$의 정수적 폐포라 하자.
$\mathfrak r \subset C$라 하고
$\mathfrak q = B \cap \mathfrak r$로 둔다. 또한
$$
D_\mathfrak r = \{\tau \in H \mid \tau(\mathfrak r) = \mathfrak r\}
$$
및
$$
I_\mathfrak r = \{\tau \in D_\mathfrak r \mid
\tau \bmod \mathfrak r = \text{id}_{\kappa(\mathfrak r)}\}
$$
로 두고, $D_\mathfrak q$와 $I_\mathfrak q$도 마찬가지로 정의한다.
사상 $H \to G$ 아래에서 유도되는 사상
$D_\mathfrak r \to D_\mathfrak q$와
$I_\mathfrak r \to I_\mathfrak q$는 전사이다.
\end{lemma}

\begin{proof}
$\sigma \in D_\mathfrak q$라 하자. $\sigma$로 가는
$\tau \in H$를 택한다. 이는 Fields, Lemma
\ref{fields-lemma-galois-infinite}에 의해 가능하다. 그러면
$\tau(\mathfrak r)$과 $\mathfrak r$은 모두 $\mathfrak q$ 위에 놓인다.
따라서 보조정리 \ref{more-algebra-lemma-one-orbit-geometric-galois}에 의해
$\sigma'(\tau(\mathfrak r)) = \mathfrak r$인
$\sigma' \in \text{Gal}(M/L)$가 존재한다. 그러므로
$\sigma'\tau \in D_\mathfrak r$은 $\sigma$로 간다. 관성군의 경우도
자기동형사상군으로의 전사성을 사용하여 정확히 같은 방식으로 증명한다.
\end{proof}

\section{이산 값매김환의 확대}
\label{more-algebra-section-discrete-valuation-rings}

\noindent
이 절과 다음 몇 절에서는 아래 정의들을 사용한다.

\begin{definition}
\label{more-algebra-definition-extension-discrete-valuation-rings}
$A$와 $B$가 이산 값매김환이고 $A \to B$가 단사인 국소 사상이면
$A \to B$ 또는 $A \subset B$를 {\it 이산 값매김환의 확대}라 한다.
특히 $\pi_A$와 $\pi_B$가 각각 $A$와 $B$의 균등화원이면, 어떤
$e \geq 1$과 $B$의 단원 $u$에 대하여
$\pi_A = u \pi_B^e$이다. 정수 $e$는 균등화원의 선택에 의존하지
않는다. 실제로 이는
$$
\mathfrak m_A B = \mathfrak m_B^e
$$
를 만족하는 유일한 정수 $\geq 1$이기도 하다.
% P02-E0589: frozen ramification-index display lacks terminal punctuation.
정수 $e$를 $A$ 위의 $B$의 {\it 분기지수}라 한다. $e = 1$이면
$B$가 $A$ 위에서 {\it 약하게 비분기}라고 한다. 잉여체의 확대
$\kappa_A = A/\mathfrak m_A \subset \kappa_B = B/\mathfrak m_B$가
유한이면 $f = [\kappa_B : \kappa_A]$로 두고 이를 확대
$A \subset B$의 {\it 잔여차수} 또는 {\it 잉여차수}라 한다.
\end{definition}

\noindent
분수체의 확대가 유한이라고 요구하지 않음에 주의하자.

\begin{lemma}
\label{more-algebra-lemma-inequality}
$A \subset B$를 분수체가 $K \subset L$인 이산 값매김환의 확대라
하자. 확대 $L/K$가 유한이면 잉여체 확대도 유한이고
$ef \leq [L : K]$이다.
\end{lemma}

\begin{proof}
% P02-E0595: frozen English equates the finiteness claim with its citation.
잉여체 확대의 유한성은 Algebra, Lemma
\ref{algebra-lemma-finite-extension-residue-fields-dimension-1}에서
따른다. 부등식은 Algebra, Lemmas
\ref{algebra-lemma-finite-length}와
\ref{algebra-lemma-pushdown-module}에서 따른다.
\end{proof}

\begin{lemma}
\label{more-algebra-lemma-multiplicative-e-f}
$A \subset B \subset C$를 이산 값매김환의 확대들이라 하자.
$B/A$와 $C/B$의 분기지수의 곱은 $C/A$의 분기지수이다.
% P02-E0596: frozen “In a formula” lacks punctuation before the equation.
식으로 쓰면 $e_{C/A} = e_{B/A} e_{C/B}$이다. 잔여차수들이 유한인
경우에도 마찬가지이다.
\end{lemma}

\begin{proof}
이는 정의와 Fields, Lemma
\ref{fields-lemma-multiplicativity-degrees}에서 즉시 따른다.
\end{proof}

\begin{lemma}
\label{more-algebra-lemma-ramification-index-a-power-of-p}
$A \subset B$를 체 확대 $K \subset L$을 유도하는 이산 값매김환의
확대라 하자. $K$의 표수가 $p > 0$이고 $L$이 $K$ 위에서 순수
비분리이면, 분기지수 $e$는 $p$의 거듭제곱이다.
\end{lemma}

\begin{proof}
어떤 $u \in B^*$에 대하여 $\pi_A = u \pi_B^e$로 쓰자. 한편 어떤
$p$-거듭제곱 $q$에 대하여 $\pi_B^q \in K$이다. 어떤 $v \in A^*$와
$k \in \mathbf{Z}$에 대하여 $\pi_B^q = v \pi_A^k$로 쓰자. 그러면
$\pi_A^q = u^q \pi_B^{qe} = u^q v^e \pi_A^{ke}$이다. $B$에서
값매김을 취하면 $ke = q$를 얻는다.
\end{proof}

\noindent
다음 보조정리에서는 이산 값매김환의 확대 $A \subset B$가
“비분기”라는 것, 즉 분기지수가 $1$이고 잉여체 확대가 (대수적이지
않을 수도 있는) 분리 확대라는 것의 의미를 논한다. 그러나 이미
비분기 환 사상이라는 개념이 있으므로 “비분기”라는 말 자체를
사용할 수는 없다. Algebra, Section
\ref{algebra-section-unramified}를 보라.

\begin{lemma}
\label{more-algebra-lemma-extension-dvrs-formally-smooth}
$A \subset B$를 이산 값매김환의 확대라 하자. 다음은 동치이다.
\begin{enumerate}
\item $A \to B$가 $\mathfrak m_B$-진 위상에 대하여 형식적 매끄럽고,
\item $A \to B$가 약하게 비분기이며 $\kappa_B/\kappa_A$가
분리 체 확대이다.
\end{enumerate}
\end{lemma}

\begin{proof}
이는 명제 \ref{more-algebra-proposition-fs-flat-fibre-fs}와 Algebra, Proposition
\ref{algebra-proposition-characterize-separable-field-extensions}에서
따른다.
\end{proof}

\begin{remark}
\label{more-algebra-remark-finite-separable-extension}
$A$를 분수체가 $K$인 이산 값매김환이라 하고, $L/K$를 유한 분리
체 확대라 하자. $B \subset L$을 $L$ 안에서의 $A$의 정수적 폐포라
하자. 그림은 다음과 같다.
% P02-E0590: frozen diagram display lacks terminal punctuation.
$$
\xymatrix{
B \ar[r] & L \\
A \ar[u] \ar[r] & K \ar[u]
}
$$
Algebra, Lemma
\ref{algebra-lemma-Noetherian-normal-domain-finite-separable-extension}에
의해 환 확대 $A \subset B$는 유한이므로 $B$는 뇌터환이다.
Algebra, Lemma \ref{algebra-lemma-integral-sub-dim-equal}에 의해
$B$의 차원은 $1$이므로 $B$는 데데킨트 정역이다. Algebra, Lemma
\ref{algebra-lemma-characterize-Dedekind}를 보라.
$\mathfrak m_1, \ldots, \mathfrak m_n$을 $B$의 극대아이디얼들
(즉 $\mathfrak m_A$ 위에 놓이는 소아이디얼들)이라 하자. 이산
값매김환의 확대
% P02-E0590: frozen localized-DVR display lacks its continuing comma.
$$
A \subset B_{\mathfrak m_i}
$$
들을 얻고, 따라서 분기지수 $e_i$와 잉여차수 $f_i$를 얻는다.
% P02-E0590: frozen degree-formula display lacks its continuing punctuation.
$$
[L : K] = \sum\nolimits_{i = 1, \ldots, n} e_i f_i
$$
이는 Algebra, Lemma \ref{algebra-lemma-finite-extension-dim-1}를
$A$의 균등화원에 적용한 결과이다. $A$가 헨젤환이면
$n = 1$임을 관찰하자(Algebra, Lemma
\ref{algebra-lemma-finite-over-henselian}에 의해). 예를 들어
$A$가 완비이면 그러하다.
\end{remark}

\begin{definition}
\label{more-algebra-definition-types-of-extensions}
$A$를 분수체가 $K$인 이산 값매김환이라 하고, $L/K$를 유한 분리
확대라 하자. 비고 \ref{more-algebra-remark-finite-separable-extension}에서와 같이
$B$와 $\mathfrak m_i$, $i = 1, \ldots, n$을 잡는다. 확대 $L/K$가
다음을 만족한다고 말한다.
\begin{enumerate}
\item 모든 $i$에 대하여 $e_i = 1$이고 확대
$\kappa(\mathfrak m_i)/\kappa_A$가 분리이면
{\it $A$에 관하여 비분기}이다.
% P02-E1246: frozen tame-ramification item ambiguously scopes its two characteristic alternatives.
\item $\kappa_A$의 표수가 $0$이거나, 또는 $\kappa_A$의 표수가
$p > 0$이고 체 확대들 $\kappa(\mathfrak m_i)/\kappa_A$가 분리이며
분기지수들 $e_i$가 $p$와 서로소이면 {\it $A$에 관하여 순분기}이다.
\item $n = 1$이고 잉여체 확대
$\kappa(\mathfrak m_1)/\kappa_A$가 자명하면
{\it $A$에 관하여 완전 분기}이다.
\end{enumerate}
이산 값매김환 $A$가 문맥에서 명백하면, 줄여서 $L/K$가 비분기,
완전 분기 또는 순분기라고 하기도 한다.
\end{definition}

\noindent
비분기 확대에 대해서는 다음 기본 보조정리가 성립한다.

\begin{lemma}
\label{more-algebra-lemma-permanence-unramified}
$A$를 분수체가 $K$인 이산 값매김환이라 하자.
\begin{enumerate}
\item $M/L/K$가 유한 분리 확대들이고 $M$이 $A$에 관하여
비분기이면, $L$도 $A$에 관하여 비분기이다.
\item $L/K$가 $A$에 관하여 비분기인 유한 분리 확대이면, $L$을
포함하고 $A$에 관하여 비분기인 갈루아 확대 $M/K$가 존재한다.
\item $L_1/K$, $L_2/K$가 $A$에 관하여 비분기인 유한 분리
확대들이면, $L_1$과 $L_2$를 포함하고 $A$에 관하여 비분기인 유한
분리 확대 $L/K$가 존재한다.
% P02-E0591: frozen part (3) duplicates the English indefinite article.
\end{enumerate}
\end{lemma}

\begin{proof}
비고 \ref{more-algebra-remark-finite-separable-extension}의 논의 결과들을 더 언급하지
않고 사용한다.

\medskip\noindent
(1)의 증명. $C/B/A$를 $M/L/K$ 안에서의 $A$의 정수적 폐포들이라
하자. $C$는 $B$의 유한 환 확대이므로
$\Spec(C) \to \Spec(B)$는 전사이다.
% P02-E0592: frozen English says “for ever maximal ideal”.
따라서 모든 극대아이디얼 $\mathfrak m \subset B$에 대하여 $\mathfrak m$ 위에
놓이는 극대아이디얼 $\mathfrak m' \subset C$가 존재한다.
분기지수의 곱셈성(보조정리 \ref{more-algebra-lemma-multiplicative-e-f})과 가정에
의해 $A$ 위의 $B_\mathfrak m$의 분기지수는 $1$이다.
$\kappa(\mathfrak m')/\kappa_A$가 유한 분리이므로
$\kappa(\mathfrak m)/\kappa_A$도 그러하다.

\medskip\noindent
(2)의 증명. $M$을 $K$ 위의 $L$의 정규 폐포라 하자. Fields,
Definition \ref{fields-definition-normal-closure}를 보라. 그러면
Fields, Lemma \ref{fields-lemma-normal-closure-galois}에 의해
$M/K$는 갈루아이다. 한편 Fields, Lemma
\ref{fields-lemma-normal-closure-tensor-product}에 의해 $K$-대수의
전사
% P02-E0593: frozen tensor-product source leaves its finite arity unbound.
$$
L \otimes_K \ldots \otimes_K L \longrightarrow M
$$
가 존재한다. $B$를 비고 \ref{more-algebra-remark-finite-separable-extension}에서와
같이 $L$ 안에서의 $A$의 정수적 폐포라 하자. $L$이 $A$에 관하여
비분기라는 조건은 정확히 $A \to B$가 에탈 환 사상이라는 뜻이다.
Algebra, Lemma \ref{algebra-lemma-characterize-etale}를 보라.
에탈 환 사상의 영속성에 의해
$$
B \otimes_A \ldots \otimes_A B
$$
는 $A$ 위에서 에탈이다. Algebra, Lemma
\ref{algebra-lemma-etale}를 보라. 따라서 표시된 환은 데데킨트
정역들의 곱이다. 보조정리 \ref{more-algebra-lemma-Dedekind-etale-extension}을
보라. 그러므로 $M$은 $A$ 위 유한 에탈인 데데킨트 정역의 분수체이다.
이는 원하는 대로 $M$이 $A$에 관하여 비분기라는 뜻이다.

\medskip\noindent
(3)의 증명.
% P02-E0594: frozen indexed closure notation does not bind the index or respective fields.
$B_i \subset L_i$를 $A$의 각 정수적 폐포라 하자. 위와 같은 방식으로
$B_1 \otimes_A B_2$가 $A$ 위 유한 에탈임을 보인다.
% P02-E1250: frozen proof omits selection of a product factor and fraction field.
세부사항은 생략한다.
\end{proof}

\begin{lemma}
\label{more-algebra-lemma-composition-unramified}
$A$를 분수체가 $K$인 이산 값매김환이라 하고, $M/L/K$를 유한 분리
확대들이라 하자. $B$를 $L$ 안에서의 $A$의 정수적 폐포라 하자.
$L/K$가 $A$에 관하여 비분기이고, $B$의 모든 극대아이디얼
$\mathfrak m$에 대하여 $M/L$이 $B_\mathfrak m$에 관하여
비분기이면, $M/K$는 $A$에 관하여 비분기이다.
\end{lemma}

\begin{proof}
$C$를 $M$ 안에서의 $A$의 정수적 폐포라 하자. $C$의 모든
극대아이디얼 $\mathfrak m'$은 $B$의 어떤 극대아이디얼
$\mathfrak m$ 위에 놓인다. 그러면 보조정리
\ref{more-algebra-lemma-multiplicative-e-f}의 분기지수 곱셈성과 유한 체 확대들의 탑
$\kappa(\mathfrak m')/\kappa(\mathfrak m)/\kappa_A$에 의해 결론이
따른다.
\end{proof}

\section{갈루아 확대와 분기}
\label{more-algebra-section-ramification}

\noindent
갈루아 확대의 경우에는 Section
\ref{more-algebra-section-discrete-valuation-rings}의 논의를 더 상세히 전개할 수 있다.

\begin{lemma}
\label{more-algebra-lemma-galois}
$A$를 분수체가 $K$인 이산 값매김환이라 하자.
$L/K$를 갈루아군이 $G$인 유한 갈루아 확대라 하자.
그러면 $G$는 비고 \ref{more-algebra-remark-finite-separable-extension}의 환 $B$에
작용하고, $B$의 극대아이디얼들의 집합에 추이적으로 작용한다.
\end{lemma}

\begin{proof}
$A$는 $K$ 안에서 정수적으로 닫혀 있고 $K = L^G$이므로
$A = B^G$임을 관찰하자. 따라서 이 보조정리는 보조정리
\ref{more-algebra-lemma-one-orbit}의 특수한 경우이다.
\end{proof}

\begin{lemma}
\label{more-algebra-lemma-galois-conclusion}
$A$를 분수체가 $K$인 이산 값매김환이라 하자.
$L/K$를 유한 갈루아 확대라 하자. 그러면 모든 $i$에 대하여
$e_i = e$이고 $f_i = f$가 되는 $e \geq 1$과 $f \geq 1$이 존재한다
(표기는 비고 \ref{more-algebra-remark-finite-separable-extension}과 같다). 특히
$[L : K] = n e f$이다.
\end{lemma}

\begin{proof}
보조정리 \ref{more-algebra-lemma-galois}와 정의들에서 즉시 따른다.
\end{proof}

\begin{definition}
\label{more-algebra-definition-decomposition-inertia}
$A$를 분수체가 $K$인 이산 값매김환이라 하자.
$L/K$를 갈루아군이 $G$인 유한 갈루아 확대라 하자.
$B$를 $L$ 안에서의 $A$의 정수적 폐포라 하자.
$\mathfrak m \subset B$를 극대아이디얼이라 하자.
\begin{enumerate}
\item $\mathfrak m$의 {\it 분해군}은 부분군
$D = \{\sigma \in G \mid \sigma(\mathfrak m) = \mathfrak m\}$이다.
\item $\mathfrak m$의 {\it 관성군}은 사상
$D \to \text{Aut}(\kappa(\mathfrak m)/\kappa_A)$의 핵 $I$이다.
\end{enumerate}
\end{definition}

\noindent
체 $\kappa(\mathfrak m)$은 $\kappa_A$ 위에서 비분리일 수도 있음에
유의하자. 특히 체 확대 $\kappa(\mathfrak m)/\kappa_A$는 갈루아일
필요가 없다. $\kappa_A$가 완전체이면 갈루아이다.

\begin{lemma}
\label{more-algebra-lemma-galois-galois}
$A$를 분수체가 $K$이고 잉여체가 $\kappa$인 이산 값매김환이라 하자.
$L/K$를 갈루아군이 $G$인 유한 갈루아 확대라 하자.
$B$를 $L$ 안에서의 $A$의 정수적 폐포라 하고, $\mathfrak m$을
$B$의 극대아이디얼이라 하자. 그러면
\begin{enumerate}
\item 체 확대 $\kappa(\mathfrak m)/\kappa$는 정규이고,
\item $D \to \text{Aut}(\kappa(\mathfrak m)/\kappa)$는 전사이다.
\end{enumerate}
어떤(동치로 모든) 극대아이디얼 $\mathfrak m \subset B$에 대하여
체 확대 $\kappa(\mathfrak m)/\kappa$가 분리가능이면,
\begin{enumerate}
\item[(3)] $\kappa(\mathfrak m)/\kappa$는 갈루아이고,
\item[(4)] $D \to \text{Gal}(\kappa(\mathfrak m)/\kappa)$는 전사이다.
\end{enumerate}
여기서 $D \subset G$는 $\mathfrak m$의 분해군이다.
\end{lemma}

\begin{proof}
$A$는 $K$ 안에서 정수적으로 닫혀 있고 $K = L^G$이므로
$A = B^G$임을 관찰하자. 따라서 (1)과 (2)는 보조정리
\ref{more-algebra-lemma-one-orbit-geometric}에서 따른다. 보조정리의 “동치로
모든” 부분은 보조정리 \ref{more-algebra-lemma-galois}에서 따른다.
$\kappa(\mathfrak m)/\kappa$가 분리가능이라고 가정하자.
그러면 (3)과 (4)는 (1)과 (2)에서 즉시 따른다.
\end{proof}

\begin{lemma}
\label{more-algebra-lemma-galois-inertia}
$A$를 분수체가 $K$인 이산 값매김환이라 하자.
$L/K$를 갈루아군이 $G$인 유한 갈루아 확대라 하자.
$B$를 $L$ 안에서의 $A$의 정수적 폐포라 하자.
$\mathfrak m \subset B$를 극대아이디얼이라 하자.
$\mathfrak m$의 관성군 $I$는 표준 완전열
$$
1 \to P \to I \to I_t \to 1
$$
에 들어가며, 다음이 성립한다.
\begin{enumerate}
\item $D$가 분해군이면
$$
P = \{\sigma \in I \mid \sigma|_{\mathfrak m/\mathfrak m^2} =
\text{id}_{\mathfrak m/\mathfrak m^2}\}
=
\{\sigma \in D \mid \sigma \text{ acts trivially on }\text{Gr}_\mathfrak m(B)\}
$$
이다.
\item $P$는 $D$의 정규부분군이다.
\item $\kappa_A$의 표수가 $p > 0$이면 $P$는 $p$-군이고,
$\kappa_A$의 표수가 영이면 $P = \{1\}$이다.
% P02-E0597: frozen source uses p in the characteristic-zero branch without defining a convention.
\item $I_t$는 정수 $e$의 $p$와 서로소인 부분을 위수로 갖는 순환군이다.
\item 표준 동형
$\theta : I_t \to \mu_e(\kappa(\mathfrak m))$가 존재한다.
% P02-E0598: frozen source writes kappa(m), not kappa(mathfrak m).
\item $\kappa(m)$이 $A$의 잉여체 위에서 분리가능이면
$P =
\{\sigma \in D \mid \sigma|_{B/\mathfrak m^2} = \text{id}_{B/\mathfrak m^2}\}$
이다.
\end{enumerate}
여기서 $e$는 보조정리 \ref{more-algebra-lemma-galois-conclusion}의 정수이다.
\end{lemma}

\begin{proof}
보조정리 \ref{more-algebra-lemma-galois-conclusion}에서
$|G| = [L : K] = nef$임을 상기하자.
$G$는 $B$의 극대아이디얼들의 집합
$\{\mathfrak m_1, \ldots, \mathfrak m_n\}$에 추이적으로 작용하고
(보조정리 \ref{more-algebra-lemma-galois}), $D$는 한 원소의 안정자이므로
$|D| = ef$이다. 보조정리 \ref{more-algebra-lemma-galois-galois}에 의해
$$
ef = |D| = |I| \cdot |\text{Aut}(\kappa(\mathfrak m)/\kappa)|
$$
이다. 여기서 $\kappa$는 $A$의 잉여체이다.
$\kappa(\mathfrak m)$은 $\kappa$ 위에서 정규이므로
$\text{Aut}(\kappa(\mathfrak m)/\kappa)$의 위수와 $f$의 차이는
$p$의 거듭제곱 인자뿐이다(Fields, Lemma
\ref{fields-lemma-normal-and-automorphisms}와 Fields, Definition
\ref{fields-definition-insep-degree} 뒤의 논의를 보라).
따라서 $|I|$의 $p$와 서로소인 부분은 $e$의 $p$와 서로소인 부분과 같다.

\medskip\noindent
$C = B_\mathfrak m$이라 두자. 그러면 $I$는 $A$ 위에서 $C$에
작용하고 $C$의 잉여체에는 자명하게 작용한다.
$\pi_A \in A$와 $\pi_C \in C$를 균등화원들이라 하자.
어떤 $C$의 단위원 $u$에 대하여 $\pi_A = u \pi_C^e$라 쓰자.
$\sigma \in I$에 대하여 어떤 $C$의 단위원 $\theta_\sigma$를 써서
$\sigma(\pi_C) = \theta_\sigma \pi_C$라 쓰자. 그러면
% P02-E0599: frozen source omits punctuation after both displayed formulas below.
$$
\pi_A = \sigma(\pi_A) = \sigma(u) (\theta_\sigma \pi_C)^e
= \sigma(u) \theta_\sigma^e \pi_C^e = \frac{\sigma(u)}{u} \theta_\sigma^e \pi_A
$$
이다. $\sigma \in I$이므로
$\sigma(u) \equiv u \bmod \mathfrak m_C$이다. 따라서
$\theta_\sigma$의 $\kappa_C = \kappa(\mathfrak m)$ 안에서의 상
$\overline{\theta}_\sigma$는 1의 $e$제곱근이다.
사상
\begin{equation}
\label{more-algebra-equation-inertia-character}
\theta : I \longrightarrow \mu_e(\kappa(\mathfrak m)),\quad
\sigma \mapsto \overline{\theta}_\sigma
\end{equation}
을 얻는다. $\theta$는 군 준동형이고, 균등화원 $\pi_C$의 선택에
무관하다고 주장한다. 실제로 $\tau$가 $I$의 둘째 원소이면
$\tau(\sigma(\pi_C)) = \tau(\theta_\sigma \pi_C) =
\tau(\theta_\sigma) \theta_\tau \pi_C$이므로
$\theta_{\tau \sigma} = \tau(\theta_\sigma) \theta_\tau$이다.
$\tau \in I$이므로
$\overline{\theta}_{\tau \sigma} =
\overline{\theta}_\sigma \overline{\theta}_\tau$라고 결론내린다.
$\pi'_C$가 둘째 균등화원이면 어떤 $C$의 단위원 $w$에 대하여
$\pi'_C = w \pi_C$이고
$\sigma(\pi'_C) = w^{-1}\sigma(w)\theta_\sigma \pi'_C$임을 알 수 있다.
따라서 $\theta'_\sigma = w^{-1}\sigma(w)\theta_\sigma$이고,
앞서와 같이 $\theta'_\sigma$와 $\theta_\sigma$는 잉여체의 같은 원소로 간다.

\medskip\noindent
% P02-E0597: frozen source again uses p unconditionally here.
$\kappa(\mathfrak m)$의 표수는 $p$이므로 군
$\mu_e(\kappa(\mathfrak m))$는 $e$의 $p$와 서로소인 부분 이하의
위수를 갖는 순환군이다(Fields, Section
\ref{fields-section-roots-of-1}을 보라).

\medskip\noindent
$P = \Ker(\theta)$라 두자. 구성에 의해 $P$의 원소들은 정확히
$\mathfrak m/\mathfrak m^2 = \text{Gr}_\mathfrak m^1(B)$에 자명하게
작용하는 $I$의 원소들이다. 즉,
$P = \{\sigma \in I \mid \sigma|_{\mathfrak m/\mathfrak m^2} =
\text{id}_{\mathfrak m/\mathfrak m^2}\}$이다. 또한 $I$는
$\kappa(\mathfrak m) = B/\mathfrak m$에 자명하게 작용하는
$D$의 원소들로 이루어진다. 등급환 $\text{Gr}_\mathfrak m(B)$는
$\text{Gr}^0_\mathfrak m(B) = B/\mathfrak m$ 위에서
$\text{Gr}^1_\mathfrak m(B) = \mathfrak m/\mathfrak m^2$에 의해
생성되므로
$P = \{\sigma \in D \mid \sigma \text{ acts trivially on }
\text{Gr}_\mathfrak m(B)\}$라고 결론내린다.
따라서 (1)이 성립한다. $P$는 준동형
$D \to \text{Aut}(\text{Gr}_\mathfrak m(B))$의 핵이므로 이는 (2)를
함의한다. (3)을 증명할 수 있다면, 위의 논의가
% P02-E1252: frozen source reverses the image/subgroup order inequality.
$|I_t| \geq |\mu_e(\kappa(\mathfrak m))|$임을 보이므로
$I_t$는 $\mu_e(\kappa(\mathfrak m))$와 동형이 되고 (4)와 (5)가 따른다.

\medskip\noindent
따라서 $\theta$의 핵 $P$가 $p$-군임을 증명하면 충분하다.
핵의 자명하지 않은 원소를 $\sigma$라 하자. 그러면 모든 $i$에 대하여
$\sigma - \text{id}$는 $\mathfrak m_C^i$를
$\mathfrak m_C^{i + 1}$ 안으로 보낸다. $\sigma$의 위수를 $m$이라
하자. $\sigma(c) \not = c$인 $c \in C$를 택하자. 그러면 어떤 $i$에
대하여 $\sigma(c) - c \in \mathfrak m_C^i$이고
$\sigma(c) - c \not \in \mathfrak m_C^{i + 1}$이며,
\begin{align*}
0
& =
\sigma^m(c) - c \\
& =
\sigma^m(c) - \sigma^{m - 1}(c) + \ldots + \sigma(c) - c \\
& =
\sum\nolimits_{j = 0, \ldots, m - 1} \sigma^j(\sigma(c) - c) \\
& \equiv
m(\sigma(c) - c) \bmod \mathfrak m_C^{i + 1}
\end{align*}
를 얻는다.
% P02-E0600: frozen characteristic-zero alternative incorrectly asserts m = 0.
$p | m$이다(또는 $p = 1$이면 $m = 0$이다). 따라서 $\theta$의 핵의
모든 원소의 위수는 $p$로 나누어지며, 즉 $\Ker(\theta)$는 $p$-군이다.

\medskip\noindent
(6)의 증명. $\kappa(\mathfrak m)/\kappa$가 분리가능이라고 가정하자.
이 경우 보조정리 \ref{more-algebra-lemma-galois-galois}에 의해 $f = |D/I|$이고
$I$의 위수는 $e$이다. $e = 1$이면 $P = I = \{\text{id}\}$이고
결론이 성립한다. $e > 1$이면
$B/\mathfrak m^2 = C/\pi_C^2C$는 $\kappa$-대수이고
$\kappa(\mathfrak m)$의 작은 확대이다.
$\kappa(\mathfrak m)$은 $\kappa$ 위에서 형식적으로 에탈이므로
(Algebra, Lemma
\ref{algebra-lemma-characterize-separable-algebraic-field-extensions}),
$B/\mathfrak m^2 \to \kappa(\mathfrak m)$의 오른쪽 역인 유일한
$\kappa$-대수 사상
$\kappa(\mathfrak m) \to B/\mathfrak m^2$가 존재한다.
달리 말하면, $\kappa$-대수의 $D$-등변 동형
$B/\mathfrak m^2 = \mathfrak m/\mathfrak m^2 \oplus \kappa(\mathfrak m)$
이 존재한다. 이는 이 경우
$P =
\{\sigma \in D \mid \sigma|_{B/\mathfrak m^2} = \text{id}_{B/\mathfrak m^2}\}$
임을 즉시 보인다.
\end{proof}

\begin{example}
\label{more-algebra-example-counter-description-P}
보조정리 \ref{more-algebra-lemma-galois-inertia}의 (6)의 등식은
$\kappa(\mathfrak m)$에 대한 가정 없이는 거짓이다.
반례는 $A = Z_2[t]_{(2)}$의 확대 $B = A[x]/(x^2 - t)$로 주어진다.
실제로 이때 $\sigma(x) = -x = x - 2x$이고
$2x$는 $\mathfrak m^2$에 속하지 않는다.
\end{example}

\begin{definition}
\label{more-algebra-definition-wild-inertia}
% P02-E0601: frozen source opens this definition with a sentence fragment.
가정과 표기는 보조정리 \ref{more-algebra-lemma-galois-inertia}와 같다고 하자.
\begin{enumerate}
\item $\mathfrak m$의 {\it 야생 관성군}은 부분군 $P$이다.
% P02-E1255: frozen source defines the tame inertia group as a quotient map, not the quotient group.
\item $\mathfrak m$의 {\it 순관성군}은 몫 $I \to I_t$이다.
\end{enumerate}
% P02-E0602: frozen source omits “by” in the notation assignment.
핵이 $P$이고 동형 $I_t \to \mu_e(\kappa(\mathfrak m))$를 유도하는
전사 (\ref{more-algebra-equation-inertia-character})를
$\theta : I \to \mu_e(\kappa(\mathfrak m))$로 나타낸다.
\end{definition}

\begin{lemma}
\label{more-algebra-lemma-inertia-character}
가정과 표기는 보조정리 \ref{more-algebra-lemma-galois-inertia}와 같다고 하자.
관성 지표 $\theta : I \to \mu_e(\kappa(\mathfrak m))$는 다음 성질을
만족한다.
% P02-E0603: frozen source does not punctuate the display introduction or conclusion.
$$
\theta(\tau \sigma \tau^{-1}) = \tau(\theta(\sigma))
$$
여기서 $\tau \in D$이고 $\sigma \in I$이다.
\end{lemma}

\begin{proof}
$I$는 $D$의 정규부분군이고, 예를 들어 보조정리
\ref{more-algebra-lemma-galois-galois}에서 논의한 사상
$D \to \text{Aut}(\kappa(\mathfrak m))$을 통하여
$\tau$가 $\kappa(\mathfrak m)$에 작용하므로 공식은 의미가 있다.
$\theta$의 구성을 상기하자. $B_\mathfrak m$의 균등화원 $\pi$를
택하고, $\sigma \in I$에 대하여
$\sigma(\pi) = \theta_\sigma \pi$라 쓰자. 그러면
$\theta(\sigma)$는 $\theta_\sigma$의 잉여체 안에서의 상
$\overline{\theta}_\sigma$이다. 모든 $\tau \in D$에 대하여 어떤
단위원 $\theta_\tau$를 써서 $\tau(\pi) = \theta_\tau \pi$라 쓸 수
있다. 그러면 $\theta_{\tau^{-1}} = \tau^{-1}(\theta_\tau^{-1})$이다.
계산하면
% P02-E0604: frozen aligned computation lacks terminal punctuation.
\begin{align*}
\theta_{\tau \sigma \tau^{-1}}
& =
\tau(\sigma(\tau^{-1}(\pi)))/\pi \\
& =
\tau(\sigma(\tau^{-1}(\theta_\tau^{-1}) \pi))/\pi \\
& =
\tau(\sigma(\tau^{-1}(\theta_\tau^{-1})) \theta_\sigma \pi)/\pi \\
& =
\tau(\sigma(\tau^{-1}(\theta_\tau^{-1}))) \tau(\theta_\sigma) \theta_\tau
\end{align*}
를 얻는다. 그러나 $\sigma$는 $\pi$를 법으로 자명하게 작용하므로
곱 $\tau(\sigma(\tau^{-1}(\theta_\tau^{-1}))) \theta_\tau$는
잉여체에서 $1$로 간다. 이로써 보조정리가 증명된다.
\end{proof}

\noindent
다음 보조정리는 Fundamental Groups, Lemma
\ref{pione-lemma-inertial-invariants-unramified}에서 일반화할 것이다.

\begin{lemma}
\label{more-algebra-lemma-inertial-invariants-unramified}
$A$를 분수체가 $K$인 이산 값매김환이라 하자.
$L/K$를 유한 갈루아 확대라 하자. $\mathfrak m \subset B$를
$L$ 안에서의 $A$의 정수적 폐포의 극대아이디얼이라 하자.
$I \subset G$를 $\mathfrak m$의 관성군이라 하자.
그러면 $B^I$는 $L^I$ 안에서의 $A$의 정수적 폐포이고,
$A \to (B^I)_{B^I \cap \mathfrak m}$는 에탈이다.
\end{lemma}

\begin{proof}
$B' = B^I$라 쓰자. 정의들로부터 $B' = B^I$는 $L^I$ 안에서의
$A$의 정수적 폐포이다. 또한
$\mathfrak m' = B^I \cap \mathfrak m = B' \cap \mathfrak m \subset B'$라
쓰자. 보조정리 \ref{more-algebra-lemma-one-orbit}에 의해 $\mathfrak m$은
$\mathfrak m'$ 위에 놓이는 $B$의 유일한 소아이디얼이다.
$I$가 $\kappa(\mathfrak m)$에 자명하게 작용하므로 보조정리
\ref{more-algebra-lemma-invariants-modulo}에서 확대
$\kappa(\mathfrak m)/\kappa(\mathfrak m')$가 순수 비분리임을 알 수
있다(아마 더 쉬운 대안은 보조정리
\ref{more-algebra-lemma-one-orbit-geometric}의 결과를 적용하는 것이다).
$D/I$는 $\kappa(\mathfrak m')$에 충실하게 작용하므로
$D/I$가 $\kappa(\mathfrak m)$에 충실하게 작용한다고 결론내린다.
물론 $A$의 잉여체 $\kappa$의 원소들은 이 작용에 의해 고정된다.
갈루아 이론에 의해
$[\kappa(\mathfrak m') : \kappa] \geq |D/I|$이다. Fields, Lemma
\ref{fields-lemma-galois-over-fixed-field}를 보라.

\medskip\noindent
$\pi$를 $A$의 균등화원이라 하자.
$\text{Norm}_{L/K}(\pi) = \pi^{[L : K]}$이므로 Algebra, Lemma
\ref{algebra-lemma-finite-extension-dim-1}에서
% P02-E0605: frozen source omits punctuation after this and two later valuation displays.
$$
|G| = [L : K] = [L : K]\ \text{ord}_A(\pi) =
|G/D|\ [\kappa(\mathfrak m) : \kappa]\ \text{ord}_{B_\mathfrak m}(\pi)
$$
를 얻는다. 이는 $B$에는 모두 $G$ 아래에서 켤레인
$n = |G/D|$개의 극대아이디얼이 있기 때문이다. 비고
\ref{more-algebra-remark-finite-separable-extension}와 보조정리
\ref{more-algebra-lemma-galois}를 보라.
% P02-E0606: frozen source duplicates “the finite extension”.
같은 논의를 차수가 $|I|$인 유한 확대 $L/L^I$에 적용하면
$$
|I|\ \text{ord}_{B'_{\mathfrak m'}}(\pi) =
[\kappa(\mathfrak m) : \kappa(\mathfrak m')]\ \text{ord}_{B_\mathfrak m}(\pi)
$$
를 얻는다. 따라서
% P02-E0605: frozen source again omits punctuation at the display boundary.
$$
\text{ord}_{B'_{\mathfrak m'}}(\pi) =
\frac{|D/I|}{[\kappa(\mathfrak m') : \kappa]}
$$
이다. 왼쪽은 양의 정수이고 위의 결과에 의해 오른쪽은 $\leq 1$이므로,
등식이 성립하고 $\text{ord}_{B'_{\mathfrak m'}}(\pi) = 1$이며
$\kappa(\mathfrak m')/\kappa$의 차수는 $|D/I|$이다.
% P02-E0607: frozen source localizes B, not B prime, at m prime on the right.
따라서 $\pi B'_{\mathfrak m'} = \mathfrak m' B_\mathfrak m'$이고,
$\kappa(\mathfrak m')$은 $\kappa$ 위 갈루아이며 그 갈루아군은
$D/I$이다. 특히 분리가능이다. Fields, Lemma
\ref{fields-lemma-finite-Galois}를 보라. Algebra, Lemma
\ref{algebra-lemma-characterize-etale}에 의해 원하는 대로
$A \to B'_{\mathfrak m'}$는 에탈이다.
\end{proof}

\begin{remark}
\label{more-algebra-remark-tower-of-rings}
$A$를 분수체가 $K$인 이산 값매김환이라 하자.
$L/K$를 유한 갈루아 확대라 하자. $\mathfrak m \subset B$를
$L$ 안에서의 $A$의 정수적 폐포의 극대아이디얼이라 하자.
% P02-E0608: frozen source names three groups with singular grammar and leaves both displays unpunctuated.
$\mathfrak m$의 야생 관성군, 관성군, 분해군을 각각
$$
P \subset I \subset D \subset G
$$
라 하자. 도식
$$
\xymatrix{
\mathfrak m \ar@{-}[d] \ar@{-}[r] &
\mathfrak m^P \ar@{-}[d] \ar@{-}[r] &
\mathfrak m^I \ar@{-}[d] \ar@{-}[r] &
\mathfrak m^D \ar@{-}[d] \ar@{-}[r] &
A \cap \mathfrak m \ar@{-}[d] \\
B & B^P \ar[l] & B^I \ar[l] & B^D \ar[l] & A \ar[l]
}
$$
을 생각하자. $B^P, B^I, B^D$는 각각 체 $L^P$, $L^I$, $L^D$
안에서의 $A$의 정수적 폐포이다. 따라서 $B^P$가 $L^P$ 안에서의
$B^I$의 정수적 폐포라는 것 등도 알 수 있다.
$\mathfrak m^P = \mathfrak m \cap B^P$,
$\mathfrak m^I = \mathfrak m \cap B^I$, 그리고
$\mathfrak m^D = \mathfrak m \cap B^D$임을 관찰하자.
따라서 도식의 윗줄은 스펙트럼들의 유도된 사상 아래에서
$\mathfrak m \in \Spec(B)$의 상들에 대응한다.
이 모든 것을 말했으므로 다음이 성립한다.
\begin{enumerate}
\item 확대 $L^I/L^D$는 군 $D/I$를 갖는 갈루아 확대이다.
\item 확대 $L^P/L^I$는 군 $I_t = I/P$를 갖는 갈루아 확대이다.
\item 확대 $L^P/L^D$는 군 $D/P$를 갖는 갈루아 확대이다.
\item $\mathfrak m^I$는 $\mathfrak m^D$ 위에 놓이는 $B^I$의 유일한
소아이디얼이다.
\item $\mathfrak m^P$는 $\mathfrak m^I$ 위에 놓이는 $B^P$의 유일한
소아이디얼이다.
\item $\mathfrak m$은 $\mathfrak m^P$ 위에 놓이는 $B$의 유일한
소아이디얼이다.
\item $\mathfrak m^P$는 $\mathfrak m^D$ 위에 놓이는 $B^P$의 유일한
소아이디얼이다.
\item $\mathfrak m$은 $\mathfrak m^I$ 위에 놓이는 $B$의 유일한
소아이디얼이다.
\item $\mathfrak m$은 $\mathfrak m^D$ 위에 놓이는 $B$의 유일한
소아이디얼이다.
\item $A \to B^D_{\mathfrak m^D}$는 에탈이고 자명한 잉여체 확대를
유도한다.
\item $B^D_{\mathfrak m^D} \to B^I_{\mathfrak m^I}$는 에탈이고,
갈루아군이 $D/I$인 잉여체들의 갈루아 확대를 유도한다.
\item $A \to B^I_{\mathfrak m^I}$는 에탈이다.
% P02-E0609: frozen remark uses p without defining it, including in characteristic zero.
\item $B^I_{\mathfrak m^I} \to B^P_{\mathfrak m^P}$는 $p$와 서로소인
분기지수 $|I/P|$를 갖고 자명한 잉여체 확대를 유도한다.
\item $B^D_{\mathfrak m^D} \to B^P_{\mathfrak m^P}$는 $p$와 서로소인
분기지수 $|I/P|$를 갖고 분리가능한 잉여체 확대를 유도한다.
\item $A \to B^P_{\mathfrak m^P}$는 $p$와 서로소인 분기지수
$|I/P|$를 갖고 분리가능한 잉여체 확대를 유도한다.
\end{enumerate}
(1), (2), (3)은 갈루아 이론(Fields, Section
\ref{fields-section-galois-theory})과 보조정리
\ref{more-algebra-lemma-galois-inertia}에서 즉시 따른다.
(4)--(9)는 보조정리 \ref{more-algebra-lemma-galois}에서 명백하다.
(12)는 보조정리 \ref{more-algebra-lemma-inertial-invariants-unramified}이다.
분해 $A \to B^D_{\mathfrak m^D} \to B^I_{\mathfrak m^I}$가 있으므로
(10)과 (11)의 에탈성은 그 결과로 얻어진다.
(10)의 잉여체 확대는 자명해야 한다. 보조정리
\ref{more-algebra-lemma-galois-galois}에서 보였듯이 이 확대는 분리가능이고
$D/I$가 $\text{Aut}(\kappa(\mathfrak m)/\kappa_A)$로 전사하기
때문이다. 같은 논증이 (11)의 잉여체에 관한 명제를 증명한다.
(13), (14), (15)를 보이려면 (13)을 증명하는 것으로 충분하다.
위에서 본 대로 확대 $L^P/L^I$는 $p$와 서로소인 위수의 순환
갈루아군을 갖는 갈루아 확대이고, 소아이디얼 $\mathfrak m^P$는
$\mathfrak m^I$ 위에 놓이는 유일한 소아이디얼이며,
$I/P$의 잉여체 위 작용은 자명하다. 따라서 이 확대와 이산 값매김환
$B^I_{\mathfrak m^I}$에 보조정리 \ref{more-algebra-lemma-galois-inertia}를 적용하면
(13)이 성립함을 알 수 있다.
\end{remark}

\begin{lemma}
\label{more-algebra-lemma-compare-inertia}
$A$를 분수체가 $K$인 이산 값매김환이라 하자.
$M/L/K$를 $M/K$와 $L/K$가 유한 갈루아인 탑이라 하자.
% P02-E0610: frozen source uses singular “integral closure” for two respective rings.
$C$와 $B$를 각각 $M$과 $L$ 안에서의 $A$의 정수적 폐포라 하자.
$\mathfrak m' \subset C$를 극대아이디얼이라 하고
$\mathfrak m = \mathfrak m' \cap B$라 두자.
$$
P \subset I \subset D \subset \text{Gal}(L/K)
\quad\text{and}\quad
P' \subset I' \subset D' \subset \text{Gal}(M/K)
$$
를 각각 $\mathfrak m$과 $\mathfrak m'$의 야생 관성군, 관성군,
분해군이라 하자. 그러면 표준 전사
$\text{Gal}(M/K) \to \text{Gal}(L/K)$는 전사
$P' \to P$, $I' \to I$, 그리고 $D' \to D$를 유도한다. 더욱이 이들은
가환도식들
$$
\vcenter{
\xymatrix{
D' \ar[r] \ar[d] &
\text{Aut}(\kappa(\mathfrak m')/\kappa_A) \ar[d] \\
D \ar[r] &
\text{Aut}(\kappa(\mathfrak m)/\kappa_A)
}
}
\quad\text{and}\quad
\vcenter{
\xymatrix{
I' \ar[r]_-{\theta'} \ar[d] &
\mu_{e'}(\kappa(\mathfrak m')) \ar[d]^{(-)^{e'/e}} \\
I \ar[r]^-\theta &
\mu_e(\kappa(\mathfrak m))
}
}
$$
에 들어간다. 여기서 $e'$과 $e$는
$A \to C_{\mathfrak m'}$와 $A \to B_\mathfrak m$의 분기지수이다.
\end{lemma}

\begin{proof}
사상 $\text{Gal}(M/K) \to \text{Gal}(L/K)$ 아래에서 군
$P', I', D'$이 각각 $P, I, D$ 안으로 가는 것은 이 군들의 정의에서
즉시 따른다. 첫째 도식의 가환성은 명백하다
($\kappa(\mathfrak m)/\kappa_A$가 정규이므로
$\kappa_A$ 위 $\kappa(\mathfrak m')$의 모든 자기동형이 실제로
$\kappa_A$ 위 $\kappa(\mathfrak m)$의 자기동형을 유도한다.
따라서 첫째 도식의 오른쪽 수직 화살표를 얻는다. 보조정리
\ref{more-algebra-lemma-galois-galois}와 Fields, Lemma
\ref{fields-lemma-lift-maps}를 보라).

\medskip\noindent
사상 $I' \to I$와 $D' \to D$는 보조정리
\ref{more-algebra-lemma-one-orbit-geometric-galois-compare}에 의해 전사이다.
$P'$과 $P$는 각각 $I'$과 $I$의 p-실로우 부분군이므로
$P' \to P$도 전사이다.

\medskip\noindent
둘째 도식의 가환성을 보이기 위해 $C_{\mathfrak m'}$의 균등화원
$\pi'$과 $B_\mathfrak m$의 균등화원 $\pi$를 택하자.
그러면 어떤 $C_{\mathfrak m'}$의 단위원 $c'$에 대하여
$\pi = c' (\pi')^{e'/e}$이다. $\sigma' \in I'$에 대하여
상 $\sigma \in I$는 단순히 $\sigma'$의 $L$에 대한 제한이다.
단위원 $c \in C_{\mathfrak m'}$에 대하여
$\sigma'(\pi') = c \pi'$라 쓰고, $B_\mathfrak m$의 단위원 $b$에
대하여 $\sigma(\pi) = b \pi$라 쓰자. 그러면
$\sigma'(\pi) = b \pi$이고
% P02-E0611: frozen displayed computation lacks terminal punctuation.
$$
b \pi = \sigma'(\pi) = \sigma'(c' (\pi')^{e'/e}) =
\sigma'(c') c^{e'/e} (\pi')^{e'/e} =
\frac{\sigma'(c')}{c'} c^{e'/e} \pi
$$
를 얻는다. $\sigma' \in I'$이므로 $b$와 $c^{e'/e}$는 잉여체에서
같은 상을 가지며, 이것이 원하는 바를 증명한다.
\end{proof}

\begin{remark}
\label{more-algebra-remark-canonical-inertia-character}
무한 갈루아 확대에 관성 지표
$\theta : I \to \mu_e(\kappa(\mathfrak m))$를 사용하려면 이를
스케일하는 것이 편리하다.
$A, K, L, B, \mathfrak m, G, P, I, D, e, \theta$가 보조정리
\ref{more-algebra-lemma-galois-inertia}와 정의 \ref{more-algebra-definition-wild-inertia}에서와
같다고 하자.
% P02-E0612: frozen source says “with q is a power” rather than “where q is”.
그러면 $e = q |I_t|$이며, $\kappa(\mathfrak m)$의 표수 $p$가 양수이면
$q$는 그 거듭제곱이고 영이면 $1$이다.
$\kappa(\mathfrak m)$의 표수가 $p$이므로
$\mu_e(\kappa(\mathfrak m)) = \mu_{|I_t|}(\kappa(\mathfrak m))$임에
유의하자. 사상
% P02-E0613: frozen source leaves this map display and the final diagram unpunctuated.
$$
\theta_{can} = q\theta : I \longrightarrow \mu_{|I_t|}(\kappa(\mathfrak m))
$$
을 생각하자. 이 사상은 동형
$\theta_{can} : I_t \to \mu_{|I_t|}(\kappa(\mathfrak m))$를 유도한다.
보조정리 \ref{more-algebra-lemma-inertia-character}에 의해
$\tau \in D$와 $\sigma \in I$에 대하여
$\theta_{can}(\tau \sigma \tau^{-1}) = \tau(\theta_{can}(\sigma))$이다.
마지막으로 $M/L$이 $M/K$가 갈루아인 확대이고,
$\mathfrak m'$이 $\mathfrak m$ 위에 놓이는 $M$ 안에서의 $A$의 정수적
폐포의 소아이디얼이면 가환도식
$$
\xymatrix{
I' \ar[r]_-{\theta'_{can}} \ar[d] &
\mu_{|I'_t|}(\kappa(\mathfrak m')) \ar[d]^{(-)^{|I'_t|/|I_t|}} \\
I \ar[r]^-{\theta_{can}} &
\mu_{|I_t|}(\kappa(\mathfrak m))
}
$$
을 보조정리 \ref{more-algebra-lemma-compare-inertia}에서 얻는다.
\end{remark}

\section{크라스너 보조정리}
\label{more-algebra-section-krasner}

\noindent
이산 값매김체의 경우의 크라스너 보조정리는 다음과 같다.

\begin{lemma}[크라스너 보조정리]
\label{more-algebra-lemma-krasner}
$A$를 차원이 $1$인 완비 국소 정역이라 하자.
$P(t) \in A[t]$를 계수가 $A$에 속하는 다항식이라 하자.
$\alpha \in A$가 $P$의 근이지만 도함수
$P' = \text{d}P/\text{d}t$의 근은 아니라고 하자.
모든 $c \geq 0$에 대하여 다음 성질을 갖는 정수 $n$이 존재한다.
계수들이 $\mathfrak m_A^n$에 속하는 모든 $Q \in A[t]$에 대하여
다항식 $P + Q$는
$\beta - \alpha \in \mathfrak m_A^c$인 근 $\beta \in A$를 갖는다.
\end{lemma}

\begin{proof}
% P02-E0614: frozen proof uses unindexed mathfrak m although the statement defines mathfrak m_A.
영이 아닌 $\pi \in \mathfrak m$을 택하자. $A$의 차원은 $1$이므로
$\mathfrak m = \sqrt{(\pi)}$이다. 가정에 의해 어떤 $m \geq 0$과
$a \in A$에 대하여 $P'(\alpha)^{-1} = \pi^{-m} a$라 쓸 수 있다.
$c \geq m + 1$이라고 가정해도 좋고 그렇게 하자.
$\mathfrak m_A^n \subset (\pi^{c + m})$이 되는 $n$을 택하자.
명제에서와 같은 임의의 $Q$를 택하자. 뒤에서 사용하기 위해 어떤
$R(x, y) \in A[x, y]$에 대하여
% P02-E0618: frozen source leaves the five Taylor/congruence displays unpunctuated.
$$
P(x + y) = P(x) + P'(x)y + R(x, y)y^2
$$
로 쓸 수 있음을 관찰하자.
% P02-E0615: frozen induction starts at alpha_m although the stated invariants need the stronger c-indexed start.
귀납법으로 다음을 만족하는 수열
$\alpha_m, \alpha_{m + 1}, \alpha_{m + 2}, \ldots$를 찾을 수 있음을
보이겠다.
\begin{enumerate}
\item $\alpha_k \equiv \alpha \bmod \pi^c$,
\item $\alpha_{k + 1} - \alpha_k \in (\pi^k)$,
\item $(P + Q)(\alpha_k) \in (\pi^{m + k})$.
\end{enumerate}
$\beta = \lim \alpha_k$라 두면 증명이 끝난다.

\medskip\noindent
기초 단계. $Q$의 계수들은 $(\pi^{c + m})$에 속하므로
$(P + Q)(\alpha) \in (\pi^{c + m})$이다. 따라서
$\alpha_m = \alpha$로 둘 수 있다. 이 선택은 모든 $k \geq m$에
대하여 $\alpha_k \equiv \alpha \bmod \pi^c$임을 보장한다.

\medskip\noindent
귀납 단계. $\alpha_k$가 주어졌다고 하자. 어떤
$\delta \in (\pi^k)$에 대하여
$\alpha_{k + 1} = \alpha_k + \delta$라 쓰자. 그러면
$$
(P + Q)(\alpha_{k + 1}) =
P(\alpha_k + \delta) + Q(\alpha_k + \delta)
$$
이다. $Q$의 계수들이 $(\pi^{c + m})$에 속하므로
$Q(\alpha_k + \delta) \equiv Q(\alpha_k) \bmod \pi^{c + m + k}$이다.
한편
$$
P(\alpha_k + \delta) =
P(\alpha_k) + P'(\alpha_k)\delta + R(\alpha_k, \delta)\delta^2
$$
이다. $\alpha_k \equiv \alpha \bmod \pi^{m + 1}$이므로
$P'(\alpha_k) \equiv P'(\alpha) \bmod (\pi^{m + 1})$임에 유의하자.
따라서
$$
P(\alpha_k + \delta) \equiv P(\alpha_k) + P'(\alpha) \delta
\bmod \pi^{k + m + 1}
$$
를 얻는다. 두 항을 다시 합치면
$$
(P + Q)(\alpha_{k + 1}) \equiv (P + Q)(\alpha_k) + P'(\alpha) \delta
\bmod \pi^{k + m + 1}
$$
임을 알 수 있다. 따라서
$\delta = -P'(\alpha)^{-1} (P + Q)(\alpha_k) =  - \pi^{-m} a (P + Q)(\alpha_k)$
로 택하면 된다. 귀납 가정에 의해 이는 $(\pi^k)$에 속한다.
\end{proof}

\begin{lemma}
\label{more-algebra-lemma-approximate-separable-extension}
$A$를 분수체가 $K$인 이산 값매김환이라 하자.
$A^\wedge$를 $A$의 완비화라 하고 그 분수체를 $K^\wedge$라 하자.
$M/K^\wedge$가 유한 분리 확대이면, 어떤 유한 분리 확대 $L/K$가
존재하여 $M = K^\wedge \otimes_K L$이다.
\end{lemma}

\begin{proof}
$A^\wedge$도 이산 값매김환임에 유의하자(보조정리
\ref{more-algebra-lemma-completion-regular}와 \ref{more-algebra-lemma-completion-dimension}에
의한다). 특히 $A^\wedge$는 정역이다. 이 증명은 더 일반적으로
$A^\wedge$가 차원 $1$인 국소 정역인 뇌터 국소환 $A$에 대해서도
성립한다.

\medskip\noindent
$\theta \in M$을 $K^\wedge$ 위에서 $M$을 생성하는 원소라 하자
(원시원소 정리). $P(t) \in K^\wedge[t]$를 $K^\wedge$ 위
$\theta$의 최소다항식이라 하자.
$\pi \in \mathfrak m_A$를 영이 아닌 원소라 하자.
$\theta$를 $\pi^n\theta$로 바꾼 뒤에는 $P(t)$의 계수들이
$A^\wedge$에 속한다고 가정할 수 있다.
$B = A^\wedge[\theta] = A^\wedge[t]/(P(t))$라 하자.
% P02-E0616: frozen source says B is finite over A, not over A-wedge.
$B$는 $A$ 위에서 유한이고 $M$에 포함되므로 차원 $1$인 완비 국소
정역임에 유의하자.
% P02-E0617: frozen source invokes separability over K, not the stated K-wedge.
$M$은 $K$ 위에서 분리가능이므로 원소 $\theta$는 $P$의 도함수의
근이 아니다. 모든 정수 $n$에 대하여 $P - P_1$의 계수들이
$\pi^nA^\wedge[t]$에 속하는 일계수 다항식 $P_1 \in A[t]$를 찾을 수
있다. 크라스너 보조정리(보조정리 \ref{more-algebra-lemma-krasner})에 의해
$n$이 충분히 크면 $P_1$은 $B$ 안에 근 $\beta$를 갖는다.
더욱이 $n$을 충분히 크게 택하면
$\theta - \beta \in \pi B$라고 가정할 수 있다.
$t$를 $\beta$로 보내는 $A^\wedge$-대수 사상
$\Phi : A^\wedge[t]/(P_1) \to B$를 생각하자.
$B = \pi B + \sum_{i < \deg(P)} A^\wedge \theta^i$이므로 나카야마
보조정리에 의해 사상 $\Phi$는 전사이다.
$\deg(P_1) = \deg(P)$이므로 $\Phi$는 동형이다. 따라서 환 확대
$L = K[t]/(P_1(t))$는 $K^\wedge \otimes_K L \cong M$을 만족한다.
이는 $L$이 체임을 함의하며 증명이 끝난다.
\end{proof}

\begin{definition}
\label{more-algebra-definition-mixed}
$A$를 이산 값매김환이라 하자. $A$의 잉여체의 표수가 $p > 0$이고
$A$의 분수체의 표수가 $0$이면 $A$가 {\it 혼합 표수}를 갖는다고 한다.
이 경우 이산 값매김환의 확대 $\mathbf{Z}_{(p)} \subset A$를 얻으며,
$A$의 {\it 절대 분기지수}는 이 확대의 분기지수이다.
\end{definition}

\section{아브얀카르 보조정리와 순분기}
\label{more-algebra-section-abhyankar-tame}

\noindent
이 절에서는 이산 값매김환에 대한 아브얀카르 보조정리의 가장 일반적인
형태라고 생각되는 것을 증명한다. 그 뒤 이를 적용하여 이산 값매김환의
분수체의 순분기 확대에 관한 몇 가지 결과를 증명한다.

\begin{remark}
\label{more-algebra-remark-construction}
$A \to B$를 분수체들이 $K \subset L$인 이산 값매김환의 확대라 하자.
$K_1/K$를 유한 체 확대라 하자.
$A_1 \subset K_1$을 $K_1$ 안에서의 $A$의 정수적 폐포라 하자.
한편 $L_1 = (L \otimes_K K_1)_{red}$라 하자. 그러면 $L_1$은
$L$의 유한 체 확대들의 공집합이 아닌 유한곱이다.
$B_1$을 $L_1$ 안에서의 $B$의 정수적 폐포라 하자.
서로 호환되는 가환도식들
% P02-E0619: frozen source leaves this and three later displays unpunctuated.
$$
\vcenter{
\xymatrix{
L \ar[r] & L_1 \\
K \ar[u] \ar[r] & K_1 \ar[u]
}
}
\quad\text{and}\quad
\vcenter{
\xymatrix{
B \ar[r] & B_1 \\
A \ar[u] \ar[r] & A_1 \ar[u]
}
}
$$
을 얻는다. 이 상황에서는 다음이 성립한다.
\begin{enumerate}
\item Algebra, Lemma \ref{algebra-lemma-integral-closure-Dedekind}에 의해
환 $A_1$은 데데킨트 정역이고 $B_1$은 데데킨트 정역들의 유한곱이다.
\item $\pi \in A$가 균등화원이면
$L \otimes_K K_1 = (B \otimes_A A_1)_\pi$이고,
$\pi$는 $B \otimes_A A_1$ 위 영인자가 아님에 유의하자.
따라서 환 사상 $B \otimes_A A_1 \to B_1$은 정수적이고 그 핵은
멱영원소들로 이루어진다. 그러므로
$\Spec(B_1) \to \Spec(B \otimes_A A_1)$은 스펙트럼 위에서 전사이다
(Algebra, Lemma
\ref{algebra-lemma-integral-overring-surjective}).
$A_1/\mathfrak m_A A_1$은 아르틴이고 환 사상
$A_1/\mathfrak m_A A_1 \to
B/\mathfrak m_AB \otimes_{\kappa_A} A_1/\mathfrak m_A A_1$이
단사이므로 사상 $\Spec(B \otimes_A A_1) \to \Spec(A_1)$은 전사이다.
따라서 $\Spec(B_1) \to \Spec(A_1)$이 전사라고 결론내린다.
% P02-E0626: frozen index range omits the comma in “1, \ldots n”.
\item $\mathfrak m_i$, $i = 1, \ldots n$을 $A_1$의 극대아이디얼들이라
하자. 여기서 $n \geq 1$이다. 각 $i = 1, \ldots, n$에 대하여
$\mathfrak m_{ij}$, $j = 1, \ldots, m_i$를 $\mathfrak m_i$ 위에 놓이는
$B_1$의 극대아이디얼들이라 하자. 여기서 $m_i \geq 1$이다.
이산 값매김환 확대들의 도식
% P02-E0619: frozen source leaves this diagram unpunctuated.
$$
\xymatrix{
B \ar[r] & (B_1)_{\mathfrak m_{ij}} \\
A \ar[u] \ar[r] & (A_1)_{\mathfrak m_i} \ar[u]
}
$$
을 얻는다.
\item $A$가 헨젤이면(예를 들어 완비이면) $A_1$은 이산 값매김환이다.
즉, $n = 1$이다. 실제로 $A_1$은 정역인 $A$의 유한 확대들의 합집합이므로
Algebra, Lemma \ref{algebra-lemma-finite-over-henselian}에 의해
국소환이다.
\item $B$가 헨젤이면(예를 들어 완비이면) $B_1$은 이산 값매김환들의
곱이다. 즉, $i = 1, \ldots, n$에 대하여 $m_i = 1$이다.
\item $K \subset K_1$이 순수 비분리이면 $A_1$과 $B_1$은 모두 이산
값매김환이다. 즉, $n = 1$이고 $m_1 = 1$이다.
% P02-E0620: frozen source duplicates the word “power”.
이는 모든 $b \in B_1$에 대하여 $b$의 어떤 $p$-거듭제곱이 $B$에
속하고, 따라서 $B_1$은 하나의 극대아이디얼만 가질 수 있기 때문이다.
\item $K \subset K_1$이 유한 분리이면
$L_1 = L \otimes_K K_1$이고, 이것도 유한 분리 확대들의 유한곱이다.
따라서 Algebra, Lemma
\ref{algebra-lemma-Noetherian-normal-domain-finite-separable-extension}에
의해 $A \subset A_1$과 $B \subset B_1$은 유한이다.
\item $A$가 나가타이면 $A \subset A_1$은 유한이다.
\item $B$가 나가타이면 $B \subset B_1$은 유한이다.
\end{enumerate}
\end{remark}

\begin{lemma}
\label{more-algebra-lemma-pull-root-uniformizer}
$A$를 균등화원이 $\pi$인 이산 값매김환이라 하고 $n \geq 2$라 하자.
$K_1 = K[\pi^{1/n}]$라 하자. 그러면
\begin{enumerate}
\item $[K_1 : K] = n$이다.
\item $K_1$ 안에서의 $A$의 정수적 폐포 $A_1$은 환
$A[\pi^{1/n}]$이다.
\item $A_1$은 이산 값매김환이다.
\item $A$ 위 $A_1$의 분기지수는 $n$이다.
\item $K_1$은 $A$에 관하여 완전 분기이다.
\item $n$이 $A$의 잉여표수와 서로소이면, $K_1/K$는 순분기이고
$K_1/K$의 모든 부분확대는 $n$의 어떤 약수 $d$에 대한
$\pi^{1/d}$에 의해 생성된다.
\end{enumerate}
\end{lemma}

\begin{proof}
환 $A' = A[x]/(x^n - \pi)$를 생각하자.
% P02-E0621: frozen source omits “by” in the notation assignment.
% P02-E0622: frozen source places the image of x in A_1 before proving A'=A_1.
$A_1$ 안에서의 $x$의 상을 $\pi'$이라 나타내자.
$A'$은 $A$-가군으로 계수 $n$의 유한 자유이므로 원소 $\pi$는
$A'$ 안에서 영인자가 아니며, 따라서 $\pi'$도 그러하다.
더욱이 $A'/\pi'A'$은 체인 $A/\pi A$와 동형이다.
그러므로 $A'$은 균등화원이 $\pi'$인 이산 값매김환이다.
$A'$의 분수체 $K'$은 차수 $n$의 확대이다.
% P02-E0623: frozen source says “by adjoint” instead of “by adjoining”.
명백히 $K'$은 $K$에 $\pi$의 $n$제곱근
$\pi' = \pi^{1/n}$을 첨가하여 얻는다. 즉, $K' \cong K_1$이다.
$A'$은 정규이므로 $A'$은 $K'$ 안에서의 $A$의 정수적 폐포이다.
즉, $A' = A_1$이다. 이로써 (1)--(5)가 증명된다.

\medskip\noindent
$n$이 잉여체에서 가역이면 $K_1/K$가 순분기라는 주장은 정의
\ref{more-algebra-definition-types-of-extensions}와 (3), (4)에서 즉시 따른다.
마지막 명제를 증명하려면 1의 모든 $n$제곱근 $\zeta \in K'$가
$K$에 속함을 보이면 충분하다. Fields, Lemma
\ref{fields-lemma-subfields-kummer}를 보라.
명백히 $\zeta \in (A')^*$이다. 어떤 $a \in A$와
$b \in A \pi' \oplus \ldots \oplus A (\pi')^{n - 1}$에 대하여
$\zeta = a + b$라 쓸 수 있다. $b$가 영이 아니면
$b = u(\pi')^t$라 쓸 수 있다. 여기서 $t \geq 1$이고,
% P02-E0624: frozen source says b, not u, is a unit of A prime.
$b$는 $A'$의 단위원이다. 그러면
% P02-E0619: frozen source leaves this congruence display unpunctuated.
$$
1 = \zeta^n \equiv a^n + n a^{n - 1} u (\pi')^t \bmod (\pi')^{t + 1}
$$
이다. $n$과 $a$는 $A$의 단위원이므로 이는
$b = u(\pi')^t$가 $(\pi')^{t + 1}$을 법으로
$(1 - a^n)/(n a^{n - 1}) \in A$와 합동임을 함의한다.
직합분해 $A' = A \oplus \ldots \oplus A(\pi')^{n - 1}$에 의해 이는
불가능하다. $(\pi')^{t + 1}$에 의한 곱셈의 상은 대응하는 직합이고,
$A(\pi')^t$와의 교집합은 $(\pi A)(\pi')^t$이기 때문이다.
따라서 원하는 대로 $b = 0$이고 $\zeta \in A$이다.
\end{proof}

\begin{lemma}
\label{more-algebra-lemma-formally-smooth-goes-up}
$A \to B$를 분수체들이 $K \subset L$인 이산 값매김환의 확대라 하자.
$A \to B$가 $\mathfrak m_B$-진 위상에서 형식적으로 매끄럽다고
가정하자. 그러면 모든 유한 확대 $K_1/K$에 대하여
$L_1 = L \otimes_K K_1$, $B_1 = B \otimes_A A_1$이고,
각 확대 $(A_1)_{\mathfrak m_i} \subset (B_1)_{\mathfrak m_{ij}}$
(비고 \ref{more-algebra-remark-construction}를 보라)는
$\mathfrak m_{ij}$-진 위상에서 형식적으로 매끄럽다.
\end{lemma}

\begin{proof}
보조정리 \ref{more-algebra-lemma-extension-dvrs-formally-smooth}의 동치성을 더
언급하지 않고 사용하겠다. $\pi \in A$와
$\pi_i \in (A_1)_{\mathfrak m_i}$를 균등화원들이라 하자.
$\kappa_A \subset \kappa_B$가 분리가능이므로 환
% P02-E0619: frozen source leaves this quotient display unpunctuated.
$$
(B \otimes_A (A_1)_{\mathfrak m_i})/\pi_i (B \otimes_A (A_1)_{\mathfrak m_i}) =
B/\pi B \otimes_{A/\pi A} (A_1)_{\mathfrak m_i}/\pi_i (A_1)_{\mathfrak m_i}
$$
은 각각 $\kappa_{\mathfrak m_i}$ 위에서 분리가능인 체들의 곱이다.
따라서 $B \otimes_A (A_1)_{\mathfrak m_i}$ 안의 원소 $\pi_i$는
영인자가 아니고, 이 원소에 의한 몫은 체들의 곱이다.
% P02-E0625: frozen source claims a single Dedekind domain although the base change may split.
그러므로 $B \otimes_A A_1$은 특히 기약인 데데킨트 정역이다.
따라서 $B \otimes_A A_1 \subset B_1$은 등식이다.
\end{proof}

\noindent
다음 보조정리는 이산 값매김환에 대한 우리의 아브얀카르 보조정리이다.
$\kappa_B/\kappa_A$가 대수적 체 확대라고 가정하지 않음에 유의하자.

\begin{lemma}[아브얀카르 보조정리]
\label{more-algebra-lemma-abhyankar}
$A \subset B$를 이산 값매김환의 확대라 하자.
$A$의 잉여표수가 $0$이거나, 잉여표수가 $p$이고 분기지수 $e$가
$p$와 서로소이며 $\kappa_B/\kappa_A$가 분리가능 체 확대라고
가정하자. $K_1/K$를 유한 확대라 하자. 비고
\ref{more-algebra-remark-construction}의 표기를 사용하고, 어떤 $i$에 대하여
$e$가 $A \subset (A_1)_{\mathfrak m_i}$의 분기지수를 나눈다고
가정하자. 그러면 모든 $j = 1, \ldots, m_i$에 대하여
$(A_1)_{\mathfrak m_i} \subset (B_1)_{\mathfrak m_{ij}}$는
$\mathfrak m_{ij}$-진 위상에서 형식적으로 매끄럽다.
\end{lemma}

\begin{proof}
$\pi \in A$를 균등화원이라 하자.
$\pi_1$을 $(A_1)_{\mathfrak m_i}$의 균등화원이라 하자.
$(A_1)_{\mathfrak m_i}$의 단위원 $u$와
$A \subset (A_1)_{\mathfrak m_i}$의 분기지수 $e_1$에 대하여
$\pi = u \pi_1^{e_1}$이라 쓰자.

\medskip\noindent
주장: $u$가 $K_1$ 안에서 $e$제곱이라고 가정할 수 있다.
실제로 $x^e = u$의 근을 첨가하여 얻는 $K_1$의 확대를 $K_2$라 하자.
즉, $K_2$는 $K_1[x]/(x^e - u)$의 한 인자이다.
$e$에 대한 가정에 의해 $K_2/K_1$은 유한 분리 확대이고,
따라서 $A_1 \subset A_2$는 유한이다.
% P02-E0627: frozen source localizes A_2 and B_2 at ideals belonging to A_1 and B_1.
$(A_1)_{\mathfrak m_i} \to
(A_1)_{\mathfrak m_i}[x]/(x^e - u)$는 유한 에탈이므로
($e$는 잉여표수와 서로소이고 $u$는 단위원이기 때문이다),
$(A_2)_{\mathfrak m_i}$는 $(A_1)_{\mathfrak m_i}$의 유한 에탈 확대의
한 인자이고, 따라서 그 자체가 $(A_1)_{\mathfrak m_i}$ 위 유한
에탈이라고 결론내린다. 같은 논의로 $B_1 \subset B_2$가 유한 에탈
확대들 $(B_1)_{\mathfrak m_{ij}} \subset
(B_2)_{\mathfrak m_{ij}}$를 유도함을 알 수 있다.
$\mathfrak m_{ij} \subset B_1$ 위에 놓이는 극대아이디얼
$\mathfrak m'_{ij} \subset B_2$를 택하고(물론 하나보다 많을 수도 있다)
도식
% P02-E0630: frozen source leaves this and two later commutative diagrams unpunctuated.
$$
\xymatrix{
(B_1)_{\mathfrak m_{ij}} \ar[r] & (B_2)_{\mathfrak m'_{ij}} \\
(A_1)_{\mathfrak m_i} \ar[u] \ar[r] &
(A_2)_{\mathfrak m'_i} \ar[u]
}
$$
을 생각하자.
% P02-E0628: frozen source calls the contracted prime an image.
여기서 $\mathfrak m'_i \subset A_2$는 그 상이다.
이제 수평 화살표들의 분기지수는 $1$이고 이들은 유한 분리 잉여체
확대들을 유도한다. 따라서 보조정리
\ref{more-algebra-lemma-extension-dvrs-formally-smooth}의 동치성을 사용하면,
오른쪽 수직 화살표가 $\mathfrak m'_{ij}$-진 위상에서 형식적으로
매끄러움을 보이는 것으로 충분하다. $u$는 $K_2$ 안에서
% P02-E0629: frozen source twice uses the wrong English article before e-th root.
$e$제곱근을 가지므로 주장을 얻는다.

\medskip\noindent
$u$가 $K_1$ 안에서 $e$제곱근을 갖는다고 가정하자.
$e | e_1$이고 $u$가 $K_1$ 안에서 $e$제곱근을 가지므로 어떤
$\theta \in K_1$에 대하여 $\pi = \theta^e$이다.
$K'_1 = K[\theta] \subset K_1$을 $\theta$에 의해 생성되는 부분체라
하자. 보조정리 \ref{more-algebra-lemma-pull-root-uniformizer}에 의해 $K[\theta]$
안에서의 $A$의 정수적 폐포 $A'_1$은 이산 값매김환
$A'_1 = A[\theta]$이고, $A$ 위 분기지수는 $e$이다.
$K'_1/K$에 대하여 보조정리를 증명할 수 있다면 모든
$j = 1, \ldots, m_i$에 대하여 도식
% P02-E0630: frozen source leaves this commutative diagram unpunctuated.
$$
\xymatrix{
(B'_1)_{B'_1 \cap \mathfrak m_{ij}} \ar[r] & (B_1)_{\mathfrak m_{ij}} \\
A'_1 \ar[u] \ar[r] &
(A_1)_{\mathfrak m_i} \ar[u]
}
$$
에 보조정리 \ref{more-algebra-lemma-formally-smooth-goes-up}을 적용하여 결론을
얻는다. 이로써 다음 문단에서 다룰 경우로 환원된다.

\medskip\noindent
$K_1 = K[\pi^{1/e}]$라고 가정하고 $\theta = \pi^{1/e}$라 두자.
$\pi_B$를 $B$의 균등화원이라 하고, 어떤 $B$의 단위원 $w$에 대하여
$\pi = w \pi_B^e$라 쓰자. 그러면 $L_1 = L \otimes_K K_1$은
$w$의 $e$제곱근인 $\theta / \pi_B$를 첨가하여 얻어진다.
따라서 $B \subset B_1$은 유한 에탈이다.
그러므로 모든 극대아이디얼 $\mathfrak m \subset B_1$에 대하여
가환도식
% P02-E0630: frozen source leaves this commutative diagram unpunctuated.
$$
\xymatrix{
B \ar[r]_1 & (B_1)_{\mathfrak m} \\
A \ar[u]^e \ar[r]^e & A_1 \ar[u]_{e_\mathfrak m}
}
$$
을 생각하자. 여기서 화살표들을 따른 수들은 분기지수이다.
분기지수의 곱셈성(보조정리 \ref{more-algebra-lemma-multiplicative-e-f})에 의해
$e_\mathfrak m = 1$이라고 결론내린다.
잉여체 확대들을 보면 $\kappa(\mathfrak m)$은 $\kappa_B$의 유한 분리
확대이고, 이는 $\kappa_A$ 위에서 분리가능이다.
따라서 $\kappa(\mathfrak m)$은 $\kappa_A$ 위에서 분리가능이다.
이는 $A_1$의 잉여체와 같으므로 보조정리
\ref{more-algebra-lemma-extension-dvrs-formally-smooth}에 의해 결론을 얻는다.
\end{proof}

\begin{lemma}
\label{more-algebra-lemma-composition-tame}
$A$를 분수체가 $K$인 이산 값매김환이라 하자.
$M/L/K$를 유한 분리 확대들이라 하자.
$B$를 $L$ 안에서의 $A$의 정수적 폐포라 하자.
$L/K$가 $A$에 관하여 순분기이고, $B$의 모든 극대아이디얼
$\mathfrak m$에 대하여 $M/L$이 $B_\mathfrak m$에 관하여
순분기이면, $M/K$는 $A$에 관하여 순분기이다.
\end{lemma}

\begin{proof}
$C$를 $M$ 안에서의 $A$의 정수적 폐포라 하자.
$C$의 모든 극대아이디얼 $\mathfrak m'$은 $B$의 어떤 극대아이디얼
$\mathfrak m$ 위에 놓인다. 그러면 분기지수의 곱셈성(보조정리
\ref{more-algebra-lemma-multiplicative-e-f})과 유한 체 확대들의 탑
$\kappa(\mathfrak m')/\kappa(\mathfrak m)/\kappa_A$가 있다는 사실에서
보조정리가 따른다.
\end{proof}

\begin{lemma}
\label{more-algebra-lemma-subextension-tame}
$A$를 분수체가 $K$인 이산 값매김환이라 하자.
$M/L/K$가 유한 분리 확대들이고 $M$이 $A$에 관하여 순분기이면,
$L$은 $A$에 관하여 순분기이다.
\end{lemma}

\begin{proof}
비고 \ref{more-algebra-remark-finite-separable-extension}의 논의 결과들을 더
언급하지 않고 사용하겠다.
% P02-E0631: frozen source uses slash notation without binding the respective closures explicitly.
$C/B/A$를 $M/L/K$ 안에서의 $A$의 정수적 폐포들이라 하자.
$C$는 $B$의 유한 환 확대이므로
$\Spec(C) \to \Spec(B)$는 전사이다.
% P02-E0632: frozen source says “for ever maximal ideal”.
따라서 모든 극대아이디얼 $\mathfrak m \subset B$에 대하여
$\mathfrak m$ 위에 놓이는 극대아이디얼 $\mathfrak m' \subset C$가
존재한다. 분기지수의 곱셈성(보조정리
\ref{more-algebra-lemma-multiplicative-e-f})과 가정에 의해
$A$ 위 $B_\mathfrak m$의 분기지수는 잉여표수와 서로소이다.
$\kappa(\mathfrak m')/\kappa_A$가 유한 분리이므로
$\kappa(\mathfrak m)/\kappa_A$도 그러하다.
\end{proof}

\begin{lemma}
\label{more-algebra-lemma-characterize-tame}
$A$를 분수체가 $K$인 이산 값매김환이라 하고,
$\pi \in A$를 균등화원이라 하자.
$L/K$를 유한 분리 확대라 하자. 다음은 동치이다.
\begin{enumerate}
\item $L$은 $A$에 관하여 순분기이다.
\item $\kappa_A$에서 가역인 $e \geq 1$과
$A' = A[\pi^{1/e}]$에 관하여 비분기인 확대
$L'/K' = K[\pi^{1/e}]$가 존재하여 $L$이 $L'$에 포함된다.
\item $\kappa_A$에서 가역인 $e_0 \geq 1$이 존재하여,
$\kappa_A$에서 가역인 모든 $d \geq 1$에 대하여
$e = de_0$로 두었을 때 (2)가 성립한다.
\end{enumerate}
\end{lemma}

\begin{proof}
$A'$은 분수체가 $K'$인 이산 값매김환이다. 보조정리
\ref{more-algebra-lemma-pull-root-uniformizer}를 보라.
물론 $A$ 위 $A'$의 분기지수는 $e$이다.
따라서 (2)가 성립하면 보조정리 \ref{more-algebra-lemma-composition-tame}에 의해
$L'$은 $A$에 관하여 순분기이다. 그러므로 보조정리
\ref{more-algebra-lemma-subextension-tame}에 의해 $L$은 $A$에 관하여 순분기이다.

\medskip\noindent
(3) $\Rightarrow$ (2)는 즉시 성립한다.

\medskip\noindent
(1)이 성립한다고 가정하자. $B$를 $L$ 안에서의 $A$의 정수적
폐포라 하고, 그 극대아이디얼들을
$\mathfrak m_1, \ldots, \mathfrak m_n$이라 하자.
% P02-E0633: frozen source omits “by” when assigning the ramification indices.
$e_i$를 $A \to B_{\mathfrak m_i}$의 분기지수라 하자.
% P02-E0634: frozen source ends the lcm at undefined e_r instead of e_n.
$e_0$를 $e_1, \ldots, e_r$의 최소공배수라 하자.
가정 (1)에 의해 이는 $\kappa_A$에서 가역이다.
(3)에서와 같이 $e = de_0$라 하자. $A' = A[\pi^{1/e}]$라 두자.
그러면 $A \to A'$은 분수체가 $K' = K[\pi^{1/e}]$인 이산
값매김환의 확대이다. 보조정리 \ref{more-algebra-lemma-pull-root-uniformizer}를 보라.
곱분해
% P02-E0636: frozen product-decomposition display lacks its sentence comma.
$$
L \otimes_K K' = \prod L'_j
$$
를 택하자. 여기서 $L'_j$들은 체이다. $B'_j$를 $L'_j$ 안에서의
$A$의 정수적 폐포라 하자. $\mathfrak m_{ijk}$를
$\mathfrak m_i$ 위에 놓이는 $B'_j$의 극대아이디얼들이라 하자.
% P02-E0635: frozen source localizes B'_j at m_i, an ideal of B, rather than m_ijk.
$(B'_j)_{\mathfrak m_i}$는 $L'_j$ 안에서의
$B_{\mathfrak m_i}$의 정수적 폐포임을 관찰하자.
$A \subset B_{\mathfrak m_i}$와 확대 $K'/K$에 아브얀카르 보조정리
(보조정리 \ref{more-algebra-lemma-abhyankar})를 적용하면
$A' \to (B'_j)_{\mathfrak m_{ijk}}$가
$\mathfrak m_{ijk}$-진 위상에서 형식적으로 매끄러움을 알 수 있다.
이는 분기지수가 $1$이고 잉여체 확대가 분리가능임을 함의한다
(보조정리 \ref{more-algebra-lemma-extension-dvrs-formally-smooth}).
이와 같이 $L'_j$가 $A'$에 관하여 비분기임을 알 수 있다.
이로써 증명이 끝난다. 어떤 $j$에 대하여 $L' = L'_j$로 택하면 된다.
\end{proof}

\begin{lemma}
\label{more-algebra-lemma-permanence-tame}
$A$를 분수체가 $K$인 이산 값매김환이라 하자.
\begin{enumerate}
\item $L/K$가 $A$에 관하여 순분기인 유한 분리 확대이면,
$L$을 포함하고 $A$에 관하여 순분기인 갈루아 확대 $M/K$가 존재한다.
\item $L_1/K$, $L_2/K$가 $A$에 관하여 순분기인 유한 분리
확대들이면, $L_1$과 $L_2$를 포함하고 $A$에 관하여 순분기인 유한
분리 확대 $L/K$가 존재한다.
\end{enumerate}
\end{lemma}

\begin{proof}
(2)의 증명. 균등화원 $\pi \in A$를 택하자.
$\kappa_A$에서 가역인 정수 $e$와,
$A' = A[\pi^{1/e}]$에 관하여 비분기인 확대들
$L_i'/K' = K[\pi^{1/e}]$을 택할 수 있다.
% P02-E0637: frozen source does not state the needed containment in “with L'_i/L_i”.
이들은 $K$의 확대들로서 $L'_i/L_i$를 갖는다. 보조정리
\ref{more-algebra-lemma-characterize-tame}를 보라.
보조정리 \ref{more-algebra-lemma-permanence-unramified}에 의해
$A'$에 관하여 비분기인 확대 $L'/K'$을 찾아서
% P02-E0638: frozen source writes L'_i/K as a subextension of L'/K'.
$i = 1, 2$에 대하여 $L'_i/K$가 $L'/K'$의 부분확대와 동형이 되게
할 수 있다.
% P02-E0639: frozen proof of item (2) says this completes item (3).
$L'/K$는 순분기이므로 이로써 (3)의 증명이 끝난다
(위와 같은 보조정리를 사용하라).

\medskip\noindent
(1)의 증명. 먼저 $L$을 더 큰 확대로 바꾸어,
$L$이 $A' = A[\pi^{1/e}]$에 관하여 비분기인
$K' = K[\pi^{1/e}]$의 확대라고 가정할 수 있다. 여기서 $e$는
$\kappa_A$에서 가역이다. 보조정리 \ref{more-algebra-lemma-characterize-tame}를
보라. $M$을 $K$ 위 $L$의 정규 폐포라 하자. Fields, Definition
\ref{fields-definition-normal-closure}를 보라. 그러면 Fields, Lemma
\ref{fields-lemma-normal-closure-galois}에 의해 $M/K$는 갈루아이다.
한편 $K$-대수의 전사
% P02-E0640: frozen source leaves the finite tensor arity unbound.
% P02-E0642: frozen source leaves this and four later displays unpunctuated.
$$
L \otimes_K \ldots \otimes_K L \longrightarrow M
$$
가 존재한다. Fields, Lemma
\ref{fields-lemma-normal-closure-tensor-product}를 보라.
$B$를 비고 \ref{more-algebra-remark-finite-separable-extension}에서와 같이
$L$ 안에서의 $A$의 정수적 폐포라 하자.
$L$이 $A' = A[\pi^{1/e}]$에 관하여 비분기라는 조건은 정확히
$A' \to B$가 에탈 환 사상이라는 뜻이다. Algebra, Lemma
\ref{algebra-lemma-characterize-etale}를 보라.
주장:
% P02-E0640: frozen tensor arity is again unbound here.
% P02-E0642: frozen product display lacks sentence punctuation.
$$
K' \otimes_K \ldots \otimes_K K' = \prod K'_i
$$
는 $A$에 관하여 순분기인 체 확대들 $K'_i/K$의 곱이다.
그러면 $A'_i$를 $K'_i$ 안에서의 $A$의 정수적 폐포라 할 때,
% P02-E0640: frozen tensor arities remain unbound in this base change.
% P02-E0642: frozen tensor-product display lacks sentence punctuation.
$$
\prod A'_i \otimes_{(A' \otimes_A \ldots \otimes_A A')}
(B \otimes_A \ldots \otimes_A B)
$$
는 $\prod A'_i$ 위 유한 에탈이고, 따라서 데데킨트 정역들의
곱이다(보조정리 \ref{more-algebra-lemma-Dedekind-etale-extension}).
$M$은 이 데데킨트 정역들 가운데 하나의 분수체이고, 그 정역은 어떤
$i$에 대하여 $A'_i$ 위 유한 에탈이라고 결론내린다.
따라서 $M/K'_i$는 $A'_i$의 모든 극대아이디얼에 관하여 비분기이고,
보조정리 \ref{more-algebra-lemma-composition-tame}에 의해 $M/K$는 순분기이다.

\medskip\noindent
% P02-E0641: frozen source says “It remains the prove”.
주장을 증명하는 일만 남는다. 이를 위해
$A' = A[x]/(x^e - \pi)$라 쓰면
% P02-E0640: frozen source first binds r only inside this display.
% P02-E0642: frozen normalization display lacks sentence punctuation.
$$
A' \otimes_A \ldots \otimes_A A' =
A'[x_1, \ldots, x_r]/(x_1^e - \pi, \ldots, x_r^e - \pi)
$$
이다. 이 환의 정규화는 분명히 $i = 2, \ldots, r$에 대한 원소
$y_i = x_i/x_1$을 포함한다. 이 원소들은 관계
$y_i^e - 1 = 0$을 만족하며,
% P02-E0642: frozen final normalization display lacks sentence punctuation.
$$
A[x_1, y_2, \ldots, y_r]/(x_1^e - \pi, y_2^e - 1, \ldots, y_r^e - 1) =
A'[y_2, \ldots, y_r]/(y_2^e - 1, \ldots, y_r^e - 1)
$$
를 얻는다. 이 환은 $e$가 $A'$에서 가역이므로 $A'$ 위 유한
에탈이다. 따라서 원하는 대로 각각 $A'$ 위 비분기인 데데킨트
정역들의 곱이다(혼동이 있으면 위에 든 참고문헌들을 보라).
\end{proof}

\begin{lemma}
\label{more-algebra-lemma-tame-goes-up}
$A \subset B$를 이산 값매김환의 확대라 하자.
대응하는 분수체 확대를 $L/K$로 나타내자.
$K'/K$를 유한 분리 확대라 하자. 그러면
% P02-E0644: frozen source leaves this and two later displays unpunctuated.
$$
K' \otimes_K L = \prod L'_i
$$
는 체들의 유한곱이고 다음이 성립한다.
\begin{enumerate}
\item $K'$이 $A$에 관하여 비분기이면 각 $L'_i$는 $B$에 관하여
비분기이다.
\item $K'$이 $A$에 관하여 순분기이면 각 $L'_i$는 $B$에 관하여
순분기이다.
\end{enumerate}
\end{lemma}

\begin{proof}
대수 $K' \otimes_K L$은 $L$ 위 유한 에탈 대수이므로 체들의 유한곱이다.
$A'$을 $K'$ 안에서의 $A$의 정수적 폐포라 하자.

\medskip\noindent
(1)의 경우 환 사상 $A \to A'$은 유한 에탈이다.
따라서 $B' = B \otimes_A A'$은 $B$ 위 유한 에탈이고 데데킨트
정역들의 유한곱이다(보조정리
\ref{more-algebra-lemma-Dedekind-etale-extension}).
그러므로 $B'$은 $K' \otimes_K L$ 안에서의 $B$의 정수적 폐포이다.
각 $L'_i$가 $B$에 관하여 비분기임이 즉시 따른다.

\medskip\noindent
균등화원 $\pi \in A$를 택하자. (2)를 증명하기 위해 $K'$을
% P02-E0643: frozen source says “a larger extension tame ramified”.
$A$에 관하여 순분기인 더 큰 확대로 바꾸어도 좋다
(세부사항은 생략한다. 힌트: 보조정리 \ref{more-algebra-lemma-subextension-tame}를
사용하라). 따라서 보조정리 \ref{more-algebra-lemma-characterize-tame}에 의해
$\kappa_A$에서 가역인 어떤 $e \geq 1$이 존재하여,
$K'$은 $K[\pi^{1/e}]$를 포함하고
$K'$이 $A[\pi^{1/e}]$에 관하여 비분기라고 가정할 수 있다.
곱분해
% P02-E0644: frozen source leaves this product display unpunctuated.
$$
K[\pi^{1/e}] \otimes_K L = \prod L_{e, j}
$$
를 택하자. 모든 $i$에 대하여
$L'_i/L_{e, j_i}$가 유한 분리 확대가 되는 $j_i$가 존재한다.
$B_{e, j}$를 $L_{e, j}$ 안에서의 $B$의 정수적 폐포라 하자.
$K'/K[\pi^{1/e}]$와
$A[\pi^{1/e}] \subset (B_{e, j_i})_\mathfrak m$에 (1)을 적용하면,
모든 극대아이디얼 $\mathfrak m \subset B_{e, j_i}$에 대하여
$L'_i$가 $(B_{e, j_i})_\mathfrak m$에 관하여 비분기임을 알 수 있다.
따라서 보조정리 \ref{more-algebra-lemma-composition-tame}에 의해
$L_{e, j}$가 $B$에 관하여 순분기임을 보이면 증명이 끝난다.

\medskip\noindent
$B$ 안에서 균등화원 $\theta$를 택하자.
$B$의 단위원 $u$와 $t \geq 1$에 대하여
$\pi = u \theta^t$라 쓰자. 그러면
% P02-E0644: frozen source leaves this ring display unpunctuated.
$$
A[\pi^{1/e}] \otimes_A B = B[x]/(x^e - u \theta^t)
\subset B[y, z]/(y^{e'} - \theta, z^e - u)
$$
이다. 여기서 $e' = e/\gcd(e, t)$이다. 사상은 $x$를
$z y^{t/\gcd(e, t)}$로 보낸다. 오른쪽은 각각 $B$ 위에서 순분기인
데데킨트 정역들의 곱이므로 증명이 끝난다(세부사항은 생략한다).
\end{proof}

\section{분기 제거}
\label{more-algebra-section-eliminating-ramification}

\noindent
이 절에서는 Helmut Epp의 결과 하나를 논한다. \cite{Epp}를 보라.
독자에게 원 논문을 읽기를 강력히 권한다. 혼합 표수와 등표수의 경우를
같은 방법으로 다루려 하기 때문에 여기서의 접근법은 조금 다르다.
관련 결과는 \cite{Ponomarev}, \cite{Ponomarev-Abhyankar},
\cite{Kuhlmann}, \cite{ZK}도 보라.

\medskip\noindent
$A \subset B$를 분수체가 $K \subset L$인 이산 값매김환의 확대라 하자.
이 절의 목표는 다음 두 그림이 평주 \ref{more-algebra-remark-construction}에서와 같을 때
% P02-E0645: frozen source leaves this display boundary unpunctuated.
$$
\vcenter{
\xymatrix{
L \ar[r] & L_1 \\
K \ar[u] \ar[r] & K_1 \ar[u]
}
}
\quad\text{and}\quad
\vcenter{
\xymatrix{
B \ar[r] & B_1 \ar[r] & (B_1)_{\mathfrak m_{ij}} \\
A \ar[u] \ar[r] & A_1 \ar[r] \ar[u] & (A_1)_{\mathfrak m_i} \ar[u]
}
}
$$
모든 확대
$(A_1)_{\mathfrak m_i} \subset (B_1)_{\mathfrak m_{ij}}$가
약하게 비분기이거나, 더 나아가 관련 진 위상에서 형식적으로 매끄럽게 되는
유한 확대 $K_1/K$를 찾는 것이다. 가장 간단하지만 자명하지는 않은 예는
아브얀카르 보조정리이다. 보조정리 \ref{more-algebra-lemma-abhyankar}를 보라.

\begin{definition}
\label{more-algebra-definition-solution}
$A \to B$를 분수체가 $K \subset L$인 이산 값매김환의 확대라 하자.
\begin{enumerate}
\item 평주 \ref{more-algebra-remark-construction}의 모든 확대
$(A_1)_{\mathfrak m_i} \subset (B_1)_{\mathfrak m_{ij}}$가
약하게 비분기이면, 유한 체 확대 $K_1/K$를
{\it $A \subset B$의 약한 해}라고 한다.
\item 평주 \ref{more-algebra-remark-construction}의 각 확대
$(A_1)_{\mathfrak m_i} \subset (B_1)_{\mathfrak m_{ij}}$가
$\mathfrak m_{ij}$-진 위상에서 형식적으로 매끄러우면, 유한 체 확대
$K_1/K$를 {\it $A \subset B$의 해}라고 한다.
\end{enumerate}
해 $K_1/K$에서 $K_1/K$가 분리 확대이면 이를 {\it 분리 해}라고 한다.
\end{definition}

\noindent
일반적으로 (약한) 해는 존재하지 않으며 \cite{Epp}에 한 예가 있다.
잉여체 확대에 대한 온화한 가정 아래에서는 \cite{Epp}를 따라 정리
\ref{more-algebra-theorem-epp}에서 약한 해의 존재를 증명한다. 다음 절에서는 기하적으로
의미 있는 경우 해의 존재, 때로는 분리 해의 존재도 도출한다. 명제
\ref{more-algebra-proposition-epp-essentially-finite-type}와 보조정리
\ref{more-algebra-lemma-epp-essentially-finite-type-separable}를 보라.
그러나 다음 예는 일반적으로 약한 해조차 얻으려면 비분리 확대가 필요함을
보여 준다.

\begin{example}
\label{more-algebra-example-inseparable-necessary}
$k$를 표수가 $p > 0$인 완전체라 하자. $A = k[[x]]$ 및
$K = k((x))$라 하자. $B = A[x^{1/p}]$라 하자. $A \to B$의 임의의
약한 해 $K_1/K$는 비분리이다(그리고 $K$의 임의의 유한 비분리 확대는
해이다). 증명은 생략한다.
\end{example}

\noindent
해는 더 큰 확대 아래에서 안정적이다. 보조정리
\ref{more-algebra-lemma-solution-goes-up}을 보라. 약한 해에는 이것이 성립하지 않을 수
있다. 약한 해는 어떤 의미에서 완전 분기 확대 아래에서 안정적이다.
보조정리 \ref{more-algebra-lemma-weakly-unramified-goes-up-along-totally-ramified}를 보라.

\begin{lemma}
\label{more-algebra-lemma-weakly-unramified-goes-up-along-totally-ramified}
$A \to B$를 분수체가 $K \subset L$인 이산 값매김환의 확대라 하자.
$A \to B$가 약하게 비분기라고 가정하자. 그러면 $A$에 관하여 완전
분기인 임의의 유한 분리 확대 $K_1/K$에 대하여
$L_1 = L \otimes_K K_1$은 체이고, $A_1$과
$B_1 = B \otimes_A A_1$은 이산 값매김환이며, 확대
$A_1 \subset B_1$은 약하게 비분기이다(평주
\ref{more-algebra-remark-construction}을 보라).
\end{lemma}

\begin{proof}
$\pi \in A$와 $\pi_1 \in A_1$을 균등화원이라 하자. $K_1/K$가
$A$에 관하여 완전 분기이므로 $A_1$의 어떤 단위원 $u_1$에 대하여
$\pi_1^e = u_1 \pi$이다. 따라서 $A_1$은 $A$ 위에서 $\pi_1$로
생성되고, $\pi_1$의 $K$ 위 최소다항식 $P(t)$는 다음 꼴이다.
% P02-E0645: frozen source leaves this display boundary unpunctuated.
$$
P(t) = t^e + a_{e - 1} t^{e - 1} + \ldots + a_0
$$
여기서 $a_i \in (\pi)$이고 $A$의 어떤 단위원 $u$에 대하여
$a_0 = u\pi$이다. 또한 $e = [K_1 : K]$이다. $A \to B$가 약하게
비분기이므로 $\pi$는 $B$의 균등화원이다. 따라서
$B_1 = B[t]/(P(t))$는 $t$의 잉여류를 균등화원으로 갖는 이산
값매김환이다. 이제 보조정리는 명백하다.
\end{proof}

\begin{lemma}
\label{more-algebra-lemma-solutions-go-down}
$A \to B \to C$를 분수체가 $K \subset L \subset M$인 이산
값매김환의 확대라 하고, $K_1/K$를 유한 확대라 하자.
\begin{enumerate}
\item $K_1$이 $A \to C$의 (약한) 해이면 $K_1$은
$A \to B$의 (약한) 해이다.
\item $K_1$이 $A \to B$의 (약한) 해이고
$L_1 = (L \otimes_K K_1)_{red}$가 각각 $B \to C$의 (약한) 해인
체들의 곱이면, $K_1$은 $A \to C$의 (약한) 해이다.
\end{enumerate}
\end{lemma}

\begin{proof}
$L_1 = (L \otimes_K K_1)_{red}$ 및
$M_1 = (M \otimes_K K_1)_{red}$라 하고, $B_1 \subset L_1$과
$C_1 \subset M_1$을 각각 $B$와 $C$의 정수적 폐포라 하자.
% P02-E0646: frozen source's claimed tensor identity is retained for canon review.
$M_1 = (M \otimes_L L_1)_{red}$이고 $L_1$은 $L$의 유한 확대들의
공집합이 아닌 유한 곱임에 주의하자. 따라서 환 사상 $B_1 \to C_1$은
평주 \ref{more-algebra-remark-construction}에서 논한 꼴의 환 사상들의 유한 곱이다.
특히 사상 $\Spec(C_1) \to \Spec(B_1)$은 전사이다. 극대아이디얼
$\mathfrak m \subset C_1$을 택하고 다음 이산 값매김환의 확대들을
생각하자.
% P02-E0645: frozen source leaves this display boundary unpunctuated.
$$
(A_1)_{A_1 \cap \mathfrak m} \to
(B_1)_{B_1 \cap \mathfrak m} \to
(C_1)_\mathfrak m
$$
합성이 약하게 비분기이면 사상
$(A_1)_{A_1 \cap \mathfrak m} \to
(B_1)_{B_1 \cap \mathfrak m}$도 약하게 비분기이다.
잉여체 확대
$\kappa_{A_1 \cap \mathfrak m} \to \kappa_\mathfrak m$가 분리이면
부분 확대
$\kappa_{A_1 \cap \mathfrak m} \to
\kappa_{B_1 \cap \mathfrak m}$도 분리이다. 보조정리
\ref{more-algebra-lemma-extension-dvrs-formally-smooth}를 함께 고려하면 (1)이
증명된다. (2)에도 비슷한 논증이 적용된다.
\end{proof}

\begin{lemma}
\label{more-algebra-lemma-solution-after-strict-henselization}
$A \to B$를 이산 값매김환의 확대라 하자. 다음 조건을 만족하는 이산
값매김환 확대의 가환 그림이 존재한다.
% P02-E0645: frozen source leaves this display boundary unpunctuated.
$$
\xymatrix{
B \ar[r] & B' \\
A \ar[r] \ar[u] & A' \ar[u]
}
$$
\begin{enumerate}
\item 분수체의 확대 $K'/K$와 $L'/L$은 분리 대수적이다.
\item $A'$과 $B'$의 잉여체는 각각 $A$와 $B$의 잉여체의 분리
대수적 폐포이다.
\item $A' \to B'$에 해, 약한 해 또는 분리 해가 존재하면
$A \to B$에도 각각 해, 약한 해 또는 분리 해가 존재한다.
\end{enumerate}
\end{lemma}

\begin{proof}
Algebra, 보조정리
\ref{algebra-lemma-colimit-finite-etale-given-residue-field}에 의해,
$A'$의 잉여체가 $A$의 잉여체의 분리 대수적 폐포이고 유한
\'etale 확대들의 여과 쌍대극한인 확대 $A \subset A'$이 존재한다.
그러면 $A \subset A'$은 이산 값매김환의 확대이고, 유도된 분수체 확대
$K'/K$는 분리 대수적이다.

\medskip\noindent
$B \subset B'$을 $B$의 엄밀 헨젤화라 하자. 그러면
$B \subset B'$은 분수체 확대가 분리 대수적인 이산 값매김환의 확대이다.
Algebra, 보조정리
\ref{algebra-lemma-strictly-henselian-functorial-prepare}에 의해 보조정리의
명제에 나온 가환 그림이 존재한다. 보조정리의 (1)과 (2)는 명백하다.

\medskip\noindent
$K'_1/K'$를 $A' \to B'$의 (약한) 해라 하자. $A'$이 쌍대극한이므로
유한 \'etale 확대 $A \subset A_1'$과 $A_1'$의 분수체 $F$의 유한 확대
$K_1$을 찾아 $K'_1 = K' \otimes_F K_1$이 되게 할 수 있다.
$A \subset A_1'$은 유한 \'etale이고 $B'$은 엄밀 헨젤이므로
$B' \otimes_A A_1'$은 $B'$과 동형인 환들의 유한 곱이다. 따라서
% P02-E0645: frozen source leaves this display boundary unpunctuated.
$$
L' \otimes_K K_1 = L' \otimes_K F \otimes_F K_1
$$
은 $L' \otimes_{K'} K'_1$과 동형인 환들의 유한 곱이다. 그러므로
$K_1/K$는 $A \to B'$의 (약한) 해이다. 보조정리
\ref{more-algebra-lemma-solutions-go-down}에 의해 이는 $A \to B$의 (약한) 해이기도
하다.
% P02-E0647: frozen source does not separately justify descent of the separable-solution case.
\end{proof}

\begin{lemma}
\label{more-algebra-lemma-galois-relative}
$A \to B$를 분수체가 $K \subset L$인 이산 값매김환의 확대라 하자.
$K_1/K$를 정규 확대라 하고 $G = \text{Aut}(K_1/K)$라 하자. 그러면
$G$는 평주 \ref{more-algebra-remark-construction}의 환 $K_1$, $L_1$, $A_1$, $B_1$에
작용하고 $B_1$의 극대아이디얼들의 집합에 추이적으로 작용한다.
% P02-E0648: frozen source's normal-extension hypotheses and automorphism-group conclusion are retained.
\end{lemma}

\begin{proof}
마지막 단언을 제외하면 모두 명백하다. 작용에 두 개 이상의 궤도가 있으면,
한 궤도의 모든 극대아이디얼에서 영이 되고 다른 궤도의 모든
극대아이디얼에서 잉여류가 $1$인 원소 $b \in B_1$을 찾을 수 있다.
그러면 $b' = \prod_{\sigma \in G} \sigma(b)$는
$B_1 \subset L_1 = (L \otimes_K K_1)_{red}$의 $G$-불변 원소이며,
$B_1$의 어떤 극대아이디얼들에는 속하지만 $B_1$의 모든 극대아이디얼에는
않는다. 이를 $L \otimes_K K_1$의 원소로 올리고 충분히 높은 거듭제곱을
취하면, 어떤 $N > 0$에 대하여 $(b')^N$으로 사상되는
$L \otimes_K K_1$의 $G$-불변 원소 $b''$을 얻는다. 실제로 이것은 표수가
$p > 0$일 때만 필요하고, 이 경우 충분히 큰 $p$의 거듭제곱 $q$를 취하면
표준 사상
$(L \otimes_K K_1)_{red} \to L \otimes_K K_1$을 얻는다.
$K = (K_1)^G$이므로 $b'' \in L$이다. $b''$이 $B_1$의 원소로 사상되므로
$b'' \in B$이다($B$가 정규이기 때문이다). 이제 한편으로 $b'$이
$B_1$의 어떤 극대아이디얼에 속하므로 $b'' \in \mathfrak m_B$이어야 하고,
다른 한편으로 $b'$이 $B_1$의 모든 극대아이디얼에 속하지는 않으므로
$b'' \not \in \mathfrak m_B$이어야 한다. 이 모순으로 증명이 끝난다.
\end{proof}

\begin{lemma}
\label{more-algebra-lemma-make-degree-q-extension}
$A$를 균등화원이 $\pi$인 이산 값매김환이라 하자. $A$의 잉여 표수가
$p > 0$이면, 모든 $n > 1$ 및 $p$의 거듭제곱 $q$에 대하여 $A$에 관해
완전 분기인 차수 $q$의 분리 확대 $L/K$가 존재한다. 또한 $L$ 안에서의
$A$의 정수적 폐포 $B$는 분기지수가 $q$이고, 어떤 $b, b' \in B$에 대하여
$\pi_B^q = \pi + \pi^n b$ 및
$\pi_B^q = \pi + (\pi_B)^{nq}b'$을 만족하는 균등화원 $\pi_B$를 갖는다.
\end{lemma}

\begin{proof}
$K$의 표수가 영이면 $\pi_B^q = \pi$로 주어지는 확대를 취할 수 있다.
보조정리 \ref{more-algebra-lemma-pull-root-uniformizer}를 보라. $K$의 표수가
$p > 0$이면 $z^q - \pi^n z = \pi^{1 - q}$로 주어지는 $K$의 확대를
취할 수 있다. 실제로 $y = \pi z$이면
$y^q - \pi^{n + q - 1} y = \pi$이다. $\pi_B = y$로 두면 원하는
결과를 얻는다.
\end{proof}

\begin{lemma}
\label{more-algebra-lemma-pre-purely-inseparable-case}
% P02-E0649: frozen source's "reside field" typo is translated by its mathematical sense.
$A$를 이산 값매김환이라 하자. 잉여체 $\kappa_A$의 표수가 $p > 0$이고
$a \in A$의 잉여류가 $\kappa_A$에서 $p$제곱이 아니라고 가정하자.
그러면 $a$는 $K$에서도 $p$제곱이 아니고, $K[a^{1/p}]$ 안에서 $A$의
정수적 폐포는 환 $A[a^{1/p}]$이다. 이 환은 $A$ 위에서 약하게 비분기인
이산 값매김환이다.
\end{lemma}

\begin{proof}
이 보조정리는 자명하다.
\end{proof}

\begin{lemma}
\label{more-algebra-lemma-purely-inseparable-case}
% P02-E0650: frozen source's "fractions fields" wording is translated by its mathematical sense.
$A \subset B \subset C$를 분수체가 $K \subset L \subset M$인 이산
값매김환의 확대라 하자. $\pi \in A$를 균등화원이라 하자. 다음을 가정한다.
\begin{enumerate}
\item $B$는 나가타 환이다.
\item $A \subset B$는 약하게 비분기이다.
\item $M$은 $L$의 차수 $p$인 순수 비분리 확대이다.
\end{enumerate}
그러면 다음 중 하나가 성립한다.
\begin{enumerate}
\item $A \to C$는 약하게 비분기이다.
\item $C = B[\pi^{1/p}]$이다.
\item $A$에 관하여 완전 분기인 차수 $p$의 분리 확대 $K_1/K$가 존재하여
$L_1 = L \otimes_K K_1$과 $M_1 = M \otimes_K K_1$은 체이고,
정수적 폐포의 사상 $A_1 \to B_1 \to C_1$은 이산 값매김환의 약하게
비분기인 확대이다.
\end{enumerate}
\end{lemma}

\begin{proof}
$e$를 $B$ 위에서 $C$의 분기지수라 하자. $e = 1$이면 증명이 끝난다.
그렇지 않으면 보조정리 \ref{more-algebra-lemma-inequality}와
\ref{more-algebra-lemma-ramification-index-a-power-of-p}에 의해 $e = p$이다.
이는 다시 $B$와 $C$의 잉여체가 같음을 뜻한다. $C$의 균등화원
$\pi_C$를 택하자. $C$의 어떤 단위원 $u$에 대하여
$\pi_C^p = u \pi$라 쓰자. $\pi_C^p \in L$이므로 $u \in B^*$이다.
또한 $M = L[\pi_C]$이다.

\medskip\noindent
다음과 같이 쓸 수 있는 정수 $m \geq 0$이 존재한다고 하자.
$$
u = \sum\nolimits_{0 \leq i < m} b_i^p \pi^i + b \pi^m
$$
여기서 $b_i \in B$이고 $b \in B$의 $\kappa_B$에서의 상은 $p$제곱이
아니다. 보조정리 \ref{more-algebra-lemma-make-degree-q-extension}에서 $n = m + 2$로
두어 확대 $K_1/K$를 택하고, $K_1$ 안에서 $A$의 정수적 폐포 $A_1$의
균등화원 가운데 어떤 $a \in A_1$에 대하여
$\pi = (\pi')^p + (\pi')^{np} a$를 만족하는 것을 $\pi'$로 쓰자.
$B_1$을 $L \otimes_K K_1$ 안에서의 $B$의 정수적 폐포라 하자.
보조정리 \ref{more-algebra-lemma-weakly-unramified-goes-up-along-totally-ramified}에 의해
$A_1 \to B_1$은 약하게 비분기이다. $B_1$ 안에서 다음이 성립한다.
$$
u \pi =
\left(\sum\nolimits_{0 \leq i < m} b_i (\pi')^{i + 1}\right)^p +
b (\pi')^{(m + 1)p} + (\pi')^{np} b_1
$$
여기서 어떤 $b_1 \in B_1$이 존재한다(계산은 생략한다). 따라서
$M_1$은 다음 원소의 $p$제곱근을 $L_1$에 첨가하여 얻는다.
$$
b + (\pi')^{n - m - 1} b_1
$$
$B_1$의 잉여체는 $B$의 잉여체와 같으므로 보조정리
\ref{more-algebra-lemma-pre-purely-inseparable-case}에 의해 $M_1/L_1$의 차수는
$p$이고, $B_1$의 정수적 폐포 $C_1$은 $B_1$ 위에서 약하게 비분기이다.
따라서 이 경우에는 결론을 얻었다.

\medskip\noindent
% P02-E0651: frozen source switches between B and B_1 in the completion/Nagata argument.
앞 단락과 같은 정수 $m$이 존재하지 않으면 $u$는 $B_1$의 $\pi$-진
완비화에서 $p$제곱이다. $B$가 나가타이므로 Algebra, 보조정리
\ref{algebra-lemma-nagata-pth-roots}에 의해 이는 $u$가 $B_1$에서
$p$제곱이라는 뜻이다. 따라서 보조정리의 명제에서 둘째 경우가 성립한다.
\end{proof}

\begin{lemma}
\label{more-algebra-lemma-cohen}
$A$를 소수 $p$에 의해 영이 되고 극대아이디얼이 멱영인 국소환이라 하자.
잉여 사상 $A \to \kappa_A$의 절단인 환 사상
$\sigma : \kappa_A \to A$가 존재한다. $A \to A'$이 국소환 사이의
국소 준동형이면, 확대 $\kappa_{A'}/\kappa_A$가 분리인 경우
$\sigma$와 호환되는 비슷한 환 사상
$\sigma' : \kappa_{A'} \to A'$을 택할 수 있다.
\end{lemma}

\begin{proof}
Algebra, 명제
\ref{algebra-proposition-characterize-separable-field-extensions}에 의해
분리 확대는 형식적으로 매끄럽다. 따라서 $\sigma$의 존재는
$\mathbf{F}_p \to \kappa_A$가 분리라는 사실에서 따른다.
$\sigma$와 호환되는 $\sigma'$의 존재도 마찬가지이다.
\end{proof}

\begin{lemma}
\label{more-algebra-lemma-pre-characteristic-p-case}
$A$를 분수체 $K$의 표수가 $p > 0$인 이산 값매김환이라 하자.
$\xi \in K$라 하자. $L$을 $z^p - z = \xi$의 근 하나를 첨가하여 얻는
$K$의 확대라 하자. 그러면 $L/K$는 갈루아 확대이고 다음 중 하나가
성립한다.
\begin{enumerate}
\item $L = K$이다.
\item $L/K$는 $A$에 관하여 비분기이고 차수가 $p$이다.
\item $L/K$는 $A$에 관하여 완전 분기이고 분기지수가 $p$이다.
\item $L$ 안에서 $A$의 정수적 폐포 $B$는 이산 값매김환이고,
$A \subset B$는 약하게 비분기이며, $A \to B$는 차수 $p$의 순수 비분리
잉여체 확대를 유도한다.
\end{enumerate}
$\pi$를 $A$의 균등화원이라 하자. 다음 함의들이 성립한다.
\begin{enumerate}
\item[(A)] $\xi \in A$이면 (1) 또는 (2)의 경우이다.
\item[(B)] $\xi = \pi^{-n}a$이고 $n > 0$이 $p$로 나누어떨어지지 않으며
$a$가 $A$의 단위원이면 (3)의 경우이다.
\item[(C)] $\xi = \pi^{-n} a$이고 $n > 0$이 $p$로 나누어떨어지며
$a$의 $\kappa_A$에서의 상이 $p$제곱이 아니면 (4)의 경우이다.
\end{enumerate}
\end{lemma}

\begin{proof}
Fields, 절 \ref{fields-section-Artin-Schreier}의 논의에 의해 이 확대는
위수가 $p$의 약수인 갈루아 확대이다. 절 \ref{more-algebra-section-ramification}의
논의에서 보조정리에 나열한 (1)--(4) 가운데 한 경우가 성립함이 즉시 따른다.

\medskip\noindent
경우 (A). 여기서는 $A \to A[x]/(x^p - x - \xi)$가 유한 \'etale 환
확대임을 알 수 있다. 따라서 (1) 또는 (2)의 경우이다.

\medskip\noindent
경우 (B). $p$가 $n$을 나누지 않도록 $\xi = \pi^{-n}a$라 쓰자.
$B \subset L$을 $L$ 안에서의 $A$의 정수적 폐포라 하자. 어떤
극대아이디얼 $\mathfrak m$에 대하여 $C = B_\mathfrak m$이면
$p \text{ord}_C(z) = -n \text{ord}_C(\pi)$임이 명백하다.
특히 $A \subset C$의 분기지수는 $p$의 배수이다. 따라서 그 지수는
$p$이고 $B = C$이다.

\medskip\noindent
경우 (C). $k = n/p$라 두자. 그러면 방정식을 다음과 같이 다시 쓸 수 있다.
$$
(\pi^kz)^p - \pi^{n - k} (\pi^kz) = a
$$
$A[y]/(y^p - \pi^{n - k}y - a)$가 $A$ 위에서 약하게 비분기인 이산
값매김환이므로 보조정리가 따른다.
\end{proof}

\begin{lemma}
\label{more-algebra-lemma-characteristic-p-case}
$A \subset B \subset C$를 분수체가 $K \subset L \subset M$인 이산
값매김환의 확대라 하자. 다음을 가정한다.
\begin{enumerate}
\item $A \subset B$는 약하게 비분기이다.
\item $K$의 표수는 $p$이다.
\item $M$은 $L$의 차수 $p$인 갈루아 확대이다.
\item $\kappa_A = \bigcap_{n \geq 1} \kappa_B^{p^n}$이다.
\end{enumerate}
그러면 $A$에 관하여 완전 분기이고 $A \to C$의 약한 해인 유한 갈루아
확대 $K_1/K$가 존재한다.
\end{lemma}

\begin{proof}
$L$의 표수가 $p$이므로 $M$은 $L$의 아르틴--슈라이어 확대이다
(Fields, 보조정리 \ref{fields-lemma-Artin-Schreier}). 따라서
$z \in M$, $z \not \in L$이며 $\xi = z^p - z \in L$이 되도록
택할 수 있다. $\pi^n\xi \in B$가 되게 하는 $n \geq 0$을
택한다. 이 가운데 $n$이 최소가 되게 $z$를 택한다. $n = 0$이면
$M/L$은 $B$에 관하여 비분기이므로(보조정리
\ref{more-algebra-lemma-pre-characteristic-p-case}) 증명이 끝난다. 따라서 $n > 0$이다.

\medskip\noindent
가정 (4)는 $\kappa_A$가 완전체임을 뜻한다. 따라서 보조정리
\ref{more-algebra-lemma-cohen}에서와 같이 호환되는 환 사상
% P02-E0652: frozen source assigns the same barred-sigma notation to both maps.
$\overline{\sigma} : \kappa_A \to A/\pi^n A$와
$\overline{\sigma} : \kappa_B \to B/\pi^n B$를 택할 수 있다.
이 가운데 둘째 사상을 집합의 사상
$\sigma : \kappa_B \to B$로 올리자\footnote{$B$가 완비이면
$\sigma$를 환 사상으로 택할 수 있다. $A$도 완비이고 $\sigma$가 환
사상이면 $\sigma$는 $\kappa_A$를 $A$ 안으로 보낸다.}. 그러면 어떤
$\lambda_i \in \kappa_B$와 $b \in B$에 대하여 다음과 같이 쓸 수 있다.
$$
\xi = \sum\nolimits_{i = n, \ldots, 1} \sigma(\lambda_i) \pi^{-i} + b
$$
다음과 같이 두자.
$$
I = \{i \in \{n, \ldots, 1\} \mid \lambda_i \in \kappa_A\}
$$
그리고
% P02-E0653: frozen source uses lambda_i rather than the bound lambda_j in J.
$$
J = \{j \in \{n, \ldots, 1\} \mid \lambda_i \not \in \kappa_A\}
$$
유한 집합 $J$의 크기에 관하여 귀납하겠다.

\medskip\noindent
$J = \emptyset$인 경우. 이때 모든 $i \in \{n, \ldots, 1\}$에 대하여
$\sigma$의 선택에 의해 어떤 $a_i \in A$와 $b_i \in B$가 존재하여
$\sigma(\lambda_i) = a_i + \pi^n b_i$가 된다. 따라서 어떤
$a \in A$와 $b \in B$에 대하여 $\xi = \pi^{-n} a + b$이다.
$p | n$이면 $A$의 잉여체가 완전체이므로 어떤 $a_0, a_1 \in A$에
대하여 $a = a_0^p + \pi a_1$이라 쓴다. 계산하면 어떤 $b' \in B$에
대하여
$$
(z - \pi^{-n/p}a_0)^p - (z - \pi^{-n/p}a_0) =
\pi^{-(n - 1)}(a_1 + \pi^{n - 1 - n/p}a_0) + b'
$$
을 얻는다. 이는 $n$의 최소성에 모순이다. 따라서 $p$는 $n$을 나누지
않는다. $w^p - w = \pi^{-n}a$로 주어지는 $K$의 차수 $p$ 확대
$K_1$을 생각하자. 보조정리 \ref{more-algebra-lemma-pre-characteristic-p-case}에 의해
이 확대는 갈루아이고 $A$에 관하여 완전 분기이다. 따라서
$L_1 = L \otimes_K K_1$은 체이고 $A_1 \subset B_1$은 약하게 비분기이다
(보조정리 \ref{more-algebra-lemma-weakly-unramified-goes-up-along-totally-ramified}).
보조정리 \ref{more-algebra-lemma-pre-characteristic-p-case}에 의해 환
$M_1 = M \otimes_K K_1$은 $L_1$의 $p$개 복사본의 곱이거나(이 경우
증명이 끝난다), $L_1$의 차수 $p$인 체 확대이다. 또한 둘째 경우에는
$C_1$이 $B_1$ 위에서 약하게 비분기이거나(이 경우 증명이 끝난다),
$M_1/L_1$이 차수 $p$이고 갈루아이며 $B_1$에 관하여 완전 분기이다.
마지막 경우 확대 $M_1/L_1$은 원소 $z - w$로 생성되고
$$
(z - w)^p - (z - w) = z^p - z - (w^p - w) = b
$$
이다. 여기서 $b \in B$이다(위 참조). 보조정리
\ref{more-algebra-lemma-pre-characteristic-p-case}를 다시 적용하면 확대 $M_1/L_1$은
$B_1$에 관하여 비분기이다. 따라서 $K_1$이 $A \to C$의 약한 해라고
결론 내린다. 이제부터 $J \not = \emptyset$이라 가정한다.

\medskip\noindent
어떤 $r > 0$에 대하여 $j' = p^r j$가 되는 $j', j \in J$가 있다고
하자. $z$의 선택을 다음으로 바꾼다.
$$
z' = z -
(\sigma(\lambda_j) \pi^{-j} + \sigma(\lambda_j^p) \pi^{-pj} + \ldots +
\sigma(\lambda_j^{p^{r - 1}}) \pi^{-p^{r - 1}j})
$$
그러면 $\xi$는 다음과 같이 $\xi' = (z')^p - (z')$로 바뀐다.
$$
\xi' =
\xi - \sigma(\lambda_j) \pi^{-j} + \sigma(\lambda_j^{p^r}) \pi^{-j'}
+ \text{something in }B
$$
앞에서와 같이
$\xi' = \sum\nolimits_{i = n, \ldots, 1} \sigma(\lambda'_i) \pi^{-i} + b'$
이라 쓰면 $i \not = j, j'$에 대하여 $\lambda'_i = \lambda_i$이고
$\lambda'_j = 0$임을 얻는다. 따라서 집합 $J$가 작아졌다.
$J$의 크기에 관한 귀납에 의해 그런 쌍 $j, j'$이 없다고 가정할 수 있다.
(이 과정에서 위에서 다룬 $J = \emptyset$인 경우로 되돌아갈 수도 있음에
주의하라.)

\medskip\noindent
$j \in J$에 대하여 어떤 $r_j \geq 0$ 및 $p$제곱이 아닌
$\mu_j \in \kappa_B$를 사용하여
$\lambda_j = \mu_j^{p^{r_j}}$라 쓰자. 이는 가정 (4)에 의해 가능하다.
$j p^{-r_j}$가 최대가 되는 유일한 지수를 $j \in J$라 하자.
(앞 단락의 결과에 의해 이 지수는 유일하다.)
% P02-E0654: frozen source leaves the indexing of max(r_j + 1) implicit.
$r > \max(r_j + 1)$이고 모든 $j \in J$에 대하여
$j p^{r - r_j} > n$이 되도록 정수를 택하자. $A$에 관하여 완전 분기인
차수 $p^r$의 분리 확대 $K_1/K$를 택하되, 이에 대응하는 이산 값매김환
$A_1 \subset K_1$이 어떤 $a \in A_1$에 대하여
$(\pi')^{p^r} = \pi + \pi^{n + 1}a$를 만족하는 균등화원 $\pi'$를 갖게
하자(보조정리 \ref{more-algebra-lemma-make-degree-q-extension}).
$L_1 = L \otimes_K K_1$이 체이고 $L_1/L$이 $B$에 관하여 완전
분기임에 주의하자(보조정리
\ref{more-algebra-lemma-weakly-unramified-goes-up-along-totally-ramified}).
정수적 폐포 $B_1$ 안에서 계산하면
$$
\xi = \sum\nolimits_{i \in I} \sigma(\lambda_i) (\pi')^{-i p^r} +
\sum\nolimits_{j \in J} \sigma(\mu_j)^{p^{r_j}} (\pi')^{-j p^r} + b_1
$$
을 얻는다. 여기서 어떤 $b_1 \in B_1$이 존재한다. $i \in I$에 대한
% P02-E0655: frozen source calls this a qth power although q is not defined here.
$\sigma(\lambda_i)$는 $\pi^n$을 법으로 $q$제곱, 즉
$(\pi')^{n p^r}$을 법으로도 같은 꼴의 거듭제곱임에 주의하자. 따라서 위 식을 다음과
같이 다시 쓸 수 있다.
$$
\xi = \sum\nolimits_{i \in I} x_i^{p^r} (\pi')^{-i p^r} +
\sum\nolimits_{j \in J} \sigma(\mu_j)^{p^{r_j}} (\pi')^{- j p^r}
+ b_1
$$
앞 단락에서처럼 $z$의 선택을 다음으로 바꾼다.
% P02-E0656: frozen source's second telescoping sum is retained exactly.
\begin{align*}
z' & = z \\
& -
\sum\nolimits_{i \in I}
\left(x_i (\pi')^{-i} + \ldots + x_i^{p^{r - 1}} (\pi')^{-i p^{r - 1}}\right)
\\
& -
\sum\nolimits_{j \in J}
\left(
\sigma(\mu_j) (\pi')^{- j p^{r - r_j}}
+ \ldots +
\sigma(\mu_j)^{p^{r_j - 1}} (\pi')^{- j p^{r - 1}}
\right)
\end{align*}
그러면
$$
(z')^p - z' =
\sum\nolimits_{i \in I} x_i (\pi')^{-i} +
\sum\nolimits_{j \in J} \sigma(\mu_j) (\pi')^{- j p^{r - r_j}} + b_1'
$$
을 얻는다. 여기서 어떤 $b'_1 \in B_1$이 존재한다. 정확히 하나의 $j$에
대하여 $j p^{r - r_j}$가 최대이고, $j p^{r - r_j}$는 모든
$i \in I$보다 크며 $p$로 나누어떨어진다. 따라서 $M_1 / L_1$은
보조정리 \ref{more-algebra-lemma-pre-characteristic-p-case}의 경우 (C)에 해당한다.
이로써 증명이 끝난다.
\end{proof}

\begin{lemma}
\label{more-algebra-lemma-prepare}
$A$를 원시 $p$제곱근 $\zeta$를 포함하는 환이라 하자.
$w = 1 - \zeta$라 두자. 그러면
$$
P(z) = \frac{(1 + wz)^p - 1}{w^p} =
z^p - z + \sum\nolimits_{0 < i < p} a_i z^i
$$
은 $A[z]$의 원소이고, 실제로 $a_i \in (w)$이다. 또한 다항식환
$A[z_1, z_2]$ 안에서
$$
P(z_1 + z_2 + w z_1 z_2) = P(z_1) + P(z_2) + w^p P(z_1) P(z_2)
$$
가 성립한다.
\end{lemma}

\begin{proof}
원분체의 정수환이
$$
A = \mathbf{Z}[\zeta] = \mathbf{Z}[x]/(x^{p - 1} + \ldots + x + 1)
$$
인 경우만 증명하면 충분하다. 다항식 항등식
$t^p - 1 = (t - 1)(t - \zeta) \ldots (t - \zeta^{p - 1})$
(양변의 근을 살펴보면 증명된다)은
$t^{p - 1} + \ldots + t + 1 = (t - \zeta) \ldots (t - \zeta^{p - 1})$
임을 보여 준다. $t = 1$을 대입하면
$p = (1 - \zeta)(1 - \zeta^2) \ldots (1 - \zeta^{p - 1})$을 얻는다.
극대아이디얼 $(p, w) = (w)$는 $p$ 위에 놓이는 $A$의 유일한
소아이디얼이다(표수가 $p$인 체에는 자명하지 않은 $1$의 $p$제곱근이
없기 때문이다). 따라서 어떤 단위원 $u$에 대하여
$p = u w^{p - 1}$이다. 이는 $p > i > 1$에 대하여
$$
a_i = \frac{1}{p} {p \choose i} u w^{i - 1}
$$
이고 $- 1 + a_1 = pw/w^p = u$임을 뜻한다. $P(-1) = 0$이므로
$(w)$를 법으로 $0 = (-1)^p - u$이다. 따라서 $a_1 \in (w)$이고
% P02-E0657: frozen source's "proof if" typo is translated by its intended sense.
첫째 부분의 증명이 끝났다. 둘째 부분은 생략한 직접 계산에서 따른다.
\end{proof}

\begin{lemma}
% P02-E0658: frozen source label misspells polynomial as "polynial".
\label{more-algebra-lemma-extension-defined-by-nice-polynial}
$A$를 $1$의 원시 $p$제곱근을 포함하는 혼합 표수 $(0, p)$의 이산
값매김환이라 하자. $P(t) \in A[t]$를 보조정리 \ref{more-algebra-lemma-prepare}의
다항식이라 하고 $\xi \in K$라 하자. $L$을 $P(z) = \xi$의 근 하나를
첨가하여 얻는 $K$의 확대라 하자. 그러면 $L/K$는 갈루아 확대이고 다음
중 하나가 성립한다.
\begin{enumerate}
\item $L = K$이다.
\item $L/K$는 $A$에 관하여 비분기이고 차수가 $p$이다.
\item $L/K$는 $A$에 관하여 완전 분기이고 분기지수가 $p$이다.
\item $L$ 안에서 $A$의 정수적 폐포 $B$는 이산 값매김환이고,
$A \subset B$는 약하게 비분기이며, $A \to B$는 차수 $p$의 순수 비분리
잉여체 확대를 유도한다.
\end{enumerate}
$\pi$를 $A$의 균등화원이라 하자. 다음 함의들이 성립한다.
\begin{enumerate}
\item[(A)] $\xi \in A$이면 (1) 또는 (2)의 경우이다.
\item[(B)] $\xi = \pi^{-n}a$이고 $n > 0$이 $p$로 나누어떨어지지 않으며
$a$가 $A$의 단위원이면 (3)의 경우이다.
\item[(C)] $\xi = \pi^{-n} a$이고 $n > 0$이 $p$로 나누어떨어지며
$a$의 $\kappa_A$에서의 상이 $p$제곱이 아니면 (4)의 경우이다.
\end{enumerate}
\end{lemma}

\begin{proof}
% P02-E0659: frozen source's Kummer-equation equivalence is retained for canon review.
$P(z) = \xi$의 근을 첨가하는 것은 $y^p = w^p(1 + \xi)$의 근을
첨가하는 것과 같다. $K$가 $1$의 원시 $p$제곱근을 포함하므로 Fields,
절 \ref{fields-section-Kummer}의 논의에 의해 이 확대는 위수가 $p$의
약수인 갈루아 확대이다. 절 \ref{more-algebra-section-ramification}의 논의에서
보조정리에 나열한 (1)--(4) 가운데 한 경우가 성립함이 즉시 따른다.

\medskip\noindent
경우 (A). 여기서는 $A \to A[x]/(P(x) - \xi)$가 유한 \'etale 환
확대임을 알 수 있다. 따라서 (1) 또는 (2)의 경우이다.

\medskip\noindent
경우 (B). $p$가 $n$을 나누지 않도록 $\xi = \pi^{-n}a$라 쓰자.
$B \subset L$을 $L$ 안에서의 $A$의 정수적 폐포라 하자. 어떤
극대아이디얼 $\mathfrak m$에 대하여 $C = B_\mathfrak m$이면
$p \text{ord}_C(z) = -n \text{ord}_C(\pi)$임이 명백하다.
특히 $A \subset C$의 분기지수는 $p$의 배수이다. 따라서 그 지수는
$p$이고 $B = C$이다.

\medskip\noindent
경우 (C). $k = n/p$라 두자. 그러면 방정식을 다음과 같이 다시 쓸 수 있다.
$$
(\pi^kz)^p - \pi^{n - k} (\pi^kz) + \sum a_i \pi^{n - ik} (\pi^kz)^i = a
$$
% P02-E0660: frozen source's sign pattern in the defining quotient is retained.
$A[y]/(y^p - \pi^{n - k}y - \sum  a_i \pi^{n - ik} y^i - a)$가
$A$ 위에서 약하게 비분기인 이산 값매김환이므로 보조정리가 따른다.
\end{proof}

\noindent
$A$를 $1$의 원시 $p$제곱근을 포함하는 혼합 표수 $(0, p)$의 이산
값매김환이라 하자. $w \in A$와 $P(t) \in A[t]$를 보조정리
\ref{more-algebra-lemma-prepare}에서와 같이 두고, $L$을 $K$의 유한 확대라 하자.
$L$이 $K$의 차수 $p$인 확대이고, $w^p \xi \in \mathfrak m_A$인
어떤 $\xi \in K$에 대하여 방정식 $P(z) = \xi$의 근을 첨가하여
그 확대를 얻으면 $L/K$를 {\it 유한 층위의 차수 $p$ 확대}라고 한다.

\medskip\noindent
다음의 직접적인 보조정리 때문에 이 정의가 이 절의 논의와 관련된다.

\begin{lemma}
\label{more-algebra-lemma-make-finite-level}
$A \subset B \subset C$를 분수체가 $K \subset L \subset M$인 이산
값매김환의 확대라 하자. 다음을 가정한다.
\begin{enumerate}
\item $A$의 표수는 혼합 표수 $(0, p)$이다.
\item $A \subset B$는 약하게 비분기이다.
\item $B$는 $1$의 원시 $p$제곱근을 포함한다.
\item $M/L$은 차수 $p$인 갈루아 확대이다.
\end{enumerate}
그러면 $A$에 관하여 완전 분기인 유한 갈루아 확대 $K_1/K$가 존재하며,
이는 $A \to C$의 약한 해이거나 $M_1/L_1$을 유한 층위의 차수 $p$
확대로 만든다.
\end{lemma}

\begin{proof}
$\pi \in A$를 균등화원이라 하자. 쿠머 이론에 의해(Fields, 보조정리
\ref{fields-lemma-Kummer}) 어떤 $b \in L$에 대한 $y^p = b$의 근을
$L$에 첨가하여 $M$을 얻는다.

\medskip\noindent
$\text{ord}_B(b)$가 $p$와 서로소이면 $A$에 관하여 완전 분기인 차수
$p$의 분리 확대 $K_1/K$를 택한다(예를 들어 보조정리
\ref{more-algebra-lemma-make-degree-q-extension}을 쓴다). $A_1$을 $K_1$ 안에서의
$A$의 정수적 폐포라 하자. 보조정리
\ref{more-algebra-lemma-weakly-unramified-goes-up-along-totally-ramified}에 의해
$L_1 = L \otimes_K K_1$ 안에서의 $B$의 정수적 폐포 $B_1$은
$A_1$ 위에서 약하게 비분기인 이산 값매김환이다. $K_1/K$가
$A \to C$의 약한 해가 아니면, $M_1 = M \otimes_K K_1$ 안에서의
$C$의 정수적 폐포 $C_1$은 이산 값매김환이고 $B_1 \to C_1$의
분기지수는 $p$이다. 이 경우 $M_1$은 $p$로 나누어떨어지는
$\text{ord}_{B_1}(b)$를 갖는 $b$의 $p$제곱근을 $L_1$에 첨가하여 얻는다.
$A$를 $A_1$ 등으로 바꾸면 $u \in B$가 단위원이고 $n$이 $p$로
나누어떨어지는 $b = \pi^n u$의 경우라고 가정할 수 있다. 물론 이 경우
확대 $M$은 어떤 단위원의 $p$제곱근을 $L$에 첨가하여 얻는다.

\medskip\noindent
$M$이 $B$의 어떤 단위원 $u$에 대한 $y^p = u$의 근을 $L$에 첨가하여
얻어진다고 하자. $\kappa_B$에서 $u$의 잉여류가 $p$제곱이 아니면
$B \subset C$는 약하게 비분기이고(보조정리
\ref{more-algebra-lemma-pre-purely-inseparable-case}) 증명이 끝난다. 그렇지 않으면
$v^p$와 $u$의 $\kappa_B$에서의 상이 같도록 $y$의 선택을 $y/v$로
바꿀 수 있다. 이렇게 바꾸고 나면 어떤 $b \in B$에 대하여
$$
y^p = 1 + \pi b
$$
이다. 그러면 $z = (y - 1)/w$일 때 $P(z) = \pi b/ w^p$이다.
따라서 이 확대는 $\xi = \pi b / w^p$인 유한 층위의 차수 $p$
확대임을 알 수 있다.
\end{proof}

\noindent
$A$를 $1$의 원시 $p$제곱근을 포함하는 혼합 표수 $(0, p)$의 이산
값매김환이라 하자. $w \in A$와 $P(t) \in A[t]$를 보조정리
\ref{more-algebra-lemma-prepare}에서와 같이 두자. $L$을 $K$의 유한 층위의 차수
$p$ 확대라 하자. $L$을 $K$ 위에서 생성하고 $\xi = P(z) \in K$를
만족하는 $z \in L$을 택하자. $A$의 균등화원 $\pi$를 택하고,
어떤 정수 $e_1 = \text{ord}_A(w)$와 단위원 $u \in A$에 대하여
$w = u \pi^{e_1}$이라 쓰자. 마지막으로 다음이 성립하도록
$n \geq 0$을 택하자.
$$
\pi^n \xi \in A
$$
% P02-E0661: frozen source's "taking over" wording is translated by its intended sense.
$L/K$를 생성하고 $\xi = P(z) \in K$를 만족하는 모든 $z$에 대하여
양 $n/e_1$이 취하는 최소값을 $L/K$의 {\it 층위}라고 한다.

\medskip\noindent
몇 가지를 지적하자. 이 확대는 유한 층위이므로 $n < pe_1$이 되도록
$z$를 택할 수 있다. 따라서 층위는 $[0, p)$에 속하는 유리수이다.
층위가 영이면 보조정리 \ref{more-algebra-lemma-extension-defined-by-nice-polynial}에
의해 $L/K$는 $A$에 관하여 비분기이다. 다음 목표는 층위를 낮추는 것이다.

\begin{lemma}
\label{more-algebra-lemma-lowering-the-level}
$A \subset B \subset C$를 분수체가 $K \subset L \subset M$인 이산
값매김환의 확대라 하자. 다음을 가정한다.
\begin{enumerate}
\item $A$의 표수는 혼합 표수 $(0, p)$이다.
\item $A \subset B$는 약하게 비분기이다.
\item $B$는 $1$의 원시 $p$제곱근을 포함한다.
\item $M/L$은 유한 층위 $l > 0$의 차수 $p$ 확대이다.
\item $\kappa_A = \bigcap_{n \geq 1} \kappa_B^{p^n}$이다.
\end{enumerate}
그러면 $A$에 관하여 완전 분기인 $K$의 유한 분리 확대 $K_1$이 존재하여,
$K_1$은 $A \to C$의 약한 해이거나 확대 $M_1/L_1$은 층위가
$\leq \max(0, l - 1, 2l - p)$인 유한 층위의 차수 $p$ 확대이다.
\end{lemma}

\begin{proof}
$\pi \in A$를 균등화원이라 하자. $w \in B$와 $P \in B[t]$를
$B$에 대한 보조정리 \ref{more-algebra-lemma-prepare}에서와 같이 두자.
$e_1 = \text{ord}_B(w)$라 두면 $w$와 $\pi^{e_1}$은 $B$에서 서로
동반이다. $M$을 $L$ 위에서 생성하고 $\xi = P(z) \in K$를 만족하는
$z \in M$을 택하고, $M$의 $L$ 위 층위의 정의에서처럼
$\pi^n\xi \in B$ 및 $l = n/e_1$이 되도록 $n$을 택하자.

\medskip\noindent
이 보조정리의 증명은 보조정리 \ref{more-algebra-lemma-characteristic-p-case}의 증명과
완전히 비슷하다. 무슨 일이 일어나는지 설명하기 위해, 다음에 주의하자.
% P02-E0663: frozen source's hypothesis uses pi^{-n}P(z); retained exactly.
\begin{equation}
\label{more-algebra-equation-first-congruence}
P(z) \equiv z^p - z \bmod \pi^{-n + e_1}B
\end{equation}
이는 $\pi^{-n} P(z) \in B$를 만족하는 모든 $z \in L$에 대하여
성립한다($z$의 값매김이 최악의 경우 $-n/p$이고 다항식 $P$가 위의 꼴임을
사용하라). 또한 $\xi_1, \xi_2 \in \pi^{-n}B$에 대하여
\begin{equation}
\label{more-algebra-equation-second-congruence}
\xi_1 + \xi_2 + w^p \xi_1 \xi_2 \equiv \xi_1 + \xi_2 \bmod \pi^{-2n + pe_1}B 
\end{equation}
이다. 마지막으로 다음에 주의하자.
% P02-E0664: frozen source's level identity is dimensionally inconsistent; retained exactly.
$n - e_1 = (l - 1)/e_1$ 및 $-2n + pe_1 = -(2l - p)e_1$이다.
$m = n - e_1 \max(0, l - 1, 2l - p)$라 쓰자. 위 사실들은
$\pi^{-n}B / \pi^{-n + m}B$ 안에서 계산할 때 다항식 $P$가 다항식
$z^p -  z$와 정확히 똑같이 행동함을 보여 준다. 이것으로 보조정리가
참인 까닭을 설명할 수 있지만, 아래에 세부사항도 제시한다.

\medskip\noindent
% P02-E0662: frozen source refers to assumption (4), while the perfectness hypothesis is separately numbered.
가정 (4)는 $\kappa_A$가 완전체임을 뜻한다. $m \leq e_1$이고, 따라서
$A/\pi^m$은 $w$에 의해, 그러므로 $p$에 의해 영이 됨에 주의하자.
보조정리 \ref{more-algebra-lemma-cohen}에서와 같이 호환되는 환 사상
$\overline{\sigma} : \kappa_A \to A/\pi^mA$와
$\overline{\sigma} : \kappa_B \to B/\pi^mB$를 택할 수 있다.
이 가운데 둘째 사상을 집합의 사상 $\sigma : \kappa_B \to B$로 올리자.
그러면 어떤 $\lambda_i \in \kappa_B$와 $b \in B$에 대하여
% P02-E0665: frozen source's unmatched parenthesis in the exponent is retained.
$$
\xi =
\sum\nolimits_{i = n, \ldots, n - m + 1} \sigma(\lambda_i) \pi^{-i} +
\pi^{-n + m)} b
$$
라고 쓸 수 있다. 다음과 같이 두자.
$$
I = \{i \in \{n, \ldots, n - m + 1\} \mid \lambda_i \in \kappa_A\}
$$
그리고
% P02-E0666: frozen source uses lambda_i rather than the bound lambda_j in J.
$$
J = \{j \in \{n, \ldots, n - m + 1\} \mid \lambda_i \not \in \kappa_A\}
$$
유한 집합 $J$의 크기에 관하여 귀납하겠다.

\medskip\noindent
$J = \emptyset$인 경우. 이때 모든
$i \in \{n, \ldots, n - m + 1\}$에 대하여 $\overline{\sigma}$의 선택에
의해 어떤 $a_i \in A$와 $b_i \in B$가 존재하여
% P02-E0667: frozen source's lifting exponent is retained for canon review.
$\sigma(\lambda_i) = a_i + \pi^{n - m}b_i$가 된다. 따라서 어떤
$a \in A$와 $b \in B$에 대하여
$\xi = \pi^{-n} a + \pi^{-n + m} b$이다. $p | n$이면 $A$의 잉여체가
완전체이므로 어떤 $a_0, a_1 \in A$에 대하여
$a = a_0^p + \pi a_1$이라 쓰자. $z_1 = - \pi^{-n/p} a_0$라 두자.
$P(z_1) \in \pi^{-n}B$이고 $z + z_1 + w z z_1$이 $M$을 $L$ 위에서
생성하는 원소임에 주의하자($n < pe_1$이므로 $wz_1 \not = -1$임에
주의하라). 또한 보조정리 \ref{more-algebra-lemma-prepare}에 의해
$$
P(z + z_1 + w z z_1) = P(z) + P(z_1) + w^p P(z) P(z_1) \in K
$$
이고, 식 (\ref{more-algebra-equation-first-congruence})와
(\ref{more-algebra-equation-second-congruence})에 의해 어떤 $b' \in B$에 대하여
$$
P(z) + P(z_1) + w^p P(z) P(z_1)
\equiv
\xi + z_1^p - z_1 \bmod \pi^{-n + m}B
$$
이다.
% P02-E0668: frozen source introduces b' without it appearing in the congruence and has a grammar defect.
이는 $n$의 최소성에 모순이다! 따라서 $p$는 $n$을 나누지 않는다.
$P(y) = -\pi^{-n}a$로 주어지는 $K$의 차수 $p$ 확대 $K_1$을 생각하자.
보조정리 \ref{more-algebra-lemma-extension-defined-by-nice-polynial}에 의해 이 확대는
분리이고 $A$에 관하여 완전 분기이다. 따라서
$L_1 = L \otimes_K K_1$은 체이고 $A_1 \subset B_1$은 약하게
비분기이다(보조정리
\ref{more-algebra-lemma-weakly-unramified-goes-up-along-totally-ramified}).
보조정리 \ref{more-algebra-lemma-extension-defined-by-nice-polynial}에 의해 환
$M_1 = M \otimes_K K_1$은 $L_1$의 $p$개 복사본의 곱이거나(이 경우
증명이 끝난다), $L_1$의 차수 $p$인 체 확대이다. 또한 둘째 경우에는
$C_1$이 $B_1$ 위에서 약하게 비분기이거나(이 경우 증명이 끝난다),
$M_1/L_1$이 차수 $p$이고 갈루아이며 $B_1$에 관하여 완전 분기이다.
마지막 경우 확대 $M_1/L_1$은 원소 $z + y + wzy$로 생성되고,
$P(z + y + wzy) \in L_1$이며
\begin{align*}
P(z + y + wzy)
& = P(z) + P(y) + w^p P(z) P(y) \\
& \equiv
\xi - \pi^{-n}a \bmod \pi^{-n + m}B_1 \\
& \equiv
0 \bmod \pi^{-n + m}B_1
\end{align*}
이다. 이는 위와 정확히 같은 방식으로 얻는다. $m$의 선택에 의해 이것은
$M_1/L_1$의 층위가 $\max(0, l - 1, 2l - p)$ 이하라는 뜻이다.
이제부터 $J \not = \emptyset$이라 가정한다.

\medskip\noindent
어떤 $r > 0$에 대하여 $j' = p^r j$가 되는 $j', j \in J$가 있다고
하자. 다음과 같이 두자.
$$
z_1 = - \sigma(\lambda_j) \pi^{-j} - \sigma(\lambda_j^p) \pi^{-pj} -
\ldots - \sigma(\lambda_j^{p^{r - 1}}) \pi^{-p^{r - 1}j}
$$
그리고 $z$를 $z' = z + z_1 + wzz_1$로 바꾼다. $z' \in M$이 $M$을
$L$ 위에서 생성하고,
% P02-E0669: frozen source uses w rather than w^p in this polynomial identity.
$\xi' = P(z') = P(z) + P(z_1) + wP(z)P(z_1) \in L$이며
$$
\xi' \equiv
\xi - \sigma(\lambda_j) \pi^{-j} + \sigma(\lambda_j^{p^r}) \pi^{-j'}
\bmod \pi^{-n + m}B
$$
임에 주의하자. 이는 위에서처럼 식 (\ref{more-algebra-equation-first-congruence})와
(\ref{more-algebra-equation-second-congruence})를 사용하여 얻는다. 앞에서와 같이
$$
\xi' = \sum\nolimits_{i = n, \ldots, n - m + 1} \sigma(\lambda'_i) \pi^{-i}
+ \pi^{-n + m}b'
$$
이라 쓰면 $i \not = j, j'$에 대하여 $\lambda'_i = \lambda_i$이고
$\lambda'_j = 0$임을 얻는다. 따라서 집합 $J$가 작아졌다.
$J$의 크기에 관한 귀납에 의해 $j'/j$가 $p$의 거듭제곱이 되는
$J$의 쌍 $j, j'$이 없다고 가정할 수 있다. (이 과정에서 위에서 다룬
$J = \emptyset$인 경우로 되돌아갈 수도 있음에 주의하라.)

\medskip\noindent
$j \in J$에 대하여 어떤 $r_j \geq 0$ 및 $p$제곱이 아닌
$\mu_j \in \kappa_B$를 사용하여
$\lambda_j = \mu_j^{p^{r_j}}$라 쓰자. 이는 가정 (4)에 의해 가능하다.
$j p^{-r_j}$가 최대가 되는 유일한 지수를 $j \in J$라 하자.
(앞 단락의 결과에 의해 이 지수는 유일하다.)
% P02-E0670: frozen source leaves the indexing of max(r_j + 1) implicit.
$r > \max(r_j + 1)$이고 모든 $j \in J$에 대하여
$j p^{r - r_j} > n$이 되도록 택하자. $K_1/K$를
$(\pi')^{p^r} = \pi$로 정의되는, $A$에 관하여 완전 분기인 차수
$p^r$의 확대라 하자. $\pi'$은 이에 대응하는 이산 값매김환
$A_1 \subset K_1$의 균등화원임에 주의하자.
$L_1 = L \otimes_K K_1$은 체이고 $L_1/L$은 $B$에 관하여 완전
분기임에 주의하자(보조정리
\ref{more-algebra-lemma-weakly-unramified-goes-up-along-totally-ramified}).
정수적 폐포 $B_1$ 안에서 계산하면
$$
\xi = \sum\nolimits_{i \in I} \sigma(\lambda_i) (\pi')^{-i p^r} +
\sum\nolimits_{j \in J} \sigma(\mu_j)^{p^{r_j}} (\pi')^{-j p^r} +
\pi^{-n + m} b_1
$$
을 얻는다. 여기서 어떤 $b_1 \in B_1$이 존재한다. $i \in I$에 대한
% P02-E0671: frozen source calls this a qth power although q is not defined here.
$\sigma(\lambda_i)$는 $\pi^m$을 법으로 $q$제곱, 즉
$(\pi')^{m p^r}$을 법으로도 같은 꼴의 거듭제곱이다. 따라서 위 식을
다음과 같이 다시 쓸 수 있다.
$$
\xi = \sum\nolimits_{i \in I} x_i^{p^r} (\pi')^{-i p^r} +
\sum\nolimits_{j \in J} \sigma(\mu_j)^{p^{r_j}} (\pi')^{- j p^r}
+ \pi^{-n + m}b_1
$$
앞 단락의 선택과 비슷하게 다음과 같이 둔다.
% P02-E0672: frozen source's aligned display omits an equality sign after z_1.
% P02-E0673: frozen source's second telescoping sum is retained exactly.
\begin{align*}
z_1 & - \sum\nolimits_{i \in I}
\left(x_i (\pi')^{-i} + \ldots + x_i^{p^{r - 1}} (\pi')^{-i p^{r - 1}}\right)
\\
& - \sum\nolimits_{j \in J}
\left(
\sigma(\mu_j) (\pi')^{- j p^{r - r_j}}
+ \ldots +
\sigma(\mu_j)^{p^{r_j - 1}} (\pi')^{- j p^{r - 1}}
\right)
\end{align*}
그리고 $z$의 선택을 $z' = z + z_1 + wzz_1$로 바꾼다. 그러면 $z'$은
$M_1$을 $L_1$ 위에서 생성하고
$\xi' = P(z') = P(z) + P(z_1) + w^p P(z) P(z_1) \in L_1$이며,
계산하면
$$
\xi' \equiv
\sum\nolimits_{i \in I} x_i (\pi')^{-i} +
\sum\nolimits_{j \in J} \sigma(\mu_j) (\pi')^{- j p^{r - r_j}} +
(\pi')^{(-n + m)p^r}b'_1
$$
을 얻는다. 여기서 어떤 $b'_1 \in B_1$이 존재한다. 정확히 하나의 $j$에
대하여 $j p^{r - r_j}$가 최대이고, $j p^{r - r_j}$는 모든
$i \in I$보다 크다. $j p^{r - r_j} \leq (n - m)p^r$이면 확대
$M_1/L_1$의 층위는
% P02-E0674: frozen source says "less than" although the lemma states a weak inequality.
$\max(0, l - 1, 2l - p)$보다 작다. 그렇지 않으면 $p$가
$j p^{r - r_j}$를 나누므로 $M_1 / L_1$은 보조정리
\ref{more-algebra-lemma-extension-defined-by-nice-polynial}의 경우 (C)에 해당한다.
이로써 증명이 끝난다.
\end{proof}

\begin{lemma}
\label{more-algebra-lemma-special-case}
$A \subset B \subset C$를 분수체가 $K \subset L \subset M$인 이산
값매김환의 확대라 하자. 다음을 가정한다.
\begin{enumerate}
\item $A$의 잉여체 $k$는 표수가 $p > 0$인 대수적으로 닫힌 체이다.
\item $A$와 $B$는 완비이다.
\item $A \to B$는 약하게 비분기이다.
\item $M$은 $L$의 유한 확대이다.
\item $k = \bigcap\nolimits_{n \geq 1} \kappa_B^{p^n}$이다.
\end{enumerate}
그러면 $A \to C$의 약한 해인 유한 확대 $K_1/K$가 존재한다.
\end{lemma}

\begin{proof}
$M'$을 $L$의 임의의 유한 확대라 하고, $M'$ 안에서의 $B$의 정수적 폐포
$C'$을 생각하자. Algebra, 보조정리
\ref{algebra-lemma-Noetherian-complete-local-Nagata}에 의해 $B$는
나가타 환이므로 $C'$은 $B$ 위에서 유한이다. 또한 $C'$은 이산
값매김환이다. 평주 \ref{more-algebra-remark-construction}의 논의를 보라. 또 $C'$은
$B$-가군으로서 완비이므로 이산 값매김환으로서 완비이다. Algebra, 절
\ref{algebra-section-completion}을 보라. 특히 $C$는 $M$ 안에서의 $B$의
정수적 폐포이다(지배 관계에 관하여 값매김환이 극대라는 정의에 의해).

\medskip\noindent
$M \subset M'$을 유한 확대라 하고 $C' \subset M'$을 위와 같은 $B$의
정수적 폐포라 하자. 보조정리 \ref{more-algebra-lemma-solutions-go-down}에 의해
$A \to B \to C'$에 대하여 결과를 증명하면 충분하다. 따라서 Fields,
보조정리 \ref{fields-lemma-normal-closure}에 의해 $M/L$이 정규라고
가정할 수 있다.

\medskip\noindent
$M / L$이 정규이면 다음과 같은 유한 확대의 사슬을 찾을 수 있다.
$$
L = L^0 \subset L^1 \subset L^2 \subset \ldots \subset L^r = M
$$
여기서 각 확대 $L^{j + 1}/L^j$는 다음 중 하나이다.
\begin{enumerate}
\item[(a)] 차수가 $p$인 순수 비분리 확대이다.
\item[(b)] $B^j$에 관하여 완전 분기이고 차수가 $p$인 갈루아 확대이다.
\item[(c)] $B^j$에 관하여 완전 분기이고 위수가 $p$와 서로소인 순환
갈루아 확대이다.
\item[(d)] $B^j$에 관하여 비분기인 갈루아 확대이다.
\end{enumerate}
여기서 $B^j$는 $L^j$ 안에서의 $B$의 정수적 폐포이다. 실제로
$M/L$이 정규이므로 이를 갈루아 확대와 순수 비분리 확대의 합성체로 쓸 수
있다(Fields, 보조정리 \ref{fields-lemma-normal-case}). 순수 비분리
확대의 경우 필터의 존재는 명백하다. 갈루아 경우에는 $G$가 ``그''
분해군임에 주의하고 $I \subset G$를 관성군이라 하자. 한편 보조정리
\ref{more-algebra-lemma-galois-inertia}에 의해 $I$는 가해군이고, 다른 한편 보조정리
\ref{more-algebra-lemma-inertial-invariants-unramified}에 의해 확대 $M^I/L$은 $B$에
관하여 비분기이다. 이로써 말한 필터가 존재함을 증명했다.

\medskip\noindent
정수 $r$에 관하여 귀납하겠다. $B^1$이 $L^1$ 안에서의 $B$의 정수적
폐포일 때 $A \to B^1$의 약한 해인 유한 확대 $K_1/K$를 찾을 수 있다고
하자. $K'_1$을 $K_1/K$의 정규 폐포라 하자(Fields, 보조정리
\ref{fields-lemma-normal-closure}). $A$가 완비이고 $A$의 잉여체가
대수적으로 닫혀 있으므로 $K'_1/K_1$은 $A_1$에 관하여 분리이고 완전
분기이다(세부사항 일부는 생략한다). 따라서 보조정리
\ref{more-algebra-lemma-weakly-unramified-goes-up-along-totally-ramified}에 의해
$K'_1/K$도 $A \to B^1$의 약한 해이다. 다시 말해 $K_1$이 $K$의 정규
확대라고 가정할 수 있고, 그렇게 가정한다. 이제 다음 열을 생각하자.
$$
L^0_1 = (L^0 \otimes_K K_1)_{red} \subset
L^1_1 = (L^1 \otimes_K K_1)_{red} \subset \ldots \subset
L^r_1 = (L^r \otimes_K K_1)_{red}
$$
그리고 이에 대응하는 정수적 폐포들을 $B^i_1$이라 하자.
$C_1 = B^r_1$은 이산 값매김환들의 곱이고, 보조정리
\ref{more-algebra-lemma-galois-relative}에 의해 $G = \text{Aut}(K_1/K)$가 이들을
추이적으로 치환한다. 특히 이산 값매김환의 모든 확대
$A_1 \to (C_1)_\mathfrak m$은 서로 동형이고, 그중 하나에 대한 약한
해는 모두에 대한 약한 해이다. 다음 열에 귀납 가설을 적용할 수 있다.
$$
A_1 \to (B^1_1)_{B^1_1 \cap \mathfrak m} \to
(B^2_1)_{B^2_1 \cap \mathfrak m} \to
\ldots \to
(B^r_1)_{B^r_1 \cap \mathfrak m} =
(C_1)_\mathfrak m
$$
그러면 $A_1 \to (C_1)_\mathfrak m$의 약한 해 $K_2/K_1$을 얻는다.
앞에서 말한 바에 의해 확대 $K_2/K$는 $A \to C$의 약한 해이다.
귀납 가설이 적용됨에 주의하자. $K_1$의 선택에 의해 환 사상
$A_1 \to (B^1_1)_{B^1_1 \cap \mathfrak m}$은 약하게 비분기이고,
분수체 확대의 열에 있는 각 확대에는 여전히 위에 나열한 (a), (b), (c),
(d) 가운데 한 성질이 있다. 또한 임의의 유한 확대
$\kappa_B \subset \kappa$에 대하여도 여전히
$k = \bigcap \kappa^{p^n}$임에 주의하자.

\medskip\noindent
따라서 모든 것은 체 확대 $M/L$이 성질 (a), (b), (c), (d) 가운데
하나를 만족할 때 $A \subset C$의 약한 해를 찾는 문제로 귀착된다.

\medskip\noindent
경우 (d). 이 경우에는 이미 $B \to C$가 비분기이므로 자명하다.

\medskip\noindent
경우 (c). $M/L$이 $p$와 서로소인 위수 $n$의 순환 확대라 하자.
$M/L$이 $B$에 관하여 완전 분기이므로 $B \subset C$의 분기지수는
$n$이고, 따라서 $A \subset C$의 분기지수도 $n$이다. 균등화원
$\pi \in A$를 택하고 $K_1 = K[\pi^{1/n}]$라 두자. 아브얀카르
보조정리에 의해(보조정리 \ref{more-algebra-lemma-abhyankar}) $K_1/K$는
$A \subset C$의 해이다.

\medskip\noindent
경우 (b). 이 경우를 혼합 표수의 경우와 등표수의 경우로 나눈다.
등표수의 경우는 보조정리 \ref{more-algebra-lemma-characteristic-p-case}이다.
혼합 표수의 경우에는 먼저 보조정리 \ref{more-algebra-lemma-make-finite-level}을
사용하여 $K$를 유한 확대 하나로 바꾸고, $M/L$이 유한 층위의 차수
$p$ 확대인 상황으로 간다. 그러면 층위는 유리수 $l \in [0, p)$이다.
보조정리 \ref{more-algebra-lemma-lowering-the-level} 앞의 논의를 보라. 층위가
$0$이면 $B \to C$는 약하게 비분기이므로 증명이 끝난다. 그렇지 않으면
보조정리 \ref{more-algebra-lemma-lowering-the-level}에 의해 체 $K$를 유한 확대 하나로
% P02-E0675: frozen source's "can replacing" grammar is translated by its intended sense.
바꾸어 층위가 $l' \leq \max(0, l - 1, 2l - p)$인 새 상황을 얻을 수 있다.
$\epsilon < 1$인 $l = p - \epsilon$이면
$l' \leq p - 2\epsilon$이다. 따라서 유한 번 바꾸고 나면 층위가
$\leq p - 1$인 경우를 얻는다. 그런 교체를 많아야 $p - 1$번 더 하면
층위가 영인 상황에 도달한다.

\medskip\noindent
경우 (a)는 보조정리 \ref{more-algebra-lemma-purely-inseparable-case}이다. 이 경우에만
$K$의 순수 비분리 확대가 필요할 수 있다. 즉 보조정리의 명제에서 경우
% P02-E0676: frozen source says adjoining a pth power of pi, rather than a pth root.
(2)에는 원소 $\pi$의 $p$제곱을 첨가하면 이긴다. 이로써 보조정리의
증명이 끝난다.
\end{proof}

\noindent
이제 이 절의 주 결과를 증명하는 데 필요한 모든 보조정리를 모았다.

\begin{theorem}[Epp]
\label{more-algebra-theorem-epp}
$A \subset B$를 분수체가 $K \subset L$인 이산 값매김환의 확대라 하자.
$\kappa_A$의 표수가 $p > 0$이면
$$
\bigcap\nolimits_{n \geq 1} \kappa_B^{p^n}
$$
의 모든 원소가 $\kappa_A$ 위에서 분리 대수적이라고 가정하자.
그러면 정의 \ref{more-algebra-definition-solution}에서 정의한 $A \to B$의 약한 해인
유한 확대 $K_1/K$가 존재한다.
\end{theorem}

\begin{proof}
$\kappa_A$의 표수가 영인 경우, 또는 잉여 표수가 $p$이고 분기지수가
$p$와 서로소이며 잉여체 확대가 분리인 경우에는 아브얀카르 보조정리
(보조정리 \ref{more-algebra-lemma-abhyankar})에서 결론이 따른다. 실제로 분기지수가
$e$라고 하자. 균등화원 $\pi \in A$를 택하자. $K_1/K$를 $\pi$의
$e$제곱근 하나를 첨가하여 얻는 확대라 하자. 보조정리
\ref{more-algebra-lemma-pull-root-uniformizer}에 의해 $K_1$ 안에서의 $A$의 정수적
% P02-E0677: frozen source omits the value after "ramification index over A".
폐포 $A_1$은 $A$ 위에서 분기지수를 갖는 이산 값매김환이다. 따라서
보조정리 \ref{more-algebra-lemma-abhyankar}에 의해 $B_1$의 모든 극대아이디얼
$\mathfrak m$에 대하여 $A_1 \to (B_1)_\mathfrak m$은
$\mathfrak m$-진 위상에서 형식적으로 매끄럽고, 특히 이들은 이산
값매김환의 약하게 비분기인 확대이다.

\medskip\noindent
이제부터 $p$를 소수라 하고 $\kappa_A$의 표수가 $p$라고 가정한다.
먼저 보조정리 \ref{more-algebra-lemma-solution-after-strict-henselization}을 적용하여
$A$와 $B$의 잉여체가 분리적으로 닫힌 경우로 환원한다. 이 과정에서
$\kappa_A$와 $\kappa_B$가 각각 그 분리 대수적 폐포로 바뀌므로 정리의
조건으로부터
$$
\kappa_A \supset \bigcap\nolimits_{n \geq 1} \kappa_B^{p^n}
$$
을 얻는다.

\medskip\noindent
$\pi \in A$를 균등화원이라 하자. $A^\wedge$와 $B^\wedge$를 각각
$A$와 $B$의 완비화라 하자. 이산 값매김환의 확대들로 이루어진 가환 그림
$$
\xymatrix{
B \ar[r] & B^\wedge \\
A \ar[u] \ar[r] & A^\wedge \ar[u]
}
$$
이 있다. $K^\wedge$를 $A^\wedge$의 분수체라 하자. (a)
$A^\wedge \to B^\wedge$의 약한 해이고, (b) 분리 확대와 $\pi$의
$p$의 거듭제곱근을 첨가하여 얻는 확대의 합성체인 유한 확대
$M/K^\wedge$를 찾을 수 있다고 하자. 그러면 보조정리
\ref{more-algebra-lemma-approximate-separable-extension}에 의해
$K^\wedge \otimes_K K_1 = M$이 되게 하는 유한 확대 $K_1/K$를 찾을 수
있다. $A_1$과 $A_1^\wedge$를 각각 $K_1$ 안에서의 $A$의 정수적 폐포와
$M$ 안에서의 $A^\wedge$의 정수적 폐포라 하자.
$A \to A^\wedge$가 $\mathfrak m^\wedge$-진 위상에서 형식적으로
매끄러우므로(보조정리 \ref{more-algebra-lemma-extension-dvrs-formally-smooth}),
$A_1 \to A_1^\wedge$는 $\mathfrak m_1^\wedge$-진 위상에서 형식적으로
매끄럽다(보조정리 \ref{more-algebra-lemma-formally-smooth-goes-up}; 평주
\ref{more-algebra-remark-construction}의 논의에 의해 $A_1$과 $A_1^\wedge$는 이산
값매김환이다). 보조정리 \ref{more-algebra-lemma-solutions-go-down}의 (2)에 의해
$K_1/K$는 $A \to B^\wedge$의 약한 해이다. 보조정리
\ref{more-algebra-lemma-solutions-go-down}의 (1)을 적용하면 $K_1/K$가 $A \to B$의
약한 해임을 알 수 있다.

\medskip\noindent
따라서 $A$와 $B$가 표수 $p$인 분리적으로 닫힌 잉여체를 갖는 완비
이산 값매김환이고
$\kappa_A \supset \bigcap\nolimits_{n \geq 1} \kappa_B^{p^n}$이라고
가정할 수 있다. 또한 균등화원 $\pi \in A$가 주어졌고, 분리 확대와
$\pi$의 $p$의 거듭제곱근을 취하여 얻는 체의 합성체인 $A \to B$의
약한 해를 찾아야 한다. $A$가 혼합 표수이면 둘째 조건은 자동임에
주의하자.

\medskip\noindent
$k = \bigcap\nolimits_{n \geq 1} \kappa_B^{p^n}$라 두자. $k$는 표수가
$p$인 대수적으로 닫힌 체임에 주의하자. $A$가 혼합 표수이면
$\Lambda$를 $k$의 코언 환이라 하고, 등표수의 경우에는
$\Lambda = k[[t]]$라 두자. 등표수의 경우 $t$를 $\pi$로 보내는 환 사상
$\Lambda \to A$를 택할 수 있다. 등표수의 경우 이는 코언 구조 정리에서
따르고(Algebra, 정리 \ref{algebra-theorem-cohen-structure-theorem}),
혼합 표수의 경우에는 $\mathbf{Z}_p \to \Lambda$가 진 위상에서
형식적으로 매끄럽다는 사실에서 따른다(보조정리
\ref{more-algebra-lemma-extension-dvrs-formally-smooth}와
\ref{more-algebra-lemma-lift-continuous}). 보조정리 \ref{more-algebra-lemma-solutions-go-down}을
적용하면 $\Lambda \to B$의 약한 해가 존재함을 증명하는 것으로 충분하다.
등표수 $p$의 경우에는 이 해가 분리 확대와 $t$의 $p$의 거듭제곱근을
취하여 얻는 체의 합성체여야 한다. 그러나 등표수의 경우
$\Lambda = k[[t]]$이고 $k((t))$의 임의의 확대가 그런 합성체이므로 이제
이 요구를 버릴 수 있다!

\medskip\noindent
따라서 $A$와 $B$가 완비이고, $A$의 잉여체 $k$가 표수 $p > 0$인
대수적으로 닫힌 체이며, $k = \bigcap \kappa_B^{p^n}$이고, 혼합 표수의
경우 $p$가 $A$의 균등화원인 상황에 이르렀다(즉 $A$는 $k$의 코언
환이다). $A$가 혼합 표수이면 $\Lambda$를 $\kappa_B$의 코언 환이라
하고, 등표수의 경우 $\Lambda = \kappa_B[[t]]$라 두자. 위와 같이
논증하여 $k \to \kappa_B$를 올리고 균등화원을 균등화원으로 보내는
환 사상 $A \to \Lambda$를 택할 수 있다. $k \subset \kappa_B$가
분리이므로 환 사상 $A \to \Lambda$는 진 위상에서 형식적으로 매끄럽다
(보조정리 \ref{more-algebra-lemma-extension-dvrs-formally-smooth}). 따라서 합성
$A \to \Lambda \to B$가 주어진 환 사상 $A \to B$가 되도록 하는
환 사상 $\Lambda \to B$를 찾을 수 있다(보조정리
\ref{more-algebra-lemma-lift-continuous}를 보라). $\Lambda$와 $B$는 잉여체가 같은
완비 이산 값매김환이므로 $B$는 $\Lambda$ 위에서 유한이다(Algebra,
보조정리 \ref{algebra-lemma-finite-over-complete-ring}). 이로써 보조정리
\ref{more-algebra-lemma-special-case}에서 논한 특수한 경우로 환원된다.
\end{proof}

\section{분기 제거, II}
\label{more-algebra-section-eliminating-ramification-bis}

\noindent
이 절에서는 절 \ref{more-algebra-section-eliminating-ramification}의 결과를 사용하여
몇 가지 경우에 (분리) 해를 얻는다.

\begin{lemma}
\label{more-algebra-lemma-solution-goes-up}
$A \to B$를 분수체가 $K \subset L$인 이산 값매김환의 확대라 하자.
$K_1/K$가 $A \subset B$의 해이면, 임의의 유한 확대 $K_2/K_1$에 대하여
확대 $K_2/K$는 $A \subset B$의 해이다.
\end{lemma}

\begin{proof}
이는 보조정리 \ref{more-algebra-lemma-formally-smooth-goes-up}에서 따른다.
세부사항은 생략한다.
\end{proof}

\begin{lemma}
\label{more-algebra-lemma-nagata-goes-down}
$A \subset B$를 이산 값매김환의 확대라 하자. $B$가 나가타이고 분수체의
확대 $L/K$가 분리이면 $A$는 나가타이다.
\end{lemma}

\begin{proof}
이산 값매김환이 나가타인 것과 N-2인 것은 동치이다. $K_1/K$를 유한
순수 비분리 체 확대라 하자. $K_1$ 안에서 $A$의 정수적 폐포 $A_1$이
$A$ 위에서 유한임을 보이면 된다. Algebra, 보조정리
\ref{algebra-lemma-domain-char-p-N-1-2}를 보라. $L/K$는 분리이고
$K_1/K$는 순수 비분리이므로 대수 $L \otimes_K K_1$은 체이다(Algebra,
보조정리 \ref{algebra-lemma-separable-extension-preserves-reducedness}와
\ref{algebra-lemma-radicial-integral-bijective}에 의해). $B_1$을
$L \otimes_K K_1$ 안에서의 $B$의 정수적 폐포라 하자. $B$가
나가타이므로 $B_1$은 $B$ 위에서 유한이다. 또
$B \otimes_A A_1 \subset B_1$이고 $B$가 뇌터이므로
$B \otimes_A A_1$은 $B$ 위에서 유한이다. $A \to B$가 충실하게
평탄하므로 이는 $A_1$이 $A$ 위에서 유한임을 뜻한다. Algebra, 보조정리
\ref{algebra-lemma-descend-properties-modules}를 보라.
\end{proof}

\begin{lemma}
\label{more-algebra-lemma-construct-extension}
$A' \subset A$를 환의 확대라 하고 $f \in A'$라 하자. (a) $A$가
$A'$ 위에서 유한이고, (b) $f$가 $A$ 위의 영인자가 아니며, (c)
$A'_f = A_f$라고 가정하자. 그러면 어떤 정수 $n_0 > 0$이 존재하여
모든 $n \geq n_0$에 대하여 다음이 성립한다. 환 $B'$, 영인자가 아닌
$g \in B'$, 그리고 $\varphi'(f) \equiv g$를 만족하는 동형
$\varphi' : A'/f^n A' \to B'/g^n B'$가 주어지면, 유한 확대
$B' \subset B$와 $\varphi'$와 호환되는 동형
$\varphi : A/fA \to B/gB$가 존재한다.
\end{lemma}

\begin{proof}
% P02-E0678: frozen source's "Cchoose" typo is translated by its intended sense.
$A$가 $A'$ 위에서 유한이고 $A'_f = A_f$이므로
$f^t A \subset A'$이 되게 하는 $t > 0$을 택할 수 있다.
$n_0 = 2t$라 두자. 보조정리의 명제에 있는 $n, B', g, \varphi'$가
주어졌을 때, $b \bmod g^nB' \in \varphi'(f^tA)$를 만족하는
$b \in B'$들의 집합을 $N \subset B'$로 쓰자. $B = g^{-t}N$이라 두자.
$f^tA' \subset f^tA$이고 $\varphi'$가 $f$를 $g$로 보내므로
$g^tB' \subset N$이고, 따라서 $B' \subset B$이다.
$f^tA \cdot f^tA \subset f^t \cdot f^tA$이고 $\varphi'$가 $f$를
$g$로 보내므로 $N \cdot N \subset g^t N$이다. 따라서 $B'$의 곱셈을
연장하는 $B$ 위의 곱셈을 얻는다. 다음 $A'/f^nA'$-가군 동형이 있다.
% P02-E0679: frozen source's displayed module chain is retained for canon review.
$$
A/f^tA' \xrightarrow{f^t} f^tA/f^nA' \xrightarrow{\varphi'}
g^tB/g^nB' \xrightarrow{g^{-t}} B/g^tB'
$$
여기서 오른쪽의 가군 구조는 $\varphi'$를 사용하여 정의한다.
$A/f^tA'$은 유한 $A'$-가군이므로 $B/g^tB'$은 유한 $B'$-가군이고,
따라서 $B' \to B$가 유한임을 알 수 있다. 마지막으로, 표시된 가군
동형이 $fA$를 $gB$ 안으로 보내고 $\varphi'$와 호환되는 환 동형
$\varphi : A/fA \to B/gB$를 유도함은 독자에게 맡긴다(실제로 동형
$A/f^tA \to B/g^tB$도 유도하지만 이것은 필요하지 않다).
\end{proof}

\begin{remark}
\label{more-algebra-remark-functoriality-construct-extension}
보조정리 \ref{more-algebra-lemma-construct-extension}의 구성은 다음 ``함자성''을
만족한다. 수평 화살표들이 단사인 가환 그림
$$
\xymatrix{
A'_2 \ar[r] & A_2 \\
A'_1 \ar[r] \ar[u] & A_1 \ar[u]
}
$$
이 있다고 하자. 원소 $f \in A'_1$이 주어져
$(A'_1 \subset A_1, f)$와 $(A'_2 \subset A_2, f)$가 보조정리
\ref{more-algebra-lemma-construct-extension}의 성질 (a), (b), (c)를 만족한다고 하자.
$n_{0, 1}$과 $n_{0, 2}$를 이 두 상황에서 보조정리가 주는 정수들이라 하자.
마지막으로 환 사상 $B'_1 \to B'_2$가 주어지고,
% P02-E0680: frozen source places g in B'_1 but states regularity on B_1 and B_2.
$g \in B'_1$이 $B_1$과 $B_2$ 위의 영인자가 아니며,
$n \geq \max(n_{0, 1}, n_{0, 2})$이고, 수평 화살표들이 동형이며
$\varphi'_1(f) \equiv g$인 가환 그림
% P02-E0681: frozen source repeats B'_2 in both right-hand vertices.
$$
\xymatrix{
A'_2/f^nA'_2 \ar[r]_{\varphi'_2} & B'_2/g^nB'_2 \\
A'_1/f^nA'_1 \ar[r]^{\varphi'_1} \ar[u] & B'_2/g^nB'_2 \ar[u]
}
$$
이 주어졌다고 하자. 그러면 다음 가환 그림들을 얻는다.
% P02-E0682: frozen source repeats B_2 in both lower and upper right quotient vertices.
$$
\vcenter{
\xymatrix{
B'_2 \ar[r] & B_2 \\
B'_1 \ar[r] \ar[u] & B_1 \ar[u]
}
}
\quad\text{and}\quad
\vcenter{
\xymatrix{
A_2/fA_2 \ar[r]_{\varphi_2} & B_2/gB_2 \\
A_1/fA_1 \ar[r]^{\varphi_1} \ar[u] & B_2/gB_2 \ar[u]
}
}
$$
여기서 $(B'_1 \subset B_1, \varphi_1)$과
$(B'_2 \subset B_2, \varphi_2)$는 보조정리
\ref{more-algebra-lemma-construct-extension}의 증명에서와 같이 구성한다. 자세한 검증은
생략한다.
\end{remark}


\begin{lemma}
\label{more-algebra-lemma-approximate-solution}
$p$를 소수라 하자. $A \subset B$를 분수체의 확대가 $L/K$인 이산
값매김환의 확대라 하자. $K_2/K_1/K$를 유한 체 확대의 탑이라 하자.
다음을 가정한다.
\begin{enumerate}
\item $K$의 표수는 $p$이다.
\item $L/K$는 분리이다.
\item $B$는 나가타이다.
\item $K_2$는 $A \subset B$의 해이다.
\item $K_2/K_1$은 차수가 $p$인 순수 비분리 확대이다.
\end{enumerate}
그러면 $A \subset B$의 해인 분리 확대 $K_3/K_1$이 존재한다.
\end{lemma}

\begin{proof}
평주 \ref{more-algebra-remark-construction}에서와 같은 기호를 쓰며, 거기서 한 모든
관찰을 사용하겠다. $L/K$가 분리이므로 대수
$L_1 = L \otimes_K K_1$은 기약이다(Algebra, 보조정리
\ref{algebra-lemma-separable-extension-preserves-reducedness}).
$B$가 나가타이므로 환 확대 $B \subset B_1$은 유한이다. 여기서 $B_1$은
$L_1$ 안에서의 $B$의 정수적 폐포이고 $B_1$은 나가타 환이다. 마찬가지로 보조정리
\ref{more-algebra-lemma-nagata-goes-down}에 의해 환 $A$는 나가타이므로
$A \subset A_1$은 유한이고 $A_1$도 나가타 환이다. 또한 같은 단언들이
$K_2$에 대해서도 참이다. 즉 $L_2 = L \otimes_K K_2$는 기약이고,
환 확대 $A_1 \subset A_2$와 $B_1 \subset B_2$는 유한이다. 여기서
$A_2$와 $B_2$는 각각 $K_2$와 $L_2$ 안에서의 $A$와 $B$의 정수적
폐포이다.

\medskip\noindent
$\pi \in A$를 균등화원이라 하자. $\pi$는 $K_1$, $K_2$, $A_1$,
$A_2$, $L_1$, $L_2$, $B_1$, $B_2$ 위의 영인자가 아니며,
$K_1 = (A_1)_\pi$, $K_2 = (A_2)_\pi$, $L_1 = (B_1)_\pi$,
$L_2 = (B_2)_\pi$임에 주의하자. Fields, 보조정리
\ref{fields-lemma-finite-purely-inseparable}에 의해
$\alpha^p = a_1 \in K_1$인 $K_2 = K_1(\alpha)$로 쓸 수 있다.
$\alpha$에 $\pi$의 거듭제곱을 곱하여 $a_1 \in A_1$이라고 가정할 수
있고, 그렇게 가정한다. 증명의 나머지에서는 편의상
$K_2 = K_1[x]/(x^p - a_1)$ 및 $L_2 = L_1[x]/(x^p - a_1)$로 쓰겠다.
다음 환의 확대들을 생각하자.
$$
A'_2 = A_1[x]/(x^p - a_1) \subset A_2
\quad\text{and}\quad
B'_2 = B_1[x]/(x^p - a_1) \subset B_2
$$
보조정리 \ref{more-algebra-lemma-construct-extension}을 $A'_2 \subset A_2$와
$f = \pi^2$에, 그리고 $B'_2 \subset B_2$와 $f = \pi^2$에 적용할 수
있다. 이 둘에 모두 적용될 만큼 충분히 큰 정수 $n$을 택하자.

\medskip\noindent
다음 대수들을 생각하자.
% P02-E0685: frozen source's 2n/n perturbation pattern is retained for canon review.
$$
K_3 = K_1[x]/(x^p - \pi^{2n} x - a_1)
\quad\text{and}\quad
L_3 = L_1[x]/(x^p - \pi^{2n} x - a_1)
$$
$K_3/K_1$과 $L_3/L_1$이 차수 $p$인 유한 \'etale 대수 확대임에
주의하자. 다음 $K_3 = (A'_2)_\pi$와 $L_3 = (B'_3)_\pi$의 부분환들을
생각하자.
% P02-E0683: frozen source changes the perturbation exponent from 2n to n.
$$
A'_3 = A_1[x]/(x^p - \pi^n x - a_1)
\quad\text{and}\quad
B'_3 = B_1[x]/(x^p - \pi^n x - a_1)
$$
% P02-E0684: frozen source identifies K_3 with (A'_2)_pi while pairing L_3 with (B'_3)_pi.
다음 가환 그림을 구성하겠다.
$$
\xymatrix{
B'_2/\pi^{2n} B'_2 \ar[r]_{\psi'} & B'_3/\pi^{2n} B'_3 \\
A'_2/\pi^{2n} A'_2 \ar[r]^{\varphi'} \ar[u] & A'_3/\pi^{2n} A'_3 \ar[u]
}
$$
실제로 $\varphi'$는 $x$의 잉여류를 $x$의 잉여류로 보내는 유일한
$A_1$-대수 동형이다. 마찬가지로 $\psi'$는 $x$의 잉여류를 $x$의
잉여류로 보내는 유일한 $B_1$-대수 동형이다. $n$의 선택에 따라 보조정리
\ref{more-algebra-lemma-construct-extension}과 평주
\ref{more-algebra-remark-functoriality-construct-extension}을 통하여 유한 환 확대
$A'_3 \subset A_3$과 $B'_3 \subset B_3$을 얻는다. 이때
$A'_3 \to B'_3$은 환 사상 $A_3 \to B_3$으로 연장되고, 다음 가환
그림을 얻는다.
$$
\xymatrix{
B_2/\pi^2 B_2 \ar[r]_\psi & B_3/\pi^2B_3 \\
A_2/\pi^2 A_2 \ar[r]^\varphi \ar[u] & A_3/\pi^2A_3 \ar[u]
}
$$
여기에는 위에서 인용한 결과들이 단언한 모든 성질이 있다(특히
$\varphi$와 $\psi$는 동형이다).

\medskip\noindent
이 모든 데이터로 증명을 끝낼 수 있다. 먼저 $A_3$과 $B_3$은
$\pi$가 모든 극대아이디얼에 들어 있는 데데킨트 정역들의 유한 곱임에
주의하자. 실제로 $\mathfrak p \subset A_3$이 극대아이디얼이면
$A \to A_3$이 유한이므로 $\pi \in \mathfrak p$이다. 그러면
$\mathfrak p/\pi^2 A_3$은 $\varphi$를 통하여
$A_2 / \pi^2A_2$의 극대아이디얼에 대응한다. $A_2$가 데데킨트 정역들의
유한 곱이므로 이 아이디얼은 주아이디얼이다. 따라서
$\mathfrak p/\pi^2 A_3$이 주아이디얼이고 나카야마 보조정리에 의해
$\mathfrak p (A_3)_\mathfrak p$가 주아이디얼임을 알 수 있다.
$B_3$에도 같은 논증이 적용된다. 따라서 $A_3$은 $K_3$ 안에서의 $A$의
정수적 폐포이고 $B_3$은 $L_3$ 안에서의 $B$의 정수적 폐포이다.
$\mathfrak q \subset B_3$를 $\mathfrak p \subset A_3$ 위에 놓이는
극대아이디얼이라 하자. 증명을 끝내려면
$(A_3)_\mathfrak p \to (B_3)_\mathfrak q$가 $\mathfrak q$-진 위상에서
형식적으로 매끄러움을 보여야 한다. 다음의 판정법에 의해
보조정리 \ref{more-algebra-lemma-extension-dvrs-formally-smooth}, 다음을 보이면
충분하다.
$\mathfrak p (B_3)_\mathfrak q = \mathfrak q (B_3)_\mathfrak q$이고
체 확대 $\kappa(\mathfrak q)/\kappa(\mathfrak p)$가 분리이다.
두 단언 모두 환 사상 $A_3/\pi^2 A_3 \to B_3/\pi^2 B_3$을 살펴보아
검사할 수 있고, 이 사상은 환 사상
$A_2/\pi^2 A_2 \to B_2/\pi^2 B_2$와 동형이므로 참이다. 뒤 사상에
대해서는 $K_2$가 $A \subset B$의 해라는 가정으로부터 대응하는 단언이
성립한다. 세부사항 일부는 생략한다.
\end{proof}

\begin{lemma}
\label{more-algebra-lemma-separable-solution-separable-solution}
$A \subset B$를 이산 값매김환의 확대라 하자. 다음을 가정한다.
\begin{enumerate}
\item 분수체의 확대 $L/K$는 분리이다.
\item $B$는 나가타이다.
\item $A \subset B$의 해가 존재한다.
\end{enumerate}
그러면 $A \subset B$의 분리 해가 존재한다.
\end{lemma}

\begin{proof}
$K$의 표수가 영이면 보조정리는 자명하다. 따라서 $K$의 표수가
$p > 0$이라고 가정할 수 있고, 그렇게 가정한다.

\medskip\noindent
$K_2/K$를 $A \to B$의 해라 하자. $K_2/K$의 비분리 차수
$[K_2 : K]_i$에 관하여 귀납하겠다(Fields, 정의
\ref{fields-definition-insep-degree}). $[K_2 : K]_i = 1$이면 $K_2$는
$K$ 위에서 분리이므로 증명이 끝난다. 그렇지 않으면 $K_2/K_1$이 차수
$p$인 순수 비분리 확대가 되는 부분체의 탑 $K_2/K_1/K$가 존재한다
(Fields, 보조정리 \ref{fields-lemma-separable-first}와
\ref{fields-lemma-finite-purely-inseparable}). 보조정리
\ref{more-algebra-lemma-approximate-solution}에 의해 $A \subset B$의 해인 분리 확대
$K_3/K_1$이 존재한다. 그러면
$[K_3 : K]_i = [K_1 : K]_i = [K_2 : K]_i/p$이고(Fields, 보조정리
\ref{fields-lemma-multiplicativity-all-degrees}), 이는 더 작다. 귀납으로
결론을 얻는다.
\end{proof}

\begin{lemma}
\label{more-algebra-lemma-big-extension-is-ok}
$A \to B$를 분수체가 $K \subset L$인 이산 값매김환의 확대라 하자.
$B$가 $A$ 위에서 본질적 유한형이라고 가정하자. $K'/K$를 대수적 체
확대라 하되, $K'$ 안에서의 $A$의 정수적 폐포 $A'$이 뇌터라고 하자.
그러면 $L' = (L \otimes_K K')_{red}$ 안에서의 $B$의 정수적 폐포
$B'$도 뇌터이다. 또한 사상 $\Spec(B') \to \Spec(A')$은 전사이고,
대응하는 잉여체 확대들은 유한 생성 체 확대이다.
\end{lemma}

\begin{proof}
$A \to C$를 $B$가 소아이디얼 $\mathfrak p$에서의 $C$의 국소화가 되는
유한형 환 사상이라 하자. 그러면 $C' = C \otimes_A A'$은 유한형
$A'$-대수이고, 특히 뇌터이다. $A \to A'$이 정수적이므로
$C \to C'$도 정수적이다. 따라서
$B = C_\mathfrak p \subset C'_\mathfrak p$도 정수적이다.
그러므로 $C'_\mathfrak p$의 차원은 $1$이다(Algebra, 보조정리
\ref{algebra-lemma-integral-sub-dim-equal}). 물론
$C'_\mathfrak p$는 뇌터이다. $\mathfrak q_1, \ldots, \mathfrak q_n$을
$C'_\mathfrak p$의 극소 소아이디얼들이라 하자. $B'_i$를
$B = C_\mathfrak p$의 정수적 폐포, 또는 위에서 본 바와 동치로
% P02-E0686: frozen source uses the undefined prime p' in this fraction-field locus.
$C'_{\mathfrak p'}/\mathfrak q_i$의 분수체 안에서의 $C'_\mathfrak p$의
정수적 폐포라 하자. 크룰--아키즈키 정리에 의해($C'_\mathfrak p$의
극대아이디얼에서의 유한 개 국소화에 Algebra, 보조정리
\ref{algebra-lemma-krull-akizuki}를 적용한다) 각 $B'_i$는 뇌터이다.
또한 Algebra, 보조정리
\ref{algebra-lemma-finite-extension-residue-fields-dimension-1}에 의해
$C'_\mathfrak p \to B'_i$에서의 잉여체 확대들은 유한이다. 마지막으로
$B' = \prod B'_i$는 $L' = (L \otimes_K K')_{red}$ 안에서의 $B$의
정수적 폐포임에 주의하자.
\end{proof}

\begin{proposition}
\label{more-algebra-proposition-epp-essentially-finite-type}
\begin{reference}
특수한 경우는 \cite[Lemma 2.13]{alterations}를 보라.
\end{reference}
$A \to B$를 분수체가 $K \subset L$인 이산 값매김환의 확대라 하자.
$B$가 $A$ 위에서 본질적 유한형이면, 정의 \ref{more-algebra-definition-solution}에서
정의한 $A \to B$의 해인 유한 확대 $K_1/K$가 존재한다.
\end{proposition}

\begin{proof}
$A$의 잉여체가 완전체이면 약한 해는 해임에 주의하자. 보조정리
\ref{more-algebra-lemma-extension-dvrs-formally-smooth}를 보라. 따라서 $A$의 잉여
표수가 $0$이면 명제는 정리 \ref{more-algebra-theorem-epp}에서 즉시 따른다(실제로는
$A \to B$가 본질적 유한형이라는 가정이 필요하지 않다). $A$의 잉여
표수가 $p > 0$이면 역시 Epp 정리로부터 도출하겠다.

\medskip\noindent
$x_i \in A$, $i \in I$를 $A$의 잉여체 $\kappa$의 $p$-기저로 사상되는
원소들의 집합이라 하자. 다음과 같이 두자.
% P02-E0688: frozen source's indexed union presentation is retained for canon review.
$$
A' = \bigcup\nolimits_{n \geq 1} A[t_{i, n}]/(t_{i, n}^{p^n} - x_i)
$$
여기서 전이 사상은 $t_{i, n + 1}$을 $t_{i, n}^p$로 보낸다. $A'$은
$A$ 위의 이산 값매김환의 약하게 비분기인 유한 확대들의 여과 쌍대극한임에
주의하자. 따라서 $A'$은 이산 값매김환이고 $A \to A'$은 약하게
비분기이다. 구성에 의해 잉여체
$\kappa' = A'/\mathfrak m_A A'$은 $\kappa$의 완전체화이다.

\medskip\noindent
$K'$을 $A'$의 분수체라 하자. 확대 $K'/K$에 보조정리
\ref{more-algebra-lemma-big-extension-is-ok}를 적용할 수 있다. 따라서 $B'$은 데데킨트
정역들의 유한 곱이다. $\mathfrak m_1, \ldots, \mathfrak m_n$을 $B'$의
극대아이디얼들이라 하자. Epp 정리(정리 \ref{more-algebra-theorem-epp})를 사용하여
각 확대 $A' \subset B'_{\mathfrak m_i}$의 약한 해 $K'_i/K'$를 찾는다.
$A'$의 잉여체가 완전체이므로 이들은 실제로 해이다. $K'_1/K'$를 각
$K'_i$를 포함하는 유한 확대라 하자. 보조정리
\ref{more-algebra-lemma-solution-goes-up}에 의해 $K'_1/K'$는 여전히 모든
$A' \subset B'_{\mathfrak m_i}$의 해이다.

\medskip\noindent
$A'_1$을 $K'_1$ 안에서의 $A$의 정수적 폐포라 하자. $K' \subset K'_1$에
평주 \ref{more-algebra-remark-construction}의 논의를 적용하면 $A'_1$은 데데킨트
정역임에 주의하자. 따라서 $K'_1/K$에 보조정리
\ref{more-algebra-lemma-big-extension-is-ok}가 적용된다. 그러므로
$L'_1 = (L \otimes_K K'_1)_{red}$ 안에서의 $B$의 정수적 폐포 $B'_1$은
데데킨트 정역이고, $K'_1/K'$가 모든
$A' \subset B'_{\mathfrak m_i}$의 해이므로
% P02-E0687: frozen source's localization locus is retained for canon review.
각 극대아이디얼 $\mathfrak m \subset B'_1$에 대하여
$(A'_1)_{A'_1 \cap \mathfrak m} \to (B'_1)_{\mathfrak m}$은
$\mathfrak m$-진 위상에서 형식적으로 매끄럽다.

\medskip\noindent
구성에 의해 체 $K'_1$은 $K$의 유한 확대들의 여과 쌍대극한이다.
$K'_1 = \colim_{i \in I} K_i$라 쓰자. 각 $i$에 대하여 $A_i$와 $B_i$를
각각 $K_i$ 안에서의 $A$의 정수적 폐포와
$L_i = (L \otimes_K K_i)_{red}$ 안에서의 $B$의 정수적 폐포라 하자.
그러면 다음은 명백하다.
$$
A'_1 =  \colim A_i\quad\text{and}\quad B'_1 = \colim B_i
$$
환 사상 $A_i \to A'_1$과 $B_i \to B'_1$은 단사 정수적 환 사상이고
$A'_1$과 $B'_1$의 스펙트럼은 유한하므로, 충분히 큰 모든 $i$에 대하여
환 사상 $A_i \to A'_1$과 $B_i \to B'_1$은 스펙트럼 위에서 전단사이다.
일단 이것이 참이면, 충분히 큰 모든 $i$에 대하여 사상
$A_i \to A'_1$과 $B_i \to B'_1$은 약하게 비분기이다(균등화원이 상에
들어오고 나면). 분기지수의 곱셈성에 의해 그런 $i$에 대하여
$A_i \to B_i$는 $B_i$의 모든 극대아이디얼에서의 국소화 위에 약하게
비분기인 사상을 유도한다. $i$를 조금 더 키우면
$$
B_i \otimes_{A_i} A'_1 \longrightarrow B'_1
$$
이 잉여체 위에 전사인 사상을 유도함을 알 수 있다(보조정리
\ref{more-algebra-lemma-big-extension-is-ok}에 의해 $B'_1$의 잉여체들이 $A'_1$의
잉여체들 위에서 유한 생성이기 때문이다). $B'_1$의 주어진 극대아이디얼
아래에 놓이는 극대아이디얼들에서의 잉여체 그림은 다음과 같다.
$$
\xymatrix{
\kappa_{B_i} \ar[r] &
\kappa_{B_{i'}} \ar[r] &
 & \ldots &
\kappa_{B'_1} \\
\kappa_{A_i} \ar[r] \ar[u] &
\kappa_{A_{i'}} \ar[r] \ar[u] &
 & \ldots &
\kappa_{A'_1} \ar[u]
}
$$
따라서 $\kappa_{B_i}$는 $\kappa_{A_i}$의 유한 생성 확대이고,
$\kappa_{B_i}$와 $\kappa_{A'_1}$의 $\kappa_{B'_1}$ 안에서의 합성체는
$\kappa_{A'_1}$ 위에서 분리이다. 그러면 이는 이미 어떤 유한 단계에서
일어난다. 예를 들어 $\kappa_{B'_1}$이
$\kappa_{A'_1}(x_1, \ldots, x_n)$ 위에서 유한 분리라고 하자. 그러면
$x_1, \ldots, x_n$이 $\kappa_{B_i}$에 들어가고 모든 생성원이
$\kappa_{A_i}(x_1, \ldots, x_n)$ 위의 분리 다항방정식을 만족하도록
$i$를 키우기만 하면 된다. 이는 $B_i$의 모든 극대아이디얼
$\mathfrak m$에 대하여 $A_i \to (B_i)_\mathfrak m$가
$\mathfrak m$-진 위상에서 형식적으로 매끄럽다는 뜻이다. 증명이 끝났다.
\end{proof}

\begin{lemma}
\label{more-algebra-lemma-epp-essentially-finite-type-separable}
$A \to B$를 분수체가 $K \subset L$인 이산 값매김환의 확대라 하자.
다음을 가정한다.
\begin{enumerate}
\item $B$는 $A$ 위에서 본질적 유한형이다.
\item $A$ 또는 $B$가 나가타 환이다.
\item $L/K$는 분리이다.
\end{enumerate}
그러면 $A \to B$의 분리 해가 존재한다(정의
\ref{more-algebra-definition-solution}).
\end{lemma}

\begin{proof}
$A$가 나가타이면 $B$도 나가타임에 주의하자(Algebra, 보조정리
\ref{algebra-lemma-nagata-localize}와 명제
\ref{algebra-proposition-nagata-universally-japanese}). 따라서 명제
\ref{more-algebra-proposition-epp-essentially-finite-type}와 보조정리
\ref{more-algebra-lemma-separable-solution-separable-solution}를 합치면 보조정리가
따른다.
\end{proof}

\section{환의 Picard 군}
\label{more-algebra-section-picard}

\noindent
먼저 다음과 같이 가역 가군을 정의한다.

\begin{definition}
\label{more-algebra-definition-invertible}
$R$을 환이라 하자. $R$-가군 $M$이 {\it 가역}이라는 것은 함자
$$
\text{Mod}_R \longrightarrow \text{Mod}_R,\quad
N \longmapsto M \otimes_R N
$$
가 범주의 동치라는 뜻이다. 가역 $R$-가군이 $R$-가군으로서
$R$과 동형이면 {\it 자명}하다고 한다.
\end{definition}

\begin{lemma}
\label{more-algebra-lemma-invertible}
$R$을 환이라 하자. $M$을 $R$-가군이라 하자. 다음 조건들은 동치이다.
\begin{enumerate}
\item $M$은 계수 $1$의 유한 국소 자유 가군이다.
\item $M$은 가역이다.
\item $M \otimes_R N \cong R$이 되는 $R$-가군 $N$이 존재한다.
\end{enumerate}
더욱이 이 경우 (3)의 가군 $N$은 $\Hom_R(M, R)$과 동형이다.
\end{lemma}

\begin{proof}
(1)을 가정하자. 가군 $N = \Hom_R(M, R)$과 평가 사상
$M \otimes_R N = M \otimes_R \Hom_R(M, R) \to R$을 생각하자.
$M_f \cong R_f$가 되는 $f \in R$을 택하면, $f$에서 국소화한 뒤
평가 사상은 동형이 된다(세부 사항은 생략한다). 따라서 Algebra,
보조정리 \ref{algebra-lemma-cover}에 의해 평가 사상은 동형이다.
그러므로 (1) $\Rightarrow$ (3)이다.

\medskip\noindent
(3)을 가정하자. 그러면 함자 $K \mapsto K \otimes_R N$은 함자
$K \mapsto K \otimes_R M$의 준역함자이다. 따라서 (3) $\Rightarrow$ (2)이다.
거꾸로 (2)가 성립하면 $K \mapsto K \otimes_R M$은 본질적으로 전사이므로
(3)이 성립함을 알 수 있다.

\medskip\noindent
동치인 조건 (2)와 (3)이 성립한다고 가정하자. (3)의 동형을
$\psi : M \otimes_R N \to R$로 나타내자.
$\psi(\xi) = 1$이 되도록 원소
$\xi = \sum_{i = 1, \ldots, n} x_i \otimes y_i$를 택한다. 동형
$$
M \to M \otimes_R M \otimes_R N \to M
$$
을 생각하자. 여기서 첫째 화살표는 $x$를
$\sum x_i \otimes x \otimes y_i$로 보내고, 둘째 화살표는
$x \otimes x' \otimes y$를 $\psi(x' \otimes y)x$로 보낸다.
따라서 $x \mapsto \sum \psi(x \otimes y_i)x_i$는 $M$의 자기동형이다.
이 자기동형은
$$
M \to R^{\oplus n} \to M
$$
으로 분해된다. 여기서 첫째 화살표는
$x \mapsto (\psi(x \otimes y_1), \ldots, \psi(x \otimes y_n))$로 주어지고,
둘째 화살표는 $(a_1, \ldots, a_n) \mapsto \sum a_i x_i$로 주어진다.
이로부터 $M$이 유한 자유 $R$-가군의 직합 인자임을 얻는다. 이는
$M$이 유한 국소 자유라는 뜻이다(Algebra, 보조정리
\ref{algebra-lemma-finite-projective}). 대칭성에 의해 $N$도 마찬가지이고,
$M \otimes_R N \cong R$이므로 $M$과 $N$은 모두 계수 $1$이어야 한다.
\end{proof}

\noindent
이러한 가군들의 동형류 집합은 흔히 $R$의 {\it 유군} 또는
{\it Picard 군}이라고 한다. 군 구조는 가역 가군 $L$과 $L'$의 동형류에
$L \otimes_R L'$의 동형류를 대응시켜 정한다. 가역 가군 $L$의 역원은 가군
$$
L^{\otimes -1} = \Hom_R(L, R),
$$
이다. 실제로 보조정리 \ref{more-algebra-lemma-invertible}의 증명에서 보았듯이 평가 사상
$L \otimes_R L^{\otimes -1} \to R$은 동형이다. $R$의 Picard 군을
$\Pic(R)$로 나타내자.

\begin{lemma}
\label{more-algebra-lemma-UFD-Pic-trivial}
$R$을 UFD라 하자. 그러면 $\Pic(R)$은 자명하다.
\end{lemma}

\begin{proof}
$L$을 가역 $R$-가군이라 하자. 보조정리 \ref{more-algebra-lemma-invertible}에 의해
$L$은 유한 국소 자유 $R$-가군이다. 특히 $L$은 비틀림이 없고 $R$ 위에서
유한하다. 쌍대 가역 가군의 영이 아닌 원소
$\varphi \in \Hom_R(L, R)$를 택하자. 그러면
$I = \varphi(L) \subset R$은 가역 가군이기도 한 아이디얼이다. 영이 아닌
$f \in I$를 택하고
$$
f = u p_1^{e_1} \ldots p_r^{e_r}
$$
를 서로 다른 소원소 $p_i$들로의 분해라 하자. $L$이 유한 국소 자유이므로,
어떤 $g_i \in R_{a_i}$에 대하여 $I_{a_i} = (g_i)$가 되도록 하는
$a_i \in R$, $a_i \not \in (p_i)$가 존재한다. 그러면 $p_i$는 UFD
$R_{a_i}$에서도 여전히 소원소이며, $p_i$로 나누어지지 않는 어떤
$g'_i \in R_{a_i}$에 대하여 $g_i = p_i^{c_i} g'_i$로 쓸 수 있다.
$f \in I_{a_i}$이므로 $e_i \geq c_i$임을 알 수 있다. $I$가
$h = p_1^{c_1} \ldots p_r^{c_r}$로 생성된다고 주장한다. 이로써 증명이 끝난다.

\medskip\noindent
주장을 증명하려면 $I_a$가 주아이디얼이 되게 하는 모든 $a \in R$에 대해
$I_a$가 $h$로 생성됨을 보이면 충분하다(Algebra, 보조정리
\ref{algebra-lemma-cover}). $I_a = (g)$라 하자.
$J \subset \{1, \ldots, r\}$를 $p_i$가 $R_a$에서 비단위원
(따라서 소원소)이 되는 $i$들의 집합이라 하자. $f \in I_a = (g)$이므로
다음 소인수분해를 얻는다.
% P02-E0689: source line 36214 binds i in the product but uses p_j and b_j.
$g = v \prod_{i \in J} p_j^{b_j}$
여기서 $v$는 단위원이고 $b_j \leq e_j$이다. 각 $j \in J$에 대하여
$I_{aa_j} = g R_{aa_j} = g_j R_{aa_j}$이다. 다시 말해 $g$와 $g_j$는
$R_{aa_j}$에서 서로 동반인 원소로 간다. 분해의 유일성으로부터
$b_j = c_j$가 따르며 증명이 끝난다.
\end{proof}

\section{행렬식}
\label{more-algebra-section-determinants}

\noindent
$R$을 환이라 하자. $M$을 유한 사영 $R$-가군이라 하자.
곱분해 $R = R_0 \times \ldots \times R_t$가 존재하여, 이에 대응하는
$M$의 분해 $M = M_0 \times \ldots \times M_t$에서 $M_i$는 $R_i$ 위의
계수 $i$인 유한 국소 자유 가군이 된다. 이는 Algebra, 보조정리
\ref{algebra-lemma-finite-projective}(계수가 국소 상수임을 보기 위해)와
Algebra, 보조정리들 \ref{algebra-lemma-disjoint-decomposition} 및
\ref{algebra-lemma-disjoint-implies-product}($R$을 곱으로 분해하기 위해)에서
따른다. 이 상황에서 $R$-가군
$$
\det(M) = \wedge^0_{R_0}(M_0) \times \ldots \times \wedge^t_{R_t}(M_t)
$$
을 정의한다. 각 항이 계수 $1$인 유한 국소 자유 가군이므로, 이것도 계수
$1$인 유한 국소 자유 가군이다.
유한 사영 $R$-가군의 동형 $\varphi : M \to N$이 주어지면, 계수 $1$의
국소 자유 가군 사이의 표준 동형
$$
\det(\varphi) : \det(M) \longrightarrow \det(N)
$$
을 얻는다. 더 일반적으로 $R$의 모든 소아이디얼 $\mathfrak p$에 대하여
자유 가군 $M_\mathfrak p$와 $N_\mathfrak p$의 계수가 같다면, 임의의
$R$-가군 준동형 $\varphi : M \to N$은 $R$-가군 사상
$\det(\varphi) : \det(M) \to \det(N)$을 유도한다. 마지막으로 $M = N$이면
$\det(\varphi) : \det(M) \to \det(M)$은 가역 $R$-가군의 자기준동형이다.
가역 $R$-가군에 대해서는 $R = \Hom_R(L, L)$이므로, $\det(\varphi)$를
$R$의 원소로 볼 수 있고 그렇게 하기로 한다. 이로써 승법적 사상인
{\it 행렬식}
$$
\det : \Hom_R(M, M) \longrightarrow R
$$
을 얻는다.

\begin{remark}
\label{more-algebra-remark-determinant-as-socle}
$R$을 환이라 하자. $M$을 유한 사영 $R$-가군이라 하자. 그러면 $R$ 위에서
$M$으로 생성되는 등급 가환 $R$-대수인 외대수 $\wedge^*_R(M)$을 생각할
수 있다. $\det(M)$의 한 공식은 $\det(M) \subset \wedge^*_R(M)$이
$M \subset \wedge^*_R(M)$의 소멸자라는 것이다. 이 표현은 $R$의
곱분해를 언급하지 않으므로 때로 유용하다. 물론 이것이 원하는 성질들을
만족함을 증명하려면 $R$을 (위에서처럼) 곱으로 분해하거나, $R$의
소아이디얼들에서의 국소화를 살펴보아야 한다.
\end{remark}

\noindent
% P02-E0690: source line 36279 says “give” where the grammar requires “given”.
다음으로 유한 사영 가군들의 짧은 완전열이 주어졌을 때 행렬식에 무슨 일이
일어나는지 살펴본다.

\begin{lemma}
\label{more-algebra-lemma-det-ses}
$R$을 환이라 하자.
$$
0 \to M' \to M \to M'' \to 0
$$
을 유한 사영 $R$-가군들의 짧은 완전열이라 하자. 그러면 표준 동형
$$
\gamma : \det(M') \otimes \det(M'') \longrightarrow \det(M)
$$
이 존재한다.
\end{lemma}

\begin{proof}[첫 번째 증명]
첫 번째 증명. $R$을 환들의 곱 $R_{ij}$로 분해하여
$M' = \prod M'_{ij}$이고 $M'' = \prod M''_{ij}$이며, $M'_{ij}$의 계수는
$i$이고 $M''_{ij}$의 계수는 $j$가 되게 하자. 물론 그러면
$M = \prod M_{ij}$이고 $M_{ij}$의 계수는 $i + j$이다. 따라서
$M'$과 $M''$의 계수가 각각 상수 $i$와 $j$인 경우로 환원된다.
이 경우 표준 사상
$$
\wedge^i(M') \otimes \wedge^j(M'') \longrightarrow \wedge^{i + j}(M)
$$
을 구성해야 한다. $M'$에서 $m'_1, \ldots, m'_i$를, $M''$에서
$m''_1, \ldots, m''_j$를 택한다. $m'_1, \ldots, m'_i$의 상들을
$m_1, \ldots, m_i \in M$로 나타내고, $M''$에서
$m''_1, \ldots, m''_j$로 가는 원소들을
$m_{i + 1}, \ldots , m_{i + j} \in M$로 나타내자. 규칙을
$$
m'_1 \wedge \ldots \wedge m'_i \otimes
m''_1 \wedge \ldots \wedge m''_j
\longmapsto
m_1 \wedge \ldots \wedge m_{i + j}
$$
로 정한다. 이것이 잘 정의되고 동형이라는 상세한 증명은 생략한다.
\end{proof}

\begin{proof}[두 번째 증명]
주의 \ref{more-algebra-remark-determinant-as-socle}에서 주어진 $\det(M)$,
$\det(M')$, $\det(M'')$의 표현을 사용하자. $R$-대수 사상
$\wedge^*_R(M') \to \wedge^*_R(M)$과
$\wedge^*_R(M) \to \wedge^*_R(M'')$을 생각하자. 첫째 사상은 단사이고
둘째 사상은 전사이다. 원소
$x' \in \det(M') \subset \wedge^*_R(M')$와 원소
$x'' \in \det(M'') \subset \wedge^*_R(M'')$를 택하자.
$x''$로 가는 원소 $y'' \in \wedge^*(M)$를 택하고
$$
\gamma(x' \otimes x'') = x' \wedge y'' \in \det(M) \subset \wedge^*_R(M)
$$
로 놓는다.
% P02-E0691: source line 36332 omits “is” in “this is well defined”.
소아이디얼들에서의 국소화를 살펴보면 이것이 잘 정의되고 동형임을 독자가
쉽게 확인할 수 있다. 더욱이 이 구성은 첫 번째 증명의 구성과 같은 사상을
준다.
\end{proof}

\begin{lemma}
\label{more-algebra-lemma-det-ses-functorial}
$R$을 환이라 하자.
$$
\xymatrix{
0 \ar[r] &
M' \ar[r] \ar[d]^u &
M \ar[r] \ar[d]^v &
M'' \ar[r] \ar[d]^w &
0 \\
0 \ar[r] &
K' \ar[r] &
K \ar[r] &
K'' \ar[r] &
0
}
$$
을 수직 화살표들이 동형인 유한 사영 $R$-가군들의 가환 그림이라 하자.
그러면 동형들의 가환 그림
$$
\xymatrix{
\det(M') \otimes \det(M'') \ar[r]_-\gamma \ar[d]_{\det(u) \otimes \det(w)} &
\det(M) \ar[d]^{\det(v)} \\
\det(K') \otimes \det(K'') \ar[r]^-\gamma & \det(K)
}
$$
을 얻는다. 여기서 수평 화살표들은 보조정리 \ref{more-algebra-lemma-det-ses}에서
구성한 것들이다.
\end{lemma}

\begin{proof}
생략한다. 힌트: 보조정리 \ref{more-algebra-lemma-det-ses}에서 사상 $\gamma$를 구성한
두 번째 방법을 사용하라.
\end{proof}

\begin{lemma}
\label{more-algebra-lemma-det-filtration}
$R$을 환이라 하자.
$$
K \subset L \subset M
$$
을 $K$, $L/K$, $M/L$이 유한 사영 $R$-가군인 $R$-가군들이라 하자.
그러면 그림
$$
\xymatrix{
\det(K) \otimes \det(L/K) \otimes \det(M/L) \ar[r] \ar[d] &
\det(L) \otimes \det(M/L) \ar[d] \\
\det(K) \otimes \det(M/K) \ar[r] &
\det(M)
}
$$
은 가환한다. 여기서 사상들은 보조정리 \ref{more-algebra-lemma-det-ses}의 사상들이다.
\end{lemma}

\begin{proof}
생략한다. 힌트: $R$의 한 소아이디얼에서 국소화한 뒤
$K \subset L \subset M$이
$R^{\oplus a} \subset R^{\oplus a + b}  \subset R^{\oplus a + b + c}$와
동형이라고 가정할 수 있으며, 이 경우 결과는 자명한 계산이다.
\end{proof}

\begin{lemma}
\label{more-algebra-lemma-det-direct-sum}
$R$을 환이라 하자. $M'$과 $M''$을 두 유한 사영 $R$-가군이라 하자.
그러면 그림
$$
\xymatrix{
\det(M') \otimes \det(M'') \ar[r]
\ar[d]_{\epsilon \cdot (\text{switch tensors})} &
\det(M' \oplus M'') \ar[d]^{\det(\text{switch summands})} \\
\det(M'') \otimes \det(M') \ar[r] &
\det(M'' \oplus M')
}
$$
은 가환한다. 여기서
$\epsilon = \det( -\text{id}_{M' \otimes M''}) \in R^*$이고 수평 화살표들은
보조정리 \ref{more-algebra-lemma-det-ses}의 사상들이다.
\end{lemma}

\begin{proof}
생략한다.
\end{proof}

\begin{lemma}
\label{more-algebra-lemma-det-switch}
$R$을 환이라 하자. $M$, $N$을 유한 사영 $R$-가군이라 하자.
$a : M \to N$과 $b : N \to M$을 $R$-선형 사상이라 하자. 그러면
$$
\det(\text{id} + a \circ b) = \det(\text{id} + b \circ a)
$$
가 $R$의 원소로서 성립한다.
\end{lemma}

\begin{proof}
$R$을 한 소아이디얼에서의 국소화로 바꾼 뒤 명제를 증명하면 충분하다.
따라서 $R$이 국소환이고 $M$과 $N$이 유한 자유라고 가정해도 된다.
이 경우 $A$의 크기가 $n \times m$이고 $B$의 크기가 $m \times n$인
행렬들의 보통 행렬식에 대한 등식
$$
\det(I_n + AB) = \det(I_m + BA)
$$
을 증명해야 한다. 이는
$R = \mathbf{Z}[a_{ij}, b_{ji}; 1 \leq i \leq n, 1 \leq j \leq m]$이고,
% P02-E0693: source line 36440 changes b_{ji} to b_{ij} in the prose.
$a_{ij}$와 $b_{ij}$가 변수이자 행렬 $A$와 $B$의 성분인 경우로 환원된다.
분수체를 취하면 표수 영인 체의 경우로 환원된다.
% P02-E0692: the source trace reduction compares different matrix sizes and is invalid for n != m.
표수 영에서는 크기가 $\leq N$인 행렬의 행렬식을 그 행렬의 거듭제곱들의
대각합으로 표현하는 보편 다항식이 존재한다. 따라서 모든 $k \geq 1$에 대해
$$
\text{Trace}((I_n + AB)^k) = \text{Trace}((I_m + BA)^k)
$$
임을 증명하면 충분하다. 전개하면 모든 $k \geq 0$에 대해
$\text{Trace}((AB)^k) = \text{Trace}((BA)^k)$임을 증명하면 충분함을
알 수 있다. $k = 1$일 때 이는 잘 알려진 사실
$\text{Trace}(AB) = \text{Trace}(BA)$이다. $k > 1$일 때는
$(AB)^k = A(BA)^{k - 1}B$와 $(BA)^k = (BA)^{k - 1} A B$로 써서
이 사실로부터 따른다.
\end{proof}

\noindent
Algebra, 절 \ref{algebra-section-K-groups}에서 유한 사영 $R$-가군들의
동형류 위의 자유군을 관계식
$[M'] + [M''] = [M' \oplus M'']$으로 나누어 군 $K_0(R)$을 정의했음을
상기하자.

\begin{lemma}
\label{more-algebra-lemma-det}
$R$을 환이라 하자. 사상
$$
\det : K_0(R) \longrightarrow \Pic(R)
$$
이 존재하며, $M$이 계수 $n$인 유한 국소 자유 가군이면 이 사상은
$[M]$을 가역 가군 $\wedge^n(M)$의 동형류로 보낸다.
\end{lemma}

\begin{proof}
이는 위의 구성들에서 즉시 따르며, 특히 관계식들이 $0$으로 가는 것은
보조정리 \ref{more-algebra-lemma-det-ses}에서 따른다.
\end{proof}

\section{퍼펙트 복합체와 K-군}
\label{more-algebra-section-perfect-K-group}

\noindent
환 $R$의 퍼펙트 복합체들의 유도 범주의 영차 K-군이 Algebra, 절
\ref{algebra-section-K-groups}에서 정의한 $K_0(R)$과 같음을 간단히
보이자.

\begin{lemma}
\label{more-algebra-lemma-perfect-to-K-group}
$R$을 환이라 하자. 다음 성질들을 갖는 사상
$$
c : \text{perfect complexes over }R \longrightarrow K_0(R)
$$
이 존재한다.
\begin{enumerate}
\item 퍼펙트 복합체 $K$에 대하여 $c(K[n]) = (-1)^nc(K)$이다.
\item 퍼펙트 복합체들의 완전 삼각형 $K \to L \to M \to K[1]$이
주어지면 $c(L) = c(K) + c(M)$이다.
\item $K$가 유한 사영 가군들로 이루어진 유한 복합체 $M^\bullet$로
표현되면
% P02-E0694: source line 36503 writes M_i although the complex uses M^i.
$c(K) = \sum (-1)^i[M_i]$이다.
\end{enumerate}
\end{lemma}

\begin{proof}
$K$를 $D(R)$의 퍼펙트 대상이라 하자. 정의에 의해 $K$를 유한 사영
$R$-가군들의 유한 복합체 $M^\bullet$로 표현할 수 있다. 다음과 같이
놓아 $c$를 정의한다.
$$
c(K) = \sum (-1)^n[M^n]
$$
이는 $K_0(R)$의 원소이다. 물론 이것이 잘 정의됨을 보여야 하지만, 일단
잘 정의되면 (1)과 (3)은 즉시 따른다. 당분간 사상 $c$를 유한 사영
$R$-가군들의 복합체 위에서 정의된 것으로 보자.

\medskip\noindent
유한 사영 $R$-가군들의 유한 복합체 사이의 전사 사상
$L^\bullet \to M^\bullet$이 주어졌다고 하자. 그 핵을 $K^\bullet$이라
하자. 그러면 $R$-가군들의 짧은 완전열
$$
0 \to K^n \to L^n \to M^n \to 0
$$
을 얻으며, 이는 $M^n$이 사영이므로 분할된다. 따라서 $K^\bullet$도
유한 사영 $R$-가군들의 유한 복합체이고,
$c(L^\bullet) = c(K^\bullet) + c(M^\bullet)$가 $K_0(R)$에서 성립한다.

\medskip\noindent
유한 사영 $R$-가군들로 이루어진 비순환 유한 복합체 $M^\bullet$이
주어졌다고 하자. $n \not \in [a, b]$에 대하여 $M^n = 0$이라고 하자.
그러면 $M^\bullet$을 다음 짧은 완전열들로 나눌 수 있다.
% P02-E0695: source line 36538 skips from N^{a+1} to N^{a+3}.
$$
\begin{matrix}
0 \to M^a \to M^{a + 1} \to N^{a + 1} \to 0, \\
0 \to N^{a + 1} \to M^{a + 2} \to N^{a + 3} \to 0, \\
\ldots \\
0 \to N^{b - 3} \to M^{b - 2} \to N^{b - 2} \to 0, \\
0 \to N^{b - 2} \to M^{b - 1} \to M^b \to 0
\end{matrix}
$$
내림 귀납법으로 $N^{b - 2}, \ldots, N^{a + 1}$이 유한 사영
$R$-가군이고 이 열들이 분할 완전이며 다음이 성립함을 알 수 있다.
% P02-E0696: source line 36547 omits the exponent n from the first (-1).
$$
c(M^\bullet) = \sum (-1)[M^n] = \sum (-1)^n([N^{n - 1}] + [N^n]) = 0
$$
따라서 우리의 구성은 비순환 복합체에서 영이 된다.

\medskip\noindent
앞 두 문단의 결과로부터 $c$가 잘 정의되고 (2)를 만족함이 형식적으로
따른다. 구체적으로 유한 사영 $R$-가군들의 유한 복합체 $M^\bullet$과
$L^\bullet$이 $D(R)$의 같은 대상을 표현한다고 하자. 그러면 이 동형을
복합체 사상 $f : M^\bullet \to L^\bullet$로 표현할 수 있다. 유도 범주,
보조정리 \ref{derived-lemma-morphisms-from-projective-complex}를 보라.
복합체들의 짧은 완전열
% P02-E0697: source lines 36561-36570 use undefined K^\bullet instead of M^\bullet.
$$
0 \to L^\bullet \to C(f)^\bullet \to K^\bullet[1] \to 0
$$
을 얻는다. 유도 범주, 정의 \ref{derived-definition-cone}를 보라.
$f$가 준동형사상이므로 원뿔 $C(f)^\bullet$은 비순환이다(예를 들어 유도
범주, 절 \ref{derived-section-canonical-delta-functor}의 논의에서 따른다).
따라서 원하는 대로
$$
0 = c(C(f)^\bullet) = c(L^\bullet) + c(K^\bullet[1]) =
c(L^\bullet) - c(K^\bullet)
$$
이다. 비슷한 (2)의 증명은 생략한다.
\end{proof}

\noindent
다음 보조정리는 $K_0(R)$이 $K_0(D_{perf}(R))$과 같음을 보인다.

\begin{lemma}
\label{more-algebra-lemma-perfect-to-K-group-universal}
$R$을 환이라 하자. $D_{perf}(R)$을 퍼펙트 대상들의 유도 범주라 하자.
보조정리 \ref{more-algebra-lemma-perfect-ring-classical-generator}를 보라.
보조정리 \ref{more-algebra-lemma-perfect-to-K-group}의 사상 $c$는 동형
$K_0(D_{perf}(R)) = K_0(R)$을 준다.
\end{lemma}

\begin{proof}
$K_0(D_{perf}(R))$의 정의(유도 범주, 정의
\ref{derived-definition-K-zero})로부터 $c$가 준동형
$K_0(D_{perf}(R)) \to K_0(R)$을 유도함이 따른다.

\medskip\noindent
$R$ 위의 유한 사영 가군 $M$이 주어지면, $M$이 차수 $0$에 놓이고 다른
차수에서는 영인 $R$ 위의 퍼펙트 복합체를 $M[0]$으로 나타내자.
유한 사영 가군들의 짧은 완전열 $0 \to M \to M' \to M'' \to 0$이
주어지면 완전 삼각형 $M[0] \to M'[0] \to M''[0] \to M[1]$을 얻는다.
유도 범주, 절 \ref{derived-section-canonical-delta-functor}를 보라.
이로써 $[M]$을 $[M[0]]$으로 보내는 사상
$K_0(R) \to K_0(D_{perf}(R))$을 얻는다. 이 끔찍한 표기에 양해를 구한다.

\medskip\noindent
$K_0(R) \to K_0(D_{perf}(R)) \to K_0(R)$이 항등임은 명백하다.
한편 $M^\bullet$이 유한 사영 $R$-가군들의 유계 복합체라면,
% P02-E0698: source line 36605 duplicates “the the”.
“stupid 절단”들의 완전 삼각형
(Homology, 절 \ref{homology-section-truncations}를 보라)
$$
\sigma_{\geq n}M^\bullet \to \sigma_{\geq n - 1}M^\bullet \to
M^{n - 1}[-n + 1] \to (\sigma_{\geq n}M^\bullet)[1]
$$
이 존재하므로 귀납법에 의해
$$
[M^\bullet] = \sum (-1)^i[M^i[0]]
$$
을 $K_0(D_{perf}(R))$에서 얻는다(표기에 대해 다시 양해를 구한다).
따라서 사상 $K_0(R) \to K_0(D_{perf}(R))$은 전사이고, 이로써 증명이
끝난다.
\end{proof}

\section{유한 길이 가군의 자기준동형의 행렬식}
\label{more-algebra-section-determinants-finite-length}

\noindent
$(R, \mathfrak m, \kappa)$를 국소환이라 하자. 유한 길이 $R$-가군과
자기준동형 $\varphi : M \to M$으로 이루어진 쌍 $(M, \varphi)$들의
범주를 생각하자. 이 범주는 아벨 범주이고 모든 대상은 아르틴 대상인 동시에
뇌터 대상이다. 정의는 Homology, 절
\ref{homology-section-jordan-holder}를 보라.

\medskip\noindent
$(M, \varphi)$가 이 범주의 단순 대상이면, 그렇지 않을 경우
% P02-E0699: source line 36642 misspells “subobject” as “suboject”.
$(\mathfrak m M, \varphi|_{\mathfrak m M})$가 자명하지 않은 부분대상이
되므로 $M$은 $\mathfrak m$에 의해 소멸된다. 또한
$\dim_\kappa(M) = \text{length}_R(M)$은 유한하다. 따라서 선형대수를
사용하여 행렬식과 대각합
$$
\det\nolimits_\kappa(\varphi),\quad \text{Trace}_\kappa(\varphi)
$$
을 $\kappa$의 원소로 정의할 수 있다. 이 경우 $\varphi$의 특성다항식도
마찬가지로 정의한다.

\medskip\noindent
Homology, 보조정리 \ref{homology-lemma-finite-length}에 의해 우리 범주의
임의의 대상 $(M, \varphi)$에는 유한 여과
$$
0 \subset M_1 \subset \ldots \subset M_n = M
$$
가 존재한다. 여기서 부분가군들은 $\varphi$ 아래 안정적이고,
$(M_i/M_{i - 1}, \varphi_i)$는 이 범주의 단순 대상이며,
$\varphi_i : M_i/M_{i - 1} \to M_i/M_{i - 1}$는 유도된 사상이다.
$(M, \varphi)$의 $\kappa$ 위 {\it 행렬식}을
$$
\det\nolimits_\kappa(\varphi) = \prod \det\nolimits_\kappa(\varphi_i)
$$
로 정의한다. 여기서 $\det_\kappa(\varphi_i)$는 앞 문단에서 정의한
것이다. $(M, \varphi)$의 $\kappa$ 위 {\it 대각합}을
$$
\text{Trace}_\kappa(\varphi) = \sum \text{Trace}_\kappa(\varphi_i)
$$
로 정의한다. 여기서 $\text{Trace}_\kappa(\varphi_i)$는 앞 문단에서
정의한 것이다. $\kappa$ 위에서의 $\varphi$의 특성다항식도 마찬가지로
앞 문단에서 정의한 $\varphi_i$들의 특성다항식들의 곱으로 정의할 수 있다.
Jordan--H\"older(Homology, 보조정리
\ref{homology-lemma-jordan-holder})에 의해 이는 잘 정의된다.

\begin{lemma}
\label{more-algebra-lemma-ses}
$(R, \mathfrak m, \kappa)$를 국소환이라 하자. 위에서 논의한 범주의 짧은
완전열
$0 \to (M, \varphi) \to (M', \varphi') \to (M'', \varphi'') \to 0$
이 주어졌다고 하자. 그러면
$$
\det\nolimits_\kappa(\varphi') =
\det\nolimits_\kappa(\varphi)\det\nolimits_\kappa(\varphi''),\quad
\text{Trace}_\kappa(\varphi') = \text{Trace}_\kappa(\varphi) + 
\text{Trace}_\kappa(\varphi'')
$$
이다. 또한 $\kappa$ 위에서의 $\varphi'$의 특성다항식은 $\varphi$와
$\varphi''$의 특성다항식들의 곱이다.
\end{lemma}

\begin{proof}
연습문제로 남긴다.
\end{proof}

\begin{lemma}
\label{more-algebra-lemma-multiplication}
$(R, \mathfrak m, \kappa) \to (R', \mathfrak m', \kappa')$를 국소환들의
국소 준동형이라 하자. $\kappa'/\kappa$가 유한 확대라고 가정하자.
$u \in R'$라 하자. 그러면 임의의 유한 길이 $R'$-가군 $M'$에 대하여
$$
\det\nolimits_\kappa(u : M' \to M') =
\text{Norm}_{\kappa'/\kappa}(u \bmod \mathfrak m')^m
$$
이다. 여기서 $m = \text{length}_{R'}(M')$이다.
\end{lemma}

\begin{proof}
$\text{length}_R(M') = \text{length}_{R'}(M') [\kappa' : \kappa]$이므로
명제의 의미가 통함에 유의하자. $M' = \kappa'$이면, 등식은
$u$를 곱하는 선형 작용소의 행렬식으로 주어진 노름의 정의에 의해 성립한다.
일반적인 경우에는 적절한 여과를 택하고 보조정리 \ref{more-algebra-lemma-ses}의
승법성을 사용하여 이 경우로 환원한다. 몇 가지 세부 사항은 생략한다.
\end{proof}

\begin{lemma}
\label{more-algebra-lemma-flat-base-change-det}
$(R, \mathfrak m, \kappa) \to (R', \mathfrak m', \kappa')$를 국소환들의
평탄 국소 준동형이라 하고
$m = \text{length}_{R'}(R'/\mathfrak mR') < \infty$라 하자. 위와 같은
임의의 $(M, \varphi)$에 대하여 원소 $\det_\kappa(\varphi)^m$은
$\kappa'$에서
$\det_{\kappa'}(\varphi \otimes 1 : M \otimes_R R' \to M \otimes_R R')$로
간다.
\end{lemma}

\begin{proof}
$R \to R'$의 평탄성은 보조정리 \ref{more-algebra-lemma-ses}에서와 같은 짧은 완전열이
기저변환 뒤에도 $R'$ 위의 짧은 완전열이 됨을 보장한다. 따라서 보조정리
\ref{more-algebra-lemma-ses}의 승법성에 의해 $(M, \varphi)$가 우리 범주의 단순
대상이라고 가정해도 된다(이 절의 도입부를 보라). 단순한 경우 $M$은
$\mathfrak m$에 의해 소멸된다. 이어지는 몫들이 $R'$-가군으로
$\kappa'$와 동형인 여과
$$
0 \subset I_1 \subset I_2 \subset \ldots \subset I_{m - 1} \subset
R'/\mathfrak mR'
$$
를 택하자. 그러면 이어지는 몫들이 $M \otimes_\kappa \kappa'$와
동형인 여과
$$
0 \subset
M \otimes_\kappa I_1 \subset
M \otimes_\kappa I_2 \subset
\ldots \subset
M \otimes_\kappa I_{m - 1} \subset
M \otimes_\kappa R'/\mathfrak mR' = M \otimes_R R'
$$
를 얻는다. 또한 이 부분가군들은 $\varphi \otimes 1$ 아래 불변이다.
보조정리 \ref{more-algebra-lemma-ses}에 의해
$$
\det\nolimits_{\kappa'}(\varphi \otimes 1 : M \otimes_R R' \to M \otimes_R R')
=
\det\nolimits_{\kappa'}(\varphi \otimes 1 :
M \otimes_\kappa \kappa' \to M \otimes_\kappa \kappa')^m =
\det\nolimits_\kappa(\varphi)^m
$$
을 얻는다. 마지막 등식은 선형 사상들의 행렬식과 체 확대의 양립성에 의해
성립한다. 이것이 보조정리를 증명한다.
\end{proof}

\section{정칙 국소환은 UFD이다}
\label{more-algebra-section-regular-local-UFD}

\noindent
절 제목에서 언급한 결과를 증명한다.

\begin{lemma}
\label{more-algebra-lemma-regular-local-Pic-zero}
$R$을 정칙 국소환이라 하자. $f \in R$라 하자. 그러면
$\Pic(R_f) = 0$이다.
\end{lemma}

\begin{proof}
$L$을 가역 $R_f$-가군이라 하자. 특히 $L$은 유한 $R_f$-가군이다.
$M_f \cong L$이 되는 유한 $R$-가군 $M$이 존재한다. Algebra, 보조정리
\ref{algebra-lemma-construct-fp-module}를 보라. Algebra, 명제
\ref{algebra-proposition-regular-finite-gl-dim}에 의해 $M$은 $R$ 위의
유한 자유 분해 $F_\bullet$을 갖는다. 따라서 $L$은 {\it 자유}
$R_f$-가군들의 유한 복합체와 준동형적이다. 그러므로 보조정리
\ref{more-algebra-lemma-perfect-to-K-group}에 의해 어떤 $n \in \mathbf{Z}$에 대하여
% P02-E0700: source line 36795 writes [L_f] although L is already over R_f.
$[L_f] = n[R_f]$가 $K_0(R_f)$에서 성립한다. 보조정리
\ref{more-algebra-lemma-det}의 사상을 적용하면 $L$이 자명함을 알 수 있다.
\end{proof}

\begin{lemma}
\label{more-algebra-lemma-regular-local-UFD}
정칙 국소환은 UFD이다.
\end{lemma}

\begin{proof}
정칙 국소환이 정역임을 상기하자. Algebra, 보조정리
\ref{algebra-lemma-regular-domain}를 보라. 정칙 국소환 $R$의 차원에 관한
귀납법으로 유일인수분해 성질을 증명하자. $\dim(R) = 0$이면 $R$은
체이고, 특히 UFD이다. $\dim(R) > 0$이라 가정하자.
$x \in \mathfrak m$, $x \not \in \mathfrak m^2$라 하자. 그러면 Algebra,
보조정리 \ref{algebra-lemma-regular-ring-CM}에 의해 $R/(x)$는 정칙이고,
따라서 Algebra, 보조정리 \ref{algebra-lemma-regular-domain}에 의해
정역이다. 그러므로 $x$는 소원소이다.
높이 $1$인 소아이디얼 $\mathfrak p \subset R$을 택하자. Algebra,
보조정리 \ref{algebra-lemma-characterize-UFD-height-1}에 의해
$\mathfrak p$가 주아이디얼임을 보여야 한다.
$x \in \mathfrak p$이면 $\mathfrak p = (x)$이어서 끝나므로,
$x \not \in \mathfrak p$라고 가정해도 된다. 모든 비극대 소아이디얼
$\mathfrak q \subset R$에 대하여 국소환 $R_\mathfrak q$는 정칙
국소환이다. Algebra, 보조정리
\ref{algebra-lemma-localization-of-regular-local-is-regular}를 보라.
귀납가정으로부터 $\mathfrak pR_\mathfrak q$가 주아이디얼임을 얻는다.
특히 $R_x$-가군 $\mathfrak p_x = \mathfrak pR_x \subset R_x$는 유한
표시 $R_x$-가군이고, 모든 소아이디얼에서의 국소화가 계수 $1$인 자유
가군이다.
Algebra, 보조정리 \ref{algebra-lemma-finite-projective}에 의해
$\mathfrak p_x$는 가역 $R_x$-가군이다. 보조정리
\ref{more-algebra-lemma-regular-local-Pic-zero}에 의해 어떤 $y \in R_x$에 대하여
$\mathfrak p_x = (y)$이다. 어떤 $f \in \mathfrak p$와
$e \in \mathbf{Z}$에 대하여 $y = x^e f$로 쓸 수 있다.
$f = a_1 \ldots a_r$을 $R$의 기약원소들로 인수분해하자(Algebra,
보조정리 \ref{algebra-lemma-factorization-exists}). $\mathfrak p$가
소아이디얼이므로 어떤 $i$에 대하여 $a_i \in \mathfrak p$이다.
$\mathfrak p_x = (y)$가 소아이디얼이고 $R_x$에서 $a_i | y$이므로,
$\mathfrak p_x$는 $R_x$에서 $a_i$로 생성된다. 즉 $a_i$의 $R_x$에서의
상은 소원소이다. $x$가 소원소이므로 Algebra, 보조정리
\ref{algebra-lemma-invert-prime-elements}에 의해 $a_i$는 $R$에서
소원소이다. $(a_i) \subset \mathfrak p$이고 $\mathfrak p$의 높이가
$1$이므로 원하는 대로 $(a_i) = \mathfrak p$라고 결론내린다.
\end{proof}

\begin{lemma}
\label{more-algebra-lemma-picard-group-generic-fibre-regular}
$R$을 분수체가 $K$이고 잉여체가 $\kappa$인 값매김환이라 하자.
$R \to A$를 다음 조건들을 만족하는 환 준동형이라 하자.
\begin{enumerate}
\item $A$는 국소환이고 $R \to A$는 국소 준동형이다.
\item $A$는 $R$ 위에서 평탄하고 본질적 유한형이다.
% P02-E0701: source line 36853 lacks the predicate “is”.
\item $A \otimes_R \kappa$는 정칙이다.
\end{enumerate}
그러면 $\Pic(A \otimes_R K) = 0$이다.
\end{lemma}

\begin{proof}
$L$을 가역 $A \otimes_R K$-가군이라 하자. 특히 $L$은 유한 가군이다.
$M \otimes_R K \cong L$이 되는 유한 $A$-가군 $M$이 존재한다.
Algebra, 보조정리 \ref{algebra-lemma-construct-fp-module}를 보라.
$M$이 $R$-가군으로서 비틀림이 없다고 가정해도 된다. 따라서 $M$은
$R$-가군으로서 평탄하다(보조정리
\ref{more-algebra-lemma-valuation-ring-torsion-free-flat}). 보조정리
\ref{more-algebra-lemma-flat-finite-type-valuation-ring-finite-presentation}로부터
$M$이 $A$-가군으로서 유한 표시이고 $A$가 $R$-대수로서 본질적 유한
표시임을 얻는다. 보조정리 \ref{more-algebra-lemma-structure-relatively-perfect}에
의해 $M$은 $R$에 대해 상대적으로 퍼펙트이고, 특히 $M$은
$A$-가군으로서 유사연접이다. 보조정리
\ref{more-algebra-lemma-perfect-over-regular-local-ring}에
의해 $M$은 퍼펙트이다. 따라서 $M$은 $A$ 위의 유한 자유 분해
$F_\bullet$을 갖는다. 그러므로 $L$은 {\it 자유}
$A \otimes_R K$-가군들의 유한 복합체와 준동형적이다. 따라서 보조정리
\ref{more-algebra-lemma-perfect-to-K-group}에 의해 어떤 $n \in \mathbf{Z}$에 대하여
$[L] = n[A \otimes_R K]$가 $K_0(A \otimes_R K)$에서 성립한다.
보조정리 \ref{more-algebra-lemma-det}의 사상을 적용하면 $L$이 자명함을 알 수 있다.
\end{proof}

\section{복합체의 행렬식}
\label{more-algebra-section-determinants-complexes}

\noindent
절 \ref{more-algebra-section-perfect-K-group}에서 환 $R$ 위의 퍼펙트 복합체 $K$에
가역 $R$-가군의 동형류, 즉 $\Pic(R)$의 원소가 어떻게 대응하는지 보았다.
실제로 절 \ref{more-algebra-section-determinants}와 마찬가지로 함자
$$
\det :
\left\{
\begin{matrix}
\text{category of perfect complexes} \\
\text{morphisms are isomorphisms}
\end{matrix}
\right\}
\longrightarrow
\left\{
\begin{matrix}
\text{category of invertible modules} \\
\text{morphisms are isomorphisms}
\end{matrix}
\right\}
$$
가 존재함을 알 수 있다. 더욱이 $R$의 여과 유도 범주 $DF(R)$의 대상
$(L, F)$가 주어지고 그 여과가 유한하며 그 등급 조각들이 퍼펙트
복합체이면, 표준 동형 $\det(\text{gr}L) \to \det(L)$이 존재한다.
원래의 설명은 \cite{determinant}를 보라.
% P02-E0702: source line 36918 retains the editorial placeholder “insert future reference”.
이 내용은 나중에 추가할 것이다(향후 참고문헌 삽입).

\medskip\noindent
지금은 $D(R)$의 퍼펙트 대상 $L$이 $[-1, 0]$에 Tor 진폭을 갖는 경우에
임시방편의 구성을 제시한다. 이러한 대상은 복합체
$$
L^\bullet = \ldots \to 0 \to L^{-1} \to L^0 \to 0 \to \ldots
$$
로 표현할 수 있다. 여기서 $L^{-1}$과 $L^0$은 유한 사영 $R$-가군이다.
보조정리 \ref{more-algebra-lemma-perfect}를 보라. 이 경우
$$
\det(L^\bullet) = \det(L^0) \otimes_R \det(L^{-1})^{\otimes -1} =
\Hom_R(\det(L^{-1}), \det(L^0))
$$
로 놓는다. 이 꼴의 복합체에서 $R$의 모든 소아이디얼에 대해
$L^{-1}_\mathfrak p$와 $L^0_\mathfrak p$의 계수가 같으면, 이 복합체가
{\it 계수 $0$}을 갖는다고 하자. $L^\bullet$의 계수가 $0$이면 절
\ref{more-algebra-section-determinants}에서 보았듯이 표준 원소
$$
\delta(L^\bullet) \in \det(L^\bullet)
$$
가 존재한다.
% P02-E0703: source line 36941 misspells “determinant” as “determininant”.
이는 단지 $d : L^{-1} \to L^0$의 행렬식이다. $\delta(L^\bullet)$이
$\det(L^\bullet)$의 자명화일 필요충분조건은 $L^\bullet$이 비순환인
것임에 유의하자.

\medskip\noindent
다음 조건들을 만족하는 복합체 사상
$a^\bullet : K^\bullet \to L^\bullet$을 생각하자.
\begin{enumerate}
\item $a^\bullet$은 준동형사상이다.
\item 모든 $n$에 대하여 $a^n : K^n \to L^n$은 전사이다.
\item $K^n$, $L^n$은 유한 사영 $R$-가군이고,
$n \in \{-1, 0\}$이 아닐 때는 영이다.
\end{enumerate}
이 상황에서 동형
$$
\det(a^\bullet) : \det(K^\bullet) \longrightarrow \det(L^\bullet)
$$
을 구성하자. 완전열 $0 \to \Ker(a^i) \to K^i \to L^i \to 0$을 사용하면
보조정리 \ref{more-algebra-lemma-det-ses}에 의해 $i = -1, 0$에 대하여 동형
$$
\gamma^i : \det(\Ker(a^i)) \otimes \det(L^i) \to \det(K^i)
$$
을 얻는다. $a^\bullet$이 준동형사상이므로 복합체
$\Ker(a^\bullet)$은 비순환이고 계수 $0$을 갖는다. 따라서 표준 원소
$\delta(\Ker(a^\bullet))$은 위에서 본 가역 $R$-가군
$\det(\Ker(a^\bullet))$의 자명화이다.
$\det(a^\bullet) : \det(K^\bullet) \to \det(L^\bullet)$을 그림
$$
\xymatrix{
\det(K^\bullet)
\ar[rr]_{\det(a^\bullet)}
& &
\det(L^\bullet)
\ar[ld]^{\delta(\Ker(a^\bullet))}
\\
& \det(L^\bullet) \otimes \det(\Ker(a^\bullet))
\ar[lu]^{\gamma^0 \otimes (\gamma^{-1})^{\otimes -1}}
}
$$
을 가환하게 하는 유일한 동형으로 정의한다.

\begin{lemma}
\label{more-algebra-lemma-canonical-element-well-defined}
$R$을 환이라 하자. $a^\bullet : K^\bullet \to L^\bullet$을 위의
(1), (2), (3)을 만족하는 $R$-가군의 복합체 사상이라 하자.
$L^\bullet$의 계수가 $0$이면 $\det(a^\bullet)$은 표준 원소
$\delta(K^\bullet)$을 $\delta(L^\bullet)$로 보낸다.
\end{lemma}

\begin{proof}
$M^i = \Ker(a^i)$로 쓰자. 그러면 짧은 완전열들의 사상
$$
\xymatrix{
0 \ar[r] &
M^{-1} \ar[r] \ar[d]_{d_M} &
K^{-1} \ar[r] \ar[d]_{d_K} &
L^{-1} \ar[r] \ar[d]_{d_L} &
0 \\
0 \ar[r] &
M^0 \ar[r] &
K^0 \ar[r] &
L^0 \ar[r] &
0
}
$$
을 얻는다. 보조정리 \ref{more-algebra-lemma-det-ses-functorial}에 의해 사상으로서
$\det(d_K)$가 $\det(d_M) \otimes \det(d_L)$에 대응함을 안다.
정의들을 풀어 쓰면 필요한 등식을 얻는다.
\end{proof}

\begin{lemma}
\label{more-algebra-lemma-homotopic-surjections}
$R$을 환이라 하자. $a^\bullet : K^\bullet \to L^\bullet$을 위의
(1), (2), (3)을 만족하는 $R$-가군의 복합체 사상이라 하자.
$h : K^0 \to L^{-1}$을
$b^0 = a^0 + d \circ h$와 $b^{-1} = a^{-1} + h \circ d$가 전사가 되게
하는 사상이라 하자. 그러면
$\det(a^\bullet) = \det(b^\bullet)$가 사상
$\det(K^\bullet) \to \det(L^\bullet)$으로서 성립한다.
\end{lemma}

\begin{proof}
$h = a^{-1} \circ \tilde h$이고
$k^0 = \text{id} + d \circ \tilde h : K^0 \to K^0$ 및
% P02-E0704: source line 37029 names this degree -1 map k^1.
$k^1 = \text{id} + \tilde h \circ d : K^{-1} \to K^{-1}$가 동형이 되게 하는
사상 $\tilde h : K^0 \to K^{-1}$이 존재한다고 하자. 그러면 복합체들의
가환 그림
$$
\xymatrix{
0 \ar[r] &
\Ker(b^\bullet) \ar[r] \ar[d]_{c^\bullet} &
K^\bullet \ar[r]_{b^\bullet}
\ar[d]_{k^\bullet} &
L^\bullet \ar[r] \ar[d]^{\text{id}} &
0 \\
0 \ar[r] &
\Ker(a^\bullet) \ar[r] &
K^\bullet \ar[r]^{a^\bullet} &
L^\bullet \ar[r] &
0
}
$$
을 얻는다. 여기서 $c^\bullet$은 유도된 핵들의 동형이다.
보조정리 \ref{more-algebra-lemma-det-ses-functorial}을 사용하면 그림
$$
\xymatrix{
\det(\Ker(b^i)) \otimes \det(L^i) \ar[r] \ar[d]_{\det(c^i) \otimes 1} &
\det(K^i) \ar[d]^{\det(k^i)} \\
\det(\Ker(a^i)) \otimes \det(L^i) \ar[r] &
\det(K^i)
}
$$
이 가환함을 알 수 있다.
% P02-E0705: source lines 37056-37058 reverse the canonical-trivialization direction and omit det around the target kernel.
$\det(c^\bullet)$가 $\det(\Ker(a^\bullet))$의 표준 자명화를
$\Ker(b^\bullet)$의 표준 자명화로 보내므로(보조정리
\ref{more-algebra-lemma-canonical-element-well-defined}), 다음이 성립할 때 그리고
그때에만 결론을 얻는다.
% P02-E0706: source line 37059 has the malformed phrase “we see that we conclude”.
$$
\det(k^0) = \det(k^{-1})
$$
이는 $R$의 원소들 사이의 등식이며 보조정리 \ref{more-algebra-lemma-det-switch}에서
따른다.

\medskip\noindent
두 사상 $a^{-1}|_U : U \to L^{-1}$와
$b^{-1}|_U : U \to L^{-1}$가 모두 동형이 되게 하는 직합 인자
$U \subset K^{-1}$이 존재한다고 하자. $\tilde h$를 $h$와
$a^{-1}|_U$의 역의 합성으로 정의한다. $\tilde h$가 증명의 첫 문단과
같은 사상이라고 주장한다. 실제로 구성에 의해
$h = a^{-1} \circ \tilde h$이다.
$k^{-1} : K^{-1} \to K^{-1}$가 동형임을 보이려면 전사임을 보이면
충분하다(Algebra, 보조정리 \ref{algebra-lemma-fun}).
$u \in U$라 하자. $b^{-1}(u') = a^{-1}(u)$가 되도록 $u' \in U$를
택할 수 있다. 그러면 $u = k^{-1}(u')$이다. 실제로 $u$와
$k^{-1}(u')$는 모두 $U$에 속하고, 다음 계산에 의해
$a^{-1}(u) = a^{-1}(k^{-1}(u'))$이다.\footnote{$a^{-1}(k^{-1}(u')) =
a^{-1}(u') + a^{-1}(\tilde h(d(u'))) =
a^{-1}(u') + h(d(u')) = b^{-1}(u') = a^{-1}(u)$} $a^{-1}|_U$가
동형이므로 등식을 얻는다. 따라서 $U \subset \Im(k^{-1})$이다.
한편 $x \in \Ker(a^{-1})$이면 $x = k^{-1}(x) \bmod U$이다.
$K^{-1} = \Ker(a^{-1}) + U$이므로 $k^{-1}$은 전사라고 결론내린다.
마지막으로 $k^0 : K^0 \to K^0$이 전사임을 보이자. 먼저
$a^0 \circ k^0 = b^0$이므로 $a^0 \circ k^0$이 전사임을 알 수 있다.
$x \in \Ker(a^0)$이면 어떤 $y \in \Ker(a^{-1})$에 대하여 $x = d(y)$이다.
위에서 본 바에 의해 어떤 $z \in K^{-1}$에 대하여
$y = k^{-1}(z)$로 쓸 수 있다. 그러면 $x = k^0(d(z))$이고 결론을 얻는다.

\medskip\noindent
증명의 마지막 단계. 앞 문단과 같은 $U$를 찾으면 충분하지만, 항상 가능하지는
않다. 그러나 두 $R$-가군 사상이 같음을 보이려면 $R$의 소아이디얼들에서
국소화한 뒤 이를 보이면 충분하다. 따라서 $R$이 국소환이라고 가정할 수 있다.
그러면 다음 문제가 된다. 두 전사 $R$-선형 사상
$$
\alpha, \beta : R^{\oplus n} \longrightarrow R^{\oplus m}
$$
이 주어졌다고 하자. $\alpha|_U$와 $\beta|_U$가 모두 동형이 되게 하는
직합 인자 $U \subset R^{\oplus n}$을 찾아야 한다. $R$이 체이면 이는
선형대수로 가능하다. 일반적인 경우에는 잉여체 위에서 해를 택하고 이를
국소환 $R$ 위의 해로 올린다. 몇 가지 세부 사항은 생략한다.
\end{proof}

\begin{lemma}
\label{more-algebra-lemma-compose-surjections}
$R$을 환이라 하자. $a^\bullet : K^\bullet \to L^\bullet$과
$b^\bullet : L^\bullet \to M^\bullet$을 위의 (1), (2), (3)을 만족하는
$R$-가군의 복합체 사상이라 하자. 그러면
$\det(b^\bullet) \circ \det(a^\bullet) =
\det(b^\bullet \circ a^\bullet)$가 사상
% P02-E0707: source line 37114 reverses the stated direction of the determinant map.
$\det(M^\bullet) \to \det(K^\bullet)$으로서 성립한다.
\end{lemma}

\begin{proof}
생략한다. 힌트: 보조정리들 \ref{more-algebra-lemma-det-ses},
\ref{more-algebra-lemma-det-ses-functorial}, \ref{more-algebra-lemma-det-filtration}에서 바로
따른다.
\end{proof}

\begin{lemma}
\label{more-algebra-lemma-determinant-two-term-complexes}
$R$을 환이라 하자. 위의 구성들은 함자
$$
\det :
\left\{
\begin{matrix}
\text{category of perfect complexes} \\
\text{with tor amplitude in }[-1, 0] \\
\text{morphisms are isomorphisms}
\end{matrix}
\right\}
\longrightarrow
\left\{
\begin{matrix}
\text{category of invertible modules} \\
\text{morphisms are isomorphisms}
\end{matrix}
\right\}
$$
를 정한다. 더욱이 $D(R)$에서 $[-1, 0]$에 Tor 진폭을 갖는 계수 $0$의
퍼펙트 대상 $L$이 주어지면, 모든 $D(R)$의 동형 $a : L \to K$에 대하여
$\det(a)(\delta(L)) = \delta(K)$가 되게 하는 표준 원소
$\delta(L) \in \det(L)$이 존재한다.
\end{lemma}

\begin{proof}
보조정리 \ref{more-algebra-lemma-perfect}에 의해 원천 범주의 모든 대상은 복합체
$$
L^\bullet = \ldots \to 0 \to L^{-1} \to L^0 \to 0 \to \ldots
$$
로 표현할 수 있다. 여기서 $L^{-1}$과 $L^0$은 유한 사영 $R$-가군이다.
이런 유형의 복합체를 잠시 좋은 복합체라고 부르자. 유도 범주, 보조정리
\ref{derived-lemma-morphisms-from-projective-complex}에 의해 유도 범주에서
좋은 복합체 사이의 사상들은 복합체 사상들의 호모토피류이다. 따라서 좋은
복합체들로 작업할 수 있고, 위에서 살펴본 행렬식
$\det(L^\bullet) = \det(L^0) \otimes \det(L^{-1})^{\otimes -1}$을
사용할 수 있다.

\medskip\noindent
$a^\bullet : L^\bullet \to K^\bullet$을 $D(R)$에서 동형인 좋은 복합체
사이의 사상, 즉 준동형사상이라 하자. 그림
$$
\xymatrix{
L^\bullet \ar[rr]_{a^\bullet} & & K^\bullet \\
& M^\bullet \ar[lu]^{b^\bullet} \ar[ru]_{c^\bullet}
}
$$
이 호모토피까지 가환하고 $b^\bullet$과 $c^\bullet$이 위의 조건
(1), (2), (3)을 만족하면 이를 좋은 그림이라고 한다. 이러한 그림이
주어지면
$$
\det(a^\bullet) = \det(c^\bullet) \circ \det(b^\bullet)^{-1}
$$
로 정의할 수 있다. 여기서 $\det(c^\bullet)$과 $\det(b^\bullet)$은
위의 글에서 구성한 동형들이다. 좋은 그림이 항상 존재하고, 그로부터 얻는
사상 $\det(a^\bullet)$이 좋은 그림의 선택과 무관함을 보이겠다.

\medskip\noindent
좋은 복합체들의 준동형사상 $a^\bullet : L^\bullet \to K^\bullet$에
대한 좋은 그림의 존재성. 전사 $p : R^{\oplus n} \to K^{-1}$을 택하자.
그러면 새로운 좋은 복합체
$$
M^\bullet = \ldots \to 0 \to
L^{-1} \oplus R^{\oplus n} \xrightarrow{d \oplus 1}
L^0 \oplus R^{\oplus n} \to 0 \to \ldots
$$
를 생각할 수 있다. 차수 $-1$에서 사영 사상
$b^\bullet : M^\bullet \to L^\bullet$과 $a^{-1} \oplus p$를 사용하고,
차수 $0$에서 $a^0 \oplus d \circ p$를 사용하는 사상
$c^\bullet : M^\bullet \to K^\bullet$을 생각하자. 사상
$b^\bullet : M^\bullet \to L^\bullet$과
$c^\bullet : M^\bullet \to K^\bullet$은 위의 조건 (1), (2), (3)을
만족하며 좋은 그림을 얻는다.

\medskip\noindent
좋은 그림
$$
\xymatrix{
L^\bullet \ar[rr]_{\text{id}^\bullet} & & L^\bullet \\
& M^\bullet \ar[lu]^{b^\bullet} \ar[ru]_{c^\bullet}
}
$$
이 주어졌다고 하자. 보조정리 \ref{more-algebra-lemma-homotopic-surjections}에 의해
$\det(c^\bullet) = \det(b^\bullet)$임을 알 수 있다. 따라서
$\det(\text{id}^\bullet) = \text{id}$는 좋은 그림의 선택과 무관하다.

\medskip\noindent
일반적인 선택 무관성을 증명하기 전에 합성을 생각하자. 좋은 복합체들의
준동형사상 $L_1^\bullet \to L_2^\bullet$과
$L_2^\bullet \to L_3^\bullet$ 및 좋은 그림
$$
\vcenter{
\xymatrix{
L_1^\bullet \ar[rr] & &
L_2^\bullet \\
& M_{12}^\bullet \ar[lu] \ar[ru]
}
}
\quad\text{and}\quad
\vcenter{
\xymatrix{
L_2^\bullet \ar[rr] & &
L_3^\bullet \\
& M_{23}^\bullet \ar[lu] \ar[ru]
}
}
$$
들이 주어졌다고 하자. 이를 그림
$$
\xymatrix{
L_1^\bullet \ar[rr] & &
L_2^\bullet \ar[rr] & &
L_3^\bullet \\
& M_{12}^\bullet \ar[lu] \ar[ru]
& & M_{23}^\bullet \ar[lu] \ar[ru] \\
& &
M_{123}^\bullet \ar[lu] \ar[ru]
}
$$
으로 확장할 수 있다. 여기서 $M_{123}^\bullet \to M_{12}^\bullet$과
$M_{123}^\bullet \to M_{23}^\bullet$은 성질 (1), (2), (3)을 갖고,
그림의 사각형은 가환한다. 실제로
$M_{123}^n = M_{12}^n \times_{L_2^n} M_{23}^n$으로 두면 된다.
그러면 보조정리 \ref{more-algebra-lemma-compose-surjections}는 그림
$$
\xymatrix{
\det(L_2^\bullet) &
\det(M_{23}^\bullet) \ar[l] \\
\det(M_{12}^\bullet) \ar[u] &
\det(M_{123}^\bullet) \ar[l] \ar[u]
}
$$
이 가환함을 보인다. 그림 추적으로부터 $M_{12}^\bullet$과
$M_{23}^\bullet$을 사용하는 두 좋은 그림에 대응하는 사상들의 합성
$\det(L_1^\bullet) \to \det(L_2^\bullet) \to \det(L_3^\bullet)$이
좋은 그림
$$
\xymatrix{
L_1^\bullet \ar[rr] & & L_3^\bullet \\
& M_{123}^\bullet \ar[lu] \ar[ru]
}
$$
에 대응하는 사상과 같음을 알 수 있다. 따라서 이 사상들이 선택과 무관함을
보이면 합성 법칙도 만족하여 원하는 함자를 얻는다.

\medskip\noindent
선택 무관성. 좋은 복합체들의 준동형사상
$a^\bullet : L^\bullet \to K^\bullet$이 주어졌다고 하자. 역
준동형사상 $b^\bullet : K^\bullet \to L^\bullet$을 택하자.
% P02-E0708: source lines 37280-37282 drop the bullet on L and omit the grammatical subject.
$L_1^\bullet = L$, $L_2^\bullet = K^\bullet$,
$L_3^\bullet = L^\bullet$로 놓고, $b^\bullet$에 대한 좋은 그림을
하나 고정한 뒤 $a^\bullet$에 대한 좋은 그림을 변화시켜 생각하자.
앞 문단들의 결과에 따르면 어떤 선택을 하더라도 합성은 항상
$\det(L^\bullet)$ 위의 항등사상과 같다. 이는 이 선택들로부터의
무관성을 분명히 증명한다.

\medskip\noindent
표준 원소에 관한 명제는 보조정리
\ref{more-algebra-lemma-canonical-element-well-defined}와 우리의 구성에서 즉시
따른다.
\end{proof}

\section{값매김환의 확대}
\label{more-algebra-section-valuation-rings}

\noindent
이 절은 일반 값매김환에 대하여 절 \ref{more-algebra-section-discrete-valuation-rings}와
대응하는 내용을 다룬다.

\begin{definition}
\label{more-algebra-definition-extension-valuation-rings}
$A$와 $B$가 값매김환이고 $A \to B$가 단사인 국소 준동형이면,
$A \to B$ 또는 $A \subset B$를 {\it 값매김환의 확대}라고 한다.
이러한 확대는 가환 그림
$$
\xymatrix{
A \setminus \{0\} \ar[r] \ar[d]_v & B \setminus \{0\} \ar[d]^v \\
\Gamma_A \ar[r] & \Gamma_B
}
$$
을 유도한다. 여기서 $\Gamma_A$와 $\Gamma_B$는 값군이다. 아래쪽 수평
화살표가 전단사이면 $B$가 $A$ 위에서 {\it 약하게 비분기}라고 한다.
잉여체 확대
$\kappa_A = A/\mathfrak m_A \subset \kappa_B = B/\mathfrak m_B$가
유한이면 $f = [\kappa_B : \kappa_A]$로 놓고, 이를 확대
$A \subset B$의 {\it 잔여차수} 또는 {\it 잉여차수}라고 한다.
\end{definition}

\noindent
% P02-E0709: source lines 37328-37329 say “inverse” rather than “inverse image”.
$A$의 단위원들은 사상 $A \to B$ 아래에서 $B$의 단위원들의 역이므로
$\Gamma_A \to \Gamma_B$는 단사임에 유의하자. 또한 분수체의 확대가
유한이라고 요구하지 않음에도 유의하자.

\begin{lemma}
\label{more-algebra-lemma-inequality-general}
$A \subset B$를 분수체가 $K \subset L$인 값매김환의 확대라 하자.
확대 $L/K$가 유한이면 잉여체 확대가 유한이고, $\Gamma_B$에서
$\Gamma_A$의 지수가 유한이며,
$$
[\Gamma_B : \Gamma_A] [\kappa_B : \kappa_A] \leq [L : K].
$$
이다.
\end{lemma}

\begin{proof}
$b_1, \ldots, b_n \in B$를 그 $\kappa_B$에서의 상들이 $\kappa_A$ 위에서
선형독립인 단위원들이라 하자. $c_1, \ldots, c_m \in B$를 그
$\Gamma_B/\Gamma_A$에서의 상들이 서로 다른 영이 아닌 원소들이라 하자.
$b_i c_j$들이 $L$에서 $K$-선형독립이라고 주장한다. 즉, 모두 영은 아닌
$a_{ij} \in K$를 계수로 하는 합
$$
\sum a_{ij} b_i c_j
$$
은 영일 수 없다고 주장한다. $v(a_{i_0j_0}b_{i_0}c_{j_0})$가 최소가
되도록 $(i_0, j_0)$를 택하자. $a_{ij}$를 $a_{ij}/a_{i_0j_0}$로
바꾸어서 $a_{i_0 j_0} = 1$이 되게 한다.
$$
P = \{(i, j) \mid
v(a_{ij}b_ic_j) = v(a_{i_0j_0}b_{i_0}c_{j_0}) \}
$$
로 놓자. $c_1, \ldots, c_m$의 선택에 의해 $(i, j) \in P$이면
$j = j_0$임을 알 수 있다. 따라서 $(i, j) \in P$이면
$v(a_{ij}) = v(a_{i_0j_0}) = 0$이고, 즉 $a_{ij}$는 단위원이다.
$b_1, \ldots, b_n$의 선택에 의해
$$
\sum\nolimits_{(i, j) \in P} a_{ij}b_i
$$
는 $B$의 단위원이다. 따라서
$\sum\nolimits_{(i, j) \in P} a_{ij}b_ic_j$의 값매김은
$v(c_{j_0}) = v(a_{i_0j_0}b_{i_0}c_{j_0})$이다.
첫 번째 표시식에서 $(i, j) \not \in P$인 항들은 값매김이 엄밀히 더
크므로, 그 합은 영일 수 없다. 이로써 보조정리가 증명된다.
\end{proof}

\begin{lemma}
\label{more-algebra-lemma-valuation-ring-purely-inseparable}
$A$를 분수체 $K$의 표수가 $p > 0$인 값매김환이라 하자.
$L/K$를 순수 비분리 확대라 하자. 그러면 $L$에서 $A$의 정수적 폐포
$B$는 분수체가 $L$인 값매김환이고, $A \subset B$는 값매김환의
확대이다.
\end{lemma}

\begin{proof}
생략한다. 힌트: 예를 들어 Algebra, 보조정리들
\ref{algebra-lemma-x-or-x-inverse-valuation-ring}과
\ref{algebra-lemma-integral-overring-surjective}를 사용하라.
\end{proof}

\begin{lemma}
\label{more-algebra-lemma-extension-normal-domains-and-roots}
$A \to B$를 뇌터 국소 정규 정역들의 평탄 국소 준동형이라 하자.
$f \in A$와 $h \in B$가 어떤 $n > 1$과 $B$의 어떤 단위원 $w$에 대하여
$f = w h^n$을 만족한다고 하자. 높이 $1$인 모든 소아이디얼
$\mathfrak p \subset A$에 대하여, $\mathfrak p$ 위에 놓이고 확대
$A_\mathfrak p \subset B_\mathfrak q$가 약하게 비분기인 높이 $1$의
소아이디얼 $\mathfrak q \subset B$가 존재한다고 가정하자. 그러면 어떤
$g \in A$와 $A$의 단위원 $u$에 대하여 $f = u g^n$이다.
\end{lemma}

\begin{proof}
$A$와 $B$를 높이 $1$인 소아이디얼에서 국소화한 국소환들은 이산
값매김환이다(Algebra, 보조정리
\ref{algebra-lemma-characterize-dvr}). 따라서 이 가정은 정의
\ref{more-algebra-definition-extension-discrete-valuation-rings}를 통해 의미가 있다.
$\mathfrak p_1, \ldots, \mathfrak p_r$을 $f$ 위에서 극소인 $A$의
소아이디얼들이라 하자. Algebra, 보조정리
\ref{algebra-lemma-minimal-over-1}에 의해 이들의 높이는 $1$이다.
각 $i$에 대해 $\mathfrak q_{i, j} \subset B$, $j = 1, \ldots, r_i$를
$\mathfrak p_i$ 위에 놓인 높이 $1$의 $B$의 소아이디얼들이라 하자.
$A_{\mathfrak p_i} \to B_{\mathfrak q_{i, 1}}$이 약하게 비분기가 되도록
번호를 붙였다고 하자. $f$가 $B_{\mathfrak q_{i, 1}}$에서 단위원과
$n$제곱의 곱으로 가므로, $A_{\mathfrak p_i}$에서 $f$의 값매김 $v_i$는
$n$으로 나누어떨어진다. 어떤 $w_i \geq 0$에 대하여
$v_i = n w_i$라고 하자. 완전열
$$
0 \to I \to A \to
\prod\nolimits_{i = 1, \ldots, r}
A_{\mathfrak p_i}/\mathfrak p_i^{w_i}A_{\mathfrak p_i}
$$
이 아이디얼 $I$를 정의한다고 하자. 완전함자 $- \otimes_A B$를 적용하면
완전열
$$
0 \to I \otimes_A B \to B \to
\prod\nolimits_{i = 1, \ldots, r}
(A_{\mathfrak p_i}/\mathfrak p_i^{w_i}A_{\mathfrak p_i}) \otimes_A B
$$
을 얻는다. $i$를 고정하자. 표준 사상
$$
(A_{\mathfrak p_i}/\mathfrak p_i^{w_i}A_{\mathfrak p_i}) \otimes_A B
\to
\prod\nolimits_{j = 1, \ldots, r_i}
B_{\mathfrak q_{i, j}}/\mathfrak q_{i, j}^{e_{i, j}w_i}B_{\mathfrak q_{i, j}}
$$
이 단사라고 주장한다. 여기서 $e_{i, j}$는
$A_{\mathfrak p_i} \to B_{\mathfrak q_{i, j}}$의 분기지수이다.
이 주장은 $\mathfrak p_i^{w_i}B_{\mathfrak p_i}$가
$\mathfrak q_{i, j}$에서의 값매김이 $\geq e_{i, j}w_i$인
$B_{\mathfrak p_i}$의 원소 $b$들의 집합과 같다는 말이다. 주아이디얼
$\mathfrak p_i^{w_i}$의 생성원 $a \in A_{\mathfrak p_i}$를 택하자.
그러면 $\mathfrak q_{i, j}$에서 $a$의 값매김은 $e_{i, j}w_i$와 같다.
$B_{\mathfrak p_i}$는 높이 1의 소아이디얼들이
$\mathfrak q_{i, j}$, $j = 1, \ldots, r_i$인 정규 정역이므로, 위와 같은
$b$에 대하여 Algebra, 보조정리
\ref{algebra-lemma-normal-domain-intersection-localizations-height-1}에 의해
$b/a \in B_{\mathfrak p_i}$임을 알 수 있다. 이로써 주장이 증명된다.

\medskip\noindent
이 주장과 위의 두 번째 완전열을 합치면 완전열
$$
0 \to I \otimes_A B \to B \to
\prod\nolimits_{i = 1, \ldots, r}
\prod\nolimits_{j = 1, \ldots, r_i}
B_{\mathfrak q_{i, j}}/\mathfrak q_{i, j}^{e_{i, j}w_i}B_{\mathfrak q_{i, j}}
$$
을 얻는다. 따라서 $I \otimes_A B$는 $\mathfrak q_{i, j}$에서의
값매김이 $\geq e_{i, j}w_i$인 $B$의 원소 $h'$들의 집합이다.
$B$에서 $f = wh^n$이므로 $h$의 $\mathfrak q_{i, j}$에서의 값매김은
$e_{i, j}w_i$이다. 따라서 Algebra, 보조정리
\ref{algebra-lemma-normal-domain-intersection-localizations-height-1}에 의해
$h'/h \in B$이다. 그러므로 $I \otimes_A B$는 $h$로 생성되는 계수
$1$의 자유 $B$-가군이다. 따라서 $I$는 계수 $1$의 자유 $A$-가군이다.
Algebra, 보조정리 \ref{algebra-lemma-finite-projective-descends}를 보라.
생성원 $g \in I$를 택하자. 그러면 $g$와 $h$는 $B$에서 단위원만큼
차이 난다. 거꾸로 따라가면 $A_{\mathfrak p_i}$에서 $g$의 값매김은
$w_i = v_i/n$이라고 결론내린다. 따라서 원하는 대로 Algebra, 보조정리
\ref{algebra-lemma-normal-domain-intersection-localizations-height-1}에 의해
$g^n$과 $f$는 $A$에서 단위원만큼 차이 난다.
\end{proof}

\begin{lemma}
\label{more-algebra-lemma-etale-extension-valuation-ring}
$A$를 값매김환이라 하자. $A \to B$를 에탈 환 사상이라 하고,
$\mathfrak m \subset B$를 $A$의 극대아이디얼 위에 놓인 소아이디얼이라
하자. 그러면 $A \subset B_\mathfrak m$은 약하게 비분기인 값매김환의
확대이다.
\end{lemma}

\begin{proof}
보조정리 \ref{more-algebra-lemma-weak-dimension-at-most-1}에 의해 환 $A$의 약한 차원은
$\leq 1$이다. 그러면 보조정리들 \ref{more-algebra-lemma-weak-dimension-goes-up}과
\ref{more-algebra-lemma-when-weakly-etale}에 의해 $B$의 약한 차원도 $\leq 1$이다.
% P02-E0710: source lines 37491-37496 begin a sentence with lowercase “hence” and use a spurious article before Gamma_A.
따라서 보조정리 \ref{more-algebra-lemma-weak-dimension-at-most-1}에 의해 국소환
$B_\mathfrak m$은 값매김환이다. 확대 $A \subset B_\mathfrak m$은
분수체의 유한 확대를 유도하므로, $\Gamma_A$는 $B_{\mathfrak m}$의
값군에서 유한 지수를 갖는다. 따라서 모든 $h \in B_\mathfrak m$에
대하여 어떤 $n > 0$, 원소 $f \in A$, 단위원 $w \in B_\mathfrak m$가
존재하여 $B_\mathfrak m$에서 $f = w h^n$이다. 이것이 어떤
$g \in A$와 단위원 $u \in A$에 대하여 $f = ug^n$임을 함의함을 보이겠다.
그러면 보조정리에서 주장한 대로 $A$와 $B_\mathfrak m$의 값군이 같음이
따른다.

\medskip\noindent
이를 $\mathbf{Z}$ 위에서 본질적 유한형인 국소 부분환들의 쌍대극한
$A = \colim A_i$로 쓰자. $A$는 정규 정역이므로(Algebra, 보조정리
\ref{algebra-lemma-valuation-ring-normal}), 각 $A_i$가 정규라고 가정해도
된다. 여기서는 정규화를 취할 때에도 국소환들이 $\mathbf{Z}$ 위에서
본질적 유한형임을 사용한다(Algebra, 명제
\ref{algebra-proposition-ubiquity-nagata}). 어떤 $i$에 대하여
$B = A \otimes_{A_i} B_i$가 되는 에탈 확대 $A_i \to B_i$를 찾을 수
있다. Algebra, 보조정리 \ref{algebra-lemma-etale}를 보라.
$\mathfrak m_i$를 $B_i$와 $\mathfrak m$의 교집합이라 하자. 그러면 환
사상 $A_i \to (B_i)_{\mathfrak m_i}$에 보조정리
\ref{more-algebra-lemma-extension-normal-domains-and-roots}를 적용하여 결론을 얻을
수 있다. 보조정리의 가정들은 다음 이유로 만족된다.
\begin{enumerate}
\item $A_i$와 $(B_i)_{\mathfrak m_i}$는 $\mathbf{Z}$ 위에서 본질적
유한형이므로 뇌터이다.
\item $A_i \to B_i$가 에탈이므로 $A_i \to (B_i)_{\mathfrak m_i}$는
평탄하다.
\item $A_i \to B_i$가 에탈이므로 $B_i$는 정규이다. Algebra, 보조정리
\ref{algebra-lemma-normal-goes-up}을 보라.
\item Algebra, 보조정리 \ref{algebra-lemma-finite-in-codim-1}과
$\Spec((B_i)_{\mathfrak m_i}) \to \Spec(A_i)$가 전사라는 사실에 의해,
$A_i$의 높이 $1$인 모든 소아이디얼 위에는
$(B_i)_{\mathfrak m_i}$의 높이 $1$인 소아이디얼이 존재한다.
\item $A_i \to B_i$가 에탈이므로, 소아이디얼 $\mathfrak p$ 위에 놓인
모든 소아이디얼 $\mathfrak q$에 대하여 유도된 확대
$(A_i)_\mathfrak p \to (B_i)_\mathfrak q$는 비분기이다.
\end{enumerate}
이로써 보조정리의 증명이 끝난다.
\end{proof}

\begin{lemma}
\label{more-algebra-lemma-henselization-valuation-ring}
$A$를 값매김환이라 하자. $A^h$와 $A^{sh}$를 각각 그 헨젤화와 엄밀
헨젤화라 하자. 그러면
$$
A \subset A^h \subset A^{sh}
$$
는 값군 위에서 전단사를 유도하는, 즉 약하게 비분기인 값매김환의 확대들이다.
\end{lemma}

\begin{proof}
$A^h = \colim (B_i)_{\mathfrak q_i}$로 쓰자. 여기서 $A \to B_i$는
에탈이고 $\mathfrak q_i \subset B_i$는 $\mathfrak m_A$ 위에 놓인
소아이디얼이다. Algebra, 보조정리
\ref{algebra-lemma-henselization-different}를 보라. 그러면 보조정리
\ref{more-algebra-lemma-etale-extension-valuation-ring}에 의해
$(B_i)_{\mathfrak q_i}$는 값매김환이고 유도된 사상
$$
(A \setminus \{0\})/A^* \longrightarrow
((B_i)_{\mathfrak q_i} \setminus \{0\}) / (B_i)_{\mathfrak q_i}^*
$$
은 전단사이다. Algebra, 보조정리
\ref{algebra-lemma-colimit-valuation-rings}에 의해 $A^h$가 값매김환임을
얻는다. 또한 $(A \setminus \{0\})/A^* \to
(A^h \setminus \{0\})/(A^h)^*$가 전단사임도 따른다. 이는 포함
$A \subset A^h$에 대한 보조정리를 증명한다. $A \subset A^{sh}$에
대해서는 Algebra, 보조정리
\ref{algebra-lemma-henselization-different}를 Algebra, 보조정리
\ref{algebra-lemma-strict-henselization-different}로 바꾸어 정확히 같은
논증을 사용할 수 있다. $A^{sh} = (A^h)^{sh}$이므로 이는 포함
$A^h \subset A^{sh}$에 대한 보조정리의 주장들도 증명한다.
\end{proof}

\section{PID 위 가군의 구조}
\label{more-algebra-section-structure-modules}

\noindent
증명들이 값매김환 위에서도 성립하도록 Warfield의 논문
\cite{Warfield-Purity}와 \cite{Warfield-Decomposition}을 따라
조금 더 일반적인 상황에서 논의한다.

\begin{lemma}
\label{more-algebra-lemma-characterize-PD-modules}
\begin{reference}
\cite[Corollary 1]{Warfield-Purity}
\end{reference}
$P$를 환 $R$ 위 가군이라 하자. 다음은 서로 동치이다.
\begin{enumerate}
\item $P$는 $f \in R$가 변할 때 꼴 $R/fR$인 가군들의
직합의 직합인자이다.
\item 모든 $f \in R$에 대해 $fA = A \cap fB$를 만족하는
$R$-가군의 임의의 짧은 완전열 $0 \to A \to B \to C \to 0$에 대해
사상 $\Hom_R(P, B) \to \Hom_R(P, C)$가 전사이다.
\end{enumerate}
\end{lemma}

\begin{proof}
(2)에서와 같은 완전열 $0 \to A \to B \to C \to 0$을 잡자.
(1)이 (2)를 함의함을 증명하려면 모든 $f \in R$에 대해
$\Hom_R(R/fR, B) \to \Hom_R(R/fR, C)$가 전사임을 증명하면 충분하다.
사상 $\psi : R/fR \to C$를 잡자. $\psi(1)$이 어떤 $b \in B$의
상이라고 하자. 그러면 $fb \in A$이다. 따라서 $fa = fb$인
$a \in A$가 존재한다. 이제 $f(b - a) = 0$이므로 $1$을 $b - a$로
보내는 사상 $\varphi : R/fR \to B$를 얻으며, 이는 $\psi$를 올린다.

\medskip\noindent
반대로 (2)가 성립한다고 하자. $I$를 $f \in R$이고
$\varphi : R/fR \to P$인 쌍 $(f, \varphi)$들의 집합이라 하자.
$i \in I$에 대해 대응하는 쌍을 $(f_i, \varphi_i)$로 나타낸다.
직합인자 $R/f_iR$의 원소 $r$을 $P$의 $\varphi_i(r)$로 보내는 사상
$$
B = \bigoplus\nolimits_{i \in I} R/f_iR \longrightarrow P
$$
을 생각하자. $A = \Ker(B \to P)$라 하자. 그러면 열
$$
0 \to A \to B \to P \to 0
$$
이 (2)에서와 같은 완전열이면 (1)이 참임을 알 수 있다. 이를 보이기 위해
$f \in R$이고 $a \in A$가 $B$에서 $f b$로 간다고 하자.
거의 모든 $r_i = 0$이 되도록 $b = (r_i)_{i \in I}$로 쓰자.
그러면 $P$에서
$$
f\sum \varphi_i(r_i) = 0
$$
이다. 따라서 $f_{i_0} = f$이고
$\varphi_{i_0}(1) = \sum \varphi_i(r_i)$인 $i_0 \in I$가 존재한다.
$x_{i_0} \in R/f_{i_0}R$를 $1$의 유라 하자. 그러면
$$
a' = (r_i)_{i \in I} - (0, \ldots, 0, x_{i_0}, 0, \ldots )
$$
은 $A$의 원소이고, 원하는 대로 $fa' = a$이다.
\end{proof}

\begin{lemma}[일반화 값매김환]
\label{more-algebra-lemma-generalized-valuation-ring}
\begin{reference}
\cite{Warfield-Decomposition}
\end{reference}
$R$을 영이 아닌 환이라 하자. 다음은 서로 동치이다.
\begin{enumerate}
\item $a, b \in R$에 대해 $a$가 $b$를 나누거나 $b$가 $a$를 나눈다.
\item 모든 유한 생성 아이디얼이 주 아이디얼이고 $R$은 국소환이다.
\item $R$의 아이디얼들의 집합은 포함관계에 의해 전순서가 매겨진다.
\end{enumerate}
특히 $R$이 값매김환이면 이 조건들이 성립한다.
\end{lemma}

\begin{proof}
(2)를 가정하고 $a, b \in R$이라 하자. 그러면 $(a, b) = (c)$이다.
$c = 0$이면 $a = b = 0$이고 $a$가 $b$를 나눈다.
$c \not = 0$이라 하자. $c = ua + vb$, $a = wc$, $b = zc$로 쓰면
$c(1 - uw - vz) = 0$이다. $R$이 국소환이므로 이는
$1 - uw - vz \in \mathfrak m$을 함의한다. 따라서 $w$와 $z$ 중
하나는 단위원이므로 $a$가 $b$를 나누거나 $b$가 $a$를 나눈다.
따라서 (2)는 (1)을 함의한다.

\medskip\noindent
(1)을 가정하자. $R$에 두 극대 아이디얼 $\mathfrak m_i$가 있으면
$a \in \mathfrak m_1$이면서 $a \not \in \mathfrak m_2$인 원소와
$b \in \mathfrak m_2$이면서 $b \not \in \mathfrak m_1$인 원소를
고를 수 있다. 그러면 $a$는 $b$를 나누지 않고 $b$도 $a$를 나누지 않는다.
따라서 $R$은 유일한 극대 아이디얼을 가지므로 국소환이다.
조건 (1)과 귀납법으로부터 모든 유한 생성 아이디얼이 주 아이디얼임이
쉽게 따른다. 따라서 (1)은 (2)를 함의한다.

\medskip\noindent
(1)과 (3)이 동치임은 직접 증명할 수 있다. 마지막 명제는 Algebra,
Lemma \ref{algebra-lemma-valuation-ring-x-or-x-inverse}이다.
\end{proof}

\begin{lemma}
\label{more-algebra-lemma-generalized-valuation-ring-modules}
\begin{reference}
\cite[Theorem 1]{Warfield-Decomposition}
\end{reference}
$R$을 Lemma \ref{more-algebra-lemma-generalized-valuation-ring}의 동치 조건들을
만족하는 환이라 하자. 그러면 모든 유한 표시 $R$-가군은 꼴 $R/fR$인
가군들의 유한 직합과 동형이다.
\end{lemma}

\begin{proof}
$M$을 유한 표시 $R$-가군이라 하자. 이제부터는 따로 언급하지 않고
Lemma \ref{more-algebra-lemma-generalized-valuation-ring}에 나오는 $R$의 동치인
성질들을 모두 사용한다. $\mathfrak m \subset R$을 극대 아이디얼,
$\kappa = R/\mathfrak m$을 잉여체라 하자. $I \subset R$을
$M$의 소멸자라 하자. 유한 차원 $\kappa$-벡터공간
$M/\mathfrak m M$의 기저 $y_1, \ldots, y_n$을 고른다.
$n$에 대한 귀납법으로 논증하겠다.

\medskip\noindent
나카야마 보조정리에 의해 $M/\mathfrak m M$의 원소
$y_1, \ldots, y_n$을 올리는 임의의 원소들
$x_1, \ldots, x_n \in M$은 $M$을 생성한다. Algebra, Lemma
\ref{algebra-lemma-NAK}를 보라. 이는 귀납법의 기초 단계 $n = 0$을
즉시 증명한다.

\medskip\noindent
다음 성질을 만족하는 지표 $i$가 존재한다고 주장한다. 즉,
$y_i$로 가는 $x_i \in M$을 어떻게 고르든 $x_i$의 소멸자는 $I$이다.
그렇지 않다면 모든 $i$에 대해
$I_i = \text{Ann}(x_i) \not = I$인 $x_1, \ldots, x_n$을 고를 수 있다.
하지만 모든 $i$에 대해 $I \subset I_i$이고 아이디얼들이 전순서로
정렬되어 있으므로 $i = 1, \ldots, n$에 대해 $I_i$는 $I$보다
엄밀히 크다. 전순서성을 한 번 더 사용하면
$\text{Ann}(M) = I_1 \cap \ldots \cap I_n$이 $I$보다 크다는 결론을
얻는데, 이는 모순이다. 번호를 다시 매겨 $y_1$이 다음 성질을 가진다고
가정해도 된다. 즉, $y_1$을 올리는 임의의 $x_1 \in M$에 대해
$x_1$의 소멸자는 $I$이다.

\medskip\noindent
$A = Rx_1 \subset M$이라 놓자. 완전열
$0 \to A \to M \to M/A \to 0$을 생각한다. $A$가 유한이므로
$M/A$는 생성원 수가 더 적은 유한 표시 $R$-가군이다
(Algebra, Lemma \ref{algebra-lemma-extension}). 귀납가정에 의해
$M/A \cong \bigoplus_{j = 1, \ldots, m} R/f_jR$이다.
한편 $A \to M$이 다음 성질을 만족한다고 주장한다. $f \in R$이면
$fA = A \cap fM$이다. 포함관계 $fA \subset A \cap fM$은 자명하다.
반대로 $x \in A \cap fM$이면 어떤 $g \in R$와 $y \in M$에 대해
$x = gx_1 = f y$이다. $f$가 $g$를 나누면 원하는 대로 $x \in fA$이다.
그렇지 않으면 어떤 $h \in \mathfrak m$에 대해 $f = hg$로 쓸 수 있다.
앞 단락에 의해 원소 $x'_1 = x_1 - hy$의 소멸자는 $I$이다.
따라서 $g \in I$이고 원하는 대로 $x = 0$이다. 이 주장과
Lemma \ref{more-algebra-lemma-characterize-PD-modules}에 의해 완전열
$0 \to A \to M \to M/A \to 0$은 분할되고
$M \cong A \oplus \bigoplus_{j = 1, \ldots, m} R/f_jR$을 얻는다.
그러면 $A = R/I$는 $M$의 직합인자이므로 유한 표시이고,
따라서 $I$는 유한 생성이어서 주 아이디얼이다. 이로써 증명이 끝난다.
\end{proof}

\begin{lemma}
\label{more-algebra-lemma-warfield}
\begin{reference}
\cite[Theorem 3]{Warfield-Decomposition}
\end{reference}
$R$을 각 극대 아이디얼에서의 $R$의 모든 국소환이 Lemma
\ref{more-algebra-lemma-generalized-valuation-ring}의 동치 조건들을 만족하는 환이라 하자.
그러면 모든 유한 표시 $R$-가군은 $f$가 $R$ 안에서 변할 때 꼴
$R/fR$인 가군들의 유한 직합의 직합인자이다.
\end{lemma}

\begin{proof}
$M$을 유한 표시 $R$-가군이라 하자. 먼저 $M$이 꼴 $R/fR$인
가군들의 직합의 직합인자임을 보이고, 마지막에 이 직합을 유한하게
잡을 수 있음을 논증하겠다. 모든 $f \in R$에 대해
$fA = A \cap fB$를 만족하는 $R$-가군의 짧은 완전열
$$
0 \to A \to B \to C \to 0
$$
을 잡자. Lemma \ref{more-algebra-lemma-characterize-PD-modules}에 의해
$\Hom_R(M, B) \to \Hom_R(M, C)$가 전사임을 보여야 한다.
극대 아이디얼 $\mathfrak m$에서 국소화한 뒤 이를 증명하면 충분하다.
Algebra, Lemma \ref{algebra-lemma-characterize-zero-local}를 보라.
국소화된 열
$0 \to A_\mathfrak m \to B_\mathfrak m \to C_\mathfrak m \to 0$은
모든 $f \in R_\mathfrak m$에 대해
$fA_\mathfrak m = A_\mathfrak m \cap fB_\mathfrak m$을 만족함에 유의하라.
실제로 $u \in R_\mathfrak m$은 단위원이고 $f' \in R$이 되도록
$f = uf'$로 쓸 수 있으며, 국소화가 완전하기 때문이다.
$M$이 유한 표시이므로 Algebra, Lemma
\ref{algebra-lemma-hom-from-finitely-presented}에 의해
$$
\Hom_R(M, B)_\mathfrak m = \Hom_{R_\mathfrak m}(M_\mathfrak m, B_\mathfrak m)
\quad\text{and}\quad
\Hom_R(M, C)_\mathfrak m = \Hom_{R_\mathfrak m}(M_\mathfrak m, C_\mathfrak m)
$$
이다. 가군 $M_\mathfrak m$은 유한 표시 $R_\mathfrak m$-가군이다.
Lemma \ref{more-algebra-lemma-generalized-valuation-ring-modules}에 의해
$M_\mathfrak m$은 꼴 $R_\mathfrak m/fR_\mathfrak m$인 가군들의
직합이다. 따라서 Lemma \ref{more-algebra-lemma-characterize-PD-modules}에 의해
국소화된 사상은 전사이다.

\medskip\noindent
이제 $M$이 $\bigoplus_{i \in I} R/f_i R$의 직합인자임을 안다.
사상 $M \to \bigoplus_{i \in I} R/f_i R$을 생각하자.
$M$은 유한 $R$-가군이므로 어떤 유한 부분집합 $I' \subset I$에 대해
그 상은 $\bigoplus_{i \in I'} R/f_i R$에 포함된다.
이로써 증명이 끝난다.
\end{proof}

\begin{definition}
\label{more-algebra-definition-bezout}
$R$을 정역이라 하자.
\begin{enumerate}
\item $R$의 모든 유한 생성 아이디얼이 주 아이디얼이면
$R$을 {\it B\'ezout 정역}이라 한다.
\item 모든 $n , m \geq 1$과 모든 $n \times m$ 행렬 $A$에 대해
크기가 각각 $n \times n, m \times m$인 가역행렬 $U, V$가 존재하여
$$
U A V =
\left(
\begin{matrix}
f_1 & 0 & 0 & \ldots \\
0 & f_2 & 0 & \ldots \\
0 & 0 & f_3 & \ldots \\
\ldots & \ldots & \ldots & \ldots
\end{matrix}
\right)
$$
이고, $f_1, \ldots, f_{\min(n, m)} \in R$이며
$f_1 | f_2 | \ldots$이면 $R$을 {\it 기본인자 정역}이라 한다.
\end{enumerate}
\end{definition}

\noindent
모든 B\'ezout 정역 $R$이 기본인자 정역인지 아닌지는 여전히
미해결 문제인 것으로 보인다. 이는 $R$ 위 모든 유한 표시 가군이
순환가군들의 직합인지 묻는 문제와 동치이다. 역방향 함의는 참이다.

\begin{lemma}
\label{more-algebra-lemma-elementary-divisor-is-bezout}
기본인자 정역은 B\'ezout 정역이다.
\end{lemma}

\begin{proof}
$a, b \in R$을 영이 아닌 원소라 하자. $1 \times 2$ 행렬
$A = (a\ b)$를 생각한다. 그러면 $u \in R$은 가역이고
$V = (g_{ij})$는 가역 $2 \times 2$ 행렬인
$u(a\ b)V = (f\ 0)$임을 알 수 있다. 이제
$f = u a g_{11} + u b g_{2 1}$이고 $(g_{11}, g_{2 1}) = R$이다.
따라서 $(a, b) = (f)$이다. 그러면 귀납논증(생략함)에 의해
$R$의 임의의 유한 생성 아이디얼은 한 원소로 생성된다.
\end{proof}

\begin{lemma}
\label{more-algebra-lemma-localize-bezout}
B\'ezout 정역의 국소화는 B\'ezout 정역이다.
B\'ezout 정역의 모든 국소환은 값매김환이다.
국소 정역이 B\'ezout 정역일 필요충분조건은 값매김환인 것이다.
\end{lemma}

\begin{proof}
국소화에 관한 명제의 증명은 생략한다. 마지막 명제는 Algebra, Lemma
\ref{algebra-lemma-characterize-valuation-ring}이다.
두 번째 명제는 나머지 두 명제로부터 따른다.
\end{proof}

\begin{lemma}
\label{more-algebra-lemma-split-off-free-part}
$R$을 B\'ezout 정역이라 하자.
\begin{enumerate}
\item 자유 가군의 모든 유한 부분가군은 유한 자유 가군이다.
\item 모든 유한 표시 $R$-가군 $M$은 유한 자유 가군과 비틀림 가군
$M_{tors}$의 직합이다. 여기서 이 가군은 영이 아닌
$f_1, \ldots, f_n \in R$에 대해 꼴
$\bigoplus_{i = 1, \ldots, n} R/f_iR$인 가군의 직합인자이다.
\end{enumerate}
\end{lemma}

\begin{proof}
(1)을 증명한다. $M \subset F$를 자유 가군 $F$의 유한 부분가군이라 하자.
$M$이 유한이므로 $F$가 유한 자유 가군이라고 가정해도 된다
(자세한 내용은 생략한다). $F = R^{\oplus n}$이라 하자.
$n$에 대한 귀납법으로 논증한다. $n = 1$이면 $M$은 유한 생성
아이디얼이므로, $R$이 B\'ezout 정역이라는 가정에 의해 주 아이디얼이다.
$n > 1$이면 마지막 직합인자로의 사영
$R^{\oplus n} \to R$ 아래에서 $M$의 상 $I$를 생각한다.
$I = (0)$이면 $M \subset R^{\oplus n - 1}$이고 귀납가정으로 끝난다.
$I \not = 0$이면 $I = (f) \cong R$이다. 따라서
$M \cong R \oplus \Ker(M \to I)$이고 역시 귀납가정으로 끝난다.

\medskip\noindent
$M$을 유한 표시 $R$-가군이라 하자.
% P02-E0711
$R$의 극대 아이디얼에서의 국소화들이 값매김환이므로
(Lemma \ref{more-algebra-lemma-localize-bezout}) Lemma \ref{more-algebra-lemma-warfield}를 적용할 수 있다.
따라서 $M$은 $f_i \not = 0$인
$R^{\oplus r} \oplus \bigoplus_{i = 1, \ldots, n} R/f_iR$
꼴 가군의 직합인자이다. 비틀림 부분가군을 취하는 것은 함자이므로
$M_{tors}$는 $\bigoplus_{i = 1, \ldots, n} R/f_iR$의 직합인자이고
$M/M_{tors}$는 $R^{\oplus r}$의 직합인자이다.
증명의 첫 부분에 의해 $M/M_{tors}$는 유한 자유 가군이다.
따라서 원하는 대로 $M \cong M_{tors} \oplus M/M_{tors}$이다.
\end{proof}

\begin{lemma}
\label{more-algebra-lemma-modules-PID}
$R$을 PID라 하자. 모든 유한 $R$-가군 $M$은
% P02-E0712
$$
R^{\oplus r} \oplus \bigoplus\nolimits_{i = 1, \ldots, n} R/f_iR
$$
꼴 가군과 동형이다. 여기서 $r, n \geq 0$이고
$f_1, \ldots, f_n \in R$은 영이 아니다.
\end{lemma}

\begin{proof}
PID는 Noether B\'ezout 환이다. Lemma \ref{more-algebra-lemma-split-off-free-part}에 의해
$M$이 비틀림 가군일 때 결과를 증명하면 충분하다. $M$이 유한이므로
이는 $M$의 소멸자가 영이 아님을 뜻한다. 어떤 영이 아닌 $f \in R$에 대해
$fM = 0$이라 하자. 그러면 $M$을 $R/fR$ 위 가군으로 볼 수 있다.
$R/fR$은 차원 $0$인 Noether 환이므로(작은 세부사항은 생략한다)
% P02-E0713
$R/fR = \prod R_j$는 Artin 국소환 $R_i$들의 유한 곱이다
(Algebra, Proposition \ref{algebra-proposition-dimension-zero-ring}).
각 $R_i$는 국소환이고 PID의 몫환이므로 Lemma
\ref{more-algebra-lemma-generalized-valuation-ring}의 의미에서 일반화 값매김환이다
(작은 세부사항은 생략한다). $R_j$에 대응하는 멱등원을
$e_j \in R/fR$라 하여 $M_j = e_j M$이고
$M = \prod M_j$가 되도록 쓰자. Lemma
\ref{more-algebra-lemma-generalized-valuation-ring-modules}에 의해 어떤
$\overline{f}_{ji} \in R_j$에 대해
$M_j = \bigoplus_{i = 1, \ldots, n_j} R_j/\overline{f}_{ji}R_j$이다.
올림 $f_{ji} \in R$을 고르고
% P02-E0714
$(g_{ji}) = (f_j, f_{ji})$인 $g_{ji} \in R$을 고른다.
그러면 $R$-가군으로
$$
M \cong \bigoplus R/g_{ji}R
$$
임을 얻으며, 이로써 증명이 끝난다.
\end{proof}

\noindent
% P02-E0715
PID가 기본인자 정역이라는 것도 증명할 수 있다(참조를 추후 여기에 삽입).
그 방법 하나는 다음과 비슷한 보조정리들을 증명하는 것이다.

\begin{lemma}
\label{more-algebra-lemma-unimodular-vector}
$R$을 B\'ezout 정역이라 하자. $n \geq 1$이고
$f_1, \ldots, f_n \in R$이 단위 아이디얼을 생성한다고 하자.
첫째 행이 $f_1 \ldots f_n$인 $R$ 위 가역
$n \times n$ 행렬이 존재한다.
\end{lemma}

\begin{proof}
이는 Lemma \ref{more-algebra-lemma-split-off-free-part}에서 따르지만 다음과 같이
직접 증명할 수도 있다. $n$에 대한 귀납법을 사용한다.
$n = 1$이면 결과가 성립한다. $n > 1$이라 하자. 번호를 다시 매겨
$f_1 \not = 0$이라 가정해도 된다.
$(f) = (f_1, \ldots, f_{n - 1})$인 $f \in R$을 고른다.
% P02-E1338
$A$를 첫째 행이 $f_1/f, \ldots, f_{n - 1}/f$인
$(n - 1) \times (n - 1)$ 행렬이라 하자. 다음이 가능하도록
$a, b \in R$을 고르자.
$1 \in (f_1, \ldots, f_n) = (f, f_n)$이므로
$af - bf_n = 1$이다. 그러면 한 해는 다음 행렬이다.
$$
\left(
\begin{matrix}
f & 0 & \ldots & 0 & f_n \\
0 & 1 & \ldots & 0 & 0 \\
  &   & \ldots \\
0 & 0 & \ldots & 1 & 0 \\
b & 0 & \ldots & 0 & a
\end{matrix}
\right)
\left(
\begin{matrix}
  & & & 0 \\
  & A \\
  & & & 0 \\
0 & \ldots & 0 & 1
\end{matrix}
\right)
$$
이다. 왼쪽 행렬은 행렬식이 $1$이므로 가역임에 유의하라.
\end{proof}

\section{주 근기 아이디얼}
\label{more-algebra-section-principal-radical-ideals}

\noindent
% P02-E0716
이 절에서는 카테너리 Noether 정규 국소 정역 안에 자명하지 않은
주 근기 아이디얼이 존재함을 증명한다. 이 결과는
\cite{Artin-Lipman}에서 찾을 수 있다.

\begin{lemma}
\label{more-algebra-lemma-polypoly}
$(R,\mathfrak m)$을 차원 1인 Noether 국소환이라 하고,
$x\in\mathfrak m$을 $R$의 어느 극소 소아이디얼에도 들어 있지 않은
원소라 하자. 그러면
\begin{enumerate}
\item 함수 $P : n \mapsto \text{length}_R(R/x^n R)$는
$n \geq 0$에 대해 $P(n) \leq n P(1)$을 만족한다.
\item $x$가 영인자가 아니면 $n \geq 0$에 대해 $P(n) = nP(1)$이다.
\end{enumerate}
\end{lemma}

\begin{proof}
$\dim(R) = 1$이므로 $\dim(R/x^n R) = 0$이다. 따라서 각 $n$에 대해
$\text{length}_R(R/x^n R)$은 유한하다
(Algebra, Lemma \ref{algebra-lemma-support-point}).
보조정리를 보이기 위해 $n$에 대해 귀납하겠다. $x^0 R = R$이므로
$P(0) = \text{length}_R(R/x^0R) = \text{length}_R 0 = 0$이다.
$n = 1$일 때도 명제가 성립한다. 이제 $n \geq 2$이고
$n - 1$에 대해 명제가 성립한다고 하자. 다음 열은 완전하다.
$$
R/x^{n-1}R \xrightarrow{x} R/x^nR \to R/xR \to 0
$$
여기서 $x$는 $x$를 곱하는 사상을 나타낸다. 길이가 가법적이므로
(Algebra, Lemma \ref{algebra-lemma-length-additive})
$P(n) \leq P(n - 1) + P(1)$이다. 귀납가정에 의해
$P(n - 1) \leq (n - 1)P(1)$이므로 $P(n) \leq nP(1)$이다.
이는 귀납 단계를 증명한다.

\medskip\noindent
$x$가 영인자가 아니면 위에 표시한 완전열은 왼쪽에서도 완전하다.
따라서 모든 $n \geq 1$에 대해
$P(n) = P(n - 1) + P(1)$을 얻는다.
\end{proof}

\begin{lemma}
\label{more-algebra-lemma-minprimespoly}
$(R, \mathfrak m)$을 차원 $1$인 Noether 국소환이라 하자.
$x \in \mathfrak m$을 $R$의 어느 극소 소아이디얼에도 들어 있지 않은
원소라 하자. $t$를 $R$의 극소 소아이디얼 개수라 하자. 그러면
$t \leq \text{length}_R(R/xR)$이다.
\end{lemma}

\begin{proof}
$\mathfrak p_1, \ldots, \mathfrak p_t$를 $R$의 극소 소아이디얼들이라 하자.
$R' = R/\sqrt{0} = R/(\bigcap_{i = 1}^t \mathfrak p_i)$라 놓자.
$R'$에 대해 보조정리를 증명하면 충분하다고 주장한다. 실제로
$R'$ 역시 $t$개의 극소 소아이디얼을 가지며,
$\text{length}_{R'}(R'/xR') = \text{length}_R(R'/xR')$은
$R/xR \to R'/xR'$가 전사이므로 $\text{length}_R(R/xR)$보다 작음이
분명하다. 따라서 $R$이 축약환이라 가정해도 된다.

\medskip\noindent
$R$이 극소 소아이디얼 $\mathfrak p_1, \ldots, \mathfrak p_t$를
가지는 축약환이라고 하자. 이는 다음 완전열이 있음을 뜻한다.
$$
0 \to R \to
\prod\nolimits_{i = 1}^t R/\mathfrak p_i \to Q \to 0
$$
여기서 $Q$는 첫째 사상의 여핵이다.
$M = \prod_{i = 1}^t R/\mathfrak p_i$라 쓰자.
$\mathfrak p_j$에서 국소화하면
$$
R_{\mathfrak p_j} \to M_{\mathfrak p_j} =
\left(\prod\nolimits_{i=1}^t R/\mathfrak p_i\right)_{\mathfrak p_j} =
(R/\mathfrak p_j)_{\mathfrak p_j}
$$
가 전사임을 알 수 있다. 따라서 모든 $j$에 대해
$Q_{\mathfrak p_j} = 0$이다. $\mathfrak m$이
$\mathfrak p_i$들과 다른 $R$의 유일한 소아이디얼이므로
$\text{Supp}(Q) = \{\mathfrak m\}$이라고 결론 내린다.
따라서 $Q$는 유한 길이를 가진다
(Algebra, Lemma \ref{algebra-lemma-support-point}).
$\text{Supp}(Q) = \{\mathfrak m\}$이므로 $x^n$이 $Q$ 위에서
$0$으로 작용하는 $n \gg 0$을 고를 수 있다
(Algebra, Lemma \ref{algebra-lemma-Noetherian-power-ideal-kills-module}).
이제 다음 그림을 생각하자.
$$
\xymatrix{
0 \ar[r] & R \ar[r] \ar[d]^-{x^n} & M
\ar[r] \ar[d]^-{x^n} & Q \ar[r] \ar[d]^-{x^n} & 0 \\
0 \ar[r] & R \ar[r] & M \ar[r] & Q \ar[r] & 0
}
$$
여기서 세로 사상들은 $x^n$을 곱하는 사상이다. $x$가 어느
$\mathfrak p_i$에도 들어 있지 않으므로 이 사상은 $R$과 $M$ 위에서
단사이다. 뱀 보조정리(Algebra, Lemma \ref{algebra-lemma-snake})에 의해
다음 열은 완전하다.
$$
0 \to Q \to R/x^nR \to M/x^nM \to Q \to 0
$$
따라서 충분히 큰 $n$에 대해
$\text{length}_R(R/x^nR) = \text{length}_R(M/x^nM)$이다.
$R_i = R/\mathfrak p_i$로 쓰면
$\text{length}(M/x^nM) =
\sum_{i = 1}^t \text{length}_R(R_i/x^nR_i)$임을 알 수 있다.
Lemma \ref{more-algebra-lemma-polypoly}와 $x$가 $R$과 $R_i$ 위에서 영인자가
아니라는 사실을 적용하면
$$
n \text{length}_R(R/xR) =
\sum\nolimits_{i = 1}^t n \text{length}_{R_i}(R_i/x R_i)
$$
을 얻는다. $\text{length}_{R_i}(R_i/x R_i) \geq 1$이므로
보조정리가 증명되었다.
\end{proof}

\begin{lemma}
\label{more-algebra-lemma-sopexists}
$(R,\mathfrak m)$을 차원 $d > 1$인 Noether 국소환이라 하고,
$f \in \mathfrak m$을 $R$의 어느 극소 소아이디얼에도 들어 있지 않은
원소라 하며, $k\in\mathbf{N}$이라 하자. 그러면
$f, g_1, \ldots, g_{d - 1}$이 매개변수계를 이루는 원소
$g_1, \ldots, g_{d - 1} \in \mathfrak m^k$가 존재한다.
\end{lemma}

\begin{proof}
Algebra, Lemma \ref{algebra-lemma-one-equation}에 의해
$\dim(R/fR) = d - 1$이다. $R/fR$에서 매개변수계
$\overline{g}_1, \ldots, \overline{g}_{d - 1}$을 고르고
(Algebra, Proposition \ref{algebra-proposition-dimension})
$R$에서 올림 $g_1, \ldots, g_{d - 1}$을 취한다.
$f, g_1, \ldots, g_{d - 1}$이 $R$의 매개변수계임은 직접 확인된다.
그러면 $f, g_1^k, \ldots, g_{d - 1}^k$ 역시 매개변수계이므로
증명이 끝난다.
\end{proof}

\begin{lemma}
\label{more-algebra-lemma-syspar}
$(R,\mathfrak m)$을 차원 2인 Noether 국소환이라 하고,
$f \in \mathfrak m$을 $R$의 어느 극소 소아이디얼에도 들어 있지 않은
원소라 하자. 그러면 다음을 만족하는
$g \in \mathfrak m$과 $N \in \mathbf{N}$이 존재한다.
\begin{enumerate}
\item[(a)] $f,g$는 $R$의 매개변수계를 이룬다.
\item[(b)] $h \in \mathfrak m^N$이면 $f + h, g$는 매개변수계이고
$\text{length}_R (R/(f, g)) = \text{length}_R(R/(f + h, g))$이다.
\end{enumerate}
\end{lemma}

\begin{proof}
Lemma \ref{more-algebra-lemma-sopexists}에 의해 $f, g$가 $R$의 매개변수계를 이루는
$g \in \mathfrak m$이 존재한다. 그러면
$\mathfrak m = \sqrt{(f, g)}$이다. 따라서 어떤 $n$에 대해
$\mathfrak m^n \subset (f, g)$이다. Algebra, Lemma
\ref{algebra-lemma-Noetherian-power}를 보라. $N = n + 1$이면 된다고
주장한다. 실제로 $h \in \mathfrak m^N$이라 하자.
$N$의 선택에 의해 $a, b \in \mathfrak m$인
$h = af + bg$로 쓸 수 있다. 따라서
$$
(f + h, g) = (f + af + bg, g) = ((1 + a)f, g) = (f, g)
$$
이다. $1 + a$가 $R$의 단위원이기 때문이다. 이는 길이의 등식과
$f + h, g$가 매개변수계라는 사실을 증명한다.
\end{proof}

\begin{lemma}
\label{more-algebra-lemma-radical-element}
$R$을 차원 $2$인 Noether 정규 국소 정역이라 하자.
$\mathfrak p_1, \ldots, \mathfrak p_r$를 서로 다른 높이 $1$의
소아이디얼들이라 하자. $R/fR$이 축약환이 되도록 하는 영이 아닌 원소
$f \in \mathfrak p_1 \cap \ldots \cap \mathfrak p_r$가 존재한다.
\end{lemma}

\begin{proof}
$f \in \mathfrak p_1 \cap \ldots \cap \mathfrak p_r$를 영이 아닌
원소라 하자. 근기 아이디얼을 생성하는 원소를 얻기 위해 $f$를
조금 변형하겠다. 각 높이 1의 소아이디얼 $\mathfrak p$에서 $R$을
국소화한 $R_\mathfrak p$는 이산 값매김환이다. Algebra, Lemma
\ref{algebra-lemma-characterize-dvr} 또는 Algebra, Lemma
\ref{algebra-lemma-criterion-normal}을 보라. 이에 대응하는
$R_{\mathfrak p}$에서의 $f$의 값매김을
$\text{ord}_\mathfrak p(f)$로 나타낸다.
$\mathfrak q_1, \ldots, \mathfrak q_s$를 $f$를 포함하는 서로 다른
높이 1의 소아이디얼들이라 하자. 각 $j$에 대해
$\text{ord}_{\mathfrak q_j}(f) = m_j \geq 1$로 쓰자.
그런 다음 정수 계수를 가지는 높이 1 소아이디얼들의 형식 선형결합으로
$\text{div}(f) = \sum_{j = 1}^s m_j\mathfrak q_j$를 정의한다.
나중에 사용할 사실로, 각 소아이디얼 $\mathfrak p_i$는
소아이디얼 $\mathfrak q_j$들 가운데 나타남에 유의하라.
$R/fR$이 축약환일 필요충분조건은 $j = 1, \ldots, s$에 대해
$m_j = 1$인 것이다. 실제로 $m_j$가 $1$이면
% P02-E0717
$(R/fR)\mathfrak q_j$는 축약환이고
$R/fR \subset \prod (R/fR)_{\mathfrak q_j}$이다.
% P02-E0718
이는 $\mathfrak q_1, \ldots, \mathfrak q_j$가 $R/fR$의 연관
소아이디얼들이기 때문이다. Algebra, Lemma
\ref{algebra-lemma-zero-at-ass-zero}와
\ref{algebra-lemma-normal-domain-intersection-localizations-height-1}을 보라.

\medskip\noindent
Lemma \ref{more-algebra-lemma-syspar}에서처럼 $g$와 $N$을 골라 고정한다.
영이 아닌 $y \in R$에 대해 $t(y)$를 $y$ 위의 극소 소아이디얼 개수라
하자. $R$은 정규 정역이므로 이 소아이디얼들은 높이 1이고,
$R/yR$의 극소 소아이디얼들과 $1$-대-$1$로 대응한다
(Algebra, Lemma \ref{algebra-lemma-minimal-over-1}과
\ref{algebra-lemma-normal-domain-intersection-localizations-height-1}).
예를 들어 $t(f) = s$는 $\text{div}(f)$에 나타나는
소아이디얼 $\mathfrak q_j$들의 개수이다. $h \in \mathfrak m^N$이라 하자.
Lemma \ref{more-algebra-lemma-minprimespoly}에 의해
\begin{align*}
t(f + h) & \leq \text{length}_{R/(f + h)}(R/(f + h, g)) \\
& = \text{length}_R(R/(f + h, g)) \\
& = \text{length}_R(R/(f, g))
\end{align*}
이다. 첫째 등식은 Algebra, Lemma
\ref{algebra-lemma-length-independent}를 보라. 따라서
$t(f + h)$는 $h \in \mathfrak m^N$과 무관하게 유계임을 알 수 있다.

\medskip\noindent
위에서 증명한 유계성에 의해
$h \in \mathfrak m^N \cap \mathfrak p_1 \cap \ldots \cap \mathfrak p_r$
중 $t(f + h)$가 최대가 되도록 $h$를 고를 수 있다.
$f' = f + h$라 놓자. 주어진
$h' \in \mathfrak m^N \cap \mathfrak p_1 \cap \ldots \cap \mathfrak p_r$
에 대해 $t(f' + h') \leq t(f + h)$임을 알 수 있다.
따라서 $f$를 $f'$으로 바꾸어 모든
$h \in \mathfrak m^N \cap \mathfrak p_1 \cap \ldots \cap \mathfrak p_r$
에 대해 $t(f + h) \leq s$라고 가정해도 된다.

\medskip\noindent
이제 각 $j$에 대해 $\text{ord}_{\mathfrak q_j}(h) \geq 1$이고
$\text{ord}_{\mathfrak q_j}(h) = 1 \Leftrightarrow m_j > 1$인 원소
$h \in \mathfrak m^N$을 찾을 수 있다고 하자.
$h \in \mathfrak m^N \cap \mathfrak p_1 \cap \ldots \cap \mathfrak p_r$
임에 유의하라. 그러면 값매김의 기본 성질에 의해 모든 $j$에 대해
$\text{ord}_{\mathfrak q_j}(f + h) = 1$이다. 따라서
$$
\text{div}(f + h) = \sum\nolimits_{j = 1}^s \mathfrak q_j +
\sum\nolimits_{k = 1}^v e_k \mathfrak r_k
$$
이다. 여기서 $\mathfrak r_1, \ldots, \mathfrak r_v$는 서로 다른
높이 1의 소아이디얼들이고 $e_k \geq 1$이다. 그러나
$s = t(f) \geq t(f + h)$이므로 $v = 0$임을 알 수 있고,
원하는 원소를 찾았다.

\medskip\noindent
이제 위 조건들을 만족하는 $h$를 고르겠다. 소아이디얼 회피
(Algebra, Lemma \ref{algebra-lemma-silly})에 의해
각 $1 \leq j \leq s$에 대해 $j' \not = j$일 때
$a_j \not \in \mathfrak q_{j'}$이고
$a_j \not \in \mathfrak q_j^{(2)}$인
$a_j \in \mathfrak q_j$를 찾을 수 있다. 여기서
$\mathfrak q_j^{(2)} = \{x \in R \mid \text{ord}_{\mathfrak q_j}(x) \geq 2\}$
는 $\mathfrak q_j$의 둘째 심볼릭 거듭제곱이다. 이제
$$
h = \prod\nolimits_{m_j = 1} a_j^2 \times
\prod\nolimits_{m_j > 1} a_j
$$
로 잡는다. 그러면 $h$는 위에서 부과한 값매김 조건들을 분명히 만족한다.
$h \not \in \mathfrak m^N$이면 모든 $j$에 대해
$\mathfrak q_j$에 들어 있지 않은 $\mathfrak m^N$의 원소를 곱한다.
\end{proof}

\begin{lemma}
\label{more-algebra-lemma-divides-radical}
$(A, \mathfrak m, \kappa)$를 차원 $2$인 Noether 정규 국소 정역이라 하자.
$a \in \mathfrak m$이 영이 아니면, $A/cA$가 축약환이고 어떤 $n$에 대해
$a$가 $c^n$을 나누도록 하는 원소 $c \in A$가 존재한다.
\end{lemma}

\begin{proof}
Lemma \ref{more-algebra-lemma-radical-element}의 증명에서와 같은 표기법으로
$\text{div}(a) = \sum_{i = 1, \ldots, r} n_i \mathfrak p_i$라 쓰자.
Lemma \ref{more-algebra-lemma-radical-element}에 의해 $A/cA$가 축약환인
$c \in \mathfrak p_1 \cap \ldots \cap \mathfrak p_r$를 고른다.
$n \geq \max(n_i)$이면
$-\text{div}(a) + \text{div}(c^n)$이 유효 약수이다
(모든 계수가 음이 아니다). 따라서 Algebra, Lemma
\ref{algebra-lemma-normal-domain-intersection-localizations-height-1}에 의해
$c^n/a \in A$이다.
\end{proof}

\noindent
이 절의 나머지에서는 차원 $> 2$일 때의 결과를 증명한다.

\begin{lemma}
\label{more-algebra-lemma-multiplicity}
$(R, \mathfrak m)$을 차원 $d$인 Noether 국소환이라 하고,
$g_1, \ldots, g_d$를 매개변수계라 하며,
$I = (g_1, \ldots, g_d)$라 하자. 수치 다항식
$n \mapsto \text{length}_R(R/I^{n+1})$의 최고차항 계수가
$e_I/d!$이면 $e_I \leq \text{length}_R(R/I)$이다.
\end{lemma}

\begin{proof}
이 함수가 수치 다항식임은 Algebra, Proposition
\ref{algebra-proposition-hilbert-function-polynomial}에서 따른다.
Algebra, Proposition \ref{algebra-proposition-dimension}에 의해
그 차수는 $d$이다. $d = 0$이면 결과는 자명하다.
$d = 1$이면 결과는 Lemma \ref{more-algebra-lemma-polypoly}이다.
일반적인 경우를 증명하기 위해
$$
\bigoplus\nolimits_{i_1, \ldots, i_d \geq 0,\ \sum i_j = n} R/I
\longrightarrow
I^n/I^{n + 1}
$$
라는 전사가 있음에 유의하라. 이 전사는
$i_1, \ldots, i_d$에 대응하는 기저원을 $I^n/I^{n + 1}$에서
$g_1^{i_1} \ldots g_d^{i_d}$의 유로 보낸다. 따라서
$$
\text{length}_R(R/I^{n + 1}) - \text{length}_R(R/I^n)
\leq  \text{length}_R(R/I) {n + d - 1 \choose d - 1}
$$
이다. $d \geq 2$이므로 왼쪽의 수치 다항식은 차수가 $d - 1$이고
최고차항 계수가 $e_I / (d - 1)!$이다. 오른쪽 다항식은 차수가
$d - 1$이고 최고차항 계수가
$\text{length}_R(R/I)/ (d - 1)!$이다. 이는 보조정리를 증명한다.
\end{proof}

\begin{lemma}
\label{more-algebra-lemma-minprimespolyhigher}
$(R, \mathfrak m)$을 차원 $d$인 Noether 국소환이라 하고,
$t$를 차원 $d$인 $R$의 극소 소아이디얼 개수라 하며,
$(g_1,\ldots,g_d)$를 매개변수계라 하자. 그러면
% P02-E0719
$t \leq \text{length}_R(R/(g_1,\ldots,g_n))$이다.
\end{lemma}

\begin{proof}
$d = 0$이면 보조정리는 자명하다. $d = 1$이면 보조정리는
Lemma \ref{more-algebra-lemma-minprimespoly}이다. 따라서 $d > 1$이라고 가정해도 된다.
$\mathfrak p_1, \ldots, \mathfrak p_s$를 $R$의 극소 소아이디얼들이라
하되, 처음 $t$개는 차원이 $d$라고 하자. 그리고
% P02-E0719
$I = (g_1, \ldots, g_n)$라 놓자. Lemma
\ref{more-algebra-lemma-minprimespoly}의 증명과 정확히 같은 방식으로 논증하여
$R$이 축약환이라고 가정해도 된다.

\medskip\noindent
% P02-E0720
$R$이 극소 소아이디얼 $\mathfrak p_1, \ldots, \mathfrak p_t$를
가지는 축약환이라고 하자. 이는 다음 완전열이 있음을 뜻한다.
$$
0 \to R \to
\prod\nolimits_{i = 1}^t R/\mathfrak p_i \to Q \to 0
$$
여기서 $Q$는 첫째 사상의 여핵이다.
$M = \prod_{i = 1}^t R/\mathfrak p_i$라 쓰자.
$\mathfrak p_j$에서 국소화하면
$$
R_{\mathfrak p_j} \to M_{\mathfrak p_j} =
\left(\prod\nolimits_{i=1}^t R/\mathfrak p_i\right)_{\mathfrak p_j} =
(R/\mathfrak p_j)_{\mathfrak p_j}
$$
가 전사임을 알 수 있다. 따라서 모든 $j$에 대해
$Q_{\mathfrak p_j} = 0$이다. 그러므로 높이 $0$인 $R$의
어느 소아이디얼도 $Q$의 지지에 들어 있지 않다. 따라서 수치 다항식
$n \mapsto \text{length}_R(Q/I^nQ)$의 차수는
$\dim(\text{Supp}(Q)) < d$와 같다. Algebra, Lemma
\ref{algebra-lemma-support-dimension-d}를 보라.
Algebra, Lemma \ref{algebra-lemma-hilbert-ses-chi}에 의해
($R$이 유한 길이를 가지지 않으므로 적용할 수 있다) 다항식
$$
n \longmapsto
\text{length}_R(M/I^nM) - \text{length}_R(R/I^n) - \text{length}_R(Q/I^nQ)
$$
의 차수는 $< d$이다. $M = \prod R/\mathfrak p_i$이고,
% P02-E0721
$n \to \text{length}_R(R/\mathfrak p_i + I^n)$는
$i = 1, \ldots, t$에 대해 차수가 정확히(!) $d$인 수치 다항식이므로
(Algebra, Lemma \ref{algebra-lemma-support-dimension-d})
$n \mapsto \text{length}_R(M/I^nM)$의 최고차항 계수는 적어도
$t/d!$임을 알 수 있다. 따라서 Lemma \ref{more-algebra-lemma-multiplicity}에 의해
결론을 얻는다.
\end{proof}

\begin{lemma}
\label{more-algebra-lemma-sysparhigher}
$(R, \mathfrak m)$을 차원 $d$인 Noether 국소환이라 하고,
$f \in \mathfrak m$을 $R$의 어느 극소 소아이디얼에도 들어 있지 않은
원소라 하자. 그러면 다음을 만족하는
$g_1, \ldots, g_{d - 1} \in \mathfrak m$과
$N \in \mathbf{N}$이 존재한다.
\begin{enumerate}
\item $f, g_1, \ldots, g_{d - 1}$은 $R$의 매개변수계를 이룬다.
\item $h \in \mathfrak m^N$이면
$f + h, g_1, \ldots, g_{d - 1}$은 매개변수계를 이루며
$\text{length}_R R/(f, g_1, \ldots, g_{d-1}) =
\text{length}_R R/(f + h, g_1, \ldots, g_{d-1})$이다.
\end{enumerate}
\end{lemma}

\begin{proof}
Lemma \ref{more-algebra-lemma-sopexists}에 의해
$f, g_1, \ldots, g_{d - 1}$이 $R$의 매개변수계를 이루는
$g_1, \ldots, g_{d - 1} \in \mathfrak m$이 존재한다.
그러면 $\mathfrak m = \sqrt{(f, g_1, \ldots, g_{d - 1})}$이다.
따라서 어떤 $n$에 대해
% P02-E0722
$\mathfrak m^n \subset (f, g)$이다. Algebra, Lemma
\ref{algebra-lemma-Noetherian-power}를 보라. $N = n + 1$이면 된다고
주장한다. 실제로 $h \in \mathfrak m^N$이라 하자. $N$의 선택에 의해
$a, b_i \in \mathfrak m$인 $h = af + \sum b_ig_i$로 쓸 수 있다.
따라서
\begin{align*}
(f + h, g_1, \ldots, g_{d - 1})
& =
(f + af + \sum b_ig_i, g_1, \ldots, g_{d - 1}) \\
& =
((1 + a)f, g_1, \ldots, g_{d - 1}) \\
& =
(f, g_1, \ldots, g_{d - 1})
\end{align*}
이다. $1 + a$가 $R$의 단위원이기 때문이다. 이는 길이의 등식과
$f + h, g_1, \ldots, g_{d - 1}$이 매개변수계라는 사실을 증명한다.
\end{proof}

\begin{proposition}
\label{more-algebra-proposition-propdimd}
\begin{reference}
\cite[Lemma 3.14]{Artin-Lipman}에는 환이 카테너리라는 가정 없이
이 결과가 제시되어 있다.
\end{reference}
$R$을 카테너리 Noether 정규 국소 정역이라 하자.
% P02-E1347
$J \subset R$을 근기 아이디얼이라 하자. 그러면 $R/fR$이 축약환이
되도록 하는 영이 아닌 원소 $f \in J$가 존재한다.
\end{proposition}

\begin{proof}
증명은 Lemma \ref{more-algebra-lemma-radical-element}의 증명과 같다.
다만 Lemma \ref{more-algebra-lemma-minprimespolyhigher}를
Lemma \ref{more-algebra-lemma-minprimespoly} 대신 사용하고,
Lemma \ref{more-algebra-lemma-sysparhigher}를 Lemma \ref{more-algebra-lemma-syspar} 대신 사용한다.
$R$은 카테너리 정역이므로 $R$의 높이 1인 모든 소아이디얼은
차원이 $d - 1$이다. 따라서 $R/(f + h)$의 스펙트럼은 등차원이다.
그러므로 Lemma \ref{more-algebra-lemma-minprimespolyhigher}를 사용할 수 있다.
독자의 편의를 위해 자세한 내용을 적겠다.

\medskip\noindent
$f \in J$를 영이 아닌 원소라 하자. 근기 아이디얼을 생성하는 원소를
얻기 위해 $f$를 조금 변형하겠다. 각 높이 1의 소아이디얼
$\mathfrak p$에서 $R$을 국소화한 $R_\mathfrak p$는
이산 값매김환이다. Algebra, Lemma \ref{algebra-lemma-characterize-dvr}
또는 Algebra, Lemma \ref{algebra-lemma-criterion-normal}을 보라.
이에 대응하는 $R_{\mathfrak p}$에서의 $f$의 값매김을
$\text{ord}_\mathfrak p(f)$로 나타낸다.
$\mathfrak q_1, \ldots, \mathfrak q_s$를 $f$를 포함하는 서로 다른
높이 1의 소아이디얼들이라 하자. 각 $j$에 대해
$\text{ord}_{\mathfrak q_j}(f) = m_j \geq 1$로 쓰자. 그리고 정수 계수를
가지는 높이 1 소아이디얼들의 형식 선형결합으로
$\text{div}(f) = \sum_{j = 1}^s m_j\mathfrak q_j$를 정의한다.
$R/fR$이 축약환일 필요충분조건은 $j = 1, \ldots, s$에 대해
$m_j = 1$인 것이다. 실제로 $m_j$가 $1$이면
% P02-E0717
$(R/fR)\mathfrak q_j$는 축약환이고
$R/fR \subset \prod (R/fR)_{\mathfrak q_j}$이다.
% P02-E0718
이는 $\mathfrak q_1, \ldots, \mathfrak q_j$가 $R/fR$의 연관
소아이디얼들이기 때문이다. Algebra, Lemma
\ref{algebra-lemma-zero-at-ass-zero}와
\ref{algebra-lemma-normal-domain-intersection-localizations-height-1}을 보라.

\medskip\noindent
% P02-E0723
Lemma \ref{more-algebra-lemma-sysparhigher}에서처럼 $g_2, \ldots, g_{d - 1}$과
$N$을 골라 고정한다. 영이 아닌 $y \in R$에 대해 $t(y)$를
$y$ 위의 극소 소아이디얼 개수라 하자. $R$은 정규 정역이므로
이 소아이디얼들은 높이 1이고 $R/yR$의 극소 소아이디얼들과
$1$-대-$1$로 대응한다
(Algebra, Lemma \ref{algebra-lemma-minimal-over-1}과
\ref{algebra-lemma-normal-domain-intersection-localizations-height-1}).
예를 들어 $t(f) = s$는 $\text{div}(f)$에 나타나는
소아이디얼 $\mathfrak q_j$들의 개수이다. $h \in \mathfrak m^N$이라 하자.
$R$이 카테너리이므로 $R$의 높이 1인 각 소아이디얼 $\mathfrak p$에 대해
% P02-E0724
$\dim(R/\mathfrak p) = d$이다. 따라서 Lemma
\ref{more-algebra-lemma-minprimespolyhigher}에 의해
\begin{align*}
t(f + h) & \leq \text{length}_{R/(f + h)}(R/(f + h, g_1, \ldots, g_{d - 1})) \\
& = \text{length}_R(R/(f + h, g_1, \ldots, g_{d - 1})) \\
& = \text{length}_R(R/(f, g_1, \ldots, g_{d - 1}))
\end{align*}
이다. 첫째 등식은 Algebra, Lemma
\ref{algebra-lemma-length-independent}를 보라. 따라서
$t(f + h)$는 $h \in \mathfrak m^N$과 무관하게 유계임을 알 수 있다.

\medskip\noindent
위에서 증명한 유계성에 의해 $h \in \mathfrak m^N \cap J$ 중
$t(f + h)$가 최대가 되도록 $h$를 고를 수 있다.
$f' = f + h$라 놓자. 주어진 $h' \in \mathfrak m^N \cap J$에 대해
$t(f' + h') \leq t(f + h)$임을 알 수 있다. 따라서 $f$를 $f'$으로
바꾸어 모든 $h \in \mathfrak m^N \cap J$에 대해
$t(f + h) \leq s$라고 가정해도 된다.

\medskip\noindent
이제 각 $j$에 대해 $\text{ord}_{\mathfrak q_j}(h) \geq 1$이고
$\text{ord}_{\mathfrak q_j}(h) = 1 \Leftrightarrow m_j > 1$인 원소
$h \in \mathfrak m^N \cap J$를 찾을 수 있다고 하자.
그러면 값매김의 기본 성질에 의해 모든 $j$에 대해
$\text{ord}_{\mathfrak q_j}(f + h) = 1$이다. 따라서
$$
\text{div}(f + h) = \sum\nolimits_{j = 1}^s \mathfrak q_j +
\sum\nolimits_{k = 1}^v e_k \mathfrak r_k
$$
이다. 여기서 $\mathfrak r_1, \ldots, \mathfrak r_v$는 서로 다른
높이 1의 소아이디얼들이고 $e_k \geq 1$이다. 그러나
$s = t(f) \geq t(f + h)$이므로 $v = 0$임을 알 수 있고,
원하는 원소를 찾았다.

\medskip\noindent
이제 위 조건들을 만족하는 $h$를 고르겠다. 소아이디얼 회피
(Algebra, Lemma \ref{algebra-lemma-silly})에 의해
각 $1 \leq j \leq s$에 대해 $j' \not = j$이면
$a_j \not \in \mathfrak q_{j'}$가 되도록 하는
$a_j \in \mathfrak q_j \cap J$를 찾을 수 있다. 다음으로
% P02-E0725
$b_j \in J \cap \mathfrak q_1 \cap \ldots \cap q_s$이면서
$b_j \not \in \mathfrak q_j^{(2)}$인 원소를 고를 수 있다. 여기서
$\mathfrak q_j^{(2)} = \{x \in R \mid \text{ord}_{\mathfrak q_j}(x) \geq 2\}$
는 $\mathfrak q_j$의 둘째 심볼릭 거듭제곱이다. 소아이디얼 회피를
적용할 수 있는 이유는 아이디얼
% P02-E0725
$J' = J \cap \mathfrak q_1 \cap \ldots \cap q_s$이 근기 아이디얼이고,
따라서 $R/J'$이 축약환이며, 따라서 $(R/J')_{\mathfrak q_j}$가
축약환이고, 따라서 $J'$이 $\text{ord}_{\mathfrak q_j}(x) = 1$인
원소 $x$를 포함하며, 따라서
$J' \not \subset \mathfrak q_j^{(2)}$이기 때문이다. 그러면 원소
$$
c = \sum\nolimits_{j = 1, \ldots, s}
b_j \times \prod\nolimits_{j' \not = j} a_{j'}
$$
는 $J$의 원소이고, 값매김의 기본 성질에 의해 모든
$j = 1, \ldots, s$에 대해
$\text{ord}_{\mathfrak q_j}(c) = 1$이다. 마지막으로
$$
h = c \times \prod\nolimits_{m_j = 1} a_j \times y
$$
라 놓는다. 여기서 $y \in \mathfrak m^N$은 모든 $j$에 대해
$\mathfrak q_j$에 들어 있지 않은 원소이다.
\end{proof}

\section{유도 범주의 가역 대상}
\label{more-algebra-section-invertible-D-or-R}

\noindent
환의 유도 범주에 있는 가역 대상들을 특징짓는다.

\begin{lemma}
\label{more-algebra-lemma-symmetric-monoidal-derived}
$R$을 환이라 하자. $R$의 유도 범주 $D(R)$은 유도 텐서곱으로
주어지는 텐서곱과 Section \ref{more-algebra-section-sign-rules}에서와 같은
결합 및 교환 제약을 가진 대칭 모노이드 범주이다.
\end{lemma}

\begin{proof}
생략한다. 힌트: 결합 제약은 Lemma \ref{more-algebra-lemma-triple-tensor-product}의
동형이고, 교환 제약은 Lemma
\ref{more-algebra-lemma-flip-douoble-tensor-product}의 동형이다.
이렇게 말하고 나면 여러 그림의 가환성은 $R$-가군 복합체들의 범주에
대한 대응하는 결과에서 따른다. Section
\ref{more-algebra-section-symmetric-monoidal}을 보라.
\end{proof}

\noindent
따라서 $D(R)$의 대상이 (왼쪽) 쌍대를 가진다는 것과 가역이라는 것의
의미를 안다. 이것이 무엇을 뜻하는지 규명하기 전에 간단한 보조정리가
필요하다.

\begin{lemma}
\label{more-algebra-lemma-complex-bounded-above-free-colim-bounded-finite-free}
$R$을 환이라 하자. $F^\bullet$을 자유 $R$-가군들로 이루어진
위로 유계인 복합체라 하자. $n_i \in \mathbf{Z}$이고
$f_i \in F^{n_i}$인 쌍들 $(n_i, f_i)$, $i = 1, \ldots, N$이
주어졌다고 하자. 모든 $f_i$를 포함하고 유계이며 유한 자유
$R$-가군들로 이루어진 부분복합체
$G^\bullet \subset F^\bullet$이 존재한다.
\end{lemma}

\begin{proof}
$a = \min(n_i; i = 1, \ldots, N)$에 대한 내림귀납법을 사용한다.
$n \geq a$에 대해 $F^n = 0$이면 영 복합체를
$G^\bullet$로 택하여 결과를 얻는다. 일반적으로 번호를 다시 매겨
% P02-E0726
$n_1 = \ldots = n_r = a$이고 $i > r$에 대해 $n_i > a$인
$1 \leq r \leq N$이 존재한다고 가정해도 된다.
$b_j, j \in J$를 $F^a$의 기저라 하자.
$i = 1, \ldots, r$에 대해
$f_i \in \bigoplus_{j \in J'} Rb_j$가 되는 유한 부분집합
$J' \subset J$를 고를 수 있다. $c_k, k \in K$를 $F^{a + 1}$의
기저라 하자. $j \in J'$에 대해
$\text{d}_F^a(b_j) \in \bigoplus_{k \in K'} Rc_k$가 되는 유한 부분집합
$K' \subset K$를 고를 수 있다. 이제 귀납가정을 적용하여
$k \in K'$에 대한 $c_k \in F^{a + 1}$과 $i > r$에 대한
$f_i \in F^{n_i}$를 포함하는 부분복합체
$H^\bullet \subset F^\bullet$을 찾을 수 있다.
$G^\bullet$은 차수 $> a$에서 $H^\bullet$과 같고,
차수 $a$에서 $\bigoplus_{j \in J'} Rb_j$와 같게 잡는다.
\end{proof}

\begin{lemma}
\label{more-algebra-lemma-have-dual-derived}
$R$을 환이라 하자. $M$을 $D(R)$의 대상이라 하자. 다음은 서로 동치이다.
\begin{enumerate}
\item $M$은 Categories, Definition \ref{categories-definition-dual}에서와
같이 $D(R)$에서 왼쪽 쌍대를 가진다.
\item $M$은 $D(R)$의 퍼펙트 대상이다.
\end{enumerate}
더욱이 이 경우 $M$의 왼쪽 쌍대는 Lemma
\ref{more-algebra-lemma-dual-perfect-complex}의 대상 $M^\vee$이다.
\end{lemma}

\begin{proof}
$M$이 퍼펙트이면 유한 사영 $R$-가군들로 이루어진 유계 복합체
$M^\bullet$로 $M$을 나타낼 수 있다. 이 경우 Lemma
\ref{more-algebra-lemma-left-dual-complex}에 의해 $M^\bullet$은 복합체들의 범주에서
왼쪽 쌍대를 가지며, 이는 당연히 $D(R)$에서도 왼쪽 쌍대이다.

\medskip\noindent
(1)을 가정하자. Categories, Definition
\ref{categories-definition-dual}에서와 같은 왼쪽 쌍대를
% P02-E0727
$N$, $\eta : R \to M \otimes_R^\mathbf{L} N$,
$\epsilon : M \otimes_R^\mathbf{L} N \to R$라 하자.
$M$을 나타내는 복합체 $M^\bullet$을 고른다.
% P02-E0728
$N$을 나타내며 항들이 평탄한 K-평탄 복합체 $N^\bullet$을 고른다.
Lemma \ref{more-algebra-lemma-K-flat-resolution}을 보라. 그러면 $\eta$는 복합체 사상
$$
\eta : R \longrightarrow \text{Tot}(M^\bullet \otimes_R N^\bullet)
$$
으로 주어진다. $1$의 상을 유한합
$$
\eta(1) = \sum\nolimits_n \sum\nolimits_i m_{n, i} \otimes n_{-n, i}
$$
으로 쓸 수 있다. 여기서 $m_{n, i} \in M^n$이고
$n_{-n, i} \in N^{-n}$이다. $K^\bullet \subset M^\bullet$을 모든 원소
$m_{n, i}$와 $\text{d}(m_{n, i})$로 생성되는 부분복합체라 하자.
$N^\bullet$의 선택에 의해
$\text{Tot}(K^\bullet \otimes_R N^\bullet) \subset
\text{Tot}(M^\bullet \otimes_R N^\bullet)$이고,
위 선택에 의해 $\eta(1)$은 이 부분복합체 안에 있다.
$K$를 $K^\bullet$로 나타내어지는 $D(R)$의 대상이라 하자.
그러면 $\eta$가 사상
$\tilde \eta : R \longrightarrow K \otimes_R^\mathbf{L} N$을
거쳐 인수분해됨을 알 수 있다. 다음이 성립하므로
$(1 \otimes \epsilon) \circ (\eta \otimes 1) = \text{id}_M$
가환 그림
$$
\xymatrix{
M \ar[rr]_-{\eta \otimes 1} \ar[rrd]_{\tilde \eta \otimes 1} & &
M \otimes_R^\mathbf{L} N \otimes_R^\mathbf{L} M
\ar[r]_-{1 \otimes \epsilon} &
M \\
& &
K \otimes_R^\mathbf{L} N \otimes_R^\mathbf{L} M
\ar[u] \ar[r]^-{1 \otimes \epsilon} &
K \ar[u]
}
$$
에 의해 $M$ 위 항등사상은 $K$를 거쳐 인수분해된다고 결론 내린다.
$K$가 위로 유계이므로 $M \in D^-(R)$이다. 따라서 예를 들어
Derived Categories, Lemma
\ref{derived-lemma-subcategory-left-resolution}에 의해 자유
$R$-가군들로 이루어진 위로 유계인 복합체 $M^\bullet$로 $M$을
나타낼 수 있다. 이전처럼
$\eta(1) = \sum\nolimits_n \sum\nolimits_i m_{n, i} \otimes n_{-n, i}$로
쓰자. Lemma
\ref{more-algebra-lemma-complex-bounded-above-free-colim-bounded-finite-free}에 의해
모든 원소 $m_{n, i}$를 포함하고 유계이며 유한 자유
$R$-가군들로 이루어진 부분복합체
$K^\bullet \subset M^\bullet$을 찾을 수 있다.
위와 같이 $M$ 위 항등사상이 $K$를 거쳐 인수분해됨을 알 수 있다.
$K$가 퍼펙트이므로 $M$ 역시 퍼펙트이다. Lemma
\ref{more-algebra-lemma-summands-perfect}를 보라.
\end{proof}

\begin{lemma}
\label{more-algebra-lemma-invertible-derived}
$R$을 환이라 하자. $M$을 $D(R)$의 대상이라 하자. 다음은 서로 동치이다.
\begin{enumerate}
\item $M$은 $D(R)$에서 가역이다. Categories, Definition
\ref{categories-definition-invertible}을 보라.
\item 모든 소아이디얼 $\mathfrak p \subset R$에 대해
$f \not \in \mathfrak p$이고 어떤 $n \in \mathbf{Z}$에 대해
$M_f \cong R_f[-n]$이 되는 $f \in R$가 존재한다.
\end{enumerate}
더욱이 이 경우
\begin{enumerate}
\item[(a)] $M$은 $D(R)$의 퍼펙트 대상이다.
\item[(b)] $D(R)$에서 $M = \bigoplus H^n(M)[-n]$이다.
\item[(c)] 각 $H^n(M)$은 유한 사영 $R$-가군이다.
\item[(d)] $H^n(M)$이 가역 $R_n$-가군에 대응하도록
$R = \prod_{a \leq n \leq b} R_n$로 쓸 수 있다.
\end{enumerate}
\end{lemma}

\begin{proof}
(2)를 가정하자. 대상 $R\Hom_R(M, R)$과 합성사상
$$
R\Hom(M, R) \otimes_R^\mathbf{L} M \to R
$$
을 생각하자. 국소적으로 확인하면 이것이 동형임을 알 수 있다.
자세한 내용은 생략한다. $D(R)$은 대칭 모노이드이므로 $M$은 가역이다.

\medskip\noindent
(1)을 가정하자. 모노이드 범주의 가역 대상은 왼쪽 쌍대, 즉 그 역을
가짐에 유의하라. 따라서 Lemma \ref{more-algebra-lemma-have-dual-derived}에 의해
$M$은 퍼펙트이다. 잉여체가 $\kappa$인 소아이디얼
$\mathfrak p \subset R$을 생각하자. 그러면
$M \otimes_R^\mathbf{L} \kappa$는 $D(\kappa)$의 가역 대상임을 알 수 있다.
이는 분명 $\dim H^i(M \otimes_R^\mathbf{L} \kappa)$가 정확히 하나의
$i$에서만 영이 아니고, 그 경우 $1$과 같음을 함의한다.
Lemma \ref{more-algebra-lemma-lift-perfect-from-residue-field}에 의해 (2)를 얻는다.

\medskip\noindent
위 증명에서 (a)가 성립함을 보았다. $U_n \subset \Spec(R)$을
$M_f \cong R_f[-n]$인 꼴 $D(f)$의 열린집합들의 합집합이라 하자.
$n \not = n'$이면 분명 $U_n \cap U_{n'} = \emptyset$이다.
$M$의 Tor 진폭이 $[a, b]$ 안에 있으면
$n \not \in [a, b]$일 때 $U_n = \emptyset$이다.
따라서 (d)에서와 같이 $U_n$이 $\Spec(R_n)$에 대응하는 곱분해
$R = \prod_{a \leq n \leq b} R_n$을 가진다. Algebra, Lemma
\ref{algebra-lemma-disjoint-implies-product}를 보라.
% P02-E0729
$D(R) = \prod_{a \leq n \leq b} D(R_n)$이고 가군들의 범주에 대해서도
마찬가지이므로 (b), (c), (d)가 즉시 따른다.
\end{proof}

\section{자유 가군을 직합인자로 떼어내기}
\label{more-algebra-section-serre-splitting}

\noindent
이 절의 논증들은 Serre에 의한 것이다. \cite{Serre-projective}를 보라.

\begin{situation}
\label{more-algebra-situation-splitting}
여기서 $R$은 환이고 $M$은 유한 표시 $R$-가군이다.
$\Omega \subset \Spec(R)$를 유도된 위상을 가진 닫힌점들의 집합이라 하자.
$x \in \Omega$에 대해 $M(x) = M/xM$를 $x$에서의 $M$의 올이라 하자.
이는 $x$에서의 잉여체 $\kappa(x)$ 위 유한 차원 벡터공간이다.
주어진 $s \in M$에 대해 $s(x)$로 $M(x)$ 안에서의 $s$의 상을 나타낸다.
\end{situation}

\begin{lemma}
\label{more-algebra-lemma-which-elements-split}
Situation \ref{more-algebra-situation-splitting}에서 $x \in \Omega$라 하자.
% P02-E0730
다음 성질을 가진 $\kappa(x)$-벡터공간들의 표준 짧은 완전열
$$
0 \to B(x) \to M(x) \to V(x) \to 0
$$
이 존재한다. $s_1, \ldots, s_r \in M$에 대해 다음은 서로 동치이다.
\begin{enumerate}
\item 어떤 $f \in R$, $f \not \in x$가 존재하여 $f$를 가역화한 뒤
사상 $s_1, \ldots, s_r : R^{\oplus r} \to M$이 직합인자의 포함이 된다.
\item $s_1(x), \ldots, s_r(x)$는 $V(x)$의 선형독립인 원소들로 간다.
\end{enumerate}
\end{lemma}

\begin{proof}
$B(x) \subset M(x)$를 사상
$$
\Hom_R(M, R) \to \Hom_{\kappa(x)}(M(x), \kappa(x))
$$
의 상의 직교공간으로 정의하고 $V(x) = M(x)/B(x)$라 놓는다.
그러면 임의의 $R$-선형사상 $\varphi : M \to R$은 사상
$\overline{\varphi} : V(x) \to \kappa(x)$를 유도하고, 반대로 임의의
$\kappa(x)$-선형사상 $\lambda : V(x) \to \kappa(x)$는 어떤
$\varphi$에 대한 $\overline{\varphi}$와 같다.
$s_1, \ldots, s_r \in M$이라 하자.

\medskip\noindent
$s_1, \ldots, s_r$가 $V(x)$의 선형독립인 원소들로 간다고 하자.
$\varphi_i(s_j)$가 $\kappa(x)$에서
$\delta_{ij}$\footnote{Kronecker 델타.}로 가는
$\varphi_1, \ldots, \varphi_r \in \Hom_R(M, R)$을 찾을 수 있다.
따라서 합성
$$
R^{\oplus r}
\xrightarrow{s_1, \ldots, s_r}
M
\xrightarrow{\varphi_1, \ldots, \varphi_r}
R^{\oplus r}
$$
의 행렬은 $\kappa(x)$에서 $1$로 가는 행렬식 $f \in R$를 가진다.
% P02-E0731
그러므로 $f$를 가역화한 뒤
$s_1, \ldots, s_r : R^{\oplus r} \to M$이 직합인자의 포함이 됨은
분명하다.

\medskip\noindent
반대로 어떤 $f \in R$, $f \not \in x$가 존재하여 $f$를 가역화한 뒤
$s_1, \ldots, s_r : R^{\oplus r} \to M$이 직합인자의 포함이 된다고 하자.
그러면 $\varphi_i(s_j) = \delta_{ij} \in R_f$인
$R_f$-선형사상 $\varphi_i : M_f \to R_f$를 찾을 수 있다.
Algebra, Lemma \ref{algebra-lemma-hom-from-finitely-presented}에 의해
$\Hom_R(M, R)_f = \Hom_{R_f}(M_f, R_f)$이므로, 어떤 $n \geq 0$과
$\varphi'_i \in \Hom_R(M, R)$을 찾아
$\varphi'_i(s_j) = f^n\delta_{ij} \in R$이 되게 할 수 있다.
따라서
% P02-E0732
$\overline{\varphi}'_i(s_j) = f^n\delta_{ij}$이므로
$s_1, \ldots, s_r$는 $V(x)$의 선형독립인 원소들로 간다.
\end{proof}

\noindent
Situation \ref{more-algebra-situation-splitting}에서
$s_1, \ldots, s_r \in M$이 주어졌을 때
$Z(s_1, \ldots, s_r) \subset \Omega$로
$s_1(x), \ldots, s_r(x)$가 $V(x)$의 선형종속인 원소들로 가는
$x \in \Omega$들의 집합을 나타낸다. 보조정리에 의해 이는
$\Omega$의 닫힌 부분집합이다.

\begin{lemma}
\label{more-algebra-lemma-choose-values}
Situation \ref{more-algebra-situation-splitting}에서
$x_1, \ldots, x_n \in \Omega$를 서로 다른 원소들이라 하자.
$v_i \in V(x_i)$라 하자. 그러면 $i = 1, \ldots, n$에 대해
$s(x_i)$가 $v_i$로 가는 $s \in M$이 존재한다.
\end{lemma}

\begin{proof}
$x_i$가 $R$의 극대 아이디얼이므로 Algebra, Lemma
\ref{algebra-lemma-chinese-remainder}를 사용하여
$M(x_1) \oplus \ldots \oplus M(x_n)$이 $M$의 몫임을 알 수 있다.
\end{proof}

\begin{proposition}
\label{more-algebra-proposition-splitting}
\begin{reference}
\cite[Theorem 2]{Serre-projective}
\end{reference}
Situation \ref{more-algebra-situation-splitting}에서 $\Omega$가 Noether 위상공간이라
가정하자. $s_1, \ldots, s_h \in M$이라 하자.
$Z(s_1, \ldots, s_h) \subset F \subset \Omega$가 닫혔다고 하자.
$x_1, \ldots, x_n \in F$를 서로 다른 원소들이라 하고,
$v_i  \in V(x_i)$라 하자. 다음을 만족하는 정수 $k \geq 0$을 잡자.
$$
(*)\quad h + k \leq \dim_{\kappa(x)} V(x)\text{ for all }x \in \Omega
$$
그러면 다음을 만족하는 $s \in M$과 닫힌
$F' \subset \Omega$가 존재한다.
\begin{enumerate}
\item[(a)] $s(x_i)$는 $v_i$로 간다.
\item[(b)] $Z(s_1, \ldots, s_h, s) \subset F \cup F'$이다.
\item[(c)] $F'$의 모든 기약성분은 $\Omega$에서 여차원
$\geq k$를 가진다.
\end{enumerate}
\end{proposition}

\begin{proof}
여차원은 Topology, Section \ref{topology-section-catenary-spaces}에서
정의되었고, Topology, Section \ref{topology-section-noetherian}에
포함된 Noether 위상공간에 관한 몇 가지 결과를 사용할 것임에 유의하라.

\medskip\noindent
$k$에 대한 귀납법으로 증명한다. $k = 0$이면
Lemma \ref{more-algebra-lemma-choose-values}에서와 같이 $s \in M$을 고르고
$F' = \Omega$로 고른다.

\medskip\noindent
$k > 0$이라 하자. 귀납가정에 의해
$s_1, \ldots, s_h$, $F$, $x_1, \ldots, x_n$,
$v_1, \ldots, v_n$, $k - 1$에 대해 (a), (b), (c)를 만족하는
$u \in M$과 닫힌 $G \subset \Omega$를 고를 수 있다.

\medskip\noindent
$G = G_1 \cup \ldots \cup G_m$를 $G$의 기약성분 분해라 하자.
$G_j \subset F$이면 목록에서 이를 제거할 수 있다. 따라서
$j = 1, \ldots, m$에 대해 $G_j$가 $F$에 포함되지 않는다고 가정해도 된다.
각 $j = 1, \ldots, m$에 대해 $y_j \not \in F$이고
$j' \not = j$일 때 $y_j \not \in G_{j'}$인 $y_j \in G_j$를 고른다.
$G$의 기약성분들 사이에 포함관계가 없으므로 이것이 가능하다.
$s_1(y_j), \ldots, s_h(y_j)$의 상들이 생성하는 공간에 들어 있지 않은
$w_j \in V(y_j)$를 고른다. $h + k \leq \dim V(y_j)$이고
$k > 0$이므로 이것이 가능하다.

\medskip\noindent
$h + 1$개의 단면 $s_1, \ldots, s_h, u$, 닫힌집합 $F \cup G$,
점들 $x_1, \ldots, x_n, y_1, \ldots, y_m \in F \cup G$,
원소들 $0 \in V(x_i)$와 $w_j \in V(y_j)$, 정수 $k - 1$에
귀납가정을 적용한다. $h$를 $1$만큼 늘리고 $k$를 $1$만큼 줄였으므로
명제의 가정 $(*)$는 여전히 성립함에 유의하라. 이로부터
$s_1, \ldots, s_h, u$, $F \cup G$,
$x_1, \ldots, x_n, y_1, \ldots, y_m$,
$0, \ldots, 0, w_1, \ldots, w_m$, $k - 1$에 대해
(a), (b), (c)를 만족하는 $t \in M$과 닫힌
$H \subset \Omega$를 얻는다.

\medskip\noindent
$H_1, \ldots, H_p \subset H$를 $F \cup G$에 포함되지 않는
$H$의 기약성분들이라 하자. 이전처럼
$z_l \not \in F \cup G$이고 $l' \not = l$일 때
$z_l \not \in H_{l'}$인 $z_l \in H_l$을 고른다.
Algebra, Lemma \ref{algebra-lemma-chinese-remainder}를 사용하여
$j = 1, \ldots, m$에 대해 $f(y_j) = 1$이고
$l = 1, \ldots, p$에 대해 $f(z_l) = 0$인 $f \in R$을 고를 수 있다.
주장: 원소 $s = u + f t$가 원하는 조건을 만족한다.

\medskip\noindent
먼저
% P02-E1358
$t(x_i) = 0$이므로 $s(x_i)$의 값은 $u(x_i)$와 같고,
따라서 $s(x_i)$가 $v_i$로 감을 알 수 있다. 이는 (a)를 증명한다.
증명을 끝내려면 $F$에 포함되지 않는
$Z(s_1, \ldots, s_h, s)$의 모든 기약성분 $Z$가 $\Omega$에서
여차원 $\geq k$임을 보이면 충분하다. 실제로 $F'$을 이들의
합집합으로 놓으면 (b)와 (c)를 얻는다. 다음과 같이
$Z(s_1, \ldots, s_h, s)$의 여차원 $\leq k - 1$인 기약성분
$Z$는 존재하지 않음을 알 수 있다.
\begin{enumerate}
\item
% P02-E0733
$s = u + ft$이므로
$Z(s_1, \ldots, s_h, s) \subset Z(s_1, \ldots, s_h, u, t) = F \cup H$
임에 유의하라. 따라서 $Z \subset H$이다.
\item $H$의 기약성분들은 여차원 $\geq k - 1$을 가진다.
따라서 $Z$는 $H$의 한 기약성분과 같다.
이는 $Z$가 여차원 $\leq k - 1$이기 때문이다.
따라서 어떤 $l \in \{1, \ldots, p\}$에 대해 $Z = H_l$이거나,
% P02-E0734
어떤 $j \in \{1, \ldots m\}$에 대해 $Z = G_j$이다.
\item 그러나 $Z = G_j$는 불가능하다.
$s_1(y_j), \ldots, s_h(y_j)$는 $V(y_j)$의 선형독립인 원소들로 가고,
$s(y_j) = u(y_j) + f(y_j) t(y_j) = u(y_j) + t(y_j)$는
% P02-E0735
$$
\text{linear combination images of }s_i(y_j) + w_j
$$
꼴 원소로 가기 때문이다. $w_j$의 선택에 의해 이는 $V(y_j)$에서
$s_1(y_j), \ldots, s_h(y_j)$의 상들과 선형독립이다.
\item $Z = H_l$도 불가능하다. 실제로 다시
$s_1(z_l), \ldots, s_h(z_l)$는 $V(z_l)$의 선형독립인 원소들로 가고,
$s(z_l) = u(z_l) + f(z_l) t(z_l) = u(z_l)$은
$z_l \not \in F \cup G$이므로 $V(z_l)$에서 이들과 선형독립인
원소로 간다.
\end{enumerate}
이로써 증명이 끝난다.
\end{proof}

\begin{theorem}
\label{more-algebra-theorem-splitting}
\begin{reference}
\cite[Theorem 1]{Serre-projective}
\end{reference}
$R$을 그 극대 스펙트럼 $\Omega \subset \Spec(R)$가 차원
$d < \infty$인 Noether 위상공간인 환이라 하자.
$M$을 유한 표시 $R$-가군이라 하고, 모든 $\mathfrak m \in \Omega$에 대해
$R_\mathfrak m$-가군 $M_\mathfrak m$이 계수 $> d$인 자유
직합인자를 가진다고 하자. 그러면 $M \cong R \oplus M'$이다.
\end{theorem}

\begin{proof}
$\mathfrak m \in \Omega$에 대해 $R_\mathfrak m^{\oplus r}$이
$M_\mathfrak m$의 직합인자라고 하자. 그러면 Algebra, Lemma
\ref{algebra-lemma-localization-colimit}와
\ref{algebra-lemma-colimit-category-fp-modules}에 의해 어떤
$f \in R$, $f \not \in \mathfrak m$에 대해
$R_f^{\oplus r}$이 $M_f$의 직합인자임을 알 수 있다.
따라서 이 가정은 Lemma \ref{more-algebra-lemma-which-elements-split}에서와 같은
$V(x)$에 대해 모든 $x \in \Omega$에서 $\dim V(x) > d$임을 뜻한다.
Proposition \ref{more-algebra-proposition-splitting}을
$F = \emptyset$, $h = 0$이고 $s_i$가 없으며, $n = 0$이고
$x_i, v_i$가 없으며, $k = d + 1$인 경우에 적용하면
$F'$의 모든 기약성분이 여차원 $\geq d + 1$을 가지며
$Z(s) \subset F'$인 $s \in M$과 $F' \subset \Omega$를 얻는다.
$d = \dim(\Omega)$이므로 이는 $F' = \emptyset$을 강제한다.
따라서 $s : R \to M$은 모든 극대 아이디얼에서 직합인자의 포함이다.
그러므로 $s$는 보편 단사이다. Algebra, Lemma
\ref{algebra-lemma-universally-injective-check-stalks}를 보라.
이제 Algebra, Lemma \ref{algebra-lemma-universally-exact-split}에 의해
$s$는 분할 단사이다.
\end{proof}


\section{큰 사영 가군은 자유이다}
\label{more-algebra-section-big-projective-free}

\noindent
이 절에서는 \cite{Bass}의 결과 중 하나를 논한다. 독자에게 원 논문을
살펴볼 것을 권한다. 우리의 논증은 논문 \cite{Akasaki}와
\cite{Hinohara}에 주어진 약간 단순화된 증명을 사용한다.

\begin{lemma}[Eilenberg의 보조정리]
\label{more-algebra-lemma-eilenberg-swindle}
\begin{slogan}
Eilenberg 속임수
\end{slogan}
\begin{history}
\cite{Bass}에는 다음과 같이 쓰여 있다. “…… Eilenberg가 몇 년 전에
발견했으며 Barry Mazur의 이마에서 솟아났다고 해도 좋을 만큼 우아한
작은 속임수이다.”
\end{history}
\begin{reference}
\cite[Eilenberg's lemma]{Bass}
\end{reference}
$P \oplus Q \cong F$이고 $F$가 유한 생성이 아닌 자유 가군이면
$P \oplus F \cong F$이다.
\end{lemma}

\begin{proof}
$$
F \cong F \oplus F \oplus \ldots \cong
P \oplus Q \oplus P \oplus Q \oplus \ldots \cong
P \oplus F \oplus F \oplus \ldots \cong P \oplus F
$$
\end{proof}

\begin{lemma}
\label{more-algebra-lemma-projective-plus-free-is-free}
$R$을 환이라 하고 $P$를 사영 가군이라 하자.
$P \oplus F$가 자유인 자유 가군 $F$가 존재한다.
\end{lemma}

\begin{proof}
$P$가 사영이므로 어떤 가군 $Q$에 대해
$F_0 = P \oplus Q$가 자유 가군이다.
$F = \bigoplus_{n \geq 1} F_0$로 놓자. 그러면 Lemma
\ref{more-algebra-lemma-eilenberg-swindle}에 의해 $P \oplus F \cong F$이다.
\end{proof}

\begin{lemma}
\label{more-algebra-lemma-element-projective}
$R$을 환이라 하고 $P$를 사영 가군이라 하자.
$s \in P$라 하자. $(0, s) \in K$를 만족하는 유한 자유 가군 $F$와
유한 자유 직합인자 $K \subset F \oplus P$가 존재한다.
\end{lemma}

\begin{proof}
Lemma \ref{more-algebra-lemma-projective-plus-free-is-free}에 의해
$F \oplus P$가 자유가 되는 (무한일 수도 있는) 자유 가군 $F$를
찾을 수 있다. 그러면 물론 $(0, s)$는 유한 자유 직합인자
$K \subset F \oplus P$에 포함된다. 다시 $K$는
유한 자유 직합인자 $F' \subset F$에 대한 $F' \oplus P$에 포함된다.
\end{proof}

\begin{lemma}
\label{more-algebra-lemma-trick-to-find-good-element}
$R$을 Jacobson 근기 $J$를 가지며 $R/J$가 Noether인 환이라 하자.
$P$를 $R$의 모든 극대 아이디얼 $\mathfrak m$에 대해
$P_\mathfrak m$이 무한 계수를 가지는 사영 $R$-가군이라 하자.
$s \in P$이고 $M \subset P$이며 $Rs + M = P$라 하자.
그러면 $R(s + m)$이 $P$의 자유 직합인자가 되는
$m \in M$을 찾을 수 있다.
\end{lemma}

\begin{proof}
Algebra, Theorem
\ref{algebra-theorem-projective-free-over-local-ring}에 의해
$P_\mathfrak m$이 자유이므로 이 명제는 의미가 있다.

\medskip\noindent
$M$의 상을 $M' \subset P/JP$로, $s$의 상을 $s' \in P/JP$로 나타내자.
% P02-E0736
$R/J s' + M' = P/JP$임에 유의하라.
$R/J(s' + m')$이 $M'$의 자유 직합인자가 되는
% P02-E0737
$m' \in M'$을 찾을 수 있다고 하자.
분할을 주는 $\varphi' : P/JP \to R/J$, 즉 $R/J$에서
$\varphi'(s' + m') = 1$인 것을 고르자.
$P$는 사영 $R$-가군이므로 $\varphi'$의 올림
$\varphi : P \to R$을 찾을 수 있다. $m'$으로 가는 $m \in M$을 고르자.
그러면 $\varphi(s + m) \in R$은 $J$를 법으로 $1$과 합동이므로
$R$의 단위원이다(Algebra, Lemma
\ref{algebra-lemma-contained-in-radical}).
따라서 $R(s + m)$은 $P$의 자유 직합인자이다.
이로써 다음 문단에서 논하는 경우로 환원된다.

\medskip\noindent
$R$이 Noether라고 하자. $m \in M$을 원소라 하고
$\varphi_1, \ldots, \varphi_n : P \to R$을 $R$-선형 사상들이라 하자.
$$
Z(s + m, \varphi_1, \ldots, \varphi_n) \subset \Spec(R)
$$
로 $\varphi_1(s + m), \ldots, \varphi_n(s + m) \in R$의
영점 집합을 나타내자.

\medskip\noindent
$\mathfrak m$이 $R$의 극대 아이디얼이고
$\mathfrak m \in Z(s + m, \varphi_1, \ldots, \varphi_n)$이라고 하자.
$K = M \cap \bigcap \Ker(\varphi_i)$로 놓자. 사상
$$
K/\mathfrak mK \to P/\mathfrak m P
$$
의 상이 무한 차원이라고 주장한다. 실제로 $P/K$는
$P/M \oplus R^{\oplus n}$의 부분가군과 동형이므로 유한 $R$-가군이다.
따라서 위 사상의 핵은
$\text{Tor}_1^R(P/K, \kappa(\mathfrak m))$의 몫이며,
Algebra, Lemma \ref{algebra-lemma-tor-noetherian}에 의해 이는 유한이다.
$P/\mathfrak mP$가 무한 차원이라는 사실과 합치면 주장을 얻는다.
따라서 $P/\mathfrak mP$의 벡터 $\overline{t}$으로 가는
$t \in K$를 찾을 수 있으며, 이는 $s + m$의 상
$\overline{s + m}$과 선형독립이다. 선형대수에 의해
% P02-E0738
$\overline{\varphi} : P \to \kappa(\mathfrak m)$인 $R$-선형 사상으로서
$\overline{\varphi}(\overline{t}) = 1$이고
$\overline{\varphi}(\overline{s + m}) = 0$인 것을 찾을 수 있다.
$P$가 사영이므로 $\overline{\varphi}$를 올리는 $R$-선형 사상
$\varphi : P \to R$을 찾을 수 있다. 그러면 영점 집합
$Z(s + m + t, \varphi_1, \ldots, \varphi_n, \varphi)$는
$Z(s + m, \varphi_1, \ldots, \varphi_n)$에 포함되지만
$\mathfrak m$을 포함하지 않는다. 즉,
$Z(s + m, \varphi_1, \ldots, \varphi_n)$보다 진부분집합이다.

\medskip\noindent
$\Spec(R)$은 Noether 위상공간이므로 위 논증으로부터
닫힌 부분집합 $Z(s + m, \varphi_1, \ldots, \varphi_n)$이
$\Spec(R)$의 닫힌 점을 하나도 포함하지 않게 하는
$m \in M$과 $\varphi_1, \ldots, \varphi_n : P \to R$을 찾을 수 있다.
따라서 $Z(s + m, \varphi_1, \ldots, \varphi_n) = \emptyset$이다.
따라서 $\sum r_i\varphi_i(s + m) = 1$이 되게 하는
$r_1, \ldots, r_n \in R$을 찾을 수 있다. 그러면
$$
R \xrightarrow{s + m} P \xrightarrow{\sum r_i \varphi_i} R
$$
은 원하는 분할이다.
\end{proof}

\begin{lemma}
\label{more-algebra-lemma-element-in-free-summand}
$R$을 Jacobson 근기 $J$를 가지며 $R/J$가 Noether인 환이라 하자.
$P$를 $R$의 모든 극대 아이디얼 $\mathfrak m$에 대해
$P_\mathfrak m$이 무한 계수를 가지는 사영 $R$-가군이라 하자.
$s \in P$라 하자. $s \in M$인 유한 안정 자유 직합인자
$M \subset P$를 찾을 수 있다.
\end{lemma}

\begin{proof}
Lemma \ref{more-algebra-lemma-element-projective}에 의해 $(0, s) \in K$인
유한 자유 가군 $F$와 유한 자유 직합인자
$K \subset F \oplus P$를 찾을 수 있다.
$F$의 계수에 대한 귀납법으로 다음 문단에서 논하는 경우로 환원된다.

\medskip\noindent
$(0, s) \in K$인 유한 안정 자유 직합인자
$K \subset R \oplus P$가 존재한다고 하자. $K$의 여가군 $K'$,
즉 $R \oplus P = K \oplus K'$인 것을 고르자. 사영
% P02-E0739
$\pi : R \oplus P \to K'$는 전사이므로 Lemma
\ref{more-algebra-lemma-trick-to-find-good-element}에 의해
$\pi(1, p) \in K'$가 자유 직합인자를 생성하게 하는
$p \in P$를 찾을 수 있다. 따라서
$K' = R\pi(1, p) \oplus K''$로 쓴다. 그러면
$$
R \oplus P = K \oplus K' = K \oplus R\pi(1, p) \oplus K''
$$
이다. 사영 $\pi' : P \to K''$는 전사이다\footnote{실제로
$k'' \in K''$이면, $K'$의 원소로 본 $k''$은 어떤
$\lambda \in R$과 $q \in P$에 대해
$k'' = \lambda \pi(1, 0) + \pi(0, q)$로 쓸 수 있다.
이는 $k'' = \lambda \pi(1, p) + \pi(0, q - \lambda p)$임을 뜻한다.
다시 이는 $K' \to K''$가 $\pi(1, p)$을 영으로 보내므로
$q - \lambda p$가 합성
$P \to R \oplus P \xrightarrow{\pi} K' \to K''$ 아래에서
$k''$으로 감을 뜻한다.}
% P02-E0740
따라서 이는 분할된다($K''$이 사영이므로).
그러므로 $\Ker(\pi') \subset P$는 $s$를 포함하는 직합인자이다.
마지막으로 구성에 의해 동형
$$
R \oplus \Ker(\pi') \cong K \oplus R\pi(1, p)
$$
이 있다. 따라서 $K$가 유한 안정 자유이므로 $\Ker(\pi')$도 그러하다.
\end{proof}

\begin{theorem}
\label{more-algebra-theorem-countable-free}
\begin{reference}
\cite[Theorem 3.1]{Bass}의 가환인 경우
\end{reference}
$R$을 Jacobson 근기 $J$를 가지며 $R/J$가 Noether인 환이라 하자.
$P$를 가산 생성 사영 $R$-가군이라 하고, $R$의 모든 극대 아이디얼
$\mathfrak m$에 대해 $P_\mathfrak m$이 무한 계수를 가진다고 하자.
그러면 $P$는 자유이다.
\end{theorem}

\begin{proof}
먼저 $P$가 유한 안정 자유 가군들의 가산 직합임을 증명하자.
$x_1, x_2, \ldots$를 $P$의 가산 생성계라 하자.
모든 $n$에 대해 $F_1 \oplus \ldots \oplus F_n$이
$x_1, \ldots, x_n$을 포함하는 $P$의 직합인자가 되게 하는
$P$의 유한 안정 자유 직합인자 $F_1, F_2, \ldots$를 귀납적으로 구성한다.
실제로 원하는 성질을 가진 $F_1, \ldots, F_n$이 주어지면
$$
P = F_1 \oplus \ldots \oplus F_n \oplus P'
$$
로 쓰고, $x_{n + 1}$의 상을 $s \in P'$이라 하자.
Lemma \ref{more-algebra-lemma-element-in-free-summand}에 의해 $s$를 포함하는
유한 안정 자유 직합인자 $F_{n + 1} \subset P'$를 찾을 수 있다.
그러면 $P = \bigoplus_{i = 1}^{\infty} F_i$이다.

\medskip\noindent
$P$가 영이 아닌 유한 안정 자유 가군들의 무한 직합
$P = \bigoplus_{i = 1}^{\infty} F_i$라고 하자.
가군 $F_i$들의 안정 자유성을 다음과 같이 사용한다.
각 $F_i$의 계수는 일정하다(그리고 양수이다). 따라서 $R$의 각
극대 아이디얼 $\mathfrak m$에 대해 $P_\mathfrak m$은 가산 무한 계수의
자유 가군이다. Lemma \ref{more-algebra-lemma-trick-to-find-good-element}를
$s = 0$과 $M = P$에 적용하면 $Rt_1$이 $P$의 자유 직합인자가 되게 하는
$t_1 \in P$를 찾을 수 있다. 그러면 어떤 $n_1 > n_0 = 0$에 대해
$t_1$은 $F_1 \oplus \ldots \oplus F_{n_1}$에 포함된다.
$\bigoplus_{n > n_1} F_n$에 같은 논증을 적용하면
$n_1 < n_2$와, 자유 직합인자를 생성하는
$t_2 \in F_{n_1 + 1} \oplus \ldots \oplus F_{n_2}$를 얻는다.
이와 같이 계속하면 무한 계수의 자유 직합인자
$$
\bigoplus\nolimits_{i \geq 1} t_i :
\bigoplus\nolimits_{i \geq 1} R
\longrightarrow
\bigoplus\nolimits_{i \geq 1}
\bigoplus\nolimits_{n_i \geq n > n_{i - 1}} F_n = P
$$
를 얻는다. 따라서 어떤 가산 계수 자유 $R$-가군 $F$에 대해
$P \cong Q \oplus F$이다. $Q$가 가산 생성이므로 어떤 가군 $Q'$에 대해
$Q \oplus Q' \cong F$이다. 그러면 Eilenberg 속임수
(Lemma \ref{more-algebra-lemma-eilenberg-swindle})로부터
$Q \oplus F \cong F$이고 $P$가 자유임을 얻는다.
\end{proof}
